\documentclass[10pt,oneside]{book}

\usepackage[T1]{fontenc}
\usepackage[utf8]{inputenc}
\usepackage[french]{babel}
\usepackage{amsmath,amssymb,mathrsfs}
\usepackage[all]{xy}
\usepackage{tikz-cd}
\usepackage[shortlabels]{enumitem}
\usepackage{xparse}
\usepackage{xstring}
\usepackage[a4paper,margin=22mm]{geometry}
\usepackage[
  colorlinks=true,
  linkcolor=blue,
  filecolor=blue,
  hypertexnames=false,
  unicode=true,
  pdfencoding=auto
]{hyperref}
\hypersetup{
  pdftitle={Éléments de géométrie algébrique IV, quatrième partie -- édition française diplomatique},
  pdfauthor={Alexandre Grothendieck et Jean Dieudonné},
  pdfsubject={Transcription diplomatique française depuis les exemplaires NUMDAM}
}

\renewcommand{\thesection}{\arabic{section}}
\renewcommand{\thesubsection}{\thesection.\arabic{subsection}}
%
% DEBUT SOURCE preamble-base.tex
% Notation of the EGA editions: the abbreviations the transcription uses for alphabets, symbols, operators and
% arrows. Written for these editions and dedicated to the public domain under CC0 1.0.
\makeatletter
% a command that may or may not exist already is (re)defined in one step
\newcommand{\ega@define}[2]{\providecommand{#1}{}\renewcommand{#1}{#2}}
\newcommand{\ega@operator}[2]{\providecommand{#1}{}\renewcommand{#1}{\operatorname{#2}}}

% alphabets: sheaves, bold letters, categories
\ega@define{\sh}{\mathcal}
\ega@define{\bb}{\mathbf}
\ega@define{\cat}{\mathtt}
\ega@define{\C}{\cat{C}}
\ega@define{\Set}{\cat{Set}}
\ega@define{\CHom}{\cat{Hom}}
\ega@define{\op}{^{\cat{op}}}
\ega@define{\GL}{\bb{GL}}
\ega@define{\bbGamma}{\boldsymbol{\Gamma}}

% letters and symbols
\renewcommand{\leq}{\leqslant}
\renewcommand{\geq}{\geqslant}
\ega@define{\setmin}{-}
\ega@define{\emp}{\varnothing}
\ega@define{\vphi}{\phi}
\ega@define{\kres}{k}
\ega@define{\rad}{\mathfrak{r}}
\ega@define{\nilrad}{\mathfrak{N}}
\ega@define{\HH}{\mathrm{H}}
\ega@define{\CHH}{\check{\mathrm{H}}}
\ega@define{\RR}{\mathrm{R}}
\ega@define{\LL}{\mathrm{L}}
\ega@define{\red}{\mathrm{red}}
\providecommand{\dual}{}\renewcommand{\dual}[1]{{#1}^{\vee}}
% the tilde of the sheaf associated with a module, written after it: M\supertilde
\ega@define{\supertilde}{{\,\widetilde{\,}\,}}

% operators
\ega@operator{\rg}{rg}
\ega@operator{\gr}{gr}
\ega@operator{\Gr}{Gr}
\ega@operator{\Hom}{Hom}
\ega@operator{\Proj}{Proj}
\ega@operator{\Tor}{Tor}
\ega@operator{\Ext}{Ext}
\ega@operator{\Supp}{Supp}
\ega@operator{\Ker}{Ker}
\ega@operator{\Im}{Im}
\ega@operator{\Coker}{Coker}
\ega@operator{\Spec}{Spec}
\ega@operator{\Spm}{Spm}
\ega@operator{\Spf}{Spf}
\ega@operator{\grad}{grad}
\ega@operator{\dimc}{dimc}
\ega@operator{\codim}{codim}
\ega@operator{\id}{id}
\ega@operator{\Der}{Der}
\ega@operator{\Diff}{Diff}
\ega@operator{\Hyp}{Hyp}
\ega@operator{\Tr}{Tr}

% arrows: long in displays, short in text, scripts and subscripts
\let\ega@shortmapsto\mapsto
\renewcommand{\to}{\mathchoice{\longrightarrow}{\rightarrow}{\rightarrow}{\rightarrow}}
\providecommand{\from}{}\renewcommand{\from}{\mathchoice{\longleftarrow}{\leftarrow}{\leftarrow}{\leftarrow}}
\renewcommand{\mapsto}{\mathchoice{\longmapsto}{\ega@shortmapsto}{\ega@shortmapsto}{\ega@shortmapsto}}
\ega@define{\isoto}{\simeq}
\ega@define{\simto}{\xrightarrow{\sim}}

% wide accents
\renewcommand{\hat}[1]{\widehat{#1}}
\renewcommand{\tilde}[1]{\widetilde{#1}}
\makeatother
% FIN SOURCE preamble-base.tex

\makeatletter
\newcommand{\EGAUnavailableHyperref}[2]{%
  \begingroup
  \let\@latex@warning\@gobble
  \let\@latex@warning@no@line\@gobble
  \let\G@refundefinedtrue\relax
  \hyperref[#1]{#2}%
  \endgroup
}
\NewDocumentCommand{\sref}{o m o}{%
  \@ifundefined{r@#2}{\EGAUnavailableHyperref{#2}}{\hyperref[#2]}%
  {\normalfont{(\IfValueT{#1}{\textbf{#1},~}\StrBehind{#2}{.}\IfValueT{#3}{,~{#3}})}}%
}
\newcommand{\EGAIVIndexExternalLabel}[3]{%
  \expandafter\gdef\csname r@canon-index.#1\endcsname{{#1}{0}{}{ega-fr:#2:#3}{EGA_French_Diplomatic_Cumulative_Reader.pdf}}}
\EGAIVIndexExternalLabel{ega4.c04.cor.13-1-4}{iv3}{ega4.c04.cor.13-1-4}
\EGAIVIndexExternalLabel{ega4.c04.cor.14-4-4}{iv3}{ega4.c04.cor.14-4-4}
\EGAIVIndexExternalLabel{ega4.c04.cor.14-5-11}{iv3}{ega4.c04.cor.14-5-11}
\EGAIVIndexExternalLabel{ega4.c04.cor.15-5-2}{iv3}{ega4.c04.cor.15-5-2}
\EGAIVIndexExternalLabel{ega4.c04.cor.15-6-2}{iv3}{ega4.c04.cor.15-6-2}
\EGAIVIndexExternalLabel{ega4.c04.cor.3-1-5}{iv2}{ega4.c04.cor.3-1-5}
\EGAIVIndexExternalLabel{ega4.c04.cor.4-5-11}{iv2}{ega4.c04.cor.4-5-11}
\EGAIVIndexExternalLabel{ega4.c04.cor.5-11-2}{iv2}{ega4.c04.cor.5-11-2}
\EGAIVIndexExternalLabel{ega4.c04.cor.5-11-4}{iv2}{ega4.c04.cor.5-11-4}
\EGAIVIndexExternalLabel{ega4.c04.cor.5-2-3}{iv2}{ega4.c04.cor.5-2-3}
\EGAIVIndexExternalLabel{ega4.c04.cor.5-5-4}{iv2}{ega4.c04.cor.5-5-4}
\EGAIVIndexExternalLabel{ega4.c04.cor.6-1-3}{iv2}{ega4.c04.cor.6-1-3}
\EGAIVIndexExternalLabel{ega4.c04.cor.6-7-2}{iv2}{ega4.c04.cor.6-7-2}
\EGAIVIndexExternalLabel{ega4.c04.cor.7-2-9}{iv2}{ega4.c04.cor.7-2-9}
\EGAIVIndexExternalLabel{ega4.c04.cor.8-12-10}{iv3}{ega4.c04.cor.8-12-10}
\EGAIVIndexExternalLabel{ega4.c04.cor.8-12-11}{iv3}{ega4.c04.cor.8-12-11}
\EGAIVIndexExternalLabel{ega4.c04.cor.8-2-11}{iv3}{ega4.c04.cor.8-2-11}
\EGAIVIndexExternalLabel{ega4.c04.cor.8-5-8}{iv3}{ega4.c04.cor.8-5-8}
\EGAIVIndexExternalLabel{ega4.c04.def.7-8-2}{iv2}{ega4.c04.def.7-8-2}
\EGAIVIndexExternalLabel{ega4.c04.item.12-1-2}{iv3}{ega4.c04.item.12-1-2}
\EGAIVIndexExternalLabel{ega4.c04.item.15-6-10}{iv3}{ega4.c04.item.15-6-10}
\EGAIVIndexExternalLabel{ega4.c04.item.6-5-5}{iv2}{ega4.c04.item.6-5-5}
\EGAIVIndexExternalLabel{ega4.c04.item.7-5-4}{iv2}{ega4.c04.item.7-5-4}
\EGAIVIndexExternalLabel{ega4.c04.item.7-8-4}{iv2}{ega4.c04.item.7-8-4}
\EGAIVIndexExternalLabel{ega4.c04.lem.1-8-7-1}{iv1}{ega4.c04.lem.1-8-7-1}
\EGAIVIndexExternalLabel{ega4.c04.lem.2-3-1}{iv2}{ega4.c04.lem.2-3-1}
\EGAIVIndexExternalLabel{ega4.c04.lem.4-4-2}{iv2}{ega4.c04.lem.4-4-2}
\EGAIVIndexExternalLabel{ega4.c04.lem.5-6-11-1}{iv2}{ega4.c04.lem.5-6-11-1}
\EGAIVIndexExternalLabel{ega4.c04.prop.1-10-3}{iv1}{ega4.c04.prop.1-10-3}
\EGAIVIndexExternalLabel{ega4.c04.prop.10-2-1}{iv3}{ega4.c04.prop.10-2-1}
\EGAIVIndexExternalLabel{ega4.c04.prop.11-3-4}{iv3}{ega4.c04.prop.11-3-4}
\EGAIVIndexExternalLabel{ega4.c04.prop.14-3-8}{iv3}{ega4.c04.prop.14-3-8}
\EGAIVIndexExternalLabel{ega4.c04.prop.15-4-2}{iv3}{ega4.c04.prop.15-4-2}
\EGAIVIndexExternalLabel{ega4.c04.prop.15-6-6}{iv3}{ega4.c04.prop.15-6-6}
\EGAIVIndexExternalLabel{ega4.c04.prop.2-3-4}{iv2}{ega4.c04.prop.2-3-4}
\EGAIVIndexExternalLabel{ega4.c04.prop.2-4-3}{iv2}{ega4.c04.prop.2-4-3}
\EGAIVIndexExternalLabel{ega4.c04.prop.2-4-5}{iv2}{ega4.c04.prop.2-4-5}
\EGAIVIndexExternalLabel{ega4.c04.prop.2-5-2}{iv2}{ega4.c04.prop.2-5-2}
\EGAIVIndexExternalLabel{ega4.c04.prop.2-6-1}{iv2}{ega4.c04.prop.2-6-1}
\EGAIVIndexExternalLabel{ega4.c04.prop.3-1-7}{iv2}{ega4.c04.prop.3-1-7}
\EGAIVIndexExternalLabel{ega4.c04.prop.3-3-1}{iv2}{ega4.c04.prop.3-3-1}
\EGAIVIndexExternalLabel{ega4.c04.prop.4-2-1}{iv2}{ega4.c04.prop.4-2-1}
\EGAIVIndexExternalLabel{ega4.c04.prop.4-5-19}{iv2}{ega4.c04.prop.4-5-19}
\EGAIVIndexExternalLabel{ega4.c04.prop.4-7-3}{iv2}{ega4.c04.prop.4-7-3}
\EGAIVIndexExternalLabel{ega4.c04.prop.5-10-2}{iv2}{ega4.c04.prop.5-10-2}
\EGAIVIndexExternalLabel{ega4.c04.prop.5-6-1}{iv2}{ega4.c04.prop.5-6-1}
\EGAIVIndexExternalLabel{ega4.c04.prop.5-6-5}{iv2}{ega4.c04.prop.5-6-5}
\EGAIVIndexExternalLabel{ega4.c04.prop.5-9-8}{iv2}{ega4.c04.prop.5-9-8}
\EGAIVIndexExternalLabel{ega4.c04.prop.6-1-5}{iv2}{ega4.c04.prop.6-1-5}
\EGAIVIndexExternalLabel{ega4.c04.prop.6-14-1}{iv2}{ega4.c04.prop.6-14-1}
\EGAIVIndexExternalLabel{ega4.c04.prop.6-14-4}{iv2}{ega4.c04.prop.6-14-4}
\EGAIVIndexExternalLabel{ega4.c04.prop.6-15-7}{iv2}{ega4.c04.prop.6-15-7}
\EGAIVIndexExternalLabel{ega4.c04.prop.6-2-3}{iv2}{ega4.c04.prop.6-2-3}
\EGAIVIndexExternalLabel{ega4.c04.prop.6-3-1}{iv2}{ega4.c04.prop.6-3-1}
\EGAIVIndexExternalLabel{ega4.c04.prop.6-5-1}{iv2}{ega4.c04.prop.6-5-1}
\EGAIVIndexExternalLabel{ega4.c04.prop.7-8-6}{iv2}{ega4.c04.prop.7-8-6}
\EGAIVIndexExternalLabel{ega4.c04.prop.8-10-1}{iv3}{ega4.c04.prop.8-10-1}
\EGAIVIndexExternalLabel{ega4.c04.prop.8-10-3}{iv3}{ega4.c04.prop.8-10-3}
\EGAIVIndexExternalLabel{ega4.c04.prop.8-12-3}{iv3}{ega4.c04.prop.8-12-3}
\EGAIVIndexExternalLabel{ega4.c04.prop.8-9-1}{iv3}{ega4.c04.prop.8-9-1}
\EGAIVIndexExternalLabel{ega4.c04.rem.14-5-7}{iv3}{ega4.c04.rem.14-5-7}
\EGAIVIndexExternalLabel{ega4.c04.statement.1-7-4}{iv1}{ega4.c04.statement.1-7-4}
\EGAIVIndexExternalLabel{ega4.c04.statement.10-4-11}{iv3}{ega4.c04.statement.10-4-11}
\EGAIVIndexExternalLabel{ega4.c04.statement.11-7-5}{iv3}{ega4.c04.statement.11-7-5}
\EGAIVIndexExternalLabel{ega4.c04.statement.13-3-9}{iv3}{ega4.c04.statement.13-3-9}
\EGAIVIndexExternalLabel{ega4.c04.statement.2-1-1}{iv2}{ega4.c04.statement.2-1-1}
\EGAIVIndexExternalLabel{ega4.c04.statement.6-10-1}{iv2}{ega4.c04.statement.6-10-1}
\EGAIVIndexExternalLabel{ega4.c04.statement.6-15-2}{iv2}{ega4.c04.statement.6-15-2}
\EGAIVIndexExternalLabel{ega4.c04.statement.6-15-4}{iv2}{ega4.c04.statement.6-15-4}
\EGAIVIndexExternalLabel{ega4.c04.statement.7-3-19}{iv2}{ega4.c04.statement.7-3-19}
\EGAIVIndexExternalLabel{ega4.c04.statement.7-8-3}{iv2}{ega4.c04.statement.7-8-3}
\EGAIVIndexExternalLabel{ega4.c04.statement.8-1-2}{iv3}{ega4.c04.statement.8-1-2}
\EGAIVIndexExternalLabel{ega4.c04.subsec.10-4}{iv3}{ega4.c04.subsec.10-4}
\EGAIVIndexExternalLabel{ega4.c04.thm.11-2-9}{iv3}{ega4.c04.thm.11-2-9}
\EGAIVIndexExternalLabel{ega4.c04.thm.11-3-10}{iv3}{ega4.c04.thm.11-3-10}
\EGAIVIndexExternalLabel{ega4.c04.thm.12-1-1}{iv3}{ega4.c04.thm.12-1-1}
\EGAIVIndexExternalLabel{ega4.c04.thm.14-2-1}{iv3}{ega4.c04.thm.14-2-1}
\EGAIVIndexExternalLabel{ega4.c04.thm.14-4-1}{iv3}{ega4.c04.thm.14-4-1}
\EGAIVIndexExternalLabel{ega4.c04.thm.5-10-5}{iv2}{ega4.c04.thm.5-10-5}
\EGAIVIndexExternalLabel{ega4.c04.thm.5-8-6}{iv2}{ega4.c04.thm.5-8-6}
\EGAIVIndexExternalLabel{ega4.c04.thm.7-6-4}{iv2}{ega4.c04.thm.7-6-4}
\EGAIVIndexExternalLabel{ega4.c04.thm.7-7-2}{iv2}{ega4.c04.thm.7-7-2}
\EGAIVIndexExternalLabel{ega4.c04.thm.8-10-5}{iv3}{ega4.c04.thm.8-10-5}
\EGAIVIndexExternalLabel{ega4.c04.thm.8-11-1}{iv3}{ega4.c04.thm.8-11-1}
\EGAIVIndexExternalLabel{ega4.c04.thm.8-8-2}{iv3}{ega4.c04.thm.8-8-2}
\newcommand{\EGAIVFourExternalLabel}[3]{%
  \expandafter\gdef\csname r@#1\endcsname{{#1}{0}{}{ega-fr:#2:canon-ref.#3}{EGA_French_Diplomatic_Cumulative_Reader.pdf}}}
\newcommand{\EGAIVFourExternalCitation}[2]{%
  \expandafter\gdef\csname b@#1\endcsname{#2}}
\makeatother
\newcommand{\oldpage}[2][]{%
  \setcounter{footnote}{0}%
  \marginpar{\scriptsize\textbf{#1}~|~#2}\ignorespaces}
\newenvironment{env}[1][]{\par\medskip\noindent\textbf{(#1)}\ }{\par}
\newenvironment{definition}[1][]{%
  \par\medskip\noindent\itshape Définition (#1). ---\ }{\par}
\newenvironment{examples}[1][]{%
  \par\medskip\noindent\itshape Exemples (#1). ---\ }{\par}
\newenvironment{proposition}[1][]{%
  \par\medskip\noindent\itshape Proposition (#1). ---\ }{\par}
\newenvironment{theorem}[1][]{%
  \par\medskip\noindent\itshape Théorème (#1). ---\ }{\par}
\newenvironment{corollary}[1][]{%
  \par\medskip\noindent\itshape Corollaire (#1). ---\ }{\par}
\newenvironment{remark}[1][]{%
  \par\medskip\noindent\itshape Remarque (#1). ---\ }{\par}
\newenvironment{remarks}[1][]{%
  \par\medskip\noindent\itshape Remarques (#1). ---\ }{\par}
\newenvironment{lemma}[1][]{%
  \par\medskip\noindent\itshape Lemme (#1). ---\ }{\par}
\newenvironment{proof}{%
  \par\medskip\noindent\textit{Démonstration. ---\ }}{\par}

\makeatletter
\newcommand{\EGAReaderTarget}[2]{\hypertarget{#1}{}\@bsphack\begingroup\def\@currentHref{#1}\label{#2}\endgroup\@esphack}
\makeatother
\begin{document}
% BEGIN RIGHTS AND AI TRANSCRIPTION NOTICE
\clearpage
\begingroup
\small
\noindent\textbf{Droits, disponibilité et provenance}

Les droits sur les œuvres originales et sur toute contribution protégée à ces traductions et transcriptions dérivées restent ceux de leurs titulaires respectifs. Ce projet ne place pas ces œuvres ou contributions dans le domaine public et n'accorde aucune licence sur elles. Les mentions de droits de tiers et les autorisations déjà accordées restent applicables.

Ces éditions sont mises à disposition gratuitement, à des fins non commerciales et exclusivement pédagogiques, dans la présomption que les titulaires de droits ne s'opposeraient pas à la libre disponibilité de ces traductions et transcriptions aux côtés des originaux français sur NUMDAM (www.numdam.org/journals/PMIHES/). Cette présomption ne signifie pas qu'une autorisation a été obtenue ; le libre accès aux originaux n'autorise pas, à lui seul, leur traduction, leur adaptation ou leur republication.

Si vous détenez des droits sur un élément présenté ici et vous opposez à sa disponibilité, veuillez ouvrir une demande sur GitHub (github.com/KokunoYumeto/ega-fr/issues). Le responsable de la publication retirera l'élément concerné à la demande du titulaire des droits.

La transcription moderne, les corrections, la composition et les notes mathématiques sont produites par IA. Les travaux documentés comprennent ceux d'OpenAI Codex, GPT-6 Astra, effort Ultra. Les indications de provenance des contributions antérieures restent applicables. Les textes originaux restent l'œuvre d'Alexandre Grothendieck et Jean Dieudonné ; ils ne sont pas attribués à l'IA. Cette édition n'est approuvée ni par les auteurs, ni par l'IHÉS, ni par NUMDAM.
\par
\endgroup
\clearpage
% END RIGHTS AND AI TRANSCRIPTION NOTICE

\EGAIVFourExternalLabel{0.4.3.3}{base}{0.4.3.3-fr}
\EGAIVFourExternalLabel{0.5.4.3}{base}{0.5.4.3-fr}
\EGAIVFourExternalLabel{0.5.4.5}{base}{0.5.4.5-fr}
\EGAIVFourExternalLabel{0.5.4.7}{base}{0.5.4.7-fr}
\EGAIVFourExternalLabel{0.6.1.2}{base}{0.6.1.2-fr}
\EGAIVFourExternalLabel{0.16.4.6}{iv1}{0.16.4.6-fr}
\EGAIVFourExternalCitation{IV-37}{37}
\setcounter{section}{15}
%
% DEBUT SOURCE ega4/ega4-16-fr.tex
\section{Invariants différentiels. Morphismes différentiellement lisses}
\label{section:IV.16}

\oldpage[IV-4]{5}
Dans ce paragraphe, nous présentons, sous forme globale, quelques notions de
calcul différentiel particulièrement utiles en Géométrie algébrique. Nous passons sous
silence de nombreux développements, classiques en Géométrie différentielle (connexions,
transformations infinitésimales associées à un champ de vecteurs, jets, etc.), bien que ces
notions s'écrivent de façon particulièrement naturelle dans le cadre des schémas. Nous
passons également sous silence ici les phénomènes spéciaux à la caractéristique $p>0$
(dont certains sont étudiés, dans le cadre affine, dans (0, 21)). Pour certains compléments
sur le formalisme différentiel dans les préschémas, le lecteur pourra consulter les
exposés II et VII de [42], ainsi que les chapitres ultérieurs de ce Traité.

\subsection{Invariants normaux d'une immersion}
\EGAReaderTarget{ega4.c04.subsec.16-1}{IV.16.1}

\begin{env}[16.1.1]
\label{IV.16.1.1}
Soient $(X,\mathcal{O}_X)$, $(Y,\mathcal{O}_Y)$ deux espaces annelés,
$f=(\psi,\theta):Y\to X$ un morphisme d'espaces annelés (0\textsubscript{I}, 4.1.1)
tel que l'homomorphisme
\[
  \theta^\#:\psi^*(\mathcal{O}_X)\longrightarrow\mathcal{O}_Y
\]
soit surjectif, de sorte que $\mathcal{O}_Y$ s'identifie à un faisceau d'anneaux quotient
$\psi^*(\mathcal{O}_X)/\mathscr{I}_f$. On peut alors munir
$\psi^*(\mathcal{O}_X)$ de la filtration $\mathscr{I}_f$-préadique.
\end{env}

\begin{definition}[16.1.2]
\EGAReaderTarget{ega4.c04.def.16-1-2}{IV.16.1.2}
Le faisceau d'anneaux $\mathcal{O}_Y$-augmenté
$\psi^*(\mathcal{O}_X)/\mathscr{I}_f^{n+1}$ est appelé le $n$-ème invariant normal
de $f$; l'espace annelé
$(Y,\psi^*(\mathcal{O}_X)/\mathscr{I}_f^{n+1})$ est appelé le $n$-ème voisinage
infinitésimal de $Y$ pour le morphisme $f$, et noté $Y_f^{(n)}$ ou simplement
$Y^{(n)}$. Le faisceau d'anneaux gradués associé au faisceau d'anneaux filtrés
$\psi^*(\mathcal{O}_X)$
\[
  \mathscr{G}r_\bullet(f)=\bigoplus_{n\geq 0}
  (\mathscr{I}_f^n/\mathscr{I}_f^{n+1})
  \tag{16.1.2.1}
  \label{IV.16.1.2.1}
\]
est appelé le faisceau d'anneaux gradués associé à $f$. Le faisceau
$\mathscr{G}r_1(f)=\mathscr{I}_f/\mathscr{I}_f^2$ est appelé le faisceau conormal
de $f$ (que l'on note aussi $\mathscr{N}_{Y/X}$ s'il n'en résulte pas de confusion).
\end{definition}

Il est clair que les
$\mathcal{O}_{Y^{(n)}}=\psi^*(\mathcal{O}_X)/\mathscr{I}_f^{n+1}$
(qu'on note aussi $\mathcal{O}_{Y_f^{(n)}}$) forment un
\oldpage[IV-4]{6}
système projectif de faisceaux d'anneaux sur $Y$, l'homomorphisme de transition
$\varphi_{nm}:\mathcal{O}_{Y^{(m)}}\to\mathcal{O}_{Y^{(n)}}$ pour $n\leq m$
identifiant $\mathcal{O}_{Y^{(n)}}$ au quotient de $\mathcal{O}_{Y^{(m)}}$ par la
puissance $(\mathscr{I}_f/\mathscr{I}_f^{n+1})^m$ de l'Idéal d'augmentation de
$\mathcal{O}_{Y^{(n)}}$, noyau de
$\varphi_{0n}:\mathcal{O}_{Y^{(n)}}\to\mathcal{O}_Y$. Les $Y^{(n)}$ forment par
suite un système inductif d'espaces annelés, ayant tous pour espace sous-jacent
l'espace $Y$, et on a des morphismes canoniques d'espaces annelés
$h_n:Y^{(n)}\to X$ égaux à $(\psi,\theta_n)$, où $\theta_n^\#$ est le morphisme
canonique
$\psi^*(\mathcal{O}_X)\to\psi^*(\mathcal{O}_X)/\mathscr{I}_f^{n+1}$.
Il est clair que le faisceau $\mathscr{G}r_\bullet(f)$ est un faisceau d'algèbres
graduées sur le faisceau d'anneaux
$\mathcal{O}_Y=\mathscr{G}r_0(f)$, et les $\mathscr{G}r_k(f)$ des
$\mathcal{O}_Y$-Modules.

Comme pour tout faisceau d'anneaux filtré, on a un homomorphisme canonique
surjectif de $\mathcal{O}_Y$-Algèbres graduées
\[
  \mathscr{S}_{\mathcal{O}_Y}^{\bullet}(\mathscr{G}r_1(f))
  \longrightarrow\mathscr{G}r_\bullet(f)
  \tag{16.1.2.2}
  \label{IV.16.1.2.2}
\]
coïncidant en degrés 0 et 1 avec les homomorphismes identiques.

\begin{examples}[16.1.3]
\label{IV.16.1.3}
\begin{enumerate}
\item[(i)] Supposons que $X$ soit un espace annelé en anneaux locaux, que $Y$ soit
réduit à un seul point $y$ (muni d'un anneau $\mathcal{O}_y$) et que, si
$x=\psi(y)$, $\theta^\#:\mathcal{O}_x\to\mathcal{O}_y$ soit un homomorphisme
surjectif d'anneaux ayant pour noyau l'idéal maximal $\mathfrak{m}_x$ de
$\mathcal{O}_x$. Alors les $\mathcal{O}_{Y^{(n)}}$ s'identifient aux anneaux
$\mathcal{O}_x/\mathfrak{m}_x^{n+1}$, et $\mathscr{G}r_\bullet(f)$ à l'anneau
gradué associé à l'anneau local $\mathcal{O}_x$ muni de sa filtration
$\mathfrak{m}_x$-préadique.
\item[(ii)] Supposons que $Y$ soit une partie fermée d'un sous-espace ouvert $U$
de $X$ et que $\mathcal{O}_Y$ soit induit sur $Y$ par un faisceau quotient
$\mathcal{O}_U/\mathscr{I}$, où $\mathscr{I}$ est un Idéal de $\mathcal{O}_U$
tel que $\mathscr{I}_x=\mathcal{O}_x$ pour tout $x\notin Y$; si $X$ est un
espace annelé en anneaux locaux, nous supposerons de plus que
$\mathscr{I}_x\neq\mathcal{O}_x$ pour $x\in Y$, de sorte que
$(Y,\mathcal{O}_Y)$ est encore un espace annelé en anneaux locaux.

Soit $\psi_0:Y\to U$ l'injection canonique, et notons
$\theta_0:\mathcal{O}_U\to(\psi_0)_*(\mathcal{O}_Y)$ l'homomorphisme tel que
$\theta_0^\#$ soit l'homomorphisme canonique
$\psi_0^*(\mathcal{O}_U)=\mathcal{O}_U|Y\to
(\mathcal{O}_U/\mathscr{I})|Y$, de sorte que
$j_0=(\psi_0,\theta_0):Y\to U$ est un morphisme d'espaces annelés (et d'espaces
annelés en anneaux locaux si $X$ est un espace annelé en anneaux locaux); si
$i:U\to X$ est l'injection canonique (morphisme d'espaces annelés),
$j=i\circ j_0$ est le morphisme $(\psi,\theta)$ de $Y$ dans $X$, où
$\psi:Y\to X$ est l'injection canonique, et
$\theta:\mathcal{O}_X\to\psi_*(\mathcal{O}_Y)$ est l'homomorphisme tel que
$\theta^\#=\theta_0^\#$. Comme $\theta^\#$ est surjectif, on peut appliquer les
définitions précédentes; $\mathcal{O}_{Y^{(n)}}$ est égal à
$\psi_0^*(\mathcal{O}_U/\mathscr{I}^{n+1})$, et l'on a
$(\psi_0)_*(\mathcal{O}_{Y^{(n)}})=\mathcal{O}_U/\mathscr{I}^{n+1}$, et
$\mathscr{G}r_n(j)=\mathscr{G}r_n(j_0)=
\psi_0^*(\mathscr{I}^n/\mathscr{I}^{n+1})=
j_0^*(\mathscr{I}^n/\mathscr{I}^{n+1})$.
\end{enumerate}
\end{examples}

\begin{env}[16.1.4]
\label{IV.16.1.4}
L'exemple (16.1.3, (ii)) montre qu'en général les
$\mathcal{O}_{Y^{(n)}}$ ne sont pas munis canoniquement d'une structure de
$\mathcal{O}_Y$-Module, ni a fortiori d'une structure de
$\mathcal{O}_Y$-Algèbre. La donnée d'une telle structure équivaut à celle d'un
homomorphisme de faisceaux d'anneaux
$\lambda_n:\mathcal{O}_Y\to\mathcal{O}_{Y^{(n)}}$, inverse à droite de
l'homomorphisme d'augmentation $\varphi_{0n}$; cela revient aussi à la donnée
d'un morphisme d'espaces annelés
$(1_Y,\lambda_n):Y^{(n)}\to Y$, inverse à gauche du morphisme canonique
$(1_Y,\varphi_{0n}):Y\to Y^{(n)}$.
\end{env}

\begin{proposition}[16.1.5]
\EGAReaderTarget{ega4.c04.prop.16-1-5}{IV.16.1.5}
Soit $f=(\psi,\theta):Y\to X$ une immersion de préschémas. Alors:
\begin{enumerate}
\item[(i)] $\mathscr{G}r_\bullet(f)$ est une $\mathcal{O}_Y$-Algèbre graduée
quasi-cohérente.
\oldpage[IV-4]{7}
\item[(ii)] Les $Y^{(n)}$ sont des préschémas, canoniquement isomorphes à des
sous-préschémas de $X$.
\item[(iii)] Tout homomorphisme de faisceaux d'anneaux
$\lambda_n:\mathcal{O}_Y\to\mathcal{O}_{Y^{(n)}}$, inverse à droite de
l'homomorphisme d'augmentation $\varphi_{0n}$, fait de
$\mathcal{O}_{Y^{(n)}}$ et des $\mathcal{O}_{Y^{(k)}}$ pour $k\leq n$, des
$\mathcal{O}_Y$-Algèbres quasi-cohérentes; les structures de
$\mathcal{O}_Y$-Module déduites des structures précédentes sur les
$\mathscr{G}r_k(f)$ pour $k\leq n$ coïncident avec celles définies dans
(16.1.2).
\end{enumerate}
\end{proposition}

\noindent(i) La question étant locale sur $X$ et sur $Y$, on peut se borner au cas
où $Y$ est un sous-préschéma fermé de $X$ défini par un Idéal quasi-cohérent
$\mathscr{I}$ de $\mathcal{O}_X$; comme $\mathcal{O}_Y$ est la restriction à
$Y$ de $\mathcal{O}_X/\mathscr{I}$, l'assertion (i) est évidente, et $Y^{(n)}$
est le sous-préschéma fermé de $X$ défini par l'Idéal quasi-cohérent
$\mathscr{I}^{n+1}$ de $\mathcal{O}_X$. Enfin, pour prouver (iii), notons que la
donnée de $\lambda_n$ fait de l'Idéal
$\mathscr{I}/\mathscr{I}^{n+1}$ de l'augmentation $\varphi_{0n}$ et de ses
quotients $\mathscr{I}/\mathscr{I}^{k+1}$ ($1\leq k\leq n$) des
$\mathcal{O}_Y$-Modules, et il suffit de prouver par récurrence sur $k$ que les
$\mathscr{I}/\mathscr{I}^{k+1}$ sont des $\mathcal{O}_Y$-Modules
quasi-cohérents et que la structure de $\mathcal{O}_Y$-Module quotient qu'on en
déduit sur $\mathscr{I}^k/\mathscr{I}^{k+1}$ est la même que celle définie en
(16.1.2). La seconde assertion est immédiate,
$\mathscr{I}^k/\mathscr{I}^{k+1}$ étant annulé par
$\mathscr{I}/\mathscr{I}^{n+1}$; la première résulte, par récurrence sur $k$,
de ce qu'elle est triviale pour $k=1$ et de ce que
$\mathscr{I}/\mathscr{I}^{k+1}$ est une extension de
$\mathscr{I}/\mathscr{I}^k$ par
$\mathscr{I}^k/\mathscr{I}^{k+1}$ (III, 1.4.17).

\begin{corollary}[16.1.6]
\label{IV.16.1.6}
Sous les hypothèses générales de (16.1.5), si l'immersion $f$ est localement de
présentation finie, les $\mathscr{G}r_n(f)$ sont des $\mathcal{O}_Y$-Modules
quasi-cohérents de type fini.
\end{corollary}

En effet, avec les notations de la démonstration de (16.1.5), $\mathscr{I}$ est un
Idéal de type fini de $\mathcal{O}_X$ (1.4.7), donc les
$\mathscr{I}^n/\mathscr{I}^{n+1}$ sont des $\mathcal{O}_Y$-Modules de type fini,
d'où la conclusion.

\begin{corollary}[16.1.7]
\label{IV.16.1.7}
Sous les hypothèses générales de (16.1.5), soit $g:X\to Y$ un morphisme de
préschémas inverse à gauche de $f$. Alors, pour tout $n$, le morphisme composé
$(1,\lambda_n):Y^{(n)}\xrightarrow{h_n}X\xrightarrow{g}Y$ définit un
homomorphisme de faisceaux d'anneaux
$\lambda_n:\mathcal{O}_Y\to\mathcal{O}_{Y^{(n)}}$ inverse à droite de
l'augmentation $\varphi_{0n}$, faisant de $\mathcal{O}_{Y^{(n)}}$ une
$\mathcal{O}_Y$-Algèbre quasi-cohérente; pour ces homomorphismes, les
homomorphismes de transition
$\varphi_{nm}:\mathcal{O}_{Y^{(m)}}\to\mathcal{O}_{Y^{(n)}}$ ($n\leq m$) sont
des homomorphismes de $\mathcal{O}_Y$-Algèbres. En outre, si $g$ est localement
de type fini, les $\mathcal{O}_{Y^{(n)}}$ sont des $\mathcal{O}_Y$-Modules
quasi-cohérents de type fini.
\end{corollary}

La première assertion résulte aussitôt des définitions et de (16.1.5). D'autre
part, si $g$ est localement de type fini, alors $f$ est localement de
présentation finie (1.4.3, (v)); les $\mathscr{G}r_n(f)$ étant alors des
$\mathcal{O}_Y$-Modules quasi-cohérents de type fini par (16.1.6), il en est de
même des $\mathcal{O}_Y$-Modules
$\mathscr{I}/\mathscr{I}^{n+1}$, qui sont des extensions d'un nombre fini de
$\mathscr{G}r_k(f)$ (III, 1.4.17).

\begin{proposition}[16.1.8]
\label{IV.16.1.8}
Soient $X$ un préschéma localement noethérien, $j:Y\to X$ une immersion. Alors
les $Y^{(n)}$ sont des préschémas localement noethériens, les
$\mathscr{G}r_n(j)$ sont des $\mathcal{O}_Y$-Modules cohérents, et
$\mathscr{G}r_\bullet(j)$ est un faisceau d'anneaux cohérent sur l'espace $Y$.
\end{proposition}

Tout étant local sur $X$ et $Y$, on est ramené au cas où $X$ est affine et $j$
une immersion fermée, et alors toutes les assertions sont évidentes sauf la
dernière, qui résulte de ce que si $A$ est un anneau noethérien et
$\mathfrak{I}$ un idéal de $A$, $\operatorname{gr}_{\mathfrak{I}}^\bullet(A)$
est un anneau noethérien, compte tenu de l'exactitude du foncteur $\psi^*$ et de
(0\textsubscript{I}, 5.3.7).

\begin{proposition}[16.1.9]
\label{IV.16.1.9}
\oldpage[IV-4]{8}
Soient $X$ un préschéma, $j:Y\to X$ une immersion localement de présentation
finie, $y$ un point de $Y$. Les conditions suivantes sont équivalentes:
\begin{enumerate}
\item[(a)] Il existe un voisinage ouvert $U$ de $y$ dans $Y$ tel que $j|U$ soit
un homéomorphisme de $U$ sur un ouvert de $X$.
\item[(b)] Il existe un entier $n>0$ tel que l'homomorphisme canonique
\[
  (\varphi_{n,n-1})_y:
  \mathcal{O}_{Y^{(n)},y}\longrightarrow\mathcal{O}_{Y^{(n-1)},y}
\]
soit bijectif.
\item[(c)] Il existe un entier $n>0$ tel que
$(\mathscr{G}r_n(j))_y=0$.
\end{enumerate}
De plus, si l'entier $n$ vérifie b) ou c), il existe un voisinage $V$ de $y$
dans $Y$ tel que $\mathscr{G}r_m(j)|V=0$ pour $m\geq n$ et que
$\varphi_{nm}|V:\mathcal{O}_{Y^{(m)}}|V\to\mathcal{O}_{Y^{(n)}}|V$ soit
bijectif pour $m\geq n$.
\end{proposition}

La question étant locale sur $Y$, on peut se borner au cas où $j$ est une
immersion fermée, $Y$ étant défini par un Idéal quasi-cohérent de type fini
$\mathscr{I}$ de $\mathcal{O}_X$. L'équivalence de b) et c), pour un $n$ donné,
est alors immédiate; en outre, comme
$\mathscr{I}^n/\mathscr{I}^{n+1}$ est un $\mathcal{O}_X$-Module de type fini,
il existe un voisinage ouvert $U$ de $y$ dans $X$ tel que
$\mathscr{I}^n|U=\mathscr{I}^{n+1}|U$ (0\textsubscript{I}, 5.2.2), donc aussi
$\mathscr{I}^n|U=\mathscr{I}^m|U$ pour $m\geq n$, ce qui prouve les dernières
assertions. Pour prouver que a) entraîne b), on peut se borner au cas où
l'espace sous-jacent à $Y$ est égal à l'espace sous-jacent à $X$, et où
$\mathscr{I}$ est engendré par un nombre fini de ses sections au-dessus de $X$:
comme $\mathscr{I}$ est alors contenu dans le Nilradical $\mathscr{N}$ de
$\mathcal{O}_X$ (I, 5.1.2), il est nilpotent, ce qui prouve b). Enfin, pour
prouver que b) entraîne a), on peut aussi se borner au cas où
$\mathscr{I}^n=\mathscr{I}^{n+1}$; alors, pour tout $y\in Y$, comme
$\mathscr{I}_y\subset\mathfrak{m}_y$, idéal maximal de
$\mathcal{O}_{X,y}$, on a nécessairement $\mathscr{I}_y^n=0$ en vertu du
lemme de Nakayama, puisque $\mathscr{I}_y$ est un idéal de type fini.
L'ensemble des $x\in X$ tels que $\mathscr{I}_x^n=0$ est donc un ouvert $U$ de
$X$ contenant $Y$ (0\textsubscript{I}, 5.2.2); comme d'autre part
$\mathscr{I}_x\neq 0$ pour $x\notin Y$, on a nécessairement $U=Y$.

\begin{corollary}[16.1.10]
\label{IV.16.1.10}
Pour que la restriction de l'immersion $j$ à un voisinage de $y$ dans $Y$ soit
une immersion ouverte (autrement dit pour que $j$ soit un isomorphisme local au
point $y$), il faut et il suffit que
$(\mathscr{G}r_1(j))_y=(\mathscr{N}_{Y/X})_y=0$.
\end{corollary}

La condition est évidemment nécessaire, et le raisonnement précédent, appliqué
pour $n=1$, prouve qu'elle est suffisante.

\begin{remarks}[16.1.11]
\EGAReaderTarget{ega4.c04.rem.16-1-11}{IV.16.1.11}
\begin{enumerate}
\item[(i)] Sous les conditions de la définition (16.1.1), la limite projective
du système projectif
$(\mathcal{O}_{Y^{(n)}},\varphi_{nm})$ de faisceaux d'anneaux sur $Y$ est
appelé l'invariant normal d'ordre infini de $f$, et parfois noté
$\mathcal{O}_{Y^{(\infty)}}$. Lorsque $X$ est un préschéma localement
noethérien, $j:Y\to X$ une immersion fermée, $Y$ étant donc un sous-préschéma
fermé de $X$ défini par un Idéal cohérent $\mathscr{I}$,
$\mathcal{O}_{Y^{(\infty)}}$ n'est autre que le complété formel de
$\mathcal{O}_X$ le long de $Y$ (I, 10.8.4), et
$Y^{(\infty)}=(Y,\mathcal{O}_{Y^{(\infty)}})$ le préschéma formel complété de
$X$ le long de $Y$ (I, 10.8.5). Dans tous les cas, on pourra dire que
$Y^{(\infty)}$ est le voisinage formel de $Y$ dans $X$ (pour le morphisme
$f$). Dans le cas particulier qu'on vient d'envisager, c'est donc le préschéma
formel limite inductive des voisinages infinitésimaux d'ordre $n$.
\item[(ii)] On notera que pour un morphisme de préschémas
$f=(\psi,\theta):Y\to X$, il peut se faire que l'homomorphisme
$\theta^\#:\psi^*(\mathcal{O}_X)\to\mathcal{O}_Y$ soit surjectif sans que $f$
soit une immersion
\oldpage[IV-4]{9}
locale, et sans que $f$ soit injectif. On en a un exemple en prenant pour $Y$
une somme de préschémas $Y_\lambda$ tous isomorphes à
$\operatorname{Spec}(\mathcal{O}_x)$, où $x\in X$, et pour $f$ le morphisme
égal au morphisme canonique dans chacun des $Y_\lambda$.
\end{enumerate}
\end{remarks}

\subsection{Propriétés fonctorielles des invariants normaux d'une immersion}
\EGAReaderTarget{ega4.c04.subsec.16-2}{IV.16.2}

\begin{env}[16.2.1]
\label{IV.16.2.1}
Soient $f=(\psi,\theta):Y\to X$,
$f'=(\psi',\theta'):Y'\to X'$ deux morphismes d'espaces annelés tels que les
homomorphismes $\theta^\#$ et $\theta'^\#$ soient surjectifs; considérons un
diagramme commutatif de morphismes d'espaces annelés
\[
  \xymatrix{
    Y \ar[r]^f & X \\
    Y' \ar[r]_{f'} \ar[u]^u & X' \ar[u]_v
  }
  \tag{16.2.1.1}
  \label{IV.16.2.1.1}
\]
Posons $u=(\rho,\lambda)$, $v=(\sigma,\mu)$. On a
$\rho^*(\psi^*(\mathcal{O}_X))=
\psi'^*(\sigma^*(\mathcal{O}_X))$ et on a par suite un diagramme commutatif
d'homomorphismes de faisceaux d'anneaux sur $Y'$
\[
  \xymatrix{
    \rho^*(\psi^*(\mathcal{O}_X))
      =\psi'^*(\sigma^*(\mathcal{O}_X))
      \ar[r]^-{\psi'^*(\mu^\#)}
      \ar[d]_{\rho^*(\theta^\#)}
      & \psi'^*(\mathcal{O}_{X'}) \ar[d]^{\theta'^\#} \\
    \rho^*(\mathcal{O}_Y) \ar[r]_-{\lambda^\#}
      & \mathcal{O}_{Y'}
  }
\]
d'où l'on conclut, si $\mathscr{I}$ et $\mathscr{I}'$ sont les noyaux de
$\theta^\#$ et $\theta'^\#$, que l'on a
$\psi'^*(\mu^\#)(\rho^*(\mathscr{I}))\subset\mathscr{I}'$, compte tenu de
l'exactitude du foncteur $\rho^*$. On en déduit aussitôt que pour tout entier
$n$,
$\psi'^*(\mu^\#)(\rho^*(\mathscr{I}^n))\subset\mathscr{I}'^n$, ce qui montre
que $\psi'^*(\mu^\#)$ définit, par passage aux quotients, un homomorphisme de
faisceaux d'anneaux
\[
  \nu_n:
  \rho^*(\psi^*(\mathcal{O}_X)/\mathscr{I}^{n+1})
  \longrightarrow
  \psi'^*(\mathcal{O}_{X'})/\mathscr{I}'^{n+1}
  \tag{16.2.1.2}
  \label{IV.16.2.1.2}
\]
et par suite un morphisme d'espaces annelés
$w_n=(\rho,\nu_n):Y'^{(n)}\to Y^{(n)}$ (qui, pour $n=0$, n'est autre que
$u$). Il résulte aussitôt de cette définition que les diagrammes
\[
  \xymatrix@R=1pc{
    Y^{(n)} \ar[r]^-{h_{mn}} & Y^{(m)} \ar[r]^-{h_m} & X \\
    & & & (n\leq m) \\
    Y'^{(n)} \ar[r]_-{h'_{mn}} \ar[uu]^-{w_n}
      & Y'^{(m)} \ar[r]_-{h'_m} \ar[uu]^-{w_m}
      & X' \ar[uu]_-v
  }
\]
(où les flèches horizontales sont les morphismes canoniques (16.1.2)) sont
commutatifs.

Par passage aux quotients à partir des homomorphismes (16.2.1.2), et en tenant
\oldpage[IV-4]{10}
compte de l'exactitude du foncteur $\rho^*$, on obtient un di-homomorphisme
d'Algèbres graduées (relatif à l'homomorphisme
$\lambda^\#:\rho^*(\mathcal{O}_Y)\to\mathcal{O}_{Y'}$)
\[
  \operatorname{gr}(u):
  \rho^*(\mathscr{G}r_\bullet(f))
  \longrightarrow\mathscr{G}r_\bullet(f')
  \tag{16.2.1.3}
  \label{IV.16.2.1.3}
\]
(ou, si on veut, un $\rho$-morphisme (0\textsubscript{I}, 3.5.1)
$\mathscr{G}r_\bullet(f)\to\mathscr{G}r_\bullet(f')$), et en particulier un
di-homomorphisme des faisceaux conormaux
\[
  \operatorname{gr}_1(u):
  \rho^*(\mathscr{G}r_1(f))\longrightarrow\mathscr{G}r_1(f').
\]
Il est immédiat en outre que ces homomorphismes donnent lieu à un diagramme
commutatif
\[
  \xymatrix{
    \rho^*(\mathscr{S}_{\mathcal{O}_Y}^{\bullet}
      (\mathscr{G}r_1(f)))
      \ar[r] \ar[d]_-{\mathscr{S}(\operatorname{gr}_1(u))}
      & \rho^*(\mathscr{G}r_\bullet(f))
        \ar[d]^-{\operatorname{gr}(u)} \\
    \mathscr{S}_{\mathcal{O}_{Y'}}^{\bullet}
      (\mathscr{G}r_1(f'))
      \ar[r]
      & \mathscr{G}r_\bullet(f')
  }
  \tag{16.2.1.4}
  \label{IV.16.2.1.4}
\]
où les flèches horizontales sont les homomorphismes canoniques (16.1.2.2).

Enfin, si l'on a un diagramme commutatif de morphismes d'espaces annelés
\[
  \xymatrix{
    Y \ar[r]^f & X \\
    Y' \ar[r]_{f'} \ar[u]^u & X' \ar[u]_v \\
    Y'' \ar[r]_{f''} \ar[u]^{u'} & X'' \ar[u]_{v'}
  }
\]
où $f''=(\psi'',\theta'')$ est tel que $\theta''^\#$ soit surjectif, et si
$w'_n$ et $w''_n$ sont définis à partir de $u'$, $v'$ d'une part, et de
$u''=u\circ u'$, $v''=v\circ v'$ de l'autre, alors on a
$w''_n=w_n\circ w'_n$, comme il résulte aussitôt des définitions et de
(0\textsubscript{I}, 3.5.5); on a de même
$\operatorname{gr}(u'')=
\operatorname{gr}(u')\circ\rho'^*(\operatorname{gr}(u))$ si
$u'=(\rho',\lambda')$. On peut donc dire que les $Y^{(n)}$ et
$\mathscr{G}r_\bullet(f)$ dépendent fonctoriellement de $f$.
\end{env}

\begin{proposition}[16.2.2]
\label{IV.16.2.2}
Avec les notations et hypothèses de (16.2.1), supposons de plus que $f$, $f'$,
$u$ et $v$ soient des morphismes de préschémas. Alors:
\begin{enumerate}
\item[(i)] Les morphismes $w_n:Y'^{(n)}\to Y^{(n)}$ sont des morphismes de
préschémas.
\item[(ii)] Si $Y'=Y\times_X X'$, $u$ et $f'$ étant les projections
canoniques, et si $f$ est une immersion ou si $v$ est plat, on a
$Y'^{(n)}=Y^{(n)}\times_X X'$.
\item[(iii)] Si $Y'=Y\times_X X'$ et si $v$ est plat (resp. si $f$ est une
immersion), l'homomorphisme
\[
  \operatorname{Gr}(u)=\operatorname{gr}(u)\otimes 1:
  \mathscr{G}r_\bullet(f)\otimes_{\mathcal{O}_Y}\mathcal{O}_{Y'}
  \longrightarrow\mathscr{G}r_\bullet(f')
\]
est bijectif (resp. surjectif).
\end{enumerate}
\end{proposition}

\noindent(i) Les hypothèses entraînent aussitôt que pour tout $y'\in Y'$,
$\rho_{y'}^*(\theta_{\psi'(y')}^\#)$ est un homomorphisme local (I, 1.6.2),
donc $w_n$ est un morphisme de préschémas (I, 2.2.1).

\oldpage[IV-4]{11}
\noindent(ii) et (iii) Si $f$ est une immersion, on peut se borner au cas où
$f$ est une immersion fermée, $Y$ étant donc défini par l'Idéal
quasi-cohérent $\mathscr{I}$ de $\mathcal{O}_X$, et $Y^{(n)}$ par l'Idéal
$\mathscr{I}^{n+1}$; les assertions résultent alors de (I, 4.4.5).

Supposons en second lieu que $v$ soit plat; on peut se borner au cas où
$X=\operatorname{Spec}(A)$, $Y=\operatorname{Spec}(B)$,
$X'=\operatorname{Spec}(A')$ sont affines, $A'$ étant un $A$-module plat;
alors $Y'=\operatorname{Spec}(B')$ avec $B'=B\otimes_A A'$; en outre, si
$\mathfrak{I}$ est le noyau de l'homomorphisme $A\to B$, le noyau
$\mathfrak{I}'$ de $A'\to B'$ s'identifie à $\mathfrak{I}\otimes_A A'$ par
platitude, et
\[
  \mathfrak{I}'^n/\mathfrak{I}'^{n+1}
  =(\mathfrak{I}^n/\mathfrak{I}^{n+1})\otimes_A A'.
\]
On en déduit aussitôt, compte tenu de (0\textsubscript{I}, 4.3.3), que le
$\mathcal{O}_{Y'}$-Module
$\mathscr{I}'^n/\mathscr{I}'^{n+1}$ est égal à
\begin{align*}
&\psi'^*\!\left(\sigma^*
  \left((\mathfrak{I}^n/\mathfrak{I}^{n+1})^{\sim}\right)\right)
  \otimes_{\sigma^*(\mathcal{O}_X)}\mathcal{O}_{X'} \\
={}&
\psi'^*\!\left(\sigma^*
  \left((\mathfrak{I}^n/\mathfrak{I}^{n+1})^{\sim}\right)\right)
  \otimes_{\psi'^*(\sigma^*(\mathcal{O}_X))}
  \psi'^*(\mathcal{O}_{X'}) \\
={}&
\rho^*(\mathscr{I}^n/\mathscr{I}^{n+1})
  \otimes_{\rho^*(\psi^*(\mathcal{O}_X))}
  \psi'^*(\mathcal{O}_{X'}),
\end{align*}
et comme en particulier pour $n=0$, on a
\[
  \mathcal{O}_{Y'}=
  \rho^*(\mathcal{O}_Y)
  \otimes_{\rho^*(\psi^*(\mathcal{O}_X))}
  \psi'^*(\mathcal{O}_{X'}),
\]
on obtient un isomorphisme canonique de
$\mathscr{I}'^n/\mathscr{I}'^{n+1}$ sur
\[
  \rho^*(\mathscr{I}^n/\mathscr{I}^{n+1})
  \otimes_{\rho^*(\mathcal{O}_Y)}\mathcal{O}_{Y'}
  =
  (\mathscr{I}^n/\mathscr{I}^{n+1})
  \otimes_{\mathcal{O}_Y}\mathcal{O}_{Y'},
\]
ce qui prouve (iii). Posons maintenant
$C_n=\Gamma(Y,\mathcal{O}_{Y^{(n)}})$,
$C'_n=\Gamma(Y',\mathcal{O}_{Y'^{(n)}})$. Puisque $Y^{(n)}$ et
$Y'^{(n)}$ sont des schémas affines (16.1.5), le noyau
$\mathfrak{K}_n$ (resp. $\mathfrak{K}'_n$) de l'homomorphisme
$C_n\to C_{n-1}$ (resp. $C'_n\to C'_{n-1}$) est
$\Gamma(Y,\mathscr{I}^n/\mathscr{I}^{n+1})$ (resp.
$\Gamma(Y',\mathscr{I}'^n/\mathscr{I}'^{n+1})$); on déduit donc de ce qui
précède que $\mathfrak{K}'_n=\mathfrak{K}_n\otimes_A A'$. Or, on a un
diagramme commutatif
\[
  \xymatrix{
    0 \ar[r] &
    \mathfrak{K}_n\otimes_A A' \ar[r]\ar[d]^r &
    C_n\otimes_A A' \ar[r]\ar[d]^{s_n} &
    C_{n-1}\otimes_A A' \ar[r]\ar[d]^{s_{n-1}} & 0 \\
    0 \ar[r] &
    \mathfrak{K}'_n \ar[r] &
    C'_n \ar[r] &
    C'_{n-1} \ar[r] & 0
  }
\]
où la flèche verticale de gauche est bijective et les deux lignes sont exactes
($A'$ étant un $A$-module plat). On en déduit par récurrence que $s_n$ est
bijectif pour tout $n$, car c'est vrai par hypothèse pour $n=0$, et se déduit
par application du lemme des cinq pour $n$ quelconque. Cela prouve la seconde
assertion de (ii).

\begin{corollary}[16.2.3]
\label{IV.16.2.3}
Soient $g:X\to Y$, $u:Y'\to Y$ deux morphismes de préschémas,
$X'=X\times_Y Y'$, $g':X'\to Y'$ et $v:X'\to X$ les projections canoniques.
Soient $f:Y\to X$ une $Y$-section de $X$ (donc une immersion),
$f'=f_{(Y')}:Y'\to X'$ la $Y'$-section de $X'$ déduite de $f$ par le
changement de base $u$. Alors:
\begin{enumerate}
\item[(i)] Le morphisme
$w_n:Y_{f'}'^{(n)}\to Y_f^{(n)}$ correspondant à $f$, $f'$, $u$, $v$
(16.2.1) et le morphisme canonique
$h'_n:Y_{f'}'^{(n)}\to X'$ identifient $Y_{f'}'^{(n)}$ au produit
$Y_f^{(n)}\times_X X'$.
\item[(ii)] Si l'on munit $\mathcal{O}_{Y_f^{(n)}}$ (resp.
$\mathcal{O}_{Y_{f'}'^{(n)}}$) de la structure de $\mathcal{O}_Y$-Algèbre
définie par $g$ (resp. de la structure de $\mathcal{O}_{Y'}$-Algèbre définie
par $g'$) (16.1.6), l'homomorphisme de $\mathcal{O}_{Y'}$-Algèbres
\[
  \rho^*(\mathcal{O}_{Y_f^{(n)}})
  \otimes_{\mathcal{O}_Y}\mathcal{O}_{Y'}
  \longrightarrow
  \mathcal{O}_{Y_{f'}'^{(n)}}
  \tag{16.2.3.1}
  \label{IV.16.2.3.1}
\]
\oldpage[IV-4]{12}
déduit de l'homomorphisme $\nu_n$ (16.2.1.2) est bijectif. En outre,
l'homomorphisme de $\mathcal{O}_{Y'}$-Modules
\[
  \operatorname{Gr}_1(u):
  \mathscr{G}r_1(f)\otimes_{\mathcal{O}_Y}\mathcal{O}_{Y'}
  \longrightarrow\mathscr{G}r_1(f')
  \tag{16.2.3.2}
  \label{IV.16.2.3.2}
\]
est bijectif.
\end{enumerate}
\end{corollary}

\noindent(i) Notons d'abord que les morphismes $f':Y'\to X'$ et $u:Y'\to Y$
identifient $Y'$ au produit $Y\times_X X'$ (pour les morphismes structuraux
$f:Y\to X$ et $v:X'\to X$) (14.5.12.1). La conclusion de (i) résulte alors de
(16.2.2, (ii)), le morphisme $g$ étant une immersion.

\noindent(ii) Le diagramme commutatif
\[
  \xymatrix{
    Y_f^{(n)} \ar[d]^{h_n}
      & Y_{f'}'^{(n)} \ar[d]^-{h'_n} \ar[l]^{w_n} \\
    X \ar[d]^g
      & X' \ar[d]^{g'} \ar[l]^v \\
    Y
      & Y' \ar[l]^u
  }
\]
identifie $Y_{f'}'^{(n)}$ au produit $Y_f^{(n)}\times_X X'$ et $X'$ au produit
$X\times_Y Y'$, donc (I, 3.3.9) il identifie (pour les morphismes
$g'\circ h'_n$ et $w_n$) $Y_{f'}'^{(n)}$ au produit
$Y_f^{(n)}\times_Y Y'$. Comme $Y_f^{(n)}$ (resp. $Y_{f'}'^{(n)}$) est le
préschéma affine au-dessus de $Y$ (resp. $Y'$) associé à la
$\mathcal{O}_Y$-Algèbre $\mathcal{O}_{Y_f^{(n)}}$ (resp. à la
$\mathcal{O}_{Y'}$-Algèbre $\mathcal{O}_{Y_{f'}'^{(n)}}$), le fait que
l'homomorphisme canonique (16.2.3.1) soit bijectif résulte de (II, 1.5.2).
Enfin, l'homomorphisme canonique (16.2.3.1) est compatible avec les
augmentations
$\mathcal{O}_{Y_f^{(n)}}\to\mathcal{O}_Y$ et
$\mathcal{O}_{Y_{f'}'^{(n)}}\to\mathcal{O}_{Y'}$; comme
$\mathcal{O}_{Y_f^{(n)}}$ est somme directe (en tant que
$\mathcal{O}_Y$-Module) de $\mathcal{O}_Y$ et de l'Idéal d'augmentation
$\mathscr{I}/\mathscr{I}^{n+1}$, on voit donc que l'homomorphisme canonique
(16.2.3.1), restreint à
$(\mathscr{I}/\mathscr{I}^{n+1})\otimes_{\mathcal{O}_Y}\mathcal{O}_{Y'}$,
est une bijection de ce dernier sur
$\mathscr{I}'/\mathscr{I}'^{n+1}$. Pour $n=1$, cela montre que
$\operatorname{Gr}_1(u)$ est bijectif.

On notera que, sous les hypothèses de (16.2.3), les homomorphismes
$\operatorname{Gr}_n(u)$ sont surjectifs en vertu de ce qui précède, mais ne
sont pas bijectifs en général pour $n\geq 2$. Toutefois:

\begin{corollary}[16.2.4]
\label{IV.16.2.4}
Sous les hypothèses de (16.2.3), supposons que $u:Y'\to Y$ soit un morphisme
plat (resp. que les $\mathscr{G}r_n(f)$ soient des $\mathcal{O}_Y$-Modules
plats pour $n\leq m$). Alors l'homomorphisme
\[
  \operatorname{Gr}_n(u):
  \mathscr{G}r_n(f)\otimes_{\mathcal{O}_Y}\mathcal{O}_{Y'}
  \longrightarrow\mathscr{G}r_n(f')
\]
est bijectif pour tout $n$ (resp. pour $n\leq m$).
\end{corollary}

Si $u$ est plat, il en est de même de $v:X'\to X$ qui s'en déduit par
changement de base, et on sait déjà dans ce cas que $\operatorname{Gr}(u)$ est
bijectif (16.2.2, (iii)). Si les $\mathscr{G}r_n(f)$ sont plats pour
$n\leq m$, on voit d'abord par récurrence sur $n$ qu'il en est de même des
$\mathscr{I}/\mathscr{I}^{n+1}$ pour $n\leq m$, en vertu des suites exactes
\[
  0\longrightarrow\mathscr{I}^n/\mathscr{I}^{n+1}
  \longrightarrow\mathscr{I}/\mathscr{I}^{n+1}
  \longrightarrow\mathscr{I}/\mathscr{I}^n
  \longrightarrow 0
\]
\oldpage[IV-4]{13}
(0\textsubscript{I}, 6.1.2); en outre, on a alors des diagrammes commutatifs
\[
  \xymatrix{
    0 \ar[r] &
    (\mathscr{I}^n/\mathscr{I}^{n+1})
      \otimes_{\mathcal{O}_Y}\mathcal{O}_{Y'}
      \ar[r]\ar[d] &
    (\mathscr{I}/\mathscr{I}^{n+1})
      \otimes_{\mathcal{O}_Y}\mathcal{O}_{Y'}
      \ar[r]\ar[d] &
    (\mathscr{I}/\mathscr{I}^n)
      \otimes_{\mathcal{O}_Y}\mathcal{O}_{Y'}
      \ar[r]\ar[d] & 0 \\
    0 \ar[r] &
    \mathscr{I}'^n/\mathscr{I}'^{n+1}
      \ar[r] &
    \mathscr{I}'/\mathscr{I}'^{n+1}
      \ar[r] &
    \mathscr{I}'/\mathscr{I}'^n
      \ar[r] & 0
  }
\]
dans lesquels les lignes sont exactes (la première par platitude
(0\textsubscript{I}, 6.1.2)) et les deux dernières flèches verticales sont
bijectives en vertu de (16.2.3, (ii)); d'où la conclusion.

\begin{remarks}[16.2.5]
\label{IV.16.2.5}
\begin{enumerate}
\item[(i)] Le raisonnement de (16.2.2, (i)) s'applique encore lorsque dans
(16.2.1.1) il s'agit de morphismes d'espaces annelés en anneaux locaux
($\mathbf{Err}_{\mathrm{II}}$, (1.8.2)).
\item[(ii)] Dans (16.2.2, (ii)), la conclusion n'est plus nécessairement valable
lorsqu'on suppose seulement que $v$ et $f$ sont des morphismes de préschémas
($f$ vérifiant la condition de (16.1.1)). Par exemple (avec les notations de la
démonstration de (16.2.2, (ii)), il peut se faire que $\mathfrak{I}=0$ mais que
le noyau $\mathfrak{I}'$ de $A'\to B'=B\otimes_A A'$ ne soit pas nul et que
$B'\neq 0$, auquel cas on a $Y^{(n)}=Y$ pour tout $n$, mais
$Y'^{(n)}\neq Y'$. On a un exemple de ce fait en prenant
$A=\mathbb{Z}$, $B=\mathbb{Q}$,
$A'=\prod_{h=1}^{\infty}(\mathbb{Z}/m^h\mathbb{Z})$ où $m>1$.
\end{enumerate}
\end{remarks}

\begin{env}[16.2.6]
\label{IV.16.2.6}
Considérons le cas particulier du diagramme (16.2.1.1) où $X'=X$, $v$ étant
l'identité, $X$ est un préschéma, $Y$ un sous-préschéma de $X$, $Y'$ un
sous-préschéma de $Y$, $f$, $u$ et $f'=f\circ u$ les injections canoniques; le
di-homomorphisme (16.2.1.3) donne, par tensorisation avec
$\mathcal{O}_{Y'}$ au-dessus de $\rho^*(\mathcal{O}_Y)$, un homomorphisme de
$\mathcal{O}_{Y'}$-Algèbres graduées
\[
  u^*(\mathscr{G}r_\bullet(f))
  \longrightarrow\mathscr{G}r_\bullet(f')
  \tag{16.2.6.1}
  \label{IV.16.2.6.1}
\]
D'autre part, on identifie $\mathcal{O}_Y$ à
$\psi^*(\mathcal{O}_X)/\mathscr{I}_f$, $\mathcal{O}_{Y'}$ à
$\rho^*(\mathcal{O}_Y)/\mathscr{I}_u$; comme $\rho^*$ est un foncteur exact,
on a
$\rho^*(\mathcal{O}_Y)=
\rho^*(\psi^*(\mathcal{O}_X))/\rho^*(\mathscr{I}_f)=
\psi'^*(\mathcal{O}_X)/\rho^*(\mathscr{I}_f)$, et comme
$\mathcal{O}_{Y'}$ s'identifie par ailleurs à
$\psi'^*(\mathcal{O}_X)/\mathscr{I}_{f'}$, on voit que l'on a
$\mathscr{I}_u=\mathscr{I}_{f'}/\rho^*(\mathscr{I}_f)$. On en déduit pour
tout entier $n$ un homomorphisme canonique
$\mathscr{I}_{f'}^n/\mathscr{I}_{f'}^{n+1}\to
\mathscr{I}_u^n/\mathscr{I}_u^{n+1}$, d'où un homomorphisme canonique de
$\mathcal{O}_{Y'}$-Algèbres graduées
\[
  \mathscr{G}r_\bullet(f')
  \longrightarrow\mathscr{G}r_\bullet(u).
  \tag{16.2.6.2}
  \label{IV.16.2.6.2}
\]
\end{env}

\begin{proposition}[16.2.7]
\label{IV.16.2.7}
Soient $X$ un préschéma, $Y$ un sous-préschéma de $X$, $Y'$ un
sous-préschéma de $Y$, $j:Y'\to Y$ l'injection canonique. On a alors une suite
exacte de faisceaux conormaux ($\mathcal{O}_{Y'}$-Modules)
\[
  \xymatrix{
    j^*(\mathscr{N}_{Y/X}) \ar[r] &
    \mathscr{N}_{Y'/X} \ar[r] &
    \mathscr{N}_{Y'/Y} \ar[r] & 0
  }
  \tag{16.2.7.1}
  \label{IV.16.2.7.1}
\]
où les flèches sont les composantes de degré 1 des homomorphismes canoniques
(16.2.6.1) et (16.2.6.2).
\end{proposition}

La question étant locale, on peut se borner au cas où
$X=\operatorname{Spec}(A)$, $Y=\operatorname{Spec}(A/\mathfrak{I})$,
$Y'=\operatorname{Spec}(A/\mathfrak{K})$, $\mathfrak{I}$ et $\mathfrak{K}$
étant deux idéaux de $A$ tels que $\mathfrak{I}\subset\mathfrak{K}$; tout
revient à voir alors
\oldpage[IV-4]{14}
que la suite d'homomorphismes canoniques
$\mathfrak{I}/\mathfrak{K}\mathfrak{I}\to
\mathfrak{K}/\mathfrak{K}^2\to
(\mathfrak{K}/\mathfrak{I})/(\mathfrak{K}/\mathfrak{I})^2\to 0$
est exacte, ce qui est immédiat puisque l'image de
$\mathfrak{I}/\mathfrak{K}\mathfrak{I}$ dans
$\mathfrak{K}/\mathfrak{K}^2$ est
$(\mathfrak{I}+\mathfrak{K}^2)/\mathfrak{K}^2$ et que
$(\mathfrak{K}/\mathfrak{I})/(\mathfrak{K}/\mathfrak{I})^2$ s'identifie à
$\mathfrak{K}/(\mathfrak{I}+\mathfrak{K}^2)$.

Il est facile de donner des exemples où la suite (16.2.7.1) prolongée à gauche
par un 0 n'est plus exacte; avec les notations précédentes, il suffit de
prendre $A=k[T]$, $\mathfrak{I}=AT^2$, $\mathfrak{K}=AT$, car on a alors
$(\mathfrak{I}+\mathfrak{K}^2)/\mathfrak{K}^2=0$ et
$\mathfrak{I}/\mathfrak{K}\mathfrak{I}\neq 0$. Voir cependant (16.9.13) et
(19.1.5) pour des cas utiles où la suite prolongée reste encore exacte.

\subsection{Invariants différentiels fondamentaux d'un morphisme de préschémas}
\EGAReaderTarget{ega4.c04.subsec.16-3}{IV.16.3}

\begin{definition}[16.3.1]
\EGAReaderTarget{ega4.c04.def.16-3-1}{IV.16.3.1}
Soient $f:X\to S$ un morphisme de préschémas, $\Delta_f:X\to X\times_S X$
le morphisme diagonal correspondant, qui est une immersion (I, 5.3.9 et
$\mathbf{Err}_{\mathrm{III}}$, 10). On désigne par $\mathscr{P}_f^n$ ou
$\mathscr{P}_{X/S}^n$, et l'on appelle faisceau des parties principales d'ordre $n$
du $S$-préschéma $X$, le faisceau d'anneaux $\mathcal{O}_X$-augmentés, $n$-ème
invariant normal de $\Delta_f$ (16.1.2). On pose
$\mathscr{P}_f^\infty=\mathscr{P}_{X/S}^\infty=
\varprojlim_n\mathscr{P}_{X/S}^n$,
$\mathscr{G}r_n(\mathscr{P}_f)=\mathscr{G}r_n(\mathscr{P}_{X/S})=
\mathscr{G}r_n(\Delta_f)$ (16.1.2); le $\mathcal{O}_X$-Module
$\mathscr{G}r_1(\Delta_f)$, Idéal d'augmentation de $\mathscr{P}_{X/S}^1$, est
aussi noté $\Omega_f^1$ ou $\Omega_{X/S}^1$, et appelé le $\mathcal{O}_X$-Module
des 1-différentielles de $f$, ou de $X$ par rapport à $S$, ou du $S$-préschéma $X$.
\end{definition}

Il résulte de cette définition que $\mathscr{P}_{X/S}^0$ s'identifie
canoniquement à $\mathcal{O}_X$ (16.1.2).

On a (16.1.2.2) un homomorphisme canonique surjectif de
$\mathcal{O}_X$-Algèbres graduées
\[
  \mathscr{S}_{\mathcal{O}_X}^{\bullet}(\Omega_{X/S}^1)
  \longrightarrow\mathscr{G}r_\bullet(\mathscr{P}_{X/S}).
  \tag{16.3.1.1}
  \label{IV.16.3.1.1}
\]
Il résulte aussi de la définition (16.3.1) que pour tout ouvert $U$ de $X$, on a
$\mathscr{P}_{f|U}^n=\mathscr{P}_f^n|U$,
$\mathscr{P}_{f|U}^\infty=\mathscr{P}_f^\infty|U$,
$\mathscr{G}r_n(\mathscr{P}_{f|U})=\mathscr{G}r_n(\mathscr{P}_f)|U$,
$\Omega_{f|U}^1=\Omega_f^1|U$ (autrement dit, les notions introduites sont
locales sur $X$).

\begin{env}[16.3.2]
\label{IV.16.3.2}
Notons $p_1$, $p_2$ les deux projections canoniques du produit $X\times_S X$;
comme $\Delta_f$ est une $X$-section de $X\times_S X$ à la fois pour les
morphismes $p_1$ et $p_2$, chacun de ces morphismes définit, pour tout $n$, un
homomorphisme de faisceaux d'anneaux $\mathcal{O}_X\to\mathscr{P}_{X/S}^n$,
inverse à droite de l'augmentation $\mathscr{P}_{X/S}^n\to\mathcal{O}_X$
(16.1.7); on peut encore dire que l'on définit ainsi sur
$\mathscr{P}_{X/S}^n$ deux structures de $\mathcal{O}_X$-Algèbres augmentées
quasi-cohérentes; les structures correspondantes de $\mathcal{O}_X$-Module sur
$\mathscr{G}r_n(\mathscr{P}_{X/S})$ sont les mêmes. On a de même, par passage à
la limite, deux structures de $\mathcal{O}_X$-Algèbre sur
$\mathscr{P}_{X/S}^\infty$.
\end{env}

\begin{env}[16.3.3]
\EGAReaderTarget{ega4.c04.statement.16-3-3}{IV.16.3.3}
Le morphisme $s=(p_2,p_1)_S:X\times_S X\to X\times_S X$ est un automorphisme
involutif de $X\times_S X$, appelé la symétrie canonique, tel que
\[
  p_1\circ s=p_2,\qquad p_2\circ s=p_1,\qquad s\circ\Delta_f=\Delta_f.
  \tag{16.3.3.1}
  \label{IV.16.3.3.1}
\]
Si l'on pose $s=(\rho,\lambda)$, $p_i=(\pi_i,\mu_i)$ ($i=1,2$),
$\Delta_f=(\delta,\nu)$, $\lambda^\#$ est donc un isomorphisme de
$\rho^*(\pi_1^*(\mathcal{O}_X))$ sur $\pi_2^*(\mathcal{O}_X)$, et
$\delta^*(\lambda^\#)$ laisse invariant $\delta^*(\mathcal{O}_{X\times_S X})$
et le noyau $\mathscr{I}$ de l'homomorphisme
$\nu^\#:\delta^*(\mathcal{O}_{X\times_S X})\to\mathcal{O}_X$. Par suite:
\end{env}

\begin{proposition}[16.3.4]
\EGAReaderTarget{ega4.c04.prop.16-3-4}{IV.16.3.4}
L'homomorphisme $\sigma=\delta^*(\lambda^\#)$ déduit de $s$ (et appelé encore
symétrie canonique) est un automorphisme involutif du système projectif
$(\mathscr{P}_{X/S}^n)$ de faisceaux
\oldpage[IV-4]{15}
d'anneaux $\mathcal{O}_X$-augmentés, et par suite aussi de leur limite projective
$\mathscr{P}_{X/S}^\infty$. Cet automorphisme permute les deux structures de
$\mathcal{O}_X$-Algèbre sur les $\mathscr{P}_{X/S}^n$ et sur
$\mathscr{P}_{X/S}^\infty$.
\end{proposition}

\begin{env}[16.3.5]
\label{IV.16.3.5}
Dans ce qui suit, les deux structures de $\mathcal{O}_X$-Algèbre définies sur les
$\mathscr{P}_{X/S}^n$ et sur $\mathscr{P}_{X/S}^\infty$ joueront des rôles très
différents: nous conviendrons désormais, sauf mention expresse du contraire,
que lorsque $\mathscr{P}_{X/S}^n$ ou $\mathscr{P}_{X/S}^\infty$ sera considéré
comme une $\mathcal{O}_X$-Algèbre, c'est de la structure d'Algèbre définie par
$p_1$ qu'il s'agira.
\end{env}

Pour tout ouvert $U$ de $X$ et toute section $t\in\Gamma(U,\mathcal{O}_X)$, on
notera simplement $t.1$ ou même $t$ l'image de $t$ par l'homomorphisme
structural $\Gamma(U,\mathcal{O}_X)\to\Gamma(U,\mathscr{P}_{X/S}^n)$ (resp.
$\Gamma(U,\mathcal{O}_X)\to\Gamma(U,\mathscr{P}_{X/S}^\infty)$)
(c'est-à-dire l'homomorphisme correspondant à $p_1$).

\begin{definition}[16.3.6]
\EGAReaderTarget{ega4.c04.def.16-3-6}{IV.16.3.6}
On désigne par $d_f^n$, ou $d_{X/S}^n$ (resp. $d_f^\infty$, ou
$d_{X/S}^\infty$), ou simplement $d^n$ (resp. $d^\infty$), l'homomorphisme de
faisceaux d'anneaux $\mathcal{O}_X\to\mathscr{P}_f^n=\mathscr{P}_{X/S}^n$
(resp. $\mathcal{O}_X\to\mathscr{P}_f^\infty=\mathscr{P}_{X/S}^\infty$)
déduit de $p_2$. Pour tout ouvert $U$ de $X$, et tout
$t\in\Gamma(U,\mathcal{O}_X)$, $d^nt$ (resp. $d^\infty t$) est appelé la partie
principale d'ordre $n$ (resp. partie principale d'ordre infini) de $t$. On pose
$dt=d^1t-t$, et on dit que $dt$ est la différentielle de $t$ (élément de
$\Gamma(U,\Omega_{X/S}^1)$, aussi noté $d_{X/S}(t)$).
\end{definition}

Il résulte aussitôt de cette définition que l'on a
\[
  d(t_1t_2)=t_1dt_2+t_2dt_1
  \tag{16.3.6.1}
  \label{IV.16.3.6.1}
\]
quels que soient $t_1,t_2$ dans $\Gamma(U,\mathcal{O}_X)$, autrement dit, $d$
est une dérivation de l'anneau $\Gamma(U,\mathcal{O}_X)$ dans le
$\Gamma(U,\mathcal{O}_X)$-module $\Gamma(U,\Omega_{X/S}^1)$.

Dans toutes les notations introduites dans (16.3.1) et (16.3.6), on remplacera
parfois $S$ par $A$ lorsque $S=\operatorname{Spec}(A)$.

\begin{env}[16.3.7]
\label{IV.16.3.7}
Supposons en particulier que $S=\operatorname{Spec}(A)$ et
$X=\operatorname{Spec}(B)$ soient des schémas affines, $B$ étant donc une
$A$-algèbre. Alors $\Delta_f$ correspond à l'homomorphisme canonique surjectif
$\pi:B\otimes_A B\to B$ tel que $\pi(b\otimes b')=bb'$, de noyau
$\mathfrak{I}=\mathfrak{I}_{B/A}$ (0, 20.4.1); $\mathscr{P}_f^n$ est le faisceau
structural du préschéma $\operatorname{Spec}(P_{B/A}^n)$, où
\[
  P_{B/A}^n=(B\otimes_A B)/\mathfrak{I}^{n+1};
\]
$\mathscr{G}r_\bullet(\mathscr{P}_f)$ est le $\mathcal{O}_X$-Module
quasi-cohérent correspondant au $B$-module gradué
\[
  \operatorname{gr}_{\mathfrak{I}}^\bullet(B\otimes_A B)
  =\bigoplus_{n\geq 0}(\mathfrak{I}^n/\mathfrak{I}^{n+1});
\]
en particulier $\Omega_f^1=\Omega_{X/S}^1$ est le $\mathcal{O}_X$-Module
quasi-cohérent correspondant au $B$-module des 1-différentielles de $B$ par
rapport à $A$, $\Omega_{B/A}^1$ (0, 20.4.3). Les morphismes projections
$p_1:X\times_S X\to X$, $p_2:X\times_S X\to X$ correspondent aux deux
homomorphismes d'anneaux $j_1:B\to B\otimes_A B$, $j_2:B\to B\otimes_A B$
tels que $j_1(b)=b\otimes 1$, $j_2(b)=1\otimes b$, de sorte que (par la
convention de (16.3.5)), $P_{B/A}^n$ est toujours considérée comme une
$B$-algèbre pour l'homomorphisme composé
$B\xrightarrow{j_1}B\otimes_A B\to P_{B/A}^n$; l'homomorphisme d'anneaux
$B\xrightarrow{j_2}B\otimes_A B\to P_{B/A}^n$ se note $d_{B/A}^n$ et correspond
à $d_{X/S}^n$ opérant sur $\Gamma(X,\mathcal{O}_X)$; pour tout $t\in B$, $dt$
est égal à $d_{B/A}t$, défini dans (0, 20.4.6) (cf. $\mathbf{Err}_{\mathrm{IV}}$,
11).

Si $\pi_n:B\otimes_A B\to P_{B/A}^n$ est l'homomorphisme canonique, on a donc,
en vertu des définitions précédentes
\[
  \pi_n(b\otimes b')=b.\pi_n(1\otimes b')=b.d_{B/A}^n(b')
  \quad\text{pour }b\in B,\ b'\in B.
  \tag{16.3.7.1}
  \label{IV.16.3.7.1}
\]
\end{env}

\oldpage[IV-4]{16}
\begin{proposition}[16.3.8]
\label{IV.16.3.8}
L'image de l'homomorphisme $d_{X/S}^n:\mathcal{O}_X\to\mathscr{P}_{X/S}^n$
engendre le $\mathcal{O}_X$-Module $\mathscr{P}_{X/S}^n$.
\end{proposition}

On se ramène aussitôt au cas où $X=\operatorname{Spec}(B)$ et
$S=\operatorname{Spec}(A)$ sont affines et la proposition résulte de (16.3.7.1)
puisque $\pi_n$ est surjectif. On notera qu'en général $d_{X/S}^n$ n'est pas
surjectif (même déjà pour $n=1$).

\begin{proposition}[16.3.9]
\label{IV.16.3.9}
Supposons que $f:X\to S$ soit un morphisme localement de type fini. Alors les
$\mathscr{P}_f^n$ et les $\mathscr{G}r_n(\mathscr{P}_f)$ sont des
$\mathcal{O}_X$-Modules quasi-cohérents de type fini.
\end{proposition}

Cela résulte de (16.1.6) et du fait que $\Delta_f$ est localement de présentation
finie (I, 4.3.1).

\subsection{Propriétés fonctorielles des invariants différentiels}
\EGAReaderTarget{ega4.c04.subsec.16-4}{IV.16.4}

\begin{env}[16.4.1]
\label{IV.16.4.1}
Considérons un diagramme commutatif de morphismes de préschémas
\[
  \xymatrix{
    X \ar[d]_-{f} & X' \ar[l]_-u \ar[d]^-{f'} \\
    S & S' \ar[l]^-w
  }
  \tag{16.4.1.1}
  \label{IV.16.4.1.1}
\]
On en déduit un diagramme commutatif
\[
  \xymatrix{
    X \ar[d]_-{\Delta_f} & X' \ar[l]_-u \ar[d]^-{\Delta_{f'}} \\
    X\times_S X & X'\times_{S'}X' \ar[l]^-v
  }
\]
où $v$ est l'homomorphisme composé (I, 5.3.5 et 5.3.15)
\[
  X'\times_{S'}X'
  \xrightarrow{(p'_1,p'_2)_S}X'\times_S X'
  \xrightarrow{u\times_S u}X\times_S X.
  \tag{16.4.1.2}
  \label{IV.16.4.1.2}
\]
On déduit donc de $u$ et $v$, comme il a été expliqué dans (16.2.1), des
homomorphismes de faisceaux d'anneaux augmentés
\[
  \nu_n:\rho^*(\mathscr{P}_{X/S}^n)\longrightarrow
  \mathscr{P}_{X'/S'}^n
  \tag{16.4.1.3}
  \label{IV.16.4.1.3}
\]
(où l'on a posé $u=(\rho,\lambda)$); ces homomorphismes forment un système
projectif, et donnent donc à la limite un homomorphisme de faisceaux d'anneaux
augmentés
\[
  \nu_\infty:\rho^*(\mathscr{P}_{X/S}^\infty)\longrightarrow
  \mathscr{P}_{X'/S'}^\infty;
  \tag{16.4.1.4}
  \label{IV.16.4.1.4}
\]
d'autre part, par passage aux quotients, les homomorphismes $\nu_n$ donnent un
di-homomorphisme d'Algèbres graduées (relatif à $\lambda^\#$):
\[
  \operatorname{gr}(u):\rho^*(\mathscr{G}r_\bullet(\mathscr{P}_{X/S}))
  \longrightarrow\mathscr{G}r_\bullet(\mathscr{P}_{X'/S'}).
  \tag{16.4.1.5}
  \label{IV.16.4.1.5}
\]
\end{env}

\begin{env}[16.4.2]
\label{IV.16.4.2}
\oldpage[IV-4]{17}
Si l'on a un diagramme commutatif
\[
  \xymatrix{
    X \ar[d]_-{f} & X' \ar[l]_-u \ar[d]^-{f'} &
      X'' \ar[l]_-{u'} \ar[d]^-{f''} \\
    S & S' \ar[l]^-w & S'' \ar[l]^-{w'}
  }
\]
on en déduit un diagramme commutatif
\[
  \xymatrix{
    X \ar[d]_-{\Delta_f} & X' \ar[l]_-u \ar[d]^-{\Delta_{f'}} &
      X'' \ar[l]_-{u'} \ar[d]^-{\Delta_{f''}} \\
    X\times_S X & X'\times_{S'}X' \ar[l]^-v &
      X''\times_{S''}X'' \ar[l]^-{v'}
  }
\]
où $v'$ est défini à partir de $u'$, $w'$, $f'$, $f''$ comme $v$ à partir de
$u$, $w$, $f$, $f'$. On vérifie aussitôt que si $u''=u\circ u'$,
$w''=w\circ w'$, alors le morphisme composé $v\circ v'$ est égal au morphisme
$v''$ défini à partir de $u''$, $w''$, $f$, $f''$ comme $v$ à partir de $u$,
$w$, $f$, $f'$. Si l'on pose $u'=(\rho',\lambda')$,
$u''=(\rho'',\lambda'')$ il résulte alors de (16.2.1) que l'homomorphisme
$\nu_n'':\rho''^*(\mathscr{P}_{X/S}^n)\to\mathscr{P}_{X''/S''}^n$ est égal au
composé
\[
  \rho'^*(\rho^*(\mathscr{P}_{X/S}^n))
  \xrightarrow{\rho'^*(\nu_n^\#)}
  \rho'^*(\mathscr{P}_{X'/S'}^n)
  \xrightarrow{\nu_n'}\mathscr{P}_{X''/S''}^n
\]
et l'on a des propriétés de transitivité analogues pour les homomorphismes
(16.4.1.4) et (16.4.1.5), ce qui permet de dire que les
$\mathscr{P}_{X/S}^n$, $\mathscr{P}_{X/S}^\infty$ et
$\mathscr{G}r_\bullet(\mathscr{P}_{X/S})$ dépendent fonctoriellement de $f$.
\end{env}

\begin{env}[16.4.3]
\label{IV.16.4.3}
On vérifie aussitôt (par exemple en se ramenant au cas affine à l'aide de
(16.3.7)) qu'avec les notations de (16.4.1), le diagramme
\[
  \xymatrix{
    \rho^*(\mathcal{O}_X) \ar[r]^{\lambda^\#} \ar[d] &
      \mathcal{O}_{X'} \ar[d] \\
    \rho^*(\mathscr{P}_{X/S}^n) \ar[r]_{\nu_n} &
      \mathscr{P}_{X'/S'}^n
  }
  \tag{16.4.3.1}
  \label{IV.16.4.3.1}
\]
où les flèches verticales sont celles définissant les structures d'Algèbre
choisies dans (16.3.5) (c'est-à-dire celles provenant des premières projections)
est commutatif; il en est de même du diagramme
\[
  \xymatrix{
    \rho^*(\mathcal{O}_X) \ar[r]^{\lambda^\#}
      \ar[d]_{\rho^*(d_{X/S}^n)} &
      \mathcal{O}_{X'} \ar[d]^{d_{X'/S'}^n} \\
    \rho^*(\mathscr{P}_{X/S}^n) \ar[r]_{\nu_n} &
      \mathscr{P}_{X'/S'}^n
  }
  \tag{16.4.3.2}
  \label{IV.16.4.3.2}
\]
\oldpage[IV-4]{18}
les flèches verticales définissant ici les structures d'Algèbre provenant des
secondes projections; d'ailleurs, si $\sigma$ et $\sigma'$ sont les symétries
canoniques correspondant à $f$ et $f'$ (16.3.4), on a
\[
  \nu_n\circ\rho^*(\sigma)=\sigma'\circ\nu_n
\]
qui fait passer d'un des diagrammes précédents à l'autre. On déduit donc de
(16.4.3.1) un homomorphisme canonique de $\mathcal{O}_{X'}$-Algèbres augmentées
\[
  P^n(u):u^*(\mathscr{P}_{X/S}^n)
  =\mathscr{P}_{X/S}^n\otimes_{\mathcal{O}_X}\mathcal{O}_{X'}
  \longrightarrow\mathscr{P}_{X'/S'}^n
  \tag{16.4.3.3}
  \label{IV.16.4.3.3}
\]
et il résulte de (16.4.3.2) que le diagramme
\[
  \xymatrix{
    \mathcal{O}_{X'} \ar[r]^{\operatorname{id}}
      \ar[d]_{u^*(d_{X/S}^n)} &
      \mathcal{O}_{X'} \ar[d]^{d_{X'/S'}^n} \\
    u^*(\mathscr{P}_{X/S}^n) \ar[r]_{P^n(u)} &
      \mathscr{P}_{X'/S'}^n
  }
  \tag{16.4.3.4}
  \label{IV.16.4.3.4}
\]
est commutatif. On en déduit un homomorphisme de $\mathcal{O}_{X'}$-Algèbres
graduées
\[
  \operatorname{Gr}_\bullet(u):u^*(\mathscr{G}r_\bullet(\mathscr{P}_{X/S}))
  \longrightarrow\mathscr{G}r_\bullet(\mathscr{P}_{X'/S'})
  \tag{16.4.3.5}
  \label{IV.16.4.3.5}
\]
et en particulier un homomorphisme de $\mathcal{O}_{X'}$-Modules
\[
  \operatorname{Gr}_1(u):
  \Omega_{X/S}^1\otimes_{\mathcal{O}_X}\mathcal{O}_{X'}
  \longrightarrow\Omega_{X'/S'}^1
  \tag{16.4.3.6}
  \label{IV.16.4.3.6}
\]
donnant lieu à un diagramme commutatif
\[
  \xymatrix{
    \mathcal{O}_{X'} \ar[r]^{\operatorname{id}}
      \ar[d]_{d_{X/S}\otimes 1} &
      \mathcal{O}_{X'} \ar[d]^{d_{X'/S'}} \\
    \Omega_{X/S}^1\otimes_{\mathcal{O}_X}\mathcal{O}_{X'} \ar[r] &
      \Omega_{X'/S'}^1
  }
  \tag{16.4.3.7}
  \label{IV.16.4.3.7}
\]
\end{env}

\begin{env}[16.4.4]
\label{IV.16.4.4}
Lorsque $S=\operatorname{Spec}(A)$, $S'=\operatorname{Spec}(A')$,
$X=\operatorname{Spec}(B)$, $X'=\operatorname{Spec}(B')$ sont affines, de sorte
que l'on a un diagramme commutatif d'homomorphismes d'anneaux
\[
  \xymatrix{
    B \ar[r] & B' \\
    A \ar[u] \ar[r] & A' \ar[u]
  }
\]
l'image de $\mathfrak{I}_{B/A}$ dans $B'\otimes_{A'}B'$ est contenue dans
$\mathfrak{I}_{B'/A'}$, et l'homomorphisme $\nu_n$ correspond à
l'homomorphisme d'anneaux $P_{B/A}^n\to P_{B'/A'}^n$ déduit de l'homomorphisme
$B\otimes_A B\to B'\otimes_{A'}B'$ par passage aux quotients. L'homomorphisme
(16.4.3.6) correspond à l'homomorphisme défini dans (0, 20.5.4.1), et le
diagramme commutatif (16.4.3.7) au diagramme (0, 20.5.4.2).
\end{env}

\oldpage[IV-4]{19}
\begin{proposition}[16.4.5]
\label{IV.16.4.5}
Supposons que $X'=X\times_S S'$, $f'$ et $u$ étant les projections canoniques.
Alors les homomorphismes canoniques $P^n(u)$ (16.4.3.3) et
$\operatorname{Gr}_1(u)$ (16.4.3.6) sont bijectifs.
\end{proposition}

On a en effet alors $X'\times_{S'}X'=(X\times_S X)\times_S S'$, et on peut donc
appliquer (16.2.3, (ii)) en remplaçant $g$ par la première projection
$p_1:X\times_S X\to X$ et $f$ par la diagonale $\Delta_f$.

On notera aussi que sous les hypothèses de (16.4.5) l'homomorphisme
$\operatorname{Gr}_\bullet(u)$ (16.4.3.5) est surjectif, mais non bijectif en
général. Toutefois (16.2.4):

\begin{corollary}[16.4.6]
\label{IV.16.4.6}
Les hypothèses étant celles de (16.4.5), supposons de plus que $w:S'\to S$ soit
plat (resp. que les $\mathscr{G}r_n(\mathscr{P}_{X/S})$ soient des
$\mathcal{O}_X$-Modules plats pour $n\leq m$); alors l'homomorphisme
\[
  \operatorname{Gr}_n(u):u^*(\mathscr{G}r_n(\mathscr{P}_{X/S}))
  \longrightarrow\mathscr{G}r_n(\mathscr{P}_{X'/S'})
\]
est bijectif pour tout $n$ (resp. pour $n\leq m$).
\end{corollary}

En effet, si $w$ est plat, il en est de même de
$v:X'\times_{S'}X'\to X\times_S X$, donc la conclusion résulte de (16.2.4).

\begin{env}[16.4.7]
\label{IV.16.4.7}
Soient $S$ un préschéma, $\mathscr{E}$ un $\mathcal{O}_S$-Module
quasi-cohérent, et posons $X=\mathbf{V}(\mathscr{E})$ (II, 1.7.8), fibré
vectoriel associé à $\mathscr{E}$, égal à
$\operatorname{Spec}(\mathscr{S}_{\mathcal{O}_S}(\mathscr{E}))$. Soit
$f:X\to S$ le morphisme structural. Pour tout ouvert $U$ de $S$, et toute
section $t\in\Gamma(U,\mathscr{E})$, $t$ s'identifie à une section de
$\mathscr{S}_{\mathcal{O}_S}(\mathscr{E})$ au-dessus de $U$; soit $t'$ son
image dans
$\Gamma(f^{-1}(U),\mathcal{O}_X)=\Gamma(U,f_*(\mathcal{O}_X))=
\Gamma(U,\mathscr{S}_{\mathcal{O}_S}(\mathscr{E}))$, et posons
\[
  \delta(t)=d_{X/S}^n(t')-t'\in
  \Gamma(f^{-1}(U),\mathscr{P}_{X/S}^n);
  \tag{16.4.7.1}
  \label{IV.16.4.7.1}
\]
il est clair que $\delta$ est un di-homomorphisme de modules (correspondant à
l'homomorphisme d'anneaux
$\Gamma(U,\mathcal{O}_S)\to\Gamma(f^{-1}(U),\mathcal{O}_X)$) de
$\Gamma(U,\mathscr{E})$ dans $\Gamma(f^{-1}(U),\mathscr{P}_{X/S}^n)$, dont
l'image appartient d'ailleurs à l'idéal d'augmentation de
$\Gamma(f^{-1}(U),\mathscr{P}_{X/S}^n)$. On en déduit (en faisant varier $U$)
un homomorphisme canonique de $\mathcal{O}_X$-Algèbres
\[
  f^*(\mathscr{S}_{\mathcal{O}_S}(\mathscr{E}))
  \longrightarrow\mathscr{P}_{X/S}^n
  \tag{16.4.7.2}
  \label{IV.16.4.7.2}
\]
et d'après la remarque précédente, si $\mathscr{K}$ est l'Idéal noyau de
l'augmentation $\mathscr{S}_{\mathcal{O}_S}(\mathscr{E})\to\mathcal{O}_S$,
l'image de $\mathscr{K}^{n+1}$ par (16.4.7.2) est nulle, si bien qu'en
factorisant par $\mathscr{K}^{n+1}$, on a finalement un homomorphisme canonique
\[
  \delta_n:f^*(\mathscr{S}_{\mathcal{O}_S}(\mathscr{E})/
  \mathscr{K}^{n+1})\longrightarrow\mathscr{P}_{X/S}^n.
  \tag{16.4.7.3}
  \label{IV.16.4.7.3}
\]
\end{env}

\begin{proposition}[16.4.8]
\label{IV.16.4.8}
Sous les conditions de (16.4.7), les homomorphismes $\delta_n$ sont bijectifs et
forment un système projectif d'isomorphismes; on en déduit un isomorphisme de
$\mathcal{O}_X$-Algèbres graduées
\[
  f^*(\mathscr{S}_{\mathcal{O}_S}^{\bullet}(\mathscr{E}))
  \longrightarrow\mathscr{G}r_\bullet(\mathscr{P}_{X/S}).
  \tag{16.4.8.1}
  \label{IV.16.4.8.1}
\]
\end{proposition}

Le fait que les homomorphismes (16.4.7.3) forment un système projectif résulte
aussitôt de leur définition. Pour prouver que ce sont des isomorphismes, il
suffira de
\oldpage[IV-4]{20}
démontrer que (16.4.8.1) est un isomorphisme, les filtrations des deux membres de
(16.4.7.3) étant finies (Bourbaki, \emph{Alg. comm.}, chap. III, \S~2,
n\textsuperscript{o} 8, cor. 3 du th. 1). Pour cela, considérons la suite exacte
scindée de $\mathcal{O}_S$-Modules
\[
  \xymatrix{
    0 \ar[r] & \mathscr{E} \ar[r]^{u} &
    \mathscr{E}\oplus\mathscr{E} \ar[r]^{v} &
    \mathscr{E} \ar[r] & 0
  }
  \tag{16.4.8.2}
  \label{IV.16.4.8.2}
\]
où, pour tout couple de sections $s$, $t$ de $\mathscr{E}$ au-dessus d'un ouvert
$U$ de $S$, on prend $u(s)=(-s,s)$ et $v(s,t)=s+t$. On a
\[
  X\times_S X=\operatorname{Spec}
  (\mathscr{S}_{\mathcal{O}_S}(\mathscr{E})\otimes_{\mathcal{O}_S}
  \mathscr{S}_{\mathcal{O}_S}(\mathscr{E}))
  =\operatorname{Spec}(\mathscr{S}_{\mathcal{O}_S}
  (\mathscr{E}\oplus\mathscr{E}))
\]
(II, 1.4.6 et 1.7.11), et le morphisme diagonal $X\to X\times_S X$ correspond
(II, 1.2.7) à l'homomorphisme de $\mathcal{O}_S$-Algèbres
$\mathscr{S}(v):\mathscr{S}_{\mathcal{O}_S}(\mathscr{E}\oplus\mathscr{E})
\to\mathscr{S}_{\mathcal{O}_S}(\mathscr{E})$ (II, 1.7.4), de sorte que si
$\mathscr{I}$ est le noyau de cet homomorphisme, on a
\[
  \mathscr{P}_{X/S}^n=f^*(\mathscr{S}_{\mathcal{O}_S}
  (\mathscr{E}\oplus\mathscr{E})/\mathscr{I}^{n+1}).
\]

La proposition sera conséquence du lemme suivant:

\begin{lemma}[16.4.8.3]
\label{IV.16.4.8.3}
Soient $Y$ un espace annelé,
$0\to\mathscr{F}'\xrightarrow{u}\mathscr{F}\xrightarrow{v}\mathscr{F}''\to0$
une suite exacte de $\mathcal{O}_Y$-Modules telle que tout point $y\in Y$
possède un voisinage ouvert $V$ tel que la suite
$0\to\mathscr{F}'|V\to\mathscr{F}|V\to\mathscr{F}''|V\to0$ soit scindée.
Soit $\mathscr{I}$ l'Idéal noyau de
$\mathscr{S}(v):\mathscr{S}_{\mathcal{O}_Y}(\mathscr{F})\to
\mathscr{S}_{\mathcal{O}_Y}(\mathscr{F}'')$, et soit
$\operatorname{gr}_{\mathscr{I}}^\bullet
(\mathscr{S}_{\mathcal{O}_Y}(\mathscr{F}))$ la $\mathcal{O}_Y$-Algèbre graduée
associée à la $\mathcal{O}_Y$-Algèbre
$\mathscr{S}_{\mathcal{O}_Y}(\mathscr{F})$ munie de la filtration
$\mathscr{I}$-préadique. Alors l'homomorphisme de $\mathcal{O}_Y$-Algèbres
graduées
\[
  \mathscr{S}_{\mathcal{O}_Y}^{\bullet}(\mathscr{F}')
  \otimes_{\mathcal{O}_Y}
  \mathscr{S}_{\mathcal{O}_Y}^{\bullet}(\mathscr{F}'')
  \longrightarrow
  \operatorname{gr}_{\mathscr{I}}^\bullet
  (\mathscr{S}_{\mathcal{O}_Y}(\mathscr{F}))
  \tag{16.4.8.4}
  \label{IV.16.4.8.4}
\]
(où le premier membre est le produit tensoriel gradué des
$\mathcal{O}_Y$-Algèbres symétriques munies de leur graduation canonique
(II, 1.7.4 et 2.1.2)), provenant de l'injection canonique
\[
  \mathscr{F}'\longrightarrow\mathscr{I}
  =\operatorname{gr}_{\mathscr{I}}^1
  (\mathscr{S}_{\mathcal{O}_Y}(\mathscr{F})),
\]
est bijectif.
\end{lemma}

L'injection $\mathscr{F}'\to\mathscr{I}$ donne en effet canoniquement un
homomorphisme de $\mathcal{O}_Y$-Algèbres graduées\linebreak
$\mathscr{S}_{\mathcal{O}_Y}^{\bullet}(\mathscr{F}')\to\allowbreak
\operatorname{gr}_{\mathscr{I}}^\bullet
(\mathscr{S}_{\mathcal{O}_Y}(\mathscr{F}))$, et comme le second membre est par
définition une $\mathscr{S}_{\mathcal{O}_Y}^{\bullet}(\mathscr{F}'')$-Algèbre
graduée, on en déduit l'homomorphisme canonique (16.4.8.4) par tensorisation du
précédent avec $\mathscr{S}_{\mathcal{O}_Y}^{\bullet}(\mathscr{F}'')$. Pour
prouver le lemme, on peut, la question étant locale, se borner au cas où
$\mathscr{F}=\mathscr{F}'\oplus\mathscr{F}''$, $u$ et $v$ étant les
homomorphismes canoniques. Alors l'Algèbre graduée
$\mathscr{S}_{\mathcal{O}_Y}^{\bullet}(\mathscr{F})$ s'identifie canoniquement au
produit tensoriel gradué
$\mathscr{S}_{\mathcal{O}_Y}^{\bullet}(\mathscr{F}')
\otimes_{\mathcal{O}_Y}
\mathscr{S}_{\mathcal{O}_Y}^{\bullet}(\mathscr{F}'')$ (II, 1.7.4), et il est
immédiat que $\mathscr{I}$ est alors l'Idéal
$\mathscr{I}'\otimes_{\mathcal{O}_Y}
\mathscr{S}_{\mathcal{O}_Y}^{\bullet}(\mathscr{F}'')$, où $\mathscr{I}'$ est
l'Idéal d'augmentation de
$\mathscr{S}_{\mathcal{O}_Y}^{\bullet}(\mathscr{F}')$, c'est-à-dire la somme
(directe) des $\mathscr{S}_{\mathcal{O}_Y}^m(\mathscr{F}')$ pour $m\geq1$. On en
conclut que $\mathscr{I}^n=\mathscr{I}'^n\otimes_{\mathcal{O}_Y}
\mathscr{S}_{\mathcal{O}_Y}^{\bullet}(\mathscr{F}'')$, où cette fois
$\mathscr{I}'^n$ est la somme (directe) des
$\mathscr{S}_{\mathcal{O}_Y}^m(\mathscr{F}')$ pour $m\geq n$; on a par suite
$\mathscr{I}^n/\mathscr{I}^{n+1}=
\mathscr{S}_{\mathcal{O}_Y}^n(\mathscr{F}')\otimes_{\mathcal{O}_Y}
\mathscr{S}_{\mathcal{O}_Y}^{\bullet}(\mathscr{F}'')$, ce qui prouve que
(16.4.8.4) est bijectif.

\oldpage[IV-4]{21}
Ce lemme étant démontré, il reste à voir que l'homomorphisme (16.4.8.1) est
bien l'image par $f^*$ de l'homomorphisme (16.4.8.4) correspondant à la suite
exacte (16.4.8.2); on constate aisément que cela résulte de la définition de
$u$ (16.4.8.2) et de celle de $\delta$ (16.4.7.1), compte tenu de la définition
de la structure de $\mathcal{O}_X$-Algèbre sur $\mathscr{P}_{X/S}^n$ et de
celle de $d_{X/S}^n$ (16.3.5 et 16.3.6).

En particulier:

\begin{corollary}[16.4.9]
\label{IV.16.4.9}
Sous les conditions de (16.4.7), on a un isomorphisme canonique
\[
  \operatorname{gr}_1(\delta):f^*(\mathscr{E})
  \xrightarrow{\sim}\Omega_{X/S}^1.
  \tag{16.4.9.1}
  \label{IV.16.4.9.1}
\]
\end{corollary}

\begin{corollary}[16.4.10]
\label{IV.16.4.10}
Si $S=\operatorname{Spec}(A)$, $\mathscr{E}=\mathcal{O}_S^m$, de sorte que
\[
  X=\operatorname{Spec}(A[T_1,\ldots,T_m]),
\]
$\mathscr{P}_{X/S}^n$ s'identifie canoniquement à la
$\mathcal{O}_X$-Algèbre correspondant à la $A[T_1,\ldots,T_m]$-algèbre quotient
\[
  A[T_1,\ldots,T_m,U_1,\ldots,U_m]/\mathfrak{K}^{n+1},
\]
où les $U_i$ ($1\leq i\leq m$) sont $m$ nouvelles indéterminées et
$\mathfrak{K}$ est l'idéal engendré par $U_1,\ldots,U_m$.
\end{corollary}

On retrouve ainsi en particulier la structure de $\Omega_{X/S}^1$ dans ce cas
(0, 20.5.13).

Notons en outre que $d_{X/S}^n$ fait alors correspondre à un polynôme
$F(T_1,\ldots,T_m)$ la classe mod.\ $\mathfrak{K}^{n+1}$ de
$F(T_1+U_1,\ldots,T_m+U_m)$, comme il résulte de la définition (16.4.7.1).

\begin{proposition}[16.4.11]
\label{IV.16.4.11}
Soient $f:X\to S$ un morphisme, $g:S\to X$ une $S$-section de $X$,
$S_g^{(n)}$ le $n$-ème voisinage infinitésimal de $S$ pour l'immersion $g$
(16.1.2). Il existe alors un isomorphisme et un seul de
$\mathcal{O}_S$-Algèbres
\[
  \varpi_n:g^*(\mathscr{P}_{X/S}^n)\longrightarrow
  \mathcal{O}_{S_g^{(n)}}
  \tag{16.4.11.1}
  \label{IV.16.4.11.1}
\]
(pour la structure de $\mathcal{O}_S$-Algèbre sur
$\mathcal{O}_{S_g^{(n)}}$ définie par $f$ (16.1.7)), rendant commutatif le
diagramme
\[
  \xymatrix{
    \mathcal{O}_S=g^*(\mathcal{O}_X)
      \ar[rr]^-{\lambda_n}\ar[dr]_-{g^*(d_{X/S}^n)}
      &&\mathcal{O}_{S_g^{(n)}}\\
    &g^*(\mathscr{P}_{X/S}^n)\ar[ur]_-{\varpi_n}
  }
  \tag{16.4.11.2}
  \label{IV.16.4.11.2}
\]
(où $\lambda_n$ est l'homomorphisme structural).
\end{proposition}

En vertu de (I, 5.3.7), où l'on remplace $X$, $Y$, $S$ par $S$, $X$, $S$
respectivement et $f$ par $g$, les diagrammes
\[
  \xymatrix@=3pc{
    S\ar[r]^-g\ar[d]_-g&X\ar[d]^-{\Delta_f}
      &S\ar[r]^-g\ar[d]_-g&X\ar[d]^-{\Delta_f}\\
    X\ar[r]_{(g\circ f,1_X)_S}&X\times_S X
      &X\ar[r]_{(1_X,g\circ f)_S}&X\times_S X
  }
  \tag{16.4.11.3}
  \label{IV.16.4.11.3}
\]
\oldpage[IV-4]{22}
identifient $S$ au produit des $(X\times_S X)$-préschémas $X$ et $X$ pour les
morphismes $\Delta_f$ et $(g\circ f,1_X)_S$ (resp.\ $(1_X,g\circ f)_S$).
D'autre part, les diagrammes
\[
  \xymatrix@=3pc{
    X\ar[r]^{(g\circ f,1_X)_S}\ar[d]_-f&X\times_S X\ar[d]^-{p_1}
      &X\ar[r]^{(1_X,g\circ f)_S}\ar[d]_-f&X\times_S X\ar[d]^-{p_2}\\
    S\ar[r]_{g}&X&S\ar[r]_{g}&X
  }
  \tag{16.4.11.4}
  \label{IV.16.4.11.4}
\]
identifient $X$ au produit des $X$-préschémas $S$ et $X\times_S X$ pour les
morphismes $g$ et $p_1$ (resp.\ $p_2$) (cas particulier de la formule
d'associativité (I, 3.3.9.1)). On peut dire que $\Delta_f$, considéré comme
$X$-section de $X\times_S X$ (relativement à $p_1$ ou $p_2$) joue le rôle
d'une \emph{section universelle} pour les $S$-sections de $X$: chacune de ces
sections $g$ s'en déduit en effet par le \emph{changement de base}
$(g\circ f,1_X)_S:X\to X\times_S X$. La définition de l'homomorphisme
$\varpi_n$ et le fait qu'il est bijectif résultent donc de ces remarques et de
(16.2.3, (ii)) appliqué au premier diagramme (16.4.11.4). La commutativité du
diagramme (16.4.11.2) résulte de même de (16.2.3, (ii)) appliqué cette fois au
second diagramme (16.4.11.4). Pour expliciter $\varpi_n$, on peut se borner au
cas où $g$ est une immersion \emph{fermée}: en effet, pour tout $s\in S$, il y
a un voisinage ouvert $W$ de $s$ dans $S$ tel que $g(W)$ soit fermé dans un
ouvert $U$ de $X$, et il est clair que $g|W$ est une $W$-section du morphisme
$U\cap f^{-1}(W)\to W$, restriction de $f$, et $g(W)$ est \emph{a fortiori}
fermé dans $U\cap f^{-1}(W)$. On peut donc supposer que $S$ est un
sous-préschéma fermé de $X$ défini par un Idéal quasi-cohérent $\mathscr{K}$.
Alors les définitions précédentes montrent que si $W$ est un ouvert de $S$,
$t$ une section de $\mathcal{O}_X$ au-dessus de $f^{-1}(W)$,
$\varpi_n(d^n t|W)$ est égal à l'image canonique de $t$ dans
$\Gamma(W,(\mathcal{O}_X/\mathscr{K}^{n+1})|W)$. L'unicité de $\varpi_n$
résulte donc de ce que l'image de $\mathcal{O}_X$ par $d_{X/S}^n$ engendre le
$\mathcal{O}_X$-Module $\mathscr{P}_{X/S}^n$ (16.3.8).

\begin{corollary}[16.4.12]
\label{IV.16.4.12}
Soient $k$ un corps, $X$ un $k$-préschéma, $x$ un point de $X$ rationnel sur
$k$. Alors $(\mathscr{P}_{X/k}^n)_x\otimes_{\mathcal{O}_x}k(x)$ est
canoniquement isomorphe (en tant que $k(x)$-algèbre augmentée) à
$\mathcal{O}_x/\mathfrak{m}_x^{n+1}$.
\end{corollary}

Il suffit de considérer l'unique $k$-section $g$ de $X$ telle que
$g(\operatorname{Spec}(k))=\{x\}$.

\begin{corollary}[16.4.13]
\label{IV.16.4.13}
Soient $f:X\to S$ un morphisme, $s$ un point de $S$,
$X_s=X\times_S\operatorname{Spec}(k(s))$ la fibre de $f$ en $s$. Si
$x\in X_s$ est rationnel sur $k(s)$,
$(\mathscr{P}_{X/S}^n)_x\otimes_{\mathcal{O}_s}k(s)$ est canoniquement
isomorphe à $\mathcal{O}_{X_s,x}/\mathfrak{m}'_x{}^{\,n+1}$, où
$\mathfrak{m}'_x$ est l'idéal maximal de $\mathcal{O}_{X_s,x}$; de façon
précise, cet isomorphisme fait correspondre à $(d^n t)_x\otimes1$ (où $t$ est
une section de $\mathcal{O}_X$ au-dessus d'un voisinage ouvert de $x$ dans
$X$) la classe de $t_x\otimes1$ modulo $\mathfrak{m}'_x{}^{\,n+1}$.
\end{corollary}

Cela résulte de (16.4.5) et (16.4.12).

Les corollaires précédents justifient la terminologie de «faisceau des parties
principales d'ordre $n$».

\begin{proposition}[16.4.14]
\label{IV.16.4.14}
Soit $\rho:A\to B$ un homomorphisme d'anneaux, $S$ une partie
multiplicativement stable de $B$. Alors les homomorphismes canoniques
\[
  S^{-1}P_{B/A}^n\longrightarrow P_{S^{-1}B/A}^n
  \tag{16.4.14.1}
  \label{IV.16.4.14.1}
\]
\oldpage[IV-4]{23}
déduits des homomorphismes canoniques
$P_{B/A}^n\to P_{S^{-1}B/A}^n$ (16.4.4), forment un système projectif et sont
bijectifs.
\end{proposition}

Il suffit de remarquer que
\[
  S^{-1}((B\otimes_A B)/\mathfrak{I}^{n+1})
  =S^{-1}(B\otimes_A B)/(S^{-1}\mathfrak{I})^{n+1}
\]
par platitude, et que
$S^{-1}(B\otimes_A B)=(S^{-1}B)\otimes_A(S^{-1}B)$ (I, 1.3.4).

\begin{corollary}[16.4.15]
\label{IV.16.4.15}
Les notations étant celles de (16.4.14), soit $R$ une partie multiplicative de
$A$ telle que $\rho(R)\subset S$. Alors on a des isomorphismes canoniques
\[
  S^{-1}P_{B/A}^n\xrightarrow{\sim}P_{S^{-1}B/R^{-1}A}^n
  \tag{16.4.15.1}
  \label{IV.16.4.15.1}
\]
formant un système projectif.
\end{corollary}

Il suffit évidemment de définir des isomorphismes canoniques
\[
  P_{S^{-1}B/A}^n\xrightarrow{\sim}P_{S^{-1}B/R^{-1}A}^n,
  \tag{16.4.15.2}
  \label{IV.16.4.15.2}
\]
c'est-à-dire que l'on est ramené au cas où $\rho(R)$ est formé d'éléments
\emph{inversibles} de $B$. Mais alors l'isomorphisme (16.4.15.2) est simplement
déduit de l'isomorphisme canonique
$B\otimes_A B\to B\otimes_{R^{-1}A}B$ par passage aux quotients
(0$_{\rm I}$, 1.5.3).

\begin{corollary}[16.4.16]
\label{IV.16.4.16}
Soient $f:X\to S$ un morphisme de préschémas, $x$ un point de $X$,
$s=f(x)$. Alors on a des isomorphismes canoniques
\[
  (\mathscr{P}_{X/S}^n)_x\xrightarrow{\sim}
  P_{\mathcal{O}_x/\mathcal{O}_s}^n
  \tag{16.4.16.1}
  \label{IV.16.4.16.1}
\]
formant un système projectif.
\end{corollary}

On déduit de là des isomorphismes pour les gradués associés, et en particulier
un isomorphisme canonique
\[
  (\Omega_{X/S}^1)_x\xrightarrow{\sim}
  \Omega_{\mathcal{O}_x/\mathcal{O}_s}^1.
  \tag{16.4.16.2}
  \label{IV.16.4.16.2}
\]

\begin{corollary}[16.4.17]
\label{IV.16.4.17}
Soient $k$ un corps, $K$ le corps des fractions rationnelles
$k(T_1,\ldots,T_r)$. Alors, pour tout entier $n$, l'homomorphisme de
$K[U_1,\ldots,U_r]$ ($U_i$ indéterminées) dans $P_{K/k}^n$ qui, à tout $U_i$,
fait correspondre $d^nT_i-T_i\cdot1$, est surjectif et définit un isomorphisme
du quotient $K[U_1,\ldots,U_r]/\mathfrak{m}^{n+1}$ (où $\mathfrak{m}$ est
l'idéal engendré par les $U_i$) sur $P_{K/k}^n$.
\end{corollary}

Cela résulte de (16.4.8), (16.4.10) et (16.4.14), où l'on fait $A=k$,
$B=k[T_1,\ldots,T_r]$ et $S=B-\{0\}$.

On retrouve ainsi le fait que les $dT_i$ forment une base du $K$-espace
vectoriel $\Omega_{K/k}^1$ (0, 20.5.10).

\begin{proposition}[16.4.18]
\label{IV.16.4.18}
Soient $f:X\to Y$, $g:Y\to Z$ deux morphismes de préschémas; et considérons
les homomorphismes canoniques de $\mathcal{O}_X$-Algèbres augmentées
(16.4.3.3)
\[
  g_{X/Y/Z}:\mathscr{P}_{X/Z}^n\longrightarrow\mathscr{P}_{X/Y}^n
  \tag{16.4.18.1}
  \label{IV.16.4.18.1}
\]
\[
  f_{X/Y/Z}:f^*(\mathscr{P}_{Y/Z}^n)
  \longrightarrow\mathscr{P}_{X/Z}^n.
  \tag{16.4.18.2}
  \label{IV.16.4.18.2}
\]
Alors $g_{X/Y/Z}$ est surjectif, et son noyau est l'Idéal engendré par l'image
par $f_{X/Y/Z}$ de l'Idéal d'augmentation de
$f^*(\mathscr{P}_{Y/Z}^n)$.
\end{proposition}

\oldpage[IV-4]{24}
Notons d'abord que $g_{X/Y/Z}$ correspond au cas où dans (16.4.3.3) on fait
$X'=X$, $S'=Y$, $S=Z$ et $u=1_X$, et $f_{X/Y/Z}$ au cas où l'on remplace
$X'$, $X$, $S$, $S'$ par $X$, $Y$, $Z$, $Z$ respectivement et $u$, $f$ par
$f$, $g$ respectivement.

On a un diagramme commutatif (I, 5.3.5)
\[
  \xymatrix{
    X\ar[r]^-{\Delta_f}\ar[dr]_-f
      &X\times_Y X\ar[d]^-p\ar[r]^-j
      &X\times_Z X\ar[d]^-{f\times_Z f}\\
    &Y\ar[r]_-{\Delta_g}&Y\times_Z Y
  }
  \tag{16.4.18.3}
  \label{IV.16.4.18.3}
\]
où $j=(1_X,1_X)_Z$ est une immersion, $j\circ\Delta_f=\Delta_{g\circ f}$ et
$p$ est le morphisme structural. Comme on peut se borner au cas où $X$, $Y$,
$Z$ sont affines, on peut supposer les immersions $\Delta_f$, $\Delta_g$ et
$j$ fermées, de sorte que $\mathcal{O}_X$ et
$\mathcal{O}_{X\times_Y X}$ s'identifient respectivement à
$\mathcal{O}_{X\times_Z X}/\mathscr{J}$ et
$\mathcal{O}_{X\times_Z X}/\mathscr{L}$, où
$\mathscr{J}\supset\mathscr{L}$ sont les deux Idéaux quasi-cohérents
correspondant respectivement aux immersions $\Delta_{g\circ f}$ et $j$. La
$\mathcal{O}_X$-Algèbre $\mathscr{P}_{X/Z}^n$ s'identifie donc à
$\mathcal{O}_{X\times_Z X}/\mathscr{J}^{n+1}$, et
$\mathscr{P}_{X/Y}^n$ s'identifie à
$\mathcal{O}_{X\times_Y X}/(\mathscr{J}/\mathscr{L})^{n+1}$, c'est-à-dire à
$\mathcal{O}_{X\times_Z X}/(\mathscr{J}^{n+1}+\mathscr{L})$, et par suite au
quotient de $\mathscr{P}_{X/Z}^n$ par
$(\mathscr{J}^{n+1}+\mathscr{L})/\mathscr{J}^{n+1}$. Mais on sait
(\emph{loc. cit.}) que $p$ et $j$ font de $X\times_Y X$ le produit des
$(Y\times_Z Y)$-préschémas $Y$ et $X\times_Z X$, donc si $\mathcal{O}_Y$ est
identifié à $\mathcal{O}_{Y\times_Z Y}/\mathscr{K}$, où $\mathscr{K}$ est
l'Idéal correspondant à $\Delta_g$, $\mathscr{L}$ est égal à
$(f\times_Z f)^*(\mathscr{K})\cdot\mathcal{O}_{X\times_Z X}$ (I, 4.4.5).
Comme $(\mathscr{J}^{n+1}+\mathscr{L})/\mathscr{J}^{n+1}$ est l'Idéal de
$\mathscr{P}_{X/Z}^n$ engendré par l'image de $\mathscr{L}$, on en déduit la
proposition.

\begin{corollary}[16.4.19]
\label{IV.16.4.19}
Avec les notations de (16.4.18), on a une suite exacte de
$\mathcal{O}_X$-Modules quasi-cohérents
\[
  \xymatrix{
    f^*(\Omega_{Y/Z}^1)\ar[r]^-{f_{X/Y/Z}}&
    \Omega_{X/Z}^1\ar[r]^-{g_{X/Y/Z}}&
    \Omega_{X/Y}^1\ar[r]&0.
  }
  \tag{16.4.19.1}
  \label{IV.16.4.19.1}
\]
\end{corollary}

Lorsque $X$, $Y$, $Z$ sont affines, on retrouve ainsi la suite exacte
(0, 20.5.7.1).

\begin{proposition}[16.4.20]
\label{IV.16.4.20}
Soient $f:Y\to Z$ un morphisme, $j:X\to Y$ une immersion fermée,
$\mathscr{K}$ l'Idéal quasi-cohérent de $\mathcal{O}_Y$ correspondant à $j$.
Alors on a
$\mathscr{P}_{X/Y}^n=\mathcal{O}_X=\mathcal{O}_Y/\mathscr{K}$,
l'homomorphisme canonique
$j_{X/Y/Z}:j^*(\mathscr{P}_{Y/Z}^n)\to\mathscr{P}_{X/Z}^n$ est surjectif, et
son noyau est l'Idéal de $j^*(\mathscr{P}_{Y/Z}^n)$ engendré par
$j^*(\mathcal{O}_Y\cdot d_{Y/Z}^n(\mathscr{K}))$ (il faut noter que
$d_{Y/Z}^n(\mathscr{K})$ est un sous-faisceau de groupes commutatifs de
$\mathscr{P}_{Y/Z}^n$, mais non un $\mathcal{O}_Y$-Module en général).
\end{proposition}

On sait (I, 5.3.8) que la diagonale
$\Delta_j:X\to X\times_Y X$ est un isomorphisme, d'où la première assertion.
Si $\varpi_1$ et $\varpi_2$ sont les deux homomorphismes d'Algèbres
$\mathcal{O}_Y\to\mathscr{P}_{Y/Z}^n$ correspondant respectivement aux deux
projections canoniques $p_1$, $p_2$ de $Y\times_Z Y$ sur $Y$, rappelons que
par définition (16.3.5 et 16.3.6) $\varpi_1$ est l'homomorphisme structural de
la $\mathcal{O}_Y$-Algèbre $\mathscr{P}_{Y/Z}^n$ et
$\varpi_2=d_{Y/Z}^n$. La $\mathcal{O}_X$-Algèbre
$j^*(\mathscr{P}_{Y/Z}^n)$ s'identifie donc à
$\mathscr{P}_{Y/Z}^n/\varpi_1(\mathscr{K})\mathscr{P}_{Y/Z}^n$ et son quotient
par l'Idéal engendré par $j^*(d_{Y/Z}^n(\mathscr{K}))$ à
\[
  \mathscr{P}_{Y/Z}^n/
  (\varpi_1(\mathscr{K})+\varpi_2(\mathscr{K}))\mathscr{P}_{Y/Z}^n.
\]
Notons maintenant qu'on a un diagramme commutatif
\oldpage[IV-4]{25}
\[
  \xymatrix{
    Y\ar[d]_-{\Delta_f}&X\ar[l]_-j\ar[d]^-{\Delta_{f\circ j}}\\
    Y\times_Z Y&X\times_Z X\ar[l]^-{j\times_Z j}
  }
\]
identifiant $X$ au produit des $(Y\times_Z Y)$-préschémas $Y$ et
$X\times_Z X$ (I, 5.3.7). Comme $j\times_Z j$ est une immersion, on déduit
donc de cette remarque et de (16.2.2) que si $\Delta_{Y/Z}^n$ et
$\Delta_{X/Z}^n$ désignent les voisinages infinitésimaux d'ordre $n$ de $Y$
et $X$ pour les immersions canoniques $\Delta_f$ et $\Delta_{f\circ j}$
respectivement, on a un diagramme
\[
  \xymatrix{
    \Delta_{Y/Z}^n\ar[d]&\Delta_{X/Z}^n\ar[l]\ar[d]\\
    Y\times_Z Y&X\times_Z X\ar[l]^-{j\times_Z j}
  }
\]
faisant de $\Delta_{X/Z}^n$ le produit des $(Y\times_Z Y)$-préschémas
$\Delta_{Y/Z}^n$ et $X\times_Z X$. On peut encore dire que
$\mathscr{P}_{X/Z}^n$ s'identifie au faisceau d'anneaux
$\mathscr{P}_{Y/Z}^n\otimes_{\mathcal{O}_{Y\times_Z Y}}
\mathcal{O}_{X\times_Z X}$. Mais on voit aussitôt (par exemple en se ramenant
au cas affine) que
\[
  \mathcal{O}_{X\times_Z X}=
  \mathcal{O}_{Y\times_Z Y}/
  (p_1^*(\mathscr{K})+p_2^*(\mathscr{K}))
  \mathcal{O}_{Y\times_Z Y}.
\]
Donc $\mathscr{P}_{X/Z}^n$ s'identifie au quotient de
$\mathscr{P}_{Y/Z}^n$ par l'Idéal engendré par l'image dans
$\mathscr{P}_{Y/Z}^n$ de $p_1^*(\mathscr{K})+p_2^*(\mathscr{K})$. Mais par
définition cet Idéal est aussi engendré par
$\varpi_1(\mathscr{K})+\varpi_2(\mathscr{K})$. C.Q.F.D.

\begin{corollary}[16.4.21]
\label{IV.16.4.21}
Soient $f:Y\to Z$ un morphisme, $j:X\to Y$ une immersion. On a une suite
exacte de $\mathcal{O}_X$-Modules quasi-cohérents
\[
  \xymatrix{
    \mathscr{N}_{X/Y}\ar[r]&j^*(\Omega_{Y/Z}^1)\ar[r]&
    \Omega_{X/Z}^1\ar[r]&0.
  }
  \tag{16.4.21.1}
  \label{IV.16.4.21.1}
\]
\end{corollary}

Lorsque $X$, $Y$ et $Z$ sont affines, on retrouve ainsi la suite exacte
(0, 20.5.12.1).

\begin{corollary}[16.4.22]
\label{IV.16.4.22}
Si $f:X\to S$ est un morphisme localement de présentation finie,
$\mathscr{P}_{X/S}^n$ et $\Omega_{X/S}^1$ sont des
$\mathcal{O}_X$-Modules quasi-cohérents de présentation finie.
\end{corollary}

On est aussitôt ramené au cas où $S=\operatorname{Spec}(A)$ est affine,
$X=\operatorname{Spec}(B)$, où
$B=A[T_1,\ldots,T_r]/\mathfrak{K}$, $\mathfrak{K}$ étant un idéal de type
fini de $C=A[T_1,\ldots,T_r]$. On applique alors (16.4.20) avec $Z=S$,
$Y=\operatorname{Spec}(C)$ et $\mathscr{K}=\widetilde{\mathfrak{K}}$. Alors
$j^*(\mathscr{P}_{Y/Z}^n)$ est un $\mathcal{O}_X$-Module libre de rang fini
(16.4.10) et l'hypothèse sur $\mathfrak{K}$ entraîne que
$j^*(\mathcal{O}_Y\cdot d_{Y/Z}^n(\mathscr{K}))$ engendre un
$\mathcal{O}_X$-Module quasi-cohérent de type fini; d'où la conclusion.

\begin{proposition}[16.4.23]
\label{IV.16.4.23}
Soient $X$, $Y$ deux $S$-préschémas, $Z=X\times_S Y$ leur produit,
$p:X\times_S Y\to X$ et $q:X\times_S Y\to Y$ les projections canoniques.
Alors l'homomorphisme canonique
\[
  p_{Z/X/S}\oplus q_{Z/Y/S}:
  p^*(\Omega_{X/S}^1)\oplus q^*(\Omega_{Y/S}^1)
  \longrightarrow\Omega_{(X\times_S Y)/S}^1
  \tag{16.4.23.1}
  \label{IV.16.4.23.1}
\]
est bijectif.
\end{proposition}

\oldpage[IV-4]{26}
Le diagramme commutatif
\[
  \xymatrix{
    Y\ar[d]_-g&X\times_S Y\ar[l]_-q\ar[d]_-h&
      X\times_S Y\ar[l]_-{\operatorname{id}}\ar[d]^-p\\
    S&S\ar[l]^-{\operatorname{id}}&X\ar[l]^-f
  }
\]
donne une factorisation de l'\emph{isomorphisme} canonique $P^n(p)$ (16.4.5)
\[
  p^*(\mathscr{P}_{X/S}^n)\longrightarrow
  \mathscr{P}_{Z/S}^n\longrightarrow\mathscr{P}_{Z/Y}^n
\]
et de même, en intervertissant les rôles de $X$ et $Y$, on a une factorisation
de l'\emph{isomorphisme} $P^n(q)$
\[
  q^*(\mathscr{P}_{Y/S}^n)\longrightarrow
  \mathscr{P}_{Z/S}^n\longrightarrow\mathscr{P}_{Z/X}^n.
\]
Ceci prouve que l'homomorphisme canonique (16.4.18.1)
\[
  p_{Z/X/S}:p^*(\mathscr{P}_{X/S}^n)\longrightarrow\mathscr{P}_{Z/S}^n
  \quad\text{(resp.\
  $q_{Z/Y/S}:q^*(\mathscr{P}_{Y/S}^n)\to\mathscr{P}_{Z/S}^n$)}
\]
est \emph{injectif}, et que le noyau de l'homomorphisme canonique surjectif
(16.4.18.2)
\[
  \mathscr{P}_{Z/S}^n\longrightarrow\mathscr{P}_{Z/Y}^n
  \quad\text{(resp.\
  $\mathscr{P}_{Z/S}^n\to\mathscr{P}_{Z/X}^n$)}
\]
est supplémentaire de l'image de $p_{Z/X/S}$ (resp.\ $q_{Z/Y/S}$). Mais
d'autre part, ce noyau est, en vertu de (16.4.18), engendré par l'image par
$q_{Z/Y/S}$ (resp.\ $p_{Z/X/S}$) de l'Idéal d'augmentation de
$q^*(\mathscr{P}_{Y/S}^n)$ (resp.\ $p^*(\mathscr{P}_{X/S}^n)$). On en conclut
la proposition en considérant le cas $n=1$.

On généralise aussitôt (16.4.23) au cas d'un produit d'un nombre fini
quelconque de $S$-préschémas.

\begin{remarks}[16.4.24]
\label{IV.16.4.24}
\begin{enumerate}
\item[(i)] Nous verrons (17.2.3) que lorsque le morphisme $f:X\to Y$ dans
(16.4.18) est \emph{lisse}, l'homomorphisme $f_{X/Y/Z}$ dans (16.4.19.1) est
localement \emph{inversible à gauche} et en particulier injectif. De même,
lorsque le morphisme $f\circ j:X\to Z$ de (16.4.20) est \emph{lisse},
l'homomorphisme de gauche dans (16.4.21.2) est localement
\emph{inversible à gauche} et \emph{a fortiori} injectif (17.2.5). Au chap. V,
nous donnerons aussi une variante, dans le cas des Modules sur les
préschémas, des «modules d'imperfection» étudiés dans (0, 20.6), et des suites
exactes où ils figurent.

\item[(ii)] Soient $X$ un espace topologique, $\mathscr{A}$ un faisceau
d'anneaux sur $X$ et $\mathscr{B}$ une $\mathscr{A}$-Algèbre sur $X$. Alors il
est clair que
\[
  U\longmapsto
  P_{\Gamma(U,\mathscr{B})/\Gamma(U,\mathscr{A})}^n
  \qquad (U\text{ ouvert dans }X)
\]
est un préfaisceau de $\Gamma(U,\mathscr{B})$-algèbres augmentées, donc le
faisceau associé $\mathscr{P}_{\mathscr{B}/\mathscr{A}}^n$ est une
$\mathscr{B}$-Algèbre augmentée. Dans le cas particulier où $X$ est un
préschéma, $f=(\psi,\theta):X\to S$ un morphisme de préschémas, il résulte
aisément de (16.4.16) et de l'exactitude du foncteur $\varinjlim$ que
$\mathscr{P}_{X/S}^n$ est canoniquement isomorphe à
$\mathscr{P}_{\mathcal{O}_X/\psi^*(\mathcal{O}_S)}^n$. Il s'ensuit que le
formalisme développé dans le présent paragraphe pourrait être considéré comme
un
\oldpage[IV-4]{27}
cas particulier d'un formalisme différentiel pour des espaces annelés munis
d'un faisceau d'algèbres sur le faisceau structural. Cependant, nous n'avons
pas voulu partir de ce point de vue, moins intuitif et moins commode pour les
applications. Il semble d'ailleurs que, pour les diverses espèces de
«variétés», la construction «globale» des $\mathscr{P}^n$ analogue à celle que
nous utilisons ici soit également mieux adaptée aux applications.
\end{enumerate}
\end{remarks}

\subsection{Faisceaux et fibrés tangents relatifs; dérivations.}
\EGAReaderTarget{ega4.c04.subsec.16-5}{IV.16.5}

\begin{env}[16.5.1]
\EGAReaderTarget{ega4.c04.statement.16-5-1}{IV.16.5.1}
Soit $f=(\psi,\theta):X\to S$ un morphisme d'espaces annelés. Pour tout
$\mathcal{O}_X$-Module $\mathscr{F}$, on appelle \emph{$S$-dérivation}
(ou \emph{$(X/S)$-dérivation}, ou \emph{$f$-dérivation}) \emph{de
$\mathcal{O}_X$ dans $\mathscr{F}$} tout homomorphisme de \emph{faisceaux de
groupes additifs} $D:\mathcal{O}_X\to\mathscr{F}$, vérifiant les conditions
suivantes:
\begin{enumerate}
\item[(a)] pour tout ouvert $V$ de $X$, et tout couple de sections
$(t_1,t_2)$ de $\mathcal{O}_X$ au-dessus de $V$, on a
\[
  D(t_1t_2)=t_1D(t_2)+D(t_1)t_2;
  \tag{16.5.1.1}
  \label{IV.16.5.1.1}
\]
\item[(b)] pour tout ouvert $V$ de $X$, toute section $t$ de $\mathcal{O}_X$
au-dessus de $V$ et toute section $s$ de $\mathcal{O}_S$ au-dessus d'un ouvert
$U$ de $S$ tel que $V\subset f^{-1}(U)$, on a
\[
  D((s|V)t)=(s|V)D(t).
  \tag{16.5.1.2}
  \label{IV.16.5.1.2}
\]
\end{enumerate}
Il est clair qu'il revient au même de dire que pour tout $x\in X$,
l'homomorphisme de groupes additifs
$D_x:\mathcal{O}_x\to\mathscr{F}_x$ est une
$\mathcal{O}_{f(x)}$-dérivation.

Une autre interprétation consiste à considérer la $\mathcal{O}_X$-Algèbre
$\mathscr{D}_{\mathcal{O}_X}(\mathscr{F})$ égale à
$\mathcal{O}_X\oplus\mathscr{F}$, la structure d'Algèbre étant définie par la
condition que pour tout ouvert $V$ de $X$, le produit de deux sections de
$\mathcal{O}_X$ (resp. d'une section de $\mathcal{O}_X$ et d'une section de
$\mathscr{F}$) au-dessus de $V$ est défini par la structure d'anneau de
$\Gamma(V,\mathcal{O}_X)$ (resp. la structure de
$\Gamma(V,\mathcal{O}_X)$-module sur $\Gamma(V,\mathscr{F})$), et le produit de
deux sections de $\mathscr{F}$ au-dessus de $V$ est pris égal à 0; alors,
$\mathscr{F}$ est un Idéal de
$\mathscr{D}_{\mathcal{O}_X}(\mathscr{F})$, noyau de l'augmentation canonique
$\mathscr{D}_{\mathcal{O}_X}(\mathscr{F})\to\mathcal{O}_X$, et dire que $D$
est une $S$-dérivation de $\mathcal{O}_X$ dans $\mathscr{F}$ signifie que
$1_{\mathcal{O}_X}+D$ est un $\mathcal{O}_S$-homomorphisme d'Algèbres de
$\mathcal{O}_X$ dans $\mathscr{D}_{\mathcal{O}_X}(\mathscr{F})$, qui, composé
avec l'augmentation, donne $1_{\mathcal{O}_X}$.

Les $S$-dérivations de $\mathcal{O}_X$ dans $\mathscr{F}$ forment évidemment
un $\Gamma(X,\mathcal{O}_X)$-module
$\operatorname{D\acute{e}r}_S(\mathcal{O}_X,\mathscr{F})$.

Lorsque $\mathscr{F}=\mathcal{O}_X$, une $S$-dérivation de $\mathcal{O}_X$ dans
lui-même est simplement appelée une \emph{$S$-dérivation de $\mathcal{O}_X$}.
\end{env}

\begin{proposition}[16.5.2]
\label{IV.16.5.2}
Soient $A$ un anneau, $B$ une $A$-algèbre, $L$ un $B$-module; on pose
$S=\operatorname{Spec}(A)$, $X=\operatorname{Spec}(B)$,
$\mathscr{F}=\widetilde{L}$. Alors l'application $D\mapsto\Gamma(D)$ qui, à
toute $S$-dérivation $D$ de $\mathcal{O}_X$ dans $\mathscr{F}$ fait
correspondre l'application $\Gamma(D):t\mapsto D(t)$ de $B$ dans $L$, est un
isomorphisme de $B$-modules de
$\operatorname{D\acute{e}r}_S(\mathcal{O}_X,\mathscr{F})$ sur
$\operatorname{D\acute{e}r}_A(B,L)$ (cf. (0, 20.1.2)).
\end{proposition}

Cela résulte aussitôt de l'interprétation donnée ci-dessus des $S$-dérivations
en termes
\oldpage[IV-4]{28}
d'homomorphismes d'Algèbres, de l'interprétation analogue donnée dans
(0, 20.1.6), et de la correspondance canonique entre homomorphismes de
$\mathcal{O}_X$-Algèbres et homomorphismes de $B$-algèbres (I, 1.3.13 et
1.3.8).

\begin{proposition}[16.5.3]
\label{IV.16.5.3}
Soit $f=(\psi,\theta):X\to S$ un morphisme de préschémas.
\begin{enumerate}
\item[(i)] La différentielle
$d_{X/S}:\mathcal{O}_X\to\Omega_{X/S}^1$ (16.3.6) est une $S$-dérivation.
\item[(ii)] Pour tout $\mathcal{O}_X$-Module $\mathscr{F}$, l'application
$u\mapsto u\circ d_{X/S}$ est un isomorphisme de
$\Gamma(X,\mathcal{O}_X)$-modules
\[
  \operatorname{Hom}_{\mathcal{O}_X}(\Omega_{X/S}^1,\mathscr{F})
  \xrightarrow{\sim}
  \operatorname{D\acute{e}r}_S(\mathcal{O}_X,\mathscr{F}).
  \tag{16.5.3.1}
  \label{IV.16.5.3.1}
\]
\end{enumerate}
\end{proposition}

L'assertion (i) a déjà été notée (16.3.6). D'autre part, il est immédiat (en
vertu de (0, 20.4.8)) que $u\mapsto u\circ d_{X/S}$ est injectif, en
considérant les restrictions à une fibre $\mathcal{O}_x$ des deux membres et
utilisant (16.4.16.2). Pour voir que l'homomorphisme (16.5.3.1) est surjectif,
considérons une $S$-dérivation $D:\mathcal{O}_X\to\mathscr{F}$; pour tout
ouvert affine $V=\operatorname{Spec}(B)$ de $X$, tel que $f(V)$ soit contenu
dans un ouvert affine $U=\operatorname{Spec}(A)$ de $S$,
$D_V:B\to\Gamma(V,\mathscr{F})$ est une $A$-dérivation, donc il existe un
\emph{unique} $B$-homomorphisme
$u_V:\Omega_{B/A}^1\to\Gamma(V,\mathscr{F})$ tel que
$D_V=u_V\circ d_{B/A}$ (0, 20.4.8); en outre l'unicité de $u_V$ montre
aussitôt que pour tout ouvert affine $W\subset V$, on a $u_W=u_V|W$, donc les
$u_V$ définissent un homomorphisme de $\mathcal{O}_X$-Modules
$u:\mathcal{O}_X\to\mathscr{F}$ répondant à la question.

\begin{env}[16.5.4]
\EGAReaderTarget{ega4.c04.statement.16-5-4}{IV.16.5.4}
Avec les notations de (16.5.1), pour tout ouvert $U$ de $X$,
$\operatorname{D\acute{e}r}_S(\mathcal{O}_U,\mathscr{F}|U)$ est un
$\Gamma(U,\mathcal{O}_X)$-module et il est clair que l'application
$U\mapsto\operatorname{D\acute{e}r}_S(\mathcal{O}_U,\mathscr{F}|U)$ est un
préfaisceau; en fait, c'est même un \emph{faisceau} (donc un
$\mathcal{O}_X$-Module), en vertu de la caractérisation ponctuelle des
$S$-dérivations, vue en (16.5.1). Cet $\mathcal{O}_X$-Module se note
$\mathscr{D}\!\mathrm{\acute{e}r}_S(\mathcal{O}_X,\mathscr{F})$ et est appelé
\emph{faisceau des $S$-dérivations de $\mathcal{O}_X$ dans $\mathscr{F}$}, et
ce qu'on vient de voir s'exprime encore par le corollaire suivant:
\end{env}

\begin{corollary}[16.5.5]
\label{IV.16.5.5}
Pour tout $\mathcal{O}_X$-Module $\mathscr{F}$, l'homomorphisme de
$\mathcal{O}_X$-Modules déduit de $u\mapsto u\circ d_{X/S}$
\[
  \mathscr{H}\!om_{\mathcal{O}_X}(\Omega_{X/S}^1,\mathscr{F})
  \longrightarrow
  \mathscr{D}\!\mathrm{\acute{e}r}_S(\mathcal{O}_X,\mathscr{F})
  \tag{16.5.5.1}
  \label{IV.16.5.5.1}
\]
est bijectif.
\end{corollary}

\begin{corollary}[16.5.6]
\label{IV.16.5.6}
\begin{enumerate}
\item[(i)] Si le morphisme $f:X\to S$ est localement de présentation finie et
si $\mathscr{F}$ est un $\mathcal{O}_X$-Module quasi-cohérent, alors
$\mathscr{D}\!\mathrm{\acute{e}r}_S(\mathcal{O}_X,\mathscr{F})$ est un
$\mathcal{O}_X$-Module quasi-cohérent.
\item[(ii)] Si de plus $S$ est localement noethérien et si $\mathscr{F}$ est
cohérent, alors $\mathscr{D}\!\mathrm{\acute{e}r}_S(\mathcal{O}_X,\mathscr{F})$ est un
$\mathcal{O}_X$-Module cohérent.
\end{enumerate}
\end{corollary}

L'assertion (i) résulte de l'isomorphisme (16.5.5.1), de (16.4.22) et
(I, 1.3.12); l'assertion (ii) résulte de (0$_{\rm I}$, 5.3.5).

\begin{env}[16.5.7]
\EGAReaderTarget{ega4.c04.statement.16-5-7}{IV.16.5.7}
On pose
\[
  \mathfrak{G}_{X/S}=
  \mathscr{H}\!om_{\mathcal{O}_X}(\Omega_{X/S}^1,\mathcal{O}_X)
  =\mathscr{D}\!\mathrm{\acute{e}r}_S(\mathcal{O}_X,\mathcal{O}_X)
  \tag{16.5.7.1}
  \label{IV.16.5.7.1}
\]
et on dit que c'est le \emph{faisceau des $S$-dérivations de
$\mathcal{O}_X$}, ou encore le \emph{faisceau tangent de $X$ relativement à
$S$}: c'est donc le \emph{dual} du $\mathcal{O}_X$-Module
$\Omega_{X/S}^1$. Si $f$ est localement de présentation
\oldpage[IV-4]{29}
finie, $\mathfrak{G}_{X/S}$ est un $\mathcal{O}_X$-Module quasi-cohérent; si
de plus $S$ est localement noethérien, $\mathfrak{G}_{X/S}$ est cohérent
(16.5.6).
\end{env}

\begin{env}[16.5.8]
\label{IV.16.5.8}
Supposons plus particulièrement que $\Omega_{X/S}^1$ soit un
$\mathcal{O}_X$-Module \emph{localement libre} (de rang fini) (ce qui sera le
cas lorsque $f$ est \emph{lisse} (17.2.3)); alors $\mathfrak{G}_{X/S}$ est un
$\mathcal{O}_X$-Module localement libre et a même rang que $\Omega_{X/S}^1$
en chaque point. De façon plus précise, supposons que $\Omega_{X/S}^1$ soit de
rang $n$ en un point $x$; il y a alors $n$ sections $s_i$ ($1\leq i\leq n$)
de $\mathcal{O}_X$ au-dessus d'un voisinage affine $U$ de $x$ telles que les
images canoniques des $ds_i$ dans
$\Omega_{X/S}^1\otimes_{\mathcal{O}_X}k(x)$ forment une base de ce
$k(x)$-espace vectoriel; en vertu du lemme de Nakayama, les germes
$(ds_i)_x=d(s_i)_x$ des $ds_i$ au point $x$ forment une base du
$\mathcal{O}_x$-module $(\Omega_{X/S}^1)_x$, donc, en restreignant $U$, on
peut supposer que les $ds_i$ forment une \emph{base} du
$\Gamma(U,\mathcal{O}_X)$-module $\Gamma(U,\Omega_{X/S}^1)$. Alors le
$\Gamma(U,\mathcal{O}_X)$-module $\Gamma(U,\mathfrak{G}_{X/S})$ est dual du
précédent; on note $(D_i)_{1\leq i\leq n}$ ou
$\left(\frac{\partial}{\partial s_i}\right)_{1\leq i\leq n}$ la base duale de
$(ds_i)_{1\leq i\leq n}$, de sorte que, par (16.5.3), on a
\[
  D_is_j=\langle D_i,ds_j\rangle
  =\left\langle\frac{\partial}{\partial s_i},ds_j\right\rangle
  =\delta_{ij}\qquad\text{(indice de Kronecker).}
  \tag{16.5.8.1}
  \label{IV.16.5.8.1}
\]
Toute $\Gamma(S,\mathcal{O}_S)$-dérivation de la
$\Gamma(S,\mathcal{O}_S)$-algèbre $\Gamma(U,\mathcal{O}_X)$ s'écrit donc
d'une seule manière
\[
  D=\sum_{i=1}^n a_iD_i=\sum_{i=1}^n a_i\frac{\partial}{\partial s_i}
\]
où les $a_i$ ($1\leq i\leq n$) sont des sections de $\mathcal{O}_X$ au-dessus
de $U$. Pour toute section $g\in\Gamma(U,\mathcal{O}_X)$, si l'on pose
$dg=\sum_{i=1}^n c_i ds_i$, on a
$c_i=\langle D_i,dg\rangle=D_ig$ en vertu de (16.5.8.1), autrement dit
\[
  dg=\sum_{i=1}^n(D_ig)ds_i
  =\sum_{i=1}^n\frac{\partial g}{\partial s_i}ds_i.
  \tag{16.5.8.2}
  \label{IV.16.5.8.2}
\]
\end{env}

\begin{env}[16.5.9]
\label{IV.16.5.9}
Soient $D_1$, $D_2$ deux $S$-dérivations de $\mathcal{O}_X$. Pour tout ouvert
$U$ de $X$, si $D_1^U$, $D_2^U$ sont les dérivations correspondantes de
l'anneau $\Gamma(U,\mathcal{O}_X)$, le \emph{crochet}
\[
  [D_1^U,D_2^U]=D_1^U\circ D_2^U-D_2^U\circ D_1^U
\]
est aussi une dérivation de cet anneau, donc le
$\psi^*(\mathcal{O}_S)$-endomorphisme de $\mathcal{O}_X$
\[
  [D_1,D_2]=D_1\circ D_2-D_2\circ D_1
  \tag{16.5.9.1}
  \label{IV.16.5.9.1}
\]
est encore une $S$-dérivation; comme on vérifie aussitôt que ce crochet
vérifie l'identité de Jacobi, on voit qu'on a ainsi défini sur
$\operatorname{D\acute{e}r}_S(\mathcal{O}_X,\mathcal{O}_X)$ une structure de
$\Gamma(S,\mathcal{O}_S)$-algèbre de Lie. Comme la définition de cette
structure commute à la restriction à un ouvert de $X$, on voit que
$\mathfrak{G}_{X/S}$ est canoniquement muni d'une structure de
$\psi^*(\mathcal{O}_S)$-Algèbre de Lie. On notera que l'application
$(D_1,D_2)\mapsto[D_1,D_2]$ n'est pas
$\Gamma(X,\mathcal{O}_X)$-bilinéaire.
\end{env}

\begin{env}[16.5.10]
\label{IV.16.5.10}
Pour tout changement de base $g:S'\to S$, si l'on pose $X'=X\times_S S'$, on
a vu (16.4.5) que l'on a un isomorphisme canonique
\[
  \Omega_{X/S}^1\otimes_S S'\xrightarrow{\sim}\Omega_{X'/S'}^1
  \tag{16.5.10.1}
  \label{IV.16.5.10.1}
\]
\oldpage[IV-4]{30}
d'où l'on déduit, en vertu de (16.5.10.1), un homomorphisme canonique
(Bourbaki, \emph{Alg.}, chap. II, 3\textsuperscript{e} éd., \S~5,
n\textsuperscript{o} 3)
\[
  \mathfrak{G}_{X/S}\otimes_{\mathcal{O}_S}\mathcal{O}_{S'}
  \longrightarrow\mathfrak{G}_{X'/S'}
  \tag{16.5.10.2}
  \label{IV.16.5.10.2}
\]
qui en général n'est ni injectif ni surjectif. Toutefois:
\end{env}

\begin{proposition}[16.5.11]
\label{IV.16.5.11}
\begin{enumerate}
\item[(i)] Si $g:S'\to S$ est un morphisme plat et si $f$ est localement de
type fini (resp.\ localement de présentation finie), l'homomorphisme
(16.5.10.2) est injectif (resp.\ bijectif).
\item[(ii)] Si $\Omega_{X/S}^1$ est un $\mathcal{O}_X$-Module localement libre
de type fini, l'homomorphisme (16.5.10.2) est bijectif.
\end{enumerate}
\end{proposition}

En effet, l'assertion (ii) résulte de Bourbaki, \emph{Alg.}, chap. II,
3\textsuperscript{e} éd., \S~5, n\textsuperscript{o} 3, prop. 7. L'assertion
(i) résulte de même de Bourbaki, \emph{Alg. comm.}, chap. I, \S~2,
n\textsuperscript{o} 10, prop. 11 et du fait que si $f$ est localement de type
fini (resp.\ localement de présentation finie), $\Omega_{X/S}^1$ est un
$\mathcal{O}_X$-Module de type fini (resp.\ de présentation finie) ((16.3.9)
et (16.4.22)).

\begin{env}[16.5.12]
\EGAReaderTarget{ega4.c04.statement.16-5-12}{IV.16.5.12}
Puisque $\Omega_{X/S}^1$ est un $\mathcal{O}_X$-Module quasi-cohérent, on peut
considérer le fibré vectoriel sur $X$ défini par $\Omega_{X/S}^1$
(II, 1.7.8)
\[
  T_{X/S}=\mathbf{V}(\Omega_{X/S}^1)
  \tag{16.5.12.1}
  \label{IV.16.5.12.1}
\]
qu'on appelle \emph{fibré tangent de $X$ relativement à $S$}. On a donc une
bijection canonique (II, 1.7.9)
\[
  \Gamma(T_{X/S}/S)\xrightarrow{\sim}
  \operatorname{Hom}_{\mathcal{O}_X}(\Omega_{X/S}^1,\mathcal{O}_X)
  =\Gamma(X,\mathfrak{G}_{X/S})
\]
par définition de $\mathfrak{G}_{X/S}$, et on peut dans cet isomorphisme
remplacer $X$ par un ouvert quelconque $U$ de $X$; on peut donc dire que le
\emph{faisceau tangent} de $X$ relativement à $S$ est isomorphe au
\emph{faisceau des germes de $S$-sections} du fibré tangent de $X$
relativement à $S$. Si $f:X\to Y$ est un $S$-morphisme, on a vu (16.4.19)
qu'on a un homomorphisme canonique
$f_{X/Y/S}:f^*(\Omega_{Y/S}^1)\to\Omega_{X/S}^1$, ce qui donne, compte tenu de
ce que
\[
  \mathbf{V}(f^*(\Omega_{Y/S}^1))
  =\mathbf{V}(\Omega_{Y/S}^1)\times_Y X\qquad\text{(II, 1.7.11),}
\]
un $X$-morphisme
$T_{X/S}(f):T_{X/S}\to T_{Y/S}\times_Y X$. Si $g:Y\to Z$ est un second
$S$-morphisme, on a
$T_{X/S}(g\circ f)=(T_{Y/S}(g)\times1_X)\circ T_{X/S}(f)$ (0, 20.5.4.1).

Il résulte de (16.5.10.1) et de (II, 1.7.11) que pour tout changement de base
$g:S'\to S$, on a un isomorphisme canonique
\[
  T_{X'/S'}\xrightarrow{\sim}T_{X/S}\times_S S'
  =T_{X/S}\times_X X'.
  \tag{16.5.12.2}
  \label{IV.16.5.12.2}
\]
\end{env}

\begin{env}[16.5.13]
\EGAReaderTarget{ega4.c04.statement.16-5-13}{IV.16.5.13}
Pour tout point $x\in X$, on appelle \emph{espace tangent à $X$ au point $x$}
(relativement à $S$) l'ensemble des points de la fibre
$T_{X/S}\times_X\operatorname{Spec}(k(x))$ \emph{rationnels sur $k(x)$}, donc
l'ensemble
\[
  T_{X/S}(x)=
  \operatorname{Hom}_{k(x)}
  (\Omega_{X/S}^1\otimes_{\mathcal{O}_x}k(x),k(x))
  \tag{16.5.13.1}
  \label{IV.16.5.13.1}
\]
qui est le \emph{dual} du $k(x)$-espace vectoriel
$\Omega_{\mathcal{O}_x/\mathcal{O}_s}^1/
\mathfrak{m}_x\cdot\Omega_{\mathcal{O}_x/\mathcal{O}_s}^1$. Lorsque
$\Omega_{X/S}^1$ est un $\mathcal{O}_X$-Module de \emph{type fini},
$T_{X/S}(x)$ est donc un espace vectoriel de rang fini sur $k(x)$, et pour
tout
\oldpage[IV-4]{31}
changement de base $g:S'\to S$, et tout point
$x'\in X'=X\times_S S'$ au-dessus de $x$, on a un isomorphisme canonique
\[
  T_{X'/S'}(x')\xrightarrow{\sim}
  T_{X/S}(x)\otimes_{k(x)}k(x').
  \tag{16.5.13.2}
  \label{IV.16.5.13.2}
\]

Si $x$ est \emph{rationnel sur $k(s)$}, où $s=f(x)$ (de sorte que
$k(s)\to k(x)$ est un isomorphisme), il résulte de (16.4.13) que l'on a un
isomorphisme canonique
\[
  T_{X/S}(x)=T_{X_s/k(s)}(x)
  =\operatorname{Hom}_{k(s)}
  (\mathfrak{m}'_x/\mathfrak{m}'_x{}^2,k(s))
  \tag{16.5.13.3}
  \label{IV.16.5.13.3}
\]
où $\mathfrak{m}'_x$ est l'idéal maximal de
$\mathcal{O}_{X_s,x}=\mathcal{O}_{X,x}/
\mathfrak{m}_s\mathcal{O}_{X,x}$. Dans le cas où $S$ est le spectre d'un corps
$k$, on retrouve ainsi la définition de Zariski de l'espace tangent en un
point $x\in X$ \emph{rationnel sur $k$}, comme dual de
$\mathfrak{m}_x/\mathfrak{m}_x^2$.

Soit $Y$ un second $S$-préschéma et soit $g:Y\to X$ un $S$-morphisme; on a
alors un homomorphisme canonique de $\mathcal{O}_Y$-Modules (16.4.19)
\[
  g_{Y/X/S}:g^*(\Omega_{X/S}^1)\longrightarrow\Omega_{Y/S}^1.
  \tag{16.5.13.4}
  \label{IV.16.5.13.4}
\]
Notons maintenant que si $y\in Y$ et $x=g(y)$, on a
\[
  g^*(\Omega_{X/S}^1)\otimes_{\mathcal{O}_Y}k(y)
  =(\Omega_{X/S}^1\otimes_{\mathcal{O}_X}k(x))
  \otimes_{k(x)}k(y)
\]
et par suite, si $\Omega_{X/S}^1$ est un $\mathcal{O}_X$-Module
\emph{de type fini}, on peut identifier
\[
  \operatorname{Hom}_{k(y)}
  (g^*(\Omega_{X/S}^1)\otimes_{\mathcal{O}_Y}k(y),k(y))
\]
à $T_{X/S}(x)\otimes_{k(x)}k(y)$. On déduit donc alors de l'homomorphisme
(16.5.13.4) un homomorphisme de \emph{$k(y)$-espaces vectoriels}
\[
  T_y(g):T_{Y/S}(y)\longrightarrow
  T_{X/S}(x)\otimes_{k(x)}k(y)
  \tag{16.5.13.5}
  \label{IV.16.5.13.5}
\]
dit \emph{application linéaire tangente à $g$ au point $y$}. Lorsque $y$ est
\emph{rationnel sur $k(s)$}, on peut identifier $k(s)$, $k(y)$ et $k(x)$ et
$T_y(g)$ est alors un homomorphisme de $k(s)$-espaces vectoriels
$T_{Y/S}(y)\to T_{X/S}(x)$; on notera d'ailleurs que dans ce cas,
$g^*(\Omega_{X/S}^1)\otimes_{\mathcal{O}_Y}k(y)$ s'identifie à
$\Omega_{X/S}^1\otimes_{\mathcal{O}_X}k(x)$, et l'homomorphisme précédent est
donc défini sans aucune condition de finitude sur $\Omega_{X/S}^1$ et n'est
autre que l'homomorphisme $T_{Y/S}(g)$ (16.5.12) restreint à la fibre au point
$y$ de $T_{Y/S}$.
\end{env}

\begin{env}[16.5.14]
\label{IV.16.5.14}
L'interprétation des dérivations d'une $A$-algèbre $B$ dans un $B$-module
$L$, donnée dans (0, 20.1.1), se traduit dans le langage des préschémas de la
façon suivante.

Considérons deux morphismes de préschémas $f:X\to S$, $g:Y\to S$, et un
sous-préschéma fermé $Y_0$ de $Y$ défini par un Idéal de carré nul
$\mathscr{J}$ de $\mathcal{O}_Y$ (de sorte que $Y$ et $Y_0$ ont
\emph{même espace topologique sous-jacent}). Supposons donné un $S$-morphisme
$u_0:Y_0\to X$, de sorte qu'on a un diagramme commutatif
\[
  \xymatrix{
    X\ar[d]_-f & Y_0\ar[l]_-{u_0}\ar[d]^j\\
    S & Y\ar@{-->}[ul]_-u\ar[l]^-g
  }
  \tag{16.5.14.1}
  \label{IV.16.5.14.1}
\]
\oldpage[IV-4]{32}
et proposons-nous de chercher s'il existe des \emph{$S$-morphismes}
$u:Y\to X$ tels que $u_0=u\circ j$ (autrement dit, s'il est possible de
compléter le diagramme précédent par la flèche pointillée $u$, de façon à le
laisser \emph{commutatif}).

Pour cela, considérons un ouvert affine $U=\operatorname{Spec}(C)$ de $Y$;
son image réciproque $j^{-1}(U)$ est l'ouvert affine
$U_0=\operatorname{Spec}(C/\mathfrak{L})$, où
$\mathfrak{L}=\Gamma(U,\mathscr{J})$, idéal \emph{de carré nul} dans $C$;
nous supposerons $U$ assez petit pour que $u_0(U_0)$ soit contenu dans un
ouvert affine $V=\operatorname{Spec}(B)$ de $X$, et
$g(U)=f(u_0(U_0))$ contenu dans un ouvert affine
$W=\operatorname{Spec}(A)$ de $S$, de sorte que $B$ et $C$ sont des
$A$-algèbres et que $u_0|U_0$ correspond à un $A$-homomorphisme $\psi$ de
$B$ dans $C/\mathfrak{L}$; soit $P(U_0)$ l'ensemble des restrictions $u|U$
des homomorphismes cherchés, qui correspondent canoniquement aux
$A$-homomorphismes d'algèbres $\varphi:B\to C$ tels que le \emph{composé}
\[
  B\xrightarrow{\varphi}C\longrightarrow C/\mathfrak{L}
\]
\emph{soit égal à $\psi$}. On sait donc (0, 20.1.1) que l'ensemble de ces
homomorphismes est vide ou de la forme
$\varphi_1+\operatorname{D\acute{e}r}_A(B,\mathfrak{L})$; lorsque $P(U_0)$ n'est pas
vide, le groupe additif $\operatorname{D\acute{e}r}_A(B,\mathfrak{L})$ opère par
addition dans $P(U_0)$, qui est donc un \emph{espace affine} pour le groupe
additif $\operatorname{D\acute{e}r}_A(B,\mathfrak{L})$ (ou encore un
\emph{espace homogène principal} (ou \emph{torseur}) \emph{sous}
$\operatorname{D\acute{e}r}_A(B,\mathfrak{L})$).

Remarquons maintenant que, puisque $\mathfrak{L}$ est muni d'une structure de
$B$-module au moyen de $\psi$, on a un \emph{isomorphisme}
$v\mapsto v\circ d_{B/A}$ de
$\operatorname{Hom}_B(\Omega_{B/A}^1,\mathfrak{L})$ sur
$\operatorname{D\acute{e}r}_A(B,\mathfrak{L})$ (0, 20.4.8). D'ailleurs, comme
$\mathfrak{L}$ est de carré nul, donc un $(C/\mathfrak{L})$-module, tout
$B$-homomorphisme $v:\Omega_{B/A}^1\to\mathfrak{L}$ peut être considéré
comme un $(C/\mathfrak{L})$-homomorphisme
$\Omega_{B/A}^1\otimes_B(C/\mathfrak{L})\to\mathfrak{L}$. Comme
$\mathscr{J}$ est de carré nul, il peut être considéré comme un
$\mathcal{O}_{Y_0}$-Module quasi-cohérent; introduisons le
$\mathcal{O}_{Y_0}$-Module
\[
  \mathscr{G}=\mathscr{H}om(u_0^*(\Omega_{X/S}^1),\mathscr{J});
  \tag{16.5.14.2}
  \label{IV.16.5.14.2}
\]
il résulte alors du fait que
$\Omega_{B/A}^1=\Gamma(V,\Omega_{X/S}^1)$ (16.3.7) que l'on peut écrire
$\operatorname{D\acute{e}r}_A(B,\mathfrak{L})=\Gamma(U_0,\mathscr{G})$.

Comme $P(U_0)$ est défini comme ensemble de $S$-morphismes $U\to X$, il est
clair que $U_0\mapsto P(U_0)$ est un \emph{faisceau d'ensembles}
$\mathscr{P}$ sur $Y_0$. Utilisons ce fait pour prouver que l'application
\[
  h:\Gamma(U_0,\mathscr{G})\times P(U_0)\longrightarrow P(U_0)
\]
définissant la structure de torseur sur $P(U_0)$ est aussi indépendante du
choix de $V$ et $W$, et en outre que, si $U'\subset U$ est un second ouvert
affine de $Y$, $U'_0$ son image réciproque dans $Y_0$, le diagramme
\[
  \xymatrix{
    \Gamma(U_0,\mathscr{G})\times P(U_0)\ar[d]\ar[r]^h
      & P(U_0)\ar[d]\\
    \Gamma(U'_0,\mathscr{G})\times P(U'_0)\ar[r]_{h'}
      & P(U'_0)
  }
  \tag{16.5.14.3}
  \label{IV.16.5.14.3}
\]
est commutatif (les flèches verticales étant les opérateurs de restriction).
En vertu de la remarque précédente, on est ramené à prouver la commutativité
du diagramme précédent lorsque $h$ est défini comme ci-dessus à partir des
ouverts affines $V$, $W$ et $h'$ à partir d'ouverts
\oldpage[IV-4]{33}
affines $V'\subset V$ et $W'\subset W$. Mais en vertu de la description
précédente de $h$, cela résulte de la commutativité du diagramme (0, 20.5.4.2).

Les applications
$\Gamma(U_0,\mathscr{G})\times P(U_0)\to P(U_0)$ définissent donc un
homomorphisme de \emph{faisceaux d'ensembles}
$m:\mathscr{G}\times\mathscr{P}\to\mathscr{P}$ tel que, pour tout ouvert
$U_0$ pour lequel $\Gamma(U_0,\mathscr{P})\neq\varnothing$,
\[
  m_{U_0}:\Gamma(U_0,\mathscr{G})\times\Gamma(U_0,\mathscr{P})
  \longrightarrow\Gamma(U_0,\mathscr{P})
\]
soit une loi externe définissant sur $\Gamma(U_0,\mathscr{P})$ une structure
de torseur sous le groupe $\Gamma(U_0,\mathscr{G})$.
\end{env}

\begin{env}[16.5.15]
\EGAReaderTarget{ega4.c04.statement.16-5-15}{IV.16.5.15}
D'une façon générale, lorsqu'on se donne sur un espace topologique $Z$ un
faisceau d'ensembles $\mathscr{P}$, un faisceau de groupes $\mathscr{G}$ (non
nécessairement commutatif) et un homomorphisme de faisceaux d'ensembles
$m:\mathscr{G}\times\mathscr{P}\to\mathscr{P}$ tel que, pour tout ouvert
$U\subset Z$ tel que $\Gamma(U,\mathscr{P})\neq\varnothing$,
\[
  m_U:\Gamma(U,\mathscr{G})\times\Gamma(U,\mathscr{P})
  \longrightarrow\Gamma(U,\mathscr{P})
\]
fasse de $\Gamma(U,\mathscr{P})$ un \emph{torseur} sous le groupe
$\Gamma(U,\mathscr{G})$, on dit que $\mathscr{P}$ est un
\emph{pseudo-torseur} (ou \emph{faisceau formellement principal homogène})
sous le faisceau de groupes $\mathscr{G}$. On dit que $\mathscr{P}$ est un
\emph{torseur} (ou \emph{faisceau principal homogène}) sous $\mathscr{G}$ si
de plus $\Gamma(U,\mathscr{P})\neq\varnothing$ pour tout ouvert
$U\neq\varnothing$ d'une base convenable de la topologie de $Z$.

Pour la théorie générale des torseurs, nous renvoyons à [42]; nous nous
bornerons ici à rappeler la correspondance canonique entre les classes
d'isomorphie de ces torseurs (pour un $\mathscr{G}$ donné) et les éléments de
l'ensemble de cohomologie $H^1(Z,\mathscr{G})$. Considérons en effet un
torseur $\mathscr{P}$ sous $\mathscr{G}$ et un recouvrement ouvert
$(U_\lambda)$ de $Z$ tel que
$\Gamma(U_\lambda,\mathscr{P})\neq\varnothing$ pour tout $\lambda$;
désignons par $p_\lambda$ un élément de
$\Gamma(U_\lambda,\mathscr{P})$. Pour tout couple d'indices $\lambda$, $\mu$
tel que $U_\lambda\cap U_\mu\neq\varnothing$, il existe alors un élément et
un seul $\gamma_{\lambda\mu}$ de
$\Gamma(U_\lambda\cap U_\mu,\mathscr{G})$ tel que
\[
  \gamma_{\lambda\mu}\cdot
  (p_\mu|U_\lambda\cap U_\mu)=p_\lambda|U_\lambda\cap U_\mu;
\]
en outre, si $\lambda$, $\mu$, $\nu$ sont trois indices tels que
$U_\lambda\cap U_\mu\cap U_\nu\neq\varnothing$, les restrictions
$\gamma'_{\lambda\mu}$, $\gamma'_{\mu\nu}$, $\gamma'_{\lambda\nu}$ de
$\gamma_{\lambda\mu}$, $\gamma_{\mu\nu}$, $\gamma_{\lambda\nu}$ à
$U_\lambda\cap U_\mu\cap U_\nu$ vérifient la condition
$\gamma'_{\lambda\nu}=\gamma'_{\lambda\mu}\gamma'_{\mu\nu}$; autrement dit
$(\lambda,\mu)\mapsto\gamma_{\lambda\mu}$ est un \emph{1-cocycle} du
recouvrement $(U_\lambda)$ à valeurs dans $\mathscr{G}$. Si, pour tout
$\lambda$, $p'_\lambda$ est un second élément de
$\Gamma(U_\lambda,\mathscr{P})$, il existe un élément et un seul
$\beta_\lambda\in\Gamma(U_\lambda,\mathscr{G})$ tel que
$p'_\lambda=\beta_\lambda\cdot p_\lambda$ et le 1-cocycle
$(\gamma'_{\lambda\mu})$ correspondant à la famille $(p'_\lambda)$ est donné
par
\[
  \gamma'_{\lambda\mu}=\beta_\lambda\gamma_{\lambda\mu}\beta_\mu^{-1},
\]
autrement dit, est \emph{cohomologue} à $(\gamma_{\lambda\mu})$.
Inversement, la donnée d'un 1-cocycle $(\gamma_{\lambda\mu})$ définit, pour
tout couple $(\lambda,\mu)$, un automorphisme $\theta_{\lambda\mu}$ du
faisceau \emph{d'ensembles}
$\mathscr{G}|U_\lambda\cap U_\mu$, savoir la translation à droite par
$\gamma_{\lambda\mu}$, et le fait qu'il s'agisse d'un cocycle montre que l'on
peut \emph{recoller} les faisceaux d'ensembles $\mathscr{G}|U_\lambda$ à
l'aide des automorphismes $\theta_{\lambda\mu}$ (0, 3.3.1); on obtient ainsi
de façon évidente un torseur sous $\mathscr{G}$, soit $\mathscr{P}$, et si
l'on prend pour $p_\lambda$ la section unité au-dessus de $U_\lambda$, le
1-cocycle correspondant n'est autre que le 1-cocycle donné
$(\gamma_{\lambda\mu})$; en outre, si l'on remplace
$(\gamma_{\lambda\mu})$ par un 1-cocycle
$\gamma'_{\lambda\mu}=\beta_\lambda\gamma_{\lambda\mu}\beta_\mu^{-1}$ qui
lui est cohomologue, on vérifie aussitôt que le torseur obtenu est isomorphe
à $\mathscr{P}$.

En particulier, si $(\gamma_{\lambda\mu})$ est un \emph{1-cobord}, autrement
dit de la forme
$\gamma_{\lambda\mu}=\beta_\lambda\beta_\mu^{-1}$, le torseur obtenu
$\mathscr{P}$ est \emph{isomorphe à $\mathscr{G}$} (considéré comme torseur
sous lui-même pour les translations à gauche); on dit dans ce cas que
$\mathscr{P}$ est \emph{trivial}, et la réciproque est évidente.

Plus particulièrement, il résulte de (III, 1.3.1) que l'on a:
\end{env}

\begin{proposition}[16.5.16]
\label{IV.16.5.16}
Soient $Z$ un schéma affine, $\mathscr{G}$ un $\mathcal{O}_Z$-Module
quasi-cohérent; alors tout torseur sous $\mathscr{G}$ est trivial.
\end{proposition}

\oldpage[IV-4]{34}
Revenant au problème considéré en (16.5.13), on obtient donc:

\begin{proposition}[16.5.17]
\label{IV.16.5.17}
Soient $X$, $Y$ deux $S$-préschémas, $Y_0$ un sous-préschéma fermé de $Y$
défini par un Idéal quasi-cohérent $\mathscr{J}$ de $\mathcal{O}_Y$ tel que
$\mathscr{J}^2=0$, $j:Y_0\to Y$ l'injection canonique. Soit
$u_0:Y_0\to X$ un $S$-morphisme, et $\mathscr{P}$ le faisceau d'ensembles sur
$Y$ tel que, pour tout ouvert $U$ de $Y$, $\Gamma(U,\mathscr{P})$ soit
l'ensemble des $S$-morphismes $u:U\to X$ tels que
$u_0|U_0=u\circ(j|U_0)$, où $U_0=j^{-1}(U)$. Alors il existe sur
$\mathscr{P}$ une structure de pseudo-torseur sous le
$\mathcal{O}_{Y_0}$-Module
\[
  \mathscr{G}=\mathscr{H}om_{\mathcal{O}_{Y_0}}
  (u_0^*(\Omega_{X/S}^1),\mathscr{J}).
\]
\end{proposition}

En particulier:

\begin{corollary}[16.5.18]
\label{IV.16.5.18}
Avec les notations de (16.5.16), supposons $Y$ affine et $\Omega_{X/S}^1$ de
présentation finie; s'il existe un recouvrement ouvert $(U_\alpha)$ de $Y$,
et, pour chaque indice $\alpha$, un $S$-morphisme $v_\alpha:U_\alpha\to X$
tel que, si $U_\alpha^0=j^{-1}(U_\alpha)$, on ait
$v_\alpha\circ(j|U_\alpha^0)=u_0|U_\alpha^0$, alors il existe un
$S$-morphisme $u:Y\to X$ tel que $u\circ j=u_0$.
\end{corollary}

En effet, $\mathscr{G}$ est alors un $\mathcal{O}_{Y_0}$-Module
quasi-cohérent (I, 1.3.12); en vertu de (16.5.16) et du fait que $Y_0$ est
alors affine, le faisceau $\mathscr{P}$, qui est par hypothèse un torseur sous
$\mathscr{G}$, et non seulement un pseudo-torseur, est \emph{trivial}; mais si
$w$ est un isomorphisme de $\mathscr{G}$ sur $\mathscr{P}$ (en tant que
torseurs sous $\mathscr{G}$), l'image par $w$ de la section 0 de
$\mathscr{G}$ est le $S$-morphisme $u$ cherché.

\subsection{\texorpdfstring{Faisceaux de $p$-différentielles et différentielle extérieure.}{Faisceaux de p-différentielles et différentielle extérieure.}}
\EGAReaderTarget{ega4.c04.subsec.16-6}{IV.16.6}

\begin{env}[16.6.1]
\EGAReaderTarget{ega4.c04.statement.16-6-1}{IV.16.6.1}
Soit $f:X\to S$ un morphisme de préschémas. On appelle \emph{faisceau des
$p$-différentielles de $X$ relativement à $S$} ($p$ entier) la
\emph{puissance extérieure $p$-ème} (0, 4.1.5) du $\mathcal{O}_X$-Module
$\Omega_{X/S}^1$, notée
\[
  \Omega_{X/S}^p=\bigwedge^p(\Omega_{X/S}^1).
  \tag{16.6.1.1}
  \label{IV.16.6.1.1}
\]
On a donc $\Omega_{X/S}^0=\mathcal{O}_X$, et $\Omega_{X/S}^p=0$ pour $p<0$;
les $\Omega_{X/S}^p$ sont les composants homogènes de l'Algèbre extérieure
sur $\Omega_{X/S}^1$
\[
  \Omega_{X/S}^{\bullet}=\bigwedge(\Omega_{X/S}^1)
  =\bigoplus_{p\in\mathbf{Z}}\bigwedge^p(\Omega_{X/S}^1),
  \tag{16.6.1.2}
  \label{IV.16.6.1.2}
\]
qui est donc une $\mathcal{O}_X$-Algèbre graduée quasi-cohérente,
anticommutative et dont les éléments de degré 1 sont de carré nul. Pour tout
ouvert affine $U$ de $X$, on a
\[
  \Gamma(U,\Omega_{X/S}^{\bullet})
  =\bigwedge(\Gamma(U,\Omega_{X/S}^1)),
\]
où $\Gamma(U,\Omega_{X/S}^1)$ est considéré comme
$\Gamma(U,\mathcal{O}_X)$-module.

Lorsque $S=\operatorname{Spec}(A)$ et $X=\operatorname{Spec}(B)$ sont
affines, $B$ étant donc une $A$-algèbre, on a (0, 4.1.5)
$\Omega_{X/S}^p=(\Omega_{B/A}^p)^{\sim}$, en posant
$\Omega_{B/A}^p=\bigwedge^p\Omega_{B/A}^1$.
\end{env}

\begin{theorem}[16.6.2]
\label{IV.16.6.2}
Il existe un endomorphisme et un seul $d$ du faisceau de groupes additifs
$\Omega_{X/S}^{\bullet}$, ayant les propriétés suivantes:
\begin{enumerate}
\item[(i)] $d\circ d=0$.
\item[(ii)] Pour tout ouvert $U$ de $X$ et toute section
$f\in\Gamma(U,\mathcal{O}_X)$, on a $df=d_{X/S}f$.
\oldpage[IV-4]{35}
\item[(iii)] Pour tout ouvert $U$ de $X$, tout couple d'entiers $p$, $q$ et
tout couple de sections
$\omega'_p\in\Gamma(U,\Omega_{X/S}^p)$,
$\omega''_q\in\Gamma(U,\Omega_{X/S}^q)$, on a
\[
  d(\omega'_p\wedge\omega''_q)
  =(d\omega'_p)\wedge\omega''_q
  +(-1)^p\omega'_p\wedge d\omega''_q.
  \tag{16.6.2.1}
  \label{IV.16.6.2.1}
\]
\end{enumerate}
En outre, $d$ est un endomorphisme de $\psi^*(\mathcal{O}_X)$-Modules
gradués, de degré $+1$.
\end{theorem}

Supposons prouvée l'existence de l'endomorphisme $d$. Pour tout ouvert affine
$U$ de $X$, toute section de $\Omega_{X/S}^p$ au-dessus de $U$ est (en vertu
de (i)) combinaison linéaire d'un nombre fini d'éléments de la forme
$g(df_1\wedge df_2\wedge\cdots\wedge df_p)$, où $g$ et les $f_i$ sont des
sections de $\mathcal{O}_X$ au-dessus de $U$ (0, 20.4.7). Les conditions (i)
et (iii) montrent alors, par récurrence sur $p$, que l'on a nécessairement
\[
  d(g(df_1\wedge df_2\wedge\cdots\wedge df_p))
  =dg\wedge df_1\wedge df_2\wedge\cdots\wedge df_p.
  \tag{16.6.2.2}
  \label{IV.16.6.2.2}
\]
Cela prouve donc l'\emph{unicité} de $d$ et la dernière assertion du théorème.
En vertu de cette propriété d'unicité, pour démontrer l'existence de $d$, on
peut se borner au cas où $S=\operatorname{Spec}(A)$ et
$X=\operatorname{Spec}(B)$ sont affines. Or (Bourbaki, \emph{Alg.}, chap.
III, 3\textsuperscript{e} éd., \S~10) pour définir une $A$-antidérivation
$D$ de degré $+1$ d'une algèbre extérieure $\bigwedge(M)$ (où $M$ est un
$B$-module et $B$ une $A$-algèbre), cette antidérivation prenant ses valeurs
dans une $A$-algèbre \emph{graduée anticommutative}
$C=\bigoplus_{n=0}^{\infty}C_n$, dont les éléments de degré 1 sont de carré
nul, il suffit de se donner arbitrairement une $A$-dérivation $D_0$ de $B$
dans $C_1$ et un $A$-homomorphisme $D_1$ de $M$ dans $C_2$; il existe alors
une $A$-antidérivation $D$ et une seule de $\bigwedge(M)$ dans $C$
coïncidant avec $D_0$ dans $B$ et avec $D_1$ dans $M$.

Dans le cas actuel, $D_0$ est nécessairement égal à $d_{B/A}$ en vertu de
(ii); tout revient à voir, compte tenu de (16.6.2.2), qu'il y a un
$A$-homomorphisme $u$ de $\Omega_{B/A}^1$ dans $\Omega_{B/A}^2$ tel que
l'on ait
\[
  u(g.df)=dg\wedge df
  \tag{16.6.2.3}
  \label{IV.16.6.2.3}
\]
quels que soient $f$, $g$ dans $A$; il suffira pour cela de montrer qu'il
existe un $A$-homomorphisme
$v:B\otimes_A\Omega_{B/A}^1\to\Omega_{B/A}^2$ tel que
\[
  v(g.\omega)=dg\wedge\omega
  \tag{16.6.2.4}
  \label{IV.16.6.2.4}
\]
pour $g\in B$ et $\omega\in\Omega_{B/A}^1$. Finalement, comme
$\Omega_{B/A}^1=\mathfrak{I}/\mathfrak{I}^2$ (où
$\mathfrak{I}=\mathfrak{I}_{B/A}$ est le noyau de l'homomorphisme canonique
$B\otimes_A B\to B$) et que $\Omega_{B/A}^1$ est engendré par les éléments
de la forme $g.df$, il suffit de définir un $A$-homomorphisme
$w:B\otimes_A(B\otimes_A B)\to\Omega_{B/A}^2$ tel que
\[
  w(g'\otimes g\otimes f)=dg'\wedge(g.df)
  \tag{16.6.2.5}
  \label{IV.16.6.2.5}
\]
et tel que $w$ soit nul dans l'image de $B\otimes_A\mathfrak{I}^2$. Or,
comme le second membre de (16.6.2.5) est $A$-trilinéaire en $g'$, $g$, $f$,
l'existence de $w$ vérifiant (16.6.2.5) est immédiate. Comme, d'autre part,
$\mathfrak{I}$ est engendré par les éléments $1\otimes x-x\otimes1$
($x\in B$), on est ramené à vérifier que lorsque
$z=(1\otimes x-x\otimes1)(1\otimes y-y\otimes1)$, on a
$w(g'\otimes z)=0$. Or, comme
$z=1\otimes(xy)+(xy)\otimes1-x\otimes y-y\otimes x$, la formule (16.6.2.4)
montre qu'il
\oldpage[IV-4]{36}
suffit de voir que l'on a $d(xy)-x.dy-y.dx=0$, ce qui exprime que $d$ est une
dérivation.

Il reste à prouver que $d$ vérifie la condition (i). Or, le carré d'une
antidérivation est une dérivation (Bourbaki, \emph{loc. cit.}), et comme
$\Omega_{B/A}^{\bullet}$ est engendré par $\Omega_{B/A}^1$ comme
$B$-algèbre, il suffit de vérifier que $d(dz)=0$ pour $z\in B$ et
$z\in\Omega_{B/A}^1$; dans le premier cas, cela résulte de la formule
(16.6.2.3) avec $g=1$; dans le second, on peut se borner au cas où $z=g.df$
avec $f$, $g$ dans $B$, et alors, en vertu de (16.6.2.1) et (16.6.2.3), on a
\[
  d(d(g.df))=d(dg\wedge df)
  =(d(dg))\wedge(df)-(dg)\wedge(d(df))=0.
\]
\hfill C.Q.F.D.

\begin{definition}[16.6.3]
\EGAReaderTarget{ega4.c04.def.16-6-3}{IV.16.6.3}
L'antidérivation $d$ définie dans (16.6.2) (aussi notée $d_{X/S}$) est
appelée \emph{différentielle extérieure sur $X$} (relativement à $S$).
\end{definition}

\begin{proposition}[16.6.4]
\label{IV.16.6.4}
Pour tout changement de base $g:S'\to S$, si l'on pose
$X'=X\times_S S'$, l'homomorphisme canonique
\[
  \Omega_{X/S}^{\bullet}\otimes_S S'
  \longrightarrow\Omega_{X'/S'}^{\bullet}
  \tag{16.6.4.1}
  \label{IV.16.6.4.1}
\]
déduit de l'isomorphisme (16.5.9.1) est bijectif. En outre, si $s$ est une
section de $\Omega_{X/S}^{\bullet}$ au-dessus d'un ouvert $U$ de $X$,
$s\otimes1$ son image réciproque, section de $\Omega_{X'/S'}^{\bullet}$
au-dessus de l'image réciproque $U'$ de $U$ dans $X'$, on a
$d_{X'/S'}(s\otimes1)=d_{X/S}(s)\otimes1$.
\end{proposition}

La première assertion est immédiate, la formation de l'algèbre extérieure
d'un module permutant avec toute extension de l'anneau des scalaires. Pour
prouver la seconde, on peut, en vertu de (16.6.2.2), se borner au cas où
$s\in\Gamma(U,\mathcal{O}_X)$, et dans ce cas l'assertion a déjà été prouvée
(16.4.3.7).

\begin{env}[16.6.5]
\label{IV.16.6.5}
Supposons que $\Omega_{X/S}^1$ soit un $\mathcal{O}_X$-Module localement libre
de rang $n$ en un point $x$, de sorte qu'il y a $n$ sections
$s_i\in\Gamma(U,\mathcal{O}_X)$ telles que les $ds_i$ forment une base du
$\Gamma(U,\mathcal{O}_X)$-module $\Gamma(U,\Omega_{X/S}^1)$ (16.5.8).
Alors, pour tout entier $p\geq1$, les $p$-différentielles
$ds_{i_1}\wedge ds_{i_2}\wedge\cdots\wedge ds_{i_p}$ (pour
$i_1<i_2<\cdots<i_p$, éléments de $[1,n]$) forment une base de
$\binom{n}{p}$ éléments de $\Gamma(U,\Omega_{X/S}^p)$ sur
$\Gamma(U,\mathcal{O}_X)$. En outre la formule (16.6.2.2) montre que pour
toute section $g\in\Gamma(U,\mathcal{O}_X)$, on a
\[
  d(g.ds_{i_1}\wedge ds_{i_2}\wedge\cdots\wedge ds_{i_p})
  =\sum_k(-1)^r\frac{\partial g}{\partial s_k}
  ds_{i_1}\wedge\cdots\wedge ds_{i_r}\wedge ds_k\wedge
  ds_{i_{r+1}}\wedge\cdots\wedge ds_{i_p}
  \tag{16.6.5.1}
  \label{IV.16.6.5.1}
\]
où, au second membre, $k$ parcourt l'ensemble des $n-p$ indices distincts des
$i_h$, $i_r$ étant le plus grand indice $<k$.

On notera que la relation $d(dg)=0$ pour toute section
$g\in\Gamma(U,\mathcal{O}_X)$ s'exprime sous la forme
\[
  D_i(D_jg)=D_j(D_ig)\qquad\text{pour $i\neq j$};
\]
autrement dit, les dérivations $D_i$ définies dans (16.5.7) sont deux à deux
permutables.
\end{env}

\subsection{\texorpdfstring{Les $\mathscr{P}_{X/S}^n(\mathscr{F})$.}{Les P\string_\string{X/S\string}\string^n(F).}}
\EGAReaderTarget{ega4.c04.subsec.16-7}{IV.16.7}

\begin{env}[16.7.1]
\label{IV.16.7.1}
Soient $f:X\to S$ un morphisme de préschémas, $\mathscr{F}$ un
$\mathcal{O}_X$-Module. Désignons par $X_{\Delta_f}^{(n)}$ le \emph{$n$-ème
voisinage infinitésimal de $X$} pour le morphisme diagonal
\oldpage[IV-4]{37}
$\Delta_f:X\to X\times_S X$, par
$h_n:X_{\Delta_f}^{(n)}\to X\times_S X$ le morphisme canonique (16.1.2), et
considérons les deux morphismes composés
\[
  p_1^{(n)}:X_{\Delta_f}^{(n)}\xrightarrow{h_n}X\times_S X
  \xrightarrow{p_1}X,qquad
  p_2^{(n)}:X_{\Delta_f}^{(n)}\xrightarrow{h_n}X\times_S X
  \xrightarrow{p_2}X
\]
de sorte que, par définition, $p_1^{(n)}$ correspond à l'homomorphisme de
faisceaux d'anneaux $\mathcal{O}_X\to\mathscr{P}_{X/S}^n$ que nous avons
choisi pour définir la structure de $\mathcal{O}_X$-Algèbre sur
$\mathscr{P}_{X/S}^n$ (16.3.5), et $p_2^{(n)}$ à l'homomorphisme de faisceaux
d'anneaux $d_{X/S}^n:\mathcal{O}_X\to\mathscr{P}_{X/S}^n$ (16.3.6). Comme
$X_{\Delta_f}^{(n)}$ et $X$ ont même espace sous-jacent, on peut écrire
\[
  \mathscr{P}_{X/S}^n
  =(p_1^{(n)})_*((p_2^{(n)})^*(\mathcal{O}_X)).
  \tag{16.7.1.1}
  \label{IV.16.7.1.1}
\]
Plus généralement, nous poserons
\[
  \mathscr{P}_{X/S}^n(\mathscr{F})
  =(p_1^{(n)})_*((p_2^{(n)})^*(\mathscr{F}))
  \tag{16.7.1.2}
  \label{IV.16.7.1.2}
\]
de sorte que
$\mathscr{P}_{X/S}^n=\mathscr{P}_{X/S}^n(\mathcal{O}_X)$; par définition,
$\mathscr{P}_{X/S}^n(\mathscr{F})$ est donc un $\mathcal{O}_X$-Module.
\end{env}

\begin{env}[16.7.2]
\label{IV.16.7.2}
Si l'on revient aux définitions des images réciproques de Modules sur les
espaces annelés (0, 4.3.1) et que l'on tienne compte de ce que
$X_{\Delta_f}^{(n)}$ et $X$ ont même espace sous-jacent, on voit que l'on peut
aussi écrire la définition (16.7.1.2) sous la forme
\[
  \mathscr{P}_{X/S}^n(\mathscr{F})
  =\mathscr{P}_{X/S}^n\otimes_{\mathcal{O}_X}\mathscr{F}
  \tag{16.7.2.1}
  \label{IV.16.7.2.1}
\]
mais où il faut prendre garde que, dans l'interprétation du signe $\otimes$,
$\mathscr{P}_{X/S}^n$ est muni de sa structure de $\mathcal{O}_X$-Module
définie par l'homomorphisme de faisceaux d'anneaux
$d_{X/S}^n:\mathcal{O}_X\to\mathscr{P}_{X/S}^n$. Il résulte aussitôt de cette
formule (ou directement de (16.7.1.2)) que
$\mathscr{P}_{X/S}^n(\mathscr{F})$ est canoniquement muni d'une structure de
$\mathscr{P}_{X/S}^n$-Module.
\end{env}

\begin{proposition}[16.7.3]
\label{IV.16.7.3}
\begin{enumerate}
\item[(i)] Le foncteur
$\mathscr{F}\mapsto\mathscr{P}_{X/S}^n(\mathscr{F})$ de la catégorie des
$\mathcal{O}_X$-Modules dans celle des $\mathscr{P}_{X/S}^n$-Modules est
exact à droite, et commute aux limites inductives quelconques; il est exact
lorsque $\mathscr{P}_{X/S}^n$ est un $\mathcal{O}_X$-Module plat.
\item[(ii)] Si $\mathscr{F}$ est un $\mathcal{O}_X$-Module quasi-cohérent
(resp. de type fini, resp. de présentation finie),
$\mathscr{P}_{X/S}^n(\mathscr{F})$ est un $\mathscr{P}_{X/S}^n$-Module
quasi-cohérent (resp. de type fini, resp. de présentation finie).
\end{enumerate}
\end{proposition}

Les assertions de (i) résultent aussitôt de la formule (16.7.2.1) et de la
considération de la symétrie de $\mathscr{P}_{X/S}^n$ (16.3.4). Les assertions
de (ii) découlent de l'exactitude à droite du foncteur
$\mathscr{F}\mapsto\mathscr{P}_{X/S}^n(\mathscr{F})$.

\begin{env}[16.7.4]
\label{IV.16.7.4}
Les deux structures de $\mathcal{O}_X$-Module de $\mathscr{P}_{X/S}^n$
définissent sur $\mathscr{P}_{X/S}^n(\mathscr{F})$ deux structures de
$\mathcal{O}_X$-Module, d'ailleurs permutables, donc une structure de
$\mathcal{O}_X$-Bimodule. Il est commode de noter \emph{à gauche} celle de
ces structures provenant de l'homomorphisme structural
$\mathcal{O}_X\to\mathscr{P}_{X/S}^n$ (choisi dans (16.3.5)) et
\emph{à droite} celle provenant de l'homomorphisme
$d_{X/S}^n:\mathcal{O}_X\to\mathscr{P}_{X/S}^n$. Autrement dit, pour tout
ouvert $U$ de $X$, et tout triplet d'éléments
$a\in\Gamma(U,\mathcal{O}_X)$, $b\in\Gamma(U,\mathscr{P}_{X/S}^n)$,
$t\in\Gamma(U,\mathscr{F})$, on a par définition
\[
  a(b\otimes t)=(ab)\otimes t,qquad
  (b\otimes t)a=(b.d^na)\otimes t=b\otimes(at)
  =(d^na).(b\otimes t).
  \tag{16.7.4.1}
  \label{IV.16.7.4.1}
\]
La structure de $\mathcal{O}_X$-Module provenant de la définition (16.7.1.2)
est donc, avec ces conventions, la structure de $\mathcal{O}_X$-Module à
gauche.

\oldpage[IV-4]{38}
Si $\mathscr{F}$ est un $\mathcal{O}_X$-Module quasi-cohérent, il en est de
même de $\mathscr{P}_{X/S}^n(\mathscr{F})$ pour l'une ou l'autre de ses
structures de $\mathcal{O}_X$-Module. Si de plus $\mathscr{F}$ est de type
fini (resp. de présentation finie) et $f:X\to S$ localement de type fini
(resp. localement de présentation finie),
$\mathscr{P}_{X/S}^n(\mathscr{F})$ est (pour l'une ou l'autre de ses
structures de $\mathcal{O}_X$-Module) de type fini (resp. de présentation
finie), comme il résulte de (16.3.9) et de (16.4.22).
\end{env}

\begin{env}[16.7.5]
\label{IV.16.7.5}
La définition (16.7.2.1) entraîne l'existence d'un homomorphisme de faisceaux
de groupes commutatifs
\[
  d_{X/S,\mathscr{F}}^n:\mathscr{F}
  \longrightarrow\mathscr{P}_{X/S}^n(\mathscr{F})
  \qquad\text{(aussi noté $d_{X/S}^n$)}
  \tag{16.7.5.1}
  \label{IV.16.7.5.1}
\]
tel qu'avec les notations de (16.7.4), on ait
\[
  d_{X/S,\mathscr{F}}^n(t)=1\otimes t
  \tag{16.7.5.2}
  \label{IV.16.7.5.2}
\]
et par suite, en vertu de (16.7.4.1)
\[
  d_{X/S,\mathscr{F}}^n(at)
  =(1\otimes t)a=(d_{X/S,\mathscr{F}}^n(t)).a
  \tag{16.7.5.3}
  \label{IV.16.7.5.3}
\]
\[
  d_{X/S,\mathscr{F}}^n(at)
  =(d_{X/S}^n(a)).(1\otimes t)
  =(d_{X/S}^n(a))(d_{X/S,\mathscr{F}}^n(t)).
  \tag{16.7.5.4}
  \label{IV.16.7.5.4}
\]
Il est donc $\mathcal{O}_X$-linéaire pour la structure de
$\mathcal{O}_X$-Module à droite sur $\mathscr{P}_{X/S}^n(\mathscr{F})$, et
\emph{semi-linéaire} (relativement à l'automorphisme $\sigma$ (16.3.4)) pour
la structure de $\mathcal{O}_X$-Module à gauche.
\end{env}

\begin{proposition}[16.7.6]
\label{IV.16.7.6}
Le $\mathcal{O}_X$-Module à gauche $\mathscr{P}_{X/S}^n(\mathscr{F})$ est
engendré par l'image de $\mathscr{F}$ par l'homomorphisme canonique
$d_{X/S,\mathscr{F}}^n$.
\end{proposition}

Cela résulte aussitôt de (16.7.5.3) et du cas particulier correspondant à
$\mathscr{F}=\mathcal{O}_X$ (16.3.8).

\begin{env}[16.7.7]
\label{IV.16.7.7}
Les homomorphismes canoniques de faisceaux d'anneaux
\[
  \varphi_{nm}:\mathscr{P}_{X/S}^m\longrightarrow\mathscr{P}_{X/S}^n
\]
pour $n\leq m$ (16.1.2) définissent, en vertu de (16.7.2.1) des
homomorphismes canoniques
\[
  \mathscr{P}_{X/S}^m(\mathscr{F})
  \longrightarrow\mathscr{P}_{X/S}^n(\mathscr{F})qquad(n\leq m)
\]
qui sont des homomorphismes de $\mathcal{O}_X$-Bimodules en vertu de (16.1.6)
et (16.7.4.1); en outre on a des diagrammes commutatifs
\[
  \xymatrix{
    \mathscr{P}_{X/S}^m(\mathscr{F})\ar[rr]
      &&\mathscr{P}_{X/S}^n(\mathscr{F})\\
    &\mathscr{F}\ar[ul]^-{d_{X/S,\mathscr{F}}^m}
      \ar[ur]_-{d_{X/S,\mathscr{F}}^n}
  }
\]
On a donc ainsi un système projectif de $\mathcal{O}_X$-Bimodules
$(\mathscr{P}_{X/S}^n(\mathscr{F}))$, et l'on pose
\[
  \mathscr{P}_{X/S}^{\infty}(\mathscr{F})
  =\varprojlim\mathscr{P}_{X/S}^n(\mathscr{F}).
  \tag{16.7.7.1}
  \label{IV.16.7.7.1}
\]
En outre, ce qui précède montre que les homomorphismes (16.7.5.1) forment un
système projectif d'homomorphismes, et définissent donc un homomorphisme
canonique
\[
  d_{X/S,\mathscr{F}}^{\infty}:\mathscr{F}
  \longrightarrow\mathscr{P}_{X/S}^{\infty}(\mathscr{F}).
  \tag{16.7.7.2}
  \label{IV.16.7.7.2}
\]
\end{env}

\oldpage[IV-4]{39}

\begin{env}[16.7.8]
\label{IV.16.7.8}
Soient $\mathscr{F}$, $\mathscr{G}$ deux $\mathcal{O}_X$-Modules; il résulte
aussitôt de la définition (16.7.2.1) que l'on a un isomorphisme canonique de
$\mathscr{P}_{X/S}^n$-Modules
\[
  \mathscr{P}_{X/S}^n(\mathscr{F}\otimes_{\mathcal{O}_X}\mathscr{G})
  \xrightarrow{\sim}
  \mathscr{P}_{X/S}^n(\mathscr{F})
  \otimes_{\mathscr{P}_{X/S}^n}\mathscr{P}_{X/S}^n(\mathscr{G})
  \tag{16.7.8.1}
  \label{IV.16.7.8.1}
\]
(Bourbaki, \emph{Alg.}, chap.~II, 3\textsuperscript{e} éd., §~5,
n\textsuperscript{o}~1, prop.~3).

On en conclut en particulier (ou l'on voit directement sur la définition
(16.7.2.1)) que si $\mathscr{F}$ est muni d'une structure de
$\mathcal{O}_X$-Algèbre (non nécessairement associative),
$\mathscr{P}_{X/S}^n(\mathscr{F})$ est canoniquement muni d'une structure de
$\mathscr{P}_{X/S}^n$-Algèbre; cette dernière est associative (resp.
commutative, resp. unitaire, resp. une Algèbre de Lie) lorsqu'il en est ainsi
de $\mathscr{F}$. En outre les homomorphismes canoniques
$\mathscr{P}_{X/S}^m(\mathscr{F})\to\mathscr{P}_{X/S}^n(\mathscr{F})$ pour
$n\leq m$ (16.7.7) sont alors des di-homomorphismes d'Algèbres; de même
(16.7.5.1) est alors un homomorphisme de $\mathcal{O}_X$-Algèbres lorsque
$\mathscr{P}_{X/S}^n(\mathscr{F})$ est muni de sa structure de
$\mathcal{O}_X$-Algèbre provenant de sa structure de $\mathcal{O}_X$-Module
\emph{à droite}.

Avec les mêmes notations, on a également un homomorphisme canonique de
$\mathscr{P}_{X/S}^n$-Modules
\[
  \mathscr{P}_{X/S}^n
  (\mathscr{H}\!om_{\mathcal{O}_X}(\mathscr{F},\mathscr{G}))
  \longrightarrow
  \mathscr{H}\!om_{\mathscr{P}_{X/S}^n}
  (\mathscr{P}_{X/S}^n(\mathscr{F}),\mathscr{P}_{X/S}^n(\mathscr{G}))
  \tag{16.7.8.2}
  \label{IV.16.7.8.2}
\]
(Bourbaki, \emph{Alg.}, 3\textsuperscript{e} éd., §~5,
n\textsuperscript{o}~3), qui est bijectif lorsque $\mathscr{P}_{X/S}^n$ est
un $\mathcal{O}_X$-Module \emph{localement libre de type fini}
(\emph{loc. cit.}, prop.~7).
\end{env}

\begin{env}[16.7.9]
\label{IV.16.7.9}
Supposons qu'on soit dans la situation décrite dans (16.4.1); alors, de
l'homomorphisme canonique $P^n(u)$ (16.4.3.3), on déduit aussitôt un
homomorphisme canonique de $\mathcal{O}_{X'}$-Bimodules
\[
  u^*(\mathscr{P}_{X/S}^n(\mathscr{F}))
  \longrightarrow
  \mathscr{P}_{X'/S'}^n(u^*(\mathscr{F})).
  \tag{16.7.9.1}
  \label{IV.16.7.9.1}
\]
Nous laissons au lecteur le soin d'étendre à cet homomorphisme les propriétés
vues dans (16.4) pour le cas $\mathscr{F}=\mathcal{O}_X$.
\end{env}

\begin{remark}[16.7.10]
\label{IV.16.7.10}
La définition de $\mathscr{P}_{X/S}^n(\mathscr{F})$ sous la forme (16.7.1.2)
garde un sens lorsque $\mathscr{F}$ est un faisceau d'ensembles quelconque
(l'image réciproque d'un faisceau d'ensembles par $p_2^{(n)}$ étant définie
dans (0\textsubscript{I}, 3.7.1)); une variante de cette définition permet de
définir le «schéma des jets» (relativement à $S$) d'un $X$-préschéma
quelconque.
\end{remark}

\subsection{Opérateurs différentiels.}
\footnote{Pour un formalisme plus général, voir l'Exposé VII de [42]
(dû à P. Gabriel).}
\EGAReaderTarget{ega4.c04.subsec.16-8}{IV.16.8}

\begin{definition}[16.8.1]
\EGAReaderTarget{ega4.c04.def.16-8-1}{IV.16.8.1}
Soient $f=(\psi,\theta):X\to S$ un morphisme de préschémas,
$\mathscr{F}$, $\mathscr{G}$ deux $\mathcal{O}_X$-Modules, $n$ un entier
$\geq0$. On dit qu'un homomorphisme de faisceaux de groupes additifs
$D:\mathscr{F}\to\mathscr{G}$ est un \emph{opérateur différentiel d'ordre
$\leq n$} (relativement à $S$) s'il existe un homomorphisme de
$\mathcal{O}_X$-Modules
$u:\mathscr{P}_{X/S}^n(\mathscr{F})\to\mathscr{G}$ (où
$\mathscr{P}_{X/S}^n(\mathscr{F})$ est muni de sa structure de
$\mathcal{O}_X$-Module à gauche (16.7.4)) tel que l'on ait
$D=u\circ d_{X/S,\mathscr{F}}^n$.
\end{definition}

Il est clair, en vertu de l'existence des homomorphismes canoniques
\[
  \mathscr{P}_{X/S}^m(\mathscr{F})
  \longrightarrow\mathscr{P}_{X/S}^n(\mathscr{F})
\]
\oldpage[IV-4]{40}
pour $n\leq m$ (16.7.7) qu'un opérateur différentiel d'ordre $\leq n$ est
aussi un opérateur différentiel d'ordre $\leq m$ pour tout $m\geq n$. Si
$D:\mathscr{F}\to\mathscr{G}$ est un opérateur différentiel d'ordre
$\leq n$, alors, pour tout ouvert $U$ de $X$,
$D|U:\mathscr{F}|U\to\mathscr{G}|U$ est aussi un opérateur différentiel
d'ordre $\leq n$.

On dit qu'un homomorphisme $D:\mathscr{F}\to\mathscr{G}$ des faisceaux de
groupes additifs sous-jacents à $\mathscr{F}$ et à $\mathscr{G}$ est un
\emph{opérateur différentiel} (relativement à $S$) si, pour tout $x\in X$, il
existe un voisinage ouvert $U$ de $x$ et un entier $n\geq0$ tels que
$D|U:\mathscr{F}|U\to\mathscr{G}|U$ soit un opérateur différentiel d'ordre
$\leq n$. L'\emph{ordre} d'un opérateur différentiel
$D:\mathscr{F}\to\mathscr{G}$ est la borne inférieure des entiers $n$ tels
que $D$ soit un opérateur différentiel d'ordre $\leq n$ (et donc $+\infty$
s'il n'y a pas de tels entiers); cet ordre est toujours fini si $X$ est
\emph{quasi-compact}. Les opérateurs différentiels d'ordre $0$ ne sont autres
que les homomorphismes de $\mathcal{O}_X$-Modules
$\mathscr{F}\to\mathscr{G}$; on convient que tout opérateur différentiel
d'ordre $<0$ est \emph{nul}. Pour $n\geq0$, un opérateur différentiel n'est
pas en général un homomorphisme de $\mathcal{O}_X$-Modules mais est toujours
un homomorphisme de $\psi^*(\mathcal{O}_S)$-Modules.

Lorsque $\mathscr{F}=\mathcal{O}_X$, un opérateur différentiel d'ordre
$\leq1$ de $\mathcal{O}_X$ dans $\mathscr{G}$ s'écrit d'une seule manière
sous la forme $v+D$, où $v:\mathcal{O}_X\to\mathscr{G}$ est un
$\mathcal{O}_X$-homomorphisme, et $D$ une $S$-\emph{dérivation} (16.5.1) de
$\mathcal{O}_X$ dans $\mathscr{G}$: cela résulte de la structure de
$P_{B/A}^1$ (0, 20.4.8).

\begin{env}[16.8.2]
\label{IV.16.8.2}
Pour décrire de façon plus explicite un opérateur différentiel d'ordre
$\leq n$, $D:\mathscr{F}\to\mathscr{G}$, il suffit, pour tout ouvert affine
$U$ de $X$, dont l'image dans $S$ est contenue dans un ouvert affine $V$, de
caractériser l'homomorphisme
$D=D_U:\Gamma(U,\mathscr{F})\to\Gamma(U,\mathscr{G})$. Si l'on pose
$\Gamma(V,\mathcal{O}_S)=A$, $\Gamma(U,\mathcal{O}_X)=B$, de sorte que $B$
est une $A$-algèbre, on a
\[
  \Gamma(U,\mathscr{P}_{X/S}^n)
  =(B\otimes_A B)/\mathfrak{I}^{n+1},
\]
où l'on a posé pour abréger $\mathfrak{I}=\mathfrak{I}_{B/A}$. Posons en
outre $M=\Gamma(U,\mathscr{F})$, $N=\Gamma(U,\mathscr{G})$; alors la
définition de $D$ signifie que pour chaque couple $(U,V)$ vérifiant les
conditions précédentes, le $A$-homomorphisme $D:M\to N$ se factorise en
\[
  M\longrightarrow
  ((B\otimes_A B)/\mathfrak{I}^{n+1})\otimes_B M
  \xrightarrow{v}N
\]
où la première flèche est l'homomorphisme canonique $t\mapsto1\otimes t$, et
$v$ est un $B$-homomorphisme, la structure de $B$-module de
$((B\otimes_A B)/\mathfrak{I}^{n+1})\otimes_B M$ provenant du premier
facteur $B$ (alors qu'on rappelle que dans la formation du produit tensoriel
sur $B$, la structure de $B$-module de
$(B\otimes_A B)/\mathfrak{I}^{n+1}$ provient du second facteur $B$). Notons
maintenant que le $B$-module
$((B\otimes_A B)/\mathfrak{I}^{n+1})\otimes_B M$ est isomorphe à
$(B\otimes_A M)/\mathfrak{I}^{n+1}(B\otimes_A M)$, où $B\otimes_A M$ est
considéré comme $(B\otimes_A B)$-module et sa structure de $B$-module
provient de l'homomorphisme $b\mapsto b\otimes1$ de $B$ dans $B\otimes_A B$.
Soit alors $D'$ le $B$-homomorphisme de $B\otimes_A M$ dans $N$ tel que
$D'(b\otimes t)=bD(t)$; la condition de factorisation sur $D$ s'exprime encore
en disant que $D'$ doit \emph{être nul dans le $B$-module}
$\mathfrak{I}^{n+1}(B\otimes_A M)$.
\end{env}

\begin{env}[16.8.3]
\label{IV.16.8.3}
Il est clair que l'ensemble des opérateurs différentiels d'ordre $\leq n$ de
$\mathscr{F}$ dans $\mathscr{G}$ est un groupe additif, noté
$\operatorname{Diff}_{X/S}^n(\mathscr{F},\mathscr{G})$; lorsque
$\mathscr{F}=\mathscr{G}=\mathcal{O}_X$, on écrit aussi
$\operatorname{Diff}_{X/S}^n$ au lieu de
$\operatorname{Diff}_{X/S}^n(\mathcal{O}_X,\mathcal{O}_X)$.

On a vu (16.8.1) que l'on a pour deux ouverts $U\supset V$ de $X$, un
homomorphisme canonique de restriction
\[
  \operatorname{Diff}_{U/S}^n(\mathscr{F}|U,\mathscr{G}|U)
  \longrightarrow
  \operatorname{Diff}_{V/S}^n(\mathscr{F}|V,\mathscr{G}|V)
\]
\oldpage[IV-4]{41}
donc $U\mapsto\operatorname{Diff}_{U/S}^n(\mathscr{F}|U,\mathscr{G}|U)$ est
un préfaisceau de groupes additifs; en fait c'est même un \emph{faisceau}, car
pour $U$ ouvert variable dans $X$, les homomorphismes
$u\mapsto u\circ d_{U/S,\mathscr{F}|U}^n$ sont des \emph{isomorphismes} de
groupes additifs
\[
  \mathscr{H}\!om_{\mathcal{O}_U}
  (\mathscr{P}_{U/S}^n(\mathscr{F}|U),\mathscr{G}|U)
  \xrightarrow{\sim}
  \operatorname{Diff}_{U/S}^n(\mathscr{F}|U,\mathscr{G}|U),
  \tag{16.8.3.1}
  \label{IV.16.8.3.1}
\]
en vertu du fait que l'image de $\mathscr{F}$ par
$d_{X/S,\mathscr{F}}^n$ engendre $\mathscr{P}_{X/S}^n(\mathscr{F})$
(16.7.6). On note ce faisceau
$\mathscr{D}\!iff_{X/S}^n(\mathscr{F},\mathscr{G})$, et on a donc:
\end{env}

\begin{proposition}[16.8.4]
\label{IV.16.8.4}
Les isomorphismes (16.8.3.1) définissent un isomorphisme de faisceaux de
groupes additifs
\[
  \mathscr{H}\!om_{\mathcal{O}_X}
  (\mathscr{P}_{X/S}^n(\mathscr{F}),\mathscr{G})
  \xrightarrow{\sim}
  \mathscr{D}\!iff_{X/S}^n(\mathscr{F},\mathscr{G}).
  \tag{16.8.4.1}
  \label{IV.16.8.4.1}
\]
\end{proposition}

Lorsque $\mathscr{F}=\mathscr{G}=\mathcal{O}_X$, on écrit aussi
$\mathscr{D}\!iff_{X/S}^n$ au lieu de
$\mathscr{D}\!iff_{X/S}^n(\mathcal{O}_X,\mathcal{O}_X)$; il résulte de
(16.8.4) que $\mathscr{D}\!iff_{X/S}^n$ s'identifie canoniquement au
\emph{dual} du $\mathcal{O}_X$-Module $\mathscr{P}_{X/S}^n$; aussi écrit-on
$\langle t,D\rangle$ au lieu de $u(t)$ si $t$ est une section de
$\mathscr{P}_{X/S}^n$ au-dessus d'un ouvert et si $u$ est l'homomorphisme de
$\mathscr{P}_{X/S}^n$ dans $\mathcal{O}_X$ correspondant à $D$.

\begin{env}[16.8.5]
\label{IV.16.8.5}
Puisque $\mathscr{P}_{X/S}^n(\mathscr{F})$ est muni d'une structure de
$\mathcal{O}_X$-Bimodule (16.7.4), on en déduit canoniquement une structure de
$\mathcal{O}_X$-Bimodule sur
$\mathscr{H}\!om_{\mathcal{O}_X}
(\mathscr{P}_{X/S}^n(\mathscr{F}),\mathscr{G})$, donc aussi sur
$\mathscr{D}\!iff_{X/S}^n(\mathscr{F},\mathscr{G})$ en vertu de (16.8.4.1).
De façon précise, à la structure de $\mathcal{O}_X$-Module à gauche de
$\mathscr{P}_{X/S}^n(\mathscr{F})$ correspond, en vertu de la définition
(16.8.1), la structure de $\mathcal{O}_X$-Module à gauche sur
$\mathscr{D}\!iff_{X/S}^n(\mathscr{F},\mathscr{G})$ explicitée comme suit:
pour tout ouvert $U$ de $X$, toute section
$a\in\Gamma(U,\mathcal{O}_X)$ et tout opérateur différentiel
$D:\mathscr{F}|U\to\mathscr{G}|U$, $aD$ est l'opérateur différentiel qui, à
toute section $t\in\Gamma(U,\mathscr{F})$, fait correspondre la section
\[
  (aD)(t)=a(D(t))
  \tag{16.8.5.1}
  \label{IV.16.8.5.1}
\]
de $\Gamma(U,\mathscr{G})$. De même, à la structure de
$\mathcal{O}_X$-Module à droite de $\mathscr{P}_{X/S}^n(\mathscr{F})$
correspond la structure de $\mathcal{O}_X$-Module à droite sur
$\mathscr{D}\!iff_{X/S}^n(\mathscr{F},\mathscr{G})$ explicitée comme suit:
avec les mêmes notations que ci-dessus, $Da$ est l'opérateur différentiel qui,
à $t\in\Gamma(U,\mathscr{F})$, fait correspondre la section
\[
  (Da)(t)=D(at).
  \tag{16.8.5.2}
  \label{IV.16.8.5.2}
\]
\end{env}

\begin{proposition}[16.8.6]
\label{IV.16.8.6}
Si $f:X\to S$ est un morphisme localement de présentation finie,
$\mathscr{F}$ un $\mathcal{O}_X$-Module quasi-cohérent de présentation finie
et $\mathscr{G}$ un $\mathcal{O}_X$-Module quasi-cohérent, alors
$\mathscr{D}\!iff_{X/S}^n(\mathscr{F},\mathscr{G})$ est un
$\mathcal{O}_X$-Module quasi-cohérent pour l'une ou l'autre structure définies
dans (16.8.5).
\end{proposition}

La proposition résulte de ce que, sous les hypothèses faites,
$\mathscr{P}_{X/S}^n(\mathscr{F})$ est un $\mathcal{O}_X$-Module
quasi-cohérent de présentation finie (16.7.4), et de (I, 1.3.12).

\begin{env}[16.8.7]
\label{IV.16.8.7}
L'ensemble des opérateurs différentiels de $\mathscr{F}$ dans $\mathscr{G}$
(d'ordre non précisé (16.8.1)) se note
$\operatorname{Diff}_{X/S}(\mathscr{F},\mathscr{G})$; on voit encore comme
dans (16.8.3) que
$U\mapsto\operatorname{Diff}_{U/S}(\mathscr{F}|U,\mathscr{G}|U)$ est un
faisceau de groupes additifs, que nous désignerons par
$\mathscr{D}\!iff_{X/S}(\mathscr{F},\mathscr{G})$. Il est immédiat que
$\mathscr{D}\!iff_{X/S}(\mathscr{F},\mathscr{G})$ est réunion de la famille
filtrante croissante de ses sous-faisceaux
$\mathscr{D}\!iff_{X/S}^n(\mathscr{F},\mathscr{G})$; si $X$ est
quasi-compact, $\operatorname{Diff}_{X/S}(\mathscr{F},\mathscr{G})$
\oldpage[IV-4]{42}
est de même réunion de ses sous-groupes
$\operatorname{Diff}_{X/S}^n(\mathscr{F},\mathscr{G})$ (16.8.1). Les
structures de $\mathcal{O}_X$-Bimodule sur les
$\mathscr{D}\!iff_{X/S}^n(\mathscr{F},\mathscr{G})$ définissent donc une
structure de $\mathcal{O}_X$-Bimodule sur
$\mathscr{D}\!iff_{X/S}(\mathscr{F},\mathscr{G})$, explicitée encore par
(16.8.5.1) et (16.8.5.2).

Notons que, pour $n\leq m$, on a un diagramme commutatif
\[
  \xymatrix{
    \mathscr{H}\!om_{\mathcal{O}_X}
      (\mathscr{P}_{X/S}^n(\mathscr{F}),\mathscr{G})
      \ar[r]^-{\sim}\ar[d]
      &\mathscr{D}\!iff_{X/S}^n(\mathscr{F},\mathscr{G})\ar[d]\\
    \mathscr{H}\!om_{\mathcal{O}_X}
      (\mathscr{P}_{X/S}^m(\mathscr{F}),\mathscr{G})
      \ar[r]^-{\sim}
      &\mathscr{D}\!iff_{X/S}^m(\mathscr{F},\mathscr{G})
  }
  \tag{16.8.7.1}
  \label{IV.16.8.7.1}
\]
où les flèches horizontales sont les isomorphismes (16.8.4.2) et la flèche
verticale de gauche provient de l'homomorphisme canonique
$\mathscr{P}_{X/S}^m(\mathscr{F})\to
\mathscr{P}_{X/S}^n(\mathscr{F})$ (16.7.7). Pour tout ouvert $U$ de $X$,
munissons alors
\[
  \Gamma(U,\mathscr{P}_{X/S}^{\infty}(\mathscr{F}))
  =\varprojlim\Gamma(U,\mathscr{P}_{X/S}^n(\mathscr{F}))
\]
de la topologie limite projective des topologies discrètes sur les
$\Gamma(U,\mathscr{P}_{X/S}^n(\mathscr{F}))$, ce qui définit
$\Gamma(U,\mathscr{P}_{X/S}^{\infty}(\mathscr{F}))$ comme un
$\Gamma(U,\mathcal{O}_X)$-bimodule topologique, de sorte que
$\mathscr{P}_{X/S}^{\infty}(\mathscr{F})$ apparaît comme un faisceau à valeurs
dans la catégorie des groupes commutatifs topologiques
(0\textsubscript{I}, 3.2.6). Alors (G, II, 1.11) la limite du système inductif
de faisceaux de groupes commutatifs
\[
  (\mathscr{H}\!om_{\mathcal{O}_X}
  (\mathscr{P}_{X/S}^n(\mathscr{F}),\mathscr{G}))
\]
n'est autre que le faisceau des germes d'homomorphismes \emph{continus} de
$\mathscr{P}_{X/S}^{\infty}(\mathscr{F})$ dans $\mathscr{G}$ (ce dernier
étant muni de la topologie discrète): les homomorphismes continus de
$\Gamma(U,\mathscr{P}_{X/S}^{\infty}(\mathscr{F}))$ dans le groupe discret
$\Gamma(U,\mathscr{G})$ correspondent en effet biunivoquement aux systèmes
inductifs d'homomorphismes de groupes
$\Gamma(U,\mathscr{P}_{X/S}^n(\mathscr{F}))\to\Gamma(U,\mathscr{G})$. On peut
donc encore exprimer (16.8.4) en disant qu'on a un isomorphisme canonique
\[
  \mathscr{H}\!om.\mathrm{cont}_{\mathcal{O}_X}
  (\mathscr{P}_{X/S}^{\infty}(\mathscr{F}),\mathscr{G})
  \xrightarrow{\sim}
  \mathscr{D}\!iff_{X/S}(\mathscr{F},\mathscr{G})
\]
où le premier membre désigne le faisceau des germes d'homomorphismes continus
de $\mathscr{P}_{X/S}^{\infty}(\mathscr{F})$ dans $\mathscr{G}$.
\end{env}

\begin{proposition}[16.8.8]
\label{IV.16.8.8}
Soient $\mathscr{F}$, $\mathscr{G}$ deux $\mathcal{O}_X$-Modules,
$D:\mathscr{F}\to\mathscr{G}$ un homomorphisme de
$\psi^*(\mathcal{O}_S)$-Modules, $n$ un entier $\geq0$. Les conditions
suivantes sont équivalentes:
\begin{enumerate}
  \item[(a)] $D$ est un opérateur différentiel d'ordre $\leq n$.
  \item[(b)] Pour toute section $a$ de $\mathcal{O}_X$ au-dessus d'un ouvert
  $U$, l'homomorphisme $D_a:\mathscr{F}|U\to\mathscr{G}|U$ tel que, pour
  toute section $t$ de $\mathscr{F}$ au-dessus d'un ouvert $V\subset U$, on
  ait
  \[
    D_a(t)=D(at)-aD(t)
    \tag{16.8.8.1}
    \label{IV.16.8.8.1}
  \]
  est un opérateur différentiel d'ordre $\leq n-1$.
  \item[(c)] Pour tout ouvert $U$ de $X$, toute famille
  $(a_i)_{1\leq i\leq n+1}$ de $n+1$ sections de $\mathcal{O}_X$ au-dessus
  de $U$ et toute section $t$ de $\mathscr{F}$ au-dessus de $U$, on a
  l'identité
  \[
    \sum_{H\subset I_{n+1}}(-1)^{\operatorname{Card}(H)}
    \left(\prod_{i\in H}a_i\right)
    D\!\left(\left(\prod_{i\notin H}a_i\right)t\right)=0
    \tag{16.8.8.2}
    \label{IV.16.8.8.2}
  \]
  (où $I_{n+1}$ est l'intervalle $1\leq i\leq n+1$ de $\mathbb{N}$).
\end{enumerate}
\oldpage[IV-4]{43}
\end{proposition}

Prouvons d'abord l'équivalence de a) et c). Par définition, pour prouver que
$D$ est un opérateur différentiel d'ordre $\leq n$, il suffit de montrer qu'il
en est ainsi pour la restriction
$D|U:\mathscr{F}|U\to\mathscr{G}|U$ à tout ouvert affine $U$ de $X$, et
d'autre part la propriété c) est valable pour tout ouvert $U$ de $X$ si elle
l'est pour tout ouvert affine. On peut donc se borner au cas où
$S=\operatorname{Spec}(A)$ et $X=\operatorname{Spec}(B)$ sont affines. En
vertu de (16.8.2) (dont on garde les notations), la condition a) signifie
alors que le $A$-homomorphisme $D':B\otimes_A M\to N$ tel que
$D'(b\otimes t)=bD(t)$ s'annule dans
$\mathfrak{I}^{n+1}(B\otimes_A M)$, ce qui, en vertu de (0, 20.4.4), équivaut
à dire que $D'$ s'annule pour tous les éléments de la forme
\[
  \left(\prod_{i=1}^{n+1}(a_i\otimes1-1\otimes a_i)\right).(1\otimes t)
\]
où $a_i\in B$ et $t\in M$. Or, cet élément s'écrit
\[
  \sum_{H\subset I_{n+1}}
  \left(\prod_{i\in H}a_i\right)\otimes
  \left(\left(\prod_{i\notin H}a_i\right)t\right),
\]
et la valeur de $D'$ pour cet élément n'est autre que le premier membre de
(16.8.8.2), ce qui prouve l'équivalence de a) et c).

Prouvons maintenant l'équivalence de b) et c). Raisonnons par récurrence sur
$n$, l'assertion étant triviale pour $n=0$. Écrivant $a_{n+1}$ au lieu de $a$
dans la condition b), on voit, en vertu de l'hypothèse de récurrence, que la
condition b) signifie que pour toute famille $(a_i)_{1\leq i\leq n}$ de $n$
sections de $\mathcal{O}_X$ au-dessus de $U$ et toute section $t$ de
$\mathscr{F}$ au-dessus de $U$, on a
\[
  \sum_{H'\subset I_n}(-1)^{\operatorname{Card}(H')}
  \left(\prod_{i\in H'}a_i\right)
  D_{a_{n+1}}\!\left(\left(\prod_{i\notin H'}a_i\right)t\right)=0.
\]
Mais si l'on remplace dans cette relation $D_{a_{n+1}}$ par sa définition
(16.8.8.1), on constate aussitôt qu'on obtient, au signe près, le premier
membre de (16.8.8.2); d'où la conclusion.

\begin{proposition}[16.8.9]
\label{IV.16.8.9}
Si $D:\mathscr{F}\to\mathscr{G}$ est un opérateur différentiel d'ordre
$\leq n$, et $D':\mathscr{G}\to\mathscr{H}$ un opérateur différentiel d'ordre
$\leq n'$, alors $D'\circ D:\mathscr{F}\to\mathscr{H}$ est un opérateur
différentiel d'ordre $\leq n+n'$.
\end{proposition}

Par hypothèse, on peut écrire $D=u\circ d_{X/S,\mathscr{F}}^n$ et
$D'=v\circ d_{X/S,\mathscr{G}}^{n'}$, où
$u:\mathscr{P}_{X/S}^n\otimes_{\mathcal{O}_X}\mathscr{F}\to\mathscr{G}$ et
$v:\mathscr{P}_{X/S}^{n'}\otimes_{\mathcal{O}_X}\mathscr{G}\to\mathscr{H}$
sont des $\mathcal{O}_X$-homomorphismes. Tout revient à montrer que
l'homomorphisme composé de faisceaux de groupes additifs
\[
  \xymatrix{
    \mathscr{F}\ar[r]^-{d_{X/S,\mathscr{F}}^n}
      &\mathscr{P}_{X/S}^n\otimes_{\mathcal{O}_X}\mathscr{F}
        \ar[r]^-u
      &\mathscr{G}\ar[r]^-{d_{X/S,\mathscr{G}}^{n'}}
      &\mathscr{P}_{X/S}^{n'}\otimes_{\mathcal{O}_X}\mathscr{G}
  }
\]
se factorise en
\[
  \xymatrix{
    \mathscr{F}\ar[r]^-{d_{X/S,\mathscr{F}}^{n+n'}}
      &\mathscr{P}_{X/S}^{n+n'}\otimes_{\mathcal{O}_X}\mathscr{F}
        \ar[r]^-w
      &\mathscr{P}_{X/S}^{n'}\otimes_{\mathcal{O}_X}\mathscr{G}
  }
\]
où $w$ est un $\mathcal{O}_X$-homomorphisme. Il suffira de prouver le

\begin{lemma}[16.8.9.1]
\label{IV.16.8.9.1}
Il existe un $\mathcal{O}_X$-homomorphisme et un seul
\[
  \delta:\mathscr{P}_{X/S}^{n+n'}
  \longrightarrow
  \mathscr{P}_{X/S}^{n'}(\mathscr{P}_{X/S}^n)
  =\mathscr{P}_{X/S}^{n'}\otimes_{\mathcal{O}_X}\mathscr{P}_{X/S}^n
  \tag{16.8.9.2}
  \label{IV.16.8.9.2}
\]
\oldpage[IV-4]{44}
rendant commutatif le diagramme
\[
  \xymatrix{
    \mathcal{O}_X\ar[r]^-{d_{X/S}^{n+n'}}\ar[d]_-{d_{X/S}^n}
      &\mathscr{P}_{X/S}^{n+n'}\ar[d]^-\delta\\
    \mathscr{P}_{X/S}^n
      \ar[r]_-{d_{X/S,\mathscr{P}_{X/S}^n}^{n'}}
      &\mathscr{P}_{X/S}^{n'}(\mathscr{P}_{X/S}^n)
  }
  \tag{16.8.9.3}
  \label{IV.16.8.9.3}
\]
\end{lemma}

On aura alors en effet un diagramme commutatif déduit de (16.8.9.3) par
tensorisation avec $\mathscr{F}$
\[
  \xymatrix{
    \mathscr{F}\ar[r]^-{d_{X/S,\mathscr{F}}^{n+n'}}
      \ar[d]_-{d_{X/S,\mathscr{F}}^n}
      &\mathscr{P}_{X/S}^{n+n'}(\mathscr{F})
        \ar[d]^-{\delta\otimes1}\\
    \mathscr{P}_{X/S}^n(\mathscr{F})
      \ar[r]_-{d_{X/S,\mathscr{P}_{X/S}^n(\mathscr{F})}^{n'}}
      &\mathscr{P}_{X/S}^{n'}(\mathscr{P}_{X/S}^n(\mathscr{F}))
  }
\]
et d'autre part, on vérifie aussitôt sur la définition (16.7.5) que le
diagramme
\[
  \xymatrix{
    \mathscr{P}_{X/S}^n(\mathscr{F})\ar[r]^-u
      \ar[d]_-{d_{X/S,\mathscr{P}_{X/S}^n(\mathscr{F})}^{n'}}
      &\mathscr{G}\ar[d]^-{d_{X/S,\mathscr{G}}^{n'}}\\
    \mathscr{P}_{X/S}^{n'}(\mathscr{P}_{X/S}^n(\mathscr{F}))
      \ar[r]_-{1\otimes u}
      &\mathscr{P}_{X/S}^{n'}(\mathscr{G})
  }
\]
est commutatif. On répondra donc à la question en prenant pour $w$ le
$\mathcal{O}_X$-homomorphisme composé
\[
  \xymatrix{
    \mathscr{P}_{X/S}^{n+n'}(\mathscr{F})
      \ar[r]^-{\delta\otimes1}
      &\mathscr{P}_{X/S}^{n'}(\mathscr{P}_{X/S}^n(\mathscr{F}))
        \ar[r]^-{1\otimes u}
      &\mathscr{P}_{X/S}^{n'}(\mathscr{G}).
  }
\]

Reste à prouver le lemme (16.8.9.1). Compte tenu de (16.7.6), qui prouve
l'unicité de $\delta$, on est ramené au cas où $S=\operatorname{Spec}(A)$ et
$X=\operatorname{Spec}(B)$ sont affines; posant
$\mathfrak{I}=\mathfrak{I}_{B/A}$, il s'agit de définir un homomorphisme
canonique de $B$-modules
\[
  \varphi:(B\otimes_A B)/\mathfrak{I}^{n+n'+1}
  \longrightarrow
  ((B\otimes_A B)/\mathfrak{I}^{n'+1})
  \otimes_B((B\otimes_A B)/\mathfrak{I}^{n+1})
\]
les structures de $B$-module des deux membres provenant du premier facteur
$B$; rappelons que dans le produit tensoriel du second membre,
$(B\otimes_A B)/\mathfrak{I}^{n'+1}$ doit être considéré
\oldpage[IV-4]{45}
comme $B$-module à droite par son second facteur $B$ et
$(B\otimes_A B)/\mathfrak{I}^{n+1}$ comme $B$-module à gauche par son premier
facteur $B$ (16.7.2). Il revient au même de définir un homomorphisme de
$B$-modules
\[
  \varphi_0:B\otimes_A B\longrightarrow
  ((B\otimes_A B)/\mathfrak{I}^{n'+1})
  \otimes_B((B\otimes_A B)/\mathfrak{I}^{n+1})
\]
et de prouver qu'il s'annule dans $\mathfrak{I}^{n+n'+1}$. Or, on définit
aussitôt un tel homomorphisme par la condition que
\[
  \varphi_0(b\otimes b')
  =\pi_{n'}(b\otimes1)\otimes\pi_n(1\otimes b')
  \qquad\text{pour $b,b'$ dans $B$}
\]
avec les notations de (16.3.7). En outre, il est immédiat que $\varphi_0$ est
un homomorphisme d'\emph{anneaux}. Or, on peut écrire
\[
  \begin{aligned}
  \varphi_0(b\otimes1-1\otimes b)
  ={}&\pi_{n'}(b\otimes1-1\otimes b)\otimes\pi_n(1\otimes1)\\
  &+\pi_{n'}(1\otimes b)\otimes\pi_n(1\otimes1)
  -\pi_{n'}(1\otimes1)\otimes\pi_n(1\otimes b)
  \end{aligned}
\]
et on a
\[
  \begin{aligned}
  \pi_{n'}(1\otimes b)\otimes\pi_n(1\otimes1)
  &=\pi_{n'}(1\otimes1)b\otimes\pi_n(1\otimes1)\\
  &=\pi_{n'}(1\otimes1)\otimes b\pi_n(1\otimes1)\\
  &=\pi_{n'}(1\otimes1)\otimes\pi_n(b\otimes1)
  \end{aligned}
\]
d'où finalement
\[
  \begin{aligned}
  \varphi_0(b\otimes1-1\otimes b)
  ={}&\pi_{n'}(b\otimes1-1\otimes b)\otimes\pi_n(1\otimes1)\\
  &+\pi_{n'}(1\otimes1)\otimes\pi_n(b\otimes1-1\otimes b).
  \end{aligned}
  \tag{16.8.9.4}
  \label{IV.16.8.9.4}
\]
Un produit de $n+n'+1$ termes de la forme (16.8.9.4) est donc
nécessairement nul, car il en est ainsi d'un produit de $n+1$ termes de la
forme $\pi_n(b\otimes1-1\otimes b)$ et d'un produit de $n'+1$ termes de la
forme $\pi_{n'}(b\otimes1-1\otimes b)$. La conclusion résulte donc de
(0, 20.4.4).

\begin{corollary}[16.8.10]
\label{IV.16.8.10}
Le faisceau $\mathscr{D}\!iff_{X/S}(\mathcal{O}_X,\mathcal{O}_X)$ (aussi noté
$\mathscr{D}\!iff_{X/S}$) est canoniquement muni d'une structure de faisceau
d'anneaux, les $\mathscr{D}\!iff_{X/S}^n$ formant une filtration croissante
compatible avec cette structure.
\end{corollary}

En particulier, $\mathscr{D}\!iff_{X/S}^0$ est un faisceau de sous-anneaux de
$\mathscr{D}\!iff_{X/S}$, qui s'identifie canoniquement à $\mathcal{O}_X$
(16.8.1). Les formules (16.8.5.1) et (16.8.5.2) montrent que la structure de
$\mathcal{O}_X$-Bimodule de $\mathscr{D}\!iff_{X/S}$ provient de la
multiplication à gauche et à droite par les sections de $\mathcal{O}_X$
considéré comme faisceau de sous-anneaux de $\mathscr{D}\!iff_{X/S}$.

\begin{remarks}[16.8.11]
\label{IV.16.8.11}
\begin{enumerate}
  \item[(i)] Supposons que
  $\mathscr{F}=\bigoplus_{\lambda\in L}\mathscr{F}_\lambda$; alors il est
  clair (16.7.2.1) que
  $\mathscr{P}_{X/S}^n(\mathscr{F})=
  \bigoplus_{\lambda\in L}\mathscr{P}_{X/S}^n(\mathscr{F}_\lambda)$; comme le
  foncteur $\mathscr{F}\mapsto\Gamma(U,\mathscr{F})$ commute à la formation
  des sommes directes quelconques, $d_{X/S,\mathscr{F}}^n$ est
  l'homomorphisme dont la restriction à chaque $\mathscr{F}_\lambda$ est
  \[
    d_{X/S,\mathscr{F}_\lambda}^n:\mathscr{F}_\lambda
    \longrightarrow\mathscr{P}_{X/S}^n(\mathscr{F}_\lambda);
  \]
  on en conclut aussitôt que l'on a
  \[
    \operatorname{Diff}_{X/S}^n(\mathscr{F},\mathscr{G})
    =\prod_{\lambda\in L}
    \operatorname{Diff}_{X/S}^n(\mathscr{F}_\lambda,\mathscr{G}),
  \]
  et par suite aussi (0\textsubscript{I}, 3.2.6)
  \[
    \mathscr{D}\!iff_{X/S}^n(\mathscr{F},\mathscr{G})
    =\prod_{\lambda\in L}
    \mathscr{D}\!iff_{X/S}^n(\mathscr{F}_\lambda,\mathscr{G}).
  \]
  Par ailleurs, si $\mathscr{G}=\prod_{\mu\in M}\mathscr{G}_\mu$
  (0\textsubscript{I}, 3.2.6), on a
  \[
    \operatorname{Hom}_{\mathcal{O}_X}
    (\mathscr{P}_{X/S}^n(\mathscr{F}),\mathscr{G})
    =\prod_{\mu\in M}\operatorname{Hom}_{\mathcal{O}_X}
    (\mathscr{P}_{X/S}^n(\mathscr{F}),\mathscr{G}_\mu),
  \]
  \oldpage[IV-4]{46}
  tout homomorphisme $u$ de $\mathscr{P}_{X/S}^n(\mathscr{F})$ dans
  $\mathscr{G}$ correspondant biunivoquement à la famille de ses composés
  \[
    u_\mu:\mathscr{P}_{X/S}^n(\mathscr{F})
    \longrightarrow\mathscr{G}\longrightarrow\mathscr{G}_\mu.
  \]
  On a donc
  \[
    \operatorname{Diff}_{X/S}^n(\mathscr{F},\mathscr{G})
    =\prod_{\mu\in M}
    \operatorname{Diff}_{X/S}^n(\mathscr{F},\mathscr{G}_\mu),
  \]
  et par suite aussi
  \[
    \mathscr{D}\!iff_{X/S}^n(\mathscr{F},\mathscr{G})
    =\prod_{\mu\in M}
    \mathscr{D}\!iff_{X/S}^n(\mathscr{F},\mathscr{G}_\mu).
  \]
  \item[(ii)] Jusqu'ici, on n'a guère rencontré d'opérateurs différentiels
  $\mathscr{F}\to\mathscr{G}$ que lorsque $\mathscr{F}$ et $\mathscr{G}$ sont
  des $\mathcal{O}_X$-Modules localement libres de rang fini, auquel cas, leur
  structure se ramène localement, en vertu de (i), à celle du faisceau
  $\mathscr{D}\!iff_{X/S}$; cette dernière va être étudiée plus loin (16.11)
  dans un cas particulier.
\end{enumerate}
\end{remarks}

\subsection{Immersions régulières et quasi-régulières.}
\EGAReaderTarget{ega4.c04.subsec.16-9}{IV.16.9}

\begin{definition}[16.9.1]
\EGAReaderTarget{ega4.c04.def.16-9-1}{IV.16.9.1}
Soit $X$ un espace annelé. On dit qu'un Idéal $\mathscr{I}$ de
$\mathcal{O}_X$ est \emph{régulier} (resp. \emph{quasi-régulier}) si, pour
tout point $x\in\operatorname{Supp}(\mathcal{O}_X/\mathscr{I})$, il existe un
voisinage ouvert $U$ de $x$ dans $X$ et une suite régulière (0, 15.2.2) (resp.
quasi-régulière (0, 15.2.2)) d'éléments de $\Gamma(U,\mathcal{O}_X)$ qui
engendre $\mathscr{I}|U$.
\end{definition}

On dira qu'une suite régulière (resp. quasi-régulière) de sections de
$\mathcal{O}_X$ au-dessus de $U$ qui engendre $\mathscr{I}|U$ est un
\emph{système régulier} (resp. \emph{quasi-régulier}) \emph{de générateurs}
de $\mathscr{I}|U$.

\begin{definition}[16.9.2]
\EGAReaderTarget{ega4.c04.def.16-9-2}{IV.16.9.2}
Soit $j:Y\to X$ une immersion de préschémas et soit $U$ un ouvert de $X$ tel
que $j(Y)\subset U$ et que $j$ soit une immersion fermée de $Y$ dans $U$. On
dit que $j$ est \emph{régulière} (resp. \emph{quasi-régulière}) si le
sous-préschéma fermé $j(Y)$ de $U$ associé à $j$ est défini par un Idéal
régulier (resp. quasi-régulier) de $\mathcal{O}_U$ (condition indépendante de
l'ouvert $U$ choisi).
\end{definition}

On dit qu'un sous-préschéma $Y$ d'un préschéma $X$ est \emph{régulièrement
immergé} (resp. \emph{quasi-régulièrement immergé}) si l'injection canonique
$j:Y\to X$ est une immersion régulière (resp. quasi-régulière). Si $Y$ est un
sous-préschéma fermé de $X$ et $\mathscr{I}$ l'Idéal de $\mathcal{O}_X$ qui
définit $Y$, il revient au même de dire que $\mathscr{I}$ est \emph{régulier}
(resp. \emph{quasi-régulier}).

Par exemple, si $A$ est un anneau \emph{intègre}, $f$ un élément $\neq0$ de
$A$, le sous-préschéma fermé $V(f)$ de $\operatorname{Spec}(A)$ (isomorphe à
$\operatorname{Spec}(A/fA)$) est \emph{régulièrement immergé} dans
$\operatorname{Spec}(A)$.

Tout Idéal régulier est quasi-régulier (0, 15.2.2); toute immersion régulière
est quasi-régulière (cf. (16.9.11) pour une réciproque).

\begin{proposition}[16.9.3]
\label{IV.16.9.3}
Soient $X$ un espace annelé, $\mathscr{I}$ un Idéal de $\mathcal{O}_X$,
$(f_i)_{1\leq i\leq m}$ une suite finie de sections de $\mathcal{O}_X$
au-dessus de $X$ engendrant $\mathscr{I}$. Pour que $(f_i)$ soit une suite
quasi-régulière (0, 15.2.2), il faut et il suffit que les conditions suivantes
soient vérifiées:
\begin{enumerate}
  \item[(i)] Les images canoniques des $f_i$ dans
  $\mathscr{I}/\mathscr{I}^2$ forment une base de cet
  $\mathcal{O}_X/\mathscr{I}$-Module.
  \item[(ii)] L'homomorphisme surjectif canonique (16.1.2.2)
  \[
    \mathbf{S}_{\mathcal{O}_X/\mathscr{I}}^{\bullet}
    (\mathscr{I}/\mathscr{I}^2)
    \longrightarrow\mathscr{G}r_{\mathscr{I}}^{\bullet}(\mathcal{O}_X)
  \]
  est bijectif.
\end{enumerate}
En outre, s'il en est ainsi, toute suite $(f'_i)_{1\leq i\leq n}$ de $n$
sections de $\mathscr{I}$ au-dessus de $X$ qui engendre $\mathscr{I}$ est
quasi-régulière.
\end{proposition}

\oldpage[IV-4]{47}

Les deux conditions de l'énoncé ne font que traduire la définition donnée dans
(0, 15.2.2), compte tenu de la définition des homomorphismes canoniques
(0, 15.2.1.1). La dernière assertion résulte de ce que, si un module $M$ sur
un anneau commutatif $A$ admet une base de $n$ éléments, tout système de
générateurs de $M$ ayant $n$ éléments est une base de $M$ (Bourbaki,
\emph{Alg. comm.}, chap.~II, §~3, cor.~5 du th.~1).

\begin{corollary}[16.9.4]
\label{IV.16.9.4}
Soient $X$ un espace annelé en anneaux locaux, $\mathscr{I}$ un Idéal de
$\mathcal{O}_X$. Pour que $\mathscr{I}$ soit quasi-régulier, il faut et il
suffit qu'il vérifie les conditions suivantes:
\begin{enumerate}
  \item[(i)] $\mathscr{I}$ est de type fini.
  \item[(ii)] $\mathscr{I}/\mathscr{I}^2$ est un
  $(\mathcal{O}_X/\mathscr{I})$-Module localement libre.
  \item[(iii)] L'homomorphisme canonique
  \[
    \mathbf{S}_{\mathcal{O}_X/\mathscr{I}}^{\bullet}
    (\mathscr{I}/\mathscr{I}^2)
    \longrightarrow\mathscr{G}r_{\mathscr{I}}^{\bullet}(\mathcal{O}_X)
    \tag{16.9.4.1}
    \label{IV.16.9.4.1}
  \]
  est bijectif.
\end{enumerate}
\end{corollary}

La nécessité des conditions résulte aussitôt de (16.9.3). Pour voir que ces
conditions sont suffisantes, il faut montrer, en vertu de (16.9.3), que si, en
un point $x\in\operatorname{Supp}(\mathcal{O}_X/\mathscr{I})$, il existe un
voisinage ouvert $U$ de $x$ dans $X$ et $n$ sections $f_i$ ($1\leq i\leq n$)
de $\mathscr{I}$ au-dessus de $U$ dont les images canoniques dans
$\mathscr{I}/\mathscr{I}^2$ forment une base de
$(\mathscr{I}/\mathscr{I}^2)|U$ sur
$(\mathcal{O}_X/\mathscr{I})|U$, alors il y a un voisinage ouvert
$V\subset U$ de $x$ tel que les $f_i|V$ engendrent $\mathscr{I}|V$. Or, par
hypothèse, on a $\mathscr{I}_x\neq\mathcal{O}_x$, donc $\mathscr{I}_x$ est
contenu dans l'idéal maximal de $\mathcal{O}_x$; comme $\mathscr{I}_x$ est un
$\mathcal{O}_x$-module de type fini et que les classes des $(f_i)_x$ dans
$\mathscr{I}_x/\mathscr{I}_x^2$ engendrent cet
$(\mathcal{O}_x/\mathscr{I}_x)$-module, le lemme de Nakayama montre que les
$(f_i)_x$ engendrent $\mathscr{I}_x$. Comme $\mathscr{I}$ est de type fini,
on conclut par (0\textsubscript{I}, 5.2.2).

\begin{corollary}[16.9.5]
\label{IV.16.9.5}
Soient $X$ un espace annelé en anneaux locaux, $\mathscr{I}$ un Idéal
quasi-régulier de $\mathcal{O}_X$, $(f_i)_{1\leq i\leq n}$ une suite de
sections de $\mathscr{I}$ au-dessus de $X$,
$x$ un point de $\operatorname{Supp}(\mathcal{O}_X/\mathscr{I})$. Les
conditions suivantes sont équivalentes:
\begin{enumerate}
  \item[(a)] Il existe un voisinage ouvert $U$ de $x$ dans $X$ tel que les
  $f_i|U$ forment une suite quasi-régulière d'éléments de
  $\Gamma(U,\mathcal{O}_X)$ engendrant $\mathscr{I}|U$.
  \item[(b)] Les $(f_i)_x$ forment un système de générateurs de
  $\mathscr{I}_x$ dont le nombre d'éléments est le plus petit possible.
  \item[(b')] Les $(f_i)_x$ forment un système minimal de générateurs de
  $\mathscr{I}_x$.
  \item[(c)] Si $\bar f_i$ est l'image canonique de $f_i$ dans
  $\Gamma(X,\mathscr{I}/\mathscr{I}^2)$, les $(\bar f_i)_x$ forment une base
  du $(\mathcal{O}_x/\mathscr{I}_x)$-module
  $\mathscr{I}_x/\mathscr{I}_x^2$.
\end{enumerate}
\end{corollary}

Par hypothèse, $\mathcal{O}_x$ est un anneau local, $\mathscr{I}_x$ un idéal
de type fini de $\mathcal{O}_x$ contenu dans l'idéal maximal de
$\mathcal{O}_x$; l'équivalence de b), b') et c) résulte donc du lemme de
Nakayama (Bourbaki, \emph{Alg. comm.}, chap.~II, §~3,
n\textsuperscript{o}~2, prop.~5). Il est clair que a) entraîne c) en vertu de
(16.9.3); d'autre part, il résulte de (0\textsubscript{I}, 5.2.2) que si la
condition c) est vérifiée (donc aussi b)), il y a un voisinage $U$ de $x$ dans
$X$ tel que $(\mathscr{I}/\mathscr{I}^2)|U$ ait un rang constant égal à $n$,
et que les $f_i|U$ engendrent $\mathscr{I}|U$; il suffit donc d'appliquer dans
$U$ la dernière assertion de (16.9.3).

\begin{remarks}[16.9.6]
\label{IV.16.9.6}
\begin{enumerate}
  \item[(i)] Sous les hypothèses générales de (16.9.5), il ne suffit pas que
  les $(\bar f_i)_y$ forment une base du
  $(\mathcal{O}_y/\mathscr{I}_y)$-module
  $\mathscr{I}_y/\mathscr{I}_y^2$ pour tout $y\in X$ pour que la suite
  $(f_i)$ engendre $\mathscr{I}$. On en a un exemple en
\oldpage[IV-4]{48}
  prenant $X=\operatorname{Spec}(A)$, où $A$ est un anneau de Dedekind, et
  $\mathscr{I}=\widetilde{\mathfrak{I}}$, où $\mathfrak{I}$ est un idéal
  premier non principal de $A$; on a alors en effet
  $\mathscr{I}_y/\mathscr{I}_y^2=0$ en tout point $y$ distinct du point
  $x\in X$ correspondant à $\mathfrak{I}$, et
  $\mathscr{I}_x/\mathscr{I}_x^2$ est de rang $1$ sur le corps
  $\mathcal{O}_x/\mathscr{I}_x$; en outre $\mathscr{I}$ est évidemment un
  Idéal régulier.
  \item[(ii)] Dans (16.9.5), on ne peut remplacer « quasi-régulière » par
  « régulière », même lorsque $X$ est un préschéma (cf. (16.9.12)).
  Désignons en effet par $B$ l'anneau des germes de fonctions indéfiniment
  différentiables au point $0$ dans $\mathbb{R}$; il a un idéal maximal
  $\mathfrak{m}$ engendré par le germe $t$ de l'application identique de
  $\mathbb{R}$ au point $0$, et l'intersection $\mathfrak{n}$ des
  $\mathfrak{m}^k$ pour $k>0$ n'est pas réduite à $0$. Soit maintenant $A$
  l'anneau quotient $B[T]/\mathfrak{n}TB[T]$, et soient $f_1,f_2$ les images
  canoniques dans $A$ des éléments $t$ et $T$ de $B[T]$. La suite
  $(f_1,f_2)$ est régulière dans $A$: en effet, $f_1$ n'est pas diviseur de
  $0$ dans $A$, car la relation $tP[T]\in\mathfrak{n}TB[T]$, pour un
  polynôme $P\in B[T]$, entraîne que les produits de $t$ par les coefficients
  de $P$ appartiennent à l'idéal $\mathfrak{n}$, et il en résulte aussitôt que
  ces coefficients sont eux-mêmes dans $\mathfrak{n}$, donc
  $P[T]\in\mathfrak{n}TB[T]$. Comme $B/tB$ est isomorphe à $\mathbb{R}$,
  $A/f_1A$ est isomorphe à l'anneau de polynômes $\mathbb{R}[T]$, donc
  intègre, et l'image de $f_2$ dans $A/f_1A$, étant égale à $T$, n'est donc
  pas diviseur de $0$, ce qui prouve notre assertion. Cependant $f_2$ est
  diviseur de $0$ dans $A$, car pour tout élément $x\in\mathfrak{n}$ non nul,
  l'image de $x$ dans $A$ est $\ne0$, mais l'image de $xT$ est nulle. On en
  conclut que la suite $(f_2,f_1)$ n'est pas régulière dans $A$; d'autre part
  l'idéal $\mathfrak{I}=f_1A+f_2A$ est distinct de $A$, donc les conditions
  b), b') et c) de (16.9.5) n'entraînent pas la condition a) lorsqu'on y
  remplace « quasi-régulière » par « régulière ».
\end{enumerate}
\end{remarks}

\begin{env}[16.9.7]
\EGAReaderTarget{ega4.c04.statement.16-9-7}{IV.16.9.7}
Si $X=\operatorname{Spec}(A)$ est un schéma affine, nous dirons qu'un idéal
$\mathfrak{I}$ de $A$ est régulier (resp. quasi-régulier) si l'Idéal
$\mathscr{I}=\widetilde{\mathfrak{I}}$ de $\mathcal{O}_X$ est régulier
(resp. quasi-régulier); on notera que cette notion est locale et n'implique
nullement l'existence d'un système de générateurs de $\mathfrak{I}$ formant
dans $A$ une suite régulière (resp. quasi-régulière) comme le montre l'exemple
(16.9.6, (i)); il en est toutefois ainsi lorsque $A$ est local (16.9.5).

La proposition (16.9.4) se traduit en termes d'immersions quasi-régulières de
la façon suivante:
\end{env}

\begin{proposition}[16.9.8]
\label{IV.16.9.8}
Soit $j:Y\to X$ un morphisme de préschémas; pour que $j$ soit une immersion
quasi-régulière, il faut et il suffit que $j$ satisfasse aux conditions
suivantes:
\begin{enumerate}
  \item[(i)] $j$ est une immersion localement de présentation finie.
  \item[(ii)] Le faisceau conormal
  $\operatorname{Gr}^1(j)=\mathcal{N}_{Y/X}$ (16.1.2) est un
  $\mathcal{O}_Y$-Module localement libre.
  \item[(iii)] L'homomorphisme canonique
  \[
    \operatorname{S}_{\mathcal{O}_Y}^{\bullet}
      (\operatorname{Gr}^1(j))\longrightarrow\operatorname{Gr}^{\bullet}(j)
  \]
  (16.1.2.2) est bijectif.
\end{enumerate}
\end{proposition}

La question étant locale sur $Y$, on peut se borner au cas où $j$ est
l'injection canonique d'un sous-préschéma fermé $Y$ de $X$, et alors la
traduction de (16.9.4) en (16.9.8) résulte de l'explicitation de
$\operatorname{Gr}^1(j)$ et $\operatorname{Gr}^{\bullet}(j)$ en termes de
l'Idéal $\mathscr{I}$ de $\mathcal{O}_X$ définissant le sous-préschéma $Y$
(16.1.3, (ii)).

\begin{corollary}[16.9.9]
\label{IV.16.9.9}
Soient $Y$ un préschéma, $X$ un $Y$-préschéma, $j:Y\to X$ une $Y$-section
de $X$, de sorte que le $n$-ème invariant normal $\mathscr{A}^{(n)}$ de $j$
(16.1.2) est une $\mathcal{O}_Y$-Algèbre augmentée (16.1.7); posons
$\mathscr{A}^{(\infty)}=\varprojlim\mathscr{A}^{(n)}$. Pour que $j$ soit une
immersion quasi-régulière, il faut et il suffit que $j$ soit localement de
présentation finie, et que tout $y\in Y$ admette un voisinage ouvert affine
$U$ d'anneau $C$ tel que $\mathscr{A}^{(\infty)}|U$ soit isomorphe comme
$\mathcal{O}_U$-Algèbre topologique augmentée à
$\mathcal{O}_U[[T_1,\ldots,T_n]]$.
\end{corollary}

On peut se borner au cas où $j$ est une immersion fermée en se restreignant
à un voisinage assez petit de $y$ (voir le raisonnement de (16.4.11)), et
alors $\mathcal{O}_Y$ s'identifie à une Algèbre quotient
$\mathcal{O}_X/\mathscr{I}$ et l'homomorphisme surjectif canonique
$\mathcal{O}_X\to\mathcal{O}_Y$ admet un
\oldpage[IV-4]{49}
inverse à droite (16.1.7). On peut donc supposer
$X=\operatorname{Spec}(B)$ et $Y=\operatorname{Spec}(A)$ affines, $B$ étant
une $A$-algèbre augmentée, et l'idéal d'augmentation $\mathfrak{I}$ étant de
type fini. Comme $\mathscr{A}^{(n)}$ s'identifie alors à
$(B/\mathfrak{I}^{n+1})^{\sim}$, le corollaire résulte de l'équivalence de b)
et c) dans (0, 19.5.4) puisque $B/\mathfrak{I}=A$.

On notera que, dans le cas affine envisagé, le fait que $j$ soit une immersion
quasi-régulière équivaut encore, en vertu de (0, 19.5.4), à dire que $B$ est
une $A$-algèbre formellement lisse pour la topologie
$\mathfrak{I}$-pré-adique.

On notera aussi que la condition que $j$ soit une immersion localement de
présentation finie est toujours vérifiée lorsque le morphisme $X\to Y$ est
localement de type fini (1.4.3, (v)).

\begin{proposition}[16.9.10]
\EGAReaderTarget{ega4.c04.prop.16-9-10}{IV.16.9.10}
Soient $X$ un préschéma localement noethérien, $Y$ un sous-préschéma de $X$,
$j:Y\to X$ l'injection canonique, $y$ un point de $Y$.
\begin{enumerate}
  \item[(i)] Pour qu'il existe un voisinage ouvert $U$ de $y$ dans $X$ tel que
  la restriction $Y\cap U\to U$ de $j$ soit une immersion régulière, il faut
  et il suffit que le noyau $\mathscr{I}_y$ de l'homomorphisme surjectif
  $\mathcal{O}_{X,y}\to\mathcal{O}_{Y,y}$ soit engendré par une suite régulière
  d'éléments de $\mathcal{O}_{X,y}$.
  \item[(ii)] Pour que l'immersion $j$ soit régulière, il faut et il suffit
  qu'elle soit quasi-régulière.
\end{enumerate}
\end{proposition}

\begin{enumerate}
  \item[(i)] On peut se borner au cas où $Y$ est un sous-préschéma fermé de
  $X$ défini par un Idéal \emph{cohérent} $\mathscr{I}$ de $\mathcal{O}_X$.
  La condition est évidemment nécessaire. Inversement, si $\mathscr{I}_y$
  est engendré par une suite régulière $(s_i)_y$, où les $s_i$ sont des
  sections de $\mathscr{I}$ au-dessus d'un voisinage ouvert $U$ de $y$ dans
  $X$, on peut supposer que les $s_i$ engendrent $\mathscr{I}|U$
  (0\textsubscript{I}, 5.2.2) et forment une suite régulière (0, 15.2.4),
  d'où l'assertion.
  \item[(ii)] Le fait qu'une immersion quasi-régulière soit régulière résulte
  de (i) et de l'identité des suites quasi-régulières et des suites régulières
  dans $\mathcal{O}_{X,y}$, formées d'éléments de l'idéal maximal
  (0, 15.1.11).
\end{enumerate}

Lorsque (sans hypothèse noethérienne sur $X$) le noyau $\mathscr{I}_y$ de
$\mathcal{O}_{X,y}\to\mathcal{O}_{Y,y}$ est engendré par une suite régulière
d'éléments de $\mathcal{O}_{X,y}$, on dit que l'immersion $j$ est
\emph{régulière au point $y$}.

\begin{corollary}[16.9.11]
\label{IV.16.9.11}
Soit $X$ un préschéma localement noethérien; alors tout Idéal quasi-régulier de
$\mathcal{O}_X$ est régulier.
\end{corollary}

\begin{remarks}[16.9.12]
\label{IV.16.9.12}
\begin{enumerate}
  \item[(i)] On notera qu'une immersion régulière n'est pas en général un
  morphisme plat, ni \emph{a fortiori} un morphisme \emph{régulier} au sens de
  (6.8.1).
  \item[(ii)] Soit $A$ un anneau local noethérien; il résulte aussitôt de
  (16.9.4) et de (0, 17.1.1) que pour que $A$ soit \emph{régulier}, il faut et
  il suffit que son idéal maximal $\mathfrak{m}$ soit \emph{quasi-régulier}
  (ou \emph{régulier}, ce qui revient au même puisque $A$ est noethérien).
  Pour qu'un schéma affine noethérien $X$ soit \emph{régulier}, il faut et il
  suffit que pour tout point fermé $x\in X$, l'injection canonique
  $\operatorname{Spec}(k(x))\to X$ soit une immersion \emph{régulière}.
\end{enumerate}
\end{remarks}

\begin{proposition}[16.9.13]
\label{IV.16.9.13}
Soient $X$ un préschéma localement noethérien, $Y$ un sous-préschéma de $X$,
$Y'$ un sous-préschéma de $Y$, tel que l'injection canonique $j:Y'\to Y$ soit
régulière. Alors la suite de $\mathcal{O}_{Y'}$-Modules
\[
  \xymatrix{
    0\ar[r] & j^*(\mathcal{N}_{Y/X})\ar[r]
      & \mathcal{N}_{Y'/X}\ar[r] & \mathcal{N}_{Y'/Y}\ar[r] & 0
  }
  \tag{16.9.13.1}
  \label{IV.16.9.13.1}
\]
\oldpage[IV-4]{50}
est exacte; en outre, pour tout $x\in X$, il existe un voisinage ouvert $U$
de $x$ tel que les restrictions à $U$ des homomorphismes de (16.9.13.1)
forment une suite exacte et scindée.
\end{proposition}

Prouvons d'abord le lemme suivant:

\begin{lemma}[16.9.13.2]
\label{IV.16.9.13.2}
Soient $A$ un anneau, $\mathfrak{I}$ un idéal de $A$,
$A'=A/\mathfrak{I}$, $(f_i)_{1\leq i\leq r}$ une suite d'éléments de $A$ qui
est $A'$-régulière, $\mathfrak{K}=\sum_i f_iA$,
$\mathfrak{L}=\mathfrak{I}+\mathfrak{K}$,
$\mathfrak{K}'=\sum_i f_iA'$, de sorte que $C=A/\mathfrak{L}$ est isomorphe à
$A'/\mathfrak{K}'$. Alors pour tout entier $n>0$ et tout entier $N\geq n$,
on a la relation
\[
  \mathfrak{I}\cap\mathfrak{K}^n
  =\mathfrak{I}\mathfrak{K}^n+(\mathfrak{I}\cap\mathfrak{K}^N).
  \tag{16.9.13.3}
  \label{IV.16.9.13.3}
\]
\end{lemma}

Il suffit évidemment de prouver que tout élément du premier membre est
contenu dans le second, et, par récurrence sur $n$, on se ramène au cas où
$N=n+1$. Un élément du premier membre de (16.9.13.3), étant dans
$\mathfrak{K}^n$, s'écrit $P(f_1,\ldots,f_r)$, où
$P\in A[T_1,\ldots,T_r]$ est homogène et de degré $n$. Si $f_i'$ est l'image
canonique de $f_i$ dans $A'$, l'hypothèse
$P(f_1,\ldots,f_r)\in\mathfrak{I}$ signifie que
$P(f_1',\ldots,f_r')=0$. Mais
$P(f_1',\ldots,f_r')\in\mathfrak{K}'^n$, donc l'image canonique de
$P(f_1',\ldots,f_r')$ dans
$\mathfrak{K}'^n/\mathfrak{K}'^{n+1}$ est nulle. Or l'hypothèse que la suite
$(f_i)$ est $A'$-régulière implique que l'homomorphisme canonique
\[
  \mathbf{S}_C^n(\mathfrak{K}'/\mathfrak{K}'^2)
  \longrightarrow\mathfrak{K}'^n/\mathfrak{K}'^{n+1}
\]
est bijectif (0, 15.1.9); on conclut que les coefficients de $P$ appartiennent
à $\mathfrak{L}=\mathfrak{I}+\mathfrak{K}$. Il en résulte aussitôt que l'on a
$P(f_1,\ldots,f_r)\in\mathfrak{I}\mathfrak{K}^n+\mathfrak{K}^{n+1}$, et comme
$P(f_1,\ldots,f_r)\in\mathfrak{I}$, on a finalement
\[
  P(f_1,\ldots,f_r)
  \in\mathfrak{I}\mathfrak{K}^n
    +(\mathfrak{I}\cap\mathfrak{K}^{n+1}),
\]
ce qui prouve le lemme.

En prenant les quotients des deux membres de (16.9.13.3) par
$\mathfrak{I}\mathfrak{K}^n$, on voit que les relations (16.9.13.3) pour
$N\geq n$ entraînent
\[
  (\mathfrak{I}\cap\mathfrak{K}^n)/\mathfrak{I}\mathfrak{K}^n
  \subset\bigcap_{N\geq n}\mathfrak{K}^N
    \cdot(A/(\mathfrak{I}\mathfrak{K}^n)).
  \tag{16.9.13.4}
  \label{IV.16.9.13.4}
\]
On en déduit le

\begin{corollary}[16.9.13.5]
\label{IV.16.9.13.5}
Supposons vérifiées les hypothèses de (16.9.13.2) et supposons en outre que
l'anneau $A$ soit noethérien et que $\mathfrak{K}$ soit contenu dans le
radical de $A$. Alors on a pour tout entier $n>0$,
\[
  \mathfrak{I}\cap\mathfrak{K}^n=\mathfrak{I}\mathfrak{K}^n.
  \tag{16.9.13.6}
  \label{IV.16.9.13.6}
\]
\end{corollary}

En effet, le second membre de (16.9.13.4) est alors nul, puisque
$A/\mathfrak{I}\mathfrak{K}^n$ est un $A$-module de type fini (Bourbaki,
\emph{Alg. comm.}, chap.~III, §~3, n\textsuperscript{o}~3, prop.~6).

Prenons en particulier $n=2$ dans (16.9.13.6), et remarquons que l'on a
$\mathfrak{L}^2=\mathfrak{I}^2+\mathfrak{I}\mathfrak{K}+\mathfrak{K}^2
=\mathfrak{I}\mathfrak{L}+\mathfrak{K}^2$; puisque
$\mathfrak{I}\mathfrak{L}\subset\mathfrak{L}^2$, on en déduit que
\[
  \mathfrak{I}\cap\mathfrak{L}^2
  =\mathfrak{I}\mathfrak{L}+(\mathfrak{I}\cap\mathfrak{K}^2)
  =\mathfrak{I}\mathfrak{L}+\mathfrak{I}\mathfrak{K}^2
  =\mathfrak{I}\mathfrak{L},
\]
autrement dit
\[
  \mathfrak{I}\cap\mathfrak{L}^2=\mathfrak{I}\mathfrak{L},
  \tag{16.9.13.7}
  \label{IV.16.9.13.7}
\]
ce qu'on peut encore exprimer en disant que l'homomorphisme canonique
\[
  \mathfrak{I}/\mathfrak{I}\mathfrak{L}
  \longrightarrow(\mathfrak{I}+\mathfrak{L}^2)/\mathfrak{L}^2
\]
est bijectif.

Ces lemmes étant démontrés, prouvons la première assertion de (16.9.13): il
suffit
\oldpage[IV-4]{51}
évidemment de prouver que la suite des fibres des faisceaux figurant dans
(16.9.13.1), en un point $x\in Y'$, est exacte. Or, si l'on pose
$A=\mathcal{O}_{X,x}$, on peut écrire
$\mathcal{O}_{Y,x}=A'=A/\mathfrak{I}$ où $\mathfrak{I}$ est un idéal contenu
dans l'idéal maximal de $A$, puis
$\mathcal{O}_{Y',x}=A'/\mathfrak{K}'$, où $\mathfrak{K}'$ est engendré par une
suite $A'$-régulière d'éléments de $A'$, eux-mêmes images des éléments d'une
suite $A'$-régulière d'éléments de $A$ appartenant à l'idéal maximal de $A$.
Si $\mathfrak{K}$ est l'idéal engendré par ces derniers et
$\mathfrak{L}=\mathfrak{I}+\mathfrak{K}$, on a
$\mathcal{O}_{Y',x}=A/\mathfrak{L}$, et comme on est dans la situation de
(16.9.13.5), l'homomorphisme canonique
\[
  \mathfrak{I}/\mathfrak{I}\mathfrak{L}
  \longrightarrow(\mathfrak{I}+\mathfrak{L}^2)/\mathfrak{L}^2
\]
est bijectif. Mais cela montre que la suite
\[
  0\longrightarrow\mathfrak{I}/\mathfrak{I}\mathfrak{L}
  \longrightarrow\mathfrak{L}/\mathfrak{L}^2
  \longrightarrow(\mathfrak{L}/\mathfrak{I})/
    (\mathfrak{L}/\mathfrak{I})^2\longrightarrow0
\]
est exacte (voir la démonstration de (16.2.7)), et les modules figurant dans
cette suite sont précisément les fibres en $x$ des faisceaux de (16.9.13.1).
La seconde assertion résulte de ce que $\mathcal{N}_{Y'/Y}$ est un
$\mathcal{O}_{Y'}$-Module localement \emph{libre} (16.9.8) et de Bourbaki,
\emph{Alg.}, chap.~II, 3\textsuperscript{e} éd., §~1, n\textsuperscript{o}~11,
prop.~21.

\subsection{Morphismes différentiellement lisses}
\EGAReaderTarget{ega4.c04.subsec.16-10}{IV.16.10}

\begin{definition}[16.10.1]
\EGAReaderTarget{ega4.c04.def.16-10-1}{IV.16.10.1}
On dit qu'un morphisme de préschémas $f:X\to S$ est
\emph{différentiellement lisse} (ou que $X$ est \emph{différentiellement
lisse sur $S$}) s'il vérifie les conditions suivantes:
\begin{enumerate}
  \item[(i)] $\Omega_{X/S}^1$ est un $\mathcal{O}_X$-Module localement
  projectif, c'est-à-dire que tout point de $X$ admet un voisinage ouvert
  affine $U$ tel que $\Gamma(U,\Omega_{X/S}^1)$ soit un
  $\Gamma(U,\mathcal{O}_X)$-module projectif (non nécessairement de type
  fini).
  \item[(ii)] L'homomorphisme canonique (16.3.1.1)
  \[
    \mathbf{S}_{\mathcal{O}_X}^{\bullet}(\Omega_{X/S}^1)
    \longrightarrow\mathscr{G}r_{\bullet}(\mathscr{P}_{X/S})
  \]
  est bijectif.
\end{enumerate}
En particulier, si $\Omega_{X/S}^1$ est localement libre de rang fini, les
$\mathscr{P}_{X/S}^n$ sont des $\mathcal{O}_X$-Modules localement libres de
rang fini (étant des extensions de tels Modules).
\end{definition}

On dit que $f$ est \emph{différentiellement lisse en un point $x\in X$} (ou
que $X$ est différentiellement lisse sur $S$ au point $x$) s'il existe un
voisinage ouvert $U$ de $x$ dans $X$ tel que $f|U$ soit différentiellement
lisse.

Nous verrons plus loin (17.12.4) qu'un morphisme lisse est différentiellement
lisse, ce qui justifie la terminologie; mais la réciproque est inexacte; en
effet, un \emph{monomorphisme} $f:X\to S$ est différentiellement lisse,
puisque $\Omega_{X/S}^1=0$ en vertu de (I, 5.3.8), et par suite
l'homomorphisme surjectif (16.3.1.1) est évidemment bijectif; or un
monomorphisme n'est même pas nécessairement plat, ni \emph{a fortiori} lisse.
Bornons-nous ici à noter la proposition suivante:

\begin{proposition}[16.10.2]
\label{IV.16.10.2}
Soient $A$ un anneau, $B$ une $A$-algèbre formellement lisse pour les
topologies discrètes (0, 19.3.1). Alors $\operatorname{Spec}(B)$ est
différentiellement lisse sur $\operatorname{Spec}(A)$.
\end{proposition}

En effet, $B\otimes_A B$ est alors (pour les topologies discrètes) une
$B$-algèbre formellement lisse (pour l'un ou l'autre des homomorphismes
canoniques $b\mapsto b\otimes1$, $b\mapsto1\otimes b$ de $B$
\oldpage[IV-4]{52}
dans $B\otimes_A B$) (0, 19.3.5, (iii)); donc $B\otimes_A B$ est aussi une
$A$-algèbre formellement lisse pour les topologies discrètes (0, 19.3.5,
(ii)). Posant $\mathfrak{I}=\mathfrak{I}_{B/A}$, il en résulte que
$B\otimes_A B$ est aussi une $A$-algèbre formellement lisse pour la topologie
$\mathfrak{I}$-préadique (0, 19.3.8); comme par hypothèse
$B=(B\otimes_A B)/\mathfrak{I}$ est une $A$-algèbre formellement lisse pour
les topologies discrètes, la proposition résulte de l'équivalence de a) et b)
dans (0, 19.5.4).

\begin{proposition}[16.10.3]
\label{IV.16.10.3}
Pour qu'un morphisme $f:X\to S$ soit différentiellement lisse, il faut et il
suffit que pour tout $x\in X$, il existe un voisinage ouvert affine $U$ de
$x$, d'anneau $A$, tel que $\Gamma(U,\mathscr{P}_{X/S}^{\infty})$ soit une
$A$-algèbre topologique augmentée isomorphe à l'algèbre complétée $\widehat B$,
où $B=\mathbf{S}_A(V)$, $V$ étant un $A$-module projectif et $B$ étant muni de
la topologie $B^+$-préadique (où $B^+$ est l'idéal d'augmentation). Si
$\Omega_{X/S}^1$ est localement libre de rang fini, on peut remplacer
$\widehat B$ par l'algèbre de séries formelles $A[[T_1,\ldots,T_n]]$.
\end{proposition}

La notion de morphisme différentiellement lisse étant évidemment locale sur
$X$, on peut se borner au cas où $S=\operatorname{Spec}(B)$,
$X=\operatorname{Spec}(C)$. Considérons $C\otimes_B C$ comme une $C$-algèbre
(pour le premier facteur); posons $\mathfrak{I}=\mathfrak{I}_{C/B}$ et
munissons $C\otimes_B C$ de la topologie $\mathfrak{I}$-préadique; on peut
appliquer à la $C$-algèbre topologique $C\otimes_B C$ et à l'idéal
$\mathfrak{I}$ de $C\otimes_B C$ l'équivalence de b) et c) dans (0, 19.5.4),
puisque $(C\otimes_B C)/\mathfrak{I}=C$ est évidemment une $C$-algèbre
formellement lisse pour les topologies discrètes. La topologie sur
$\Gamma(U,\mathscr{P}_{X/S}^{\infty})$ est évidemment la topologie de limite
projective de cet anneau (16.1.11).

On notera que l'entier $n$ de l'énoncé de (16.10.3) est le \emph{rang} de
$\Omega_{X/S}^1$ au point $x$. Nous verrons plus loin (17.13.5) que lorsque
$f$ est différentiellement lisse et localement de type fini, $n$ est aussi
égal à la dimension de la fibre $f^{-1}(f(x))$ au point $x$.

\begin{proposition}[16.10.4]
\label{IV.16.10.4}
Soient $f:X\to S$, $g:S'\to S$ deux morphismes, et posons
$X'=X\times_S S'$, $f'=f_{(S')}:X'\to S'$.
\begin{enumerate}
  \item[(i)] Si $f$ est différentiellement lisse, il en est de même de $f'$.
  \item[(ii)] Inversement, si $g$ est fidèlement plat et quasi-compact, et si
  $f'$ est différentiellement lisse et $\Omega_{X'/S'}^1$ un
  $\mathcal{O}_{X'}$-Module de type fini, $f$ est différentiellement lisse et
  $\Omega_{X/S}^1$ est un $\mathcal{O}_X$-Module de type fini.
\end{enumerate}
\end{proposition}

En effet, si $f$ est différentiellement lisse, les
$\mathscr{G}r_n(\mathscr{P}_{X/S})$ sont des $\mathcal{O}_X$-Modules
\emph{plats}; par suite (16.4.6), l'homomorphisme
\[
  \mathscr{G}r_n(\mathscr{P}_{X/S})\otimes_{\mathcal{O}_X}\mathcal{O}_{X'}
  \longrightarrow\mathscr{G}r_n(\mathscr{P}_{X'/S'})
\]
est bijectif pour tout $n$, et en vertu de la commutativité du diagramme
(16.2.1.3), il résulte de la définition (16.10.1) que $f'$ est
différentiellement lisse. D'autre part, si $g$ est fidèlement plat et
quasi-compact, il résulte encore de (16.4.6) que
\[
  \mathscr{G}r_n(\mathscr{P}_{X/S})\otimes_{\mathcal{O}_X}\mathcal{O}_{X'}
  \longrightarrow\mathscr{G}r_n(\mathscr{P}_{X'/S'})
\]
est bijectif pour tout $n$. Supposons alors $f'$ différentiellement lisse et
$\Omega_{X'/S'}^1$ de rang fini. Comme la projection canonique $X'\to X$ est
un morphisme fidèlement plat et quasi-compact, il résulte d'abord de (2.5.2)
que $\Omega_{X/S}^1$ est un $\mathcal{O}_X$-Module localement libre de rang
fini, puis de (2.2.7) que l'homomorphisme canonique (16.3.1.1) est bijectif,
donc $f$ est différentiellement lisse.

\begin{proposition}[16.10.5]
\label{IV.16.10.5}
Pour qu'un morphisme localement de type fini $f:X\to S$ soit
différentiellement lisse, il faut et il suffit que l'immersion diagonale
$\Delta_f:X\to X\times_S X$ soit quasi-régulière.
\end{proposition}

\oldpage[IV-4]{53}

La question étant locale, on peut se borner au cas où $S$ et $X$ sont
affines, et par suite le sous-préschéma diagonal de $X\times_S X$ est fermé.
L'hypothèse que $f$ est localement de type fini entraîne que $\Delta_f$ est
localement de présentation finie (I, 4.3.1), donc le sous-préschéma diagonal
de $X\times_S X$ est défini par un Idéal de type fini $\mathscr{I}$, et
$\Omega_{X/S}^1=\mathscr{I}/\mathscr{I}^2$ est un $\mathcal{O}_X$-Module de
type fini. La proposition résulte alors aussitôt de la comparaison des
conditions dans (16.10.1) et (16.9.4).

\begin{remark}[16.10.6]
\label{IV.16.10.6}
Soit $f:X\to S$ un morphisme tel que le $\mathcal{O}_X$-Module
$\Omega_{X/S}^1$ soit localement libre de rang fini. Il résulte alors de
(0, 20.4.7) que tout $x\in X$ possède un voisinage ouvert $U$ tel qu'il existe
une famille finie $(z_\lambda)_{\lambda\in L}$ de sections de
$\mathcal{O}_X$ au-dessus de $U$ pour lesquelles $(dz_\lambda)_{\lambda\in L}$
forme une \emph{base} du $\Gamma(U,\mathcal{O}_X)$-module
$\Gamma(U,\Omega_{X/S}^1)$.
\end{remark}

\subsection{\texorpdfstring{Opérateurs différentiels sur un $S$-préschéma différentiellement
lisse}{Opérateurs différentiels sur un S-préschéma différentiellement lisse}}
\EGAReaderTarget{ega4.c04.subsec.16-11}{IV.16.11}

\begin{env}[16.11.1]
\label{IV.16.11.1}
Soient $f:X\to S$ un morphisme, $U$ un ouvert de $X$,
$(z_\lambda)_{\lambda\in L}$ une famille de sections de $\mathcal{O}_X$
au-dessus de $U$ telle que les $dz_\lambda$ forment un système de générateurs
de $\Omega_{X/S}^1|U=\Omega_{U/S}^1$. Soit $m$ un entier ou le symbole
$\infty$, et posons, pour tout $\lambda$
\[
  \zeta_\lambda=\delta z_\lambda=d^m z_\lambda-z_\lambda
  \in\Gamma(U,\mathscr{P}_{X/S}^m).
  \tag{16.11.1.1}
  \label{IV.16.11.1.1}
\]

Nous utiliserons d'autre part les notations usuelles de l'analyse; pour tout
$\mathbf{p}=(p_\lambda)\in\mathbb{N}^{(L)}$ (avec $p_\lambda=0$ sauf pour un
nombre fini d'indices), nous poserons
\[
  |\mathbf{p}|=\sum_\lambda p_\lambda,
  \qquad \mathbf{p}!=\prod_\lambda(p_\lambda!)
  \tag{16.11.1.2}
  \label{IV.16.11.1.2}
\]
\[
  \binom{\mathbf{p}}{\mathbf{q}}
  =\mathbf{p}!/(\mathbf{q}!(\mathbf{p}-\mathbf{q})!)
  \quad\text{pour $\mathbf{p},\mathbf{q}$ dans $\mathbb{N}^{(L)}$,
  $\mathbf{q}\leq\mathbf{p}$}
  \tag{16.11.1.3}
  \label{IV.16.11.1.3}
\]
et on convient que $\binom{\mathbf{p}}{\mathbf{q}}=0$ si
$\mathbf{q}\nleq\mathbf{p}$,
\[
  \mathbf{z}^{\mathbf{p}}=\prod_\lambda(z_\lambda)^{p_\lambda},
  \qquad
  \boldsymbol{\zeta}^{\mathbf{p}}=\prod_\lambda(\zeta_\lambda)^{p_\lambda}.
  \tag{16.11.1.4}
  \label{IV.16.11.1.4}
\]

On aura donc, avec ces notations
\[
  d^m(\mathbf{z}^{\mathbf{p}})
  =(d^m(\mathbf{z}))^{\mathbf{p}}
  =(\boldsymbol{\zeta}+\mathbf{z})^{\mathbf{p}}
  =\sum_{\mathbf{q}\leq\mathbf{p}}
    \binom{\mathbf{p}}{\mathbf{q}}
    \mathbf{z}^{\mathbf{p}-\mathbf{q}}\boldsymbol{\zeta}^{\mathbf{q}}
  \tag{16.11.1.5}
  \label{IV.16.11.1.5}
\]
\[
  \boldsymbol{\zeta}^{\mathbf{p}}
  =(d^m\mathbf{z}-\mathbf{z})^{\mathbf{p}}
  =\sum_{\mathbf{q}\leq\mathbf{p}}
    (-1)^{|\mathbf{p}-\mathbf{q}|}
    \binom{\mathbf{p}}{\mathbf{q}}
    \mathbf{z}^{\mathbf{p}-\mathbf{q}}d^m(\mathbf{z}^{\mathbf{q}}).
  \tag{16.11.1.6}
  \label{IV.16.11.1.6}
\]

Comme les $dz_\lambda$ engendrent $\Omega_{X/S}^1$ et sont les images des
$\delta z_\lambda$, et que l'homomorphisme canonique (16.3.1.1) est
surjectif, on en conclut que pour $m$ fini, les $\delta z_\lambda$ engendrent
la $\mathcal{O}_U$-Algèbre $\mathscr{P}_{U/S}^m$ (Bourbaki,
\emph{Alg. comm.}, chap.~III, §~2, n\textsuperscript{o}~8, cor.~2 du th.~1).
Donc les $\boldsymbol{\zeta}^{\mathbf{p}}$ (pour $|\mathbf{p}|\leq m$)
engendrent le $\mathcal{O}_U$-Module $\mathscr{P}_{U/S}^m$. Un opérateur
différentiel $D\in\operatorname{Diff}_{U/S}^m$ est par suite entièrement
déterminé par les valeurs des
$\langle\boldsymbol{\zeta}^{\mathbf{p}},D\rangle$ pour
$|\mathbf{p}|\leq m$, ou, ce qui revient au même par (16.11.1.5) et
(16.11.1.6), par les valeurs
\oldpage[IV-4]{54}
des $\langle d^m(\mathbf{z}^{\mathbf{p}}),D\rangle=D(\mathbf{z}^{\mathbf{p}})$
pour $|\mathbf{p}|\leq m$; de façon précise, il résulte de (16.11.1.5) que
l'on a
\[
  D(\mathbf{z}^{\mathbf{p}})
  =\langle d^m(\mathbf{z}^{\mathbf{p}}),D\rangle
  =\sum_{\mathbf{q}\leq\mathbf{p}}
    \binom{\mathbf{p}}{\mathbf{q}}
    \langle\boldsymbol{\zeta}^{\mathbf{q}},D\rangle
    \mathbf{z}^{\mathbf{p}-\mathbf{q}}.
  \tag{16.11.1.7}
  \label{IV.16.11.1.7}
\]
\end{env}

\begin{theorem}[16.11.2]
\label{IV.16.11.2}
Soient $f:X\to S$ un morphisme, $U$ un ouvert de $X$,
$(z_\lambda)_{\lambda\in L}$ une famille de sections de $\mathcal{O}_X$
au-dessus de $U$ telle que la famille $(dz_\lambda)_{\lambda\in L}$ engendre
$\Omega_{X/S}^1|U=\Omega_{U/S}^1$. Les conditions suivantes sont
équivalentes:
\begin{enumerate}
  \item[a)] $f|U$ est différentiellement lisse et $(dz_\lambda)$ est une base
  du $\mathcal{O}_U$-Module $\Omega_{U/S}^1$.
  \item[b)] Il existe une famille
  $(D_{\mathbf{p}})_{\mathbf{p}\in\mathbb{N}^{(L)}}$ d'opérateurs
  différentiels de $\mathcal{O}_U$ dans lui-même vérifiant les conditions
  \[
    D_{\mathbf{p}}(\mathbf{z}^{\mathbf{q}})
    =\binom{\mathbf{q}}{\mathbf{p}}
      \mathbf{z}^{\mathbf{q}-\mathbf{p}}
    \quad(\mathbf{p},\mathbf{q}\text{ dans }\mathbb{N}^{(L)}).
    \tag{16.11.2.1}
    \label{IV.16.11.2.1}
  \]
\end{enumerate}
De plus, lorsque ces conditions sont vérifiées, la famille
$(D_{\mathbf{p}})$ est déterminée de façon unique par les conditions
(16.11.2.1) et vérifie les relations
\[
  D_{\mathbf{p}}\circ D_{\mathbf{q}}
  =D_{\mathbf{q}}\circ D_{\mathbf{p}}
  =\frac{(\mathbf{p}+\mathbf{q})!}{\mathbf{p}!\mathbf{q}!}
    D_{\mathbf{p}+\mathbf{q}}
  \quad(\mathbf{p},\mathbf{q}\text{ dans }\mathbb{N}^{(L)}).
  \tag{16.11.2.2}
  \label{IV.16.11.2.2}
\]
Enfin, si $L$ est fini, pour tout entier $m$, les $D_{\mathbf{p}}$ tels que
$|\mathbf{p}|\leq m$ forment une base du $\mathcal{O}_U$-Module
$\operatorname{Diff}_{U/S}^m$, autrement dit, tout opérateur différentiel
d'ordre $\leq m$ sur $U$ s'écrit d'une seule façon sous la forme
\[
  D=\sum_{|\mathbf{p}|\leq m}a_{\mathbf{p}}D_{\mathbf{p}}
\]
où les $a_{\mathbf{p}}$ sont des sections de $\mathcal{O}_X$ au-dessus de
$U$.
\end{theorem}

Notons d'abord qu'en vertu de (16.11.1.6) et (16.11.1.5), on vérifie aussitôt
que les conditions (16.11.2.1) sont équivalentes à
\[
  \langle\boldsymbol{\zeta}^{\mathbf{p}},D_{\mathbf{q}}\rangle
  =\delta_{\mathbf{p}\mathbf{q}}
  \quad\text{(indice de Kronecker)}.
  \tag{16.11.2.3}
  \label{IV.16.11.2.3}
\]
L'existence de la famille $(D_{\mathbf{p}})$ vérifiant ces relations entraîne
donc d'abord (en prenant $|\mathbf{p}|=1$) que les $dz_\lambda$ sont
linéairement indépendants, donc forment une \emph{base} du
$\mathcal{O}_U$-Module $\Omega_{U/S}^1$. Puis, pour tout entier $m\geq1$, on
déduit de même de (16.11.2.3) que les $\boldsymbol{\zeta}^{\mathbf{p}}$ tels
que $|\mathbf{p}|\leq m$ sont linéairement indépendants; par suite
l'homomorphisme canonique (16.3.1.1) est injectif, donc bijectif, et cela
prouve que b) entraîne a). La réciproque découle aussitôt de la définition
(16.10.1), le fait que les $\boldsymbol{\zeta}^{\mathbf{p}}$ forment une base
de $\mathscr{P}_{U/S}^m$ pour $|\mathbf{p}|\leq m$ entraînant l'existence et
l'unicité d'une famille d'homomorphismes
$u_{\mathbf{q},m}:\mathscr{P}_{U/S}^m\to\mathcal{O}_U$
($|\mathbf{q}|\leq m$) tels que
$\langle\boldsymbol{\zeta}^{\mathbf{p}},u_{\mathbf{q},m}\rangle
=\delta_{\mathbf{p}\mathbf{q}}$ pour $|\mathbf{p}|\leq m$,
$|\mathbf{q}|\leq m$. Pour une valeur de $\mathbf{q}$ donnée, les opérateurs
différentiels correspondant aux $u_{\mathbf{q},m}$ pour
$m\geq|\mathbf{q}|$ sont identifiés à un même opérateur $D_{\mathbf{q}}$.
Cela prouve donc que a) entraîne b), et en outre que la famille
$(D_{\mathbf{p}})$ est déterminée de façon unique et que, si $L$ est fini,
pour $|\mathbf{p}|\leq m$, les $D_{\mathbf{p}}$ forment une base du dual
$\operatorname{Diff}_{U/S}^m$ de $\mathscr{P}_{U/S}^m$. Enfin, les relations
(16.11.2.2) découlent aussitôt de l'expression des valeurs des trois
opérateurs considérés pour les $\mathbf{z}^{\mathbf{r}}$, et du fait que les
$\boldsymbol{\zeta}^{\mathbf{r}}$ pour $|\mathbf{r}|\leq m$ engendrent
$\mathscr{P}_{U/S}^m$.

\oldpage[IV-4]{55}

\begin{remarks}[16.11.3]
\EGAReaderTarget{ega4.c04.statement.16-11-3}{IV.16.11.3}
\begin{enumerate}
  \item[(i)] Le fait que les $D_{\mathbf{p}}$ sont deux à deux permutables en
  vertu de (16.11.2.2) n'implique naturellement pas que la
  $\mathcal{O}_U$-Algèbre $\operatorname{Diff}_{U/S}$ soit commutative, les
  $D_{\mathbf{p}}$ ne permutant aux produits par les sections de
  $\mathcal{O}_U$ que si $n=0$.
  \item[(ii)] Les indices $\mathbf{p}$ tels que $|\mathbf{p}|=1$ sont les
  $\boldsymbol{\varepsilon}_\lambda
  =(\varepsilon_{\lambda\mu})_{\mu\in L}$ avec
  $\varepsilon_{\lambda\mu}=0$ si $\mu\neq\lambda$ et
  $\varepsilon_{\lambda\lambda}=1$; lorsque $L$ est fini, les opérateurs
  $D_{\boldsymbol{\varepsilon}_\lambda}$ ne sont autres que les
  $S$-dérivations $D_i$ introduites dans (16.5.7). On notera qu'en général
  (et contrairement à ce qui se passe en Analyse classique), il n'est pas
  vrai qu'un opérateur différentiel d'ordre quelconque puisse s'écrire comme
  combinaison linéaire de puissances des $D_i$ (cf. (16.12)).
  \item[(iii)] Pour tout entier $r\geq1$, on peut définir la notion de
  morphisme \emph{différentiellement lisse jusqu'à l'ordre $r$} en remplaçant
  dans (16.10.1) la condition (ii) par la condition que les homomorphismes
  \[
    \mathbf{S}_{\mathcal{O}_X}^m(\Omega_{X/S}^1)
    \longrightarrow\mathscr{G}r_m(\mathscr{P}_{X/S})
  \]
  soient bijectifs \emph{pour tout $m\leq r$}. Le raisonnement de (16.11.2)
  prouve alors que si, dans la condition a), on remplace «différentiellement
  lisse» par «différentiellement lisse jusqu'à l'ordre $r$», cette condition
  est équivalente à la condition b) dans laquelle on se borne aux
  $\mathbf{p}\in\mathbb{N}^{(L)}$, $\mathbf{q}\in\mathbb{N}^{(L)}$ tels que
  $|\mathbf{p}|\leq r$, $|\mathbf{q}|\leq r$.
\end{enumerate}
\end{remarks}

\subsection{Cas de la caractéristique nulle: critère jacobien pour les
morphismes différentiellement lisses}
\EGAReaderTarget{ega4.c04.subsec.16-12}{IV.16.12}

\begin{env}[16.12.1]
\EGAReaderTarget{ega4.c04.statement.16-12-1}{IV.16.12.1}
On dit qu'un préschéma $X$ est \emph{de caractéristique $p$} ($p$ égal à $0$
ou à un nombre premier) si, pour tout ouvert affine $U$ de $X$, l'anneau
$\Gamma(U,\mathcal{O}_X)$ est de caractéristique $p$ (0, 21.1.1). Il résulte
de (0, 21.1.3) que pour que $X$ soit de caractéristique $0$, il faut et il
suffit que pour tout point \emph{fermé} $x$ de $X$, le corps résiduel $k(x)$
soit de caractéristique $0$, ou encore que $X$ puisse être muni d'une
structure de $\mathbb{Q}$-préschéma (nécessairement unique).
\end{env}

\begin{theorem}[16.12.2]
\label{IV.16.12.2}
Soient $X$ un préschéma de caractéristique $0$, $f:X\to S$ un morphisme. Si
$\Omega_{X/S}^1$ est un $\mathcal{O}_X$-Module localement libre (non
nécessairement de type fini), $f$ est différentiellement lisse.
\end{theorem}

La question étant locale sur $X$, on peut supposer qu'il existe une famille
$(z_\lambda)$ de sections de $\mathcal{O}_X$ au-dessus de $X$ telle que
$(dz_\lambda)$ soit une base du $\mathcal{O}_X$-Module $\Omega_{X/S}^1$.
Appliquant le critère (16.11.2), il suffit de remarquer que les opérateurs
\[
  D_{\mathbf{p}}=(\mathbf{p}!)^{-1}
  \prod_\lambda D_\lambda^{p_\lambda}
\]
(où les $D_\lambda$ sont les formes coordonnées correspondant à la base
$(dz_\lambda)$) vérifient les relations (16.11.2.1), ce qui est une
conséquence du fait que les $D_\lambda$ sont des dérivations.

\begin{env}[16.12.3]
\label{IV.16.12.3}
Le théorème précédent n'est plus exact lorsqu'on abandonne l'hypothèse que
$X$ est de caractéristique $0$. Par exemple, si
$S=\operatorname{Spec}(k)$, où $k$ est un corps de caractéristique $p>0$,
$X=\operatorname{Spec}(K)$ où $K=k(\alpha)$ avec
$\alpha\notin k$, $\alpha^p\in k$, on vérifie aussitôt
\oldpage[IV-4]{56}
que $\Omega_{X/S}^1$ est de rang $1$, et que le morphisme $X\to S$ est
différentiellement lisse jusqu'à l'ordre $p-1$ (16.11.3, (iii)), mais non
jusqu'à l'ordre $p$. Cependant, la démonstration de (16.12.2) prouve que si
$\Omega_{X/S}^1$ est localement libre, et si $n!1_{\mathcal{O}_X}$ est
inversible dans $\Gamma(X,\mathcal{O}_X)$, alors $X$ est différentiellement
lisse sur $S$ jusqu'à l'ordre $n$.
\end{env}

% FIN SOURCE ega4/ega4-16-fr.tex

%
% DEBUT SOURCE ega4/ega4-17-fr.tex
\section[Morphismes lisses, morphismes non ramifiés (ou nets) morphismes étales]{Morphismes lisses, morphismes non ramifiés (ou nets)\\morphismes étales}
\label{section:IV.17}

Dans le présent paragraphe, nous reprenons les notions étudiées dans (0, 19), exprimées à l'aide du langage géométrique des schémas et du point de vue global, pour les préschémas localement de présentation finie sur un préschéma de base donné. La plupart des résultats (à l'exception des n\textsuperscript{os} 17.7, 17.8, 17.9, 17.13 et 17.16) se réduisent d'ailleurs à des variantes de propriétés déjà rencontrées dans (0, 19). Pour des résultats plus spéciaux sur les morphismes étales, le lecteur consultera le § 18.

\subsection{Morphismes formellement lisses, morphismes formellement non ramifiés, morphismes formellement étales}
\EGAReaderTarget{ega4.c04.subsec.17-1}{subsection:IV.17.1}

\begin{definition}[17.1.1]
\EGAReaderTarget{ega4.c04.def.17-1-1}{IV.17.1.1}
Soit $f:X\to Y$ un morphisme de préschémas. On dit que $f$ est \emph{formellement lisse} (resp. \emph{formellement non ramifié} ou \emph{formellement net}, resp. \emph{formellement étale}) si, pour tout schéma affine $Y'$, tout sous-schéma fermé $Y'_0$ de $Y'$ défini par un Idéal nilpotent $\sh{J}$ de $\sh{O}_{Y'}$, et tout morphisme $Y'\to Y$, l'application
\[
\label{IV.17.1.1.1}
  \Hom_Y(Y',X)\longrightarrow\Hom_Y(Y'_0,X)
  \tag{17.1.1.1}
\]
déduite de l'injection canonique $Y'_0\to Y'$, est \emph{surjective} (resp. \emph{injective}, resp. \emph{bijective}).
\end{definition}

On dit encore alors que $X$ est \emph{formellement lisse} (resp. \emph{formellement non ramifié} ou \emph{formellement net}, resp. \emph{formellement étale}) sur $Y$.

Il est clair que dire que $f$ est formellement étale signifie qu'il est \emph{à la fois} formellement lisse et formellement non ramifié.

\begin{remarks}[17.1.2]
\label{IV.17.1.2}
\begin{enumerate}
  \item[(i)] Supposons que $Y=\Spec(A)$ et $X=\Spec(B)$ soient affines, de sorte que $f$ provient d'un homomorphisme d'anneaux $\varphi:A\to B$. En vertu de (0, 19.3.1 et 0, 19.10.1), dire que $f$ est formellement lisse (resp. formellement non ramifié, resp. formellement étale) signifie que $\varphi$ fait de $B$ une $A$-algèbre \emph{formellement lisse} (resp. \emph{formellement non ramifiée}, resp. \emph{formellement étale}) pour les topologies \emph{discrètes} sur $A$ et $B$.
  \item[(ii)] Pour vérifier que $f$ est formellement lisse (resp. formellement non ramifié, resp. formellement étale), on peut, dans la définition (17.1.1) se borner au cas où $\sh{J}^2=0$. En effet, si $f$ vérifie la condition correspondante de la définition (17.1.1) dans ce cas particulier, et si l'on a $\sh{J}^n=0$, on considère le sous-schéma fermé $Y'_j$ de $Y'$ défini par l'Idéal $\sh{J}^{j+1}$ pour $0\leq j\leq n-1$, de sorte que $Y'_j$ est un sous-schéma fermé de $Y'_{j+1}$ défini par un Idéal de carré nul; l'hypothèse entraîne que chacune des applications
  \[
    \Hom_Y(Y'_{j+1},X)\longrightarrow\Hom_Y(Y'_j,X)qquad(0\leq j\leq n-1)
  \]
  \oldpage[IV-4]{57}
  est surjective (resp. injective, resp. bijective); par composition, on en conclut qu'il en est de même de (17.1.1.1).
  \item[(iii)] On notera que les propriétés du morphisme $f$ définies dans (17.1.1) sont des propriétés du \emph{foncteur représentable} $(0_{\mathrm{III}}, 8.1.8)$
  \[
    Y'\longmapsto\Hom_Y(Y',X)
  \]
  de la catégorie des $Y$-préschémas dans la catégorie des ensembles; elles gardent un sens pour \emph{n'importe quel} foncteur contravariant ayant mêmes catégories de départ et d'arrivée, représentable ou non.
  \item[(iv)] Supposons que le morphisme $f$ soit formellement non ramifié (resp. formellement étale); considérons un $Y$-préschéma \emph{quelconque} $Z$ et un sous-préschéma fermé $Z_0$ de $Z$ défini par un Idéal \emph{localement nilpotent} $\sh{J}$ de $\sh{O}_Z$. Alors l'application
  \[
  \label{IV.17.1.2.1}
    \Hom_Y(Z,X)\longrightarrow\Hom_Y(Z_0,X)
    \tag{17.1.2.1}
  \]
  déduite de l'injection canonique $Z_0\to Z$, est encore injective (resp. bijective). En effet, soit $(U_\alpha)$ un recouvrement ouvert affine de $Z$ tel que les Idéaux $\sh{J}|U_\alpha$ soient nilpotents, et pour tout $\alpha$, soit $U^0_\alpha$ l'image réciproque de $U_\alpha$ dans $Z_0$, qui est le sous-schéma fermé de $U_\alpha$ défini par $\sh{J}|U_\alpha$. Soit $f_0:Z_0\to X$ un $Y$-morphisme; par hypothèse, pour tout $\alpha$, il y a au plus un (resp. un et un seul) $Y$-morphisme $f_\alpha:U_\alpha\to X$ dont la restriction à $U^0_\alpha$ soit égale à $f_0|U^0_\alpha$. On en conclut aussitôt que si $f_\alpha$ et $f_\beta$ sont définis, alors, pour tout ouvert affine $V\subset U_\alpha\cap U_\beta$, on a $f_\alpha|V=f_\beta|V$, puisque les restrictions de ces morphismes à l'image réciproque $V_0$ de $V$ dans $Z_0$ coïncident. Il y a donc au plus un (resp. un et un seul) $Y$-morphisme $f:Z\to X$ dont la restriction à $Z_0$ coïncide avec $f_0$.
\end{enumerate}
\end{remarks}

\begin{proposition}[17.1.3]
\label{IV.17.1.3}
\begin{enumerate}
  \item[(i)] Un monomorphisme de préschémas est formellement non ramifié; une immersion ouverte est formellement étale.
  \item[(ii)] Le composé de deux morphismes formellement lisses (resp. formellement non ramifiés, resp. formellement étales) est formellement lisse (resp. formellement non ramifié, resp. formellement étale).
  \item[(iii)] Si $f:X\to Y$ est un $S$-morphisme formellement lisse (resp. formellement non ramifié, resp. formellement étale), il en est de même de $f_{(S')}:X_{(S')}\to Y_{(S')}$ pour toute extension $S'\to S$ du préschéma de base.
  \item[(iv)] Si $f:X\to X'$ et $g:Y\to Y'$ sont deux $S$-morphismes formellement lisses (resp. formellement non ramifiés, resp. formellement étales), il en est de même de $f\times_S g:X\times_S Y\to X'\times_S Y'$.
  \item[(v)] Soient $f:X\to Y$, $g:Y\to Z$ deux morphismes; si $g\circ f$ est formellement non ramifié, il en est de même de $f$.
  \item[(vi)] Si $f:X\to Y$ est un morphisme formellement non ramifié, il en est de même de $f_{\mathrm{red}}:X_{\mathrm{red}}\to Y_{\mathrm{red}}$.
\end{enumerate}
\end{proposition}

En vertu de (I, 5.5.12), il suffit de prouver (i), (ii) et (iii). Les deux assertions de (i) sont triviales. Pour prouver (ii), considérons deux morphismes $f:X\to Y$, $g:Y\to Z$, un schéma affine $Z'$, un sous-schéma fermé $Z'_0$ de $Z'$ défini par un Idéal nilpotent et un morphisme $Z'\to Z$. Supposons $f$ et $g$ formellement lisses, et considérons un
\oldpage[IV-4]{58}
$Z$-morphisme $u_0:Z'_0\to X$; l'hypothèse sur $g$ entraîne qu'il existe un $Z$-morphisme $v:Z'\to Y$ tel que $f\circ u_0=v\circ j$ (où $j:Z'_0\to Z'$ est l'injection canonique); l'hypothèse sur $f$ entraîne alors qu'il existe un morphisme $u:Z'\to X$ tel que $f\circ u=v$ et $u\circ j=u_0$, donc $(g\circ f)\circ u$ est égal au morphisme donné $Z'\to Z$ et $u\circ j=u_0$, ce qui prouve que $g\circ f$ est formellement lisse; on raisonne de même lorsque l'on suppose $f$ et $g$ formellement non ramifiés.

Enfin, pour démontrer (iii), posons $X'=X_{(S')}$, $Y'=Y_{(S')}$, $f'=f_{(S')}$; considérons un schéma affine $Y''$, un sous-schéma fermé $Y''_0$ de $Y''$ défini par un Idéal nilpotent et un morphisme $g:Y''\to Y'$ faisant de $Y''$ un $Y'$-préschéma; on sait alors (I, 3.3.8) que $\Hom_{Y'}(Y'',X')$ s'identifie canoniquement à $\Hom_Y(Y'',X)$ et $\Hom_{Y'}(Y''_0,X')$ à $\Hom_Y(Y''_0,X)$, et la conclusion résulte alors immédiatement de la définition (17.1.1).

On notera qu'une \emph{immersion fermée} n'est pas nécessairement un morphisme formellement lisse.

\begin{proposition}[17.1.4]
\label{IV.17.1.4}
Soient $f:X\to Y$, $g:Y\to Z$ deux morphismes, et supposons $g$ formellement non ramifié. Alors, si $g\circ f$ est formellement lisse (resp. formellement étale), il en est de même de $f$.
\end{proposition}

En effet, soient $Y'$ un schéma affine, $Y'_0$ un sous-schéma fermé de $Y'$ défini par un Idéal nilpotent, $h:Y'\to Y$ un morphisme, $j:Y'_0\to Y'$ l'injection canonique, $u_0:Y'_0\to X$ un $Y$-morphisme, donc tel que $f\circ u_0=h\circ j$. Supposons $g\circ f$ formellement lisse; alors il existe un morphisme $u:Y'\to X$ tel que $u\circ j=u_0$ et $(g\circ f)\circ u=g\circ h$. Mais ces relations entraînent que $f\circ u$ et $h$ sont deux $Z$-morphismes de $Y'$ dans $Y$ tels que $(f\circ u)\circ j=h\circ j$; en vertu de l'hypothèse que $g$ est formellement non ramifié, on en tire que $f\circ u=h$, autrement dit $u$ est un $Y$-morphisme; donc $f$ est formellement lisse. Compte tenu de (17.1.3, (v)), cela démontre la proposition.

\begin{corollary}[17.1.5]
\label{IV.17.1.5}
Supposons $g$ formellement étale; alors, pour que $g\circ f$ soit formellement lisse (resp. formellement non ramifié, resp. formellement étale), il faut et il suffit que $f$ le soit.
\end{corollary}

Cela résulte de (17.1.4) et de (17.1.3, (ii) et (v)).

\begin{proposition}[17.1.6]
\label{IV.17.1.6}
Soit $f:X\to Y$ un morphisme de préschémas.
\begin{enumerate}
  \item[(i)] Soit $(U_\alpha)$ un recouvrement ouvert de $X$ et, pour tout $\alpha$, soit $i_\alpha:U_\alpha\to X$ l'injection canonique. Pour que $f$ soit formellement lisse (resp. formellement non ramifié, resp. formellement étale), il faut et il suffit que chacun des morphismes $f\circ i_\alpha$ le soit.
  \item[(ii)] Soit $(V_\lambda)$ un recouvrement ouvert de $Y$. Pour que $f$ soit formellement lisse (resp. formellement non ramifié, resp. formellement étale), il faut et il suffit que chacune des restrictions $f^{-1}(V_\lambda)\to V_\lambda$ de $f$ le soit.
\end{enumerate}
\end{proposition}

Notons d'abord que (ii) est conséquence de (i): en effet, si $j_\lambda:V_\lambda\to Y$ et $i_\lambda:f^{-1}(V_\lambda)\to X$ sont les injections canoniques, la restriction $f_\lambda:f^{-1}(V_\lambda)\to V_\lambda$ de $f$ est telle que $j_\lambda\circ f_\lambda=f\circ i_\lambda$; si $f$ est formellement lisse (resp. formellement non ramifié), il en est de même de $f\circ i_\lambda$ puisque $i_\lambda$ est formellement étale (17.1.3); mais comme $j_\lambda$ est formellement étale, cela entraîne que $f_\lambda$ est formellement lisse (resp. formellement non ramifié) en vertu de (17.1.5). Inversement, si tous les $f_\lambda$ sont formellement lisses (resp. formellement non ramifiés), il en est de même des $j_\lambda\circ f_\lambda$ (17.1.3), donc aussi de $f$ en vertu de (i).

\oldpage[IV-4]{59}
Si l'on tient compte de ce que les $i_\alpha$ sont formellement étales, tout revient donc à prouver que si les $f\circ i_\alpha$ sont formellement lisses (resp. formellement non ramifiés), il en est de même de $f$.

Soient donc $Y'$ un schéma affine, $Y'_0$ un sous-schéma fermé de $Y'$ défini par un Idéal nilpotent $\sh{J}$, que l'on peut supposer tel que $\sh{J}^2=0$ (17.1.2, (ii)), et enfin soit $g:Y'\to Y$ un morphisme. Supposons donné un $Y$-morphisme $u_0:Y'_0\to X$; désignons par $W_\alpha$ (resp. $W^0_\alpha$) le préschéma induit par $Y'$ (resp. $Y'_0$) sur l'ouvert $u_0^{-1}(U_\alpha)$ (on rappelle que $Y'$ et $Y'_0$ ont \emph{même espace topologique sous-jacent}). Supposons d'abord que les $f\circ i_\alpha$ soient \emph{formellement non ramifiés}, et montrons que, si $u'$ et $u''$ sont deux $Y$-morphismes de $Y'$ dans $X$ dont les restrictions à $Y'_0$ coïncident, alors on a $u'=u''$. En effet, compte tenu de (17.1.2, (iv)), l'hypothèse que les $f\circ i_\alpha$ sont non ramifiés entraîne que pour tout $\alpha$, on a $u'|W_\alpha=u''|W_\alpha$, puisque les restrictions de ces deux $Y$-morphismes à $W^0_\alpha$ coïncident. D'où la conclusion dans ce cas.

Supposons maintenant tous les $f\circ i_\alpha$ \emph{formellement lisses} et prouvons qu'il existe un $Y$-morphisme $u:Y'\to X$ dont $u_0$ est la restriction à $Y'_0$. Or, puisque $Y'$ est un \emph{schéma affine}, on peut appliquer (16.5.17) dont les hypothèses sont satisfaites, et dont la conclusion démontre précisément l'existence de $u$.

On peut donc dire que les notions introduites dans (17.1.1) sont \emph{locales} sur $X$ et sur $Y$, ce qui permet toujours, en vertu de (17.1.2, (i)), de se ramener à l'étude des \emph{algèbres} formellement lisses (resp. formellement non ramifiées, resp. formellement étales).

\subsection{Propriétés différentielles générales}
\EGAReaderTarget{ega4.c04.subsec.17-2}{subsection:IV.17.2}

\begin{proposition}[17.2.1]
\label{IV.17.2.1}
Pour qu'un morphisme $f:X\to Y$ soit formellement non ramifié, il faut et il suffit que $\Omega_f^1=0$ (ce qu'on écrit encore $\Omega_{X/Y}^1=0$ (16.3.1)).
\end{proposition}

Compte tenu de (17.1.6), on est ramené au cas où $Y=\Spec(A)$ et $X=\Spec(B)$ sont affines, et la conclusion résulte alors de (0, 20.7.4) et de l'interprétation de $\Omega_{X/Y}^1$ dans ce cas (16.3.7).

\begin{corollary}[17.2.2]
\label{IV.17.2.2}
Soient $f:X\to Y$, $g:Y\to Z$ deux morphismes. Pour que $f$ soit formellement non ramifié, il faut et il suffit que l'homomorphisme canonique (16.4.19)
\[
  f^*(\Omega_{Y/Z}^1)\longrightarrow\Omega_{X/Z}^1
\]
soit surjectif.
\end{corollary}

C'est une conséquence immédiate de (17.2.1) et de la suite exacte (16.4.19.1).

\begin{proposition}[17.2.3]
\label{IV.17.2.3}
Soit $f:X\to Y$ un morphisme formellement lisse.
\begin{enumerate}
  \item[(i)] Le $\sh{O}_X$-Module $\Omega_{X/Y}^1$ est localement projectif (16.10.1). Si $f$ est localement de type fini, $\Omega_{X/Y}^1$ est localement libre de type fini.
  \item[(ii)] Pour tout morphisme $g:Y\to Z$, la suite (16.4.19) de $\sh{O}_X$-Modules
  \[
  \label{IV.17.2.3.1}
    0\longrightarrow f^*(\Omega_{Y/Z}^1)\longrightarrow\Omega_{X/Z}^1\longrightarrow\Omega_{X/Y}^1\longrightarrow0
    \tag{17.2.3.1}
  \]
  est exacte; en outre, pour tout $x\in X$, il existe un voisinage ouvert $U$ de $x$ tel que les restrictions à $U$ des homomorphismes de (17.2.3.1) forment une suite exacte et scindée.
\end{enumerate}
\end{proposition}

\oldpage[IV-4]{60}
\begin{enumerate}
  \item[(i)] On sait (16.3.9) que si $f$ est localement de type fini, $\Omega_f^1$ est un $\sh{O}_X$-Module de type fini. Pour prouver que, dans tous les cas, il est localement projectif, on peut se borner, en vertu de (17.1.6), au cas où $Y=\Spec(A)$ et $X=\Spec(B)$ sont affines, et cela résulte de l'hypothèse sur $f$ et de (0, 20.4.9 et 0, 19.2.1).
  \item[(ii)] Ici encore, on peut se borner au cas où $X$, $Y$ et $Z$ sont affines (17.1.6) et la conclusion résulte dans ce cas de l'interprétation des Modules figurant dans la suite (17.2.3.1) et de (0, 20.5.7).
\end{enumerate}

\begin{corollary}[17.2.4]
\label{IV.17.2.4}
Si $f:X\to Y$ est un morphisme formellement étale, alors, pour tout morphisme $g:Y\to Z$, l'homomorphisme canonique de $\sh{O}_X$-Modules
\[
  f^*(\Omega_{Y/Z}^1)\longrightarrow\Omega_{X/Z}^1
\]
est bijectif.
\end{corollary}

Cela résulte de l'exactitude de la suite (17.2.3.1) et du fait que l'on a alors $\Omega_{X/Y}^1=0$ (17.2.1).

\begin{proposition}[17.2.5]
\label{IV.17.2.5}
Soient $f:X\to Y$ un morphisme, $X'$ un sous-préschéma de $X$ tel que le morphisme composé $X'\xrightarrow{j}X\to Y$ (où $j$ est l'injection canonique) soit formellement lisse. Alors la suite de $\sh{O}_{X'}$-Modules (16.4.21)
\[
\label{IV.17.2.5.1}
  0\longrightarrow\mathcal{N}_{X'/X}\longrightarrow\Omega_{X/Y}^1\otimes_{\sh{O}_X}\sh{O}_{X'}\longrightarrow\Omega_{X'/Y}^1\longrightarrow0
  \tag{17.2.5.1}
\]
est exacte; en outre, pour tout $x\in X$, il existe un voisinage ouvert $U$ de $x$ tel que les restrictions à $U$ des homomorphismes de (17.2.5.1) forment une suite exacte et scindée.
\end{proposition}

Toujours en vertu de (17.1.6), on peut se borner au cas où $Y=\Spec(A)$ et $X=\Spec(B)$ sont affines, et $X'=\Spec(B/\mathfrak{J})$, où $\mathfrak{J}$ est un idéal de $B$. Alors le faisceau conormal $\mathcal{N}_{X'/X}$ correspond au $B$-module $\mathfrak{J}/\mathfrak{J}^2$ (16.1.3), et la conclusion découle de (0, 20.5.14).

\begin{proposition}[17.2.6]
\label{IV.17.2.6}
Soient $X$, $Y$ deux préschémas, $f:X\to Y$ un morphisme localement de type fini. Les conditions suivantes sont équivalentes:
\begin{enumerate}
  \item[a)] $f$ est un monomorphisme.
  \item[b)] $f$ est radiciel et formellement non ramifié.
  \item[c)] Pour tout $y\in Y$, la fibre $f^{-1}(y)$ est vide ou $k(y)$-isomorphe à $\Spec(k(y))$ (autrement dit, est réduite à un seul point $z$ tel que $k(y)\to\sh{O}_z/\mathfrak{m}_y\sh{O}_z$ soit un isomorphisme).
\end{enumerate}
\end{proposition}

Le fait que a) entraîne c) résulte de (8.11.5.1). Il est clair que c) entraîne que $f$ est radiciel; montrons qu'il résulte aussi de c) que $\Omega_{X/Y}^1=0$, ce qui prouvera que c) entraîne b) (17.2.1). Notons que le $\sh{O}_X$-Module $\Omega_{X/Y}^1$ est quasi-cohérent de type fini (16.3.9). Il résulte donc de (I, 9.1.13.1) que, pour que $(\Omega_{X/Y}^1)_x=0$, il faut et il suffit que si l'on pose $Y_1=\Spec(k(y))$, $X_1=f^{-1}(y)=X\times_Y Y_1$, on ait $(\Omega_{X_1/Y_1}^1)_x=0$; mais comme le morphisme $f_1:X_1\to Y_1$ déduit de $f$ est formellement non ramifié en vertu de l'hypothèse c) (17.1.3), la conclusion résulte de (17.2.1). Prouvons enfin que b) entraîne a); pour cela, considérons le morphisme diagonal $g=\Delta_f:X\to X\times_Y X$; puisque $f$ est radiciel, $g$ est surjectif (I, 8.7.1); d'autre part, $\Omega_{X/Y}^1$ est par définition le faisceau conormal $\mathop{Gr}\nolimits_1(g)$ de l'immersion $g$ (16.3.1), et dire que $f$ est formellement non ramifié signifie donc que

\oldpage[IV-4]{61}
$\mathop{Gr}\nolimits_1(g)=0$ (17.2.1). En outre, $g$ est localement de présentation finie (I, 4.3.1); donc l'hypothèse $\mathop{Gr}\nolimits_1(g)=0$ entraîne que $g$ est une immersion ouverte (16.1.10); étant surjective, cette immersion est un isomorphisme, donc $f$ est un monomorphisme (I, 5.3.8).

\subsection{Morphismes lisses, morphismes non ramifiés, morphismes étales}
\EGAReaderTarget{ega4.c04.subsec.17-3}{subsection:IV.17.3}

\begin{definition}[17.3.1]
\EGAReaderTarget{ega4.c04.def.17-3-1}{IV.17.3.1}
On dit qu'un morphisme $f:X\to Y$ est \emph{lisse} (resp. \emph{non ramifié}, ou \emph{net}\footnote{Les mots « net » et « formellement net » paraissent bien préférables à la terminologie consacrée de « non ramifié » (resp. « formellement non ramifié ») et seront employés à peu près exclusivement à partir du chap. V. Dans ce chapitre, nous avons conservé l'ancienne terminologie afin de ne pas entrer en conflit avec (0, 19.10).}, resp. \emph{étale}) s'il est localement de présentation finie et formellement lisse (resp. formellement non ramifié, resp. formellement étale).
\end{definition}

On dit encore alors que $X$ est \emph{lisse} (resp. \emph{non ramifié} ou \emph{net}, resp. \emph{étale}) sur $Y$.

Nous verrons plus loin (17.5.2) que cette définition d'un morphisme lisse coïncide avec celle déjà donnée dans (6.8.1); jusque-là, c'est la définition de (17.3.1) que nous utiliserons exclusivement.

Il est clair que dire que $f$ est étale signifie qu'il est \emph{à la fois} lisse et non ramifié.

\begin{remarks}[17.3.2]
\EGAReaderTarget{ega4.c04.subsec.17-3-2}{IV.17.3.2}
\begin{enumerate}
  \item[(i)] On notera que la définition (17.3.1) peut s'exprimer uniquement à l'aide du foncteur
  \[
    Y'\longmapsto\Hom_Y(Y',X)
  \]
  considéré dans (17.1.2, (iii)), car dire que $f$ est localement de présentation finie équivaut à dire que le foncteur précédent \emph{commute aux limites projectives de schémas affines} (8.14.2).
  \item[(ii)] Soient $A$ un anneau, $B$ une $A$-algèbre. On dit que $B$ est une $A$-algèbre \emph{lisse} (resp. \emph{non ramifiée}, resp. \emph{étale}) si le morphisme correspondant $\Spec(B)\to\Spec(A)$ est lisse (resp. non ramifié, resp. étale). Il revient au même de dire que $B$ est une $A$-algèbre \emph{de présentation finie} (I, 4.4.6) et formellement lisse (resp. formellement non ramifiée, resp. formellement étale) pour les topologies discrètes.
  \item[(iii)] Il résulte de (17.1.6) et de la définition d'un morphisme localement de présentation finie (I, 4.4.2) que les notions de morphisme lisse, non ramifié et étale sont \emph{locales sur $X$ et sur $Y$}.
\end{enumerate}
\end{remarks}

\begin{proposition}[17.3.3]
\label{IV.17.3.3}
\begin{enumerate}
  \item[(i)] Une immersion ouverte est étale. Pour qu'une immersion soit non ramifiée, il faut et il suffit qu'elle soit localement de présentation finie.
  \item[(ii)] Le composé de deux morphismes lisses (resp. non ramifiés, resp. étales) est lisse (resp. non ramifié, resp. étale).
  \item[(iii)] Si $f:X\to Y$ est un $S$-morphisme lisse (resp. non ramifié, resp. étale), il en est de même de $f_{(S')}:X_{(S')}\to Y_{(S')}$ pour toute extension $S'\to S$ du préschéma de base.
  \item[(iv)] Si $f:X\to X'$ et $g:Y\to Y'$ sont deux $S$-morphismes lisses (resp. non ramifiés, resp. étales), il en est de même de $f\times_S g:X\times_S Y\to X'\times_S Y'$.
\end{enumerate}

\oldpage[IV-4]{62}

\begin{enumerate}
  \item[(v)] Soient $f:X\to Y$, $g:Y\to Z$ deux morphismes; si $g$ est localement de type fini et si $g\circ f$ est non ramifié, alors $f$ est non ramifié.
\end{enumerate}
\end{proposition}

Cela résulte aussitôt de (I, 4.3) et de (17.1.3).

\begin{proposition}[17.3.4]
\label{IV.17.3.4}
Soient $f:X\to Y$, $g:Y\to Z$ deux morphismes, et supposons $g$ non ramifié. Alors, si $g\circ f$ est lisse (resp. non ramifié, resp. étale), il en est de même de $f$.
\end{proposition}

En effet, comme $g$ et $g\circ f$ sont localement de présentation finie, il en est de même de $f$ (I, 4.4.3, (v)); la conclusion résulte donc de (17.1.4) et (17.1.3, (v)).

\begin{corollary}[17.3.5]
\label{IV.17.3.5}
Supposons $g$ étale; alors, pour que $f$ soit lisse (resp. non ramifié, resp. étale), il faut et il suffit que $g\circ f$ le soit.
\end{corollary}

Cela résulte de (17.3.4) et de (17.3.3, (iii)).

\begin{proposition}[17.3.6]
\label{IV.17.3.6}
Soient $g:Y\to S$, $h:X\to S$ deux morphismes localement de présentation finie. Pour qu'un $S$-morphisme $f:X\to Y$ soit non ramifié, il faut et il suffit que l'homomorphisme canonique (16.4.19)
\[
  f^*(\Omega_{Y/S}^1)\longrightarrow\Omega_{X/S}^1
\]
soit surjectif.
\end{proposition}

Comme $f$ est alors localement de présentation finie (I, 4.4.3, (v)), la proposition résulte aussitôt de (17.2.2).

\begin{definition}[17.3.7]
\EGAReaderTarget{ega4.c04.def.17-3-7}{IV.17.3.7}
Soit $f:X\to Y$ un morphisme. On dit que $f$ est \emph{lisse} (resp. \emph{non ramifié}, resp. \emph{étale}) \emph{en un point} $x\in X$ s'il existe un voisinage ouvert $U$ de $x$ dans $X$ tel que la restriction $f|U$ soit un morphisme lisse (resp. non ramifié, resp. étale) de $U$ dans $Y$.
\end{definition}

On dit encore alors que $X$ est lisse (resp. non ramifié, resp. étale) sur $Y$ au point $x$.

Compte tenu de la remarque (17.3.2, (iii)), il revient au même de dire que $f$ est un morphisme lisse (resp. non ramifié, resp. étale) ou qu'il est lisse (resp. non ramifié, resp. étale) en chaque point de $X$.

Il est clair que l'ensemble des points de $X$ où un morphisme $f:X\to Y$ est lisse (resp. non ramifié, resp. étale) est \emph{ouvert} dans $X$.

\begin{proposition}[17.3.8]
\label{IV.17.3.8}
Pour tout préschéma $Y$ et tout $\sh{O}_Y$-Module localement libre $\sh{E}$ de type fini, le préschéma fibré vectoriel $\mathbf{V}(\sh{E})$ (II, 1.7.8) est un $Y$-préschéma lisse.
\end{proposition}

En effet (17.3.2, (iii)), on peut se borner au cas où $Y=\Spec(A)$ est affine et $\mathbf{V}(\sh{E})=\Spec(A[T_1,\ldots,T_r])$; comme $A[T_1,\ldots,T_r]$ est une $A$-algèbre formellement lisse pour les topologies discrètes (0, 19.3.2) et de présentation finie, cela démontre la proposition (17.3.2, (ii)).

\begin{corollary}[17.3.9]
\label{IV.17.3.9}
Sous les hypothèses de (17.3.8), le préschéma fibré projectif $\mathbf{P}(\sh{E})$ (II, 4.1.1) est un $Y$-préschéma lisse.
\end{corollary}

On peut encore se borner au cas où $Y=\Spec(A)$ est affine et $\mathbf{P}(\sh{E})=\mathbf{P}^r_Y$. On sait alors (II, 2.3.14) qu'on a un recouvrement ouvert fini de $\mathbf{P}^r_A$ en prenant les $D_+(T_i)$ ($0\leq i\leq r$), égaux respectivement aux spectres des anneaux $S_{(f)}$, où l'on remplace $S$ par $A[T_0,T_1,\ldots,T_r]$ et $f$ par $T_i$; mais il résulte aussitôt de la définition de $S_{(f)}$ (II, 2.2.1) que cet anneau, dans le cas considéré, est isomorphe à $A[T_0,\ldots,T_{i-1},T_{i+1},\ldots,T_r]$; donc le corollaire résulte de (17.3.8).

\oldpage[IV-4]{63}

\subsection{Caractérisations des morphismes non ramifiés}
\EGAReaderTarget{ega4.c04.subsec.17-4}{subsection:IV.17.4}

\begin{theorem}[17.4.1]
\label{IV.17.4.1}
Soient $f:X\to Y$ un morphisme localement de présentation finie, $x$ un point de $X$. Les propriétés suivantes sont équivalentes:
\begin{enumerate}
  \item[a)] $f$ est non ramifié au point $x$.
  \item[b)] Le morphisme diagonal $\Delta_f:X\to X\times_Y X$ est un isomorphisme local au point $x$.
  \item[b')] Si l'on pose $\Delta_f=(\psi,\theta)$, $Z=X\times_Y X$ et $z=\psi(x)$, l'homomorphisme $\theta_z^\#:\sh{O}_{Z,z}\to\sh{O}_{X,x}$ est bijectif.
  \item[b'')] Pour tout morphisme $g:Y'\to Y$, et tout point $y'\in Y'$ au-dessus de $y=f(x)$, toute $Y'$-section $s'$ de $X'=X\times_Y Y'$ telle que $x'=s'(y')$ soit au-dessus de $x$ est un isomorphisme local au point $y'$.
  \item[c)] On a $(\Omega^1_{X/Y})_x=0$.
  \item[d)] Le $k(y)$-préschéma $f^{-1}(y)$ est non ramifié sur $k(y)$ au point $x$.
  \item[d')] Le point $x$ est isolé dans $X_y=f^{-1}(y)$ (autrement dit (Err$_{\mathrm{III}}$, 20), le morphisme $f$ est quasi-fini au point $x$) et l'anneau $\sh{O}_{X_y,x}$ est un corps, extension séparable de $k(y)$.
  \item[d'')] L'anneau $\sh{O}_{X_y,x}=\sh{O}_{X,x}/\mathfrak{m}_y\sh{O}_{X,x}$ est un corps, extension finie séparable de $k(y)$.
  \item[e)] L'anneau $\sh{O}_{X,x}$ est une $\sh{O}_{Y,y}$-algèbre formellement non ramifiée pour les topologies discrètes.
\end{enumerate}
\end{theorem}

Comme $f$ est localement de type fini, le $\sh{O}_X$-Module $\Omega^1_{X/Y}$ est de type fini (16.3.9), donc il revient au même de dire que $(\Omega^1_{X/Y})_x=0$ ou qu'il existe un voisinage ouvert $U$ de $x$ tel que $\Omega^1_{X/Y}|U=0$. Compte tenu de (17.2.1), cela prouve l'équivalence de a) et c). D'autre part, si l'on pose $A=\sh{O}_{Y,y}$, $B=\sh{O}_{X,x}$, on a $(\Omega^1_{X/Y})_x=\Omega^1_{B/A}$ (16.4.15), et l'équivalence de c) et e) résulte donc de (0, 20.7.4).

Comme d') ne fait intervenir que des propriétés du morphisme $X_y\to\Spec(k(y))$, l'équivalence de a) et d') entraînera \emph{ipso facto} celle de d) et d'). D'autre part d') et d'') sont équivalentes, car il revient au même de dire que $\sh{O}_{X_y,x}$ est une $k(y)$-algèbre \emph{finie} ou que $x$ est un point isolé de $X_y$, puisque $X_y$ est un $k(y)$-préschéma localement de type fini (I, 6.4.4).

Prouvons maintenant l'équivalence de b) et b'). On peut se limiter au cas où $Y=\Spec(R)$ et $X=\Spec(S)$ sont affines et $f$ de présentation finie; alors on a $Z=\Spec(S\otimes_R S)$ et $\Delta_f$ correspond à l'homomorphisme canonique surjectif $S\otimes_R S\to S$, dont on sait que le noyau $\mathfrak{J}$ est un idéal de type fini (0, 20.4.4). Si l'on pose $\sh{J}=\widetilde{\mathfrak{J}}$, le $\sh{O}_Z$-Module $\psi_*(\sh{O}_X)=\sh{O}_Z/\sh{J}$ est donc de présentation finie, et l'hypothèse que l'homomorphisme $\theta_z^\#: \sh{O}_{Z,z}\to(\psi_*(\sh{O}_X))_z$ est bijectif entraîne qu'en remplaçant au besoin $X$ par un voisinage ouvert de $x$, l'homomorphisme $\theta:\sh{O}_Z\to\psi_*(\sh{O}_X)$ est lui-même bijectif ($0_{\mathrm I}$, 5.2.7). Cela montre donc que b') entraîne b); la réciproque est évidente.

D'autre part, l'équivalence de b) et b'') résulte de (I, 5.3.7) sans hypothèse de finitude sur $f$: la donnée d'une $Y'$-section $s':Y'\to X'$ équivaut en effet à celle d'un

\oldpage[IV-4]{64}
$Y$-morphisme $h=g'\circ s':Y'\to X$ (où $g':X'\to X$ est la projection canonique), de sorte que $s'=(1_{Y'},h)_X$, et alors le diagramme
\[
\tag{17.4.1.1}
\label{IV.17.4.1.1}
\xymatrix{
Y' \ar[r]^{s'} \ar[d]_h & X'=Y'\times_Y X \ar[d]^{h\times 1_X}\\
X \ar[r]_{\Delta_f} & X\times_Y X
}
\]
identifie $Y'$ au produit des $(X\times_Y X)$-préschémas $X$ et $X'$. Par suite (I, 4.3.2) si $\Delta_f$ est un isomorphisme local au point $x$, $s'$ est un isomorphisme local au point $y'$ (puisque $x=h(y')$), ce qui prouve que b) implique b''). La réciproque s'obtient en appliquant b'') au cas où l'on prend $Y'=X$, $y'=x$, $g=f$ et $s'=\Delta_f$.

Pour achever la démonstration de (17.4.1), il suffit de prouver les implications
\[
d'')\Rightarrow c)\Rightarrow b)\Rightarrow d'').
\]

$d'')\Rightarrow c)$: Comme $\Omega^1_{X/Y}$ est un $\sh{O}_X$-Module de type fini, il résulte du lemme de Nakayama que la condition c) équivaut à $(\Omega^1_{X/Y})_x/\mathfrak{m}_y(\Omega^1_{X/Y})_x=0$, c'est-à-dire (16.4.5) à $(\Omega^1_{X_y/\Spec(k(y))})_x=0$. On est donc ramené au cas où $Y$ est le spectre d'un corps $k$ et $X$ un $k$-préschéma algébrique. L'hypothèse que $\sh{O}_{X,x}$ est un corps $k'$, extension finie de $k$, entraîne d'abord que $x$ est \emph{fermé} dans $X$ (I, 6.4.2), puis que $x$ est point \emph{maximal} du préschéma noethérien $X$, donc est un point isolé de $X$. Remplaçant $X$ par l'ouvert $\{x\}$ de $X$, on peut donc supposer que $X=\Spec(k')$; mais alors l'hypothèse que $k'$ est extension \emph{finie séparable} de $k$ entraîne $\Omega^1_{k'/k}=0$ (0, 20.6.20), ce qui prouve c).

$c)\Rightarrow b)$: On a vu plus haut que l'on a alors $\Omega^1_{X/Y}|U=0$ pour un voisinage ouvert $U$ de $x$ dans $X$; l'assertion b) résulte alors de la définition de $\Omega^1_{X/Y}$ (16.3.1) et de (16.1.9).

$b)\Rightarrow d'')$: Remplaçant $X$ par un voisinage ouvert de $x$, on peut supposer que $\Delta_f$ est une immersion ouverte; si l'on désigne par $f_y:X_y\to\Spec(k(y))$ le morphisme déduit de $f$ par changement de base, $\Delta_{f_y}$ est alors aussi une immersion ouverte (I, 5.3.4), et comme la condition d'') ne concerne que le préschéma $X_y$, on voit qu'on peut se borner au cas où $Y$ est le spectre d'un corps $k$, $X$ le spectre d'une $k$-algèbre $A$ de type fini; la propriété d'') sera établie si l'on prouve que $A$ est une $k$-algèbre \emph{finie et séparable}, une telle algèbre étant composée directe d'extensions finies séparables de $k$. Si $K$ est une extension algébriquement close de $k$, il revient au même de dire que $A\otimes_k K$ est une $K$-algèbre finie et séparable (4.6.1), donc on voit qu'on peut se borner au cas où $k$ est algébriquement clos. Montrons d'abord que $A$ est une $k$-algèbre \emph{finie}: il suffira de montrer que tout point \emph{fermé} $x$ de $X$ est isolé, car alors l'ensemble de ces points est ouvert dans $X$ et discret, donc fini puisqu'il est quasi-compact ($X$ étant noethérien), ce qui établira notre assertion en vertu de (I, 6.4.4). Or, on a alors $k(x)=k$ puisque $k$ est algébriquement clos (I, 6.4.3), donc il y a une $Y$-section $s$ de $X$ telle que $s(Y)=\{x\}$, et en vertu de (17.4.1.1), $\{x\}$ est l'image réciproque de la diagonale $\Delta_X(X)$ par un morphisme $X\to X\times_Y X$, donc $\{x\}$ est ouvert dans $X$ en vertu de l'hypothèse b). On a ainsi montré que $A$ est une

\oldpage[IV-4]{65}
$k$-algèbre \emph{finie}, composée directe de $k$-algèbres finies locales. Pour exprimer que $\Delta_f$ est une immersion ouverte, on peut donc se borner au cas où $A$ est une $k$-algèbre locale finie, $X=\Spec(A)$ étant donc réduit à un seul point; le corps résiduel de $A$, étant une extension finie de $k$, est nécessairement identique à $k$, et par suite (I, 3.4.9) $X\times_k X$ est réduit à un seul point et $\Delta_f$ est donc nécessairement un \emph{isomorphisme}. Or, puisque $A$ est une $k$-algèbre, l'homomorphisme canonique $A\otimes_k A\to A$ ne peut être bijectif que si $A=k$. C.Q.F.D.

\begin{remark}[17.4.1.2]
\label{IV.17.4.1.2}
Supposons seulement que $f$ soit \emph{localement de type fini}. Alors $\Omega^1_{X/Y}$ est encore un $\sh{O}_X$-Module de type fini (16.3.9), et $\Delta_f$ un morphisme localement de présentation finie (I, 4.3.1). Toute la démonstration de (17.4.1) est donc valable, à condition de remplacer a) par: \emph{la restriction de $f$ à un voisinage convenable de $x$ est un morphisme formellement non ramifié}. On voit en outre que dans ce cas la restriction de $f$ à un voisinage convenable de $x$ est un morphisme \emph{localement quasi-fini}.
\end{remark}

\begin{corollary}[17.4.2]
\label{IV.17.4.2}
Soit $f:X\to Y$ un morphisme localement de présentation finie. Les propriétés suivantes sont équivalentes:
\begin{enumerate}
  \item[a)] $f$ est non ramifié.
  \item[b)] Le morphisme diagonal $\Delta_f:X\to X\times_Y X$ est une immersion ouverte.
  \item[b')] Pour tout morphisme $Y'\to Y$, toute $Y'$-section de $X'=X\times_Y Y'$ est une immersion ouverte.
  \item[c)] On a $\Omega^1_{X/Y}=0$.
  \item[d)] Pour tout $y\in Y$, le $k(y)$-préschéma $f^{-1}(y)$ est non ramifié sur $k(y)$.
  \item[d')] Pour tout $y\in Y$, le $k(y)$-préschéma $f^{-1}(y)$ est isomorphe à un préschéma de la forme $\coprod_{\lambda\in L}\Spec(K_\lambda)$, où, pour tout $\lambda\in L$, $K_\lambda$ est une extension finie et séparable de $k(y)$.
  \item[e)] Pour tout $x\in X$, l'anneau $\sh{O}_{X,x}$ est une $\sh{O}_{Y,f(x)}$-algèbre formellement non ramifiée pour les topologies discrètes.
\end{enumerate}
\end{corollary}

\begin{corollary}[17.4.3]
\label{IV.17.4.3}
Si $f:X\to Y$ est non ramifié, alors $f$ est localement quasi-fini (Err$_{\mathrm{III}}$, 20).
\end{corollary}

\begin{proposition}[17.4.4]
\label{IV.17.4.4}
Soient $Y$ un préschéma localement noethérien, $f:X\to Y$ un morphisme localement de type fini, $x$ un point de $X$, $y=f(x)$. Posons $A=\sh{O}_{Y,y}$, $B=\sh{O}_{X,x}$, qui sont des anneaux locaux noethériens, et soit $k$ le corps résiduel de $A$. Alors les conditions équivalentes a) à e) du théorème (17.4.1) sont aussi équivalentes à chacune des suivantes:
\begin{enumerate}
  \item[f)] $\widehat{B}\otimes_{\widehat{A}}k$ est un corps, extension finie séparable de $k$ (ce qui entraîne que $\widehat{B}$ est une $\widehat{A}$-algèbre finie).
  \item[f')] $\widehat{B}$ est une $\widehat{A}$-algèbre formellement non ramifiée pour les topologies adiques.
\end{enumerate}
Si de plus $k(x)=k(y)$, ou si $k$ est séparablement clos, ces conditions équivalent aussi à:
\begin{enumerate}
  \item[f'')] L'homomorphisme $\widehat{A}\to\widehat{B}$ est surjectif.
\end{enumerate}
\end{proposition}

Notons d'abord, par le même raisonnement que dans (0, 19.3.6), qu'il revient au même de dire que $B$ est une $A$-algèbre formellement non ramifiée pour les topologies préadiques, ou que $\widehat{B}$ est une $\widehat{A}$-algèbre formellement non ramifiée pour les topologies adiques. D'autre part, l'hypothèse que $f$ est localement de type fini entraîne que $\Omega^1_{B/A}$

\oldpage[IV-4]{66}
est un $B$-module de type fini (16.3.9), donc \emph{séparé} pour la topologie $\mathfrak n$-préadique (où $\mathfrak n$ est l'idéal maximal de $B$) ($0_{\mathrm I}$, 7.3.5); il revient par suite au même de dire que $\Omega^1_{B/A}=0$ ou que $\widehat{\Omega}^1_{B/A}=0$; donc (0, 20.7.4), il revient au même de dire que $B$ est une $A$-Algèbre formellement non ramifiée pour les topologies \emph{discrètes}, ou pour les topologies \emph{préadiques}. Ceci prouve l'équivalence des conditions e) et f'). Si $\mathfrak m$ est l'idéal maximal de $A$, on a
\[
k=A/\mathfrak m=\widehat A/\mathfrak m\widehat A,
\qquad
\widehat B\otimes_{\widehat A}k=\widehat B/\mathfrak m\widehat B
=\widehat B\otimes_B(B/\mathfrak m B),
\]
et par suite ($0_{\mathrm I}$, 7.3.5) $\widehat B/\mathfrak m\widehat B$ est le complété de $B/\mathfrak m B=B\otimes_A k$ pour la topologie $\mathfrak n$-préadique; ceci prouve l'équivalence de d'') et de f). Enfin, lorsque $k(x)=k(y)$ ou lorsque $k$ est séparablement clos, la condition f) entraîne que l'homomorphisme $\widehat A/\mathfrak m\widehat A\to\widehat B/\mathfrak m\widehat B$ est bijectif; la condition f) implique d'autre part que $\widehat B$ est une $\widehat A$-algèbre \emph{quasi-finie} ($0_{\mathrm I}$, 7.4.4), donc \emph{finie} puisque $\widehat A$ est complet et $\widehat B$ séparé pour la topologie $\mathfrak m$-préadique, $\mathfrak m\widehat B$ étant un idéal de définition de $\widehat B$ ($0_{\mathrm I}$, 7.4.1). L'homomorphisme $\widehat A\to\widehat B$ est donc surjectif en vertu du lemme de Nakayama. Donc f) entraîne alors f''), et la réciproque est évidente.

\begin{env}[17.4.5]
\EGAReaderTarget{ega4.c04.statement.17-4-5}{IV.17.4.5}
Étant donnés un $S$-préschéma $Y$ et deux $S$-morphismes $f:X\to Y$, $g:X\to Y$, on en déduit canoniquement un $S$-morphisme $(f,g)_S:X\to Y\times_S Y$. Nous appellerons \emph{préschéma des coïncidences} de $f$ et $g$ l'image réciproque par $(f,g)_S$ de la diagonale $\Delta_{Y|S}$; c'est donc un sous-préschéma de $X$, qui est \emph{fermé} lorsque $Y$ est un $S$-schéma (I, 5.4.1).
\end{env}

\begin{proposition}[17.4.6]
\label{IV.17.4.6}
Soient $h:Y\to S$ un morphisme non ramifiée, et soient $f:X\to Y$, $g:X\to Y$ deux $S$-morphismes. Alors le préschéma $C$ des coïncidences de $f$ et $g$ est un sous-préschéma induit sur un ouvert de $X$; si de plus $Y$ est un $S$-schéma (I, 5.4.1), $C$ est un sous-préschéma fermé de $X$.
\end{proposition}

En effet, puisque $\Delta_h:Y\to Y\times_S Y$ est une immersion ouverte (17.4.2), l'image réciproque par $(f,g)_S$ de $\Delta_{Y|S}$ est un sous-préschéma induit sur un ouvert de $X$ (I, 4.4.1). La dernière assertion résulte de (17.4.5).

\begin{corollary}[17.4.7]
\label{IV.17.4.7}
Sous les hypothèses de (17.4.6), soit $x$ un point de $X$ tel que les deux morphismes composés $\Spec(k(x))\to X\xrightarrow{f}Y$ et $\Spec(k(x))\to X\xrightarrow{g}Y$ soient égaux. Alors il existe un voisinage ouvert $U$ de $x$ tel que $f|U=g|U$. Si de plus $Y$ est un $S$-schéma, il existe un voisinage ouvert et fermé $X'$ de $x$ dans $X$ tel que $f|X'=g|X'$. Si enfin on suppose en outre que $X$ est connexe, on a $f=g$.
\end{corollary}

Cela résulte de (17.4.6) et de (I, 5.3.17).

\begin{corollary}[17.4.8]
\label{IV.17.4.8}
Sous les hypothèses de (17.4.6), supposons que le morphisme structural $\varphi=h\circ f=h\circ g$ de $X$ dans $S$ soit fermé. Soit $s$ un point de $S$; soit $X_s$ le $k(s)$-préschéma $\varphi^{-1}(s)$ et supposons que les deux morphismes composés $X_s\to X\xrightarrow{f}Y$ et $X_s\to X\xrightarrow{g}Y$ soient égaux. Alors il existe un voisinage ouvert $V$ de $s$ dans $S$ tel que $f|\varphi^{-1}(V)=g|\varphi^{-1}(V)$. Si de plus $Y$ est un $S$-schéma et si $\varphi$ est ouvert, on peut prendre $V$ ouvert et fermé. Si enfin on suppose de plus $S$ connexe, on a $f=g$.
\end{corollary}

Il résulte de (17.4.7) que le préschéma $C$ des coïncidences de $f$ et $g$ est induit sur un ouvert de $X$ et contient $X_s$. Comme $\varphi$ est fermé, il existe un voisinage ouvert $V$ de $s$ tel que $\varphi^{-1}(V)\subset C$. Si de plus $Y$ est un $S$-schéma, $C$ est fermé, donc $\varphi(X-C)$ est à la

\oldpage[IV-4]{67}
fois ouvert et fermé dans $S$, et son complémentaire $V$ dans $S$ est donc un voisinage ouvert et fermé de $s$ tel que $\varphi^{-1}(V)\subset C$.

\begin{proposition}[17.4.9]
\label{IV.17.4.9}
Soient $Y$ un préschéma connexe, $f:X\to Y$ un morphisme non ramifié et séparé. Alors toute $Y$-section $g$ de $X$ est un isomorphisme de $Y$ sur une composante connexe ouverte de $X$, et l'application $g\mapsto g(Y)$ est une bijection de $\Gamma(X/Y)$ sur l'ensemble des composantes connexes $Z$ de $X$ (nécessairement ouvertes dans $X$) telles que la restriction de $f$ à $Z$ soit un isomorphisme de $Z$ sur $Y$. En particulier, si $g'$ et $g''$ sont deux $Y$-sections de $X$ telles que $g'(y)=g''(y)$ pour un $y\in Y$, on a $g'=g''$.
\end{proposition}

Il résulte en effet de (17.4.1, b'')) qu'une $Y$-section $s$ de $X$ est une immersion ouverte, et comme $X$ est un $Y$-schéma, $s$ est aussi une immersion fermée (I, 5.4.7); il en résulte que $s$ est un isomorphisme de $Y$ sur un sous-préschéma de $X$ induit sur une partie ouverte et fermée de $X$, et comme $s(Y)$ est connexe, c'est nécessairement une composante connexe de $X$. Le reste de la proposition est immédiat.

\begin{remark}[17.4.10]
\label{IV.17.4.10}
Compte tenu de la remarque (17.4.1.2), on voit que, dans les énoncés (17.4.6) à (17.4.9), on peut remplacer partout les mots « non ramifié » par « formellement non ramifié et localement de type fini ».
\end{remark}

\subsection{Caractérisations des morphismes lisses}
\EGAReaderTarget{ega4.c04.subsec.17-5}{subsection:IV.17.5}

\begin{theorem}[17.5.1]
\label{IV.17.5.1}
Soient $f:X\to Y$ un morphisme localement de présentation finie, $x$ un point de $X$, $y=f(x)$. Les conditions suivantes sont équivalentes:
\begin{enumerate}
  \item[a)] $f$ est lisse au point $x$.
  \item[b)] $f$ est plat au point $x$ et le $k(y)$-préschéma $f^{-1}(y)$ est lisse sur $k(y)$ au point $x$.
  \item[b')] $f$ est régulier au point $x$ (6.8.1).
  \item[c)] L'anneau $\sh{O}_{X,x}$ est une $\sh{O}_{Y,y}$-algèbre formellement lisse pour les topologies discrètes.
\end{enumerate}
\end{theorem}

On peut se borner au cas où $Y=\Spec(A)$, $X=\Spec(C)$, où $C=B/\mathfrak{J}$, $B=A[T_1,\ldots,T_n]$ étant une algèbre de polynômes et $\mathfrak{J}$ un idéal de $B$ de type fini. L'équivalence de a) et c) résulte alors de l'équivalence de a) et c) dans (0, 22.6.4). D'autre part, appliquant ce résultat au morphisme localement de type fini $f^{-1}(y)\to\Spec(k(y))$, on voit que l'équivalence de b) et b') résulte de l'équivalence de a) et b) dans (6.8.6). Reste donc à démontrer l'équivalence de a) et b).

Montrons d'abord que a) entraîne b); notons $\mathfrak p$ l'idéal premier $\mathfrak i_x$ dans $C$, $\mathfrak r$ l'idéal premier $\mathfrak i_y$ dans $A$; on a $\mathfrak p=\mathfrak q/\mathfrak J$, où $\mathfrak q$ est un idéal premier de $B$ et $\mathfrak r$ est l'image réciproque de $\mathfrak q$ dans $A$. L'hypothèse a) entraîne d'abord que $f^{-1}(y)$ est lisse sur $k(y)$ au point $x$ par (17.3.3), et il s'agit de montrer en outre que $C_{\mathfrak p}=\sh{O}_{X,x}$ est un $A_{\mathfrak r}$-module \emph{plat}. Comme $C_{\mathfrak p}$ est une $A_{\mathfrak r}$-algèbre formellement lisse et que $B$ est une $A$-algèbre formellement lisse (pour les topologies discrètes), le critère jacobien (0, 22.6.4), joint à (0, 19.1.12), entraîne qu'il existe dans $\mathfrak J$ un système de $r$ polynômes $u_i$ ($1\leq i\leq r$) et $r$ indices $j_i$ ($1\leq i\leq r$) tels que les images des $u_i$ dans $\mathfrak J_{\mathfrak q}/\mathfrak J_{\mathfrak q}^2$ engendrent ce $B_{\mathfrak q}$-module et que l'on ait
\[
\label{IV.17.5.1.1}
\det(\partial u_i/\partial T_{j_k})\notin\mathfrak q.
\tag{17.5.1.1}
\]

\oldpage[IV-4]{68}
Notons maintenant que si $g:\Spec(B)\to\Spec(A)$ est le morphisme structural, la fibre $g^{-1}(y)$ est le spectre de l'anneau régulier $k(y)[T_1,\ldots,T_n]$ (0, 17.3.7), donc l'anneau local noethérien $B_{\mathfrak q}/\mathfrak rB_{\mathfrak q}$ en un point de cette fibre est régulier. Or, la condition (17.5.1.1) entraîne que les images canoniques $v_i$ des $u_i$ dans l'idéal maximal $\mathfrak m$ de $B_{\mathfrak q}/\mathfrak rB_{\mathfrak q}$ sont \emph{linéairement indépendantes} mod. $\mathfrak m^2$: sinon, en effet, il existerait des polynômes $w_i\in B$ ($1\leq i\leq r$) n'appartenant pas tous à $\mathfrak q$ et tels que $\sum_{i=1}^r w_i u_i\in\mathfrak q^2$. En dérivant par rapport aux $T_{j_k}$, on en conclurait que $\sum_{i=1}^r w_i(\partial u_i/\partial T_{j_k})\in\mathfrak q$ pour $1\leq k\leq r$, ce qui contredirait (17.5.1.1) puisque $\mathfrak q$ est premier. On conclut donc de (0, 17.1.7) que $(v_i)$ est une \emph{suite régulière} dans $B_{\mathfrak q}/\mathfrak rB_{\mathfrak q}$. Mais comme le morphisme $g$ est localement de présentation finie et que $B$ est un $A$-module plat, il résulte de (11.3.8) que les images canoniques $u'_i$ des $u_i$ dans $B_{\mathfrak q}$ forment aussi une \emph{suite régulière} et que $B_{\mathfrak q}/(\sum_i u'_iB_{\mathfrak q})$ est un $A_{\mathfrak r}$-module \emph{plat}. Comme les images des $u'_i$ dans $\mathfrak J_{\mathfrak q}/\mathfrak J_{\mathfrak q}^2$ engendrent ce $B_{\mathfrak q}$-module, il résulte du lemme de Nakayama que l'on a $\sum_i u'_iB_{\mathfrak q}=\mathfrak J_{\mathfrak q}$, et $C_{\mathfrak p}=B_{\mathfrak q}/\mathfrak J_{\mathfrak q}$ est donc bien un $A_{\mathfrak r}$-module plat.

Prouvons enfin que b) implique a). Avec les mêmes notations, l'hypothèse que $B_{\mathfrak q}/\mathfrak J_{\mathfrak q}$ est un $A_{\mathfrak r}$-module plat entraîne que l'homomorphisme canonique $\mathfrak J_{\mathfrak q}/\mathfrak r\mathfrak J_{\mathfrak q}\to B_{\mathfrak q}/\mathfrak rB_{\mathfrak q}$ est \emph{injectif} ($0_{\mathrm I}$, 6.1.2), de sorte que $\mathfrak J_{\mathfrak q}/\mathfrak r\mathfrak J_{\mathfrak q}$ est identifié à un idéal de $B_{\mathfrak q}/\mathfrak rB_{\mathfrak q}$. Comme $B_{\mathfrak q}/\mathfrak rB_{\mathfrak q}$ est une $A_{\mathfrak r}$-algèbre formellement lisse pour les topologies discrètes, on peut appliquer à $C_{\mathfrak p}=(B_{\mathfrak q}/\mathfrak rB_{\mathfrak q})/(\mathfrak J_{\mathfrak q}/\mathfrak r\mathfrak J_{\mathfrak q})$, le critère jacobien (0, 22.6.4); joint à (0, 19.1.12), ce dernier prouve, en vertu de l'hypothèse b), l'existence de $r$ polynômes $v_i\in k(y)[T_1,\ldots,T_n]$ tels que leurs images dans $(\mathfrak J_{\mathfrak q}/\mathfrak r\mathfrak J_{\mathfrak q})/(\mathfrak J_{\mathfrak q}/\mathfrak r\mathfrak J_{\mathfrak q})^2$ engendrent ce $(B_{\mathfrak q}/\mathfrak rB_{\mathfrak q})$-module et que l'on ait
\[
\label{IV.17.5.1.2}
\det(\partial v_i/\partial T_{j_k})\notin\mathfrak qB_{\mathfrak q}/\mathfrak rB_{\mathfrak q}.
\tag{17.5.1.2}
\]
Si, pour tout $i$, on désigne alors par $u_i$ un élément de $\mathfrak J$ dont $v_i$ est l'image canonique, il résulte de (17.5.1.2) que les $u_i$ vérifient la condition (17.5.1.1); d'autre part, en vertu du lemme de Nakayama, les images $u'_i$ des $u_i$ dans $\mathfrak J_{\mathfrak q}$ engendrent ce $B_{\mathfrak q}$-module. Le critère jacobien (0, 22.6.4) joint à (0, 19.1.12), prouve alors que $C_{\mathfrak p}=B_{\mathfrak q}/\mathfrak J_{\mathfrak q}$ est une $A_{\mathfrak r}$-algèbre formellement lisse pour les topologies discrètes. C.Q.F.D.

\begin{corollary}[17.5.2]
\label{IV.17.5.2}
Soit $f:X\to Y$ un morphisme localement de présentation finie. Pour que $f$ soit lisse (au sens de (17.3.1)), il faut et il suffit que $f$ soit régulier (6.8.1), autrement dit que $f$ soit plat et que pour tout $y\in Y$, $f^{-1}(y)$ soit un $k(y)$-préschéma géométriquement régulier (6.7.6).
\end{corollary}

On a ainsi établi l'équivalence des deux définitions de « morphisme lisse » données dans (6.8.1) et (17.3.1).

\begin{proposition}[17.5.3]
\label{IV.17.5.3}
Soient $Y$ un préschéma localement noethérien, $f:X\to Y$ un morphisme localement de type fini, $x$ un point de $X$, $y=f(x)$. Posons $A=\sh{O}_{Y,y}$, $B=\sh{O}_{X,x}$, qui sont des anneaux locaux noethériens. Alors les conditions équivalentes a) à c) de (17.5.1) sont aussi équivalentes à chacune des suivantes:
\begin{enumerate}
  \item[d)] $B$ est une $A$-algèbre formellement lisse pour les topologies préadiques.
  \item[d')] $\widehat B$ est une $\widehat A$-algèbre formellement lisse pour les topologies adiques.
\end{enumerate}
Si de plus $k(x)=k(y)$, ces conditions équivalent aussi à:
\begin{enumerate}
  \item[d'')] $\widehat B$ est une $\widehat A$-algèbre isomorphe à une algèbre de séries formelles $\widehat A[[T_1,\ldots,T_n]]$.
\end{enumerate}
\end{proposition}

\oldpage[IV-4]{69}
L'équivalence de la condition c) de (17.5.1) et de d) résulte de l'équivalence de a) et d) dans le critère jacobien (0, 22.6.4), et l'équivalence de d) et d') résulte de (0, 19.3.6). D'autre part, d'') implique d') sans hypothèse sur les corps résiduels (0, 19.3.4). Enfin, si $\mathfrak m$ désigne l'idéal maximal de $\widehat A$, l'hypothèse d') entraîne que $\widehat B/\mathfrak m\widehat B$ est une $k(y)$-algèbre locale noethérienne complète, formellement lisse pour sa topologie adique (0, 19.3.5); l'hypothèse $k(y)=k(x)$ entraîne donc que $\widehat B/\mathfrak m\widehat B$ est $k(y)$-isomorphe à une algèbre de séries formelles $k(y)[[T_1,\ldots,T_n]]$ (0, 19.6.4). Comme d'autre part, $\widehat A[[T_1,\ldots,T_n]]$ est un $\widehat A$-module plat et une $\widehat A$-algèbre locale noethérienne complète, on conclut de (0, 19.7.1.5) que cette algèbre est isomorphe à $\widehat B$. Donc d') entraîne d'') sous l'hypothèse additionnelle $k(x)=k(y)$.

\begin{remark}[17.5.4]
\label{IV.17.5.4}
Supposons que $Y$ soit un préschéma localement noethérien, et $f:X\to Y$ un morphisme localement de type fini. Le critère (17.5.3, d)), joint à (0, 22.1.4) montre que pour démontrer que $f$ est lisse, on peut appliquer la définition (17.1.1), en se restreignant au cas où le schéma affine $Y'$ est le spectre d'un anneau local artinien.
\end{remark}

\begin{proposition}[17.5.5]
\label{IV.17.5.5}
Soient $Y$ un préschéma localement noethérien, $f:X\to Y$ un morphisme localement de type fini, $x$ un point de $X$, $y=f(x)$. Supposons que $Y$ soit réduit au point $y$. Alors, pour que $f$ soit lisse au point $x$, il faut et il suffit que $f$ soit universellement ouvert dans un voisinage de $x$ dans $f^{-1}(y)$ et que $f^{-1}(y)$ soit un $k(y)$-préschéma géométriquement régulier au point $x$.
\end{proposition}

Compte tenu de (17.5.1), tout revient à voir que, si $f^{-1}(y)$ est un $k(y)$-préschéma géométriquement régulier au point $x$, il est équivalent de dire que $f$ est plat au point $x$ ou universellement ouvert dans un voisinage de $x$ dans $f^{-1}(y)$. Or, si $f$ est plat au point $x$, il l'est dans un voisinage de $x$ dans $X$ (11.1.1) et par suite est universellement ouvert dans ce voisinage (2.4.6). Réciproquement, l'hypothèse que $f$ est universellement ouvert dans un voisinage de $x$ dans $f^{-1}(y)$ et que $f^{-1}(y)$ est un $k(y)$-préschéma géométriquement régulier au point $x$ entraîne que $f$ est plat au point $x$, puisque $\sh{O}_{Y,y}$ est réduit (15.2.2).

\begin{corollary}[17.5.6]
\label{IV.17.5.6}
Soient $Y$ un préschéma localement noethérien, $f:X\to Y$ un morphisme localement de type fini, $x$ un point de $X$, $y=f(x)$. On suppose que $Y$ soit réduit et géométriquement unibranche (6.15.1) au point $y$. Alors, pour que $f$ soit lisse au point $x$, il faut et il suffit que $f$ soit équidimensionnel au point $x$ et que $f^{-1}(y)$ soit un $k(y)$-préschéma géométriquement régulier au point $x$.
\end{corollary}

En remarquant que l'ensemble des points où $f$ est équidimensionnel est ouvert (13.3.2), on voit que le corollaire résulte de (17.5.5) et du critère de Chevalley (14.4.4).

Le fait que $f$ soit un morphisme lisse en un point entraîne en particulier que $f$ vérifie

\oldpage[IV-4]{70}
en ce point toutes les propriétés définies dans (6.8.1). On a donc les propriétés suivantes, que nous rappelons pour la commodité des références:

\begin{proposition}[17.5.7]
\label{IV.17.5.7}
Soit $f:X\to Y$ un morphisme localement de présentation finie, lisse en un point $x\in X$; posons $y=f(x)$. Alors, pour que l'anneau $\sh{O}_{X,x}$ soit réduit (resp. intégralement clos, resp. géométriquement unibranche), il faut et il suffit que $\sh{O}_{Y,y}$ le soit.
\end{proposition}

Cela a en effet été prouvé dans (11.3.13) et (11.3.14) complété par Err$_{\mathrm{IV}}$, 30.

\begin{proposition}[17.5.8]
\label{IV.17.5.8}
Soient $Y$ un préschéma localement noethérien, $f:X\to Y$ un morphisme localement de type fini, lisse en un point $x\in X$; posons $y=f(x)$. Alors:
\begin{enumerate}
  \item[(i)] On a $\dim(\sh{O}_{X,x})=\dim(\sh{O}_{Y,y})+\dim(\sh{O}_{X,x}\otimes_{\sh{O}_{Y,y}}k(y))$.
  \item[(ii)] On a $\operatorname{coprof}(\sh{O}_{X,x})=\operatorname{coprof}(\sh{O}_{Y,y})$.
  \item[(iii)] Pour que l'anneau $\sh{O}_{X,x}$ possède la propriété $(S_n)$ (5.7.2) (resp. $(R_n)$ (5.8.2)), il faut et il suffit que l'anneau $\sh{O}_{Y,y}$ la possède. En particulier, pour que $\sh{O}_{X,x}$ soit régulier, il faut et il suffit que $\sh{O}_{Y,y}$ le soit.
\end{enumerate}
\end{proposition}

Ce sont des cas particuliers de (6.1.2), (6.3.2), (6.4.1) et (6.5.3).

\subsection{Caractérisations des morphismes étales}
\EGAReaderTarget{ega4.c04.subsec.17-6}{subsection:IV.17.6}

\begin{theorem}[17.6.1]
\label{IV.17.6.1}
Soient $f:X\to Y$ un morphisme localement de présentation finie, $x$ un point de $X$, $y=f(x)$. Les conditions suivantes sont équivalentes:
\begin{enumerate}
  \item[a)] $f$ est étale au point $x$.
  \item[a')] $f$ est lisse au point $x$ et non ramifié au point $x$.
  \item[b)] $f$ est lisse au point $x$ et quasi-fini au point $x$ (Err$_{\mathrm{III}}$, 20).
  \item[c)] $f$ est plat au point $x$ et non ramifié au point $x$.
  \item[c')] $f$ est plat au point $x$ et l'anneau $\sh{O}_{X,x}/\mathfrak m_y\sh{O}_{X,x}$ est un corps, extension finie séparable de $k(y)$.
  \item[d)] L'anneau $\sh{O}_{X,x}$ est une $\sh{O}_{Y,y}$-algèbre formellement étale pour les topologies discrètes.
\end{enumerate}
\end{theorem}

L'équivalence de a) et a') résulte aussitôt des définitions; celle de a) et d) résulte de l'équivalence de a) et e) dans (17.4.1) et de l'équivalence de a) et d) dans (17.5.1). L'équivalence de c) et c') résulte de l'équivalence de a) et d') dans (17.4.1). Le fait que a') implique c') découle de (17.5.1); inversement, si c') est vérifiée, $f$ est régulier (donc lisse par (17.5.1)) au point $x$, car si $K$ est une extension finie séparable d'un corps $k$, alors, pour toute extension $k'$ de $k$, $\Spec(K\otimes_k k')$ est régulier, étant somme d'un nombre fini de spectres de corps. Le fait que a') implique b) résulte de (17.4.1, d')) et de (17.5.1, b)). Il reste donc à voir que b) entraîne c), et comme on sait déjà que $f$ est plat au point $x$, par (17.5.1), il suffit de montrer que $f^{-1}(y)$ est un $k(y)$-préschéma non ramifié sur $k(y)$; autrement dit, on est ramené à prouver que b) entraîne c) lorsque $Y=\Spec(k)$ est le spectre d'un corps $k$. Comme la question est locale sur $X$, on peut se borner au cas où $X=\Spec(A)$, où $A$ est une $k$-algèbre finie et locale ($0_{\mathrm I}$, 7.4.1). En vertu de l'hypothèse b), $A$ est une $k$-algèbre formellement lisse pour les topologies discrètes, qui coïncident ici avec les topologies préadiques; donc (0, 19.6.5), $A$ est un anneau local régulier, donc

\oldpage[IV-4]{71}
un corps puisqu'il est artinien, et il résulte alors de (0, 19.6.5.1) que $A$ doit être une extension finie et séparable de $k$, ce qui achève la démonstration (17.4.1).

\begin{corollary}[17.6.2]
\label{IV.17.6.2}
Soit $f:X\to Y$ un morphisme localement de présentation finie. Les conditions suivantes sont équivalentes:
\begin{enumerate}
  \item[a)] $f$ est étale.
  \item[a')] $f$ est lisse et non ramifié.
  \item[b)] $f$ est lisse et localement quasi-fini (Err$_{\mathrm{III}}$, 20).
  \item[c)] $f$ est plat et non ramifié.
  \item[c')] $f$ est plat, et toute fibre $f^{-1}(y)$ est un $k(y)$-schéma somme de spectres de corps, extensions finies et séparables de $k(y)$.
  \item[c'')] $f$ est plat, et pour tout $y\in Y$ et toute extension algébriquement close $k'$ de $k(y)$, la « fibre géométrique » $f^{-1}(y)\otimes_{k(y)}k'$ est somme de spectres de corps isomorphes à $k'$.
\end{enumerate}
\end{corollary}

Le seul point qui reste à démontrer est l'équivalence de c') et c''). Il est clair que c') entraîne c'') par changement de base (17.3.3). D'autre part, comme le morphisme projection $f^{-1}(y)\otimes_{k(y)}k'\to f^{-1}(y)$ est ouvert (2.4.10), l'hypothèse c'') entraîne que l'espace $f^{-1}(y)$ est discret, donc, pour tout point $x\in X_y=f^{-1}(y)$, l'anneau local $\sh{O}_{X_y,x}=A$ est une $k(y)$-algèbre finie, donc un anneau local artinien; en outre il résulte de c'') que $\Spec(A\otimes_{k(y)}k')$ est somme de spectres de corps isomorphes à $k'$, ce qui n'est possible que si $A$ est un corps, extension finie séparable de $k(y)$ (4.6.1).

\begin{proposition}[17.6.3]
\label{IV.17.6.3}
Soient $Y$ un préschéma localement noethérien, $f:X\to Y$ un morphisme localement de type fini, $x$ un point de $X$, $y=f(x)$. Posons $A=\sh{O}_{Y,y}$, $B=\sh{O}_{X,x}$, qui sont des anneaux locaux noethériens, et soit $k$ le corps résiduel de $A$. Alors les conditions équivalentes a) à d) de (17.6.1) sont aussi équivalentes à chacune des suivantes:
\begin{enumerate}
  \item[e)] $\widehat B$ est une $\widehat A$-algèbre formellement étale pour les topologies adiques.
  \item[e')] $\widehat B$ est un $\widehat A$-module libre et $\widehat B\otimes_{\widehat A}k$ est un corps, extension finie et séparable de $k$ (ce qui entraîne que $\widehat B$ est une $\widehat A$-algèbre finie).
\end{enumerate}
Si de plus $k(x)=k(y)$, ou si $k$ est séparablement clos, ces conditions équivalent aussi à:
\begin{enumerate}
  \item[e'')] L'homomorphisme canonique $\widehat A\to\widehat B$ est bijectif.
\end{enumerate}
\end{proposition}

L'équivalence de e) et de chacune des conditions de (17.6.1) résulte aussitôt de (17.4.4, f')) et (17.5.3, d')). Le fait que e) entraîne e') résulte de (17.4.4, f)) et de (0, 19.7.1), compte tenu de ce que $\widehat B$ est alors une $\widehat A$-algèbre finie (17.4.4) et qu'il revient donc au même de dire que $\widehat B$ est un $\widehat A$-module plat ou un $\widehat A$-module libre ($0_{\mathrm{III}}$, 10.1.3). Inversement, le fait que e') entraîne e) résulte de (17.4.4) et de (0, 19.7.1). Enfin, e') entraîne que l'homomorphisme $\widehat A\to\widehat B$ est injectif, et si $k(x)=k(y)$ ou si $k$ est séparablement clos, cet homomorphisme est surjectif par (17.4.4). La réciproque est immédiate.

\begin{proposition}[17.6.4]
\label{IV.17.6.4}
Sous les hypothèses de (17.6.3), si $f$ est étale au point $x$, on a $\dim(\sh{O}_{X,x})=\dim(\sh{O}_{Y,y})$.
\end{proposition}

C'est un cas particulier de (17.5.8, (i)) puisque $x$ est isolé dans sa fibre $f^{-1}(y)$.

\oldpage[IV-4]{72}

\subsection{Propriétés de descente, de passage à la limite et de constructibilité}
\EGAReaderTarget{ega4.c04.subsec.17-7}{subsection:IV.17.7}
\label{IV.17.7}

\begin{proposition}[17.7.1]
\label{IV.17.7.1}
Soient $f:X\to Y$ un morphisme localement de présentation finie, $g:Y'\to Y$ un morphisme, $X'=X\times_Y Y'$, $f'=f_{(Y')}:X'\to Y'$ et $g':X'\to X$ les projections canoniques. Soit $x'$ un point de $X'$ et posons $x=g'(x')$, $y'=f'(x')$.
\begin{enumerate}
  \item[(i)] Si $f'$ est non ramifié au point $x'$, alors $f$ est non ramifié au point $x$.
  \item[(ii)] Supposons de plus que $g$ soit plat au point $y'$. Alors, si $f'$ est lisse (resp. étale) au point $x'$, $f$ est lisse (resp. étale) au point $x$.
\end{enumerate}
\end{proposition}

Posons $y=f(x)=g(y')$, de sorte que l'on a $f'^{-1}(y')=f^{-1}(y)\otimes_{k(y)}k(y')$; notons que $f'$ est localement de présentation finie.

(i) Comme la propriété pour un morphisme (localement de présentation finie) d'être non ramifié en un point ne fait intervenir que la fibre du morphisme en ce point (17.4.1, $d$)), on peut se borner au cas où $Y$ et $Y'$ sont des spectres de corps. Mais alors il revient au même de dire que $f$ (resp. $f'$) est non ramifié en $x$ (resp. $x'$) ou qu'il est étale en ce point (17.6.1, $c$)), donc (i) est une conséquence de (ii).

(ii) Comme $f'$ est plat au point $x'$ (17.5.1), l'hypothèse que $g$ est plat au point $y'$ entraîne que $f$ est plat au point $x$, car la projection $g':X'\to X$ est un morphisme plat au point $x'$, et $f\circ g'=g\circ f'$ est un morphisme plat au point $X'$, d'où la conclusion (2.2.11, (iv)). Le fait que $f$ soit lisse (resp. étale) au point $x$ ne fait plus alors intervenir que la fibre $f^{-1}(y)$ (17.5.1 et 17.6.1), et on est donc encore ramené au cas où $Y=\Spec(k)$ et $Y'=\Spec(k')$ sont des spectres de corps. Dire que $f'$ est lisse au point $x'$ signifie alors (17.5.1) que $X'$ est un $k'$-préschéma géométriquement régulier au point $x'$, et cela entraîne (6.7.8) que $X$ est un $k$-préschéma géométriquement régulier au point $x$, donc que $f$ est lisse au point $x$. Supposons de plus que $f'$ soit étale au point $x'$, de sorte que $x'$ est isolé dans $f'^{-1}(y')$ (17.6.1); comme la projection $f'^{-1}(y')\to f^{-1}(y)$ est un morphisme ouvert (2.4.10), $x$ est isolé dans $f^{-1}(y)$; comme on sait déjà que $f$ est lisse au point $x$, il est étale en ce point (17.6.1).

\begin{corollary}[17.7.2]
\label{IV.17.7.2}
\begin{enumerate}
  \item[(i)] Avec les notations de (17.7.1), soit $U$ (resp. $U'$) l'ensemble des points de $X$ (resp. $X'$) où $f$ (resp. $f'$) est non ramifié; alors on a $U'=g'^{-1}(U)$.
  \item[(ii)] Supposons de plus que $g$ soit plat, et soit $V$ (resp. $V'$) l'ensemble des points de $X$ où $f$ (resp. $f'$) est lisse; alors $V'=g'^{-1}(V)$.
\end{enumerate}
\end{corollary}

\begin{corollary}[17.7.3]
\label{IV.17.7.3}
\begin{enumerate}
  \item[(i)] Supposons $g$ surjectif; alors, pour que $f$ soit non ramifié, il faut et il suffit que $f'$ le soit.
  \item[(ii)] Soient $f:X\to Y$ un $S$-morphisme, $g:S'\to S$ un morphisme fidèlement plat. Supposons que $f$ soit localement de présentation finie, ou que $g$ soit quasi-compact. Alors, pour que $f$ soit lisse (resp. non ramifié, resp. étale), il faut et il suffit que $f'$ le soit.
\end{enumerate}
\end{corollary}

Dans (ii), le cas où $f$ est localement de présentation finie résulte de (17.7.1). Si $g$ est quasi-compact et $f'$ lisse (resp. non ramifié, resp. étale), $f'$ est localement de présentation finie par (2.7.1, (iv)) et on est ramené au premier cas.

\begin{proposition}[17.7.4]
\label{IV.17.7.4}
Soient $f:X\to Y$ un morphisme, $g:Y'\to Y$ un morphisme plat et localement de présentation finie, $X'=X\times_Y Y'$, $f'=f_{(Y')}:X'\to Y'$ et $g':X'\to X$ les

\oldpage[IV-4]{73}

projections canoniques. Soit $V$ (resp. $V'$) l'ensemble des $x\in X$ (resp. $x'\in X'$) où $f$ possède l'une des propriétés suivantes (resp. où $f'$ possède la même propriété): être:
\begin{enumerate}
  \item[(i)] localement de type fini;
  \item[(ii)] localement de présentation finie;
  \item[(iii)] plat;
  \item[(iv)] non ramifié;
  \item[(v)] lisse;
  \item[(vi)] étale.
\end{enumerate}
Alors on a $V'=g'^{-1}(V)$ (autrement dit, pour que $f'$ ait la propriété considérée en un point $x'$, il faut et il suffit que $f$ ait cette propriété au point $x$).
\end{proposition}

La propriété (iii) n'est mise que pour mémoire, et ne nécessite pas l'hypothèse que $g$ soit localement de présentation finie ((2.2.11, (iv)), compte tenu du fait que la projection $g':X'\to X$ est un morphisme plat). En vertu de (17.7.2), les assertions relatives aux propriétés (iv), (v) et (vi) sont des conséquences de l'assertion relative à (ii). Il suffit donc de considérer les cas (i) et (ii). Il est clair que $V'\supset g'^{-1}(V)$; reste donc à montrer que $V'\subset g'^{-1}(V)$; notons que les ensembles $V$ et $V'$ sont ouverts, et $W=g'(V')$ est aussi ouvert dans $X$ en vertu de (2.4.6). Il s'agit donc de prouver que le morphisme $f|W:W\to Y$ est localement de type fini (resp. localement de présentation finie); comme par hypothèse le composé
\[
V'\xrightarrow{g''}W\xrightarrow{f|W}Y
\]
(où $g''$ est la restriction de $g'$), égal à $V'\xrightarrow{f'}Y'\xrightarrow{g}Y$, est localement de type fini (resp. localement de présentation finie) et que $g''$ est surjectif (donc fidèlement plat), on est ramené à prouver le lemme suivant, qui améliore (11.3.16):

\begin{lemma}[17.7.5]
\label{IV.17.7.5}
Soient $f:X\to Y$ un morphisme fidèlement plat et localement de présentation finie, $g:Y\to Z$ un morphisme tel que $g\circ f:X\to Z$ ait l'une des propriétés suivantes: être:
\begin{enumerate}
  \item[(i)] localement de type fini;
  \item[(ii)] localement de présentation finie;
  \item[(iii)] de type fini.
\end{enumerate}
Alors $g$ a la même propriété. Si en outre $f$ est quasi-compact ou $g$ quasi-séparé, la même conclusion est valable pour la propriété:
\begin{enumerate}
  \item[(iv)] être de présentation finie.
\end{enumerate}
\end{lemma}

Dans les cas (i) et (ii), il s'agit de voir que pour tout $y\in Y$, il y a un voisinage ouvert affine $V$ de $y$ dans $Y$ et un voisinage ouvert affine $W$ de $z=g(y)$ dans $Z$ contenant $g(V)$, tels que le morphisme $V\to W$, restriction de $g$, soit de type fini (resp. de présentation finie). Or, il existe par hypothèse un $x\in X$ tel que $f(x)=y$, un voisinage ouvert affine $U$ de $x$ dans $X$, un voisinage ouvert affine $V'$ de $y$ dans $Y$ contenant $f(U)$ et un voisinage ouvert affine $W$ de $z$ dans $Z$ contenant $g(V')$ tels que le morphisme $f_1:U\to V'$, restriction de $f$ soit plat et de présentation finie, et le morphisme $g_1:V'\to W$ restriction de $g$, tel que $g_1\circ f_1$ soit de type fini (resp. de présentation finie). Alors $f(U)$ est ouvert dans $Y$ (2.4.6), et si $V\subset f(U)$ est un voisinage affine de $y$, le morphisme $f_2:f_1^{-1}(V)\to V$,

\oldpage[IV-4]{74}

restriction de $f_1$, est encore de présentation finie et est en outre fidèlement plat; de plus, si $g_2=g_1|V$, $g_2\circ f_2$ est de type fini (resp. de présentation finie), l'ouvert $f_1^{-1}(V)$ étant quasi-compact (resp. quasi-compact et quasi-séparé). On est alors ramené aux hypothèses de (11.3.16), d'où l'on conclut que $g_2$ est un morphisme de type fini (resp. de présentation finie).

Dans le cas (iii) la question est locale sur $Z$, donc on peut supposer $Z$ affine, et il en résulte alors que $X$ est quasi-compact, donc aussi $Y=f(X)$. Le cas (iii) est donc conséquence de (i).

Dans le cas (iv) (avec les hypothèses supplémentaires sur $f$ ou $g$), on peut aussi supposer $Z$ affine, donc $X$ et $Y$ quasi-compacts; on sait en outre déjà que $g$ est localement de présentation finie et quasi-compact (1.1.3), donc tout revient à voir que $g$ est quasi-séparé, et il suffit donc de montrer que cette propriété est vraie lorsque l'on suppose $f$ quasi-compact. Or, puisque $g\circ f$ est quasi-séparé, il en est de même de $f$ (1.2.2, (v)); comme $f$ est quasi-compact et localement de présentation finie, il est de présentation finie (1.6.1); il suffit alors de répéter le raisonnement du premier alinéa de la démonstration de (11.3.16).

\begin{remark}[17.7.6]
\label{IV.17.7.6}
Dans l'assertion relative à la propriété (iv), on ne peut supprimer l'hypothèse que $f$ est quasi-compact; sans quoi, cela impliquerait que tout morphisme $g:Y\to Z$ quasi-compact et localement de présentation finie serait de présentation finie, conclusion que l'on sait être erronée (1.6.4). En effet, on peut se borner au cas où $Z$ est affine, donc $Y$ quasi-compact; il y a par suite un recouvrement fini $(U_i)$ de $Y$ par des ouverts affines tels que les restrictions $g|U_i$ soient de présentation finie; il suffirait alors de prendre pour $X$ le préschéma \emph{somme} des $U_i$, pour $f:X\to Y$ le morphisme canonique, qui est évidemment fidèlement plat et localement de présentation finie (1.4.3); $g\circ f$ serait de présentation finie en vertu du choix des $U_i$ et de (1.6.5), d'où notre assertion.
\end{remark}

\begin{proposition}[17.7.7]
\label{IV.17.7.7}
Soient $f:Y\to S$ et $h:X\to S$ deux morphismes localement de présentation finie, $g:X\to Y$ un $S$-morphisme, $x$ un point de $X$, $y=g(x)$. Supposons que $g$ soit plat au point $x$. Alors, si $h$ est lisse (resp. non ramifié, resp. étale) au point $x$, $f$ est lisse (resp. non ramifié, resp. étale) au point $y$.
\end{proposition}

Posons $s=h(x)=f(y)$. Dire que $h$ est non ramifié au point $x$ (resp. que $f$ est non ramifié au point $y$) équivaut à dire que $h^{-1}(s)$ est étale sur $k(s)$ au point $x$ (resp. que $f^{-1}(s)$ est étale sur $k(s)$ au point $y$) (17.4.1 et 17.6.1); comme le morphisme $g_s:h^{-1}(s)\to f^{-1}(s)$ déduit de $g$ est plat au point $x$, on voit que l'on peut se borner à prouver la proposition lorsque $h$ est lisse ou étale au point $x$. En outre, comme $h$ est alors plat au point $x$ (17.5.1), $f$ est plat au point $y$, comme il résulte de (2.2.11, (iv)). Il revient donc au même (17.5.1) de dire que $f$ est lisse (resp. étale) au point $y$, ou que $f^{-1}(s)$ est lisse (resp. étale) sur $k(s)$ au point $x$. On est ainsi ramené au cas où $S=\Spec(k)$ est le spectre d'un corps.

(i) \emph{Cas des morphismes lisses}. Comme $g$ est un morphisme localement de présentation finie (1.4.3, (v)), il y a un voisinage ouvert $U$ de $x$ dans $X$ dans lequel $g$ est plat (11.3.1) et $h$ lisse. En outre $g(U)$ est un voisinage ouvert de $y$ dans $Y$ (2.4.6); remplaçant $X$ par $U$ et $Y$ par $g(U)$, on peut donc supposer que $g$ est \emph{fidèlement plat} et que $h$ est lisse, et on est ramené à prouver que $f$ est alors lisse. Si $k'$ est une extension algébriquement

\oldpage[IV-4]{75}

close de $k$, $X\otimes_k k'$ est alors lisse sur $k'$, et en vertu de (17.7.3, (ii)), il suffit de prouver que $Y\otimes_k k'$ est lisse sur $k'$ (puisque $\Spec(k')\to\Spec(k)$ est fidèlement plat); on peut donc se borner au cas où $k$ est \emph{algébriquement clos}. Comme l'ensemble des points de $Y$ rationnels sur $k$ est alors très dense dans $Y$ (10.4.8) et comme l'ensemble des points de $Y$ où $Y$ est lisse sur $k$ est ouvert, on voit qu'il suffit de prouver que $Y$ est lisse sur $k$ en tout point rationnel sur $k$. Mais en un tel point $y$, dire que $Y$ est lisse sur $k$ en ce point équivaut à dire que $Y$ est \emph{régulier} en $y$ (17.5.1 et 6.7.8). Or on a $y=g(x)$ pour un $x\in X$ et par hypothèse (17.5.1) $X$ est régulier au point $x$; comme $X$ et $Y$ sont alors localement noethériens et que $g$ est plat, $Y$ est bien régulier au point $y$ (6.5.1, (i)).

(ii) \emph{Cas des morphismes étales}. Par (i) on sait déjà que $f$ est lisse au point $y$; en vertu de (17.6.1), il suffit donc de montrer que $f$ est quasi-fini au point $y$, ou encore que $\sh{O}_{Y,y}$ est une $k$-algèbre finie. Or, comme $\sh{O}_{X,x}$ est un $\sh{O}_{Y,y}$-module fidèlement plat ($0_{\mathrm I}$, 6.6.2), $\sh{O}_{Y,y}$ s'identifie à un sous-$k$-module de $\sh{O}_{X,x}$ et comme $\sh{O}_{X,x}$ est par hypothèse une $k$-algèbre finie, il en est de même de $\sh{O}_{Y,y}$.

\begin{proposition}[17.7.8]
\label{IV.17.7.8}
Les notations étant celles de (8.8.1), on suppose $X_\alpha$ et $Y_\alpha$ localement de présentation finie sur $S_\alpha$. Soient $f_\alpha:X_\alpha\to Y_\alpha$ un $S_\alpha$-morphisme, $f:X\to Y$ le $S$-morphisme correspondant.
\begin{enumerate}
  \item[(i)] Soient $x$ un point de $X$, $x_\lambda$ sa projection canonique dans $X_\lambda$. Pour que $f$ soit lisse (resp. non ramifié, resp. étale) au point $x$, il faut et il suffit qu'il existe $\lambda\geq\alpha$ tel que $f_\lambda$ soit lisse (resp. non ramifié, resp. étale) au point $x_\lambda$.
  \item[(ii)] Supposons de plus $X_\alpha$ quasi-compact. Pour que $f$ soit lisse (resp. non ramifié, resp. étale), il faut et il suffit qu'il existe $\lambda\geq\alpha$ tel que $f_\lambda$ soit lisse (resp. non ramifié, resp. étale).
\end{enumerate}
\end{proposition}

(i) Si $y=f(x)$, $y_\lambda=f_\lambda(x_\lambda)$ est la projection canonique de $y$ dans $Y_\lambda$, et l'on a $f^{-1}(y)=f_\lambda^{-1}(y_\lambda)\otimes_{k(y_\lambda)}k(y)$; la partie de l'énoncé concernant les morphismes non ramifiés résulte donc de (17.7.1, (i)), et il suffit donc de considérer le cas des morphismes lisses. Comme $f$ et $f_\lambda$ sont localement de présentation finie, il revient au même de dire que $f^{-1}(y)$ est géométriquement régulier au point $x$ ou que $f_\lambda^{-1}(y_\lambda)$ est géométriquement régulier au point $x_\lambda$ (6.7.8). La proposition résulte donc de (17.5.1) et de (11.2.6).

(ii) Pour tout $\lambda$, soit $U_\lambda$ l'ensemble ouvert des $x_\lambda\in X_\lambda$ tels que $f_\lambda$ soit lisse (resp. non ramifié, resp. étale) au point $x_\lambda$; soit $V_\lambda$ son image réciproque dans $X$. Comme, par hypothèse, pour tout $x\in X$ il existe, en vertu de (i), un $\lambda$ tel que $f_\lambda$ soit lisse (resp. non ramifié, resp. étale) au point $x_\lambda$, $X$ est réunion des $V_\lambda$. D'ailleurs (17.3.3), pour $\lambda\leq\mu$ on a $V_\lambda\subset V_\mu$, donc, comme $X$ est quasi-compact, il existe un indice $\mu$ tel que $X=V_\mu$. Comme les $X_\lambda$ sont quasi-compacts, il résulte alors de (8.3.4) qu'il existe un indice $\nu\geq\mu$ tel que l'image réciproque de $U_\mu$ dans $X_\nu$ soit $X_\nu$ tout entier, ce qui signifie que $f_\nu$ est lisse (resp. non ramifié, resp. étale) par (17.3.3).

\begin{corollary}[17.7.9]
\label{IV.17.7.9}
Soient $S=\Spec(A)$ un schéma affine, $f:X\to S$ un morphisme. Les conditions suivantes sont équivalentes:
\begin{enumerate}
  \item[a)] $f$ est un morphisme de présentation finie et lisse (resp. non ramifié, resp. étale).

  \oldpage[IV-4]{76}

  \item[b)] Il existe un schéma affine noethérien $S_0=\Spec(A_0)$, un morphisme de type fini $f_0:X_0\to S_0$, et un morphisme $S\to S_0$ tels que le $S$-préschéma $X_0\otimes_{S_0}S$ soit $S$-isomorphe à $X$ et que $f_0$ soit lisse (resp. non ramifié, resp. étale).
  \item[c)] Les conditions de b) sont vérifiées, et en outre $A_0$ est une sous-$\mathbb{Z}$-algèbre de type fini de $A$, le morphisme $S\to S_0$ correspondant à l'injection canonique $A_0\to A$.
\end{enumerate}
\end{corollary}

La démonstration à partir de (17.7.8) est la même que celle de (11.2.7) à partir de (11.2.6).

\begin{proposition}[17.7.10]
\label{IV.17.7.10}
Soient $f:Y\to S$, $h:X\to S$ deux morphismes localement de présentation finie, $g:X\to Y$ un $S$-morphisme, $x$ un point de $X$, $y=g(x)$. Supposons que $h$ soit plat au point $x$ et que $f$ soit non ramifié au point $y$. Alors $f$ est étale au point $y$ et $g$ est plat au point $x$.
\end{proposition}

(On verra plus loin (18.4.9) que l'on peut en fait se dispenser de l'hypothèse que $h$ est localement de présentation finie).

La question étant locale sur $S$, $X$ et $Y$, on peut supposer que $S$, $X$ et $Y$ sont affines, $f,g,h$ des morphismes de présentation finie, $h$ plat et $f$ non ramifié. Compte tenu de (11.2.7) et (17.7.9), on peut supposer en outre que $S$, $X$ et $Y$ sont noethériens; enfin on peut se borner au cas où $S=\Spec(\sh{O}_{S,s})$ (où $s=f(y)=h(x)$). L'anneau $A=\sh{O}_{S,s}$ étant un anneau local noethérien, il existe un anneau local noethérien $B$ complet, de corps résiduel algébriquement clos et un homomorphisme local $A\to B$ faisant de $B$ un $A$-module fidèlement plat ($0_{\mathrm{III}}$, 10.3.1). Remplaçant $X$ et $Y$ par $X\times_S\Spec(B)$ et $Y\times_S\Spec(B)$, on conclut par (2.5.1) et (17.7.1) qu'on peut se borner au cas où $A$ est complet et a un corps résiduel algébriquement clos. L'hypothèse que $f$ est non ramifié au point $y$ entraîne alors (17.4.4) que $\sh{O}_{Y,y}$ est une $\sh{O}_{S,s}$-algèbre finie et un anneau local noethérien complet ($0_{\mathrm I}$, 7.4.2), et comme le corps résiduel de $\sh{O}_{S,s}$ est algébriquement clos, l'homomorphisme $\sh{O}_{S,s}\to\sh{O}_{Y,y}$ est surjectif (17.4.4). Mais, d'autre part, l'hypothèse que $h$ est plat au point $x$ entraîne que l'homomorphisme composé $\sh{O}_{S,s}\to\sh{O}_{Y,y}\to\sh{O}_{X,x}$ est injectif ($0_{\mathrm I}$, 6.5.1), donc l'homomorphisme $\sh{O}_{S,s}\to\sh{O}_{Y,y}$ est bijectif, ce qui montre que $f$ est étale au point $y$ (17.6.3); en outre, $\sh{O}_{X,x}$ est un $\sh{O}_{Y,y}$-module plat, donc $g$ est plat au point $x$.

\begin{proposition}[17.7.11]
\label{IV.17.7.11}
Soient $S$ un préschéma, $X$, $Y$ deux $S$-préschémas localement de présentation finie sur $S$, $f:X\to Y$ un $S$-morphisme. Pour tout $s\in S$, on note $X_s$, $Y_s$, $f_s$ les préschémas et le morphisme déduits de $X$, $Y$, $f$ par le changement de base $\Spec(k(s))\to S$. Alors:
\begin{enumerate}
  \item[(i)] L'ensemble des $x\in X$ tels que, si $s$ est l'image de $x$ dans $S$, $f_s:X_s\to Y_s$ soit lisse (resp. non ramifié, resp. étale, resp. différentiellement lisse) au point $x$, est localement constructible.
  \item[(ii)] Supposons $f$ de présentation finie. L'ensemble des $y\in Y$ tels que, si $s$ est l'image de $y$ dans $S$, $f_s$ soit lisse (resp. non ramifié, resp. étale, resp. différentiellement lisse) en tous les points de $f_s^{-1}(y)$, est localement constructible.
  \item[(iii)] Supposons $X$ et $Y$ de présentation finie sur $S$. L'ensemble des $s\in S$ tels que $f_s$ soit lisse (resp. non ramifié, resp. étale, resp. différentiellement lisse) est localement constructible.
\end{enumerate}
\end{proposition}

Soit $E$ l'ensemble des $x\in X$ vérifiant la propriété considérée dans (i). Alors l'ensemble $F$ des $y\in Y$ vérifiant la propriété correspondante considérée dans (ii) n'est

\oldpage[IV-4]{77}

autre que $Y-f(X-E)$, donc (ii) résulte de (i) et du théorème de Chevalley lorsque $f$ est de présentation finie (1.8.4). De même, si $h:Y\to S$ est le morphisme structural, l'ensemble des $s\in S$ vérifiant la propriété correspondante considérée dans (iii) est $S-h(Y-F)$, donc (iii) découle encore de (ii) et du théorème de Chevalley lorsque $f$ et $h$ sont de présentation finie. Il suffit donc de prouver les assertions de (i).

Prouvons d'abord (i) lorsqu'il s'agit de la propriété d'être \emph{lisse}. La question étant locale sur $X$, on peut se borner au cas où $S=\Spec(A)$, $X=\Spec(B)$ et $Y=\Spec(C)$ sont affines, $B$ et $C$ étant des $A$-algèbres de présentation finie. Raisonnant comme au début de (9.9.1) et utilisant (17.7.2, (ii)), on se ramène au cas où $A$ est noethérien. En vertu de ($0_{\mathrm{III}}$, 9.2.3), on est ramené à voir que si $x\in E$ (resp. si $x\notin E$), il existe un voisinage $V$ de $x$ dans $\overline{\{x\}}$ contenu dans $E$ (resp. dans $X-E$). Désignant par $g:X\to S$ et $h:Y\to S$ les morphismes structuraux, on peut d'abord remplacer $S$ par le sous-préschéma réduit $S'$ de $S$ ayant $\overline{\{g(x)\}}$ pour espace sous-jacent, $X$ par $g^{-1}(S')$, $Y$ par $h^{-1}(S')$, les fibres de $X$ et $X'$ (resp. $Y$ et $Y'$) aux points de $S'$ étant les mêmes. Autrement dit on peut se borner au cas où $S$ est \emph{intègre} et où $\eta=g(x)=h(y)$ (où $y=f(x)$) est son point générique.

$1^\circ$ Supposons d'abord que $x\in E$. Les anneaux locaux $\sh{O}_{X_\eta,x}$ et $\sh{O}_{Y_\eta,y}$ sont respectivement égaux à $\sh{O}_{X,x}$ et $\sh{O}_{Y,y}$; comme la propriété de lissité d'un morphisme de présentation finie en un point ne dépend que de l'anneau local de ce point et de l'anneau local de son image (17.5.1), on voit que l'hypothèse $x\in E$ revient à dire que le morphisme $f$ est lisse au point $x$; il possède alors encore cette propriété aux points d'un voisinage ouvert de $x$ dans $X$, et il suffit d'appliquer (17.3.3, (iii)) pour obtenir alors la conclusion.

$2^\circ$ Supposons en second lieu que $x\in X-E$, et que le morphisme $f_\eta$ ne soit pas plat au point $x$. La conclusion résulte alors du lemme suivant qui précise (11.2.8):

\begin{lemma}[17.7.11.1]
\label{IV.17.7.11.1}
Soient $g:X\to S$, $h:Y\to S$ deux morphismes localement de présentation finie, $f:X\to Y$ un $S$-morphisme, $\sh{F}$ un $\sh{O}_X$-module quasi-cohérent de présentation finie. Alors l'ensemble $E$ des $x\in X$ tels que $\sh{F}_{g(x)}$ soit $f_{g(x)}$-plat au point $x$ est localement constructible.
\end{lemma}

Raisonnant encore comme au début de (9.9.1) et utilisant (2.5.1), on se ramène au cas où $S$, $X$ et $Y$ sont noethériens; puis on se ramène comme ci-dessus au cas où $S$ est intègre, où $\eta=g(x)=h(f(x))$ est son point générique, et il s'agit de montrer que si $x\in E$ (resp. $x\notin E$) il y a un voisinage $V$ de $x$ dans $\overline{\{x\}}$ contenu dans $E$ (resp. dans $X-E$). Le cas où $x\in E$ est conséquence immédiate de (11.1.1). Pour traiter le cas où $x\notin E$, on raisonne comme dans (9.4.7.1), dont nous conservons les notations, de sorte que l'on peut supposer, en remplaçant éventuellement $Y$ et $X$ par des voisinages de $f(x)$ et $x$ respectivement, qu'il existe deux $\sh{O}_Y$-modules cohérents $\sh{G}$, $\sh{H}$ et un $\sh{O}_Y$-homomorphisme $u:\sh{G}\to\sh{H}$ tels que pour tout $s\in S$, $u_s:\sh{G}_s\to\sh{H}_s$ soit injectif, mais que l'homomorphisme $1\otimes u_\eta:\sh{F}_\eta\otimes_{\sh{O}_{Y_\eta}}\sh{G}_\eta\to\sh{F}_\eta\otimes_{\sh{O}_{Y_\eta}}\sh{H}_\eta$ ne soit pas injectif au point $x$, autrement dit $x\in\operatorname{Supp}(\operatorname{Ker}(1\otimes u_\eta))$; si l'on pose $T=\overline{\{x\}}$, on a donc $\operatorname{Supp}(\operatorname{Ker}(1\otimes u_\eta))\supset T_\eta$. Mais on peut supposer que pour tout $s\in S$, on a $\operatorname{Ker}(1\otimes u_s)=(\operatorname{Ker}(1\otimes u))_s$ (9.4.2), donc

\oldpage[IV-4]{78}

$\operatorname{Supp}(\operatorname{Ker}(1\otimes u_s))=(\operatorname{Supp}(\operatorname{Ker}(1\otimes u)))_s$ (I, 9.1.13.1); il résulte enfin de (9.5.2) que pour $s$ dans un voisinage de $\eta$, on a $(\operatorname{Supp}(\operatorname{Ker}(1\otimes u)))_s\supset T_s$, ce qui établit le lemme.

$3^\circ$ Supposons maintenant que $x\in X-E$, que le morphisme $f_\eta$ soit plat au point $x$, mais que $f_\eta$ ne soit pas lisse au point $x$. Remarquons que dire que $f_\eta$ est plat au point $x$ équivaut à dire que $f$ lui-même est plat au point $x$ et en remplaçant $X$ par un voisinage de $x$, on peut supposer que $f$ est plat (11.1.1); on en conclut qu'il en est de même de $f_s$ pour tout $s\in S$, et comme pour tout $y\in Y$, $f^{-1}(y)=f_{h(y)}^{-1}(y)$, il revient au même de dire que $f_{g(x')}$ est lisse au point $x'$ ou de dire que $f$ est lisse au point $x'$. Mais l'ensemble des $x'\in X$ où $f$ est lisse est \emph{ouvert} dans $X$ (12.1.7), donc l'ensemble des $x'\in X$ où $f$ n'est pas lisse est fermé, et puisqu'il contient $x$ par hypothèse, il contient aussi $\overline{\{x\}}$, ce qui achève la démonstration pour la première propriété considérée dans (i).

Prouvons en second lieu (i) lorsqu'il s'agit de la propriété d'être \emph{étale}. Notons pour cela que cette propriété pour $f_s$ au point $x$ équivaut à dire que $f_s$ est à la fois lisse et quasi-fini au point $x$ (17.6.1). Or, il revient au même de dire que $f_{g(x)}$ est quasi-fini au point $x$ ou que $f$ est lui-même quasi-fini en ce point; il résulte donc de (13.1.4) que l'ensemble des points $x$ tels que $f_{g(x)}$ soit quasi-fini au point $x$ est ouvert dans $X$, et \emph{a fortiori} localement constructible; la conclusion résulte donc de ce que l'ensemble des $x$ tels que $f_{g(x)}$ soit lisse au point $x$ est lui aussi localement constructible.

Passons à la preuve de (i) lorsqu'il s'agit de la propriété d'être \emph{différentiellement lisse}. Soit $p:X\times_Y X\to X$ la seconde projection canonique; la seconde projection $X_s\times_{Y_s}X_s\to X_s$ n'est autre que $p_s$ pour tout $s\in S$, et il résulte donc de (17.12.5.1)\footnote{Le lecteur vérifiera que le résultat de (17.7.11) n'est pas utilisé dans la suite du § 17 et qu'il n'y a donc pas de cercle vicieux.} que pour que $f_{g(x)}$ soit différentiellement lisse au point $x$, il faut et il suffit que $p_{g(x)}$ soit lisse au point $\Delta_f(x)$. Comme $p$ est localement de présentation finie, l'ensemble des points $z\in X\times_Y X$ tels que $p_{g(p(z))}$ soit lisse au point $z$ est localement constructible, et l'intersection de cet ensemble avec l'ensemble localement fermé $\Delta_f(X)$ est donc aussi localement constructible dans $\Delta_f(X)$ (1.8.2); la restriction de $p$ à $\Delta_f(X)$ étant un isomorphisme sur $X$, on voit que l'ensemble des $x\in X$ tels que $f_{g(x)}$ soit différentiellement lisse au point $x$ est localement constructible.

Considérons enfin la propriété d'être non ramifié; notons que le morphisme diagonal $\Delta_f:X\to X\times_Y X$ est une immersion localement de présentation finie (1.4.3.1) et pour tout $s\in S$, le morphisme diagonal $\Delta_{f_s}:X_s\to X_s\times_{Y_s}X_s$ n'est autre que $(\Delta_f)_s$; dire que $f_{g(x)}$ est non ramifié au point $x$ équivaut à dire que $(\Delta_f)_{g(x)}$ est un isomorphisme local au point $x$ (17.4.1), et puisque c'est une immersion localement de présentation finie, il revient au même (17.9.1)\footnote{Le lecteur vérifiera que le résultat de (17.7.11) n'est pas utilisé dans le reste du § 17 et qu'il n'y a donc pas de cercle vicieux.} de dire que $(\Delta_f)_{g(x)}$ est étale au point $x$; il suffit donc d'appliquer ce qui a été vu plus haut pour la propriété d'être étale.

\oldpage[IV-4]{79}

\subsection{Critères de lissité et non-ramification par fibres}
\EGAReaderTarget{ega4.c04.subsec.17-8}{subsection:IV.17.8}
\label{IV.17.8}

\begin{proposition}[17.8.1]
\label{IV.17.8.1}
Soient $g:Y\to S$, $h:X\to S$ deux morphismes localement de présentation finie. Pour qu'un $S$-morphisme $f:X\to Y$ soit non ramifié, il faut et il suffit que pour tout $s\in S$, le morphisme $f_s:h^{-1}(s)\to g^{-1}(s)$ déduit de $f$ par le changement de base $\Spec(k(s))\to S$, soit non ramifié.
\end{proposition}

Cela résulte aussitôt du fait que, pour un morphisme localement de présentation finie, le fait d'être non ramifié est une propriété des \emph{fibres} de ce morphisme (17.4.1, $d$)), et de ce que, pour tout $y\in Y$, on a $f^{-1}(y)=f_s^{-1}(y)$ si $s=g(y)$.

\begin{proposition}[17.8.2]
\label{IV.17.8.2}
Soient $g:Y\to S$, $h:X\to S$ deux morphismes localement de présentation finie; on suppose en outre que $h$ est plat. Pour qu'un $S$-morphisme $f:X\to Y$ soit lisse (resp. étale), il faut et il suffit que, pour tout $s\in S$, le morphisme $f_s:h^{-1}(s)\to g^{-1}(s)$ déduit de $f$ par le changement de base $\Spec(k(s))\to S$ soit lisse (resp. étale). Lorsqu'il en est ainsi, le morphisme $g$ est plat aux points de $f(X)$.
\end{proposition}

On sait en effet (11.3.10) que pour que $f$ soit plat, il faut et il suffit que $f_s$ le soit pour tout $s\in S$, et qu'alors $g$ est plat aux points de $f(X)$. Mais pour un morphisme plat localement de présentation finie, le fait d'être lisse est une propriété des fibres de ce morphisme (17.5.1, $b$)), et pour tout $y\in Y$, on a $f^{-1}(y)=f_s^{-1}(y)$ si $s=g(y)$.

\begin{remark}[17.8.3]
\label{IV.17.8.3}
Les démonstrations précédentes montrent (compte tenu de (11.3.10)) que si les hypothèses sur $g$ et $h$ sont les mêmes que ci-dessus, alors, pour que $f$ soit non ramifié (resp. lisse, resp. étale) \emph{en un point} $x\in X$, il suffit que, si l'on pose $s=h(x)$, $f_s$ soit non ramifié (resp. lisse, resp. étale) au point $x$.
\end{remark}

\subsection{Morphismes étales et immersions ouvertes}
\EGAReaderTarget{ega4.c04.subsec.17-9}{subsection:IV.17.9}
\label{IV.17.9}

\begin{theorem}[17.9.1]
\label{IV.17.9.1}
Soit $f:X\to Y$ un morphisme. Les propriétés suivantes sont équivalentes:
\begin{enumerate}
  \item[a)] $f$ est une immersion ouverte.
  \item[b)] $f$ est un monomorphisme plat localement de présentation finie.
  \item[c)] $f$ est étale et radiciel.
\end{enumerate}
\end{theorem}

Il résulte de (1.4.3, (i)) que a) implique b). La condition b) implique que pour tout $y\in Y$, la fibre $f^{-1}(y)$ est vide ou isomorphe à $\Spec(k(y))$ (8.11.5.1), donc b) entraîne c), en vertu de (17.6.2, $c$)). Reste à voir que c) implique a).

La question étant locale sur $X$ et sur $Y$ (puisque $f$ est injectif), on peut se borner au cas où $Y$ est affine et $f$ de présentation finie. Comme $f$ est plat, c'est un morphisme \emph{ouvert} (2.4.6), donc, en remplaçant $Y$ par $f(X)$, on peut supposer que $f$ est \emph{surjectif}. Pour tout morphisme $Y'\to Y$, $f'=f_{(Y')}:X_{(Y')}\to Y'$ est encore étale, radiciel, surjectif et de présentation finie, donc ouvert, et par suite un homéomorphisme; autrement dit $f$ est un \emph{homéomorphisme universel}, et étant de type fini et séparé (1.8.7.1), $f$ est \emph{propre}. L'hypothèse que $f$ est radiciel et de type fini entraîne alors que $f$ est quasi-fini; donc (8.11.1) $f$ est un morphisme \emph{fini}. Pour prouver que $f$ est un isomorphisme, on peut se borner au cas

\oldpage[IV-4]{80}

où $Y=\Spec(A)$, $A$ étant un anneau local. Comme $f$ est de présentation finie, on a $X=\Spec(B)$, où $B$ est un $A$-module plat de présentation finie (1.4.7), donc \emph{libre} (Bourbaki, \emph{Alg. comm.}, chap. II, § 5, n\textsuperscript{o} 2, cor. 2 du th. 1). En outre, si $\mathfrak m$ est l'idéal maximal de $A$ et $k$ son corps résiduel, $B/\mathfrak mB$ est par hypothèse un corps, à la fois extension radicielle et extension finie séparable de $k$, puisque $f$ est étale et radiciel (17.6.1); donc $B/\mathfrak mB$ est isomorphe à $k$, et comme $B$ est un $A$-module libre, $B$ est isomorphe à $A$. C.Q.F.D.

\begin{corollary}[17.9.2]
\label{IV.17.9.2}
Soit $X$ un préschéma connexe; si $f:X\to Y$ est une immersion fermée étale, alors $f$ est un isomorphisme de $X$ sur une composante connexe ouverte de $Y$.
\end{corollary}

En effet $f$ est étale et radiciel, donc un isomorphisme de $X$ sur un préschéma induit sur une partie ouverte de $Y$; mais par hypothèse $f(X)$ est fermé dans $Y$, donc à la fois ouvert et fermé, et puisque $f(X)$ est connexe, c'est une composante connexe de $Y$.

\begin{corollary}[17.9.3]
\label{IV.17.9.3}
Soit $f:X\to Y$ un morphisme étale (resp. étale et séparé). Alors toute $Y$-section $g:Y\to X$ de $X$ est une immersion ouverte (resp. ouverte et fermée); en outre l'application $g\mapsto g(Y)$ est une bijection de l'ensemble $\Gamma(X/Y)$ des $Y$-sections de $X$ sur l'ensemble des parties ouvertes $Z$ (resp. ouvertes et fermées) de $X$ telles que la restriction de $f$ à $Z$ soit un morphisme surjectif et radiciel de $Z$ sur $Y$.
\end{corollary}

En effet, le fait que $g$ soit une immersion ouverte résulte déjà de ce que $f$ est non ramifié (17.4.1, $b''$)), et la restriction de $f$ à l'ouvert $g(Y)$ de $X$ est un isomorphisme. Inversement, si $Z$ est un ouvert de $X$ tel que $f|Z$ soit un morphisme surjectif et radiciel de $Z$ sur $Y$, $f|Z$ est un isomorphisme en vertu de (17.9.1), puisqu'il est étale. Si $f$ est en outre séparé, on sait que $g$ est une immersion fermée (1, 5.4.6), ce qui achève de prouver le corollaire.

\begin{corollary}[17.9.4]
\label{IV.17.9.4}
Soient $Y$ un préschéma connexe, $f:X\to Y$ un morphisme étale et séparé. Alors toute $Y$-section $g$ de $X$ est un isomorphisme de $Y$ sur une composante connexe ouverte de $X$, et l'application $g\mapsto g(Y)$ est une bijection de $\Gamma(X/Y)$ sur l'ensemble des composantes connexes ouvertes $Z$ de $X$ telles que la restriction de $f$ à $Z$ soit un morphisme surjectif et radiciel de $Z$ sur $Y$.
\end{corollary}

\begin{corollary}[17.9.5]
\label{IV.17.9.5}
Soient $g:Y\to S$, $h:X\to S$ deux morphismes localement de présentation finie; on suppose en outre que $h$ est plat. Pour qu'un $S$-morphisme $f:X\to Y$ soit une immersion ouverte (resp. un isomorphisme), il faut et il suffit que, pour tout $s\in S$, le morphisme $f_s:h^{-1}(s)\to g^{-1}(s)$ déduit de $f$ par le changement de base $\Spec(k(s))\to S$, soit une immersion ouverte (resp. un isomorphisme).
\end{corollary}

En effet, si $f_s$ est une immersion ouverte pour tout $s\in S$, il résulte de (17.8.3) que $f$ est un morphisme étale; comme pour tout $y\in Y$, on a $f^{-1}(y)=f_s^{-1}(y)$ avec $s=g(y)$, $f$ est radiciel; donc $f$ est une immersion ouverte en vertu de (17.9.1). Si de plus $f_s$ est surjectif pour tout $s\in S$, $f$ est surjectif, donc un isomorphisme.

La proposition suivante précise (10.4.11):

\begin{proposition}[17.9.6]
\EGAReaderTarget{ega4.c04.prop.17-9-6}{IV.17.9.6}
Soient $S$ un préschéma, $X$ un $S$-préschéma de présentation finie. Tout $S$-endomorphisme de $X$ qui est un monomorphisme est un automorphisme de $X$.
\end{proposition}

Soient $f:X\to S$ le morphisme structural, $g$ le $S$-endomorphisme considéré. La question est locale sur $S$, et on peut par suite supposer que $S$ est affine. Utilisant (8.9.1) et (8.10.5, (i \emph{bis})), on est ramené au cas où $S=\Spec(A)$, où $A$ est une $\mathbb{Z}$-algèbre de type fini, et par suite $X$ est un $\mathbb{Z}$-préschéma de type fini. Il résulte déjà de (10.4.11) que $g$ est un morphisme bijectif, puisqu'un monomorphisme est radiciel (8.11.5.1); il suffira donc de montrer

\oldpage[IV-4]{81}

que $g$ est une immersion ouverte, et puisque $g$ est radiciel, il suffira, en vertu de (17.9.1), de prouver que $g$ est étale. En outre, comme l'ensemble des points où $g$ est étale est ouvert et que $X$ est un préschéma de Jacobson (10.4.7), il suffira de montrer que $g$ est étale en tout point \emph{fermé} $z$ de $X$ (10.3.1). Posons $z'=g(z)$; il résulte de (10.4.11.1, (i)) que $z'$ est aussi un point fermé de $X$ et puisque $g$ est un monomorphisme, l'application $k(z')\to k(z)$ déduite de $g$ est un \emph{isomorphisme}. Pour que $g$ soit étale au point $z$, il faut et il suffit donc que l'homomorphisme canonique $\widehat{\sh{O}}_{X,z'}\to\widehat{\sh{O}}_{X,z}$ soit bijectif (17.6.3, $e''$)). Nous allons pour cela prouver que pour tout entier $n$, l'homomorphisme canonique $\sh{O}_{X,z'}/\mathfrak m_{z'}^{n+1}\to\sh{O}_{X,z}/\mathfrak m_z^{n+1}$ est bijectif, d'où résultera aussitôt la conclusion. On peut supposer que $z$ appartient à l'ensemble fini $T_{p,d}$ des points fermés $t$ de $X$ tels que $k(t)$ soit une extension de $\mathbb{F}_p$ dont le degré divise $d$; auquel cas il en est de même de $z'$. Désignons par $\sh{J}$ l'idéal cohérent de $\sh{O}_X$ tel que $\sh{J}_x=\sh{O}_x$ pour $x\notin T_{p,d}$, et $\sh{J}_x=\mathfrak m_x^{n+1}$ pour $x\in T_{p,d}$. Soient $T_{p,d,n}$ le sous-préschéma fermé de $X$ défini par $\sh{J}$, $j:T_{p,d,n}\to X$ l'injection canonique; le morphisme composé $g\circ j:T_{p,d,n}\to X$ applique l'ensemble $T_{p,d}$ dans lui-même, et pour tout $x\in T_{p,d}$, l'homomorphisme $\sh{O}_{X,g(x)}\to\sh{O}_{X,x}/\mathfrak m_x^{n+1}$ se factorise en
\[
\sh{O}_{X,g(x)}\to\sh{O}_{X,g(x)}/\mathfrak m_{g(x)}^{n+1}\to\sh{O}_{X,x}/\mathfrak m_x^{n+1}.
\]
Donc (I, 4.9), il existe un unique endomorphisme $g'$ de $T_{p,d,n}$ tel que $j\circ g'=g\circ j$. Comme $j$ est un monomorphisme, et qu'il en est de même de $g$ par hypothèse, on en déduit que $g'$ est aussi un monomorphisme, et la proposition sera prouvée si l'on montre que $g'$ est un \emph{automorphisme} de $T_{p,d,n}$. Or, pour tout $x\in T_{p,d,n}$, on a vu que $k(x)$ est un corps fini, et comme $\sh{O}_{X,x}$ est noethérien, chacun des $\mathfrak m_x^h/\mathfrak m_x^{h+1}$ est un $k(x)$-espace vectoriel de rang fini; d'où l'on conclut aussitôt que l'anneau local $\sh{O}_{X,x}/\mathfrak m_x^{n+1}$ de $T_{p,d,n}$ au point $x$ a un nombre \emph{fini} d'éléments, et \emph{a fortiori} est un $\mathbb{Z}$-module de type fini. Comme $T_{p,d,n}$ est somme des préschémas $\Spec(\sh{O}_{X,x}/\mathfrak m_x^{n+1})$ en nombre fini, c'est un $\mathbb{Z}$-préschéma \emph{fini}; l'endomorphisme $g'$ est donc lui aussi \emph{fini} (II, 6.1.5, (v)), donc propre. Mais un monomorphisme propre est une \emph{immersion fermée} (8.11.5); si $T_{p,d,n}=\Spec(B)$, $g'$ correspond donc à un endomorphisme surjectif $\varphi:B\to B$ de l'anneau $B$. Or, l'ensemble $B$ est fini, donc $\varphi$ est nécessairement bijectif. C.Q.F.D.

\subsection{Dimension relative d'un préschéma lisse sur un autre}
\EGAReaderTarget{ega4.c04.subsec.17-10}{subsection:IV.17.10}

\begin{definition}[17.10.1]
\EGAReaderTarget{ega4.c04.def.17-10-1}{IV.17.10.1}
Soit $f:X\to Y$ un morphisme localement de type fini. On appelle \emph{dimension relative de $f$ au point $x\in X$} (ou \emph{dimension relative de $X$ sur $Y$ au point $x$}) et l'on note $\dim_x f$ l'entier positif $\dim_x(f^{-1}(f(x)))$.
\end{definition}

Dire que $f$ est \emph{quasi-fini} au point $x$ (Err$_{\mathrm{III}}$, 20) équivaut donc à dire que $\dim_x f=0$. On a vu (13.1.3) que la fonction $x\mapsto\dim_x f$ est semi-continue supérieurement. On notera que, même lorsque le morphisme $f$ a la propriété $(S_1)$ (autrement dit (6.8.1) est plat et tel que ses fibres n'aient pas de cycle premier associé immergé), la fonction $x\mapsto\dim_x f$ n'est pas nécessairement continue, comme le montre l'exemple où $Y=\Spec(k)$, où $k$ est un corps, et $X=\Spec(k[U,V,W]/\mathfrak p\mathfrak q)$, où $\mathfrak p=(W)$ et $\mathfrak q=(U)+(V-W)$, idéaux premiers de $k[U,V,W]$ ($X$ étant donc la réunion, dans l'espace à 3 dimensions, d'un plan et d'une droite non parallèle à ce plan).

Dire qu'un morphisme localement de présentation finie $f:X\to Y$ est étale au point $x\in X$ signifie encore que $f$ est lisse au point $x$ et que $\dim_x f=0$ (17.6.1).

\begin{proposition}[17.10.2]
\label{IV.17.10.2}
Soit $f:X\to Y$ un morphisme lisse. Pour tout $x\in X$, le $\sh{O}_X$-Module localement libre $\Omega_f^1$ (17.2.3) a un rang en $x$ égal à $\dim_x f$ (ce qui entraîne que $x\mapsto\dim_x f$ est une fonction continue dans $X$).
\end{proposition}

En effet, si $y=f(x)$, $X_y=f^{-1}(y)$ est lisse sur $k(y)$, et si $f_y:X_y\to\Spec(k(y))$ est le morphisme structural, on a $\Omega_{f_y}^1=\Omega_f^1\otimes_{\sh{O}_X}\sh{O}_{X_y}$ (16.4.5); on peut donc se borner au cas où $Y=\Spec(k)$ est le spectre d'un corps; en outre, en vertu de (16.4.5) et (17.7.1), on peut remplacer $k$ par une extension algébriquement close, autrement dit

\oldpage[IV-4]{82}

supposer $k$ algébriquement clos. L'ensemble des points de $X$ rationnels sur $k$ étant alors dense dans $X$ (10.4.8), on peut se borner au cas où $x$ est rationnel sur $k$. Or (16.4.12) $(\Omega^1_{X/k})_x\otimes_{\sh{O}_{X,x}}k(x)$ est alors $k$-isomorphe à $\mathfrak m_x/\mathfrak m_x^2$, et puisque $\sh{O}_x$ est un anneau local régulier (17.5.1), $\operatorname{rg}_{k(x)}(\mathfrak m_x/\mathfrak m_x^2)=\dim(\sh{O}_x)$ (0, 17.1.1); donc $\operatorname{rg}_{\sh{O}_x}(\Omega^1_{X/k})_x=\dim(\sh{O}_x)$. Mais puisque $k(x)=k$, on a $\dim_x f=\dim(\sh{O}_x)$ (5.2.3). C.Q.F.D.

Nous démontrerons plus loin une réciproque de ce résultat (17.15.5).

\begin{corollary}[17.10.3]
\label{IV.17.10.3}
Soient $f:X\to Y$, $g:Y\to Z$ deux morphismes lisses. Alors, pour tout $x\in X$, on a
\begin{equation}
\label{IV.17.10.3.1}
\tag{17.10.3.1}
\dim_x(g\circ f)=\dim_x f+\dim_{f(x)}g.
\end{equation}
\end{corollary}

En effet, $g\circ f$ est lisse (17.3.3), donc les trois $\sh{O}_X$-Modules $\Omega^1_{X/Y}$, $\Omega^1_{X/Z}$ et $f^*(\Omega^1_{Y/Z})$ sont localement libres (17.2.3 et $0_{\mathrm I}$, 5.4.5); en outre le rang en $x$ de $f^*(\Omega^1_{Y/Z})$ est égal au rang en $y=f(x)$ de $\Omega^1_{Y/Z}$. L'égalité (17.10.3.1) est donc conséquence de (17.10.2) et de l'exactitude de la suite (17.2.3.1).

\begin{corollary}[17.10.4]
\label{IV.17.10.4}
Soient $f:X\to Y$ un morphisme lisse, $X'$ un sous-préschéma de $X$ tel que le morphisme composé $X'\xrightarrow{j}X\xrightarrow{f}Y$ (où $j$ est l'injection canonique) soit lisse. Alors le faisceau conormal $\sh{N}_{X'/X}$ est un $\sh{O}_{X'}$-Module localement libre, et pour tout $x\in X'$, on a
\begin{equation}
\label{IV.17.10.4.1}
\tag{17.10.4.1}
\dim_x f=\dim_x(f\circ j)+\operatorname{rg}_{\sh{O}_{X',x}}(\sh{N}_{X'/X})_x.
\end{equation}
\end{corollary}

En effet, $\Omega^1_{X/Y}\otimes_{\sh{O}_X}\sh{O}_{X'}$ et $\Omega^1_{X'/Y}$ sont alors localement libres et la suite exacte (17.2.5.1) est scindée dans un voisinage convenable de chaque point de $X'$, donc $\sh{N}_{X'/X}$ est localement libre (Bourbaki, \emph{Alg. comm.}, chap. II, § 5, n\textsuperscript{o} 2, th. 1), et la relation (17.10.4.1) résulte aussitôt de l'exactitude de la suite (17.2.5.1).

\subsection{Morphismes lisses de préschémas lisses}
\EGAReaderTarget{ega4.c04.subsec.17-11}{subsection:IV.17.11}

\begin{theorem}[17.11.1]
\label{IV.17.11.1}
Soient $f:Y\to S$ et $h:X\to S$ deux morphismes localement de présentation finie, $g:X\to Y$ un $S$-morphisme, $x$ un point de $X$; posons $y=g(x)$, $s=f(y)=h(x)$. Les conditions suivantes sont équivalentes:
\begin{enumerate}
  \item[a)] $f$ est lisse au point $y$ et $g$ est lisse au point $x$.
  \item[b)] $g$ et $h$ sont lisses au point $x$.
  \item[c)] $h$ est lisse au point $x$, et l'homomorphisme canonique (16.4.18)
  \begin{equation}
  \label{IV.17.11.1.1}
  \tag{17.11.1.1}
  (g^*(\Omega^1_{Y/S}))_x\to(\Omega^1_{X/S})_x
  \end{equation}
  est inversible à gauche (autrement dit est un isomorphisme sur un \emph{facteur direct} de $(\Omega^1_{X/S})_x$).
  \item[c')] $h$ est lisse au point $x$, et l'homomorphisme canonique
  \begin{equation}
  \label{IV.17.11.1.2}
  \tag{17.11.1.2}
  ((\Omega^1_{Y/S})_y\otimes_{\sh{O}_{Y,y}}k(y))\otimes_{k(y)}k(x)\to(\Omega^1_{X/S})_x\otimes_{\sh{O}_{X,x}}k(x)
  \end{equation}
  est injectif.
\end{enumerate}
Supposons en outre que l'homomorphisme $k(y)\to k(x)$ soit bijectif. Alors les conditions précédentes sont aussi équivalentes à la suivante:
\begin{enumerate}
  \item[d)] $h$ est lisse au point $x$, et l'application canonique $\mathrm{T}_{X/S}(x)\to\mathrm{T}_{Y/S}(y)$ de l'espace vectoriel tangent en $x$ à $X$ dans l'espace vectoriel tangent en $y$ à $Y$ (16.5.12) est surjective.
\end{enumerate}
\end{theorem}

Le fait que a) entraîne b) résulte trivialement de (17.3.3, (iii)); b) entraîne c) par application de (17.2.3, (ii)). Pour voir que c) est équivalente à c'), notons qu'en vertu de (17.2.3, (i)), $\Omega^1_{X/S}$ est un $\sh{O}_X$-Module localement libre de type fini; il suffit alors d'appliquer (0, 19.1.12) à l'anneau local $\sh{O}_{X,x}$ et à l'homomorphisme (17.11.1.1) de $\sh{O}_{X,x}$-modules de type fini. Lorsque $k(x)=k(y)$, l'application linéaire tangente $\mathrm{T}_{X/S}(x)\to\mathrm{T}_{Y/S}(y)$ est transposée de (17.11.1.2), par (16.5.12), ce qui établit dans ce cas l'équivalence de c') et d).

\oldpage[IV-4]{83}

Il reste donc à prouver que c) entraîne a). On peut se borner au cas où $S=\Spec(A)$, $Y=\Spec(B)$, $X=\Spec(C)$ sont affines. L'hypothèse implique (17.2.3, (i)) que $(\Omega^1_{X/S})_x=\Omega^1_{C_{\mathfrak m}/A_{\mathfrak s}}$ est un $C_{\mathfrak m}$-module libre: en effet (16.10.6) il existe des éléments $t_i$ ($1\leq i\leq r$) de $B_{\mathfrak p}=\sh{O}_{Y,y}$ tels que les différentielles $d_{B_{\mathfrak p}/A_{\mathfrak s}}(t_i)$ engendrent le $B_{\mathfrak p}$-module $\Omega^1_{B_{\mathfrak p}/A_{\mathfrak s}}$ et que leurs images dans $\Omega^1_{C_{\mathfrak m}/A_{\mathfrak s}}$ fassent partie d'une base de ce $C_{\mathfrak m}$-module libre. Comme $\Omega^1_{Y/S}$ (resp. $\Omega^1_{X/S}$) est un $\sh{O}_Y$-Module (resp. un $\sh{O}_X$-Module) de présentation finie (16.4.22), on peut, en remplaçant au besoin $X$ et $Y$ par des voisinages ouverts affines convenables de $x$ et $y$ respectivement, supposer que $\Omega^1_{C/A}$ est un $C$-module libre et que les $t_i$ sont les images d'éléments $s_i$ ($1\leq i\leq r$) de $B$ tels que les $d_{B/A}(s_i)$ engendrent le $B$-module $\Omega^1_{B/A}$ et que leurs images dans $\Omega^1_{C/A}$ fassent partie d'une base de ce $C$-module (Bourbaki, \emph{Alg. comm.}, chap. II, § 5, n\textsuperscript{o} 1, prop. 2). Soit $\varphi$ le $A$-homomorphisme de $B'=A[T_1,\ldots,T_r]$ dans $B$ tel que $\varphi(T_i)=s_i$ pour tout $i$; le di-homomorphisme correspondant $\Omega^1_{B'/A}\to\Omega^1_{B/A}$ (0, 20.5.2) transforme les $d_{B'/A}(T_i)$, qui forment une base de $\Omega^1_{B'/A}$ (0, 20.4.13), en les $d_{B/A}(s_i)$ et est par suite \emph{surjectif}; si $Y'=\Spec(B')$, et si $u:Y\to Y'$ est le $S$-morphisme correspondant à $\varphi$, on conclut de (17.2.2) que $u$ est \emph{non ramifié}. Si l'on prouve que le morphisme composé $u\circ g:X\to Y'$ est lisse au point $x$, il résultera donc de (17.7.10) que $u$ est étale au point $y$, puis de (17.3.5) que $g$ est lisse au point $x$; enfin, comme le morphisme structural $f':Y'\to S$ est lisse (17.3.8), $f=f'\circ u$ sera lisse au point $y$. Comme les images canoniques des $d_{B'/A}(T_i)$ dans $\Omega^1_{C/A}$ sont celles des $d_{B/A}(s_i)$, elles font partie d'une base de $\Omega^1_{C/A}$. On voit ainsi que pour prouver que c) implique a), on peut se borner au cas où $Y'=Y$, et par suite supposer que $f$ est un morphisme \emph{lisse}.

En vertu de (17.8.2), on peut alors se borner au cas où $S=\Spec(k)$ est le spectre d'un corps; de plus, grâce à (17.7.1, (ii)), on peut supposer que $k$ est algébriquement clos; enfin, remplaçant au besoin $X$ et $Y$ par des voisinages ouverts de $x$ et $y$ respectivement, on peut supposer que $f$ et $h$ sont tous deux lisses, et que l'homomorphisme canonique
\[
g^*(\Omega^1_{Y/S})\to\Omega^1_{X/S}
\]
est inversible à gauche (0, 19.1.12). Alors l'ensemble des points de $X$ rationnels sur $k$ est très dense dans $X$ (10.4.8), donc, pour prouver que $g$ est lisse, il suffit de prouver que $g$ est lisse en tout point $x\in X$ rationnel sur $k$, ou encore qu'en un tel point, $\sh{O}_{X,x}$ est un $\sh{O}_{Y,y}$-module plat et $\sh{O}_{X,x}/\mathfrak m_y\sh{O}_{X,x}$ un anneau régulier (17.5.1). Or, $y=g(x)$ est

\oldpage[IV-4]{84}

\emph{a fortiori} rationnel sur $k$, et comme $\sh{O}_{Y,y}$ est alors une $k$-algèbre formellement lisse pour les topologies discrètes et que $\sh{O}_{Y,y}/\mathfrak m_y=k$, il résulte de (0, 20.5.14) que l'homomorphisme canonique $\mathfrak m_y/\mathfrak m_y^2\to(\Omega^1_{Y/S})_y\otimes_{\sh{O}_{Y,y}}k(y)$ est bijectif; de même, l'homomorphisme canonique $\mathfrak m_x/\mathfrak m_x^2\to(\Omega^1_{X/S})_x\otimes_{\sh{O}_{X,x}}k(x)$ est bijectif. L'hypothèse c'), équivalente à c), signifie donc ici que l'homomorphisme canonique
\[
(\mathfrak m_y/\mathfrak m_y^2)\otimes_{k(y)}k(x)\to\mathfrak m_x/\mathfrak m_x^2
\]
est injectif. Comme l'anneau $\sh{O}_{X,x}$ est régulier, la conclusion résulte de (0, 17.3.3). C.Q.F.D.

\begin{corollary}[17.11.2]
\label{IV.17.11.2}
Sous les hypothèses générales de (17.11.1), les conditions suivantes sont équivalentes:
\begin{enumerate}
  \item[a)] $f$ est lisse au point $y$ et $g$ est étale au point $x$.
  \item[b)] $h$ est lisse au point $x$ et $g$ est étale au point $x$.
  \item[c)] $h$ est lisse au point $x$, et l'homomorphisme canonique
  \[
  (g^*(\Omega^1_{Y/S}))_x\to(\Omega^1_{X/S})_x
  \]
  est bijectif.
  \item[c')] $h$ est lisse au point $x$, et l'homomorphisme canonique
  \[
  ((\Omega^1_{Y/S})_y\otimes_{\sh{O}_{Y,y}}k(y))\otimes_{k(y)}k(x)\to(\Omega^1_{X/S})_x\otimes_{\sh{O}_{X,x}}k(x)
  \]
  est bijectif.
\end{enumerate}
Supposons en outre que l'homomorphisme $k(y)\to k(x)$ soit bijectif. Alors les conditions précédentes sont aussi équivalentes à la suivante:
\begin{enumerate}
  \item[d)] $h$ est lisse au point $x$, et l'application canonique $\mathrm{T}_{X/S}(x)\to\mathrm{T}_{Y/S}(y)$ (16.5.12) est bijective.
\end{enumerate}
\end{corollary}

Chacune des conditions a), b), c) de (17.11.2) est équivalente à la conjonction de la condition correspondante de (17.11.1) et du fait que $g$ est non ramifié au point $x$, compte tenu de (17.2.2) en ce qui concerne la condition c); d'où l'équivalence de a), b) et c). L'équivalence de c) et c') résulte de l'équivalence des conditions correspondantes de (17.11.1) et du lemme de Nakayama; l'équivalence de c') et d) lorsque $k(x)=k(y)$ est immédiate par transposition (16.5.12).

\begin{corollary}[17.11.3]
\label{IV.17.11.3}
Soient $h:X\to S$ un morphisme lisse, $s_i$ ($1\leq i\leq r$) des sections de $\sh{O}_X$ au-dessus de $X$ (qui sont aussi des sections de $h_*(\sh{O}_X)$ au-dessus de $S$), $g:X\to S[T_1,\ldots,T_r]=\mathbf{V}_S^r$ le $S$-morphisme correspondant à l'homomorphisme de $\sh{O}_S$-Modules $\sh{O}_S^r\to h_*(\sh{O}_X)$ défini par ces sections (II, 1.2.7). Pour que $g$ soit lisse (resp. étale) en un point $x\in X$ il faut et il suffit que les $(d_{X/S}(s_i))_x$ fassent partie d'une base (resp. constituent une base) du $\sh{O}_{X,x}$-module $(\Omega^1_{X/S})_x$.
\end{corollary}

Il suffit d'appliquer (17.11.1) (resp. (17.11.2)) en prenant pour $f$ le morphisme structural $S[T_1,\ldots,T_r]\to S$.

\begin{corollary}[17.11.4]
\label{IV.17.11.4}
Pour qu'un morphisme $h:X\to S$ soit lisse en un point $x\in X$, il faut et il suffit qu'il existe un voisinage ouvert $U$ de $x$, un entier $r$, et un $S$-morphisme étale $U\to S[T_1,\ldots,T_r]$.
\end{corollary}

\oldpage[IV-4]{85}

La condition est évidemment suffisante puisque le morphisme structural $S[T_1,\ldots,T_r]\to S$ est lisse (17.3.8). Pour montrer qu'elle est nécessaire, notons que puisque $h$ est lisse au point $x$, il existe un voisinage ouvert $U$ de $x$ tel que $\Omega^1_{X/S}|U$ soit localement libre (17.2.3); il suffit alors d'utiliser (16.10.6) pour obtenir (en restreignant au besoin $U$) des sections $s_i$ de $\sh{O}_X$ au-dessus de $U$ telles que les $d_{X/S}(s_i)$ forment une base de $\Omega^1_{X/S}|U$; on conclut en appliquant (17.11.3).

\begin{proposition}[17.11.5]
\label{IV.17.11.5}
Soient $f:Y\to S$, $h:X\to S$ deux morphismes lisses. Pour qu'un $S$-morphisme $g:X\to Y$ soit une immersion ouverte, il faut et il suffit que $g$ soit un monomorphisme de préschémas et que pour tout $x\in X$, si l'on pose $y=g(x)$, on ait $\dim_y(f)=\dim_x(h)$ (pour une généralisation de cette proposition, voir (18.10.5)).
\end{proposition}

La condition est évidemment nécessaire; prouvons qu'elle est suffisante.

En vertu de (17.9.5), on est aussitôt ramené au cas où $S=\Spec(k)$ est le spectre d'un corps, et en vertu de (2.7.1, (x)), on peut supposer $k$ algébriquement clos. Compte tenu de (17.9.1), il suffit alors de prouver que $g$ est étale, et comme l'ensemble des points où $g$ est étale est ouvert, il suffit de montrer que $g$ est étale aux points \emph{fermés} (ou encore, rationnels sur $k$) de $X$ (10.4.8). Soit donc $x$ un tel point, et posons $y=g(x)$, qui est aussi rationnel sur $k$; par hypothèse, l'anneau $A=\sh{O}_{Y,y}$ est régulier et de corps résiduel $k$; soient $d=\dim(A)$ et $(t_i)_{1\leq i\leq d}$ un système régulier de paramètres pour $A$. Posons $B=\sh{O}_{X,x}$, $C=B/\mathfrak m_yB$. Comme $g$ est un monomorphisme, il en est de même du morphisme $\Spec(C)\to\Spec(k(y))=\Spec(k)$ déduit de $g$ par changement de base (I, 3.3.12); mais cela signifie que l'homomorphisme correspondant $u:k\to C$ est \emph{surjectif} (donc bijectif), car $u$ admet un inverse à gauche $v:C\to k$, et $u\circ v$ et l'identité de $C$, composés avec $u$, donnent le même morphisme $u:k\to C$. Comme par hypothèse $B$ est un anneau régulier de dimension $d$, les images des $t_i$ dans $B$ forment un système régulier de paramètres pour $B$ (0, 17.1.7); la condition (17.6.3, $e''$) est donc vérifiée par $g$ au point $x$ (0, 17.1.1 et Bourbaki, \emph{Alg. comm.}, chap. III, § 2, n\textsuperscript{o} 8, cor. 3 du th. 1), ce qui achève la démonstration.

\subsection{Sous-préschémas lisses d'un préschéma lisse. Morphismes lisses et morphismes différentiellement lisses}
\EGAReaderTarget{ega4.c04.subsec.17-12}{subsection:IV.17.12}

\begin{theorem}[17.12.1]
\label{IV.17.12.1}
Soient $f:X\to S$, $h:Y\to S$ deux morphismes localement de présentation finie, $j:Y\to X$ une immersion, $y$ un point de $Y$, $x=j(y)$. Les conditions suivantes sont équivalentes:
\begin{enumerate}
  \item[a)] $h$ est lisse au point $y$ et $f$ est lisse au point $x$.
  \item[b)] $f$ est lisse au point $x$, et l'homomorphisme canonique (16.4.21)
  \[
  (\sh{N}_{Y/X})_y\to(j^*(\Omega^1_{X/S}))_y
  \]
  est inversible à gauche.
  \item[c)] $h$ est lisse au point $y$ et il existe un voisinage ouvert $U$ de $y$ dans $Y$ tel que $j|U:U\to X$ soit une immersion régulière (16.9.2).

\oldpage[IV-4]{86}

  \item[c')] $h$ est lisse au point $y$ et il existe un voisinage ouvert $U$ de $y$ dans $Y$ tel que $j|U:U\to X$ soit une immersion quasi-régulière (autrement dit (16.9.8), il existe un voisinage $U$ de $y$ dans $Y$ dans lequel $\sh{N}_{Y/X}$ est un $\sh{O}_Y$-Module localement libre et l'homomorphisme canonique
  \[
  \mathbf{S}_{\sh{O}_Y}(\sh{N}_{Y/X})\to\operatorname{gr}^{\bullet}(j)
  \]
  est bijectif).
\end{enumerate}
\end{theorem}

Pour démontrer l'équivalence de a) et b), on peut se borner au cas où $S=\Spec(A)$ et $X=\Spec(B)$ sont affines, avec $Y=\Spec(B/\mathfrak J)$, où $\mathfrak J$ est un idéal de type fini de $B$, et $B$ est une $A$-algèbre lisse (17.3.2, (ii)). Il suffit alors d'appliquer le critère jacobien (0, 22.6.1) ainsi que (0, 19.1.14), compte tenu de ce que $\Omega^1_{X/S}$ est un $\sh{O}_X$-Module localement libre dans un voisinage de $x$ et $\sh{N}_{Y/X}$ un $\sh{O}_Y$-Module de type fini (16.1.6).

En second lieu, prouvons que a) entraîne c'). On peut encore se borner au cas où $S$, $X$, $Y$ sont affines, $B$ et $B/\mathfrak J$ des $A$-algèbres lisses, et il résulte alors de (17.7.9) et (8.10.5, (iv)) qu'il existe un sous-anneau noethérien $A'$ de $A$, une $A'$-algèbre de type fini $B'$ et un idéal $\mathfrak J'$ de $B'$ tels que $B=B'\otimes_{A'}A$, $\mathfrak J=\mathfrak J'B$ et que $B'$ et $B'/\mathfrak J'$ soient des $A'$-algèbres lisses. Notons que, d'après (0, 19.3.8), $B'$ est encore une $A'$-algèbre formellement lisse lorsqu'on munit $A'$ de la topologie discrète et $B'$ de la topologie $\mathfrak J'$-préadique. Par suite (0, 19.5.4), $\mathfrak J'/\mathfrak J'^2$ est un $(B'/\mathfrak J')$-module projectif et l'homomorphisme canonique
\[
\mathbf{S}_{B'/\mathfrak J'}(\mathfrak J'/\mathfrak J'^2)\to\operatorname{gr}_{\mathfrak J'}^{\bullet}(B')
\]
est bijectif. En d'autres termes (16.9.8), l'immersion $\Spec(B'/\mathfrak J')\to\Spec(B')$ est quasi-régulière, donc régulière puisque $B'$ est noethérien (16.9.10). Mais comme par hypothèse $B'/\mathfrak J'$ est un $A'$-module plat (17.5.1) et $B'$ une $A'$-algèbre de présentation finie, on peut appliquer (11.3.8) en remplaçant $\sh{F}$ par $\widetilde{\mathfrak J'}$, et on voit donc (par changement de base) que $j$ est une immersion régulière, et \emph{a fortiori} quasi-régulière. Le fait que $B$ est une $A$-algèbre de présentation finie et un $A$-module plat montre d'ailleurs, en vertu de (11.3.8), que les conditions c') et c) sont équivalentes, et sont stables par changement de base.

Montrons enfin que c) entraîne a). Du fait que, dans (11.3.8), la condition b) entraîne c), on voit déjà que si c) est vérifiée, $f$ est \emph{plat} au point $x$; en vertu de (17.5.1), tout revient donc à voir que sous l'hypothèse c), $f^{-1}(s)$ est lisse sur $k(s)$ au point $x$, en posant $s=h(x)$; comme on a remarqué que la condition c) est stable par changement de base, on voit qu'on est ramené au cas où $S=\Spec(k)$ est le spectre d'un corps, et compte tenu de (17.7.1, (ii)) on peut supposer que $k$ est \emph{algébriquement clos}. Il suffit alors de prouver que $\sh{O}_{X,x}$ est un anneau \emph{régulier} (17.5.1). Or, par hypothèse on a $\sh{O}_{Y,y}=\sh{O}_{X,x}/\mathfrak J_x$, où $\mathfrak J_x$ est un idéal engendré par une suite $\sh{O}_{X,x}$-régulière $(t_i)$, et $\sh{O}_{Y,y}$ est un anneau régulier; comme les $t_i$ font partie d'un système de paramètres pour $\sh{O}_{X,x}$ (0, 16.4.1), la conclusion résulte de (0, 17.1.7).

\begin{corollary}[17.12.2]
\label{IV.17.12.2}
Avec les notations de (17.12.1), supposons que $f$ soit lisse au point $x$, et soit $Y$ un sous-préschéma fermé de $X$ défini par un Idéal $\sh{J}$ de $\sh{O}_X$, et $S$-lisse au point $x$. Soient $(g_i)_{1\leq i\leq r}$ des sections de $\sh{J}$ au-dessus de $X$. Les conditions suivantes sont équivalentes:
\begin{enumerate}
  \item[a)] Les $(g_i)_x$ forment un système de générateurs du $\sh{O}_{X,x}$-module $\sh{J}_x$ dont le nombre d'éléments est le plus petit possible.
  \item[b)] Si $g_i'$ est l'image canonique de $g_i$ dans $\Gamma(X,\sh{J}/\sh{J}^2)$, les $(g_i')_x$ forment une base du $(\sh{O}_{X,x}/\sh{J}_x)$-module $\sh{J}_x/\sh{J}_x^2$.

\oldpage[IV-4]{87}

  \item[c)] Les images des $d_{X/S}g_i$ dans $(\Omega^1_{X/S})_x\otimes_{\sh{O}_{X,x}}k(x)$ sont des éléments linéairement indépendants sur $k(x)$, et les $(g_i)_x$ engendrent $\sh{J}_x$.
  \item[d)] Il existe un voisinage ouvert $U$ de $x$ dans $X$ et des sections $g_j$ ($r+1\leq j\leq n$) de $\sh{O}_X$ au-dessus de $U$ telles que le $S$-morphisme $u:U\to S[T_1,\ldots,T_n]$ correspondant à l'homomorphisme $\sh{O}_S^n\to f_*(\sh{O}_U)$ défini par les sections $g_i|U$ ($1\leq i\leq r$) et $g_j$ ($r+1\leq j\leq n$) (II, 1.2.7) soit un morphisme étale, et que l'image réciproque du sous-préschéma $Y'=S[T_{r+1},\ldots,T_n]$ par ce morphisme soit le préschéma induit par $j(Y)$ sur l'ouvert $U\cap j(Y)$.
\end{enumerate}
\end{corollary}

Comme $(g_i')_x$ est l'image canonique de $(g_i)_x$, l'équivalence de a) et b) résulte du lemme de Nakayama, $\sh{J}_x$ étant de type fini et $\sh{J}_x/\sh{J}_x^2$ un $(\sh{O}_{X,x}/\sh{J}_x)$-module libre (17.10.4) (Bourbaki, \emph{Alg. comm.}, chap. II, § 3, n\textsuperscript{o} 2, prop. 5). En vertu de (17.12.1, b)), $\sh{J}_x/\sh{J}_x^2$ s'identifie canoniquement à un facteur direct du $(\sh{O}_{X,x}/\sh{J}_x)$-module libre de rang $n$, $(\Omega^1_{X/S})_x\otimes_{\sh{O}_{X,x}}(\sh{O}_{X,x}/\sh{J}_x)$, et l'équivalence de b) et c) résulte de Bourbaki, \emph{loc. cit.} En outre, si a) est vérifiée, $(g_i')_x$ est ainsi identifié à $(d_{X/S}g_i)_x\otimes1$ pour $1\leq i\leq r$; le $\sh{O}_{X,x}$-module $(\Omega^1_{X/S})_x$ étant libre de rang $n$, il résulte encore de Bourbaki, \emph{loc. cit.}, qu'en remplaçant au besoin $X$ par un voisinage ouvert de $x$, on peut supposer qu'il existe $n-r$ sections $g_i$ ($r+1\leq i\leq n$) de $\sh{O}_X$ au-dessus de $X$, telles que les $(d_{X/S}g_i)_x$ pour $1\leq i\leq n$ forment une base de $(\Omega^1_{X/S})_x$; le fait que le morphisme correspondant $u$ soit étale au point $x$ résulte alors de (17.11.3); pour la même raison, le morphisme $u^{-1}(Y')\to Y'$, restriction de $u$, est étale au point $x$; en remplaçant $X$ par un voisinage de $x$, on peut supposer ces deux morphismes étales. En outre, il est immédiat que $Y$ (identifié, pour simplifier, au sous-préschéma fermé $j(Y)$ de $X$) est un sous-préschéma fermé de $u^{-1}(Y')$, et en vertu du choix des $g_i$ pour $i>r$ et de (17.2.5) et (17.11.2), la restriction à $Y$ de $u$ peut encore être supposée étale. On en déduit que pour tout $y'\in Y'$, l'immersion $Y\cap u^{-1}(y')\to u^{-1}(y')$ est ouverte, d'où l'on conclut à l'aide de (17.9.6) que $Y\to u^{-1}(Y')$ est une immersion ouverte; cela prouve que a) entraîne d). La réciproque découle aussitôt de (17.11.3).

\begin{corollary}[17.12.3]
\label{IV.17.12.3}
Soient $f:X\to S$ un morphisme localement de présentation finie, $u:S\to X$ une $S$-section de $X$. Pour que $f$ soit lisse en un point $x\in X$, il faut et il suffit que $u$ soit une immersion quasi-régulière dans un voisinage de $s=f(x)$.
\end{corollary}

La conclusion résulte de (17.12.1), puisque $f\circ u=1_S$ est lisse.

\begin{proposition}[17.12.4]
\label{IV.17.12.4}
Tout morphisme lisse $f:X\to S$ est différentiellement lisse (autrement dit (16.10.5) $\Delta_f:X\to X\times_SX$ est une immersion quasi-régulière). \emph{En particulier, les $\sh{P}_f^n=\sh{P}_{X/S}^n$ et les $\operatorname{gr}_n(\sh{P}_{X/S})$ sont des $\sh{O}_X$-Modules localement libres de type fini.}
\end{proposition}

Il suffit de remarquer que le morphisme structural $p_1:X\times_SX\to X$ est lisse (17.3.3, (iii)) et que $\Delta_f$ est une $X$-section pour $p_1$; la conclusion résulte de (17.12.3).

\begin{proposition}[17.12.5]
\label{IV.17.12.5}
Soit $f:X\to Y$ un morphisme localement de présentation finie. Les conditions suivantes sont équivalentes:
\begin{enumerate}
  \item[a)] $f$ est différentiellement lisse.
  \item[b)] Pour tout morphisme $g:Y'\to Y$, si l'on pose $X'=X\times_Y Y'$ et $f'=f_{(Y')}:X'\to Y'$, alors, pour toute $Y'$-section $s'$ de $X'$, $f'$ est lisse en tous les points de $s'(Y')$.
  \item[c)] La seconde projection $p_2:X\times_YX\to X$ est lisse en tous les points de la diagonale $\Delta_f(X)$.
\end{enumerate}
\end{proposition}

\oldpage[IV-4]{88}

La condition c) est un cas particulier de b): il suffit en effet de prendre $Y'=X$ et $g=f$ dans b), car alors $f'$ n'est autre que la seconde projection $p_2$, et $\Delta_f$ est une $X$-section de $X\times_YX$. En second lieu, c) entraîne a), car $p_2$ est un morphisme localement de présentation finie, et si $p_2$ est lisse aux points de $\Delta_f(X)$, $\Delta_f$ est une immersion quasi-régulière (17.12.3), donc $f$ est différentiellement lisse. Enfin, montrons que a) entraîne b); si $\Delta_f$ est une immersion quasi-régulière, $p_2$ est lisse en tous les points de $\Delta_f(X)$ (17.12.3).

Or, si $g':X'\to X$ est la projection canonique, et $v=(g',g')_Y:X'\to X\times_YX$, on vérifie aussitôt que le diagramme
\[
\xymatrix{
  X\times_Y X \ar[d]_{p_2} & X' \ar[l]_-v \ar[d]^{f'}\\
  X & Y' \ar[l]^{h=g'\circ s'}
}
\]
est commutatif et identifie $X'$ au produit $(X\times_YX)\times_XY'$ (I, 3.3.9); tenant compte de ce que le diagramme (17.4.1.1) identifie $Y'$ au produit des $(X\times_YX)$-préschémas $X$ et $X'$, on conclut, de ce que $p_2$ est lisse en tout point de $\Delta_f(X)$, que $f'$ est lisse en tout point de $s'(Y')$ (17.3.3, (iii)).

\begin{corollary}[17.12.6]
\label{IV.17.12.6}
Les notations étant celles de (8.8.1), on suppose $X_\alpha$ quasi-compact, et $f_\alpha:X_\alpha\to S_\alpha$ localement de présentation finie. Pour que $f:X\to S$ soit différentiellement lisse, il faut et il suffit qu'il existe $\lambda\geq\alpha$ tel que $f_\lambda:X_\lambda\to S_\lambda$ soit différentiellement lisse (auquel cas $f_\mu:X_\mu\to S_\mu$ est différentiellement lisse pour $\mu\geq\lambda$).
\end{corollary}

La suffisance de la condition et la dernière assertion résultent de (16.10.4). Pour prouver que la condition est nécessaire, notons que $\Delta_f:X\to X\times_SX$ et $p_2:X\times_SX\to X$ se déduisent par changement de base de $\Delta_{f_\lambda}:X_\lambda\to X_\lambda\times_{S_\lambda}X_\lambda$ et $p_{2,\lambda}:X_\lambda\times_{S_\lambda}X_\lambda\to X_\lambda$. En vertu de (17.12.5), pour tout $z\in\Delta_f(X)$, il existe un voisinage ouvert affine $U(z)$ de $z$ dans $X\times_SX$ tel que $p_2$ soit lisse dans $U(z)$; en vertu de (8.2.11) et (17.7.8), il existe un indice $\lambda(z)$ et un voisinage $U_{\lambda(z)}$ de la projection de $z$ dans $X_{\lambda(z)}\times_{S_{\lambda(z)}}X_{\lambda(z)}$ tels que $U(z)$ soit l'image réciproque de $U_{\lambda(z)}$, et que $p_{2,\lambda(z)}$ soit lisse dans $U_{\lambda(z)}$. Comme $X$ est quasi-compact, on peut recouvrir $\Delta_f(X)$ par un nombre fini de voisinages $U(z_i)$; en vertu de (8.3.4), il existe un $\lambda$ supérieur à tous les $\lambda(z_i)$ tels que les images réciproques des $U_{\lambda(z_i)}$ dans $X_\lambda\times_{S_\lambda}X_\lambda$ forment un recouvrement de $\Delta_{f_\lambda}(X_\lambda)$; il résulte alors de (17.12.5) que cet indice $\lambda$ répond à la question.

\begin{examples}[17.12.7]
\label{IV.17.12.7}
Soient $S$ un préschéma, $G$ un $S$-\emph{préschéma en groupes}, localement de présentation finie sur $S$. Alors, pour que $G$ soit différentiellement lisse, il faut et il suffit que $G$ soit lisse sur $S$ aux points de la « section unité » $e$ de $G$. La condition est en effet nécessaire par (17.12.5). Inversement, supposons-la remplie; pour tout morphisme $S'\to S$, $G'=G\times_SS'$ est alors un $S'$-préschéma en groupes localement de présentation finie sur $S'$; on peut donc, en vertu de (17.12.5), se borner à prouver que pour toute $S$-section $s$ de $G$, $f$ est lisse aux points de $s(S)$. Or, si $m:G\times_SG\to G$ est le morphisme qui définit la structure de préschéma en groupes de $G$, $m\circ(s\times1_G):S\times_SG\to G\times_SG\to G$ est un

\oldpage[IV-4]{89}

$S$-isomorphisme du préschéma $G$, la « translation à gauche » $\tau_s$, transformant les points de la section unité de $G$ en les points de $s(S)$; on en conclut que $G$ est lisse sur $S$ aux points de $s(S)$.
\end{examples}

\begin{remark}[17.12.8]
\label{IV.17.12.8}
Un $S$-préschéma en groupes $G$, de type fini (et même fini) sur un préschéma localement noethérien $S$, peut être différentiellement lisse sur $S$ sans être lisse (ni même \emph{plat}) sur $S$. Prenons par exemple pour $S$ le spectre de l'algèbre des nombres duaux $D=k[T]/T^2k[T]$ sur un corps $k$, pour $G$ le spectre de la $D$-algèbre composée directe $E=D\oplus k$. On définit sur $G$ une structure de $S$-schéma en groupes en définissant une « application diagonale » $\delta$, homomorphisme de $E$ dans $E\otimes_DE$: si $e'$, $e''$ sont les idempotents, images canoniques de l'unité $1$ de $D$ dans les facteurs $D$ et $k$ de $E$, on vérifie sans peine qu'on obtient une telle application diagonale en prenant
\[
\delta(e')=e'\otimes e'+e''\otimes e'',\qquad
\delta(e'')=e'\otimes e''+e''\otimes e'.
\]
La section unité du schéma en groupes $G$ correspond à l'homomorphisme $E\to D$ qui est l'identité dans $D$ et $0$ dans $k$; c'est un isomorphisme de $S$ sur une composante connexe de $G$, et \emph{a fortiori} $G$ est différentiellement lisse sur $S$ (17.12.6), mais il est clair que $G$ n'est pas $S$-plat.
\end{remark}

\subsection{Morphismes transversaux.}
\EGAReaderTarget{ega4.c04.subsec.17-13}{subsection:IV.17.13}
\label{IV.17.13}

\begin{env}[17.13.1]
\label{IV.17.13.1}
Soient $S$ un préschéma, $X$, $Y$, $X'$ trois $S$-préschémas, $i:Y\to X$ une $S$-immersion, $f:X'\to X$ un $S$-morphisme. Posons $Y'=Y\times_XX'$, et soient $g:Y'\to Y$, $j:Y'\to X'$ les projections canoniques, de sorte qu'on a le diagramme commutatif
\[
\label{IV.17.13.1.1}
\tag{17.13.1.1}
\xymatrix{
  Y \ar[d]_i & Y' \ar[l]_g \ar[d]^j\\
  X & X' \ar[l]^f
}
\]
et que $j$ est une $S$-immersion. On a alors un diagramme commutatif de $\sh{O}_{Y'}$-Modules quasi-cohérents
\[
\label{IV.17.13.1.2}
\tag{17.13.1.2}
\begin{tikzcd}
  g^*(\sh{N}_{Y/X}) \ar{r} \ar{d}{\operatorname{gr}_1(g)} & f^*(\Omega^1_{X/S})\otimes_{\sh{O}_{X'}}\sh{O}_{Y'} \ar{r} \ar{d} & g^*(\Omega^1_{Y/S}) \ar{r} \ar{d} & 0\\
  \sh{N}_{Y'/X'} \ar{r} & \Omega^1_{X'/S}\otimes_{\sh{O}_{X'}}\sh{O}_{Y'} \ar{r} & \Omega^1_{Y'/S} \ar{r} & 0
\end{tikzcd}
\]
où la ligne inférieure est la suite exacte (16.4.21) appliquée à $j$, la ligne supérieure provient de la même suite exacte pour $i$, par application du foncteur exact à droite $g^*$ (qui la laisse donc exacte); $\operatorname{gr}_1(g)$ est défini en (16.2.1), et la commutativité des deux carrés résulte de (0, 20.5.7.3) et (0, 20.5.11.3).

\oldpage[IV-4]{90}

On a vu en outre (16.2.2, (iii)) que $\operatorname{gr}_1(g)$ est ici surjectif, donc on déduit de (17.13.1.2) la suite exacte
\[
\label{IV.17.13.1.3}
\tag{17.13.1.3}
g^*(\sh{N}_{Y/X})\xrightarrow{\alpha}\Omega^1_{X'/S}\otimes_{\sh{O}_{X'}}\sh{O}_{Y'}\to\Omega^1_{Y'/S}\to0.
\]
\end{env}

\begin{proposition}[17.13.2]
\label{IV.17.13.2}
Les notations étant celles de (17.13.1), soit $x'$ un point de $Y'$, $x=g(x')$ son image dans $Y$; supposons $X$ et $Y$ lisses sur $S$ au point $x$, $X'$ lisse sur $S$ au point $x'$. Soient $m$, $m-c$ les dimensions relatives de $X$ et $Y$ sur $S$ au point $x$ (17.10.1), $n$ celle de $X'$ sur $S$ au point $x'$. Les conditions suivantes sont équivalentes:
\begin{enumerate}
  \item[a)] $Y'$ est lisse sur $S$ au point $x'$ et de dimension relative $n-c$.
  \item[b)] L'homomorphisme $\alpha\otimes1:g^*(\sh{N}_{Y/X})\otimes_{\sh{O}_{Y'}}k(x')\to\Omega^1_{X'/S}\otimes_{\sh{O}_{X'}}k(x')$ déduit de l'homomorphisme $\alpha$ de (17.13.1.3) est injectif.
\end{enumerate}
Soit $s$ l'image canonique de $x$ dans $S$, et supposons que $x'$ soit rationnel sur $k(s)$. Alors les conditions a) et b) sont aussi équivalentes à la suivante:
\begin{enumerate}
  \item[b')] L'homomorphisme transposé de $\alpha\otimes1$
  \[
  \mathrm{T}_{X'/S}(x')\to(\sh{N}_{Y/X}\otimes_{\sh{O}_Y}k(x'))^\vee=\mathrm{T}_{X/S}(x)/\mathrm{T}_{Y/S}(x)
  \]
  (cf. (16.5.12)) est surjectif.
\end{enumerate}
De plus, lorsque les conditions a) et b) sont vérifiées en $x'$, elles le sont dans un voisinage de $x'$ dans $Y'$; en remplaçant au besoin $X'$ par un voisinage de $x'$, l'homomorphisme
\[
\operatorname{gr}_1(g):g^*(\sh{N}_{Y/X})\to\sh{N}_{Y'/X'}
\]
est bijectif, et la suite
\[
\label{IV.17.13.2.1}
\tag{17.13.2.1}
0\to g^*(\sh{N}_{Y/X})\xrightarrow{\alpha}\Omega^1_{X'/S}\otimes_{\sh{O}_{X'}}\sh{O}_{Y'}\to\Omega^1_{Y'/S}\to0
\]
obtenue en ajoutant un $0$ à (17.13.1.3), est exacte.
\end{proposition}

Le fait que si les conditions équivalentes a) et b) sont vérifiées au point $x'$, elles le sont aussi au voisinage de $x'$ dans $Y'$, résulte de ce que l'ensemble des points où un morphisme est lisse est ouvert (17.3.7) et de (17.10.2).

En vertu de (17.12.1) appliqué à $X'$ et $Y'$, dire que $Y'$ est lisse sur $S$ au point $x'$ équivaut à dire que l'homomorphisme
\[
\delta\otimes1:\sh{N}_{Y'/X'}\otimes_{\sh{O}_{Y'}}k(x')\to\Omega^1_{X'/S}\otimes_{\sh{O}_{X'}}k(x')
\]
est injectif, compte tenu de (0, 19.1.12) et de ce que $\Omega^1_{X'/S}$ est un $\sh{O}_{X'}$-Module localement libre au point $x'$ (17.2.3); comme alors $\Omega^1_{Y'/S}$ est aussi un $\sh{O}_{Y'}$-Module localement libre au voisinage de $x'$ et que la suite
\[
\label{IV.17.13.2.2}
\tag{17.13.2.2}
0\to\sh{N}_{Y'/X'}\to\Omega^1_{X'/S}\otimes_{\sh{O}_{X'}}\sh{O}_{Y'}\to\Omega^1_{Y'/S}\to0
\]
est exacte (17.2.5), dire que $Y'$ est de dimension relative $n-c$ sur $S$ au point $x'$ signifie, par (17.10.2), que $\sh{N}_{Y'/X'}$ (qui est localement libre au voisinage de $x'$) est de rang $c$ au point $x'$. Mais en vertu des hypothèses sur $X$ et $Y$ au point $x$, et par le même raisonnement, $\sh{N}_{Y/X}$ est localement libre et de rang $c$ au point $x$, donc $g^*(\sh{N}_{Y/X})$ est aussi localement libre et de rang $c$ au point $x'$. Comme l'homomorphisme $\operatorname{gr}_1(g):g^*(\sh{N}_{Y/X})\to\sh{N}_{Y'/X'}$ est surjectif, les conditions précédentes équivalent à dire que cet homomorphisme est

\oldpage[IV-4]{91}

bijectif au point $x'$ (Bourbaki, \emph{Alg. comm.}, chap. II, § 3, n\textsuperscript{o} 2, cor. de la prop. 6), donc aussi dans un voisinage de $x'$ (0$_{\mathrm I}$, 5.2.7); cela entraîne évidemment b), ainsi que la dernière assertion de l'énoncé, en vertu de l'exactitude de (17.13.2.2). Inversement, comme $\alpha\otimes1$ se factorise en
\[
g^*(\sh{N}_{Y/X})\otimes k(x')\xrightarrow{\operatorname{gr}_1(g)\otimes1}\sh{N}_{Y'/X'}\otimes k(x')\xrightarrow{\delta\otimes1}\Omega^1_{X'/S}\otimes k(x')
\]
et que $\operatorname{gr}_1(g)$ est surjectif, dire que $\alpha\otimes1$ est injectif entraîne que $\delta\otimes1$ l'est et que $\operatorname{gr}_1(g)\otimes1$ est bijectif. On en conclut (17.12.1) que $Y'$ est lisse sur $S$ au point $x'$, et que $\operatorname{gr}_1(g)$ est bijectif dans un voisinage de $x'$ (Bourbaki, \emph{loc. cit.}). En outre, comme la suite (17.13.2.2) est alors exacte, et que $\sh{N}_{Y'/X'}$, isomorphe à $g^*(\sh{N}_{Y/X})$, est de rang $c$ au point $x'$, $\Omega^1_{Y'/S}$ est de rang $n-c$ en ce point, ce qui achève de prouver l'équivalence de a) et b), en vertu de (17.10.2).

Reste à montrer l'équivalence de b) et b') lorsque $x'$ est rationnel sur $k(s)$ (ce qui implique qu'il en est de même de $x$); alors $g^*(\sh{N}_{Y/X})\otimes_{\sh{O}_{Y'}}k(x')$ s'identifie à $\sh{N}_{Y/X}\otimes_{\sh{O}_Y}k(x')$, et puisque la suite
\[
0\to\sh{N}_{Y/X}\to\Omega^1_{X/S}\otimes_{\sh{O}_X}\sh{O}_Y\to\Omega^1_{Y/S}\to0
\]
est exacte au point $x$ et formée de $\sh{O}_Y$-Modules localement libres, le dual de $\sh{N}_{Y/X}\otimes_{\sh{O}_Y}k(x)$ s'identifie à l'espace quotient $\mathrm{T}_{X/S}(x)/\mathrm{T}_{Y/S}(x)$ (16.5.12); d'où l'équivalence de b) et b').

\begin{definition}[17.13.3]
\EGAReaderTarget{ega4.c04.def.17-13-3}{IV.17.13.3}
Avec les notations de (17.13.1), on dit que le morphisme $f$ est \emph{transversal} à $Y$ au point $x'$, relativement à $S$, si $X$ et $Y$ sont lisses sur $S$ au point $x$, $X'$ lisse sur $S$ au point $x'$, et si les conditions équivalentes a), b) de (17.13.2) sont satisfaites.
\end{definition}

Lorsque aucune confusion n'est à craindre, on supprime la mention du préschéma $S$ et on dit simplement que $f$ est transversal à $Y$ au point $x'$.

\begin{remarks}[17.13.4]
\label{IV.17.13.4}
(i) Supposons que $X$, $Y$ et $X'$ soient \emph{plats et localement de présentation finie} sur $S$. Pour tout $s\in S$, notons $X_s$, $Y_s$, $X'_s$, $Y'_s$ les fibres de $X$, $Y$, $X'$, $Y'$ au point $s$, $f_s:X'_s\to X_s$ le morphisme déduit de $f$ par changement de base. Il résulte alors de (17.5.9) que pour que $f$ soit un morphisme transversal à $Y$ au point $x'$, il faut et il suffit que, si $s$ est l'image de $x$ dans $S$, $f_s$ soit transversal à $Y_s$ au point $x'$, relativement à $k(s)$.

(ii) On a vu dans (17.13.2) que l'ensemble des points $x'\in Y'$ où $f$ est transversal à $Y$ est \emph{ouvert} dans $Y'$; si $Y'$ est de plus \emph{propre} sur $S$, on en déduit que l'ensemble des $s\in S$ tels que $f$ soit transversal à $Y$ en tous les points de $Y'_s$, relativement à $S$, est \emph{ouvert} dans $S$.

Lorsque $X$, $Y$ et $X'$ sont plats et localement de présentation finie sur $S$, il résulte de (i) que l'ensemble des $x'\in Y'$ tels que $f_s$ soit transversal à $Y_s$ au point $x'$ ($s$ image de $x'$ dans $S$), relativement à $k(s)$, est \emph{ouvert} dans $Y'$. Si de plus $Y'$ est propre sur $S$, l'ensemble des $s\in S$ tels que $f_s$ soit transversal à $Y_s$ en tous les points de $Y'_s$ (relativement à $k(s)$) est \emph{ouvert} dans $S$.

(iii) La propriété d'être lisse sur $S$, ainsi que la notion de dimension relative à $S$ en un point, étant stables par tout changement de base $S'\to S$ ((17.3.3) et (4.2.7)), il en est de même de la propriété pour un morphisme $f:X'\to X$ d'être transversal à un sous-préschéma de $X$ en un point, relativement à $S$.

\oldpage[IV-4]{92}

(iv) La condition b) de (17.13.2) exprime encore que l'homomorphisme $\alpha:g^*(\sh{N}_{Y/X})\to\Omega^1_{X'/S}\otimes_{\sh{O}_{X'}}\sh{O}_{Y'}$ est \emph{universellement injectif} relativement à $Y'$ (II, 9.18).
\end{remarks}

\begin{env}[17.13.5]
\label{IV.17.13.5}
Considérons maintenant un préschéma $S$, trois $S$-préschémas $X$, $Y$, $Z$, deux $S$-morphismes $f:Y\to X$, $g:Z\to X$; posons $T=Y\times_XZ$; on sait alors (I, 5.3.5) qu'on a un diagramme commutatif
\[
\label{IV.17.13.5.1}
\tag{17.13.5.1}
\xymatrix{
  X \ar[d]^{\Delta} & T=Y\times_X Z \ar[l] \ar[d]^{j}\\
  X\times_S X & Y\times_S Z \ar[l]^{u}
}
\]
faisant de $T$ le produit des $(X\times_SX)$-préschémas $X$ et $Y\times_SZ$, où $u=f\times_Sg$. Comme $\Delta$ est une $S$-immersion (I, 5.3.9), on est dans la situation du diagramme (17.13.1.1); ce qui correspond à $\Omega^1_{X'/S}$ dans (17.13.1) est alors $(\Omega^1_{Y/S}\otimes_{\sh{O}_Y}\sh{O}_{Y\times_SZ})\oplus(\Omega^1_{Z/S}\otimes_{\sh{O}_Z}\sh{O}_{Y\times_SZ})$ en vertu de (16.4.23). D'autre part, ce qui correspond au $\sh{O}_Y$-Module $\sh{N}_{Y/X}$ dans (17.13.1) est ici par définition $\Omega^1_{X/S}$ (16.3.1); il correspond donc à (17.13.1.3) une suite exacte
\[
\label{IV.17.13.5.2}
\tag{17.13.5.2}
\Omega^1_{X/S}\otimes_{\sh{O}_X}\sh{O}_T\xrightarrow{\rho}
(\Omega^1_{Y/S}\otimes_{\sh{O}_Y}\sh{O}_T)\oplus
(\Omega^1_{Z/S}\otimes_{\sh{O}_Z}\sh{O}_T)
\xrightarrow{\sigma}\Omega^1_{T/S}\to0
\]
où il reste à préciser les homomorphismes $\rho$ et $\sigma$. Tenant compte d'abord de (16.4.23) et de (0, 20.5.2), on voit que si $p:T\to Y$, $q:T\to Z$ sont les projections canoniques, on a, avec les notations de (16.4.19),
\[
\label{IV.17.13.5.3}
\tag{17.13.5.3}
\sigma=p_{T/Y/S}+q_{T/Z/S}.
\]
D'autre part, pour évaluer $\rho$, utilisons la commutativité du carré de gauche dans (17.13.1.2), ce qui, dans le cas présent, ramène d'abord à expliciter l'homomorphisme canonique
\[
\rho':\Omega^1_{X/S}\to\Omega^1_{(X\times_SX)/S}\otimes_{\sh{O}_{X\times_SX}}\sh{O}_X
\]
défini dans (16.4.21) appliqué à l'immersion $\Delta$. On peut se borner au cas où $S=\Spec(A)$, $X=\Spec(B)$ sont affines, et alors $\rho'$ correspond à l'homomorphisme $\delta$ de (0, 20.5.11.2), où il faut remplacer $B$ par $B\otimes_AB$ et $C$ par $B=(B\otimes_AB)/\mathfrak J_{B/A}$. On voit alors que $\delta$ fait correspondre à la classe de $x\otimes1-1\otimes x\pmod{\mathfrak J_{B/A}^2}$ (pour un $x\in B$) l'image de
\[
(x\otimes1-1\otimes x)\otimes(1\otimes1)-(1\otimes1)\otimes(x\otimes1-1\otimes x)
\]
dans $\Omega^1_{(B\otimes_AB)/A}\otimes_{B\otimes_AB}B$; mais l'élément précédent peut s'écrire
\[
\begin{split}
&((x\otimes1)\otimes(1\otimes1)-(1\otimes1)\otimes(x\otimes1))\\
&\quad-((1\otimes x)\otimes(1\otimes1)-(1\otimes1)\otimes(1\otimes x))
\end{split}
\]
et on voit donc que $\rho'$ est la différence des deux homomorphismes $\pi^{(1)}$ et $\pi^{(2)}$ de $\Omega^1_{X/S}$ dans $\Omega^1_{(X\times_SX)/S}\otimes_{\sh{O}_{X\times_SX}}\sh{O}_X$, correspondant respectivement à la première et à la seconde projection de $X\times_SX$ dans $X$, par (16.4.3.3). Pour obtenir $\rho$, il faut d'abord considérer l'homomorphisme
\[
\rho'':\Omega^1_{(X\times_SX)/S}\otimes_{\sh{O}_{X\times_SX}}\sh{O}_{Y\times_SZ}\to\Omega^1_{(Y\times_SZ)/S}
\]
correspondant par (16.4.3.3) au morphisme $u$ de (17.13.5.1), puis, après tensorisation par $\sh{O}_T$, former le composé $(\rho''\otimes1)\circ(\rho'\otimes1)$; il résulte de ce qui précède que l'on a, avec les notations de (16.4.18),
\[
\label{IV.17.13.5.4}
\tag{17.13.5.4}
\rho=(f_{Y/X/S}\otimes1_{\sh{O}_T})\oplus(-g_{Z/X/S}\otimes1_{\sh{O}_T}).
\]

\oldpage[IV-4]{93}

Ceci dit, l'application de (17.13.2) à la situation du diagramme (17.13.5.1) (compte tenu de (17.3.3, (iv)), qui implique que $X\times_SX$ est lisse sur $S$ en $x$ si $X$ l'est) donne le corollaire suivant.
\end{env}

\begin{corollary}[17.13.6]
\EGAReaderTarget{ega4.c04.cor.17-13-6}{IV.17.13.6}
Soient $S$ un préschéma, $X$, $Y$, $Z$ trois $S$-préschémas, $f:Y\to X$, $g:Z\to X$ deux $S$-morphismes; posons $T=Y\times_XZ$, et soient $p:T\to Y$, $q:T\to Z$ les projections canoniques. Soient $t$ un point de $T$, $y=p(t)$, $z=q(t)$, $x=f(y)=g(z)$. Supposons que $X$ soit lisse sur $S$ au point $x$, de dimension relative $m$, $Y$ (resp. $Z$) lisse sur $S$ au point $y$ (resp. $z$), de dimension relative $m+a$ (resp. $m+b$), $a$ et $b$ étant positifs ou négatifs. Alors les conditions suivantes sont équivalentes:
\begin{enumerate}
  \item[a)] $T$ est lisse sur $S$ au point $t$, de dimension relative $m+a+b$.
  \item[b)] L'homomorphisme
  \[
  \rho\otimes1:\Omega^1_{X/S}\otimes_{\sh{O}_X}k(t)\to
  (\Omega^1_{Y/S}\otimes_{\sh{O}_Y}k(t))\oplus
  (\Omega^1_{Z/S}\otimes_{\sh{O}_Z}k(t))
  \]
  où $\rho$ est donné par (17.13.5.4), est injectif.
  \item[c)] Le morphisme $T=Y\times_SZ\to X\times_SX$ est transversal à la diagonale de $X\times_SX$ au point $t$.
\end{enumerate}
Si $t$ est rationnel sur le corps résiduel de son image $s$ dans $S$, ces conditions sont aussi équivalentes à la suivante:
\begin{enumerate}
  \item[b')] L'homomorphisme
  \[
  \label{IV.17.13.6.1}
  \tag{17.13.6.1}
  \mathrm{T}_y(f)-\mathrm{T}_z(g):\mathrm{T}_{Y/S}(y)\oplus\mathrm{T}_{Z/S}(z)\to\mathrm{T}_{X/S}(x)
  \]
  (cf. (16.5.12.5)) est surjectif.
\end{enumerate}
De plus, lorsque les conditions a) et b) sont vérifiées en $t$, elles le sont dans un voisinage de $t$ dans $T$, et en restreignant $T$ à un tel voisinage, la suite
\[
\label{IV.17.13.6.2}
\tag{17.13.6.2}
0\to\Omega^1_{X/S}\otimes_{\sh{O}_X}\sh{O}_T\xrightarrow{\rho}
(\Omega^1_{Y/S}\otimes_{\sh{O}_Y}\sh{O}_T)\oplus
(\Omega^1_{Z/S}\otimes_{\sh{O}_Z}\sh{O}_T)
\xrightarrow{\sigma}\Omega^1_{T/S}\to0
\]
(où $\sigma$ est donné par (17.13.5.3)) est exacte.

Le seul point qui reste à prouver dans (17.13.6) concerne b'), et résulte de ce que, sous l'hypothèse que $t$ est rationnel sur $k(s)$, l'homomorphisme $\rho\otimes1$ est \emph{transposé} de $\mathrm{T}_y(f)-\mathrm{T}_z(g)$, en vertu de (17.13.5.4).

Lorsque les conditions équivalentes de (17.13.6) sont satisfaites, on dit que $f$ et $g$ forment un \emph{couple de morphismes transversaux au point $t$, relativement} à $S$; ici encore, on supprime souvent la mention de $S$.
\end{corollary}

\begin{env}[17.13.7]
\EGAReaderTarget{ega4.c04.statement.17-13-7}{IV.17.13.7}
Considérons en particulier le cas où $Y$ et $Z$ sont des \emph{sous-préschémas} de $X$, $f$ et $g$ étant les injections canoniques; $T$ est alors le sous-préschéma « intersection »

\oldpage[IV-4]{94}

$\inf(Y,Z)$ de $X$ (I, 4.4.3), et l'on a $t=y=z=x$. Au lieu de dire que le couple $(f,g)$ est transversal au point $x$, on dit alors que $Y$ et $Z$ \emph{se coupent transversalement au point $x$} (relativement à $S$). Notant $X_s$, $Y_s$, $Z_s$, $T_s$ les fibres de $X$, $Y$, $Z$, $T$ au point $s$, image de $x$ dans $S$, et tenant compte de (5.2.3), (5.1.9) et (0, 16.5.12), on voit que pour qu'il en soit ainsi (lorsque $X$, $Y$, $Z$ sont lisses sur $S$ au point $x$), il faut et il suffit que $T$ soit lisse sur $S$ au point $x$, et que l'on ait la relation
\[
\label{IV.17.13.7.1}
\tag{17.13.7.1}
\operatorname{codim}_x(T_s,X_s)=\operatorname{codim}_x(Y_s,X_s)+\operatorname{codim}_x(Z_s,X_s).
\]
\end{env}

\begin{proposition}[17.13.8]
\label{IV.17.13.8}
Soient $S$ un préschéma, $X$ un $S$-préschéma localement de présentation finie sur $S$, $Y$, $Z$ deux sous-préschémas de $X$, $T=\inf(Y,Z)$ le sous-préschéma « intersection » de $Y$ et $Z$, $x$ un point de $T$. Les conditions suivantes sont équivalentes:
\begin{enumerate}
  \item[a)] L'injection canonique $i:Y\to X$ est un morphisme transversal à $Z$ au point $x$, relativement à $S$ (17.13.3).
  \item[a')] L'injection canonique $j:Z\to X$ est un morphisme transversal à $Y$ au point $x$, relativement à $S$.
  \item[a'')] $Y$ et $Z$ se coupent transversalement au point $x$, relativement à $S$.
  \item[b)] $X$, $Y$, $Z$ sont lisses sur $S$ au point $x$, et l'homomorphisme $\rho$ (17.13.5.4) est tel que
  \[
  \rho\otimes1:\Omega^1_{X/S}\otimes_{\sh{O}_X}k(x)\to
  (\Omega^1_{Y/S}\otimes_{\sh{O}_Y}k(x))\oplus
  (\Omega^1_{Z/S}\otimes_{\sh{O}_Z}k(x))
  \]
  est injectif.
\end{enumerate}
Lorsque $x$ est rationnel sur $k(s)$ ($s$ image de $x$ dans $S$) ces conditions sont aussi équivalentes à la suivante:
\begin{enumerate}
  \item[b')] L'homomorphisme
  \[
  \mathrm{T}_x(i)-\mathrm{T}_x(j):\mathrm{T}_{Y/S}(x)\oplus\mathrm{T}_{Z/S}(x)\to\mathrm{T}_{X/S}(x)
  \]
  est surjectif.
\end{enumerate}
En outre, lorsque les conditions équivalentes a) à b) sont vérifiées au point $x$, elles le sont dans un voisinage de $x$ dans $T$, et en restreignant $X$ à un voisinage de $x$, la suite (17.13.6.2) est exacte, et l'on a un isomorphisme canonique
\[
\label{IV.17.13.8.1}
\tag{17.13.8.1}
\sh{N}_{T/X}\xrightarrow{\sim}
(\sh{N}_{Y/X}\otimes_{\sh{O}_Y}\sh{O}_T)\oplus
(\sh{N}_{Z/X}\otimes_{\sh{O}_Z}\sh{O}_T).
\]
\end{proposition}

Les conditions a), a'), a'') impliquent toutes que $X$, $Y$, $Z$ sont lisses sur $S$ au point $x$. Soient en outre, $m$, $m-a$, $m-b$ les dimensions relatives de $X$, $Y$, $Z$ sur $S$ au point $x$. Il résulte alors de (17.13.2) appliqué en remplaçant $X'$ par $Z$ et $Y'$ par $T=Y\times_XZ$, que les conditions a), a') équivalent toutes deux à dire que $T$ est lisse sur $S$ au point $x$ et de dimension relative $m-a-b$ en ce point; mais en vertu de (17.13.7), cela signifie précisément que la condition a'') est vérifiée, d'où l'équivalence de a), a') et a''). L'équivalence de a'') et de b) (ou b') lorsque $x$ est rationnel sur $k(s)$) a été prouvée dans (17.13.6), ainsi que le fait que si ces conditions sont satisfaites au point $x$, elles le sont dans un voisinage de $x$, et l'exactitude de la suite (17.13.6.2) dans un tel voisinage. Reste à définir l'isomorphisme canonique (17.13.8.1).

\oldpage[IV-4]{95}

Notons $\alpha$ et $-\beta$ les homomorphismes figurant au second membre de (17.13.5.4), $\gamma$ et $\delta$ ceux figurant au second membre de (17.13.5.2). Dire que la suite (17.13.6.2)
\[
0\to\Omega^1_{X/S}\otimes_{\sh{O}_X}\sh{O}_T
\xrightarrow{\alpha\oplus(-\beta)}
(\Omega^1_{Y/S}\otimes_{\sh{O}_Y}\sh{O}_T)\oplus
(\Omega^1_{Z/S}\otimes_{\sh{O}_Z}\sh{O}_T)
\xrightarrow{\gamma+\delta}\Omega^1_{T/S}\to0
\]
est exacte signifie que, dans la catégorie des $\sh{O}_T$-Modules, $\Omega^1_{X/S}\otimes_{\sh{O}_X}\sh{O}_T$ s'identifie canoniquement au \emph{produit fibré} de $\Omega^1_{Y/S}\otimes_{\sh{O}_Y}\sh{O}_T$ et $\Omega^1_{Z/S}\otimes_{\sh{O}_Z}\sh{O}_T$ sur $\Omega^1_{T/S}$, pour les homomorphismes $\gamma$ et $\delta$. Le même raisonnement que dans (0, 18.1.2 et 18.1.3), où l'on remplace anneaux et idéaux bilatères respectivement par Modules et sous-Modules, fournit un diagramme commutatif
{\scriptsize
\[
\begin{tikzcd}[column sep=small,row sep=large]
0 \ar[dr] & & 0 \ar[d] & 0 \ar[d] & \\
& (\sh{N}_{Y/X}\otimes\sh{O}_T)\oplus(\sh{N}_{Z/X}\otimes\sh{O}_T)
  \ar[r] \ar[d] \ar[dr,"\iota" description]
& \sh{N}_{Z/X}\otimes\sh{O}_T \ar[r,"\sim"] \ar[d]
& \sh{N}_{T/Y} \ar[d] & \\
0 \ar[r] & \sh{N}_{Y/X}\otimes\sh{O}_T \ar[r] \ar[d]
& \Omega^1_{X/S}\otimes\sh{O}_T \ar[r,"\alpha"] \ar[d,"\beta"] \ar[dr,"\varepsilon" description]
& \Omega^1_{Y/S}\otimes\sh{O}_T \ar[r] \ar[d,"\gamma"] & 0\\
0 \ar[r] & \sh{N}_{T/Z} \ar[r]
& \Omega^1_{Z/S}\otimes\sh{O}_T \ar[r,"\delta"] \ar[d]
& \Omega^1_{T/S} \ar[r] \ar[d] \ar[dr] & 0\\
& & 0 & 0 & 0
\end{tikzcd}
\]
}
où les 3\textsuperscript{e} et 4\textsuperscript{e} lignes et colonnes sont exactes en vertu des hypothèses de lissité. Le fait que l'homomorphisme composé $\varepsilon=\gamma\circ\alpha=\delta\circ\beta$ soit l'homomorphisme correspondant à l'injection canonique $T\to X$ résulte de (0, 20.5.2.7). Comme la diagonale du diagramme précédent est exacte, $\operatorname{Ker}(\varepsilon)$ s'identifie canoniquement avec $\sh{N}_{T/X}$, compte tenu de ce que $T$ est lisse sur $S$ (17.2.5); donc on a un isomorphisme canonique (17.13.8.1) réciproque de $\iota$.

\begin{env}[17.13.9]
\label{IV.17.13.9}
On peut généraliser les résultats de (17.13.5) et (17.13.6) à un nombre fini quelconque de $S$-morphismes $f_i:Y_i\to X$ ($i\in I$, $I$ ensemble fini). Pour cela, notons $X^I$ le produit de la famille $(X_i)_{i\in I}$ de $S$-préschémas tous identiques à $X$ (I, 3.3.5), et soient $p_i$ ($i\in I$) les projections canoniques. Le \emph{morphisme diagonal} $\Delta:X\to X^I$ est défini (comme dans (I, 5.3.1)) comme l'unique $S$-morphisme tel que $p_i\circ\Delta=1_X$ pour tout $i\in I$. C'est une $X$-section de $X^I$ pour le morphisme $p_1$, donc (I, 5.3.11) une \emph{immersion}.

\oldpage[IV-4]{96}

Soit alors $Y$ le produit des $S$-préschémas $Y_i$ (pour les morphismes composés $Y_i\xrightarrow{f_i}X\to S$), $T$ le produit des $X$-préschémas $Y_i$ (pour les morphismes $f_i$). On prouve comme dans (I, 5.3.5) que l'on a un diagramme commutatif
\[
\label{IV.17.13.9.1}
\tag{17.13.9.1}
\xymatrix{
  X \ar[d]^{\Delta} & T \ar[l]_-v \ar[d]^{j}\\
  X^I & Y \ar[l]^{u}
}
\]
(où $u$ est le produit (sur $S$) des $f_i$), qui fait de $T$ le produit des $X^I$-préschémas $X$ et $Y$.

Par récurrence sur le nombre d'éléments de $I$, il résulte de (16.4.23) que $\Omega^1_{X^I/S}$ (resp. $\Omega^1_{Y/S}$) s'identifie canoniquement à la somme directe des $p_i^*(\Omega^1_{X/S})$ (resp. des $q_i^*(\Omega^1_{Y_i/S})$, si $q_i:Y\to Y_i$ sont les projections canoniques).

Supposons maintenant que $X$ soit \emph{lisse} sur $S$ en un point $x$ (donc dans un voisinage de ce point). Il en est alors de même de $X^I$ en ce point (17.3.3, (iv)), et restreignant $X$ à un voisinage de $x$, on a donc (17.2.5) une suite exacte de $\sh{O}_X$-Modules localement libres
\[
0\to\operatorname{gr}_1(\Delta)\to(\Omega^1_{X/S})^I\to\Omega^1_{X/S}\to0.
\]
Posons $\sh{J}=\operatorname{gr}_1(\Delta)$; on déduit de la suite précédente une suite exacte de $\sh{O}_T$-Modules localement libres
\[
\label{IV.17.13.9.1-seq}
\tag{17.13.9.1}
0\to\sh{J}\otimes_{\sh{O}_X}\sh{O}_T\to(\Omega^1_{X/S})^I\otimes_{\sh{O}_X}\sh{O}_T\xrightarrow{\sigma}\Omega^1_{X/S}\otimes_{\sh{O}_X}\sh{O}_T\to0
\]
qui correspond à la première ligne du diagramme (17.13.1.2), et où $\sigma$ n'est autre que l'homomorphisme canonique qui, à chaque famille de sections $(t_i)_{i\in I}$ de $\Omega^1_{X/S}\otimes_{\sh{O}_X}\sh{O}_T$ au-dessus d'un ouvert de $T$, fait correspondre sa \emph{somme}.

D'autre part, ce qui correspond ici à la seconde flèche verticale du diagramme (17.13.1.2) est l'homomorphisme
\[
\label{IV.17.13.9.2}
\tag{17.13.9.2}
\tau:(\Omega^1_{X/S})^I\otimes_{\sh{O}_X}\sh{O}_T\to\bigoplus_{i\in I}(\Omega^1_{Y_i/S}\otimes_{\sh{O}_{Y_i}}\sh{O}_T)
\]
qui, à toute famille $(t'_i\otimes1)_{i\in I}$, où ici les $t'_i$ sont des sections au-dessus d'un ouvert de $X$ de $\Omega^1_{X/S}$, associe la somme des $((f_i)_{Y_i/X/S}(t'_i))\otimes1$, avec la notation de (16.4.18). L'homomorphisme $\rho$ qui correspond à l'homomorphisme $\alpha$ de (17.13.1.3) est donc la \emph{restriction} à $\sh{J}\otimes_{\sh{O}_X}\sh{O}_T$ de l'homomorphisme $\tau$ précédent.

Ces hypothèses et notations étant précisées, on peut appliquer à la situation de (17.13.9.1) la prop. (17.13.2), ce qui donne le
\end{env}

\begin{corollary}[17.13.10]
\EGAReaderTarget{ega4.c04.cor.17-13-10}{IV.17.13.10}
Sous les hypothèses générales de (17.13.9), soit $t$ un point de $T$, $y_i$ ($i\in I$) ses projections dans les $Y_i$, $x$ le point de $X$ égal à chacun des $f_i(y_i)$. Supposons que $X$ soit lisse sur $S$ au point $x$ et de dimension relative $d$, et que chacun des $Y_i$ soit lisse sur $S$ au point $y_i$; soit $d+c_i$ ($c_i$ entier positif ou négatif) la dimension relative de $Y_i$ au point $y_i$. Alors les conditions suivantes sont équivalentes:

\oldpage[IV-4]{97}

\begin{enumerate}
  \item[a)] $T$ est lisse sur $S$ au point $t$, de dimension relative $d+\sum_{i\in I}c_i$.
  \item[b)] L'homomorphisme
  \[
  \rho\otimes1:\sh{J}\otimes_{\sh{O}_X}k(t)\to\bigoplus_{i\in I}(\Omega^1_{Y_i/S}\otimes_{\sh{O}_{Y_i}}k(t))
  \]
  est injectif.
\end{enumerate}
Lorsque $t$ est rationnel sur le corps $k(s)$ (où $s$ est l'image de $t$ dans $S$), ces conditions sont aussi équivalentes à la suivante:
\begin{enumerate}
  \item[b')] L'homomorphisme transposé de $\rho\otimes1$
  \[
  \bigoplus_{i\in I}\mathrm{T}_{Y_i/X}(y_i)\to(\mathrm{T}_{X/S}(x))^I/\delta(\mathrm{T}_{X/S}(x))
  \]
  (où $\delta$ est l'application diagonale) est surjectif.
\end{enumerate}
Il suffit de remarquer que $X^I$ est lisse sur $S$ et de dimension relative $d\cdot\operatorname{Card}(I)$, et $Y$ lisse sur $S$ de dimension relative $d\cdot\operatorname{Card}(I)+\sum_{i\in I}c_i$.

Lorsque les conditions équivalentes de (17.13.10) sont vérifiées, on dit encore que les $f_i$ forment une \emph{famille de morphismes transversaux au point $t$, relativement} à $S$. On voit encore que l'ensemble des points $t\in T$ où cela a lieu est \emph{ouvert} dans $T$. La remarque (17.13.4, (iii)) montre alors que si $X$ et les $Y_i$ sont plats et localement de présentation finie sur $S$, et de plus si $T$ est \emph{propre} sur $S$, l'ensemble des $s\in S$ tels que les $f_i$ forment une famille de morphismes transversaux en \emph{tous} les points de $T_s$, relativement à $S$, est \emph{ouvert} dans $S$.
\end{corollary}

\begin{env}[17.13.11]
\EGAReaderTarget{ega4.c04.statement.17-13-11}{IV.17.13.11}
Considérons en particulier le cas, généralisant (17.13.7), où les $Y_i$ ($i\in I$) sont des \emph{sous-préschémas} de $X$, les $f_i$ étant les injections canoniques, de sorte que $T=\inf_{i\in I}(Y_i)$ est encore le sous-préschéma « intersection » des $Y_i$, et $t=y_i=x$ pour tout $i$; au lieu de dire que les $f_i$ forment une famille de morphismes transversaux au point $x$, on dit encore que \emph{les $Y_i$ se coupent transversalement au point $x$} (relativement à $S$). La condition a) de (17.13.10) s'exprime encore en la relation qui généralise (17.13.7.1)
\[
\label{IV.17.13.11.1}
\tag{17.13.11.1}
\operatorname{codim}_x(T_s,X_s)=\sum_{i\in I}\operatorname{codim}_x((Y_i)_s,X_s).
\]
En outre, on a la propriété suivante, qui étend (17.13.8.1), et en donne une autre démonstration quand $I$ a 2 éléments:
\end{env}

\begin{corollary}[17.13.12]
\label{IV.17.13.12}
Lorsque les $Y_i$ se coupent transversalement au point $x$, on a un isomorphisme canonique
\[
\label{IV.17.13.12.1}
\tag{17.13.12.1}
\sh{N}_{T/X}\xrightarrow{\sim}\bigoplus_{i\in I}(\sh{N}_{Y_i/X}\otimes_{\sh{O}_{Y_i}}\sh{O}_T).
\]
\end{corollary}

On peut se borner au cas où les $Y_i$ sont des sous-préschémas fermés de $X$, définis par des Idéaux quasi-cohérents $\sh{J}_i$, de sorte que $T$ est défini par l'Idéal $\sh{J}=\sum_i\sh{J}_i$. Par définition du faisceau conormal à une immersion (16.1.3), l'homomorphisme canonique $\bigoplus_i\sh{J}_i\to\sh{J}$ donne, par passage aux quotients, un homomorphisme \emph{surjectif}
\[
\label{IV.17.13.12.2}
\tag{17.13.12.2}
\bigoplus_{i\in I}(\sh{N}_{Y_i/X}\otimes_{\sh{O}_{Y_i}}\sh{O}_T)\to\sh{N}_{T/X}.
\]

\oldpage[IV-4]{98}

Mais ici les $\sh{O}_T$-Modules des deux membres de (17.13.12.2) sont localement libres et \emph{de même rang} $\sum_i c_i$ (si $c_i$ est le rang de $\sh{N}_{Y_i/X}$), en vertu de (17.2.5) et de la condition a) de (17.13.10); on conclut donc de Bourbaki, \emph{Alg. comm.}, chap. II, § 3, n\textsuperscript{o} 2, cor. de la prop. 6, que (17.13.12.2) est bijectif, et (17.13.12.1) est l'isomorphisme réciproque.

\subsection{Caractérisations locales et infinitésimales des morphismes lisses, des morphismes non ramifiés et des morphismes étales.}
\EGAReaderTarget{ega4.c04.subsec.17-14}{subsection:IV.17.14}

\begin{proposition}[17.14.1]
\label{IV.17.14.1}
Soit $f:X\to Y$ un morphisme localement de présentation finie, $x$ un point de $X$, $y=f(x)$. Pour que $f$ soit lisse (resp. non ramifié, resp. étale) au point $x$, il faut et il suffit que la condition suivante soit vérifiée:

Pour tout schéma local $Y'=\operatorname{Spec}(A')$ de point fermé $y'$, tout morphisme $h:Y'\to Y$ tel que $h(y')=y$, tout sous-schéma fermé $Y'_0=\operatorname{Spec}(A'/\mathfrak{J}')$ de $Y'$, où l'idéal $\mathfrak{J}'$ est de carré nul, et tout $Y$-morphisme $g_0:Y'_0\to X$ tel que $g_0(y')=x$, il existe au moins un (resp. au plus un, resp. un et un seul) $Y$-morphisme $g:Y'\to X$ dont $g_0$ soit la restriction à $Y'_0$.
\end{proposition}

Compte tenu des définitions (17.1.1 et 17.3.1), il s'agit de montrer que la condition de l'énoncé est \emph{suffisante} pour que $f$ soit lisse (resp. non ramifié) au point $x$.

\emph{(i) Cas des morphismes lisses.} On peut se borner au cas où $Y=\operatorname{Spec}(A)$ et $X=\operatorname{Spec}(C)$ sont affines, avec $C=B/\mathfrak{R}$, où $B=A[T_1,\ldots,T_n]$ est une $A$-algèbre de polynômes et $\mathfrak{R}$ un idéal de type fini de $B$. Pour prouver que $f$ est lisse au point $x$, il suffit d'établir que $\sh{O}_{X,x}$ est une $\sh{O}_{Y,y}$-algèbre formellement lisse pour les topologies discrètes (17.5.1). Or on a $\sh{O}_{X,x}=B_{\mathfrak q}/\mathfrak{R}B_{\mathfrak q}$, où $\mathfrak q$ est un idéal premier de $B$, et si $\mathfrak p$ est l'image réciproque de $\mathfrak q$ dans $A$, $B_{\mathfrak q}$ est une $A_{\mathfrak p}$-algèbre formellement lisse pour les topologies discrètes (0, 19.3.2 et 19.3.5). Par application du critère jacobien (0, 22.6.1 et 20.5.12) il suffit donc de voir que $B_{\mathfrak q}/\mathfrak{R}^2B_{\mathfrak q}$ est une extension $A_{\mathfrak p}$-triviale de $B_{\mathfrak q}/\mathfrak{R}B_{\mathfrak q}$ par $\mathfrak{R}B_{\mathfrak q}/\mathfrak{R}^2B_{\mathfrak q}$. Mais cela résulte précisément de l'hypothèse appliquée à $A'=B_{\mathfrak q}/\mathfrak{R}^2B_{\mathfrak q}$ et $\mathfrak{J}'=\mathfrak{R}B_{\mathfrak q}/\mathfrak{R}^2B_{\mathfrak q}$ et au morphisme $g_0$ correspondant à l'automorphisme identique de $A'/\mathfrak{J}'=\sh{O}_{X,x}$.

\emph{(ii) Cas des morphismes non ramifiés.} Posons ici $A=\sh{O}_{Y,y}$, $B=\sh{O}_{X,x}$, et notons qu'en vertu de (17.4.1, c')), il suffit de montrer que $\Omega^1_{B/A}=(\Omega^1_{X/Y})_x=0$. Or, avec les notations de (0, 20.4.1), l'anneau $C=\sh{P}^1_{B/A}=(B\otimes_A B)/\mathfrak{J}^2$ est local puisqu'il en est ainsi de $C/(\mathfrak{J}/\mathfrak{J}^2)$ isomorphe à $B$. L'hypothèse appliquée à $A'=C$ et $\mathfrak{J}'=\mathfrak{J}/\mathfrak{J}^2$ montre par définition que l'application $d_{B/A}:B\to\Omega^1_{B/A}$ est nulle (0, 20.4.6), donc que $\Omega^1_{B/A}=0$ par (0, 20.4.7).

\begin{proposition}[17.14.2]
\label{IV.17.14.2}
Soient $Y$ un préschéma localement noethérien, $f:X\to Y$ un morphisme localement de type fini, $x$ un point de $X$, $y=f(x)$. Pour que $f$ soit lisse (resp. non ramifié, resp. étale) au point $x$, il faut et il suffit qu'il vérifie la condition de (17.14.1) où l'on suppose de plus que l'anneau local $A'$ est artinien, que $\mathfrak m'\mathfrak{J}'=0$, où $\mathfrak m'$ est l'idéal maximal de $A'$, et que le corps résiduel $A'/\mathfrak m'$ de $A'$ est égal à $k(x)$.
\end{proposition}

Il s'agit encore de montrer que ces conditions sont suffisantes dans le cas des morphismes lisses et le cas des morphismes non ramifiés.

\oldpage[IV-4]{99}

\emph{(i) Cas des morphismes lisses.} Si l'on pose $A=\sh{O}_{Y,y}$ et $B=\sh{O}_{X,x}$, il s'agit ici de prouver, compte tenu de (17.5.3), que $B$ est une $A$-algèbre formellement lisse pour les topologies préadiques. Or, en vertu de (0, 22.1.4), cela découle de l'hypothèse lorsqu'on remplace dans celle-ci la condition $\mathfrak m'\mathfrak{J}'=0$ par $\mathfrak{J}'\subset\mathfrak m'$. Mais comme $\mathfrak m'$ est nilpotent, on voit que la condition de l'énoncé est déjà suffisante en considérant les anneaux $A'/\mathfrak m'^{,j}\mathfrak{J}'$ et en raisonnant par récurrence sur $j$ comme dans (0, 19.4.3).

\emph{(ii) Cas des morphismes non ramifiés.} Reprenons les notations de la démonstration de (17.14.1, (ii)); pour prouver que l'idéal $\mathfrak{J}/\mathfrak{J}^2$ de $C$ est nul, notons que $C$ est alors un anneau local noethérien, étant anneau local de l'anneau d'un ouvert affine d'un sous-préschéma de $X\times_YX$, qui est localement de type fini sur $X$, donc aussi sur $Y$. Si $\mathfrak r$ est l'idéal maximal de $C$, il suffit donc de vérifier que l'on a $\mathfrak{J}/\mathfrak{J}^2\subset\mathfrak r^n$ pour tout $n$ (0$_{\mathrm I}$, 7.3.5), ou encore que $\mathfrak r^n=\mathfrak r^n+(\mathfrak{J}/\mathfrak{J}^2)$. Avec les notations de (0, 20.4.6), cela signifie encore que pour tout $n$, les deux applications composées $B\xrightarrow{p_1}C\to C/\mathfrak r^n$ et $B\xrightarrow{p_2}C\to C/\mathfrak r^n$ sont identiques; mais comme $C/\mathfrak r^n$ est artinien, cette identité résulte de l'hypothèse appliquée, par récurrence sur $n$, à $A'=C/\mathfrak r^n$ et $\mathfrak{J}'=(\mathfrak r^n+(\mathfrak{J}/\mathfrak{J}^2))/\mathfrak r^n$.

\begin{remark}[17.14.3]
\label{IV.17.14.3}
On notera que, sous les hypothèses de (17.14.2), si de plus le corps $k(x)$ est une extension \emph{finie} de $k(y)$, alors les $A'$-modules $\mathfrak m'^{,j}/\mathfrak m'^{,j+1}$ sont des $k(y)$-espaces vectoriels de rang fini, donc \emph{a fortiori} $A'$ est une $\sh{O}_{Y,y}$-algèbre finie.
\end{remark}

\subsection{Cas des préschémas sur un corps de base.}
\EGAReaderTarget{ega4.c04.subsec.17-15}{subsection:IV.17.15}

Rappelons d'abord (6.7.7, 6.7.8 et 6.8.1) la

\begin{proposition}[17.15.1]
\label{IV.17.15.1}
Soient $k$ un corps, $X$ un préschéma localement de type fini sur $k$. Pour que $X$ soit lisse en un point $x$, il faut et il suffit que pour toute extension radicielle $k'$ de $k$ (ou seulement pour toute extension finie radicielle de $k$), l'anneau local $\sh{O}_{X,x}\otimes_k k'$ soit régulier. Si $k(x)$ est une extension séparable de $k$ (en particulier si $k$ est parfait), il revient au même de dire que $\sh{O}_{X,x}$ est régulier.
\end{proposition}

\begin{corollary}[17.15.2]
\label{IV.17.15.2}
Sous les conditions de (17.15.1), pour que $X$ soit lisse sur $k$, il faut et il suffit que $X$ soit géométriquement régulier sur $k$ (6.7.6), ce qui entraîne que $X$ est régulier. Si $k$ est parfait, alors $X$ est lisse sur $k$ si et seulement si $X$ est régulier.
\end{corollary}

\begin{proposition}[17.15.3]
\label{IV.17.15.3}
Soient $k$ un corps, $X$ un préschéma localement de type fini sur $k$, $x$ un point de $X$, $n=\dim_x(X)$. Soit $(g_i)_{1\leq i\leq n}$ une famille de $n$ sections de $\sh{O}_X$ au-dessus de $X$, et soit $f:X\to Y=\operatorname{Spec}(k[T_1,\ldots,T_n])$ le morphisme correspondant au $k$-homomorphisme $k[T_1,\ldots,T_n]\to\Gamma(X,\sh{O}_X)$ transformant $T_i$ en $g_i$ (I, 2.2.4). Les conditions suivantes sont équivalentes (et impliquent que $X$ est lisse sur $k$ au point $x$ et par suite régulier au point $x$):
\begin{enumerate}
  \item[a)] $f$ est étale au point $x$.
  \item[b)] Les images des $d_{X/k}g_i$ forment une base du $\sh{O}_{X,x}$-module $(\Omega^1_{X/k})_x$.
  \item[c)] Les images des $d_{X/k}g_i$ engendrent le $\sh{O}_{X,x}$-module $(\Omega^1_{X/k})_x$.
\end{enumerate}
\end{proposition}

Si $f$ est étale au point $x$, $X$ est lisse sur $k$ au point $x$ puisque $Y=\operatorname{Spec}(k[T_1,\ldots,T_n])$ est lisse sur $k$ (17.3.8); le fait que a) entraîne b) est donc un cas particulier de (17.11.3). Comme b) entraîne trivialement c), il reste à voir que c) implique a). Notons d'abord

\oldpage[IV-4]{100}

que l'hypothèse entraîne que l'homomorphisme $(f^*(\Omega^1_{Y/k}))_x\to(\Omega^1_{X/k})_x$ est surjectif, donc, en remplaçant $X$ par un voisinage ouvert de $x$, on peut supposer que l'homomorphisme $f^*(\Omega^1_{Y/k})\to\Omega^1_{X/k}$ est surjectif (Bourbaki, \emph{Alg. comm.}, chap. II, § 5, n\textsuperscript{o} 1, prop. 2); par suite $f$ est \emph{non ramifié} (17.2.2). Nous allons voir d'abord qu'on peut se borner au cas où $x$ est \emph{rationnel sur $k$}. En effet, si l'on pose $k'=k(x)$, et $X'=X\otimes_k k'$, $Y'=Y\otimes_k k'=\operatorname{Spec}(k'[T_1,\ldots,T_n])$, il existe un point $x'\in X'$ au-dessus de $x$, tel que $k(x')=k'$. Pour prouver que $f$ est étale au point $x$, il suffit de montrer que $f'=f_{(k')}:X'\to Y'$ est étale au point $x'$ (17.7.1, (ii)); d'ailleurs $f'$ est non ramifié (17.3.3, (iii)) et l'on a $\dim_{x'}(X')=n$ (4.2.7). De la même manière on peut, en remplaçant $k'$ par une extension algébriquement close de $k'$, supposer que $k$ est \emph{algébriquement clos}. Posons alors $y=f(x)$, $A=\sh{O}_{Y,y}$, $B=\sh{O}_{X,x}$; comme le corps résiduel de $B$ est égal à $k$, il en est de même de celui de $A$, donc $x$ (resp. $y$) est un point fermé de $X$ (resp. $Y$) (I, 6.4.2) et l'on a par suite $\dim(A)=\dim(B)=n$. Puisque $f$ est non ramifié, l'homomorphisme $\widehat A\to\widehat B$ est surjectif (17.4.4, f'')); mais comme $\dim(\widehat A)=\dim(\widehat B)=n$, et que $\widehat A$, étant un anneau régulier, est intègre, l'homomorphisme $\widehat A\to\widehat B$ est aussi injectif (0, 16.3.10). Cet homomorphisme est donc bijectif, ce qui achève de prouver que $f$ est étale au point $x$ (17.6.3, e'')).

\begin{corollary}[17.15.4]
\label{IV.17.15.4}
Sous les hypothèses générales de (17.15.3), supposons en outre que $k(x)$ soit une extension finie et séparable de $k$ et que les germes $(g_i)_x$ appartiennent à $\mathfrak m_x$. Alors les conditions a), b), c) de (17.15.3) équivalent aussi à chacune des suivantes:
\begin{enumerate}
  \item[d)] Les $n$ germes $(g_i)_x$ engendrent l'idéal maximal $\mathfrak m_x$ de $\sh{O}_{X,x}$.
  \item[d')] L'anneau $\sh{O}_{X,x}$ est régulier et les $(g_i)_x$ forment un système régulier de paramètres pour cet anneau (0, 17.1.6).
\end{enumerate}
\end{corollary}

En effet, d') entraîne trivialement d). D'autre part, comme $k(x)$ est une extension finie et séparable de $k$, on a $\Omega^1_{k(x)/k}=0$ (0, 20.6.20) et $k(x)=\sh{O}_{X,x}/\mathfrak m_x$ est une $k$-algèbre formellement lisse pour les topologies discrètes (0, 19.6.1), donc la suite exacte (0, 20.5.14.1) s'applique à $A=k$, $B=\sh{O}_{X,x}$, $\mathfrak{R}=\mathfrak m_x$, et fournit un isomorphisme canonique
\[
\delta:\mathfrak m_x/\mathfrak m_x^2\xrightarrow{\sim}(\Omega^1_{X/k})_x\otimes_{\sh{O}_{X,x}}k(x).
\]
De là on déduit d'abord l'équivalence des conditions c) et d), compte tenu du lemme de Nakayama. D'autre part, si $f$ est étale au point $x$, l'anneau $\sh{O}_{X,x}$ est régulier et de dimension $n$, puisque $x$ est un point fermé de $X$ (I, 6.4.2), et $n$ éléments de $\mathfrak m_x$ qui engendrent cet idéal maximal forment alors nécessairement un système régulier de paramètres pour $\sh{O}_{X,x}$ (0, 17.1.6), ce qui prouve que a) entraîne d').

\begin{proposition}[17.15.5]
\label{IV.17.15.5}
Soient $k$ un corps, $X$ un préschéma localement de type fini sur $k$, $x$ un point de $X$, $n=\dim_x(X)$. Les conditions suivantes sont équivalentes:
\begin{enumerate}
  \item[a)] $X$ est lisse sur $k$ au point $x$.
  \item[b)] $X$ est différentiellement lisse sur $k$ au point $x$.
  \item[c)] $(\Omega^1_{X/k})_x$ est un $\sh{O}_{X,x}$-module libre de rang $n$.
  \item[d)] $(\Omega^1_{X/k})_x$ est un $\sh{O}_{X,x}$-module admettant un système de $n$ générateurs.

\oldpage[IV-4]{101}

  \item[e)] Il existe une extension parfaite $k'$ de $k$ et un voisinage ouvert $U$ de $x$ dans $X$ tels que le préschéma $U\otimes_k k'$ soit régulier.
\end{enumerate}
\end{proposition}

Le fait que $X$ soit lisse sur $k$ au point $x$ entraîne l'existence d'un voisinage ouvert $U$ de $x$ qui est lisse sur $k$, donc $U\otimes_k k'$ est régulier pour toute extension $k'$ de $k$ (17.15.2), ce qui prouve que a) implique e). Inversement e) implique a), car alors $U\otimes_k k'$ est lisse sur $k'$ (17.15.2), donc $U$ est lisse sur $k$ (17.7.1, (ii)).

On a déjà prouvé que a) entraîne c) (17.10.2); c) implique d) trivialement et le fait que d) implique a) résulte de (0, 20.4.7) et (17.15.3, c)).

Enfin, on a déjà vu que a) implique b) (17.12.4). Inversement, pour prouver que b) entraîne a), on peut se borner au cas où $x$ est rationnel sur $k$, en considérant, comme dans la démonstration de (17.15.3), un point $x'$ de $X'=X\otimes_k k'$ au-dessus de $x$ et tel que $k(x')=k'=k(x)$, utilisant encore (17.7.1, (ii)) et le fait que si $X$ est différentiellement lisse au point $x$, $X'$ est différentiellement lisse au point $x'$ (16.10.4). Supposant donc $x$ rationnel sur $k$, le fait que $f$ soit lisse en $x$ résulte alors de l'hypothèse b) et de (17.12.5) appliqué à la $k$-section $u:\operatorname{Spec}(k)\to X$ telle que $u(\operatorname{Spec}(k))=\{x\}$.

\begin{corollary}[17.15.6]
\label{IV.17.15.6}
Soit $X$ un préschéma de type fini sur un corps $k$. Pour que $X$ soit lisse sur $k$, il faut et il suffit que le $\sh{O}_X$-Module $\Omega^1_{X/k}$ soit localement libre, et que les anneaux locaux aux points maximaux de $X$ soient des corps, extensions séparables de $k$ (cette dernière condition étant automatiquement vérifiée si $k$ est un corps parfait et $X$ un préschéma réduit).
\end{corollary}

Les conditions sont nécessaires, car si $X$ est lisse sur $k$, il résulte de (17.10.2) que $\Omega^1_{X/k}$ est localement libre; d'autre part, $X$ est régulier, donc \emph{a fortiori} réduit, donc en tout point maximal $x$ de $X$, $\sh{O}_{X,x}$ est un corps, qui doit être une $k$-algèbre formellement lisse pour les topologies discrètes (17.5.1), donc une extension séparable de $k$ (0, 19.6.1).

Les conditions sont suffisantes. On peut en effet se borner au cas où $X$ est connexe, donc $\Omega^1_{X/k}$ localement libre de rang constant $n$. Pour tout point maximal $x$ de $X$, on a alors $\operatorname{rg}_{k(x)}\Omega^1_{k(x)/k}=n$ puisque $\sh{O}_{X,x}=k(x)$; comme par hypothèse $k(x)$ est une extension séparable de $k$, on a $\Upsilon_{k(x)/k}=0$ (0, 20.6.3), donc $\deg.\operatorname{tr}_k k(x)=n$ par l'égalité de Cartier (0, 21.7.1). Toutes les composantes irréductibles de $X$ ont donc même dimension $n$ (5.2.1), et l'on conclut que $X$ est lisse sur $k$ en tout point en vertu de ce que c) entraîne a) dans (17.15.5).

\begin{corollary}[17.15.7]
\label{IV.17.15.7}
Si $k$ est de caractéristique 0, alors, pour que $X$ soit lisse sur $k$ au point $x$, il faut et il suffit que $(\Omega^1_{X/k})_x$ soit un $\sh{O}_{X,x}$-module libre.
\end{corollary}

Cela résulte de (16.12.2) et de l'équivalence de b) et a) dans (17.15.5).

\begin{proposition}[17.15.8]
\label{IV.17.15.8}
Soient $k$ un corps, $X$ un préschéma localement de type fini sur $k$, $x$ un point de $X$; posons $n=\dim_x(X)$, $r=\dim(\sh{O}_{X,x})$, de sorte que $n-r=\deg.\operatorname{tr}_k k(x)$ (5.2.3). Soit $(g_i)_{1\leq i\leq n}$ une famille de $n$ sections de $\sh{O}_X$ au-dessus de $X$, telles que $(g_i)_x\in\mathfrak m_x$ pour $1\leq i\leq r$; soit $f:X\to Y=\operatorname{Spec}(k[T_1,\ldots,T_n])$ le morphisme correspondant au $k$-homomorphisme $k[T_1,\ldots,T_n]\to\Gamma(X,\sh{O}_X)$ transformant $T_i$ en $g_i$. Les conditions suivantes sont équivalentes (et impliquent que $X$ est lisse sur $k$ au point $x$ et que $k(x)$ est une extension séparable de $k$):
\begin{enumerate}
  \item[a)] $f$ est étale au point $x$.

\oldpage[IV-4]{102}

  \item[b)] Les $(g_i)_x$ tels que $1\leq i\leq r$ engendrent $\mathfrak m_x$ (et par suite forment un système régulier de paramètres pour $\sh{O}_{X,x}$ (0, 17.1.6)) et les images dans $\Omega^1_{k(x)/k}$ des éléments $(d_{X/k}g_i)_x$ pour $r+1\leq i\leq n$ engendrent $\Omega^1_{k(x)/k}$.
\end{enumerate}
\end{proposition}

On a en effet (0, 20.5.12.1) la suite exacte de $k(x)$-modules
\[
\label{IV.17.15.8.1}
\tag{17.15.8.1}
\mathfrak m_x/\mathfrak m_x^2\to(\Omega^1_{X/k})_x\otimes_{\sh{O}_{X,x}}k(x)\to\Omega^1_{k(x)/k}\to0
\]
et la condition b) entraîne par suite que les $(d_{X/k}g_i)_x$ engendrent $(\Omega^1_{X/k})_x$ compte tenu du lemme de Nakayama; le fait que b) implique a) résulte donc de (17.15.3). Inversement, si a) est vérifiée, les $(d_{X/k}g_i)_x$ pour $1\leq i\leq n$ forment une base de $(\Omega^1_{X/k})_x$ en vertu de (17.15.3). Si $t_i$ ($1\leq i\leq n$) est l'image de $(g_i)_x$ dans $k(x)$, on conclut de ce qui précède que les $dt_i$ engendrent $\Omega^1_{k(x)/k}$ pour $1\leq i\leq n$, et comme par hypothèse $t_i=0$ pour $1\leq i\leq r$, les $dt_i$ pour $r+1\leq i\leq n$ engendrent déjà $\Omega^1_{k(x)/k}$. Comme $\deg.\operatorname{tr}_k k(x)=n-r$, il suit de l'égalité de Cartier (0, 21.7.1) que $\Upsilon_{k(x)/k}=0$, donc $k(x)$ est une extension séparable de $k$ (0, 20.6.3), et la suite de $k(x)$-espaces vectoriels
\[
\label{IV.17.15.8.2}
\tag{17.15.8.2}
0\to\mathfrak m_x/\mathfrak m_x^2\to(\Omega^1_{X/k})_x\otimes_{\sh{O}_{X,x}}k(x)\to\Omega^1_{k(x)/k}\to0
\]
est exacte ((0, 20.5.14) et (0, 19.6.1)); en outre les $dt_i$ pour $r+1\leq i\leq n$ forment une base de $\Omega^1_{k(x)/k}$, donc aucun des $t_i$ tels que $r+1\leq i\leq n$ ne peut être nul. Cela montre que les images des $(g_i)_x$ dans $\mathfrak m_x/\mathfrak m_x^2$ pour $1\leq i\leq r$ engendrent nécessairement $\mathfrak m_x/\mathfrak m_x^2$, donc les $(g_i)_x$ pour $1\leq i\leq r$ engendrent $\mathfrak m_x$ en vertu du lemme de Nakayama, ce qui achève de prouver que a) implique b).

\begin{corollary}[17.15.9]
\label{IV.17.15.9}
Soient $X$ un préschéma localement de type fini sur un corps $k$, $x$ un point de $X$ et posons $n=\dim_x(X)$, $r=\dim(\sh{O}_{X,x})$. Les conditions suivantes sont équivalentes:
\begin{enumerate}
  \item[a)] L'anneau $\sh{O}_{X,x}$ est régulier et $k(x)$ est une extension séparable de $k$.
  \item[b)] $X$ est lisse sur $k$ au point $x$, et l'homomorphisme canonique
  \[
  \mathfrak m_x/\mathfrak m_x^2\to(\Omega^1_{X/k})_x\otimes_{\sh{O}_{X,x}}k(x)
  \]
  est injectif.
  \item[c)] Il existe $n$ sections $g_i$ de $\sh{O}_X$ au-dessus d'un voisinage ouvert $U$ de $x$ telles que $(g_i)_x\in\mathfrak m_x$ pour $1\leq i\leq r$ et que le morphisme $U\to\operatorname{Spec}(k[T_1,\ldots,T_n])$ correspondant aux $g_i$ (cf. (17.15.8)) soit étale au point $x$.
  \item[d)] Il existe $n$ sections $g_i$ ($1\leq i\leq n$) de $\sh{O}_X$ au-dessus d'un voisinage ouvert $U$ de $x$, telles que les $(g_i)_x$ pour $1\leq i\leq r$ engendrent $\mathfrak m_x$ et que les images dans $\Omega^1_{k(x)/k}$ des $(d_{U/k}g_i)_x$ pour $r+1\leq i\leq n$ engendrent $\Omega^1_{k(x)/k}$.
\end{enumerate}
\end{corollary}

L'équivalence de c) et d) résulte de (17.15.8). De plus, on a vu dans la démonstration de (17.15.8) que la condition c) entraîne que $X$ est lisse sur $k$ au point $x$ et que la suite (17.15.8.2) est exacte, donc c) entraîne b). La condition b) entraîne que l'anneau $\sh{O}_{X,x}$ est régulier, donc $\mathfrak m_x/\mathfrak m_x^2$ est de rang $r$; en outre, $(\Omega^1_{X/k})_x$ est alors un $\sh{O}_{X,x}$-module libre de rang $n$ (17.10.2); comme la suite (17.15.7.2) est exacte par hypothèse, $\Omega^1_{k(x)/k}$ est de rang $n-r=\deg.\operatorname{tr}_k k(x)$ et l'égalité de Cartier montre que $\Upsilon_{k(x)/k}=0$, donc $k(x)$ est séparable sur $k$ (0, 20.6.3); ainsi b) entraîne a). Enfin, si a)

\oldpage[IV-4]{103}

est vérifiée, on déduit encore de l'égalité de Cartier que $\Omega^1_{k(x)/k}$ est de rang $n-r$. Comme d'autre part l'anneau $\sh{O}_{X,x}$ est régulier, l'existence des $g_i$ vérifiant les conditions de d) est immédiate.

\begin{remarks}[17.15.10]
\label{IV.17.15.10}
\begin{enumerate}
  \item[(i)] Pour qu'un préschéma de type fini $X$ sur $k$ soit lisse sur $k$, il ne suffit pas que $\Omega^1_{X/k}$ soit localement libre, comme le montre l'exemple où $X=\operatorname{Spec}(K)$, $K$ étant une extension finie non séparable de $k$.
  \item[(ii)] Lorsque $k$ n'est pas parfait, il peut se faire que $X$ soit lisse sur $k$ sans que $k(x)$ soit séparable sur $k$. On en a un exemple en prenant $X=\operatorname{Spec}(k[T])$ et pour $x$ le point correspondant à l'idéal premier principal $(T^p-\lambda)$ ($p>0$ caractéristique de $k$, $\lambda\in k-k^p$).
  \item[(iii)] Toutefois, si $X$ est un préschéma lisse sur $k$, l'ensemble des points fermés de $X$ tels que $k(x)$ soit séparable (et fini) sur $k$ est dense dans $X$. En effet, soit $f:X\to\operatorname{Spec}(k)$ le morphisme structural; pour tout $x_0\in X$, il y a un voisinage ouvert $U$ de $x_0$ dans $X$ et une factorisation de $f|U:U\xrightarrow{g}\operatorname{Spec}(k[T_1,\ldots,T_n])\xrightarrow{h}\operatorname{Spec}(k)$ où $g$ est étale (17.11.4). Comme on peut se borner au cas où $k$ est non parfait, donc infini, l'ensemble des points de l'ouvert $g(U)$ qui sont rationnels sur $k$ est non vide; si $y$ est un tel point et $x\in U$ un point au-dessus de $y$, $x$ est fermé dans $X$ et $k(x)$ est séparable sur $k(y)=k$ (17.6.2).
\end{enumerate}
\end{remarks}

\begin{proposition}[17.15.11]
\label{IV.17.15.11}
Soit $X$ un préschéma de type fini sur un corps $k$. Les conditions suivantes sont équivalentes:
\begin{enumerate}
  \item[a)] $X$ est étale sur $k$.
  \item[b)] $X$ est non ramifié sur $k$.
  \item[c)] $X$ est isomorphe à $\operatorname{Spec}(A)$, où $A$ est une $k$-algèbre finie et séparable.
\end{enumerate}
\end{proposition}

Cela résulte de (17.6.2) et (17.4.2), compte tenu de ce que la condition b) implique ici que $X$ est fini sur $k$.

\begin{proposition}[17.15.12]
\label{IV.17.15.12}
Soient $k$ un corps, $X$ un préschéma localement de type fini sur $k$. Pour qu'il existe un ouvert partout dense $U$ de $X$ qui soit lisse sur $k$, il faut et il suffit que pour tout point maximal $x$ de $X$, $X$ soit réduit en ce point et que $k(x)$ soit une extension séparable de $k$. Pour qu'il existe un ouvert partout dense $U$ de $X$ tel que $U_{\mathrm{red}}$ soit lisse sur $k$, il faut et il suffit que pour tout point maximal $x$ de $X$, $k(x)$ soit une extension séparable de $k$.
\end{proposition}

La seconde assertion résulte évidemment de la première. Dire qu'il existe un ouvert partout dense $U$ lisse sur $k$ signifie que $X$ est lisse sur $k$ en chacun de ses points maximaux $x$. Il est nécessaire pour cela que $\sh{O}_{X,x}$ soit régulier (17.15.1) et \emph{a fortiori} réduit, donc un corps, égal à $k(x)$; en outre (17.15.1), $k(x)\otimes_k k'$ doit être régulier pour toute extension radicielle $k'$ de $k$, donc $k(x)$ doit être une extension séparable de $k$ (4.3.5); la réciproque est immédiate, en vertu de (17.15.1).

\begin{corollary}[17.15.13]
\label{IV.17.15.13}
Soit $X$ un préschéma algébrique sur un corps $k$. Alors il existe un ouvert partout dense $U$ dans $X$ et une extension finie radicielle $k'$ de $k$ tels que $(U_{(k')})_{\mathrm{red}}$ soit lisse sur $k'$.
\end{corollary}

En effet, il y a une telle extension $k'$ telle que $(X_{(k')})_{\mathrm{red}}$ soit géométriquement réduit sur $k'$ (4.6.6), ce qui équivaut à dire (4.6.1) que pour tout point maximal $x'$ de $X_{(k')}$, $k(x')$ est une extension séparable de $k'$. On conclut donc de (17.15.12) qu'il y a un ouvert partout dense $U'$ de $X_{(k')}$ tel que $U'_{\mathrm{red}}$ soit lisse sur $k'$. Mais le \oldpage[IV-4]{104}morphisme $X_{(k')}\to X$ est un homéomorphisme (2.4.5), donc $U'$ est de la forme $U_{(k')}$, où $U$ est un ouvert partout dense dans $X$.

\begin{proposition}[17.15.14]
\label{IV.17.15.14}
Soit $X$ un préschéma algébrique sur $k$, de dimension $\leq1$. Alors il existe une extension finie radicielle $k'$ de $k$ telle que le normalisé (II, 6.3.8) de $(X_{(k')})_{\mathrm{red}}$ soit lisse sur $k'$.
\end{proposition}

Démontrons d'abord les deux lemmes suivants:

\begin{lemma}[17.15.14.1]
\label{IV.17.15.14.1}
Soient $X$ un préschéma réduit dont l'ensemble des composantes irréductibles est localement fini, $f:Y\to X$ un morphisme. Pour que $Y$ soit $X$-isomorphe au normalisé $X'$ de $X$ (II, 6.3.8), il faut et il suffit que les conditions suivantes soient vérifiées:
\begin{enumerate}
  \item[(i)] $Y$ est normal;
  \item[(ii)] $f$ est un morphisme entier et birationnel.
\end{enumerate}
Lorsqu'il existe dans $X$ un ouvert dense et normal (ce qui est toujours le cas lorsque $X$ est excellent (7.8.3), en particulier lorsque $X$ est localement de type fini sur un corps), on peut remplacer la condition (ii) par la suivante:
\begin{enumerate}
  \item[(ii bis)] $f$ est entier et il existe un ouvert dense $U$ dans $X$ tel que $f^{-1}(U)$ soit dense dans $Y$ et que la restriction $f^{-1}(U)\to U$ de $f$ soit un isomorphisme.
\end{enumerate}
\end{lemma}

La question étant locale sur $X$, on peut supposer que $X$ n'a qu'un nombre fini de composantes irréductibles; soient $X_i$ ($1\leq i\leq m$) les sous-préschémas réduits de $X$ ayant ces composantes pour espaces sous-jacents. On sait (II, 6.3.6) que $X'$ est somme des normalisés $X'_i$ des $X_i$; chacun des morphismes structuraux $f_i:X'_i\to X_i$ est donc entier et birationnel, donc $f$ est entier et birationnel; il est immédiat que (ii) entraîne (ii bis) lorsqu'il existe un ouvert $V$ dense et normal dans $X$, en prenant $U=V$. Inversement, si les conditions (i) et (ii) sont remplies, $Y$ n'a qu'un nombre fini de composantes irréductibles $Y_i$ ($1\leq i\leq m$), $Y_i$ dominant $X_i$ pour chaque indice $i$ (6.15.4); $Y$ est somme des $Y_i$ puisqu'il est normal, et le morphisme $g_i:Y_i\to X_i$ en lequel se factorise la restriction $Y_i\to X$ de $f$ (I, 5.2.2) est birationnel; comme $X_i$ et $Y_i$ sont intègres, il résulte alors du fait que $g_i$ est entier et $Y_i$ normal que, pour tout ouvert affine $V$ de $X_i$, $\Gamma(g_i^{-1}(V),\sh{O}_{Y_i})$ est la clôture intégrale de $\Gamma(V,\sh{O}_{X_i})$, donc $Y$ s'identifie canoniquement à $X'$ (II, 6.3.4). Enfin, si l'on suppose vérifiée la condition (ii bis), on peut se borner au cas où $U$ est réunion d'ouverts irréductibles disjoints $U_i$ ($1\leq i\leq m$), donc $f^{-1}(U)$ réunion disjointe des $V_i=f^{-1}(U_i)$. Comme les $V_i$ sont irréductibles et $f^{-1}(U)$ dense dans $Y$, les $\overline{V_i}$ sont les composantes irréductibles de $Y$ et par suite $f$ est birationnel.

\begin{lemma}[17.15.14.2]
\label{IV.17.15.14.2}
Soient $k$ un corps, $X$ un préschéma algébrique sur $k$. Alors il existe une extension finie radicielle $k'$ de $k$ telle que $(X_{(k')})_{\mathrm{red}}$ soit géométriquement réduit sur $k'$ et que son normalisé soit géométriquement normal sur $k'$.
\end{lemma}

Compte tenu de (4.6.6) et de (I, 5.1.8), on peut se borner au cas où $X$ est déjà géométriquement réduit sur $k$. Soient $p$ l'exposant caractéristique de $k$ et $k_1=k^{p^{-\infty}}$ la plus petite extension parfaite de $k$; $k_1$ est donc limite inductive des extensions finies radicielles de $k$. Posons $X_1=X_{(k_1)}$, qui est réduit par hypothèse, et soit $Y_1$ son normalisé; si $f_1:Y_1\to X_1$ est le morphisme structural, $f_1$ est donc fini (7.8.3, (vi)) et surjectif, et il existe un ouvert dense $U_1$ dans $X_1$ tel que $f_1^{-1}(U_1)=V_1$ soit dense dans $Y_1$ et que \oldpage[IV-4]{105}le morphisme $V_1\to U_1$ restriction de $f_1$ soit un isomorphisme (17.15.14.1). Appliquant (8.9.1) et (8.10.5, (vi) et (x)), on voit donc d'abord qu'il existe une extension finie radicielle $k''$ de $k$ et un morphisme fini surjectif $Y''\to X''=X_{(k'')}$ tels que $Y_1=Y''\times_{X''}X_1=Y''\otimes_{k''}k_1$; puisque $Y_1$ est normal et $k_1$ parfait, il résulte de (6.7.7) que $Y''$ est géométriquement normal sur $k''$. Comme les projections $X_1\to X''$ et $Y_1\to Y''$ sont des morphismes entiers, surjectifs et radiciels, ce sont des homéomorphismes (2.4.5), et si $U''$ et $V''$ sont les images de $U_1$ et $V_1$ dans $X''$ et $Y''$ respectivement, ce sont des ouverts denses dans $X''$ et $Y''$ respectivement tels que $V''=f''^{-1}(U'')$, où $f'':Y''\to X''$ est le morphisme structural. Comme on a $U_1=U''\otimes_{k''}k_1$, il résulte de (8.10.5, (i)) qu'il existe une extension finie radicielle $k'$ de $k''$ telle que si l'on pose $X'=X\otimes_k k'=X''\otimes_{k''}k'$, $Y'=Y''\otimes_{k''}k'$, $f'=f''_{(k')}:Y'\to X'$, et si $U'$ et $V'$ sont les images de $U''$ et $V''$ dans $X'$ et $Y'$ respectivement, la restriction $V'\to U'$ de $f'$ soit un isomorphisme. Comme $Y'$ est normal et $f'$ entier et birationnel, on conclut de (17.15.14.1) que $Y'$ est isomorphe au normalisé de $X'$, ce qui prouve le lemme puisque $Y'$ est géométriquement normal.

Revenons maintenant à la démonstration de (17.15.14), et appliquons à $k$ et $X$ le lemme (17.15.14.2); puisque $\dim(X)\leq1$, on a aussi $\dim(X_{(k')})\leq1$ (4.1.4), et le normalisé $Y'$ de $X_{(k')}$ est aussi de dimension $\leq1$ (5.4.2 et II, 7.4.6). Dire que $Y'$ est géométriquement normal sur $k'$ revient alors à dire, en vertu des définitions et de (II, 7.4.5), que $Y'$ est géométriquement régulier sur $k'$, donc lisse sur $k'$ (17.5.1), ce qui prouve (17.15.14).

\begin{proposition}[17.15.15]
\label{IV.17.15.15}
Soit $f:X\to Y$ un morphisme localement de présentation finie. Pour que $f$ soit lisse en un point $x\in X$, il faut et il suffit que $f$ soit plat au point $x$ et que $\Omega^1_f$ soit un $\sh{O}_X$-Module localement libre dans un voisinage de $x$, de rang en $x$ égal à $\dim_x f$.
\end{proposition}

La nécessité des conditions résulte de (17.5.1) et (17.10.2). Inversement, si ces conditions sont vérifiées, et si $y=f(x)$, il suffit de montrer (17.5.1) que $f^{-1}(y)$ est lisse sur $k(y)$ au point $x$; mais cela résulte de la définition de $\dim_x f$ (17.10.1), de (16.4.5) et de (17.15.5).

\subsection{Quasi-sections de morphismes plats ou lisses.}
\EGAReaderTarget{ega4.c04.subsec.17-16}{subsection:IV.17.16}

Les énoncés de ce numéro complètent ceux de (14.5), moyennant des hypothèses de platitude.

\begin{proposition}[17.16.1]
\label{IV.17.16.1}
Soit $f:X\to S$ un morphisme plat localement de présentation finie. Soient $s\in S$, $x$ un point fermé de $X_s$ tel que $\sh{O}_{X_s,x}$ soit un anneau de Cohen--Macaulay; alors il existe un voisinage ouvert $U$ de $s$ dans $S$ et un sous-préschéma $X'\subset f^{-1}(U)$ de $X$ tel que $x\in X'$ et que le morphisme $X'\to U$, restriction de $f$, soit plat, quasi-fini et de présentation finie.
\end{proposition}

La question étant locale sur $X$ et $Y$, on peut supposer $f$ de présentation finie. Soit $(t_i)_{1\leq i\leq m}$ un système de paramètres de l'anneau local $\sh{O}_{X_s,x}$; il existe un voisinage ouvert affine $V=\Spec(A)$ de $x$ dans $X$ et $m$ sections $g_i$ de $\sh{O}_X$ au-dessus de $V$ telles que

\oldpage[IV-4]{106}

les images des $(g_i)_x\in\sh{O}_{X,x}$ dans $\sh{O}_{X_s,x}=\sh{O}_{X,x}/\mathfrak m_s\sh{O}_{X,x}$ soient égales aux $t_i$. Soit $X'$ le sous-préschéma fermé de $V$ défini par l'idéal de $A$ engendré par les $g_i$; la suite $(t_i)$ étant par hypothèse régulière (0, 16.5.7) il résulte de (11.3.8) qu'en remplaçant au besoin $V$ par un voisinage ouvert plus petit, on peut supposer que le morphisme $X'\to S$, restriction de $f$, est plat et de présentation finie. D'autre part, puisque les $t_i$ forment un système de paramètres de $\sh{O}_{X_s,x}$, l'anneau $\sh{O}_{X'_s,x}$ est par définition artinien, et comme $x$ est fermé dans $X'_s$, on en conclut qu'il est isolé dans $X'_s$. En remplaçant $V$ au besoin par un voisinage plus petit de $x$, on en conclut, grâce à (13.1.4) que le morphisme $X'\to S$ est quasi-fini.

\begin{corollary}[17.16.2]
\label{IV.17.16.2}
Soit $f:X\to S$ un morphisme fidèlement plat et localement de présentation finie. Alors il existe un morphisme $g:S'\to S$, fidèlement plat, localement de présentation finie et localement quasi-fini, tel qu'il existe un $S$-morphisme $S'\to X$ (autrement dit, tel qu'il existe une $S'$-section de $X'=X\times_S S'$ (I, 3.3.14)). Si $S$ est quasi-compact (resp. quasi-compact et quasi-séparé), on peut supposer $S'$ affine (resp. $S'$ affine et le morphisme $g$ quasi-fini).
\end{corollary}

Pour tout $s\in S$, la fibre $X_s$ est non vide par hypothèse et est un préschéma localement de type fini sur $k(s)$; l'ensemble $U_s$ des points $x$ de $X_s$ où $\sh{O}_{X_s,x}$ est un anneau de Cohen--Macaulay est ouvert dans $X_s$ (6.11.3) et est non vide, puisqu'il contient les points maximaux de $X_s$ (0, 16.5.1); il contient par suite un point $x_s$ fermé dans $X_s$ (10.4.7). Soit $X'(s)$ un sous-préschéma affine de $X$ contenant $x_s$ et vérifiant les conditions de (17.16.1). Pour obtenir un préschéma $S'$ vérifiant les conditions de l'énoncé, il suffit de prendre la somme des $X'(s)$, où $s$ parcourt $S$. Puisque le morphisme $X'(s)\to S$ est plat et localement de présentation finie, l'image $U(s)$ de $X'(s)$ est ouverte dans $S$ (2.4.6); lorsque $S$ est quasi-compact, on peut donc recouvrir $S$ par un nombre fini de $U(s_i)$ et le préschéma $S'$ somme des $X'(s_i)$ répond encore à la question et est affine. Si de plus $S$ est quasi-séparé, on peut supposer les immersions ouvertes $U(s_i)\to S$ quasi-compactes (I, 2.7), donc de présentation finie (I, 6.2), et alors les morphismes $X'(s_i)\to S$ sont de présentation finie (I, 6.2) et il en est par suite de même du morphisme $S'\to S$.

\begin{corollary}[17.16.3]
\label{IV.17.16.3}
\begin{enumerate}
  \item[(i)] Soit $f:X\to S$ un morphisme lisse. Soient $s\in S$, $x$ un point fermé de $X_s$ tel que le corps résiduel $k(x)$ soit séparable sur $k(s)$; alors, dans la conclusion de (17.16.1), on peut prendre $X'$ tel que le morphisme $X'\to U$, restriction de $f$, soit étale.
  \item[(ii)] Soit $f:X\to S$ un morphisme lisse surjectif. Alors, dans les conclusions de (17.16.2), on peut supposer en outre que $g:S'\to S$ soit étale.
\end{enumerate}
\end{corollary}

Il est clair que (ii) se déduit de (i) comme (17.16.2) de (17.16.1), compte tenu de ce que, pour tout $s\in S$, comme $X_s$ est non vide et lisse sur $k(s)$, il existe un point fermé $x\in X_s$ tel que $k(x)$ soit séparable sur $k(s)$ (17.15.10, (iii)). Il suffit donc de prouver (i). On notera que l'anneau $\sh{O}_{X_s,x}$ est alors régulier; si l'on répète la construction faite dans (17.16.1) en prenant pour $(t_i)$ un système régulier de paramètres de $\sh{O}_{X_s,x}$, l'anneau $\sh{O}_{X'_s,x}$ est un corps isomorphe à $k(x)$, donc séparable sur $k(s)$ par hypothèse. La conclusion résulte alors de (17.6.1, $c'$)).

\begin{proposition}[17.16.4]
\label{IV.17.16.4}
Soient $S$ un préschéma quasi-compact et quasi-séparé, $f:X\to S$ un morphisme surjectif localement de présentation finie. Alors il existe une famille finie $(S_i)_{i\in I}$

\oldpage[IV-4]{107}

de sous-préschémas affines de $S$, de présentation finie sur $S$, deux à deux disjoints, de réunion $S$, et ayant la propriété suivante: pour chaque $i\in I$, il existe deux morphismes finis de présentation finie et surjectifs
\[
  S_i''\xrightarrow{g_i}S_i'\xrightarrow{h_i}S_i,
\]
où $h_i$ est étale et $g_i$ plat et radiciel, et un $S$-morphisme $S_i''\to X$ (autrement dit une $S_i''$-section de $X_i''=X\times_S S_i''$).
\end{proposition}

Il y a un recouvrement fini $(U_i)_{1\leq i\leq n}$ de $S$ par des ouverts affines; pour tout $i$, l'intersection
\[
  W_i=U_i\cap\left(\bigcup_{j<i}U_j\right)
\]
est quasi-compacte (I, 2.7) donc il existe un sous-préschéma fermé $V_i$ de $U_i$ ayant pour espace sous-jacent $U_i-W_i$ et défini par un idéal de type fini de l'anneau de $U_i$; donc $V_i$ est un sous-préschéma affine de $S$ qui est de présentation finie sur $S$ (I, 6.2). Il est clair qu'on peut se borner à prouver la proposition en y remplaçant $S$ par $V_i$ et $X$ par $X\times_S V_i$. Autrement dit, on peut déjà supposer $S$ affine. D'autre part, $X$ est réunion d'ouverts affines $W_\alpha$ et les $f(W_\alpha)$ sont constructibles dans $S$ (I, 8.4) et forment un recouvrement de $S$; par suite (I, 9.9) il existe une sous-famille finie $(W_{\alpha_j})_{1\leq j\leq m}$ telle que les $f(W_{\alpha_j})$ forment déjà un recouvrement de $S$. Soit alors $X'$ le préschéma somme des sous-préschémas induits sur les ouverts $W_{\alpha_j}$ de $X$; il est immédiat qu'il suffit de prouver la proposition en remplaçant $X$ par $X'$ puisqu'un $S$-morphisme $S_i''\to X'$ donne par composition un $S$-morphisme $S_i''\to X'\to X$.

On peut donc supposer $X$ affine et $f$ de présentation finie. On peut alors (8.9.1) écrire $X$ sous la forme $X_0\times_{S_0}S$, où $S_0$ est noethérien, $X_0\to S_0$ un morphisme de type fini, qu'on peut supposer surjectif (8.10.5, (vi)). Si la proposition est prouvée pour ce morphisme, elle en résultera aussitôt pour $X$ par changement de base $S\to S_0$. On peut donc supposer de plus $S$ noethérien et $f$ de type fini.

Soient $s_h$ $(1\leq h\leq r)$ les points maximaux de $S$. Montrons qu'il suffit de prouver l'énoncé en remplaçant $S$ par un voisinage ouvert assez petit $U_h$ de chacun des $s_h$, $X$ par l'image réciproque de $U_h$ dans $X$. En effet, supposons la proposition établie dans ce cas, et raisonnons par récurrence noethérienne ($0_{\mathrm I}$, 2.2.2) en supposant l'énoncé établi pour tout sous-préschéma fermé de $S$ ayant un espace sous-jacent $\neq S$. On peut supposer les $U_h$ deux à deux sans point commun; si $T$ est un sous-préschéma fermé ayant pour espace sous-jacent $S-(\bigcup_h U_h)$, l'hypothèse de récurrence entraîne que l'énoncé est vrai pour $T$; comme il est vrai aussi pour chacun des $U_h$, il l'est évidemment pour $S$.

Comme on peut évidemment (en remplaçant $S$ par $S_{\mathrm{red}}$ et $X$ par $X\times_S S_{\mathrm{red}}$) supposer $S$ réduit, chacun des $\sh{O}_{S,s_h}$ est un corps $k(s_h)$. Supposons que l'on ait prouvé la proposition lorsque $S$ est le spectre d'un corps et $f$ est de type fini. Alors l'existence des $U_h$ résulte de la méthode de (8.1.2, $a$)) et de (8.8.2, (i) et (ii)), (8.10.5, (vi), (vii) et (x)), (11.2.6, (ii)) et (17.7.8, (ii)).

Supposons donc $S=\Spec(k)$, où $k$ est un corps, $X$ étant de type fini sur $k$ et $X\neq\varnothing$. Comme $X$ est noethérien, il existe dans $X$ un point fermé $x$ ($0_{\mathrm I}$, 2.1.3), donc $k(x)$ est une extension finie de $k$ (I, 6.4.2). Il y a par suite une extension finie séparable $k'$ de $k$ telle que $k''=k(x)$ soit extension finie radicielle de $k'$. On répond alors à la question en prenant $I$ réduit à un élément, $S_i'=\Spec(k')$, $S_i''=\Spec(k'')$.

\oldpage[IV-4]{108}

\begin{corollary}[17.16.5]
\label{IV.17.16.5}
Soit $f:X\to S$ un morphisme surjectif localement de présentation finie. Alors il existe un morphisme $g:S'\to S$ surjectif, localement de présentation finie et localement quasi-fini, tel qu'il existe un $S$-morphisme $S'\to X$ (autrement dit tel qu'il existe une $S'$-section de $X'=X\times_S S'$). Si $S$ est quasi-compact (resp. quasi-compact et quasi-séparé), on peut supposer $S'$ affine (resp. $S'$ affine et le morphisme $g$ de présentation finie et quasi-fini).
\end{corollary}

Il suffit de prouver le corollaire en supposant $S$ affine: $S$ est en effet réunion d'une famille $(S_\alpha)$ d'ouverts affines, et si pour chaque $\alpha$, $S'_\alpha$ est affine et le morphisme $g_\alpha:S'_\alpha\to S_\alpha$ répond à la question et est de présentation finie et quasi-fini, alors en prenant pour $S'$ le préschéma somme des $S'_\alpha$, $g:S'\to S$ coïncidant avec $g_\alpha$ dans chacun des $S'_\alpha$, ce morphisme répond à la question; en outre, si $S$ est quasi-compact, on peut supposer la famille $(S_\alpha)$ finie, donc $S'$ affine; si de plus $S$ est quasi-séparé, les immersions $S_\alpha\to S$ sont de présentation finie (I, 6.2), donc il en est de même de $g$.

Considérons donc le cas où $S$ est affine: on forme alors les morphismes finis et de présentation finie $S_i''\to S_i$ de (17.16.4); comme les immersions $S_i\to S$ sont de présentation finie (les $S_i$ étant affines), les morphismes $g_i'':S_i''\to S_i\to S$ sont de présentation finie et quasi-finis, et on répond à la question en prenant $S'$ égal à la somme des $S_i''$ et $g$ coïncidant avec $g_i''$ dans chacun des $S_i''$.

\begin{proposition}[17.16.6]
\label{IV.17.16.6}
Soient $S$ un préschéma quasi-compact et quasi-séparé, $f:X\to S$ un morphisme de présentation finie tel que, pour tout $s\in S$, $X_s=f^{-1}(s)$ soit un préschéma propre sur $k(s)$. Alors il existe une famille finie $(S_i)_{i\in I}$ de sous-préschémas affines de $S$, de présentation finie sur $S$, deux à deux disjoints, de réunion $S$ et tels que, pour tout $i$, le morphisme $X_i=X\times_S S_i\to S_i$ soit propre et plat. Si de plus, pour tout $s\in S$, $X_s$ est fini sur $k(s)$, on peut prendre les $S_i$ tels que chacun des morphismes $X_i\to S_i$ se factorise en
\[
  X_i\xrightarrow{u_i}X_i'\xrightarrow{h_i}S_i
\]
où $u_i$ et $h_i$ sont finis et localement libres (18.2.7), $h_i$ est étale, $u_i$ est radiciel et surjectif.
\end{proposition}

On suit une marche analogue à celle de (17.16.4). On se ramène d'abord au cas où $S$ est affine et où $f$ est plat, en utilisant le théorème de platitude générique (8.9.5) (on observera que les sous-préschémas $S_i'$ définis dans la démonstration de (8.9.5) sont affines). Puisque $f$ est de présentation finie, on peut ensuite écrire $X=X_0\times_{S_0}S$, où $S_0$ est noethérien, $X_0\to S_0$ un morphisme de type fini et plat (11.2.7); en outre chaque fibre $(X_0)_{s_0}$ est encore propre sur $k(s_0)$, comme il résulte de (2.7.1, (vii)) et on est donc ramené au cas où $S$ est noethérien. Par récurrence noethérienne et application du procédé de (8.1.2, $a$)) (en utilisant aussi cette fois (8.10.5, (xii))), on est finalement ramené, pour prouver la première assertion, au cas où $S$ est le spectre d'un corps, qui résulte trivialement de l'hypothèse (avec $S_i=S$). On traite de même le cas où $X_s$ est supposé fini sur $k(s)$ pour tout $s\in S$ (il faut cette fois utiliser (2.7.1, (xv)), (8.10.5, (x), (iv), (vi) et (vii)), (17.7.8) et (2.1.12)), et on est ramené à prouver la dernière assertion de l'énoncé lorsque $S=\Spec(k)$ est le spectre d'un corps. Mais alors $X=\Spec(A)$, où $A$ est une $k$-algèbre finie (I, 6.4.4); donc $A$ est composée directe de $k$-algèbres finies $A_j$ qui sont des anneaux locaux $(1\leq j\leq m)$; puisque $A_j$ est un anneau artinien et une

\oldpage[IV-4]{109}

$k$-algèbre, il contient un sous-corps $K'_j$ isomorphe canoniquement au corps résiduel de $A_j$ (0, 19.6.2). On prend alors $X_i=X$, $X'_i$ somme des $\Spec(K'_j)$ et il est clair que l'on répond ainsi à la question puisque pour tout $j$ on a deux homomorphismes $K'_j\to A_j\to k(A_j)$ dont le composé est un isomorphisme.

% FIN SOURCE ega4/ega4-17-fr.tex

%
% DEBUT SOURCE ega4/ega4-18-fr.tex
\section[Compléments sur les morphismes étales. Anneaux locaux henséliens et anneaux strictement locaux]{Compléments sur les morphismes étales\\Anneaux locaux henséliens et anneaux strictement locaux}
\label{section:IV.18}

Dans le présent paragraphe, nous étudions diverses propriétés spéciales aux morphismes étales. De plus, la notion de morphisme étale permet de développer de façon fort naturelle la théorie de Nagata des anneaux henséliens, ainsi que celle des anneaux strictement locaux. Ces anneaux jouent un rôle important dans de nombreux développements récents, en s'introduisant chaque fois que l'on a besoin d'un procédé de « localisation » plus fin que celui fourni par la topologie de Zariski (cf. par exemple [43], en attendant la parution du chapitre de notre Traité consacré à l'étude de la « topologie étale »).

\subsection{Une équivalence remarquable de catégories.}
\EGAReaderTarget{ega4.c04.subsec.18-1}{subsection:IV.18.1}

\begin{proposition}[18.1.1]
\label{IV.18.1.1}
Soient $S$ un préschéma, $S_0$ un sous-préschéma fermé de $S$, $X_0$ un $S_0$-préschéma lisse (resp. étale) sur $S_0$, $x_0$ un point de $X_0$. Il existe alors un voisinage ouvert $U_0$ de $x_0$ dans $X_0$, un $S$-préschéma $U$ lisse (resp. étale) sur $S$ et un $S_0$-isomorphisme
\[
  U\times_S S_0\xrightarrow{\sim}U_0.
\]
\end{proposition}

Notons que si $X_0$ est étale sur $S_0$ au point $x_0$, il est a fortiori non ramifié sur $S_0$ en ce point; si l'on a construit un $S$-préschéma $U$ lisse sur $S$ tel que $U\times_S S_0$ soit isomorphe à $U_0$, comme les fibres des morphismes $U_0\to S_0$ et $U\to S$ contenant $x_0$ sont alors isomorphes, il en résultera que $U$ est non ramifié sur $S$ au point $x_0$ (17.4.1, $d$)), donc aussi dans un voisinage de $x_0$; en remplaçant $U$ par ce voisinage, on en conclut que $U$ sera étale sur $S$. Il suffit donc de prouver la proposition lorsqu'on suppose seulement $X_0$ lisse sur $S_0$.

La question étant locale sur $S$ et sur $X_0$, on peut supposer que $S=\Spec(A)$ et $X_0=\Spec(C_0)$ sont affines, de sorte que $S_0=\Spec(A_0)$, où $A_0$ est un anneau quotient de $A$, $C_0=B_0/\mathfrak J_0$, où $B_0=A_0[T_1,\ldots,T_n]$ et $\mathfrak J_0$ est un idéal de type fini de $B_0$; enfin, $C_0$ est une $A_0$-algèbre formellement lisse pour les topologies discrètes. Soit $\mathfrak p_0$ l'idéal $\mathfrak j_{x_0}$ dans $C_0$; on a $\mathfrak p_0=\mathfrak q_0/\mathfrak J_0$, où $\mathfrak q_0$ est un idéal premier de $B_0$. Le critère jacobien (0, 22.6.4) joint à (0, 19.1.12) montre qu'il existe dans $\mathfrak J_0$ une famille de $r$ polynômes $u_i$ $(1\leq i\leq r)$ et $r$ indices $j_h$ $(1\leq h\leq r)$ tels que les images des $u_i$ dans $(\mathfrak J_0)_{\mathfrak q_0}/(\mathfrak J_0^2)_{\mathfrak q_0}$ engendrent ce $(B_0)_{\mathfrak q_0}$-module et que l'on ait
\begin{equation}
  \det\left(\frac{\partial u_i}{\partial T_{j_h}}\right)\notin\mathfrak q_0.
  \tag{18.1.1.1}\label{IV.18.1.1.1}
\end{equation}

Comme $(B_0)_{\mathfrak q_0}$ est un anneau local, il résulte du lemme de Nakayama que l'on peut supposer que les images des $u_i$ dans $(\mathfrak J_0)_{\mathfrak q_0}$ engendrent ce $(B_0)_{\mathfrak q_0}$-module, puis, en remplaçant au besoin $X_0$ par un voisinage ouvert affine de $x_0$, que les $u_i$ engendrent $\mathfrak J_0$ ($0_{\mathrm I}$, 5.2.2). Posons alors $B=A[T_1,\ldots,T_n]$; $B_0$ est donc un anneau quotient

\oldpage[IV-4]{110}

de $B$, $\mathfrak q_0$ est l'image d'un idéal premier $\mathfrak q$ de $B$ et $\mathfrak q$ l'image réciproque de $\mathfrak q_0$. Pour tout $i$, soit $v_i\in B$ un élément dont $u_i$ est l'image dans $B_0$, et soit $\mathfrak J$ l'idéal de $B$ engendré par les $v_i$, de sorte que $\mathfrak J_0$ est l'image de $\mathfrak J$ dans $B_0$. La proposition sera établie en prenant pour $U$ un voisinage ouvert du point de $\Spec(B/\mathfrak J)$ correspondant à l'idéal premier $\mathfrak p=\mathfrak q/\mathfrak J$, pourvu que l'on prouve que $B_{\mathfrak q}/\mathfrak J_{\mathfrak q}$ est une $A$-algèbre formellement lisse pour les topologies discrètes. Or, cela résulte du critère jacobien, car les images des $v_i$ dans $\mathfrak J_{\mathfrak q}/\mathfrak J_{\mathfrak q}^2$ engendrent ce $B_{\mathfrak q}$-module, et il résulte de (18.1.1.1) que l'on a $\det(\partial v_i/\partial T_{j_h})\notin\mathfrak q$.

\begin{theorem}[18.1.2]
\label{IV.18.1.2}
Soient $S$ un préschéma, $S_0$ un sous-préschéma fermé de $S$ dont l'espace sous-jacent est identique à celui de $S$. Alors le foncteur
\[
  X\longmapsto X\times_S S_0
\]
de la catégorie des $S$-préschémas étales sur $S$, dans la catégorie des $S_0$-préschémas étales sur $S_0$, est une équivalence de catégories.
\end{theorem}

Montrons d'abord que ce foncteur est pleinement fidèle. Soient $X$, $Y$ deux $S$-préschémas étales sur $S$, et posons $X_0=X\times_S S_0$, $Y_0=Y\times_S S_0$. Si $Z=X\times_S Y$, l'ensemble $\Hom_S(X,Y)$ est en correspondance biunivoque canonique avec l'ensemble des $X$-sections $\Gamma(Z/X)$, et de même $\Hom_{S_0}(X_0,Y_0)$ est en correspondance biunivoque canonique avec $\Gamma(Z_0/X_0)$, où $Z_0=Z\times_S S_0=X_0\times_{S_0}Y_0$. Or $Z$ est étale sur $X$, $Z_0$ étale sur $X_0$, et $X_0$ (resp. $Z_0$) est un sous-préschéma fermé de $X$ (resp. $Z$) ayant même espace sous-jacent. Les parties ouvertes de $Z$ telles que la restriction du morphisme $Z\to X$ soit surjective et radicielle sont donc les mêmes que les parties de $Z_0$ ayant les propriétés correspondantes et notre assertion découle par suite de (17.9.3).

Pour achever la démonstration, il suffit de voir que pour tout $S_0$-préschéma $X_0$ étale sur $S_0$, il existe un $S$-préschéma $X$ étale sur $S$ et un $S_0$-isomorphisme $X_0\xrightarrow{\sim}X\times_S S_0$. En vertu de la prop. (18.1.1), il existe un recouvrement ouvert $(U_\alpha)$ de $X_0$ et pour tout $\alpha$, un $S$-préschéma $V_\alpha$ qui est étale sur $S$, et enfin un $S_0$-isomorphisme $\theta_\alpha:U_\alpha\xrightarrow{\sim}V_\alpha\times_S S_0$. En outre, en raison de la première partie de la démonstration, il existe un unique $S_0$-isomorphisme $\varphi_{\alpha\beta}$ de $\theta_\alpha(U_\alpha\cap U_\beta)$ sur $\theta_\beta(U_\alpha\cap U_\beta)$, correspondant à l'automorphisme identique de $U_\alpha\cap U_\beta$, et il est immédiat, pour la même raison, que ces isomorphismes vérifient la condition de recollement ($0_{\mathrm I}$, 4.1.7). Il y a par suite un $S$-préschéma $X$ tel que les $V_\alpha$ s'identifient canoniquement à des sous-préschémas induits sur des ouverts de $X$, les $\theta_\alpha$ s'identifiant à des $S_0$-isomorphismes qui coïncident dans les intersections $U_\alpha\cap U_\beta$, et définissent donc un $S_0$-isomorphisme $X_0\xrightarrow{\sim}X\times_S S_0$. Il est clair que $X$ est étale sur $S$ (17.3.2), ce qui achève la démonstration.

\begin{corollary}[18.1.3]
\label{IV.18.1.3}
Soient $S$ un préschéma, $X$ un $S$-préschéma étale sur $S$, $S'$ un $S$-préschéma, $S'_0$ un sous-préschéma fermé de $S'$ ayant même espace sous-jacent. Alors l'application canonique $X(S')_S\to X(S'_0)_S$ (I, 3.4.3) est bijective.
\end{corollary}

Posons en effet $X'=X\times_S S'$, $X'_0=X\times_S S'_0$, de sorte que $X(S')_S=\Gamma(X'/S')$ et $X(S'_0)_S=\Gamma(X'_0/S'_0)$; le corollaire résulte du fait que le foncteur défini dans (18.1.2) (avec $S$ et $S_0$ remplacés par $S'$ et $S'_0$ respectivement) est pleinement fidèle (ou directement de (17.9.3)).

\subsection{Revêtements étales.}
\EGAReaderTarget{ega4.c04.subsec.18-2}{subsection:IV.18.2}

\oldpage[IV-4]{111}

\begin{env}[18.2.1]
\EGAReaderTarget{ega4.c04.statement.18-2-1}{IV.18.2.1}
Étant donnés un anneau $A$ et une $A$-algèbre commutative $B$ qui est \emph{finie} et est un $A$-module \emph{libre}, rappelons (Bourbaki, \emph{Alg.}, chap. VIII, \textsection~12, n\textsuperscript{o}~2) que l'on définit sur $B$ une \emph{forme $A$-linéaire} $\operatorname{Tr}_{B/A}$, la « forme trace »; on en déduit la définition d'une \emph{forme $A$-bilinéaire symétrique} (dite encore « forme trace »)
\begin{equation}
  (x,y)\longmapsto\operatorname{Tr}_{B/A}(xy)
  \tag{18.2.1.1}\label{IV.18.2.1.1}
\end{equation}
dont la donnée est équivalente à celle de l'application $A$-linéaire associée $\operatorname{astr}_{B/A}:B\to\check B$ du $A$-module $B$ dans son \emph{dual} $\check B$, égale à sa transposée. Lorsque $A$ est un \emph{corps}, il est équivalent de dire que cette forme bilinéaire est \emph{non dégénérée} ou que $B$ est une $A$-algèbre \emph{séparable} (Bourbaki, \emph{Alg.}, chap. IX, \textsection~2, prop. 5).

Soit $f:A\to A'$ un homomorphisme d'anneaux; posons $B'=B\otimes_A A'$, et soit $g:B\to B'$ l'homomorphisme canonique; l'image par $g$ d'une base du $A$-module $B$ est alors une base du $A'$-module $B'$, et il résulte des définitions que l'on a, pour tout $x\in B$,
\begin{equation}
  \operatorname{Tr}_{B'/A'}(g(x))=f(\operatorname{Tr}_{B/A}(x)).
  \tag{18.2.1.2}\label{IV.18.2.1.2}
\end{equation}
\end{env}

\begin{env}[18.2.2]
\EGAReaderTarget{ega4.c04.statement.18-2-2}{IV.18.2.2}
Considérons maintenant un \emph{espace annelé} $(X,\sh{O}_X)$ et soit $\sh{B}$ une $\sh{O}_X$-Algèbre qui, en tant que $\sh{O}_X$-Module, est \emph{localement libre de rang fini}; alors, pour tout ouvert $U\subset X$ tel que $\sh{B}|U$ soit (en tant que $\sh{O}_U$-Module) isomorphe à $\sh{O}_U^n$ (pour un $n$ dépendant de $U$), $\Gamma(U,\sh{B})$ est une $\Gamma(U,\sh{O}_X)$-algèbre qui, en tant que $\Gamma(U,\sh{O}_X)$-module, est libre de rang fini et définit donc une forme $\Gamma(U,\sh{O}_X)$-linéaire $\operatorname{Tr}_{\Gamma(U,\sh{B})/\Gamma(U,\sh{O}_X)}$, que nous noterons aussi $\operatorname{Tr}_{\sh{B}|\sh{O}_X,U}$; on en déduit une application linéaire associée
\[
  \operatorname{astr}_{\sh{B}|\sh{O}_X,U}:\Gamma(U,\sh{B})\longrightarrow\Gamma(U,\sh{B})^\vee=\Gamma(U,\check{\sh{B}}).
\]
En outre, il résulte de (18.2.1.2) que ces applications linéaires sont compatibles avec les opérations de restriction de $U$ à un ouvert plus petit, et définissent donc, d'une part un \emph{homomorphisme de $\sh{O}_X$-Modules}, appelé encore \emph{homomorphisme trace}:
\begin{equation}
  \operatorname{Tr}_{\sh{B}|\sh{O}_X}:\sh{B}\longrightarrow\sh{O}_X
  \tag{18.2.2.1}\label{IV.18.2.2.1}
\end{equation}
et d'autre part un \emph{homomorphisme de $\sh{O}_X$-Modules}
\begin{equation}
  \operatorname{astr}_{\sh{B}|\sh{O}_X}:\sh{B}\longrightarrow\check{\sh{B}}
  \tag{18.2.2.2}\label{IV.18.2.2.2}
\end{equation}
dit \emph{associé à la trace}, et égal à son transposé. Il résulte aussi de (18.2.1.2) que pour tout $x\in X$, on a
\begin{equation}
  (\operatorname{Tr}_{\sh{B}|\sh{O}_X})_x=\operatorname{Tr}_{\sh{B}_x|\sh{O}_{X,x}}
  \tag{18.2.2.3}\label{IV.18.2.2.3}
\end{equation}
\begin{equation}
  (\operatorname{astr}_{\sh{B}|\sh{O}_X})_x=\operatorname{astr}_{\sh{B}_x|\sh{O}_{X,x}}.
  \tag{18.2.2.4}\label{IV.18.2.2.4}
\end{equation}
Enfin, sous les conditions de (18.2.1), si l'on pose $X=\Spec(A)$ et si $\sh{B}=\widetilde B$ est la $\sh{O}_X$-Algèbre correspondant à $B$, la forme $\operatorname{Tr}_{\sh{B}|\sh{O}_X}$ (resp. l'homomorphisme $\operatorname{astr}_{\sh{B}|\sh{O}_X}$) correspond à la forme $\operatorname{Tr}_{B/A}$ (resp. à l'homomorphisme de $A$-modules $\operatorname{astr}_{B/A}$), comme il résulte encore de (18.2.1.2).
\end{env}


\oldpage[IV-4]{112}

\begin{proposition}[18.2.3]
\EGAReaderTarget{ega4.c04.prop.18-2-3}{IV.18.2.3}
Soit $f:X\to Y$ un morphisme fini de préschémas et soit $\sh{B}=f_*(\sh{O}_X)$. Les conditions suivantes sont équivalentes:
\begin{enumerate}
  \item[a)] $f$ est étale.
  \item[a$'$)] $f$ est un morphisme plat de présentation finie, et pour tout $x\in X$, si l'on pose $y=f(x)$, $\sh{O}_{X,x}/\mathfrak m_y\sh{O}_{X,x}$ est un corps, extension finie séparable de $\kappa(y)$.
  \item[b)] $\sh{B}$ est un $\sh{O}_Y$-Module localement libre et pour tout $y\in Y$, $\sh{B}_y\otimes_{\sh{O}_{Y,y}}\kappa(y)$ est une $\kappa(y)$-algèbre finie séparable (donc le composé direct d'un nombre fini de corps, extensions finies séparables de $\kappa(y)$).
  \item[c)] $\sh{B}$ est un $\sh{O}_Y$-Module localement libre, et l'homomorphisme $\operatorname{astr}_{\sh{B}|\sh{O}_Y}:\sh{B}\to\check{\sh{B}}$ (18.2.2) est bijectif.
\end{enumerate}
\end{proposition}

Compte tenu de ce que $f$ est quasi-compact, l'équivalence de $a)$ et $a')$ a déjà été démontrée (17.6.2). Pour prouver le reste de la proposition, on peut se borner au cas où $Y=\Spec(A)$ et $X=\Spec(B)$ sont affines, $B$ étant une $A$-algèbre finie et $\sh{B}=\widetilde B$. Dire que $f$ est un morphisme de présentation finie équivaut alors à dire que $B$ est un $A$-module de présentation finie (1.4.7). Si en outre $f$ est plat, donc $B$ un $A$-module plat, on sait (Bourbaki, \emph{Alg. comm.}, chap. II, \textsection~5, n\textsuperscript{o}~2, cor. 2 du th. 1) que $B$ est un $A$-module projectif, donc $\sh{B}$ un $\sh{O}_Y$-Module localement libre (\emph{loc. cit.}, n\textsuperscript{o}~2, th. 1), et la réciproque est immédiate. D'autre part $f^{-1}(y)$ n'est autre que le spectre de la $\kappa(y)$-algèbre $\sh{B}(y)=\sh{B}_y\otimes_{\sh{O}_{Y,y}}\kappa(y)$, ce qui achève de prouver l'équivalence de $a')$ et $b)$. Pour voir que $b)$ est équivalente à $c)$, notons que la seconde assertion de $b)$ équivaut au fait que l'homomorphisme $\operatorname{astr}_{\sh{B}(y)|\kappa(y)}:\sh{B}(y)\to\check{\sh{B}(y)}$ est bijectif; comme $\sh{B}(y)=\sh{B}_y/\mathfrak m_y\sh{B}_y$ et $\check{\sh{B}(y)}=\check{\sh{B}}_y/\mathfrak m_y\check{\sh{B}}_y$ et que $\sh{B}_y$ et $\check{\sh{B}}_y$ sont des $\sh{O}_{Y,y}$-modules libres, il résulte de (18.2.2.4) et de Bourbaki, \emph{Alg. comm.}, chap. II, \textsection~3, n\textsuperscript{o}~3, cor. de la prop. 6, que l'homomorphisme $\operatorname{astr}_{\sh{B}_y|\sh{O}_{Y,y}}:\sh{B}_y\to\check{\sh{B}}_y$ est lui aussi bijectif; la réciproque étant évidente, cela achève la démonstration.

Lorsqu'une $\sh{O}_X$-Algèbre $\sh{B}$ vérifie les conditions équivalentes $b)$ et $c)$ de (18.2.3), on dit que $\sh{B}$ est une $\sh{O}_X$-Algèbre \emph{finie et étale}. Lorsque $X=\Spec(A)$ est affine et que l'on a donc $\sh{B}=\widetilde B$, où $B$ est une $A$-algèbre, il revient au même, en vertu de (18.2.3), de dire que $\sh{B}$ est une $\sh{O}_X$-Algèbre finie et étale, ou que $B$ est une $A$-algèbre finie et étale (au sens de (17.3.2)).

\begin{corollary}[18.2.4]
\label{IV.18.2.4}
Soit $f:X\to Y$ un morphisme fini et de présentation finie, et posons $\sh{B}=f_*(\sh{O}_X)$. Soit $y$ un point de $Y$. Les conditions suivantes sont équivalentes:
\begin{enumerate}
  \item[a)] Il existe un voisinage ouvert $U$ de $y$ dans $Y$ tel que la restriction $f^{-1}(U)\to U$ de $f$ soit un morphisme étale.
  \item[b)] $\sh{B}_y$ est un $\sh{O}_{Y,y}$-module libre de type fini et $\sh{B}_y\otimes_{\sh{O}_{Y,y}}\kappa(y)$ est une $\kappa(y)$-algèbre séparable.
\end{enumerate}
\end{corollary}

Il est clair que $a)$ entraîne $b)$ en vertu de (18.2.3). D'autre part, $\sh{B}$ est un $\sh{O}_Y$-Module de présentation finie (1.6.3 et 1.4.7), donc, si $\sh{B}_y$ est un $\sh{O}_{Y,y}$-module libre, il existe un voisinage ouvert $U$ de $y$ dans $Y$ tel que $\sh{B}|U$ soit un $\sh{O}_U$-Module localement libre ($0_{\mathrm I}$, 5.2.7); en outre, par hypothèse, l'homomorphisme $\operatorname{astr}_{\sh{B}_y|\sh{O}_{Y,y}}:\sh{B}_y\to\check{\sh{B}}_y$ étant bijectif, il résulte aussi de ($0_{\mathrm I}$, 5.2.7) que l'on peut supposer $U$ choisi de sorte que

\oldpage[IV-4]{113}

l'homomorphisme $\operatorname{astr}_{\sh{B}|U/\sh{O}_U}$ soit bijectif. Le fait que $b)$ implique $a)$ résulte alors de (18.2.3).

\begin{corollary}[18.2.5]
\label{IV.18.2.5}
Soient $Y$ un préschéma quasi-compact ou localement noethérien, $f:X\to Y$ un morphisme fini et de présentation finie; posons $\sh{B}=f_*(\sh{O}_X)$. Supposons que pour tout point fermé $y$ de $Y$, $\sh{B}_y$ soit un $\sh{O}_{Y,y}$-module libre et $\sh{B}_y\otimes_{\sh{O}_{Y,y}}\kappa(y)$ une $\kappa(y)$-algèbre séparable. Alors $f$ est étale.
\end{corollary}

En effet, il résulte de (18.2.4) que tout point fermé $y$ de $Y$ possède un voisinage ouvert $U$ tel que la restriction $f^{-1}(U)\to U$ de $f$ soit étale; la conclusion résulte de ce que dans les deux cas considérés, toute partie fermée non vide de $Y$ contient un point fermé (5.1.11 et $0_{\mathrm I}$, 2.1.3).

\begin{corollary}[18.2.6]
\label{IV.18.2.6}
Si $f:X\to Y$ est un morphisme fini et étale, et si l'on pose $\sh{B}=f_*(\sh{O}_X)$, alors le $\sh{O}_Y$-homomorphisme $\operatorname{Tr}_{\sh{B}|\sh{O}_Y}:\sh{B}\to\sh{O}_Y$ (18.2.2) (aussi noté $\operatorname{Tr}_f$) est surjectif.
\end{corollary}

La question étant locale, on peut, en vertu de (18.2.3) supposer que $Y=\Spec(A)$, $X=\Spec(B)$ avec $\sh{B}=\widetilde B$, $B$ étant un $A$-module libre; comme en vertu de (18.2.3) la forme bilinéaire (18.2.1.1) est non dégénérée, cela entraîne en particulier que la forme linéaire $\operatorname{Tr}_{B/A}$ est surjective.

\begin{remarks}[18.2.7]
\EGAReaderTarget{ega4.c04.item.18-2-7}{IV.18.2.7}
\begin{enumerate}
  \item[(i)] Lorsque $f:X\to Y$ est un morphisme fini tel que $f_*(\sh{O}_X)$ soit un $\sh{O}_Y$-Module localement libre (resp. localement libre de rang $n$), on dit encore que $f$ est un morphisme \emph{fini localement libre} (resp. \emph{localement libre de rang $n$}). Cette condition, en vertu de (18.2.3), est vérifiée si $f$ est un morphisme fini étale, mais n'implique pas à elle seule que $f$ soit étale, comme le montre l'exemple où $X=\Spec(K)$ et $Y=\Spec(k)$ sont des spectres de corps, $K$ étant une extension finie non séparable de $k$. Lorsque $f:X\to Y$ est un morphisme fini étale, on dit encore que $X$ est un \emph{revêtement étale} de $Y$.

On notera que dans ce cas, $f$ est universellement ouvert et universellement fermé, et en particulier $f(X)$ est une partie à la fois ouverte et fermée de $Y$.

On dit qu'un revêtement étale $X$ de $Y$ est \emph{trivial} si $X$ est \emph{somme} d'un nombre fini de préschémas \emph{isomorphes} à $Y$. On dit qu'un revêtement étale $X$ de $Y$ est \emph{localement trivial} si le morphisme $f:X\to Y$ est tel que tout point $y\in Y$ ait un voisinage ouvert $U$ pour lequel le revêtement $f^{-1}(U)$ de $U$ soit trivial.

  \item[(ii)] Soit $f:X\to Y$ un morphisme fini, localement libre de rang $n$; posons $\sh{B}=f_*(\sh{O}_X)$ et soit $u=\operatorname{astr}_{\sh{B}|\sh{O}_Y}:\sh{B}\to\check{\sh{B}}$; on en déduit un homomorphisme puissance extérieure $n$-ème
\[
  \bigwedge^n u:\bigwedge^n\sh{B}\longrightarrow\bigwedge^n\check{\sh{B}}=(\bigwedge^n\sh{B})^\vee
\]
de $\sh{O}_Y$-Modules \emph{inversibles}, et par suite ($0_{\mathrm I}$, 5.4.2) un élément
\begin{equation}
  d_{X/Y}\in\Gamma\left(Y,(\bigwedge^n\check{\sh{B}})\otimes_{\sh{O}_Y}(\bigwedge^n\check{\sh{B}})\right)
  \tag{18.2.7.1}\label{IV.18.2.7.1}
\end{equation}
que l'on appelle le \emph{discriminant de $X$ sur $Y$}. D'ailleurs, comme $(\bigwedge^n\check{\sh{B}})\otimes_{\sh{O}_Y}(\bigwedge^n\check{\sh{B}})$ est le dual de $(\bigwedge^n\sh{B})\otimes_{\sh{O}_Y}(\bigwedge^n\sh{B})$, $d_{X/Y}$ peut aussi être identifié à un homomorphisme
\begin{equation}
  (\bigwedge^n\sh{B})\otimes_{\sh{O}_Y}(\bigwedge^n\sh{B})\longrightarrow\sh{O}_Y
  \tag{18.2.7.2}\label{IV.18.2.7.2}
\end{equation}
et on note $\mathscr{D}_{X/Y}$ l'Idéal quasi-cohérent de type fini de $\sh{O}_Y$, \emph{image} de l'homomorphisme (18.2.7.2), encore appelé \emph{Idéal discriminant de $X$ sur $Y$}.
\end{enumerate}
\end{remarks}

\oldpage[IV-4]{114}

Cela étant, pour que l'homomorphisme $u$ soit bijectif, il faut et il suffit que $\bigwedge^n u$ soit bijectif, ou encore que la section $d_{X/Y}$ ait un germe \emph{inversible} en tout point $y\in Y$, ce qui s'écrit aussi $d_{X/Y}(y)\neq 0$ pour tout $y$ ($0_{\mathrm I}$, 5.5.2). Il revient aussi au même de dire que l'Idéal discriminant $\mathscr{D}_{X/Y}$ est \emph{égal à $\sh{O}_Y$}.

La terminologie de « revêtement » introduite dans (18.2.7, (i)), se justifie par la proposition suivante:

\begin{proposition}[18.2.8]
\label{IV.18.2.8}
Soit $f:X\to Y$ un morphisme étale, séparé et de type fini et pour tout $y\in Y$, soit $n(y)$ le nombre géométrique des points de $f^{-1}(y)$. Alors la fonction $y\mapsto n(y)$ est semi-continue inférieurement dans $Y$. Pour qu'elle soit continue en un point $y$ (donc constante dans un voisinage de $y$), il faut et il suffit qu'il existe un voisinage ouvert $U$ de $y$ tel que la restriction $f^{-1}(U)\to U$ de $f$ soit un morphisme fini (étale).
\end{proposition}

Comme $f$ est quasi-fini (17.6.1) et localement de présentation finie, il revient au même de dire que $f$ est fini ou que $f$ est propre (8.11.1): en outre, chaque fibre $f^{-1}(y)$ est géométriquement réduite sur $\kappa(y)$. Les conclusions résultent donc de (15.5.9, (i) et (ii)) et de ce que $f$ est plat.

\begin{corollary}[18.2.9]
\label{IV.18.2.9}
Soient $Y$ un préschéma connexe, $f:X\to Y$ un morphisme étale, séparé et de type fini. Pour que $f$ soit fini (autrement dit, que $X$ soit un revêtement étale de $Y$), il faut et il suffit que toutes les fibres de $f$ aient même nombre géométrique de points.
\end{corollary}

\begin{remarks}[18.2.10]
\label{IV.18.2.10}
\begin{enumerate}
  \item[(i)] L'exemple de la « droite affine avec un point dédoublé » (I, 5.5.11) montre qu'un morphisme étale de type fini de préschémas noethériens peut être \emph{non séparé}; pour cet exemple, la première assertion de (18.2.8) n'est plus exacte.

  \item[(ii)] Pour qu'un morphisme séparé, de type fini et étale $f:X\to Y$ fasse de $X$ un revêtement \emph{localement trivial}, il faut et il suffit que pour tout $x\in X$ il existe un voisinage ouvert $U$ de $y=f(x)$ et une $U$-section $g$ de $f^{-1}(U)$ telle que $g(y)=x$. En effet, la condition est évidemment nécessaire; le fait qu'elle est suffisante résulte de ce que toute fibre $f^{-1}(y)$ est finie (17.6.1) et (I, 6.2.2), de la caractérisation des sections d'un $Y$-schéma étale (17.9.3) et de la prop. (18.2.8).
\end{enumerate}
\end{remarks}

\subsection{Algèbres finies et étales.}
\EGAReaderTarget{ega4.c04.subsec.18-3}{subsection:IV.18.3}

\begin{proposition}[18.3.1]
\label{IV.18.3.1}
Soient $A$ un anneau, $B$ une $A$-algèbre de présentation finie.
\begin{enumerate}
  \item[(i)] Pour que $B$ soit une $A$-algèbre non ramifiée, il faut et il suffit que $B$ soit un $A$-module de présentation finie et que $B$ soit un $(B\otimes_A B)$-module projectif.
  \item[(ii)] Supposons de plus que $B$ soit une $A$-algèbre finie. Pour que $B$ soit une $A$-algèbre étale, il faut et il suffit que $B$ soit un $A$-module projectif et un $(B\otimes_A B)$-module projectif.
\end{enumerate}
\end{proposition}

Bien entendu, la structure de $(B\otimes_A B)$-module sur $B$ est celle provenant de la structure de $(B\otimes_A B)$-algèbre sur $B$ correspondant au $A$-homomorphisme canonique d'anneaux $B\otimes_A B\to B$, qui est surjectif et de noyau $\mathfrak J_{B/A}$ (0, 20.4.1).

(i) Dire que le morphisme $\Spec(B)\to\Spec(A)$ est localement de présentation finie équivaut à dire que $B$ est une $A$-algèbre de présentation finie (1.4.6). Dire

\oldpage[IV-4]{115}

que $B$ est une $A$-algèbre non ramifiée signifie alors (17.4.2) que $\Spec((B\otimes_A B)/\mathfrak J_{B/A})$ est un sous-schéma induit sur une partie \emph{ouverte et fermée} de $\Spec(B\otimes_A B)$ et on sait que pour qu'il en soit ainsi, il faut et il suffit que $\mathfrak J_{B/A}$ soit un idéal \emph{facteur direct} de $B\otimes_A B$ (Bourbaki, \emph{Alg. comm.}, chap. II, \textsection~4, n\textsuperscript{o}~3, prop. 15); mais il revient au même de dire que le $(B\otimes_A B)$-module quotient $(B\otimes_A B)/\mathfrak J_{B/A}$ est projectif (Bourbaki, \emph{Alg.}, chap. II, 3\textsuperscript{e} éd., \textsection~2, n\textsuperscript{o}~2, prop. 4).

(ii) Si on se rappelle qu'un $A$-module plat de présentation finie est projectif et réciproquement (Bourbaki, \emph{Alg. comm.}, chap. II, \textsection~5, n\textsuperscript{o}~2, cor. 2 du th. 1), l'assertion de (ii) résulte de celle de (i) et de (17.6.2).

\begin{proposition}[18.3.2]
\label{IV.18.3.2}
Soient $A$ un anneau, $\mathfrak J$ un idéal de $A$ tel que, pour la topologie $\mathfrak J$-pré-adique, $A$ soit séparé et complet; on pose $A_0=A/\mathfrak J$. Alors le foncteur
\[
  B\longmapsto B\otimes_A A_0
\]
est une équivalence de la catégorie des $A$-algèbres finies et étales, dans la catégorie des $A_0$-algèbres finies et étales.
\end{proposition}

Nous démontrerons d'abord le lemme suivant:

\begin{lemma}[18.3.2.1]
\label{IV.18.3.2.1}
Soient $A$ un anneau, $\mathfrak J$ un idéal de $A$ tel que, pour la topologie $\mathfrak J$-pré-adique, $A$ soit séparé et complet.
\begin{enumerate}
  \item[(i)] Tout $A$-module projectif de type fini $M$ est séparé et complet pour la topologie $\mathfrak J$-pré-adique, donc limite projective des $(A/\mathfrak J^{n+1})$-modules projectifs $M/\mathfrak J^{n+1}M$.
  \item[(ii)] Inversement, posons $A_n=A/\mathfrak J^{n+1}$, et soit $(M_n)$ un système projectif de $A_n$-modules, tel que, pour tout $n$, l'homomorphisme $M_{n+1}\otimes_{A_{n+1}}A_n\to M_n$ déduit du di-homomorphisme de transition $M_{n+1}\to M_n$ soit bijectif. Supposons en outre les $M_n$ projectifs et $M_0$ de type fini. Alors $M=\varprojlim M_n$ est un $A$-module projectif de type fini tel que l'homomorphisme canonique $M\otimes_A A_0\to M_0$ soit bijectif.
\end{enumerate}
\end{lemma}

(i) Il existe un $A$-module libre de type fini $L$ tel que $M$ soit isomorphe à un facteur direct de $L$; comme $L$ est séparé pour la topologie $\mathfrak J$-pré-adique, il en est de même de tout sous-module $N$ de $L$ puisque $\mathfrak J^{n+1}N\subset\mathfrak J^{n+1}L$; en particulier $M$ est séparé pour cette topologie, et comme l'homomorphisme surjectif $f:L\to M$ est continu pour la topologie $\mathfrak J$-pré-adique, son noyau $N$ est fermé pour la topologie induite par celle de $L$; puisque $L$ est complet et $f$ un morphisme strict, on en conclut que $M$ est complet (Bourbaki, \emph{Top. gén.}, chap. IX, 2\textsuperscript{e} éd., \textsection~3, n\textsuperscript{o}~1, prop. 4).

(ii) Il résulte du lemme de Nakayama que si $M_0$ est engendré par une famille finie $(x_{i0})$ de $r$ éléments et si pour tout $n$, $x_{in}$ est un élément de $M_n$ dont $x_{i,n-1}$ est l'image dans $M_{n-1}$, alors $(x_{in})$ $(1\leq i\leq r)$ est un système de générateurs de $M_n$ (Bourbaki, \emph{Alg. comm.}, chap. II, \textsection~3, cor. 2 de la prop. 4). Cela étant, pour tout $n$, posons $L_n=A_n^r$; si $(e_{in})_{1\leq i\leq r}$ est la base canonique de $L_n$, soit $u_n:L_n\to M_n$ l'application $A_n$-linéaire telle que $u_n(e_{in})=x_{in}$ pour tout $i$. Par hypothèse on a une suite exacte scindée
\[
  0\longrightarrow N_n\xrightarrow{v_n}L_n\xrightarrow{u_n}M_n\longrightarrow 0
\]

\oldpage[IV-4]{116}

et comme $L_n=L_{n+1}/\mathfrak J^{n+1}L_{n+1}$ et $M_n=M_{n+1}/\mathfrak J^{n+1}M_{n+1}$, les flèches verticales dans le diagramme commutatif
\[
\xymatrix{
0\ar[r] & N_{n+1}\ar[r]^{v_{n+1}}\ar[d] & L_{n+1}\ar[r]^{u_{n+1}}\ar[d] & M_{n+1}\ar[r]\ar[d] & 0\\
0\ar[r] & N_n\ar[r]_{v_n} & L_n\ar[r]_{u_n} & M_n\ar[r] & 0
}
\]
sont toutes trois surjectives. Or on a $M=\varprojlim M_n$ et $L=A^r=\varprojlim L_n$; si l'on pose $N=\varprojlim N_n$, $u=\varprojlim u_n$, $v=\varprojlim v_n$, on a, en vertu de ($0_{\mathrm{III}}$, 13.2.2), la suite exacte
\begin{equation}
  0\longrightarrow N\xrightarrow{v}L\xrightarrow{u}M\longrightarrow 0.
  \tag{18.3.2.2}\label{IV.18.3.2.2}
\end{equation}
En outre, comme pour tout $n$, $v_n$ est inversible à gauche et que $M_n$ est un $A_n$-module projectif, il résulte de (0, 19.1.8) que la suite exacte (18.3.2.2) est scindée, ce qui prouve le lemme.

Cela étant, montrons d'abord que le foncteur de l'énoncé de (18.3.2) est \emph{pleinement fidèle}. Posons, comme dans le lemme, $A_n=A/\mathfrak J^{n+1}$; soient $B$, $C$ deux $A$-algèbres finies et étales, et posons, pour tout $n$, $B_n=B\otimes_A A_n$, $C_n=C\otimes_A A_n$; en vertu de (18.3.1) et (18.3.2.1), $B$ et $C$ sont séparés et complets pour la topologie $\mathfrak J$-pré-adique, et l'on a $B=\varprojlim B_n$, $C=\varprojlim C_n$; en outre tout homomorphisme de $A$-algèbres $u:B\to C$ est continu pour les topologies $\mathfrak J$-pré-adiques, et donne donc un système projectif d'homomorphismes de $A_n$-algèbres $u_n=u\otimes 1:B_n\to C_n$, dont il est la limite projective; la réciproque étant évidente, on a donc une bijection canonique
\[
  \Hom_{A\text{-alg.}}(B,C)\xrightarrow{\sim}\varprojlim\Hom_{A_n\text{-alg.}}(B_n,C_n).
\]
Mais comme $B$ et $C$ sont des $A$-algèbres étales, il résulte aussitôt de (18.1.2) que l'application canonique
\[
  \Hom_{A_{n+1}\text{-alg.}}(B_{n+1},C_{n+1})\longrightarrow\Hom_{A_n\text{-alg.}}(B_n,C_n)
\]
est bijective pour $n\geq 0$, ce qui achève de prouver que l'application canonique
\[
  \Hom_{A\text{-alg.}}(B,C)\longrightarrow\Hom_{A_0\text{-alg.}}(B_0,C_0)
\]
est bijective.

Pour achever de démontrer (18.3.2), il suffit de voir que pour toute $A_0$-algèbre finie et étale $B_0$, il existe une $A$-algèbre finie et étale $B$ et un $A_0$-isomorphisme $B_0\xrightarrow{\sim}B\otimes_A A_0$. Or, il résulte de (18.1.2) qu'il y a un système projectif $(B_n)$ tel que $B_n$ soit une $A_n$-algèbre étale et que les homomorphismes $B_{n+1}\otimes_{A_{n+1}}A_n\to B_n$ soient bijectifs. Il résulte de (18.3.1) et (18.3.2.1) que la $A$-algèbre $B=\varprojlim B_n$ est un $A$-module projectif de type fini et que $B_0$ est isomorphe à $B\otimes_A A_0$. Pour prouver que $B$ est une $A$-algèbre étale, il suffit, en vertu de (18.2.5), de montrer que pour tout idéal maximal $\mathfrak m$ de $A$, $B_{\mathfrak m}/\mathfrak mB_{\mathfrak m}$ est une $(A/\mathfrak m)$-algèbre séparable. Or, puisque $\mathfrak J$ est contenu dans le radical de $A$ ($0_{\mathrm I}$, 7.1.10),

\oldpage[IV-4]{117}

on a $\mathfrak J\subset\mathfrak m$ et si $\mathfrak m_0=\mathfrak m/\mathfrak J$, on a $A_0/\mathfrak m_0=A/\mathfrak m$ et $B_{\mathfrak m}/\mathfrak mB_{\mathfrak m}=(B_0)_{\mathfrak m_0}/\mathfrak m_0(B_0)_{\mathfrak m_0}$; la conclusion résulte donc du fait que $B_0$ est une $A_0$-algèbre étale (18.2.5).

\begin{env}[18.3.3]
\label{IV.18.3.3}
\emph{Exemple.} --- La prop. (18.3.2) s'applique en particulier lorsque $A$ est un \emph{anneau local} séparé et complet, $\mathfrak J$ étant l'idéal maximal de $A$, de sorte que $A_0$ est un \emph{corps} et la catégorie des $A_0$-algèbres finies et étales identique à celle des $A_0$-algèbres finies et séparables, donc isomorphes à des composées directes de corps, extensions séparables et finies de $A_0$. En particulier, si le corps $A_0$ est \emph{séparablement clos}, ces extensions sont toutes identiques à $A_0$, et par suite tout revêtement étale de $\Spec(A)$ est \emph{trivial} (18.2.7) en vertu de (18.3.2).
\end{env}

\begin{theorem}[18.3.4]
\label{IV.18.3.4}
Soient $A$ un anneau noethérien, $\mathfrak J$ un idéal de $A$ tel que $A$ soit séparé et complet pour la topologie $\mathfrak J$-pré-adique, $A_0=A/\mathfrak J$. Posons $S=\Spec(A)$, $S_0=\Spec(A_0)$. Soit $X$ un $S$-schéma propre sur $S$, et posons $X_0=X\times_S S_0$. Alors le foncteur
\[
  Z\longmapsto Z\times_X X_0
\]
de la catégorie des $X$-schémas finis et étales sur $X$ dans la catégorie des $X_0$-schémas finis et étales sur $X_0$, est une équivalence de catégories.
\end{theorem}

Montrons d'abord que ce foncteur est \emph{pleinement fidèle}. Soient donc $Z'$ et $Z''$ deux $X$-schémas finis et étales sur $X$. Posons $S_n=\Spec(A/\mathfrak J^{n+1})$, $X_n=X\times_S S_n$, $Z'_n=Z'\times_S S_n$, $Z''_n=Z''\times_S S_n$ pour tout $n\geq 0$. Il résulte de (III, 5.4.1) que l'on a une bijection canonique $\Hom_X(Z',Z'')\xrightarrow{\sim}\varprojlim\Hom_{X_n}(Z'_n,Z''_n)$. Or, en vertu de (18.1.2), l'application canonique $\Hom_{X_{n+1}}(Z'_{n+1},Z''_{n+1})\to\Hom_{X_n}(Z'_n,Z''_n)$ est bijective, ce qui achève de prouver notre assertion.

Il reste à prouver que si $\sh{B}_0$ est une $\sh{O}_{X_0}$-Algèbre finie et étale, il existe une $\sh{O}_X$-Algèbre finie et étale $\sh{B}$ et un isomorphisme $\sh{B}_0\xrightarrow{\sim}\sh{B}\otimes_{\sh{O}_X}\sh{O}_{X_0}$. Il résulte de (18.1.2) qu'il y a un système projectif $(\sh{B}_n)$, où $\sh{B}_n$ est une $\sh{O}_{X_n}$-Algèbre finie et étale, et le deuxième théorème de comparaison (III, 5.1.4) prouve l'existence d'un $\sh{O}_X$-Module cohérent $\sh{B}$ et d'un système projectif d'isomorphismes $\sh{B}_n\xrightarrow{\sim}\sh{B}\otimes_{\sh{O}_X}\sh{O}_{X_n}$. La donnée d'une structure d'Algèbre sur un Module $\sh{F}$ étant équivalente à celle d'un homomorphisme $\sh{F}\otimes\sh{F}\to\sh{F}$ rendant commutatifs des diagrammes où n'interviennent que des puissances tensorielles de $\sh{F}$, il résulte de (III, 5.1.3), (I, 10.11.4) et (I, 10.11.7) que $\sh{B}$ est muni de façon naturelle d'une structure de $\sh{O}_X$-Algèbre pour laquelle l'isomorphisme $\sh{B}_0\xrightarrow{\sim}\sh{B}\otimes_{\sh{O}_X}\sh{O}_{X_0}$ est un isomorphisme d'Algèbres. De plus, $\sh{B}$ est un $\sh{O}_X$-Module \emph{localement libre}; cela résulte encore de (III, 5.1.3), (I, 10.11.4), (I, 10.11.7) et du fait que, dans la catégorie des $\sh{O}_X$-Modules cohérents, les $\sh{O}_X$-Modules localement libres $\sh{F}$ peuvent être définis comme ceux pour lesquels le foncteur $\sh{G}\mapsto\mathscr{H}\!om_{\sh{O}_X}(\sh{F},\sh{G})$ est \emph{exact}. Enfin, pour voir que $\sh{B}$ est une $\sh{O}_X$-Algèbre étale, il suffit (18.2.5) de montrer que pour tout point \emph{fermé} $x\in X$, $\sh{B}_x\otimes_{\sh{O}_{X,x}}\kappa(x)$ est une $\kappa(x)$-algèbre séparable. Mais comme le morphisme structural $f:X\to S$ est \emph{propre}, $f(x)$ est un point fermé de $S$, donc appartient à $S_0$, puisque $\mathfrak J$ est contenu dans le radical de $A$ ($0_{\mathrm I}$, 7.1.10); la conclusion résulte donc de ce que $X_0=f^{-1}(S_0)$ et de ce que $\sh{B}_0$ est une $\sh{O}_{X_0}$-Algèbre finie et étale.

\subsection{Structure locale des morphismes non ramifiés et des morphismes étales.}
\EGAReaderTarget{ega4.c04.subsec.18-4}{subsection:IV.18.4}

\oldpage[IV-4]{118}

\begin{lemma}[18.4.1]
\label{IV.18.4.1}
Soient $A$ un anneau, $B$ une $A$-algèbre finie monogène, $u$ un générateur de la $A$-algèbre $B$, $F\in A[T]$ un polynôme tel que $F(u)=0$, $F'$ le polynôme dérivé; posons $u'=F'(u)$. Alors l'idéal de $B$, annulateur de $\Omega^1_{B/A}$, contient $u'B$; il est égal à $u'B$ si l'idéal $\mathfrak J$ de $A[T]$ formé des polynômes $G$ tels que $G(u)=0$, est engendré par $F$, autrement dit si l'homomorphisme surjectif canonique $\varphi:A[T]/F.A[T]\to B$ transformant l'image de $T$ en $u$, est bijectif.
\end{lemma}

Posons $C=A[T]$, de sorte que $B=C/\mathfrak J$. On a la suite exacte (0, 20.5.12.1)
\[
  \mathfrak J/\mathfrak J^2\longrightarrow\Omega^1_{C/A}\otimes_C B\longrightarrow\Omega^1_{B/A}\longrightarrow 0
\]
et $\Omega^1_{B/A}$ s'identifie donc au quotient $B/\mathfrak J'$, $\mathfrak J'$ étant l'idéal engendré par les éléments $G'(u)$, où $G$ parcourt un système de générateurs de l'idéal $\mathfrak J$ (0, 20.5.13); d'où aussitôt le lemme.

\begin{proposition}[18.4.2]
\label{IV.18.4.2}
Les notations étant celles de (18.4.1), soit $\mathfrak q$ un idéal premier de $B$. Alors:
\begin{enumerate}
  \item[(i)] Si $\mathfrak q$ ne contient pas $u'$, $B_{\mathfrak q}$ est une $A_{\mathfrak p}$-algèbre formellement non ramifiée ($\mathfrak p$ étant l'image réciproque de $\mathfrak q$ dans $A$); en d'autres termes, $\Spec(B_{u'})$ est formellement non ramifié sur $\Spec(A)$.
  \item[(ii)] Supposons de plus que $F$ soit unitaire et engendre $\mathfrak J$. Alors, pour que $\Spec(B)$ soit étale sur $\Spec(A)$ au point $\mathfrak q$, il faut et il suffit que $u'\notin\mathfrak q$.
\end{enumerate}
\end{proposition}

L'hypothèse que $u'\notin\mathfrak q$ entraîne que $\Omega^1_{B_{\mathfrak q}/A_{\mathfrak p}}=0$ (0, 20.5.9), donc (i) résulte de (17.2.1). En outre, sous les hypothèses de (ii), $B$ est un $A$-module libre en vertu de la division euclidienne; comme l'annulateur $\mathfrak J'$ de $\Omega^1_{B/A}$ est alors égal à $u'B$ en vertu de (18.4.1) et que $\Omega^1_{B/A}$ est un $B$-module de présentation finie (16.4.22), l'annulateur de $\Omega^1_{B_{\mathfrak q}/A_{\mathfrak p}}$ est égal à $u'B_{\mathfrak q}$ (Bourbaki, \emph{Alg. comm.}, chap. II, \textsection~2, n\textsuperscript{o}~4, formule (9)), et (ii) résulte donc de (i) et de l'implication $c)\Rightarrow a)$ dans (17.6.1).

\begin{corollary}[18.4.3]
\EGAReaderTarget{ega4.c04.cor.18-4-3}{IV.18.4.3}
Les notations étant celles de (18.4.2), supposons que $F$ soit unitaire et engendre $\mathfrak J$. Alors, pour que $B$ soit une $A$-algèbre étale, il faut et il suffit que $u'$ soit inversible dans $B$ (ou, ce qui revient au même, que l'idéal de $A[T]$ engendré par $F$ et $F'$ soit égal à $A[T]$).
\end{corollary}

Compte tenu de (18.4.2), (ii), dire que $\Spec(B)$ est étale sur $\Spec(A)$ signifie en effet que $u'$ n'appartient à aucun idéal premier de $B$, i.e. qu'il est inversible dans $B$.

On dit qu'un polynôme unitaire $F\in A[T]$ tel que l'idéal de $A[T]$ engendré par $F$ et $F'$ soit égal à $A[T]$ lui-même, est \emph{séparable}; il est immédiat que cette définition coïncide avec la définition usuelle (Bourbaki, \emph{Alg.}, chap. V, \textsection~7, n\textsuperscript{o}~6) lorsque $A$ est un corps.

\begin{lemma}[18.4.4]
\label{IV.18.4.4}
Soient $A$ un anneau local, $B$ une $A$-algèbre finie monogène, $u$ un générateur de la $A$-algèbre $B$. Soient $\mathfrak n_i$ ($1\leq i\leq r$) des idéaux maximaux de $B$ tels que $\Spec(B)$ soit formellement non ramifié sur $\Spec(A)$ aux points $\mathfrak n_i$. Alors il existe un polynôme unitaire $F\in A[T]$ tel que $F(u)=0$ et $F'(u)\notin\mathfrak n_i$ pour tout indice $i$ ($1\leq i\leq r$). En outre, si $k$ est le corps

\oldpage[IV-4]{119}

résiduel de $A$, $f\in k[T]$ le polynôme minimal de l'image de $u$ dans $B\otimes_A k$, il existe un $F\in A[T]$ dont $f$ est l'image canonique et tel que $F(u)=0$; un tel polynôme $F$ vérifie les conditions $F'(u)\notin\mathfrak n_i$ pour $1\leq i\leq r$.
\end{lemma}

L'idéal maximal $\mathfrak m$ de $A$ est l'image réciproque de chacun des $\mathfrak n_i$ (II, 6.1.10); soit $L=B\otimes_A k$, qui est une $k$-algèbre finie. Soit $\xi$ l'image de $u$ dans $L$, et soit $n$ le rang de $L$ sur $k$; le polynôme minimal $f\in k[T]$ de $\xi$ sur $k$ est donc de degré $n$, et $L$ est isomorphe à $k[T]/f.k[T]$. Si $\mathfrak n'_i=\mathfrak n_i/\mathfrak mB$, l'hypothèse entraîne que $\Spec(L)$ est étale sur $k$ aux points $\mathfrak n'_i$ ($1\leq i\leq r$) en vertu de (17.6.1, $c)$), donc $f'(\xi)\notin\mathfrak n'_i$ en vertu de (18.4.2, (ii)). Notons maintenant que $\mathfrak mB$ est contenu dans le radical de $B$; comme $1,\xi,\ldots,\xi^{n-1}$ forment une base de $L$ sur $k$, il résulte du lemme de Nakayama que $1,u,\ldots,u^{n-1}$ engendrent le $A$-module $B$, et par suite il existe un polynôme unitaire $F\in A[T]$ de degré $n$, tel que $F(u)=0$; en outre, comme $\xi$ est racine de l'image canonique de $F$ dans $k[T]$, cette image est nécessairement égale à $f$. Mais alors l'image de $F'(u)$ dans $L$ est $f'(\xi)$, et comme $f'(\xi)\notin\mathfrak n'_i$, on a $F'(u)\notin\mathfrak n_i$ pour tout $i$.

\begin{proposition}[18.4.5]
\label{IV.18.4.5}
Soient $A$ un anneau local, $k$ son corps résiduel, $B$ une $A$-algèbre finie (resp. finie et de présentation finie). On suppose en outre, ou bien que le corps $k$ est infini, ou bien que $B$ est un anneau local. Soit $n$ le rang de $L=B\otimes_A k$ sur $k$. Pour que $B$ soit une $A$-algèbre formellement non ramifiée (resp. étale), il faut et il suffit qu'il existe un polynôme unitaire séparable $F\in A[T]$ (18.4.3) tel que $B$ soit isomorphe à un quotient de $A[T]/F.A[T]$ (resp. isomorphe à $A[T]/F.A[T]$). En outre, on peut supposer que $F$ est de degré $n$ (resp. $F$ est nécessairement de degré $n$).
\end{proposition}

Les conditions sont suffisantes en vertu de (18.4.2), sans supposer $k$ infini ni $B$ local. Pour voir que les conditions sont nécessaires, notons que si $B$ est une $A$-algèbre formellement non ramifiée, $L$ est une algèbre finie et séparable sur $k$, donc composée directe d'un nombre fini d'extensions finies et séparables $k_j$ de $k$ ($1\leq j\leq r$), $k_j$ étant donc engendrée par un élément $\xi_j$, de polynôme minimal $f_j$ ($1\leq j\leq r$) (Bourbaki, \emph{Alg.}, chap. V, \textsection~11, n\textsuperscript{o}~4, prop. 4). Montrons qu'en vertu des hypothèses faites sur $k$ ou $B$, il existe un élément $\xi$ de $L$ engendrant la $k$-algèbre $L$. C'est immédiat si $B$ est local, puisque alors $r=1$. Sinon, $k$ étant supposé infini, on peut supposer que les polynômes irréductibles $f_j\in k[T]$ sont tous \emph{distincts} (en remplaçant au besoin chaque $\xi_j$ par $\xi_j+a_j$, pour un élément convenable $a_j\in k$); si l'on pose $f=f_1f_2\cdots f_r$, il est clair que $L$ est isomorphe à $k[T]/f.k[T]$ dans les deux cas envisagés, donc est engendrée par un élément $\xi$ de polynôme minimal $f\in k[T]$ de degré $n$. Si $u\in B$ est un élément d'image $\xi$ dans $L$, le lemme de Nakayama montre que les éléments $1,u,\ldots,u^{n-1}$ engendrent le $A$-module $B$; ceci montre déjà qu'il existe un polynôme unitaire $F\in A[T]$ de degré $n$ tel que $F(u)=0$, $u$ engendrant la $A$-algèbre $B$, qui est par suite isomorphe à une algèbre quotient de $A[T]/F.A[T]$; en outre $B$ est un anneau semi-local et en chacun de ses idéaux maximaux $\mathfrak n_i$, on a $F'(u)\notin\mathfrak n_i$ d'après (18.4.4), ce qui prouve que $F'(u)$ est inversible dans $B$, donc que $F$ est un polynôme \emph{séparable}. Enfin, si $B$ est une $A$-algèbre étale, $B$ étant un $A$-module plat de présentation finie (1.4.7) est un $A$-module libre, et $1,u,\ldots,u^{n-1}$ forment une base du $A$-module $B$ (Bourbaki, \emph{Alg. comm.}, chap. II, \textsection~3, n\textsuperscript{o}~2, prop. 5), autrement dit la $A$-algèbre $B$ est

\oldpage[IV-4]{120}

isomorphe à $A[T]/F.A[T]$, et pour tout autre polynôme unitaire $G\in A[T]$ tel que $B$ soit isomorphe à $A[T]/G.A[T]$, $G$ est nécessairement de degré $n$.

\begin{theorem}[18.4.6]
\label{IV.18.4.6}
(Chevalley). ---
\begin{enumerate}
  \item[(i)] Soient $f:X\to Y$ un morphisme localement de type fini, $x$ un point de $X$, $y=f(x)$, et posons $A=\sh{O}_{Y,y}$. Pour que $\sh{O}_{X,x}$ soit une $A$-algèbre formellement étale (resp. formellement non ramifiée), il faut et il suffit qu'il existe un polynôme unitaire $F\in A[T]$ et un idéal maximal $\mathfrak n$ de $B=A[T]/F.A[T]$ (resp. d'une algèbre quotient $B$ de $A[T]/F.A[T]$) tels que $\sh{O}_{X,x}$ soit $A$-isomorphe à $B_{\mathfrak n}$ et que, si $u$ est l'image de $T$ dans $B$, on ait $F'(u)\notin\mathfrak n$.
  \item[(ii)] Supposons en outre que $f$ soit localement de présentation finie. Alors, pour que $f$ soit étale au point $x$, il faut et il suffit de plus que l'on puisse prendre $B=A[T]/F.A[T]$.
\end{enumerate}
\end{theorem}

Les conditions sont suffisantes en vertu de (18.4.2). Pour voir qu'elles sont nécessaires, on peut évidemment se borner au cas où $X$ et $Y$ sont affines, et, compte tenu de la remarque (17.4.1.2), au cas où $f$ est formellement non ramifié et quasi-fini. Comme $f$ est affine, il résulte de (8.12.8) qu'il existe une $A$-algèbre finie $C$ et un idéal maximal $\mathfrak r$ de $C$ (nécessairement au-dessus de l'idéal maximal $\mathfrak m$ de $A$) tels que $\sh{O}_{X,x}$ soit $A$-isomorphe à $C_{\mathfrak r}$. En outre (17.4.1.2) le corps résiduel $C/\mathfrak r=C_{\mathfrak r}/\mathfrak rC_{\mathfrak r}$ est une extension finie séparable de $k=A/\mathfrak m$, donc de la forme $k[v]$, où $v$ est \emph{séparable} sur $k$. Soient $\mathfrak r_i$ ($1\leq i\leq h$) les idéaux maximaux de l'anneau semi-local $C$ \emph{autres} que $\mathfrak r$; il existe un élément $u\in C$ appartenant à tous les $\mathfrak r_i$ et tel que son image dans $C/\mathfrak r$ soit égale à $v$ (Bourbaki, \emph{Alg. comm.}, chap. II, \textsection~1, n\textsuperscript{o}~2, prop. 5). Nous allons montrer que la sous-$A$-algèbre $B=A[u]$ de $C$ et l'idéal (nécessairement maximal puisque $C$ est finie sur $B$) $\mathfrak n=\mathfrak r\cap B$ de $B$ répondent à la question.

Pour traiter le cas où $\sh{O}_{X,x}$ est une $A$-algèbre formellement ramifiée, il suffira de prouver que $B_{\mathfrak n}$ est isomorphe à $C_{\mathfrak r}$; en effet, $B_{\mathfrak n}$ sera alors formellement non ramifié sur $A$ et l'existence du polynôme $F\in A[T]$ ayant les propriétés de l'énoncé résultera de (18.4.4). Notons maintenant que, puisque $f$ est formellement non ramifié, $C/\mathfrak r$ est isomorphe à la $A$-algèbre $C_{\mathfrak r}/\mathfrak mC_{\mathfrak r}$ (17.4.1.2). On est ainsi ramené à prouver le lemme suivant:

\begin{lemma}[18.4.6.1]
\label{IV.18.4.6.1}
Soient $A$ un anneau local, $\mathfrak m$ son idéal maximal, $C$ une $A$-algèbre finie, $\mathfrak r$ un idéal maximal de $C$. Soit $u$ un élément de $C$ appartenant à tous les idéaux maximaux de $C$ distincts de $\mathfrak r$, n'appartenant pas à $\mathfrak r$, et tel que $C_{\mathfrak r}/\mathfrak mC_{\mathfrak r}$ soit une algèbre monogène sur $k=A/\mathfrak m$, engendrée par l'image de $u$ dans $C_{\mathfrak r}/\mathfrak mC_{\mathfrak r}$. Posons $B=A[u]$, $\mathfrak n=\mathfrak r\cap B$. Alors l'homomorphisme canonique $B_{\mathfrak n}\to C_{\mathfrak r}$ est un isomorphisme.
\end{lemma}

Posons $R=B-\mathfrak n$, $S=C-\mathfrak r$, de sorte que $B_{\mathfrak n}=R^{-1}B$ et $C_{\mathfrak r}=S^{-1}C$; l'homomorphisme canonique $B_{\mathfrak n}\to C_{\mathfrak r}$ peut s'écrire comme le composé
\[
  R^{-1}B\xrightarrow{g}R^{-1}C\xrightarrow{h}S^{-1}C
\]
et il suffit de montrer que chacun de ces deux homomorphismes est bijectif.

Montrons d'abord que $h:R^{-1}C\to S^{-1}C=C_{\mathfrak r}$ est bijectif; il suffit de voir que les images dans $R^{-1}C$ des éléments de $S$ sont inversibles, ou encore que tout idéal maximal $\mathfrak p$ de $R^{-1}C$ a une image réciproque dans $C$ ne rencontrant pas $S$, donc nécessairement égale à $\mathfrak r$. Or, comme $R^{-1}C$ est une $(R^{-1}B)$-algèbre finie, l'image réciproque de $\mathfrak p$ dans

\oldpage[IV-4]{121}

$R^{-1}B=B_{\mathfrak n}$ est l'unique idéal maximal $\mathfrak nB_{\mathfrak n}$ de cet anneau, donc l'image réciproque de $\mathfrak p$ dans $B$ est égale à $\mathfrak n$. Mais d'autre part, si $\mathfrak q$ est l'image réciproque de $\mathfrak p$ dans $C$, on a $\mathfrak q\cap B=\mathfrak n$, et comme $\mathfrak n$ est maximal dans $B$ et $C$ une $B$-algèbre finie, $\mathfrak q$ est nécessairement un des idéaux maximaux de $C$; en outre, on a $u\notin\mathfrak q$ puisque $u\notin\mathfrak r$ et que $u\in B$, donc par hypothèse on a nécessairement $\mathfrak q=\mathfrak r$.

D'autre part, comme $B\subset C$, l'homomorphisme $g:R^{-1}B\to R^{-1}C$ est injectif ($0_{\mathrm I}$, 1.3.2); pour voir qu'il est surjectif, notons que $R^{-1}C$ est un $(R^{-1}B)$-module de type fini, et d'autre part que $\mathfrak mR^{-1}B$ est contenu dans l'idéal maximal de l'anneau local $R^{-1}B$; en vertu du lemme de Nakayama, il suffit de prouver que l'homomorphisme
\[
  R^{-1}B/\mathfrak mR^{-1}B\longrightarrow R^{-1}C/\mathfrak mR^{-1}C
\]
est surjectif. Mais en vertu de la première partie de la démonstration, $R^{-1}C/\mathfrak mR^{-1}C$ s'identifie à $C_{\mathfrak r}/\mathfrak mC_{\mathfrak r}$; par hypothèse cette $k$-algèbre est engendrée par l'image de $u$, et \emph{a fortiori} elle est égale à l'image de $R^{-1}B/\mathfrak mR^{-1}B$.

Considérons en second lieu le cas où $f$ est étale au point $x$. En remplaçant $X$ par un voisinage de $x$, on peut supposer que $X$ est un voisinage de $\mathfrak n$ dans $\Spec(B)$ (I, 7.2). Posons $B'=\Spec(A[T]/F.A[T])$ et soit $\mathfrak n'$ l'image réciproque de $\mathfrak n$ dans $B'$; comme l'image de $F'(T)$ dans $B'$ n'appartient pas à $\mathfrak n'$ par hypothèse, le morphisme $\Spec(B')\to\Spec(A)$ est étale au point $\mathfrak n'$ par (18.4.2, (ii)). Comme par hypothèse $\Spec(B)\to\Spec(A)$ est étale au point $\mathfrak n$, on en conclut (17.3.4) que $\Spec(B)\to\Spec(B')$ est étale au point $\mathfrak n$; mais comme ce morphisme est une immersion, il ne peut être étale en un point que si c'est un isomorphisme local en ce point (17.9.1), donc $B_{\mathfrak n}$ et $B'_{\mathfrak n'}$ sont isomorphes.

Enfin, supposons que $\sh{O}_{X,x}$ soit une $A$-algèbre formellement étale; avec les notations précédentes, il résulte de (17.1.5) que $B_{\mathfrak n}$ est une $B'_{\mathfrak n'}$-algèbre formellement étale; mais puisque l'homomorphisme $B'_{\mathfrak n'}\to B_{\mathfrak n}$ est surjectif, cela ne peut avoir lieu que si cet homomorphisme est \emph{bijectif} (0, 19.10.3, (i)). Ceci achève la démonstration de (18.4.6).

\begin{corollary}[18.4.7]
\label{IV.18.4.7}
Soient $f:X\to Y$ un morphisme localement de type fini, $x$ un point de $X$. Pour que $f$ soit formellement non ramifié au point $x$, il faut et il suffit qu'il existe un voisinage ouvert $U$ de $x$ tel que $f|U$ se factorise en $U\xrightarrow{j}X'\xrightarrow{h}Y$, où $h$ est un morphisme étale et $j$ une immersion fermée.
\end{corollary}

On peut évidemment se borner au cas où $Y=\Spec(R)$ est affine et $f$ de type fini. Si $A=\sh{O}_{Y,y}$, avec $y=f(x)$, la condition pour que $f$ soit formellement non ramifié au point $x$ équivaut, en vertu de (17.4.1.2), à dire que $\sh{O}_{X,x}$ est une $A$-algèbre formellement non ramifiée. S'il en est ainsi, on peut appliquer (18.4.6, (ii)); remplaçant au besoin $Y$ par un voisinage affine de $y$, on peut supposer (avec les notations de (18.4.6)) que le polynôme $F$ est l'image dans $A[T]$ d'un polynôme unitaire $G\in R[T]$. On pose alors $X'=\Spec(R[T]/G.R[T])$; soit $x'$ l'image du point $\mathfrak n$ de $\Spec(B)$ par le morphisme correspondant à l'homomorphisme composé $R[T]/G.R[T]\to A[T]/F.A[T]\to B$. Il résulte de (18.4.2) que le morphisme $h:X'\to Y$ correspondant à l'homomorphisme canonique $R\to R[T]/G.R[T]$ est étale au point $x'$, donc, en restreignant au besoin $X'$ et $Y$ à des voisinages ouverts de $x'$ et $y$ respectivement, on peut supposer $h$ étale. D'autre part, en vertu de (I, 6.5.1, (ii)) et de (I, 7.2), il résulte du fait que l'on a un homomorphisme local $\varphi:\sh{O}_{X',x'}\to\sh{O}_{X,x}$ que cet homomorphisme correspond à un morphisme

\oldpage[IV-4]{122}

$j:U\to X'$ d'un voisinage ouvert $U$ de $x$ dans $X$; restreignant au besoin $U$, on peut, en appliquant (I, 6.5.1, (i)), supposer que $h\circ j=f|U$, donc $j$ est de type fini (I, 6.3.4, (v)); enfin, l'homomorphisme $\varphi$ est surjectif par (18.4.6, (ii)), donc il résulte de (I, 6.5.4, (i)) qu'on peut, en restreignant encore $U$ et $X'$, supposer que $j$ est une immersion fermée. Cela prouve donc la nécessité de la condition énoncée; sa suffisance est immédiate (17.1.3, (i) et (ii)).

\begin{corollary}[18.4.8]
\label{IV.18.4.8}
Soit $f:X\to Y$ un morphisme localement de présentation finie. Pour que $f$ soit non ramifié (resp. étale), il faut et il suffit qu'il existe une famille de morphismes plats $g_\alpha:Y'_\alpha\to Y$ et, pour chaque $\alpha$, un ouvert $U_\alpha$ dans $X'_\alpha=X\times_Y Y'_\alpha$, tels que, si $g'_\alpha:X'_\alpha\to X$ et $f'_\alpha:X'_\alpha\to Y'_\alpha$ sont les projections canoniques, les $g'_\alpha(U_\alpha)$ forment un recouvrement de $X$, et que chacun des morphismes composés $U_\alpha\to X'_\alpha\xrightarrow{f'_\alpha}Y'_\alpha$ soit une immersion fermée (resp. ouverte). On peut de plus alors prendre les $g_\alpha$ étales.
\end{corollary}

La nécessité de la condition pour les morphismes étales est triviale, en prenant un seul $Y'_\alpha$ égal à $X$, le morphisme $g_\alpha$ correspondant étant égal à $f$, et l'ouvert $U\subset X\times_Y X$ étant la diagonale. Lorsque $f$ est non ramifié, la nécessité de la condition résulte de (18.4.7); on prend un recouvrement ouvert $(V_\alpha)$ de $X$ tel que pour chaque $\alpha$, $f|V_\alpha$ se factorise en $V_\alpha\xrightarrow{j_\alpha}Y'_\alpha\xrightarrow{g_\alpha}Y$, où $j_\alpha$ est une immersion fermée et $g_\alpha$ un morphisme étale. Alors $j_\alpha:V_\alpha\to Y'_\alpha$ se factorise en $V_\alpha\xrightarrow{s_\alpha}X'_\alpha\xrightarrow{f'_\alpha}Y'_\alpha$, où $s_\alpha$ est une $V_\alpha$-section de $X'_\alpha$, et comme le morphisme $g'_\alpha:X'_\alpha\to X$ est étale, $s_\alpha$ est une immersion ouverte (17.4.1), et il suffit de prendre $U_\alpha=s_\alpha(V_\alpha)$ pour répondre à la question.

La suffisance des conditions résulte de (17.7.1); on en conclut que $f$ est non ramifié (resp. étale) en chaque point de $g'(U_\alpha)$, donc dans $X$ tout entier puisque les $g'(U_\alpha)$ recouvrent $X$.

\begin{proposition}[18.4.9]
\label{IV.18.4.9}
Soient $S$ un préschéma, $f:X\to S$ un morphisme, $h:Y\to S$ un morphisme localement de présentation finie, $g:X\to Y$ un $S$-morphisme, $x$ un point de $X$, $y=g(x)$. Les conditions suivantes sont équivalentes:
\begin{enumerate}
  \item[a)] $h$ est étale au point $y$ et $g$ est plat au point $x$.
  \item[b)] $h$ est non ramifié au point $y$ et $f$ est plat au point $x$.
\end{enumerate}
\end{proposition}

Comme $f=h\circ g$, $a)$ entraîne que $f$ est plat au point $x$ (2.1.6), et évidemment que $h$ est non ramifié au point $y$, donc $a)$ implique $b)$.

Pour prouver que $b)$ entraîne $a)$, on peut d'abord supposer que $h$ est non ramifié (en remplaçant $Y$ par un voisinage de $y$); puis, en remplaçant $S$ par un voisinage ouvert de $s=h(y)=f(x)$, on peut supposer qu'il existe un morphisme étale $u:S'\to S$, un point $y'$ de $Y'=Y_{(S')}$ au-dessus de $y$ et un voisinage ouvert $V$ de $y'$ dans $Y'$ tel que, si $h'=h_{(S')}:Y'\to S'$, la restriction de $h'$ à $V$ soit une immersion fermée (18.4.8). Si l'on prouve alors que $h'$ est étale au point $y'$, il en résultera que $h$ est étale au point $y$ (17.7.1, (ii)); d'ailleurs, $f'=f_{(S')}:X_{(S')}\to S'$ est plat en tout point $x'$ au-dessus de $x$. Comme les projections $v:Y'\to Y$, $w:X'\to X$ sont des morphismes étales (donc plats), si l'on prouve que $g'=g_{(S')}:X_{(S')}\to Y_{(S')}$ est plat au point $x'$, il résultera de (2.2.11, (iv)) que $g$ sera plat au point $x$. On peut donc se borner au cas où $h$ est une \emph{immersion fermée}

\oldpage[IV-4]{123}

de présentation finie, $f$ étant supposé plat au point $x$. Soit $\mathscr J$ l'Idéal quasi-cohérent de type fini de $\sh{O}_S$ qui définit le sous-préschéma fermé $Y$ de $S$. L'hypothèse que $f$ est plat au point $x$ entraîne que l'homomorphisme $\sh{O}_{S,s}\to\sh{O}_{X,x}$ est injectif ($0_{\mathrm I}$, 6.5.1); \emph{a fortiori} l'homomorphisme $\sh{O}_{S,s}\to\sh{O}_{Y,y}$ est injectif, ce qui signifie que $\mathscr J_s=0$, puisque c'est le noyau de l'homomorphisme précédent. Puisque $\mathscr J$ est de type fini, il y a un voisinage ouvert $U$ de $s$ dans $S$ tel que $\mathscr J|U=0$ ($0_{\mathrm I}$, 5.2.2). On peut donc supposer que $h$ est une immersion ouverte, et alors il est clair que $g$ est plat au point $x$.

\begin{proposition}[18.4.10]
\label{IV.18.4.10}
Désignons par $\mathbf P(f,x)$ une propriété vérifiant les conditions suivantes:
\begin{enumerate}
  \item[$1^\circ$] Pour tout morphisme $f:X\to Y$ et tout isomorphisme local $h:Y\to Z$, $\mathbf P(f,x)$ est équivalente à $\mathbf P(h\circ f,x)$ pour $x\in X$.
  \item[$2^\circ$] Pour tout morphisme $f:X\to Y$, tout morphisme étale $g:Y'\to Y$, tout point $x\in X$, si l'on pose $X'=X\times_Y Y'$, $f'=f_{(Y')}:X'\to Y'$ et si $x'\in X'$ est au-dessus de $x$, les propriétés $\mathbf P(f,x)$ et $\mathbf P(f',x')$ sont équivalentes (« invariance par changement de base étale »).
\end{enumerate}
Soient alors $S$ un préschéma, $f:X\to S$ et $h:Y\to S$ deux morphismes, $g:X\to Y$ un $S$-morphisme, $x$ un point de $X$, $y=g(x)$, et supposons $h$ étale au point $y$. Alors les propriétés $\mathbf P(f,x)$ et $\mathbf P(g,x)$ sont équivalentes.
\end{proposition}

Avec les mêmes notations que dans la démonstration de (18.4.9) et en remplaçant $S$ (resp. $Y$) par un voisinage de $s=h(y)$ (resp. de $y$), on peut ici, en vertu de (18.4.8), trouver un morphisme étale $u:S'\to S$ tel que $h'$ soit une immersion ouverte; si $x'\in X'$ est au-dessus de $x$, $\mathbf P(f',x')$ (resp. $\mathbf P(g',x')$) est alors par hypothèse équivalente à $\mathbf P(f,x)$ (resp. $\mathbf P(g,x)$). On peut donc se borner au cas où $h$ est une immersion ouverte. Comme l'hypothèse entraîne que $\mathbf P(g,x)$ est alors équivalente à $\mathbf P(h\circ g,x)$, on en déduit la conclusion.

\begin{env}[18.4.11]
\label{IV.18.4.11}
\emph{Exemples.} --- On peut prendre pour propriété $\mathbf P(f,x)$ l'une quelconque des suivantes, compte tenu de (17.7.4, (ii)):
\begin{enumerate}
  \item[(i)] $f$ est plat au point $x$ (2.2.11, (iv));
  \item[(ii)] $f$ est localement de présentation finie et de coprofondeur $\leq n$ au point $x$ (6.8.1 et 6.7.8);
  \item[(iii)] $f$ est localement de présentation finie et de Cohen-Macaulay au point $x$ (6.8.1 et 6.7.8);
  \item[(iv)] $f$ est localement de présentation finie et possède la propriété $(S_n)$ au point $x$ (6.8.1 et 6.7.8);
  \item[(v)] $f$ est localement de présentation finie et possède la propriété $(R_n)$ au point $x$ (6.8.1 et 6.7.8);
  \item[(vi)] $f$ est localement de présentation finie et normal au point $x$ (6.8.1 et 6.7.8);
  \item[(vii)] $f$ est localement de présentation finie et réduit au point $x$ (6.8.1 et 6.7.8);
  \item[(viii)] $f$ est non ramifié au point $x$ (17.7.4);
  \item[(ix)] $f$ est lisse au point $x$ (17.7.4);
  \item[(x)] $f$ est étale au point $x$ (17.7.4).
\end{enumerate}
\end{env}

\oldpage[IV-4]{124}

\begin{corollary}[18.4.12]
\label{IV.18.4.12}
\begin{enumerate}
  \item[(i)] Soient $f:X\to Y$ un morphisme localement de type fini, $x$ un point de $X$. Si $f$ est plat et formellement non ramifié au point $x$, alors, dans toute factorisation $f|U:U\xrightarrow{j}X'\xrightarrow{h}Y$, où $U$ est un voisinage ouvert de $x$, $j$ une immersion fermée et $h$ un morphisme étale (18.4.7), l'homomorphisme $\sh{O}_{X',j(x)}\to\sh{O}_{X,x}$ correspondant à $j$ est bijectif (en particulier $\sh{O}_{X,x}$ est une $\sh{O}_{Y,f(x)}$-algèbre \emph{essentiellement étale} (18.6.1)).
  \item[(ii)] Pour qu'un morphisme $f:X\to Y$ soit étale, il faut et il suffit qu'il soit localement de type fini, formellement non ramifié et plat.
\end{enumerate}
\end{corollary}

(i) Comme l'homomorphisme $\sh{O}_{X',j(x)}\to\sh{O}_{X,x}$ est surjectif, il suffit de prouver qu'il est injectif, et pour cela qu'il fait de $\sh{O}_{X,x}$ un $\sh{O}_{X',j(x)}$-module \emph{fidèlement plat}, ou seulement \emph{plat} ($0_{\mathrm I}$, 6.5.1 et 6.6.2); autrement dit, il s'agit de montrer que $j$ est un morphisme \emph{plat au point} $x$; mais puisque $h\circ j=f$ est par hypothèse plat au point $x$ et que $h$ est étale, cela résulte de (18.4.10) et (18.4.11, (i)).

(ii) Il n'y a à prouver que la suffisance des conditions énoncées. Pour tout $x\in X$, on a donc dans un voisinage ouvert $U$ de $x$ une factorisation de $f|U$ ayant les propriétés considérées dans (i). Comme par hypothèse $f$ est plat et formellement non ramifié en \emph{tous} les points de $U$, le résultat de (i) s'applique non seulement à $x$ mais à tous les points de $U$; cela signifie que si $\mathscr J$ est l'Idéal quasi-cohérent de $\sh{O}_{X'}$ correspondant au sous-préschéma fermé de $X'$ associé à $j$, on a $\mathscr J_{j(z)}=0$ pour tout $z\in U$, donc $j$ est une immersion ouverte, et $f$ est par suite étale en tout point de $U$, donc en tout point de $X$.

\begin{corollary}[18.4.13]
\label{IV.18.4.13}
Soient $f:X\to Y$ un morphisme localement de type fini, $x$ un point de $X$, $y=f(x)$. Supposons que $y$ admette un voisinage ouvert qui soit un préschéma réduit et n'ayant qu'un nombre fini de composantes irréductibles. Alors, pour que $f$ soit étale au point $x$, il faut et il suffit que $f$ soit plat et formellement non ramifié au point $x$.
\end{corollary}

Il n'y a à démontrer que la suffisance de la condition. La question étant locale sur $X$ et $Y$, on peut supposer (18.4.7) que $f$ se factorise en $X\xrightarrow{j}X'\xrightarrow{h}Y$ où $h$ est étale et $j$ une immersion fermée, et en outre que $Y$ est réduit et n'a qu'un nombre fini de composantes irréductibles. Alors $X'$ est réduit (17.5.7) et, en remplaçant au besoin $X'$ par un voisinage ouvert de $j(x)$, on peut supposer que $X'$ n'a qu'un nombre fini de composantes irréductibles: en effet, on peut supposer $h$ quasi-fini (17.6.1) et comme les points maximaux de $X'$ sont au-dessus des points maximaux de $Y$ (2.3.4) leur nombre est fini. Tout revient à montrer, avec les notations de la démonstration de (18.4.12), que l'on a $\mathscr J_{x'}=0$ pour tous les points $x'$ d'un voisinage de $j(x)$ dans $X'$, sachant que $\mathscr J_{j(x)}=0$. Or, en remplaçant $X'$ par un voisinage affine de $j(x)$, on peut supposer que toutes les composantes irréductibles de $X'$ contiennent $j(x)$; si $X'=\Spec(A')$, et si $\mathfrak p'$ est l'idéal premier de $A'$ correspondant au point $j(x)$, le morphisme $\Spec(A'_{\mathfrak p'})\to\Spec(A')$ est dominant, donc l'homomorphisme correspondant $A'\to A'_{\mathfrak p'}$ est \emph{injectif} puisque $A'$ est réduit (I, 1.2.7). Si $\mathscr J=\widetilde{\mathfrak J}$, où $\mathfrak J$ est un idéal de $A'$, $\mathfrak J$ s'identifie donc à une partie de $\mathfrak J_{\mathfrak p'}$, et l'hypothèse $\mathfrak J_{\mathfrak p'}=0$ entraîne donc $\mathfrak J=0$.

\begin{corollary}[18.4.14]
\label{IV.18.4.14}
Soient $A$ un anneau local d'idéal maximal $\mathfrak m$, de corps résiduel $k$, $B$ une $A$-algèbre finie.

\oldpage[IV-4]{125}

\begin{enumerate}
  \item[(i)] Pour que $B$ soit une $A$-algèbre formellement non ramifiée, il faut et il suffit que $B\otimes_A k$ soit une $k$-algèbre étale.
  \item[(ii)] Pour que $B$ soit une $A$-algèbre étale, il faut et il suffit que $B\otimes_A k$ soit une $k$-algèbre étale et que $B$ soit un $A$-module plat (ce qui équivaut à dire ($0_{\mathrm I}$, 6.3.3) que pour tout idéal maximal $\mathfrak n$ de $B$ (nécessairement au-dessus de $\mathfrak m$), $B_{\mathfrak n}$ est un $A$-module plat).
\end{enumerate}
\end{corollary}

On n'a évidemment à prouver que la suffisance des conditions énoncées. Il est clair que (ii) résulte de (i) et de (18.4.12, (ii)), compte tenu de (17.1.2, (i)). Pour prouver (i) remarquons que si $B\otimes_A k$ est une $k$-algèbre étale, on a $\Omega^1_{(B\otimes_A k)/k}=0$ (17.2.1). Or on a $\Omega^1_{(B\otimes_A k)/k}=\Omega^1_{B/A}\otimes_B k$ (0, 20.5.5) et puisque $B$ est une $A$-algèbre de type fini, $\Omega^1_{B/A}$ est un $B$-module de type fini (0, 20.4.7). Mais comme $B$ est une $A$-algèbre finie, $\mathfrak mB$ est contenu dans le radical de $B$ (Bourbaki, \emph{Alg. comm.}, chap. V, \textsection~2, n\textsuperscript{o}~1, prop. 1), donc le lemme de Nakayama prouve que $\Omega^1_{B/A}=0$, et par suite $B$ est une $A$-algèbre formellement non ramifiée (17.2.1).

\subsection{Anneaux locaux henséliens\texorpdfstring{\protect\footnotemark}{}.}
\footnotetext{La notion d'anneau local hensélien est due à Azumaya, celle de hensélisation à Nagata, à qui l'on doit aussi les principaux résultats de cette théorie.}
\EGAReaderTarget{ega4.c04.subsec.18-5}{subsection:IV.18.5}

\begingroup
\newcommand{\shHom}{\mathscr{H}\!om}
\begin{env}[18.5.1]
\label{IV.18.5.1}
Soient $X$ un préschéma, $\sh{E}$ un $\sh{O}_X$-Module \emph{localement libre} de rang fini; le dual $\check{\sh{E}}=\shHom_{\sh{O}_X}(\sh{E},\sh{O}_X)$ est donc un $\sh{O}_X$-Module localement libre dont le rang en tout point de $X$ est égal à celui de $\sh{E}$ en ce point, et l'homomorphisme canonique $\sh{E}\to\shHom_{\sh{O}_X}(\check{\sh{E}},\sh{O}_X)=(\check{\sh{E}})^\vee$ est \emph{bijectif}. Pour tout morphisme $X'\to X$, posons $\sh{E}_{(X')}=\sh{E}\otimes_{\sh{O}_X}\sh{O}_{X'}$ qui est un $\sh{O}_{X'}$-Module localement libre, et considérons l'\emph{ensemble} $\Gamma(X',\sh{E}_{(X')})$ des sections de ce $\sh{O}_{X'}$-Module au-dessus de $X'$. Nous allons voir que l'on définit ainsi un \emph{foncteur contravariant représentable}
\begin{equation}
  {}^t\!\mathbf V_{\sh{E}}:X'\longmapsto\Gamma(X',\sh{E}_{(X')})
  \tag{18.5.1.1}\label{IV.18.5.1.1}
\end{equation}
de la catégorie des $X$-préschémas dans celle des \emph{ensembles} ($0_{\mathrm{III}}$, 8.1.8).

Tout d'abord, on a bien défini un foncteur, car si $f:X''\to X'$ est un $X$-morphisme de $X$-préschémas, on a $\sh{E}_{(X'')}=f^*(\sh{E}_{(X')})$, d'où ($0_{\mathrm I}$, 4.4.3.2) une application $\Gamma(X',\sh{E}_{(X')})\to\Gamma(X'',f_*(\sh{E}_{(X'')}))=\Gamma(X'',\sh{E}_{(X'')})$ qui achève de définir le foncteur ${}^t\!\mathbf V$. Montrons ensuite que le $X$-préschéma $\mathbf V(\check{\sh{E}})$ (II, 1.7.8) \emph{représente} le foncteur ${}^t\!\mathbf V$. Il est immédiat en effet que l'on a $(\sh{E}_{(X')})^\vee=(\check{\sh{E}})_{(X')}$, donc $\mathbf V(\check{\sh{E}})\times_X X'=\mathbf V((\check{\sh{E}})_{(X')})$; compte tenu de (I, 3.3.14), on est ramené à définir une bijection $\Gamma(\mathbf V(\check{\sh{E}})/X)\to\Gamma(X,\sh{E})$, et à vérifier que la bijection $\Gamma(\mathbf V((\check{\sh{E}})_{(X')})/X')\to\Gamma(X',\sh{E}_{(X')})$ est fonctorielle en $X'$. Or, on a une bijection canonique de $\Gamma(\mathbf V(\check{\sh{E}})/X)$ sur $\Hom_{\sh{O}_X}(\check{\sh{E}},\sh{O}_X)$ (II, 1.7.8), et la transposition $u\mapsto{}^t u$ est une bijection canonique de $\Hom_{\sh{O}_X}(\check{\sh{E}},\sh{O}_X)$ sur $\Hom_{\sh{O}_X}(\sh{O}_X,\sh{E})=\Gamma(X,\sh{E})$, en raison de l'identification de $(\check{\sh{E}})^\vee$ et de $\sh{E}$. La vérification de la propriété fonctorielle est immédiate.
\end{env}
\endgroup

\oldpage[IV-4]{126}

Notons encore que, conformément à la théorie générale ($0_{\mathrm{III}}$, 8.1.6), l'automorphisme identique de $\mathbf V(\check{\sh{E}})$ correspond canoniquement à une section $c$ de $\sh{E}_{(\mathbf V(\check{\sh{E}}))}$ au-dessus de $\mathbf V(\check{\sh{E}})$, c'est-à-dire (II, 1.4.1) à un homomorphisme de $\mathbf S_{\sh{O}_X}(\check{\sh{E}})$-Modules
\[
  u:\mathbf S_{\sh{O}_X}(\check{\sh{E}})\longrightarrow
  \mathbf S_{\sh{O}_X}(\check{\sh{E}})\otimes_{\sh{O}_X}\sh{E};
\]
pour tout ouvert affine $W$ de $X$, si l'on pose $\Gamma(W,\sh{O}_X)=A$, si on identifie $\Gamma(W,\mathbf S_{\sh{O}_X}(\check{\sh{E}}))$ à une algèbre de polynômes $C=A[T_1,\ldots,T_n]$, de sorte que les $T_i$ forment une base de $\Gamma(W,\check{\sh{E}})$, et si enfin on désigne par $(e_i)$ la base duale de $(T_i)$ dans $\Gamma(W,\sh{E})=\Gamma(W,\check{\sh{E}})^\vee$, on voit aussitôt que $u$ correspond à l'homomorphisme de $C$-modules tel que
\[
  u(1)=\sum_{i=1}^n T_i\otimes e_i.
\]
Si $X'=\Spec(A)$ est affine et tel que $\sh{E}_{(X')}$ soit isomorphe à $\sh{O}_{X'}^n$, $\Gamma(X',\sh{E}_{(X')})$ est un $A$-module libre de rang $n$; de façon imagée, on peut dire que l'objet $\mathbf V(\check{\sh{E}})$ représente « l'ensemble des points de l'espace affine tordu sur $X$ défini par $\sh{E}$ ».

\begin{env}[18.5.2]
\label{IV.18.5.2}
Rappelons (II, 1.7.8) que l'on a par définition $\mathbf V(\check{\sh{E}})=\Spec(\mathbf S_{\sh{O}_X}(\check{\sh{E}}))$. Soit $\mathscr J$ un Idéal quasi-cohérent de $\mathbf S_{\sh{O}_X}(\check{\sh{E}})$, de sorte que $\Spec(\mathbf S_{\sh{O}_X}(\check{\sh{E}})/\mathscr J)$ est un sous-préschéma fermé de $\mathbf V(\check{\sh{E}})$; nous allons l'interpréter comme \emph{représentant un foncteur} de la catégorie des $X$-préschémas dans celle des ensembles. Notons pour cela qu'une section $u\in\Gamma(X,\sh{E})$ s'identifie canoniquement à un $\sh{O}_X$-homomorphisme $u:\sh{O}_X\to\sh{E}$, auquel correspond par transposition un $\sh{O}_X$-homomorphisme ${}^t u:\check{\sh{E}}\to\sh{O}_X$, et par suite un homomorphisme de $\sh{O}_X$-Algèbres $v:\mathbf S_{\sh{O}_X}(\check{\sh{E}})\to\sh{O}_X$. Soit $\operatorname{Al}(X,\sh{E},\mathscr J)$ l'ensemble des $u\in\Gamma(X,\sh{E})$ tels que $\mathscr J$ soit \emph{contenu dans le noyau} de $v$; il résulte aussitôt de ces définitions que
\[
  X'\longmapsto\operatorname{Al}(X',\sh{E}_{(X')},\mathscr J\otimes_{\sh{O}_X}\sh{O}_{X'})
\]
est un foncteur représenté par $\Spec(\mathbf S_{\sh{O}_X}(\check{\sh{E}})/\mathscr J)$. Si $X'=\Spec(A)$ est affine et tel que $\sh{E}_{(X')}$ soit isomorphe à $\sh{O}_{X'}^n$, $\Gamma(X',\sh{E}_{(X')})$ peut s'identifier à l'ensemble $A^n$, et $\mathscr J\otimes_{\sh{O}_X}\sh{O}_{X'}$ à un $\sh{O}_{X'}$-Module de la forme $\widetilde{\mathfrak J}$, où $\mathfrak J$ est un idéal de l'anneau de polynômes $A[T_1,\ldots,T_n]$; l'ensemble $\operatorname{Al}(X',\sh{E}_{(X')},\mathscr J\otimes_{\sh{O}_X}\sh{O}_{X'})$ s'identifie alors à la partie de $A^n$ formée des points $(t_1,\ldots,t_n)$ tels que $P(t_1,\ldots,t_n)=0$ pour tous les polynômes $P\in\mathfrak J$; de façon imagée on peut donc dire que l'objet $\Spec(\mathbf S_{\sh{O}_X}(\check{\sh{E}})/\mathscr J)$ représente « le sous-ensemble algébrique de l'espace affine tordu $\Gamma(X,\sh{E})$ formé des points annulant l'idéal $\Gamma(X,\mathscr J)$ ». On note encore ce $X$-préschéma $\mathbf{Al}(\sh{E},\mathscr J)$. On notera que si $\mathscr J$ est un Idéal de type fini de $\mathbf S_{\sh{O}_X}(\check{\sh{E}})$, $\mathbf{Al}(\sh{E},\mathscr J)$ est un $X$-préschéma \emph{de présentation finie}, puisque $\mathbf S_{\sh{O}_X}(\check{\sh{E}})$ est une $\sh{O}_X$-Algèbre de présentation finie.
\end{env}

\begin{lemma}[18.5.3]
\label{IV.18.5.3}
Soient $S$ un préschéma, $f:X\to S$ un morphisme fini et localement libre (18.2.7). Considérons le foncteur contravariant de la catégorie des $S$-préschémas dans la catégorie des ensembles
\begin{equation}
  S'\longmapsto\operatorname{Of}(X\times_S S')
  \tag{18.5.3.1}\label{IV.18.5.3.1}
\end{equation}
où $\operatorname{Of}(X\times_S S')$ est l'ensemble des parties à la fois ouvertes et fermées de l'espace sous-jacent à $X\times_S S'$. Alors ce foncteur est représentable par un $S$-préschéma $\mathbf{Of}(X)$, qui est affine, étale et de présentation finie sur $S$.
\end{lemma}

\oldpage[IV-4]{127}

On a par hypothèse $X=\Spec(\mathscr B)$, où $\mathscr B=f_*(\sh{O}_X)$ est une $\sh{O}_S$-Algèbre finie et localement libre. Posons $X'=X\times_S S'$, $f'=f_{(S')}$, $\mathscr B'=\mathscr B\otimes_{\sh{O}_S}\sh{O}_{S'}=f'_*(\sh{O}_{X'})$, de sorte que $X'=\Spec(\mathscr B')$; il y a alors une bijection canonique, fonctorielle en $S'$, de l'ensemble $\operatorname{Of}(X')$ sur l'ensemble $\operatorname{Id}(\mathscr B')$ des \emph{idempotents} de l'anneau $\Gamma(S',\mathscr B')=\Gamma(X',\sh{O}_{X'})$. En effet, en vertu de l'\emph{équivalence} de la catégorie des $S'$-schémas affines sur $S'$ et de la catégorie opposée de la catégorie des $\sh{O}_{S'}$-Algèbres quasi-cohérentes (II, 1.2.7 et 1.3.1), il y a correspondance biunivoque canonique entre les décompositions de $X'$ en \emph{somme} $X'_1\amalg X'_2$ de deux sous-préschémas induits sur des ouverts de $X'$ et les décompositions de $\mathscr B'$ en \emph{composée directe} de deux Idéaux $\mathscr B'_1$, $\mathscr B'_2$; ces dernières à leur tour forment un ensemble en correspondance biunivoque canonique (et fonctorielle en $S'$) avec $\operatorname{Id}(\mathscr B')$. Il suffit donc de prouver le lemme pour le foncteur $S'\longmapsto\operatorname{Id}(\mathscr B')$.

Pour cela, nous allons montrer qu'il existe un Idéal de type fini $\mathscr J$ de $\mathbf S_{\sh{O}_S}(\check{\mathscr B})$ tel que $\operatorname{Id}(\mathscr B')$ soit de la forme $\operatorname{Al}(S',\mathscr B',\mathscr J\otimes_{\sh{O}_S}\sh{O}_{S'})$ (18.5.2). Notons à cet effet que puisque $\mathscr B$ est une $\sh{O}_S$-Algèbre, son image réciproque $\mathscr B_{(\mathbf V(\check{\mathscr B}))}$ est une $\sh{O}_{(\mathbf V(\check{\mathscr B}))}$-Algèbre, et l'on peut donc former le \emph{carré} $c^2$, dans cette Algèbre, de la section canonique $c$ (18.5.1); elle correspond canoniquement à un homomorphisme $u^{(2)}:\mathbf S_{\sh{O}_S}(\check{\mathscr B})\to\mathbf S_{\sh{O}_S}(\check{\mathscr B})\otimes_{\sh{O}_S}\mathscr B$ de $\mathbf S_{\sh{O}_S}(\check{\mathscr B})$-Modules, et on vérifie aussitôt que pour un ouvert affine $W$ de $S$, avec les notations de (18.5.1), $u^{(2)}$ correspond à l'homomorphisme de $C$-modules tel que
\[
  u^{(2)}(1)=\sum_k\Bigl(\sum_{i,j}c_{ijk}T_iT_j\Bigr)\otimes e_k,
\]
où $(c_{ijk})$ est la table de multiplication de l'algèbre $\Gamma(W,\mathscr B)$. Montrons que l'Idéal $\mathscr J$ de $\mathbf S_{\sh{O}_S}(\check{\mathscr B})$ engendré par le \emph{noyau} de l'homomorphisme $u-u^{(2)}$ répond à la question. Il suffit en effet de noter que l'idéal $\Gamma(W,\mathscr J)$ de $C$ est engendré par les polynômes $P_k(T_1,\ldots,T_n)=T_k-\sum_{i,j}c_{ijk}T_iT_j$ et que les idempotents de $\Gamma(W,\mathscr B)$ sont précisément les éléments $\sum_i t_i e_i$ de cette algèbre tels que $(t_1,\ldots,t_n)$ annule tous les polynômes $P_k$ ($1\leq i\leq n$). On déduit donc de (18.5.2) que le $S$-préschéma affine de présentation finie $\mathbf{Of}(X)=\mathbf{Al}(\mathscr B,\mathscr J)$ représente bien le foncteur (18.5.3.1). Il reste à prouver que $\mathbf{Of}(X)$ est étale sur $S$, ou, ce qui revient au même, qu'il est \emph{formellement étale} sur $S$. Mais si $S'$ est un $S$-préschéma, $S'_0$ un sous-préschéma fermé de $S'$ défini par un Idéal localement nilpotent de $\sh{O}_{S'}$ (et ayant donc même espace sous-jacent que $S'$), il est clair que $X'=X\times_S S'$ et $X'_0=X\times_S S'_0$ ont même espace sous-jacent, donc l'application canonique $\operatorname{Of}(X')\to\operatorname{Of}(X'_0)$ est \emph{bijective}, ce qui achève la démonstration (17.1.1).

\begin{proposition}[18.5.4]
\label{IV.18.5.4}
Soient $S$ un préschéma, $S_0$ un sous-préschéma fermé de $S$; considérons les propriétés suivantes:
\begin{enumerate}
  \item[a)] Pour tout morphisme fini $g:S'\to S$, l'application canonique
  \begin{equation}
    \operatorname{Of}(S')\longrightarrow\operatorname{Of}(S'\times_S S_0)
    \tag{18.5.4.1}\label{IV.18.5.4.1}
  \end{equation}
  (cf. (18.5.3)) est bijective.
  \item[a')] Pour tout morphisme fini et localement libre $g:S'\to S$, l'application (18.5.4.1) est bijective.

\oldpage[IV-4]{128}

  \item[b)] Pour tout morphisme étale et séparé $g:S'\to S$, l'application canonique
  \begin{equation}
    \Gamma(S'/S)\longrightarrow\Gamma(S'\times_S S_0/S_0)
    \tag{18.5.4.2}\label{IV.18.5.4.2}
  \end{equation}
  est bijective.
\end{enumerate}
La condition $b)$ entraîne $a')$; si de plus $S$ est quasi-compact et quasi-séparé, la condition $a)$ entraîne $b)$.
\end{proposition}

Prouvons d'abord que $b)$ entraîne $a')$. Supposons donc $b)$ vérifiée, et soit $g:S'\to S$ un morphisme fini et localement libre; posons $S'_0=S'\times_S S_0$, de sorte que $g_0=g_{(S_0)}:S'_0\to S_0$ est fini et localement libre. Alors il résulte de (18.5.3) que $P=\mathbf{Of}(S')$ est un $S$-préschéma \emph{étale et séparé}; en outre, la définition du foncteur $\operatorname{Of}$ montre aussitôt que si l'on pose $P_0=\mathbf{Of}(S'_0)$ (pour la catégorie des $S_0$-préschémas), on a $P_0=P\times_S S_0$. Cela étant, on a par définition le diagramme commutatif
\[
\xymatrix{
  \Gamma(P/S)\ar[r]\ar[d]^{\sim} & \Gamma(P_0/S_0)\ar[d]^{\sim}\\
  \operatorname{Of}(S')\ar[r] & \operatorname{Of}(S'_0)
}
\]
où les flèches verticales sont les bijections canoniques. Comme l'hypothèse $b)$, appliquée au morphisme $P\to S$, entraîne que la première ligne est une bijection, il en est de même de la seconde, ce qui établit notre assertion.

Avant de prouver que $a)$ entraîne $b)$ lorsque $S$ est quasi-compact et quasi-séparé, nous établirons le

\begin{lemma}[18.5.4.3]
\label{IV.18.5.4.3}
Si $S$ et $S_0$ vérifient la condition $a)$ de (18.5.4), alors, pour tout morphisme fini $g:S'\to S$, $S'$ est l'unique voisinage de $S'_0=g^{-1}(S_0)=S'\times_S S_0$ dans $S'$.
\end{lemma}

En effet, il revient au même de dire que si $T'$ est une partie fermée de $S'$ telle que $T'\cap S'_0=\varnothing$, alors $T'=\varnothing$. Or, si l'on désigne encore par $T'$ un sous-préschéma fermé de $S'$ ayant $T'$ pour espace sous-jacent, le morphisme composé $h:T'\to S'\xrightarrow{g}S$ est fini et $h^{-1}(S_0)$ est vide; la condition $a)$ appliquée au morphisme $h$ entraîne que $T'$ est nécessairement vide.

Ce lemme étant démontré, prouvons d'abord que, sous l'hypothèse $a)$, l'application (18.5.4.2) est \emph{injective}. En effet, si $u'$, $u''$ sont deux $S$-sections de $S'$, le fait que le morphisme $S'\to S$ soit non ramifié entraîne que le préschéma des coïncidences de $u'$ et $u''$ est induit sur un ouvert $U$ de $S$ (17.4.6). Si les restrictions à $S_0$ de $u'$ et $u''$ sont les mêmes, le fait que $u'$ et $u''$ sont des immersions ouvertes (17.4.1) entraîne que $U$ contient $S_0$, donc est égal à $S$ en vertu du lemme (18.5.4.3) appliqué au cas où $S'=S$.

\oldpage[IV-4]{129}

Reste à montrer que sous l'hypothèse $a)$, l'application (18.5.4.2) est \emph{surjective} ($S$ étant quasi-compact et quasi-séparé). Soit donc $u_0:S_0\to S'_0$ une $S_0$-section de $S'_0$; $u_0(S_0)$ étant quasi-compact dans $S'$, peut être recouvert par un nombre fini d'ouverts affines $V_i$ tels que la restriction $g|V_i$ soit un morphisme de type fini; si $V$ est la réunion des $V_i$, on en conclut que $g|V:V\to S$ est un morphisme \emph{de présentation finie}, étant séparé et localement de présentation finie par hypothèse (1.6.1). Remplaçant $S'$ par $V$, on peut donc supposer que $g$ est de présentation finie. Comme $S$ est quasi-compact et quasi-séparé et $g$ quasi-fini et séparé (17.6.1), il résulte du « Main theorem » (8.12.6) que $g$ se factorise en $S'\xrightarrow{j}S''\xrightarrow{f}S$, où $j$ est une immersion ouverte et $f$ un morphisme \emph{fini}. Posons $S''_0=S''\times_S S_0$, $j_0=j_{(S_0)}:S'_0\to S''_0$, qui est une immersion ouverte et $f_0=f_{(S_0)}:S''_0\to S_0$, qui est un morphisme fini. Alors $u_0$ est aussi une $S_0$-section de $S''_0$. Comme $g_0:S'_0\to S_0$ est étale, $u_0$ est une immersion ouverte de $S_0$ dans $S'_0$ (17.4.1), donc $u_0(S_0)$ est ouvert dans $S'_0$, et \emph{a fortiori} dans $S''_0$; mais d'autre part, comme $f_0$ est un morphisme fini, donc séparé, $u_0$ est une immersion fermée de $S_0$ dans $S''_0$ (I, 5.4.6), donc $X_0=u_0(S_0)$ est \emph{à la fois ouvert et fermé} dans $S''_0$. En vertu de l'hypothèse $a)$, il existe un sous-ensemble à la fois ouvert et fermé $X$ de $S''$ tel que $X\cap S''_0=X_0$. Montrons d'abord que le morphisme $f:S''\to S$ est \emph{étale aux points de $X$}: en effet, l'ensemble $U$ des points de $X$ où $f|X$ est étale est ouvert et contient par hypothèse $X_0\subset S'_0$. Mais le lemme (18.5.4.3) appliqué au morphisme fini $f|X$ prouve que $U=X$. D'autre part, $S'\cap X$ est ouvert dans $X$ et contient $X_0$ par hypothèse, donc le même raisonnement prouve que $S'\cap X=X$, c'est-à-dire $X\subset S'$. Il reste à montrer que, pour tout $s\in S$, le nombre géométrique $n(s)$ de points de $X\cap f^{-1}(s)$ est égal à $1$, car il en résultera que $f|X$ est radiciel et surjectif, et comme $f|X$ est étale, on aura montré (17.9.1) que $f|X$ est un isomorphisme de l'ouvert $X\subset S'$ sur $S$, dont l'isomorphisme réciproque $u$ sera la $S$-section cherchée prolongeant $u_0$. Mais comme $f|X$ est étale et \emph{fini}, $s\mapsto n(s)$ est continue dans $S$ (18.2.8), et comme $X\cap S''_0=X_0$, on a $n(s)=1$ dans $S_0$; l'ensemble des points $s\in S$ tels que $n(s)=1$ étant ouvert dans $S$ et contenant $S_0$, il est égal à $S$ par (18.5.4.3). C.Q.F.D.

\begin{remark}[18.5.4.4]
\label{IV.18.5.4.4}
On peut montrer que l'énoncé (18.5.4) reste valable lorsque, dans la condition $b)$, on suppose seulement le morphisme $g$ étale (mais non nécessairement séparé) [43, exp. XII].
\end{remark}

\begin{definition}[18.5.5]
\EGAReaderTarget{ega4.c04.def.18-5-5}{IV.18.5.5}
On dit qu'un préschéma $S$ et un sous-préschéma fermé $S_0$ de $S$ forment un \emph{couple hensélien} s'ils vérifient la condition $a)$ de (18.5.4).
\end{definition}

Compte tenu de (I, 5.1.8), il revient au même de dire que $(S,S_0)$ est un couple hensélien ou que $(S_{\mathrm{red}},(S_0)_{\mathrm{red}})$ en est un.

\begin{proposition}[18.5.6]
\label{IV.18.5.6}
\begin{enumerate}
  \item[(i)] Si $(S,S_0)$ est un couple hensélien, alors, pour tout morphisme fini $f:S'\to S$, si $S'_0$ est le sous-préschéma $f^{-1}(S_0)$ de $S'$, le couple $(S',S'_0)$ est hensélien.
  \item[(ii)] Soient $S=\coprod_\alpha S^{(\alpha)}$ une somme de préschémas, $S_0$ un sous-préschéma fermé de $S$, somme des sous-préschémas fermés $S_0^{(\alpha)}$ des $S^{(\alpha)}$. Pour que le couple $(S,S_0)$ soit hensélien, il faut et il suffit que chacun des couples $(S^{(\alpha)},S_0^{(\alpha)})$ le soit.
\end{enumerate}
\end{proposition}

L'assertion (i) est conséquence immédiate de la définition, puisque pour tout morphisme fini $g:S''\to S'$, $f\circ g:S''\to S$ est un morphisme fini. De même, sous les conditions de (ii), pour qu'un morphisme $g:S'\to S$ soit fini, il faut et il suffit que chacune

\oldpage[IV-4]{130}

de ses restrictions $g^{(\alpha)}:S'^{(\alpha)}=g^{-1}(S^{(\alpha)})\to S^{(\alpha)}$ le soit, et si l'on pose $S'_0=g^{-1}(S_0)$, $S'^{(\alpha)}_0=(g^{(\alpha)})^{-1}(S^{(\alpha)}_0)$, il y a correspondance biunivoque entre les parties ouvertes et fermées $U$ (resp. $U_0$) de $S'$ (resp. $S'_0$) et les familles $(U^{(\alpha)})$ (resp. $U^{(\alpha)}_0$) où $U^{(\alpha)}$ (resp. $U^{(\alpha)}_0$) est une partie ouverte et fermée de $S'^{(\alpha)}$ (resp. $S'^{(\alpha)}_0$), d'où l'assertion (ii).

\begin{remark}[18.5.7]
\label{IV.18.5.7}
Soient $S=\Spec(A)$ un schéma affine, $S_0$ un sous-schéma fermé de $S$ défini par un idéal $\mathfrak J$ de $A$. Alors, si le couple $(S,S_0)$ est hensélien, l'idéal $\mathfrak J$ est nécessairement contenu dans le \emph{radical} de $A$. En effet, si $\mathfrak m$ est un idéal maximal de $A$, $\mathfrak m$ doit appartenir à $V(\mathfrak J)=S_0$, en vertu de (18.5.4.3), autrement dit on doit avoir $\mathfrak J\subset\mathfrak m$, d'où la conclusion. En particulier, supposons que $S_0$ soit \emph{réduit à un point}, c'est-à-dire que l'idéal $\mathfrak J$ soit \emph{maximal}; alors $\mathfrak J$ doit être le radical de $A$, autrement dit $A$ doit être un anneau \emph{local}, $S_0$ étant l'unique point fermé de $\Spec(A)$.
\end{remark}

\begin{definition}[18.5.8]
\EGAReaderTarget{ega4.c04.def.18-5-8}{IV.18.5.8}
On dit qu'un anneau $A$ est \emph{hensélien} s'il est semi-local et si, en désignant par $\mathfrak r$ le radical de $A$, le couple $(\Spec(A),\Spec(A/\mathfrak r))$ est hensélien. On appelle \emph{schéma local hensélien} un schéma isomorphe au spectre d'un anneau local hensélien.
\end{definition}

\begin{proposition}[18.5.9]
\label{IV.18.5.9}
\begin{enumerate}
  \item[(i)] Pour qu'un anneau semi-local $A$ soit hensélien, il faut et il suffit qu'il soit composé direct d'anneaux locaux henséliens.
  \item[(ii)] Pour qu'un anneau local $A$ soit hensélien, il faut et il suffit que toute $A$-algèbre finie $B$ soit isomorphe à un produit d'anneaux locaux.
\end{enumerate}
\end{proposition}

(i) En effet, la définition, appliquée à $S=\Spec(A)$ et $S_0=\Spec(A/\mathfrak r)$ montre, puisque $S_0$ est un ensemble fini discret et fermé dans $S$, que $S$ est réunion d'un nombre fini de parties ouvertes et fermées $S_i$ ($1\leq i\leq n$) deux à deux disjointes, dont chacune contient exactement un des idéaux maximaux $\mathfrak m_i$ de $A$; la conclusion résulte de (18.5.6, (ii)) et de la remarque (18.5.7).

(ii) Tout morphisme fini $S'\to\Spec(A)$ est de la forme $\Spec(B)\to\Spec(A)$, où $B$ est une $A$-algèbre finie. Si $k$ est le corps résiduel de $A$, $\Spec(B\otimes_A k)$ est un spectre d'anneau artinien, donc fini et discret. Dire que le couple $(\Spec(A),\Spec(k))$ est hensélien signifie donc que $B$ est composé direct d'anneaux $A_i$ tels que $\Spec(A_i\otimes_A k)$ soit réduit à un point, c'est-à-dire que $A_i$ (qui est une $A$-algèbre \emph{finie}) ne doit avoir qu'un seul idéal maximal (Bourbaki, \emph{Alg. comm.}, chap. V, \textsection~2, n\textsuperscript{o}~1, prop. 1).

L'étude des anneaux henséliens est donc essentiellement ramenée à celle des anneaux \emph{locaux} henséliens.

\begin{proposition}[18.5.10]
\label{IV.18.5.10}
Si $A$ est un anneau hensélien, toute $A$-algèbre finie $B$ est un anneau hensélien (donc composé direct d'anneaux locaux henséliens (18.5.9)).
\end{proposition}

En effet, si $\mathfrak r'$ est le radical de $B$, l'image réciproque de $\mathfrak r'$ dans $A$ est le radical $\mathfrak r$ de $A$, et tout idéal de $B$ au-dessus d'un idéal maximal de $A$ est un idéal maximal de $B$, donc l'ensemble $V(\mathfrak r')$ dans $\Spec(B)$ est l'image réciproque de l'ensemble $V(\mathfrak r)$ dans $\Spec(A)$. La proposition est alors une conséquence de (18.5.6, (i)).

\begin{theorem}[18.5.11]
\label{IV.18.5.11}
Soient $A$ un anneau local, $\mathfrak m$ son idéal maximal. Les conditions suivantes sont équivalentes:
\begin{enumerate}
  \item[a)] $A$ est hensélien, autrement dit, toute $A$-algèbre finie $B$ est isomorphe à un produit d'anneaux locaux.

\oldpage[IV-4]{131}

  \item[a')] La condition $a)$ est satisfaite pour toutes les $A$-algèbres $B$ de la forme $A[T]/F.A[T]$, où $F\in A[T]$ est un polynôme unitaire.
  \item[b)] Soient $S=\Spec(A)$, $S_0=\Spec(k)$. Pour tout morphisme étale $g:S'\to S$, si l'on pose $S'_0=S'\otimes_A k$, toute $S_0$-section $u_0$ de $S'_0$ est la restriction d'une $S$-section $u$ de $S'$.
  \item[c)] Pour tout morphisme $f:X\to S$, séparé et localement de type fini, et tout point $x\in X$ tel que $f(x)$ soit égal au point fermé $s$ de $S$ et que $f$ soit quasi-fini au point $x$ (Err$_{\mathrm{III}}$, 20), $X$ est somme de deux préschémas $X'$, $X''$ tels que $X'=\Spec(\sh{O}_{X,x})$ et que $f|X':X'\to S$ soit un morphisme fini.
  \item[c')] Pour tout morphisme localement de type fini $f:X\to S$, et tout point $x\in X$ tel que $f$ soit quasi-fini au point $x$, et que $f(x)$ soit égal au point fermé $s$ de $S$, $\sh{O}_{X,x}$ est une algèbre finie sur $\sh{O}_{S,s}=A$.
  \item[c'')] Pour tout morphisme localement de présentation finie $f:X\to S$, et tout point $x\in X$ tel que $f$ soit quasi-fini au point $x$, et que $f(x)$ soit égal au point fermé $s$ de $S$, $\sh{O}_{X,x}$ est une algèbre finie et de présentation finie sur $\sh{O}_{S,s}=A$.
\end{enumerate}
\end{theorem}

Notons d'abord que la condition $c')$ (resp. $c'')$ est équivalente à la même condition où on suppose de plus $f$ \emph{séparé}, la question étant locale sur $X$. De même, la condition $b)$ est équivalente à la même condition où l'on suppose de plus $g$ \emph{séparé}: en effet, il suffit d'appliquer cette dernière à la restriction de $g$ à un voisinage ouvert affine du point $u_0(S_0)$ dans $S'$. Bornons-nous donc désormais au cas où, dans $b)$, $c')$ et $c'')$, les morphismes donnés sont \emph{séparés}.

Le fait que $a)$ implique $b)$ et que $b)$ implique $a')$ est un cas particulier de (18.5.4). Montrons en outre que $a')$ implique $a)$. Il s'agit de prouver que si $e_0$ est un idempotent de $C=B\otimes_A k$, il existe un idempotent $e\in B$ dont $e_0$ est l'image canonique. Si $b$ est un élément de $B$ dont l'image dans $C$ est $e_0$, la sous-$A$-algèbre $A[b]=B'$ de $B$ est finie, et l'image canonique $C'$ de $B'\otimes_A k$ dans $C$ contient $e_0$. Or, $B'\otimes_A k$ est une $k$-algèbre finie, donc composée directe de $k$-algèbres locales finies, et par suite $e_0$ est l'image dans $C'$ d'un idempotent $e'_0$ de $B'\otimes_A k$. On est ainsi ramené au cas où $B$ est monogène, et par suite isomorphe à une $A$-algèbre quotient d'une algèbre de la forme $A[T]/F.A[T]$, où $F$ est un polynôme unitaire. Or, en vertu de $a')$, $A[T]/F.A[T]$ est composée directe d'anneaux locaux, donc il en est de même de toutes ses algèbres quotients, ce qui achève de prouver l'existence de l'idempotent $e$.

Il est immédiat que $c)$ entraîne $a)$, comme on le voit en raisonnant par récurrence sur le nombre d'idéaux maximaux de l'anneau semi-local $B$. Pour voir que $a)$ implique $c)$, on peut, en vertu de (13.1.4), se borner au cas où $f$ est un morphisme affine et quasi-fini. Alors, par application du « Main theorem » (8.12.8), $f$ peut s'écrire comme un morphisme composé
\[
  X\xrightarrow{j}Y\xrightarrow{g}S,
\]
où $g$ est un morphisme fini et $j$ une immersion ouverte. Comme $Y=\Spec(B)$, où $B$ est une $A$-algèbre finie, il résulte de $a)$ que $B$ est composée directe d'anneaux locaux, qui sont évidemment des $A$-algèbres finies, et $\sh{O}_{X,x}$ s'identifie à l'un de ces anneaux locaux puisque $f(x)=s$; en outre, tout ouvert de $Y$ contenant $x$ contient nécessairement $\Spec(\sh{O}_{X,x})$.

Il est trivial que $c)$ implique $c')$, en vertu de la remarque du début. Prouvons que $c')$ implique $c'')$. Supposons en effet $c')$ vérifiée, et prouvons que sous les conditions

\oldpage[IV-4]{132}

de $c'')$, l'ensemble $Z=\Spec(\sh{O}_{X,x})$ s'identifie alors à une partie \emph{ouverte et fermée} de $X$, ce qui établira $c'')$ (I, 2.4.2). En premier lieu, le morphisme composé $Z\xrightarrow{j}X\xrightarrow{f}S$, où $j$ est le morphisme canonique (I, 2.4.1), est fini et de présentation finie par hypothèse, et comme $f$ est séparé et localement de présentation finie, $j$ est aussi un morphisme fini et de présentation finie ((II, 6.1.5) et (I, 4.3)); $j$ est par suite un morphisme fermé (II, 6.1.10), ce qui prouve que $Z$ est fermé dans $X$. Il résulte alors de (I, 2.4.2) et (I, 4.2.2) que $j$ est une immersion fermée. Mais si $\mathscr J$ est l'idéal de $\sh{O}_X$ définissant $Z$, on a alors par hypothèse $\mathscr J_x=0$, donc aussi $\mathscr J_z=0$ en tout point $z$ d'un voisinage ouvert $V$ de $x$ dans $X$, puisque par hypothèse $\mathscr J$ est un $\sh{O}_X$-Module quasi-cohérent de type fini ((1.4.7) et ($0_{\mathrm I}$, 5.2.2)). Or, un tel voisinage contient $Z$, donc $Z$ est ouvert dans $X$ puisque $\mathscr J|V=0$.

Enfin, $c'')$ implique $a')$: en effet, si $B=A[T]/F.A[T]$, le morphisme $X=\Spec(B)\to\Spec(A)=Y$ est fini et de présentation finie, et la démonstration précédente montre que $B$ est composée directe des $A$-algèbres finies $\sh{O}_{X,x_i}$, où les $x_i$ sont les points de la fibre du point fermé de $Y$.

\begin{corollary}[18.5.12]
\label{IV.18.5.12}
Soient $A$ un anneau semi-local, $\mathfrak r$ son radical; posons $S=\Spec(A)$, $S_0=\Spec(A/\mathfrak r)$. Pour que $A$ soit hensélien, il faut et il suffit que, pour tout morphisme fini $f:X\to S$ et tout morphisme étale et séparé $g:Y\to S$, si l'on pose $X_0=X\times_S S_0$ et $Y_0=Y\times_S S_0$, l'application canonique
\[
  \Hom_S(X,Y)\longrightarrow\Hom_{S_0}(X_0,Y_0)
\]
soit bijective.
\end{corollary}

La suffisance de la condition résulte de l'équivalence de $a)$ et $b)$ dans (18.5.11), en appliquant cette condition au cas où $f=1_S$. Pour voir que la condition est nécessaire, notons que si $A$ est hensélien, le couple $(X,X_0)$ est hensélien par (18.5.6, (i)); en outre $\Hom_S(X,Y)=\Gamma(X\times_S Y/X)$ et $\Hom_{S_0}(X_0,Y_0)=\Gamma(X_0\times_{S_0}Y_0/X_0)$; comme $X\times_S Y$ est étale et séparé sur $X$, la conclusion résulte de ce que $a)$ implique $b)$ dans (18.5.4).

\begin{remark}[18.5.13]
\label{IV.18.5.13}
Les conditions équivalentes $a)$, $a')$, $b)$ du théorème (18.5.11) sont encore équivalentes à la suivante (« lemme de Hensel »):
\begin{enumerate}
  \item[a'')] Pour tout polynôme unitaire $F\in A[T]$, d'image canonique $F_0\in k[T]$, et toute décomposition $F_0=G_0H_0$ de $F_0$ en un produit de deux polynômes unitaires étrangers $G_0$, $H_0$ de $k[T]$, il existe un couple unique $(G,H)$ de polynômes unitaires de $A[T]$ possédant les propriétés suivantes: $G_0$ et $H_0$ sont les images canoniques respectives de $G$ et $H$, on a $F=GH$, et l'idéal de $A[T]$ engendré par $G$ et $H$ est égal à $A[T]$.
\end{enumerate}
\end{remark}

Nous établirons d'abord le lemme suivant:

\begin{lemma}[18.5.13.1]
\label{IV.18.5.13.1}
Soient $A$ un anneau local de corps résiduel $k$, $F\in A[T]$ un polynôme unitaire, $B$ la $A$-algèbre $A[T]/F.A[T]$. Il existe une correspondance canonique entre les décompositions de $B$ en composée directe de deux $A$-algèbres quotients $B'$, $B''$ et les décompositions $F=GH$ de $F$ en produit de deux polynômes unitaires $G$, $H$ de $A[T]$, tels que l'idéal engendré par $G$ et $H$ soit égal à $A[T]$; les algèbres quotients $B'$, $B''$ correspondant à un tel couple de polynômes $G$, $H$ sont respectivement $A[T]/H.A[T]$ et $A[T]/G.A[T]$.
\end{lemma}

\oldpage[IV-4]{133}

Si $F=GH$ et si $G$ et $H$ engendrent l'idéal $A[T]$, il y a deux polynômes $P$, $Q$ de $A[T]$ tels que $1=PG+QH$. On en déduit que l'intersection des idéaux principaux $\mathfrak a=G.A[T]$ et $\mathfrak b=H.A[T]$ est égale à $\mathfrak c=F.A[T]$: en effet, si $R\in G.A[T]\cap H.A[T]$, on peut écrire $R=PRG+QRH$; or $RH$ (resp. $RG$) est un multiple de $F$ puisque $R$ est un multiple de $G$ (resp. $H$), donc $R\in F.A[T]$. Comme $A[T]=\mathfrak a+\mathfrak b$, $A[T]/\mathfrak c$ est somme directe des idéaux $\mathfrak a/(\mathfrak a\cap\mathfrak b)$ et $\mathfrak b/(\mathfrak a\cap\mathfrak b)$, canoniquement isomorphes respectivement à $A[T]/\mathfrak b$ et $A[T]/\mathfrak a$.

Inversement, supposons donnée une décomposition de $B$ en composée directe de deux $A$-algèbres $B'$, $B''$, qui s'identifient canoniquement à deux idéaux $e'B$, $e''B$ de $B$, correspondant à une décomposition $1=e'+e''$ de $1$ en idempotents orthogonaux $e'$, $e''$ de $B$. Posons encore $B_0=B\otimes_A k=k[T]/F_0.k[T]$, où $F_0$ est l'image canonique de $F$ dans $k[T]$, de même degré $n$ que $F$; si $e'_0$, $e''_0$ sont les images canoniques de $e'$, $e''$ dans $B_0$, ce sont deux idempotents orthogonaux tels que $1=e'_0+e''_0$, et $B_0$ est donc composée directe de $B'_0=e'_0B_0$ et $B''_0=e''_0B_0$. Soient $t$ et $t_0$ les images canoniques de $T$ dans $B$ et $B_0$; comme $B'_0$ (resp. $B''_0$) est une $k$-algèbre finie engendrée par $t'_0=e'_0t_0$ (resp. $t''_0=e''_0t_0$), elle admet une base de la forme $\{e'_0,t'_0,t_0'{}^2,\ldots,t_0'{}^{s-1}\}$ (resp. $\{e''_0,t''_0,t_0''{}^2,\ldots,t_0''{}^{r-1}\}$) avec $r+s=n$. D'autre part, $B$ étant un $A$-module libre (de base $\{1,t,\ldots,t^{n-1}\}$), $B'$ et $B''$ sont des $A$-modules projectifs, donc libres puisque $A$ est un anneau local (Bourbaki, \emph{Alg. comm.}, chap. II, \textsection~5, n\textsuperscript{o}~3, cor. de la prop. 5); il résulte donc de ce qui précède et de Bourbaki, \emph{Alg. comm.}, chap. II, \textsection~3, n\textsuperscript{o}~3, prop. 5, que si l'on pose $t'=e't$, $t''=e''t$, $\{e',t',\ldots,t'^{s-1}\}$ (resp. $\{e'',t'',\ldots,t''^{r-1}\}$) est une base du $A$-module $B'$ (resp. $B''$). Il y a donc un polynôme unitaire $H$ (resp. $G$) de degré $s$ (resp. $r$) de $A[T]$ tel que $e'H(t')=0$ et $e''G(t'')=0$; comme $t^h=t'^h+t''^h$ pour tout entier $h\geq1$, et $t'^h=e't'^h$, $t''^h=e''t''^h$, on a aussi $G(t)=e'G(t')$ et $H(t)=e''H(t'')$, d'où $G(t)H(t)=0$; on en conclut que le polynôme $G(T)H(T)$ est divisible par $F(T)$; mais comme les degrés de ces deux polynômes unitaires sont les mêmes, on a $GH=F$. En outre, $B'$ (resp. $B''$) est isomorphe à $A[T]/H.A[T]$ (resp. $A[T]/G.A[T]$). Enfin, il y a deux polynômes $R$, $S$ de $A[T]$ tels que $e'=R(t)$ et $e''=S(t)$; comme $R(t)=e'R(t')+e''R(t'')$, on a nécessairement $e''R(t'')=0$, et de même $e'S(t')=0$, de sorte que, par définition de $G$ et $H$, $R=QH$ et $S=PG$, où $P$, $Q$ appartiennent à $A[T]$; la relation $1=R(t)+S(t)$ dans $B$ donne donc par définition $PG+QH=1+LF$ pour un certain polynôme $L\in A[T]$, et comme $F=GH$, cela prouve que l'idéal engendré par $G$ et $H$ est $A[T]$, et achève de démontrer le lemme.

Ce lemme étant établi, il suffit de l'appliquer à l'anneau local $A$ d'une part, au corps $k$ de l'autre, pour voir aussitôt que les conditions $a')$ et $a'')$ sont équivalentes.

\begin{proposition}[18.5.14]
\label{IV.18.5.14}
Tout anneau semi-local $A$, séparé et complet pour la topologie $\mathfrak r$-préadique (où $\mathfrak r$ est le radical de $A$) est hensélien.
\end{proposition}

En effet, $A$ est composé direct d'anneaux locaux séparés et complets (Bourbaki, \emph{Alg. comm.}, chap. III, \textsection~2, n\textsuperscript{o}~13, cor. de la prop. 19), donc on est ramené au cas où $A$ est un anneau local. Vérifions le critère $a')$ de (18.5.11). Comme $B$ est un $A$-module libre de type fini, il est évidemment séparé et complet pour la topologie $\mathfrak r$-préadique,

\oldpage[IV-4]{134}

qui est aussi la topologie $\mathfrak s$-préadique, où $\mathfrak s$ est le radical de l'anneau semi-local $B$, car $B/\mathfrak rB$ est un anneau artinien de radical $\mathfrak s/\mathfrak rB$. On sait alors (Bourbaki, \emph{Alg. comm.}, chap. III, \textsection~2, n\textsuperscript{o}~13, cor. de la prop. 19) que $B$ est composé direct d'anneaux locaux.

\begin{proposition}[18.5.15]
\label{IV.18.5.15}
Soient $A$ un anneau hensélien, $\mathfrak r$ son radical. Alors le foncteur $B\longmapsto B/\mathfrak rB$ est une équivalence de la catégorie des $A$-algèbres finies et étales avec la catégorie des $(A/\mathfrak r)$-algèbres finies et étales.
\end{proposition}

Le fait que le foncteur de l'énoncé soit pleinement fidèle est un cas particulier de (18.5.12). Pour montrer que ce foncteur est une équivalence de catégories, on peut se borner au cas où $A$ est local; il suffit alors d'appliquer (18.1.1) (pour les morphismes étales), à $S=\Spec(A)$ et $S_0=\Spec(A/\mathfrak r)$, réduit à un seul point.

\begin{remarks}[18.5.16]
\label{IV.18.5.16}
\begin{enumerate}
  \item[(i)] Nous ignorons si la proposition (18.5.15) se généralise à un couple hensélien $(S,S_0)$, même lorsque $S$ est affine et noethérien.
  \item[(ii)] Soient $A$ un anneau noethérien, $\mathfrak J$ un idéal de $A$ tel que $A$ soit \emph{séparé et complet} pour la topologie $\mathfrak J$-préadique. Alors le couple $(\Spec(A),\Spec(A/\mathfrak J))$ est hensélien: en effet, pour toute $A$-algèbre finie $B$, $B$ est séparé et complet pour la topologie $\mathfrak J$-préadique ($0_{\mathrm I}$, 7.3.6). Remplaçant $B$ par $A$ et $\mathfrak J B$ par $\mathfrak J$, tout revient donc à voir que l'application qui, à tout idempotent de $A$, fait correspondre sa classe mod. $\mathfrak J$, est \emph{bijective}. Or on a $A=\varprojlim (A/\mathfrak J^n)$. Notons $\operatorname{Idem}(A)$ l'ensemble des idempotents de $A$, et pour tout homomorphisme d'anneaux $\varphi:A\to B$, soit $\operatorname{Idem}(\varphi)$ l'application de $\operatorname{Idem}(A)$ dans $\operatorname{Idem}(B)$ restriction de $\varphi$; il résulte de la définition de la limite projective que l'on a
  \[
    \operatorname{Idem}(A)=\varprojlim_n\operatorname{Idem}(A/\mathfrak J^n)
  \]
  pour les applications $\psi_{nm}:\operatorname{Idem}(A/\mathfrak J^m)\to\operatorname{Idem}(A/\mathfrak J^n)$ restrictions des applications canoniques $A/\mathfrak J^m\to A/\mathfrak J^n$. Mais puisque $\Spec(A/\mathfrak J^n)\to\Spec(A/\mathfrak J^m)$ est un homéomorphisme, les $\psi_{nm}$ sont des bijections (comme on l'a vu dans la démonstration de (18.5.3)); cela prouve donc notre assertion.
\end{enumerate}
\end{remarks}

\begin{theorem}[18.5.17]
\label{IV.18.5.17}
Soient $A$ un anneau local hensélien, $S=\Spec(A)$, $s$ le point fermé de $S$. Pour tout morphisme lisse $f:X\to S$, l'application canonique
\[
  \Gamma(X/S)\longrightarrow\Gamma(X_s/k(s))
\]
(où $X_s=f^{-1}(s)$) est surjective.
\end{theorem}

La donnée d'une $k(s)$-section de $X_s$ équivaut à celle d'un point $x\in X$ au-dessus de $s$ rationnel sur $k(s)$, et il s'agit de prouver qu'il existe une $S$-section $u:S\to X$ telle que $u(s)=x$. Compte tenu de (17.16.3, (i)), on peut supposer que $f$ est étale. Alors la conclusion résulte du critère (18.5.11, $b)$). (Le lecteur notera qu'en vertu de ce critère, la validité de (18.5.17) pour un anneau local donné $A$ est nécessaire et suffisante pour que $A$ soit hensélien).

\begin{remark}[18.5.18]
\label{IV.18.5.18}
Les conditions de (18.5.11) sont encore équivalentes à la suivante:
\begin{enumerate}
  \item[d)] Pour tout morphisme localement de type fini $f:X\to S$ et tout point $x\in X$ tel que $f(x)$ soit le point fermé $s$ de $S$ et que $\sh{O}_{X,x}/\mathfrak m_s\sh{O}_{X,x}$ soit un corps canoniquement isomorphe à $k(s)=k$, il existe un voisinage ouvert $U$ de $x$ dans $X$ tel que $f|U$ soit une immersion fermée.
\end{enumerate}
\end{remark}

\oldpage[IV-4]{135}

Il est immédiat que $d)$ implique la condition $b)$ de (18.5.11) car une immersion fermée étale est une immersion ouverte (18.9.1). Inversement, supposons vérifiées les conditions de (18.5.11) et prouvons $d)$. L'hypothèse de $d)$ implique que $f$ est quasi-fini au point $x$ (Err$_{\mathrm{III}}$, 20), en vertu de (13.1.4). Donc en vertu de la condition $c)$ de (18.5.11), $\sh{O}_{X,x}$ est une $A$-algèbre finie et $\Spec(\sh{O}_{X,x})$ un voisinage ouvert $U$ de $x$ dans $X$. En outre, comme $\sh{O}_{X,x}/\mathfrak m_s\sh{O}_{X,x}$ est isomorphe à $k=A/\mathfrak m_s$, le lemme de Nakayama prouve que l'homomorphisme $A\to\sh{O}_{X,x}$ est surjectif, donc que $f|U$ est une immersion fermée.

\begin{proposition}[18.5.19]
\label{IV.18.5.19}
Soient $A$ un anneau local noethérien et hensélien, $S=\Spec(A)$, $s$ le point fermé de $S$, $f:X\to S$ un morphisme propre; posons $X_s=f^{-1}(s)$. Alors l'application $Y\longmapsto Y\cap X_s$ est une bijection de l'ensemble des composantes connexes de $X$ sur l'ensemble des composantes connexes de $X_s$.
\end{proposition}

En considérant un sous-préschéma fermé de $X$ ayant pour espace sous-jacent une composante connexe de $X$, on est ramené à prouver que si $X$ est \emph{connexe et non vide}, alors $X_s$ est \emph{connexe et non vide}. Le fait que $X_s$ soit non vide résulte de (II, 7.2.1); pour prouver que $X_s$ est connexe, raisonnons par l'absurde, en considérant la factorisation de Stein $f:X\xrightarrow{f'}S'\xrightarrow{g}S$ du morphisme propre $f$ (III, 4.3.3); par hypothèse, l'ensemble discret fini $g^{-1}(s)$ contiendrait au moins deux points. Puisque $A$ est hensélien et $g$ séparé et de type fini, il résulterait alors de (18.5.11, $c)$ que $S'$ serait somme de deux préschémas non vides $S'_1$, $S'_2$ (puisque l'intersection de l'un d'eux avec $g^{-1}(s)$ est réduite à un seul point). Puisque $f'$ est surjectif (III, 4.3.1), on en conclurait que $X$ est somme de deux préschémas non vides, contrairement à l'hypothèse.

\subsection{Hensélisation.}
\EGAReaderTarget{ega4.c04.subsec.18-6}{subsection:IV.18.6}

\begin{env}[18.6.1]
\EGAReaderTarget{ega4.c04.statement.18-6-1}{IV.18.6.1}
Étant donné un anneau local $A$, nous dirons qu'une $A$-algèbre locale $B$ est \emph{essentiellement étale} s'il existe une $A$-algèbre étale $C$ et un idéal premier $\mathfrak n$ de $C$ tels que $B$ soit $A$-isomorphe à $C_{\mathfrak n}$ et que l'homomorphisme composé $A\to C\to C_{\mathfrak n}$ soit \emph{local} (autrement dit, que $\mathfrak n$ soit au-dessus de l'idéal maximal $\mathfrak m$ de $A$); cela entraîne (17.6.1) que $\mathfrak n$ est le \emph{seul} idéal premier de $B$ au-dessus de $\mathfrak m$ puisque $B$ est un anneau local; si $B$, $B'$ sont deux $A$-algèbres locales essentiellement étales, tout $A$-homomorphisme $B\to B'$ est donc \emph{local}.

Si $A'$ est une $A$-algèbre locale essentiellement étale et $A''$ une $A'$-algèbre locale essentiellement étale, alors $A''$ est une $A$-algèbre locale essentiellement étale. En effet, on a par hypothèse $A'=B_{\mathfrak n}$, où $B$ est une $A$-algèbre étale et $\mathfrak n$ un idéal premier de $B$ au-dessus de l'idéal maximal de $A$, et $A''=B'_{\mathfrak n'}$, où $B'$ est une $A'$-algèbre étale et $\mathfrak n'$ un idéal premier de $B'$ au-dessus de l'idéal $\mathfrak nB_{\mathfrak n}$. Si l'on pose $S=B-\mathfrak n$, de sorte que $A'=S^{-1}B$, $B'$ est de la forme $S^{-1}C$, où $C$ est une $B$-algèbre de présentation finie. Par suite $C$ est une $A$-algèbre de présentation finie et $A''$ est de la forme $C_{\mathfrak r}$, où $\mathfrak r$ est un idéal premier de $C$ au-dessus de l'idéal maximal de $A$. Comme $A'$ est une $A$-algèbre formellement étale et $A''$ une $A'$-algèbre formellement étale, $A''$ est une $A$-algèbre formellement

\oldpage[IV-4]{136}

étale (pour les topologies discrètes (17.1.3)); le morphisme $\Spec(C)\to\Spec(A)$ est donc étale au point $\mathfrak r$ (17.6.1), et il y a par suite un élément $g\in C$ tel que $C_g$ soit une $A$-algèbre étale, ce qui prouve que $A''$ est une $A$-algèbre essentiellement étale.

Étant donné un anneau local $A$, il existe un \emph{ensemble} $\mathfrak E$ de $A$-algèbres locales essentiellement étales tel que toute $A$-algèbre locale essentiellement étale soit $A$-isomorphe à une algèbre appartenant à $\mathfrak E$. Il suffit évidemment pour le voir d'observer qu'il existe un ensemble $\mathfrak F$ de $A$-algèbres de type fini tel que toute $A$-algèbre de type fini soit isomorphe à une algèbre appartenant à $\mathfrak F$; on peut prendre $\mathfrak F$ égal à l'ensemble des quotients des algèbres de polynômes $A[T_1,\ldots,T_n]$ ($n\in\mathbf N$).

Dans ce qui suit, si $A$ et $B$ sont deux anneaux locaux, nous noterons $\Hom.\operatorname{loc}(A,B)$ l'ensemble des homomorphismes \emph{locaux} de $A$ dans $B$.
\end{env}

\begin{lemma}[18.6.2]
\EGAReaderTarget{ega4.c04.lem.18-6-2}{IV.18.6.2}
Soient $A$, $A'$ deux anneaux locaux, $k$, $k'$ leurs corps résiduels respectifs, $\varphi:A\to A'$ un homomorphisme (local) faisant de $A'$ une $A$-algèbre essentiellement étale (18.6.1) et tel que l'homomorphisme correspondant $k\to k'$ soit bijectif. Alors, pour tout anneau local hensélien $B$, l'application canonique
\[
  \Hom(\varphi,1_B):\Hom.\operatorname{loc}(A',B)\longrightarrow\Hom.\operatorname{loc}(A,B)
\]
est bijective.
\end{lemma}

Posons $S=\Spec(A)$, $Y=\Spec(B)$, et soient $s$ et $y$ les points fermés de $S$ et $Y$ respectivement. Par hypothèse, $A'$ est isomorphe à un anneau local $\sh{O}_{X,x}$ d'un $S$-schéma $X$ étale sur $S$, en un point $x\in X$ au-dessus de $s$. Supposons donné un homomorphisme local $\psi:A\to B$, faisant de $Y$ un $S$-schéma; il s'agit de voir qu'il existe \emph{un et un seul} $S$-morphisme $f:Y\to X$ tel que $f(y)=x$. Posons $X'=X\times_S Y$, et notons que puisque $k(x)=k(s)$, il existe un seul point $x'\in X'$ au-dessus de $x$ et de $y$, et que $k(x')=k(y)$. Il faut prouver qu'il existe une seule $Y$-section $f'$ de $X'$ telle que $f'(y)=x'$. Or, le morphisme $g:X'\to Y$ est étale et séparé, et la fibre $X'_0=g^{-1}(y)$ a pour anneau local au point $x'$ le corps $k(x')=k(y)$. Si l'on pose $Y_0=\Spec(k(y))$, il existe donc une unique $Y_0$-section $f'_0$ de $X'_0$ telle que $f'_0(y)=x'$, et la conclusion résulte de ce que $B$ est supposé hensélien et de (18.5.11, $b)$).

Nous dirons qu'une $A$-algèbre locale $A'$ vérifiant les conditions de (18.6.2) est \emph{strictement essentiellement étale}. Notons que le critère (18.5.11, $b)$) signifie que, \emph{pour que $A$ soit hensélien, il faut et il suffit que toute $A$-algèbre strictement essentiellement étale soit $A$-isomorphe à $A$}.

\begin{lemma}[18.6.3]
\label{IV.18.6.3}
Soient $A$ un anneau local, $A_1$, $A_2$ deux $A$-algèbres locales strictement essentiellement étales.
\begin{enumerate}
  \item[(i)] Il existe au plus un $A$-homomorphisme (nécessairement local) de $A_1$ dans $A_2$.
  \item[(ii)] Il existe une $A$-algèbre locale strictement essentiellement étale $A_3$ et deux $A$-homomorphismes $A_1\to A_3$, $A_2\to A_3$.
\end{enumerate}
\end{lemma}

Posons $S=\Spec(A)$; par hypothèse il y a deux $S$-schémas étales sur $S$, $X_1$, $X_2$, et deux points $x_1\in X_1$, $x_2\in X_2$ au-dessus du point fermé $s$ de $S$ et tels que $A_1=\sh{O}_{X_1,x_1}$, $A_2=\sh{O}_{X_2,x_2}$. Posons $X_3=X_1\times_S X_2$; les hypothèses $k(x_1)=k(x_2)=k(s)$ entraînent qu'il existe un seul point $x_3\in X_3$ au-dessus de $x_1$ et $x_2$ et que $k(x_3)=k(s)$ (I, 3.4.9). En outre, $X_3$ est étale sur $S$ (17.3.3), donc $A_3=\sh{O}_{X_3,x_3}$ vérifie les conditions de (ii). Par ailleurs,

\oldpage[IV-4]{137}

on a vu qu'un $A$-homomorphisme de $A_1$ dans $A_2$ est nécessairement local; il correspond à un $S$-morphisme $f$ de $X'_2=\Spec(A_2)$ dans $X_1$ tel que $f(x_2)=x_1$, ou encore, en posant $X'_3=X_1\times_S X'_2$, à une $X'_2$-section $f'$ de $X'_3$ telle que $f'(x_2)=x_3$. Comme $k(x_3)=k(x_2)$ et que $X'_2$ est connexe, l'unicité de $f'$ résulte de (17.4.9), d'où (i).

\begin{env}[18.6.4]
\label{IV.18.6.4}
Désignons par $\mathfrak S$ le sous-ensemble de l'ensemble $\mathfrak E$ défini dans (18.6.1) formé des $A$-algèbres strictement essentiellement étales appartenant à $\mathfrak E$. Il résulte de (18.6.3) que la relation « il existe un $A$-homomorphisme de $A_1$ dans $A_2$ » est une relation de \emph{préordre} dans $\mathfrak S$ et fait de $\mathfrak S$ un ensemble \emph{filtrant croissant}. En indexant $\mathfrak S$ par lui-même au moyen de l'application identique, il résulte de (18.6.3, (i)) que si $\lambda\leq\mu$ dans $\mathfrak S$, il existe un $A$-homomorphisme \emph{unique} $\varphi_{\mu\lambda}:A_\lambda\to A_\mu$ et $(A_\lambda,\varphi_{\mu\lambda})$ est évidemment un \emph{système inductif} de $A$-algèbres locales, les $\varphi_{\mu\lambda}$ étant des homomorphismes locaux. En outre, il résulte de (17.3.5) que pour $\lambda\leq\mu$, $A_\mu$ est une $A_\lambda$-algèbre strictement essentiellement étale.
\end{env}

\begin{definition}[18.6.5]
\EGAReaderTarget{ega4.c04.def.18-6-5}{IV.18.6.5}
On appelle \emph{hensélisé} d'un anneau local $A$ et on désigne par ${}^h\!A$ la $A$-algèbre limite inductive du système inductif $(A_\lambda,\varphi_{\mu\lambda})$ défini dans (18.6.4).
\end{definition}

Cette définition ne dépend qu'en apparence du choix de $\mathfrak E$; si $\mathfrak E'$ est un autre ensemble de $A$-algèbres locales essentiellement étales ayant la même propriété que $\mathfrak E$, $\mathfrak S'$ le sous-ensemble de $\mathfrak E'$ formé des $A$-algèbres strictement essentiellement étales, il résulte de (18.6.3, (ii)) que, pour les relations de préordre envisagées, $\mathfrak S$ et $\mathfrak S'$ sont \emph{cofinaux} à $\mathfrak S''=\mathfrak S\cup\mathfrak S'$, donc donnent la même limite inductive à un isomorphisme près. Nous allons voir d'ailleurs ci-dessous (18.6.6, (i) et (ii)) que ${}^h\!A$ et l'homomorphisme canonique $A\to{}^h\!A$ forment une solution d'un problème universel, et par suite sont déterminés à isomorphisme unique près.

On remarquera que si $A$ est un corps, on a évidemment ${}^h\!A=A$. Dans le cas général, ${}^h\!A$ est aussi le hensélisé de toutes les $A$-algèbres strictement essentiellement étales $A_\lambda$ ($\lambda\in\mathfrak S$), car si $B$ est une $A_\lambda$-algèbre strictement essentiellement étale, $B$ est aussi une $A$-algèbre strictement essentiellement étale (18.6.1), donc (à un $A$-isomorphisme près) une des $A_\mu$ pour $\mu\geq\lambda$, et les $A_\mu$ telles que $\mu\geq\lambda$ et que $A_\mu$ soit une $A_\lambda$-algèbre strictement essentiellement étale forment un ensemble cofinal dans l'ensemble préordonné des $A_\lambda$, en vertu de (18.6.1) et (18.6.3).

\begin{theorem}[18.6.6]
\label{IV.18.6.6}
Soient $A$ un anneau local, ${}^h\!A$ son hensélisé.
\begin{enumerate}
  \item[(i)] ${}^h\!A$ est un anneau local hensélien et l'homomorphisme structural $A\to{}^h\!A$ est local.
  \item[(ii)] Pour tout anneau local hensélien $B$, l'application canonique
  \[
    \Hom.\operatorname{loc}({}^h\!A,B)\longrightarrow\Hom.\operatorname{loc}(A,B)
  \]
  est bijective.
  \item[(iii)] ${}^h\!A$ est un $A$-module fidèlement plat, et si $\mathfrak m$ est l'idéal maximal de $A$, $\mathfrak m.{}^h\!A$ est l'idéal maximal de ${}^h\!A$, et l'homomorphisme $A/\mathfrak mA\to{}^h\!A/\mathfrak m.{}^h\!A$ des corps résiduels est bijective.
  \item[(iv)] Si $\widehat A$ et $({}^h\!A)\widehat{\phantom{x}}$ sont les séparés complétés des anneaux locaux $A$ et ${}^h\!A$, l'homomorphisme $\widehat A\to({}^h\!A)\widehat{\phantom{x}}$ déduit de l'homomorphisme structural $A\to{}^h\!A$ par complétion est bijectif.
  \item[(v)] Pour que ${}^h\!A$ soit noethérien, il faut et il suffit que $A$ le soit.
  \item[(vi)] Si $A$ est hensélien, l'homomorphisme canonique $A\to{}^h\!A$ est bijectif.
\end{enumerate}
\end{theorem}

\oldpage[IV-4]{138}

Soit $\mathfrak m_\lambda$ l'idéal maximal de $A_\lambda$; il résulte du fait que, dans (17.6.1), $a)$ implique $c')$, que pour $\lambda\leq\mu$, on a $\mathfrak m_\mu=\mathfrak m_\lambda A_\mu$, l'homomorphisme $A_\lambda/\mathfrak m_\lambda\to A_\mu/\mathfrak m_\mu$ est bijectif et $A_\mu$ est un $A_\lambda$-module plat. Le fait que ${}^h\!A$ soit local, que l'homomorphisme $A\to{}^h\!A$ soit local, l'assertion (iii) et la suffisance de (v) résultent donc de ($0_{\mathrm{III}}$, 10.3.1.3); la nécessité de (v) résulte de ce que ${}^h\!A$ est un $A$-module fidèlement plat ($0_{\mathrm I}$, 6.5.2).

Pour prouver que ${}^h\!A$ est hensélien, appliquons le critère (18.5.11, $b)$). Posons donc $S=\Spec({}^h\!A)$, $S_0=\Spec({}^h\!A/\mathfrak m.{}^h\!A)$, et soit $g:S'\to S$ un morphisme étale; posons $S'_0=g^{-1}(S_0)$, et soit $f_0:S_0\to S'_0$ une $S_0$-section de $S'_0$. Raisonnant comme dans (18.5.4), on peut supposer que $g$ est de présentation finie. Il résulte alors de (8.8.2) et (17.7.5) qu'il existe un indice $\lambda\in\mathfrak S$, un morphisme étale $g_\lambda:S'^{(\lambda)}\to S^{(\lambda)}=\Spec(A_\lambda)$ et, si l'on pose $S_0^{(\lambda)}=\Spec(A_\lambda/\mathfrak mA_\lambda)$ et $S_0'^{(\lambda)}=g_\lambda^{-1}(S_0^{(\lambda)})$, une $S_0^{(\lambda)}$-section $f_0^{(\lambda)}:S_0^{(\lambda)}\to S_0'^{(\lambda)}$ tels que
\[
  S'=S'^{(\lambda)}\times_{S^{(\lambda)}}S,\qquad
  g=g_\lambda\times1,\qquad
  f_0=f_0^{(\lambda)}\times1.
\]
Soient $s$ le point fermé de $S$, $x=f_0(s)$, $x_\lambda$ la projection de $x$ dans $S'^{(\lambda)}$; comme $x_\lambda$ est au-dessus du point fermé $s_\lambda$ de $S^{(\lambda)}$, l'anneau local $C_\lambda$ de $S'^{(\lambda)}$ au point $x_\lambda$ est une $A_\lambda$-algèbre essentiellement étale; en outre, comme $f_0^{(\lambda)}(s_\lambda)=x_\lambda$, on a $k(x_\lambda)=k(s_\lambda)$, autrement dit, $C_\lambda$ est une $A_\lambda$-algèbre locale strictement essentiellement étale, et par suite (18.6.1) est $A_\lambda$-isomorphe à une $A$-algèbre $A_\mu$ avec $\mu\geq\lambda$. Il y a donc un $A_\lambda$-homomorphisme $C_\lambda\to{}^h\!A$, c'est-à-dire un $S^{(\lambda)}$-morphisme $h:S\to S'^{(\lambda)}$ tel que $h(s)=x_\lambda$, et par suite il existe bien une $S$-section $f$ de $S'$ telle que $f(s)=x$, ce qui achève de prouver (i).

Pour prouver (ii), il suffit de noter que l'on a
\[
  \Hom.\operatorname{loc}({}^h\!A,B)=\varprojlim_\lambda\Hom.\operatorname{loc}(A_\lambda,B)
\]
et que par (18.6.2) les homomorphismes canoniques $\Hom.\operatorname{loc}(A_\lambda,B)\leftarrow\Hom.\operatorname{loc}(A,B)$ sont bijectifs.

Pour démontrer (iv), notons que pour tout $\lambda$ et tout entier $n>0$, $\mathfrak m_\lambda^n=\mathfrak m^nA_\lambda$, et
\[
  (\mathfrak m.{}^h\!A)^n=\mathfrak m^n.{}^h\!A=\varinjlim_\lambda\mathfrak m_\lambda^n.
\]
On en conclut, par l'exactitude du foncteur $\varinjlim$ dans la catégorie des $A$-modules, que
\[
  {}^h\!A/(\mathfrak m.{}^h\!A)^n=\varinjlim_\lambda(A_\lambda/\mathfrak m_\lambda^n),
\]
et par suite, il suffit de montrer que, pour tout entier $n$, et tout indice $\lambda$, l'homomorphisme $A/\mathfrak m^n\to A_\lambda/\mathfrak m_\lambda^n$ est bijectif. Or, cela est vrai par hypothèse pour $n=1$; d'autre part, puisque $A_\lambda$ est un $A$-module plat, on a
\[
  \mathfrak m_\lambda^n/\mathfrak m_\lambda^{n+1}
  =(\mathfrak m^n/\mathfrak m^{n+1})\otimes_A A_\lambda
  =(\mathfrak m^n/\mathfrak m^{n+1})\otimes_{A/\mathfrak m}(A_\lambda/\mathfrak mA_\lambda)
\]
et comme $A/\mathfrak m\to A_\lambda/\mathfrak mA_\lambda=A_\lambda/\mathfrak m_\lambda$ est bijectif, l'homomorphisme $\mathfrak m^n/\mathfrak m^{n+1}\to\mathfrak m_\lambda^n/\mathfrak m_\lambda^{n+1}$ est lui aussi bijectif; la conclusion résulte donc de Bourbaki, \emph{Alg. comm.}, chap. III, \textsection~2, n\textsuperscript{o}~8, cor. 3 du th. 1.

Enfin, si $A$ est hensélien, il résulte de la remarque qui précède (18.6.3) que les homomorphismes $A\to A_\lambda$ sont bijectifs, ce qui prouve (vi) par définition de ${}^h\!A$.

\begin{env}[18.6.7]
\EGAReaderTarget{ega4.c04.statement.18-6-7}{IV.18.6.7}
Soient maintenant $A$ un anneau semi-local, $\mathfrak m_i$ ($1\leq i\leq r$) ses idéaux maximaux. On appelle \emph{hensélisé} de $A$ et l'on note encore ${}^h\!A$ l'anneau produit $\prod_{i=1}^r{}^h(A_{\mathfrak m_i})$ des hensélisés des anneaux locaux $A_{\mathfrak m_i}$. C'est un $A$-module fidèlement plat et une $A$-algèbre semi-locale, dont les idéaux maximaux sont les $\mathfrak m_i.{}^h\!A$ en vertu de (18.6.6, (iii)); en outre, si $\mathfrak r=\bigcap_i\mathfrak m_i$ est le radical de $A$, il résulte de ce qui précède que $\mathfrak r.{}^h\!A$

\oldpage[IV-4]{139}

est le radical de ${}^h\!A$ et que l'application canonique $A/\mathfrak r\to{}^h\!A/\mathfrak r.{}^h\!A$ est bijective. Comme le séparé complété $\widehat A$ de $A$ pour la topologie $\mathfrak r$-préadique est le produit des séparés complétés $\widehat A_{\mathfrak m_i}$ (Bourbaki, \emph{Alg. comm.}, chap. III, \textsection~2, n\textsuperscript{o}~13, prop. 18) l'homomorphisme canonique $\widehat A\to({}^h\!A)\widehat{\phantom{x}}$ est bijectif par (18.6.6, (iv)), et il est clair par (18.6.6, (v)) que pour que ${}^h\!A$ soit noethérien, il faut et il suffit que $A$ le soit.

Pour obtenir l'analogue de la propriété universelle (18.6.6, (ii)), convenons, lorsque $A$ et $B$ sont deux anneaux semi-locaux de radicaux respectifs $\mathfrak r$, $\mathfrak s$, d'appeler \emph{homomorphisme semi-local} de $A$ dans $B$ tout homomorphisme $\varphi$ tel que $\varphi(\mathfrak r)\subset\mathfrak s$. Dans le cas particulier où $B$ est un \emph{produit} d'anneaux locaux $B_j$ ($1\leq j\leq n$), la donnée d'un tel homomorphisme revient à celle de ses projections $\varphi_j:A\to B_j$, qui sont des homomorphismes soumis à la seule condition que $\varphi_j(\mathfrak r)\subset\mathfrak n_j$, où $\mathfrak n_j$ est l'idéal maximal de $B_j$. D'ailleurs, si $\mathfrak m_i$ ($1\leq i\leq m$) sont les idéaux maximaux de $A$, l'idéal premier $\varphi_j^{-1}(\mathfrak n_j)$ doit contenir l'un des $\mathfrak m_i$ puisqu'il contient leur intersection $\mathfrak r$, donc il est égal à l'un des $\mathfrak m_i$ et $\varphi_j$ se factorise en $A\to A_{\mathfrak m_i}\xrightarrow{\psi_{ij}}B_j$, où $\psi_{ij}$ est un homomorphisme \emph{local}. Si $\Hom.\operatorname{sloc}(A,B)$ désigne l'ensemble des homomorphismes semi-locaux de $A$ dans $B$, on peut donc identifier cet ensemble à
\[
  \prod_j\left(\bigcup_i\Hom.\operatorname{loc}(A_{\mathfrak m_i},B_j)\right)
\]
dans le cas considéré; comme un anneau semi-local hensélien est produit direct d'anneaux locaux henséliens, on voit que la propriété universelle (18.6.6, (ii)) est encore valable pour un anneau semi-local $A$, lorsqu'on y remplace les homomorphismes locaux par les homomorphismes semi-locaux.
\end{env}

\begin{proposition}[18.6.8]
\label{IV.18.6.8}
Soient $A$ un anneau semi-local, $B$ une $A$-algèbre semi-locale et entière sur $A$. Alors $B\otimes_A({}^h\!A)$ est une $B$-algèbre semi-locale isomorphe à ${}^h\!B$.
\end{proposition}

Comme $B\otimes_A({}^h\!A)$ est entier au-dessus de ${}^h\!A$, chacun de ses idéaux maximaux est au-dessus d'un des idéaux maximaux de ${}^h\!A$, donc au-dessus d'un idéal maximal de $A$, et comme $B$ est entière sur $A$, les idéaux premiers de $B$ au-dessus d'un idéal maximal de $A$ sont des idéaux maximaux de $B$, si bien que finalement un idéal maximal de $B\otimes_A({}^h\!A)$ a pour projections dans $\Spec({}^h\!A)$ et $\Spec(B)$ des idéaux maximaux; compte tenu de (18.6.6, (iii)) et de (I, 3.4.9), on conclut que l'anneau $B\otimes_A({}^h\!A)$ est semi-local. En outre, pour tout anneau local hensélien $C$, $\Hom.\operatorname{sloc}(B\otimes_A({}^h\!A),C)$ est en correspondance biunivoque avec l'ensemble des couples $(\varphi,\psi)$ d'homomorphismes semi-locaux $\varphi:B\to C$, $\psi:{}^h\!A\to C$ tels que les composés
\[
  A\longrightarrow B\xrightarrow{\varphi}C,\qquad
  A\longrightarrow{}^h\!A\xrightarrow{\psi}C
\]
soient égaux. Mais en vertu de la correspondance biunivoque entre $\Hom.\operatorname{sloc}(A,C)$ et $\Hom.\operatorname{sloc}({}^h\!A,C)$ on voit que pour tout $\varphi\in\Hom.\operatorname{sloc}(B,C)$ il existe un $\psi$ et un seul ayant la propriété précédente, donc l'application
\[
  \Hom.\operatorname{sloc}(B\otimes_A({}^h\!A),C)\longrightarrow\Hom.\operatorname{sloc}(B,C)
\]
est bijective, ce qui prouve la proposition en vertu de l'unicité de la solution d'un problème universel.

\begin{theorem}[18.6.9]
\label{IV.18.6.9}
Soit $A$ un anneau semi-local.
\begin{enumerate}
  \item[(i)] Pour que ${}^h\!A$ soit réduit (resp. normal), il faut et il suffit que $A$ le soit.
  \item[(ii)] Supposons $A$ noethérien. Alors, pour tout idéal premier $\mathfrak p$ de $A$, l'anneau $({}^h\!A)_{\mathfrak p}/\mathfrak p.({}^h\!A)_{\mathfrak p}$ est composé direct d'un nombre fini d'extensions algébriques séparables de $k(\mathfrak p)$ (ce qui entraîne que les fibres du morphisme canonique $\Spec({}^h\!A)\to\Spec(A)$ sont \emph{géométriquement régulières}).
\end{enumerate}
\end{theorem}

\oldpage[IV-4]{140}

(i) Comme ${}^h\!A$ est un $A$-module fidèlement plat, il résulte de (2.1.13) que si ${}^h\!A$ est réduit (resp. normal), il en est de même de $A$. Pour prouver la réciproque, on peut se borner au cas où l'anneau $A$ est local, de sorte qu'avec les notations de (18.6.4), on a ${}^h\!A=\varinjlim A_\alpha$. Or, il résulte de la définition des $A_\alpha$ et de (17.5.7) que si $A$ est réduit (resp. intègre et intégralement clos), il en est de même des $A_\alpha$; en outre, les homomorphismes $A_\alpha\to A_\beta$ pour $\alpha\leq\beta$ sont injectifs en vertu de ($0_{\mathrm I}$, 6.5.1) puisque $A_\beta$ est un $A_\alpha$-module fidèlement plat; donc les morphismes $\Spec(A_\beta)\to\Spec(A_\alpha)$ sont dominants, et on conclut de (5.13.2) (resp. (5.13.4)) que ${}^h\!A$ est réduit (resp. intègre et intégralement clos).

(ii) On peut se borner au cas où $A$ est local. En vertu du fait que le foncteur $\varinjlim$ permute au produit tensoriel, la fibre du morphisme $\Spec({}^h\!A)\to\Spec(A)$ en un point $\mathfrak p$ est limite inductive des fibres des morphismes $\Spec(A_\alpha)\to\Spec(A)$ en ce point. Comme ${}^h\!A$ est noethérien, on voit, compte tenu de (17.6.2), que l'on est ramené à prouver le lemme suivant:

\begin{lemma}[18.6.9.1]
\label{IV.18.6.9.1}
Soient $(B_\alpha)$ un système inductif d'anneaux artiniens, $B=\varinjlim B_\alpha$ sa limite inductive. Si $B$ est noethérien, alors $B$ est artinien; si de plus, pour $\alpha\leq\beta$, les homomorphismes $\varphi_{\beta\alpha}:B_\alpha\to B_\beta$ sont injectifs, alors il existe $\lambda$ tel que, pour $\alpha\geq\lambda$ le nombre des composants locaux $B_\alpha^{(i)}$ de $B_\alpha$ soit constant, que $\varphi_{\beta\alpha}(B_\alpha^{(i)})\subset B_\beta^{(i)}$ pour tout $i$ et que les $B^{(i)}=\varinjlim B_\alpha^{(i)}$ soient les composants locaux de $B$.
\end{lemma}

Soient en effet $Y_\alpha=\Spec(B_\alpha)$, $Y=\Spec(B)$, qui, en tant qu'espace topologique, est égal à $\varprojlim Y_\alpha$ (8.2.10). Les espaces $Y_\alpha$ sont finis et discrets; en outre, si les $\varphi_{\beta\alpha}$ sont injectifs, les morphismes canoniques $Y_\beta\to Y_\alpha$ sont dominants, donc surjectifs. Le lemme résultera donc du lemme purement topologique suivant:

\begin{lemma}[18.6.9.2]
\label{IV.18.6.9.2}
Soit $(Y_\alpha,\psi_{\alpha\beta})$ un système projectif d'espaces topologiques finis et discrets. Si $Y=\varprojlim Y_\alpha$ est noethérien, $Y$ est fini et discret. Si de plus les $\psi_{\alpha\beta}$ sont surjectifs, il existe $\lambda$ tel que pour $\alpha\geq\lambda$, le nombre d'éléments de $Y_\alpha$ soit constant et égal au nombre d'éléments de $Y$.
\end{lemma}

En effet, $Y$ étant compact et noethérien, toute partie de $Y$ est compacte, donc fermée, ce qui implique que $Y$ est discret, donc fini puisqu'il est compact. Si de plus les $\psi_{\alpha\beta}$ sont surjectifs, il en est de même de $\psi_\alpha:Y\to Y_\alpha$, donc $\operatorname{Card}(Y)\geq\operatorname{Card}(Y_\alpha)$, et comme $\operatorname{Card}(Y_\alpha)$ est fonction croissante de $\alpha$, il existe $\lambda$ tel que pour $\alpha\geq\lambda$, $\operatorname{Card}(Y_\alpha)=\operatorname{Card}(Y_\lambda)$; les $\psi_{\alpha\beta}$ sont alors bijectifs pour $\alpha\geq\lambda$ et on a donc $\operatorname{Card}(Y)=\operatorname{Card}(Y_\lambda)$.

\begin{corollary}[18.6.10]
\label{IV.18.6.10}
Soit $A$ un anneau local noethérien. Pour que ${}^h\!A$ possède une des propriétés suivantes:
\begin{enumerate}
  \item[$a)$] être un anneau de Cohen-Macaulay (0, 16.5.3);
  \item[$b)$] vérifier la propriété $(S_n)$ (5.7.3);
  \item[$c)$] être régulier;
  \item[$d)$] vérifier la propriété $(R_n)$ (5.8.2);
\end{enumerate}
il faut et il suffit que $A$ possède cette même propriété.
\end{corollary}

\oldpage[IV-4]{141}

Cela résulte aussitôt de ce que les fibres du morphisme $\Spec({}^h\!A)\to\Spec(A)$ sont géométriquement régulières (18.6.9) et de (6.4.1), (6.5.1) et (6.5.3).

\begin{corollary}[18.6.11]
\label{IV.18.6.11}
Soit $A$ un anneau semi-local noethérien. Pour qu'un idéal premier $\mathfrak p$ de ${}^h\!A$ appartienne à $\operatorname{Ass}({}^h\!A)$, il faut et il suffit que $\mathfrak p\cap A$ appartienne à $\operatorname{Ass}(A)$.
\end{corollary}

Compte tenu de la description des fibres du morphisme $\Spec({}^h\!A)\to\Spec(A)$, cela résulte aussitôt de (3.3.1).

\begin{proposition}[18.6.12]
\label{IV.18.6.12}
Soit $A$ un anneau local; les propriétés suivantes sont équivalentes:
\begin{enumerate}
  \item[$a)$] $A$ est unibranche (0, 23.2.1) (resp. réduit et unibranché).
  \item[$b)$] Pour toute $A$-algèbre locale strictement essentiellement étale $A'$, $\Spec(A')$ est irréductible (resp. intègre).
  \item[$c)$] $\Spec({}^h\!A)$ est irréductible (resp. intègre).
\end{enumerate}
\end{proposition}

Notons d'abord que ${}^h(A_{\mathrm{red}})=({}^h\!A)_{\mathrm{red}}$, car ${}^h(A_{\mathrm{red}})=({}^h\!A)\otimes_A A_{\mathrm{red}}$ (18.6.8), et comme ${}^h(A_{\mathrm{red}})$ est réduit (18.6.9), il est égal à $({}^h\!A)_{\mathrm{red}}$. Ceci permet, dans la démonstration, de considérer seulement le cas où $A$ est réduit, donc aussi $A'$ (17.5.7) et ${}^h\!A$ (18.6.9).

Avec les notations de (18.6.4), la condition $b)$ signifie que tous les $A_\alpha$ sont intègres; on en conclut que ${}^h\!A=\varinjlim A_\alpha$ est intègre (5.13.3), donc $b)$ entraîne $c)$.

Il est clair que si ${}^h\!A$ est intègre, il en est de même des $A_\alpha\subset{}^h\!A$, donc $c)$ entraîne $b)$.

Montrons que $c)$ entraîne $a)$. Notons d'abord que $A\subset{}^h\!A$ est intègre; soient $K$ et $L$ les corps des fractions de $A$ et ${}^h\!A$ respectivement. Il suffit de voir que pour toute $A$-algèbre \emph{finie} $B\subset K$, $B$ est un anneau local. Or $B$ est un anneau semi-local, dont le hensélisé est donc $B\otimes_A({}^h\!A)$ (18.6.8). Par platitude, $B\otimes_A({}^h\!A)$ s'identifie à un sous-anneau de $K\otimes_A({}^h\!A)$ et ce dernier à un sous-anneau de $L$, donc ${}^h\!B$ est intègre. Mais comme ${}^h\!B$ est un anneau semi-local hensélien, il est composé direct d'anneaux locaux (18.5.9), donc ne peut être intègre que s'il est local, ce qui entraîne bien que $B$ est local.

Prouvons enfin que $a)$ implique $c)$. Soit $C$ la clôture intégrale de l'anneau intègre $A$. Comme $C$ est un anneau local par hypothèse, $C\otimes_A({}^h\!A)$ est le hensélisé de $C$ (18.6.8), donc est intègre et intégralement clos (18.6.9). Comme $A$ est un sous-anneau de $C$, ${}^h\!A$ s'identifie à un sous-anneau de ${}^h\!C$ par platitude, donc est lui aussi intègre, ce qui achève la démonstration.

\begin{corollary}[18.6.13]
\label{IV.18.6.13}
Soit $A$ un anneau local hensélien. Alors, pour que $A$ soit unibranche, il faut et il suffit que $A_{\mathrm{red}}$ soit intègre.
\end{corollary}

\begin{proposition}[18.6.14]
\label{IV.18.6.14}
\begin{enumerate}
  \item[(i)] Toute limite inductive filtrante d'anneaux locaux henséliens (où les homomorphismes de transition sont locaux) est un anneau local hensélien.
  \item[(ii)] Le foncteur $A\leadsto{}^h\!A$ dans la catégorie des anneaux locaux (où les morphismes sont les homomorphismes locaux) commute aux limites inductives filtrantes.
  \item[(iii)] Tout anneau local hensélien $A$ est limite inductive d'un système inductif filtrant d'anneaux locaux henséliens et noethériens $A_\lambda$ (les homomorphismes de transition $A_\lambda\to A_\mu$ étant locaux).
\end{enumerate}
\end{proposition}

(i) Soit $(A_\lambda)$ un système inductif filtrant d'anneaux locaux henséliens, où les homomorphismes de transition sont locaux. Montrons que $A'=\varinjlim A_\lambda$, qui est local ($0_{\mathrm{III}}$, 10.3.1.3) est hensélien. Or, soient $C$ une $A'$-algèbre étale, $\mathfrak n$ un idéal premier

\oldpage[IV-4]{142}

de $C$ au-dessus de l'idéal maximal $\mathfrak m'$ de $A'$, tel que $C_{\mathfrak n}$ soit strictement essentiellement étale sur $A'$ (18.6.1), autrement dit que $k(\mathfrak m')\to k(\mathfrak n)$ soit un isomorphisme. En vertu de (8.8.2, (ii)) et de (17.7.8), il existe un indice $\lambda$ et une $A_\lambda$-algèbre étale $C_\lambda$ tels que $C=C_\lambda\otimes_{A_\lambda}A'$ à isomorphisme près; en outre, si $\mathfrak n_\lambda$ est l'image réciproque de $\mathfrak n$ dans $C_\lambda$, et $\mathfrak m_\lambda$ l'idéal maximal de $A_\lambda$, on peut, en vertu de (8.8.2.4) et de la transitivité des fibres (I, 3.6.4), supposer que $(C_\lambda)_{\mathfrak n_\lambda}$ est strictement essentiellement étale sur $A_\lambda$; puisque $A_\lambda$ est hensélien, $(C_\lambda)_{\mathfrak n_\lambda}$ est nécessairement isomorphe à $A_\lambda$ (18.5.11, $b)$); il existe par suite un voisinage de $\mathfrak n_\lambda$ dans $\Spec(C_\lambda)$ isomorphe à $\Spec(A_\lambda)$ (I, 6.5.4); on en conclut qu'il y a un voisinage de $\mathfrak n$ dans $\Spec(C)$ isomorphe à $\Spec(A')$, donc que $C_{\mathfrak n}$ est isomorphe à $A'$, ce qui prouve que $A'$ est hensélien (18.6.2) et achève la démonstration.

(ii) Supposons que $B=\varinjlim B_\lambda$, où les $B_\lambda$ sont des anneaux locaux, les homomorphismes de transition $B_\lambda\to B_\mu$ ($\lambda\leq\mu$) étant locaux. Posons $A_\lambda={}^h\!B_\lambda$; les $A_\lambda$ sont des anneaux locaux henséliens (18.6.6, (vi)), et en vertu de (18.6.6, (ii)) les homomorphismes de transition $B_\lambda\to B_\mu$ déterminent de façon unique des homomorphismes locaux $A_\lambda\to A_\mu$, de sorte que $(A_\lambda)$ est un système inductif. Tout revient à voir que $A'=\varinjlim A_\lambda$ est canoniquement isomorphe à $A={}^h\!B$. En vertu de (18.6.6, (ii)) et de la définition des limites inductives, on a, pour tout anneau local hensélien $E$,
\[
  \Hom.\operatorname{loc}(A',E)
  =\varprojlim\Hom.\operatorname{loc}(A_\lambda,E)
  =\varprojlim\Hom.\operatorname{loc}(B_\lambda,E)
  =\Hom.\operatorname{loc}(B,E).
\]
Mais comme $A'$ est hensélien par (i), cela prouve notre assertion.

(iii) L'anneau $A$ est limite inductive filtrante de ses sous-anneaux locaux noethériens $B_\lambda$, les homomorphismes de transition étant locaux (5.13.3, (iii)). Comme les ${}^h\!B_\lambda$ sont des anneaux locaux henséliens et noethériens (18.6.6, (v)) et que $A={}^h\!A$ (18.6.6, (vi)), il suffit d'appliquer (ii) au système inductif $(B_\lambda)$.

\begin{corollary}[18.6.15]
\label{IV.18.6.15}
Soient $A$ un anneau local hensélien, $X$ un $A$-préschéma de présentation finie sur $A$. Il existe alors un anneau local noethérien et hensélien $A_0$, un homomorphisme local $A_0\to A$, un préschéma $X_0$ de type fini sur $A_0$ et un $A$-isomorphisme $X_0\otimes_{A_0}A\xrightarrow{\sim}X$.
\end{corollary}

Cela résulte de (18.6.14) et de (8.8.2, (ii)).

\subsection{Hensélisation et anneaux excellents.}
\EGAReaderTarget{ega4.c04.subsec.18-7}{subsection:IV.18.7}

\begin{env}[18.7.1]
\label{IV.18.7.1}
Nous désignerons dans ce numéro par $\mathbf P(Z,k)$ une propriété de la forme considérée dans (7.3.1), où nous supposons en outre que la propriété $\mathbf Q(A,k)$ satisfait à la condition suivante:

Pour toute extension algébrique séparable $k'$ de $k$, et toute $k'$-algèbre locale noethérienne $A$, la propriété $\mathbf Q(A,k)$ est \emph{équivalente} à $\mathbf Q(A,k')$.

Il est immédiat que toutes les propriétés $\mathbf P$ considérées dans (7.3.8) vérifient la condition précédente, compte tenu du fait que si $K$ est une extension finie de $k$, $k'\otimes_k K$ est composé direct d'extensions algébriques séparables de $K$.
\end{env}

\oldpage[IV-4]{143}

\begin{proposition}[18.7.2]
\label{IV.18.7.2}
Soit $A$ un anneau local noethérien. Pour que ${}^h\!A$ soit un $\mathbf P$-anneau (7.3.13), il faut et il suffit que $A$ le soit.
\end{proposition}

En effet, en vertu de (18.6.6, (iv)), le complété $\widehat A$ de $A$ est aussi le complété de ${}^h\!A$ et l'on a donc $A\subset{}^h\!A\subset\widehat A$. D'après (18.6.9), pour tout $x\in\Spec(A)$, la fibre en $x$ du morphisme $\Spec({}^h\!A)\to\Spec(A)$ est discrète et finie, et en chacun des points $y_i$ de cette fibre, $k(y_i)$ est une extension algébrique séparable de $k(x)$. La fibre formelle de $A$ au point $x$ est donc un préschéma \emph{somme} des fibres formelles de ${}^h\!A$ aux points $y_i$. La conclusion résulte alors de l'hypothèse faite sur $\mathbf P$ dans (18.7.1).

\begin{corollary}[18.7.3]
\label{IV.18.7.3}
Soit $A$ un anneau local noethérien. Pour que ${}^h\!A$ soit universellement japonais, il faut et il suffit que $A$ le soit.
\end{corollary}

En effet, pour un anneau local noethérien $B$, il revient au même de dire que $B$ est universellement japonais, ou que les fibres formelles de $B$ sont géométriquement réduites (7.6.4 et 7.7.1); la conclusion résulte donc de (18.7.2).

\begin{corollary}[18.7.4]
\label{IV.18.7.4}
Soit $A$ un anneau local noethérien. Pour que les fibres formelles de ${}^h\!A$ soient géométriquement régulières, il faut et il suffit que celles de $A$ le soient.
\end{corollary}

C'est un cas particulier de (18.7.2).

\begin{proposition}[18.7.5]
\label{IV.18.7.5}
Soit $A$ un anneau local noethérien. Si $A$ est universellement caténaire (5.6.2) (resp. formellement caténaire (7.1.9), resp. strictement formellement caténaire (7.2.6)), il en est de même de ${}^h\!A$.
\end{proposition}

Avec les notations de (18.6.4), si $A$ est universellement caténaire, il en est de même des $A_\alpha$, qui sont des $A$-algèbres essentiellement de type fini (5.6.3). Pour prouver que dans ce cas ${}^h\!A$ est universellement caténaire, il suffira donc d'établir le lemme suivant:

\begin{lemma}[18.7.5.1]
\label{IV.18.7.5.1}
Soit $(B_\alpha,\varphi_{\beta\alpha})$ un système inductif d'anneaux noethériens et supposons que, pour $\alpha\leq\beta$, $\varphi_{\beta\alpha}$ fasse de $B_\beta$ un $B_\alpha$-module fidèlement plat et que $B=\varinjlim B_\alpha$ soit noethérien. Alors, si les $B_\alpha$ sont caténaires (resp. universellement caténaires), il en est de même de $B$.
\end{lemma}

Dire que $B$ est universellement caténaire équivaut à dire que $B[T_1,\ldots,T_n]$ est caténaire pour tout $n$ (5.6.2), et il est immédiat que le système inductif $(B_\alpha[T_1,\ldots,T_n])$ vérifie les mêmes conditions que $(B_\alpha)$; il suffit donc de prouver que si les $B_\alpha$ sont caténaires, il en est de même de $B$. Si $\mathfrak p_0\supset\mathfrak p_1\supset\cdots\supset\mathfrak p_r$ est une chaîne saturée d'idéaux premiers de $B$, les images réciproques $\mathfrak p_{0\alpha}\supset\mathfrak p_{1\alpha}\supset\cdots\supset\mathfrak p_{r\alpha}$ de ces idéaux dans $B_\alpha$ forment une chaîne d'idéaux premiers, et il suffit de montrer que pour $\alpha$ assez grand, cette chaîne est saturée. Or, chaque $\mathfrak p_i$ est un idéal de type fini, donc il existe un $\lambda$ tel que, pour $\alpha\geq\lambda$, on ait $\mathfrak p_i=\mathfrak p_{i\alpha}.B$. On a donc
\[
1=\operatorname{codim}(V(\mathfrak p_i),V(\mathfrak p_{i+1}))
=\operatorname{codim}(V(\mathfrak p_{i\alpha}),V(\mathfrak p_{i+1,\alpha}))
\]
puisque $B$ est un $B_\alpha$-module plat ((6.1.4) et ($0_{\mathrm I}$, 6.2.3)), ce qui achève de prouver le lemme.

Supposons maintenant $A$ formellement caténaire. Soient $\mathfrak p$ un idéal premier de ${}^h\!A$, $\mathfrak q$ son image réciproque dans $A$, de sorte que $\mathfrak q.{}^h\!A\subset\mathfrak p$, et par suite ${}^h\!A/\mathfrak p$ est un anneau quotient de ${}^h\!A/\mathfrak q.{}^h\!A$. Il suffit donc de prouver que ${}^h\!A/\mathfrak q.{}^h\!A$ est formellement caténaire pour tout idéal premier $\mathfrak q$ de $A$ (7.1.9), et comme ${}^h\!A/\mathfrak q.{}^h\!A={}^h(A/\mathfrak q)$ par (18.6.8), il suffit (7.1.11) de voir que si $A$ est \emph{intègre et formellement caténaire}, ${}^h\!A$ est \emph{formellement équidimensionnel}. Mais comme le complété de ${}^h\!A$ est égal à $\widehat A$, cela résulte de l'hypothèse que $A$ est formellement caténaire.

\oldpage[IV-4]{144}

Supposons enfin $A$ strictement formellement caténaire; alors on vient de voir que ${}^h\!A$ est formellement caténaire, et il reste à prouver que les fibres du morphisme $\Spec(\widehat A)\to\Spec({}^h\!A)$ vérifient la propriété $(S_1)$ (7.2.5, $b)$). Mais cela résulte de l'hypothèse sur $A$ et de (18.7.2) appliqué au cas où $\mathbf P(Z,k)$ est la propriété $(S_1)$ pour $Z$.

\begin{corollary}[18.7.6]
\label{IV.18.7.6}
Si un anneau local noethérien $A$ est excellent (7.8.2), il en est de même de ${}^h\!A$.
\end{corollary}

\begin{remark}[18.7.7]
\label{IV.18.7.7}
Il peut se faire que ${}^h\!A$ soit un anneau local excellent sans que $A$ soit universellement caténaire (ni par suite excellent). Pour le voir, reprenons l'exemple de l'anneau $A$ de (5.6.11), avec les mêmes notations. L'anneau $E$ construit dans (5.6.11) est excellent lorsque $k_0$ est de caractéristique $0$ (7.8.3, (ii) et (iii)); considérons deux anneaux $E_1$, $E_2$ isomorphes à $E$, et « recollons » $\Spec(E_1)$ et $\Spec(E_2)$ de sorte que si $\mathfrak m_i$ et $\mathfrak m'_i$ sont les idéaux maximaux de $E_i$ ($i=1,2$) correspondant à $\mathfrak m$ et $\mathfrak m'$, le point $\mathfrak m_1$ soit « recollé » à $\mathfrak m'_2$ et le point $\mathfrak m_2$ à $\mathfrak m'_1$; de façon précise, si $\varepsilon_i$ et $\varepsilon'_i$ sont les homomorphismes canoniques de $E_i$ sur $k(\mathfrak m_i)$ et $k(\mathfrak m'_i)$ ($i=1,2$), le schéma obtenu est $\Spec(R)$, où $R$ est le sous-anneau de $E_1\times E_2$ formé des couples $(x_1,x_2)$ tels que $\varepsilon'_2(x_2)=\sigma(\varepsilon_1(x_1))$ et $\varepsilon'_1(x_1)=\sigma(\varepsilon_2(x_2))$. On vérifie aisément (par exemple à l'aide de (17.6.3)) que $R$ est une $C$-algèbre finie et étale, et qu'il y a deux idéaux maximaux $\mathfrak r_1$, $\mathfrak r_2$ de $R$ au-dessus de l'idéal maximal $\mathfrak n$ de $C$, les complétés de $R_{\mathfrak r_1}$ et $R_{\mathfrak r_2}$ étant canoniquement isomorphes à celui de $A=C_{\mathfrak n}$; $R_{\mathfrak r_1}$ est donc une $A$-algèbre strictement essentiellement étale, et par suite son hensélisé est égal à ${}^h\!A$. En outre, on a vu (7.8.4, (ii)) que les fibres formelles de $A$ sont géométriquement régulières, donc il en est de même de celles de ${}^h\!A$ (18.7.2). Enfin, l'anneau $R$ est universellement caténaire, car ses quotients par ses deux idéaux premiers minimaux sont isomorphes à $E$, d'où la conclusion par (5.6.3, (iii)). L'anneau $R_{\mathfrak r_1}$ est donc excellent, et il en est par suite de même de ${}^h\!A$ (18.7.6), alors que $A$ n'est pas universellement caténaire.
\end{remark}

\subsection{Anneaux strictement locaux et hensélisation stricte.}
\EGAReaderTarget{ega4.c04.subsec.18-8}{subsection:IV.18.8}

\begin{proposition}[18.8.1]
\label{IV.18.8.1}
Soit $A$ un anneau local. Les conditions suivantes sont équivalentes:
\begin{enumerate}
  \item[$a)$] $A$ est un anneau hensélien et son corps résiduel $k$ est séparablement clos.
  \item[$b)$] $A$ est un anneau hensélien et tout revêtement étale de $S=\Spec(A)$ est trivial.
  \item[$c)$] Pour tout morphisme étale $f:X\to S$ et tout point $x\in X$ tel que $f(x)$ soit le point fermé $s$ de $S$, il existe une $S$-section $u:S\to X$ de $X$ telle que $u(s)=x$.
\end{enumerate}
\end{proposition}

L'hypothèse $c)$ implique que pour tout morphisme étale $f:X\to S$, les corps résiduels aux points de $f^{-1}(s)$ sont tous isomorphes à $k$, et en outre que la condition $b)$ de (18.5.11) est vérifiée. Donc $A$ est hensélien, et le fait que tout revêtement étale de $S$ est trivial résulte de (18.2.10, (ii)); donc $c)$ implique $b)$. Comme les revêtements étales de $\Spec(k)$ sont triviaux si et seulement si $k$ est séparablement clos, $b)$ implique $a)$ en vertu de (18.5.11, $b)$). Enfin, il résulte aussi de (18.5.11, $b)$) que $a)$ implique $c)$.

\begin{definition}[18.8.2]
\EGAReaderTarget{ega4.c04.def.18-8-2}{IV.18.8.2}
On dit qu'un anneau local est un anneau \emph{strictement local} s'il vérifie les conditions équivalentes de (18.8.1). On appelle \emph{schéma strictement local} un schéma isomorphe au spectre d'un anneau strictement local.
\end{definition}

\begin{remark}[18.8.3]
\label{IV.18.8.3}
Les conditions de (18.8.1) sont encore équivalentes à la suivante:
\begin{enumerate}
  \item[$d)$] Pour tout morphisme localement de type fini $f:X\to S$ et tout point $x\in X$ tel que $f(x)$ soit le point fermé $s$ de $S$ et que $\sh{O}_{X,x}/\mathfrak m_s\sh{O}_{X,x}$ soit un corps, extension finie séparable de $k(s)=k$, il existe un voisinage ouvert $U$ de $x$ dans $X$ tel que $f|U$ soit une immersion fermée.
\end{enumerate}
\end{remark}

\oldpage[IV-4]{145}

(On notera que l'hypothèse de $d)$ est vérifiée si $f$ est \emph{formellement non ramifié} au point $x$ (17.4.1.2).)

Nous laissons au lecteur la démonstration, qui est substantiellement la même que celle de (18.5.18).

\begin{lemma}[18.8.4]
\label{IV.18.8.4}
Soient $A$, $A'$ deux anneaux locaux, $\varphi:A\to A'$ un homomorphisme local faisant de $A'$ une $A$-algèbre essentiellement étale (18.6.1). Alors, pour tout anneau strictement local $B$, tout homomorphisme local $\psi:A\to B$ et tout homomorphisme de $k(A)$-algèbres $\alpha:k(A')\to k(B)$, il existe un et un seul homomorphisme local $\psi':A'\to B$ tel que $\psi=\psi'\circ\varphi$ et que $\alpha$ soit égal à l'homomorphisme $\overline{\psi'}$ déduit de $\psi'$ par passage aux quotients.
\end{lemma}

Avec les notations de (18.6.2), l'homomorphisme $\alpha$ correspond à un $k(s)$-morphisme $\Spec(k(y))\to\Spec(k(x))$, donc à un point $x'$ bien déterminé dans $X'$ au-dessus de $y$. L'existence de la $Y$-section $f'$ de $X'$ telle que $f'(y)=x'$ résulte donc de (18.8.1, $c)$) puisque $f'$ est étale et $B$ strictement local, et son unicité résulte de ce que $f'$ est séparé et de (17.4.9).

\begin{lemma}[18.8.5]
\label{IV.18.8.5}
Soient $A$ un anneau local, $A_1$, $A_2$ deux $A$-algèbres locales essentiellement étales, $K$ un corps, extension de $k(A)$.
\begin{enumerate}
  \item[(i)] Pour tout $k(A)$-homomorphisme $\gamma:k(A_1)\to k(A_2)$, il existe au plus un $A$-homomorphisme (nécessairement local) $\psi:A_1\to A_2$ tel que $\gamma$ soit l'homomorphisme $\overline\psi$ déduit de $\psi$ par passage aux quotients.
  \item[(ii)] Pour tout couple de $A$-homomorphismes locaux $\beta_1:A_1\to K$, $\beta_2:A_2\to K$ il existe une $A$-algèbre locale essentiellement étale $A_3$, deux $A$-homomorphismes $\varphi_1:A_1\to A_3$, $\varphi_2:A_2\to A_3$ et un $A$-homomorphisme $\beta:A_3\to K$ tels que $\beta_1=\beta\circ\varphi_1$ et $\beta_2=\beta\circ\varphi_2$.
\end{enumerate}
\end{lemma}

Avec les notations de (18.6.3), les homomorphismes $\beta_1$, $\beta_2$ dans (ii) correspondent à deux $S$-morphismes $\Spec(K)\to X_1$, $\Spec(K)\to X_2$, d'images respectives $x_1$, $x_2$; on en déduit un $S$-morphisme bien déterminé $\Spec(K)\to X_3$ (I, 3.2.1) d'image $x_3$ au-dessus de $x_1$ et $x_2$; la $A$-algèbre $A_3=\sh{O}_{X_3,x_3}$ et l'homomorphisme $A_3\to K$ correspondant à ces données répondent aux conditions de (ii). D'autre part, la donnée de $\gamma$ dans (i) correspond à un $S$-morphisme $\Spec(k(x_2))\to\Spec(k(x_1))$ ou encore à un point $x'_3$ de $X'_3$ au-dessus de $x_2$; l'unicité de $\psi$ résulte de l'unicité d'une $X_2$-section de $X'_3$ passant par $x'_3$ (17.4.9).

\begin{env}[18.8.6]
\label{IV.18.8.6}
Soient $A$ un anneau local, $\mathfrak m$ son idéal maximal, $k$ son corps résiduel, $\Omega$ une clôture séparable de $k$. Considérons l'ensemble $\mathfrak E$ de $A$-algèbres essentiellement étales défini dans (18.6.1), et désignons par $L$ l'ensemble des $A$-homomorphismes locaux $A'\to\Omega$, où $A'$ parcourt $\mathfrak E$; on notera qu'un tel homomorphisme se factorise en
\[
  \lambda:A'\longrightarrow k(A')\xrightarrow{\omega_\lambda}\Omega,
\]
où $\omega_\lambda:k(A')\to\Omega$ est un $k$-homomorphisme (puisque $\mathfrak mA'$ est l'idéal maximal de $A'$); inversement la donnée d'un tel homomorphisme $\omega_\lambda$ détermine uniquement un homomorphisme $\lambda\in L$. Pour tout $\lambda\in L$, nous désignerons par $A_\lambda$ la $A$-algèbre essentiellement étale appartenant à $\mathfrak E$ sur laquelle $\lambda$ est défini. L'ensemble $L$ est \emph{préordonné} par la relation « il existe un $A$-homomorphisme $\varphi_{\mu\lambda}:A_\lambda\to A_\mu$ tel que $\lambda=\mu\circ\varphi_{\mu\lambda}$ », et il résulte de (18.8.5, (i)) que cet homomorphisme $\varphi_{\mu\lambda}$ est \emph{unique}. En outre, $L$ est \emph{filtrant croissant} pour la relation de préordre précédente (que l'on note encore $\lambda\leq\mu$), en vertu

\oldpage[IV-4]{146}

de (18.8.5, (ii)). Il est clair que $(A_\lambda,\varphi_{\mu\lambda})$ est un \emph{système inductif} de $A$-algèbres locales, les $\varphi_{\mu\lambda}$ étant des homomorphismes locaux; en outre, il résulte de (17.3.5) que pour $\lambda\leq\mu$, $A_\mu$ est une $A_\lambda$-algèbre essentiellement étale. Si $\overline\varphi_{\mu\lambda}:k(A_\lambda)\to k(A_\mu)$ est le $k$-homomorphisme déduit de $\varphi_{\mu\lambda}$ par passage aux quotients, $(k(A_\lambda),\overline\varphi_{\mu\lambda})$ est un \emph{système inductif} d'extensions algébriques séparables de $k$, et les $\omega_\lambda:k(A_\lambda)\to\Omega$ forment un système inductif de $k$-homomorphismes. En outre, la limite inductive
\begin{equation}
  \omega:\varinjlim k(A_\lambda)\longrightarrow\Omega
  \label{IV.18.8.6.1}
  \tag{18.8.6.1}
\end{equation}
est un $k$-\emph{isomorphisme}. Il suffit en effet de prouver que pour toute extension finie séparable $k'$ de $k$, il existe une $A$-algèbre essentiellement étale $A'$ telle que $k'$ soit le corps résiduel de $A'$ et que l'homomorphisme $k\to k'$ se déduise de $A\to A'$ par passage aux quotients. Mais cela résulte de (18.1.1) appliqué aux morphismes étales, compte tenu de ce que $\Spec(A)$ est le seul voisinage de son point fermé.

Notons maintenant que si $\mathfrak m_\lambda$ est l'idéal maximal de $A_\lambda$, il résulte de (17.6.1) que pour $\lambda\leq\mu$, on a $\mathfrak m_\mu=\mathfrak m_\lambda A_\mu$ et que $A_\mu$ est un $A_\lambda$-module plat. Donc il résulte de ($0_{\mathrm{III}}$, 10.3.1.3) que l'anneau $\varinjlim A_\lambda$ est \emph{local}, que l'homomorphisme canonique $w:A\to\varinjlim A_\lambda$ est local, et que l'on a un $k$-isomorphisme canonique
\[
  \varinjlim k(A_\lambda)\xrightarrow{\sim}k(\varinjlim A_\lambda);
\]
identifiant ces deux corps, on en déduit un $k$-\emph{isomorphisme canonique}
\begin{equation}
  \omega^{-1}:\Omega\xrightarrow{\sim}k(\varinjlim A_\lambda).
  \label{IV.18.8.6.2}
  \tag{18.8.6.2}
\end{equation}
\end{env}

\begin{definition}[18.8.7]
\EGAReaderTarget{ega4.c04.def.18-8-7}{IV.18.8.7}
Soient $A$ un anneau local, $k=k(A)$ son corps résiduel, $i:k\to\Omega$ un homomorphisme de $k$ dans une clôture séparable $\Omega$ de $k$; on appelle \emph{hensélisé strict} de $A$ relativement à $i$, et on note ${}^{hs}\!A_{(i)}$ (ou ${}^{hs}\!A$ lorsque cela ne prête pas à confusion) la limite inductive du système inductif $(A_\lambda,\varphi_{\mu\lambda})$ défini dans (18.8.6).
\end{definition}

Comme on a un $k$-isomorphisme canonique $\omega^{-1}:\Omega\xrightarrow{\sim}k({}^{hs}\!A_{(i)})$, $\Omega$ est muni canoniquement d'une structure de ${}^{hs}\!A_{(i)}$-algèbre par l'homomorphisme local
\[
  {}^{hs}\!A_{(i)}\longrightarrow k({}^{hs}\!A_{(i)})\xrightarrow{\omega}\Omega.
\]

Comme dans (18.6.5) on voit que la définition donnée dans (18.8.7) ne dépend qu'en apparence du choix de l'ensemble $\mathfrak E$; nous allons voir encore ci-dessous (18.8.8, (i) et (ii)) que ${}^{hs}\!A_{(i)}$ est un objet représentant un foncteur entièrement défini par la donnée de $A$ et de $i$ ($0_{\mathrm{III}}$, 8.1.8), et est donc défini comme tel à isomorphisme unique près.

\begin{proposition}[18.8.8]
\label{IV.18.8.8}
Soient $A$ un anneau local, $k$ son corps résiduel, $i:k\to\Omega$ un homomorphisme de $k$ dans une clôture séparable de $k$, ${}^{hs}\!A_{(i)}$ le hensélisé strict de $A$ correspondant à $i$.
\begin{enumerate}
  \item[(i)] ${}^{hs}\!A_{(i)}$ est un anneau strictement local (18.8.2) et l'homomorphisme structural $w:A\to{}^{hs}\!A_{(i)}$ est local.
  \item[(ii)] Pour tout homomorphisme local $u:A\to B$, où $B$ est un anneau strictement local, et tout $k$-homomorphisme $\psi:\Omega\to k(B)$, il existe un et un seul $A$-homomorphisme $v:{}^{hs}\!A_{(i)}\to B$ tel que $\psi$ se factorise en $\Omega\xrightarrow{\omega^{-1}}k({}^{hs}\!A_{(i)})\xrightarrow{\overline v}k(B)$, où $\overline v$ est déduit de $v$ par passage aux quotients.

\oldpage[IV-4]{147}

  \item[(iii)] ${}^{hs}\!A_{(i)}$ est un $A$-module fidèlement plat, et si $\mathfrak m$ est l'idéal maximal de $A$, $\mathfrak m.{}^{hs}\!A_{(i)}$ est l'idéal maximal de ${}^{hs}\!A_{(i)}$, et $k({}^{hs}\!A_{(i)})$ est une clôture séparable de $k$.
  \item[(iv)] Pour que ${}^{hs}\!A_{(i)}$ soit noethérien, il faut et il suffit que $A$ le soit.
  \item[(v)] Si $A$ est strictement local (donc $\Omega=k$), on a ${}^{hs}\!A\xrightarrow{\sim}A$.
  \item[(vi)] Si $i':k\to\Omega'$ est un second homomorphisme de $k$ dans une clôture séparable de $k$, pour tout $k$-isomorphisme $\sigma:\Omega\xrightarrow{\sim}\Omega'$, le $A$-homomorphisme correspondant (par (ii)) $v_\sigma:{}^{hs}\!A_{(i)}\to{}^{hs}\!A_{(i')}$ est un isomorphisme.
\end{enumerate}
\end{proposition}

On a déjà vu ci-dessus que ${}^{hs}\!A_{(i)}$ est un anneau local et que $w$ est un homomorphisme local; le fait que ${}^{hs}\!A_{(i)}$ soit hensélien (et par suite strictement local) se démontre comme dans (18.6.6). Les assertions de (iii) découlent encore de ($0_{\mathrm{III}}$, 10.3.1.3), ainsi que la suffisance de la condition (iv); la nécessité de cette condition résulte de ce que ${}^{hs}\!A_{(i)}$ est un $A$-module fidèlement plat ($0_{\mathrm I}$, 6.5.2).

Pour prouver (ii), remarquons, avec les notations de (18.8.6), que pour tout $\lambda\in L$, il existe un et un seul $A$-homomorphisme local $v_\lambda:A_\lambda\to B$, tel que le composé $k(A_\lambda)\xrightarrow{\omega_\lambda}\Omega\xrightarrow{\psi}k(B)$ soit déduit de $v_\lambda$ par passage aux quotients, en vertu de (18.8.4) et de l'hypothèse que $B$ est strictement local. L'unicité de $v_\lambda$ entraîne en outre que les $v_\lambda$ forment un système inductif, d'où l'existence de l'homomorphisme $v$ répondant aux conditions de (ii); son unicité résulte de (17.4.9). L'unicité de $v$ entraîne aussitôt l'assertion (vi), en considérant les composés $v_\sigma\circ v_{\sigma^{-1}}$ et $v_{\sigma^{-1}}\circ v_\sigma$. Enfin, si $A$ est strictement local, et $A'=B_{\mathfrak n}$ une $A$-algèbre locale essentiellement étale, où $B$ est une $A$-algèbre étale, il résulte de (18.8.1) qu'il existe une $\Spec(A)$-section de $\Spec(B)$ prenant la valeur $\mathfrak n$ au point fermé de $\Spec(A)$, donc on déduit de (17.4.1) que l'homomorphisme $A\to A'$ est bijectif, ce qui prouve (v).

\begin{env}[18.8.8.1]
\label{IV.18.8.8.1}
Soit $\mathcal C$ la catégorie des anneaux strictement locaux, avec pour morphismes les homomorphismes locaux; pour tout $B\in\mathcal C$, désignons par $F(B)$ l'ensemble des couples $(u,\psi)$ formés d'un homomorphisme local $u:A\to B$ et d'un $k$-isomorphisme $\psi:\Omega\to k(B)$ tels que le composé $k(A)=k\xrightarrow{i}\Omega\xrightarrow{\psi}k(B)$ soit déduit de $u$ par passage aux quotients. On peut dire que l'objet ${}^{hs}\!A_{(i)}$ \emph{représente} le foncteur covariant $F:\mathcal C\to\mathbf{Ens}$ ($0_{\mathrm{III}}$, 8.1.11).

On remarquera qu'en vertu de (18.8.8, (vi)) les hensélisés stricts ${}^{hs}\!A_{(i)}$ sont tous $A$-isomorphes (pour les divers $k$-homomorphismes $i$ de $k$ dans des clôtures séparables de $k$), mais pour $i$ et $i'$ donnés, il y a en général une infinité de $A$-isomorphismes ${}^{hs}\!A_{(i)}\xrightarrow{\sim}{}^{hs}\!A_{(i')}$. De façon précise, le groupe des $A$-automorphismes de ${}^{hs}\!A_{(i)}$ est isomorphe au groupe de Galois de $\Omega$ sur $k$.
\end{env}

\begin{env}[18.8.9]
\EGAReaderTarget{ega4.c04.statement.18-8-9}{IV.18.8.9}
Soient maintenant $A$ un anneau \emph{semi-local}, $\mathfrak m_j$ ($1\leq j\leq r$) ses idéaux maximaux, et pour tout indice $j$, considérons un homomorphisme $i_j:k(A_{\mathfrak m_j})\to\Omega_j$ de $k(A_{\mathfrak m_j})$ dans une clôture séparable de ce corps. On appelle \emph{hensélisé strict} de $A$ relatif aux $i_j$ et l'on note ${}^{hs}\!A$ (lorsqu'il n'en résulte pas de confusion) le produit $\prod_{j=1}^r({}^{hs}(A_{\mathfrak m_j}))$ des hensélisés stricts des anneaux locaux $A_{\mathfrak m_j}$ relativement aux homomorphismes $i_j$.

\oldpage[IV-4]{148}

C'est donc, en vertu de (18.8.8), un $A$-module fidèlement plat et une $A$-algèbre semi-locale dont les idéaux maximaux sont les $\mathfrak m_j.{}^{hs}\!A$ et le radical $\mathfrak r.{}^{hs}\!A$ (si $\mathfrak r$ désigne le radical de $A$). Enfin, la propriété universelle (18.8.8, (ii)) subsiste en remplaçant « homomorphisme local » par « homomorphisme semi-local » (18.6.7), et en remplaçant $\Omega$ par $\prod_j\Omega_j$.
\end{env}

\begin{proposition}[18.8.10]
\label{IV.18.8.10}
Soient $A$ un anneau semi-local, $B$ une $A$-algèbre finie. Soient $\mathfrak n_j$ ($1\leq j\leq r$) les idéaux maximaux de $B$. Alors il y a, pour chaque $j$, un idéal maximal $\mathfrak q_j$ de $B\otimes_A({}^{hs}\!A)$ au-dessus de $\mathfrak n_j$, tel que ${}^{hs}\!B$ soit isomorphe au composé direct des anneaux locaux $(B\otimes_A({}^{hs}\!A))_{\mathfrak q_j}$ ($1\leq j\leq r$).
\end{proposition}

On peut évidemment se borner au cas où $A$ est local, en remplaçant $A$ par $A_{\mathfrak m_j}$, où $\mathfrak m_j$ est l'idéal maximal de $A$ image réciproque de $\mathfrak n_j$. Posons d'autre part $C=B_{\mathfrak n_j}$; en vertu de la définition de ${}^{hs}\!B$ (18.8.1), tout revient à voir que ${}^{hs}\!C$ est $C$-isomorphe à $(C\otimes_A({}^{hs}\!A))_{\mathfrak q}$, où $\mathfrak q$ est un des idéaux maximaux de $C\otimes_A({}^{hs}\!A)$. Soient $k$ et $K$ les corps résiduels de $A$ et $C$, $\Omega$ celui de ${}^{hs}\!A$, qui est par définition une clôture séparable de $k$; comme $K$ est une extension finie de $k$, on peut supposer $K$ et $\Omega$ contenus dans une même clôture algébrique de $k$; il est alors immédiat que $K\otimes_k\Omega$ est composé direct d'un nombre fini de corps, tous isomorphes à la clôture séparable $K\Omega$ de $K$. D'autre part, l'anneau ${}^{hs}\!A$ étant hensélien, $C\otimes_A({}^{hs}\!A)$ est composé direct des anneaux locaux $(C\otimes_A({}^{hs}\!A))_{\mathfrak q_i}$, où $\mathfrak q_i$ ($1\leq i\leq s$) parcourt les idéaux maximaux de $C\otimes_A({}^{hs}\!A)$, et ces anneaux locaux sont henséliens (18.5.10) et ont $K\Omega$ pour corps résiduel, donc sont strictement locaux (18.8.1). Posons pour abréger $C'=C\otimes_A({}^{hs}\!A)$ et $C'_i=C'_{\mathfrak q_i}$, qui est donc un anneau quotient de $C'$; soient $f:C\to C'$, $g:{}^{hs}\!A\to C'$, $\varphi_i:C'\to C'_i$ les applications canoniques.

\[
\xymatrix{
  &&& {}^{hs}\!C \ar@<0.8ex>[dl]_{w'} \\
  && C'_i \ar@<0.8ex>[ur]_{w} \\
  C \ar[uurrr]^{v} \ar[r]_{f} & C' \ar[ur] \\
  A \ar[u] \ar[r] & {}^{hs}\!A \ar[u]^{g} \ar[uuurr]_{u}
}
\]

Soit $v:C\to{}^{hs}\!C$ l'homomorphisme canonique; il existe un homomorphisme local et un seul $u:{}^{hs}\!A\to{}^{hs}\!C$ qui, par passage aux quotients, donne l'injection canonique $\Omega\to K\Omega$ et tel que le composé $A\to C\xrightarrow{v}{}^{hs}\!C$ soit égal au composé $A\to{}^{hs}\!A\xrightarrow{u}{}^{hs}\!C$ (18.8.8); donc il existe un unique homomorphisme local $w_0:C'\to{}^{hs}\!C$ tel que $w_0\circ g=u$, $w_0\circ f=v$, et puisque ${}^{hs}\!C$ est un anneau local, cet homomorphisme se factorise en $C'\xrightarrow{\varphi_i}C'_i\xrightarrow{w}{}^{hs}\!C$ pour un indice $i$ bien déterminé, $w$ étant un homomorphisme local. D'autre part, l'anneau $C'_i$ étant strictement local, il existe un homomorphisme local $w':{}^{hs}\!C\to C'_i$ qui, par passage aux quotients, donne l'automorphisme identique de $K\Omega$ et qui est tel que $w'\circ v=\varphi_i\circ f$ (18.8.8). On déduit d'abord de là que $w\circ w'$ est un endomorphisme de ${}^{hs}\!C$ qui, par passage aux quotients, donne l'automorphisme identique de $K\Omega$, donc est

\oldpage[IV-4]{149}

l'identité (18.8.8). D'autre part, on a $w'\circ w\circ\varphi_i\circ f=w'\circ v=\varphi_i\circ f$, et $w'\circ w\circ\varphi_i\circ g=w'\circ u=\varphi_i\circ g$; d'où $w'\circ w\circ\varphi_i=\varphi_i$, autrement dit $w'\circ w$ est l'automorphisme identique de $C'_i$. C.Q.F.D.

\begin{remark}[18.8.11]
\label{IV.18.8.11}
La démonstration de (18.8.10) utilise le fait que $B$ est entière sur $A$, mais n'utilise le fait que $B$ est une $A$-algèbre \emph{finie}, que pour établir que $K\otimes_k\Omega$ est composé direct d'un nombre fini de corps. Or, cette dernière propriété est encore vérifiée lorsque l'on suppose que $B$ est une $A$-algèbre \emph{entière semi-locale} dont les idéaux maximaux $\mathfrak n_j$ ont des corps résiduels de degrés séparables $[k(\mathfrak n_j):k]_s$ finis.
\end{remark}

\begin{proposition}[18.8.12]
\label{IV.18.8.12}
Soit $A$ un anneau semi-local.
\begin{enumerate}
  \item[(i)] Pour que ${}^{hs}\!A$ soit réduit (resp. normal), il faut et il suffit que $A$ le soit.
  \item[(ii)] Supposons $A$ noethérien. Alors, pour tout idéal premier $\mathfrak p$ de $A$, l'anneau $({}^{hs}\!A)_{\mathfrak p}/\mathfrak p.({}^{hs}\!A)_{\mathfrak p}$ est composé direct d'un nombre fini d'extensions algébriques séparables de $k(\mathfrak p)$ (ce qui entraîne que les fibres du morphisme canonique $\Spec({}^{hs}\!A)\to\Spec(A)$ sont \emph{géométriquement régulières} et sont des espaces \emph{discrets}).
\end{enumerate}
\end{proposition}

\begin{corollary}[18.8.13]
\label{IV.18.8.13}
Soit $A$ un anneau local noethérien. Pour que ${}^{hs}\!A$ possède une des propriétés suivantes:
\begin{enumerate}
  \item[$a)$] être un anneau de Cohen-Macaulay (0, 16.5.3);
  \item[$b)$] vérifier la condition $(S_n)$ (5.7.3);
  \item[$c)$] être régulier;
  \item[$d)$] vérifier la condition $(R_n)$ (5.8.2);
\end{enumerate}
il faut et il suffit que $A$ possède cette même propriété.
\end{corollary}

\begin{corollary}[18.8.14]
\label{IV.18.8.14}
Soit $A$ un anneau semi-local noethérien. Pour qu'un idéal premier $\mathfrak p$ de ${}^{hs}\!A$ appartienne à $\operatorname{Ass}({}^{hs}\!A)$, il faut et il suffit que $\mathfrak p\cap A$ appartienne à $\operatorname{Ass}(A)$.
\end{corollary}

Les démonstrations sont les mêmes que celles de (18.6.9), (18.6.10) et (18.6.11) respectivement.

\begin{proposition}[18.8.15]
\label{IV.18.8.15}
Soit $A$ un anneau local; les propriétés suivantes sont équivalentes:
\begin{enumerate}
  \item[$a)$] $A$ est géométriquement unibranche (resp. réduit et géométriquement unibranche).
  \item[$b)$] Pour toute $A$-algèbre locale essentiellement étale $A'$, $\Spec(A')$ est irréductible (resp. intègre).
  \item[$c)$] $\Spec({}^{hs}\!A)$ est irréductible (resp. intègre).
\end{enumerate}
\end{proposition}

On se ramène comme dans (18.6.12) au cas où $A$ est réduit, en utilisant la relation ${}^{hs}(A_{\mathrm{red}})=({}^{hs}\!A)\otimes_A A_{\mathrm{red}}$ (18.8.11), et l'équivalence de $b)$ et $c)$ se prouve comme dans (18.6.12); il en est de même du fait que $a)$ entraîne $c)$, en tenant compte de (18.8.11) et de la définition d'un anneau local intègre géométriquement unibranche. Enfin, pour démontrer que $c)$ entraîne $a)$, gardons les mêmes notations que dans (18.6.12); il s'agit encore de voir que ${}^{hs}\!B$ est intègre. Mais, par (18.8.10), ${}^{hs}\!B$ est un anneau localisé de l'anneau semi-local $B\otimes_A({}^{hs}\!A)$, qui s'identifie encore par platitude à un sous-anneau du corps $L$, donc est intègre; il en est par suite de même de ${}^{hs}\!B$, ce qui achève la démonstration.

\begin{corollary}[18.8.16]
\label{IV.18.8.16}
Soit $A$ un anneau local hensélien (resp. strictement local). Pour que $A$ soit unibranche (resp. géométriquement unibranche), il faut et il suffit que $\Spec(A)$ soit irréductible.
\end{corollary}

\oldpage[IV-4]{150}

\begin{proposition}[18.8.17]
\label{IV.18.8.17}
Soit $A$ un anneau local noethérien. Si $A$ est universellement caténaire, il en est de même de ${}^{hs}\!A$.
\end{proposition}

La démonstration est la même que celle de (18.7.5).

\begin{proposition}[18.8.18]
\label{IV.18.8.18}
\begin{enumerate}
  \item[(i)] Toute limite inductive filtrante d'anneaux strictement locaux (où les homomorphismes de transition sont locaux) est un anneau strictement local.
  \item[(ii)] Le foncteur $A\leadsto{}^{hs}\!A$ dans la catégorie des anneaux locaux (où les morphismes sont les homomorphismes locaux) commute aux limites inductives filtrantes.
  \item[(iii)] Tout anneau strictement local $A$ est limite inductive d'un système inductif filtrant d'anneaux strictement locaux noethériens (les homomorphismes de transition étant locaux).
\end{enumerate}
\end{proposition}

Les démonstrations sont calquées sur celles de (18.6.14).

\subsection{Fibres formelles des anneaux noethériens henséliens.}
\EGAReaderTarget{ega4.c04.subsec.18-9}{subsection:IV.18.9}

\begin{theorem}[18.9.1]
\label{IV.18.9.1}
Soit $A$ un anneau local noethérien hensélien, dont les fibres formelles sont géométriquement normales. Alors les fibres formelles de $A$ sont géométriquement intègres (donc géométriquement connexes).
\end{theorem}

Les deux assertions de l'énoncé sont en fait équivalentes, un préschéma localement noethérien connexe et normal étant intègre. Comme toute $A$-algèbre intègre finie est un anneau local hensélien (18.5.9 et 18.5.10), il résulte de (7.3.16.2) que l'on est ramené à prouver que si l'on suppose de plus $A$ \emph{intègre}, alors la fibre du morphisme $\Spec(\widehat A)\to\Spec(A)$ au point générique de $\Spec(A)$ est \emph{intègre}, ce qui résulte du

\begin{corollary}[18.9.2]
\label{IV.18.9.2}
Soit $A$ un anneau local noethérien hensélien intègre, dont les fibres formelles sont géométriquement normales. Alors $\widehat A$ est intègre.
\end{corollary}

En effet, il résulte de (18.8.16) que $A$ est unibranche, donc le théorème résulte de l'hypothèse faite sur les fibres formelles de $A$ et de (7.6.3).

\begin{corollary}[18.9.3]
\label{IV.18.9.3}
Sous les hypothèses de (18.9.2), le corps des fractions $L$ de $\widehat A$ est une extension séparable du corps des fractions $K$ de $A$, et $K$ est algébriquement fermé dans $L$.
\end{corollary}

Cela résulte de ce que la fibre du morphisme $\Spec(\widehat A)\to\Spec(A)$ au point générique de $\Spec(A)$ est géométriquement intègre et de (4.3.2) et (4.3.5).

\begin{corollary}[18.9.4]
\label{IV.18.9.4}
Soit $A$ un anneau local noethérien, hensélien et dont les fibres formelles sont géométriquement normales (par exemple un anneau local noethérien hensélien et excellent), et posons $A'=\widehat A$. Soient $X$ un $A$-préschéma, $X'=X\otimes_A A'$, $g:X'\to X$ la projection canonique. Alors l'application $U\mapsto g^{-1}(U)$ est une bijection de l'ensemble des parties à la fois ouvertes et fermées de $X$ sur l'ensemble des parties à la fois ouvertes et fermées de $X'$.
\end{corollary}

Comme $g$ est un morphisme fidèlement plat et quasi-compact, on sait (2.3.12) que la topologie de $X$ est \emph{quotient} de celle de $X'$ par la relation d'équivalence définie par $g$. Tout revient donc à voir qu'une partie ouverte et fermée $U'$ de $X'$ est \emph{saturée} pour cette relation. Or, comme le morphisme $\Spec(A')\to\Spec(A)$ a ses fibres géométriquement connexes (18.9.1), les fibres de $g$ sont connexes, d'où aussitôt notre assertion.

En particulier, si $X$ est localement connexe (ce qui sera le cas si $X$ est localement noethérien), alors, pour toute composante connexe $U$ de $X$ (qui est ouverte et fermée

\oldpage[IV-4]{151}

dans $X$), $g^{-1}(U)$ est connexe (Bourbaki, \emph{Top. gén.}, chap.~I\textsuperscript{er}, 3\textsuperscript{e} éd., \textsection~11, n\textsuperscript{o}~3, prop.~7). Si l'on désigne par $\pi_0(X)$ l'ensemble des composantes connexes de $X$, l'application $\pi_0(X')\to\pi_0(X)$ déduite canoniquement de $g$ est donc \emph{bijective}.

\begin{corollary}[18.9.5]
\label{IV.18.9.5}
Sous les hypothèses de (18.9.4), le foncteur
\begin{equation}
  Z\longmapsto Z\otimes_A A'
  \tag{18.9.5.1}
  \label{IV.18.9.5.1}
\end{equation}
de la catégorie des préschémas étales sur $X$, dans celle des préschémas étales sur $X'$, est pleinement fidèle.
\end{corollary}

Soient $Z_1$, $Z_2$ deux préschémas étales sur $X$, et posons $Z'_i=Z_i\otimes_A A'$ ($i=1,2$); il s'agit de prouver que tout $X'$-morphisme $Z'_1\to Z'_2$ provient, par changement de base, d'un $X$-morphisme unique $Z_1\to Z_2$. Supposons d'abord que $Z_2$ soit \emph{séparé} sur $X$; alors, comme $\Hom_X(Z_1,Z_2)$ s'identifie à $\Gamma((Z_1\times_X Z_2)/Z_1)$, et que $Z_1\times_X Z_2$ est étale et séparé sur $Z_1$, $\Hom_X(Z_1,Z_2)$ s'identifie fonctoriellement, en vertu de (17.9.3), à l'ensemble des parties ouvertes et fermées $U$ de $Z_1\times_X Z_2$ telles que la restriction à $U$ de la projection $p:Z_1\times_X Z_2\to Z_1$ soit un morphisme surjectif et radiciel. L'assertion résulte donc dans le cas considéré de (18.9.4) et de ce que $Z'_1\times_{X'}Z'_2=(Z_1\times_X Z_2)\otimes_A A'$.

Passons au cas général: il suffira de prouver que le graphe $\Gamma'$ d'un $X'$-morphisme $Z'_1\to Z'_2$, qui est ouvert (17.9.3) dans $Z'_1\times_{X'}Z'_2$, est de la forme $\Gamma\otimes_A A'$, où $\Gamma$ est induit sur une partie \emph{ouverte} de $Z=Z_1\times_X Z_2$. En effet, la restriction à $\Gamma$ de la projection $p:Z\to Z_1$ est alors un \emph{isomorphisme}, car par changement de base $A\to A'$ cette restriction devient l'isomorphisme restriction à $\Gamma'$ de la projection $p':Z'\to Z'_1$ (avec $Z'=Z\otimes_A A'$), et il suffit d'appliquer (2.7.1, (viii)); la conclusion résulte alors de la caractérisation des graphes de morphismes (I, 5.3).

Si $q:Z'\to Z$ est la projection canonique, pour prouver que $\Gamma'$ est de la forme $q^{-1}(\Gamma)$, où $\Gamma$ est ouvert dans $Z$, il suffit, puisque $q$ est un morphisme fidèlement plat et quasi-compact, de prouver qu'il existe un ensemble $U\subset Z$ tel que $\Gamma'=q^{-1}(U)$ (2.3.12). Posons $Z''=Z'\times_Z Z'=Z\times_X X''$, où $X''=X'\times_X X'$, et soient $q_1:Z''\to Z'$, $q_2:Z''\to Z'$ les projections canoniques. Appliquant (4.5.19.1), il suffit de montrer que $q_1^{-1}(\Gamma')=q_2^{-1}(\Gamma')$. Mais c'est une propriété qui est vraie si et seulement si elle l'est après chaque changement de base $\Spec(k(x))\to X$, où $x$ parcourt $X$. Autrement dit, il suffit de prouver le corollaire lorsque $X$ est le \emph{spectre d'un corps}; mais comme tout $X$-préschéma étale est alors automatiquement \emph{séparé} sur $X$ (17.6.2, $c'$)), on est ramené au cas envisagé au début de la démonstration.

\begin{remarks}[18.9.6]
\label{IV.18.9.6}
\begin{enumerate}
  \item[(i)] Il est possible que, si le corps résiduel de $A$ est de caractéristique 0, le foncteur (18.9.5.1) induise même une \emph{équivalence} de la catégorie des revêtements étales de $X$ et de la catégorie des revêtements étales de $X'$. On peut en tout cas le démontrer lorsque $A$ est \emph{excellent}, en utilisant la résolution des singularités de Hironaka (M.~Artin). Il est plausible que l'énoncé analogue soit encore vrai sans restriction sur la caractéristique du corps résiduel, à condition de se borner aux revêtements « galoisiens principaux » dont le groupe est d'ordre premier à la caractéristique résiduelle (cf. [41]).

\oldpage[IV-4]{152}

  \item[(ii)] On ignore si les fibres formelles de $A$ sont géométriquement connexes (voire géométriquement irréductibles) lorsque $A$ est un anneau local hensélien quelconque. Il suffirait que, pour tout anneau local noethérien, hensélien et \emph{intègre} $A$, $\Spec(\widehat A)$ soit irréductible. On ignore s'il en est toujours ainsi.

  \item[(iii)] On ignore si lorsque $A$ est un anneau local excellent, son hensélisé strict ${}^{hs}\!A$ est excellent. On peut voir que, pour élucider cette question, on peut se ramener au cas où $A=k[[T_1,\ldots,T_n]]$, $k$ étant un corps de caractéristique $p>0$. Compte tenu de (18.8.17), la question est de savoir si les fibres formelles de ${}^{hs}\!A$ sont géométriquement régulières. La réponse est affirmative lorsque $n=1$, mais n'est pas connue pour $n=2$. On peut montrer que la réponse est affirmative chaque fois que, pour tout schéma $Y_1$ fini sur $Y=\Spec(\widehat A)$, on peut résoudre les singularités de $Y_1$ (7.9.1).
\end{enumerate}
\end{remarks}

L'assertion de connexité faite dans (18.9.1) se généralise de la façon suivante:

\begin{theorem}[18.9.7]
\label{IV.18.9.7}
Soient $A$, $B$ deux anneaux locaux noethériens, $\rho:A\to B$ un homomorphisme local, $f:\Spec(B)=X\to\Spec(A)=Y$ le morphisme correspondant. Supposons vérifiées les conditions suivantes:
\begin{enumerate}
  \item[(i)] $f$ est plat et toutes les fibres $X_y=f^{-1}(y)$ ($y\in Y$) sont géométriquement réduites (4.6.2).
  \item[(ii)] Ou bien l'anneau $A$ est géométriquement unibranche (0, 23.2.1), ou bien l'anneau $A$ est unibranche (0, 23.2.1) et $k(B)$ est une extension primaire (4.3.1) de $k(A)$.
\end{enumerate}
Alors, si $\eta$ est le point générique de $Y$, $X_\eta$ est connexe, et pour toute partie fermée $T$ de $X$ telle que $f(T)$ soit rare dans $Y$, $X-T$ est connexe.
\end{theorem}

Les deux conclusions énoncées sont en fait équivalentes; en effet, l'hypothèse que $f(T)$ est rare entraîne que $f(T)$ ne contient pas $\eta$; d'autre part, comme $\{\eta\}$ est proconstructible (I, 9.6) et que $f$ est plat, $X_\eta$ est dense dans $X$ (2.3.10), et par suite si $X_\eta$ est connexe, il en est de même de $X-T$ qui le contient. Inversement, lorsque $U$ parcourt l'ensemble des voisinages ouverts affines de $\eta$ dans $Y$, $X_\eta$ est limite projective des sous-préschémas induits sur les ouverts $f^{-1}(U)$ de $X$ (8.1.2, $a)$); si l'on suppose que $X-T$ est connexe pour toute partie fermée $T$ de $X$ telle que $f(T)$ soit rare dans $Y$, les $f^{-1}(U)$ sont connexes, donc il en est de même de $X_\eta$ (8.4.1).

Montrons d'abord qu'on peut se limiter au cas où $A$ est intègre et $Y=\Spec(A)$ géométriquement unibranche (6.15.1) (ce qui implique que $A$ est géométriquement unibranche, mais n'est pas équivalent à cette condition (6.15.2)). Il est clair qu'on peut d'abord supposer $Y$ réduit, en considérant le morphisme $X\times_Y Y_{\mathrm{red}}\to Y_{\mathrm{red}}$, déduit de $f$ par changement de base, qui est plat, et a mêmes fibres que $f$. On peut donc supposer $A$ intègre et unibranche; si $K$ est le corps des fractions de $A$, il existe alors une sous-$A$-algèbre finie $A''$ de $K$ telle que, si $A'$ est la clôture intégrale de $A$, le morphisme $\Spec(A')\to\Spec(A'')$ soit radiciel (0, 23.2.5); on en conclut (6.15.5) que $\Spec(A'')$ est géométriquement unibranche. Comme $A$ est unibranche, $A'$ est un anneau local, donc il en est de même de $A''$; montrons que l'anneau $B''=B\otimes_A A''$ est aussi un anneau local. En effet, $B''$ est une $B$-algèbre finie, donc un anneau semi-local noethérien; si $k''$ est le corps résiduel de $A''$, les idéaux maximaux de $B''$ sont les points de $\Spec(B'')$ au-dessus du point fermé de $\Spec(B)$, donc les points de $\Spec(k(B)\otimes_{k(A)}k'')$ (I, 3.4.9)

\oldpage[IV-4]{153}

puisque le point fermé de $\Spec(A'')$ est le seul point au-dessus du point fermé de $\Spec(A)$. Or, lorsque $A$ est géométriquement unibranche, $k''$ est une extension radicielle de $k(A)$, donc $\Spec(k(B)\otimes_{k(A)}k'')$ est radiciel sur $\Spec(k(B))$ et ne comprend donc qu'un seul point. Lorsque $A$ est unibranche et $k(B)$ extension primaire de $k(A)$, $\Spec(k(B)\otimes_{k(A)}k'')$ est irréductible (4.3.2) et comme c'est un espace fini et discret (I, 6.4.4) il est encore réduit à un point, ce qui montre que dans les deux cas envisagés $B''$ est local. Si $Y''=\Spec(A'')$, $X''=\Spec(B'')=X\times_Y Y''$, le morphisme $f''=f_{(Y'')}:X''\to Y''$ est plat et a ses fibres géométriquement réduites; en outre, si $\eta''$ est le point générique de $Y''$, $k(\eta'')=k(\eta)=K$ par définition, donc les fibres $X_\eta$ et $X''_{\eta''}$ sont isomorphes. Il suffit donc de prouver que $X''_{\eta''}$ est connexe.

\begin{remark}[18.9.7.1]
\label{IV.18.9.7.1}
Lorsque l'anneau $A$ est \emph{japonais}, on peut, dans ce qui précède, prendre $A''=A'$ par définition (0, 23.1.1); on voit donc que dans ce cas on serait même ramené à prouver le théorème lorsque $A$ est \emph{intégralement clos}.
\end{remark}

\begin{env}[18.9.7.2]
\label{IV.18.9.7.2}
Supposons donc désormais que $Y$ soit intègre et géométriquement unibranche. Notons que si $T$ est une partie fermée de $X$ telle que $f(T)$ soit rare dans $Y$, $T$ ne contient aucun point maximal de $X$ puisque $f$ est plat (2.3.4), donc est rare dans $X$; comme $X$ est un schéma local, donc connexe, pour prouver que $X-T$ est connexe, il suffit, en vertu de (15.5.6.1) de montrer que pour tout $x\in X$ tel que $f(x)\neq\eta$, $\Spec(\sh{O}_{X,x})-\{x\}$ est connexe. Posons $y=f(x)$, $A_1=\sh{O}_{Y,y}$, $B_1=\sh{O}_{X,x}$, $Y_1=\Spec(A_1)$, $X_1=\Spec(B_1)$, $f_1:X_1\to Y_1$ le morphisme correspondant à l'homomorphisme local $A_1\to B_1$ déduit de $\rho$; il est clair que $f_1$ est plat, et il résulte de (4.6.1) que ses fibres sont géométriquement réduites; en outre $A_1$ est intègre et géométriquement unibranche (6.15.1) et $\dim(A_1)\geq 1$. On est ainsi ramené à prouver le
\end{env}

\begin{lemma}[18.9.7.3]
\label{IV.18.9.7.3}
Soient $A$, $B$ deux anneaux locaux noethériens, $Y=\Spec(A)$, $X=\Spec(B)$, $\rho:A\to B$ un homomorphisme local, $f:X\to Y$ le morphisme correspondant. On suppose vérifiées les conditions suivantes:
\begin{enumerate}
  \item[(i)] $f$ est plat et toutes les fibres $X_y$ ($y\in Y$) sont géométriquement réduites.
  \item[(ii)] $A$ est intègre, géométriquement unibranche et $\dim(A)\geq 1$.
\end{enumerate}
Alors, si $b$ est le point fermé de $X$, $X-\{b\}$ est connexe.
\end{lemma}

Notons d'abord qu'en vertu du théorème de Hartshorne (5.10.7), $X-\{b\}$ est connexe si l'on a $\operatorname{prof}(B)\geq 2$. Or on a (6.3.1)
\begin{equation}
  \operatorname{prof}(B)=\operatorname{prof}(A)+\operatorname{prof}(B\otimes_A k)
  \tag{18.9.7.4}
  \label{IV.18.9.7.4}
\end{equation}
en désignant par $k$ le corps résiduel de $A$. D'autre part, puisque $A$ est intègre et $\dim(A)\geq 1$, on a $\operatorname{prof}(A)\geq 1$, donc $\operatorname{prof}(B)\geq 2$ \emph{sauf} lorsque $\operatorname{prof}(B\otimes_A k)=0$; d'ailleurs $B\otimes_A k$ est réduit par (i), donc la relation $\operatorname{prof}(B\otimes_A k)=0$ signifie que $B\otimes_A k$ est un \emph{corps} (0, 16.4.7). On supposera donc désormais que cette dernière condition est vérifiée. On commencera par traiter un cas où la démonstration est très simple.

A) \emph{Cas où $A$ est intégralement clos.} Alors, si $\dim(A)\geq 2$, on a $\operatorname{prof}(A)\geq 2$ (0, 16.5.1), donc $\operatorname{prof}(B)\geq 2$. On n'a donc qu'à considérer le cas où $\dim(A)=1$ et où $B\otimes_A k$ est un corps; $A$ est alors un anneau de valuation discrète, donc régulier, et puisque $B\otimes_A k$

\oldpage[IV-4]{154}

est un corps et que $f$ est plat, $B$ est régulier (0, 17.3.3); mais en outre (6.1.1.1), on a $\dim(B)=\dim(A)+\dim(B\otimes_A k)=\dim(A)=1$, donc $B$ est aussi un anneau de valuation discrète. Mais alors $X-\{b\}$ est réduit à un seul point, d'où le lemme (18.9.7.3) dans ce cas.

On observera que cela prouve l'énoncé du théorème (18.9.7) lorsqu'on suppose en outre que, dans cet énoncé, l'anneau $A$ est \emph{japonais}, car en vertu de la remarque (18.9.7.1), on est ramené au cas où $A$ est \emph{intégralement clos}, et alors, dans la réduction qui précède (18.9.7.3), l'anneau $A_1=\sh{O}_{Y,y}$ est lui aussi intégralement clos et on peut appliquer le résultat qui vient d'être démontré.

B) \emph{Cas où $\dim(A)\geq 2$.} Le résultat de (18.9.7.3) résultera alors des deux lemmes suivants:

\begin{lemma}[18.9.7.5]
\label{IV.18.9.7.5}
Soient $A$ un anneau semi-local noethérien et réduit, $A'$ sa fermeture intégrale dans son anneau total des fractions, $X=\Spec(A)$, $X'=\Spec(A')$; soient $S$ l'ensemble des points fermés de $X$ et $U=X-S$. Supposons que pour tout point $x'\in X'$ au-dessus d'un point de $S$, on ait $\dim(\sh{O}_{X',x'})\geq 2$; alors l'anneau $A^{(1)}=\Gamma(U,\sh{O}_X)$ est entier sur $A$ et est une $A$-algèbre isomorphe à une sous-algèbre de $A'$.
\end{lemma}

Rappelons que si les $\mathfrak p_i$ ($1\leq i\leq m$) sont les idéaux minimaux de $A$, $A'$ est le composé direct des clôtures intégrales $A'_i$ des anneaux intègres $A_i=A/\mathfrak p_i$, de sorte que $X'$ est la somme des schémas $X'_i=\Spec(A'_i)$. Si $X_0$ est la somme des schémas $X_i=\Spec(A_i)$ ($1\leq i\leq m$) et $U_0$ l'image réciproque de $U$ dans $X_0$, l'homomorphisme canonique $\Gamma(U,\sh{O}_X)\to\Gamma(U_0,\sh{O}_{X_0})$ est injectif puisque $X$ est réduit et le morphisme $X_0\to X$ (donc aussi le morphisme $U_0\to U$) surjectif. Comme $\Gamma(U_0,\sh{O}_{X_0})$ est composé direct des $\Gamma(U_i,\sh{O}_{X_i})$ ($1\leq i\leq m$), où $U_i$ est l'image réciproque de $U$ dans $X_i$, tout revient à prouver que pour chaque $i$, $\Gamma(U_i,\sh{O}_{X_i})$ est isomorphe à une sous-$A_i$-algèbre de $A'_i$. Autrement dit, on est ramené au cas où $A$ est un anneau semi-local noethérien intègre. Comme $X$ est réduit et le morphisme $X'\to X$ surjectif, l'homomorphisme $\Gamma(U,\sh{O}_X)\to\Gamma(U',\sh{O}_{X'})$ (où $U'$ est l'image réciproque de $U$ dans $X'$) est injectif, et tout revient à voir que l'homomorphisme canonique $A'=\Gamma(X',\sh{O}_{X'})\to\Gamma(U',\sh{O}_{X'})$ est \emph{bijectif}. Or (I, 8.2.1.1), $\Gamma(U',\sh{O}_{X'})$ est l'intersection des anneaux locaux $A'_{\mathfrak p'}$, où $\mathfrak p'$ parcourt $U'$, et en vertu de l'hypothèse, parmi ces anneaux locaux figurent tous ceux pour lesquels $\mathfrak p'$ est de hauteur 1. Mais le raisonnement de (0, 23.2.7) s'applique aussi à un anneau semi-local intègre noethérien, donc $A'$ est un \emph{anneau de Krull semi-local}, et est donc l'intersection des anneaux locaux $A'_{\mathfrak p'}$, où $\mathfrak p'$ parcourt l'ensemble des idéaux premiers de hauteur 1 de $A'$ (Bourbaki, \emph{Alg. comm.}, chap.~VII, \textsection~1, n\textsuperscript{o}~6, th.~4); \emph{a fortiori}, $A'$ est l'intersection des $A'_{\mathfrak p'}$ pour $\mathfrak p'\in U'$, ce qui achève la démonstration du lemme.

\begin{remark}[18.9.7.6]
\label{IV.18.9.7.6}
On sait (0, 23.2.5) qu'il existe une sous-$A$-algèbre finie $A''$ du corps des fractions $K$ de $A$ telle que le morphisme $\Spec(A')\to\Spec(A'')$ soit \emph{radiciel}; comme ce morphisme est aussi entier et dominant, donc surjectif (II, 6.1.10), c'est un \emph{homéomorphisme} (2.4.5). En vertu de l'interprétation géométrique de $\dim(\sh{O}_{X',x'})$ (5.1.2), on a donc, pour tout point $x''$ de $X''=\Spec(A'')$ au-dessus d'un point $x\in X$, $\dim(\sh{O}_{X'',x''})=\dim(\sh{O}_{X',x'})\geq 2$ si $x'$ est l'unique point de $X'$ au-dessus de $x$. Cela

\oldpage[IV-4]{155}

étant, il résulte d'abord de (0, 16.1.5) et (0, 16.1.4.1) que la relation $\dim(\sh{O}_{X',x'})\geq 2$ entraîne $\dim(\sh{O}_{X,x})\geq 2$. Inversement, supposons que $\dim(\sh{O}_{X,x})\geq 2$ et que $A$ soit universellement caténaire (5.6.2); alors il en est de même de $\sh{O}_{X,x}$ (5.6.3), donc $\dim(\sh{O}_{X'',x''})=\dim(\sh{O}_{X,x})\geq 2$ en vertu de (5.6.10) et on a donc alors $\dim(\sh{O}_{X',x'})\geq 2$. La même conclusion est valable si $\sh{O}_{X,x}$ est \emph{géométriquement unibranche}, car la fibre du morphisme $X'\to X$ au point $x$ est alors réduite à un seul point, et $\sh{O}_{X',x'}$ est une $\sh{O}_{X,x}$-algèbre \emph{entière}; il suffit donc d'appliquer (0, 16.1.5).
\end{remark}

\begin{lemma}[18.9.7.7]
\label{IV.18.9.7.7}
Soient $A$ un anneau local noethérien, de dimension $\geq 2$, intègre et géométriquement unibranche, $Y=\Spec(A)$, $y$ le point fermé de $Y$. Soient $X$ un préschéma localement noethérien et connexe, $f:X\to Y$ un morphisme plat. Alors, pour toute partie fermée $T\subset f^{-1}(y)$, $X-T$ est connexe.
\end{lemma}

Le même raisonnement qu'au début de (18.9.7.2) ramène à prouver que, pour tout point $x\in f^{-1}(y)$, $\Spec(\sh{O}_{X,x})-\{x\}$ est connexe. On peut donc se borner au cas où $X=\Spec(B)$, $B$ étant un anneau local noethérien et l'homomorphisme $A\to B$ correspondant à $f$ étant local; posons $U=Y-\{y\}$, et notons que $V=f^{-1}(U)$ est dense dans $X$ (2.3.10); il suffit donc de prouver que $V$ est \emph{connexe}. Gardons les notations de (18.9.7.5), et posons donc $A^{(1)}=\Gamma(U,\sh{O}_Y)$, $Y_1=\Spec(A^{(1)})$, $X_1=X\times_Y Y_1=\Spec(B\otimes_A A^{(1)})$, de sorte qu'on a le diagramme commutatif
\[
  \xymatrix{
    X \ar[d]_{f} & X_1 \ar[l]^{g_1} \ar[d]^{f_1} \\
    Y & Y_1 \ar[l]_{g}
  }
\]

En vertu de (2.3.1) et de la définition de $A^{(1)}$, on a un isomorphisme canonique $\Gamma(V,\sh{O}_X)\xrightarrow{\sim}B\otimes_A A^{(1)}$, et il suffit donc de prouver que ce dernier anneau est \emph{local}, un tel anneau ne pouvant être le produit de deux anneaux non réduits à 0 (III, 7.8.6.1). Mais on a vu (18.9.7.5) que $A^{(1)}$ est une sous-$A$-algèbre de la clôture intégrale $A'$ de $A$, puisque $\dim(A)\geq 2$. L'hypothèse que $A$ est géométriquement unibranche entraîne que le morphisme $\Spec(A')\to\Spec(A)$ est radiciel au point $y$ (6.15.3); \emph{a fortiori} $g:Y_1\to Y$ est radiciel au point $y$, et par suite $g_1$ est radiciel au point fermé $x$ de $X$; comme en outre $g_1$ est un morphisme entier, les points fermés de $X_1$ sont au-dessus de $x$, donc il n'existe qu'un seul point fermé de $X_1$, ce qui achève la démonstration de (18.9.7.7).

Il est clair que (18.9.7.7) démontre (18.9.7.3) lorsque $\dim(A)\geq 2$ (même sans supposer les fibres $X_y$ géométriquement réduites).

C) \emph{Cas où $\dim(A)=1$.} On a vu au début de la preuve de (18.9.7.3) que l'on peut supposer $\dim(B\otimes_A k)=0$, et par platitude (6.1.1.1), on a donc
\[
  \dim(B)=\dim(A)+\dim(B\otimes_A k)=1.
\]
Puisque $A$ est intègre, $Y$ se compose de deux points, le point générique $\eta$ et le point fermé $a$. En outre, puisque $f$ est plat, toute composante irréductible de $X$ domine $Y$ (2.3.4), donc $X-\{b\}$, qui est l'ensemble des points maximaux $\xi_i$ de $X$, est égal à

\oldpage[IV-4]{156}

l'ensemble sous-jacent à $f^{-1}(\eta)$, et la fibre $f^{-1}(\eta)$ est somme des $\Spec(k(\xi_i))$. L'hypothèse sur les $X_y$ et le fait que ces fibres sont de dimension 0, montrent qu'elles sont \emph{géométriquement régulières} (6.7.6), et \emph{a fortiori} géométriquement normales, autrement dit $f$ est un \emph{morphisme normal}. Mais puisque $A$ est intègre et géométriquement unibranche, il résulte alors de (6.15.10) que $B$ est intègre et géométriquement unibranche, donc $f^{-1}(\eta)=X-\{b\}$ est réduit à un seul point, et \emph{a fortiori} connexe; ceci termine donc la démonstration de (18.9.7.3) et celle de (18.9.7).

\begin{remark}[18.9.7.8]
\label{IV.18.9.7.8}
Dans le cas C) de la démonstration de (18.9.7.3), on peut éviter de faire appel au délicat résultat (6.15.10) en raisonnant de la façon suivante: puisque $A$ est intègre et de dimension 1, il résulte du théorème de Krull-Akizuki (Bourbaki, \emph{Alg. comm.}, chap.~VII, \textsection~2, n\textsuperscript{o}~5, prop.~5) que sa clôture intégrale $A'$ est un anneau \emph{noethérien}. Le même raisonnement qu'au début de la démonstration de (18.9.7) montre alors que $B'=B\otimes_A A'$ est un anneau local; en outre, par platitude, $B'$ est contenu dans l'anneau total des fractions $R$ de l'anneau réduit $B$ (3.3.5). Or, le raisonnement qui prouve le théorème de Krull-Akizuki (Bourbaki, \emph{loc. cit.}) s'applique également à un anneau local noethérien \emph{réduit} de dimension 1, et montre que pour un tel anneau $B$, tout anneau compris entre $B$ et son anneau total des fractions est \emph{noethérien}. L'anneau $B'$ étant un anneau local noethérien et le morphisme $\Spec(B')\to\Spec(A')$ étant normal, $B'$ est un anneau intègre et intégralement clos (6.5.4), et on conclut comme au début de la preuve de (18.9.7).
\end{remark}

\begin{corollary}[18.9.8]
\label{IV.18.9.8}
Avec les notations de (18.9.7), supposons vérifiée la condition (i) de (18.9.7) et l'une des deux conditions suivantes:
\begin{enumerate}
  \item[(ii$'$)] $A$ est un anneau strictement local (18.8.2).
  \item[(ii$''$)] $A$ est un anneau local hensélien et $k(B)$ est une extension primaire de $k(A)$.
\end{enumerate}
Alors toutes les fibres $X_y$ ($y\in Y$) sont géométriquement connexes.
\end{corollary}

Soit $\mathfrak p$ l'idéal premier de $A$ correspondant au point $y$, et posons $A'=A/\mathfrak p$ et $B'=B\otimes_A A'=B/\mathfrak pB$, qui sont évidemment des anneaux locaux noethériens; il est clair que le morphisme $f':\Spec(B')\to\Spec(A')$ vérifie la condition (i) de (18.9.7); comme $y$ est le point générique de $\Spec(A')$ et que $k(A')=k(A)$, $k(B')=k(B)$, on voit qu'il suffit de vérifier que, dans le cas (ii$'$) (resp. (ii$''$)), $A'$ est géométriquement unibranche (resp. unibranche). Mais $A'$ est intègre, et dans le cas (ii$'$) (resp. (ii$''$)) il est strictement local (resp. hensélien) en vertu de (18.5.10). On en conclut que $A'$ est géométriquement unibranche (resp. unibranche) en vertu de (18.8.15) (resp. (18.6.12)).

\begin{corollary}[18.9.9]
\label{IV.18.9.9}
Soit $A$ un anneau local noethérien hensélien, universellement japonais (\emph{i.e.} ((7.6.4) et (7.7.2)), dont les fibres formelles (7.3.13) sont géométriquement réduites); alors les fibres formelles de $A$ sont géométriquement connexes.
\end{corollary}

Il suffit d'appliquer (18.9.8) au cas où $B=\widehat A$, puisque $B$ est un $A$-module plat et que $k(B)=k(A)$.

\begin{remark}[18.9.10]
\label{IV.18.9.10}
Les démonstrations de (18.9.4) et (18.9.5) n'utilisent que le fait que les fibres formelles de $A$ sont géométriquement connexes. Les conclusions

\oldpage[IV-4]{157}

de ces deux propositions sont donc encore valables lorsqu'on suppose l'anneau local $A$ noethérien, hensélien et universellement japonais.
\end{remark}

\begin{corollary}[18.9.11]
\label{IV.18.9.11}
Soient $X$, $Y$ deux préschémas localement noethériens, $Y$ étant supposé géométriquement unibranche et $X$ connexe. Soit $f:X\to Y$ un morphisme réduit (\emph{i.e.} (6.8.1) plat et à fibres géométriquement réduites). Alors, pour toute partie fermée $T$ de $X$ telle que $f(T)$ soit rare dans $Y$, $X-T$ est connexe.
\end{corollary}

Comme $f$ est plat et que $f(T)$, étant rare, ne peut contenir de point maximal de $Y$, $T$ ne peut contenir de point maximal de $X$ (2.3.4), donc est rare dans $X$. Utilisant (15.5.6.1), il suffit de montrer que, pour tout $x\in T$, $\Spec(\sh{O}_{X,x})-\{x\}$ est connexe. Posons $y=f(x)$, qui n'est pas un point maximal de $Y$. Comme l'anneau $\sh{O}_{Y,y}$ est géométriquement unibranche, on peut appliquer (18.9.7) au morphisme
\[
  f_1:\Spec(\sh{O}_{X,x})\longrightarrow\Spec(\sh{O}_{Y,y}),
\]
qui vérifie les conditions (i) et (ii) de ce théorème, et à la partie fermée $\{x\}$ de $\Spec(\sh{O}_{X,x})$, puisque $f_1(\{x\})=\{y\}$ est rare dans $\Spec(\sh{O}_{Y,y})$.

\subsection{Préschémas étales sur un préschéma géométriquement unibranche ou normal}
\EGAReaderTarget{ega4.c04.subsec.18-10}{subsection:IV.18.10}

\begin{theorem}[18.10.1]
\label{IV.18.10.1}
Soient $f:X\to Y$ un morphisme localement de présentation finie, $x$ un point de $X$, $y=f(x)$. Supposons que $Y$ soit intègre et géométriquement unibranche au point $y$. Alors, pour que $f$ soit étale au point $x$, il faut et il suffit que les deux conditions suivantes soient vérifiées:
\begin{enumerate}
  \item[(i)] $f$ est non ramifié au point $x$;
  \item[(ii)] l'homomorphisme $\sh{O}_{Y,y}\to\sh{O}_{X,x}$ est injectif.
\end{enumerate}
De plus, s'il en est ainsi, $X$ est intègre et géométriquement unibranche au point $x$.

Enfin, si $\sh{O}_{Y,y}$ est noethérien, on peut, dans le critère précédent, remplacer la condition (ii) par la condition

\emph{(ii bis)} On a $\dim\sh{O}_{Y,y}\leq\dim\sh{O}_{X,x}$.
\end{theorem}

Le fait que, si $f$ est étale, $X$ soit intègre et géométriquement unibranche au point $x$ est un cas particulier de (17.5.7). Il est clair d'autre part que les conditions (i) et (ii) sont nécessaires pour que $f$ soit étale au point $x$, puisque $\sh{O}_{X,x}$ est alors un $\sh{O}_{Y,y}$-module fidèlement plat (17.6.1). Prouvons que ces conditions sont suffisantes.

Posons $A=\sh{O}_{Y,y}$, $A'={}^{hs}\!A$ (18.8.7), $Y'=\Spec(A')$, $X'=X\times_Y Y'$, $f'=f_{(Y')}:X'\to Y'$. Soient $y'$ le point fermé de $Y'$, $x'$ un point de $X'$ au-dessus de $x$ et de $y'$; puisque le morphisme $g:Y'\to Y$ est plat (18.8.8), tout revient à voir que $f'$ est étale au point $x'$ (17.7.1, (iii)). L'hypothèse (ii) entraîne que le morphisme $\Spec(\sh{O}_{X,x})\to\Spec(\sh{O}_{Y,y})$ est dominant (I, 1.2.7); comme $\sh{O}_{Y,y}$ est intègre, $\Spec(\sh{O}_{Y,y})$ a un seul point générique $y_1$, et il y a donc une générisation $x_1$ de $x$ dans $X$ au-dessus de $y_1$. Le morphisme projection $X'\to X$ étant plat, il existe une générisation $x'_1$ de $x'$ au-dessus de $x_1$ (2.3.4); posons $y'_1=f'(x'_1)$; on a par suite $g(y'_1)=y_1$. Or, l'hypothèse sur $y$ entraîne que $A'$ est intègre (18.8.15), donc $Y'$ a un seul point générique $\eta$, nécessairement au-dessus de $y_1$ puisque $g$

\oldpage[IV-4]{158}

est plat (2.3.4); d'autre part, $\eta$ est aussi point générique de $g^{-1}(y_1)$ ($0_{\mathrm I}$, 2.1.8), et on sait par ailleurs que les fibres de $g$ sont des espaces discrets (18.8.12), donc $g^{-1}(y_1)=\{\eta\}$, et par suite $y'_1=\eta$. Le morphisme $\Spec(\sh{O}_{X',x'})\to\Spec(\sh{O}_{Y',y'})$ restriction de $f'$ est donc dominant, et puisque $\sh{O}_{Y',y'}=A'$ est intègre, l'homomorphisme correspondant $\sh{O}_{Y',y'}\to\sh{O}_{X',x'}$ est injectif (I, 1.2.7). D'autre part, $f'$ est non ramifié au point $x'$, (donc (18.8.3, $d$)), il y a un voisinage $U$ de $x'$ dans $X'$ tel que $f'|U$ soit une immersion fermée de $U$ dans $Y'$; \emph{a fortiori} l'homomorphisme $\sh{O}_{Y',y'}\to\sh{O}_{X',x'}$ est surjectif, et puisqu'il est injectif, il est bijectif. Mais alors $\sh{O}_{X',x'}$ est un $\sh{O}_{Y',y'}$-module plat, et $f'$ est donc étale au point $x'$ (17.6.1).

Enfin, lorsque $\sh{O}_{Y,y}$ est noethérien, on sait que si $f$ est étale au point $x$, on a $\dim\sh{O}_{Y,y}=\dim\sh{O}_{X,x}$ (17.6.4). Réciproquement, supposons vérifiées les conditions (i) et (ii \emph{bis}), et montrons qu'elles entraînent (ii). Soit $\mathfrak J$ le noyau de l'homomorphisme canonique $\sh{O}_{Y,y}\to\sh{O}_{X,x}$; restreignant au besoin $Y$ à un voisinage de $y$, on peut supposer que $\mathfrak J=\mathscr{J}_y$, où $\mathscr{J}$ est un Idéal de $\sh{O}_Y$; si $Y_1$ est le sous-préschéma fermé de $Y$ défini par $\mathscr{J}$, on peut, en restreignant encore au besoin $X$ et $Y$ à des voisinages de $x$ et $y$ respectivement, supposer que $f$ se factorise en $X\xrightarrow{f_1}Y_1\to Y$ (I, 6.5.1). Il est clair que $f_1$ est encore non ramifié au point $x$, donc (17.4.1) quasi-fini en ce point; la démonstration de (5.4.1, (i)) prouve alors que l'on a $\dim(\sh{O}_{X,x})\leq\dim(\sh{O}_{Y,y}/\mathfrak J)$. Mais si l'on avait $\mathfrak J\neq0$, on en conclurait que $\dim(\sh{O}_{Y,y}/\mathfrak J)<\dim(\sh{O}_{Y,y})$ puisque $Y$ est intègre ($0$, 16.1.2.2); on aurait par suite $\dim\sh{O}_{X,x}<\dim\sh{O}_{Y,y}$ contrairement à l'hypothèse, ce qui prouve (ii).

\begin{remark}[18.10.2]
\label{IV.18.10.2}
\begin{enumerate}
  \item[(i)] Lorsque $A=\sh{O}_{Y,y}$ est noethérien, et que l'on sait que son complété $\widehat A$ est intègre (par exemple si $A$ est régulier), on peut, dans la démonstration précédente, remplacer $A'$ par $\widehat A$.
  \item[(ii)] Lorsque $Y$ est localement noethérien, on peut donner de (18.10.1) une démonstration plus rapide, n'utilisant pas la hensélisation stricte, mais faisant intervenir les résultats assez délicats des \textsection\textsection~14 et 15. Comme $f$ est quasi-fini au point $x$ par hypothèse (17.4.1), on peut, en remplaçant $X$ par un voisinage ouvert de $x$, supposer que $f^{-1}(y)=\{x\}$. L'hypothèse (ii) entraîne, comme on l'a vu, l'existence d'une composante irréductible $Z$ de $X$ contenant $x$ et dominant l'unique composante irréductible $Y_0$ de $Y$ contenant $y$. Le morphisme $f|Z$ étant quasi-fini au point $x$, donc équidimensionnel en ce point (13.2.2), il résulte de l'hypothèse sur $Y$ et du critère de Chevalley (14.4.4) que $g=f|Z$ est universellement ouvert au point $x$. Par ailleurs puisque $f$ (donc aussi $g$) est non ramifié au point $x$, $g^{-1}(y)$ est géométriquement réduit sur $k(y)$ (17.4.1); comme $\sh{O}_{Y,y}$ est intègre, on conclut de (15.2.3) que $g$ est plat au point $x$, donc $f$ est étale en ce point (18.4.9).
  \item[(iii)] Les notations et les hypothèses sur $Y$ étant celles de (18.10.1), supposons seulement que $f$ soit localement de type fini (et non nécessairement localement de présentation finie). Alors, pour que $\sh{O}_{X,x}$ soit une $\sh{O}_{Y,y}$-algèbre essentiellement étale (18.6.1), il faut et il suffit que $f$ vérifie les deux conditions suivantes:

\oldpage[IV-4]{159}

\begin{enumerate}
  \item[$1^\circ$] $f$ est formellement non ramifié au point $x$;
  \item[$2^\circ$] l'homomorphisme $\sh{O}_{Y,y}\to\sh{O}_{X,x}$ est injectif.
\end{enumerate}

Il n'y a encore à prouver que la suffisance de ces conditions. Tenant compte de (18.4.12), il suffit de montrer que $f$ est plat au point $x$. Or, en reprenant la démonstration de (18.10.1) et les notations utilisées dans cette démonstration, il suffit de montrer que $f'$ est plat au point $x'$ (2.5.1); par ailleurs $f'$ est formellement non ramifié au point $x'$ (17.1.3); on peut donc se borner à faire la démonstration lorsque $Y=\Spec(A)$, où $A$ est strictement local et $y=f(x)$ le point fermé de $Y$. Utilisant alors (18.8.3) on voit que l'homomorphisme $\sh{O}_{Y,y}\to\sh{O}_{X,x}$ est surjectif, donc bijectif en vertu de l'hypothèse, et cela suffit à montrer la platitude de $f$ au point $x$.

Supposons de plus que $Y$ soit localement intègre au point $y$ (I, 2.1.8). Alors les conditions $1^\circ$ et $2^\circ$ précédentes (jointe au fait que $Y$ est géométriquement unibranche au point $y$ et $f$ localement de type fini) entraînent déjà que $f$ est étale au point $x$. En effet, cela résulte de ce qui précède et de (18.4.13).
\end{enumerate}
\end{remark}

\begin{corollary}[18.10.3]
\label{IV.18.10.3}
Soient $Y$ un préschéma intègre et géométriquement unibranche, $\eta$ son point générique, $X$ un préschéma connexe, $f:X\to Y$ un morphisme localement de type fini et tel que $f^{-1}(\eta)$ soit non vide. Alors, si $f$ est formellement non ramifié, $f$ est étale, et $X$ est intègre et géométriquement unibranche.
\end{corollary}

Il résulte de l'hypothèse et de (18.4.13) que $f$ est étale en tous les points où il est plat, et en particulier aux points de $f^{-1}(\eta)$, puisque $\sh{O}_{Y,\eta}=k(\eta)$ est un corps. L'ensemble ouvert $U$ des points de $X$ où $f$ est étale est donc non vide. Montrons d'autre part que pour tout $x\in\overline U$, l'homomorphisme $\sh{O}_{Y,f(x)}\to\sh{O}_{X,x}$ est injectif. Soit en effet $s$ une section de $\sh{O}_Y$ au-dessus d'un voisinage ouvert de $f(x)$, qu'on peut toujours supposer être $Y$ lui-même (la question étant locale sur $X$ et $Y$), et supposons que son image $t\in\Gamma(Y,f_*(\sh{O}_X))=\Gamma(X,\sh{O}_X)$ ait un germe nul au point $x$, donc s'annule dans un voisinage ouvert $V$ de $x$; alors la restriction de $t$ à l'ouvert $U\cap V\neq\varnothing$ est nulle; mais la restriction $f|(U\cap V)$ est étale, donc $W=f(U\cap V)$ est ouvert dans $Y$ et contient par suite le point générique $\eta$; l'homomorphisme $\sh{O}_{Y,\eta}\to\sh{O}_{X,z}$ étant injectif pour tout point $z\in f^{-1}(\eta)$, l'hypothèse que $t|(U\cap V)=0$ entraîne que $s_\eta=0$, et comme $Y$ est intègre, cela entraîne $s=0$, d'où notre assertion. Appliquant maintenant (18.10.2, (iii)), on conclut que $f$ est plat au point $x$, autrement dit $x\in U$, et $U$ est donc à la fois ouvert et fermé dans $X$; comme il est non vide et que $X$ est connexe, on a $U=X$.

Notons maintenant que puisque $f$ est étale, donc localement quasi-fini et que $Y$ est irréductible, les points maximaux de $X$ appartiennent à $f^{-1}(\eta)$ (2.3.4) et pour tout $x\in X$, il y a un voisinage de $x$ qui ne contient qu'un nombre fini de ces points maximaux; autrement dit l'ensemble des composantes irréductibles de $X$ est localement fini. Mais par (18.10.1) $X$ est intègre et géométriquement unibranche en tout point, donc tout point de $X$ n'appartient qu'à une seule composante irréductible, et puisque ces dernières forment un ensemble localement fini, elles sont ouvertes (et fermées) dans $X$; puisque $X$ est connexe, il est intègre.

\oldpage[IV-4]{160}

\begin{remark}[18.10.3.1]
\label{IV.18.10.3.1}
Si $f:X\to Y$ est localement de présentation finie, dominant et quasi-compact, la fibre $f^{-1}(\eta)$ est non vide en vertu de (I, 1.1.5). Par contre, si l'on ne suppose pas $f$ quasi-compact, $f$ peut être non ramifié et dominant sans que $f$ soit étale (en supposant toujours vérifiées les hypothèses de (18.10.3) pour $X$ et $Y$). C'est ce que montre l'exemple suivant: on prend pour $Y$ le plan affine $\Spec(\mathbf C[T,U])$; on considère d'autre part une famille $(X_j)_{j\in\mathbf Z}$ de préschémas isomorphes à la droite affine $\Spec(\mathbf C[T])$, et on forme le préschéma obtenu en « recollant » $X_{2j}$ et $X_{2j+1}$ au point $-1$ et $X_{2j+1}$ et $X_{2j+2}$ au point $+1$. Définissons d'autre part $f:X\to Y$ comme étant égal sur $X_{2j}$ à l'immersion fermée dont l'image est la droite d'équation $y=2j$, transformant le point $-1$ en $(2j-1,2j)$ et le point $+1$ en $(2j+1,2j)$; sur $X_{2j+1}$, $f_j$ est l'immersion fermée dont l'image est la droite d'équation $x=2j+1$, transformant le point $-1$ en $(2j+1,2j)$ et le point $-1$ en $(2j+1,2j+2)$. Il est clair que $f$ est une immersion locale (donc est non ramifié) et est dominant mais la fibre du point générique de $Y$ est vide et $f$ n'est pas étale.
\end{remark}

\begin{corollary}[18.10.4]
\label{IV.18.10.4}
Soient $g:Y\to S$, $h:X\to S$ deux morphismes localement de présentation finie, $f:X\to Y$ un $S$-morphisme, $x$ un point de $X$, $y=f(x)$, $s=g(y)=h(x)$. On suppose que $h$ est plat au point $x$ et $g^{-1}(s)$ normal au point $y$. Alors, pour que $f$ soit étale au point $x$, il faut et il suffit que $f$ soit non ramifié au point $x$ et que $\dim_xh^{-1}(s)=\dim_yg^{-1}(s)$.
\end{corollary}

Les conditions sont nécessaires en vertu du fait que si $f$ est étale au point $x$, il en est de même de $f_s:h^{-1}(s)\to g^{-1}(s)$, et la conclusion résulte alors de ce que $f_s$ est plat et quasi-fini au point $x$ (6.1.2). Inversement, si les conditions de l'énoncé sont vérifiées, pour voir que $f$ est étale au point $x$, il suffit, en vertu de (17.8.3) de voir que $f_s$ l'est. On peut donc se borner au cas où $S$ est le spectre d'un corps. Il résulte alors de l'hypothèse $\dim_xX=\dim_yY$ et du fait que $f$ est non ramifié au point $x$ que l'on a (par (5.2.3) et (17.4.1)) $\dim\sh{O}_{X,x}=\dim\sh{O}_{Y,y}$; puisque $f$ est quasi-fini au point $x$ et $\sh{O}_{Y,y}$ intègre, on déduit de ($0_{\mathrm I}$, 7.4.4) et ($0$, 16.3.10) que l'homomorphisme $\sh{O}_{Y,y}\to\sh{O}_{X,x}$ est injectif; il suffit alors pour conclure d'appliquer (18.10.1).

\begin{corollary}[18.10.5]
\label{IV.18.10.5}
Avec les notations de (18.10.4), on suppose que $h$ est plat et que $g^{-1}(s)$ est normal pour tout $s\in S$ (ce qui sera par exemple le cas si $g$ est lisse). Pour que $f$ soit une immersion ouverte, il faut et il suffit que $f$ soit un monomorphisme et que, pour tout $x\in X$, on ait $\dim_xh^{-1}(s)=\dim_yg^{-1}(s)$, où $y=f(x)$ et $s=g(y)=h(x)$.
\end{corollary}

On n'a à prouver que la suffisance des conditions; or, il résulte de (17.1.3) qu'un monomorphisme localement de présentation finie est non ramifié, donc les hypothèses entraînent que $f$ est étale (18.10.4); on conclut alors par (17.9.1).

\begin{lemma}[18.10.6]
\label{IV.18.10.6}
Soit $f:X\to Y$ un morphisme plat et localement quasi-fini ($\mathrm{Err}_{\mathrm{III}}$, 20).
\begin{enumerate}
  \item[(i)] L'ensemble $\operatorname{Max}(X)$ des points maximaux de $X$ est en correspondance biunivoque canonique avec l'ensemble $\coprod_{y\in\operatorname{Max}(Y)}f^{-1}(y)$.
  \item[(ii)] Si $Y$ est irréductible et de point générique $\eta$, l'ensemble des composantes irréductibles de $X$ est en correspondance biunivoque canonique avec $f^{-1}(\eta)$, et en particulier est fini si et seulement si $f^{-1}(\eta)$ est fini.

\oldpage[IV-4]{161}

  \item[(iii)] Si l'ensemble des composantes irréductibles de $Y$ est localement fini, il en est de même de l'ensemble des composantes irréductibles de $X$, et en particulier $X$ est localement connexe.
\end{enumerate}
\end{lemma}

L'assertion (i) résulte aussitôt de (2.3.4) et ($0_{\mathrm I}$, 2.1.6) et du fait que les fibres $f^{-1}(y)$ sont des ensembles discrets. Il est clair que (ii) est un cas particulier de (i). Enfin, pour prouver (iii), on peut, en vertu de (i), se borner au cas où $f$ est quasi-fini et où $Y$ est irréductible ($0_{\mathrm I}$, 2.1.6) et la conclusion résulte alors de (ii) et du fait que les fibres de $f$ sont des ensembles finis (II, 6.2.2).

\begin{proposition}[18.10.7]
\label{IV.18.10.7}
Soient $Y$ un préschéma géométriquement unibranche et irréductible (resp. intègre), $f:X\to Y$ un morphisme étale. Alors $X$ est isomorphe à une somme de préschémas irréductibles (resp. intègres) $X_\lambda$, où $\lambda$ parcourt la fibre $f^{-1}(\eta)$ du point générique $\eta$ de $Y$. Les $X_\lambda$ sont géométriquement unibranches; si $Y$ est normal, les $X_\lambda$ sont normaux.
\end{proposition}

En effet, $X$ est géométriquement unibranche (17.5.7), donc tout point de $X$ n'appartient qu'à une seule composante irréductible; en outre, en vertu de (18.10.6) et (17.6.1) l'ensemble des composantes irréductibles de $X$ est localement fini, donc ($0_{\mathrm I}$, 2.1.6) ce sont les composantes connexes $X_\lambda$ de $X$ ($\lambda\in f^{-1}(\eta)$), et $X$ est somme des $X_\lambda$; en vertu de (17.5.7), si $Y$ est intègre (resp. normal), les $X_\lambda$ sont intègres (resp. normaux).

\begin{proposition}[18.10.8]
\label{IV.18.10.8}
Soient $Y$ un préschéma normal et intègre, de point générique $\eta$, $K=k(\eta)$ son corps de fonctions rationnelles, $f:X\to Y$ un morphisme étale. On suppose que $f^{-1}(\eta)$ est finie et non vide, de sorte que le $K$-schéma $f^{-1}(\eta)$ est le spectre d'une $K$-algèbre finie séparable $L$, composée directe d'un nombre fini de corps $K_i$ ($1\leq i\leq n$), extensions finies séparables de $K$ (17.6.2). Soit $Y'$ la fermeture intégrale de $Y$ dans $L$, isomorphe à la somme des fermetures intégrales $Y_i$ de $Y$ dans les $K_i$ (II, 6.3.6). Alors le morphisme $f$ se factorise d'une seule manière en $f:X\xrightarrow{f'}Y'\xrightarrow{g}Y$, tel que la restriction $f'_\eta:f^{-1}(\eta)\to g^{-1}(\eta)$ soit l'isomorphisme canonique. En outre $f'$ est un isomorphisme local; pour que $f'$ soit une immersion ouverte, il faut et il suffit que $f$ soit séparé.
\end{proposition}

En vertu de (18.10.7), on peut se borner au cas où $X$ est intègre et normal, $L$ son corps de fonctions rationnelles, $f$ étant dominant. L'existence et l'unicité de la factorisation considérée de $f$ résultent alors de (II, 6.3.9). Comme $f'$ est birationnel et localement quasi-fini, les dernières assertions résultent de (8.12.10).

\begin{corollary}[18.10.9]
\label{IV.18.10.9}
Sous les hypothèses de (18.10.8), pour que $f$ soit un morphisme fini (autrement dit, pour que $X$ soit un revêtement étale de $Y$), il faut et il suffit que $X$ soit isomorphe à la fermeture intégrale $Y'$ de $Y$ dans $L$.
\end{corollary}

Si $f$ est fini, l'injection canonique $j:X\to Y'$ (qui est une immersion ouverte) est aussi un morphisme fini (II, 6.1.5, (v)), donc un morphisme fermé (II, 6.1.10), et comme $j(X)$ est dense dans $Y'$ par (18.10.8), on a $j(X)=Y'$. Inversement, comme $Y$ est normal et $L$ composé direct d'extensions finies séparables de $K$, $Y'$ est fini sur $Y$ (Bourbaki, \emph{Alg. comm.}, chap.~V, \textsection~1, n\textsuperscript{o}~6, cor.~1 de la prop.~18), d'où le corollaire.

\begin{env}[18.10.10]
\label{IV.18.10.10}
Soient $Y$ un préschéma normal et intègre, $K=R(Y)$ son corps de fonctions rationnelles. Nous dirons qu'une $K$-algèbre de rang fini $L$ est non ramifiée sur $Y$ si: $1^\circ$ $L$ est une $K$-algèbre séparable, donc composée directe d'extensions finies séparables $K_i$ de $K$ ($1\leq i\leq n$); $2^\circ$ la fermeture intégrale $Y'$ de $Y$ dans $L$ (somme des préschémas fermetures intégrales de $Y$ dans les $K_i$) est non ramifiée sur $Y$ (ce qui, par (18.10.3),

\oldpage[IV-4]{162}

équivaut à dire que $Y'$ est étale sur $Y$). Lorsque $Y=\Spec(A)$, où $A$ est un anneau intègre et intégralement clos, $K$ son corps des fractions, on dira encore (par abus de langage, et lorsqu'une confusion avec la terminologie de (17.3.2, (ii)) n'est pas à craindre) que $L$ est non ramifiée sur $A$ au lieu de « non ramifiée sur $Y$ ».
\end{env}

\begin{remark}[18.10.11]
\EGAReaderTarget{ega4.c04.rem.18-10-11}{IV.18.10.11}
Au lieu de dire que $L$ est non ramifiée sur $Y$ certains auteurs disent aussi que $L$ est « non ramifiée sur $K$ »; nous nous garderons de suivre cet usage qui prête à confusions.
\end{remark}

On peut alors exprimer (18.10.9) sous la forme suivante:

\begin{corollary}[18.10.12]
\label{IV.18.10.12}
Soient $Y$ un préschéma normal et intègre, $K=R(Y)$ son corps de fonctions rationnelles. Alors le foncteur $X\rightsquigarrow R(X)$ est une équivalence de la catégorie des revêtements étales de $Y$, et de la catégorie des $K$-algèbres finies étales qui sont non ramifiées sur $Y$. On obtient un foncteur quasi-inverse en faisant correspondre à toute $K$-algèbre finie étale $L$, non ramifiée sur $Y$, la fermeture intégrale $Y'$ de $Y$ dans $L$.
\end{corollary}

\begin{proposition}[18.10.13]
\label{IV.18.10.13}
Soient $Y$ un préschéma normal et intègre, $K=R(Y)$ son corps de fonctions rationnelles.
\begin{enumerate}
  \item[(i)] Le corps $K$ est une $K$-algèbre non ramifiée sur $Y$.
  \item[(ii)] Soit $L$ une extension finie de $K$, non ramifiée sur $Y$, et soit $Z$ la fermeture intégrale de $Y$ dans $L$. Alors, si $M$ est une $L$-algèbre non ramifiée sur $Z$, $M$ est aussi une $K$-algèbre non ramifiée sur $Y$ (« transitivité » de la non-ramification).
  \item[(iii)] Soient $Y'$ un préschéma normal et intègre, $K'$ son corps de fonctions rationnelles, $g:Y'\to Y$ un morphisme dominant. Si $L$ est une $K$-algèbre non ramifiée sur $Y$, alors $L\otimes_K K'$ est une $K'$-algèbre non ramifiée sur $Y'$ (propriété de « translation »). Si de plus $Y=\Spec(A)$ et $Y'=\Spec(A')$ sont affines, et si $C$ est la fermeture intégrale de $A$ dans $L$, alors $C'=C\otimes_A A'$ est la fermeture intégrale de $A'$ dans $L\otimes_K K'$.
\end{enumerate}
\end{proposition}

L'assertion (i) est immédiate, $Y$ étant par hypothèse son propre normalisé dans $K$. Pour prouver (ii), soit $Z'$ la fermeture intégrale de $Z$ dans $M$, qui est par hypothèse un revêtement étale de $Z$. Comme $Z$ est étale et séparé sur $Y$, il en est de même de $Z'$ (17.3.3) qui est évidemment la fermeture intégrale de $Y$ dans $M$. Comme $M$ est une $K$-algèbre finie et séparable, il résulte de (18.10.10) que $M$ est non ramifiée sur $Y$. Enfin, pour démontrer (iii), notons que si $Z$ est la fermeture intégrale de $Y$ dans $L$, $Z\times_Y Y'$ est étale et séparé sur $Y'$ (17.3.3), fini sur $Y'$ puisque $Z$ est fini sur $Y$, et admet pour anneau de fonctions rationnelles $L\otimes_K K'$ (I, 3.4.9); comme $Y'$ est normal, il en est de même de $Z\times_Y Y'$ (17.5.7), qui est donc la fermeture intégrale de $Y'$ dans $L\otimes_K K'$, ce qui achève la démonstration.

\begin{corollary}[18.10.14]
\label{IV.18.10.14}
Soient $Y$ et $Y'$ deux préschémas normaux et intègres, $K$, $K'$ leurs corps de fonctions rationnelles respectifs, $g:Y'\to Y$ un morphisme dominant, de sorte que $K'$ est une extension de $K$. Soient $L$ une extension finie de $K$ non ramifiée sur $Y$, $L_1$ une extension composée (Bourbaki, \emph{Alg.}, chap.~VIII, \textsection~8, déf.~1) de $L$ et de $K'$. Alors $L_1$ est une extension de $K'$ non ramifiée sur $Y'$. Si de plus $Y=\Spec(A)$, $Y'=\Spec(A')$ sont affines, et si $C$ est la fermeture intégrale de $A$ dans $L$, alors le sous-anneau $C_1=A[C,A']$ de $L_1$ est la fermeture intégrale de $A'$ dans $L_1$.
\end{corollary}

En effet, $L\otimes_K K'$ est composée directe d'extensions finies séparables de $K'$, et $L_1$

\oldpage[IV-4]{163}

est l'une de ces extensions; donc la fermeture intégrale de $Y$ dans $L$ est somme de préschémas dont l'un est la fermeture intégrale de $Y$ dans $L_1$. Le corollaire résulte donc aussitôt de (18.10.14).

Lorsque par exemple $A$ est l'anneau des entiers d'un corps de nombres algébriques $K$, $K'$ et $L$ des extensions algébriques de $K$, il y a des exemples classiques montrant que la relation $C_1=A[C,A']$ est en défaut lorsque $L$ n'est pas non ramifiée sur $A$ [\ ].

\begin{env}[18.10.15]
\label{IV.18.10.15}
Supposons que $Y=\Spec(A)$ soit affine, $A$ étant intègre et intégralement clos; soient $K$ son corps des fractions, $L$ une extension finie séparable de $K$, et supposons que la fermeture intégrale $C$ de $A$ dans $L$ soit un $A$-module projectif de type fini (ce qui sera par exemple le cas si $A$ est un anneau de Dedekind (Bourbaki, \emph{Alg. comm.}, chap.~VII, \textsection~4, n\textsuperscript{o}~10, prop.~22)). Dire que $L$ est non ramifiée sur $A$ signifie que $C$ est étale (18.10.10) et en vertu de (18.2.7, (ii)), cela équivaut donc à dire que le discriminant $d_{C/A}$ de $\Spec(C)$ sur $\Spec(A)$ est inversible dans $A$. Dans le cas particulier où $C$ est un $A$-module libre et $(x_i)_{1\leq i\leq n}$ une base de $C$ sur $A$, cela signifie que $\det(\operatorname{Tr}_{C/A}(x_ix_j))$ est inversible dans $A$.
\end{env}

\begin{theorem}[18.10.16]
\label{IV.18.10.16}
Soient $Y$ un préschéma intègre, $\eta$ son point générique, $f:X\to Y$ un morphisme séparé et localement de type fini. On suppose que toute composante irréductible de $X$ domine $Y$ et que la fibre générique $X_\eta=f^{-1}(\eta)$ est un préschéma fini sur $K=k(\eta)$, donc $(X_\eta)_{\mathrm{red}}$ est égal à $\Spec(L)$, où $L=\prod_iL_i$ est un produit d'extensions finies $L_i$ de $K$. On pose
\[
  n=[L:K]=\sum_i[L_i:K],\qquad n_s=\sum_i[L_i:K]_s \quad\text{(somme des degrés séparables des $L_i$)}.
\]
Soit $y$ un point géométriquement unibranche de $Y$, et désignons par $n(y)$ la somme des degrés séparables sur $k(y)$ des corps résiduels des points isolés de $f^{-1}(y)$. Alors:
\begin{enumerate}
  \item[(i)] On a $n(y)\leq n_s\leq n$.
  \item[(ii)] Supposons $X$ réduit et le point $y$ normal dans $Y$. Pour que $n(y)=n$, il faut et il suffit qu'il existe un voisinage ouvert $U$ de $y$ dans $Y$ tel que la restriction $f^{-1}(U)\to U$ de $f$ soit un morphisme étale et fini.
\end{enumerate}
\end{theorem}

\emph{A) Réduction au cas où $f$ est quasi-fini.} --- On sait (13.1.4) que l'ensemble $Z$ des points $x\in X$ isolés dans $f^{-1}(f(x))$ est ouvert dans $X$ et par hypothèse $Z$ contient $f^{-1}(\eta)$. Puisque $f^{-1}(\eta)$ est finie, $Z$ est réunion filtrante croissante d'ouverts quasi-compacts $V_\lambda$ contenant $f^{-1}(\eta)$ (donc denses dans $X$ par hypothèse). Si $n_\lambda(y)$ est la somme des degrés séparables sur $k(y)$ des corps résiduels des points de $V_\lambda\cap f^{-1}(y)$, on a $n(y)=\sup_\lambda n_\lambda(y)$; pour prouver (i), il suffira donc de montrer que $n_\lambda(y)\leq n_s$. D'autre part, si $n(y)=n$, il y a un indice $\lambda$ tel que $n_\lambda(y)=n$. Supposons démontré qu'il existe alors un voisinage ouvert $U$ de $y$ dans $Y$ tel que la restriction $V_\lambda\cap f^{-1}(U)\to U$ de $f$ soit un morphisme étale et fini. Puisque $f$ est séparé, l'injection canonique $V_\lambda\cap f^{-1}(U)\to f^{-1}(U)$ est alors aussi un morphisme fini (II, 6.1.5, (v)), donc $V_\lambda\cap f^{-1}(U)$ est fermé dans $f^{-1}(U)$ (II, 6.1.10); mais étant partout dense dans $f^{-1}(U)$, il lui est nécessairement égal.

\emph{B) Réduction au cas où $X$ est réduit, $f$ fini et $Y=\Spec(\sh{O}_{Y,y})$.} --- On notera que l'hypothèse sur $f$ entraîne qu'elle est aussi vérifiée pour $f_{\mathrm{red}}$ et que la conclusion de (i) ne concerne que $f_{\mathrm{red}}$. On peut donc remplacer partout $f$ par $f_{\mathrm{red}}$ et supposer donc $X$ réduit.

\oldpage[IV-4]{164}

D'autre part, en vertu de (18.12.13)\footnote{Le lecteur vérifiera que (18.10.16) n'est pas utilisé dans la démonstration de (18.12.13). Si $f$ est supposé localement de présentation finie (resp. quasi-projectif), on peut remplacer (18.12.13) par (8.12.8) (resp. (8.12.6)).}, on peut (après avoir remplacé $Y$ par un voisinage ouvert affine $V$ de $y$ et $f$ par sa restriction $f^{-1}(V)\to V$) factoriser $f$ en $X\xrightarrow{j}X'\xrightarrow{g}Y$, où $g$ est fini et $j$ une immersion ouverte. Puisque $X$ est réduit, on peut remplacer $X'$ par le sous-préschéma fermé réduit d'espace sous-jacent $\overline{j(X)}$ (I, 5.2.2), donc supposer $X'$ réduit et tel que toute composante irréductible de $X'$ domine $Y$. En outre, $j(X)$ est un ouvert dense dans $X'$ et n'a qu'un nombre fini de composantes irréductibles en vertu de l'hypothèse; donc ($0_{\mathrm I}$, 1.2.7) les composantes irréductibles de $X'$ sont les adhérences de celles de $j(X)$, et par suite la fibre $g^{-1}(\eta)$ s'identifie à $f^{-1}(\eta)$. Si on suppose le théorème démontré lorsqu'en outre $f$ est un morphisme fini, on pourra l'appliquer à $g$, et comme $f^{-1}(y)$ s'identifie à un sous-préschéma de $g^{-1}(y)$, la relation (i) pour $g$ entraîne la même relation pour $f$. Supposons d'autre part que $n(y)=n$; alors, en vertu de ce qui précède, on peut appliquer la conclusion de (ii) au morphisme $g$, et en remplaçant au besoin $Y$ par un voisinage ouvert de $y$, on peut supposer $g$ étale, et il en est par suite de même de $f$; en outre, on a alors $f^{-1}(y)=g^{-1}(y)$. Comme le morphisme $g$ est fermé et que $j(X)$ est ouvert dans $X'$, il existe un voisinage $U$ de $y$ dans $Y$ tel que $j(X)\cap g^{-1}(U)=j(f^{-1}(U))$, ce qui prouve que (ii) est vrai aussi pour $f$.

On est ainsi ramené au cas où $f$ est en outre un morphisme fini. Il est immédiat que, pour prouver (i), on peut se borner au cas où $Y=\Spec(\sh{O}_{Y,y})$. Montrons qu'il en est de même pour prouver (ii). En effet, supposons prouvé que la condition $n(y)=n$ entraîne que $X\times_Y\Spec(\sh{O}_{Y,y})$ est étale et fini sur $\Spec(\sh{O}_{Y,y})$; comme cela entraîne que $f$ est plat et formellement non ramifié aux points de $f^{-1}(y)$, et que $Y$ est intègre, on conclut de (18.4.13) que $f$ est étale aux points de $f^{-1}(y)$, donc aussi dans un voisinage ouvert $V$ de $f^{-1}(y)$ dans $X$; mais puisque $f$ est un morphisme fini, donc fermé, $V$ contient un ouvert de la forme $f^{-1}(U)$, où $U$ est un voisinage ouvert de $y$ dans $Y$.

\emph{C) Réduction au cas où $Y=\Spec(\sh{O}_{Y,y})$ et $\sh{O}_{Y,y}$ est strictement local.} --- Posons $A=\sh{O}_{Y,y}$, et soit $A'={}^{hs}\!A$ le hensélisé strict de $A$ (18.8.7), qui est un anneau local intègre et géométriquement unibranche (18.8.16). Soit $Y'=\Spec(A')$, et soient $y'$ le point fermé et $\eta'$ le point générique de $Y'$; $y'$ est au-dessus de $y$ et puisque le morphisme $g:Y'\to Y$ est plat (18.8.8), $\eta'$ est au-dessus de $\eta$ (2.3.4). Posons $X'=X\times_Y Y'$, $f'=f_{(Y')}:X'\to Y'$; $f'$ est fini et puisque $g$ est plat, toute composante irréductible de $X'$ domine $Y'$ (2.3.7).

Si $n'(y')$, $n'$ et $n'_s$ sont les nombres définis pour $f'$ de la même façon que $n(y)$, $n$ et $n_s$ pour $f$, on a $n'_s=n_s$ et $n'(y')=n(y)$ (I, 6.4.8), et il revient donc au même de prouver l'assertion (i) de l'énoncé pour $f$ ou pour $f'$. En outre, si $X$ est réduit et $y$ un point normal de $Y$ (donc $A$ intégralement clos), alors $A'$ est intégralement clos (18.8.12). D'autre part il résulte de la définition (18.8.7), de la remarque (8.1.2, $a$)) et du théorème de la double limite inductive, que $A'$ est limite inductive d'anneaux $A_\alpha$ tels que les morphismes $\Spec(A_\alpha)\to\Spec(A)=Y$ soient étales; par suite $X'$ est limite projective des $X_\alpha=X\otimes_A A_\alpha$, qui sont étales sur $X$, donc réduits puisque $X$ l'est (17.5.7); par passage à la limite,

\oldpage[IV-4]{165}

on en déduit que $X'$ est réduit (8.7.1). Cela étant, si $n(y)=n$, on a $n_s=n$ et par suite les corps résiduels aux points de $X_\eta$ sont nécessairement séparables sur $k(\eta)$; les corps résiduels aux points de $X'_{\eta'}$ sont par suite algébriques et séparables sur $k(\eta')$ (4.6.1), autrement dit $n'_s=n'=n$, et on a donc en vertu de ce qui précède $n'(y')=n'$. Le morphisme $f'$ a donc les mêmes propriétés que $f$, et si l'on prouve que $f'$ est étale, il en résultera que $f$ est étale par (17.7.1).

\emph{D) Fin de la démonstration.} --- Supposons donc $A$ strictement local. Puisque $A$ est hensélien et le morphisme $f$ fini, il résulte de (18.5.11) que si $x_j$ ($1\leq j\leq m$) sont les points distincts de $f^{-1}(y)$, $X$ est somme des sous-préschémas ouverts $\Spec(\sh{O}_{X,x_j})=X_j$, qui sont donc les composantes connexes de $X$. Par hypothèse, chacune d'elles domine $Y$, donc rencontre $f^{-1}(\eta)$, et par suite le nombre de points de $f^{-1}(\eta)$ est $\geq m$; \emph{a fortiori} on a $n_s\geq m$. Mais puisque $k(y)$ est séparablement clos et que les $k(x_j)$ sont algébriques sur $k(y)$, ce sont nécessairement des extensions radicielles de $k(y)$, donc $m=n(y)$, et cela prouve (i). Supposons maintenant vérifiées les hypothèses de (ii); les relations $n(y)=n_s=n$ entraînent d'abord, puisque toute composante irréductible de $X$ domine $Y$, que chacun des $X_j$ est irréductible; en outre, comme les $X_j$ sont réduits, ils sont intègres; enfin les relations $m=n_s=n$ entraînent que pour le point générique $\xi_j$ de $X_j$, $k(\xi_j)$ est une extension séparable de $k(\eta)$ de degré 1, donc isomorphe à $k(\eta)$; autrement dit, la restriction $f|X_j:X_j\to Y$ de $f$ est un morphisme fini et birationnel; en outre, $X_j$ est intègre et par hypothèse $Y$ est normal; donc (8.12.10.1) $f|X_j$ est un isomorphisme, ce qui prouve que $f$ est étale, et termine la démonstration de (18.10.16).

\begin{remark}[18.10.17]
\label{IV.18.10.17}
La méthode de «localisation étale» utilisée dans la démonstration de (18.10.16) permet aussi d'améliorer les résultats de (15.5.1) en éliminant l'hypothèse noethérienne. On a tout d'abord le résultat suivant qui généralise (15.5.2):
\end{remark}

\begin{env}[18.10.17.1]
\label{IV.18.10.17.1}
Soient $f:X\to Y$ un morphisme séparé et quasi-fini, $y$ un point de $Y$ tel que $f$ soit universellement ouvert aux points de $f^{-1}(y)$. Alors, pour toute générisation $y'$ de $y$, on a $n(y)\leq n(y')$.
\end{env}

On peut remplacer $Y$ par le sous-préschéma fermé réduit $Z$ ayant pour espace sous-jacent $\overline{\{y'\}}$ et $X$ par $f^{-1}(Z)$, donc supposer $Y$ intègre de point générique $y'$. Utilisant (18.12.13) comme dans (18.10.16, B)), on peut supposer que $X$ est ouvert dense dans un préschéma $X'$ et que $f$ est la restriction à $X$ d'un morphisme fini $g:X'\to Y$. Montrons que l'on a $g^{-1}(y')=f^{-1}(y')$. En effet, soit $i:g^{-1}(y')\to X'$ l'injection canonique, qui est un morphisme plat puisque $y'$ est point générique de $Y$ (I, 3.6.5); on peut écrire $f^{-1}(y')=i^{-1}(X)$. D'autre part, l'injection canonique $j:X\to X'$ est un morphisme de type fini puisque le composé $g\circ j=f$ est de type fini et que $g$ est séparé (1.5.4); donc $X$ est un ouvert rétrocompact dans $X'$, et par suite est pro-constructible dans $X'$ (1.9.5, (v)). On conclut donc de (2.3.10) que l'on a $i^{-1}(\overline{X})=\overline{i^{-1}(X)}$, autrement dit $f^{-1}(y')$ est dense dans $g^{-1}(y')$; mais comme $g^{-1}(y')$ est discret, on a nécessairement $f^{-1}(y')=g^{-1}(y')$.

On est ainsi ramené à prouver que la somme des nombres $[k(x):k(y)]_s$ où $x$ parcourt l'ensemble des points de $g^{-1}(y)$ où $g$ est universellement ouvert, est au plus égale

\oldpage[IV-4]{166}

au nombre $n(y')$ (pour le morphisme $g$). Utilisant ensuite le fait que la propriété d'être universellement ouvert en un point se conserve par changement de base, on montre, comme dans (18.10.16, C)), que l'on peut se ramener au cas où $Y=\Spec(\sh{O}_{Y,y})$, avec $\sh{O}_{Y,y}$ strictement local. Alors, avec les notations de (18.10.16, D)), $X$ est somme des sous-préschémas ouverts $X_j=\Spec(\sh{O}_{X,x_j})$, où $x_j$ parcourt $g^{-1}(y)$; l'hypothèse que $g$ est ouvert en un de ces points $x_j$ entraîne que le sous-préschéma $X_j$ correspondant domine $Y$ (I, 10.3); on termine comme dans (18.10.16, D)).

\emph{On tire ensuite de (18.10.17.1) que les conclusions de (15.5.1) sont encore valables lorsqu'on ne suppose plus $Y$ localement noethérien, mais seulement que $f$ est séparé, quasi-fini et de présentation finie.}

En effet, pour prouver l'assertion (i) de (15.5.1), on remarque que, puisque $f$ est de présentation finie, l'ensemble $E$ des points $y'\in Y$ tels que $n(y')\geq n(y)$ est localement constructible (9.7.9); pour montrer que $y$ est intérieur à $E$, il suffit, en vertu de (I, 10.1), de voir que toute générisation $y'$ de $y$ appartient à $E$, ce qui n'est autre que (18.10.17.1).

Pour prouver l'assertion (ii) de (15.5.1), on peut, puisque $f$ est de présentation finie, utiliser (8.10.5, (xii)) appliqué suivant la méthode de (8.1.2, a)), et on peut donc déjà supposer que $Y=\Spec(\sh{O}_{Y,y})$. Appliquant ensuite (2.7.1, (vii)), on voit qu'on peut par le changement de base fidèlement plat $Y'\to Y$, où $Y'=\Spec({}^{hs}\!\sh{O}_{Y,y})$, supposer que l'anneau $\sh{O}_{Y,y}$ est strictement local. Appliquant (18.5.11, c)), on voit alors que si $x_j$ ($1\leq j\leq n$) sont les points de $f^{-1}(y)$, $X$ est somme de $n$ sous-préschémas ouverts $X_j=\Spec(\sh{O}_{X,x_j})$ qui sont finis sur $Y$, et d'un préschéma ouvert $X''$. Si on prouve que $X''=\varnothing$, on aura démontré que $f$ est fini, donc propre. Or, les corps $k(x_j)$ étant des extensions algébriques de $k(y)$, sont radicielles, et par suite $n(y)=n$; comme $f$ est ouvert en chacun des points $x_j$, $X_j$ domine $Y$ (I, 10.3), donc la restriction de $f$ à $X_j$, étant un morphisme fini, est surjective (II, 6.1.10). L'hypothèse que $z\mapsto n(z)$ est constant dans $Y$ entraîne donc que $X''\cap f^{-1}(y')=\varnothing$ pour tout $y'\in Y$, c'est-à-dire $X''=\varnothing$.

Enfin, la preuve de (iii) reproduit sans changement celle donnée dans (15.5.1).

\begin{env}[18.10.17.2]
\label{IV.18.10.17.2}
Par des méthodes analogues utilisant la hensélisation stricte, on peut, dans la plupart des résultats des \textsection\textsection~14 et 15, se débarrasser des hypothèses noethériennes.
\end{env}

\begin{lemma}[18.10.18]
\label{IV.18.10.18}
Soient $X$, $Y$ deux préschémas, $f:X\to Y$ un morphisme birationnel\footnote{Nous entendons par là la notion définie dans (6.15.4), mais où l'on ne suppose pas $X$ et $Y$ réduits.}, $x$ un point de $X$. Pour que $f$ soit un isomorphisme local au point $x$, il faut et il suffit que $f$ soit étale au point $x$. Pour que $f$ soit une immersion ouverte, il faut et il suffit que $f$ soit étale et séparé.
\end{lemma}

Les conditions énoncées étant trivialement nécessaires, tout revient à voir qu'elles sont suffisantes. Pour la première assertion, la question est locale sur $X$ et $Y$, donc on peut supposer $f$ étale, $X$ et $Y$ affines, $f$ séparé, et on est ramené à prouver la seconde assertion. En vertu de (17.9.1), il s'agit de prouver que $f$ est radiciel. Soit donc $y\in Y$ et prouvons que $f^{-1}(y)$ est radiciel sur $k(y)$. Le changement de base $\Spec(\sh{O}_{Y,y})\to Y$ ne changeant pas le fait que $f$ est étale, séparé et birationnel, on peut se borner au cas où $Y=\Spec(A)$, avec $A=\sh{O}_{Y,y}$.

\oldpage[IV-4]{167}

Soit alors $A'={}^{hs}\!A$, qui est un anneau local et un $A$-module fidèlement plat (18.8.8), l'homomorphisme $A\to A'$ étant local. Posons $Y'=\Spec(A')$, $X'=X\times_Y Y'$, $f'=f_{(Y')}:X'\to Y'$, de sorte que $f'$ est étale et séparé; en outre, comme le morphisme $Y'\to Y$ est plat, le raisonnement de (6.15.4.1) montre que $f'$ est aussi birationnel. Si $y'$ est le point fermé de $Y'$, il suffira, en vertu de (2.6.1, (v)), de prouver que $f'^{-1}(y')$ est radiciel sur $k(y')$; puisque $f'$ est étale, il suffira de prouver que $f'^{-1}(y')$ contient au plus un point (17.6.1, c')).

Supposons donc $Y=\Spec(A)$, $A$ étant strictement local, $f$ séparé, étale et birationnel, et montrons que $f^{-1}(y)$ ne peut contenir plus d'un point. En effet, s'il existait dans $f^{-1}(y)$ deux points distincts $x_1$, $x_2$, il existerait deux $Y$-sections $u_1$, $u_2$ de $f$ telles que $u_1(y)=x_1$ et $u_2(y)=x_2$ (18.8.1); mais puisque $f$ est étale et séparé, $u_1$ et $u_2$ seraient des isomorphismes de $Y$ sur deux composantes connexes (ouvertes) de $X$ (17.9.4) et il y aurait donc deux points maximaux au moins de $X$ au-dessus de tout point maximal de $Y$, ce qui contredit l'hypothèse que $f$ est birationnel.

\begin{proposition}[18.10.19]
\label{IV.18.10.19}
Soient $Y$ un préschéma réduit dont l'ensemble des composantes irréductibles est localement fini, $f:X\to Y$ un morphisme séparé, localement de type fini et formellement non ramifié, $g$ une $Y$-section rationnelle de $f$ (c'est-à-dire une $Y$-application rationnelle de $Y$ dans $X$), $U$ le domaine de définition de $g$ (I, 7.2.1), et soit $Z$ l'adhérence de $g(U)$ dans $X$. Alors, pour tout point $y\in Y-U$ tel que $Y$ soit géométriquement unibranche au point $y$, on a $Z_y=Z\cap f^{-1}(y)=\varnothing$. En particulier, si $Y$ est géométriquement unibranche, $g(U)$ est une partie à la fois ouverte et fermée de $X$. Si, pour toute partie fermée irréductible $T$ de $X$, $f(T)$ est fermé dans $Y$, $g$ est définie en tout point géométriquement unibranche de $Y$.
\end{proposition}

Puisque $Y$ est réduit et $f$ séparé, il résulte de (I, 7.2.2) qu'il existe une $U$-section $u$ de $f^{-1}(U)$ appartenant à la classe $g$; d'ailleurs (17.4.1.2), puisque $f$ est formellement non ramifié, $u$ est un isomorphisme de $U$ sur le préschéma induit sur l'ouvert $u(U)$ de $X$. Si on note encore $Z$ le sous-préschéma réduit de $X$ ayant $Z$ pour espace sous-jacent, $Z|_U=u(U)$, étant réduit, est induit aussi par $Z$, donc la restriction $f'=f|_Z$ de $f$ est un morphisme birationnel de $Z$ dans $Y$, d'ailleurs formellement non ramifié, puisque $f$ l'est (17.1.3). La question étant locale sur $Y$, puisque $Y$ est réduit et géométriquement unibranche (donc intègre) au point $y$, on peut se borner (en remplaçant $Y$ par un voisinage ouvert de $y$) au cas où $Y$ est intègre; alors $u(U)$ est irréductible, donc aussi $Z$, et par suite $Z$ est aussi intègre. Alors (I, 8.2.7), puisque $f'$ est dominant, l'homomorphisme $\sh{O}_{Y,f(x)}\to\sh{O}_{Z,x}$ est injectif pour tout $x\in Z_y$, et on conclut de (18.10.2, (iii)) que $f'$ est un morphisme étale; étant séparé et birationnel, c'est un isomorphisme local au point $x$ en vertu de (18.10.18); mais cela entraînerait que la $U$-section $u$ serait prolongeable à un ouvert distinct de $U$, contrairement à la définition de ce dernier. On a donc nécessairement $Z_y=\varnothing$, ce qui établit la première assertion; la seconde en est une conséquence évidente. En outre, sous les hypothèses de la dernière assertion de l'énoncé, comme $Z$ est fermé et irréductible, $f(Z)$ est fermé dans $Y$, donc égal à $Y$, c'est-à-dire $Z_y\ne\varnothing$ pour tout $y\in Y$, ce qui achève la démonstration.

\begin{remark}[18.10.20]
\EGAReaderTarget{ega4.c04.rem.18-10-20}{IV.18.10.20}
Soit $Y$ un préschéma localement noethérien. On dit qu'un morphisme $f:X\to Y$ est essentiellement propre s'il est localement de type fini et si, pour tout $Y$-schéma de la forme $Y'=\Spec(A)$ où $A$ est un anneau de valuation discrète, l'application canonique (II, 7.3.2.2) relative au $Y'$-préschéma $X\times_Y Y'$ est bijective.

\oldpage[IV-4]{168}

Le raisonnement de (II, 7.3.8) prouve qu'il revient au même de dire que $f$ (supposé localement de type fini) est séparé et que, pour tout changement de base $Y''\to Y$, où $Y''$ est localement noethérien, l'image par $f_{(Y'')}:X_{(Y'')}\to Y''$ de toute partie fermée irréductible de $X''$ est fermée. Dire que $f$ est propre (pour $Y$ localement noethérien) signifie donc que $f$ est essentiellement propre et de type fini (II, 7.3.8); mais on rencontre des exemples importants de morphismes essentiellement propres qui ne sont pas de type fini (par exemple certains «préschémas de Picard» ou certains «préschémas de Néron-Severi» (chap. VI)). Il résulte évidemment de (18.10.19) que si $Y$ est localement noethérien, réduit et géométriquement unibranche, et si le morphisme $f:X\to Y$ est essentiellement propre et non ramifié, alors toute $Y$-section rationnelle de $f$ est partout définie.
\end{remark}

\begin{proposition}[18.10.21]
\label{IV.18.10.21}
Soient $f:X\to Y$ un morphisme localement de présentation finie, $x$ un point de $X$, $y=f(x)$ et supposons l'anneau $\sh{O}_{Y,y}$ intègre et géométriquement unibranche. Soit $r$ le nombre minimum de générateurs du $\sh{O}_{X,x}$-module $(\Omega^1_{X/Y})_x$ (égal au rang du $k(x)$-espace vectoriel $(\Omega^1_{X/Y})_x\otimes_{\sh{O}_{X,x}}k(x))$. Alors, pour que $f$ soit lisse au point $x$, il faut et il suffit que, si $y'$ est le point générique de l'unique composante irréductible de $Y$ contenant $y$, il existe une générisation $x'$ de $x$ telle que $f(x')=y'$ et $\dim_{x'}f^{-1}(y')\geq r$.
\end{proposition}

Si $f$ est lisse au point $x$, il l'est aussi dans un voisinage ouvert $U$ de $x$, donc en toute générisation $x'$ de $x$ et $\dim_{x'}(f^{-1}(f(x')))=r$ en chacun de ces points, en vertu de (17.10.2). Comme $f$ est en outre plat dans $U$, il existe une générisation de $x$ au-dessus de toute générisation de $y$ (2.3.4), ce qui prouve que la condition est nécessaire. Pour montrer qu'elle est suffisante, notons d'abord que par (17.5.1), il suffit de prouver que le morphisme déduit de $f$ par le changement de base $\Spec(\sh{O}_{Y,y})\to Y$ est lisse au point $x$; on peut donc supposer que $Y=\Spec(\sh{O}_{Y,y})$.

La question est locale sur $X$, donc ((16.4.22) et ($0_{\mathrm I}$, 5.2.2)) on peut supposer que $X$ est affine, et qu'il existe $r$ sections $t_i$ ($1\leq i\leq r$) de $\sh{O}_X$ au-dessus de $X$ telles que les sections $d_{X/Y}t_i$ engendrent $\Omega^1_{X/Y}$. Soit
\[
  g:X\to Z=Y[T_1,\ldots,T_r]=\mathbf V^r_Y
\]
le $Y$-morphisme correspondant à l'homomorphisme $\sh{O}_Y^r\to f_*(\sh{O}_X)$ de $\sh{O}_Y$-Modules défini par les $t_i$ (considérées comme sections de $f_*(\sh{O}_X)$ au-dessus de $Y$) (II, 1.2.7). On a la suite exacte (16.4.19)
\[
  g^*(\Omega^1_{Z/Y})\longrightarrow\Omega^1_{X/Y}\longrightarrow\Omega^1_{X/Z}\longrightarrow0.
\]
Comme les $dT_i$ ($1\leq i\leq r$) engendrent $\Omega^1_{Z/Y}$, il résulte de la définition de $g$ que l'homomorphisme $g^*(\Omega^1_{Z/Y})\to\Omega^1_{X/Y}$ est surjectif, donc $\Omega^1_{X/Z}=0$ et $g$ est par suite non ramifié (17.4.1). Nous allons voir que $g$ est étale au point $x$, et compte tenu de (17.3.8), cela prouvera que $f$ est lisse au point $x$. Comme $\sh{O}_{Y,y}$ est supposé intègre et géométriquement

\oldpage[IV-4]{169}

unibranche, il en est de même de $\sh{O}_{Z,z}$ où $z=g(x)$, puisque le morphisme structural $Z\to Y$ est lisse (17.3.8). En vertu de (18.10.1), il suffit de montrer que l'homomorphisme $\sh{O}_{Z,z}\to\sh{O}_{X,x}$ est injectif; il revient au même de voir que le morphisme $\Spec(\sh{O}_{X,x})\to\Spec(\sh{O}_{Z,z})$ correspondant est dominant, puisque $\sh{O}_{Z,z}$ est intègre (I, 1.2.7).

Soit $y'$ le point générique de $Y=\Spec(\sh{O}_{Y,y})$ et soit $z'$ le point générique de $Z=\Spec(\sh{O}_{Y,y}[T_1,\ldots,T_r])$ qui est intègre; si $h:Z\to Y$ est le morphisme structural, on a $h(z')=y'$. Il suffit de prouver qu'il existe dans $f^{-1}(y')$ une générisation $x'$ de $x$ telle que l'image de $x'$ par le morphisme $g_{y'}:f^{-1}(y')\to h^{-1}(y')$ soit égale à $z'$. D'ailleurs $z'$ est aussi le point générique de $h^{-1}(y')=\Spec(k(y')[T_1,\ldots,T_r])$; comme le morphisme $g_{y'}$ est non ramifié, donc quasi-fini (17.4.2), il en est de même de sa restriction à toute composante irréductible de $f^{-1}(y')$; comme $f^{-1}(y')$ et $h^{-1}(y')$ sont localement noethériens et que $\dim(h^{-1}(y'))=r$, il résulte de (5.4.1) que la restriction de $g_{y'}$ à toute composante irréductible de dimension $\geq r$ de $f^{-1}(y')$ est un morphisme dominant. Or, il existe par hypothèse une telle composante contenant une générisation $x'$ de $x$; \emph{a fortiori} son point générique est une générisation de $x$, d'où la conclusion.

\subsection{Application aux algèbres locales noethériennes complètes sur un corps.}
\EGAReaderTarget{ega4.c04.subsec.18-11}{subsection:IV.18.11}

Le lemme suivant généralise (0, 21.9.1) et (0, 21.9.2):

\begin{lemma}[18.11.1]
\label{IV.18.11.1}
Soient $k$ un corps; $A$ un anneau local noethérien complet qui est une $k$-algèbre, $K$ le corps résiduel de $A$.
\begin{enumerate}
  \item[(i)] Si le $K$-espace vectoriel $\Omega^1_{K/k}$ est de rang fini, le $A$-module $\hat{\Omega}^1_{A/k}$ est de type fini.
  \item[(ii)] Si de plus $A$ est une $k$-algèbre formellement lisse (pour la topologie adique), $\hat{\Omega}^1_{A/k}$ est un $A$-module libre de rang égal à
\[
    \dim(A)+\operatorname{rg}_K(\Omega^1_{K/k})-\operatorname{rg}_K(\Upsilon_{K/k}).
\]
\end{enumerate}
En outre, pour tout sous-corps $k_0$ de $k$ tel que $\Omega^1_{k/k_0}$ soit un $k$-espace vectoriel de rang fini, $\hat{\Omega}^1_{A/k_0}$ est un $A$-module libre de rang égal à
\[
  \dim(A)+\operatorname{rg}_K(\Omega^1_{K/k_0})-\operatorname{rg}_K(\Upsilon_{K/k_0})+\operatorname{rg}_k(\Omega^1_{k/k_0}).
\]
\end{lemma}

L'assertion (i) n'est mise que pour mémoire, ayant été prouvée dans (0, 20.7.15). Pour prouver (ii), notons que cette assertion a été en fait établie par la démonstration de (0, 21.9.2), l'énoncé de (0, 21.9.2) faisant seul intervenir la finitude du degré de transcendance de $K$ sur $k$ (par l'intermédiaire de l'égalité de Cartier (0, 21.7.1)).

\begin{lemma}[18.11.2]
\label{IV.18.11.2}
Soient $k$ un corps de caractéristique $p>0$, $C$ une $k$-algèbre locale noethérienne, complète, dont le corps résiduel est une extension finie de $k$, et qui est intègre; soient $n$ sa dimension, $L$ son corps des fractions. Alors $\Omega^1_{L/k}$ et $\Upsilon_{L/k}$ sont des $L$-espaces vectoriels de rang fini, et l'on a
\begin{equation}
\label{IV.18.11.2.1}
\tag{18.11.2.1}
  \operatorname{rg}_L(\Omega^1_{L/k})-\operatorname{rg}_L(\Upsilon_{L/k})\geq n.
\end{equation}

\oldpage[IV-4]{170}

Supposons en outre que $[k:k^p]<+\infty$. Alors on a l'égalité
\begin{equation}
\label{IV.18.11.2.2}
\tag{18.11.2.2}
  \operatorname{rg}_L(\Omega^1_{L/k})-\operatorname{rg}_L(\Upsilon_{L/k})=n.
\end{equation}
En outre, $\Omega^1_{C/k}$ est alors isomorphe à $\hat{\Omega}^1_{C/k}$, donc (18.11.1) est un $C$-module de type fini.
\end{lemma}

On sait en effet (0, 19.8.9) qu'il existe une sous-$k$-algèbre $C_0$ de $C$ isomorphe à une algèbre de séries formelles $k[[T_1,\ldots,T_n]]$ et telle que $C$ soit une $C_0$-algèbre finie. Par suite, $L$ est une extension finie du corps des fractions $L_0=k((T_1,\ldots,T_n))$ de $C_0$. Or, on a la suite exacte de $L$-espaces vectoriels obtenue en appliquant (0, 20.6.17.1) au sous-corps premier de $k$ et aux trois corps $k\subset L_0\subset L$ (et en tenant compte de (0, 20.6.21, (i))
\[
  0\to\Upsilon_{L_0/k}\otimes_{L_0}L\to\Upsilon_{L/k}\to\Upsilon_{L/L_0}
  \to\Omega^1_{L_0/k}\otimes_{L_0}L\to\Omega^1_{L/k}\to\Omega^1_{L/L_0}\to0.
\]
Puisque $L$ est une extension finie de $L_0$, $\Omega^1_{L/L_0}$ et $\Upsilon_{L/L_0}$ sont des $L$-espaces vectoriels de rang fini ayant même rang, en vertu de l'égalité de Cartier (0, 21.7.1). Comme $L_0$ est séparable sur $k$ (0, 21.9.6.4), on a $\Upsilon_{L_0/k}=0$ (0, 20.6.3); on déduit donc déjà de la suite exacte précédente que $\Upsilon_{L/k}$ est de rang fini et que l'on a dans tous les cas ($0_{\mathrm{III}}$, 11.10.2)
\[
  \operatorname{rg}_L(\Omega^1_{L/k})-\operatorname{rg}_L(\Upsilon_{L/k})
  =\operatorname{rg}_{L_0}(\Omega^1_{L_0/k})-\operatorname{rg}_{L_0}(\Upsilon_{L_0/k}).
\]

Pour prouver (18.11.2.1) ou (18.11.2.2), on peut donc se borner à prouver ces relations en remplaçant $C$ et $L$ par $C_0$ et $L_0$. Comme $L_0$ est séparable sur $k$ (0, 21.9.6.4), on a $\Upsilon_{L_0/k}=0$ (0, 20.6.3); d'autre part $\Omega^1_{L_0/k}=\Omega^1_{C_0/k}\otimes_{C_0}L_0$ (0, 20.5.9). On sait (0, 21.9.4) que si $C_1=k[[T_1^p,\ldots,T_n^p]]$, $\hat{\Omega}^1_{C_0/k}$ s'identifie à $\Omega^1_{C_0/C_1}$; d'autre part, on a $\Omega^1_{C_0/k}=\Omega^1_{C_0/k[C_0^p]}$ (0, 21.1.5), et comme $C_0^p=k^p[[T_1^p,\ldots,T_n^p]]$, on a $k[C_0^p]\subset C_1$, avec égalité lorsque $[k:k^p]<+\infty$. Or, on a une suite exacte $\Omega^1_{C_0/k[C_0^p]}\to\Omega^1_{C_0/C_1}\to0$ (0, 20.5.7), d'où une suite exacte $\Omega^1_{L_0/k}\to\hat{\Omega}^1_{C_0/k}\otimes_{C_0}L_0\to0$. Comme $\hat{\Omega}^1_{C_0/k}$ est un $C_0$-module libre de rang $n$ (0, 21.9.3), on voit que l'on a dans tous les cas $\operatorname{rg}_{L_0}(\Omega^1_{L_0/k})\geq n$, avec égalité lorsque $[k:k^p]<+\infty$; ceci prouve déjà (18.11.2.1) et (18.11.2.2).

Enfin, pour voir que $\Omega^1_{C/k}$ est isomorphe à $\hat{\Omega}^1_{C/k}$ lorsque $[k:k^p]<+\infty$, il suffit de prouver que $\Omega^1_{C/k}$ est un $C$-module de type fini, puisque $C$ est un anneau local noethérien complet ($0_{\mathrm I}$, 7.3.3). Puisque $\Omega^1_{C/k}$ est isomorphe à $\Omega^1_{C/k[C^p]}$ et que $k[C_0^p]\subset k[C^p]$, tout revient à prouver que $C$ est un $k[C_0^p]$-module de type fini (0, 20.4.7); mais cela résulte de ce que $C$ est un $C_0$-module de type fini et $C_0$ un $k[C_0^p]$-module de type fini en vertu de l'hypothèse $[k:k^p]<+\infty$.

\begin{lemma}[18.11.3]
\label{IV.18.11.3}
Soient $k$ un corps, $p$ son exposant caractéristique, et supposons que $[k:k^p]<+\infty$. Soit $A$ une $k$-algèbre locale noethérienne complète, dont le corps résiduel est une extension finie de $k$. Soit $\mathfrak p$ un idéal premier de $A$ tel que $A_{\mathfrak p}$ soit géométriquement régulier sur $k$ (6.7.6), de sorte qu'il existe un seul idéal premier minimal $\mathfrak q$ de $A$ contenu dans $\mathfrak p$. Alors $(\hat{\Omega}^1_{A/k})_{\mathfrak p}$ est un $A_{\mathfrak p}$-module libre de rang égal à $\dim(A/\mathfrak q)$.
\end{lemma}

Comme $A$ est noethérien, et $\Spec(A)$ régulier (et \,\emph{a fortiori} intègre) au point $\mathfrak p$, $\mathfrak p$ n'appartient qu'à une seule composante irréductible de $\Spec(A)$, donc ne contient

\oldpage[IV-4]{171}

qu'un seul idéal minimal $\mathfrak q$ de $A$, et en outre on a $\mathfrak q_{\mathfrak p}=0$. Si l'on pose $B=A/\mathfrak q$, la suite (0, 20.7.20)
\[
  \mathfrak q/\mathfrak q^2\longrightarrow
  \hat{\Omega}^1_{A/k}\widehat{\otimes}_A B
  \longrightarrow\hat{\Omega}^1_{B/k}\longrightarrow0
\]
est exacte; en effet, $\hat{\Omega}^1_{A/k}$ est un $A$-module de type fini (18.11.1), donc on a $\hat{\Omega}^1_{A/k}\widehat{\otimes}_A B=\hat{\Omega}^1_{A/k}\otimes_A B=\hat{\Omega}^1_{A/k}/\mathfrak q\hat{\Omega}^1_{A/k}$, et comme ce $B$-module est de type fini, tout sous-$B$-module en est fermé ($0_{\mathrm I}$, 7.3.5), compte tenu de (0, 20.4.5). Comme l'image de $\mathfrak q/\mathfrak q^2$ est dense dans le noyau de l'homomorphisme $\hat{\Omega}^1_{A/k}\widehat{\otimes}_A B\to\hat{\Omega}^1_{B/k}$ (0, 20.7.20), elle est nécessairement égale à ce noyau, d'où notre assertion. En localisant en $\mathfrak p$ la suite exacte précédente, on voit en outre que l'homomorphisme canonique $(\hat{\Omega}^1_{A/k})_{\mathfrak p}\to(\hat{\Omega}^1_{B/k})_{\mathfrak p/\mathfrak q}$ est bijectif. On voit ainsi qu'on peut se borner au cas où $\mathfrak q=0$, autrement dit au cas où $A$ est intègre. Distinguons alors deux cas:

\emph{I) $p>1$.} Alors on peut appliquer (18.11.2) à l'algèbre quotient $C=A/\mathfrak p$, dont le corps des fractions $K$ n'est autre que le corps résiduel de $A_{\mathfrak p}$; on voit donc que l'on a
\begin{equation}
\label{IV.18.11.3.1}
\tag{18.11.3.1}
  \operatorname{rg}_K(\Omega^1_{K/k})-\operatorname{rg}_K(\Upsilon_{K/k})=\dim(A/\mathfrak p).
\end{equation}

Notons maintenant (0, 19.6.6) que $A_{\mathfrak p}$ est une $k$-algèbre formellement lisse pour sa topologie $\mathfrak pA_{\mathfrak p}$-préadique; par suite $\Omega^1_{A_{\mathfrak p}/k}$ est formellement projectif (0, 20.4.9) pour la topologie $\mathfrak p$-préadique (0, 20.4.5); d'autre part, $\Omega^1_{A_{\mathfrak p}/k}=(\Omega^1_{A/k})_{\mathfrak p}$ (0, 20.5.9) est un $A_{\mathfrak p}$-module de type fini, en vertu de (18.11.2) appliqué à $A$. Pour tout entier $j$, $\Omega^1_{A_{\mathfrak p}/k}/\mathfrak p^{j+1}\Omega^1_{A_{\mathfrak p}/k}$ est donc un $(A_{\mathfrak p}/\mathfrak p^{j+1}A_{\mathfrak p})$-module projectif de rang $m=\operatorname{rg}_K(\Omega^1_{A_{\mathfrak p}/k}\otimes_{A_{\mathfrak p}}K)$ (0, 19.2.4); on conclut donc de ($0_{\mathrm{III}}$, 10.2.1 et 10.1.3) que $\hat{\Omega}^1_{A_{\mathfrak p}/k}$ est un $A_{\mathfrak p}$-module libre de rang $m$. Soit $A'=(A_{\mathfrak p})^\wedge$ l'algèbre complétée de $A_{\mathfrak p}$ pour sa topologie $\mathfrak pA_{\mathfrak p}$-préadique; $A'$ est encore une $k$-algèbre formellement lisse (0, 19.3.6) pour sa topologie adique, et il résulte de (0, 20.7.14) et de (0, 20.4.5) que $\hat{\Omega}^1_{A'/k}=\hat{\Omega}^1_{A_{\mathfrak p}/k}$; on en conclut que $\hat{\Omega}^1_{A'/k}$ est un $A'$-module libre de rang $m$. Mais il résulte alors de (18.11.1) et de ce que $\dim(A')=\dim(A_{\mathfrak p})$ (0, 16.2.4) que l'on a
\begin{equation}
\label{IV.18.11.3.2}
\tag{18.11.3.2}
  m=\dim(A_{\mathfrak p})+
  \bigl(\operatorname{rg}_K(\Omega^1_{K/k})-\operatorname{rg}_K(\Upsilon_{K/k})\bigr),
\end{equation}
d'où, en vertu de (18.11.3.1)
\[
  m=\dim(A/\mathfrak p)+\dim(A_{\mathfrak p}).
\]
Mais puisque $A$ est un anneau local noethérien complet, c'est un quotient d'anneau régulier (0, 19.8.8), donc (0, 16.5.12) on a $\dim(A)=\dim(A/\mathfrak p)+\dim(A_{\mathfrak p})$, ce qui achève la démonstration dans ce cas.

\emph{II) $p=1$.} On a, comme ci-dessus $\dim(A)=\dim(A/\mathfrak p)+\dim(A_{\mathfrak p})$; posons $n=\dim(A)$, $r=\dim(A/\mathfrak p)$, $s=\dim(A_{\mathfrak p})$; nous allons voir que $(\hat{\Omega}^1_{A/k})_{\mathfrak p}$ est un $A_{\mathfrak p}$-module libre de rang $n$. Nous allons d'abord prouver le lemme suivant:

\begin{lemma}[18.11.3.3]
\label{IV.18.11.3.3}
Soient $A$ un anneau local noethérien, $\mathfrak p$ un idéal premier de $A$ tel que $A_{\mathfrak p}$ soit un anneau de Cohen-Macaulay. Pour tout système de paramètres $(t_i)_{1\leq i\leq s}$ de $A_{\mathfrak p}$, il existe une suite $(x_i)_{1\leq i\leq s}$ d'éléments de $\mathfrak p$ tels que les $x_i$ fassent partie d'un système de paramètres de $A$

\oldpage[IV-4]{172}

et que, pour tout $i$, $x_i/1$ soit congru à $t_i$ modulo l'idéal $t_1A_{\mathfrak p}+\cdots+t_{i-1}A_{\mathfrak p}$. En particulier, si $A_{\mathfrak p}$ est régulier et si les $t_i$ forment un système régulier de paramètres de $A_{\mathfrak p}$, il en est de même des $x_i/1$.
\end{lemma}

Ce lemme sera lui-même une conséquence du suivant:

\begin{lemma}[18.11.3.4]
\label{IV.18.11.3.4}
Soient $A$ un anneau local noethérien, $\mathfrak p$ un idéal premier de $A$, $x$ un élément de $A$; alors il existe un élément $x'\in A$ tel que $x'/1=x/1$ dans $A_{\mathfrak p}$ et que $x'$ n'appartienne à aucun des idéaux premiers minimaux de $A$ non contenus dans $\mathfrak p$.
\end{lemma}

Montrons d'abord comment (18.11.3.4) entraîne (18.11.3.3). En multipliant au besoin les $t_i$ par des éléments inversibles de $A_{\mathfrak p}$, on peut supposer qu'ils sont de la forme $x_i/1$, avec $x_i\in\mathfrak p$ pour $1\leq i\leq s$. Raisonnons par récurrence sur $s$; comme $t_1$ fait partie d'un système de paramètres de $A_{\mathfrak p}$, il n'appartient à aucun des idéaux premiers minimaux de $A_{\mathfrak p}$ (0, 16.5.5), donc aucun $x'_1\in A$ tel que $x'_1/1=t_1$ ne peut appartenir à un idéal premier minimal de $A$ contenu dans $\mathfrak p$. En vertu de (18.11.3.4), on peut en outre choisir $x'_1$ tel que $x'_1/1=t_1$ (donc $x'_1\in\mathfrak p$) et que $x'_1$ n'appartienne à aucun idéal premier minimal de $A$, donc fasse partie d'un système de paramètres de $A$ (0, 16.3.4 et 16.3.7). On raisonne ensuite par récurrence, en considérant l'anneau $A'=A/x'_1A$ et l'idéal premier $\mathfrak p'=\mathfrak p/x'_1A$ de cet anneau; comme $A'_{\mathfrak p'}=A_{\mathfrak p}/t_1A_{\mathfrak p}$, $A'_{\mathfrak p'}$ est aussi un anneau de Cohen-Macaulay (0, 16.5.5), et les images $t'_i$ ($2\leq i\leq s$) des $t_i$ dans $A'_{\mathfrak p'}$ forment un système de paramètres de cet anneau; il suffit d'appliquer à $A'$ et aux $t'_i$ ($2\leq i\leq s$) l'hypothèse de récurrence. La dernière assertion de (18.11.3.3) résulte de ce que dans $A_{\mathfrak p}$, un système de paramètres qui engendre l'idéal maximal est un système régulier de paramètres (0, 17.1.1).

Prouvons donc (18.11.3.4). Soit $(\mathfrak p'_k)_{1\leq k\leq r}$ la suite des idéaux premiers minimaux de $A$ non contenus dans $\mathfrak p$ et tels que $x\in\mathfrak p'_k$, et soit $(\mathfrak p_j)_{1\leq j\leq n}$ la suite des idéaux premiers minimaux de $A$ autres que les $\mathfrak p'_k$; on peut supposer que $r\geq1$. Comme $\mathfrak p'_k$ ne contient pas $\bigcap_{1\leq j\leq n}\mathfrak p_j$, il existe un élément $y\in\bigcap_{1\leq j\leq n}\mathfrak p_j$ qui n'est contenu dans aucun des $\mathfrak p'_k$ (Bourbaki, \emph{Alg. comm.}, chap. II, \textsection~1, n\textsuperscript{o}~1, prop. 2); d'ailleurs $y/1$ appartient à tous les idéaux premiers minimaux de $A_{\mathfrak p}$, donc est nilpotent; si $(y/1)^h=0$, l'élément $x'=x+y^h$ répondra aux conditions de l'énoncé, car d'une part $y^h\notin\mathfrak p'_k$ pour $1\leq k\leq r$, et par définition des $\mathfrak p'_k$ on a bien $x'\notin\mathfrak p'_k$ pour $1\leq k\leq r$, et d'autre part, si $\mathfrak p_j$ est un idéal premier minimal de $A$ non contenu dans $\mathfrak p$ mais tel que $x\notin\mathfrak p_j$, on a aussi $x'\notin\mathfrak p_j$ puisque $y\in\mathfrak p_j$.

Revenons à la preuve de (18.11.3) lorsque $p=1$. En vertu de (18.11.3.3), il existe un système régulier de paramètres $(t_i)_{1\leq i\leq s}$ de $A_{\mathfrak p}$ tel que $t_i=x_i/1$, où les $x_i$ ($1\leq i\leq s$) appartiennent à $\mathfrak p$ et font partie d'un système de paramètres $(x_j)_{1\leq j\leq n}$ de $A$. Posons $A_0=k[[T_1,\ldots,T_n]]$; comme $A$ est une $k$-algèbre complète et que les $x_j$ appartiennent à l'idéal maximal $\mathfrak m$ de $A$, il existe un $k$-homomorphisme local $u:A_0\to A$ tel que $u(T_j)=x_j$ pour $1\leq j\leq n$ (Bourbaki, \emph{Alg. comm.}, chap. III, \textsection~4, n\textsuperscript{o}~5, prop. 6); si $\mathfrak n$ est l'idéal de $A$ engendré par les $x_j$ ($1\leq j\leq n$), c'est par hypothèse un idéal de définition de $A$; on déduit donc de ($0_{\mathrm I}$, 7.4.4 et 7.4.3) que $u$ fait de $A$ une $A_0$-algèbre finie.

Posons $\mathfrak p_0=\sum_{j=1}^s A_0T_j$, $B_0=(A_0)_{\mathfrak p_0}$ et $B=A\otimes_{A_0}B_0$, de sorte que $u_{\mathfrak p_0}:B_0\to B$ fait

\oldpage[IV-4]{173}

de $B$ une $B_0$-algèbre finie; en outre, si $\mathfrak p'$ est l'idéal de $B$ engendré par $\mathfrak p$, on a $B_{\mathfrak p'}=A_{\mathfrak p}$; comme $\mathfrak p'$ contient $\mathfrak p_0B$ par construction, il est au-dessus de l'idéal maximal $\mathfrak p_0B_0$ de $B_0$. Montrons que le morphisme $\Spec(B)\to\Spec(B_0)$ est non ramifié au point $\mathfrak p'$: cela résulte en effet de ce que $k(\mathfrak p')$ est une extension finie du corps $k(\mathfrak p_0)$ de caractéristique 0, donc est nécessairement séparable, et que l'on a $B_{\mathfrak p'}/\mathfrak p_0B_{\mathfrak p'}=k(\mathfrak p')$ en vertu du choix des $x_j$ pour $1\leq j\leq s$ (17.4.1).

Notons maintenant le lemme suivant:

\begin{lemma}[18.11.3.5]
\label{IV.18.11.3.5}
Soient $k$ un corps, $R$, $S$ deux $k$-algèbres locales noethériennes complètes, telles que, si $K$ est le corps résiduel de $S$, $\Omega^1_{K/k}$ soit de rang fini sur $K$; soit $u:R\to S$ un $k$-homomorphisme faisant de $S$ une $R$-algèbre finie; alors on a $\hat{\Omega}^1_{S/R}=\Omega^1_{S/R}$, et la suite
\begin{equation}
\label{IV.18.11.3.6}
\tag{18.11.3.6}
  \hat{\Omega}^1_{R/k}\otimes_R S\xrightarrow{v}\hat{\Omega}^1_{S/k}\xrightarrow{w}\Omega^1_{S/R}\to0
\end{equation}
(cf. (0, 20.7.17.3)) est exacte.
\end{lemma}

La première assertion résulte de ce que $\Omega^1_{S/R}$ est un $S$-module de type fini (0, 20.4.7) et de ($0_{\mathrm I}$, 7.3.6). D'autre part, on sait (0, 20.7.17.3) que l'image de $v$ est dense dans $\operatorname{Ker}(w)$ et que $w$ est surjectif; mais il résulte de (18.11.1) que $\hat{\Omega}^1_{S/k}$ est un $S$-module de type fini; on en conclut que tout sous-$S$-module de $\hat{\Omega}^1_{S/k}$ est fermé ($0_{\mathrm I}$, 7.3.5), d'où le lemme.

Appliquons ce lemme au cas où $R=A_0$, $S=A$, et notons que $\hat{\Omega}^1_{A_0/k}$ est un $A_0$-module libre de rang $n$ (0, 21.9.3); donc on a $\hat{\Omega}^1_{A_0/k}\widehat{\otimes}_{A_0}A=\hat{\Omega}^1_{A_0/k}\otimes_{A_0}A$ et ce $A$-module est libre de rang $n$. Puisque $\Spec(B)$ est non ramifié sur $\Spec(B_0)$ au point $\mathfrak p'$, on a $(\Omega^1_{A/A_0})_{\mathfrak p}=\Omega^1_{B_{\mathfrak p'}/B_0}=0$ (16.4.15 et 17.4.1). Si on localise la suite exacte (18.11.3.6) (appliquée à $A_0$ et $A$) en $\mathfrak p$, on obtient donc un homomorphisme surjectif
\[
  (\hat{\Omega}^1_{A_0/k}\otimes_{A_0}A)_{\mathfrak p}\longrightarrow(\hat{\Omega}^1_{A/k})_{\mathfrak p},
\]
d'où l'on conclut que le $A_{\mathfrak p}$-module $(\hat{\Omega}^1_{A/k})_{\mathfrak p}$ admet un système de $n$ générateurs. Mais le $A$-module $\hat{\Omega}^1_{A/k}$ est de rang $n$, en vertu de (0, 21.9.5), qui est applicable à l'anneau complet et intègre $A$ en raison du théorème de Cohen (0, 19.8.8, (ii)) et du fait que le corps des fractions de $A$ est de caractéristique 0. Le $A_{\mathfrak p}$-module $(\hat{\Omega}^1_{A/k})_{\mathfrak p}$ est donc aussi de rang $n$, et comme son quotient par son sous-module de torsion admet un système de $n$ générateurs, ce quotient est nécessairement libre; on en déduit aussitôt que les $n$ générateurs de $(\hat{\Omega}^1_{A/k})_{\mathfrak p}$ obtenus ci-dessus forment un système libre, d'où la conclusion.

\begin{lemma}[18.11.4]
\label{IV.18.11.4}
Soient $k$ un corps, $A$ une $k$-algèbre locale noethérienne complète, dont le corps résiduel est une extension finie de $k$. Soit $k'$ une extension quelconque de $k$, et posons $A'=A\widehat{\otimes}_k k'$. Alors:
\begin{enumerate}
  \item[(i)] $A'$ est un anneau semi-local noethérien complet, composé direct d'anneaux locaux complets $A'_i$ ($1\leq i\leq r$) qui sont des $A$-modules fidèlement plats, et dont les corps résiduels sont des extensions finies de $k'$; si $\mathfrak m$ est l'idéal maximal de $A$, $\mathfrak mA'$ est un idéal de définition de $A'$; on a $\dim(A'_i)=\dim(A)$ pour tout $i$.
  \item[(ii)] Pour tout $i$, $\hat{\Omega}^1_{A'_i/k'}$ est canoniquement isomorphe à $\hat{\Omega}^1_{A/k}\otimes_A A'_i$.
\end{enumerate}
\end{lemma}

\oldpage[IV-4]{174}

(i) Les premières assertions résultent aussitôt de (7.5.5) et de ($0_{\mathrm I}$, 6.6.2); le fait que $\mathfrak mA'$ soit un idéal de définition de $A'$ résulte aussi de (7.5.5), car si $K$ est le corps résiduel de $A$, $K\otimes_k k'$ est une $k'$-algèbre finie. Enfin $A'_i/\mathfrak mA'_i$ est un des anneaux locaux artiniens composants directs de $K\otimes_k k'$ (7.5.5), donc est de dimension 0; comme $A'_i$ est un $A$-module plat, l'égalité des dimensions de $A$ et de $A'_i$ résulte de (6.1.2).

(ii) Comme $\hat{\Omega}^1_{A/k}$ est, en vertu de l'hypothèse sur $K$, un $A$-module de type fini (18.11.1), $\hat{\Omega}^1_{A/k}\otimes_A A'$ est complet et s'identifie au produit tensoriel complété $\hat{\Omega}^1_{A/k}\widehat{\otimes}_A A'$; en vertu de l'associativité du produit tensoriel complété ($0_{\mathrm I}$, 7.7.4), $\hat{\Omega}^1_{A/k}\widehat{\otimes}_A A'=\hat{\Omega}^1_{A/k}\widehat{\otimes}_A(A\widehat{\otimes}_k k')$ s'identifie à $\hat{\Omega}^1_{A/k}\widehat{\otimes}_k k'$. Mais $\Omega^1_{A/k}\otimes_k k'$ s'identifie à $\Omega^1_{A''/k'}$, où $A''=A\otimes_k k'$ (0, 20.5.5), et comme $A'$ est par définition le séparé complété de $A''$, le séparé complété de $\Omega^1_{A''/k'}$ s'identifie par construction à celui de $\Omega^1_{A'/k'}$ (0, 20.7.4); autrement dit, on a un isomorphisme canonique
\[
  \hat{\Omega}^1_{A/k}\widehat{\otimes}_k k'\xrightarrow{\sim}\hat{\Omega}^1_{A'/k'}.
\]
La conclusion de (ii) résulte maintenant de ce que $\Omega^1_{A'/k'}$ est somme directe des $\Omega^1_{A'_i/k'}$ (0, 20.4.13), et si $\mathfrak r$ est le radical de $A'$, la topologie $\mathfrak r$-préadique sur $\Omega^1_{A'/k'}$ s'identifie au produit des topologies $\mathfrak m'_i$-préadiques sur les $\Omega^1_{A'_i/k'}$ (où $\mathfrak m'_i$ est l'idéal maximal de $A'_i$); finalement, il en résulte que le séparé complété $\hat{\Omega}^1_{A'/k'}$ pour la topologie $\mathfrak r$-préadique s'identifie au produit des séparés complétés $\hat{\Omega}^1_{A'_i/k'}$ pour les topologies $\mathfrak m'_i$-préadiques, et il suffit d'utiliser (0, 20.4.5).

\begin{proposition}[18.11.5]
\label{IV.18.11.5}
Soient $k$ un corps, $A$ une $k$-algèbre locale noethérienne et complète dont le corps résiduel est une extension finie de $k$, $\mathfrak p$ un idéal premier de $A$, distinct de l'idéal maximal $\mathfrak m$, tel qu'il existe un idéal premier minimal $\mathfrak q\subset\mathfrak p$ pour lequel $\dim(A/\mathfrak q)=\dim(A)$ (ce qui aura lieu en particulier lorsque $A$ est équidimensionnel), $m$ un entier $\geq0$. Les conditions suivantes sont équivalentes:
\begin{enumerate}
  \item[$a)$] Le $A_{\mathfrak p}$-module $(\hat{\Omega}^1_{A/k})_{\mathfrak p}$ admet un système de $m$ générateurs.
  \item[$b)$] Il existe un $k$-homomorphisme local $u:B\to A$, où $B=k[[T_1,\ldots,T_m]]$, faisant de $A$ une $B$-algèbre finie, et tel que le morphisme correspondant $\Spec(A)\to\Spec(B)$ soit non ramifié au point $\mathfrak p$.
\end{enumerate}
\end{proposition}

Pour prouver que $b)$ entraîne $a)$, notons qu'en vertu du lemme (18.11.3.5), on a alors la suite exacte
\[
  \hat{\Omega}^1_{B/k}\otimes_B A\longrightarrow\hat{\Omega}^1_{A/k}\longrightarrow\Omega^1_{A/B}\longrightarrow0
\]
puisque $A$ est une $B$-algèbre finie; localisant en $\mathfrak p$ et notant que par hypothèse on a alors $(\Omega^1_{A/B})_{\mathfrak p}=0$ (17.4.1), on obtient un homomorphisme surjectif $(\hat{\Omega}^1_{B/k}\otimes_B A)_{\mathfrak p}\to(\hat{\Omega}^1_{A/k})_{\mathfrak p}$, et la conclusion résulte de ce que $\hat{\Omega}^1_{B/k}$ est un $B$-module libre de rang $m$ (0, 21.9.3).

Pour prouver que $a)$ implique $b)$, prouvons d'abord les lemmes suivants.

\begin{lemma}[18.11.5.1]
\label{IV.18.11.5.1}
Soient $k$ un corps, $A$ une $k$-algèbre locale noethérienne et complète, dont le corps résiduel est une extension finie de $k$, $\mathfrak p$ un idéal premier de $A$, $\mathfrak q$ un idéal minimal de $A$ contenu dans $\mathfrak p$. Alors le nombre minimum de générateurs du $A_{\mathfrak p}$-module $(\hat{\Omega}^1_{A/k})_{\mathfrak p}$ est au moins égal à $\dim(A/\mathfrak q)$.
\end{lemma}

\oldpage[IV-4]{175}

Posons $n=\dim(A/\mathfrak q)$, et soit $m$ le nombre minimum de générateurs du $A_{\mathfrak p}$-module $(\hat{\Omega}^1_{A/k})_{\mathfrak p}$, qui est égal à $\operatorname{rg}_{k(\mathfrak p)}((\hat{\Omega}^1_{A/k})_{\mathfrak p}\otimes_{A_{\mathfrak p}}k(\mathfrak p))$ (Bourbaki, \emph{Alg. comm.}, chap. II, \textsection~3, n\textsuperscript{o}~2, cor. 2 de la prop. 4). Notons que l'on a $(\hat{\Omega}^1_{A/k})_{\mathfrak q}=(\hat{\Omega}^1_{A/k})_{\mathfrak p}\otimes_{A_{\mathfrak p}}A_{\mathfrak q}$, donc le nombre minimum de générateurs du $A_{\mathfrak q}$-module $(\hat{\Omega}^1_{A/k})_{\mathfrak q}$ est au plus égal à $m$. Il suffit par suite de considérer le cas où $\mathfrak p=\mathfrak q$ est minimal. En second lieu, montrons qu'on peut se borner au cas où $k$ est algébriquement clos. En effet, soit $k'$ une clôture algébrique de $k$, et posons $A'=A\widehat{\otimes}_k k'$, qui est composé direct de $k'$-algèbres locales $A'_i$ ($1\leq i\leq r$), dont les corps résiduels sont isomorphes à $k'$, et qui sont des $A$-modules fidèlement plats (18.11.4). Appliquant (2.3.4) et (6.1.1), on voit que dans un $A'_i$ donné, il existe un idéal minimal $\mathfrak q'_i$ au-dessus de $\mathfrak q$ tel que $n'=\dim(A'_i/\mathfrak q'_i)\geq\dim(A/\mathfrak q)=n$. D'autre part (18.11.4), on a $(\hat{\Omega}^1_{A'_i/k'})_{\mathfrak q'_i}=(\hat{\Omega}^1_{A/k})_{\mathfrak q}\otimes_{A_{\mathfrak q}}(A'_i)_{\mathfrak q'_i}$, donc le nombre minimum $m'$ de générateurs du $(A'_i)_{\mathfrak q'_i}$-module $(\hat{\Omega}^1_{A'_i/k'})_{\mathfrak q'_i}$ est au plus égal à $m$; d'où notre assertion.

Supposons donc $k$ algébriquement clos, et montrons qu'on peut en outre se borner au cas où $\mathfrak q=0$. En effet, on a $\Omega^1_{(A/\mathfrak q)/A}=0$ (0, 20.4.12), donc la suite exacte (18.11.3.6) fournit un homomorphisme surjectif $\hat{\Omega}^1_{A/k}\otimes_A(A/\mathfrak q)\to\hat{\Omega}^1_{(A/\mathfrak q)/k}$, d'où on déduit aussitôt que le nombre minimum de générateurs du $K$-module $\hat{\Omega}^1_{(A/\mathfrak q)/k}\otimes_{A/\mathfrak q}K$, où $K$ est le corps des fractions de $A/\mathfrak q$, est au plus égal à $m$.

Mais si $A$ est intègre et $K$ son corps des fractions, $K$ est séparable sur le corps algébriquement clos $k$, donc une $k$-algèbre géométriquement régulière (6.7.6); on peut donc appliquer (18.11.3) pour $\mathfrak p=0$ en observant que $k=k^p$ puisque $k$ est algébriquement clos, et dans ce cas on a $m=n$.

\begin{lemma}[18.11.5.2]
\label{IV.18.11.5.2}
Soient $k$ un corps, $A$ une $k$-algèbre locale noethérienne complète, dont le corps résiduel est une extension finie de $k$, $\mathfrak m$ son idéal maximal, $\mathfrak p\subset\mathfrak m$ un idéal premier de $A$. Soit $m$ le nombre minimum de générateurs du $A_{\mathfrak p}$-module $(\hat{\Omega}^1_{A/k})_{\mathfrak p}$, et soit $(x_i)_{1\leq i\leq m}$ une famille d'éléments de $\mathfrak m$ n'appartenant pas à $\mathfrak p$. Alors il existe $m$ éléments inversibles $u_i$ ($1\leq i\leq m$) de $A$ tels que si l'on pose $y_i=u_ix_i$, les images canoniques des $dy_i$ dans $(\hat{\Omega}^1_{A/k})_{\mathfrak p}$ engendrent cet $A_{\mathfrak p}$-module.
\end{lemma}

Pour tout $x\in A$, désignons par $\delta(x)$ l'image canonique de $dx$ ($=d_{A/k}x$) dans $(\hat{\Omega}^1_{A/k})_{\mathfrak p}\otimes_{A_{\mathfrak p}}k(\mathfrak p)=\hat{\Omega}^1_{A/k}\otimes_A k(\mathfrak p)$; comme ce $k(\mathfrak p)$-espace vectoriel est par hypothèse de rang $m$, il suffit (en vertu du lemme de Nakayama) de prouver qu'on peut déterminer les $u_i$ tels que les $\delta(u_ix_i)$ forment un système libre. Raisonnons par récurrence et supposons que, pour un entier $r<m$, on ait déterminé les $u_i$ ($1\leq i\leq r$) tels que les $\delta(u_ix_i)$ pour $1\leq i\leq r$ soient linéairement indépendants sur $k(\mathfrak p)$; si $\delta(x_{r+1})$ n'est pas combinaison linéaire des $\delta(u_ix_i)$ pour $1\leq i\leq r$, il suffit de prendre $u_{r+1}=1$ pour poursuivre la récurrence. Dans le cas contraire, notons que pour $u\in A$, $\delta(ux_{r+1})$ est l'image canonique de $(du)x_{r+1}+u(dx_{r+1})$; comme $u(dx_{r+1})$ a une image canonique combinaison linéaire des $\delta(u_ix_i)$ pour $1\leq i\leq r$, et que d'autre part l'image canonique de $x_{r+1}$ dans $k(\mathfrak p)$ est $\neq0$, on voit qu'il suffit de prouver qu'il existe un élément inversible $u\in A$ tel que $\delta(u)$ ne soit pas combinaison linéaire des $\delta(u_ix_i)$ pour $1\leq i\leq r$.

\oldpage[IV-4]{176}

Or, il résulte de (0, 20.7.15) et de ($0_{\mathrm I}$, 7.2.9) que les $\delta(x)$ engendrent le $k(\mathfrak p)$-espace vectoriel $\hat{\Omega}^1_{A/k}\otimes_A k(\mathfrak p)$; comme par hypothèse $r<m$, il existe donc $z\in A$ tel que $\delta(z)$ ne soit pas combinaison linéaire des $\delta(u_ix_i)$ pour $1\leq i\leq r$. Si $z\notin\mathfrak m$, on prendra $u=z$; sinon, $1+z$ est inversible et on a $\delta(1+z)=\delta(z)$ puisque $d(1+z)=dz$; on prendra alors $u=1+z$, ce qui achève de prouver (18.11.5.2).

Revenons maintenant à la démonstration de l'implication $a)\Rightarrow b)$ dans (18.11.5). Notons d'abord que puisque $\mathfrak p\neq\mathfrak m$, il existe un système de paramètres $(x_i)_{1\leq i\leq n}$ de $A$ (avec $n=\dim(A)$) tel que $x_i\notin\mathfrak p$ pour $1\leq i\leq n$: en effet, on ne peut avoir $x_i\in\mathfrak p$ pour tout $i$, sans quoi $A/\mathfrak p$ serait de longueur finie, et $\mathfrak p$ serait maximal, contrairement à l'hypothèse. Mais si par exemple $x_1\notin\mathfrak p$, il suffit, pour chaque indice $i$ tel que $x_i\in\mathfrak p$, de remplacer $x_i$ par $x_1+x_i$, pour avoir un système de paramètres dont aucun élément n'appartient à $\mathfrak p$. L'hypothèse $a)$ et la relation $\dim(A/\mathfrak q)=n$ entraînent, en vertu de (18.11.5.1), que $m\geq n$; on peut donc considérer une famille $(x_i)_{1\leq i\leq m}$ d'éléments $x_i\notin\mathfrak p$, dont les $n$ premiers forment un système de paramètres de $A$. Multipliant en outre les $x_i$ par des éléments inversibles $u_i$ de $A$, on peut, grâce à (18.11.5.2) supposer que les images des $dx_i$ ($1\leq i\leq m$) dans $(\hat{\Omega}^1_{A/k})_{\mathfrak p}$ engendrent cet $A_{\mathfrak p}$-module, et la multiplication par les $u_i$ n'a pas altéré le fait que les $x_i$ pour $1\leq i\leq n$ forment un système de paramètres.

Considérons alors le $k$-homomorphisme local $u:B\to A$ tel que $u(T_i)=x_i$ pour $1\leq i\leq m$ (Bourbaki, \emph{Alg. comm.}, chap. III, \textsection~4, n\textsuperscript{o}~5, prop. 6); comme les $x_i$ engendrent un idéal de définition de $A$, il résulte de ($0_{\mathrm I}$, 7.4.4 et 7.4.3) que $u$ fait de $A$ une $B$-algèbre finie. On a donc (18.11.3.5) la suite exacte
\[
  \hat{\Omega}^1_{B/k}\otimes_B A\xrightarrow{v}\hat{\Omega}^1_{A/k}\longrightarrow\Omega^1_{A/B}\longrightarrow0.
\]
Mais les $d_Ax_i$ sont les images canoniques par $v$ des éléments $d_BT_i\otimes1$ (0, 20.5.2.6). Si on localise la suite exacte précédente en $\mathfrak p$, on voit donc que $v_{\mathfrak p}:(\hat{\Omega}^1_{B/k}\otimes_B A)_{\mathfrak p}\to(\hat{\Omega}^1_{A/k})_{\mathfrak p}$ est surjectif, et par suite on a $0=(\Omega^1_{A/B})_{\mathfrak p}=\Omega^1_{A_{\mathfrak p}/B_{\mathfrak r}}$ (où $\mathfrak r$ est l'image réciproque de $\mathfrak p$ dans $A$, cf. (16.4.15)). En vertu de (17.4.1), cela entraîne la propriété $b)$ de (18.11.5). C.Q.F.D.

\begin{remark}[18.11.6]
\label{IV.18.11.6}
Dans l'énoncé de (18.11.5), on ne peut supprimer l'hypothèse $\mathfrak p\neq\mathfrak m$. En effet, pour tout $k$-homomorphisme local $u:B\to A$, où $B=k[[T_1,\ldots,T_r]]$ ($r$ entier quelconque) faisant de $A$ une $B$-algèbre finie, $\mathfrak m$ est le seul point de $\Spec(A)$ au-dessus de l'idéal maximal $\mathfrak n$ de $B$; le morphisme $\Spec(A)\to\Spec(B)$ ne peut être non ramifié en $\mathfrak m$ que si $A/\mathfrak nA$ est un corps, extension séparable de $k$ (17.4.1), ce qui implique que le corps résiduel $K$ de $A$ est une extension séparable de $k$. Si cette condition n'est pas remplie, la conclusion de (18.11.5) ne peut jamais être vérifiée pour $\mathfrak p=\mathfrak m$, quel que soit l'entier $m$.
\end{remark}

\begin{corollary}[18.11.7]
\label{IV.18.11.7}
Soient $k$ un corps, $A$ une $k$-algèbre locale noethérienne complète équidimensionnelle, dont le corps résiduel est une extension finie de $k$, $\mathfrak p$ un idéal premier de $A$. Alors le nombre minimum de générateurs du $A_{\mathfrak p}$-module $(\hat{\Omega}^1_{A/k})_{\mathfrak p}$ est au moins égal à $\dim(A)$.
\end{corollary}

C'est un cas particulier de (18.11.5.1).

Plus particulièrement:

\oldpage[IV-4]{177}

\begin{corollary}[18.11.8]
\label{IV.18.11.8}
Soient $k$ un corps, $A$ une $k$-algèbre locale noethérienne complète intègre, dont le corps résiduel est une extension finie de $k$. Si $K$ est le corps des fractions de $A$, on a
\[
  \operatorname{rg}_K\bigl((\hat{\Omega}^1_{A/k})\otimes_A K\bigr)\geq\dim(A).
\]
\end{corollary}

Il suffit de faire $\mathfrak p=0$ dans (18.11.7).

\begin{corollary}[18.11.9]
\label{IV.18.11.9}
Soient $k$ un corps, $A$ une $k$-algèbre locale noethérienne complète de dimension $n$, dont le corps résiduel est une extension finie de $k$. Soit $\mathfrak p$ un idéal premier de $A$ distinct de l'idéal maximal, et contenant un idéal premier minimal $\mathfrak q$ tel que $\dim(A/\mathfrak q)=n$. Supposons que le $A_{\mathfrak p}$-module $(\hat{\Omega}^1_{A/k})_{\mathfrak p}$ admette $n$ générateurs. Alors:
\begin{enumerate}
  \item[(i)] Le $A_{\mathfrak p}$-module $(\hat{\Omega}^1_{A/k})_{\mathfrak p}$ est libre de rang $n$.
  \item[(ii)] Il existe un $k$-homomorphisme local $u:B\to A$, où $B=k[[T_1,\ldots,T_n]]$, faisant de $A$ une $B$-algèbre finie, et tel que le morphisme correspondant $\Spec(A)\to\Spec(B)$ soit étale au point $\mathfrak p$.
  \item[(iii)] La $k$-algèbre $A_{\mathfrak p}$ est géométriquement régulière.
\end{enumerate}
\end{corollary}

Prouvons d'abord (ii); en vertu de (18.11.5), il existe un homomorphisme local
\[
  u:B=k[[T_1,\ldots,T_n]]\longrightarrow A
\]
faisant de $A$ une $B$-algèbre finie, et tel que le morphisme correspondant $\Spec(A)\to\Spec(B)$ soit non ramifié au point $\mathfrak p$; posons $\mathfrak r=u^{-1}(\mathfrak p)$. L'hypothèse sur $\mathfrak q\subset\mathfrak p$, et le fait que $A$ soit un quotient d'anneau régulier (0, 19.8.8) impliquent que l'on a $\dim(A_{\mathfrak p})=n-\dim(A/\mathfrak p)$ (0, 16.5.12). On a de même $\dim(B_{\mathfrak r})=n-\dim(B/\mathfrak r)$. Enfin, puisque le morphisme $\Spec(A)\to\Spec(B)$ est non ramifié au point $\mathfrak p$, la fibre de ce morphisme au point $\mathfrak r$ est de dimension 0, donc (0, 16.3.9) on a $\dim(A/\mathfrak p)\leq\dim(B/\mathfrak r)$ et par suite $\dim(B_{\mathfrak r})\leq\dim(A_{\mathfrak p})$. On conclut donc de (18.10.1) que le morphisme $\Spec(A)\to\Spec(B)$ est étale au point $\mathfrak p$.

L'assertion (iii) résulte de ce que $A_{\mathfrak p}$ est formellement lisse sur $B_{\mathfrak r}$ et que $B_{\mathfrak r}$ est formellement lisse sur $k$ (pour les topologies préadiques) (0, 19.3.4 et 19.3.5), donc $A_{\mathfrak p}$ est formellement lisse sur $k$ pour sa topologie préadique, et par suite est géométriquement régulier sur $k$ (0, 22.5.8).

Prouvons enfin (i). Soit $k'$ une extension algébriquement close de $k$, et considérons la $k'$-algèbre semi-locale $A'=A\widehat{\otimes}_k k'$. Si l'on considère $A$ comme $B$-module de type fini au moyen de l'homomorphisme $u$, on peut donc écrire, à un isomorphisme canonique près, $A'=A\otimes_B(B\widehat{\otimes}_k k')$ ((7.5.7.1), dans l'énoncé duquel on rappelle qu'il n'est pas nécessaire de supposer que le corps résiduel de $B$ soit une extension finie de $k$). Posons $B'=B\widehat{\otimes}_k k'$, qui s'identifie canoniquement à l'algèbre de séries formelles $k'[[T_1,\ldots,T_n]]$. Puisque le morphisme $\Spec(A)\to\Spec(B)$ est fini et étale au point $\mathfrak p$, le morphisme $\Spec(A')\to\Spec(B')$ est fini, et étale en tout point $\mathfrak p'$ au-dessus de $\mathfrak p$ (17.3.3); d'ailleurs, $A'$ est composé direct d'anneaux locaux $A'_i$ de dimension $n$ ($1\leq i\leq r$) (18.11.4) et $\mathfrak p'$ s'identifie à un idéal premier $\mathfrak p'_i$ de l'un des $A'_i$. Le même raisonnement que ci-dessus prouve alors que $(A'_i)_{\mathfrak p'_i}$ est géométriquement régulier sur $k'$, et comme $k'$ est parfait, il résulte de (18.11.3) que $(\hat{\Omega}^1_{A'_i/k'})_{\mathfrak p'_i}$ est un $(A'_i)_{\mathfrak p'_i}$-module libre de type fini. Mais puisque $(\hat{\Omega}^1_{A'_i/k'})_{\mathfrak p'_i}$ est isomorphe à $(\hat{\Omega}^1_{A/k})_{\mathfrak p}\otimes_{A_{\mathfrak p}}(A'_i)_{\mathfrak p'_i}$ (18.11.4) et que $(A'_i)_{\mathfrak p'_i}$ est un $A_{\mathfrak p}$-module fidèlement plat, on voit que $(\hat{\Omega}^1_{A/k})_{\mathfrak p}$ est un $A_{\mathfrak p}$-module libre de type fini (2.5.2);

\oldpage[IV-4]{178}

d'ailleurs, en vertu de (18.11.5.1) et de l'hypothèse $\dim(A/\mathfrak q)=n$, on conclut que $(\hat{\Omega}^1_{A/k})_{\mathfrak p}$ est un $A_{\mathfrak p}$-module libre de rang $n$. C.Q.F.D.

\begin{theorem}[18.11.10]
\label{IV.18.11.10}
Soient $k$ un corps d'exposant caractéristique $p$, $A$ une $k$-algèbre locale noethérienne complète, dont le corps résiduel est une extension finie de $k$, $\mathfrak p$ un idéal premier de $A$ distinct de l'idéal maximal. Les conditions suivantes sont équivalentes:
\begin{enumerate}
  \item[$a)$] Pour toute extension $k'$ de $k$, et tout idéal premier $\mathfrak p'$ de $A'=A\widehat{\otimes}_k k'$ au-dessus de $\mathfrak p$, $A'_{\mathfrak p'}$ est un anneau régulier.
  \item[$a')$] Il existe une extension parfaite $k'$ de $k$ et un idéal premier $\mathfrak p'$ de $A'=A\widehat{\otimes}_k k'$ au-dessus de $\mathfrak p$, tels que $A'_{\mathfrak p'}$ soit régulier.
  \item[$b)$] Soit $n$ la plus grande des dimensions des composantes irréductibles de $\Spec(A)$ auxquelles appartient $\mathfrak p$; alors $(\hat{\Omega}^1_{A/k})_{\mathfrak p}$ est un $A_{\mathfrak p}$-module libre de rang $n$.
\end{enumerate}
Supposons en outre que $n=\dim(A)$ (ce qui sera le cas si $A$ est équidimensionnel); alors les conditions précédentes sont aussi équivalentes à:
\begin{enumerate}
  \item[$c)$] Il existe un $k$-homomorphisme local $u:B\to A$, où $B=k[[T_1,\ldots,T_n]]$, faisant de $A$ une $B$-algèbre finie, et tel que le morphisme correspondant $\Spec(A)\to\Spec(B)$ soit étale au point $\mathfrak p$.
\end{enumerate}
Chacune des conditions $a)$, $a')$, $b)$ entraîne la suivante:
\begin{enumerate}
  \item[$d)$] L'anneau $A_{\mathfrak p}$ est géométriquement régulier sur $k$.
\end{enumerate}
Si de plus $[k:k^p]<+\infty$, la condition $d)$ est équivalente à $a)$, $a')$ et $b)$.
\end{theorem}

Le fait que $d)$ entraîne $b)$ lorsque $[k:k^p]<+\infty$ n'est autre que (18.11.3). Il est clair que $a)$ entraîne trivialement $a')$; montrons que $c)$ entraîne $a)$ lorsque $\dim(A)=n$. Avec les notations de $a)$, on a alors
\[
  A'=A\otimes_B B',\qquad B'=B\widehat{\otimes}_k k'=k'[[T_1,\ldots,T_n]]
\]
(7.5.7.1). Le morphisme $\Spec(A')\to\Spec(B')$ est alors fini et étale au point $\mathfrak p'$, et le raisonnement qui prouve (18.11.9, (iii)) montre que $A'_{\mathfrak p'}$ est régulier.

Le fait que $b)$ entraîne $c)$ lorsque $\dim(A)=n$ résulte de (18.11.9). Montrons que $a')$ entraîne $b)$ lorsque $\dim(A)=n$. On sait que $A'$ est un $A$-module plat (18.11.4) et $\dim(A')=\dim(A)=n$; il résulte de (2.3.4) et de (6.1.1) qu'il existe un idéal premier minimal $\mathfrak q'$ de $A'$ contenu dans $\mathfrak p'$, au-dessus de $\mathfrak q$ et tel que $\dim(A'/\mathfrak q')=\dim(A/\mathfrak q)=n$. Par ailleurs, il résulte de (18.11.4) que
\[
  (\hat{\Omega}^1_{A'/k'})_{\mathfrak p'}=(\hat{\Omega}^1_{A/k})_{\mathfrak p}\otimes_{A_{\mathfrak p}}A'_{\mathfrak p'};
\]
comme $k'$ est parfait, l'hypothèse que $A'_{\mathfrak p'}$ est régulier entraîne qu'il est géométriquement régulier sur $k'$ (6.7.7). Puisque $k'^p=k'$, on peut appliquer à $A'$ et $\mathfrak p'$ le fait que $d)$ entraîne $b)$, donc $(\hat{\Omega}^1_{A'/k'})_{\mathfrak p'}$ est un $A'_{\mathfrak p'}$-module libre de rang $n$; par fidèle platitude (18.11.4 et 2.5.2), on en conclut que $(\hat{\Omega}^1_{A/k})_{\mathfrak p}$ est un $A_{\mathfrak p}$-module libre de rang $n$.

Ceci montre l'équivalence de $a)$, $a')$, $b)$ et $c)$ lorsque $\dim(A)=n$. Il reste à prouver que $a)$, $a')$ et $b)$ sont encore équivalentes dans le cas général. Soit $\mathfrak J$ l'idéal de $A$, noyau de l'homomorphisme canonique $A\to A_{\mathfrak p}$, et posons $A_1=A/\mathfrak J$; on a nécessairement $\mathfrak J\subset\mathfrak p$, et si l'on pose $\mathfrak p_1=\mathfrak p/\mathfrak J$, l'homomorphisme canonique $A_{\mathfrak p}\to(A_1)_{\mathfrak p_1}$ est bijectif; on en conclut (I, 6.5.4) que l'injection canonique $\Spec(A_1)\to\Spec(A)$ est un isomorphisme local au point $\mathfrak p_1$, et l'on a $\mathfrak J_{\mathfrak p}=0$. On voit comme dans (18.11.3) que l'on a une suite exacte
\[
  \mathfrak J/\mathfrak J^2\longrightarrow\hat{\Omega}^1_{A/k}\otimes_A A_1\longrightarrow\hat{\Omega}^1_{A_1/k}\longrightarrow0.
\]

\oldpage[IV-4]{179}

Et en localisant en $\mathfrak p$, il vient un isomorphisme
\[
  (\hat{\Omega}^1_{A/k})_{\mathfrak p}\xrightarrow{\sim}(\hat{\Omega}^1_{A_1/k})_{\mathfrak p_1}.
\]
Ceci montre que la condition $b)$ pour l'anneau $A$ et l'idéal $\mathfrak p$ est équivalente à la condition $b)$ pour l'anneau $A_1$ et l'idéal $\mathfrak p_1$. D'autre part, avec les notations de $a)$, on a $A'_1=A_1\widehat{\otimes}_k k'=A'/\mathfrak JA'$ à isomorphisme près ((7.5.7.1), où l'hypothèse sur le corps résiduel de $B$ est superflue); si $\mathfrak p'$ est un idéal premier de $A'$ au-dessus de $\mathfrak p$, tout élément de $\mathfrak JA'$ annule un élément de $A'-\mathfrak p'$, donc $\mathfrak JA'\subset\mathfrak p'$, et si l'on pose $\mathfrak p'_1=\mathfrak p'/\mathfrak JA'$, $\mathfrak p'_1$ est au-dessus de $\mathfrak p_1$ et $(A'_1)_{\mathfrak p'_1}$ s'identifie canoniquement à $A'_{\mathfrak p'}$; ceci montre donc que la condition $a)$ (resp. $a')$) pour l'anneau $A$ et l'idéal $\mathfrak p$ est équivalente à la condition $a)$ (resp. $a')$) pour l'anneau $A_1$ et l'idéal $\mathfrak p_1$. Or, tous les idéaux premiers minimaux de $A$ contenus dans $\mathfrak p$ contiennent $\mathfrak J$ puisque $\Spec(A_1)\to\Spec(A)$ est un isomorphisme local au point $\mathfrak p_1$; d'autre part les idéaux de $\operatorname{Ass}_A(A/\mathfrak J)$ sont les idéaux de $\operatorname{Ass}(A)$ qui sont contenus dans $\mathfrak p$ (Bourbaki, \emph{Alg. comm.}, chap. IV, \textsection~1, n\textsuperscript{o}~2, prop. 6); donc les idéaux premiers minimaux de $A_1$ sont tous contenus dans $\mathfrak p_1$, et on a par suite $\dim(A_1)=n$. Il suffit alors d'appliquer à $A_1$ et à $\mathfrak p_1$ ce qui a été prouvé plus haut.

\begin{remarks}[18.11.11]
\label{IV.18.11.11}
\begin{enumerate}
  \item[(i)] L'équivalence des conditions $d)$ et $b)$ dans (18.11.10) n'est plus valable lorsqu'on ne suppose plus que $[k:k^p]<+\infty$. En effet, dans l'exemple de (0, 22.7.7, (ii)), l'anneau $B=A/\mathfrak q$ est intègre et de dimension 1; d'autre part, la suite
  \[
    (\mathfrak q/\mathfrak q^2)\otimes_A L\xrightarrow{j}\hat{\Omega}^1_{A/k}\otimes_A L\longrightarrow\hat{\Omega}^1_{B/k}\otimes_B L\longrightarrow0
  \]
  est exacte (même démonstration que dans (18.11.3)), et on sait que $j$ n'est pas injectif, donc $j=0$ puisque $(\mathfrak q/\mathfrak q^2)\otimes_A L$ est de rang 1. On en conclut que $\operatorname{rg}_L(\hat{\Omega}^1_{B/k}\otimes_B L)=2$. Comme $B_{\mathfrak p}$ est une $k$-algèbre géométriquement régulière, on voit qu'ici la condition $d)$ n'entraîne pas $b)$.
  \item[(ii)] Les notations étant celles de (18.11.10), supposons que $n=\dim(A)$ et soit $(x_i)_{1\leq i\leq n}$ un système de paramètres de $A$ n'appartenant pas à $\mathfrak p$; on en déduit un $k$-homomorphisme local $u:B\to A$ tel que $u(T_i)=x_i$ pour $1\leq i\leq n$, faisant de $A$ une $B$-algèbre finie. Pour que le morphisme correspondant $\Spec(A)\to\Spec(B)$ soit étale au point $\mathfrak p$, il faut et il suffit que les images des $d_Ax_i$ dans $(\hat{\Omega}^1_{A/k})_{\mathfrak p}$ forment un système de générateurs de cet $A_{\mathfrak p}$-module. En effet, on a vu dans la démonstration de (18.11.5) que si le morphisme $\Spec(A)\to\Spec(B)$ est non ramifié au point $\mathfrak p$, l'homomorphisme canonique $(\hat{\Omega}^1_{B/k}\otimes_B A)_{\mathfrak p}\to(\hat{\Omega}^1_{A/k})_{\mathfrak p}$ est surjectif et les images des éléments $d_BT_i\otimes1$, qui engendrent $(\hat{\Omega}^1_{B/k}\otimes_B A)_{\mathfrak p}$, sont les $d_Ax_i$; d'où la nécessité de la condition. Inversement, le même raisonnement montre que si cette condition est vérifiée, le morphisme $\Spec(A)\to\Spec(B)$ est non ramifié en $\mathfrak p$, et le raisonnement de (18.11.9) prouve en fait que ce morphisme est étale en $\mathfrak p$.
\end{enumerate}
\end{remarks}

\begin{corollary}[18.11.12]
\label{IV.18.11.12}
Soient $k$ un corps d'exposant caractéristique $p$ tel que $[k:k^p]<+\infty$, $A$ une $k$-algèbre locale noethérienne complète et intègre, qui n'est pas un corps, et dont le corps résiduel est une extension finie de $k$. Il existe alors une extension finie radicielle $k'$ de $k$ telle qu'en posant $A'=A\otimes_k k'$, la $k'$-algèbre $A'_{\mathrm{red}}$ et l'idéal premier $0$ de cette algèbre vérifient les conditions équivalentes $a)$, $a')$, $b)$, $c)$ et $d)$ de (18.11.10). En particulier, si $n=\dim(A)=\dim(A')$,

\oldpage[IV-4]{180}

il existe un $k'$-homomorphisme local $B'=k'[[T_1,\ldots,T_n]]\to A'_{\mathrm{red}}$ faisant de $A'_{\mathrm{red}}$ une $B'$-algèbre finie, et tel que le corps des fractions $K'$ de $A'_{\mathrm{red}}$ soit une extension finie séparable du corps des fractions $k'((T_1,\ldots,T_n))$ de $B'$.
\end{corollary}

Le morphisme $\Spec(A')\to\Spec(A)$ étant radiciel et fini, $A'$ est un anneau local complet et son nilradical est le seul idéal premier au-dessus de l'idéal $0$ de $A$; par platitude, $A'$ s'identifie à un sous-anneau de $K\otimes_k k'$, qui est d'ailleurs l'anneau total des fractions de $A'$; on a par suite $K'=(K\otimes_k k')_{\mathrm{red}}$. Il s'agit de prouver, vu l'équivalence $d)\Leftrightarrow c)$ de (18.11.10), qu'il existe une extension radicielle finie $k'$ de $k$ telle que $(K\otimes_k k')_{\mathrm{red}}$ soit une extension séparable de $k'$ (6.7.6). Or on sait (0, 19.8.9) que sous les hypothèses faites, il existe une sous-$k$-algèbre $C=k[[T_1,\ldots,T_m]]$ de $A$ telle que $A$ soit une $C$-algèbre finie; $K$ est donc une extension finie du corps des fractions $K_1=k((T_1,\ldots,T_m))$ de $C$, qui est séparable sur $k$ (0, 21.9.6.4). Si $p$ est l'exposant caractéristique de $k$, on peut donc écrire
\[
  K\otimes_k k^{p^{-\infty}}=K\otimes_{K_1}(K_1\otimes_k k^{p^{-\infty}})
\]
et $K_1\otimes_k k^{p^{-\infty}}$ est un corps, extension radicielle de $K_1$; on en conclut que $K\otimes_k k^{p^{-\infty}}$ est une algèbre finie sur le corps $K_1\otimes_k k^{p^{-\infty}}$, donc un anneau artinien. La conclusion résulte donc du lemme plus général suivant:

\begin{lemma}[18.11.12.1]
\label{IV.18.11.12.1}
Soient $k$ un corps de caractéristique $p>0$, $K$ une extension de $k$. Les conditions suivantes sont équivalentes:
\begin{enumerate}
  \item[$a)$] L'anneau $K\otimes_k k^{p^{-\infty}}$ est artinien.
  \item[$b)$] Il existe une extension $k'$ de $k$ de type fini telle que $(K\otimes_k k')_{\mathrm{red}}$ soit une $k'$-algèbre séparable.
  \item[$c)$] Il existe une extension finie radicielle $k'$ de $k$ telle que $(K\otimes_k k')_{\mathrm{red}}$ soit une extension séparable de $k'$.
\end{enumerate}
\end{lemma}

Montrons d'abord que $c)$ entraîne $a)$. L'anneau $A=K\otimes_k k'$ est alors un anneau artinien local, et si $\mathfrak N$ est son nilradical, le corps résiduel $L=A/\mathfrak N$ est séparable sur $k'$; par suite $L\otimes_{k'}k^{p^{-\infty}}$ est un corps, et comme il est égal à
\[
  (A\otimes_{k'}k^{p^{-\infty}})/(\mathfrak N\otimes_{k'}k^{p^{-\infty}}),
\]
on voit que $\mathfrak N\otimes_{k'}k^{p^{-\infty}}$ est le nilradical de $A\otimes_{k'}k^{p^{-\infty}}=K\otimes_k k^{p^{-\infty}}$; comme $\mathfrak N$ est un idéal de type fini, le nilradical de $K\otimes_k k^{p^{-\infty}}$ est donc de type fini, ce qui entraîne que l'anneau $K\otimes_k k^{p^{-\infty}}$ est artinien.

Inversement, prouvons que $a)$ entraîne $c)$. Soit $\mathfrak N$ le nilradical de l'anneau local artinien $B=K\otimes_k k^{p^{-\infty}}$, qui est par hypothèse engendré par un nombre fini d'éléments de la forme
\[
  y_i=\sum_j x_{ij}\otimes\xi_{ij},\qquad x_{ij}\in K,\quad \xi_{ij}\in k^{p^{-\infty}}.
\]
Soient $k'$ l'extension radicielle finie de $k$ engendrée par les $\xi_{ij}$, $\mathfrak N_0$ l'idéal de $B_0=K\otimes_k k'$ engendré par les $y_i$; il est clair que les $y_i$ sont nilpotents dans $B_0$; d'autre part, on a $\mathfrak N_0\otimes_{k'}k^{p^{-\infty}}=\mathfrak N$, et par suite $\mathfrak N\cap B_0=\mathfrak N_0$, $\mathfrak N_0$ contient le nilradical de $B_0$, donc il lui est égal. Comme $B/\mathfrak N=(B_0/\mathfrak N_0)\otimes_{k'}k^{p^{-\infty}}$ est réduit, on en conclut que $B_0/\mathfrak N_0=(K\otimes_k k')_{\mathrm{red}}$ est séparable sur $k$ (4.6.1).

Il est clair que $c)$ entraîne $b)$. Inversement, supposons $b)$ vérifiée, et notons qu'il existe une extension séparable $k_1$ de $k$ telle que $k'$ soit une extension radicielle finie de $k_1$; posons $K_1=K\otimes_k k_1$, qui est un corps. Appliquant l'équivalence de $a)$ et $c)$ à l'extension $K_1$ de $k_1$, on voit que $K_1\otimes_{k_1}k_1^{p^{-\infty}}$ est un anneau artinien; mais cet anneau est égal

\oldpage[IV-4]{181}

à
\[
  K\otimes_k k_1^{p^{-\infty}}=(K\otimes_k k^{p^{-\infty}})\otimes_{k^{p^{-\infty}}}k_1^{p^{-\infty}},
\]
donc $K\otimes_k k^{p^{-\infty}}$ est aussi artinien (Bourbaki, \emph{Alg. comm.}, chap. I, \textsection~3, n\textsuperscript{o}~5, cor. de la prop. 8); on a ainsi prouvé que $b)$ entraîne $a)$, ce qui achève la démonstration de (18.11.12.1) et de (18.11.12).

\subsection{Applications de la localisation étale aux morphismes quasi-finis (généralisations de résultats antérieurs)}
\EGAReaderTarget{ega4.c04.subsec.18-12}{subsection:IV.18.12}

Les résultats de ce numéro nous ont été communiqués par P.~Deligne.

\begin{theorem}[18.12.1]
\label{IV.18.12.1}
Soient $f:X\to Y$ un morphisme localement de type fini, $x$ un point de $X$, $y=f(x)$. Supposons que $x$ soit un point isolé de l'espace $f^{-1}(y)$. Alors, il existe un morphisme étale $Y'\to Y$, un point $x'$ de $X'=X\times_Y Y'$ au-dessus de $x$ et un voisinage ouvert $V'$ de $x'$ dans $X'$ tel que, si $f'=f_{(Y')}:X'\to Y'$, $f'|V'$ soit un morphisme fini. Si de plus $f$ est séparé, $V'$ est à la fois ouvert et fermé, et $X'$ est donc somme de deux sous-préschémas induits sur des ouverts de $X'$, dont l'un est fini sur $Y'$ et contient $x'$.
\end{theorem}

La dernière assertion résulte de ce que, si $f$ est séparé, il en est de même de $f'$; donc, si $j':V'\to X'$ est l'injection canonique, le fait que le morphisme $f'|V'=f'\circ j'$ soit fini entraîne qu'il en est de même de $j'$ (II, 6.1.5), et par suite $V'=j'(V')$ est fermé dans $X'$ (II, 6.1.10).

La question étant locale sur $X$, on peut supposer $Y=\Spec(A)$ et $X=\Spec(B)$ affines, $B$ étant une $A$-algèbre de type fini, donc de la forme $C/\mathfrak J$, où $C=A[T_1,\ldots,T_r]$ et $\mathfrak J$ est un idéal de $C$; $f$ est donc séparé, et on peut en outre supposer que $f^{-1}(y)=\{x\}$. Soit $(\mathfrak J_\lambda)$ la famille des idéaux de type fini de $C$ contenus dans $\mathfrak J$, de sorte que $\mathfrak J$ est la réunion filtrante des $\mathfrak J_\lambda$. Si $X$ et les $X_\lambda=\Spec(C/\mathfrak J_\lambda)$ sont considérés comme des sous-préschémas fermés de $Z=\Spec(C)$, on a donc, pour les espaces sous-jacents, $X=\bigcap_\lambda X_\lambda$; si $f_\lambda:X_\lambda\to Y$ est le morphisme structural, on en déduit que
\[
  f^{-1}(y)=\bigcap_\lambda f_\lambda^{-1}(y),
\]
et comme les ensembles $f_\lambda^{-1}(y)$ sont fermés dans l'espace noethérien $Z_y=\Spec(k(y)[T_1,\ldots,T_r])$, il existe un indice $\lambda$ tel que $f_\lambda^{-1}(y)=f^{-1}(y)=\{x\}$. On peut donc supposer que les $X_\lambda$ vérifient la même condition que $X$ au point $x$; si l'on prouve la proposition pour un tel $X_\lambda$, elle en résultera pour $X$, car si $x'$ est un point au-dessus de $x$ dans $Z'=Z\times_Y Y'$, $X'_\lambda=p^{-1}(X_\lambda)$, où $p:Z'\to Z$ est la projection (de sorte qu'on a aussi $X'=p^{-1}(X)$), s'il existe un voisinage ouvert $V'_\lambda$ de $x'$ dans $X'_\lambda$ qui est fini sur $Y'$, \emph{a fortiori} la restriction de $f'_\lambda$ au sous-préschéma fermé $X'\cap V'_\lambda$ de $V'_\lambda$ sera un morphisme fini. On est donc ramené au cas où $f$ est de présentation finie.

Notons en outre que l'ensemble des points de $X$ isolés dans leur fibre est ouvert dans $X$ (13.1.4), donc $f$ est quasi-fini au point $x$, et on peut par suite se borner au cas où $f$ est quasi-fini. Soit $A''$ l'anneau local hensélisé de $\sh{O}_{Y,y}$, et posons $Y''=\Spec(A'')$. Si $X''=X\times_Y Y''$, $f''=f_{(Y'')}:X''\to Y''$ est séparé, quasi-fini et de présentation finie; comme $A''$ est hensélien, il y a pour tout $x''\in X''$ au-dessus de $x$ un voisinage ouvert et fermé $V''$ de $x''$ dans $X''$ qui est fini sur $Y''$ (18.5.11, c)); comme l'immersion $V''\to X''$ est ouverte et fermée, elle est quasi-compacte, donc de présentation finie (1.6.2), donc $f''|V''$ est de présentation finie.

Cela étant, par définition, $A''$ est limite inductive d'une famille filtrante $(A_\lambda)$ de $\sh{O}_{Y,y}$-algèbres strictement essentiellement étales (18.6.5); chaque $A_\lambda$ est lui-même limite inductive d'une famille filtrante $(B_{\lambda\mu})$ de $\sh{O}_{Y,y}$-algèbres étales ((18.6.1) et (8.1.2, a)). Enfin, comme $B_{\lambda\mu}$ peut être supposée une $\sh{O}_{Y,y}$-algèbre de présentation finie, il résulte encore de (8.1.2, a)) et de (17.7.8) que $B_{\lambda\mu}$ est limite inductive d'une famille filtrante $(C_{\lambda\mu\nu})$ de $A$-algèbres étales. Finalement, on voit, en utilisant le théorème de la double limite inductive, que $Y''$ est limite projective d'une famille filtrante $(Y'_\alpha)$ de schémas affines étales sur $Y$. Si $x'_\alpha$ est la projection de $x''$ sur $X'_\alpha=X\times_Y Y'_\alpha$, on peut supposer que le voisinage $V''$, qui est quasi-compact, est de la forme $V'_\alpha\times_{Y'_\alpha}Y''$ où $V'_\alpha$ est un voisinage ouvert de $x'_\alpha$ dans $X'_\alpha$ (8.2.11). Enfin, $V''$ étant de présentation finie sur $Y''$, on conclut de (8.10.5, (x)) que pour un $\alpha$ convenable, $V'_\alpha$ est fini sur $Y'_\alpha$. C.Q.F.D.

\begin{remark}[18.12.2]
\label{IV.18.12.2}
Dans la démonstration précédente, on voit (compte tenu de (18.6.2)) que l'on a construit un $Y'$ répondant à la question, et tel en plus que si $y'=f'(X')$, l'homomorphisme $k(y)\to k(y')$ soit bijectif.
\end{remark}

\begin{corollary}[18.12.3]
\label{IV.18.12.3}
Soient $f:X\to Y$ un morphisme localement de type fini et séparé, $y$ un point de $Y$ tel que le sous-espace $f^{-1}(y)$ soit fini et discret. Il existe alors un morphisme étale $Y'\to Y$, un point $y'\in Y'$ au-dessus de $y$, tel que $k(y')=k(y)$, et une décomposition de $X'$ en somme de deux sous-préschémas $X'_1$, $X'_2$, induits sur des ouverts de $X'$, tels que la restriction de $f'=f_{(Y')}:X'\to Y'$ à $X'_1$ soit un morphisme fini et que l'on ait $X'_2\cap f'^{-1}(y')=\varnothing$.
\end{corollary}

Si $n$ est le nombre de points de $f^{-1}(y)$, on raisonne par récurrence sur $n$, le corollaire étant trivial pour $n=0$. Soit $x$ un point de $f^{-1}(y)$; en vertu de (18.12.1) et (18.12.2) il y a un morphisme étale $Y_1\to Y$, un point $y_1$ de $Y_1$ au-dessus de $y$ tel que $k(y_1)=k(y)$ et si l'on pose $S_1=X\times_Y Y_1$, $f_1=f_{(S_1)}:S_1\to Y_1$, il existe un point $x_1$ de $f_1^{-1}(y_1)$ tel que $S_1$ soit somme de deux ouverts $V_1$, $X_1$, $V_1$ étant fini sur $Y_1$ et un voisinage de $x_1$.

\oldpage[IV-4]{182}

En vertu de la relation $k(y_1)=k(y)$, la fibre $f_1^{-1}(y_1)$ dans $S_1$ est isomorphe à $f^{-1}(y)$, donc $X_1\cap f_1^{-1}(y_1)$ est fini, discret et a $n-1$ points. Comme $f_1|X_1$ est localement de type fini et séparé, on applique l'hypothèse de récurrence à ce morphisme: il y a un morphisme étale $Y'\to Y_1$, un point $y'\in Y'$ au-dessus de $y_1$ tel que $k(y')=k(y_1)$, et, si $X'=X\times_Y Y'$ et si $p:X'\to S_1$ est la projection canonique, on a une décomposition de $p^{-1}(X_1)$ en somme de deux sous-préschémas $U'$ et $X'_2$ induits sur des ouverts de $X'$ tels que $X'_2\cap f'^{-1}(y')=\varnothing$ et que $U'$ soit fini sur $Y'$. Par ailleurs $V'=p^{-1}(V_1)$ est fini sur $Y'$ et $X'$ est somme de $U'$, $V'$ et $X'_2$; on répondra donc à la question en prenant $X'_1$ somme de $U'$ et $V'$.

Le corollaire suivant améliore (8.11.1):

\begin{corollary}[18.12.4]
\label{IV.18.12.4}
Pour qu'un morphisme $f:X\to Y$ soit fini, il faut et il suffit qu'il soit séparé, universellement fermé, localement de type fini et que, pour tout $y\in Y$, $f^{-1}(y)$ soit un espace fini discret; en particulier, un morphisme propre et quasi-fini est fini.
\end{corollary}

En effet, on peut appliquer le corollaire (18.12.3) à un point quelconque $y$ de $Y$. Avec les notations de ce corollaire, $f'$ est un morphisme fermé, donc, puisque $X'_2$ est fermé dans $X'$, $f'(X'_2)$ est fermé dans $Y'$ et ne contient pas $y'$; il existe par suite un voisinage ouvert $U'$ de $y'$ dans $Y'$ tel que $U'\to Y$ soit de type fini et que $f'^{-1}(U')=X\times_Y U'$ soit fini sur $U'$. Soit $U$ l'image de $U'$ dans $Y$, qui est un ouvert de $Y$ en vertu de (II, 3.1) puisque le morphisme $Y'\to Y$ est étale, donc plat et localement de présentation finie; on a évidemment encore $X\times_Y U'=X\times_U U'=f^{-1}(U)\times_U U'$. Puisque maintenant le morphisme $U'\to U$ est fidèlement plat et quasi-compact, on déduit de (2.7.1, (xv)) que le morphisme $f^{-1}(U)\to U$, restriction de $f$, est fini, ce qui prouve le corollaire.

\begin{remark}[18.12.5]
\label{IV.18.12.5}
On notera que dans la démonstration de (18.12.4) on n'a pas utilisé l'hypothèse que $f$ est universellement fermé sous sa forme générale, mais seulement que $f_{(Y')}$ est un morphisme fermé pour tout morphisme étale $Y'\to Y$.
\end{remark}

\begin{corollary}[18.12.6]
\label{IV.18.12.6}
Soit $f:X\to Y$ un morphisme de préschémas. Les conditions suivantes sont équivalentes:
\begin{enumerate}
  \item[$a)$] $f$ est une immersion fermée.
  \item[$b)$] $f$ est un monomorphisme propre.
  \item[$c)$] $f$ est propre et, pour tout $y\in Y$, le $k(y)$-préschéma $X_y=f^{-1}(y)$ est radiciel et géométriquement réduit sur $k(y)$ (c'est-à-dire vide ou $k(y)$-isomorphe à $\Spec(k(y))$).
\end{enumerate}
\end{corollary}

Cela se déduit de (18.12.4) comme (8.11.5) se déduit de (8.11.1).

\begin{proposition}[18.12.7]
\label{IV.18.12.7}
Soient $f:X\to Y$ un morphisme localement de type fini, $y$ un point de $Y$. Pour qu'il existe un voisinage ouvert $U$ de $y$ tel que la restriction $f^{-1}(U)\to U$ de $f$ soit une immersion fermée, il faut et il suffit que $X_y$ soit un $k(y)$-préschéma radiciel et géométriquement réduit sur $k(y)$ (c'est-à-dire vide ou $k(y)$-isomorphe à $\Spec(k(y))$) et qu'il existe un voisinage ouvert $V$ de $y$ tel que la restriction $f^{-1}(V)\to V$ de $f$ soit un morphisme universellement fermé.
\end{proposition}

Les conditions étant évidemment nécessaires, prouvons qu'elles sont suffisantes. On peut se borner au cas où $X_y\neq\varnothing$. En vertu de la seconde condition, on peut déjà supposer $f$ universellement fermé. Montrons en outre qu'on peut supposer $f$ affine, donc séparé; cela résulte du lemme suivant:

\begin{lemma}[18.12.7.1]
\label{IV.18.12.7.1}
Soient $f:X\to Y$ un morphisme fermé, $y\in Y$ un point tel que tout voisinage de $X_y$ contienne un voisinage affine de $X_y$ (condition toujours vérifiée lorsque $X_y$ est vide ou réduite à un seul point). Alors il existe un voisinage ouvert affine $U$ de $y$ tel que la restriction $f^{-1}(U)\to U$ de $f$ soit un morphisme affine.
\end{lemma}

En effet, soit $U_0$ un voisinage ouvert affine de $y$ dans $Y$ et soit $V$ un voisinage ouvert affine de $X_y$ contenu dans $f^{-1}(U_0)$. Puisque $f$ est fermé, il existe un voisinage ouvert affine $U\subset U_0$ de $y$ tel que $f^{-1}(U)\subset V$. Comme la restriction $g:V\to U_0$ de $f$ est un morphisme affine, il en est de même de sa restriction $f^{-1}(U)\to U$ (II, 1.2.5).

Supposons donc $f$ affine et universellement fermé, et montrons qu'il existe un voisinage ouvert $U$ de $y$ tel que la restriction $f^{-1}(U)\to U$ soit un morphisme fini; comme on peut supposer $Y$ et $X$ affines, on peut supposer $f$ de type fini, donc propre, et il suffit de prouver qu'il existe un voisinage $U$ de $y$ tel que $f^{-1}(U)\to U$ soit un morphisme quasi-fini (18.12.4). Mais il existe un voisinage $V$ de l'ensemble $X_y$ (réduit à un seul point) tel que $f|V$ soit quasi-fini (13.1.4), et comme $f$ est fermé il y a un voisinage $U$ de $y$ tel que $f^{-1}(U)\subset V$, ce qui achève la démonstration, compte tenu de (18.12.6).

\begin{proposition}[18.12.8]
\label{IV.18.12.8}
Soit $f:X\to Y$ un morphisme de préschémas. Pour que $f$ soit entier, il faut et il suffit que $f$ soit affine et universellement fermé.
\end{proposition}

Les conditions sont nécessaires (II, 6.1.10); prouvons qu'elles sont suffisantes. On peut supposer $Y=\Spec(A)$ affine, donc aussi $X=\Spec(B)$, et il faut prouver que tout élément $b\in B$ est entier sur $A$ (II, 6.1.1). Soit $B'$ la sous-$A$-algèbre de $B$ engendrée par $b$, qui est une $A$-algèbre de type fini; posons $X'=\Spec(B')$, de sorte que $f:X\to Y$ se factorise en
\[
  X\xrightarrow{g}X'\xrightarrow{h}Y,
\]
où $h$ est de type fini, et $g$ est dominant puisque l'homomorphisme $B'\to B$ est injectif (I, 1.2.7). Comme $h$ est séparé et $h\circ g$ universellement fermé, $g$ est universellement fermé (II, 5.4.3 et

\oldpage[IV-4]{183}

5.4.9), donc surjectif puisqu'il est dominant; on en conclut (II, 5.4.3 et 5.4.9) que $h$ est aussi universellement fermé, donc propre puisqu'il est de type fini et séparé. Mais alors, pour tout $y\in Y$, le morphisme $h^{-1}(y)\to\Spec(k(y))$ déduit de $h$ est propre et affine, donc fini (III, 4.4.2), autrement dit $h$ est quasi-fini; on déduit donc de (18.12.4) que $h$ est fini, ce qui prouve que $B'$ est une $A$-algèbre finie, et par suite que $b$ est entier sur $A$. C.Q.F.D.

\begin{remark}[18.12.9]
\label{IV.18.12.9}
Il se peut qu'un morphisme $f:X\to Y$ soit entier lorsqu'on le suppose seulement séparé, universellement fermé et tel que pour tout $y\in Y$, le morphisme $f^{-1}(y)\to\Spec(k(y))$ déduit de $f$ soit entier. Il faudrait pour cela prouver que ces conditions impliquent que $f$ est affine, ou encore que toute fibre $X_y$ est contenue dans un voisinage ouvert affine.
\end{remark}

\begin{corollary}[18.12.10]
\label{IV.18.12.10}
Un morphisme $f:X\to Y$ qui est injectif et universellement fermé est entier.
\end{corollary}

Il suffit, par (18.12.8), de prouver que $f$ est affine, ce qui résulte du lemme (18.12.7.1) et de l'hypothèse.

\begin{corollary}[18.12.11]
\label{IV.18.12.11}
Soit $f:X\to Y$ un morphisme de préschémas (resp. un morphisme localement de type fini). Pour que $f$ soit un homéomorphisme universel (2.4.2), il faut et il suffit que $f$ soit entier (resp. fini), radiciel et surjectif.
\end{corollary}

On sait que les conditions sont suffisantes (2.4.5); elles sont nécessaires en vertu de (18.12.10) et (18.12.4).

La proposition suivante améliore (8.11.2):

\begin{proposition}[18.12.12]
\label{IV.18.12.12}
Tout morphisme $f:X\to Y$ qui est quasi-fini et séparé est quasi-affine.
\end{proposition}

Posons $\mathscr A=f_*(\sh{O}_X)$, qui est une $\sh{O}_Y$-Algèbre quasi-cohérente puisque $f$ est quasi-compact et séparé (1.7.4); soit $Z=\Spec(\mathscr A)$ (II, 1.3.1) de sorte que $f$ se factorise canoniquement en
\[
  X\xrightarrow{g}Z\xrightarrow{h}Y,
\]
$h$ étant affine et $g$ correspondant à l'automorphisme identique de $\mathscr A$ (II, 1.2.7); il suffira de prouver que $g$ est une immersion ouverte, ou, ce qui revient au même (17.9.1) que $g$ est étale et radiciel. Il suffira évidemment de montrer que, pour tout $y\in Y$, le morphisme $g$ est étale et radiciel en chaque point de $f^{-1}(y)$. Or, pour tout $y\in Y$, on peut appliquer à $f$ le résultat de (18.12.3), dont nous conservons les notations; puisque le morphisme $Y'\to Y$ est plat, et que $f$ est quasi-compact et séparé, on a, à un isomorphisme canonique près,
\[
  f'_*(\sh{O}_{X'})=\mathscr A\otimes_{\sh{O}_Y}\sh{O}_{Y'}=\mathscr A'
\]
(2.3.1), donc $Z'=Z\times_Y Y'$ s'identifie à $\Spec(f'_*(\sh{O}_{X'}))$, et dans la factorisation canonique
\[
  X'\xrightarrow{g'}Z'\xrightarrow{h'}Y'
\]
de $f'$ (II, 1.2.7), on a $h'=h_{(Y')}$ et $g'=g_{(Y')}$. Cela étant, la décomposition de $X'$ en somme de deux sous-préschémas $X'_1$, $X'_2$ entraîne la décomposition de $\mathscr A'$ en produit direct des deux $\sh{O}_{Y'}$-Algèbres quasi-cohérentes $\mathscr A'_1$, $\mathscr A'_2$, images directes respectives de $\sh{O}_{X'_1}$ et $\sh{O}_{X'_2}$, de sorte que $Z'$ s'identifie à la somme $Z'_1\amalg Z'_2$, où $Z'_i=\Spec(\mathscr A'_i)$ et $g'(X'_i)\subset Z'_i$ pour $i=1,2$. Comme $X'_1$ est fini sur $Y'$, $g'_1=g'|X'_1$ est un isomorphisme de $X'_1$ sur $Z'_1$, puisque $f'|X'_1$ est affine; comme $X'_2\cap f'^{-1}(y')=\varnothing$, on voit que $g'$ est étale et radiciel en chaque point de $f'^{-1}(y')$. Le morphisme $Y'\to Y$ étant plat et localement de présentation finie, on déduit donc d'abord de (17.7.4) que $g$ est étale en tous les points de $X$ projections de points de $f'^{-1}(y')$, c'est-à-dire en tous les points de $f^{-1}(y)$ (I, 3.5.2). D'autre part, le morphisme
\[
  g'_{y'}:f'^{-1}(y')\longrightarrow h'^{-1}(y')
\]
déduit de $g'$ est radiciel; comme $k(y')=k(y)$, le morphisme $g_y:f^{-1}(y)\to h^{-1}(y)$ est aussi radiciel, autrement dit $g$ est radiciel en chaque point de $f^{-1}(y)$, ce qui achève de prouver la proposition.

L'énoncé suivant améliore de même (8.12.6):

\begin{corollary}[18.12.13]
\label{IV.18.12.13}
(« Main Theorem » de Zariski). Soient $Y$ un préschéma quasi-compact et quasi-séparé, $f:X\to Y$ un morphisme quasi-fini et séparé. Alors il existe une factorisation de $f$
\[
  X\xrightarrow{g}Z\xrightarrow{u}Y
\]
où $g$ est une immersion ouverte (nécessairement quasi-compacte) et $u$ un morphisme fini.
\end{corollary}

En effet, il résulte de (18.12.12) que $f$ est quasi-affine, donc quasi-projectif. On peut alors appliquer (8.12.8), où l'hypothèse de l'existence d'un $\sh{O}_Y$-Module ample peut en fait être remplacée par la seule hypothèse que $Y$ est quasi-compact et quasi-séparé: en effet, il résulte de (8.12.3) que l'existence de la factorisation annoncée dans (8.12.8) est une propriété locale sur $Y$, et qu'il suffit donc de la prouver lorsque $Y$ est affine.

\begin{remark}[18.12.14]
\label{IV.18.12.14}
On peut donner de (18.12.13) une démonstration analogue à celle de (18.12.12) n'utilisant pas (8.12.8) (mais utilisant (18.12.3), donc (18.5.11), qui lui-même se sert du « Main Theorem » sous sa forme locale (8.12.9)). Gardons en effet les notations de la démonstration de (18.12.12), et soit $\mathscr B$ la $\sh{O}_Y$-Algèbre quasi-cohérente, fermeture intégrale de $\sh{O}_Y$ dans $\mathscr A$ (II, 6.3.2). Si l'on pose $T=\Spec(\mathscr B)$, il suffit, en vertu de (8.12.3), de prouver que le $Y$-morphisme $g:X\to T$ correspondant à l'injection canonique $\mathscr B\to\mathscr A$ est une immersion, et pour cela il suffit de montrer (17.9.1) que $g$ est étale et radiciel. Avec les notations de (18.12.12), on peut supposer que $Y=\Spec(C)$ et $Y'=\Spec(C')$ sont affines, $C'$ étant une $C$-algèbre étale (17.3.2); donc $\mathscr A=\widetilde A$, où $A$ est une $C$-algèbre, $\mathscr A'=\widetilde{A'}$ avec $A'=A\otimes_C C'$, et $\mathscr B=\widetilde B$, où $B$ est la fermeture intégrale de $C$ dans $A$. L'algèbre $A'$ est isomorphe au produit $A'_1\times A'_2$, où $A'_1$ est une $C'$-algèbre finie. Il suffira de prouver que $B'=B\otimes_C C'$, qui

\oldpage[IV-4]{184}

s'identifie par platitude à une sous-$C'$-algèbre de $A'$, contient $A'_1$; en effet, $B'$ se décomposera alors en produit $A'_1\times A''_2$, où $A''_2$ est une sous-$C'$-algèbre de $A'_2$, et si l'on pose $T''_2=\Spec(A''_2)$, $T'=T\times_Y Y'$ sera somme de $Z'_1=\Spec(A'_1)$ et de $T''_2$, et $g'|X'_1$ sera un isomorphisme de $X'_1$ sur $Z'_1$, ce qui permettra de conclure comme dans (18.12.12).

Pour prouver que $B'$ contient $A'_1$, il suffit évidemment de prouver la proposition suivante, qui étend partiellement (6.14.4) lorsqu'on ne fait plus d'hypothèses noethériennes:
\end{remark}

\begin{proposition}[18.12.15]
\label{IV.18.12.15}
Soient $C$ un anneau, $C'$ une $C$-algèbre étale, $A$ une $C$-algèbre, $B$ la fermeture intégrale de $C$ dans $A$; on pose $A'=A\otimes_C C'$, $B'=B\otimes_C C'$, $B'$ s'identifiant à une sous-algèbre de $A'$; alors $B'$ est la fermeture intégrale de $C'$ dans $A'$.
\end{proposition}

En considérant la famille filtrante croissante des sous-$C$-algèbres de type fini de $A$ et raisonnant comme dans (6.14.4, II)), on peut d'abord supposer que $A$ est une $C$-algèbre de type fini, donc de la forme $E/\mathfrak J$, où $E=C[T_1,\ldots,T_r]$ et $\mathfrak J$ est un idéal de $E$; on a alors $A'=E'/\mathfrak JE'$, où $E'=C'[T_1,\ldots,T_r]$. Soit $(\mathfrak J_\lambda)$ la famille filtrante croissante des idéaux de type fini de $E$ contenus dans $\mathfrak J$, de sorte que $A$ est limite inductive des $E/\mathfrak J_\lambda$; si $B_\lambda$ est la fermeture intégrale de $C$ dans $E/\mathfrak J_\lambda$, $B$ est limite inductive des $B_\lambda$, comme le montre le raisonnement de (5.13.4). De même la fermeture intégrale de $C'$ dans $E'/\mathfrak JE'$ est limite inductive de la fermeture intégrale de $C'$ dans $E'/\mathfrak J_\lambda E'$; si on prouve que cette dernière est égale à $B_\lambda\otimes_C C'$, il en résultera que
\[
  B'=\varinjlim_\lambda(B_\lambda\otimes_C C')
\]
sera la fermeture intégrale de $C'$ dans $A'$. On est ainsi ramené à prouver la proposition lorsque $A$ est une $C$-algèbre de présentation finie.

Montrons ensuite qu'on peut se ramener au cas où $C$ est noethérien. En effet, $C$ est réunion filtrante de ses sous-$\mathbf Z$-algèbres de type fini $C_\alpha$, donc il résulte de (17.7.8) qu'il existe un indice $\alpha$ et une $C_\alpha$-algèbre étale $C'_\alpha$ tels que $C'=C'_\alpha\otimes_{C_\alpha}C$; en outre $C'$ est limite inductive des $C'_\beta=C'_\alpha\otimes_{C_\alpha}C_\beta$ pour $\beta\geq\alpha$. On peut en outre supposer, puisque $A$ est une $C$-algèbre de présentation finie, que $A=A_\alpha\otimes_{C_\alpha}C$, où $A_\alpha$ est une $C_\alpha$-algèbre de type fini, et $A$ est limite inductive des $A_\beta=A_\alpha\otimes_{C_\alpha}C_\beta$ pour $\beta\geq\alpha$ (8.9.1). Le raisonnement de (5.13.4) montre alors que $B$ est la limite inductive des fermetures intégrales $B_\beta$ des $C_\beta$ dans $A_\beta$. De même $A'$ est limite inductive des
\[
  A'_\beta=A_\beta\otimes_{C_\beta}C'_\beta=A_\alpha\otimes_{C_\alpha}C'_\beta
\]
pour $\beta\geq\alpha$, donc le même raisonnement que ci-dessus montre qu'il suffira de prouver que $B'_\alpha=B_\alpha\otimes_{C_\alpha}C'_\alpha$ est la fermeture intégrale de $C'_\alpha$ dans $A'_\alpha$.

Une fois ramené au cas où $C$ est noethérien, la proposition devient un cas particulier de (6.14.4). Il faut toutefois observer que, lorsque $C'$ est supposé étale sur $C$ noethérien, la démonstration de (6.14.4) n'exige plus le délicat théorème (6.14.1). En effet, les réductions successives de la preuve de (6.14.4) ramènent sans utiliser (6.14.1), au cas où $C$ est intègre, $A$ le corps des fractions de $C$. Comme alors $B$ est un anneau normal et que le morphisme $\Spec(B')\to\Spec(B)$ est étale, il résulte de (17.5.7) que $B'$ est un anneau normal; le raisonnement se termine en faisant appel seulement au lemme élémentaire (6.14.1.1), mais non à la partie difficile de la preuve de (6.14.1).

\begin{proposition}[18.12.16]
\label{IV.18.12.16}
Soient $g:Y\to S$ un morphisme quasi-compact, $f:X\to Y$ un morphisme séparé et quasi-fini. Pour tout $\sh{O}_Y$-Module inversible $\mathscr L$, ample relativement à $S$ (II, 4.6.1), le $\sh{O}_X$-Module $f^*(\mathscr L)$ est ample relativement à $S$.
\end{proposition}

En effet, en vertu de (18.12.12), le morphisme $f$ est quasi-affine, donc (II, 5.1.6) le $\sh{O}_X$-Module $\sh{O}_X$ est ample pour $f$. On en déduit (II, 4.6.13, (ii)) que
\[
  \sh{O}_X\otimes_{\sh{O}_X}f^*(\mathscr L^{\otimes n})=f^*(\mathscr L^{\otimes n})
\]
est ample relativement à $S$ pour un $n$ assez grand. Mais comme $f^*(\mathscr L^{\otimes n})=(f^*(\mathscr L))^{\otimes n}$, cela entraîne que $f^*(\mathscr L)$ est ample relativement à $S$ (II, 4.6.9, (i)).

\begin{corollary}[18.12.17]
\label{IV.18.12.17}
Soient $Z$ un schéma affine, $h:X\to Z$ un morphisme de type fini, $\mathscr L$ un $\sh{O}_X$-Module inversible. Avec les notations de (II, 4.5.2), les conditions de (II, 4.5.2) (qui équivalent au fait que $\mathscr L$ est ample) sont aussi équivalentes à chacune des suivantes:
\begin{enumerate}
  \item[$b'')$] $h$ est séparé, on a $G(\varepsilon)=X$ et le morphisme canonique $u:G(\varepsilon)\to\operatorname{Proj}(S)$ a ses fibres finies et discrètes.
  \item[$b''')$] On a $G(\varepsilon)=X$ et le morphisme canonique $u$ est radiciel.
\end{enumerate}
\end{corollary}

Si $Z=\Spec(A)$, l'anneau $S$ est canoniquement muni d'une structure de $A$-algèbre graduée, donc on a un morphisme structural $g:\operatorname{Proj}(S)\to Z$ qui est séparé (II, 2.4.2), et on a $h=g\circ u$. Comme tout morphisme radiciel est séparé (I, 1.8.7.1), $b''')$ entraîne que $h$ est séparé, donc entraîne $b'')$. Comme la condition $b)$ de (II, 4.5.2) entraîne évidemment $b''')$, il reste à voir que $b'')$ entraîne que $\mathscr L$ est ample. Puisque $h=g\circ u$ est de type fini par hypothèse et que $g$ est séparé, $u$ est aussi de type fini (I, 6.3.4), donc l'hypothèse $b'')$ entraîne que $u$ est quasi-fini et séparé (I, 5.5.1). Pour prouver que $\mathscr L$ est ample, nous allons appliquer (18.12.16). Posons $Y=\operatorname{Proj}(S)$. Comme $X$ est quasi-compact, il existe un entier $n>0$ et un nombre fini d'éléments $f_j\in S_n$ tels que les images réciproques $u^{-1}(D_+(f_j))$ recouvrent $X$; on peut donc considérer $u$ comme un morphisme quasi-fini et séparé de $X$ dans l'ouvert
\[
  Y'=\bigcup_j D_+(f_j)
\]
de $Y$. Considérons alors le $\sh{O}_{Y'}$-Module inversible
\[
  \mathscr L'=\sh{O}_Y(n)|Y'
\]
(II, 2.5.8); il est très ample relativement à $Z$ (II, 4.4.3), et $u^*(\mathscr L')$ n'est autre par définition que $\mathscr L^{\otimes n}$

\oldpage[IV-4]{185}

(II, 3.7.9, (ii)). En vertu de (18.12.16), $\mathscr L^{\otimes n}$ est donc ample relativement à $Z$ (ce qui équivaut à dire qu'il est ample (II, 4.6.6)), et par suite il en est de même de $\mathscr L$ (II, 4.5.6, (i)).

Nous laissons au lecteur le soin d'énoncer de même des conditions équivalentes à celles de (II, 4.6.3) pour les $\sh{O}_X$-Modules relativement amples.

\section{Immersions régulières et platitude normale}
\label{section:IV.19}

Le présent paragraphe est consacré, d'une part à une étude des immersions régulières (16.9.2) $Y\to X$ de préschémas, notamment dans le cas où $X$ et $Y$ sont plats sur un même préschéma $S$; d'autre part, il donne des compléments sur les suites $M$-régulières et la platitude normale, en généralisant notamment des résultats de platitude établis par H. Hironaka dans sa théorie de la résolution des singularités [35].

\subsection{Propriétés des immersions régulières.}
\EGAReaderTarget{ega4.c04.subsec.19-1}{subsection:IV.19.1}

\begin{proposition}[19.1.1]
\label{IV.19.1.1}
Soient $X$ un préschéma localement noethérien, $j:Y\to X$ une immersion, $y$ un point régulier de $Y$. Pour que $j$ soit une immersion régulière au point $y$, il faut et il suffit que $X$ soit régulier au point $j(y)$.
\end{proposition}

Compte tenu de (16.9.10), on est ramené à prouver le

\begin{corollary}[19.1.2]
\EGAReaderTarget{ega4.c04.cor.19-1-2}{IV.19.1.2}
Soient $A$ un anneau local noethérien, $\mathfrak J$ un idéal de $A$; on suppose que l'anneau quotient $B=A/\mathfrak J$ soit régulier. Pour que $A$ soit régulier, il faut et il suffit que l'idéal $\mathfrak J$ soit régulier.
\end{corollary}

En effet, si $A$ est régulier, on sait ($0$, 17.1.9) que $\mathfrak J$ est engendré par une partie d'un système régulier de paramètres de $A$, donc $\mathfrak J$ est régulier. Inversement, si $\mathfrak J$ est régulier, et si $(x_i)_{1\leq i\leq r}$ est une suite $A$-régulière engendrant $\mathfrak J$, $(x_i)$ fait partie d'un système de paramètres de $A$ ($0$, 16.4.1), donc on déduit de ($0$, 17.1.7) que $A$ est régulier.

\begin{definition}[19.1.3]
\EGAReaderTarget{ega4.c04.def.19-1-3}{IV.19.1.3}
Soient $X$ un préschéma, $Y$ un sous-préschéma de $X$, $U$ un ouvert de $X$ contenant $Y$ et tel que $Y$ soit défini par un Idéal $\mathscr J$ de $\sh{O}_U$; on suppose que $\mathscr J/\mathscr J^2$ soit un $(\sh{O}_U/\mathscr J)$-Module localement libre (ce qui est le cas si $Y$ est quasi-régulièrement immergé dans $X$). Pour tout $y\in Y$, on appelle \emph{codimension transversale de $Y$ dans $X$ au point $y$} et on note $\operatorname{codim}_y^*(Y,X)$ le rang du $(\sh{O}_{U,y}/\mathscr J_y)$-module libre $\mathscr J_y/\mathscr J_y^2$. Si $f:Z\to X$ est une immersion d'image $Y$, on appelle de même \emph{codimension transversale de $f$ au point $z\in Z$} la codimension transversale dans $X$ au point $f(z)$ du sous-préschéma $Y$.
\end{definition}

\begin{proposition}[19.1.4]
\label{IV.19.1.4}
Soient $X$ un préschéma localement noethérien, $Y$ un sous-préschéma de $X$ régulièrement immergé; alors, pour tout $y\in Y$, on a
\begin{equation}
  \operatorname{codim}_y^*(Y,X)=\operatorname{codim}_y(Y,X).
  \tag{19.1.4.1}\label{IV.19.1.4.1}
\end{equation}
\end{proposition}

En vertu de (5.1.3.2), $\operatorname{codim}_y(Y,X)$ est égal à la plus petite valeur des dimensions des anneaux locaux $A=\sh{O}_{X,z}$, où $z$ est point générique d'une composante

\oldpage[IV-4]{186}

irréductible de $Y$ contenant $y$; comme un tel point $z$ est contenu dans tout voisinage de $y$ dans $X$, $\mathscr J_z/\mathscr J_z^2$ est un $A$-module libre de rang $n=\operatorname{codim}_y^*(Y,X)$; tout revient à voir que $\dim(A)=n$. Or, puisque $z$ est point maximal du sous-préschéma $Y$ défini par $\mathscr J$, on a $\dim(\sh{O}_{X,z}/\mathscr J_z)=0$, donc $\mathscr J_z$ est un \emph{idéal de définition} de l'anneau local noethérien $A$; en outre, l'hypothèse que $\mathscr J$ est quasi-régulier entraîne que $\mathscr J_z^m/\mathscr J_z^{m+1}$ est un $(A/\mathscr J_z)$-module libre de rang $\binom{m+n-1}{n-1}$; si $r$ est la longueur de l'anneau artinien $A/\mathscr J_z$, la longueur de $\mathscr J_z^m/\mathscr J_z^{m+1}$ est donc $r\binom{m+n-1}{n-1}$; le polynôme de Hilbert-Samuel de $A$ pour la filtration $\mathscr J_z$-préadique est donc bien de degré $n$, ce qui prouve la proposition ($0$, 16.2.3).

En vertu de (19.1.4), on dira désormais « codimension » au lieu de « codimension transversale » lorsqu'il s'agira d'un sous-préschéma $Y$ régulièrement immergé dans $X$ (\emph{même lorsque $X$ n'est pas localement noethérien}).

\begin{proposition}[19.1.5]
\label{IV.19.1.5}
\begin{enumerate}
  \item[(i)] Pour qu'une immersion $j:Y\to X$ soit ouverte, il faut et il suffit qu'elle soit régulière et de codimension $0$ en tout point.
  \item[(ii)] Soient $f:Y\to X$ une immersion régulière (resp. quasi-régulière), $g:X'\to X$ un morphisme plat, $Y'=Y\times_X X'$; alors le morphisme $f'=f_{(X')}:Y'\to X'$ est une immersion régulière (resp. quasi-régulière), et pour tout $y'\in Y'$, si $y$ est la projection de $y'$ dans $Y$, la codimension de $f'$ au point $y'$ est égale à la codimension de $f$ au point $y$. Inversement, si $f$ est une immersion et $g:X'\to X$ un morphisme fidèlement plat et quasi-compact, alors, si $f'$ est une immersion quasi-régulière, il en est de même de $f$.
  \item[(iii)] Si $f:Y\to X$ et $g:Z\to Y$ sont des immersions régulières, alors $f\circ g:Z\to X$ est une immersion régulière, et pour tout $z\in Z$, la codimension de $f\circ g$ au point $z$ est la somme de la codimension de $g$ au point $z$ et de la codimension de $f$ au point $g(z)$. De plus, la suite d'homomorphismes canoniques (16.2.7.1)
\begin{equation}
  0\longrightarrow g^*(\sh{N}_{Y/X})\longrightarrow\sh{N}_{Z/X}\longrightarrow\sh{N}_{Z/Y}\longrightarrow0
  \tag{19.1.5.1}\label{IV.19.1.5.1}
\end{equation}
est exacte, et pour tout $z\in Z$, il existe un voisinage de $z$ dans $Z$ dans lequel les restrictions des homomorphismes de cette suite forment une suite scindée.
  \item[(iv)] Soient $X$ un préschéma localement noethérien, $f:Y\to X$ et $g:Z\to Y$ deux immersions, $z$ un point de $Z$. Les conditions suivantes sont équivalentes:
  \begin{enumerate}
    \item[a)] $g$ est une immersion régulière au point $z$ et $f$ une immersion régulière au point $g(z)$.
    \item[b)] $g$ et $f\circ g$ sont des immersions régulières au point $z$.
    \item[c)] $f\circ g$ est une immersion régulière au point $z$, $f$ est une immersion régulière au point $g(z)$, et la suite d'homomorphismes canoniques (19.1.5.1), restreinte à un voisinage ouvert assez petit de $z$ dans $Z$, est exacte et scindée.
  \end{enumerate}
  \item[(v)] Soient $X$ un préschéma localement noethérien, $f:Y\to X$, $g:Z\to Y$ deux immersions, $z$ un point de $Z$; supposons que $y=g(z)$ soit un point régulier de $Y$ et $x=f(g(z))$ un point régulier

\oldpage[IV-4]{187}

de $X$ (ce qui entraîne, par (19.1.1), que $f$ est une immersion régulière au point $g(z)$); alors, pour que $f\circ g$ soit une immersion régulière au point $z$, il faut et il suffit que $g$ soit une immersion régulière au point $z$.
\end{enumerate}
\end{proposition}

L'assertion (i) est une autre manière d'exprimer (16.1.10). Pour démontrer la première assertion de (ii), on peut se borner au cas où $Y$ est un sous-préschéma fermé de $X$ défini par un Idéal régulier (resp. quasi-régulier) $\mathscr J$; alors $Y'$ est le sous-préschéma fermé de $X'$ défini par l'Idéal $g^*(\mathscr J)\sh{O}_{X'}$ (I, 4.4.5) qui s'identifie ici à $g^*(\mathscr J)$ puisque $g$ est plat; en remplaçant $X$ par un ouvert affine assez petit, on peut supposer que $\mathscr J$ admet une suite régulière (resp. quasi-régulière) de générateurs $f_i$ (sections de $\sh{O}_X$ au-dessus de $X$). Si la suite $(f_i)$ est régulière, il en est de même de la suite $(f_i\otimes1)$ (qui engendre $g^*(\mathscr J)\sh{O}_{X'}$) en vertu de ($0$, 15.2.5); le résultat analogue pour les suites quasi-régulières découle aussitôt de la platitude de $g$ et du critère (16.9.4). De même, si $g$ est fidèlement plat et quasi-compact, et si $g^*(\mathscr J)$ est quasi-régulier, $\mathscr J$ est de type fini en vertu de (2.5.2); par platitude, on a $g^*(\mathscr J)/g^*(\mathscr J)^2=g^*(\mathscr J/\mathscr J^2)$, donc, puisque $g^*(\mathscr J/\mathscr J^2)$ est localement libre par hypothèse (16.9.4), il en est de même de $\mathscr J/\mathscr J^2$ (2.5.2). Enfin, la condition (iii) de (16.9.4) relative à $g^*(\mathscr J)$ entraîne la même condition pour $\mathscr J$ en vertu de la fidèle platitude de $g$ (2.2.7). On conclut donc de (16.9.4) que $\mathscr J$ est quasi-régulier.

Pour prouver (iii), on peut de même supposer que $Y$ et $Z$ sont des sous-préschémas fermés de $X$, définis respectivement par des Idéaux $\mathscr J$, $\mathscr K$ de $\sh{O}_X$ tels que $\mathscr J\subset\mathscr K$; de plus, on peut supposer qu'il existe des sections $f_i$ ($1\leq i\leq m$) de $\mathscr J$ au-dessus de $X$ engendrant $\mathscr J$ et telles que pour $1\leq i\leq m$, l'endomorphisme de $\sh{O}_X/(\sum_{j=1}^{i-1}f_j\sh{O}_X)$ défini par $f_i$ soit injectif, et des sections $\overline f_i$ ($m+1\leq i\leq n$) de $\mathscr K/\mathscr J$ au-dessus de $X$ engendrant $\mathscr K/\mathscr J$ et telles que pour $m+1\leq i\leq n$, l'endomorphisme de
\[
  (\sh{O}_X/\mathscr J)\big/\Bigl(\sum_{j=m+1}^{i-1}\overline f_j(\sh{O}_X/\mathscr J)\Bigr)
\]
défini par $\overline f_i$ soit injectif. Pour tout $z\in Z$, soit $U$ un voisinage ouvert de $z$ dans $X$ tel qu'il existe une section $f_i$ de $\mathscr K$ au-dessus de $U$, dont $\overline f_i|U$ soit la classe dans $\Gamma(U,\mathscr K/\mathscr J)$ pour $m+1\leq i\leq n$; alors
\[
  (\sh{O}_X/\mathscr J)\big/\Bigl(\sum_{j=m+1}^{i-1}\overline f_j(\sh{O}_X/\mathscr J)\Bigr),
\]
restreint à $U$ est isomorphe à
\[
  \sh{O}_U\big/\Bigl((\mathscr J|U)+\sum_{j=m+1}^{i-1}f_j\sh{O}_U\Bigr)
  =\sh{O}_U\big/\Bigl(\sum_{j=1}^{i-1}f_j\sh{O}_U\Bigr),
\]
et on voit donc que la suite $(f_i)_{1\leq i\leq n}$ est régulière dans $\sh{O}_U$; comme elle engendre $\mathscr K|U$, cela prouve la première assertion de (iii). Pour prouver la dernière assertion de (iii), on peut se borner (avec les mêmes notations) au cas où les images $t_i$ des $f_i$ dans $\mathscr K/\mathscr K^2$ (pour $1\leq i\leq n$) forment une base de cet $(\sh{O}_X/\mathscr K)$-Module (16.9.5); pour la même raison, on peut supposer que les $t_i$ tels que $1\leq i\leq m$ sont les images canoniques des éléments d'une base du $(\sh{O}_X/\mathscr K)$-Module $g^*(\mathscr J/\mathscr J^2)$, et que les images canoniques $\overline t_i$ des $t_i$ dans $(\mathscr K/\mathscr J)/(\mathscr K/\mathscr J)^2$ pour $m+1\leq i\leq n$ forment une base de cet $(\sh{O}_X/\mathscr K)$-Module. Cela achève de prouver (iii).

Passons à la démonstration de (iv). Le fait que $a)$ implique $c)$ résulte de (iii). Prouvons que $c)$ implique $a)$. Posons encore $A=\sh{O}_{X,f(g(z))}$, de sorte que $\sh{O}_{Y,g(z)}=A/\mathfrak J$,

\oldpage[IV-4]{188}

$\sh{O}_{Z,z}=A/\mathfrak R$ avec $\mathfrak J\subset\mathfrak R$. Par hypothèse, $\mathfrak J$ et $\mathfrak R$ sont des idéaux réguliers de $A$ et on a une suite exacte scindée
\[
  0\longrightarrow\mathfrak J/\mathfrak J\mathfrak R\longrightarrow\mathfrak R/\mathfrak R^2\longrightarrow(\mathfrak R/\mathfrak J)/(\mathfrak R/\mathfrak J)^2\longrightarrow0
\]
(cf. (16.2.7)). Cela entraîne qu'il existe des éléments $f_j$ ($1\leq j\leq r+s$) de $\mathfrak R$ dont les images $f'_j$ dans $\mathfrak R/\mathfrak R^2$ forment une base de cet $(A/\mathfrak R)$-module, et tels en outre que les $f_j$ d'indice $j\leq r$ appartiennent à $\mathfrak J$ et sont tels que leurs images dans $\mathfrak J/\mathfrak J\mathfrak R$ forment une base de cet $(A/\mathfrak R)$-module. Les $f_j$ pour $1\leq j\leq r+s$ engendrent donc $\mathfrak R$, et les $f_j$ pour $1\leq j\leq r$ engendrent $\mathfrak J$ puisque $A$ est noethérien (Bourbaki, \emph{Alg. comm.}, chap. II, \textsection~3, n\textsuperscript{o}~2, prop. 5); comme $\mathfrak R$ est un idéal régulier, il résulte de (16.9.5) que la suite $(f_j)_{1\leq j\leq r+s}$ est régulière. Par définition, les images de $f_{r+1},\ldots,f_{r+s}$ dans $\mathfrak R/\mathfrak J$ forment donc une suite régulière dans cet $(A/\mathfrak J)$-module, ce qui achève de prouver que $c)$ entraîne $a)$ dans (iv).

Enfin, compte tenu de (16.9.10), il est immédiat que (v), ainsi que l'équivalence de $a)$ et $b)$ dans (iv), résultent du

\begin{corollary}[19.1.6]
\label{IV.19.1.6}
Soient $A$ un anneau local noethérien, $\mathfrak J$, $\mathfrak R$ deux idéaux de $A$ contenus dans son idéal maximal $\mathfrak m$ et tels que $\mathfrak J\subset\mathfrak R$, $A'=A/\mathfrak J$, $\mathfrak R'=\mathfrak R/\mathfrak J$. Alors:
\begin{enumerate}
  \item[(i)] Si $\mathfrak J$ et $\mathfrak R'$ sont des idéaux réguliers, il en est de même de $\mathfrak R$.
  \item[(ii)] Si $\mathfrak R$ et $\mathfrak R'$ sont des idéaux réguliers, il en est de même de $\mathfrak J$.
  \item[(iii)] Supposons les anneaux $A$ et $A'=A/\mathfrak J$ réguliers. Pour que l'idéal $\mathfrak R$ soit régulier, il faut et il suffit que $\mathfrak R'$ le soit.
\end{enumerate}
\end{corollary}

L'assertion (i) est un cas particulier de (19.1.5, (iii)). Pour démontrer (ii), considérons l'homomorphisme surjectif canonique $u:\mathfrak R/\mathfrak R^2\to\mathfrak R'/\mathfrak R'^2=\mathfrak R/(\mathfrak R^2+\mathfrak J)$. Comme par hypothèse $\mathfrak R/\mathfrak R^2$ et $\mathfrak R'/\mathfrak R'^2$ sont des $(A/\mathfrak R)$-modules libres, dont nous désignerons les rangs respectifs par $p+q$ et $q$, le noyau de $u$ est un $(A/\mathfrak R)$-module \emph{projectif} (Bourbaki, \emph{Alg.}, chap. II, 3\textsuperscript{e} éd., \textsection~2, n\textsuperscript{o}~2, prop. 4), donc \emph{libre} de rang $p$ puisque $A/\mathfrak R$ est un anneau local noethérien ($0_{\mathrm{III}}$, 10.1.3); autrement dit, il existe une base de $\mathfrak R/\mathfrak R^2$ dont les $p$ premiers éléments forment une base de $(\mathfrak R^2+\mathfrak J)/\mathfrak R^2$. Cela signifie encore (en vertu du lemme de Nakayama) qu'il existe un système de générateurs $(f_i)_{1\leq i\leq p+q}$ de $\mathfrak R$, formant une suite régulière dans $A$, telle que la suite $(f_i)_{1\leq i\leq p}$ soit formée d'éléments de $\mathfrak J$; il s'agit de montrer que cette dernière suite engendre $\mathfrak J$. Pour cela, désignons par $\mathfrak J'\subset\mathfrak J$ l'idéal qu'elle engendre; en considérant l'anneau $A/\mathfrak J'$, on voit qu'on est ramené à prouver le lemme suivant:

\begin{lemma}[19.1.6.1]
\label{IV.19.1.6.1}
Soient $B$ un anneau local noethérien d'idéal maximal $\mathfrak n$, $(g_j)_{1\leq j\leq q}$ une suite régulière dans $B$ d'éléments de $\mathfrak n$, $\mathfrak b$ l'idéal qu'elle engendre, $\mathfrak c\subset\mathfrak b$ un idéal tel que les images canoniques des $g_j$ dans $B/\mathfrak c$ forment encore une suite régulière dans cet anneau. Alors on a $\mathfrak c=0$.
\end{lemma}

Le lemme étant trivial pour $q=0$, raisonnons par récurrence sur $q$. L'hypothèse de récurrence, appliquée à l'anneau $B/g_1B$ et aux images canoniques des $g_j$ ($2\leq j\leq q$)

\oldpage[IV-4]{189}

et de $\mathfrak c$ dans cet anneau quotient, montre que $\mathfrak c\subset g_1B$. Soit $\mathfrak d$ l'idéal de $B$ formé des $x\in B$ tels que $g_1x\in\mathfrak c$; on a donc $\mathfrak c=g_1\mathfrak d$; mais comme par hypothèse $g_1$ est régulier dans $B/\mathfrak c$, on a nécessairement $\mathfrak d=\mathfrak c$, donc $\mathfrak c=g_1\mathfrak c$. Comme $\mathfrak c$ est de type fini et que $g_1$ est contenu dans le radical de $B$, le lemme de Nakayama montre que $\mathfrak c=0$.

Démontrons enfin l'assertion (iii) de (19.1.6). En vertu de (i) et de (19.1.2), il suffit de voir que si $\mathfrak R$ est régulier, il en est de même de $\mathfrak R'$. Notons que $\mathfrak J$ est régulier (19.1.2), et raisonnons par récurrence sur le rang $q$ du $(A/\mathfrak J)$-module libre $\mathfrak J/\mathfrak J^2$. Si $q=0$, on a $\mathfrak J=0$ par le lemme de Nakayama et l'assertion est triviale. Comme $A$ et $A/\mathfrak J$ sont réguliers, on sait ($0$, 17.1.9) que $\mathfrak J$ est engendré par une partie à $q$ éléments d'un système régulier de paramètres de $A$, donc, si $f$ est un des éléments de ce système de générateurs de $\mathfrak J$, $A_1=A/fA$ est régulier et $f\notin\mathfrak m^2$ ($0$, 17.1.8). Soient $\mathfrak J_1$, $\mathfrak R_1$ les images canoniques de $\mathfrak J$, $\mathfrak R$ dans $A_1$; on a $\mathfrak J_1\subset\mathfrak R_1$, $A_1/\mathfrak J_1=A/\mathfrak J$, donc $A_1/\mathfrak J_1$ est régulier et par suite $\mathfrak J_1$ est un idéal régulier (19.1.2). Montrons que $\mathfrak R_1$ est un idéal régulier dans $A_1$; comme $\mathfrak R$ est un idéal régulier dans $A$, il suffit de montrer que $f$ fait partie d'un système régulier de générateurs de $\mathfrak R$, et pour cela (16.9.5) il suffit de montrer que l'image de $f$ dans $\mathfrak R/(\mathfrak R^2+\mathfrak m\mathfrak R)=\mathfrak R/\mathfrak m\mathfrak R$ fait partie d'une base de cet $(A/\mathfrak m)$-espace vectoriel, autrement dit que $f\notin\mathfrak m\mathfrak R$, ce qui résulte en effet de $f\notin\mathfrak m^2$. Alors, comme $\mathfrak J_1/\mathfrak J_1^2$ est un $(A_1/\mathfrak J_1)$-module libre de rang $q-1$, l'hypothèse de récurrence montre que $\mathfrak R'_1=\mathfrak R_1/\mathfrak J_1$ est un idéal régulier de $A_1/\mathfrak J_1$, et comme $\mathfrak R'$ s'identifie à $\mathfrak R'_1$ dans l'isomorphisme canonique entre $A/\mathfrak J$ et $A_1/\mathfrak J_1$, $\mathfrak R'$ est régulier. C.Q.F.D.

\begin{remarks}[19.1.7]
\label{IV.19.1.7}
\begin{enumerate}
  \item[(i)] Nous ignorons si le composé de deux immersions quasi-régulières est quasi-régulière.
  \item[(ii)] Dans un espace annelé en anneaux locaux $X$, pour tout Idéal $\mathscr J$ de $\sh{O}_X$, $\sh{O}_X/\mathscr J$ est de type fini et a donc un support fermé dans $X$ ($0_{\mathrm I}$, 5.2.2). On peut par suite définir comme dans (I, 4.1.3 et 4.2.1) les notions de sous-espace annelé (localement fermé) de $X$ défini par un Idéal $\mathscr J$ d'un faisceau $\sh{O}_U$, où $U$ est un ouvert de $X$, et la notion d'immersion. Les définitions d'immersions quasi-régulières et régulières, de codimension transversale s'appliquent alors sans modification, et les résultats de (19.1.4, (i) à (iv)) sont encore valables, en supposant dans la seconde assertion de (i) et dans (iv) que le faisceau $\sh{O}_X$ soit cohérent et les anneaux locaux $\sh{O}_{X,x}$ noethériens; quant à (ii), il faut \emph{définir} $Y'$ comme le sous-espace annelé défini par le $\sh{O}_{X'}$-Idéal $g^*(\mathscr J)$; les démonstrations sont inchangées.
\end{enumerate}
\end{remarks}

\begin{env}[19.1.8]
\label{IV.19.1.8}
Nous avons déjà signalé (notamment dans (16.9.6, (ii))) que dans les anneaux non noethériens, la notion de « suite régulière » d'éléments de l'anneau ne possède pas les propriétés qui la rendraient utilisable. De récentes recherches semblent montrer que la notion qui doit dans ce cas se substituer à celle de suite régulière est celle de suite $\mathbf f=(f_1,\ldots,f_n)$ ayant la propriété que le complexe de l'algèbre extérieure $\mathrm K_0(\mathbf f,M)$ (III, 1.1.3) ait \emph{son homologie nulle en degrés $>0$} (cf. (III, 1.1.4)); voir en particulier A. Grothendieck, Séminaire de géométrie algébrique, 1966, à paraître prochainement.
\end{env}

\subsection{Immersions transversalement régulières.}
\EGAReaderTarget{ega4.c04.subsec.19-2}{subsection:IV.19.2}

\oldpage[IV-4]{190}

\begin{definition}[19.2.1]
\EGAReaderTarget{ega4.c04.def.19-2-1}{IV.19.2.1}
Soient $u:X\to S$ un morphisme de préschémas, $\mathscr F$ un $\sh{O}_X$-Module quasi-cohérent. On dit qu'une suite $(f_i)_{1\leq i\leq n}$ de sections de $\mathscr F$ au-dessus de $X$ est \emph{transversalement $\mathscr F$-régulière relativement à $S$} (ou relativement à $u$) si la suite $(f_i)$ est $\mathscr F$-régulière et telle que les $\sh{O}_X$-Modules
\[
  \mathscr F\big/\Bigl(\sum_{j\leq i}f_j\mathscr F\Bigr)
\]
soient $S$-plats pour $0\leq i\leq n$. On dit qu'un Idéal quasi-cohérent $\mathscr J$ de $\sh{O}_X$ est \emph{transversalement régulier relativement à $S$} (ou relativement à $u$) en un point $x\in\operatorname{Supp}(\sh{O}_X/\mathscr J)$ s'il existe un voisinage ouvert $U$ de $x$ et une suite finie $(f_i)$ de sections de $\mathscr J$ au-dessus de $U$ qui est transversalement $\sh{O}_U$-régulière relativement à $S$ et qui engendre $\mathscr J|U$.

Dans une $A$-algèbre $B$, on dit qu'un idéal $\mathfrak J$ est \emph{transversalement régulier relativement à $A$} si $\widetilde{\mathfrak J}$ est transversalement régulier relativement à $S=\Spec(A)$ en tout point de $V(\mathfrak J)$ dans $\Spec(B)$.
\end{definition}

\begin{definition}[19.2.2]
\EGAReaderTarget{ega4.c04.def.19-2-2}{IV.19.2.2}
Soient $S$ un préschéma, $X$, $Y$ deux $S$-préschémas, $u:Y\to X$ un $S$-morphisme qui est une immersion. On dit que $u$ est une \emph{immersion transversalement régulière relativement à $S$ en un point $y\in Y$} s'il existe un voisinage $V$ de $u(y)$ dans $X$ tel que le sous-préschéma induit sur $u(Y)\cap V$ par le sous-préschéma de $X$ associé à $u$ soit défini par un Idéal de $\sh{O}_U$ transversalement régulier relativement à $S$ au point $u(y)$.

Il est clair que l'ensemble des points de $Y$ où une immersion est transversalement régulière relativement à $S$ est \emph{ouvert} dans $Y$; s'il est égal à $Y$, on dit simplement que $u$ est une \emph{immersion transversalement régulière relativement à $S$}.

Si $S\to T$ est un morphisme \emph{plat}, alors, si $\mathscr J$ est un Idéal de $\sh{O}_X$ transversalement régulier relativement à $S$ en $x\in X$, il l'est aussi relativement à $T$ ($0_{\mathrm I}$, 6.2.1). Une $S$-immersion $Y\to X$ transversalement régulière en un point $y\in Y$ relativement à $S$ est donc aussi une $T$-immersion transversalement régulière relativement à $T$ en ce point.
\end{definition}

\begin{proposition}[19.2.3]
\label{IV.19.2.3}
Soient $S$ un préschéma, $X$, $Y$ deux $S$-préschémas, $u:Y\to X$ un $S$-morphisme qui est une immersion transversalement régulière relativement à $S$ en un point $y\in Y$. Alors, pour tout changement de base $S'\to S$, si l'on pose $X'=X\times_S S'$, $Y'=Y\times_S S'$, l'immersion $u'=u_{(S')}:Y'\to X'$ est transversalement régulière relativement à $S'$ en tout point $y'\in Y'$ au-dessus de $y$.
\end{proposition}

Cela résulte aussitôt de la définition et de ($0$, 15.1.15).

\begin{proposition}[19.2.4]
\label{IV.19.2.4}
Soient $S$ un préschéma, $X$, $Y$ deux $S$-préschémas, $u:Y\to X$ un $S$-morphisme qui est une immersion. On suppose, soit que $S$ et $X$ sont localement noethériens, soit que les morphismes structuraux $g:X\to S$, $h:Y\to S$ soient localement de présentation finie. Soient $y$ un point de $Y$, $s$ son image dans $S$. Les conditions suivantes sont équivalentes:
\begin{enumerate}
  \item[a)] $u$ est transversalement régulière relativement à $S$ au point $y$.
  \item[b)] Les morphismes $g$ et $h$ sont plats au voisinage des points $u(y)$ et $y$ respectivement, et si l'on pose $X_s=g^{-1}(s)$, $Y_s=h^{-1}(s)$, l'immersion $u_s:Y_s\to X_s$ est régulière au point $y$.
  \item[b$'$)] Les morphismes $g$ et $h$ sont plats au voisinage des points $u(y)$ et $y$ respectivement, et, pour tout changement de base $S'\to S$, si l'on pose $X'=X\times_S S'$, $Y'=Y\times_S S'$, l'immersion $u'=u_{(S')}:Y'\to X'$ est régulière en tout point de $Y'$ au-dessus de $y$.
  \item[c)] Le morphisme $h$ est plat au voisinage du point $y$, et l'immersion $u$ est régulière au voisinage du point $y$.

\oldpage[IV-4]{191}

  \item[c$'$)] Le morphisme $h$ est plat au voisinage du point $y$, et l'immersion $u$ est quasi-régulière au voisinage du point $y$.
\end{enumerate}
\end{proposition}

On a déjà prouvé dans (19.2.3) que $a)$ entraîne $b')$ sans hypothèse de finitude, et il est trivial que $b')$ entraîne $b)$. Montrons que $b)$ entraîne $c)$. On peut pour cela (la question étant locale sur $X$ et sur $S$) supposer que $Y$ est un sous-préschéma fermé de $X$ défini par un Idéal quasi-cohérent $\mathscr J$ de $\sh{O}_X$, et $u$ l'injection canonique. L'hypothèse que $\sh{O}_X/\mathscr J$ est $g$-plat entraîne que la suite
\[
  0\longrightarrow\mathscr J\otimes_{\sh{O}_s}\mathbf k(s)\longrightarrow\sh{O}_X\otimes_{\sh{O}_s}\mathbf k(s)\longrightarrow(\sh{O}_X/\mathscr J)\otimes_{\sh{O}_s}\mathbf k(s)\longrightarrow0
\]
est exacte ($0_{\mathrm I}$, 6.1.2), donc $Y_s$ est le sous-préschéma fermé de $X_s$ défini par l'Idéal $\mathscr J_s$ de $\sh{O}_{X_s}=\sh{O}_X\otimes_{\sh{O}_s}\mathbf k(s)$ qui s'identifie à $\mathscr J\otimes_{\sh{O}_s}\mathbf k(s)=\mathscr J/\mathfrak m_s\mathscr J$. Considérons une suite finie $(f_i)$ de sections de $\mathscr J$ au-dessus d'un voisinage de $y$ dans $X$, dont les images dans
\[
  \mathscr J_y/(\mathfrak m_s\mathscr J_y+\mathscr J_y^2)
\]
forment une base de ce $(\sh{O}_{X,y}/(\mathfrak m_s\sh{O}_{X,y}+\mathscr J_y))$-module (qui est libre en vertu de l'hypothèse $b)$). Puisque $X_s$ est localement noethérien, il résulte de (16.9.11), (16.9.5) et de l'hypothèse $b)$ que les images $(f_i)_s=f_i\otimes1$, sections de $\mathscr J_s=\mathscr J\otimes_{\sh{O}_s}\mathbf k(s)$ au-dessus d'un voisinage de $y$ dans $X_s$, forment une suite régulière engendrant la restriction de $\mathscr J_s$ à ce voisinage. Comme $\mathfrak m_s\sh{O}_{X,y}$ est contenu dans le radical de $\sh{O}_{X,y}$, et que dans les deux cas envisagés $\mathscr J_y$ est un idéal de type fini de $\sh{O}_{X,y}$, il résulte du lemme de Nakayama que les $(f_i)_y$ engendrent aussi $\mathscr J_y$; comme $\mathscr J$ est un Idéal de type fini de $\sh{O}_X$ dans les deux cas envisagés, il y a un voisinage $U$ de $y$ dans $X$ tel que les $f_i|U$ engendrent $\mathscr J|U$ ($0_{\mathrm I}$, 5.2.2). Le fait que $b)$ entraîne $c)$ est alors conséquence de ($0$, 15.1.16) lorsque $S$ et $X$ sont localement noethériens, et de (11.3.8) lorsque $g$ et $h$ sont localement de présentation finie; dans ce dernier cas, (11.3.8) prouve aussi l'équivalence de $c)$ et $c')$ (qui est triviale lorsque $X$ et $S$ sont localement noethériens ($0$, 15.1.11)). Enfin, si $c)$ est vérifiée, pour prouver $a)$, on peut encore se borner au cas où $Y$ est un sous-préschéma fermé de $X$ défini par $\mathscr J$, et par hypothèse, il y a une suite régulière $(f_i)$ de sections de $\mathscr J$ au-dessus d'un voisinage de $y$ dans $X$ telle que les $f_i$ engendrent $\mathscr J|U$ et que $h$ soit plat dans $U$; pour voir que $c)$ entraîne $a)$, il suffit encore d'appliquer ($0$, 15.1.6) lorsque $S$ et $X$ sont localement noethériens, et (11.3.8) lorsque $g$ et $h$ sont localement de présentation finie.

\begin{remarks}[19.2.5]
\label{IV.19.2.5}
\begin{enumerate}
  \item[(i)] Lorsque dans la condition $b)$, on supprime l'hypothèse que $h$ est plat au voisinage de $y$, la condition ainsi modifiée n'entraîne plus $a)$. Il suffit de prendre pour $S$ le spectre d'un anneau local qui n'est pas un corps, pour $s$ le point fermé de $S$, et $Y=X_s$.
  \item[(ii)] Supposons que $Y$ soit un sous-préschéma fermé de $X$ défini par un Idéal quasi-cohérent $\mathscr J$ de $\sh{O}_X$. On a prouvé au cours de la démonstration de (19.2.4) que, lorsque les conditions équivalentes de cette proposition sont remplies, alors, pour tout système $(f_i)$ de sections de $\mathscr J$ au-dessus d'un voisinage de $y$ dans $X$, tel que les images des $f_i$ dans $\mathscr J_y/(\mathfrak m_s\mathscr J_y+\mathscr J_y^2)$ forment une base de ce $(\sh{O}_{X,y}/(\mathfrak m_s\sh{O}_{X,y}+\mathscr J_y))$-module, il existe un voisinage ouvert $U$ de $y$ dans $X$ tel que les $f_i|U$ vérifient les conditions de la définition (19.2.1).
\end{enumerate}
\end{remarks}

\oldpage[IV-4]{192}

\begin{corollary}[19.2.6]
\label{IV.19.2.6}
Supposons, soit que $S$ et $X$ soient localement noethériens, soit que $g$ et $h$ soient localement de présentation finie. Supposons en outre que les fibres $X_s$ et $Y_s$ soient des préschémas réguliers aux points $u(y)$ et $y$ respectivement, et que les morphismes $g$ et $h$ soient plats au voisinage de $u(y)$ et $y$ respectivement. Alors l'immersion $Y\to X$ est transversalement régulière au point $x$.
\end{corollary}

En effet, il résulte de (19.1.1) que l'immersion $Y_s\to X_s$ est régulière au point $x$, et il suffit d'appliquer (19.2.4, $b)$).

\begin{proposition}[19.2.7]
\label{IV.19.2.7}
\begin{enumerate}
  \item[(i)] Toute $S$-immersion ouverte est transversalement régulière relativement à $S$.
  \item[(ii)] Soient $f:Y\to X$ une $S$-immersion, $g:S'\to S$ un morphisme et posons $X'=X_{(S')}$, $Y'=Y_{(S')}$, $f'=f_{(S')}:Y'\to X'$. Si $f$ est une immersion transversalement régulière relativement à $S$, $f'$ est une immersion transversalement régulière relativement à $S'$. La réciproque est vraie si $g$ est fidèlement plat, et si de plus $X$ et $Y$ sont localement de présentation finie sur $S$.
  \item[(iii)] Si $f:Y\to X$ et $g:Z\to Y$ sont des immersions transversalement régulières relativement à $S$, il en est de même de $f\circ g:Z\to X$.
  \item[(iv)] Soient $X$ un $S$-préschéma, $f:Y\to X$, $g:Z\to Y$ deux $S$-immersions. On suppose, soit que $S$ et $X$ soient localement noethériens, soit que $X$, $Y$ et $Z$ soient localement de présentation finie sur $S$. Alors, si $g$ et $f\circ g$ sont transversalement régulières au point $z\in Z$ relativement à $S$, $f$ est transversalement régulière au point $g(z)$ relativement à $S$.
  \item[(v)] Sous les conditions générales de (iv), supposons que $X$ et $Y$ soient $S$-plats, $Y_s$ régulier au point $g(z)$ et $X_s$ régulier au point $g(z)$; alors, si $f\circ g$ est transversalement régulière au point $z$, il en est de même de $g$.
\end{enumerate}
\end{proposition}

L'assertion (i) est immédiate, et la première assertion de (ii) résulte aussitôt de (19.2.3). Pour démontrer la seconde assertion de (ii), appliquons (19.2.4): $X'$ et $Y'$ étant localement de présentation finie sur $S'$, il résulte d'abord de ce critère que $X'$ et $Y'$ sont plats sur $S'$, donc $X$ et $Y$ sont plats sur $S$ (2.5.1); d'autre part, pour tout $s\in S$, si $s'\in S'$ est un point au-dessus de $s$, $Y'_{s'}\to X'_{s'}$ est par hypothèse une immersion régulière, et on déduit donc de (19.1.5, (ii)) qu'il en est de même de $Y_s\to X_s$, puisque ces préschémas sont localement noethériens et qu'il revient alors au même de dire que l'immersion $Y_s\to X_s$ est régulière ou quasi-régulière. Pour prouver (iii), on se ramène au cas où $Z$ est un sous-préschéma fermé de $Y$ et $Y$ un sous-préschéma fermé de $X$, et on forme comme dans (19.1.5, (iii)) un système de générateurs de l'Idéal de $\sh{O}_X$ définissant $Z$. Pour prouver (iv), on remarque que l'hypothèse jointe à (19.2.4) montre d'abord que $X$, $Y$ et $Z$ sont $S$-plats dans un voisinage de $z$, et il suffit en outre, si $s$ est l'image de $z$ dans $S$, de prouver que l'immersion $f_s:Y_s\to X_s$ est régulière au point $z$, sachant qu'il en est de même de $g_s:Z_s\to Y_s$ et de $f_s\circ g_s$; mais cela résulte de (19.1.5, (iv)). Enfin, pour (v), on raisonne de la même manière que pour (iv); l'hypothèse sur $f\circ g$ implique déjà par (19.2.4, $b)$) que $Z$ est $S$-plat dans un voisinage de $z$, et il en est de même de $X$ et $Y$ par hypothèse; on est alors ramené à prouver que $g_s$ est régulière au point $z$, sachant qu'il en est de même de $f_s\circ g_s$, ce qui découle de (19.1.5, (v)), vu les hypothèses faites sur $X_s$ et $Y_s$.

\oldpage[IV-4]{193}

\begin{proposition}[19.2.8]
\label{IV.19.2.8}
Soient $X$ un $S$-préschéma, $Y$ un sous-préschéma fermé de $X$ défini par un Idéal quasi-cohérent $\mathscr J$ de $\sh{O}_X$, $u:Y\to X$ l'injection canonique. Si $u$ est transversalement régulière relativement à $S$, les $\sh{O}_X$-Modules $\mathscr J^n$ et $\mathscr J^n/\mathscr J^m$ ($0\leq n<m$) sont $S$-plats aux points de $Y$. Pour tout changement de base $S'\to S$, si $X'=X_{(S')}$, $Y'=Y_{(S')}$ et $u'=u_{(S')}$, on a
\[
  \sh{G}r_n(u')=\sh{G}r_n(u)\otimes_{\sh{O}_S}\sh{O}_{S'}
\]
à isomorphisme près pour tout $n>0$.
\end{proposition}

En effet, $\sh{O}_X$ et $\sh{O}_X/\mathscr J$ sont par hypothèse $S$-plats aux points de $Y$, et $\mathscr J^n/\mathscr J^{n+1}$ est un $(\sh{O}_X/\mathscr J)$-module localement libre (16.9.3), donc est $S$-plat aux points de $Y$. La première assertion résulte alors de ($0_{\mathrm I}$, 6.1.2) et la seconde de (11.2.9.2).

La proposition suivante généralise et précise (17.16.1):

\begin{proposition}[19.2.9]
\label{IV.19.2.9}
Soient $f:X\to S$ un morphisme, $x$ un point de $X$, $s=f(x)$. Supposons que $f$ soit de présentation finie au point $x$, et soit un morphisme de Cohen-Macaulay au point $x$ (6.8.1). Alors il existe un sous-préschéma $Y$ de $X$ contenant $x$ et ayant les propriétés suivantes:
\begin{enumerate}
  \item[(i)] L'injection canonique $j:Y\to X$ est transversalement régulière relativement à $S$.
  \item[(ii)] Le morphisme $g=f\circ j:Y\to S$ est de Cohen-Macaulay (ce qui équivaut à dire, compte tenu de (i), que toutes les fibres $g^{-1}(g(y))$ de $g$ sont des préschémas de Cohen-Macaulay).
  \item[(iii)] Le point $x$ est point maximal de $Y_s=g^{-1}(s)$.
\end{enumerate}
\end{proposition}

Posons $X_s=f^{-1}(s)$; par hypothèse, l'anneau $\sh{O}_{X_s,x}$ est de Cohen-Macaulay; soit $(t_i)_{1\leq i\leq n}$ un système de paramètres de cet anneau, qui est donc une suite régulière ($0$, 16.5.7). Il y a un voisinage ouvert $U$ de $x$ dans $X$ et des sections $h_i$ ($1\leq i\leq n$) de $\sh{O}_X$ au-dessus de $U$ telles que les $t_i$ soient les images canoniques des $(h_i)_x$ dans l'anneau quotient $\sh{O}_{X_s,x}$ de $\sh{O}_{X,x}$. Soit $\mathscr J$ l'Idéal de $\sh{O}_U$ engendré par les $h_i$, et prenons pour $Y$ le sous-préschéma fermé de $U$ défini par $\mathscr J$. On peut supposer que $f|U$ est localement de présentation finie, et il en est donc de même de $g$ par définition de $Y$; il résulte alors de l'équivalence de $c)$ et $a)$ dans (11.3.8) que la suite $((h_i)_x)$ est régulière dans $\sh{O}_{X,x}$ et que $\sh{O}_{Y,x}$ est plat sur $\sh{O}_{S,s}$; en vertu de (11.3.1) $f$ et $g$ sont donc plats dans un voisinage de $x$ et il résulte de (19.2.4) que $j$ est une immersion transversalement régulière relativement à $S$ (en remplaçant au besoin $Y$ par $Y\cap V$ pour un voisinage ouvert $V$ de $x$ dans $X$). On a ainsi prouvé (i) et on voit qu'on peut supposer que $g:Y\to S$ est un morphisme plat. Comme $\sh{O}_{Y_s,x}$, quotient de $\sh{O}_{X_s,x}$ par l'idéal engendré par les $t_i$, est artinien par construction ($0$, 16.3.6), $x$ est point maximal de $Y_s$. Enfin, l'anneau $\sh{O}_{Y_s,x}$ étant artinien, est un anneau de Cohen-Macaulay ($0$, 16.5.1), donc, en remplaçant au besoin $Y$ par son intersection avec un voisinage ouvert de $x$ dans $X$, on peut supposer que $g$ est un morphisme de Cohen-Macaulay, en vertu de (12.1.7, (iii)).

En particulier, on retrouve l'existence des « quasi-sections » \emph{plates} $Y$ d'un morphisme $f:X\to S$ en un point $x\in X_s$ qui est \emph{fermé} dans $X_s$ et tel que $f$ soit plat au point $x$ et que $\sh{O}_{X_s,x}$ soit un anneau de Cohen-Macaulay (17.16.1) et on voit en outre que l'immersion $Y\to X$ est \emph{transversalement régulière} relativement à $S$. Ce dernier complément est aussi valable pour la quasi-section étale définie dans (17.16.3, (i)).

\subsection{Intersections complètes relatives (cas plat).}
\EGAReaderTarget{ega4.c04.subsec.19-3}{subsection:IV.19.3}

\oldpage[IV-4]{194}

\begin{definition}[19.3.1]
\EGAReaderTarget{ega4.c04.def.19-3-1}{IV.19.3.1}
Soit $A$ un anneau local noethérien. On dit que $A$ est un \emph{anneau d'intersection complète} (ou encore, \emph{anneau d'intersection complète absolue}, s'il y a danger de confusion) si le complété $\widehat A$ est isomorphe au quotient d'un anneau local noethérien complet régulier $B$ par un idéal régulier (i.e. (16.9.7) un idéal engendré par une suite régulière d'éléments de $B$).

Tout anneau local régulier est un anneau d'intersection complète.

On dit qu'un préschéma localement noethérien $X$ est \emph{intersection complète} (absolue) \emph{au point} $x\in X$ si l'anneau $\sh{O}_{X,x}$ est un anneau d'intersection complète.
\end{definition}

\begin{proposition}[19.3.2]
\label{IV.19.3.2}
Soient $B$ un anneau local noethérien régulier, $\mathfrak J$ un idéal de $B$. Pour que $A=B/\mathfrak J$ soit un anneau d'intersection complète, il faut et il suffit que l'idéal $\mathfrak J$ soit régulier (i.e. (16.9.7) engendré par une suite régulière d'éléments de $B$).
\end{proposition}

On peut écrire $\widehat A=\widehat B/\mathfrak J\widehat B$ et puisque $\widehat B$ est un $B$-module fidèlement plat et $B$ noethérien, $\mathfrak J\widehat B$ est régulier si et seulement si $\mathfrak J$ l'est (19.1.5, (ii)). Ceci montre (en vertu de ($0$, 17.1.5)) d'une part que la condition de l'énoncé est suffisante, et d'autre part que pour prouver qu'elle est nécessaire, on peut se borner au cas où $A$ et $B$ sont complets. Par hypothèse, il existe alors un anneau local noethérien complet régulier $B'$ tel que $A$ soit isomorphe à un quotient $B'/\mathfrak J'$, où $\mathfrak J'$ est un idéal régulier de $B'$. Considérons alors le \emph{produit fibré} $B\times_A B'=B''$ pour les homomorphismes surjectifs $B\to A$, $B'\to A$ ($0$, 18.1.2); nous prouverons d'abord le lemme suivant:

\begin{lemma}[19.3.2.1]
\label{IV.19.3.2.1}
Soient $A$, $E$, $F$ trois anneaux, $f:E\to A$, $g:F\to A$ deux homomorphismes, et posons $G=E\times_A F$ ($0$, 18.1.2).
\begin{enumerate}
  \item[(i)] Si $f$ est surjectif et si $E$ et $F$ sont des anneaux noethériens, il en est de même de $G$.
  \item[(ii)] Si $A$, $E$ et $F$ sont des anneaux locaux et si $f$ et $g$ sont des homomorphismes locaux, $G$ est un anneau local et les homomorphismes canoniques $G\to E$, $G\to F$ sont locaux.
  \item[(iii)] Si les conditions de (i) et (ii) sont vérifiées, si $g$ est surjectif, et si $E$ et $F$ sont des anneaux locaux noethériens complets, il en est de même de $G$.
\end{enumerate}
\end{lemma}

\begin{enumerate}
  \item[(i)] Compte tenu de ($0$, 18.1.3 et 18.1.5), $G$ est un $A$-anneau augmenté sur $F$, dont l'idéal d'augmentation $\mathfrak J'$ s'identifie canoniquement au noyau $\mathfrak J$ de $f$, la structure de $G$-module sur $\mathfrak J'$ s'identifiant à celle de $E$-module sur $\mathfrak J$ au moyen de l'homomorphisme canonique $G\to E$. Si $\mathfrak a$ est un idéal quelconque de $G$, $\mathfrak a/(\mathfrak a\cap\mathfrak J')$ est isomorphe à $(\mathfrak a+\mathfrak J')/\mathfrak J'$, donc à un idéal de $F$, et par suite est de type fini; d'autre part, $\mathfrak a\cap\mathfrak J'$ est isomorphe à un idéal de $E$ contenu dans $\mathfrak J$, donc est aussi de type fini, ce qui prouve que $\mathfrak a$ est de type fini, donc $G$ est noethérien.
  \item[(ii)] Soient $\mathfrak r$, $\mathfrak s$ les idéaux maximaux de $E$ et $F$ respectivement; l'ensemble $\mathfrak m=\mathfrak r\times_A\mathfrak s$ est évidemment un idéal de $G$, image réciproque de $\mathfrak r$ par la projection $G\to E$ et de $\mathfrak s$ par la projection $G\to F$. On en déduit aussitôt que les éléments de $\mathfrak m$ sont non inversibles dans $G$. D'autre part, si $(x,y)\notin\mathfrak m$, et si par exemple $x\notin\mathfrak r$, $f(x)$ n'appartient pas à l'idéal maximal de $A$, donc, puisque $g(y)=f(x)$, et que $g$ est local, on a aussi $y\notin\mathfrak s$; on en

\oldpage[IV-4]{195}

conclut que $x$ et $y$ sont inversibles, donc il en est de même de $(x,y)$, ce qui achève de prouver (ii).
  \item[(iii)] L'image de $\mathfrak m$ par l'homomorphisme surjectif $G\to E$ est l'idéal maximal $\mathfrak r$ de $E$, donc l'image de $\mathfrak m^n\cap\mathfrak J'$ est $\mathfrak r^n\cap\mathfrak J$, et comme $\mathfrak J$ est fermé (donc complet) pour la topologie induite par la topologie $\mathfrak r$-adique de $E$ ($0_{\mathrm I}$, 7.3.5), $\mathfrak J'$ est complet pour la topologie induite par la topologie $\mathfrak m$-préadique de $G$. D'autre part, $G/\mathfrak J'$, muni de la topologie $\mathfrak m$-préadique, est isomorphe à $F$ muni de la topologie $\mathfrak s$-adique, donc est complet pour la topologie quotient de la topologie $\mathfrak m$-préadique de $G$. On en déduit aussitôt que $G$ est complet pour la topologie $\mathfrak m$-préadique.
\end{enumerate}

Ce lemme étant établi, il résulte du théorème de Cohen ($0$, 19.8.8) que $B''$ est isomorphe à un quotient d'un anneau local noethérien complet \emph{régulier} $C$. Il résulte alors de (19.1.2) que l'immersion $\Spec(B')\to\Spec(C)$ est régulière; comme par hypothèse il en est de même de l'immersion $\Spec(A)\to\Spec(B')$, on conclut de (19.1.5, (iii)) que l'immersion $\Spec(A)\to\Spec(C)$ est régulière; mais cette immersion s'écrit aussi comme composée $\Spec(A)\to\Spec(B)\to\Spec(C)$. Or comme $B$ est régulier par hypothèse, l'immersion $\Spec(B)\to\Spec(C)$ est régulière (19.1.2); donc il en est de même de $\Spec(A)\to\Spec(B)$ par (19.1.5, (iv)).

\begin{corollary}[19.3.3]
\label{IV.19.3.3}
Soit $X$ un préschéma localement noethérien localement immersible dans un préschéma régulier. Alors l'ensemble des $x\in X$ où $X$ est une intersection complète est ouvert dans $X$ et contient l'ensemble des points réguliers de $X$.
\end{corollary}

On peut en effet se borner au cas où $X=\Spec(B/\mathfrak J)$, où $B$ est un anneau régulier et $\mathfrak J$ un idéal de $B$. L'ensemble des $\mathfrak p\supset\mathfrak J$ dans $\Spec(B)$ tels que $\mathfrak J_{\mathfrak p}$ soit régulier dans $B_{\mathfrak p}$ est alors ouvert (16.9.5), et c'est l'ensemble des points de $\Spec(B)$ où $X$ est une intersection complète en vertu de (19.3.2).

\begin{corollary}[19.3.4]
\label{IV.19.3.4}
Soient $k$ un corps, $X$ un préschéma localement de type fini sur $k$, $k'$ une extension de $k$, $X'=X\otimes_k k'$. Soient $x'$ un point de $X'$, $x$ sa projection dans $X$. Pour que $X$ soit intersection complète en $x$, il faut et il suffit que $X'$ soit intersection complète en $x'$.
\end{corollary}

La question étant locale sur $X$, on peut supposer que $X=\Spec(A)$, où $A$ est une $k$-algèbre de type fini, donc quotient d'une algèbre de polynômes $k[T_1,\ldots,T_n]$, si bien que $X$ est sous-préschéma fermé du schéma régulier $Y=\Spec(k[T_1,\ldots,T_n])$ ($0$, 17.3.7). Si l'on pose $Y'=Y\otimes_k k'=\Spec(k'[T_1,\ldots,T_n])$, $Y'$ est aussi un schéma régulier. Or, puisque $\Spec(k')$ est fidèlement plat sur $\Spec(k)$, il résulte de (19.1.5, (ii)) et du fait qu'il s'agit de préschémas noethériens que l'immersion $Y\to X$ est régulière au point $x$ si et seulement si l'immersion $Y'\to X'$ l'est au point $x'$. La conclusion résulte donc du critère (19.3.2).

\begin{corollary}[19.3.5]
\label{IV.19.3.5}
Soient $A$, $C$ deux anneaux locaux noethériens, $\rho:A\to C$ un homomorphisme local, $B=C/\mathfrak J$ un anneau quotient de $C$, $k$ le corps résiduel de $A$. On suppose que $\rho$ fait de $C$ un $A$-module plat, et que l'anneau $C\otimes_A k$ est régulier. Alors les conditions suivantes sont équivalentes:
\begin{enumerate}
  \item[a)] L'idéal $\mathfrak J$ est transversalement régulier relativement à $A$.
  \item[b)] $B$ est un $A$-module plat, et $B\otimes_A k$ est un anneau d'intersection complète.
\end{enumerate}
\end{corollary}

\oldpage[IV-4]{196}

En vertu de (19.2.4), la condition $a)$ équivaut à dire que $B$ est un $A$-module plat et que $\mathfrak J\otimes_A k$ (qui s'identifie à un idéal de $C\otimes_A k$ en vertu de la platitude de $B$ sur $A$ ($0_{\mathrm I}$, 6.1.2)) est régulier dans cet anneau. Comme l'anneau $C\otimes_A k$ est régulier, il revient au même, en vertu de (19.3.2), de dire (lorsque $B$ est un $A$-module plat) que $\mathfrak J\otimes_A k$ est un idéal régulier de $C\otimes_A k$, ou que $B\otimes_A k=(C\otimes_A k)/(\mathfrak J\otimes_A k)$ est un anneau d'intersection complète; d'où le corollaire.

\begin{definition}[19.3.6]
\EGAReaderTarget{ega4.c04.def.19-3-6}{IV.19.3.6}
Soit $f:X\to S$ un morphisme plat localement de présentation finie. On dit que $X$ est une \emph{intersection complète relativement à $S$ en un point} $x$, si la fibre $f^{-1}(f(x))$ est une intersection complète (absolue) au point $x$. On dit que $X$ est une \emph{intersection complète relativement à $S$}, ou que le morphisme $f$ est un \emph{morphisme d'intersection complète}, si $X$ est intersection complète relativement à $S$ en chacun de ses points.
\end{definition}

\begin{proposition}[19.3.7]
\label{IV.19.3.7}
Soient $g:X\to S$, $h:Y\to S$ deux morphismes plats localement de présentation finie, $f:Y\to X$ une $S$-immersion, $y$ un point de $Y$, $x=f(y)$, $s=g(x)=h(y)$. On suppose que la fibre $X_s=g^{-1}(s)$ soit régulière au point $x$. Alors, pour que $f$ soit une immersion transversalement régulière au point $y$, il faut et il suffit que $Y$ soit une intersection complète relativement à $S$ au point $y$.
\end{proposition}

Dire que $f$ est transversalement régulière au point $y$ revient, en vertu de (19.2.4), à dire que $f_s:Y_s\to X_s$ est une immersion régulière au point $y$. Mais comme $X_s$ est régulier au point $x=f(y)$, cela équivaut, en vertu de (19.3.2), à dire que $Y_s$ est intersection complète (absolue) au point $y$, d'où la proposition.

\begin{corollary}[19.3.8]
\label{IV.19.3.8}
Soit $f:X\to S$ un morphisme plat localement de présentation finie. L'ensemble $U$ des $x\in X$ en lesquels $X$ est une intersection complète relativement à $S$ est ouvert dans $X$. Si de plus $f$ est propre, l'ensemble $E$ des $s\in S$ tels que $X_s=f^{-1}(s)$ soit intersection complète en chacun de ses points est une partie ouverte de $S$.
\end{corollary}

Comme $E=Y-f(X-U)$, la seconde assertion résulte de la première et du fait que $f$ est une application fermée. La première assertion est locale sur $X$, donc on peut supposer que $X$ est un sous-préschéma fermé d'un préschéma de la forme $Z=S[T_1,\ldots,T_n]$ ($T_i$ indéterminées). Ce dernier étant lisse sur $S$ (17.3.8), ses fibres $Z_s$ sont régulières, et il revient donc au même, en vertu de (19.3.7), de dire que $x\in U$ ou que l'immersion $X\to Z$ est transversalement régulière au point $x$; on en conclut donc le corollaire par (19.2.2).

\begin{proposition}[19.3.9]
\label{IV.19.3.9}
\begin{enumerate}
  \item[(i)] Toute immersion ouverte est un morphisme d'intersection complète.
  \item[(ii)] Soit $f:X\to S$ un morphisme plat et localement de présentation finie qui est un morphisme d'intersection complète. Pour tout changement de base $g:S'\to S$, $f'=f_{(S')}:X_{(S')}\to S'$ est un morphisme d'intersection complète. La réciproque est vraie si $g$ est fidèlement plat et quasi-compact.
  \item[(iii)] Si $f:X\to Y$ et $g:Y\to Z$ sont deux morphismes plats localement de présentation finie qui sont des morphismes d'intersection complète, il en est de même de $g\circ f:X\to Z$.
\end{enumerate}
\end{proposition}

L'assertion (i) découle aussitôt de la définition (19.3.1). Pour prouver (ii), remarquons d'abord que si $f$ est plat et localement de présentation finie, il en est de

\oldpage[IV-4]{197}

même de $f'$ (2.1.4), la réciproque étant vraie si $g$ est fidèlement plat ((2.5.1) et (2.7.1)); en outre pour tout $s'\in S'$, si $s=g(s')$, il est équivalent de dire que $f^{-1}(s)$ est une intersection complète ou que $f'^{-1}(s')$ est une intersection complète (19.3.4); d'où (ii), compte tenu de ce que $g$ est surjectif lorsqu'il est fidèlement plat.

Enfin, sous les hypothèses de (iii), $g\circ f$ est plat et localement de présentation finie. Par définition (19.3.6), on est donc ramené au cas où $Z$ est le spectre d'un corps, et à prouver alors que les anneaux locaux $\sh{O}_{X,x}$ sont des anneaux d'intersection complète, ce qui est contenu dans le résultat plus général suivant:

\begin{corollary}[19.3.10]
\label{IV.19.3.10}
Soient $f:X\to S$ un morphisme plat localement de présentation finie qui soit un morphisme d'intersection complète, $x$ un point de $X$, $s=f(x)$. Si $\sh{O}_{S,s}$ est un anneau noethérien d'intersection complète, il en est de même de $\sh{O}_{X,x}$.
\end{corollary}

On peut évidemment se borner au cas où $S=\Spec(A)$ avec $A=\sh{O}_{S,s}$. Montrons en outre qu'on peut se borner au cas où l'anneau local $A$ est \emph{complet}. En effet, posons $A'=\widehat A$, $X'=X\otimes_A A'$, et soit $s'$ l'unique point fermé de $S'=\Spec(A')$, qui est l'unique point de $S'$ au-dessus de $s$; si $f'=f_{(S')}:X'\to S'$, la fibre $f'^{-1}(s')$ est isomorphe canoniquement à $f^{-1}(s)$, puisque $A$ et $A'$ ont même corps résiduel (I, 3.6.4); il y a donc un seul point $x'\in X'$ au-dessus de $x$ et de $s'$, et on a par suite $\sh{O}_{X',x'}=\sh{O}_{X,x}\otimes_A A'$. On en déduit que les anneaux locaux $\sh{O}_{X,x}$ et $\sh{O}_{X',x'}$ ont des \emph{complétés isomorphes}, car de façon générale, si $E$ est un anneau local qui est une $A$-algèbre (l'homomorphisme $A\to E$ étant local), le séparé complété $(E\otimes_A\widehat A)^{\wedge}$ de l'anneau $E\otimes_A\widehat A$ muni de la topologie produit tensoriel est isomorphe au séparé complété de $\widehat E\widehat\otimes_A\widehat A\cong\widehat E$, donc au séparé complété $\widehat E$ de $E$ ($0_{\mathrm I}$, 7.7.1). En vertu de la définition (19.3.1), il revient donc au même de dire que $\sh{O}_{X,x}$ est un anneau d'intersection complète, ou que $\sh{O}_{X',x'}$ est un anneau d'intersection complète.

Supposons donc $A$ \emph{complet}; on peut en outre supposer que $X=\Spec(B)$, $B$ étant quotient d'une algèbre de polynômes $C=A[T_1,\ldots,T_r]$. Comme les anneaux de polynômes sur un corps sont des anneaux réguliers ($0$, 17.3.7), il résulte de (19.3.7) que l'immersion $X\to Z=\Spec(C)$ peut être supposée transversalement régulière relativement à $S$; donc on peut supposer que $B=C/\mathfrak J$, où $\mathfrak J$ est engendré par une suite régulière $(g_i)_{1\leq i\leq n}$ d'éléments de $C$. D'autre part, puisque $A$ est complet, on peut par hypothèse l'écrire $A'/\mathfrak R$, où $A'$ est un anneau local régulier et $\mathfrak R$ un idéal de $A'$ ($0$, 19.8.8), et l'hypothèse que $A$ est un anneau d'intersection complète entraîne que $\mathfrak R$ est engendré par une suite régulière $(f_j)_{1\leq j\leq m}$ (19.3.2). Dans l'anneau de polynômes $C'=A'[T_1,\ldots,T_r]$, les éléments $f_j$ ($1\leq j\leq m$) forment encore une suite régulière engendrant l'idéal $\mathfrak R'=\mathfrak R[T_1,\ldots,T_r]$ ($0$, 15.1.4); comme $C'/\mathfrak R'=C$, on voit que si, pour chaque $i$, $g'_i$ est un élément de $C'$ d'image $g_i$ dans $C$, la suite formée des $f_j$ ($1\leq j\leq m$) et des $g'_i$ ($1\leq i\leq n$) est régulière dans $C'$; si $\mathfrak J'$ est l'idéal de $C'$ qu'elle engendre, on a $B=C'/\mathfrak J'$, d'où la conclusion en vertu de (19.3.2), puisque $C'$ est un anneau régulier ($0$, 17.3.7).

\subsection{Application: critères de régularité et de lissité pour les préschémas éclatés.}
\EGAReaderTarget{ega4.c04.subsec.19-4}{subsection:IV.19.4}

\oldpage[IV-4]{198}

\begin{env}[19.4.1]
\label{IV.19.4.1}
Soient $X$ un préschéma, $Y$ un sous-préschéma fermé de $X$, défini par un Idéal quasi-cohérent $\mathscr J$ de $\sh{O}_X$. Rappelons (II, 8.1.3) que le $X$-schéma $X'$ obtenu en faisant éclater $Y$ est le préschéma
\[
  X'=\operatorname{Proj}(\mathscr S),\qquad\text{où}\qquad \mathscr S=\bigoplus_{n\geq0}\mathscr J^n.
\]
Si $\mathscr J$ est de type fini, le morphisme structural $f:X'\to X$ est projectif. Sans hypothèse sur $\mathscr J$, le sous-préschéma fermé
\[
  Y'=f^{-1}(Y)
\]
de $X'$ est défini par l'Idéal quasi-cohérent $\mathscr J\sh{O}_{X'}$ de $\sh{O}_{X'}$, canoniquement isomorphe à $\sh{O}_{X'}(1)$; de façon plus précise, on a une suite exacte
\begin{equation}
  0\longrightarrow\sh{O}_{X'}(1)\xrightarrow{\widetilde u_0}\sh{O}_{X'}\longrightarrow\sh{O}_{Y'}\longrightarrow0.
  \tag{19.4.1.1}\label{IV.19.4.1.1}
\end{equation}
où $\mathscr J\sh{O}_{X'}$ est l'image de $\widetilde u_0$ (II, 8.1.7 et 8.1.8). Comme au voisinage d'un point de $X'$, le $\sh{O}_{X'}$-Module $\sh{O}_{X'}(1)$ est libre de rang $1$, $\mathscr J\sh{O}_{X'}$ est engendré dans un tel voisinage par une section inversible de $\sh{O}_{X'}$ et $\mathscr J\sh{O}_{X'}/\mathscr J^2\sh{O}_{X'}$ est par suite un $\sh{O}_{X'}/\mathscr J\sh{O}_{X'}$-Module libre de rang $1$. Ceci prouve la première assertion du lemme suivant:
\end{env}

\begin{lemma}[19.4.2]
\label{IV.19.4.2}
L'immersion canonique $Y'\to X'$ est régulière et de codimension $1$, et $Y'$ est canoniquement isomorphe au $Y$-schéma $\operatorname{Proj}(\operatorname{gr}_{\mathscr J}^*(\sh{O}_X))$.
\end{lemma}

La seconde assertion résulte de (II, 3.3.3) appliqué à l'injection canonique $Y\to X$, compte tenu de ce que $\mathscr S\otimes_{\sh{O}_X}\sh{O}_Y=\operatorname{gr}_{\mathscr J}^*(\sh{O}_X)$ par définition, puisque $\sh{O}_Y$ est isomorphe à $\sh{O}_X/\mathscr J$ et $\mathscr J^n\otimes_{\sh{O}_X}(\sh{O}_X/\mathscr J)$ à $\mathscr J^n/\mathscr J^{n+1}$.

\begin{corollary}[19.4.3]
\label{IV.19.4.3}
Si l'immersion canonique $Y\to X$ est quasi-régulière, la restriction $g:Y'\to Y$ du morphisme $f$ est lisse.
\end{corollary}

En effet, $\sh{N}_{Y/X}=\mathscr J/\mathscr J^2$ est alors un $\sh{O}_Y$-Module localement libre de rang fini et $\operatorname{gr}_{\mathscr J}^*(\sh{O}_X)$ est isomorphe à $\mathbf S_{\sh{O}_Y}^*(\sh{N}_{Y/X})$ (16.9.8), donc $Y'$ est $Y$-isomorphe à $\mathbf P(\sh{N}_{Y/X})$ et par suite est lisse sur $Y$ (17.3.9).

Nous verrons plus tard (chap. V) que $g$ est encore lisse lorsque l'immersion $Y\to X$ « ne présente que des singularités ordinaires ».

\begin{proposition}[19.4.4]
\label{IV.19.4.4}
Avec les notations de (19.4.1), supposons que $X$ soit localement noethérien et que le morphisme $g:Y'\to Y$, restriction de $f$, soit lisse. Soient $x'$ un point de $Y'$, $x=g(x')$. Pour que $Y$ soit régulier au point $x$, il faut et il suffit que $Y'$ soit régulier au point $x'$; lorsqu'il en est ainsi, $X'$ est alors régulier au point $x'$.
\end{proposition}

La première assertion résulte de (17.5.8) et la seconde de (19.1.1) et (19.4.2).

\begin{corollary}[19.4.5]
\label{IV.19.4.5}
Sous les hypothèses générales de (19.4.4), supposons que $Y$ soit régulier au point $x$. Alors, pour toute générisation $x_1$ de $x$ dans $X$ n'appartenant pas à $Y$, $X$ est régulier au point $x_1$. Si $\operatorname{Reg}(X)$ est ouvert ((6.12), par exemple si $X$ est excellent (7.8.6, (iii)))

\oldpage[IV-4]{199}

et si $Y$ est régulier, alors il existe un voisinage ouvert $U$ de $Y$ dans $X$ tel que $X$ soit régulier aux points de $U-Y$.
\end{corollary}

Pour prouver la première assertion, on peut se borner au cas où $X$ est un schéma local de point fermé $x$; l'hypothèse entraîne alors que $Y$ est régulier ($0$, 17.3.2), donc il en est de même de $Y'$ par (19.4.4). Notons maintenant que le morphisme $f:X'\to X$ est projectif, donc propre, et par suite (II, 7.2.1), si l'on considère l'unique point $x'_1$ de $X'$ au-dessus de $x_1$ (II, 8.1.3), il existe une spécialisation de $x'_1$ appartenant à $Y'$; en lui appliquant (19.4.4) et utilisant ($0$, 17.3.2) on en conclut que $X'$ est régulier au point $x'_1$, donc $X$ régulier au point $x_1$ puisque $X'-Y'$ est isomorphe à $X-Y$ (II, 8.1.3). Si $Y$ est régulier, tout point de $X-Y$ qui est générisation d'un point de $Y$ appartient donc à $\operatorname{Reg}(X)$; si $\operatorname{Reg}(X)$ est ouvert, $U=Y\cup\operatorname{Reg}(X)$ est donc localement constructible et stable par générisation, donc ouvert ($0_{\mathrm{III}}$, 9.2.5) et répond à la question.

\begin{proposition}[19.4.6]
\label{IV.19.4.6}
Soient $g:X\to S$ un morphisme, $Y$ un sous-préschéma fermé de $X$ défini par un Idéal quasi-cohérent $\mathscr J$, $X'$ le $X$-schéma obtenu en faisant éclater $Y$, $Y'$ l'image réciproque de $Y$ dans $X'$.
\begin{enumerate}
  \item[(i)] Les conditions suivantes sont équivalentes:
  \begin{enumerate}
    \item[a)] $\operatorname{gr}_{\mathscr J}^*(\sh{O}_X)$ est une $\sh{O}_X$-Algèbre $g$-plate.
    \item[b)] Pour tout $n\geq0$, le $\sh{O}_X$-Module $\sh{O}_X/\mathscr J^{n+1}$ est $g$-plat (en d'autres termes, le $n$-ème voisinage infinitésimal $Y^{(n)}$ de $Y$ dans $X$ (16.1.2) est $S$-plat).
    \item[c)] Pour tout changement de base $S_1\to S$ et tout entier $n\geq1$, si l'on pose $X_1=X\times_S S_1$, l'homomorphisme canonique $\mathscr J^n\otimes_{\sh{O}_S}\sh{O}_{S_1}\to\mathscr J^n\sh{O}_{X_1}$ est bijectif.
  \end{enumerate}
  Lorsqu'il en est ainsi, $Y'$ est $S$-plat, et pour tout changement de base $S_1\to S$, si $Y_1$ est l'image réciproque de $Y$ dans $X_1$, le préschéma $X'_1$ obtenu en faisant éclater $Y_1$ dans $X_1$ est canoniquement isomorphe à $X'\times_S S_1$.
  \item[(ii)] Supposons que les conditions équivalentes de (i) soient satisfaites et de plus que les morphismes $X\to S$ et $Y\to S$ soient localement de présentation finie. Alors $\mathscr S=\bigoplus_{n\geq0}\mathscr J^n$ est une $\sh{O}_X$-Algèbre de présentation finie, le morphisme $X'\to X$ est de présentation finie, l'immersion canonique $Y'\to X'$ est transversalement régulière de codimension $1$ relativement à $S$, et l'ensemble des points de $X$ (resp. $X'$) où $X$ (resp. $X'$) est $S$-plat, est un voisinage de $Y$ (resp. $Y'$).
\end{enumerate}
\end{proposition}

\begin{enumerate}
  \item[(i)] L'équivalence de $a)$ et $b)$ résulte aussitôt de ($0_{\mathrm I}$, 6.1.2); pour montrer l'équivalence de $b)$ et $c)$, on peut se borner au cas où $S=\Spec(A)$ et $X=\Spec(B)$ sont affines, avec $\mathscr J=\widetilde{\mathfrak J}$, où $\mathfrak J$ est un idéal de $B$; la suite exacte des Tor montre alors que la condition $c)$ implique que $\operatorname{Tor}_1^A(B/\mathfrak J^{n+1},A')=0$ pour tout entier $n\geq0$ et toute $A$-algèbre $A'$; en prenant pour $A'$ une algèbre somme directe $A\oplus M$, où $M$ est un $A$-module quelconque et où la multiplication est définie par $(a,m)(a',m')=(aa',am'+a'm)$ on voit que la condition précédente équivaut à $\operatorname{Tor}_1^A(B/\mathfrak J^{n+1},M)=0$, donc au fait que $B/\mathfrak J^{n+1}$ est un $A$-module plat. Si l'on pose $\mathscr J_1=\mathscr J\otimes_{\sh{O}_S}\sh{O}_{S_1}$, on a alors $\bigoplus_{n\geq0}\mathscr J_1^n=(\bigoplus_{n\geq0}\mathscr J^n)\otimes_{\sh{O}_S}\sh{O}_{S_1}$ et l'assertion relative au préschéma éclaté $X'_1$ s'en déduit aussitôt. Enfin, pour prouver que $Y'$ est $S$-plat, on peut encore se borner au cas où $S$ et $X$ sont affines; si l'on pose $C=\operatorname{gr}_{\mathfrak J}^*(B)$, tout point de $Y'$ admet alors, en vertu de (19.4.2) et de (II, 2.3.6), un

\oldpage[IV-4]{200}

voisinage ouvert affine dont l'anneau est de la forme $C_{(t)}$, où $t$ est un élément homogène de degré $1$ de $C$; comme par hypothèse $C$ est un $A$-module plat, il en est de même de son anneau de fractions $C_t$ et de la composante de degré $0$, $C_{(t)}$, de ce $A$-module gradué, d'où notre assertion.
  \item[(ii)] On peut toujours se borner au cas où $S$ et $X$ sont affines, $B$ étant donc une $A$-algèbre de présentation finie, $\mathfrak J$ un idéal de type fini de $B$. En vertu de (8.9.1), (8.6.3) et (11.2.9), il existe un sous-anneau noethérien $A_0$ de $A$, une $A_0$-algèbre de type fini $B_0$, un idéal $\mathfrak J_0$ de $B_0$ tels que $B=B_0\otimes_{A_0}A$, $\mathfrak J=\mathfrak J_0B$, que $\operatorname{gr}_{\mathfrak J_0}^*(B_0)$ soit un $A_0$-module \emph{plat} et que $\operatorname{gr}_{\mathfrak J}^*(B)=\operatorname{gr}_{\mathfrak J_0}^*(B_0)\otimes_{A_0}A$. En vertu de (i), on a alors $\bigoplus_{n\geq0}\mathfrak J^n=(\bigoplus_{n\geq0}\mathfrak J_0^n)\otimes_{A_0}A=(\bigoplus_{n\geq0}\mathfrak J_0^n)\otimes_{B_0}B$; comme $\bigoplus_{n\geq0}\mathfrak J_0^n$ est une $B_0$-algèbre engendrée par l'idéal de type fini $\mathfrak J_0$ de $B_0$, donc une $B_0$-algèbre de type fini, et par suite de présentation finie puisque $B_0$ est noethérien, on en déduit que $\bigoplus_{n\geq0}\mathfrak J^n$ est une $B$-algèbre de présentation finie. De même, le morphisme $\operatorname{Proj}(\bigoplus_{n\geq0}\mathfrak J_0^n)\to\Spec(B_0)$ est de type fini (II, 2.7.1), et par suite de présentation finie; donc le morphisme $X'\to X$, qui s'en déduit par changement de base, est de présentation finie. On sait déjà (19.4.2) que l'immersion canonique $Y'\to X'$ est régulière; comme $Y'$ est $S$-plat, ainsi qu'on l'a vu dans (i), on déduit de (19.2.4) que l'immersion $Y'\to X'$ est transversalement régulière relativement à $S$ et que $X'$ est $S$-plat aux points de $Y'$, donc (11.3.1) dans un voisinage ouvert de $Y'$. D'autre part, le fait que $\operatorname{gr}_{\mathscr J}^*(\sh{O}_X)$ est une $\sh{O}_X$-Algèbre $S$-plate entraîne, par (11.3.4), que $X$ est $S$-plat aux points de $Y$, donc (11.3.1) dans un voisinage de $Y$.
\end{enumerate}

\begin{corollary}[19.4.7]
\label{IV.19.4.7}
Soient $g:X\to S$ un morphisme localement de présentation finie, $Y$ un sous-préschéma fermé de $X$ tel que le morphisme composé $Y\to X\to S$ soit localement de présentation finie; supposons en outre que $Y$ soit $S$-plat et que $\sh{O}_X$ soit normalement plat le long de $Y$ (11.3.4). Alors (avec les notations de (19.4.1)), le morphisme $X'\to X$ est de présentation finie, l'immersion canonique $Y'\to X'$ est transversalement régulière de codimension $1$ relativement à $S$, $Y'$ est $Y$-plat (donc $S$-plat) et l'ensemble des points de $X$ (resp. $X'$) où $X$ (resp. $X'$) est $S$-plat est un voisinage de $Y$ (resp. $Y'$).
\end{corollary}

En effet, par hypothèse $\operatorname{gr}_{\mathscr J}^*(\sh{O}_X)$ est un $(\sh{O}_X/\mathscr J)$-Module plat, donc est $S$-plat puisque $Y$ est supposé être $S$-plat. On peut par suite appliquer les résultats de (19.4.6, (i) et (ii)). Comme $Y'$ est $Y$-isomorphe à $\operatorname{Proj}(\operatorname{gr}_{\mathscr J}^*(\sh{O}_X))$ (19.4.2), le même raisonnement que dans (19.4.6) montre que $Y'$ est $Y$-plat.

\begin{proposition}[19.4.8]
\label{IV.19.4.8}
Supposons vérifiées les hypothèses de (19.4.7) et en outre que le morphisme $Y'\to Y$ soit lisse. Soient $x'$ un point de $Y'$, $x$ son image dans $Y$. Pour que $Y$ soit lisse sur $S$ au point $x$, il faut et il suffit que $Y'$ soit lisse sur $S$ au point $x'$; lorsqu'il en est ainsi, $X'$ est alors lisse sur $S$ au point $x'$.

Les conclusions précédentes sont en particulier vérifiées lorsque le morphisme $X\to S$ est localement de présentation finie et que l'immersion $Y\to X$ est transversalement régulière relativement à $S$.
\end{proposition}

Compte tenu de la compatibilité de la formation d'un préschéma éclaté avec les changements de base dans le cas envisagé (19.4.6, (i)) et de (17.5.1) et (17.7.1),

\oldpage[IV-4]{201}

il suffit de considérer le cas où $S$ est le spectre d'un corps algébriquement clos. Mais alors il revient au même de dire qu'un $S$-préschéma localement de type fini est lisse sur $S$ en un point ou qu'il est régulier en ce point (17.15.1), donc les conclusions résultent de (19.4.4).

Lorsque $X\to S$ est localement de présentation finie, et que l'immersion canonique $Y\to X$ est quasi-régulière, cette immersion est par définition de présentation finie puisque l'Idéal qui définit $Y$ est de type fini, donc le morphisme composé $Y\to X\to S$ est localement de présentation finie. On a déjà vu (19.4.3) que dans ces conditions le morphisme $Y'\to Y$ est lisse; en outre, on sait alors que $\sh{N}_{Y/X}$ est un $\sh{O}_Y$-Module localement libre et que $\operatorname{gr}_{\mathscr J}^*(\sh{O}_X)$ est isomorphe à $\mathbf S_{\sh{O}_Y}^*(\sh{N}_{Y/X})$, donc est un $\sh{O}_Y$-Module plat. Les hypothèses de (19.4.7) sont donc toutes vérifiées et l'on peut appliquer les conclusions de l'énoncé de (19.4.8).

\begin{corollary}[19.4.9]
\label{IV.19.4.9}
Sous les hypothèses générales de (19.4.8), supposons en outre que $Y$ soit lisse sur $S$ au point $x$. Alors, pour toute générisation $x_1$ de $x$ dans $X$ n'appartenant pas à $Y$, $X$ est lisse sur $S$ au point $x_1$. Si $Y$ est lisse sur $S$, alors il existe un voisinage ouvert $U$ de $Y$ dans $X$ tel que $X$ soit lisse sur $S$ aux points de $U-Y$.
\end{corollary}

Pour la première assertion, on peut se borner au cas où $X$ est un schéma local de point fermé $x$, et l'hypothèse entraîne alors que $Y$ est lisse sur $S$, tout voisinage du point fermé de $Y$ dans $Y$ étant nécessairement $Y$ lui-même; on conclut donc par le même raisonnement que dans (19.4.5), utilisant le fait que l'ensemble des points où $X'$ est lisse sur $S$ est ouvert dans $X'$. Pour prouver la seconde assertion, on peut se ramener au cas où $S$ est noethérien, et $X$ de type fini sur $S$, grâce à (8.9.1), (8.6.3), (11.2.9) et (17.7.6); on conclut alors comme dans (19.4.5), en utilisant le fait que l'ensemble des points où $X$ est lisse sur $S$ est ouvert dans $X$.

\begin{remark}[19.4.10]
\label{IV.19.4.10}
Dans (19.4.6, (ii)), remplaçons l'hypothèse que $X$ et $Y$ sont localement de présentation finie sur $S$ par celle que $S$ et $X$ (donc aussi $Y$) sont \emph{localement noethériens}. Alors il est clair que $\bigoplus_{n\geq0}\mathscr J^n$ est une $\sh{O}_X$-Algèbre de type fini, que le morphisme $X'\to X$ est de type fini, et il résulte encore de (19.2.4) appliqué aux préschémas localement noethériens, que l'immersion $Y'\to X'$ est transversalement régulière relativement à $S$ et que $X'$ est $S$-plat aux points de $Y'$.
\end{remark}

\begin{proposition}[19.4.11]
\EGAReaderTarget{ega4.c04.prop.19-4-11}{IV.19.4.11}
Soient $X$ un préschéma, $\mathscr E$ un $\sh{O}_X$-Module quasi-cohérent, $u:\mathscr E\to\sh{O}_X$ un homomorphisme de $\sh{O}_X$-Modules, $\mathscr J=u(\mathscr E)$ l'Idéal quasi-cohérent de $\sh{O}_X$ image de $u$, $Y$ le sous-préschéma fermé de $X$ défini par $\mathscr J$. Soit $X'$ le $X$-schéma obtenu en faisant éclater $Y$. D'autre part, soient $P=\mathbf P(\mathscr E)$ le fibré projectif sur $X$ défini par $\mathscr E$ (II, 4.1.1), $p:P\to X$ le morphisme structural, $\alpha_1^{\#}:p^*(\mathscr E)\to\sh{O}_P(1)$ l'homomorphisme canonique (II, 4.1.5.1) et $\mathscr H$ son noyau, de sorte qu'on a la suite exacte
\[
  0\longrightarrow\mathscr H\longrightarrow p^*(\mathscr E)\xrightarrow{\alpha_1^{\#}}\sh{O}_P(1)\longrightarrow0.
\]
Enfin, soit $\mathscr K$ l'Idéal quasi-cohérent de $\sh{O}_P$ image de la restriction $v:\mathscr H\to\sh{O}_P$ de $p^*(u)$, et soit $Z$ le sous-préschéma fermé de $P$ défini par $\mathscr K$.

\oldpage[IV-4]{202}

\begin{enumerate}
  \item[(i)] L'image de $X'$ par l'immersion fermée $j:X'\to\mathbf P(\mathscr E)$ correspondant à l'homomorphisme surjectif de $\sh{O}_X$-Algèbres graduées $\mathbf S_{\sh{O}_X}(\mathscr E)\to\mathbf S_{\sh{O}_X}(\mathscr J)\to\bigoplus_{n\geq0}\mathscr J^n$ (II, 3.6.2) est un sous-préschéma fermé de $Z$.
  \item[(ii)] Si $\mathscr E$ est un $\sh{O}_X$-Module localement libre de rang $m$, $\mathscr H$ est un $\sh{O}_P$-Module localement libre de rang $m-1$.
  \item[(iii)] On suppose que $X$ soit localement noethérien, que $\mathscr E$ soit localement libre de type fini, que le sous-préschéma $Y$ soit régulier, et qu'en tout point $x\in Y$, l'immersion canonique $i:Y\to X$ soit régulière et de codimension égale à $\operatorname{rg}_{\sh{O}_{X,x}}(\mathscr E_x)$ (ce que nous exprimerons en disant que l'homomorphisme $u$ est \emph{régulier} au point $x$ (cf. chap. V)). Alors l'image de $X'$ par l'immersion fermée $j:X'\to P$ est le sous-préschéma fermé $Z$ de $P$; en tous les points de $X'-Y'$ ($Y'$ image réciproque de $Y$ dans $X'$) l'immersion $j$ est transversalement régulière relativement à $X$; enfin en tous les points $x'$ de $Y'$, $X'$ est régulier, l'immersion $j$ est régulière et de codimension $\operatorname{rg}_{\sh{O}_{P,z}}(\mathscr H_z)$ où $z=j(x')$ (autrement dit l'homomorphisme $v:\mathscr H\to\sh{O}_P$ est régulier en $z$).
\end{enumerate}
\end{proposition}

Toutes les questions étant locales sur $X$, on peut se borner au cas où $X=\Spec(A)$ est affine et $\mathscr E=\widetilde E$, où $E$ est un $A$-module. Pour faire la vérification de (i), on peut en outre se borner à examiner ce qui se passe dans un ouvert affine de $P$ de la forme $D_+(t)$, où $t\in E=\mathbf S_A^1(E)$ (II, 2.3.14). Si l'on pose $S=\mathbf S_A(E)$, on a alors $\Gamma(D_+(t),\sh{O}_P)=S_{(t)}$ (II, 2.4.1), et $p^*(\mathscr E)|D_+(t)=(E\otimes_A S_{(t)})^\sim$. Si l'on se reporte à la définition de $\alpha_1^{\#}$ (II, 2.6.2.2), on voit qu'à tout élément
\[
  a=\sum_i x^{(i)}\otimes\frac{x_1^{(i)}x_2^{(i)}\cdots x_n^{(i)}}{t^n}
\]
de $E\otimes_A S_{(t)}$ (où $x^{(i)}$ et les $x_j^{(i)}$ sont dans $E$) $\alpha_1^{\#}$ fait correspondre l'élément
\[
  \sum_i\frac{x^{(i)}x_1^{(i)}\cdots x_n^{(i)}}{t^n}
\]
de $(S(1))_{(t)}$. Il s'agit (pour prouver (i)) de montrer que si ce dernier élément est nul, alors il en est de même de l'image de
\[
  a'=v(a)=\sum_i u(x^{(i)})\frac{x_1^{(i)}\cdots x_n^{(i)}}{t^n}
\]
par l'homomorphisme canonique $S_{(t)}\to(\bigoplus_{n\geq0}(u(E))^n)_{(u(t))}$ (II, 3.6.2). Or, cette image n'est autre que
\[
  \sum_i\frac{u(x^{(i)})u(x_1^{(i)})\cdots u(x_n^{(i)})}{(u(t))^n}
\]
dans l'anneau $A_{u(t)}$, c'est-à-dire le produit de l'élément $u(t)$ et de l'image canonique de $\sum_i x^{(i)}x_1^{(i)}\cdots x_n^{(i)}/t^{n+1}$ par l'homomorphisme d'algèbres, prolongement canonique à $S_{(t)}$ de l'homomorphisme $u:E\to A$ de $A$-modules. Cela prouve donc (i).

Pour établir (ii), notons que l'on peut supposer que $E$ est un $A$-module libre de rang $m$; puisque $\alpha_1^{\#}:E\otimes_A S_{(t)}\to(S(1))_{(t)}$ est surjectif et que $(S(1))_{(t)}$ est un $S_{(t)}$-module libre de rang $1$, la suite exacte
\[
  0\longrightarrow H_{(t)}\longrightarrow E\otimes_A S_{(t)}\xrightarrow{\alpha_1^{\#}}(S(1))_{(t)}\longrightarrow0
\]
est \emph{scindée}, et $H_{(t)}$ est donc un $S_{(t)}$-module projectif de rang $m-1$, d'où (ii).

Pour prouver (iii), examinons d'abord ce qui se passe au-dessus d'un point $x\in X-Y$. Alors $u:\mathscr E\to\sh{O}_X$, restreint à un voisinage de $x$, est surjectif, donc on peut se borner au cas où $\mathscr J=\sh{O}_X$, cas où $X'$ s'identifie canoniquement à $X$. Pour voir alors que le sous-préschéma fermé de $P$ image de $X'$ s'identifie à $Z$, on peut se borner comme au début au cas où $X=\Spec(A)$ est affine, et en outre $E=A^m$, de sorte que $S=A[T_0,\ldots,T_{m-1}]$. On peut en outre, avec les notations du début, supposer $u(t)\ne0$ (les valeurs de $t\in E$

\oldpage[IV-4]{203}

qui vérifient cette condition engendrant l'algèbre symétrique $\mathbf S_A(E)$); changeant au besoin de base dans $E$, on peut donc supposer que $t=T_0$, et $u(T_i)=0$ pour $1\leq i\leq m-1$, de sorte que $S_{(t)}$ s'identifie canoniquement à $A[T_1,\ldots,T_{m-1}]$. Les éléments de $H$ sont alors les éléments de la forme
\[
  \sum_\alpha L_\alpha(\mathbf T)\otimes\mathbf T^\alpha
\]
de $E\otimes_A A[T_1,\ldots,T_{m-1}]$ où les $L_\alpha(\mathbf T)=\lambda_{\alpha0}+\lambda_{\alpha1}T_1+\cdots+\lambda_{\alpha,m-1}T_{m-1}$ doivent être tels que $\sum_\alpha L_\alpha(\mathbf T)\mathbf T^\alpha=0$ dans $A[T_1,\ldots,T_{m-1}]$. On vérifie aisément que lorsque la condition $\lambda_{\mathbf0,0}=0$ est remplie, les valeurs des $\lambda_{\alpha0}$ peuvent être prises \emph{arbitrairement} pour $|\alpha|\geq1$, et on peut toujours alors déterminer les $\lambda_{\alpha i}$ pour $i\geq1$ de façon à vérifier la condition $\sum_\alpha L_\alpha(\mathbf T)\mathbf T^\alpha=0$. Les éléments de l'idéal $\mathscr K$ correspondants sont alors les polynômes $\sum_\alpha\lambda_{\alpha0}\mathbf T^\alpha$ \emph{sans terme constant}, et comme $(\mathbf S_A(A))_{(t)}$ s'identifie à $A$, ces polynômes sont exactement les éléments du noyau de l'homomorphisme $(\mathbf S_A(E))_{(t)}\to(\mathbf S_A(A))_{(t)}$. On a ainsi prouvé (19.2.1) que l'immersion $j:X'\to P$ est \emph{transversalement régulière} relativement à $X$ en tous les points de $X'-Y'$ et que $j(X')=Z$ en ces points.

Reste à examiner ce qui se passe au-dessus d'un point $x\in Y$. L'hypothèse de « régularité » de $u$ en ce point faite dans (iii) équivaut (en vertu de (19.1.1) et de ($0$, 17.1.7)) à dire que $X$ est \emph{régulier} au point $x$ et que l'homomorphisme
\[
  u_x\otimes1:\mathscr E_x\otimes_{\sh{O}_{X,x}}\mathbf k(x)\longrightarrow\mathfrak m_x/\mathfrak m_x^2
\]
est \emph{injectif}. Cela étant, comme le morphisme $p:P\to X$ est lisse (17.3.9), pour tout $z\in P$ au-dessus de $x$, $P$ est régulier au point $z$ (17.5.8); en outre ($0$, 17.3.3) l'homomorphisme canonique
\begin{equation}
  (\mathfrak m_x/\mathfrak m_x^2)\otimes_{\mathbf k(x)}\mathbf k(z)\longrightarrow\mathfrak m_z/\mathfrak m_z^2
  \tag{19.4.11.1}\label{IV.19.4.11.1}
\end{equation}
est \emph{injectif}. On en déduit que l'homomorphisme
\[
  (\mathscr E_x\otimes_{\sh{O}_{X,x}}\sh{O}_{P,z})\otimes_{\sh{O}_{P,z}}\mathbf k(z)\longrightarrow\mathfrak m_z/\mathfrak m_z^2
\]
déduit de $p^*(u)$ est aussi injectif, car il s'écrit comme le composé
\[
  (\mathscr E_x\otimes_{\sh{O}_{X,x}}\mathbf k(x))\otimes_{\mathbf k(x)}\mathbf k(z)\longrightarrow(\mathfrak m_x/\mathfrak m_x^2)\otimes_{\mathbf k(x)}\mathbf k(z)\longrightarrow\mathfrak m_z/\mathfrak m_z^2
\]
où la première flèche est injective par platitude et la seconde est l'homomorphisme injectif (19.4.11.1). Comme $\mathscr H$ est (dans un voisinage de $z$ dans $P$) facteur direct de $\mathscr E\otimes_{\sh{O}_X}\sh{O}_P=p^*(\mathscr E)$, l'homomorphisme $v_z\otimes1:\mathscr H_z\otimes_{\sh{O}_{P,z}}\mathbf k(z)\to\mathfrak m_z/\mathfrak m_z^2$ est aussi injectif, ce qui établit la « régularité » de l'homomorphisme $v$ au point $z$. Reste à voir que l'image par $j$ de $X'$ est identique au sous-préschéma fermé $Z$. Il suffira pour cela de montrer d'une part que l'image par $j$ des points de $q^{-1}(x)$ (où $q:X'\to X$ est le morphisme structural) sont exactement ceux de $Z\cap p^{-1}(x)$ et d'autre part qu'en chacun de ces points $z=j(x')$, l'homomorphisme surjectif $\sh{O}_{Z,z}\to\sh{O}_{X',x'}$ est bijectif. Mais, compte tenu de l'hypothèse et de (19.4.3), la restriction $Y'\to Y$ de $q$ est lisse, et puisque $\sh{O}_{Y,x}$ est supposé régulier, il en est de même de $\sh{O}_{Y',x'}$ (17.5.8); en outre, l'immersion $Y'\to X'$ est régulière (19.4.2), donc $\sh{O}_{X',x'}$ est aussi un anneau régulier (19.1.1). Pour prouver que l'homomorphisme $\sh{O}_{Z,z}\to\sh{O}_{X',x'}$ est bijectif, il suffit donc de montrer que les

\oldpage[IV-4]{204}

dimensions de ces deux anneaux réguliers sont égales ($0$, 17.1.9). Mais en vertu de ($0$, 17.3.3), cela équivaut aussi à l'égalité des dimensions des deux anneaux $\sh{O}_{Z,z}\otimes_{\sh{O}_{X,x}}\mathbf k(x)$ et $\sh{O}_{X',x'}\otimes_{\sh{O}_{X,x}}\mathbf k(x)$, si bien que finalement on peut remplacer $X'$ et $Z$ par les \emph{fibres} $q^{-1}(x)$ et $Z\cap p^{-1}(x)$ respectivement, et se ramener à prouver que le morphisme $q^{-1}(x)\to Z\cap p^{-1}(x)$ déduit de $j$ est un \emph{isomorphisme}. Or, on peut ici se borner au cas où $X=\Spec(B)$ avec $B=\sh{O}_{X,x}$ et $E=B^m$, d'où $S=B[T_1,\ldots,T_m]$; en outre, l'hypothèse signifie qu'il existe un système régulier de paramètres $(t_i)_{1\leq i\leq n}$ de $B$ tel que $u(T_i)=t_i$ pour $1\leq i\leq m$ ($0$, 17.1.7); soit $\mathfrak J$ l'image de $u$; si $\mathfrak m$ est l'idéal maximal de $B$, on a $\mathfrak J^n\otimes_B B/\mathfrak m=\mathfrak J^n/\mathfrak m\mathfrak J^n=(\mathfrak J^n/\mathfrak J^{n+1})\otimes_B(B/\mathfrak m)$ et par suite (16.9.4.1), $(\bigoplus_{n\geq0}\mathfrak J^n)\otimes_B\mathbf k(x)$ s'identifie à $\mathbf k(x)[T_1,\ldots,T_m]$, d'où résulte que l'immersion $q^{-1}(x)\to p^{-1}(x)$ est ici un \emph{isomorphisme}; ceci implique \emph{a fortiori} que le sous-préschéma fermé $Z\cap p^{-1}(x)$ est identique à la fibre $p^{-1}(x)$, et termine la démonstration de (19.4.11).

\subsection{Critères de M-régularité.}
\EGAReaderTarget{ega4.c04.subsec.19-5}{subsection:IV.19.5}

Nous reprenons ici, en les complétant, les critères pour qu'une suite d'éléments d'un anneau $A$ soit M-régulière ou M-quasi-régulière ($M$ étant un $A$-module), déjà étudiés dans ($0$, 15.1).

\begin{theorem}[19.5.1]
\label{IV.19.5.1}
Soient $A$ un anneau, $\mathbf f=(f_i)_{1\leq i\leq n}$ une suite d'éléments de $A$, $M$ un $A$-module. Considérons les propriétés suivantes:
\begin{enumerate}
  \item[a)] La suite $\mathbf f$ est M-régulière.
  \item[b)] On a $\operatorname{H}_i(\mathbf f,M)=0$ pour tout $i>0$ (III, 1.1.3).
  \item[b$'$)] On a $\operatorname{H}_1(\mathbf f,M)=0$.
  \item[c)] La suite $\mathbf f$ est M-quasi-régulière.
\end{enumerate}
Alors on a les implications
\[
  a)\Rightarrow b)\Rightarrow b')\qquad\text{et}\qquad a)\Rightarrow c).
\]
Posons $\mathfrak J=\sum_{i=1}^n f_iA$. Si les modules $M_i=M/(\sum_{j=1}^i f_jM)$ sont séparés pour la topologie $\mathfrak J$-préadique (resp. si tout module quotient d'un sous-module de $M$ est séparé pour la topologie $\mathfrak J$-préadique), alors on a aussi l'implication $c)\Rightarrow a)$ (resp. $b')\Rightarrow a)$.
\end{theorem}

Les implications $a)\Rightarrow c)$, et $c)\Rightarrow a)$ lorsque les $M_i$ sont séparés pour la topologie $\mathfrak J$-préadique, ont déjà été prouvées ($0$, 15.1.9), et ne sont données que pour mémoire. On a aussi montré (III, 1.1.4 et 1.1.3.3) que $a)$ entraîne $b)$, et $b)$ implique trivialement $b')$. Il reste donc à voir que lorsque tout module quotient d'un sous-module de $M$ est séparé pour la topologie $\mathfrak J$-préadique, $b')$ entraîne $a)$. Raisonnons par récurrence sur $n$. Pour $n=1$, l'assertion résulte trivialement des définitions, car alors $\operatorname{H}_1(\mathbf f,M)$ n'est autre que le noyau de l'homothétie $z\mapsto f_1z$ dans $M$ (III, 1.1.1 et 1.1.2). Supposons donc $n\geq2$, et posons $\mathbf f'=(f_i)_{1\leq i\leq n-1}$; avec les notations de (III, 1.1.2), on a $\operatorname{K}_\bullet(\mathbf f,M)=\operatorname{K}_\bullet(f_n)\otimes_A\operatorname{K}_\bullet(\mathbf f',M)$. On a donc (III, 1.1.4.1) la suite exacte
\begin{equation}
  0\longrightarrow\operatorname{H}_0(f_n,\operatorname{H}_1(\mathbf f',M))\longrightarrow\operatorname{H}_1(\mathbf f,M)\longrightarrow\operatorname{H}_1(f_n,\operatorname{H}_0(\mathbf f',M))\longrightarrow0.
  \tag{19.5.1.1}\label{IV.19.5.1.1}
\end{equation}

\oldpage[IV-4]{205}

La relation $\operatorname{H}_1(\mathbf f,M)=0$ implique donc $\operatorname{H}_0(f_n,\operatorname{H}_1(\mathbf f',M))=0$ et $\operatorname{H}_1(f_n,\operatorname{H}_0(\mathbf f',M))=0$. La première de ces deux relations signifie que l'on a $f_n\operatorname{H}_1(\mathbf f',M)=\operatorname{H}_1(\mathbf f',M)$ (III, 1.1.3.5). Or, par définition (III, 1.1.1), $\operatorname{H}_1(\mathbf f',M)$ est isomorphe à un quotient d'un sous-module $N$ de $M^{n-1}$; prenant sur $M^{n-1}$ la filtration des $M^j$ pour $j\leq n-1$, sur $N$ la filtration induite et sur $\operatorname{H}_1(\mathbf f',M)$ la filtration quotient de celle de $N$, on obtient sur $\operatorname{H}_1(\mathbf f',M)$ une filtration finie dont les quotients sont isomorphes à des quotients de sous-modules de $M$, donc sont séparés pour la topologie $\mathfrak J$-préadique par hypothèse. Comme la relation $f_n\operatorname{H}_1(\mathbf f',M)=\operatorname{H}_1(\mathbf f',M)$ implique \emph{a fortiori} $\mathfrak J\operatorname{H}_1(\mathbf f',M)=\operatorname{H}_1(\mathbf f',M)$, on déduit de l'hypothèse et des remarques précédentes que $\operatorname{H}_1(\mathbf f',M)=0$. L'hypothèse de récurrence prouve alors que la suite $\mathbf f'$ est M-régulière. D'autre part, la relation $\operatorname{H}_1(f_n,\operatorname{H}_0(\mathbf f',M))=0$ signifie que $f_n$ est $(M/(\sum_{i=1}^{n-1}f_iM))$-régulier, donc la suite $\mathbf f$ est M-régulière.

\begin{corollary}[19.5.2]
\label{IV.19.5.2}
Supposons que $A$ soit un anneau noethérien, que les $f_i$ ($1\leq i\leq n$) appartiennent au radical de $A$ et que $M$ soit un $A$-module de type fini. Alors les quatre conditions $a)$, $b)$, $b')$, $c)$ de (19.5.1) sont équivalentes.
\end{corollary}

On sait en effet alors ($0_{\mathrm I}$, 7.3.5) que tout $A$-module de type fini est séparé pour la topologie $\mathfrak J$-préadique.

\begin{proposition}[19.5.3]
\label{IV.19.5.3}
Soient
\[
  0\longrightarrow M'\longrightarrow M\longrightarrow M''\longrightarrow0
\]
une suite exacte de $A$-modules, $\mathbf f=(f_i)_{1\leq i\leq n}$ une suite M-régulière, $\mathfrak J=\sum_{i=1}^n f_iA$. Considérons les propriétés suivantes:
\begin{enumerate}
  \item[a)] La suite $\mathbf f$ est $M''$-régulière.
  \item[b)] La filtration $\mathfrak J$-préadique de $M'$ est induite sur $M'$ par la filtration $\mathfrak J$-préadique de $M$ (autrement dit, l'homomorphisme canonique $\operatorname{gr}_{\mathfrak J}^*(M')\to\operatorname{gr}_{\mathfrak J}^*(M)$ est injectif).
  \item[c)] L'homomorphisme canonique $M'/(\sum_{i=1}^n f_iM')\to M/(\sum_{i=1}^n f_iM)$ est injectif.
\end{enumerate}
Alors on a les implications $a)\Rightarrow b)\Rightarrow c)$, et $a)$ entraîne en outre que la suite $\mathbf f$ est $M'$-régulière. Si tout quotient d'un sous-module de $M''$ est séparé pour la topologie $\mathfrak J$-préadique, les conditions $a)$, $b)$ et $c)$ sont équivalentes.
\end{proposition}

Il est clair que $c)$ est conséquence de $b)$. Prouvons que sous les hypothèses de la dernière assertion, $c)$ entraîne $a)$: puisque $\mathbf f$ est M-régulière, on a, par (19.5.1), $\operatorname{H}_1(\mathbf f,M)=0$, d'où une suite exacte, portion de la suite exacte d'homologie
\[
  0\longrightarrow\operatorname{H}_1(\mathbf f,M'')\longrightarrow\operatorname{H}_0(\mathbf f,M')\longrightarrow\operatorname{H}_0(\mathbf f,M).
\]
Or, la condition $c)$ exprime que l'homomorphisme $\operatorname{H}_0(\mathbf f,M')\to\operatorname{H}_0(\mathbf f,M)$ est injectif (III, 1.1.3.5); elle équivaut donc à $\operatorname{H}_1(\mathbf f,M'')=0$, d'où, en vertu de l'hypothèse de séparation et de (19.5.1), le fait que $\mathbf f$ est $M''$-régulière.

Montrons ensuite que $a)$ entraîne $c)$ et le fait que $\mathbf f$ est $M'$-régulière, en procédant par récurrence sur $n$. Pour $n=1$, $f_1$ étant M-régulier est aussi $M'$-régulier et la propriété $c)$ n'est autre que le lemme (3.4.1.4). Pour $n>1$, l'hypothèse de récurrence montre que

\oldpage[IV-4]{206}

si l'on pose $\mathbf f'=(f_1,\ldots,f_{n-1})$, la suite $\mathbf f'$ est $M'$-régulière et l'on a, en vertu de $c)$ appliqué à $\mathbf f'$, une suite exacte
\[
  0\longrightarrow M'/(\textstyle\sum_{i=1}^{n-1}f_iM')\longrightarrow M/(\textstyle\sum_{i=1}^{n-1}f_iM)\longrightarrow M''/(\textstyle\sum_{i=1}^{n-1}f_iM'')\longrightarrow0.
\]
Par hypothèse, $f_n$ est $(M/(\sum_{i=1}^{n-1}f_iM))$-régulier et $(M''/(\sum_{i=1}^{n-1}f_iM''))$-régulier, donc le même raisonnement montre d'une part que $f_n$ est $(M'/(\sum_{i=1}^{n-1}f_iM'))$-régulier, et d'autre part, en vertu de (3.4.1.4), que la suite
\[
  0\longrightarrow M'/(\textstyle\sum_{i=1}^{n}f_iM')\longrightarrow M/(\textstyle\sum_{i=1}^{n}f_iM)\longrightarrow M''/(\textstyle\sum_{i=1}^{n}f_iM'')\longrightarrow0
\]
est exacte, d'où $c)$.

Montrons enfin que $a)$ entraîne $b)$. Comme la suite $\mathbf f$ est alors M-régulière et $M'$-régulière, donc M-quasi-régulière et $M'$-quasi-régulière, on a des isomorphismes canoniques
\[
  \operatorname{gr}_{\mathfrak J}^*(M')\xrightarrow{\sim}\operatorname{gr}_{\mathfrak J}^0(M')[T_1,\ldots,T_n],
  \qquad
  \operatorname{gr}_{\mathfrak J}^*(M)\xrightarrow{\sim}\operatorname{gr}_{\mathfrak J}^0(M)[T_1,\ldots,T_n]
\]
($0$, 15.1.7), et comme on a vu ci-dessus que $\operatorname{gr}_{\mathfrak J}^0(M')\to\operatorname{gr}_{\mathfrak J}^0(M)$ est injectif, il en est de même de $\operatorname{gr}_{\mathfrak J}^*(M')\to\operatorname{gr}_{\mathfrak J}^*(M)$.

On notera encore que l'équivalence des conditions $a)$, $b)$, $c)$ de (19.5.3) est valable en particulier lorsque $A$ est noethérien, $M''$ un $A$-module de type fini et que les $f_i$ appartiennent au radical de $A$.

\begin{env}[19.5.4]
\label{IV.19.5.4}
Dans la suite de ce numéro, nous conservons les notations $A$, $\mathbf f$, $M$, $\mathfrak J$ de (19.5.1), et nous supposons de plus $M$ muni d'une filtration décroissante $(M_k)$ formée de sous-$A$-modules, telle que $M_0=M$. Rappelons ($0$, 15.1.5) que l'on définit alors sur $M$ une seconde filtration décroissante formée des sous-modules
\[
  M_k'=M_k+\mathfrak JM_{k-1}+\cdots+\mathfrak J^{k-1}M_1+\mathfrak J^kM_0
\]
et que, si $\operatorname{gr}_\bullet(M)$ et $\operatorname{gr}_\bullet'(M)$ sont les $A$-modules gradués associés aux filtrations $(M_k)$ et $(M_k')$ respectivement, on définit un homomorphisme gradué surjectif de degré $0$ ($0$, 15.1.5.2)
\[
  \psi_M:(\operatorname{gr}_\bullet(M)\otimes_A(A/\mathfrak J))[T_1,\ldots,T_n]\longrightarrow\operatorname{gr}_\bullet'(M).
\]
Rappelons aussi que l'on a l'homomorphisme canonique surjectif ($0$, 15.1.1.1)
\[
  \varphi_{\operatorname{gr}_\bullet(M)}:(\operatorname{gr}_\bullet(M)\otimes_A(A/\mathfrak J))[T_1,\ldots,T_n]\longrightarrow\operatorname{gr}_{\mathfrak J}^*(\operatorname{gr}_\bullet(M)).
\]
\end{env}

\begin{theorem}[19.5.5]
\label{IV.19.5.5}
Avec les notations de (19.5.4), considérons les propriétés suivantes:
\begin{enumerate}
  \item[a)] La suite $\mathbf f$ est $\operatorname{gr}_\bullet(M)$-régulière.
  \item[b)] L'homomorphisme canonique $\psi_M$ est bijectif.
  \item[c)] L'homomorphisme canonique $\varphi_{\operatorname{gr}_\bullet(M)}$ est bijectif (autrement dit ($0$, 15.1.7), la suite $\mathbf f$ est $\operatorname{gr}_\bullet(M)$-quasi-régulière).
\end{enumerate}

\oldpage[IV-4]{207}

On a alors les implications\footnote{L'implication $b)\Rightarrow c)$, sans hypothèse de séparation sur les $\operatorname{gr}_k(M)$, est due à P. Deligne, dont nous reproduisons ci-dessous la démonstration.}
\[
  a)\Rightarrow b)\Rightarrow c).
\]
En outre, si, pour tout entier $k\geq0$, les modules quotients des sous-modules de $\operatorname{gr}_k(M)$ sont séparés pour la topologie $\mathfrak J$-préadique (ce qui est le cas lorsque $A$ est noethérien, que les $f_i$ sont dans le radical de $A$ et que les $\operatorname{gr}_k(M)$ sont des $A$-modules de type fini), alors les conditions $a)$, $b)$ et $c)$ sont équivalentes.
\end{theorem}

L'implication $a)\Rightarrow b)$ a été démontrée dans ($0$, 15.1.8); l'implication $c)\Rightarrow a)$ sous l'hypothèse de séparation est un cas particulier de ($0$, 15.1.9). Il reste à prouver que $b)$ implique $c)$, ce qui va se faire en plusieurs étapes.

Nous désignerons par $\mathfrak J$ l'idéal $\sum_{i=1}^n f_iA$ engendré par les $f_i$; pour toute suite $\mathbf q=(q_1,\ldots,q_n)$ de $n$ entiers $\geq0$, on posera $|\mathbf q|=\sum_{i=1}^nq_i$ et $\mathbf f^{\mathbf q}=f_1^{q_1}f_2^{q_2}\cdots f_n^{q_n}$.

\begin{env}[19.5.5.1]
\label{IV.19.5.5.1}
Pour que l'application $\psi_M$ soit injective (et par suite bijective), il faut et il suffit que pour $p\geq0$ et $q\geq0$ la condition suivante soit vérifiée:

$(\mathrm I_{q,p})$ Pour toute famille $(x_{\mathbf q}^p)_{|\mathbf q|=q}$ d'éléments de $M_p$, la relation
\[
  \sum_{|\mathbf q|=q}\mathbf f^{\mathbf q}x_{\mathbf q}^p\in\sum_{i=1}^q\mathfrak J^{q-i}M_{p+i}+M_{p+q+1}'
\]
implique $x_{\mathbf q}^p\in M_{p+1}+\mathfrak JM_p$ pour tout $\mathbf q$ tel que $|\mathbf q|=q$.

On procède comme dans ($0$, 15.1.6) en considérant le sous-$A$-module $Q_k$ (resp. $Q_k'$) des termes de degré $k$ au premier (resp. second) membre de ($0$, 15.1.5.2). On munit $Q_k$ de la filtration
\[
  (Q_k)_i=\sum_{j\leq k-i}\left(\sum_{|\mathbf l|=j}(\operatorname{gr}_{k-j}(M)\otimes_A(A/\mathfrak J))T^{\mathbf l}\right)
\]
et $Q_k'$ de la filtration image formée des $\psi_M((Q_k)_i)=(Q_k')_i$; il suffit encore de prouver que les homomorphismes $(\psi_M)_{ki}:\operatorname{gr}_i(Q_k)\to\operatorname{gr}_i(Q_k')$ sont injectifs. Or, on a
\[
  \operatorname{gr}_i(Q_k)=\sum_{|\mathbf l|=k-i}((M_i/M_{i+1})\otimes_A(A/\mathfrak J))T^{\mathbf l}
  =\sum_{|\mathbf l|=k-i}(M_i/(\mathfrak JM_i+M_{i+1}))T^{\mathbf l};
\]

\oldpage[IV-4]{208}

d'autre part, $(Q_k')_{i+1}$ est l'image de $\sum_{h=0}^{k-i-1}\mathfrak J^{k-i-h-1}M_{i+h+1}$ dans $M_k'/M_{k+1}'$. Écrire que $(\psi_M)_{ki}$ est injectif équivaut donc à écrire la condition $(\mathrm I_{k-i,i})$; d'où notre assertion.
\end{env}

\begin{env}[19.5.5.2]
\label{IV.19.5.5.2}
Pour que la suite $\mathbf f$ soit $\operatorname{gr}_\bullet(M)$-quasi-régulière, il suffit que, pour $p\geq0$ et $q\geq0$, la condition suivante soit vérifiée:

$(\mathrm S_{q,p})$ Pour toute famille $(x_{\mathbf q}^p)_{|\mathbf q|=q}$ d'éléments de $M_p$, la relation
\[
  \sum_{|\mathbf q|=q}\mathbf f^{\mathbf q}x_{\mathbf q}^p\in M_{p+1}
\]
implique $x_{\mathbf q}^p\in M_{p+1}+\mathfrak JM_p$ pour tout $\mathbf q$ tel que $|\mathbf q|=q$.

Par définition, dire que $\mathbf f$ est $\operatorname{gr}_\bullet(M)$-quasi-régulière signifie que, pour tout $p\geq0$ et tout $q\geq0$, une relation
\[
  \sum_{|\mathbf q|=q}\mathbf f^{\mathbf q}x_{\mathbf q}^p\in M_{p+1}+\mathfrak J^{q+1}M_p
\]
pour une famille $(x_{\mathbf q}^p)_{|\mathbf q|=q}$ d'éléments de $M_p$ entraîne $x_{\mathbf q}^p\in M_{p+1}+\mathfrak JM_p$ pour tout $\mathbf q$ tel que $|\mathbf q|=q$. Or l'hypothèse sur la suite $(x_{\mathbf q}^p)$ signifie qu'il existe une famille $(y_{\mathbf q}^p)_{|\mathbf q|=q}$ d'éléments de $\mathfrak JM_p$ telle que l'on ait $\sum_{|\mathbf q|=q}\mathbf f^{\mathbf q}(x_{\mathbf q}^p-y_{\mathbf q}^p)\in M_{p+1}$. La condition $(\mathrm S_{q,p})$ entraîne alors $x_{\mathbf q}^p-y_{\mathbf q}^p\in M_{p+1}+\mathfrak JM_p$ pour tout $\mathbf q$ tel que $|\mathbf q|=q$, d'où $x_{\mathbf q}^p\in M_{p+1}+\mathfrak JM_p$; on voit donc que les conditions $(\mathrm S_{q,p})$ entraînent que la suite $\mathbf f$ est $\operatorname{gr}_\bullet(M)$-quasi-régulière.
\end{env}

Il suffit donc de prouver la proposition suivante:

\begin{env}[19.5.5.3]
\label{IV.19.5.5.3}
Pour tout $r\geq0$, si les conditions $(\mathrm I_{q,p})$ sont vérifiées pour $p+q<r$, alors les conditions $(\mathrm S_{q,p})$ sont vérifiées pour $p+q<r$.

Introduisons les conditions

$(\mathrm A_{q,p})$ Pour toute famille $(x_{\mathbf q}^p)_{|\mathbf q|=q}$ d'éléments de $M_p$, la relation
\[
  \sum_{|\mathbf q|=q}\mathbf f^{\mathbf q}x_{\mathbf q}^p\in M_{p+1}
\]
implique
\[
  \sum_{|\mathbf q|=q}\mathbf f^{\mathbf q}x_{\mathbf q}^p\in\sum_{i=1}^q\mathfrak J^{q-i}M_{p+i}.
\]

Il est clair que la conjonction de $(\mathrm A_{q,p})$ et de $(\mathrm I_{q,p})$ implique $(\mathrm S_{q,p})$. D'autre part, les conditions $(\mathrm A_{1,p})$ et $(\mathrm S_{0,p})$ sont triviales. Par suite, (19.5.5.3) résultera de la proposition suivante:
\end{env}

\begin{env}[19.5.5.4]
\label{IV.19.5.5.4}
Quels que soient $p\geq0$ et $q\geq1$, la condition $(\mathrm A_{q+1,p})$ est entraînée par la conjonction des conditions $(\mathrm A_{q,p})$, $(\mathrm I_{q-1,p+1})$, $(\mathrm I_{q-2,p+2})$, $\ldots$, $(\mathrm I_{0,p+q})$.

En effet, une fois prouvée cette proposition, les conditions $(\mathrm I_{q,p})$ supposées vérifiées pour $p+q<r$ entraîneront, pour chaque $p\leq r$, la condition $(\mathrm A_{q,p})$ pour $1\leq q<r-p$, par récurrence sur $q$.

Prouvons donc (19.5.5.4). Notons que la condition $(\mathrm A_{q,p})$ équivaut aussi à
\end{env}

\begin{env}[19.5.5.5]
\label{IV.19.5.5.5}
\[
  \mathfrak J^qM_p\cap M_{p+1}\subset\sum_{i=1}^q\mathfrak J^{q-i}M_{p+i}.
\]

Soit donc $m\in\mathfrak J^{q+1}M_p\cap M_{p+1}\subset\mathfrak J^qM_p\cap M_{p+1}$. Appliquant la condition $(\mathrm A_{q,p})$, on voit qu'il existe, pour chaque $i$ tel que $1\leq i\leq q$, une famille $(y_{\mathbf q}^{p+i})_{|\mathbf q|=q-i}$ d'éléments de $M_{p+i}$ telle que
\[
  m=\sum_{i=1}^q\left(\sum_{|\mathbf q|=q-i}\mathbf f^{\mathbf q}y_{\mathbf q}^{p+i}\right).
\]
Supposons que, pour $1\leq i<j$ ($j\geq1$), on ait prouvé que
\[
  y_{\mathbf q}^{p+i}\in\mathfrak JM_{p+i}+M_{p+i+1}\qquad(\text{pour }|\mathbf q|=q-i).
\]
On en déduit par définition (19.5.4)
\[
  m-\sum_{i<j}\left(\sum_{|\mathbf q|=q-i}\mathbf f^{\mathbf q}y_{\mathbf q}^{p+i}\right)\in M_{p+q+1}'.
\]
Puisque par hypothèse $(\mathrm I_{q-j,p+j})$ est vérifiée, on en déduit
\[
  y_{\mathbf q}^{p+j}\in\mathfrak JM_{p+j}+M_{p+j+1}
\]

\oldpage[IV-4]{209}

et, par récurrence sur $j$, on voit donc que cette condition est vérifiée pour tout $j$ tel que $1\leq j\leq q$. On a donc
\[
  m\in\sum_{i=1}^q\mathfrak J^{q-i}(\mathfrak JM_{p+i}+M_{p+i+1})=\sum_{i=1}^{q+1}\mathfrak J^{q-i+1}M_{p+i}
\]
et on a ainsi prouvé $(\mathrm A_{q+1,p})$ et terminé la démonstration de (19.5.5).
\end{env}

\begin{remarks}[19.5.6]
\label{IV.19.5.6}
\begin{enumerate}
  \item[(i)] Les résultats de ce numéro gardent un sens et sont valables lorsque, au lieu d'un $A$-module $M$ et d'éléments $f_i\in A$, on prend un objet $M$ d'une catégorie abélienne quelconque, et $n$ endomorphismes $f_i$ de $M$ qui commutent deux à deux; les démonstrations s'adaptent sans peine à ce cas plus général. L'hypothèse de séparation pour un $A$-module $N$ relativement à la topologie $\mathfrak J$-préadique doit ici être remplacée par l'hypothèse que tout sous-objet de $N$ contenu dans tous les $\sum_i f_i^r(N)$ ($r$ entier arbitraire) est nécessairement nul. La même remarque s'applique aux résultats de 19.7.
  \item[(ii)] Supposons que $A$ soit un anneau \emph{gradué} à degrés $\geq0$. Si $M$ (resp. chaque $\operatorname{gr}_p(M)$) est un $A$-module gradué à degrés bornés inférieurement et si les $f_i$ ne contiennent pas de composant homogène de degré $0$, l'hypothèse de séparation de ($0$, 15.1.9) est \emph{ipso facto} vérifiée pour $M$ (resp. chaque $\operatorname{gr}_p(M)$), et on voit donc que dans ce cas on a l'implication $c)\Rightarrow a)$ dans (19.5.1) (resp. (19.5.5)).
  \item[(iii)] Rappelons que, sans hypothèse de séparation, l'implication $c)\Rightarrow b)$ n'est plus valable même pour $n=1$ ($0$, 15.1.12, (iii)).
\end{enumerate}
\end{remarks}

\subsection{Suites régulières relativement à un module filtré quotient.}
\EGAReaderTarget{ega4.c04.subsec.19-6}{subsection:IV.19.6}

\begin{env}[19.6.1]
\label{IV.19.6.1}
Soient $A$ un anneau, $\mathbf f=(f_1,\ldots,f_n)$ une suite finie d'éléments de $A$, $\mathfrak J=f_1A+\cdots+f_nA$ l'idéal qu'elle engendre, et considérons une suite exacte de $A$-modules
\[
  0\longrightarrow R\xrightarrow{j}N\xrightarrow{p}M\longrightarrow0
\]
où l'on suppose que $N$ est muni d'une filtration décroissante $(N_k)$ formée de sous-$A$-modules, avec $N_0=N$. On pose $R_k=R\cap N_k$ (filtration induite par $(N_k)$), $M_k=p(N_k)$ (filtration quotient de $(N_k)$), et on désigne par $\operatorname{gr}_\bullet(N)$, $\operatorname{gr}_\bullet(R)$ et $\operatorname{gr}_\bullet(M)$ les $A$-modules gradués associés à ces trois filtrations. On pose d'autre part pour tout $k$,
\[
  N_k'=N_k+\mathfrak JN_{k-1}+\cdots+\mathfrak J^{k-1}N_1+\mathfrak J^kN_0
\]
de sorte que $M_k'=p(N_k')=M_k+\mathfrak JM_{k-1}+\cdots+\mathfrak J^{k-1}M_1+\mathfrak J^kM_0$; on pose d'autre part $R_k'=R\cap N_k'$ (filtration \emph{induite} par $(N_k')$), et on note $\operatorname{gr}_\bullet'(N)$, $\operatorname{gr}_\bullet'(M)$ et $\operatorname{gr}_\bullet'(R)$ les $A$-modules gradués associés à ces trois filtrations; on a donc un diagramme commutatif de suites exactes
\begin{equation}
\xymatrix{
  0 \ar[r] & \operatorname{gr}_\bullet(R) \ar[r]^-{\operatorname{gr}(j)} \ar[d] & \operatorname{gr}_\bullet(N) \ar[r]^-{\operatorname{gr}(p)} \ar[d] & \operatorname{gr}_\bullet(M) \ar[r] \ar[d] & 0 \\
  0 \ar[r] & \operatorname{gr}_\bullet'(R) \ar[r]^-{\operatorname{gr}'(j)} & \operatorname{gr}_\bullet'(N) \ar[r]^-{\operatorname{gr}'(p)} & \operatorname{gr}_\bullet'(M) \ar[r] & 0
}
\tag{19.6.1.1}\label{IV.19.6.1.1}
\end{equation}

\oldpage[IV-4]{210}

où les flèches verticales sont les homomorphismes canoniques d'un module gradué associé à une filtration dans le module gradué associé à une filtration moins fine.

On notera que si l'on pose
\[
  R_k''=R_k+\mathfrak JR_{k-1}+\cdots+\mathfrak J^{k-1}R_1+\mathfrak J^kR_0
\]
on a évidemment $R_k''\subset R_k'$, mais les filtrations $(R_k'')$ et $(R_k')$ sont en général distinctes; on notera $\operatorname{gr}_\bullet''(R)$ le $A$-module gradué associé à la filtration $(R_k'')$.
\end{env}

\begin{env}[19.6.2]
\label{IV.19.6.2}
On a défini en ($0$, 15.1.5.2) les homomorphismes surjectifs gradués de degré $0$
\begin{align*}
  \psi_N&:(\operatorname{gr}_\bullet(N)\otimes_A(A/\mathfrak J))[T_1,\ldots,T_n]\longrightarrow\operatorname{gr}_\bullet'(N),\\
  \psi_M&:(\operatorname{gr}_\bullet(M)\otimes_A(A/\mathfrak J))[T_1,\ldots,T_n]\longrightarrow\operatorname{gr}_\bullet'(M),\\
  \psi_R&:(\operatorname{gr}_\bullet(R)\otimes_A(A/\mathfrak J))[T_1,\ldots,T_n]\longrightarrow\operatorname{gr}_\bullet''(R),
\end{align*}
dont les deux premiers sont déduits des deux dernières flèches verticales de (19.6.1.1).

Comme la filtration $(R_k'')$ est plus fine que $(R_k')$ sur $R$, on a un homomorphisme canonique $\gamma:\operatorname{gr}_\bullet''(R)\to\operatorname{gr}_\bullet'(R)$; nous désignerons par $\psi_R'$ l'homomorphisme composé $\gamma\circ\psi_R$. Il résulte alors aussitôt des définitions que l'on a un diagramme commutatif
\begin{equation}
\xymatrix{
  & 0 \ar[d] \\
  (\operatorname{gr}_\bullet(R)\otimes_A(A/\mathfrak J))[T_1,\ldots,T_n] \ar[r]^-{\psi_R'} \ar[d]^{j'} & \operatorname{gr}_\bullet'(R) \ar[d]^{\operatorname{gr}'(j)} \\
  (\operatorname{gr}_\bullet(N)\otimes_A(A/\mathfrak J))[T_1,\ldots,T_n] \ar[r]^-{\psi_N} \ar[d]^{p'} & \operatorname{gr}_\bullet'(N) \ar[d]^{\operatorname{gr}'(p)} \\
  (\operatorname{gr}_\bullet(M)\otimes_A(A/\mathfrak J))[T_1,\ldots,T_n] \ar[r]^-{\psi_M} \ar[d] & \operatorname{gr}_\bullet'(M) \ar[d] \\
  0 & 0
}
\tag{19.6.2.1}\label{IV.19.6.2.1}
\end{equation}
où les colonnes sont exactes.
\end{env}

\begin{proposition}[19.6.3]
\label{IV.19.6.3}
Avec les notations précédentes, supposons que la suite $\mathbf f$ soit $\operatorname{gr}_\bullet(N)$-régulière. Considérons les propriétés suivantes:
\begin{enumerate}
  \item[a)] $\mathbf f$ est $\operatorname{gr}_\bullet(M)$-régulière.
  \item[b)] $\psi_M$ est bijectif.
  \item[c)] $\psi_R'$ est surjectif (autrement dit $\operatorname{gr}_\bullet'(R)$ est engendré comme $\operatorname{gr}_{\mathfrak J}^*(A)$-module par $\operatorname{Im}(\operatorname{gr}_\bullet(R)\to\operatorname{gr}_\bullet'(R))$).
  \item[d)] $\psi_R'$ est injectif.
  \item[d$'$)] L'homomorphisme $(\psi_R')_0:\operatorname{gr}_\bullet(R)\otimes_A(A/\mathfrak J)\to\operatorname{gr}_\bullet'(R)$ obtenu en restreignant $\psi_R'$ aux polynômes de degré $0$, est injectif.
  \item[e)] $\psi_R'$ est bijectif.
  \item[f)] $j'$ est injectif.

\oldpage[IV-4]{211}

  \item[f$'$)] L'homomorphisme $\operatorname{gr}_\bullet(R)\otimes_A(A/\mathfrak J)\to\operatorname{gr}_\bullet(N)\otimes_A(A/\mathfrak J)$ obtenu en restreignant $j'$ aux polynômes de degré $0$, est injectif.
  \item[g)] La filtration $\mathfrak J$-préadique de $\operatorname{gr}_\bullet(R)$ est induite par celle de $\operatorname{gr}_\bullet(N)$.
  \item[h)] Les filtrations $(R_k')$ et $(R_k'')$ sur $R$ sont identiques.
\end{enumerate}
On a alors les implications suivantes:
\[
\xymatrix@C=2.4em@R=1.8em{
  a) \ar@{=>}[r] \ar@{=>}[dr] & e) \ar@{=>}[r] \ar@{=>}[dr] & b) \ar@{<=>}[r] & c) \ar@{<=>}[r] & h) \\
  & g) \ar@{=>}[r] & d) \ar@{<=>}[r] & d') \ar@{<=>}[r] & f) \ar@{<=>}[r] & f')
}
\]
En outre, lorsque tout module quotient d'un sous-module d'un $\operatorname{gr}_k(M)$ est séparé pour la topologie $\mathfrak J$-préadique (ce qui sera le cas lorsque $A$ est noethérien, que les $f_i$ appartiennent au radical de $A$ et que les $\operatorname{gr}_k(M)$ sont des $A$-modules de type fini), alors les conditions $a)$ à $h)$ sont toutes équivalentes.
\end{proposition}

Notons que d'après (19.5.5), de l'hypothèse que $\mathbf f$ est $\operatorname{gr}_\bullet(N)$-régulière, il résulte que $\psi_N$ est \emph{bijectif}; comme $\operatorname{gr}'(j)$ est injectif, l'équivalence de $d)$ et $f)$ résulte du diagramme (19.6.2.1), ainsi que celle de $d')$ et $f')$; par ailleurs $f)$ et $f')$ sont trivialement équivalentes, ce qui montre l'équivalence de $d)$, $d')$, $f)$ et $f')$.

Comme on sait déjà que $\psi_M$ est surjectif et $\psi_N$ bijectif, l'équivalence de $b)$ et $c)$ découle aussi du diagramme (19.6.2.1) par un cas particulier du lemme des cinq.

La relation $\psi_R'=\gamma\circ\psi_R$ (19.6.2) et le fait que $\psi_R$ est surjectif montrent que la condition $c)$ équivaut à dire que $\gamma$ est surjectif, ce qui équivaut encore à la condition $h)$; on a donc prouvé l'équivalence de $b)$, $c)$ et $h)$.

Il est trivial que $e)$ équivaut à la conjonction de $c)$ et $d)$, donc aussi de $b)$ et $d)$. Notons maintenant que $a)$ implique $b)$ et que $b)$ implique $a)$ sous l'hypothèse de séparation (19.5.5). D'autre part, les implications $a)\Rightarrow g)\Rightarrow f')$, et l'implication $f')\Rightarrow a)$ sous l'hypothèse de séparation, résultent de (19.5.3). On a donc prouvé que $a)$ entraîne $e)$ (équivalente à la conjonction de $b)$ et $f')$), et aussi que toutes les conditions sont équivalentes sous l'hypothèse de séparation. C.Q.F.D.

\begin{corollary}[19.6.4]
\label{IV.19.6.4}
Soient $\mathfrak R$ un idéal de $A$, $\mathfrak L=\mathfrak J+\mathfrak R$. Supposons que la filtration $(N_k)$ (donc aussi $(M_k)$) soit la filtration $\mathfrak R$-préadique, de sorte que les filtrations $(M_k')$ et $(N_k')$ sont les filtrations $\mathfrak L$-préadiques. Considérons $\operatorname{gr}_\bullet(R)$ (resp. $\operatorname{gr}_\bullet'(R)$) comme un module gradué sur $\operatorname{gr}_{\mathfrak R}^*(A)$ (resp. $\operatorname{gr}_{\mathfrak L}^*(A)$).
\begin{enumerate}
  \item[(i)] Soit $S$ une partie de $\operatorname{gr}_\bullet(R)$ qui engendre $\operatorname{gr}_\bullet(R)$ comme $\operatorname{gr}_{\mathfrak R}^*(A)$-module. Alors la condition $c)$ de (19.6.3) équivaut à la condition suivante:
  \item[c$'$)] L'image de $S$ dans $\operatorname{gr}_\bullet'(R)$ par $\psi_R'$ engendre $\operatorname{gr}_\bullet'(R)$ comme $\operatorname{gr}_{\mathfrak L}^*(A)$-module.
  \item[(ii)] Supposons vérifiée la condition $e)$ de (19.6.3), et de plus que les quotients des $\operatorname{gr}_k(R)$ soient séparés pour la topologie $\mathfrak J$-préadique (ce qui sera le cas lorsque $A$ est noethérien, les $f_i$ dans le radical de $A$ et les $\operatorname{gr}_k(N)$ des $A$-modules de type fini). Alors, pour qu'une partie $S$ de $\operatorname{gr}_\bullet(R)$ formée d'éléments homogènes engendre $\operatorname{gr}_\bullet(R)$ comme $\operatorname{gr}_{\mathfrak R}^*(A)$-module, il faut et il suffit que son image par $\psi_R'$ engendre $\operatorname{gr}_\bullet'(R)$ comme $\operatorname{gr}_{\mathfrak L}^*(A)$-module.
\end{enumerate}
\end{corollary}

\oldpage[IV-4]{212}

\begin{enumerate}
  \item[(i)] Considérons l'homomorphisme de degré $0$
  \[
    \psi_A:(\operatorname{gr}_{\mathfrak R}^*(A)\otimes_A(A/\mathfrak J))[T_1,\ldots,T_n]\longrightarrow\operatorname{gr}_\bullet'(A)=\operatorname{gr}_{\mathfrak L}^*(A)
  \]
  ($0$, 15.1.5.2). Le fait que cet homomorphisme soit \emph{surjectif} entraîne que la sous-$\operatorname{gr}_{\mathfrak L}^*(A)$-algèbre de $\operatorname{gr}_\bullet'(R)$ engendrée par $S'=\psi_R'(S)$ n'est autre que $\operatorname{Im}(\psi_R')$, d'où l'équivalence de $c)$ et $c')$.
  \item[(ii)] Puisque $e)$ est vérifiée, il en est de même de $c)$, et en vertu de (i), il reste à prouver la suffisance de la condition de l'énoncé. Or, comme $\psi_R'$ est bijectif, dire que $S'$ engendre $\operatorname{gr}_\bullet'(R)$ considéré comme $\operatorname{gr}_{\mathfrak L}^*(A)$-module équivaut à dire que $S$ engendre $(\operatorname{gr}_\bullet(R)\otimes_A(A/\mathfrak J))[T_1,\ldots,T_n]$ comme $\operatorname{gr}_{\mathfrak L}^*(A)$-module, et en vertu de la surjectivité de $\psi_A$, on peut dans cet énoncé remplacer l'anneau $\operatorname{gr}_{\mathfrak L}^*(A)$ par $(\operatorname{gr}_{\mathfrak R}^*(A)\otimes_A(A/\mathfrak J))[T_1,\ldots,T_n]$. Il revient évidemment au même de dire que $S$ engendre $\operatorname{gr}_\bullet(R)\otimes_A(A/\mathfrak J)$ en tant que module sur l'anneau $\operatorname{gr}_{\mathfrak R}^*(A)\otimes_A(A/\mathfrak J)$, ou aussi en tant que $\operatorname{gr}_{\mathfrak R}^*(A)$-module. Or, si $T$ est le sous-$\operatorname{gr}_{\mathfrak R}^*(A)$-module gradué de $\operatorname{gr}_\bullet(R)$ engendré par $S$, et si l'on pose $T'=\operatorname{gr}_\bullet(R)/T$, on a par hypothèse $T'/\mathfrak JT'=0$, ou $\mathfrak JT'=T'$; mais $T'$ est un module gradué, et chaque $T_p'$ est un quotient de $\operatorname{gr}_p(R)$, donc séparé pour la topologie $\mathfrak J$-préadique; l'hypothèse $\mathfrak JT_p'=T_p'$ entraîne donc $T_p'=0$ pour tout entier $p$, donc $T'=0$, ce qui achève de prouver le corollaire.
\end{enumerate}

\subsection{Critère de platitude normale de Hironaka.}
\EGAReaderTarget{ega4.c04.subsec.19-7}{subsection:IV.19.7}
\label{IV.19.7}

\begin{theorem}[19.7.1]
\label{IV.19.7.1}
\textup{(Hironaka).} Soient $A$ un anneau, $\mathfrak R$ un idéal de $A$, $\mathbf f=(f_1,\ldots,f_n)$ une suite d'éléments de $A$ qui est $(A/\mathfrak R)$-régulière, $\mathfrak J=f_1A+\cdots+f_nA$ l'idéal engendré par $\mathbf f$, $\mathfrak L=\mathfrak J+\mathfrak R$, $M$ un $A$-module.

On a dans ce cas
\[
  (\operatorname{gr}_{\mathfrak R}^*(M))\otimes_A(A/\mathfrak J)
  =(\operatorname{gr}_{\mathfrak R}^*(M))\otimes_A(A/\mathfrak L)
  =(\operatorname{gr}_{\mathfrak R}^*(M))\otimes_{A/\mathfrak R}(A/\mathfrak L),
\]
car
\[
  (\mathfrak R^nM/\mathfrak R^{n+1}M)\otimes_A(A/\mathfrak J)
  =\mathfrak R^nM/(\mathfrak R^{n+1}+\mathfrak J\mathfrak R^n)M
  =\mathfrak R^nM/(\mathfrak R^{n+1}+\mathfrak L\mathfrak R^n)M,
\]
et la dernière égalité résulte de Bourbaki, \emph{Alg.}, chap. II, 3\textsuperscript{e} éd., § 3, n\textsuperscript{o} 7, cor. 2 de la prop. 6.

Considérons les conditions suivantes:
\begin{enumerate}
  \item[a)] $\operatorname{gr}_{\mathfrak R}^*(M)$ est un $(A/\mathfrak R)$-module plat.
  \item[b)] $\operatorname{gr}_{\mathfrak L}^*(M)$ est un $(A/\mathfrak L)$-module plat et l'homomorphisme canonique ($0$, 15.1.5.2)
  \[
    \psi_M:((\operatorname{gr}_{\mathfrak R}^*(M))\otimes_A(A/\mathfrak L))[T_1,\ldots,T_n]\longrightarrow\operatorname{gr}_{\mathfrak L}^*(M)
  \]
  est bijectif.
  \item[c)] $(\operatorname{gr}_{\mathfrak R}^*(M))\otimes_{A/\mathfrak R}(A/\mathfrak L)$ est un $(A/\mathfrak L)$-module plat et la suite $\mathbf f$ est $(\operatorname{gr}_{\mathfrak R}^*(M))$-régulière.
  \item[d)] $\operatorname{gr}_{\mathfrak L}^*(M)$ est un $(A/\mathfrak L)$-module plat, et pour tout $\mathfrak p\in\operatorname{Ass}_{A/\mathfrak R}(A/\mathfrak L)$, $(\operatorname{gr}_{\mathfrak R}^*(M))_{\mathfrak p}$ est un $(A/\mathfrak R)_{\mathfrak p}$-module plat.
\end{enumerate}
On a alors les implications
\[
\xymatrix@C=3em@R=1.6em{
  a) \ar@{=>}[r] \ar@{=>}[dr] & c) \ar@{=>}[r] & b) \\
  & d)
}
\]
Lorsque tout quotient d'un sous-module d'un $\operatorname{gr}_{\mathfrak R}^p(M)$ ($p\geq0$) est idéalement séparé pour $\mathfrak J$ (Bourbaki, \emph{Alg. comm.}, chap. III, § 5, n\textsuperscript{o} 1), les conditions $a)$, $b)$ et $c)$ sont équivalentes.

\oldpage[IV-4]{213}

Enfin, lorsque $A$ est noethérien, que les $f_i$ appartiennent au radical de $A$, que la suite $\mathbf f$ est $\operatorname{gr}_{\mathfrak R}^*(A)$-régulière et que $M$ est un $A$-module de type fini, les conditions $a)$, $b)$, $c)$, $d)$ sont équivalentes.
\end{theorem}

L'implication $a)\Rightarrow c)$ est immédiate ($0$, 15.1.13). Si $c)$ est vérifiée, il résulte de (19.5.5) que $\psi_M$ est bijectif; en outre, $\operatorname{gr}_{\mathfrak R}^*(M)\otimes_A(A/\mathfrak L)=\operatorname{gr}_{\mathfrak R}^*(M)\otimes_{A/\mathfrak R}(A/\mathfrak L)$ est un $(A/\mathfrak L)$-module plat, donc il en est de même de $(\operatorname{gr}_{\mathfrak R}^*(M)\otimes_A(A/\mathfrak L))[T_1,\ldots,T_n]$, et par suite de $\operatorname{gr}_{\mathfrak L}^*(M)$ puisque $\psi_M$ est un $(A/\mathfrak L)$-isomorphisme; donc $c)$ entraîne $b)$. Le fait que $a)$ implique $b)$ montre aussitôt que $a)$ implique aussi $d)$.

Lorsque tout quotient d'un $\operatorname{gr}_{\mathfrak R}^p(M)$ est idéalement séparé pour $\mathfrak J$, il résulte de (19.5.5) que la condition $b)$ entraîne que $\mathbf f$ est $(\operatorname{gr}_{\mathfrak R}^*(M))$-régulière; en outre, comme $(\operatorname{gr}_{\mathfrak R}^*(M)\otimes_{A/\mathfrak R}(A/\mathfrak L))[T_1,\ldots,T_n]$, isomorphe à $\operatorname{gr}_{\mathfrak L}^*(M)$, est alors un $(A/\mathfrak L)$-module plat, il en est de même de $\operatorname{gr}_{\mathfrak R}^*(M)\otimes_{A/\mathfrak R}(A/\mathfrak L)$, qui en est un facteur direct; cela prouve donc que $b)$ entraîne alors $c)$. D'autre part, sous la même hypothèse, $c)$ entraîne $a)$ en vertu de ($0$, 15.1.21) (où on peut remplacer l'hypothèse noethérienne par l'hypothèse que $M$ est idéalement séparé pour $\mathfrak J$, en vertu de ($0_{\mathrm{III}}$, 10.2.1)).

Il reste donc à prouver que lorsque $A$ est noethérien, $M$ de type fini et que les $f_i$ appartiennent au radical de $A$, $d)$ entraîne $a)$.

Il existe par hypothèse un $A$-module \emph{libre de type fini} $N$ et une suite exacte $0\to R\to N\to M\to0$. Il suffit de prouver que $d)$ entraîne $b)$, puisque tout quotient d'un sous-module d'un $\operatorname{gr}_{\mathfrak R}^p(M)$ est alors idéalement séparé pour $\mathfrak J$, en vertu du fait que $\operatorname{gr}_{\mathfrak R}^p(M)$ est un $A$-module de type fini, que $\mathfrak J$ est contenu dans le radical de $A$, et que $A$ est noethérien (Bourbaki, \emph{Alg. comm.}, chap. III, § 5, n\textsuperscript{o} 1); en d'autres termes, il s'agit de voir que $\psi_M$ est \emph{bijectif}.

Comme par hypothèse la suite $\mathbf f$ est $\operatorname{gr}_{\mathfrak R}^*(A)$-régulière, elle est aussi $\operatorname{gr}_{\mathfrak R}^*(N)$-régulière; on peut par suite appliquer (19.6.3), car tout module quotient d'un sous-module d'un $\operatorname{gr}_{\mathfrak R}^k(M)$ est alors séparé pour la topologie $\mathfrak J$-préadique puisque les $f_i$ appartiennent au radical de $A$ par hypothèse. Le diagramme commutatif (19.6.2.1) s'écrit donc ici
\begin{equation}
\xymatrix{
  & 0 \ar[d] \\
  (\operatorname{gr}_{\mathfrak R}^*(R,N)\otimes_{A/\mathfrak R}(A/\mathfrak L))[T_1,\ldots,T_n] \ar[r]^-{\psi_R'} \ar[d] & \operatorname{gr}_{\mathfrak L}^*(R,N) \ar[d] \\
  (\operatorname{gr}_{\mathfrak R}^*(N)\otimes_{A/\mathfrak R}(A/\mathfrak L))[T_1,\ldots,T_n] \ar[r]^-{\psi_N} \ar[d] & \operatorname{gr}_{\mathfrak L}^*(N) \ar[d] \\
  (\operatorname{gr}_{\mathfrak R}^*(M)\otimes_{A/\mathfrak R}(A/\mathfrak L))[T_1,\ldots,T_n] \ar[r]^-{\psi_M} \ar[d] & \operatorname{gr}_{\mathfrak L}^*(M) \ar[d] \\
  0 & 0
}
\tag{19.7.1.1}\label{IV.19.7.1.1}
\end{equation}
où $\operatorname{gr}_{\mathfrak R}^*(R,N)$ (resp. $\operatorname{gr}_{\mathfrak L}^*(R,N)$) est le gradué associé à $R$ pour la filtration des $R\cap\mathfrak R^mN$ (resp. $R\cap\mathfrak L^mN$); rappelons que dans ce diagramme les colonnes sont exactes et que $\psi_N$

\oldpage[IV-4]{214}

est \emph{bijectif}. Il s'agit, en vertu de (19.6.3), de prouver que $\psi_R'$ est \emph{surjectif}, ce qui va être fait en plusieurs étapes.

\begin{env}[19.7.1.2]
\label{IV.19.7.1.2}
Posons $B=\operatorname{gr}_{\mathfrak R}^*A\otimes_{A/\mathfrak R}(A/\mathfrak L)$, $P=\operatorname{gr}_{\mathfrak R}^*(N)\otimes_{A/\mathfrak R}(A/\mathfrak L)$, $Q=\operatorname{gr}_{\mathfrak L}^*(R,N)$; puisque $\psi_N$ est bijectif, $P[T_1,\ldots,T_n]$ s'identifie à $\operatorname{gr}_{\mathfrak L}^*(N)$, et $Q$ à un sous-$B[T_1,\ldots,T_n]$-module de $P[T_1,\ldots,T_n]$; nous allons d'abord voir que l'on a
\end{env}

\begin{env}[19.7.1.3]
\label{IV.19.7.1.3}
\[
  Q=(Q\cap P)[T_1,\ldots,T_n].
\]
\end{env}

Pour cela, posons $Z=Q/((Q\cap P)[T_1,\ldots,T_n])$; c'est un sous-$B[T_1,\ldots,T_n]$-module de
\[
  P[T_1,\ldots,T_n]/((Q\cap P)[T_1,\ldots,T_n])=(P/(Q\cap P))[T_1,\ldots,T_n].
\]
En tant que $(A/\mathfrak L)$-module, il est donc isomorphe à un sous-module d'une somme directe de $(A/\mathfrak L)$-modules isomorphes à $P/(Q\cap P)=(P+Q)/Q$. Mais $(P+Q)/Q$ est un $(A/\mathfrak L)$-module isomorphe à un sous-module de $P[T_1,\ldots,T_n]/Q$, qui n'est autre que $\operatorname{gr}_{\mathfrak L}^*(M)$ en vertu du diagramme (19.7.1.1). On a donc
\begin{equation}
  \operatorname{Ass}_{A/\mathfrak L}(Z)\subset\operatorname{Ass}_{A/\mathfrak L}(\operatorname{gr}_{\mathfrak L}^*(M)).
  \tag{19.7.1.4}\label{IV.19.7.1.4}
\end{equation}

Mais chacun des $\operatorname{gr}_{\mathfrak L}^m(M)$ est par hypothèse un $(A/\mathfrak L)$-module \emph{plat} de type fini, donc projectif puisque $A/\mathfrak L$ est noethérien, et on a par suite $\operatorname{Ass}_{A/\mathfrak L}(\operatorname{gr}_{\mathfrak L}^*(M))\subset\operatorname{Ass}_{A/\mathfrak L}(A/\mathfrak L)$. Pour voir que $Z=0$, il suffit par suite (Bourbaki, \emph{Alg. comm.}, chap. IV, § 1, n\textsuperscript{o} 1, cor. 1 de la prop. 2 et n\textsuperscript{o} 3, prop. 7) de voir que pour tout $\mathfrak q\in\operatorname{Ass}_{A/\mathfrak L}(A/\mathfrak L)$, on a $Z_{\mathfrak q}=0$. Mais les idéaux de $\operatorname{Ass}_{A/\mathfrak L}(A/\mathfrak L)$ sont les idéaux de la forme $\mathfrak p/(\mathfrak L/\mathfrak R)$, où $\mathfrak p\in\operatorname{Ass}_{A/\mathfrak R}(A/\mathfrak L)$, et il revient donc au même de voir que $Z_{\mathfrak p}=0$ pour tout $\mathfrak p\in\operatorname{Ass}_{A/\mathfrak R}(A/\mathfrak L)$. Or, si $\mathfrak p=\mathfrak r/\mathfrak R$, où $\mathfrak r$ est un idéal premier de $A$, on a $(\operatorname{gr}_{\mathfrak R}^*(M))_{\mathfrak p}=\operatorname{gr}_{\mathfrak R_{\mathfrak r}}^*(M_{\mathfrak r})$ et $(A/\mathfrak L)_{\mathfrak p}=A_{\mathfrak r}/\mathfrak L_{\mathfrak r}$, et les images des $f_i$ dans $A_{\mathfrak r}$ forment une suite $\operatorname{gr}_{\mathfrak R_{\mathfrak r}}^*(A_{\mathfrak r})$-régulière ($0$, 15.1.14); appliquant à $A_{\mathfrak r}$, $\mathfrak R_{\mathfrak r}$ et $M_{\mathfrak r}$ l'implication $a)\Rightarrow b)$ de l'énoncé, on voit que l'hypothèse $d)$ entraîne que $\psi_{M_{\mathfrak r}}$ est bijectif, donc aussi $\psi_{R_{\mathfrak r}}'$ en vertu de (19.6.3); en d'autres termes, $Q_{\mathfrak r}=Q_{\mathfrak r}'[T_1,\ldots,T_n]$, où on a posé $Q'=\operatorname{gr}_{\mathfrak R}^*(R,N)\otimes_{A/\mathfrak R}(A/\mathfrak L)$, qui est ici un sous-$B$-module de $P$; en particulier $Q_{\mathfrak r}'=Q_{\mathfrak r}\cap P_{\mathfrak r}$, donc finalement on a bien $Z_{\mathfrak r}=Z_{\mathfrak p}=0$, ce qui achève de démontrer (19.7.1.3).

\begin{env}[19.7.1.5]
\label{IV.19.7.1.5}
En vertu de (19.7.1.3), pour prouver que $\psi_R'$ est surjectif, il suffit de montrer que tout élément homogène de degré $m$ de $Q\cap P$ est l'image d'un élément de même degré de $\operatorname{gr}_{\mathfrak R}^*(R,N)\otimes_{A/\mathfrak R}(A/\mathfrak L)$. Or, les éléments de degré $m$ de $P$ s'identifient (par $\psi_N$) aux éléments de $\operatorname{gr}_{\mathfrak L}^m(N)$ qui appartiennent à $(\mathfrak R^mN+\mathfrak L^{m+1}N)/\mathfrak L^{m+1}N$; ceux qui sont images d'éléments de $\operatorname{gr}_{\mathfrak R}^m(R,N)\otimes_{A/\mathfrak R}(A/\mathfrak L)$ s'identifient aux éléments de $\operatorname{gr}_{\mathfrak L}^m(N)$ appartenant à $((R\cap\mathfrak R^mN)+\mathfrak L^{m+1}N)/\mathfrak L^{m+1}N$. Il suffira donc finalement de prouver, pour tout entier $m>0$, l'inclusion
\[
  ((R\cap\mathfrak L^mN)+\mathfrak L^{m+1}N)\cap(\mathfrak R^mN+\mathfrak L^{m+1}N)\subset(R\cap\mathfrak R^mN)+\mathfrak L^{m+1}N.
\]
Comme $\mathfrak R\subset\mathfrak L$, on vérifie aussitôt que le premier membre de cette relation est égal à $(R\cap(\mathfrak R^mN+\mathfrak L^{m+1}N))+\mathfrak L^{m+1}N$, si bien que tout se réduit à prouver
\end{env}

\begin{env}[19.7.1.6]
\label{IV.19.7.1.6}
\[
  R\cap(\mathfrak R^mN+\mathfrak L^{m+1}N)\subset(R\cap\mathfrak R^mN)+\mathfrak L^{m+1}N.
\]
\end{env}

\oldpage[IV-4]{215}

Nous procéderons par récurrence sur $m$, supposant donc (19.7.1.6) vérifiée lorsqu'on remplace $m$ par un entier $m'<m$. Nous considérerons d'autre part, pour $m$ fixé et pour un entier $d\geq0$, la relation
\begin{equation}
  R\cap(\mathfrak R^mN+\mathfrak L^{m+1}N)\subset(R\cap\mathfrak R^dN\cap\mathfrak L^mN)+\mathfrak L^{m+1}N.
  \tag{$*_d$}
\end{equation}

Il est clair qu'il suffit de démontrer $(*_d)$ pour $d=m$, et d'autre part $(*_0)$ est trivialement vraie. Nous prouverons $(*_d)$ par récurrence sur $d$; autrement dit, nous supposons, pour un $d\geq0$ fixé (avec $d<m$), que $(*_d)$ est vraie, et nous voulons prouver
\begin{equation}
  R\cap(\mathfrak R^mN+\mathfrak L^{m+1}N)\subset(R\cap\mathfrak R^{d+1}N\cap\mathfrak L^mN)+\mathfrak L^{m+1}N.
  \tag{$*_{d+1}$}
\end{equation}

Considérons pour cela, pour un $h\geq0$, la relation
\begin{equation}
  R\cap(\mathfrak R^mN+\mathfrak L^{m+1}N)\subset(R\cap(\mathfrak R^{d+1}N+\mathfrak R^d\mathfrak L^hN)\cap\mathfrak L^mN)+\mathfrak L^{m+1}N.
  \tag{$**_{d+1,h}$}
\end{equation}

L'hypothèse $(*_d)$ implique que $(**_{d+1,0})$ est vraie. Nous prouverons par récurrence sur $h$ que $(**_{d+1,h})$ est vraie pour tout $h\geq0$. Mais on a le

\begin{lemma}[19.7.1.7]
\label{IV.19.7.1.7}
Soient $A$ un anneau noethérien, $N$ un $A$-module de type fini, $E$, $F$ deux sous-modules de $N$, $\mathfrak L$ un idéal de $A$. Pour tout $k>0$ il existe $h>0$ tel que
\[
  E\cap(F+\mathfrak L^hN)\subset(E\cap F)+\mathfrak L^kN.
\]
\end{lemma}

En effet, soit $\varphi:N\to N/(E\cap F)=N_1$ l'homomorphisme canonique, et posons $E_1=\varphi(E)$, $F_1=\varphi(F)$, de sorte que $E_1\cap F_1=0$ et $\varphi(E\cap(F+\mathfrak L^hN))=E_1\cap(F_1+\mathfrak L^hN_1)$. On sait ($0_{\mathrm I}$, 7.3.2) qu'il existe $h$ tel que $(E_1+F_1)\cap\mathfrak L^hN_1\subset\mathfrak L^k(E_1+F_1)$; si $h$ est ainsi choisi et si $x_1\in E_1$ est tel qu'il existe $y_1\in F_1$ pour lequel $x_1-y_1\in\mathfrak L^hN_1$, on en déduit d'après ce qui précède qu'il existe $u_1\in\mathfrak L^kE_1$ et $v_1\in\mathfrak L^kF_1$ tels que $x_1=u_1$ et $y_1=v_1$; donc $E_1\cap(F_1+\mathfrak L^hN_1)\subset\mathfrak L^kN_1$, ce qui prouve le lemme.

Il suffit alors d'appliquer ce lemme avec $E=R\cap\mathfrak L^mN$, $F=\mathfrak R^{d+1}N$ et de prendre $k=m+1$ pour en déduire qu'avec la valeur de $h$ correspondante, $(**_{d+1,h})$ entraîne $(*_{d+1})$.

Nous supposons donc dans tout ce qui suit que $(**_{d+1,h})$ est vérifiée.

\begin{env}[19.7.1.8]
\label{IV.19.7.1.8}
Démonstration de $(**_{d+1,h+1})$: Premier cas: $h\leq m-d$.
\end{env}

Partons donc d'un élément
\begin{env}[19.7.1.9]
\label{IV.19.7.1.9}
\[
  g\in R\cap(\mathfrak R^{d+1}N+\mathfrak R^d\mathfrak L^hN)\cap\mathfrak L^mN
\]
\end{env}
congru mod. $\mathfrak L^{m+1}N$ à un élément $g_1\in R\cap(\mathfrak R^mN+\mathfrak L^{m+1}N)$; on a donc
\begin{env}[19.7.1.10]
\label{IV.19.7.1.10}
\[
  g\in\mathfrak R^mN+\mathfrak L^{m+1}N.
\]
\end{env}

Il suffira de prouver que $g$ est aussi congru mod. $\mathfrak L^{m+1}N$ à un élément $g_2\in R\cap(\mathfrak R^{d+1}N+\mathfrak R^d\mathfrak L^{h+1}N)$, car cela entraînera $g_2\in\mathfrak L^mN$, puisque l'on aura $g\in\mathfrak L^mN$ et $g_2-g\in\mathfrak L^{m+1}N$. Dans le cas $h\leq m-d$ que nous considérons, il suffit de démontrer la relation
\begin{env}[19.7.1.11]
\label{IV.19.7.1.11}
\[
  (\mathfrak R^mN+\mathfrak L^{m+1}N)\cap\mathfrak R^dN\subset\mathfrak L^{m-d+1}\mathfrak R^dN+\mathfrak R^{d+1}N.
\]
\end{env}

\oldpage[IV-4]{216}

car les relations (19.7.1.9) et (19.7.1.10) entraînent que $g$ appartient au premier membre de (19.7.1.11), et comme $h\leq m-d$, on a $\mathfrak L^{m-d+1}\subset\mathfrak L^{h+1}$, ce qui établira $(**_{d+1,h+1})$. D'ailleurs on a $\mathfrak R^mN\subset\mathfrak R^dN$, et la relation
\begin{env}[19.7.1.12]
\label{IV.19.7.1.12}
\[
  \mathfrak L^{m+1}N\cap\mathfrak R^dN\subset\mathfrak L^{m-d+1}\mathfrak R^dN
\]
\end{env}
entraînera (19.7.1.11). Comme $N$ est un $A$-module libre de type fini, on voit qu'il suffit donc de prouver
\begin{env}[19.7.1.13]
\label{IV.19.7.1.13}
\[
  \mathfrak L^{m+1}\cap\mathfrak R^d\subset\mathfrak L^{m-d+1}\mathfrak R^d.
\]
\end{env}

Il suffira de prouver que, de façon générale, pour $d<m$, on a
\begin{env}[19.7.1.14]
\label{IV.19.7.1.14}
\[
  \mathfrak L^{m+1}\cap\mathfrak R^d\subset\mathfrak L^{m-d+1}\mathfrak R^d+\mathfrak R^{d+1}.
\]
\end{env}
En effet, comme $\mathfrak R\subset\mathfrak L$, donc $\mathfrak L^{m-d+1}\mathfrak R^d\subset\mathfrak L^{m+1}$, la relation précédente entraînera
\[
  \mathfrak L^{m+1}\cap\mathfrak R^d\subset\mathfrak L^{m-d+1}\mathfrak R^d+\mathfrak L^{m+1}\cap\mathfrak R^{d+1}
\]
et par récurrence sur $d\leq m$, on en conclura
\[
  \mathfrak L^{m+1}\cap\mathfrak R^d\subset\mathfrak L^{m-d+1}\mathfrak R^d+\mathfrak R^{m+1}
\]
d'où (19.7.1.13), puisque $\mathfrak R^{m+1}\subset\mathfrak L^{m-d+1}\mathfrak R^d$.

Comme $\mathfrak L=\mathfrak J+\mathfrak R$, la relation (19.7.1.14) s'écrit aussi
\begin{env}[19.7.1.15]
\label{IV.19.7.1.15}
\[
\begin{aligned}
  &(\mathfrak J^{m+1}+\mathfrak J^m\mathfrak R+\cdots+\mathfrak J\mathfrak R^m+\mathfrak R^{m+1})\cap\mathfrak R^d\\
  &\quad\subset\mathfrak J^{m-d+1}\mathfrak R^d+\mathfrak J^{m-d}\mathfrak R^{d+1}+\cdots+\mathfrak J\mathfrak R^m+\mathfrak R^{m+1}+\mathfrak R^{d+1}\\
  &\quad=\mathfrak J^{m-d+1}\mathfrak R^d+\mathfrak R^{d+1}.
\end{aligned}
\]
\end{env}

Nous allons voir que cette inclusion est elle-même conséquence de
\begin{env}[19.7.1.16]
\label{IV.19.7.1.16}
\[
  \mathfrak R^a\mathfrak J^b\cap\mathfrak R^{a+1}\subset\mathfrak R^{a+1}\mathfrak J^{b-1}
\]
\end{env}
valable pour $a\geq0$, $b\geq1$. En effet, il suffit, pour prouver (19.7.1.15) de montrer que pour $0<q\leq d$, on a
\begin{env}[19.7.1.17]
\label{IV.19.7.1.17}
\[
\begin{aligned}
  &(\mathfrak J^{m+1}+\mathfrak J^m\mathfrak R+\cdots+\mathfrak J\mathfrak R^m+\mathfrak R^{m+1})\cap\mathfrak R^d\\
  &\quad\subset(\mathfrak J^{m+1-q}\mathfrak R^q+\mathfrak J^{m-q}\mathfrak R^{q+1}+\cdots+\mathfrak R^{m+1})\cap\mathfrak R^d
\end{aligned}
\]
\end{env}
car pour $q=d$, cela entraînera (19.7.1.15). Or, pour prouver (19.7.1.17), il suffit de procéder par récurrence sur $q$ en supposant la relation vraie pour $q<d$; un élément du premier membre s'écrit donc par hypothèse $y+z$ avec $y\in\mathfrak J^{m+1-q}\mathfrak R^q$, $z\in\mathfrak J^{m-q}\mathfrak R^{q+1}+\cdots+\mathfrak R^{m+1}$; comme $y+z\in\mathfrak R^d\subset\mathfrak R^{q+1}$, on en déduit $y\in\mathfrak R^{q+1}$, puisque $z\in\mathfrak R^{q+1}$; tenant compte de (19.7.1.16), on a $y\in\mathfrak J^{m-q}\mathfrak R^{q+1}$, donc
\[
  y+z\in\mathfrak J^{m-q}\mathfrak R^{q+1}+\cdots+\mathfrak R^{m+1},\qquad\text{et}\qquad y+z\in\mathfrak R^d
\]
par hypothèse, ce qui prouve (19.7.1.17) où $q$ a été remplacé par $q+1$.

Reste donc à établir (19.7.1.16). Un élément du premier membre s'écrit $a_1f_1+\cdots+a_nf_n$, où les $a_i$ appartiennent à $\mathfrak R^a\mathfrak J^{b-1}$; tenant compte de ce que la suite $\mathbf f$ est $(\mathfrak R^a/\mathfrak R^{a+1})$-régulière par hypothèse, donc aussi $(\mathfrak R^a/\mathfrak R^{a+1})$-quasi-régulière ($0$, 15.1.9), il résulte de la définition ($0$, 15.1.7) que la relation $a_1f_1+\cdots+a_nf_n\in\mathfrak R^{a+1}$, où les $a_i\in\mathfrak R^a$, entraîne nécessairement $a_i\in\mathfrak R^{a+1}$, donc $a_i\in\mathfrak R^a\mathfrak J^{b-1}\cap\mathfrak R^{a+1}$; il suffit de raisonner par

\oldpage[IV-4]{217}

récurrence sur $b$ pour $b\geq1$ (le cas $b=1$ étant trivial) pour achever de prouver (19.7.1.16) et par suite aussi $(**_{d+1,h+1})$ pour $h\leq m-d$.

\begin{env}[19.7.1.18]
\label{IV.19.7.1.18}
Démonstration de $(**_{d+1,h+1})$; Deuxième cas: $h>m-d$.
\end{env}

Nous allons tout d'abord démontrer que, pour tout $h\geq0$, on a
\begin{env}[19.7.1.19]
\label{IV.19.7.1.19}
\[
  R\cap(\mathfrak R^{d+1}N+\mathfrak R^d\mathfrak L^hN)\subset\mathfrak J^h(R\cap\mathfrak R^dN)+\mathfrak L^{h+1}\mathfrak R^dN+\mathfrak R^{d+1}N.
\]
\end{env}

Considérons pour cela un élément $z$ du premier membre de (19.7.1.19). Notons que l'on a $\mathfrak R^{d+1}N+\mathfrak R^d\mathfrak L^hN=\mathfrak R^{d+1}N+\mathfrak R^d\mathfrak J^hN$ puisque $\mathfrak L=\mathfrak R+\mathfrak J$; nous allons considérer la classe $\bar z$ de $z$ dans $\operatorname{gr}_{\mathfrak J}^h(\operatorname{gr}_{\mathfrak R}^d(N))$, qui, en vertu de (19.5.5) et de l'hypothèse sur $\mathbf f$, s'identifie à un sous-$(A/\mathfrak L)$-module de $\operatorname{gr}_{\mathfrak L}^{h+d}(N)$. Du fait que $z\in R$, nous allons montrer que l'on a même (avec l'identification précédente)
\begin{env}[19.7.1.20]
\label{IV.19.7.1.20}
\[
  \bar z\in\operatorname{gr}_{\mathfrak L}^{h+d}(R,N).
\]
\end{env}

En effet, on a noté, dans la démonstration de (19.7.1.3) que, pour tout idéal premier $\mathfrak r$ de $A$ tel que $\mathfrak p=\mathfrak r/\mathfrak R\in\operatorname{Ass}_{A/\mathfrak R}(A/\mathfrak L)$, l'application $\psi_{R_{\mathfrak r}}'$ est bijective, et par suite la relation (19.7.1.6), où l'on remplace tous les modules par leurs localisés en $\mathfrak r$, est vraie. Comme $z$ appartient à $R\cap(\mathfrak R^dN+\mathfrak L^{d+1}N)$, son image $z/1$ dans $R_{\mathfrak r}$ appartient à $(R_{\mathfrak r}\cap\mathfrak R_{\mathfrak r}^dN_{\mathfrak r})+\mathfrak L_{\mathfrak r}^{d+1}N_{\mathfrak r}$, donc l'image $\bar z/1$ de $\bar z$ dans $\operatorname{gr}_{\mathfrak L_{\mathfrak r}}^{h+d}(N_{\mathfrak r})$ appartient à
\[
  \operatorname{gr}_{\mathfrak L_{\mathfrak r}}^{h+d}(R_{\mathfrak r},N_{\mathfrak r})
  =(\operatorname{gr}_{\mathfrak L}^{h+d}(R,N))_{\mathfrak r}
  =(\operatorname{gr}_{\mathfrak L}^{h+d}(R,N))_{\mathfrak q},
\]
où $\mathfrak q=\mathfrak r/\mathfrak L$. En d'autres termes l'image de $\bar z$ par l'application canonique
\[
  (\operatorname{gr}_{\mathfrak L}^{h+d}(N))_{\mathfrak q}\longrightarrow(\operatorname{gr}_{\mathfrak L}^{h+d}(M))_{\mathfrak q}
\]
est nulle pour tout $\mathfrak q\in\operatorname{Ass}_{A/\mathfrak L}(A/\mathfrak L)$. Comme on a vu dans la démonstration de (19.7.1.3) que
\[
  \operatorname{Ass}_{A/\mathfrak L}(\operatorname{gr}_{\mathfrak L}^{h+d}(M))\subset\operatorname{Ass}_{A/\mathfrak L}(A/\mathfrak L),
\]
on en déduit bien (Bourbaki, \emph{Alg. comm.}, chap. IV, § 1, n\textsuperscript{o} 1, cor. 1 de la prop. 2 et n\textsuperscript{o} 3, prop. 7) que l'image de $\bar z$ par l'application canonique $\operatorname{gr}_{\mathfrak L}^{h+d}(N)\to\operatorname{gr}_{\mathfrak L}^{h+d}(M)$ est nulle, c'est-à-dire que l'on a la relation (19.7.1.20).

Cela étant, on peut écrire par définition
\begin{env}[19.7.1.21]
\label{IV.19.7.1.21}
\[
  z\equiv\sum_{|\mathbf p|=h}c_{\mathbf p}\mathbf f^{\mathbf p}\pmod{\mathfrak R^{d+1}N}
\]
\end{env}
où on pose comme d'ordinaire $\mathbf p=(p_1,\ldots,p_n)$, $\mathbf f^{\mathbf p}=f_1^{p_1}f_2^{p_2}\cdots f_n^{p_n}$, $|\mathbf p|=p_1+\cdots+p_n$ et où $c_{\mathbf p}\in\mathfrak R^dN$. En outre, comme $\psi_N$ identifie $\operatorname{gr}_{\mathfrak L}^*(N)$ à $(\operatorname{gr}_{\mathfrak R}^*(N)\otimes_{A/\mathfrak R}(A/\mathfrak L))[T_1,\ldots,T_n]$, les $c_{\mathbf p}$ sont déterminés mod. $\mathfrak R^d\mathfrak LN$, et si $\bar c_{\mathbf p}$ est la classe de $c_{\mathbf p}$ mod. $\mathfrak R^d\mathfrak LN$, on a, par l'identification précédente,
\begin{env}[19.7.1.22]
\label{IV.19.7.1.22}
\[
  \bar z=\sum_{|\mathbf p|=h}\bar c_{\mathbf p}T^{\mathbf p}.
\]
\end{env}
Utilisant (19.7.1.3), on en déduit que l'on a nécessairement, pour tout $\mathbf p$ tel que $|\mathbf p|=h$ ($\bar c_{\mathbf p}$ étant cette fois identifié par $\psi_N$ à un élément de $\operatorname{gr}_{\mathfrak R}^d(N)$)
\begin{env}[19.7.1.23]
\label{IV.19.7.1.23}
\[
  \bar c_{\mathbf p}\in\operatorname{gr}_{\mathfrak R}^d(R,N).
\]
\end{env}
Mais puisque $d<m$, l'hypothèse que $(*_d)$ est vraie implique que $\bar c_{\mathbf p}$ appartient à l'image de $\psi_R'$, donc qu'on peut supposer que $c_{\mathbf p}\in(R\cap\mathfrak R^dN)+\mathfrak R^d\mathfrak LN$. La relation (19.7.1.21) donne alors
\[
  z\in\mathfrak J^h(R\cap\mathfrak R^dN)+\mathfrak J^h\mathfrak R^d\mathfrak LN+\mathfrak R^{d+1}N
\]
et comme $\mathfrak J\subset\mathfrak L$, cela prouve (19.7.1.19).

\oldpage[IV-4]{218}

En particulier, on peut donc écrire
\[
  g=t+g_2
\]
avec $t\in\mathfrak J^h(R\cap\mathfrak R^dN)$ et $g_2\in\mathfrak R^{d+1}N+\mathfrak R^d\mathfrak L^hN$; en outre, comme $t\in R$ et $g\in R$, on a $g_2\in R$. D'autre part, utilisant l'hypothèse $h\geq m-d+1$, on voit que $t\in\mathfrak L^{m+1}N$. L'élément $g_2$ répond donc à toutes les conditions énoncées au début de (19.7.1.8), ce qui achève la démonstration de (19.7.1).

\begin{remarks}[19.7.2]
\label{IV.19.7.2}
\begin{enumerate}
  \item[(i)] Lorsque l'on suppose que $\operatorname{gr}_{\mathfrak R}^*(A)$ est un $(A/\mathfrak R)$-module \emph{plat}, l'hypothèse que la suite $\mathbf f$ est $(A/\mathfrak R)$-régulière entraîne qu'elle est aussi $\operatorname{gr}_{\mathfrak R}^*(A)$-régulière ($0$, 15.1.14). On notera que ce sera le cas lorsque $A$ et $A/\mathfrak R$ sont des anneaux réguliers ($0$, 17.3.6): en effet, $\mathfrak R$ est alors un idéal régulier de $A$ (19.1.2), donc quasi-régulier, et en localisant aux idéaux maximaux de $A$ contenant $\mathfrak R$ et utilisant ($0$, 15.1.7), on voit que $\operatorname{gr}_{\mathfrak R}^*(A)$ est un $(A/\mathfrak R)$-module \emph{projectif}.
  \item[(ii)] Il serait intéressant de pouvoir prouver la dernière conclusion de (19.7.1) en supposant seulement qu'il existe une $A$-algèbre $B$ qui est un anneau noethérien, pour lequel $\mathfrak JB$ est contenu dans le radical de $B$, et que $M$ est un $B$-module de type fini.
\end{enumerate}
\end{remarks}

\begin{corollary}[19.7.3]
\label{IV.19.7.3}
Soient $A$ un anneau local noethérien, $\mathfrak m$ son idéal maximal, $\mathfrak R$ un idéal de $A$ tel que l'anneau $A/\mathfrak R$ soit régulier et de dimension $n$, $M$ un $A$-module de type fini. Les conditions suivantes sont équivalentes:
\begin{enumerate}
  \item[a)] $\operatorname{gr}_{\mathfrak R}^*(M)$ est un $(A/\mathfrak R)$-module plat.
  \item[b)] Si $P(T)$ et $Q(T)$ sont les séries de Poincaré des $(A/\mathfrak m)$-modules gradués $\operatorname{gr}_{\mathfrak m}^*(M)$ et $\operatorname{gr}_{\mathfrak R}^*(M)\otimes_{A/\mathfrak R}(A/\mathfrak m)$, on a
\end{enumerate}
\begin{equation}
  Q(T)=P(T)(1-T)^n.
  \tag{19.7.3.1}\label{IV.19.7.3.1}
\end{equation}

Soit $\mathbf f=(f_1,\ldots,f_n)$ une suite d'éléments de $A$ tels que les images des $f_i$ dans $A/\mathfrak R$ forment un système régulier de paramètres de $A/\mathfrak R$ ($0$, 17.1.6), de sorte que $\mathbf f$ est une suite $(A/\mathfrak R)$-régulière; en outre, si $\mathfrak J$ est l'idéal engendré par $\mathbf f$, on a $\mathfrak J+\mathfrak R=\mathfrak m$, et comme $A/\mathfrak m$ est un corps, $\operatorname{gr}_{\mathfrak m}^*(M)$ est un $(A/\mathfrak m)$-module plat. Comme $A$ est noethérien et $M$ un $A$-module de type fini, on peut appliquer l'équivalence de $a)$ et $b)$ dans (19.7.1). En outre, comme $A/\mathfrak m$ est un corps, et que $\psi_M$ est surjectif, le fait que $\psi_M$ soit bijectif équivaut à dire que pour tout $h\geq0$, $\operatorname{gr}_{\mathfrak m}^h(M)$ et le sous-module des éléments de degré $h$ de $(\operatorname{gr}_{\mathfrak R}^*(M)\otimes_{A/\mathfrak R}(A/\mathfrak m))[T_1,\ldots,T_n]$ ont même rang sur $A/\mathfrak m$; mais pour le second de ces modules le rang en question est évidemment égal à
\[
  \sum_{r=0}^h\binom{r+n-1}{n-1}\operatorname{rg}(\operatorname{gr}_{\mathfrak R}^{h-r}(M)\otimes_{A/\mathfrak R}(A/\mathfrak m)).
\]
Comme
\[
  (1-T)^{-n}=\sum_{r=0}^{\infty}\binom{r+n-1}{n-1}T^r,
\]
cela prouve que la relation (19.7.3.1) est nécessaire et suffisante pour que $\psi_M$ soit bijectif.
\end{corollary}

\begin{corollary}[19.7.4]
\label{IV.19.7.4}
Soient $X$ un préschéma localement noethérien, $Y$ un sous-préschéma fermé de $X$, régulier et connexe, $\sh F$ un $\sh O_X$-Module cohérent normalement plat (6.10.1) le long de $Y$. Alors il existe une série formelle $R\in\mathbb{Q}[[T]]$ (indépendante de $x\in Y$) telle que, pour tout $x\in Y$, la série de Poincaré de $\operatorname{gr}_{\mathfrak m_x}^*(\sh F_x)$ soit donnée par la formule
\begin{equation}
  P_x(\sh F)(T)=R(T)(1-T)^{-n}\qquad\text{où }n=\dim(\sh O_{Y,x}).
  \tag{19.7.4.1}\label{IV.19.7.4.1}
\end{equation}

\oldpage[IV-4]{219}

Soit $\sh J$ l'Idéal cohérent de $\sh O_X$ définissant $Y$; l'hypothèse sur $\sh F$ signifie que $\operatorname{gr}_{\sh J}^*(\sh F)$ est un $\sh O_Y$-Module plat. On peut donc, pour tout $x\in Y$, appliquer (19.7.3) en prenant $A=\sh O_{X,x}$, $\mathfrak R=\sh J_x$, $M=\sh F_x$; l'hypothèse $a)$ étant remplie par hypothèse, la série de Poincaré $P_x(\sh F)(T)$ est donc égale à $R_x(T)(1-T)^{-n}$, où $R_x(T)$ est la série de Poincaré de $\operatorname{gr}_{\sh J_x}^*(\sh F_x)\otimes_{\sh O_{Y,x}}k(x)$. Or, chacun des $\sh O_Y$-Modules $\sh J^h\sh F/\sh J^{h+1}\sh F$ est plat et cohérent par hypothèse, donc localement libre (2.1.12), et puisque $Y$ est supposé connexe, chacun des $\sh J^h\sh F/\sh J^{h+1}\sh F$ est de rang constant dans $Y$, ce qui prouve que la série de Poincaré $R_x(T)$ est indépendante de $x\in Y$.
\end{corollary}

\begin{corollary}[19.7.5]
\label{IV.19.7.5}
Soient $X$ un préschéma localement noethérien, $Y$ un sous-préschéma fermé de $X$, $Z$ un sous-préschéma fermé de $Y$; on suppose $Y$ et $Z$ réguliers. Soit $\sh F$ un $\sh O_X$-Module cohérent.
\begin{enumerate}
  \item[(i)] Si $\sh F$ est normalement plat le long de $Y$ aux points de $Z$ (cf. (11.3.4) pour la définition de la platitude normale en un point), alors $\sh F$ est normalement plat le long de $Z$.
  \item[(ii)] Supposons en outre que $Z$ soit connexe et que $\sh O_X$ soit normalement plat le long de $Y$ aux points de $Z$ (ce qui aura lieu en particulier si $X$ est régulier aux points de $Z$). Si $\sh F$ est normalement plat le long de $Z$ et s'il existe un point $z\in Z$ tel que $\sh F$ soit normalement plat le long de $Y$ au point $z$, alors $\sh F$ est normalement plat le long de $Y$ aux points de $Z$.
\end{enumerate}

Soient $\sh K$ et $\sh L\supset\sh K$ les Idéaux cohérents de $\sh O_X$ définissant respectivement $Y$ et $Z$. Pour tout point $z\in Z$, l'immersion canonique $Z\to Y$ est régulière au point $z$, puisque $Y$ et $Z$ sont réguliers en ce point (19.1.1); autrement dit, il existe une suite $(f_i)_{1\leq i\leq n}$ d'éléments de $\sh O_{X,z}$ qui est $\sh K_z$-régulière et est telle que $\sh L_z=\sh K_z+\sum_i f_i\sh O_{X,z}$.
\begin{enumerate}
  \item[(i)] L'hypothèse signifie que $\operatorname{gr}_{\sh K_z}^*(\sh F_z)$ est un $\sh O_{Y,z}$-module plat; elle entraîne donc que $\operatorname{gr}_{\sh L_z}^*(\sh F_z)$ est un $\sh O_{Z,z}$-module plat, en vertu du fait que dans (19.7.1), $a)$ entraîne $b)$.
  \item[(ii)] L'hypothèse supplémentaire sur $X$ entraîne que la suite $(f_i)$ est $\operatorname{gr}_{\sh K_z}^*(\sh F_z)$-régulière en tout point $z\in Z$ (19.7.2). D'autre part, $Z$ étant régulier et connexe, est intègre, et si $\operatorname{gr}_{\sh K_z}^*(\sh F_z)$ est $\sh O_{Y,z}$-plat en un point $z\in Z$, alors, comme le point générique $\zeta$ de $Z$ est une générisation de $z$, on en conclut que $\operatorname{gr}_{\sh K_\zeta}^*(\sh F_\zeta)$ est $\sh O_{Y,\zeta}$-plat ($0_{\mathrm I}$, 6.3.1). On en déduit que si en tout point $z'$ de $Z$, $\operatorname{gr}_{\sh L_{z'}}^*(\sh F_{z'})$ est un $\sh O_{Z,z'}$-module plat, alors $\operatorname{gr}_{\sh K_{z'}}^*(\sh F_{z'})$ est un $\sh O_{Y,z'}$-module plat, car on peut appliquer l'équivalence de $a)$ et $d)$ dans (19.7.1).
\end{enumerate}
\end{corollary}

\subsection{Propriétés de passage à la limite projective.}
\EGAReaderTarget{ega4.c04.subsec.19-8}{subsection:IV.19.8}
\label{IV.19.8}

Dans ce numéro, les notations et conventions sur les limites projectives sont celles de (8.5.1) et de (8.8.1).

\begin{proposition}[19.8.1]
\label{IV.19.8.1}
Supposons que les morphismes de transition $S_\mu\to S_\lambda$ $(\lambda\leq\mu)$ soient plats, et en outre que l'une des deux hypothèses suivantes soit vérifiée:
\begin{enumerate}
  \item[$1^\circ$] Les préschémas $S_\lambda$ sont localement noethériens.
  \item[$2^\circ$] Les morphismes de transition $S_\mu\to S_\lambda$ sont surjectifs (donc fidèlement plats).
\end{enumerate}

\oldpage[IV-4]{220}

\emph{Dans ces conditions:}
\begin{enumerate}
  \item[(i)] Supposons $S_\alpha$ quasi-compact. Soit $\sh F_\alpha$ un $\sh O_{S_\alpha}$-Module quasi-cohérent, que l'on suppose en outre de type fini quand l'hypothèse $1^\circ$ est vérifiée. Soit $(f_{i\alpha})_{1\leq i\leq n}$ une suite finie de sections de $\sh F_\alpha$ au-dessus de $S_\alpha$. Pour que la suite des sections $f_i$ du $\sh O_S$-Module $\sh F$ au-dessus de $S$, correspondant aux $f_{i\alpha}$ $(1\leq i\leq n)$, soit $\sh F$-régulière, il faut et il suffit qu'il existe $\lambda\geq\alpha$ tel que la suite des sections $f_{i\lambda}$ de $\sh F_\lambda$ au-dessus de $S_\lambda$ soit $\sh F_\lambda$-régulière.
  \item[(ii)] Supposons $X_\alpha$ quasi-compact, et soit $j_\alpha:X_\alpha\to S_\alpha$ une immersion, que l'on suppose localement de présentation finie lorsque l'hypothèse $2^\circ$ est vérifiée. Alors, pour que l'immersion correspondante $j:X\to S$ soit régulière, il faut et il suffit qu'il existe $\lambda\geq\alpha$ tel que $j_\lambda:X_\lambda\to S_\lambda$ soit régulière.
\end{enumerate}
\end{proposition}

\begin{enumerate}
  \item[(i)] Si $p_\lambda:S\to S_\lambda$ est la projection canonique, on a
  \[
    \sh F/\Bigl(\sum_{j=1}^{i-1}f_j\sh F\Bigr)
    =p_\lambda^*\Bigl(\sh F_\lambda/\Bigl(\sum_{j=1}^{i-1}f_{j\lambda}\sh F_\lambda\Bigr)\Bigr),
  \]
  donc on est ramené au cas $n=1$, auquel cas on supprime l'indice $i$. Le fait que la condition soit suffisante résulte de ce que $p_\lambda$ est plat (8.3.8) et de ($0$, 15.1.5). Dans le cas $2^\circ$, $p_\lambda$ est fidèlement plat (8.3.8) et la nécessité de la condition résulte encore de ($0$, 15.2.5), avec $\lambda=\alpha$. Dans le cas $1^\circ$, notons $\sh N_\lambda$ (resp. $\sh N$) le noyau de l'homomorphisme $\sh F_\lambda\to\sh F_\lambda$ (resp. $\sh F\to\sh F$), multiplication par $f_\lambda$ (resp. par $f$); puisque $\sh F_\lambda$ est cohérent par hypothèse, il en est de même de $\sh N_\lambda$, donc l'hypothèse $\sh N=0$ entraîne $\sh N_\lambda=0$ pour un $\lambda\geq\alpha$ en vertu de (8.5.8, (ii)).
  \item[(ii)] Comme $p_\lambda$ est plat, la suffisance de la condition résulte de (19.1.5, (ii)). Pour prouver sa nécessité, notons que, puisque $j_\alpha(X_\alpha)$ est quasi-compact, il est contenu dans un ouvert quasi-compact de $S_\alpha$ et on peut donc se limiter au cas où $S_\alpha$ est aussi quasi-compact et $j_\alpha$ une immersion fermée, de sorte que l'image de $X_\alpha$ est définie par un Idéal quasi-cohérent $\sh J_\alpha$ de $\sh O_{S_\alpha}$, qu'on peut en outre supposer de type fini dans les cas $1^\circ$ et $2^\circ$ ($j_\alpha$ étant localement de présentation finie dans le cas $2^\circ$); l'image de $X_\lambda$ (resp. $X$) dans $S_\lambda$ (resp. $S$) est alors définie par $\sh J_\lambda=\sh J_\alpha\sh O_{S_\lambda}$ (resp. $\sh J=\sh J_\alpha\sh O_S$) qui est encore de type fini. On peut en outre, compte tenu de (8.2.11), supposer que $\sh J$ est engendré par une suite régulière $(f_i)_{1\leq i\leq n}$ de sections au-dessus de $S$, qui définissent par suite un homomorphisme surjectif $u:\sh O_S^n\to\sh J$. Compte tenu de (8.5.2, (i)) et (8.5.7), il existe un $\lambda\geq\alpha$ et un homomorphisme surjectif $u_\lambda:\sh O_{S_\lambda}^n\to\sh J_\lambda$ tels que $u=p_\lambda^*(u_\lambda)$, donc les $f_i$ sont les images canoniques de sections $f_{i\lambda}$ de $\sh J_\lambda$ au-dessus de $S_\lambda$ engendrant cet Idéal. En vertu de (i), il existe donc $\mu\geq\lambda$ tel que la suite $(f_{i\mu})$ soit $\sh O_{S_\mu}$-régulière, donc l'immersion $j_\mu$ est régulière.
\end{enumerate}

\begin{proposition}[19.8.2]
\label{IV.19.8.2}
Soient $X_\alpha$ un $S_\alpha$-préschéma quasi-compact et localement de présentation finie.
\begin{enumerate}
  \item[(i)] Soient $\sh F_\alpha$ un $\sh O_{X_\alpha}$-Module quasi-cohérent de présentation finie, $(f_{i\alpha})_{1\leq i\leq n}$ une suite de sections de $\sh O_{X_\alpha}$ au-dessus de $X_\alpha$. Pour que la suite de sections $f_i$ de $\sh O_X$ au-dessus de $X$, correspondant aux $f_{i\alpha}$, soit $\sh F$-transversalement régulière relativement à $S$ (19.2.1), il faut et il suffit qu'il existe $\lambda\geq\alpha$ tel que la suite des sections $f_{i\lambda}$ de $\sh O_{X_\lambda}$ au-dessus de $X_\lambda$ soit $\sh F_\lambda$-transversalement régulière relativement à $S_\lambda$.
  \item[(ii)] Soient $Y_\alpha$ un $S_\alpha$-préschéma quasi-compact, $j_\alpha:Y_\alpha\to X_\alpha$ une $S_\alpha$-immersion localement

\oldpage[IV-4]{221}

  de présentation finie. Pour que la $S$-immersion correspondante $j:Y\to X$ soit transversalement régulière relativement à $S$, il faut et il suffit qu'il existe $\lambda\geq\alpha$ tel que $j_\lambda:Y_\lambda\to X_\lambda$ soit transversalement régulière relativement à $S_\lambda$.
\end{enumerate}
\end{proposition}

\begin{enumerate}
  \item[(i)] La suffisance de la condition résulte de ($0$, 15.1.15). Pour prouver qu'elle est nécessaire, il suffira de montrer que tout point $x\in X$ possède un voisinage ouvert $V(x)$ pour lequel il existe un indice $\lambda=\lambda(x)\geq\alpha$ et un voisinage ouvert $V_\lambda$ de l'image $x_\lambda$ de $x$ dans $X_\lambda$ tels que $V(x)$ soit l'image réciproque de $V_\lambda$ et que la suite des restrictions à $V_\lambda$ des $f_{i\lambda}$ soit $(\sh F_\lambda|V_\lambda)$-transversalement régulière relativement à $S_\lambda$: en effet, il suffira alors d'appliquer (8.3.4) à l'ensemble ouvert $U_\lambda$ de $X_\lambda$, défini par la condition d'être le plus grand ouvert tel que les restrictions à $U_\lambda$ des $f_{i\lambda}$ forment une suite $(\sh F_\lambda|U_\lambda)$-transversalement régulière relativement à $S_\lambda$.

  Soient $s$ l'image de $x$ dans $S$, et pour tout $\lambda\geq\alpha$, soit $s_\lambda$ l'image de $s$ dans $S_\lambda$. Alors la fibre $X_s$ est limite projective des fibres $(X_\lambda)_{s_\lambda}$ et est localement noethérienne puisque le morphisme $X\to S$ est localement de présentation finie; comme en outre $(X_\alpha)_{s_\alpha}$ est quasi-compact et que les morphismes de transition $(X_\mu)_{s_\mu}\to(X_\lambda)_{s_\lambda}$ sont plats (puisqu'il en est ainsi de $\Spec(k(s_\mu))\to\Spec(k(s_\lambda))$) on peut appliquer aux sections $f_i\otimes1$ de $\sh F_s=\sh F\otimes_{\sh O_S}k(s)$ au-dessus de $X_s$ le résultat de (19.8.1, (i)); d'où l'existence d'un $\lambda\geq\alpha$ tel que les sections $f_{i\lambda}\otimes1$ de $(\sh F_\lambda)_{s_\lambda}=\sh F_\lambda\otimes_{\sh O_{S_\lambda}}k(s_\lambda)$ au-dessus de $(X_\lambda)_{s_\lambda}$ forment une suite $(\sh F_\lambda)_{s_\lambda}$-régulière. En outre (11.2.6), on peut supposer que $\sh F_\lambda$ est $S_\lambda$-plat au point $x_\lambda$. Il résulte alors de (11.3.8) qu'il existe un voisinage ouvert $V_\lambda$ de $x_\lambda$ dans $X_\lambda$ vérifiant les conditions requises.

  \item[(ii)] La suffisance de la condition résulte de (19.2.7, (ii)). Pour prouver sa nécessité, il suffit encore de montrer que tout point $x\in j(Y)$ admet un voisinage ouvert $V(x)$ pour lequel il existe un indice $\lambda=\lambda(x)\geq\alpha$ et un voisinage ouvert $V_\lambda$ de l'image $x_\lambda$ de $x$ dans $X_\lambda$ tels que $V(x)$ soit l'image réciproque de $V_\lambda$ et que la restriction $j_\lambda^{-1}(V_\lambda)\to V_\lambda$ de $j_\lambda$ soit transversalement régulière relativement à $S_\lambda$. Soient $s$ l'image de $x$ dans $S$, et pour tout $\lambda\geq\alpha$, soit $s_\lambda$ l'image de $s$ dans $S_\lambda$. Comme $X_s$ (resp. $Y_s$) est limite projective des fibres $(X_\lambda)_{s_\lambda}$ (resp. $(Y_\lambda)_{s_\lambda}$), on peut comme dans (i) utiliser (19.8.1, (ii)), et comme par hypothèse l'immersion $Y_s\to X_s$ est régulière, il existe un indice $\lambda(s)$ tel que pour $\lambda\geq\lambda(s)$ l'immersion $(Y_\lambda)_{s_\lambda}\to(X_\lambda)_{s_\lambda}$ soit régulière. En outre, en vertu de (19.2.4) il existe un voisinage ouvert quasi-compact $W(x)$ de $x$ dans $X$ tel que les morphismes structuraux $W(x)\to S$ et $Y\cap W(x)\to S$ soient plats; appliquant (11.2.6) et (8.2.11), on peut supposer qu'il existe un $\lambda\geq\lambda(s)$ et un voisinage ouvert quasi-compact $W(x_\lambda)$ de $x_\lambda$ dans $X_\lambda$, tels que $W(x)$ soit l'image réciproque de $W(x_\lambda)$ et que les restrictions à $W(x_\lambda)$ et à $Y_\lambda\cap W(x_\lambda)$ respectivement des morphismes structuraux $X_\lambda\to S_\lambda$ et $Y_\lambda\to S_\lambda$ soient plats. On conclut donc de (19.2.4) qu'il existe un voisinage ouvert $V_\lambda$ de $x_\lambda$ répondant à la question.
\end{enumerate}

\begin{remark}[19.8.3]
\label{IV.19.8.3}
Les démonstrations précédentes montrent que les énoncés de (19.8.1) et (19.8.2) subsistent lorsqu'on remplace les conditions concernant la régularité ou la régularité transversale dans tout $X$ (resp. $X_\lambda$) par ces conditions dans un voisinage d'un point $x\in X$ (resp. dans un voisinage de $x_\lambda$, projection de $x$).
\end{remark}

\subsection{\texorpdfstring{Suites $\sh F$-régulières et profondeur.}{Suites F-régulières et profondeur.}}
\EGAReaderTarget{ega4.c04.subsec.19-9}{subsection:IV.19.9}
\label{IV.19.9}

\oldpage[IV-4]{222}

\begin{env}[19.9.1]
\EGAReaderTarget{ega4.c04.statement.19-9-1}{IV.19.9.1}
Rappelons que si $X$ est un préschéma localement noethérien, $\sh F$ un $\sh O_X$-Module cohérent, $T$ une partie de $X$, on pose (5.10.1)
\begin{equation}
  \operatorname{prof}_T(\sh F)=\inf_{x\in T}\operatorname{prof}(\sh F_x).
  \tag{19.9.1.1}\label{IV.19.9.1.1}
\end{equation}

On pose d'autre part, pour tout point $t\in T$
\begin{equation}
  \operatorname{prof}_{T,t}(\sh F)=\inf_{x\in T\cap\Spec(\sh O_{X,t})}\operatorname{prof}(\sh F_x),
  \tag{19.9.1.2}\label{IV.19.9.1.2}
\end{equation}
et on dit que $\operatorname{prof}_{T,t}(\sh F)$ est la \emph{$T$-profondeur de $\sh F$ au point $t$}; il est clair que l'on a
\begin{equation}
  \operatorname{prof}_T(\sh F)=\inf_{t\in T}\operatorname{prof}_{T,t}(\sh F).
  \tag{19.9.1.3}\label{IV.19.9.1.3}
\end{equation}
\end{env}

\begin{lemma}[19.9.2]
\label{IV.19.9.2}
Soient $A$ un anneau noethérien, $M$ un $A$-module de type fini, $\mathfrak J$ un idéal de $A$. Afin que, pour tout idéal premier $\mathfrak p\supset\mathfrak J$, on ait $\operatorname{prof}_{A_{\mathfrak p}}(M_{\mathfrak p})\geq r$, il faut et il suffit qu'il existe une suite $M$-régulière de $r$ éléments appartenant à $\mathfrak J$.
\end{lemma}

La suffisance de la condition résultant aussitôt de ($0$, 16.4.5), prouvons sa nécessité. Raisonnons par récurrence sur $r$ (il n'y a rien à démontrer si $r=0$): puisque $\operatorname{prof}_{A_{\mathfrak p}}(M_{\mathfrak p})\geq r-1$ pour tout $\mathfrak p\supset\mathfrak J$, il existe par hypothèse une suite $M$-régulière $(g_i)_{1\leq i\leq r-1}$ formée d'éléments de $\mathfrak J$. Posons $N=M/(g_1M+\cdots+g_{r-1}M)$; pour tout $\mathfrak p\supset\mathfrak J$, les images des $g_i$ dans $A_{\mathfrak p}$ appartiennent donc à l'idéal maximal $\mathfrak pA_{\mathfrak p}$, et forment une suite $M_{\mathfrak p}$-régulière ($0$, 15.1.14); on conclut donc de ($0$, 16.4.6) que l'on a $\operatorname{prof}_{A_{\mathfrak p}}(N_{\mathfrak p})=\operatorname{prof}_{A_{\mathfrak p}}(M_{\mathfrak p})-(r-1)\geq1$, et on voit que tout revient à prouver le lemme pour $r=1$. Or l'hypothèse signifie alors que, pour tout $\mathfrak p\supset\mathfrak J$, $\mathfrak pA_{\mathfrak p}$ n'est pas associé à $M_{\mathfrak p}$ ($0$, 16.4.6), donc (Bourbaki, \emph{Alg. comm.}, chap. IV, § 1, n\textsuperscript{o} 2, prop. 5), $\mathfrak p$ n'est pas associé à $M$. Autrement dit, $\mathfrak J$ n'est contenu dans aucun des idéaux premiers de $\operatorname{Ass}(M)$, donc il n'est pas contenu dans leur réunion (Bourbaki, \emph{Alg. comm.}, chap. II, § 1, n\textsuperscript{o} 1, prop. 2). Mais comme cette réunion est l'ensemble des éléments qui ne sont pas $M$-réguliers (Bourbaki, \emph{Alg. comm.}, chap. IV, § 1, n\textsuperscript{o} 1, cor. 2 de la prop. 2), cela prouve le lemme.

\begin{proposition}[19.9.3]
\label{IV.19.9.3}
Soient $X$ un préschéma localement noethérien, $\sh F$ un $\sh O_X$-Module cohérent, $Y$ une partie fermée de $X$, $y$ un point de $Y$, $n$ un entier $>0$. Les conditions suivantes sont équivalentes:
\begin{enumerate}
  \item[a)] Il existe un voisinage ouvert $U$ de $y$ dans $X$ tel que $\operatorname{prof}_{Y\cap U}(\sh F|U)\geq n$.
  \item[b)] Il existe un voisinage ouvert $U$ de $y$ dans $X$ et une suite $(\sh F|U)$-régulière de $n$ sections $f_i$ de $\sh O_U$ au-dessus de $U$, telles que $f_i(x)=0$ ($0_{\mathrm I}$, 5.5.1) pour $1\leq i\leq n$ en tout point $x\in U\cap Y$.
  \item[c)] On a $\operatorname{prof}_{Y,y}(\sh F)\geq n$.
\end{enumerate}
En outre, l'ensemble des $y\in Y$ vérifiant ces conditions est ouvert dans $Y$.
\end{proposition}

La dernière assertion résulte trivialement de $a)$. L'équivalence de $a)$ et $b)$ découle de (19.9.2) appliqué à un voisinage ouvert affine de $y$ dans $X$ et de ($0$, 15.1.14). Il est clair que $a)$ entraîne $c)$, tout voisinage ouvert de $y$ dans $X$ contenant $\Spec(\sh O_{X,y})$; montrons inversement que $c)$ implique $b)$. Il résulte de $c)$ et de la définition (19.9.1.1) que l'on

\oldpage[IV-4]{223}

peut appliquer (19.9.2) à l'anneau $\sh O_{X,y}$, au $\sh O_{X,y}$-module $\sh F_y$ et à l'idéal $\sh J_y$ de $\sh O_{X,y}$, donc il existe une suite $\sh F_y$-régulière $(t_i)_{1\leq i\leq n}$ de $n$ éléments de $\sh J_y$. Il y a donc un voisinage ouvert $U$ de $y$ dans $X$ et une suite $(f_i)_{1\leq i\leq n}$ de sections de $\sh J$ au-dessus de $U$ telles que les $t_i$ soient les germes des $f_i$ au point $y$; utilisant ($0$, 15.2.4), on en déduit qu'il existe un voisinage ouvert $U'\subset U$ de $y$ dans $X$ tel que les $f_i|U'$ forment une suite $(\sh F|U')$-régulière.

\begin{corollary}[19.9.4]
\label{IV.19.9.4}
Avec les notations de (19.9.3), la fonction $y\mapsto\operatorname{prof}_{Y,y}(\sh F)$ est semi-continue inférieurement dans $Y$.
\end{corollary}

\begin{proposition}[19.9.5]
\label{IV.19.9.5}
Avec les notations de (19.9.3), soient $X'$ un préschéma localement noethérien, $g:X'\to X$ un morphisme plat, $Y'=g^{-1}(Y)$, $\sh F'=g^*(\sh F)$. Alors, pour tout point $y'\in Y'$ au-dessus de $y$, on a $\operatorname{prof}_{Y',y'}(\sh F')=\operatorname{prof}_{Y,y}(\sh F)$.
\end{proposition}

En effet, pour tout $z'\in\Spec(\sh O_{X',y'})$, il résulte de (6.3.1) que l'on a $\operatorname{prof}(\sh F'_{z'})\geq\operatorname{prof}(\sh F_{g(z')})$, et si $z''$ est un point maximal de la fibre $g^{-1}(g(z'))$, générisation de $z'$ (donc appartenant aussi à $\Spec(\sh O_{X',y'})$), on a (\emph{loc. cit.}) $\operatorname{prof}(\sh F'_{z''})=\operatorname{prof}(\sh F_{g(z')})$. La proposition résulte alors de ce que l'on a $g(\Spec(\sh O_{X',y'}))=\Spec(\sh O_{X,g(y')})$ puisque $g$ est plat (2.3.4).

\begin{proposition}[19.9.6]
\label{IV.19.9.6}
Soient $f:X\to S$ un morphisme localement de présentation finie, $Y$ une partie fermée de $X$, $\sh F$ un $\sh O_X$-Module de présentation finie et $f$-plat, $n$ un entier $>0$. Pour tout point $y\in Y$, les conditions suivantes sont équivalentes (en désignant par $X_{f(y)}$ la fibre de $f$ au point $f(y)$, par $\sh F_{f(y)}$ l'image réciproque de $\sh F$ dans $X_{f(y)}$, et en posant $Y_{f(y)}=X_{f(y)}\cap Y$):
\begin{enumerate}
  \item[a)] On a $\operatorname{prof}_{Y_{f(y)},y}(\sh F_{f(y)})\geq n$.
  \item[b)] Il existe un voisinage ouvert $U$ de $y$ dans $X$ et une suite $(g_i)_{1\leq i\leq n}$ de $n$ sections de $\sh O_X$ au-dessus de $U$, qui soit transversalement $\sh F$-régulière relativement à $S$ (cf. (19.2.1) pour la terminologie) et telle que $g_i(x)=0$ pour $1\leq i\leq n$ et pour tout $x\in U\cap Y$.
\end{enumerate}
L'ensemble $V_n$ des $y\in Y$ vérifiant ces conditions est ouvert dans $Y$; si de plus $Y$ est localement constructible dans $X$, $V_n$ est rétrocompact dans $X$ ($0_{\mathrm{III}}$, 9.1.1).
\end{proposition}

En vertu de (19.9.3), la condition $a)$ équivaut à l'existence d'un voisinage ouvert $U$ de $y$ dans $X$ et d'une suite $(\sh F_{f(y)}|(U\cap X_{f(y)}))$-régulière $(h_i)_{1\leq i\leq n}$ de sections de $\sh O_{X_{f(y)}}$ au-dessus de $U\cap X_{f(y)}$ telles que $h_i(z)=0$ pour tout $z\in U\cap Y_{f(y)}$. Si on désigne encore par $Y$ un sous-préschéma fermé de $X$ ayant $Y$ pour espace sous-jacent, par $\sh J$ l'Idéal quasi-cohérent de $\sh O_X$ définissant $Y$, on peut encore dire que l'on a $h_i\in\sh J_{f(y)}$ pour $1\leq i\leq n$. Comme on peut remplacer $U$ par un voisinage affine plus petit de $y$, on peut supposer que les $h_i$ sont les images d'une suite $(g_i)_{1\leq i\leq n}$ de sections de $\sh J$ au-dessus de $U$, de sorte que les germes $(g_i)_y$ appartiennent à l'idéal maximal $\mathfrak m_y$ de $\sh O_{X,y}$. L'équivalence de $a)$ et $b)$ résulte alors de (11.3.8), qui prouve aussi que l'ensemble $V_n$ est ouvert dans $Y$.

Reste à prouver la dernière assertion; comme il s'agit encore de propriétés locales sur $X$, on peut supposer $S=\Spec(A)$ affine et $X$ de présentation finie sur $S$; il existe alors un sous-anneau noethérien $A_0$ de $A$, un préschéma $X_0$ de type fini sur $S_0=\Spec(A_0)$ et un $\sh O_{X_0}$-Module cohérent $\sh F_0$ tels que $X=X_0\times_{S_0}S$ et $\sh F=\sh F_0\otimes_{\sh O_{X_0}}\sh O_X$ (8.9.1); on peut en outre supposer que, si $f_0:X_0\to S_0$ est le morphisme structural, $\sh F_0$ est $f_0$-plat (11.2.6). Puisque $Y$ est constructible, on peut en outre (8.3.11) supposer que $Y=p^{-1}(Y_0)$, où $p:X\to X_0$ est la projection canonique. Alors, pour tout $y\in Y$, il résulte de (19.9.5) et de la transitivité des fibres que l'on a
\[
  \operatorname{prof}_{Y_{f(y)},y}(\sh F_{f(y)})
  =\operatorname{prof}_{(Y_0)_{f_0(y_0)},y_0}((\sh F_0)_{f_0(y_0)}),
\]
en posant $y_0=p(y)$; si $V_n^0$ est l'ensemble des $y_0\in Y_0$ pour lesquels le second membre de cette relation est $\geq n$, on a donc $V_n=p^{-1}(V_n^0)$. Comme $X_0$ est noethérien, $V_n^0$ est un ouvert rétrocompact dans $X_0$, et par suite (cf. démonstration de (1.8.2)) $V_n$ est rétrocompact dans $X$.

\oldpage[IV-4]{224}

\begin{corollary}[19.9.7]
\label{IV.19.9.7}
Sous les hypothèses générales de (19.9.6), la fonction $y\mapsto\operatorname{prof}_{Y_{f(y)},y}(\sh F_{f(y)})$ est semi-continue inférieurement dans $Y$. Si de plus $Y$ est localement constructible, cette fonction est localement constructible.
\end{corollary}

\begin{proposition}[19.9.8]
\label{IV.19.9.8}
Soient $f:X\to S$ un morphisme localement de présentation finie, $Z$ une partie fermée de $X$, $\sh F$ un $\sh O_X$-Module de présentation finie et $f$-plat, $\sh G$ un $\sh O_S$-Module quasi-cohérent. Supposons que, pour tout $x\in Z$, on ait $\operatorname{prof}((\sh F_{f(x)})_x)\geq1$ (resp. $\operatorname{prof}((\sh F_{f(x)})_x)\geq2$). Alors, si $j:X-Z\to X$ est l'injection canonique, et si on pose $\sh H=\sh F\otimes_{\sh O_X}f^*(\sh G)$, l'homomorphisme canonique $\rho_{\sh H}:\sh H\to j_*(j^*(\sh H))$ relatif à $j$ est injectif (resp. bijectif).
\end{proposition}

La question est évidemment locale sur $X$ et $S$, et il suffit de démontrer la proposition dans un voisinage d'un point $x\in Z$. Autrement dit, on peut se borner au cas où $X$ et $S$ sont affines. Grâce à (19.9.6), on peut supposer qu'il existe une section transversalement $\sh F$-régulière $g_1$ de $\sh O_X$ au-dessus de $X$ relativement à $S$ (resp. une suite transversalement $\sh F$-régulière $(g_1,g_2)$ de deux sections de $\sh O_X$ au-dessus de $X$ relativement à $S$) telle que si $Z'$ est l'ensemble des $x'\in X$ pour lesquels $g_1(x')=0$ (resp. $g_1(x')=g_2(x')=0$), on ait $Z\subset Z'$.

Cela étant, revenons à la démonstration de (19.9.8); en vertu de (19.9.6) on a encore $\operatorname{prof}((\sh F_{f(x)})_x)\geq1$ (resp. $\operatorname{prof}((\sh F_{f(x)})_x)\geq2$) pour tout $x\in Z'$. Si la proposition est prouvée en remplaçant le couple $(X,Z)$ par $(X,Z')$ et $(X-Z,Z'\cap(X-Z))$, il est immédiat qu'elle le sera aussi pour $(X,Z)$; on peut donc se borner à la prouver pour $X$ et $Z'$. Autrement dit, on peut se borner à démontrer (19.9.8) lorsque $S=\Spec(A)$, $X=\Spec(B)$, où $B$ est une $A$-algèbre de présentation finie, $Z=V(\mathfrak J)$, où $\mathfrak J$ est un idéal de type fini de l'anneau $B$, $\sh F=\widetilde M$ et $\sh G=\widetilde N$, où $N$ est un $A$-module et $M$ un $B$-module de présentation finie qui est un $A$-module plat. On peut en outre se ramener au cas où $N$ est un $A$-module de présentation finie. En effet, $N$ est limite inductive d'un système inductif filtrant de $A$-modules de présentation finie $N_\alpha$ (par double limite inductive) si l'on pose $\sh G_\alpha=\widetilde{N_\alpha}$, $\sh H$ est limite inductive des $\sh H_\alpha=\sh F\otimes_{\sh O_X}f^*(\sh G_\alpha)$; comme $j_*$ et $j^*$ commutent aux limites inductives et que $\varinjlim$ est un foncteur exact dans la catégorie des Modules quasi-cohérents, si la proposition est prouvée pour chacun des $\sh H_\alpha$, elle le sera pour $\sh H$.

Il suffit en outre de prouver que l'homomorphisme canonique
\[
  \Gamma(X,\sh H)\longrightarrow\Gamma(X-Z,\sh H)
\]
est injectif (resp. bijectif). Il existe alors un sous-anneau noethérien $A_0$ de $A$, une $A_0$-algèbre de type fini $B_0$, un idéal $\mathfrak J_0$ de $B_0$ et un $B_0$-module de type fini $M_0$ qui est

\oldpage[IV-4]{225}

un $A_0$-module plat et un $A_0$-module $N_0$, tels que $B=B_0\otimes_{A_0}A$, $\mathfrak J=\mathfrak J_0B$, $M=M_0\otimes_{A_0}A$ et $N=N_0\otimes_{A_0}A$ (8.9.1, 8.5.11 et 11.2.7). De plus, soit $(A_\lambda)$ la famille des sous-$A_0$-algèbres de type fini de $A$, de sorte que $A=\varinjlim A_\lambda$; posons
\[
\begin{gathered}
  B_\lambda=B_0\otimes_{A_0}A_\lambda,\qquad
  \mathfrak J_\lambda=\mathfrak J_0B_\lambda,\qquad
  M_\lambda=M_0\otimes_{A_0}A_\lambda=M_0\otimes_{B_0}B_\lambda,\\
  N_\lambda=N_0\otimes_{A_0}A_\lambda,\qquad
  S_\lambda=\Spec(A_\lambda),\quad
  X_\lambda=\Spec(B_\lambda),\quad
  Z_\lambda=V(\mathfrak J_\lambda);
\end{gathered}
\]
si $p_\lambda:X\to X_\lambda$ est la projection canonique, on a donc $Z=p_\lambda^{-1}(Z_\lambda)$, et $X-Z=p_\lambda^{-1}(X_\lambda-Z_\lambda)$. Comme $X_\lambda-Z_\lambda$ est quasi-compact et quasi-séparé, on en conclut (8.5.2) que si l'on pose $\sh F_\lambda=\widetilde{M_\lambda}$ et $\sh H_\lambda=(M_\lambda\otimes_{A_\lambda}N_\lambda)^\sim$, les homomorphismes
\[
  \varinjlim\Gamma(X_\lambda,\sh H_\lambda)\longrightarrow\Gamma(X,\sh H)
  \qquad\text{et}\qquad
  \varinjlim\Gamma(X_\lambda-Z_\lambda,\sh H_\lambda)\longrightarrow\Gamma(X-Z,\sh H)
\]
sont bijectifs; en vertu de l'exactitude du foncteur $\varinjlim$, il suffira donc de prouver que pour tout $\lambda$ assez grand, l'homomorphisme canonique $\Gamma(X_\lambda,\sh H_\lambda)\to\Gamma(X_\lambda-Z_\lambda,\sh H_\lambda)$ est injectif (resp. bijectif). Mais pour tout $x\in Z$, si $x_\lambda$ est la projection de $x$ dans $Z_\lambda$, on a
\[
  \operatorname{prof}((\sh F_{f(x)})_x)
  =\operatorname{prof}(((\sh F_\lambda)_{f_\lambda(x_\lambda)})_{x_\lambda})
\]
en vertu de (4.2.7) et (6.7.1); pour être ramené à démontrer (19.9.8) dans le cas où $A$ est noethérien, il faut prouver que, pour $\lambda$ assez grand on a $\operatorname{prof}(((\sh F_\lambda)_{f_\lambda(x_\lambda)})_{x_\lambda})\geq1$ (resp. $\geq2$) pour tout point $x_\lambda\in Z_\lambda$. Or, notons $Z_\lambda'$ l'ensemble des points $x_\lambda\in Z_\lambda$ tels que, pour toute générisation $z_\lambda$ de $x_\lambda$ dans $Z_\lambda\cap(X_\lambda)_{f_\lambda(x_\lambda)}$, on ait $\operatorname{prof}(((\sh F_\lambda)_{f_\lambda(x_\lambda)})_{z_\lambda})\geq1$; si $p_{\lambda\mu}:Z_\mu\to Z_\lambda$ est la projection canonique pour $\lambda\leq\mu$, il résulte de l'hypothèse et de (19.9.5) que l'on a $Z_\mu'=p_{\lambda\mu}^{-1}(Z_\lambda')$ et $Z=p_\lambda^{-1}(Z_\lambda')$; d'autre part, il résulte de (19.9.6) que $Z_\lambda'$ est ouvert dans $Z_\lambda$, donc notre assertion résulte de (8.3.4).

Pour terminer la démonstration lorsque $A$ est noethérien, notons qu'en vertu de l'hypothèse de platitude et de (6.3.1), on a alors $\operatorname{prof}_Z(\sh H)\geq1$ (resp. $\operatorname{prof}_Z(\sh H)\geq2$); mais la conclusion de (19.9.8) résulte dans ce cas de (5.10.2) (resp. 5.10.4)).

\begin{remark}[19.9.9]
\label{IV.19.9.9}
La même démonstration, jointe aux résultats du chap. III, 3\textsuperscript{e} partie sur la profondeur et la cohomologie locale, permet d'établir la généralisation suivante de (19.9.8):

\emph{Sous les conditions générales de (19.9.8), supposons que, pour tout $x\in Z$, on ait $\operatorname{prof}((\sh F_{f(x)})_x)\geq k$; alors l'homomorphisme canonique}
\[
  H^i(X,\sh H)\longrightarrow H^i(X-Z,\sh H|(X-Z))
\]
\emph{est bijectif pour $i\leq k-2$ et injectif pour $i=k-1$} (ce qui, exprimé à l'aide de la notion cohomologique générale de profondeur introduite au chap. III, s'écrit aussi $\operatorname{prof}_Z(\sh H)\geq k$).
\end{remark}

% FIN SOURCE ega4/ega4-18-fr.tex

\newcommand{\shHom}{\mathcal{H}\!om}
%
% DEBUT SOURCE ega4/ega4-20-fr.tex
\section{Fonctions méromorphes et pseudo-morphismes}
\label{IV.20}

\setcounter{subsection}{-1}
\subsection{Introduction}
\EGAReaderTarget{ega4.c04.subsec.20-0}{IV.20.0}

La plupart des notions et résultats des \textsection\textsection20 et 21 se rattachent directement au chap.~I, et ne dépendent guère des chap.~II à IV, sauf en ce qui concerne l'usage occasionnel de la notion de profondeur et d'anneau local régulier (dans \sref{IV.20.6}, \sref{IV.21.11}, \sref{IV.21.13} et \sref{IV.21.15}), du « Main theorem » de Zariski dans \sref{IV.20.4} et \sref{IV.21.12}, et des propriétés des immersions transversalement régulières dans \sref{IV.20.6} et \sref{IV.21.15}.

\oldpage[IV-4]{226}
Au \textsection20, nous introduisons plusieurs variantes de la notion d'application rationnelle, déjà étudiée dans \sref[I]{section:I.7} d'un point de vue encore assez proche du point de vue classique, et pour cette raison assez mal adaptée au cas des préschémas non nécessairement réduits.
Les notions et résultats du \textsection20 sont utilisés pour développer au \textsection21 (n\textsuperscript{os} \sref{IV.21.1} à \sref{IV.21.7}) la notion générale de diviseur et ses propriétés les plus élémentaires.
Cette notion est surtout commode lorsque les anneaux locaux des préschémas considérés sont noethériens et intégralement clos, et surtout lorsqu'ils sont en outre \emph{factoriels} (\sref{IV.21.6} et \sref{IV.21.7}), en raison de son identification dans ce dernier cas avec la notion de \emph{cycle $1$-codimensionnel} (combinaison linéaire de sous-préschémas irréductibles de codimension~$1$).
Dans \sref{IV.21.9}, on détermine les diviseurs sur un préschéma noethérien de dimension~$1$ mais non nécessairement normal, ce qui est utile pour diverses applications.
Les n\textsuperscript{os} \sref{IV.21.11} et \sref{IV.21.12} donnent deux théorèmes importants, dus respectivement à Auslander--Buchsbaum et Van der Waerden, et se rattachant à la notion d'anneau factoriel (les n\textsuperscript{os} \sref{IV.21.9}, \sref{IV.21.11} et \sref{IV.21.12} sont indépendants les uns des autres).
Dans les n\textsuperscript{os} \sref{IV.21.13} et \sref{IV.21.14}, également indépendants des trois précédents, nous étudions une variante utile de la notion d'anneau local factoriel, celle d'anneau local \emph{parafactoriel}, qui s'introduit notamment~\cite{IV-41} dans le développement de théorèmes de comparaison du groupe de Picard d'un préschéma projectif $X$ sur un corps $k$ et d'une « section hyperplane ».
On verra dans \sref{IV.21.14.1} (théorème de Ramanujam--Samuel) que les anneaux locaux parafactoriels sont bien plus nombreux que l'on ne pouvait s'y attendre \emph{a priori}.

Dans \sref{IV.20.5}, \sref{IV.20.6} et \sref{IV.21.15}, nous reprenons les notions précédentes mais d'un point de vue « relatif » à un préschéma de base fixé.
Pour l'instant ces notions ne sont utilisées qu'assez rarement ; en particulier la notion de diviseur relatif n'est guère employée que lorsqu'il s'agit de diviseurs positifs, et dans ce cas elle s'explicite avantageusement sans le secours de la notion de fonction méromorphe relative, à l'aide de la notion d'immersion transversalement régulière de codimension~$1$.
On aura donc avantage à omettre ces sections en première lecture.

\subsection{Fonctions méromorphes}
\EGAReaderTarget{ega4.c04.subsec.20-1}{IV.20.1}

\begin{env}[20.1.1]
\EGAReaderTarget{ega4.c04.statement.20-1-1}{IV.20.1.1}
Soit $(X,\sh{O}_X)$ un espace annelé, et soit $\sh{S}$ un sous-faisceau \emph{d'ensembles} de $\sh{O}_X$.
Pour tout ouvert $U$ de $X$, considérons l'\emph{anneau de fractions} $\Gamma(U,\sh{O}_X)[\Gamma(U,\sh{S})^{-1}]$ (Bourbaki, \emph{Alg. comm.}, chap.~II, \textsection2, n\textsuperscript{o}1).
Il est immédiat que l'application $U\mapsto\Gamma(U,\sh{O}_X)[\Gamma(U,\sh{S})^{-1}]$ est un \emph{préfaisceau d'anneaux} (\sref[0]{0.1.5.1} et \sref[0]{0.1.5.7}).
On note $\sh{O}_X[\sh{S}^{-1}]$ le \emph{faisceau d'anneaux} associé à ce préfaisceau et on dit que c'est le \emph{faisceau d'anneaux de fractions de $\sh{O}_X$, à dénominateurs dans $\sh{S}$} ; c'est un $\sh{O}_X$-module \emph{plat}.
Il est immédiat que pour tout $x\in X$, on a un isomorphisme canonique
\[
\label{IV.20.1.1.1}
  (\sh{O}_X[\sh{S}^{-1}])_x\isoto\sh{O}_{X,x}[\sh{S}_x^{-1}],
  \tag{20.1.1.1}
\]
car le raisonnement de \sref[0]{0.1.4.5} se généralise aussitôt au cas où l'on a un système inductif $(A_\alpha,\vphi_{\beta\alpha})$ d'anneaux, et pour chaque indice $\alpha$ une partie $S_\alpha$ de $A_\alpha$ telle que
\oldpage[IV-4]{227}
$\vphi_{\beta\alpha}(S_\alpha)\subset S_\beta$ pour $\alpha\leq\beta$ ; on prend alors pour $S$ la limite inductive dans $A=\varinjlim A_\alpha$ du système inductif de parties $(S_\alpha)$.
\end{env}

\begin{env}[20.1.2]
\EGAReaderTarget{ega4.c04.statement.20-1-2}{IV.20.1.2}
Soit maintenant $\sh{F}$ un $\sh{O}_X$-module.
On pose alors
\[
\label{IV.20.1.2.1}
  \sh{F}[\sh{S}^{-1}]=\sh{F}\otimes_{\sh{O}_X}\sh{O}_X[\sh{S}^{-1}]
  \tag{20.1.2.1}
\]
et on dit que c'est le \emph{faisceau de modules de fractions de $\sh{F}$ à dénominateurs dans $\sh{S}$} ; il est immédiat qu'il est associé au préfaisceau de modules $U\mapsto\Gamma(U,\sh{F})[\Gamma(U,\sh{S})^{-1}]$, et que pour tout $x\in X$, on a un isomorphisme canonique
\[
\label{IV.20.1.2.2}
  (\sh{F}[\sh{S}^{-1}])_x\isoto\sh{F}_x[\sh{S}_x^{-1}].
  \tag{20.1.2.2}
\]
\end{env}

\begin{env}[20.1.3]
\EGAReaderTarget{ega4.c04.statement.20-1-3}{IV.20.1.3}
Nous allons nous intéresser ici au cas où $\sh{S}$ est le sous-faisceau $\sh{S}(\sh{O}_X)$ de $\sh{O}_X$ tel que pour tout ouvert $U$, $\Gamma(U,\sh{S})$ soit l'\emph{ensemble des éléments réguliers} de l'anneau $\Gamma(U,\sh{O}_X)$ ; il est immédiat qu'il s'agit d'un faisceau (et non seulement d'un préfaisceau), la régularité d'une section de $\sh{O}_X$ au-dessus de $U$ se vérifiant « fibre par fibre » (i.e. signifiant que le germe de la section en $x$ est régulier dans $\sh{O}_{X,x}$ pour tout $x\in U$), autrement dit $\sh{S}(\sh{O}_X)_x$ n'est autre que l'ensemble des éléments réguliers de $\sh{O}_{X,x}$.
Le faisceau d'anneaux correspondant
\[
  \sh{M}_X=\sh{O}_X[\sh{S}^{-1}]
\]
est appelé le \emph{faisceau des germes de fonctions méromorphes sur $X$}, et les sections de $\sh{M}_X$ au-dessus de $X$ sont appelées les \emph{fonctions méromorphes sur $X$} ; elles forment un anneau que l'on note $M(X)$.
Pour tout $\sh{O}_X$-module $\sh{F}$,
\[
  \sh{F}\otimes_{\sh{O}_X}\sh{M}_X=\sh{F}[\sh{S}^{-1}]
\]
est encore noté $\sh{M}_X(\sh{F})$ et appelé le \emph{faisceau des germes de sections méromorphes de $\sh{F}$} ; ses sections au-dessus de $X$ forment un $M(X)$-module noté $M(X,\sh{F})$, dont les éléments sont appelés \emph{sections méromorphes de $\sh{F}$ au-dessus de $X$}.
Ces définitions impliquent que pour tout ouvert $U$ de $X$, on a un isomorphisme canonique $\sh{M}_X(\sh{F})|U\isoto\sh{M}_U(\sh{F}|U)$, en particulier $\sh{M}_X|U\isoto\sh{M}_U$.
\end{env}

\begin{env}[20.1.3.1]
\label{IV.20.1.3.1}
Si $X$ est un \emph{préschéma réduit}, on notera que si un élément $s\in\Gamma(U,\sh{O}_X)$ est tel que $s_\xi\neq0$ pour tout point maximal $\xi$ de $U$, alors $s$ est \emph{régulier}.
En effet, si $st=0$ pour un $t\in\Gamma(U,\sh{O}_X)$, on a $s_\xi t_\xi=0$, donc $t_\xi=0$ puisque $\sh{O}_{X,\xi}$ est un corps, et dire que $t_\xi=0$ pour tout point maximal $\xi$ de $X$ signifie que $t=0$ : en effet, on est aussitôt ramené au cas où $U$ est affine, et un élément d'un anneau réduit qui appartient à tous les idéaux premiers minimaux est nul par définition.
La réciproque est vraie si l'ensemble des composantes irréductibles de $X$ est \emph{localement fini}.
On est en effet aussitôt ramené au cas où $X=\Spec(A)$ est affine ; si $\mathfrak p_i$ ($1\leq i\leq n$) sont les idéaux premiers minimaux de $A$ et si $s\in\mathfrak p_i$ pour un indice $i$, il existe $t\in A$ tel que $t\in\mathfrak p_j$ pour $j\neq i$ et $t\notin\mathfrak p_i$ (Bourbaki, \emph{Alg. comm.}, chap.~II, \textsection1, n\textsuperscript{o}1, prop.~1) ; on a donc $st\in\mathfrak p_i$ pour tout $i$, et par suite $st=0$ puisque $A$ est réduit ; $s$ est donc non régulier.
\end{env}

\begin{env}[20.1.4]
\EGAReaderTarget{ega4.c04.statement.20-1-4}{IV.20.1.4}
Pour tout ouvert $U$ de $X$, l'homomorphisme $t\mapsto t/1$ de $\Gamma(U,\sh{O}_X)$ dans $\Gamma(U,\sh{O}_X)[\Gamma(U,\sh{S})^{-1}]$ (qui n'est autre que l'\emph{anneau total des fractions} de
\oldpage[IV-4]{228}
$\Gamma(U,\sh{O}_X)$) est injectif ; ces homomorphismes définissent donc un \emph{homomorphisme canonique injectif}
\[
\label{IV.20.1.4.1}
  i:\sh{O}_X\to\sh{M}_X
  \tag{20.1.4.1}
\]
qui permet d'identifier $\sh{O}_X$ à un sous-faisceau de $\sh{M}_X$.
Étant donnée une fonction méromorphe $\vphi\in M(X)$, on dit que $\vphi$ est \emph{définie} dans un ouvert $U$ de $X$ si $\vphi|U$ est une \emph{section de $\sh{O}_U$} au-dessus de $U$ ; les axiomes des faisceaux montrent qu'il y a, pour une section $\vphi$ donnée, un \emph{plus grand} ouvert dans lequel $\vphi$ est définie ; on l'appelle \emph{domaine de définition} de $\vphi$ et on le note $\operatorname{dom}(\vphi)$.
\end{env}

\begin{env}[20.1.5]
\EGAReaderTarget{ega4.c04.statement.20-1-5}{IV.20.1.5}
Pour tout $\sh{O}_X$-module $\sh{F}$, on déduit de \sref{IV.20.1.4.1} un di-homomorphisme formé de $i$ et de l'homomorphisme de faisceaux de groupes additifs
\[
\label{IV.20.1.5.1}
  1_{\sh{F}}\otimes i:\sh{F}\to\sh{M}_X(\sh{F})=\sh{F}\otimes_{\sh{O}_X}\sh{M}_X.
  \tag{20.1.5.1}
\]
On notera que ce dernier n'est plus injectif en général ; lorsqu'il est injectif, on dit que $\sh{F}$ est \emph{strictement sans torsion} : cela signifie que pour tout ouvert $U$ de $X$ et toute section $s\in\Gamma(U,\sh{O}_X)$ qui est un élément régulier de cet anneau, l'homothétie $z\mapsto sz$ de $\Gamma(U,\sh{F})$ est injective ; cette condition est évidemment remplie si $\sh{F}$ est \emph{localement libre}.
\end{env}

\begin{proposition}[20.1.6]
\label{IV.20.1.6}
Soient $X$ un préschéma localement noethérien, $\sh{F}$ un $\sh{O}_X$-module quasi-cohérent.
Pour que $\sh{F}$ soit strictement sans torsion, il faut et il suffit que $\operatorname{Ass}(\sh{F})\subset\operatorname{Ass}(\sh{O}_X)$.
\end{proposition}

\begin{proof}
On est en effet aussitôt ramené au cas où $X=\Spec(A)$ est affine, $\sh{F}=\widetilde M$, et on sait que les éléments $s$ de $A$ appartenant à un idéal de $\operatorname{Ass}(M)$ sont exactement ceux pour lesquels l'homothétie $z\mapsto sz$ n'est pas injective (Bourbaki, \emph{Alg. comm.}, chap.~IV, \textsection1, n\textsuperscript{o}1, cor.~2 de la prop.~2).
\end{proof}

\begin{env}[20.1.7]
\EGAReaderTarget{ega4.c04.statement.20-1-7}{IV.20.1.7}
Si $u$ est une section de $\sh{M}_X(\sh{F})$ au-dessus de $X$, on dit que $u$ est \emph{définie} en un point $x\in X$ s'il existe un voisinage ouvert $V$ de $x$ dans $X$ tel que $u|V$ soit l'image d'une section de $\sh{F}$ au-dessus de $V$ par le di-homomorphisme \sref{IV.20.1.5.1}.
On dit que $u$ est \emph{définie} dans un ouvert $U$ de $X$ si elle est définie en tout point de $U$ ; il y a encore un plus grand ouvert dans lequel $u$ est définie, appelé \emph{domaine de définition} de $u$ et noté $\operatorname{dom}(u)$.
Lorsque $\sh{F}$ est strictement sans torsion, de sorte que $\sh{F}$ s'identifie par \sref{IV.20.1.5.1} à un sous-faisceau de $\sh{M}_X(\sh{F})$, dire que $u$ est définie dans $U$ signifie que $u|U$ est une \emph{section de $\sh{F}$ au-dessus de $U$}.
\end{env}

\begin{env}[20.1.8]
\EGAReaderTarget{ega4.c04.statement.20-1-8}{IV.20.1.8}
Conformément aux notations générales \sref[0\textsubscript{I}]{0.5.4.7}, on note $\sh{M}_X^*$ le faisceau de groupes multiplicatifs tel que $\Gamma(U,\sh{M}_X^*)$ soit (pour tout ouvert $U$ de $X$) le groupe des \emph{éléments inversibles} de $\Gamma(U,\sh{M}_X)$.
Ce faisceau n'est autre que le faisceau $\sh{S}(\sh{M}_X)$ défini en \sref{IV.20.1.3} : en effet, si $s\in\Gamma(U,\sh{S}(\sh{M}_X))$, pour tout $x\in U$, il existe un voisinage ouvert $V\subset U$ de $x$ tel que $s|V$ soit un élément régulier dans l'\emph{anneau total des fractions} de $\Gamma(V,\sh{O}_X)$ et on sait qu'un tel élément est nécessairement inversible dans cet anneau de fractions.
Nous dirons que les sections de $\sh{M}_X^*$ au-dessus de $X$ sont les \emph{fonctions méromorphes régulières} (on notera que nous nous écartons ici de la terminologie suivie par certains auteurs, qui nomment fonctions méromorphes « régulières » celles qui sont des \emph{sections de $\sh{O}_X$}, identifié à un sous-faisceau de $\sh{M}_X$).

Soit $\sh{L}$ un $\sh{O}_X$-module inversible \sref[0\textsubscript{I}]{0.5.4.1} ; alors il est clair que $\sh{M}_X(\sh{L})=\sh{L}\otimes_{\sh{O}_X}\sh{M}_X$
\oldpage[IV-4]{229}
est un $\sh{M}_X$-module inversible.
Soit $U$ un ouvert tel que $\sh{L}|U$ soit isomorphe à $\sh{O}_U$ ; comme tout automorphisme de $\sh{M}_U$ est la multiplication par un élément inversible de $\Gamma(U,\sh{M}_X)$ \sref[0\textsubscript{I}]{0.5.4.7}, il revient au même de dire qu'une section $s\in\Gamma(U,\sh{M}_X(\sh{L}))$ a une image inversible dans $\Gamma(U,\sh{M}_X)$ par \emph{un} isomorphisme ou par \emph{tout} isomorphisme sur $\Gamma(U,\sh{M}_X)$ ; on dira dans ce cas que $s$ est une \emph{section méromorphe régulière de $\sh{L}$ au-dessus de $U$} ; une section $s$ de $\sh{L}$ au-dessus de $X$ sera appelée une \emph{section méromorphe régulière de $\sh{L}$} si, pour tout ouvert $U$ tel que $\sh{L}|U$ soit isomorphe à $\sh{O}_U$, $s|U$ est une section méromorphe régulière de $\sh{L}$ au-dessus de $U$.
On notera $(\sh{M}_X(\sh{L}))^*$ le sous-faisceau de $\sh{M}_X(\sh{L})$ tel que pour tout ouvert $U$, $\Gamma(U,(\sh{M}_X(\sh{L}))^*)$ soit l'ensemble des sections méromorphes régulières de $\sh{L}$ au-dessus de $U$.
Soit $s$ une section méromorphe de $\sh{L}$ au-dessus de $X$ (i.e. une section de $\sh{M}_X(\sh{L})$) ; elle définit un homomorphisme $h_s:\sh{M}_X\to\sh{M}_X(\sh{L})$ qui, à toute section $t$ de $\sh{M}_X$ au-dessus d'un ouvert $U$, fait correspondre $(s|U)t$.
Il résulte aussitôt de ce qui précède que, pour que $s$ soit \emph{régulière}, il faut et il suffit que $h_s$ soit \emph{injectif}, et en fait $h_s$ est alors un homomorphisme \emph{bijectif} de $\sh{M}_X$ sur $\sh{M}_X(\sh{L})$, et sa restriction à $\sh{M}_X^*$ est une bijection sur $(\sh{M}_X(\sh{L}))^*$.
On en conclut que l'homothétie $t\mapsto ts$ est un isomorphisme de $M(X)$ sur $M(X,\sh{L})$.
\end{env}

\begin{env}[20.1.9]
\label{IV.20.1.9}
Soit $s$ une section méromorphe régulière du $\sh{O}_X$-module inversible $\sh{L}$ au-dessus de $X$ ; alors, pour tout $\sh{O}_X$-module $\sh{F}$, $s$ définit de même un homomorphisme $h_s\otimes1_{\sh{F}}:\sh{M}_X(\sh{F})\to\sh{M}_X(\sh{F}\otimes_{\sh{O}_X}\sh{L})$, qui est encore \emph{bijectif}.
\end{env}

\begin{env}[20.1.10]
\label{IV.20.1.10}
Soit $s$ une section méromorphe d'un $\sh{O}_X$-module inversible $\sh{L}$ au-dessus de $X$ ; pour que $s$ soit régulière, il faut et il suffit qu'il existe une section méromorphe $s'$ de $\sh{L}^{-1}$ au-dessus de $X$ telle que l'image canonique de $s\otimes s'$ dans $\sh{M}_X$ \sref[0\textsubscript{I}]{0.5.4.3} soit la section unité, et cette section $s'$ est alors unique : en effet, la nécessité de l'existence locale d'une telle section est évidente, et son unicité locale entraîne son existence (et unicité) globale ; en outre, l'existence de $s'$ est trivialement suffisante pour que $s$ soit régulière.
On posera $s'=s^{-1}$.

Enfin, si $\sh{L}'$ est un second $\sh{O}_X$-module inversible, $s$ (resp. $s'$) une section méromorphe régulière de $\sh{L}$ (resp. $\sh{L}'$) au-dessus de $X$, $s\otimes s'$ est évidemment une section méromorphe régulière de $\sh{L}\otimes\sh{L}'$ au-dessus de $X$.
\end{env}

\begin{env}[20.1.11]
\EGAReaderTarget{ega4.c04.statement.20-1-11}{IV.20.1.11}
Si $f:X'\to X$ est un morphisme d'espaces annelés, il n'y a pas en général d'application naturelle associant à une fonction méromorphe sur $X$ une fonction méromorphe sur $X'$.
Par exemple, si $X$ est le spectre d'un anneau local intègre $A$, $X'$ celui de son corps résiduel $k$, il n'y a aucun homomorphisme naturel du corps des fractions $K$ de $A$ dans $k$, et on ne peut faire correspondre à un élément de $K$ un élément de $k$ que s'il est déjà dans $A$.

De façon générale, si $f=(\psi,\theta)$, notons, pour tout ouvert $U$ de $X$, $\sh{S}_f(U)$ l'ensemble des sections \emph{régulières} $s\in\Gamma(U,\sh{O}_X)$ telles que l'image de $s$ par
\[
  \Gamma(\theta^\#):\Gamma(U,\sh{O}_X)\to\Gamma(f^{-1}(U),\sh{O}_{X'})
\]
soit une section \emph{régulière}.
Il est immédiat que $U\mapsto\sh{S}_f(U)$ est un \emph{sous-faisceau} du faisceau d'ensembles $\sh{S}(\sh{O}_X)$, que l'on note $\sh{S}_f$.
On pose $\sh{M}_f=\sh{O}_X[\sh{S}_f^{-1}]$ ; c'est un sous-faisceau
\oldpage[IV-4]{230}
d'anneaux de $\sh{M}_X$, et on déduit canoniquement de $\theta^\#:\psi^*(\sh{O}_X)\to\sh{O}_{X'}$ un homomorphisme de faisceaux d'anneaux $\theta^{\prime\#}:\psi^*(\sh{M}_f)\to\sh{M}_{X'}$ prolongeant $\theta^\#$ (Bourbaki, \emph{Alg. comm.}, chap.~II, \textsection2, n\textsuperscript{o}1, prop.~2) ; d'où, en se rappelant que $f^*(\sh{M}_f)=\psi^*(\sh{M}_f)\otimes_{\psi^*(\sh{O}_X)}\sh{O}_{X'}$, un homomorphisme canonique de $\sh{O}_{X'}$-algèbres
\[
\label{IV.20.1.11.1}
  f^*(\sh{M}_f)\to\sh{M}_{X'}.
  \tag{20.1.11.1}
\]
Pour toute fonction méromorphe $\vphi$ sur $X$, qui \emph{est une section de $\sh{M}_f$}, $\Gamma(\theta^{\prime\#})(\vphi)$ est une fonction méromorphe sur $X'$, dite \emph{image réciproque de $\vphi$ par $f$}, et notée $\vphi\circ f$ si cela n'entraîne pas confusion.

De même, si $\sh{F}$ est un $\sh{O}_X$-module, on pose $\sh{M}_f(\sh{F})=\sh{F}\otimes_{\sh{O}_X}\sh{M}_f$, et on déduit aussitôt de $\theta^{\prime\#}$ un homomorphisme canonique (qu'on écrit aussi $u\mapsto u\circ f$)
\[
  \Gamma(X,\sh{M}_f(\sh{F}))\to\Gamma(X',\sh{M}_{X'}(f^*(\sh{F}))).
\]
En outre, si $u\in\Gamma(X,\sh{M}_f(\sh{F}))$ est définie \sref{IV.20.1.7} en un point $x$, $u$ coïncide, dans un voisinage $U$ de $x$, avec une section de la forme $\sum_i h_i\otimes(t_i/s_i)$, où les $h_i$ appartiennent à $\Gamma(U,\sh{F})$, les $t_i$ à $\Gamma(U,\sh{O}_X)$ et les $s_i$ à $\Gamma(U,\sh{S}_f)$.
Comme par hypothèse les images des $s_i$ dans $\Gamma(f^{-1}(U),\sh{O}_{X'})$ sont régulières, on voit que $u\circ f$ est définie en tout point de $f^{-1}(U)$ ; autrement dit, on a
\[
\label{IV.20.1.11.2}
  f^{-1}(\operatorname{dom}(u))\subset\operatorname{dom}(u\circ f).
  \tag{20.1.11.2}
\]
Nous verrons plus loin \sref{IV.20.6.5}[(i)] des exemples (avec $\sh{F}=\sh{O}_X$) où les deux membres de \sref{IV.20.1.11.2} peuvent être distincts.

Considérons en particulier le cas où $\sh{M}_f=\sh{M}_X$ ; alors, si $\sh{L}$ est un $\sh{O}_X$-module inversible, l'image dans $\sh{M}_{X'}(f^*(\sh{L}))$, par $\Gamma(\theta^{\prime\#})$, d'une section méromorphe \emph{régulière} de $\sh{L}$ au-dessus de $X$ \sref{IV.20.1.8} est une section méromorphe \emph{régulière} de $f^*(\sh{L})$ au-dessus de $X'$, comme il résulte aussitôt de la définition de ces sections, et du fait qu'un homomorphisme d'anneaux transforme un élément inversible en un élément inversible.

Soit $f':X''\to X'$ un second morphisme d'espaces annelés, et supposons que $\sh{M}_f=\sh{M}_X$ et $\sh{M}_{f'}=\sh{M}_{X'}$ ; alors, si l'on pose $f''=f\circ f'$, on a aussi $\sh{M}_{f''}=\sh{M}_X$, et l'on voit aussitôt que pour toute section méromorphe $u$ de $\sh{F}$ au-dessus de $X$, on a $u\circ f''=(u\circ f)\circ f'$.
\end{env}

\begin{proposition}[20.1.12]
\label{IV.20.1.12}
Si le morphisme $f:X'\to X$ est plat \sref[0\textsubscript{I}]{0.6.7.1}, on a $\sh{M}_f=\sh{M}_X$, et l'homomorphisme $\vphi\mapsto\vphi\circ f$ est défini dans $M(X)$ tout entier.
En outre, si $f$ est un morphisme (plat) d'espaces annelés en anneaux locaux, on a $\operatorname{dom}(\vphi\circ f)=f^{-1}(\operatorname{dom}(\vphi))$ ; si de plus $f$ est surjectif (donc fidèlement plat), l'homomorphisme $\vphi\mapsto\vphi\circ f$ est injectif.
\end{proposition}

\begin{proof}
La première assertion résulte de ce que, si $B$ est une $A$-algèbre qui est un $A$-module plat, tout élément de $A$ non diviseur de $0$ dans $A$ est non diviseur de $0$ dans $B$ \sref[0\textsubscript{I}]{0.6.3.4}.
Pour prouver les autres assertions, notons que, pour tout $x'\in X'$, si $x=f(x')$, $\sh{O}_{X',x'}$ est un $\sh{O}_{X,x}$-module plat, et comme l'homomorphisme $\sh{O}_{X,x}\to\sh{O}_{X',x'}$ est \emph{local} par hypothèse, il est injectif (\sref[0\textsubscript{I}]{0.6.5.1} et \sref[0\textsubscript{I}]{0.6.6.2}) ; si l'on pose $A=\sh{O}_{X,x}$, $B=\sh{O}_{X',x'}$, de sorte que $A$ s'identifie à un sous-anneau de $B$, $(f^*(\sh{M}_X))_{x'}$ est égal à $S^{-1}A\otimes_A B=S^{-1}B$, où $S$ est l'ensemble des éléments réguliers de $A$, $(\sh{M}_{X'})_{x'}$ est égal à $T^{-1}B$, où $T$ est l'ensemble
\oldpage[IV-4]{231}
des éléments réguliers de $B$, et comme on a vu que $S\subset T$, l'homomorphisme $S^{-1}B\to T^{-1}B$ est injectif ; autrement dit, cela prouve que l'homomorphisme \sref{IV.20.1.11.1} $f^*(\sh{M}_X)\to\sh{M}_{X'}$ est \emph{injectif} (d'où la dernière assertion de l'énoncé).
Le quotient $f^*(\sh{M}_X)/\sh{O}_{X'}$ s'identifie à un sous-$\sh{O}_{X'}$-module de $\sh{M}_{X'}/\sh{O}_{X'}$, et $(f^*(\sh{M}_X)/\sh{O}_{X'})_{x'}$ s'identifie à $(\sh{M}_X/\sh{O}_X)_x\otimes_{\sh{O}_{X,x}}\sh{O}_{X',x'}$.
Supposons alors que $x\notin\operatorname{dom}(\vphi)$ ; l'image de $\vphi_x$ dans $(\sh{M}_X/\sh{O}_X)_x$ est donc $\neq0$ ; par fidèle platitude, on en déduit qu'il en est de même de l'image de $(\vphi\circ f)_{x'}$ dans $(\sh{M}_{X'}/\sh{O}_{X'})_{x'}$, donc $x'\notin\operatorname{dom}(\vphi\circ f)$, ce qui achève la démonstration.
\end{proof}

\begin{remark}[20.1.13]
\label{IV.20.1.13}
Soit $X$ un espace analytique complexe \emph{réduit} ; alors la notion de fonction méromorphe sur $X$ définie ci-dessus coïncide avec la notion usuelle.
Considérons d'autre part un préschéma $Y$, localement de type fini sur le corps $\bb{C}$ ; on sait alors qu'on peut associer à $Y$ un espace analytique $Y^\mathrm{an}$ ayant même espace topologique sous-jacent, et que le morphisme canonique $f:Y^\mathrm{an}\to Y$ est \emph{plat}~\cite{IV-37} ; en vertu de \sref{IV.20.1.12}, l'homomorphisme canonique $u\mapsto u\circ f$ de $M(Y)$ dans $M(Y^\mathrm{an})$ est donc partout défini et est injectif ; mais il n'est pas \emph{surjectif} en général.
Par exemple, lorsque $Y=\bb{V}_{\bb{C}}^r$ (\textbf{Err}\textsubscript{\textbf{III}},~14) est l'espace affine de dimension $r$ sur $\bb{C}$, $M(Y)$ s'identifie canoniquement au corps $R(Y)$ des fonctions rationnelles sur $Y$ \sref{IV.20.2.13}[(i)], tandis que $M(Y^\mathrm{an})$ est le corps des fonctions méromorphes usuelles sur $\bb{C}^r$.
En raison de ce fait, il est souvent préférable, en Géométrie algébrique, de s'abstenir de la terminologie introduite dans cette section, et d'utiliser la terminologie équivalente de « pseudo-fonction » qui va être définie ci-dessous.
\end{remark}

\subsection{Pseudo-morphismes et pseudo-fonctions}
\EGAReaderTarget{ega4.c04.subsec.20-2}{IV.20.2}

Les seuls espaces annelés dont il sera question dans cette section sont les \emph{préschémas}.

\begin{env}[20.2.1]
\EGAReaderTarget{ega4.c04.statement.20-2-1}{IV.20.2.1}
Rappelons \sref{IV.11.10.2} que dans un préschéma $X$, on dit qu'un ouvert $U$ est \emph{schématiquement dense} si, pour tout ouvert $V$ de $X$, l'homomorphisme canonique $\Gamma(V,\sh{O}_X)\to\Gamma(V\cap U,\sh{O}_X)$ est \emph{injectif}.

Considérons deux préschémas $X$, $Y$, deux ouverts \emph{schématiquement denses} $U$, $U'$ de $X$ ; on dit que deux morphismes $u:U\to Y$, $u':U'\to Y$ sont \emph{équivalents} s'il existe un ouvert $U''\subset U\cap U'$, \emph{schématiquement dense} dans $X$, tel que $u|U''=u'|U''$.
Comme il résulte aussitôt de la définition des ouverts schématiquement denses que l'intersection de deux tels ouverts en est un autre, il est immédiat que la relation précédente est bien une relation d'équivalence.
Une classe d'équivalence suivant cette relation s'appelle un \emph{pseudo-morphisme} de $X$ dans $Y$, ou une \emph{application rationnelle stricte} de $X$ dans $Y$.

Si $S$ est un préschéma et $X$, $Y$ deux $S$-préschémas, on appelle \emph{pseudo-$S$-morphisme} la classe d'équivalence (pour la relation précédente) d'un $S$-morphisme d'un ouvert schématiquement dense de $X$ dans $Y$.
Un pseudo-morphisme n'est donc autre chose qu'un pseudo-$S$-morphisme pour $S=\Spec(\bb{Z})$.

On note $\operatorname{Ps.hom}(X,Y)$ (resp. $\operatorname{Ps.hom}_S(X,Y)$) l'ensemble des pseudo-morphismes (resp. des pseudo-$S$-morphismes) de $X$ dans $Y$.
\end{env}

\begin{env}[20.2.2]
\EGAReaderTarget{ega4.c04.statement.20-2-2}{IV.20.2.2}
\oldpage[IV-4]{232}
Il résulte de la définition rappelée ci-dessus que si $U$ est un ouvert schématiquement dense dans $X$, alors, pour tout ouvert $V$ de $X$, $U\cap V$ est schématiquement dense dans $V$.
Si deux morphismes $u:U\to Y$, $u':U'\to Y$ d'ouverts schématiquement denses de $X$ dans $Y$ sont équivalents il s'ensuit donc que leurs restrictions $u|(V\cap U)$ et $u'|(V\cap U')$ sont aussi des morphismes équivalents (pour le préschéma induit sur $V$) ; leur classe est dite la \emph{restriction} à $V$ du pseudo-morphisme $\omega$, classe de $u$, et on note ce pseudo-morphisme $\omega|V$.
Si $W\subset V$ est un ouvert de $X$, il est clair que $(\omega|V)|W=\omega|W$.
Ceci montre que les applications de restriction définissent un \emph{préfaisceau} d'ensembles $U\mapsto\operatorname{Ps.hom}(U,Y)$ ; on notera que ce préfaisceau \emph{n'est pas} un faisceau en général (cf. \sref{IV.20.2.16}) ; le faisceau associé se note $\mathscr{P}\!s.\!\operatorname{hom}(X,Y)$.
Pour les pseudo-$S$-morphismes, on voit de même que $U\mapsto\operatorname{Ps.hom}_S(U,Y)$ est un préfaisceau d'ensembles, dont on note $\mathscr{P}\!s.\!\operatorname{hom}_S(X,Y)$ le faisceau associé.

Si $V$ est un ouvert \emph{schématiquement dense} dans $X$, alors, pour tout ouvert $U$ schématiquement dense dans $X$, $U\cap V$ est aussi schématiquement dense dans $X$, donc l'application $\omega\mapsto\omega|V$ est une \emph{bijection} de l'ensemble $\operatorname{Ps.hom}(X,Y)$ (resp. $\operatorname{Ps.hom}_S(X,Y)$) sur $\operatorname{Ps.hom}(V,Y)$ (resp. $\operatorname{Ps.hom}_S(V,Y)$).
\end{env}

\begin{env}[20.2.3]
\EGAReaderTarget{ega4.c04.statement.20-2-3}{IV.20.2.3}
Étant donné un pseudo-$S$-morphisme $\omega$ de $X$ dans $Y$, on dit qu'il est \emph{défini} en un point $x\in X$ s'il existe un voisinage ouvert $V$ de $x$ dans $X$, un ouvert $U\subset V$ \emph{contenant $x$ et schématiquement dense dans $V$}, et enfin un $S$-morphisme $u:U\to Y$ dont la classe dans $\operatorname{Ps.hom}_S(V,Y)$ soit égale à $\omega|V$ \sref{IV.20.2.2} ; on appelle \emph{domaine de définition} de $\omega$ et on note $\operatorname{dom}_S(\omega)$ (ou simplement $\operatorname{dom}(\omega)$ si $S=\Spec(\bb{Z})$) l'ensemble des $x\in X$ où $\omega$ est défini ; c'est évidemment un \emph{ouvert} de $X$.
En outre, pour tout ouvert $W$ de $X$, on a
\[
\label{IV.20.2.3.1}
  \operatorname{dom}_S(\omega|W)=(\operatorname{dom}_S(\omega))\cap W
  \tag{20.2.3.1}
\]
en vertu de la propriété des ouverts schématiquement denses rappelée dans \sref{IV.20.2.2}.
\end{env}

\begin{proposition}[20.2.4]
\label{IV.20.2.4}
Supposons que $X$, $Y$ soient des $S$-préschémas tels que le morphisme structural $Y\to S$ soit séparé ; alors, si $\omega$ est un pseudo-$S$-morphisme de $X$ dans $Y$, $\operatorname{dom}_S(\omega)$ est le plus grand des ouverts schématiquement denses $U$ de $X$ tels qu'il y ait un $S$-morphisme $u:U\to Y$ appartenant à la classe $\omega$.
\end{proposition}

\begin{proof}
Il suffit évidemment de prouver que si $U$, $U'$ sont deux ouverts schématiquement denses dans $X$ tels que deux $S$-morphismes $u:U\to Y$ et $u':U'\to Y$ soient équivalents, alors les restrictions de $u$ et $u'$ à $U\cap U'$ sont égales.
Or il existe par hypothèse un ouvert $U''\subset U\cap U'$, schématiquement dense dans $X$ et dans lequel $u$ et $u'$ coïncident ; comme $U''$ est aussi schématiquement dense dans $U\cap U'$, notre assertion résulte de \sref{IV.11.10.1}[(d)].
\end{proof}

\begin{corollary}[20.2.5]
\label{IV.20.2.5}
Soient $S$ un $S_0$-schéma, $X$, $Y$ deux $S$-préschémas tels que le morphisme composé $Y\to S\to S_0$ soit séparé (ce qui implique \sref[I]{I.5.5.1} que $Y\to S$ l'est aussi).
Pour tout pseudo-$S$-morphisme $\omega$ de $X$ dans $Y$, on a alors $\operatorname{dom}_S(\omega)=\operatorname{dom}_{S_0}(\omega)$.
En particulier, si $Y$ est un schéma, on a $\operatorname{dom}_S(\omega)=\operatorname{dom}(\omega)$.
\end{corollary}

\begin{proof}
En effet, il est clair que $\operatorname{dom}_S(\omega)\subset\operatorname{dom}_{S_0}(\omega)$ sans hypothèse de séparation ; en vertu de \sref{IV.20.2.4}, il suffit de prouver que si un $S_0$-morphisme $u_0:U_0\to Y$ d'un ouvert
\oldpage[IV-4]{233}
schématiquement dense $U_0$ de $X$ dans $Y$ est tel que le composé $v_0:U_0\xrightarrow{u_0}Y\to S$ coïncide avec la restriction du morphisme structural $w_0:U_0\to S$ dans un ouvert $U\subset U_0$ schématiquement dense dans $X$, alors on a $v_0=w_0$.
Mais en vertu de l'hypothèse que le morphisme $S\to S_0$ est séparé, cela résulte encore de \sref{IV.11.10.1}[(d)].
\end{proof}

\begin{corollary}[20.2.6]
\label{IV.20.2.6}
Sous les hypothèses de \sref{IV.20.2.4}, le préfaisceau $U\mapsto\operatorname{Ps.hom}_S(U,Y)$ est un faisceau.
\end{corollary}

\begin{proof}
En effet, soient $U$ un ouvert de $X$, $(U_\alpha)$ un recouvrement de $U$ par des ouverts de $U$ ; supposons donné pour chaque $\alpha$ un pseudo-$S$-morphisme $\omega_\alpha$ de $U_\alpha$ dans $Y$, de sorte que pour tout couple d'indices $\alpha$, $\beta$, les restrictions \sref{IV.20.2.2} des pseudo-$S$-morphismes $\omega_\alpha$ et $\omega_\beta$ à $U_\alpha\cap U_\beta$ soient égales ; en vertu de \sref{IV.20.2.3.1}, cela entraîne que $\operatorname{dom}_S(\omega_\alpha)\cap U_\beta=\operatorname{dom}_S(\omega_\beta)\cap U_\alpha$.
Soit alors $W$ la réunion des ouverts $\operatorname{dom}_S(\omega_\alpha)$, et, pour tout $\alpha$, soit $u_\alpha$ le $S$-morphisme unique $\operatorname{dom}_S(\omega_\alpha)\to Y$ appartenant à la classe $\omega_\alpha$ \sref{IV.20.2.4} ; en raison de l'hypothèse et de \sref{IV.20.2.4}, les restrictions de $u_\alpha$ et $u_\beta$ à $\operatorname{dom}_S(\omega_\alpha)\cap\operatorname{dom}_S(\omega_\beta)$ sont égales, donc il existe un morphisme $u:W\to Y$ dont la restriction à chaque ouvert $\operatorname{dom}_S(\omega_\alpha)$ est égale à $u_\alpha$ ; il est clair que $W$ est schématiquement dense dans $U$, donc définit un pseudo-$S$-morphisme $\omega$ de $U$ dans $Y$ dont les restrictions aux $U_\alpha$ sont les $\omega_\alpha$.
\end{proof}

\begin{remark}[20.2.7]
\label{IV.20.2.7}
On sait \sref{IV.11.10.4} que lorsque $X$ est \emph{réduit}, il revient au même de dire qu'un ouvert de $X$ est dense dans $X$ ou qu'il est schématiquement dense dans $X$ ; la notion de pseudo-morphisme (resp. pseudo-$S$-morphisme) de $X$ dans $Y$ coïncide alors avec celle d'\emph{application rationnelle} (resp. $S$-\emph{application rationnelle}) de $X$ dans $Y$ \sref[I]{I.7.1.2}.
Dans le cas général, la notion de pseudo-morphisme semble plus utile que celle d'application rationnelle.
\end{remark}

\begin{env}[20.2.8]
\EGAReaderTarget{ega4.c04.statement.20-2-8}{IV.20.2.8}
On appelle \emph{pseudo-fonction} sur $X$ un pseudo-morphisme de $X$ dans $\Spec(\bb{Z}[T])$ ($T$ indéterminée), ou, ce qui revient au même, un $X$-pseudo-morphisme de $X$ dans $X\otimes_{\bb{Z}}\bb{Z}[T]$ ; il revient aussi au même \sref[I]{I.3.3.15} de dire qu'une pseudo-fonction sur $X$ est une \emph{classe d'équivalence de sections de $\sh{O}_X$ au-dessus d'ouverts schématiquement denses de $X$}, deux sections $g\in\Gamma(U,\sh{O}_X)$, $g'\in\Gamma(U',\sh{O}_X)$ au-dessus de tels ouverts étant équivalentes s'il existe un ouvert $U''\subset U\cap U'$, schématiquement dense dans $X$, où $g$ et $g'$ coïncident.
On peut ici appliquer \sref{IV.20.2.4} avec $S=X$ et $Y=X\otimes_{\bb{Z}}\bb{Z}[T]$ ; donc, pour toute pseudo-fonction $\vphi$ sur $X$, il existe un \emph{plus grand} ouvert schématiquement dense $\operatorname{dom}(\vphi)$ dans $X$ tel qu'il y ait une section de $\sh{O}_X$ au-dessus de $\operatorname{dom}(\vphi)$ appartenant à la classe $\vphi$.
Il est clair que le faisceau $\mathscr{P}\!s.\!\operatorname{hom}_X(X,X\otimes_{\bb{Z}}\bb{Z}[T])$ est ici un \emph{faisceau d'anneaux}, et même une $\sh{O}_X$-\emph{Algèbre}, que nous noterons $\sh{M}'_X$ ; il résulte de \sref{IV.20.2.6} que, pour tout ouvert $U$ de $X$, $\Gamma(U,\sh{M}'_X)$ est égal à l'ensemble des pseudo-morphismes de $U$ dans $\Spec(\bb{Z}[T])$ ; on pose $M'(X)=\Gamma(X,\sh{M}'_X)$.
Dire qu'une section $\vphi$ de $\sh{M}'_X$ au-dessus de $U$ est \emph{inversible} dans l'anneau $\Gamma(U,\sh{M}'_X)$ signifie qu'il existe un ouvert $U'$ schématiquement dense dans $\operatorname{dom}(\vphi)$, donc dans $U$, tel que, si $g$ est l'unique section de $\sh{O}_X$ au-dessus de $\operatorname{dom}(\vphi)$ appartenant à $\vphi$, $g|U'$ soit \emph{inversible} dans $\Gamma(U',\sh{O}_X)$.
Il résulte de \sref[I]{I.3.3.15} que, dans la correspondance canonique entre $\Gamma(V,\sh{O}_X)$ et $\Hom(V,\bb{Z}[T])$ ($V$ ouvert de $X$), les éléments inversibles de $\Gamma(V,\sh{O}_X)$ correspondent aux morphismes qui se
\oldpage[IV-4]{234}
\emph{factorisent} en $V\to\Spec(\bb{Z}[T,T^{-1}])\to\Spec(\bb{Z}[T])$.
On en conclut que le faisceau ${\sh{M}'}^*_X$ des germes de sections inversibles de $\sh{M}'_X$ s'identifie canoniquement au faisceau $\mathscr{P}\!s.\!\operatorname{hom}_X(X,X\otimes_{\bb{Z}}\bb{Z}[T,T^{-1}])$.
\end{env}

\begin{lemma}[20.2.9]
\label{IV.20.2.9}
Soient $U$ un ouvert de $X$, $s$ un élément régulier de l'anneau $\Gamma(U,\sh{O}_X)$ ; alors l'ensemble ouvert $U_s$ des $x\in U$ tels que $s(x)\neq0$ est schématiquement dense dans $U$.
\end{lemma}

\begin{proof}
En effet, soient $V$ un ouvert de $U$, $t$ une section de $\sh{O}_X$ au-dessus de $V$ telle que $t|(V\cap U_s)=0$.
Pour tout ouvert affine $W\subset V$, il existe alors un entier $n$ tel que $s^nt|W=0$ \sref[I]{I.1.4.1} ; mais puisque $s$ est un élément \emph{régulier} de $\Gamma(U,\sh{O}_X)$, cela entraîne $t|W=0$, d'où $t=0$.
\end{proof}

\begin{env}[20.2.10]
\label{IV.20.2.10}
Considérons une fonction méromorphe $f\in M(X)$ \sref{IV.20.1.4} ; alors $\operatorname{dom}(f)$ est \emph{schématiquement dense} dans $X$ : en effet, tout point de $X$ admet un voisinage ouvert $U$ tel qu'il existe une section $s\in\Gamma(U,\sh{O}_X)$ qui est un élément régulier de cet anneau, et pour laquelle $s(f|U)\in\Gamma(U,\sh{O}_X)$ ; comme $s|U_s$ est inversible, on en conclut que $f|U_s\in\Gamma(U_s,\sh{O}_X)$, donc $U_s\subset\operatorname{dom}(f)$ par définition \sref{IV.20.1.4}, et notre assertion résulte donc du lemme \sref{IV.20.2.9}.
On peut donc associer à $f$ la pseudo-fonction égale à la classe de la section de $\sh{O}_X$ au-dessus de $\operatorname{dom}(f)$, restriction de $f$ ; opérant de même en remplaçant $X$ par un ouvert quelconque de $X$, on obtient ainsi un \emph{homomorphisme canonique de $\sh{O}_X$-Algèbres}
\[
\label{IV.20.2.10.1}
  \sh{M}_X\to\sh{M}'_X
  \tag{20.2.10.1}
\]
qui, par restriction, donne évidemment un homomorphisme de faisceaux de groupes abéliens
\[
\label{IV.20.2.10.2}
  \sh{M}_X^*\to{\sh{M}'}^*_X
  \tag{20.2.10.2}
\]
pour les faisceaux de germes inversibles de sections de $\sh{M}_X$ et $\sh{M}'_X$.
\end{env}

\begin{proposition}[20.2.11]
\EGAReaderTarget{ega4.c04.prop.20-2-11}{IV.20.2.11}
\medskip
\begin{enumerate}[label=\emph{(\roman*)}]
  \item L'homomorphisme canonique \emph{\sref{IV.20.2.10.1}} (et par suite aussi \emph{\sref{IV.20.2.10.2}}) est \emph{injectif}.
  \item Supposons, soit que $X$ soit localement noethérien, soit que $X$ soit réduit et que l'ensemble de ses composantes irréductibles soit localement fini.
  Alors l'homomorphisme canonique \emph{\sref{IV.20.2.10.1}} (et par suite aussi \emph{\sref{IV.20.2.10.2}}) est \emph{bijectif}.
\end{enumerate}
\end{proposition}

\begin{proof}
Les questions envisagées étant locales sur $X$, on peut se borner à considérer le cas où $X=\Spec(A)$ est affine, et à montrer alors que l'homomorphisme canonique $M(X)\to M'(X)$ est toujours injectif, et qu'il est bijectif dans les cas considérés dans~(ii).
Compte tenu de la définition du faisceau $\sh{M}_X$ \sref{IV.20.1.3}, on peut d'ailleurs remarquer que \sref{IV.20.2.10.1} provient en fait d'un homomorphisme de \emph{préfaisceaux}
\[
  \Gamma(U,\sh{O}_X)[\Gamma(U,\sh{S})^{-1}]\to\Gamma(U,\sh{M}'_X)
\]
et il suffit de montrer que, pour $U$ affine, ce dernier est injectif (resp. bijectif sous les hypothèses de~(ii)).
Notant $S$ l'ensemble des éléments réguliers de $A$ (de sorte que $S^{-1}A$ est \emph{l'anneau total des fractions} de $A$), on doit donc considérer l'homomorphisme canonique
\[
\label{IV.20.2.11.1}
  S^{-1}A\to M'(X)
  \tag{20.2.11.1}
\]
\oldpage[IV-4]{235}
et montrer qu'il est injectif (resp. bijectif sous les conditions de~(ii)).
Or, on peut écrire $S^{-1}A=\varinjlim A_t$, où $t$ parcourt l'ensemble des éléments \emph{réguliers} de $A$, ordonné suivant la relation « $t$ est un diviseur de $t'$ » \sref[0]{0.1.4.5} ; on peut aussi écrire $A_t=\Gamma(\mathrm{D}(t),\sh{O}_X)$.
D'autre part on a par définition $M'(X)=\varinjlim\Gamma(U,\sh{O}_X)$, où $U$ parcourt l'ensemble des ouverts schématiquement denses de $X$ (ordonné par la relation $\supset$), et l'application \sref{IV.20.2.11.1} n'est autre que l'application canonique provenant du fait que les $\mathrm{D}(t)$ constituent une partie de l'ensemble des ouverts schématiquement denses dans $X$ \sref{IV.20.2.9}.
Notons que par définition d'un ouvert schématiquement dense, l'homomorphisme $\Gamma(U,\sh{O}_X)\to\Gamma(U',\sh{O}_X)$ pour deux tels ouverts $U'\subset U$ est toujours \emph{injectif}, donc il en est de même des homomorphismes $\Gamma(U,\sh{O}_X)\to M'(X)$, et on en conclut que \sref{IV.20.2.11.1} est injectif.
Pour prouver que \sref{IV.20.2.11.1} est bijectif, il suffit de montrer que l'ensemble des $\mathrm{D}(t)$ est \emph{cofinal} à l'ensemble des ouverts schématiquement denses, autrement dit que pour un tel ouvert $U$, il existe $t$ régulier dans $A$ tel que $\mathrm{D}(t)\subset U$.
Or, lorsque $X$ est réduit et que les composantes irréductibles de $X$ forment un ensemble localement fini, cet ensemble est fini puisque $X$ a été supposé affine, autrement dit $A$ n'a qu'un nombre fini d'idéaux premiers minimaux $\mathfrak{p}_i$ ; comme $A$ est réduit, l'intersection des $\mathfrak{p}_i$ est réduite à $0$, et dire que $t$ est régulier équivaut donc à dire que $t$ n'appartient à aucun des $\mathfrak{p}_i$ ; on conclut donc par le raisonnement de \sref[I]{I.7.1.9.1}.
Lorsque $A$ est noethérien, dire que $U=X-Y$ (où $Y=V(\mathfrak{J})$ est fermé dans $X$) est schématiquement dense signifie \sref{IV.5.10.2} que $Y$ \emph{ne rencontre pas} $\operatorname{Ass}(\sh{O}_X)$, et en vertu de Bourbaki, \emph{Alg.~comm.}, chap.~IV, \textsection1, n\textsuperscript{o}~4, prop.~8, cela entraîne l'existence d'un $t\in\mathfrak{J}$ tel que $t$ soit $A$-régulier, donc $U\supset\mathrm{D}(t)$.

On a en outre prouvé au cours de cette démonstration le
\end{proof}

\begin{corollary}[20.2.12]
\label{IV.20.2.12}
Si $X=\Spec(A)$, où $A$ est noethérien, ou réduit et n'ayant qu'un nombre fini d'idéaux premiers minimaux, les ouverts schématiquement denses dans $X$ sont ceux qui contiennent un ouvert de la forme $\mathrm{D}(t)$, où $t$ est régulier dans $A$, et $M(X)=M'(X)$ est l'anneau total des fractions $S^{-1}A$, où $S$ est l'ensemble des éléments réguliers de $X$.
\end{corollary}

\begin{remarks}[20.2.13]
\EGAReaderTarget{ega4.c04.statement.20-2-13}{IV.20.2.13}
\begin{enumerate}[label=(\roman*)]
  \item Soient $\vphi$ un élément de $M(X)$, $\vphi'$ son image dans $M'(X)$ ; on a évidemment, par définition (\sref{IV.20.1.4} et \sref{IV.20.2.8}) $\operatorname{dom}(\vphi)\subset\operatorname{dom}(\vphi')$ ; mais en fait on a l'égalité $\operatorname{dom}(\vphi)=\operatorname{dom}(\vphi')$, car il existe une section de $\sh{O}_X$ au-dessus de $\operatorname{dom}(\vphi')$, appartenant à la classe $\vphi'$, et la fonction méromorphe correspondante est égale à $\vphi$ \sref{IV.20.2.11}[(i)], donc $\operatorname{dom}(\vphi')\subset\operatorname{dom}(\vphi)$.

  \item On a déjà noté que lorsque $X$ est \emph{réduit}, on a $\sh{M}'_X=\sh{R}(X)$ (faisceau des fonctions rationnelles sur $X$ \sref[I]{I.7.3.2}) ; si en outre les composantes irréductibles de $X$ forment un ensemble localement fini, on a $\sh{M}_X=\sh{M}'_X=\sh{R}(X)$.
  En général, puisque tout ensemble ouvert schématiquement dense est dense, on a un homomorphisme canonique $\sh{M}'_X\to\sh{R}(X)$, mais même lorsque $X$ est localement noethérien, cet homomorphisme n'est pas nécessairement injectif.
  Par exemple, si $X=\Spec(A)$, où $A$ est un anneau noethérien tel que $\operatorname{Ass}(A)$ contienne des idéaux premiers immergés (ce qui entraîne que $A$ n'est pas réduit), on a vu que $M(X)$ et $M'(X)$ s'identifient à l'anneau total de fractions $S^{-1}A$,
  \oldpage[IV-4]{236}
  où $S$ est l'ensemble des éléments réguliers de $A$, complémentaire de la réunion des idéaux $\mathfrak{p}\in\operatorname{Ass}(A)$ ; par contre, $R(X)$ s'identifie à $Q^{-1}A$, où $Q$ est le complémentaire de la réunion des idéaux premiers \emph{minimaux} de $A$ \sref[I]{I.7.1.9}, et l'homomorphisme canonique $A\to Q^{-1}A$ (et \emph{a fortiori} $S^{-1}A\to Q^{-1}A$) n'est donc pas injectif, puisqu'il existe dans $A-Q$ des éléments $\neq0$ de $A$ annulés par des éléments de $Q$ (Bourbaki, \emph{Alg.~comm.}, chap.~IV, \textsection1, n\textsuperscript{o}~1, cor.~2 de la prop.~1).

  \item On notera que même lorsque $X$ est localement noethérien, le $\sh{O}_X$-Module $\sh{M}'_X$ \emph{n'est pas nécessairement quasi-cohérent}.
  Considérons par exemple un anneau local noethérien $A$ de dimension $\geq2$, dont l'idéal maximal $\mathfrak{m}$ est tel que $\mathfrak{m}\in\operatorname{Ass}(A)$ et que si l'on pose $X=\Spec(A)$, le schéma induit sur l'ouvert $U$ de $X$, complémentaire de $\{\mathfrak{m}\}$, soit \emph{intègre}.
  Les seuls éléments réguliers de $A$ sont alors les éléments inversibles, donc $\Gamma(X,\sh{M}'_X)=M'(X)=A$ ; si $\sh{M}'_X$ était quasi-cohérent, il serait donc égal à $\sh{O}_X$ ; mais comme $U$ est un schéma intègre, $\sh{M}'_X|U=R(U)$ est un faisceau \emph{simple} \sref[I]{I.7.3.5}, alors que $\sh{O}_U$ n'est pas un faisceau simple puisque $\dim(A)\geq2$.

  Reste à donner un exemple d'anneau $A$ ayant les propriétés précédentes.
  Il suffit de considérer un anneau local intègre $B$ de dimension $\geq2$ et de corps résiduel $k$, et de prendre $A=B\oplus k$ avec la loi de multiplication $(b,x)(b',x')=(bb',bx'+b'x)$.

  \item Si $X$ est localement noethérien, il résulte de \sref{IV.5.10.2} que les ouverts schématiquement denses dans $X$ sont ceux qui \emph{contiennent} l'ensemble $\operatorname{Ass}(\sh{O}_X)$.
\end{enumerate}
\end{remarks}

\begin{env}[20.2.14]
\EGAReaderTarget{ega4.c04.statement.20-2-14}{IV.20.2.14}
Soient $X$ un préschéma, $\sh{F}$ un $\sh{O}_X$-Module quasi-cohérent et \emph{strictement sans torsion} \sref{IV.20.1.5}, de sorte que $\sh{F}$ s'identifie à un sous-$\sh{O}_X$-Module de $\sh{M}_X(\sh{F})$.
Pour toute section méromorphe $u$ de $\sh{F}$ au-dessus de $X$, on appelle l'\emph{Idéal des dénominateurs de $u$} l'annulateur $\sh{J}$ de la section $\overline{u}$ image de $u$ dans $\sh{M}_X(\sh{F})/\sh{F}$.
L'Idéal $\sh{J}$ est \emph{quasi-cohérent} : en effet, la question étant locale sur $X$, on peut se borner au cas où $X$ est affine, et il existe une section $s\in\Gamma(X,\sh{S}(\sh{O}_X))$ telle que $v=su\in\Gamma(X,\sh{F})$.
Dire que, pour un ouvert $U\subset X$, une section $f\in\Gamma(U,\sh{O}_X)$ appartient à $\Gamma(U,\sh{J})$ signifie que $f(u|U)\in\Gamma(U,\sh{F})$, et puisque $s|U$ est un élément régulier de $\Gamma(U,\sh{O}_X)$ et que $\sh{F}$ est strictement sans torsion, la relation précédente équivaut encore à $f((su)|U)\in\Gamma(U,s\sh{F})$ ; si $\overline{v}$ est la section de $\sh{F}/s\sh{F}$, image canonique de $v$, on voit donc que $\sh{J}$ est le noyau de l'homomorphisme $\sh{O}_X\to\sh{F}/s\sh{F}$ obtenu par multiplication par la section $\overline{v}$.
Comme $\sh{F}/s\sh{F}$ est quasi-cohérent, il en est de même de $\sh{J}$.

Il résulte aussitôt de la définition précédente que $\operatorname{dom}(u)$ est \emph{l'ouvert complémentaire du sous-préschéma fermé de $X$ défini par l'Idéal des dénominateurs de $u$}.
\end{env}

\begin{proposition}[20.2.15]
\label{IV.20.2.15}
Soient $f:X'\to X$ un morphisme, $\sh{F}$ un $\sh{O}_X$-Module quasi-cohérent, $\vphi$ une section de $\sh{M}_f(\sh{F})$ au-dessus de $X$ \sref{IV.20.1.11}.
Alors $f^{-1}(\operatorname{dom}(\vphi))$ est schématiquement dense dans $X'$.
\end{proposition}

\begin{proof}
La question étant évidemment locale sur $X$ et $X'$, on peut supposer $X=\Spec(A)$, $X'=\Spec(A')$ affines, $\sh{F}=\widetilde{M}$, et $\vphi=h\otimes(1/s)$, où $h\in M$ et $s$ est un élément régulier de $A$ dont l'image $s'$ dans $A'$ est un élément régulier.
On sait alors \sref{IV.20.2.9} que $\mathrm{D}(s')$ est un ouvert schématiquement dense dans $X'$, et c'est l'image réciproque par $f$ de $\mathrm{D}(s)$, qui est contenu dans $\operatorname{dom}(\vphi)$.
\end{proof}

\begin{remark}[20.2.16]
\label{IV.20.2.16}
Lorsque $Y$ n'est pas séparé, le préfaisceau $U\mapsto\operatorname{Ps.hom}_S(U,Y)$ sur $X$ n'est pas nécessairement un faisceau, même lorsque $X$ est noethérien, comme le montre l'exemple suivant.
Nous allons prendre pour $X$ un spectre d'un anneau semi-local noethérien $A$ de dimension $\geq2$, ayant exactement deux idéaux
\oldpage[IV-4]{237}
maximaux $\mathfrak{m}'$, $\mathfrak{m}''$ (donc tel que $X$ ait exactement deux points fermés $x'$, $x''$), de sorte que $\mathfrak{m}'$ et $\mathfrak{m}''$ appartiennent à $\operatorname{Ass}(A)$ et que l'ouvert $U=X-\{x',x''\}$ soit \emph{intègre}.
Soient $U'=X-\{x'\}$, $U''=X-\{x''\}$, de sorte que $U=U'\cap U''$.
Notons que les seuls ouverts schématiquement denses dans $X$ sont ceux qui contiennent $x'$ et $x''$ \sref{IV.20.2.13}[(iv)], donc ils sont nécessairement identiques à $X$.
Pour avoir un contre-exemple, il suffira donc de définir deux $S$-morphismes $u':U'\to Y$, $u'':U''\to Y$ (avec $S=X$) dont les restrictions à $U$ appartiennent au même pseudo-morphisme de $U$ dans $Y$, mais qui soient tels que, pour tout voisinage $V'$ de $x''$ dans $U'$ et tout voisinage $V''$ de $x'$ dans $U''$, les restrictions de $u'$ et $u''$ à $V'\cap V''$ soient distinctes.
Pour cela, considérons une partie fermée irréductible $Z$ de $X$ contenant $x'$ et $x''$, distincte de $X$ ; soit $Y$ le $X$-préschéma obtenu en recollant deux préschémas $Y'$, $Y''$ isomorphes à $X$ le long de l'ouvert partout dense $X-Z$ \sref[I]{I.2.3.1}.
On prendra alors pour $u'$ et $u''$ les restrictions respectives à $U'$ et $U''$ des isomorphismes réciproques des isomorphismes structuraux $Y'\to X$, $Y''\to X$.
Comme $V'$ et $V''$ contiennent le point générique de $Z$, les restrictions de $u'$ et $u''$ à $V'\cap V''$ sont distinctes, mais les restrictions de $u'$ et $u''$ à $U-(U\cap Z)$ sont égales, et $U-(U\cap Z)$ est un ouvert dense dans $U$, donc schématiquement dense puisque $U$ est réduit.

Reste donc seulement à définir $A$ et $Z$ de façon à avoir les propriétés précédentes.
Soient $X_0=\Spec(B)$ un schéma affine intègre (par exemple le plan affine sur un corps $k$), $Y$ une partie fermée irréductible de $X_0$ (par exemple une droite affine) contenant deux points fermés distincts $x'$ et $x''$, correspondant aux idéaux maximaux $\mathfrak{n}'$, $\mathfrak{n}''$ de $B$.
On considère l'anneau $C=B\oplus(B/\mathfrak{n}'\oplus B/\mathfrak{n}'')$ avec la multiplication $(b,z)(b',z')=(bb',bz'+b'z)$.
Si $X_1=\Spec(C)$, on a $X_0=(X_1)_\mathrm{red}$ et $X_1$ est réduit sauf aux points $x'$, $x''$.
Il suffit alors de prendre $A=R^{-1}C$, où $R$ est le complémentaire de la réunion des idéaux maximaux de $C$ aux points $x'$, $x''$, et pour $Z$ la trace de $Y$ sur $X=\Spec(A)$.
\end{remark}

\subsection{Composition des pseudo-morphismes}
\EGAReaderTarget{ega4.c04.subsec.20-3}{IV.20.3}

\begin{env}[20.3.1]
\EGAReaderTarget{ega4.c04.statement.20-3-1}{IV.20.3.1}
Soient $X$, $Y$, $Z$ trois préschémas, $\omega$ un pseudo-morphisme de $X$ dans $Y$, $f:Y\to Z$ un morphisme.
Il est clair que si $U'$, $U''$ sont deux ouverts schématiquement denses dans $X$, $u':U'\to Y$, $u'':U''\to Y$ deux morphismes de la classe $\omega$, les morphismes $f\circ u'$ et $f\circ u''$ sont équivalents (pour la relation définie dans \sref{IV.20.2.1}), donc, pour tous les morphismes $u$ de la classe $\omega$, les morphismes $f\circ u$ appartiennent à une même classe d'équivalence, que l'on note $f\circ\omega$ et que l'on appelle le pseudo-morphisme de $X$ dans $Z$ \emph{composé de $f$ et de $\omega$}.
On a $\operatorname{dom}(f\circ\omega)=\operatorname{dom}(\omega)$.
Si $g:Z\to T$ est un morphisme, il est clair que l'on a $g\circ(f\circ\omega)=(g\circ f)\circ\omega$.
\end{env}

\begin{env}[20.3.2]
\EGAReaderTarget{ega4.c04.statement.20-3-2}{IV.20.3.2}
Supposons maintenant donnés un pseudo-$S$-morphisme $\omega$ de $X$ dans $Y$, où $Y$ est \emph{séparé sur $S$}, de sorte qu'il existe un $S$-morphisme $u:\operatorname{dom}_S(\omega)\to Y$ de la classe $\omega$ \sref{IV.20.2.4}.
Soit $f:X'\to X$ un $S$-morphisme tel que l'ouvert $V'=f^{-1}(\operatorname{dom}_S(\omega))$ soit schématiquement dense dans $X'$ ; alors on dit que la classe (pour la relation d'équivalence de \sref{IV.20.2.1}) du $S$-morphisme $u\circ(f|U')$ (classe qui est définie en vertu de l'hypothèse précédente) est le pseudo-$S$-morphisme \emph{composé de $\omega$ et de $f$}, et on le note $\omega\circ f$ ; il est clair que l'on a
\[
\label{IV.20.3.2.1}
  f^{-1}(\operatorname{dom}_S(\omega))\subset\operatorname{dom}_S(\omega\circ f)
  \tag{20.3.2.1}
\]

Pour $f:X'\to X$ donné, on note $\operatorname{Ps.hom}_S(X,Y)^f$ l'ensemble des pseudo-$S$-morphismes $\omega$ de $X$ dans $Y$ qui vérifient la condition précédente.
Si $\omega$ est un tel pseudo-$S$-morphisme, il est clair que pour tout ouvert $V$ de $X$,
\[
  f^{-1}(\operatorname{dom}_S(\omega|V))=f^{-1}(V\cap\operatorname{dom}_S(\omega))=f^{-1}(V)\cap f^{-1}(\operatorname{dom}_S(\omega))
\]
est schématiquement dense dans $f^{-1}(V)$, donc, si $f^V:f^{-1}(V)\to V$ est la restriction de $f$, le composé $(\omega|V)\circ f^V$ est défini et égal à $(\omega\circ f)|f^{-1}(V)$.
On définit ainsi une
\oldpage[IV-4]{238}
application canonique de restriction $\operatorname{Ps.hom}_S(X,Y)^f\to\operatorname{Ps.hom}_S(V,Y)^{f^V}$, ce qui permet de définir un \emph{préfaisceau} d'ensembles $V\mapsto\operatorname{Ps.hom}_S(V,Y)^{f^V}$ sur $X$, sous-préfaisceau de $V\mapsto\operatorname{Ps.hom}_S(V,Y)$, d'où un faisceau d'ensembles associé que l'on note $\mathscr{P}\!s.\!\operatorname{hom}_S(X,Y)^f$.
En outre, pour tout ouvert $V$ de $X$, on a une application $\omega\mapsto\omega\circ f^V$ de $\operatorname{Ps.hom}_S(V,Y)^{f^V}$ dans $\operatorname{Ps.hom}_S(f^{-1}(V),Y)$, qui définit donc un $f$-morphisme de faisceaux d'ensembles
\[
  \mathscr{P}\!s.\!\operatorname{hom}_S(X,Y)^f\to\mathscr{P}\!s.\!\operatorname{hom}_S(X',Y).
\]
\end{env}

\begin{env}[20.3.3]
\label{IV.20.3.3}
Soit maintenant $f':X''\to X'$ un $S$-morphisme tel que l'ouvert $f'^{-1}(f^{-1}(\operatorname{dom}_S(\omega)))$ soit schématiquement dense dans $X''$ ; alors $\omega\circ(f\circ f')$ est défini et $u\circ f\circ f'$ appartient à ce pseudo-$S$-morphisme ; en outre, en vertu de \sref{IV.20.3.2.1}, $f'^{-1}(\operatorname{dom}_S(\omega\circ f))$ est \emph{a fortiori} schématiquement dense dans $X''$, donc $(\omega\circ f)\circ f'$ est aussi défini et $u\circ f\circ f'$ appartient à ce pseudo-$S$-morphisme, donc on a $(\omega\circ f)\circ f'=\omega\circ(f\circ f')$.
D'autre part, pour tout $S$-morphisme $g:Y\to Z$, on a $\operatorname{dom}_S(g\circ\omega)=\operatorname{dom}_S(\omega)$ \sref{IV.20.3.1}, donc $(g\circ\omega)\circ f$ est défini, et $g\circ u\circ f$ appartient à ce pseudo-$S$-morphisme, ce qui montre que $(g\circ\omega)\circ f=g\circ(\omega\circ f)$.
\end{env}

\begin{env}[20.3.4]
\label{IV.20.3.4}
Le cas le plus important dans la définition \sref{IV.20.3.2} est celui où l'on a $\mathscr{P}\!s.\!\operatorname{hom}_S(X,Y)^f=\mathscr{P}\!s.\!\operatorname{hom}_S(X,Y)$ : il suffit pour cela que pour tout ouvert $U$ de $X$, et tout ouvert $V$ schématiquement dense dans $U$, $f^{-1}(V)$ soit schématiquement dense dans $f^{-1}(U)$ ; lorsqu'il en est ainsi, alors, pour \emph{tout} ouvert $U$ de $X$ et \emph{tout} pseudo-$S$-morphisme $\omega:U\to Y$, on peut définir le composé $\omega\circ f^U$, \emph{même si $Y$ n'est pas séparé sur $S$}.
En effet, si $u':U'\to Y$ et $u'':U''\to Y$ sont deux morphismes de la classe $\omega$, ils coïncident dans un ouvert $U_0$ schématiquement dense dans $U$, donc les morphismes composés $f^{-1}(U')\to U'\xrightarrow{u'}Y$ et $f^{-1}(U'')\to U''\xrightarrow{u''}Y$ coïncident dans $f^{-1}(U_0)$, qui est par hypothèse schématiquement dense dans $f^{-1}(U)$ ; on peut donc définir $\omega\circ f^U$ comme la classe d'un quelconque des morphismes $f^{-1}(U')\to U'\xrightarrow{u'}Y$, où $u'$ appartient à $\omega$.
\end{env}

\begin{proposition}[20.3.5]
\EGAReaderTarget{ega4.c04.prop.20-3-5}{IV.20.3.5}
Soient $X$, $X'$ deux préschémas, $f:X'\to X$ un morphisme.
Alors, dans chacun des trois cas suivants, pour tout ouvert $U$ de $X$ et tout ouvert $V$ schématiquement dense dans $U$, $f^{-1}(V)$ est schématiquement dense dans $f^{-1}(U)$ :
\begin{enumerate}[label=\emph{(\roman*)}]
  \item $f$ est plat et localement de présentation finie.
  \item $f$ est plat et l'espace sous-jacent à $X$ est localement noethérien.
  \item $X'$ est réduit, l'ensemble des composantes irréductibles de $X$ est localement fini, et toute composante irréductible de $X'$ domine une composante irréductible de $X$.
\end{enumerate}
\end{proposition}

\begin{proof}
L'assertion~(i) résulte de \sref{IV.11.10.5}[(ii), (b)] ; l'assertion~(ii) résulte de \sref{IV.11.10.5}[(ii), (a)], car alors tout ouvert $V$ dans $U$ est rétrocompact, autrement dit l'injection canonique $j:V\to U$ est un morphisme quasi-compact.
Enfin, pour prouver~(iii), notons que puisque $X'$ est réduit, il suffit de montrer que pour tout ouvert $U$ de $X$ et tout ouvert $V$ dense dans $U$, $f^{-1}(V)$ est \emph{dense} dans $f^{-1}(U)$.
Or, pour que $f^{-1}(V)$ soit dense dans $f^{-1}(U)$, il suffit que $f^{-1}(V)$ contienne tous les points maximaux de $X'$ appartenant à $f^{-1}(U)$ ; la conclusion résulte donc de l'hypothèse que l'image par $f$ de tout point
\oldpage[IV-4]{239}
maximal de $X'$ appartenant à $f^{-1}(U)$ est un point maximal de $X$ appartenant à $U$, donc à $V$ puisque l'ensemble des composantes irréductibles de $X$ est localement fini.
\end{proof}

\begin{env}[20.3.6]
\EGAReaderTarget{ega4.c04.statement.20-3-6}{IV.20.3.6}
Soient $X$, $Y$ deux $S$-préschémas ; supposons que $X$ vérifie l'une des deux hypothèses suivantes :
\begin{enumerate}[label=\emph{\alph*)}]
  \item $X$ est localement noethérien ;
  \item $X$ est réduit et l'ensemble de ses composantes irréductibles est localement fini.
\end{enumerate}
Alors, pour tout $x\in X$, le $S$-morphisme canonique $j:\Spec(\sh{O}_{X,x})\to X$ est plat et vérifie la condition~(ii) de \sref{IV.20.3.5} dans le cas~a), la condition~(iii) de \sref{IV.20.3.5} dans le cas~b).
Pour tout pseudo-$S$-morphisme $\omega$ de $X$ dans $Y$, le composé $\omega\circ j$ est donc défini et est un pseudo-$S$-morphisme de $\Spec(\sh{O}_{X,x})$ dans $Y$, appelé la \emph{restriction} de $\omega$ à $\Spec(\sh{O}_{X,x})$.
Notons maintenant que si $X$ vérifie la condition~a) (resp. b)) de \sref{IV.20.3.6}, il en est de même de tout préschéma induit sur un ouvert $U$ de $X$ contenant $x$.
Par passage à la limite inductive, on déduit donc, des applications canoniques $\operatorname{Ps.hom}_S(U,Y)\to\operatorname{Ps.hom}_S(\Spec(\sh{O}_{X,x}),Y)$ ainsi obtenues, une application canonique
\[
\label{IV.20.3.6.1}
  (\mathscr{P}\!s.\!\operatorname{hom}_S(X,Y))_x\to\operatorname{Ps.hom}_S(\Spec(\sh{O}_{X,x}),Y)
  \tag{20.3.6.1}
\]
où le premier membre est la fibre au point $x$ du faisceau $\mathscr{P}\!s.\!\operatorname{hom}_S(X,Y)$, ensemble des germes en $x$ de pseudo-$S$-morphismes de voisinages ouverts de $x$ dans $Y$.
\end{env}

\begin{proposition}[20.3.7]
\label{IV.20.3.7}
Sous les hypothèses de \sref{IV.20.3.6}, l'application canonique \sref{IV.20.3.6.1} est injective.
Si en outre $Y$ est localement de présentation finie sur $S$, l'application \sref{IV.20.3.6.1} est bijective.
\end{proposition}

\begin{proof}
Par application de la méthode de \sref{IV.8.1.2}[(a)], cette proposition sera un cas particulier de la suivante :
\end{proof}

\begin{proposition}[20.3.8]
\EGAReaderTarget{ega4.c04.prop.20-3-8}{IV.20.3.8}
Les notations étant celles de \sref{IV.8.8.1}, supposons $X_\alpha$ quasi-compact (resp. quasi-compact et quasi-séparé) et $Y_\alpha$ localement de type fini (resp. localement de présentation finie) sur $S_\alpha$.
Supposons en outre que l'une des conditions suivantes soit vérifiée :
\begin{enumerate}[label=\emph{(\roman*)}]
  \item Les morphismes de transition $S_\mu\to S_\lambda$ (pour $\lambda\leq\mu$) sont plats, et les $X_\lambda$ et $X$ sont noethériens.
  \item Les $X_\lambda$ sont réduits, l'ensemble des composantes irréductibles de $X$ et de chacun des $X_\lambda$ est fini, et, pour $\lambda\leq\mu$, toute composante irréductible de $X_\mu$ domine une composante irréductible de $X_\lambda$.
\end{enumerate}
Alors, l'application canonique
\[
\label{IV.20.3.8.1}
  \varinjlim\operatorname{Ps.hom}_{S_\lambda}(X_\lambda,Y_\lambda)\to\operatorname{Ps.hom}_S(X,Y)
  \tag{20.3.8.1}
\]
est injective (resp. bijective).
\end{proposition}

\begin{proof}
Notons d'abord que, dans le cas~(i), les morphismes $X_\mu\to X_\lambda$ (pour $\lambda\leq\mu$) et $X\to X_\lambda$ sont plats, donc il résulte de \sref{IV.20.3.4} et \sref{IV.20.3.5} que les applications canoniques
\[
  \operatorname{Ps.hom}_{S_\lambda}(X_\lambda,Y_\lambda)\to\operatorname{Ps.hom}_{S_\mu}(X_\mu,Y_\mu)
\]
\oldpage[IV-4]{240}
pour $\lambda\leq\mu$, et $\operatorname{Ps.hom}_{S_\lambda}(X_\lambda,Y_\lambda)\to\operatorname{Ps.hom}_S(X,Y)$ sont définies (sans hypothèse de séparation sur les $Y_\lambda$ ni $Y$) ; il en est par suite de même de l'application \sref{IV.20.3.8.1}.
La proposition va résulter du lemme suivant :

\begin{lemma}[20.3.8.2]
\label{IV.20.3.8.2}
Les notations étant celles de \sref{IV.8.8.1}, supposons $X_\alpha$ quasi-compact, et soit $U_\alpha$ un ouvert dans $X_\alpha$ ; soient $U_\lambda$ et $U$ ses images réciproques dans $X_\lambda$ et $X$ pour $\lambda\geq\alpha$.
Supposons que l'une des conditions~(i),~(ii) de \sref{IV.20.3.8} soit vérifiée.
Alors, pour que $U$ soit schématiquement dense dans $X$, il faut et il suffit qu'il existe $\lambda\geq\alpha$ tel que $U_\lambda$ soit schématiquement dense dans $X_\lambda$, et dans ce cas $U_\mu$ est schématiquement dense dans $X_\mu$ pour $\mu\geq\lambda$.
\end{lemma}

Supposons d'abord vérifiée la condition~(i) ; désignons par $j_\lambda:U_\lambda\to X_\lambda$ et $j:U\to X$ les injections canoniques, par $\sh{J}_\lambda$ le noyau de l'homomorphisme canonique $\sh{O}_{X_\lambda}\to(j_\lambda)_*(\sh{O}_{U_\lambda})$.
L'immersion $j_\lambda$ étant quasi-compacte et quasi-séparée, $(j_\lambda)_*(\sh{O}_{U_\lambda})$ est un $\sh{O}_{X_\lambda}$-Module quasi-cohérent, donc $\sh{J}_\lambda$ est un Idéal quasi-cohérent de $\sh{O}_{X_\lambda}$, et puisque $X_\lambda$ est noethérien, $\sh{J}_\lambda$ est cohérent, donc de type fini.
D'autre part, le morphisme de transition $p_{\lambda\mu}:X_\mu\to X_\lambda$ (resp. $p_\lambda:X\to X_\lambda$) étant plat, il résulte de \sref{IV.2.3.1} que l'on peut identifier $(j_\mu)_*(\sh{O}_{U_\mu})$ à $p_{\lambda\mu}^*((j_\lambda)_*(\sh{O}_{U_\lambda}))$ (resp. $j_*(\sh{O}_U)$ à $p_\lambda^*((j_\lambda)_*(\sh{O}_{U_\lambda}))$).
L'assertion résulte alors de la définition d'un ouvert schématiquement dense et de \sref{IV.8.5.8}[(ii)].

Pour établir \sref{IV.20.3.8.2} lorsque la condition~(ii) est vérifiée, nous prouverons d'abord deux lemmes.

\begin{lemma}[20.3.8.3]
\label{IV.20.3.8.3}
Sous les hypothèses de \sref{IV.8.2.2}, soit $M_\lambda$ (resp. $M$) l'ensemble des points maximaux de $S_\lambda$ (resp. $S$).
Supposons que pour tout $\lambda$, l'ensemble $M_\lambda$ soit fini, et que les $M_\lambda$ forment un système projectif d'ensembles.
Alors $M$ est limite projective du système des $M_\lambda$.
\end{lemma}

Montrons d'abord qu'un point $s\in\varprojlim M_\lambda$ est maximal dans $S$ : en effet, si $s'$ est une généralisation de $s$, l'image $s'_\lambda$ de $s'$ dans $S_\lambda$ est une généralisation de l'image $s_\lambda$ de $s$, donc égale à $s_\lambda$ pour tout $\lambda$, ce qui implique $s'=s$, puisque l'ensemble sous-jacent à $S$ est limite projective des ensembles sous-jacents aux $S_\lambda$ \sref{IV.8.2.9}.
Inversement, soit $s$ un point maximal de $S$ et prouvons qu'il appartient à $\varprojlim M_\lambda$.
Soient $s_\lambda$ l'image de $s$ dans $S_\lambda$, $M'_\lambda$ l'ensemble des points maximaux de $S_\lambda$ qui sont des généralisations de $s_\lambda$ ; les $M'_\lambda$ sont des ensembles finis non vides, qui forment un système projectif, donc $M'=\varprojlim M'_\lambda$ n'est pas vide et l'application $M'\to M'_\lambda$ est surjective (Bourbaki, \emph{Ens.}, chap.~III, 2\textsuperscript{e} éd., \textsection~7, n\textsuperscript{o}~4, Exemple~I).
D'autre part, on a $\Spec(\sh{O}_{S,s})=\varprojlim\Spec(\sh{O}_{S_\lambda,s_\lambda})$ en vertu de \sref{IV.8.2.12} et \sref{IV.8.2.9}, donc les points de $\varprojlim M'_\lambda$ sont aussi points maximaux de $\Spec(\sh{O}_{S,s})$ d'après la première partie du raisonnement.
Donc $M'=\varprojlim M'_\lambda$ se réduit nécessairement au point $s$ ; on en conclut que $M'_\lambda$ se réduit au point $s_\lambda$ et par suite les $s_\lambda$ sont maximaux dans les $S_\lambda$, ce qui achève de prouver le lemme.

\begin{lemma}[20.3.8.4]
\label{IV.20.3.8.4}
Les hypothèses étant celles de \sref{IV.20.3.8.3}, supposons en outre $S_\alpha$ quasi-compact ; soit $V_\alpha$ un ensemble ouvert de $S_\alpha$, et soient $V_\lambda$ et $V$ ses images réciproques dans $S_\lambda$ pour $\lambda\geq\alpha$ et dans $S$.
Si $U_\alpha$ est dense dans $S_\alpha$, $U_\lambda$ est dense dans $S_\lambda$ pour $\lambda\geq\alpha$ et $U$ est dense dans $S$.
Inversement, si $U$ est dense dans $S$ et si de plus l'ensemble $M$ des points maximaux de $S$ est fini, il existe $\lambda\geq\alpha$ tel que $U_\lambda$ soit dense dans $S_\lambda$.
\end{lemma}

En effet, puisque $M_\alpha$ est fini, l'hypothèse que $U_\alpha$ est dense dans $S_\alpha$ entraîne que
\oldpage[IV-4]{241}
$M_\alpha\subset U_\alpha$, donc $M_\lambda\subset U_\lambda$ pour $\lambda\geq\alpha$ et $M\subset U$ en vertu de \sref{IV.20.3.8.3}, ce qui prouve la première assertion.
Inversement, supposons $M$ fini et $U$ dense dans $S$, donc $M\subset U$ ; notons que les $U_\lambda$ sont ouverts, donc ind-constructibles et les $M_\lambda$ finis, donc pro-constructibles \sref{IV.1.9.6}.
La seconde assertion résulte alors de \sref{IV.8.3.2}.

\phantomsection\label{IV.20.3.8.5}
\emph{20.3.8.5. Fin de la démonstration de \sref{IV.20.3.8.2}.}
Supposons la condition~(ii) vérifiée, et notons que $X$ est alors \emph{réduit} en vertu de \sref{IV.8.7.1} ; il revient donc au même de dire que $U_\lambda$ (resp. $U$) est schématiquement dense dans $X_\lambda$ (resp. $X$) ou qu'il est dense dans $X_\lambda$ (resp. $X$), et la conclusion résulte de \sref{IV.20.3.8.4} appliqué aux $X_\lambda$ et à $X$.

\phantomsection\label{IV.20.3.8.6}
\emph{20.3.8.6. Fin de la démonstration de \sref{IV.20.3.8}.}
Pour montrer que l'application \sref{IV.20.3.8.1} est injective, considérons deux morphismes $u'_\lambda$, $u''_\lambda$ d'un ouvert schématiquement dense $U_\lambda\subset X_\lambda$ dans $Y_\lambda$, et supposons que leurs images $u'$, $u''$, morphismes de $U$ dans $Y$, coïncident dans un ouvert schématiquement dense $V$ de $U$.
En outre, dans l'une ou l'autre des hypothèses~(i),~(ii), on peut supposer $V$ \emph{quasi-compact} ; c'est évident si $X$ est noethérien ; sinon, comme $X$ n'a qu'un nombre fini de points maximaux et est réduit, il suffit de remplacer $V$ par la réunion d'un nombre fini de voisinages ouverts affines de ces points maximaux (contenus dans $V$ par hypothèse).
Alors il existe un $\lambda\geq\alpha$ tel que $V$ soit l'image réciproque d'un ouvert quasi-compact $V_\lambda$ de $X_\lambda$ \sref{IV.8.2.11}, et il résulte de \sref{IV.20.3.8.2} qu'en prenant $\lambda$ assez grand, on peut supposer que $V_\lambda$ est schématiquement dense dans $X_\lambda$.
Il résulte alors de \sref{IV.8.8.2}[(i)] qu'en prenant $\lambda$ assez grand, $u'_\lambda$ et $u''_\lambda$ coïncident dans $V_\lambda$, donc appartiennent au même pseudo-homomorphisme.

Prouvons enfin que l'application \sref{IV.20.3.8.1} est surjective.
On considère maintenant un morphisme $u$ d'un ouvert schématiquement dense $U$ de $X$, dans $Y$ ; comme ci-dessus, on peut supposer $U$ quasi-compact et quasi-séparé ($U$ pouvant être remplacé, dans le cas~(ii), par une réunion d'un nombre fini d'ouverts affines deux à deux disjoints).
On peut donc encore supposer qu'il existe un $\lambda\geq\alpha$ tel que $U$ soit l'image réciproque d'un ouvert quasi-compact $U_\lambda$ de $X_\lambda$ \sref{IV.8.2.11} qui est automatiquement quasi-séparé, et par \sref{IV.20.3.8.2} on peut supposer en outre que $U_\lambda$ soit schématiquement dense dans $X_\lambda$.
Puisque les $Y_\lambda$ sont supposés localement de présentation finie, il résulte de \sref{IV.8.8.2}[(i)] qu'il existe un $\lambda$ tel que $u$ soit l'image d'un morphisme $u_\lambda:U_\lambda\to Y_\lambda$ ; d'où la conclusion.

\begin{remarks}[20.3.8.7]
\label{IV.20.3.8.7}
\begin{enumerate}[label=(\roman*)]
  \item Pour prouver que l'application \sref{IV.20.3.8.1} est \emph{injective}, il n'est pas nécessaire, dans l'hypothèse~(i) de \sref{IV.20.3.8}, de supposer $X$ noethérien.
  En effet, le lemme \sref{IV.20.3.8.2} n'utilise pas cette hypothèse.
  Avec les notations de \sref{IV.20.3.8.6}, soit $Z_\lambda$ le sous-préschéma des coïncidences de $u'_\lambda$ et $u''_\lambda$, et soit $Z$ le sous-préschéma des coïncidences de $u'$ et $u''$ ; il résulte de la définition \sref{IV.17.4.5} et de \sref[I]{I.3.3.10.1} que $Z$ est limite projective des $Z_\mu$ pour $\mu\geq\lambda$.
  Or, par hypothèse, $Z$ majore un ouvert schématiquement dense dans $X$ ; il en résulte que $Z$ est lui-même induit sur un ouvert de $X$ en vertu du lemme suivant :

\begin{lemma}[20.3.8.8]
\label{IV.20.3.8.8}
Soit $T$ un préschéma.
Alors tout sous-préschéma $W$ de $T$ qui majore un ouvert schématiquement dense $V$ de $T$ est induit sur un ouvert (schématiquement dense) de $T$.
\end{lemma}

  En effet, le sous-espace de $T$ sous-jacent à $W$ est localement fermé, donc ouvert
  \oldpage[IV-4]{242}
  dans son adhérence, ce qui prouve déjà que l'espace sous-jacent à $W$ est ouvert dans $T$ ; la conclusion résulte alors de \sref{IV.11.10.1}[(c)].

  Ce lemme étant établi, on en conclut que pour $\mu\geq\lambda$ assez grand, $Z_\mu$ est induit sur un ouvert de $X_\mu$ en vertu de \sref{IV.8.6.3}, car $Z_\lambda$, en tant que sous-préschéma d'un préschéma noethérien, est de présentation finie sur $X_\lambda$ \sref{IV.1.6.2}, et il en est donc de même des $Z_\mu$ pour $\mu\geq\lambda$ sur $X_\mu$ et de $Z$ sur $X$.
  On peut maintenant appliquer \sref{IV.20.3.8.2} qui montre que pour $\mu\geq\lambda$ assez grand, $Z_\mu$ est schématiquement dense dans $X_\mu$, d'où la conclusion.

  \item Si, dans l'hypothèse~(i) de \sref{IV.20.3.8}, on supprime la condition que $X$ est noethérien, on voit que le raisonnement de \sref{IV.20.3.8.6} montre encore que l'image de \sref{IV.20.3.8.1} est formée des pseudo-$S$-morphismes ayant un représentant qui est un $S$-morphisme $U\to Y$, où $U$ est schématiquement dense dans $X$ et \emph{quasi-compact et quasi-séparé}.
\end{enumerate}
\end{remarks}
\end{proof}

\begin{corollary}[20.3.9]
\label{IV.20.3.9}
Supposons vérifiées l'une ou l'autre hypothèse~a),~b) de \sref{IV.20.3.6} sur $X$, et que $Y$ soit localement de présentation finie sur $S$.
Alors, pour qu'un pseudo-$S$-morphisme $\omega$ de $X$ dans $Y$ soit défini au point $x$ \sref{IV.20.2.3}, il faut et il suffit que sa restriction à $\Spec(\sh{O}_{X,x})$ soit partout définie (autrement dit, soit un $S$-morphisme de $\Spec(\sh{O}_{X,x})$ dans $Y$).
\end{corollary}

Le résultat suivant, qui nous servira dans la démonstration de \sref{IV.20.3.11}, utilise la théorie de la descente fidèlement plate du chap.~VI.
Le lecteur pourra vérifier que les résultats de \sref{IV.20.3} ne seront pas utilisés dans cette théorie.

\begin{lemma}[20.3.10]
\label{IV.20.3.10}
Soient $f:X'\to X$ un $S$-morphisme fidèlement plat et quasi-compact, $X''=X'\times_X X'$, $p_1$ et $p_2$ les projections canoniques de $X''$ dans $X'$, $Y$ un préschéma séparé sur $S$.
Soient $U$ un ouvert de $X$, $U'=f^{-1}(U)$, $U''=p_1^{-1}(U')=p_2^{-1}(U')$, et supposons que $U''$ soit schématiquement dense dans $X''$.
Soit $g:U\to Y$ un $S$-morphisme ; alors, si $g\circ(f|U')$ se prolonge en un $S$-morphisme $X'\to Y$, $g$ se prolonge en un $S$-morphisme $X\to Y$.
\end{lemma}

On notera que les hypothèses entraînent que $U$ (resp. $U'$) est schématiquement dense dans $X$ (resp. $X'$) \sref{IV.11.10.5}[(i)] ; on peut donc encore dire que si $\omega$ désigne le pseudo-$S$-morphisme classe de $g$, l'énoncé de \sref{IV.20.3.10} signifie que \emph{si $\omega\circ f$ est partout défini, il en est de même de $\omega$}.

\begin{proof}
Pour prouver \sref{IV.20.3.10}, désignons par $g'$ un $S$-morphisme $X'\to Y$ qui prolonge $g\circ(f|U')$, et posons $g''_i=g'\circ p_i:X''\to Y$ ($i=1,2$).
Si l'on pose $f''=f\circ p_1=f\circ p_2:X''\to X$, il est clair que $g''_1$ et $g''_2$ coïncident dans $U''$ avec $g\circ(f''|U'')$.
Mais comme $Y$ est séparé sur $S$ et $U''$ schématiquement dense dans $X''$, on a $g''_1=g''_2$ \sref{IV.11.10.1}[(d)].
Comme $f$ est fidèlement plat et quasi-compact, il résulte de la théorie de la descente (chap.~VI) qu'il existe un unique $S$-morphisme $h:X\to Y$ tel que $h\circ f=g'$ ; comme la restriction $U'\to U$ de $f$ est un morphisme fidèlement plat et quasi-compact et que $U''=U'\times_U U'$, le résultat d'unicité précédent, appliqué à $U$ au lieu de $X$, montre que $h|U=g$, ce qui prouve le lemme.
\end{proof}

\begin{proposition}[20.3.11]
\label{IV.20.3.11}
Soient $Y$ un $S$-préschéma séparé sur $S$, $\omega$ un pseudo-$S$-morphisme de $X$ dans $Y$, $f:X'\to X$ un $S$-morphisme.
Supposons que $f$ soit plat, et que l'une des conditions suivantes soit réalisée :
\begin{enumerate}[label=\emph{(\roman*)}]
\oldpage[IV-4]{243}
  \item $f$ est un morphisme ouvert, ou surjectif et quasi-compact, et $\operatorname{dom}_S(\omega)$ contient un ouvert $U$ schématiquement dense dans $X$ et rétrocompact dans $X$.
  \item $f$ est localement de présentation finie.
  \item $Y$ est localement de présentation finie sur $S$, et $X$ vérifie l'une des conditions~a),~b) de \sref{IV.20.3.6}.
\end{enumerate}
Alors, $f^{-1}(\operatorname{dom}_S(\omega))$ est schématiquement dense dans $X'$, de sorte que $\omega\circ f$ est défini \sref{IV.20.3.2} et l'on a
\[
\label{IV.20.3.11.1}
  \operatorname{dom}_S(\omega\circ f)=f^{-1}(\operatorname{dom}_S(\omega)).
  \tag{20.3.11.1}
\]
\end{proposition}

\begin{proof}
Prouvons d'abord que $f^{-1}(\operatorname{dom}_S(\omega))$ est schématiquement dense dans $X'$.
La question étant locale sur $X$ et $X'$, on peut supposer $X$ et $X'$ affines et puisque $f$ est plat, il suffit de voir, en vertu de \sref{IV.11.10.5}[(ii), (a)] que $\operatorname{dom}_S(\omega)$ contient un ensemble ouvert $U$ rétrocompact et schématiquement dense dans $X$.
Cela résulte de l'hypothèse dans le cas~(i), et de \sref{IV.20.2.12} dans le cas~(iii), compte tenu du fait que dans un schéma affine, tout ouvert de la forme $\mathrm{D}(t)$ est rétrocompact ; enfin, dans le cas~(ii), on voit directement que $f^{-1}(\operatorname{dom}_S(\omega))$ est schématiquement dense dans $X'$ en appliquant \sref{IV.20.3.5}[(i)].

Prouvons maintenant \sref{IV.20.3.11.1}, autrement dit que, pour tout point $x'\in\operatorname{dom}_S(\omega\circ f)$, on a $x=f(x')\in\operatorname{dom}_S(\omega)$.
Notons d'abord qu'on peut se borner au cas où $f$ est \emph{fidèlement plat et quasi-compact}.
C'est en effet l'hypothèse dans le second cas de~(i) ; dans les autres cas, la question est locale sur $X'$, et on peut donc supposer déjà $X$ et $X'$ affines, donc $f$ quasi-compact.
Dans le premier cas de~(i) et dans le cas~(ii), $f$ est un morphisme ouvert \sref{IV.11.3.1}, donc on peut, en remplaçant $X$ par l'ouvert $f(X')$, supposer $f$ surjectif, donc fidèlement plat.
Dans le cas~(iii), en utilisant \sref{IV.20.3.9}, on peut se borner à prouver que la restriction de $\omega$ à $\Spec(\sh{O}_{X,x})$ est partout définie, et on peut donc remplacer $X$ par $\Spec(\sh{O}_{X,x})$, $X'$ par $X'\times_X\Spec(\sh{O}_{X,x})$ et $f$ par sa restriction à ce dernier préschéma, qui est un morphisme surjectif \sref{IV.2.3.4}, donc fidèlement plat.

Supposons donc $f$ fidèlement plat et quasi-compact ; avec les notations du lemme \sref{IV.20.3.10}, il suffit de voir alors que $U''$ est schématiquement dense dans $X''$, en prenant pour $U$ un ouvert schématiquement dense dans $X$, contenu dans $\operatorname{dom}_S(\omega)$ ; cela sera le cas, en vertu de \sref{IV.11.10.5}[(ii), (a)] si $U$ est pris rétrocompact dans $X$ (puisque le morphisme $f'':X''\to X$ est plat).
Or l'existence d'un tel ouvert $U$ fait partie de l'hypothèse dans le cas~(i) ; dans le cas~(iii) elle résulte de \sref{IV.20.2.12} et du fait que dans un schéma affine $\Spec(A)$, tout ouvert de la forme $\mathrm{D}(t)$ est rétrocompact.
Enfin, dans le cas~(ii), prenons $U=\operatorname{dom}_S(\omega)$ et montrons directement que $U''$ est schématiquement dense dans $X''$ : il suffit pour cela de remarquer que $f'':X''\to X$ est plat et localement de présentation finie et d'appliquer \sref{IV.11.10.5}[(ii), (b)].
\end{proof}

\begin{corollary}[20.3.12]
\label{IV.20.3.12}
Soit $\vphi$ une pseudo-fonction sur un préschéma $X$.
Alors, pour tout morphisme plat et localement de présentation finie $f:X'\to X$, la pseudo-fonction $\vphi\circ f$ est définie et l'on a $\operatorname{dom}(\vphi\circ f)=f^{-1}(\operatorname{dom}(\vphi))$.
\end{corollary}

\begin{remark}[20.3.13]
\label{IV.20.3.13}
\oldpage[IV-4]{244}
Lorsque $X$ vérifie l'une des conditions~a),~b) de \sref{IV.20.3.6}, alors on a vu \sref{IV.20.2.11} que les pseudo-fonctions sur $X$ s'identifient aux fonctions méromorphes sur $X$.
En vertu de \sref{IV.20.1.12} et de \sref{IV.20.2.13}[(i)], si l'on suppose seulement que le morphisme $f:X'\to X$ est \emph{plat}, alors, pour toute pseudo-fonction $\vphi$ sur $X$, $\vphi\circ f$ est définie et l'on a $\operatorname{dom}(\vphi\circ f)=f^{-1}(\operatorname{dom}(\vphi))$.
\end{remark}

\subsection{Propriétés des domaines de définition des fonctions rationnelles}
\EGAReaderTarget{ega4.c04.subsec.20-4}{IV.20.4}

\begin{env}[20.4.1]
\EGAReaderTarget{ega4.c04.statement.20-4-1}{IV.20.4.1}
Soient $X$, $Y$ deux $S$-préschémas, $\omega$ un pseudo-$S$-morphisme de $X$ dans $Y$.
Soit $u$ un $S$-morphisme $U\to Y$ appartenant à $\omega$, où $U$ est schématiquement dense dans $X$, et considérons le \emph{graphe} $\Gamma_u$, sous-préschéma de $U\times_S Y$ \sref[I]{I.5.3.11}, donc sous-préschéma de $X\times_S Y$.
Supposons que ce sous-préschéma admette une \emph{adhérence} \sref[I]{I.9.5.11} dans $X\times_S Y$ ; si $j:\Gamma_u\to X\times_S Y$ est l'injection canonique, ce sera le cas lorsque le $\sh{O}_{X\times_S Y}$-Module $j_*(\sh{O}_{\Gamma_u})$ est quasi-cohérent ; il résulte de la définition de la classe d'équivalence $\omega$ \sref{IV.20.2.1} que $j_*(\sh{O}_{\Gamma_u})$ ne dépend pas du représentant $u$ considéré, donc le préschéma adhérence de $\Gamma_u$ est alors bien déterminé par $\omega$ ; on dit que ce sous-préschéma fermé de $X\times_S Y$ est le \emph{graphe} du pseudo-$S$-morphisme $\omega$ et on le note $\Gamma_\omega$.
On notera que $\Gamma_\omega$ est défini s'il existe un morphisme $u:U\to Y$ de la classe $\omega$ tel que $U$ soit \emph{rétrocompact} dans $X$ et si de plus $Y$ est quasi-séparé sur $S$, car alors l'injection $j$ est un morphisme quasi-compact et quasi-séparé (\sref[I]{I.2.2} et \sref[I]{I.7.4}) ; la première hypothèse sera toujours vérifiée lorsque $X$ est \emph{localement noethérien}.
Notons aussi que lorsque $X$ est \emph{réduit}, il en est de même de $\Gamma_u$, qui est isomorphe à $U$ \sref[I]{I.5.3.11} ; alors $\Gamma_\omega$ existe et n'est autre que le sous-préschéma \emph{réduit} de $X\times_S Y$ ayant pour espace sous-jacent l'adhérence dans $X\times_S Y$ de l'espace sous-jacent à $\Gamma_u$ (\sref[I]{I.5.2.1} et \sref[I]{I.5.2.2}).
Notons enfin que si $Y$ est \emph{séparé} sur $S$, $\Gamma_u$ est \emph{fermé} dans $U\times_S Y$ \sref[I]{I.5.4.3}, donc est \emph{induit par} $\Gamma_\omega$ (lorsque ce dernier existe) sur l'ouvert $\Gamma_\omega\cap(U\times_S Y)$ \sref[I]{I.9.5.10} ; par contre, ce préschéma induit est en général différent de $\Gamma_u$ lorsque $Y$ n'est pas séparé.
En particulier, si $v:X\to Y$ est un $S$-morphisme, le graphe $\Gamma_\omega$ de la classe $\omega$ de $v$ peut être distinct du graphe $\Gamma_v$.
Aussi allons-nous dans ce qui suit, lorsqu'il sera question du graphe d'un pseudo-$S$-morphisme, nous borner au cas où \emph{$Y$ est séparé sur $S$}.
\end{env}

\begin{env}[20.4.2]
\label{IV.20.4.2}
Supposons donc $\Gamma_\omega$ défini et $Y$ séparé sur $S$ ; notons $p$ et $q$ les restrictions à $\Gamma_\omega$ des projections canoniques
\[
  \xymatrix{
  & \Gamma_\omega \ar[dl]_p \ar[dr]^q & \\
  X & & Y
  }
\]
Alors, si $U\subset\operatorname{dom}_S(\omega)$, la restriction $p^{-1}(U)\to U$ de $p$ est un isomorphisme \sref[I]{I.5.3.11} ; réciproquement, si $U$ est un ouvert de $X$ ayant cette propriété, et si $u$ est l'isomorphisme réciproque de la restriction $p^{-1}(U)\to U$ de $p$, $q\circ u$ est un $S$-morphisme de $U$ dans $Y$ qui coïncide avec un morphisme de la classe $\omega$ dans $U\cap\operatorname{dom}_S(\omega)$.
On en conclut que $\operatorname{dom}_S(\omega)$ est le \emph{plus grand ouvert} $U$ de $X$ tel que la restriction $p^{-1}(U)\to U$ de $p$ soit un
\oldpage[IV-4]{245}
isomorphisme.
Soit $v:\operatorname{dom}_S(\omega)\to\Gamma_\omega$ le $S$-morphisme réciproque de la restriction $p^{-1}(\operatorname{dom}_S(\omega))\to\operatorname{dom}_S(\omega)$ de $p$ ; on note parfois $p^{-1}$ le pseudo-$S$-morphisme de $X$ dans $\Gamma_\omega$, classe de $v$ (dont l'application rationnelle associée \sref{IV.20.2.13}[(ii)] est alors \emph{birationnelle}) ; comme $p^{-1}(\operatorname{dom}_S(\omega))$ est le graphe du $S$-morphisme $u=q\circ v:\operatorname{dom}_S(\omega)\to Y$, il est schématiquement dense dans $\Gamma_\omega$ \sref{IV.11.10.3}[(iv)], donc $\omega$ peut être considéré comme le \emph{composé} $q\circ p^{-1}$ \sref{IV.20.3.2}.
Pour toute partie $M$ de l'espace sous-jacent à $X$, on pose parfois $\omega(M)=q(p^{-1}(M))$, ce qui revient alors à considérer $\omega$ comme une application de $X$ dans $\mathfrak{P}(Y)$ (ou, comme disent certains auteurs, une « fonction multiforme »).
On notera que lorsque $x\in\operatorname{dom}_S(\omega)$, $\omega(\{x\})$ est l'ensemble $\{u(x)\}$ ; en général, pour un $x\in X$ quelconque, si l'on désigne par $s$ l'image de $x$ dans $S$, par $Y_s$ la fibre au point $s$ du morphisme structural $Y\to S$, l'ensemble $\omega(\{x\})$ est une partie du préschéma $Y_s$.
\end{env}

\begin{env}[20.4.3]
\label{IV.20.4.3}
Dans tout le reste de ce numéro, nous allons nous borner au cas où $X$ est \emph{réduit}, de sorte qu'il y a alors identité entre pseudo-$S$-morphismes et $S$-\emph{applications rationnelles} \sref{IV.20.2.7}.
En outre, le \emph{graphe} de toute $S$-application rationnelle de $X$ dans $Y$ est alors défini \sref{IV.20.4.1}.
\end{env}

\begin{proposition}[20.4.4]
\label{IV.20.4.4}
Soient $X$ un $S$-préschéma localement intègre, $Y$ un $S$-préschéma localement de type fini et séparé sur $S$, $\omega$ une $S$-application rationnelle de $X$ dans $Y$, $p:\Gamma_\omega\to X$ la projection canonique.
Alors, si $x\in X$ est un point normal tel que l'ensemble $p^{-1}(x)$ contienne un point isolé, $\omega$ est définie au point $x$.
\end{proposition}

\begin{proof}
En effet, la première projection $p_1:X\times_S Y\to X$ est alors un morphisme séparé et localement de type fini, donc il en est de même de sa restriction $p:\Gamma_\omega\to X$, qui est en outre \emph{birationnel} et comme $\Gamma_\omega$ est réduit et $X$ intègre, $\Gamma_\omega$ est \emph{intègre} ; il résulte donc de (\textbf{Err}\textsubscript{\textbf{IV}}, 30) que l'hypothèse sur $x$ entraîne l'existence d'un voisinage ouvert $U$ de $x$ tel que la restriction $p^{-1}(U)\to U$ de $p$ soit un isomorphisme ; d'où la conclusion \sref{IV.20.4.2}.

On notera que l'énoncé \sref{IV.20.4.4} est la formulation originelle donnée par Zariski à son « Main theorem » (à cela près qu'il se restreignait aux schémas algébriques sur un corps de base).
\end{proof}

\begin{proposition}[20.4.5]
\label{IV.20.4.5}
Les hypothèses et notations étant celles de \sref{IV.20.4.4}, supposons en outre $X$ normal, et soient $X'$ un préschéma réduit, $f:X'\to X$ un morphisme localement de type fini et universellement ouvert.
Alors $\omega\circ f$ est défini, et l'on a $\operatorname{dom}_S(\omega\circ f)=f^{-1}(\operatorname{dom}_S(\omega))$ (en d'autres termes, si $x'\in X'$ et $x=f(x')$, alors, pour que $\omega$ soit définie au point $x$, il faut et il suffit que $\omega\circ f$ le soit au point $x'$).
\end{proposition}

\begin{proof}
Comme $X'$ est réduit et $f^{-1}(\operatorname{dom}_S(\omega))$ partout dense dans $X'$ en vertu de \sref{IV.2.4.11}, le composé $\omega\circ f$ est défini ; pour prouver que, lorsque $\omega\circ f$ est définie au point $x'$, $\omega$ l'est au point $x$, on peut évidemment remplacer $X'$ par un voisinage ouvert de $x'$, donc supposer $\omega\circ f$ \emph{partout définie} ; en outre, comme $f(X')$ est ouvert dans $X$, on peut supposer $f$ \emph{surjectif}.
En vertu de l'hypothèse sur $f$, le morphisme $f_{(Y)}:X'\times_S Y\to X\times_S Y$ est alors ouvert, donc $\overline{f_{(Y)}^{-1}(M)}=f_{(Y)}^{-1}(\overline{M})$ pour toute partie $M$ de $X\times_S Y$ ; prenant pour $M$ l'ensemble $\Gamma_u$, où $u:\operatorname{dom}_S(\omega)\to Y$ est la restriction de $\omega$ à $\operatorname{dom}_S(\omega)$, il résulte de la relation précédente et de \sref[I]{I.3.3.12} que l'ensemble sous-jacent à $\Gamma_{\omega\circ f}$ est égal à $f_{(Y)}^{-1}(\Gamma_\omega)$ ; comme on sait
\oldpage[IV-4]{246}
que $\Gamma_{\omega\circ f}$ est un préschéma réduit \sref{IV.20.4.1}, on voit que le $X'$-préschéma $\Gamma_{\omega\circ f}$ est égal à $(\Gamma_\omega\times_X X')_\mathrm{red}$.
Mais comme par hypothèse le morphisme composé $\Gamma_{\omega\circ f}\to\Gamma_\omega\times_X X'\xrightarrow{p^{(X')}}X'$ est un isomorphisme, $p^{(X')}$ est nécessairement \emph{radiciel} ; comme $f$ est surjectif, il en est de même de $p:\Gamma_\omega\to X$ \sref[I]{I.3.4.8}, donc $p^{-1}(x)$ est un ensemble réduit à un point \sref[I]{I.3.6.4} ; il résulte alors de \sref{IV.20.4.4} que $\omega$ est définie au point $x$.
\end{proof}

La proposition suivante donne un \emph{critère valuatif} pour qu'une application rationnelle soit définie en un point :

\begin{proposition}[20.4.6]
\label{IV.20.4.6}
Soient $S$ un préschéma, $X$, $Y$ deux $S$-préschémas ; on suppose $X$ localement noethérien, $Y$ localement de type fini sur $S$.
Soient $U$ un ouvert dense dans $X$, $h:U\to Y$ un $S$-morphisme, $x\in X-U$ un point normal de $X$, $h'_x:\Spec(k(x))\to Y$ un $S$-morphisme.
Afin que $h$ puisse être prolongé en un $S$-morphisme $h':U'\to Y$, où $U'$ est un ouvert de $X$ contenant $U$ et $x$, et où le morphisme composé $\Spec(k(x))\to U'\xrightarrow{h'}Y$ est le $S$-morphisme donné $h'_x$, il faut et il suffit que la condition suivante soit vérifiée :
\begin{quote}
\emph{(P) Pour tout spectre $S_1$ d'un anneau de valuation discrète, de point fermé $a$ et de point générique $b$, et tout morphisme $u:S_1\to X$ tel que $u(a)=x$, $u(b)\in U$, le morphisme composé $h_1=h\circ(u|\{b\}):\{b\}=u^{-1}(U)\to Y$, se prolonge en un morphisme $h'_1:S_1\to Y$ tel que le diagramme}
\[
  \xymatrix{
  \Spec(k(a)) \ar[r] & S_1 \ar[d]^{h'_1} \\
  \Spec(k(x)) \ar[r]_{h'_x} & Y
  }
\]
\emph{soit commutatif.}
\end{quote}
De plus, si cette condition est vérifiée, et si $h'':U''\to Y$ est un morphisme vérifiant les mêmes conditions que $h'$, alors il existe un ensemble ouvert contenant $U\cup\{x\}$, et dans lequel $h'$ et $h''$ coïncident.
\end{proposition}

\begin{proof}
Prouvons d'abord la dernière assertion ; on peut supposer $h'$ et $h''$ définis dans le même ouvert $U_0\supset U\cup\{x\}$.
Le sous-préschéma $Z$ des coïncidences de $h'$ et $h''$ \sref{IV.17.4.5} contient $U$ et $x$, donc il y a un voisinage ouvert $V$ de $x$ dans $U_0$ tel que le sous-préschéma induit par $Z$ sur l'ouvert $Z\cap V$ soit un sous-préschéma \emph{fermé} du préschéma induit par $X$ sur $V$ ; comme ce préschéma $Z\cap V$ majore le sous-préschéma induit par $V$ sur l'ouvert $U\cap V$, et que ce dernier est schématiquement dense dans $V$, $Z\cap V$ est nécessairement égal à $V$ \sref{IV.20.3.8.8}, ce qui prouve que $h'$ et $h''$ coïncident dans $U\cup V$.

La nécessité de la condition de l'énoncé étant évidente, prouvons qu'elle est suffisante.
En vertu de la correspondance biunivoque entre $S$-morphismes de $X$ dans $Y$ et $X$-sections de $X\times_S Y$ \sref[I]{I.3.3.15}, on peut se borner au cas où $S=X$ et où $h$ est donc une $U$-section de $Y$ ; on notera qu'alors $Y$ est localement noethérien, et on peut évidemment (puisque $X$ est localement intègre et localement noethérien) supposer $X=S$ irréductible.
En outre, compte tenu de \sref{IV.20.3.7}, on peut supposer que $X=\Spec(A)$, où $A$ est un anneau \emph{local} noethérien, intègre et intégralement clos.

\oldpage[IV-4]{247}
Notons que pour tout $x'\in U$, $h(x')$ est une spécialisation de $h'_x(x)=y$.
En effet, il existe un spectre $S_1$ d'anneau de valuation discrète et un morphisme $u:S_1\to X$ tel que $u(a)=x$, $u(b)=x'$ \sref[II]{II.7.1.9}.
Appliquant la condition de l'énoncé, on obtient aussitôt notre assertion, puisque $h(x')=h_1(b)=h'_1(b)$ et $y=h'_1(a)$.
Si $Y'$ est un voisinage ouvert affine de $y$, on a donc $h(X\cap U)\subset Y'$ ; on peut par suite remplacer $Y$ par $Y'$, autrement dit supposer $Y$ \emph{affine}, donc séparé sur $X$.
Soit $\omega$ la $X$-section rationnelle de $Y$ à laquelle $h$ appartient, de sorte que son graphe a ici pour ensemble sous-jacent l'adhérence de $h(U)$ dans $Y$.
Puisque $Y$ est séparé sur $X$, on peut appliquer \sref{IV.20.4.4} à $\omega$ : il suffira de prouver que, si $p:\Gamma_\omega\to X$ est la projection canonique, $p^{-1}(x)$ est réduit à un seul point $y$ et que $h'_x(x)=y$.
En effet, par \sref{IV.20.4.4} $h$ se prolongera en une section $h'$ au-dessus d'un ouvert $U'$ contenant $U\cup\{x\}$, telle que $h'(x)=y$, et comme alors il n'existe qu'un seul $X$-morphisme $\Spec(k(x))\to Y$ transformant $x$ en $y$, cela prouvera l'identité de $h'_x$ et du composé de $h'$ et de $\Spec(k(x))\to U'$.

Puisque pour $x'\in X\cap U$, $h(x')$ est une spécialisation de $y$, on a $y\in p^{-1}(x)$.
Supposons alors que $z$ soit un point quelconque de $p^{-1}(x)$.
Comme $\Gamma_\omega$ est l'adhérence de $h(U)$ et que $h(U)$ est formé de points adhérents à $h(\xi)$, où $\xi$ est le point générique de $X$, $\Gamma_\omega$ est l'adhérence de $h(\xi)$ dans $Y$.
On prend alors un spectre $S_1$ d'anneau de valuation discrète et un morphisme $v:S_1\to Y$ tels que $v(a)=z$, $v(b)=h(\xi)$ \sref[II]{II.7.1.9}, et on pose $u=p\circ v$, de sorte que $u(a)=x$, $u(b)=\xi$.
Appliquant à $u$ la condition de l'énoncé, on voit que l'on obtient un morphisme $h'_1:S_1\to Y$ tel que $h'_1(a)=y$ et $h'_1(b)=h(\xi)$ ; mais cela entraîne $z=y$ en vertu de \sref[II]{II.7.2.3}, puisque $Y$ est séparé sur $X$ et que $v$ et $h'_1$ doivent donc coïncider, puisqu'ils sont égaux au point $b$.
C.Q.F.D.
\end{proof}

\begin{corollary}[20.4.7]
\EGAReaderTarget{ega4.c04.cor.20-4-7}{IV.20.4.7}
Les hypothèses sur $S$, $X$, $Y$, $U$ et $h$ étant celles de \sref{IV.20.4.6}, soit $E$ une partie de $X-U$ telle que $X$ soit normal en tout point de $E$, et pour chaque $x\in E$, soit $h'_x:\Spec(k(x))\to Y$ un $S$-morphisme, tel que la condition~(P) de \sref{IV.20.4.6} soit vérifiée.
On suppose en outre que la réunion $F$ des graphes des $h'_x$ (identifiés à des parties de $X\times_S Y$) pour $x\in E$, soit contenue dans la réunion d'un nombre fini d'ouverts $V_i$ de $X\times_S Y$ qui sont \emph{séparés} sur $X$ (condition automatiquement vérifiée si $Y$ est \emph{séparé} sur $S$, ou si $X$ est \emph{quasi-compact} et $Y$ de \emph{type fini} sur $S$).
Alors il existe un $S$-morphisme $h':U'\to Y$, où $U'$ est un ouvert de $X$ contenant $U\cup E$, tel que, pour tout $x\in E$, le composé
\[
  \Spec(k(x))\to U'\xrightarrow{h'}Y
\]
soit égal à $h'_x$.
\end{corollary}

\begin{proof}
Notons d'abord que, dans \sref{IV.20.4.6}, si $Y$ est supposé \emph{séparé}, il y a un plus grand ouvert $U_0\supset U$, égal au domaine de la $S$-application rationnelle correspondant à $h$, dans lequel $h$ peut être prolongé, et ce prolongement est unique \sref[I]{I.7.2.2} ; d'où la conclusion dans ce cas, grâce à \sref{IV.20.4.6}.
Dans le cas général, soit $E_i$ l'ensemble des $x\in E$ tels que $(x,h'_x(x))\in V_i$.
En vertu de \sref{IV.20.4.6}, pour tout $x\in E$, il y a un prolongement $h^{(x)}$ de $h$ à $U\cup W^{(x)}$, où $W^{(x)}$ est un voisinage de $x$ dans $X$ tel que le graphe de $h^{(x)}|W^{(x)}$ soit contenu dans \emph{tous} les $V_i$ tels que $x\in E_i$.
Puisque $V_i$ est séparé sur $X$ et $X$ réduit, pour deux points $x'$, $x''$ de $E_i$, $h^{(x')}$ et $h^{(x'')}$ coïncident dans $W^{(x')}\cap W^{(x'')}$ puisqu'ils
\oldpage[IV-4]{248}
coïncident dans l'ouvert partout dense $U\cap W^{(x')}\cap W^{(x'')}$ de cet ensemble.
Il y a donc un morphisme $h_i:U\cup U_i\to Y$ qui prolonge $h$ à un ouvert $U\cup U_i$, contenant $E_i$ ; en outre, pour tout couple d'indices $i$, $j$, les graphes des restrictions $h_i|(U_i\cap U_j)$ et $h_j|(U_i\cap U_j)$ sont contenus dans $V_i\cap V_j$ ; comme $V_i\cap V_j$ est séparé sur $X$ et que les morphismes précédents coïncident dans un ouvert partout dense $U\cap U_i\cap U_j$ de $U_i\cap U_j$, ils sont égaux.
Le morphisme $h'$ égal à $h$ dans $U$, à $h_i$ dans chacun des $U_i$, répond à la question.
\end{proof}

\begin{remark}[20.4.8]
\label{IV.20.4.8}
Lorsque $E=X-U$, il est clair que si, pour tout ouvert affine $T$ de $X$, il existe un $S$-morphisme $h'_T:T\to Y$ prolongeant $h|(U\cap T)$ et tel que le composé $\Spec(k(x))\to T\xrightarrow{h'_T}Y$ soit égal à $h'_x$ pour tout $x\in T-(U\cap T)$, alors les $h'_T$ sont les restrictions d'un $S$-morphisme $h':X\to Y$ (partout défini) en vertu de l'assertion d'unicité dans \sref{IV.20.4.6}.
Pour prouver l'existence de $h'$, on est donc ramené au cas où $X$ est \emph{affine}, et alors il suffit que l'ensemble $F$, réunion des graphes des $h'_x$, soit \emph{quasi-compact} dans $X\times_S Y$ pour que l'on puisse appliquer \sref{IV.20.4.7}.
Il en sera ainsi lorsque les $h'_x$ sont de la forme $\Spec(k(x))\to Z\xrightarrow{h_0}Y$, où $Z$ est un sous-préschéma fermé de $X$ ayant $X-U$ pour espace sous-jacent, et $h_0$ un $S$-morphisme.
\end{remark}

\begin{corollary}[20.4.9]
\label{IV.20.4.9}
Soient $S$ un préschéma localement noethérien, $X$ un préschéma localement noethérien, $f:X\to S$ un morphisme plat, $g:Y\to S$ un morphisme localement de type fini.
Soient $U$ un ouvert dense dans $S$, $h:f^{-1}(U)\to Y$ un $S$-morphisme, $T=S-U$, $Z$ le sous-préschéma fermé réduit de $X$ ayant $f^{-1}(T)$ comme espace sous-jacent, $h_0:Z\to Y$ un $S$-morphisme.
Supposons $X$ normal en tous les points de $Z$.
Afin qu'il existe un $S$-morphisme (nécessairement unique) $h':X\to Y$ prolongeant $h$ et $h_0$, il faut et il suffit que la condition suivante soit satisfaite :

\emph{Pour tout spectre $S_1$ d'un anneau de valuation discrète, de point fermé $a$ et de point générique $b$, et tout morphisme $u:S_1\to S$ tel que $u(a)\in T$ et $u(b)\in U$, il existe un $S_1$-morphisme $h'_1:X_{(S_1)}\to Y_{(S_1)}$ qui prolonge $h_{(S_1)}:f_{(S_1)}^{-1}(b)\to Y_{(S_1)}$ et est tel que le diagramme}
\[
  \xymatrix{
  \Spec(k(z_1)) \ar[r] \ar[d] & Z_{(S_1)} \ar[d]^{(h_0)_{(S_1)}} \\
  X_{(S_1)} \ar[r]_{h'_1} & Y_{(S_1)}
  }
\]
\emph{soit commutatif pour tout $z_1\in Z_{(S_1)}$.}

L'hypothèse que $f$ est plat entraîne en effet que $f^{-1}(U)$ est dense dans $X$ \sref{IV.2.3.10}, et il suffit alors d'appliquer \sref{IV.20.4.8}.
\end{corollary}

\begin{corollary}[20.4.10]
\label{IV.20.4.10}
Sous les hypothèses de \sref{IV.20.4.6}, supposons en outre $Y$ séparé et localement quasi-fini sur $S$.
Soient $U$ un ouvert dense dans $X$, $h:U\to Y$ un $S$-morphisme, $\omega$ la $S$-application rationnelle correspondante, $x\in X-U$ un point normal de $X$.
Afin que $\omega$ soit définie au point $x$, il faut et il suffit que la condition suivante soit vérifiée : il existe un spectre $S_1$ d'un anneau de valuation discrète, de point fermé $a$ et de point générique $b$, et un morphisme $u:S_1\to X$
\oldpage[IV-4]{249}
tels que $u(a)=x$, $u(b)\in U$ et que la $S$-application rationnelle $\omega\circ u$ soit partout définie.

En effet, par hypothèse toutes les fibres du morphisme projection $X\times_S Y\to X$ sont formées de points isolés, et pour appliquer \sref{IV.20.4.4} il suffit de savoir que la fibre $p^{-1}(x)$ est non vide dans $\Gamma_\omega$.
Or si $h_1$ est l'unique morphisme de la classe $\omega\circ u$, l'unique point de $X\times_S Y$ au-dessus de $x$ et de $h_1(a)$ appartient à $\Gamma_\omega$, d'où la conclusion.
\end{corollary}

\begin{proposition}[20.4.11]
\label{IV.20.4.11}
Soient $X$ un préschéma localement noethérien, $Y$ un $S$-préschéma affine sur $S$, $U$ un ouvert de $X$, $Z=X-U$.
Supposons que l'on ait $\operatorname{prof}_Z(\sh{O}_X)\geq 2$ \sref{IV.5.10.1} ; alors tout $S$-morphisme $f:U\to Y$ se prolonge de façon unique en un $S$-morphisme de $X$ dans $Y$.
\end{proposition}

\begin{proof}
On peut se borner au cas où $S$ et $X$ sont affines et (en vertu de \sref[I]{I.3.3.14}) au cas où $S=X$ ; on a donc $X=\Spec(A)$, $Y=\Spec(B)$, $B$ étant une $A$-algèbre de type fini.
Comme $B$ est quotient d'une algèbre de polynômes $A[(T_\lambda)]_{\lambda\in L}$, $Y$ est un sous-préschéma fermé de $Y'=X[(T_\lambda)]_{\lambda\in L}$.
D'autre part, l'hypothèse sur $Z$ entraîne que $U$ est schématiquement dense dans $X$ en vertu de \sref{IV.20.2.13}[(iv)] et \sref{IV.5.10.2}.
Si l'on prouve que tout $X$-morphisme $u:U\to Y$ se prolonge de façon unique en un $X$-morphisme $v':X\to Y'$, il en résultera que $v'$ se factorisera en $X\xrightarrow{v}Y\to Y'$ ; en effet, le sous-préschéma $v'^{-1}(Y)$ est fermé et majore $U$ \sref[I]{I.4.4.1}, donc est identique à $X$ \sref{IV.20.3.8.8}.
Dans ces conditions, $v$ sera l'unique $S$-morphisme de $X$ dans $Y$ prolongeant $u$.
On peut donc se borner au cas où $Y=Y'$.
Mais alors il y a correspondance biunivoque entre les $X$-morphismes d'un ouvert $U\subset X$ dans $Y'$ et les familles $(s_\lambda)_{\lambda\in L}$ de sections de $\sh{O}_X$ au-dessus de $U$ \sref[II]{II.1.7.9} ; la conclusion résulte donc de \sref{IV.5.10.5}.
\end{proof}

\begin{corollary}[20.4.12]
\label{IV.20.4.12}
Soient $X$ un $S$-préschéma localement noethérien réduit et vérifiant la condition $(S_2)$ \sref{IV.5.7.2} (par exemple un $S$-préschéma localement noethérien et normal \sref{IV.5.8.6}), $Y$ un $S$-préschéma affine sur $S$, $f$ une $S$-application rationnelle de $X$ dans $Y$ ; alors toute composante irréductible de $X-\operatorname{dom}(f)$ est de codimension $1$ dans $X$.
\end{corollary}

\begin{proof}
Il revient au même de dire que si $Z^{(2)}$ est l'ensemble des $x\in X$ tels que $\dim(\sh{O}_{X,x})\geq 2$, alors, pour toute partie fermée $Z\subset Z^{(2)}$ de $X$, tout $S$-morphisme de $X-Z$ dans $Y$ se prolonge en un $S$-morphisme de $X$ dans $Y$ ; or, cela résulte de l'hypothèse sur $X$ \sref{IV.5.7.2} et de \sref{IV.20.4.11}.
\end{proof}

\subsection{Pseudo-morphismes relatifs}
\EGAReaderTarget{ega4.c04.subsec.20-5}{IV.20.5}

\begin{env}[20.5.1]
\EGAReaderTarget{ega4.c04.statement.20-5-1}{IV.20.5.1}
Soient $X$, $Y$ deux $S$-préschémas.
Il résulte des définitions \sref{IV.11.10.8} que l'intersection de deux ouverts $U$, $U'$ de $X$, \emph{universellement schématiquement denses par rapport à $S$}, possède encore cette propriété.
On définit donc une relation d'équivalence entre $S$-morphismes $u:U\to Y$ en remplaçant dans \sref{IV.20.2.1} « schématiquement dense » par « universellement schématiquement dense par rapport à $S$ ».
Une classe d'équivalence suivant cette relation s'appelle un \emph{pseudo-morphisme de $X$ dans $Y$ relativement à $S$}, et l'ensemble de ces classes se note $\operatorname{Ps.hom}_{X/S}(X,Y)$.
\end{env}

\begin{env}[20.5.2]
\label{IV.20.5.2}
Comme tout ouvert de $X$ universellement schématiquement dense par rapport à $S$ est en particulier schématiquement dense, les éléments d'un pseudo-morphisme
\oldpage[IV-4]{250}
de $X$ dans $Y$ relativement à $S$ sont équivalents au sens de \sref{IV.20.2.1}, d'où une application canonique
\[
\label{IV.20.5.2.1}
  \operatorname{Ps.hom}_{X/S}(X,Y)\to\operatorname{Ps.hom}_S(X,Y).
  \tag{20.5.2.1}
\]
\end{env}

\begin{proposition}[20.5.3]
\label{IV.20.5.3}
Supposons $Y$ séparé sur $S$.
Alors :
\begin{enumerate}[label=\emph{(\roman*)}]
  \item L'application \sref{IV.20.5.2.1} est injective et identifie $\operatorname{Ps.hom}_{X/S}(X,Y)$ au sous-ensemble de $\operatorname{Ps.hom}_S(X,Y)$ formé des pseudo-$S$-morphismes $\omega$ tels que $\operatorname{dom}_S(\omega)$ soit universellement schématiquement dense relativement à $S$.
  \item Le préfaisceau $U\mapsto\operatorname{Ps.hom}_{U/S}(U,Y)$ sur $X$ est un faisceau.
\end{enumerate}
\end{proposition}

\begin{proof}
L'assertion~(i) est immédiate, car dire que deux $S$-morphismes $u:U\to Y$, $u':U'\to Y$ sont équivalents au sens de \sref{IV.20.2.1} signifie alors que ce sont les restrictions d'un même morphisme $\operatorname{dom}_S(\omega)\to Y$ \sref{IV.20.2.4}, et si $U$ et $U'$ sont universellement schématiquement denses par rapport à $S$, il en est de même \emph{a fortiori} de $\operatorname{dom}_S(\omega)$.
Pour prouver~(ii), notons que $U\mapsto\operatorname{Ps.hom}_S(U,Y)$ est alors un faisceau \sref{IV.20.2.6} ; d'autre part, étant donné un recouvrement ouvert $(U_\alpha)$ de $U$, et un pseudo-$S$-morphisme $\omega$ de $U$ dans $Y$, pour que $\operatorname{dom}_S(\omega)$ soit universellement schématiquement dense dans $U$ relativement à $S$, il résulte aussitôt des définitions (cf.\ \sref{IV.20.2.1}) qu'il faut et il suffit que chacun des ensembles $\operatorname{dom}_S(\omega)\cap U_\alpha=\operatorname{dom}_S(\omega|U_\alpha)$ soit universellement schématiquement dense dans $U_\alpha$ relativement à $S$ ; d'où~(ii).
\end{proof}

Lorsque $Y$ est séparé, on notera $\mathscr{P}\!s.\!\operatorname{hom}_{X/S}(X,Y)$ le sous-faisceau
\[
  U\mapsto\operatorname{Ps.hom}_{U/S}(U,Y)
\]
de $\mathscr{P}\!s.\!\operatorname{hom}_S(X,Y)$.

Lorsque $X$ est \emph{plat et localement de présentation finie} sur $S$ et $Y$ séparé sur $S$, les pseudo-morphismes de $X$ dans $Y$ relativement à $S$ s'identifient encore aux pseudo-$S$-morphismes $\omega$ tels que, pour toute fibre $X_s$ du morphisme $X\to S$, $\operatorname{dom}_S(\omega)\cap X_s$ soit \emph{schématiquement dense} dans $X_s$ \sref{IV.11.10.10}.

\begin{env}[20.5.4]
\EGAReaderTarget{ega4.c04.statement.20-5-4}{IV.20.5.4}
Un cas particulièrement important où $Y$ est séparé sur $S$ est le cas où $Y=S[T]=S\otimes_{\bb{Z}}\Spec(\bb{Z}[T])$ ($T$ indéterminée).
Il y a alors correspondance biunivoque entre les pseudo-$S$-morphismes de $X$ dans $Y$ et les \emph{pseudo-fonctions} sur $X$ \sref{IV.20.2.8} en vertu de la définition d'un produit de $\bb{Z}$-préschémas.
Les pseudo-morphismes de $X$ dans $Y$ relativement à $S$ s'identifient alors, en vertu de \sref{IV.20.5.3}, aux \emph{pseudo-fonctions $\varphi$ sur $X$ telles que $\operatorname{dom}(\varphi)$ soit universellement schématiquement dense relativement à $S$}.
Le faisceau $\mathscr{P}\!s.\!\operatorname{hom}_{X/S}(X,Y)$ est un \emph{sous-faisceau d'anneaux} de $\sh{M}'_X$, que l'on note $\sh{M}'_{X/S}$.

On pose alors $\operatorname{Ps.hom}_{X/S}(X,Y)=M'(X/S)$ et on dit que ses éléments sont les \emph{pseudo-fonctions sur $X$ relatives à $S$}.
\end{env}

\begin{env}[20.5.5]
\label{IV.20.5.5}
Soient $X$, $Y$, $Z$ trois $S$-préschémas, $f:Y\to Z$ un $S$-morphisme.
On peut, dans le raisonnement de \sref{IV.20.3.1}, remplacer partout « schématiquement dense » par « universellement schématiquement dense relativement à $S$ » ; pour tout pseudo-morphisme $\omega\in\operatorname{Ps.hom}_{X/S}(X,Y)$, les morphismes $f\circ u$, où $u$ parcourt $\omega$, appartiennent donc à une même classe d'équivalence (pour la relation définie dans \sref{IV.20.5.1}), que l'on
\oldpage[IV-4]{251}
appelle le pseudo-morphisme de $X$ dans $Z$, relativement à $S$, \emph{composé} de $f$ et de $\omega$ et que l'on note $f\circ\omega$.
Si $g:Z\to T$ est un morphisme, il est immédiat que $g\circ(f\circ\omega)=(g\circ f)\circ\omega$.
\end{env}

\begin{env}[20.5.6]
\label{IV.20.5.6}
Supposons $Y$ \emph{séparé} sur $S$, et soit $\omega$ un pseudo-morphisme de $X$ dans $Y$ relativement à $S$.
Si $f:X'\to X$ est un $S$-morphisme tel que $f^{-1}(\operatorname{dom}_S(\omega))$ soit \emph{universellement schématiquement dense dans $X'$ relativement à $S$}, il résulte de \sref{IV.20.3.2} que le pseudo-$S$-morphisme $\omega\circ f$ a un domaine universellement schématiquement dense relativement à $S$, donc \sref{IV.20.5.3} peut être considéré comme un pseudo-morphisme \emph{relativement à $S$}.
Lorsque $X'$ est plat et localement de présentation finie sur $S$, la condition que $f^{-1}(\operatorname{dom}_S(\omega))$ soit universellement schématiquement dense relativement à $S$ équivaut encore à dire que pour tout $s\in S$, $(f^{-1}(\operatorname{dom}_S(\omega)))_s$ (notation de \sref{IV.11.10.10}) est schématiquement dense dans $X'_s$, ou encore, en notant $f_s:X'_s\to X_s$ le morphisme déduit de $f$ par changement de base, que l'image réciproque par $f_s$ de $(\operatorname{dom}_S(\omega))_s$ soit schématiquement dense dans $X'_s$.
Cette dernière condition sera en particulier remplie si, pour tout $s\in S$, $X_s$, $X'_s$ et $f_s$ vérifient l'une des trois conditions~(i),~(ii),~(iii) de \sref{IV.20.3.5}.
\end{env}

\begin{env}[20.5.7]
\label{IV.20.5.7}
Supposons maintenant que $X$ et $X'$ soient tous deux des $S$-préschémas \emph{plats et localement de présentation finie sur $S$}, et que $f:X'\to X$ soit un $S$-morphisme \emph{plat} (ou, ce qui revient au même \sref{IV.11.3.10}, que pour tout $s\in S$, $f_s:X'_s\to X_s$ soit un morphisme plat).
Alors, pour tout ouvert $U$ de $X$ et tout ouvert $V\subset U$ universellement schématiquement dense dans $U$ relativement à $S$, il résulte de \sref{IV.11.10.5} et \sref{IV.11.10.9} que $f^{-1}(V)$ est universellement schématiquement dense dans $f^{-1}(U)$ relativement à $S$.
Pour tout pseudo-morphisme $\omega$ de $X$ dans $Y$ relativement à $S$, il résulte de \sref{IV.20.3.4} que le pseudo-$S$-morphisme $\omega\circ f$ est défini et est un pseudo-morphisme de $X'$ dans $Y$ relativement à $S$, \emph{même lorsqu'on ne suppose pas $Y$ séparé sur $S$}.
On en déduit que dans ce cas, pour tout $S$-morphisme $g:Y\to Z$, $(g\circ\omega)\circ f$ est encore défini et égal à $g\circ(\omega\circ f)$ (avec les définitions de \sref{IV.20.5.1}), et est donc un pseudo-morphisme \emph{relativement à $S$}.
\end{env}

\begin{env}[20.5.8]
\label{IV.20.5.8}
Soient $X$ un $S$-préschéma, $S'\to S$ un morphisme quelconque, $X'=X_{(S')}$, $g:X'\to X$ la projection canonique.
Alors, pour tout ouvert $U$ de $X$ et tout ouvert $V\subset U$ universellement schématiquement dense dans $U$ relativement à $S$, $V'=g^{-1}(V)$ est universellement schématiquement dense dans $U'=g^{-1}(U)$ relativement à $S'$ par définition \sref{IV.11.10.8}.
Soit alors $\omega$ un pseudo-morphisme de $X$ dans un $S$-préschéma $Y$ relativement à $S$ ; si $u_1$, $u_2$ sont des $S$-morphismes de la classe $\omega$, définis respectivement dans des ouverts $U_1$, $U_2$ de $X$, universellement schématiquement denses dans $X$ relativement à $S$, il résulte de ce qui précède que les morphismes $u'_1=u_1\circ(g|g^{-1}(U_1))$ et $u'_2=u_2\circ(g|g^{-1}(U_2))$ coïncident dans un ouvert $U'_3$ universellement schématiquement dense relativement à $S'$.
Or, si $Y'=Y_{(S')}$ et si $h:Y'\to Y$ est la projection canonique, $u'_1$ et $u'_2$ se factorisent canoniquement en $u'_1=h\circ v_1$, $u'_2=h\circ v_2$, et $v_1$ et $v_2$ sont deux $S'$-morphismes dans $Y'$ qui coïncident dans $U'_3$.
On voit donc que lorsque $u_1$ parcourt la classe $\omega$, les $S'$-morphismes $v_1$ correspondants appartiennent à un même pseudo-morphisme relativement à $S'$, dit \emph{image réciproque} de $\omega$ par le morphisme changement de base $S'\to S$ et noté $\omega_{(S')}$.
Il est clair que si $S''\to S'$ est un second morphisme, on a $(\omega_{(S')})_{(S'')}=\omega_{(S'')}$ (pour le morphisme changement de base composé $S''\to S'\to S$).
\end{env}

\oldpage[IV-4]{252}
\subsection{Fonctions méromorphes relatives}
\EGAReaderTarget{ega4.c04.subsec.20-6}{IV.20.6}

\begin{env}[20.6.1]
\EGAReaderTarget{ega4.c04.statement.20-6-1}{IV.20.6.1}
Soient $S$ un préschéma, $X$ un $S$-préschéma qui est \emph{plat et localement de présentation finie} sur $S$ ; pour tout $s\in S$, nous noterons $X_s$ la fibre au point $s$ du morphisme structural $X\to S$.
En général, si $\varphi$ est une fonction méromorphe sur $X$, il n'est pas possible de lui associer de façon « naturelle », pour chaque $s\in S$, une fonction méromorphe « induite » sur $X_s$ \sref{IV.20.1.11}.
Pour tout ouvert $U$ de $X$, notons $T_{X/S}(U)$ l'ensemble des sections $t\in\Gamma(U,\sh{O}_X)$ telles que, pour tout $s\in S$, l'image $t_s$ de $t$ par l'homomorphisme canonique $\Gamma(U,\sh{O}_X)\to\Gamma(U\cap X_s,\sh{O}_{X_s})$ soit une section \emph{régulière} ; cela implique d'ailleurs, par l'équivalence de a) et b) dans \sref{IV.11.3.7}, que $t$ est elle-même une section \emph{régulière}.
Il est clair que $U\mapsto T_{X/S}(U)$ est un \emph{sous-faisceau} du faisceau d'ensemble $\sh{S}=\sh{S}(\sh{O}_X)$ (notation de \sref{IV.20.1.3}), que l'on note $\sh{T}_{X/S}$ ; on pose
\[
\label{IV.20.6.1.1}
  \sh{M}_{X/S}=\sh{O}_X[\sh{T}_{X/S}^{-1}]
  \tag{20.6.1.1}
\]
(notation de \sref{IV.20.1.1}) et on dit que ce faisceau d'anneaux est le \emph{faisceau des germes de fonctions méromorphes sur $X$ relatives à $S$} ; ses sections au-dessus de $X$ sont appelées les \emph{fonctions méromorphes sur $X$ relatives à $S$} et leur ensemble est noté $M(X/S)$.
Il est clair que $\sh{M}_{X/S}$ est un \emph{sous-faisceau} de $\sh{M}_X=\sh{O}_X[\sh{S}^{-1}]$ ; pour toute fonction méromorphe $\varphi\in M(X/S)$ et tout $s\in S$, l'image réciproque de $\varphi$ par le morphisme injection canonique $j_s:X_s\to X$ est alors définie \sref{IV.20.1.11}, et notée $\varphi_s$.
\end{env}

\begin{env}[20.6.2]
\EGAReaderTarget{ega4.c04.statement.20-6-2}{IV.20.6.2}
Soit maintenant $\sh{F}$ un $\sh{O}_X$-Module quasi-cohérent ; on pose
\[
\label{IV.20.6.2.1}
  \sh{M}_{X/S}(\sh{F})=\sh{F}[\sh{T}_{X/S}^{-1}]=\sh{F}\otimes_{\sh{O}_X}\sh{M}_{X/S};
  \tag{20.6.2.1}
\]
les sections de $\sh{M}_{X/S}(\sh{F})$ sont dites \emph{sections méromorphes de $\sh{F}$ au-dessus de $X$, relatives à $S$} et leur ensemble est noté $M(X/S,F)$.
L'homomorphisme canonique $\sh{F}\to\sh{M}_{X/S}(\sh{F})$ n'est pas nécessairement injectif ; lorsqu'il l'est, on dit que $\sh{F}$ est \emph{sans torsion relativement à $S$} : cela signifie que pour tout ouvert $U$ de $X$ et toute section $t\in T_{X/S}(U)$, $t$ est $(\sh{F}|U)$-régulière ; cette condition est remplie \emph{a fortiori} lorsque $\sh{F}$ est \emph{strictement sans torsion} \sref{IV.20.1.5}.
Dans ce dernier cas, il résulte aussitôt des définitions \sref{IV.20.1.2} que l'homomorphisme canonique de $\sh{O}_X$-Modules
\[
\label{IV.20.6.2.2}
  \sh{M}_{X/S}(\sh{F})\to\sh{M}_X(\sh{F})
  \tag{20.6.2.2}
\]
est \emph{injectif}, de sorte que les sections méromorphes de $\sh{F}$ relatives à $S$ sont des sections méromorphes de $\sh{F}$ au sens de \sref{IV.20.1.3}.
\end{env}

\begin{proposition}[20.6.3]
\label{IV.20.6.3}
L'image par l'homomorphisme injectif \sref{IV.20.2.10.1} du sous-faisceau $\sh{M}_{X/S}$ de $\sh{M}_X$ est le sous-faisceau $\sh{M}'_{X/S}$ des pseudo-fonctions sur $X$ relatives à $S$ (i.e.\ des pseudo-fonctions dont le domaine de définition est universellement schématiquement dense relativement à $S$ \sref{IV.20.5.4}).
\end{proposition}

\begin{proof}
On peut évidemment se borner à prouver que l'image de $M(X/S)$ par l'homomorphisme canonique $M(X)\to M'(X)$ est égale à $M'(X/S)$ ; la proposition est alors conséquence de la proposition plus générale suivante :
\end{proof}

\oldpage[IV-4]{253}
\begin{proposition}[20.6.4]
\label{IV.20.6.4}
Soit $\sh{F}$ un $\sh{O}_X$-Module quasi-cohérent de présentation finie et strictement sans torsion.
Alors, pour qu'une section méromorphe $\varphi$ de $\sh{F}$ au-dessus de $X$ soit une section méromorphe relative à $S$, il faut et il suffit que $\operatorname{dom}(\varphi)$ soit universellement schématiquement dense relativement à $S$.
\end{proposition}

\begin{proof}
La nécessité de la condition résulte de \sref{IV.20.2.15} appliqué à chaque injection canonique $X_s\to X$ ($s\in S$), compte tenu de \sref{IV.11.10.9}.
Pour voir que la condition est suffisante, il faut prouver que pour tout $x\in X$, il existe un voisinage ouvert $U$ de $x$ dans $X$ et une section de $\sh{M}_{X/S}(\sh{F})$ au-dessus de $U$ dont la restriction à $U\cap\operatorname{dom}(\varphi)$ coïncide avec $\varphi$ dans un ouvert schématiquement dense de $U\cap\operatorname{dom}(\varphi)$.
Considérons l'Idéal des dénominateurs $\sh{J}$ de $\varphi$ \sref{IV.20.2.14}, qui est quasi-cohérent, et qui définit un sous-préschéma fermé de $X$ dont l'espace sous-jacent est $X-\operatorname{dom}(\varphi)$.
Par hypothèse, si $s$ est l'image de $x$ dans $S$, $\operatorname{dom}(\varphi)\cap X_s$ est schématiquement dense dans le préschéma localement noethérien $X_s$, donc \sref{IV.20.2.13}[(iv)] contient $\operatorname{Ass}(\sh{O}_{X_s})$ ; cela implique donc que l'idéal $\sh{J}_x$ de $\sh{O}_{X,x}$ a une image dans $\sh{O}_{X_s,x}=\sh{O}_{X,x}\otimes_{\sh{O}_{S,s}}k(s)$ qui n'est contenue dans aucun des idéaux premiers $\mathfrak{p}_i\in\operatorname{Ass}(\sh{O}_{X_s,x})$ (en nombre fini) ; donc (Bourbaki, \emph{Alg.~comm.}, chap.~II, \textsection1, n\textsuperscript{o}~1, prop.~2) il existe un élément $t_x\in\sh{J}_x$ dont l'image dans $\sh{O}_{X_s,x}$ n'appartient à aucun des $\mathfrak{p}_i$, et est par suite régulière dans cet anneau noethérien.
Soit $t$ une section de $\sh{J}$ au-dessus d'un voisinage ouvert affine $U$ de $x$, dont $t_x$ soit le germe au point $x$ ; puisque $X$ est plat et localement de présentation finie sur $S$, on peut supposer \sref{IV.11.3.8} que $t$ est une section \emph{régulière} de $\sh{O}_X$ au-dessus de $U$ et que, pour tout $s'\in S$, l'image de $t$ dans $\Gamma(U\cap X_{s'},\sh{O}_{X_{s'}})$ soit aussi \emph{régulière} ; autrement dit, on a $t\in T_{X/S}(U)$.
Mais alors, par définition de $\sh{J}$, puisque $\sh{F}$ est strictement sans torsion, $t(\varphi|(U\cap\operatorname{dom}(\varphi)))$ est une section $u$ de $\sh{F}$ au-dessus de $U\cap\operatorname{dom}(\varphi)$ ; d'autre part, $U\cap\operatorname{dom}(\varphi)$ contient l'ensemble ouvert $U_t$ des points $x'\in U$ où $t(x')\neq0$, et ce dernier contient $x$ et est schématiquement dense dans $U$ \sref{IV.20.2.9}.
On voit donc que dans $U_t$, $\varphi$ coïncide avec la restriction à $U_t$ de la section $u/t$ de $\sh{M}_{X/S}(\sh{F})$ au-dessus de $U$.
C.Q.F.D.
\end{proof}

\begin{remarks}[20.6.5]
\EGAReaderTarget{ega4.c04.statement.20-6-5}{IV.20.6.5}
\begin{enumerate}[label=(\roman*)]
  \item Soit $\varphi$ une fonction méromorphe sur $X$ relative à $S$, de sorte que pour tout $s\in S$, $\varphi_s$ est une fonction méromorphe sur $X_s$ \sref{IV.20.6.1} ; en vertu de \sref{IV.20.1.11.1}, on a
  \[
  \label{IV.20.6.5.1}
    \operatorname{dom}(\varphi)\cap X_s\subset\operatorname{dom}(\varphi_s).
    \tag{20.6.5.1}
  \]
  Mais il convient de noter que même lorsque $S$ est le spectre d'un anneau de valuation discrète $A$ et $X=S[T]$ ($T$ indéterminée), les deux membres de \sref{IV.20.6.5.1} ne sont pas nécessairement égaux : par exemple, si $\pi$ est une uniformisante de $A$, il est immédiat que $\varphi=\pi/T$ est une fonction méromorphe sur $S$ relative à $S$, car si $a$ et $b$ sont le point fermé et le point générique de $S$, $T$ est régulière dans $\Gamma(X_a,\sh{O}_{X_a})=k[T]$ et dans $\Gamma(X_b,\sh{O}_{X_b})=K[T]$, $k$ et $K$ étant le corps résiduel et le corps des fractions de $A$.
  On a $\operatorname{dom}(\varphi)=\mathrm{D}(T)$ dans $X$, mais $\operatorname{dom}(\varphi_a)=X_a$ puisque $\varphi_a=0$.

  \item Pour qu'une fonction méromorphe $\varphi$ relative à $S$ soit \emph{inversible dans l'anneau} $M(X/S)$, il faut et il suffit que pour tout $s\in S$, $\varphi_s$ soit inversible dans $M(X_s)$ (autrement dit, que $\varphi_s$ soit une fonction méromorphe régulière sur $X_s$ \sref{IV.20.1.8}).
  La condition est
  \oldpage[IV-4]{254}
  en effet trivialement nécessaire.
  Inversement, si elle est vérifiée, et si $x$ est un point quelconque de $X$, $s$ son image dans $S$, il existe par hypothèse un voisinage ouvert $U$ de $x$ dans $X$ et deux sections $u$, $t$ de $\sh{O}_X$ au-dessus de $U$ telles que $t\in T_{X/S}(U)$ et $\varphi|U=u/t$ ; l'hypothèse entraîne que si $u_s$ est l'image de $u$ dans $\Gamma(U\cap X_s,\sh{O}_{X_s})$, $u_s$ est régulière au point $x$.
  En restreignant $U$, on peut donc supposer, en vertu de \sref{IV.11.3.8}, que $u\in T_{X/S}(U)$, d'où la conclusion.

  Lorsque $\varphi$ est inversible dans $M(X/S)$, on dit encore que $\varphi$ est une \emph{fonction méromorphe régulière relative à $S$}.
  On notera que $\varphi\in M(X/S)$ peut être inversible dans $M(X)$ (autrement dit, être \emph{régulière} au sens de \sref{IV.20.1.8}) sans l'être dans $M(X/S)$, comme le montre aussitôt l'exemple de~(i).

  \item Soit $\sh{L}$ un $\sh{O}_X$-Module \emph{inversible}, et soit $\varphi$ une section méromorphe \emph{régulière} de $\sh{L}$ au-dessus de $X$ \sref{IV.20.1.8} ; on dit que $\varphi$ est \emph{régulière relativement à $S$} si, pour tout ouvert $U$ de $X$ tel que $\sh{L}|U$ soit isomorphe à $\sh{O}_U$, $\varphi|U$ correspond à un élément de $\Gamma(U,\sh{M}_X)$ qui est \emph{régulière relativement à $S$} ; en vertu de~(ii), il est immédiat qu'il faut et il suffit pour cela que, pour tout $s\in S$, $\varphi_s$ soit une section méromorphe \emph{régulière} \sref{IV.20.1.8} du $\sh{O}_{X_s}$-Module inversible $\sh{L}_s=\sh{L}\otimes_{\sh{O}_X}k(s)$.
  Si $\varphi'$ est l'inverse de $\varphi$ dans $\sh{L}^{-1}$ \sref{IV.20.1.10}, $\varphi'$ est alors aussi régulière relative à $S$.
  Si $\sh{L}_1$ est un second $\sh{O}_X$-Module inversible, $\varphi_1$ une section méromorphe de $\sh{L}_1$ au-dessus de $X$, régulière relativement à $S$, alors $\varphi\otimes\varphi_1$ est une section méromorphe de $\sh{L}\otimes\sh{L}_1$ au-dessus de $X$, régulière relativement à $S$.
\end{enumerate}
\end{remarks}

\begin{proposition}[20.6.6]
\label{IV.20.6.6}
Soient $X$ un $S$-préschéma plat et localement de présentation finie sur $S$, $\sh{F}$ un $\sh{O}_X$-Module localement libre de type fini ; pour tout $s\in S$, on note $X_s$ la fibre au point $s$ du morphisme structural $f:X\to S$.
Soit $\varphi$ une section méromorphe de $\sh{F}$ au-dessus de $X$, relativement à $S$, et supposons que $\varphi$ soit définie en tous les points $x\in X$ tels que $\operatorname{prof}(\sh{O}_{X_{f(x)},x})=1$.
Alors $\varphi$ est partout définie.
\end{proposition}

\begin{proof}
Par hypothèse, $\operatorname{dom}(\varphi)\cap X_s$ est schématiquement dense dans $X_s$ pour tout $s\in S$, donc contient les points $x$ de $X_s$ tels que $\operatorname{prof}(\sh{O}_{X_s,x})=0$ \sref{IV.5.10.2} ; l'hypothèse signifie donc que si l'on pose $Z=X-\operatorname{dom}(\varphi)$, on a $\operatorname{prof}(\sh{O}_{X_{f(x)},x})\geq2$ en tous les points de $Z$.
Il suffit donc d'appliquer \sref{IV.19.9.8}.
\end{proof}

\begin{env}[20.6.7]
\EGAReaderTarget{ega4.c04.statement.20-6-7}{IV.20.6.7}
Soient $X$, $X'$ deux $S$-préschémas plats et localement de présentation finie sur $S$, $f=(\psi,\theta):X'\to X$ un $S$-morphisme.
Pour tout ouvert $U$ de $X$, notons $T_f(U)$ l'ensemble des sections $t\in T_{X/S}(U)$ dont l'image dans $\Gamma(f^{-1}(U),\sh{O}_{X'})$ appartienne à $T_{X'/S}(f^{-1}(U))$ ; il est immédiat que $U\mapsto T_f(U)$ est un sous-faisceau du faisceau d'ensembles $\sh{T}_{X/S}$, que l'on note $\sh{T}_f$.
On pose $\sh{M}_{X/S,f}=\sh{O}_X[\sh{T}_f^{-1}]$ ; c'est un sous-faisceau d'anneaux de $\sh{M}_{X/S}$, et on en déduit canoniquement de $\theta^\#: \psi^*(\sh{O}_X)\to\sh{O}_{X'}$ un homomorphisme de faisceaux d'anneaux $\theta'^\#: \psi^*(\sh{M}_{X/S,f})\to\sh{M}_{X'/S}$ prolongeant $\theta^\#$.
Si une fonction méromorphe $\varphi$ sur $X$, relative à $S$, est \emph{une section de $\sh{M}_{X/S,f}$}, $\Gamma(\theta'^\#)(\varphi)$ est une fonction méromorphe sur $X'$, dite \emph{image réciproque de $\varphi$ par $f$}, et notée $\varphi\circ f$ si cela n'entraîne pas confusion.
On étend aussitôt de même les définitions de \sref{IV.20.1.11} relatives aux $\sh{O}_X$-Modules.
\end{env}

\begin{proposition}[20.6.8]
\label{IV.20.6.8}
Avec les notations de \sref{IV.20.6.7}, si le $S$-morphisme $f:X'\to X$ est \emph{plat}, on a $\sh{M}_{X/S,f}=\sh{M}_{X/S}$, et l'homomorphisme $\varphi\mapsto\varphi\circ f$ est défini dans $M(X/S)$ tout entier.
\end{proposition}

\begin{proof}
\oldpage[IV-4]{255}
En effet, l'hypothèse entraîne, en vertu de \sref{IV.11.3.10}, que pour tout $s\in S$, $f_s:X'_s\to X_s$ est plat ; donc, pour toute section $t\in T_{X/S}(U)$, si $t'$ est son image réciproque dans $\Gamma(f^{-1}(U),\sh{O}_{X'})$, $t'_{s}$, qui est l'image réciproque de $t_s$, est une section \emph{régulière} de $\sh{O}_{X'_s}$ au-dessus de $f^{-1}(U)\cap X'_s$, en vertu de \sref{IV.20.1.12} ; on en conclut que par définition $t'\in T_{X'/S}(f^{-1}(U))$, d'où la proposition.
\end{proof}

On déduit de là, comme dans \sref{IV.20.1.12}, un homomorphisme canonique de $\sh{O}_{X'}$-Algèbres
\[
\label{IV.20.6.8.1}
  f^*(\sh{M}_{X/S})\to\sh{M}_{X'/S}.
  \tag{20.6.8.1}
\]

\begin{env}[20.6.9]
\label{IV.20.6.9}
Considérons enfin un morphisme \emph{quelconque} $S'\to S$, et posons $X'=X\times_S S'$, qui est plat et localement de présentation finie sur $S'$ ; soit $p:X'\to X$ la projection canonique.
Soient $U$ un ouvert de $X$, $t$ une section appartenant à $T_{X/S}(U)$, $t'$ son image dans $\Gamma(p^{-1}(U),\sh{O}_{X'})$ ; pour tout $s'\in S'$, si $s\in S$ est l'image de $s'$, on a $X'_{s'}=X_s\otimes_{k(s)}k(s')$, donc le morphisme $X'_{s'}\to X_s$ est plat, et par suite \sref{IV.20.1.12} l'image réciproque $t'_{s'}$ de $t_s$ dans $\Gamma(p^{-1}(U)\cap X'_{s'},\sh{O}_{X'_{s'}})$ est régulière ; ceci prouve que l'on a $t'\in T_{X'/S'}(p^{-1}(U))$.
Ceci permet de définir canoniquement, comme dans \sref{IV.20.6.8}, un homomorphisme canonique de $\sh{O}_{X'}$-Algèbres
\[
\label{IV.20.6.9.1}
  p^*(\sh{M}_{X/S})\to\sh{M}_{X'/S'}.
  \tag{20.6.9.1}
\]

Moyennant l'identification permise par \sref{IV.20.6.3}, cette notion de changement de base pour les fonctions méromorphes relatives est un cas particulier de la notion analogue pour les pseudo-morphismes relatifs \sref{IV.20.5.8}.
\end{env}

% FIN SOURCE ega4/ega4-20-fr.tex

%
% DEBUT SOURCE ega4/ega4-21a-fr.tex
\section{Diviseurs}
\label{section:IV.21}

Sur le contenu du présent paragraphe, voir les commentaires de l'introduction du \textsection~20.
Pour les propriétés \emph{globales} des diviseurs, le lecteur se reportera au paragraphe qui leur est consacré au chap.~V.

\subsection{Diviseurs sur un espace annelé}
\EGAReaderTarget{ega4.c04.subsec.21-1}{subsection:IV.21.1}
\label{IV.21.1}

\begin{env}[21.1.1]
\label{IV.21.1.1}
Soient $(X,\sh{O}_X)$ un espace annelé, $\sh{M}_X$ le faisceau des germes de fonctions méromorphes sur $X$ \sref{IV.20.1.3}, $\sh{M}_X^*$ le faisceau (de groupes multiplicatifs) des germes de fonctions méromorphes régulières sur $X$ \sref{IV.20.1.8}.
Il est clair que le faisceau (de groupes multiplicatifs) $\sh{O}_X^*$ des germes de sections inversibles de $\sh{O}_X$ s'identifie canoniquement à un sous-faisceau (de groupes multiplicatifs) de $\sh{M}_X^*$.
\end{env}

\begin{definition}[21.1.2]
\EGAReaderTarget{ega4.c04.statement.21-1-2}{IV.21.1.2}
On appelle \emph{faisceau des diviseurs sur $X$} et on note $\sh{D}iv_X$ le faisceau quotient (de groupes commutatifs) $\sh{M}_X^*/\sh{O}_X^*$;
les sections de ce faisceau au-dessus de $X$ se nomment les \emph{diviseurs sur $X$}; ils forment un groupe commutatif noté $\operatorname{Div}(X)$.
Pour toute section $f$ de $\sh{M}_X^*$ au-dessus de $X$ (autrement dit, toute fonction méromorphe régulière sur $X$ \sref{IV.20.1.8}) on appelle \emph{diviseur de $f$} et on note $\operatorname{div}(f)$ (ou $\operatorname{div}_X(f)$) le diviseur sur $X$ image de $f$ par l'homomorphisme canonique
\[
  \Gamma(X,\sh{M}_X^*)\longrightarrow\Gamma(X,\sh{D}iv_X).
\]
\end{definition}

\oldpage[IV-4]{256}

Le \emph{support} d'un diviseur $D$ est l'ensemble fermé des $x\in X$ tels que $D_x\neq0$.
On le note $\operatorname{Supp}(D)$.

Pour tout ouvert $U$ de $X$, on a évidemment $\sh{M}_X^*|U=\sh{M}_U^*$, $\sh{O}_X^*|U=\sh{O}_U^*$, donc $\sh{D}iv_X|U=\sh{D}iv_U$, et par suite le faisceau $\sh{D}iv_X$ est égal au préfaisceau $U\mapsto\operatorname{Div}(U)$.

Lorsque $X=\operatorname{Spec}(A)$ est affine, on écrit $\operatorname{Div}(A)$ au lieu de $\operatorname{Div}(\operatorname{Spec}(A))$.

\begin{env}[21.1.3]
\label{IV.21.1.3}
Nous noterons toujours \emph{additivement} le groupe $\operatorname{Div}(X)$ des diviseurs sur $X$.
Pour deux fonctions méromorphes régulières $f$, $g$ sur $X$, on a donc
\[
\label{IV.21.1.3.1}
  \operatorname{div}(fg)=\operatorname{div}(f)+\operatorname{div}(g),
  \tag{21.1.3.1}
\]
\[
\label{IV.21.1.3.2}
  \operatorname{div}(f^{-1})=-\operatorname{div}(f).
  \tag{21.1.3.2}
\]
Par définition, pour toute fonction méromorphe régulière $f$ sur $X$, on a l'équivalence
\[
\label{IV.21.1.3.3}
  \operatorname{div}(f)=0 \iff f\in\Gamma(X,\sh{O}_X^*),
  \tag{21.1.3.3}
\]
d'où, pour deux fonctions méromorphes régulières $f$, $g$ sur $X$
\[
\label{IV.21.1.3.4}
  \operatorname{div}(f)=\operatorname{div}(g) \iff fg^{-1}\in\Gamma(X,\sh{O}_X^*).
  \tag{21.1.3.4}
\]
\end{env}

\begin{env}[21.1.4]
\EGAReaderTarget{ega4.c04.statement.21-1-4}{IV.21.1.4}
Soit maintenant $\sh{L}$ un $\sh{O}_X$-Module inversible, et soit $s$ une \emph{section méromorphe régulière} \sref{IV.20.1.8} de $\sh{L}$ au-dessus de $X$.
Tout $x\in X$ possède un voisinage ouvert $U$ tel que $\sh{L}|U$ soit isomorphe à $\sh{O}_U$, donc $\sh{M}_X(\sh{L})|U$ isomorphe à $\sh{M}_X|U$;
par un de ces isomorphismes, $s|U$ correspond à une section $f\in\Gamma(U,\sh{M}_X^*)$, et puisque deux isomorphismes de $\sh{L}|U$ sur $\sh{O}_U$ ne diffèrent que par la multiplication par un élément de $\Gamma(U,\sh{O}_X^*)$ (\hyperref[0.5.4.3]{$\mathbf{0}_{\mathrm I}$, 5.4.3}), l'élément $\operatorname{div}_U(f)$ de $\Gamma(U,\sh{D}iv_X)$ est \emph{indépendant} de l'isomorphisme choisi;
il est clair que ces éléments (pour $U$ variable) sont les restrictions d'une section de $\sh{D}iv_X$ au-dessus de $X$, que l'on appelle le \emph{diviseur de $s$} et que l'on note $\operatorname{div}(s)$ (un tel diviseur n'est pas nécessairement de la forme $\operatorname{div}(g)$ pour une fonction méromorphe régulière $g$ sur $X$; voir \sref{IV.21.2.9}).
Pour $\sh{L}=\sh{O}_X$, la définition de $\operatorname{div}(s)$ coïncide avec celle de \sref{IV.21.1.2}.
Si $\sh{L}$, $\sh{L}'$ sont deux $\sh{O}_X$-Modules inversibles, $s$ (resp. $s'$) une section méromorphe régulière de $\sh{L}$ (resp. $\sh{L}'$) au-dessus de $X$, il est immédiat que l'on a
\[
\label{IV.21.1.4.1}
  \operatorname{div}(s\otimes s')=\operatorname{div}(s)+\operatorname{div}(s')
  \tag{21.1.4.1}
\]
\[
\label{IV.21.1.4.2}
  \operatorname{div}(s^{\otimes n})=n.\operatorname{div}(s) \quad\text{pour tout}\quad n\in\mathbf{Z}
  \tag{21.1.4.2}
\]
($s^{-1}$ étant la section méromorphe régulière de $\sh{L}^{-1}$ au-dessus de $X$ définie dans \sref{IV.20.1.10})
et, pour deux sections méromorphes régulières $s$, $s'$ de $\sh{L}$ au-dessus de $X$, on a la relation
\[
\label{IV.21.1.4.3}
  \operatorname{div}(s)=\operatorname{div}(s') \iff s'=ts \quad\text{avec}\quad t\in\Gamma(X,\sh{O}_X^*).
  \tag{21.1.4.3}
\]
\end{env}

\begin{env}[21.1.5]
\label{IV.21.1.5}
Le faisceau $\sh{S}(\sh{O}_X)$ \sref{IV.20.1.3} dont les sections au-dessus d'un ouvert $U$ de $X$ sont les éléments réguliers de $\Gamma(U,\sh{O}_X)$ est un \emph{sous-faisceau de monoïdes} de $\sh{M}_X^*$; on peut écrire
\[
\label{IV.21.1.5.1}
  \sh{S}(\sh{O}_X)=\sh{O}_X\cap\sh{M}_X^*.
  \tag{21.1.5.1}
\]
\end{env}

\oldpage[IV-4]{257}

Si on note $\sh{S}(\sh{O}_X)^{-1}$ le faisceau dont les sections au-dessus de $U$ sont les inverses dans $\Gamma(U,\sh{M}_X^*)$ des éléments de $\Gamma(U,\sh{S}(\sh{O}_X))$, il est clair que l'on a $\Gamma(U,\sh{S}(\sh{O}_X))\cap\Gamma(U,\sh{S}(\sh{O}_X)^{-1})=\Gamma(U,\sh{O}_X^*)$, donc
\[
\label{IV.21.1.5.2}
  \sh{S}(\sh{O}_X)\cap\sh{S}(\sh{O}_X)^{-1}=\sh{O}_X^*.
  \tag{21.1.5.2}
\]

\begin{definition}[21.1.6]
\EGAReaderTarget{ega4.c04.def.21-1-6}{IV.21.1.6}
Le sous-faisceau d'ensembles de $\sh{D}iv_X$, image canonique du sous-faisceau $\sh{S}(\sh{O}_X)$ de $\sh{M}_X^*$, est noté $\sh{D}iv_X^+$; ses sections au-dessus de $X$ sont appelées \emph{diviseurs positifs sur $X$}, et leur ensemble est noté $\operatorname{Div}^+(X)$.
\end{definition}

Puisque $\sh{S}(\sh{O}_X)$ est un faisceau de monoïdes (multiplicatifs), on a
\[
\label{IV.21.1.6.1}
  \operatorname{Div}^+(X)+\operatorname{Div}^+(X)\subset\operatorname{Div}^+(X)
  \tag{21.1.6.1}
\]
et d'autre part, en vertu de \sref{IV.21.1.5.2} et \sref{IV.21.1.3.3}
\[
\label{IV.21.1.6.2}
  \operatorname{Div}^+(X)\cap(-\operatorname{Div}^+(X))=\{0\}.
  \tag{21.1.6.2}
\]
Ces deux relations montrent que $\operatorname{Div}^+(X)$ est l'ensemble des éléments positifs pour une \emph{structure d'ordre} sur le groupe $\operatorname{Div}(X)$, compatible avec cette structure de groupe; on note cette relation d'ordre $D\leq D'$, autrement dit, on a
\[
\label{IV.21.1.6.3}
  D\geq0 \iff D\in\operatorname{Div}^+(X).
  \tag{21.1.6.3}
\]

Nous supposerons toujours dans ce qui suit que $\operatorname{Div}(X)$ est muni de cette structure d'ordre;
il est clair que $\sh{D}iv_X^+|U=\sh{D}iv_U^+$, donc $\Gamma(U,\sh{D}iv_X^+)=\operatorname{Div}^+(U)$, et on peut donc dire que $\sh{D}iv_X^+$ définit sur $\sh{D}iv_X$ une structure de faisceau de groupes ordonnés.
La fibre $(\sh{D}iv_X^+)_x$ en un point $x$ du faisceau $\sh{D}iv_X^+$ est un sous-monoïde du groupe $(\sh{D}iv_X)_x$, ensemble des éléments $\geq0$ pour une structure d'ordre compatible avec la structure de groupe;
pour un diviseur $D$ sur $X$, il revient au même de dire que $D\geq0$ ou que $D_x\geq0$ pour tout $x\in X$.

Par définition, pour toute fonction méromorphe régulière $f$ sur $X$, on a la relation
\[
\label{IV.21.1.6.4}
  \operatorname{div}(f)\geq0 \iff f\in\Gamma(X,\sh{S}(\sh{O}_X))
  \tag{21.1.6.4}
\]
autrement dit, $\operatorname{div}(f)\geq0$ signifie que $f$ est une \emph{section régulière} de $\sh{O}_X$, ou encore une section de $\sh{O}_X$ \emph{inversible dans} $M(X)$.

Plus généralement, étant donné un diviseur $D$ sur $X$, la relation $\operatorname{div}(f)\geq D$ est équivalente à la suivante : pour tout ouvert $U$ de $X$ tel que $D|U=\operatorname{div}(g)$, où $g\in\Gamma(U,\sh{M}_X^*)$, il existe un élément régulier $t$ de $\Gamma(U,\sh{O}_X)$ tel que $f|U=tg$.

\begin{env}[21.1.7]
\label{IV.21.1.7}
Soient $\sh{L}$ un $\sh{O}_X$-Module inversible, $s$ une section méromorphe régulière de $\sh{L}$ au-dessus de $X$; on a la relation
\[
\label{IV.21.1.7.1}
  \operatorname{div}(s)\geq0 \iff s\in\Gamma(X,\sh{L})\cap\Gamma(X,(\sh{M}_X(\sh{L}))^*)
  \tag{21.1.7.1}
\]
comme il résulte aussitôt des définitions \sref{IV.21.1.4} et \sref{IV.21.1.6}.
\end{env}

\begin{proposition}[21.1.8]
\EGAReaderTarget{ega4.c04.prop.21-1-8}{IV.21.1.8}
Soient $X$ un préschéma localement noethérien, $D$ un diviseur sur $X$.
Supposons que pour tout $x\in X$ tel que $\operatorname{prof}(\sh{O}_{X,x})=1$ on ait $D_x\geq0$ (resp. $D_x=0$).
Alors on a $D\geq0$ (resp. $D=0$).
\end{proposition}

\begin{proof}
La question étant locale sur $X$, on peut supposer que $D=\operatorname{div}(f)$, $f$ étant une
\oldpage[IV-4]{258}
fonction méromorphe régulière sur $X$;
la relation $D_x\geq0$ équivaut à $x\in\operatorname{dom}(f)$, donc l'hypothèse signifie que si $T=X-\operatorname{dom}(f)$, on a $\operatorname{prof}(\sh{O}_{X,x})\geq2$ pour tout $x\in T$ (car $\operatorname{dom}(f)$ contient les points maximaux de $X$).
Par suite \sref{IV.5.10.5} l'homomorphisme de restriction $\Gamma(X,\sh{O}_X)\to\Gamma(X-T,\sh{O}_X)$ est bijectif, ce qui montre qu'il existe une section $s$ de $\sh{O}_X$ au-dessus de $X$ telle que $f=s|(X-T)$.
Mais par définition de $T$, cela entraîne $T=\varnothing$, donc $f=s$ et $D\geq0$.
L'assertion relative à la relation $D_x=0$ s'en déduit aussitôt en appliquant ce qui précède à $-D$, en vertu de \sref{IV.21.1.6.2}.
\end{proof}


\begin{corollary}[21.1.9]
\label{IV.21.1.9}
Soient $X$ un préschéma localement noethérien, $D$ un diviseur sur $X$.
Soit $S$ le support de $D$.
Alors, pour tout point maximal $x$ de $S$, on a $\operatorname{prof}(\sh{O}_{X,x})=1$.
\end{corollary}

\begin{proof}
Posons en effet $X_1=\operatorname{Spec}(\sh{O}_{X,x})$;
compte tenu de \sref{IV.20.2.11} et \sref{IV.20.3.6}, le faisceau $\sh{M}_{X_1}^*$ est induit sur $X_1$ par $\sh{M}_X^*$, donc on peut se borner au cas où $X=X_1$, auquel cas, comme $x\in S$ et que $x$ est point maximal de $S$, on a nécessairement $S=\{x\}$.
Si l'on avait alors $\operatorname{prof}(\sh{O}_{X,x})\neq1$, on en conclurait, en vertu de \sref{IV.21.1.8}, que $D=0$, ce qui contredit la définition de $S$.
\end{proof}

\begin{proposition}[21.1.10]
\label{IV.21.1.10}
Soit $A$ un anneau local noethérien;
pour que $\operatorname{Div}(A)=0$, il faut et il suffit que $\operatorname{prof}(A)=0$ (autrement dit, que l'idéal maximal $\mathfrak{m}$ de $A$ soit associé à $A$ (\hyperref[0.16.4.6]{$\mathbf{0}$, 16.4.6})).
\end{proposition}

\begin{proof}
En effet, dire que $\operatorname{Div}(A)=0$ signifie que dans $A$ tout élément régulier est inversible, ou encore que tous les éléments de $\mathfrak{m}$ sont diviseurs de zéro, ce qui signifie que $\mathfrak{m}\in\operatorname{Ass}(A)$ (Bourbaki, \emph{Alg. comm.}, chap.~IV, \textsection~1, n\textsuperscript{o}~1, cor.~3 de la prop.~2).
\end{proof}

\subsection{Diviseurs et Idéaux fractionnaires inversibles}
\EGAReaderTarget{ega4.c04.subsec.21-2}{subsection:IV.21.2}

\begin{env}[21.2.1]
\EGAReaderTarget{ega4.c04.statement.21-2-1}{IV.21.2.1}
Soit $(X,\sh{O}_X)$ un espace annelé.
On appelle \emph{Idéal fractionnaire sur $X$} un sous-$\sh{O}_X$-Module du $\sh{O}_X$-Module $\sh{M}_X$ des germes de fonctions méromorphes sur $X$.
Un Idéal fractionnaire $\sh{J}$ sur $X$ qui est un $\sh{O}_X$-Module inversible est appelé \emph{Idéal fractionnaire inversible}.
\end{env}

\begin{proposition}[21.2.2]
\label{IV.21.2.2}
Pour qu'un Idéal fractionnaire $\sh{J}$ sur $X$ soit inversible, il faut et il suffit que pour tout $x\in X$, il existe un voisinage ouvert $U$ de $x$ et une section $f\in\Gamma(U,\sh{M}_X^*)$ tels que $\sh{J}|U=\sh{O}_U.f$.
\end{proposition}

\begin{proof}
La condition est évidemment suffisante, l'application $s\mapsto s(f|V)$ de $\Gamma(V,\sh{O}_X)$ dans $\Gamma(V,\sh{J})$ étant évidemment bijective pour tout ouvert $V\subset U$.
Pour voir qu'elle est nécessaire, notons qu'il existe par hypothèse un voisinage ouvert $U$ de $x$ et un isomorphisme de $\sh{O}_X$-Modules $\sh{O}_U\xrightarrow{\sim}\sh{J}|U$.
Si $f$ est l'image de la section $1\in\Gamma(U,\sh{O}_X)$ par cet isomorphisme, on peut supposer, en restreignant $U$, que l'on a $f=u/s$, où $u\in\Gamma(U,\sh{O}_X)$ et $s\in\Gamma(U,\sh{S}(\sh{O}_X))$, et l'isomorphisme considéré fait alors correspondre, à toute section $v\in\Gamma(V,\sh{O}_X)$ (où $V$ est un ouvert contenu dans $U$), la section $v(u|V)/(s|V)$;
dire que l'application ainsi définie est bijective signifie que $u|V$ est un élément régulier de $\Gamma(V,\sh{O}_X)$, donc $f\in\Gamma(U,\sh{M}_X^*)$.
\end{proof}

On notera que pour tout ouvert $U$ de $X$ tel que $\sh{J}|U=\sh{O}_U.f$ avec $f\in\Gamma(U,\sh{M}_X^*)$, la section $f$ est déterminée \emph{à la multiplication près} par un élément de $\Gamma(U,\sh{O}_X^*)$, puisque la
\oldpage[IV-4]{259}
multiplication par ces éléments fournit tous les automorphismes du $\sh{O}_U$-Module $\sh{O}_U$ (\hyperref[0.5.4.3]{$\mathbf{0}_{\mathrm I}$, 5.4.3}).

\begin{corollary}[21.2.3]
\label{IV.21.2.3}
\begin{enumerate}[label=(\roman*)]
  \item Soit $\sh{J}$ un Idéal fractionnaire inversible;
  alors le $\sh{O}_X$-Module inversible $\sh{J}^{-1}$ s'identifie canoniquement à l'Idéal fractionnaire $\sh{J}'$ (\emph{transporteur} de $\sh{J}$ dans $\sh{O}_X$) défini de la façon suivante : pour tout ouvert $U$ de $X$ tel que $\sh{J}|U=\sh{O}_U.f$, où $f\in\Gamma(U,\sh{M}_X^*)$, on a $\sh{J}'|U=\sh{O}_U.f^{-1}$.
  \item Si $\sh{J}_1$ et $\sh{J}_2$ sont deux Idéaux fractionnaires inversibles, l'application canonique $\sh{J}_1\otimes\sh{J}_2\to\sh{J}_1\sh{J}_2$ est un isomorphisme de $\sh{O}_X$-Modules.
\end{enumerate}
\end{corollary}

\begin{proof}
L'assertion (ii) résulte aussitôt de \sref{IV.21.2.2}.
D'autre part, la remarque faite à la fin de \sref{IV.21.2.2} prouve qu'il existe bien un Idéal fractionnaire $\sh{J}'$ et un seul défini par la condition de l'énoncé;
l'isomorphisme canonique de $\sh{J}'$ sur $\sh{J}^{-1}=\shHom_{\sh{O}_X}(\sh{J},\sh{O}_X)$ s'obtient en faisant correspondre à toute section $s(f^{-1}|V)$ de $\Gamma(V,\sh{J}')$ (où $V$ est un ouvert contenu dans $U$ et $s\in\Gamma(V,\sh{O}_X)$) l'homomorphisme $t(f|V)\mapsto st$ de $\Gamma(V,\sh{J})$ dans $\Gamma(V,\sh{O}_X)$.
\end{proof}

En vertu de \sref{IV.21.2.3},~(i), on identifiera généralement les $\sh{O}_X$-Modules inversibles $\sh{J}'$ et $\sh{J}^{-1}$, considérant donc $\sh{J}^{-1}$ comme un sous-$\sh{O}_X$-Module de $\sh{M}_X$.

\begin{env}[21.2.4]
\label{IV.21.2.4}
Il résulte de \sref{IV.21.2.3} que l'ensemble $\operatorname{Id.inv}(X)$ des Idéaux fractionnaires inversibles sur $X$ est muni d'une structure de \emph{groupe commutatif} pour la loi de composition $(\sh{J}_1,\sh{J}_2)\mapsto\sh{J}_1\sh{J}_2$, l'élément neutre de ce groupe étant $\sh{O}_X$.
Il est clair que pour tout ouvert $U$ de $X$, on a $(\sh{J}_1\sh{J}_2)|U=(\sh{J}_1|U)(\sh{J}_2|U)$, donc l'application de restriction $\sh{J}\mapsto\sh{J}|U$ est un homomorphisme de groupes $\operatorname{Id.inv}(X)\to\operatorname{Id.inv}(U)$;
on définit donc ainsi un \emph{préfaisceau de groupes commutatifs} $U\mapsto\operatorname{Id.inv}(U)$;
il est immédiat qu'en fait ce préfaisceau est un \emph{faisceau de groupes commutatifs}, que l'on note $\sh{J}d.inv_X$.
\end{env}

\begin{env}[21.2.5]
\label{IV.21.2.5}
Pour toute fonction méromorphe régulière $f\in\Gamma(X,\sh{M}_X^*)$, il résulte de \sref{IV.21.2.2} que $\sh{J}(f)=\sh{O}_X.f$ est un Idéal fractionnaire inversible, et on a évidemment $\sh{J}(f_1f_2)=\sh{J}(f_1)\sh{J}(f_2)$, autrement dit l'application $f\mapsto\sh{J}(f)$ est un homomorphisme du groupe commutatif $\Gamma(X,\sh{M}_X^*)$ dans le groupe commutatif $\operatorname{Id.inv}(X)$.
Remplaçant $X$ par un ouvert quelconque $U$ de $X$ et notant que les homomorphismes obtenus sont compatibles avec les opérations de restriction, on obtient un homomorphisme canonique de faisceaux de groupes commutatifs :
\[
\label{IV.21.2.5.1}
  I_0:\sh{M}_X^*\to\sh{J}d.inv_X.
  \tag{21.2.5.1}
\]
Notons que si $f\in\Gamma(X,\sh{O}_X^*)$, on a $\sh{J}(f)=\sh{O}_X$;
on en déduit aussitôt que l'homomorphisme $I_0$ se factorise en
\[
\label{IV.21.2.5.2}
  \sh{M}_X^*\to\sh{M}_X^*/\sh{O}_X^*=\sh{D}iv_X\xrightarrow{I}\sh{J}d.inv_X
  \tag{21.2.5.2}
\]
où $I$ est un homomorphisme du faisceau de groupes additifs $\sh{D}iv_X$ dans le faisceau de groupes multiplicatifs $\sh{J}d.inv_X$;
on a donc pour tout ouvert $U$ de $X$ un homomorphisme $\sh{J}_U:\operatorname{Div}(U)\to\operatorname{Id.inv}(U)$ de groupes commutatifs, tel que, pour toute section $f\in\Gamma(U,\sh{M}_X^*)$, on ait
\[
\label{IV.21.2.5.3}
  \sh{J}_U(\operatorname{div}_U(f))=\sh{O}_U.f.
  \tag{21.2.5.3}
\]

\oldpage[IV-4]{260}

On en conclut que $\sh{J}_X(D)$, pour tout diviseur $D\in\operatorname{Div}(X)$, est l'Idéal fractionnaire inversible défini de la façon suivante : pour tout ouvert $U$ de $X$ tel que $D|U=\operatorname{div}_U(f)$, où $f\in\Gamma(U,\sh{M}_X^*)$, $\sh{J}_X(D)|U$ est l'Idéal fractionnaire inversible $\sh{O}_U.f$.
On a donc, en vertu de \sref{IV.21.1.6}, pour toute fonction méromorphe régulière $f\in\Gamma(X,\sh{M}_X^*)$, la relation
\[
\label{IV.21.2.5.4}
  f\in\Gamma(X,\sh{J}_X(D)) \iff \operatorname{div}(f)\geq D.
  \tag{21.2.5.4}
\]
\end{env}

\begin{proposition}[21.2.6]
\label{IV.21.2.6}
L'homomorphisme $I:\sh{D}iv_X\to\sh{J}d.inv_X$ est bijectif.
\end{proposition}

\begin{proof}
On définit en effet un homomorphisme $I_X'$ de $\operatorname{Id.inv}(X)$ dans $\operatorname{Div}(X)$ en faisant correspondre à tout Idéal fractionnaire inversible $\sh{J}$ sur $X$ le diviseur $I_X'(\sh{J})$ défini de la façon suivante : pour tout ouvert $U$ de $X$ tel que $\sh{J}|U=\sh{O}_U.f$, où $f\in\Gamma(U,\sh{M}_X^*)$ \sref{IV.21.2.2}, on prend $I_X'(\sh{J})|U=\operatorname{div}_U(f)$;
en vertu de la remarque suivant \sref{IV.21.2.2}, cette définition est indépendante du générateur $f$ choisi dans $\sh{J}|U$, et détermine bien un diviseur sur $X$.
En outre, cette définition montre aussitôt que les homomorphismes $\sh{J}_X$ et $I_X'$ sont réciproques l'un de l'autre.
Remplaçant $X$ par un ouvert quelconque $U$, on en déduit la définition de l'isomorphisme $I':\sh{J}d.inv_X\to\sh{D}iv_X$ réciproque de $I$, d'où la proposition.
On posera $I_X'(\sh{J})=\operatorname{div}(\sh{J})$, de sorte que l'on a, pour toute fonction méromorphe régulière $f$ sur $X$
\[
\label{IV.21.2.6.1}
  \operatorname{div}(\sh{O}_X.f)=\operatorname{div}(f).
  \tag{21.2.6.1}
\]
\end{proof}

\begin{env}[21.2.7]
\EGAReaderTarget{ega4.c04.statement.21-2-7}{IV.21.2.7}
On identifiera souvent les faisceaux $\sh{D}iv_X$ et $\sh{J}d.inv_X$ (resp. les groupes $\operatorname{Div}(X)$ et $\operatorname{Id.inv}(X)$) au moyen des isomorphismes $I$ et $I'$ (resp. $\sh{J}_X$ et $I_X'$) précédents.
On notera que l'on a la relation
\[
\label{IV.21.2.7.1}
  D\geq0 \iff \sh{J}_X(D)\subset\sh{O}_X \quad\text{pour}\quad D\in\operatorname{Div}(X)
  \tag{21.2.7.1}
\]
comme il résulte aussitôt des définitions \sref{IV.21.1.6} et de \sref{IV.21.1.5.1};
en d'autres termes, l'image $\sh{J}_X(\operatorname{Div}^+(X))$ est l'ensemble des \emph{Idéaux} de $\sh{O}_X$ (dits aussi parfois \emph{Idéaux entiers}) qui sont des $\sh{O}_X$-Modules inversibles : un tel Idéal $\sh{J}$ est encore caractérisé par le fait que pour tout $x\in X$, il y a un voisinage ouvert $U$ de $x$ dans $X$ tel que $\sh{J}|U=\sh{O}_U.f$, où $f$ est un élément régulier de $\Gamma(U,\sh{O}_X)$.
L'ensemble $\sh{J}_X(\operatorname{Div}^+(X))$ de ces Idéaux est donc un sous-monoïde de $\operatorname{Id.inv}(X)$, égal à l'ensemble des éléments positifs pour une relation d'ordre compatible avec la structure de groupe de $\operatorname{Id.inv}(X)$, et il est immédiat que cette relation n'est autre que la relation \emph{opposée} à l'inclusion;
autrement dit, on a
\[
\label{IV.21.2.7.2}
  \sh{J}_1\subset\sh{J}_2 \iff \operatorname{div}(\sh{J}_1)\geq\operatorname{div}(\sh{J}_2)
  \tag{21.2.7.2}
\]
pour $\sh{J}_1$, $\sh{J}_2$ dans $\operatorname{Id.inv}(X)$.
\end{env}

\begin{env}[21.2.8]
\label{IV.21.2.8}
Pour tout diviseur $D$ sur $X$, on pose
\[
\label{IV.21.2.8.1}
  \sh{O}_X(D)=(\sh{J}_X(D))^{-1};
  \tag{21.2.8.1}
\]
$\sh{O}_X(D)$ est donc un Idéal fractionnaire inversible, défini de la façon suivante : pour tout ouvert $U$ de $X$ tel que $D|U=\operatorname{div}_U(f)$, où $f\in\Gamma(U,\sh{M}_X^*)$, $\sh{O}_X(D)|U$ est l'Idéal
\oldpage[IV-4]{261}
fractionnaire inversible $\sh{O}_U.f^{-1}$;
en vertu de \sref{IV.21.1.6}, pour toute fonction méromorphe régulière $f$, on a la relation
\[
\label{IV.21.2.8.2}
  f\in\Gamma(X,\sh{O}_X(D)) \iff \operatorname{div}(f)\geq-D.
  \tag{21.2.8.2}
\]
En outre, il est clair que l'on a des isomorphismes canoniques \sref{IV.21.2.3}
\[
\label{IV.21.2.8.3}
  \left\{\begin{aligned}
    &\sh{O}_X(0)=\sh{O}_X,\quad \sh{O}_X(D+D')=\sh{O}_X(D)\sh{O}_X(D')\xrightarrow{\sim}\sh{O}_X(D)\otimes\sh{O}_X(D')\\
    &\sh{O}_X(nD)=(\sh{O}_X(D))^n\xrightarrow{\sim}(\sh{O}_X(D))^{\otimes n}
  \end{aligned}\right.
  \tag{21.2.8.3}
\]
pour tout entier $n\in\mathbf{Z}$, et pour deux diviseurs quelconques $D$, $D'$ sur $X$.
\end{env}

\begin{env}[21.2.9]
\label{IV.21.2.9}
Soit $\sh{J}$ un Idéal fractionnaire inversible sur $X$.
L'injection canonique $\sh{J}\hookrightarrow\sh{M}_X$ définit par tensorisation un homomorphisme de $\sh{O}_X$-Modules
\[
\label{IV.21.2.9.1}
  \sh{M}_X(\sh{J})=\sh{J}\otimes_{\sh{O}_X}\sh{M}_X\to\sh{M}_X\otimes_{\sh{O}_X}\sh{M}_X=\sh{M}_X
  \tag{21.2.9.1}
\]
qui est un \emph{isomorphisme} : en effet, si $U$ est un ouvert de $X$ tel que $\sh{J}|U=\sh{O}_U.f$, où $f\in\Gamma(U,\sh{M}_X^*)$, l'homomorphisme \sref{IV.21.2.9.1} restreint à $U$ n'est autre que l'isomorphisme qui à toute section $t$ de $\sh{M}_X(\sh{J})|U=\sh{M}_X|U$ au-dessus de $V\subset U$, fait correspondre la section $t(f|V)^{-1}$ du même faisceau.
Dans l'isomorphisme \sref{IV.21.2.9.1}, aux sections méromorphes régulières de $\sh{J}$ au-dessus de $X$ correspondent les fonctions méromorphes régulières sur $X$.

Considérons en particulier le cas où $\sh{J}=\sh{O}_X(D)$, où $D$ est un diviseur sur $X$;
on a alors un isomorphisme canonique
\[
\label{IV.21.2.9.2}
  \sh{M}_X(\sh{O}_X(D))\xrightarrow{\sim}\sh{M}_X
  \tag{21.2.9.2}
\]
et on désigne par $s_D$ la section méromorphe régulière de $\sh{O}_X(D)$ au-dessus de $X$ qui correspond par cet isomorphisme à la section $1$ de $\sh{M}_X$.
Si $U$ est un ouvert de $X$ tel que $\sh{O}_X(D)|U=\sh{O}_U.f^{-1}$, où $f\in\Gamma(U,\sh{M}_X^*)$, on a $s_D|U=1$ dans $\Gamma(U,\sh{M}_X)$;
comme alors on a $D|U=\operatorname{div}_U(f)$, on en déduit \sref{IV.21.1.4} que l'on a
\[
\label{IV.21.2.9.3}
  \operatorname{div}(s_D)=D.
  \tag{21.2.9.3}
\]
D'autre part, on déduit aussitôt des isomorphismes canoniques \sref{IV.21.2.8.3} les formules
\[
\label{IV.21.2.9.4}
  s_0=1,\quad s_{D+D'}=s_D\otimes s_{D'},\quad s_{nD}=s_D^{\otimes n}\quad(n\in\mathbf{Z}).
  \tag{21.2.9.4}
\]
\end{env}

\begin{env}[21.2.10]
\label{IV.21.2.10}
Considérons, entre deux couples $(\sh{L},s)$, $(\sh{L}',s')$, où $\sh{L}$ et $\sh{L}'$ sont deux $\sh{O}_X$-Modules inversibles, $s$ (resp. $s'$) une section méromorphe régulière de $\sh{L}$ (resp. $\sh{L}'$) au-dessus de $X$, la relation : « il existe un isomorphisme $u:\sh{L}\xrightarrow{\sim}\sh{L}'$ tel que $\overline{u}(s)=s'$ », où $\overline{u}:\Gamma(X,\sh{M}_X(\sh{L}))\xrightarrow{\sim}\Gamma(X,\sh{M}_X(\sh{L}'))$ est l'isomorphisme déduit de $u$ (on notera que l'isomorphisme $u$ vérifiant cette condition est alors déterminé de façon unique).
Il est clair que c'est une relation d'équivalence, et puisqu'il existe un ensemble de $\sh{O}_X$-Modules inversibles tel que tout $\sh{O}_X$-Module inversible soit isomorphe à un élément de cet ensemble (\hyperref[0.5.4.7]{$\mathbf{0}_{\mathrm I}$, 5.4.7}), on peut parler de l'ensemble $D(X)$ des classes d'équivalence des couples $(\sh{L},s)$ pour la relation précédente.
Pour tout $\sh{O}_X$-Module inversible $\sh{L}$ et toute section
\oldpage[IV-4]{262}
méromorphe régulière $s$ de $\sh{L}$ au-dessus de $X$, on notera $\operatorname{cl}(\sh{L},s)$ l'élément de $D(X)$ correspondant au couple $(\sh{L},s)$.
Il résulte de (\hyperref[0.5.4.3]{$\mathbf{0}_{\mathrm I}$, 5.4.3}) que si $s$, $s'$ sont deux sections méromorphes régulières de $\sh{L}$ au-dessus de $X$, la relation $\operatorname{cl}(\sh{L},s)=\operatorname{cl}(\sh{L},s')$ équivaut à l'existence d'une section $t\in\Gamma(X,\sh{O}_X^*)$ telle que $s'=ts$.

Il est immédiat que si $(\sh{L},s)$ est équivalent à $(\sh{L}_1,s_1)$ et $(\sh{L}',s')$ à $(\sh{L}_1',s_1')$, les couples $(\sh{L}\otimes\sh{L}',s\otimes s')$ et $(\sh{L}_1\otimes\sh{L}_1',s_1\otimes s_1')$ sont équivalents;
on définit donc dans $D(X)$ une loi de composition en posant
\[
\label{IV.21.2.10.1}
  \operatorname{cl}(\sh{L},s)\operatorname{cl}(\sh{L}',s')=\operatorname{cl}(\sh{L}\otimes\sh{L}',s\otimes s');
  \tag{21.2.10.1}
\]
il est immédiat que c'est une loi de groupe commutatif, dont l'élément neutre est $\operatorname{cl}(\sh{O}_X,1)$ et où l'inverse de $\operatorname{cl}(\sh{L},s)$ est $\operatorname{cl}(\sh{L}^{-1},s^{\otimes(-1)})$.
\end{env}

\begin{proposition}[21.2.11]
\EGAReaderTarget{ega4.c04.prop.21-2-11}{IV.21.2.11}
Les applications
\[
\label{IV.21.2.11.1}
  D\mapsto\operatorname{cl}(\sh{O}_X(D),s_D),\qquad \operatorname{cl}(\sh{L},s)\mapsto\operatorname{div}(s)
  \tag{21.2.11.1}
\]
sont des isomorphismes réciproques de $\operatorname{Div}(X)$ sur $D(X)$ et de $D(X)$ sur $\operatorname{Div}(X)$ respectivement.
\end{proposition}

\begin{proof}
Compte tenu de \sref{IV.21.2.8.3}, \sref{IV.21.2.9.4} et \sref{IV.21.2.10.1}, il suffit de voir que les composés de ces deux applications sont les identités dans $\operatorname{Div}(X)$ et $D(X)$ respectivement.
La première assertion n'est autre que \sref{IV.21.2.9.3}.
D'autre part, soit $D=\operatorname{div}(s)$, où $s$ est une section méromorphe régulière de $\sh{L}$ au-dessus de $X$, et soit $U$ un ouvert de $X$ tel qu'il existe un isomorphisme de $\sh{L}|U$ sur $\sh{O}_U$ transformant $s|U$ en $f\in\Gamma(U,\sh{M}_X^*)$, de sorte que $D|U=\operatorname{div}_U(f)$, $\sh{O}_X(D)|U=\sh{O}_U.f^{-1}$ et $s_D|U$ est l'élément unité de $\Gamma(U,\sh{M}_X)$.
Il y a donc un isomorphisme $v_U:\sh{L}|U\to\sh{O}_U.f^{-1}=\sh{O}_X(D)|U$ tel que $\overline{v}_U$ (notation de \sref{IV.21.2.10}) transforme $s|U$ en $f.f^{-1}=1$;
on voit aussitôt que ces isomorphismes sont compatibles avec les opérations de restriction, donc définissent un isomorphisme $v:\sh{L}\to\sh{O}_X(D)$ tel que $\overline{v}(s)=s_D$.
\end{proof}

On peut transporter par le premier des isomorphismes \sref{IV.21.2.11.1} la structure de groupe ordonné de $\operatorname{Div}(X)$ à $D(X)$;
les éléments $\geq0$ de $D(X)$ sont donc les classes $\operatorname{cl}(\sh{L},s)$ telles que $\operatorname{div}(s)\geq0$, c'est-à-dire \sref{IV.21.1.7.1} telles que
\[
  s\in\Gamma(X,\sh{L})\cap\Gamma(X,(\sh{M}_X(\sh{L}))^*).
\]

\begin{env}[21.2.12]
\label{IV.21.2.12}
Soit $D$ un diviseur \emph{positif} sur un \emph{préschéma} $X$;
l'Idéal fractionnaire $\sh{J}_X(D)$ est donc un Idéal de $\sh{O}_X$, qui est un $\sh{O}_X$-Module inversible;
soit $Y(D)$ le sous-préschéma fermé de $X$ qu'il définit.
Pour tout $x\in Y(D)$, il y a par hypothèse un voisinage ouvert $U$ de $x$ dans $X$ et une section régulière $t\in\Gamma(U,\sh{O}_X)$ telle que $\sh{J}_X(D)|U=\sh{O}_U.t$ \sref{IV.21.2.7};
en d'autres termes, l'immersion canonique $Y(D)\to X$ est \emph{régulière et de codimension $1$} en tout point de $Y(D)$ \sref{IV.19.1.4}.
Inversement, si $Y$ est un sous-préschéma fermé de $X$, régulièrement immergé dans $X$ et de codimension $1$ en tout point de $Y$, il existe un diviseur positif et un seul $D$ tel que $Y(D)=Y$, car pour tout $x\in Y$, il y a un voisinage $U$ de $x$ dans $X$ tel que $Y\cap U$ soit défini par un Idéal de $\sh{O}_U$ de la forme $\sh{O}_U.t$, où $t$ est régulier dans $\Gamma(U,\sh{O}_X)$.
On notera que l'on a alors $\operatorname{Supp}(D)=Y(D)$, car dire que $D_x\neq0$ signifie que $t_x$ (avec les notations ci-dessus) n'est pas inversible, c'est-à-dire que $x\in Y(D)$.
\end{env}

\oldpage[IV-4]{263}

\subsection{Équivalence linéaire des diviseurs}
\EGAReaderTarget{ega4.c04.subsec.21-3}{subsection:IV.21.3}

\begin{env}[21.3.1]
\EGAReaderTarget{ega4.c04.statement.21-3-1}{IV.21.3.1}
On dit qu'un diviseur $D$ sur $X$ est \emph{principal} s'il est de la forme $\operatorname{div}(f)$, où $f$ est une fonction méromorphe régulière sur $X$;
les fonctions méromorphes régulières $f'$ telles que $\operatorname{div}(f')=D$ sont alors toutes celles de la forme $tf$, où $t\in\Gamma(X,\sh{O}_X^*)$ \sref{IV.21.1.3.4}.
L'ensemble des diviseurs principaux est un sous-groupe de $\operatorname{Div}(X)$, noté $\operatorname{Div.princ}(X)$, isomorphe à $\Gamma(X,\sh{M}_X^*)/\Gamma(X,\sh{O}_X^*)$.
Deux diviseurs $D$, $D'$ sont dits \emph{linéairement équivalents} si $D-D'$ est un diviseur principal;
les diviseurs principaux sont donc les diviseurs linéairement équivalents à $0$.
\end{env}

\begin{env}[21.3.2]
\label{IV.21.3.2}
Rappelons (\hyperref[0.5.4.7]{$\mathbf{0}_{\mathrm I}$, 5.4.7}) que l'on peut parler de l'ensemble des classes d'équivalence des $\sh{O}_X$-Modules inversibles pour la relation d'isomorphie;
on désigne cet ensemble par $\operatorname{Pic}(X)$, et pour tout $\sh{O}_X$-Module inversible $\sh{L}$, on note $\operatorname{cl}(\sh{L})$ la classe d'équivalence des $\sh{O}_X$-Modules isomorphes à $\sh{L}$;
en outre, $\operatorname{Pic}(X)$ est un groupe commutatif pour la multiplication définie par $\operatorname{cl}(\sh{L})\operatorname{cl}(\sh{L}')=\operatorname{cl}(\sh{L}\otimes\sh{L}')$.
Il est clair que l'application
\[
\label{IV.21.3.2.1}
  r:\operatorname{cl}(\sh{L},s)\mapsto\operatorname{cl}(\sh{L})
  \tag{21.3.2.1}
\]
est un homomorphisme du groupe $D(X)$ \sref{IV.21.2.10} dans le groupe $\operatorname{Pic}(X)$.
Par composition, on en déduit donc un homomorphisme
\[
\label{IV.21.3.2.2}
  l:\operatorname{Div}(X)\xrightarrow{\sim}D(X)\xrightarrow{r}\operatorname{Pic}(X)
  \tag{21.3.2.2}
\]
(aussi noté $l_X$) tel que, pour tout diviseur $D$, on ait
\[
\label{IV.21.3.2.3}
  l(D)=\operatorname{cl}(\sh{O}_X(D)).
  \tag{21.3.2.3}
\]

Notons enfin que, si $u:X'\to X$ est un morphisme d'espaces annelés, $\sh{L}_1$, $\sh{L}_2$ deux $\sh{O}_X$-Modules inversibles isomorphes, les $\sh{O}_{X'}$-Modules inversibles $u^*(\sh{L}_1)$ et $u^*(\sh{L}_2)$ (\hyperref[0.5.4.5]{$\mathbf{0}_{\mathrm I}$, 5.4.5}) sont isomorphes;
comme d'autre part, pour deux $\sh{O}_X$-Modules inversibles quelconques $\sh{L}_1$, $\sh{L}_2$, on a $u^*(\sh{L}_1\otimes\sh{L}_2)=u^*(\sh{L}_1)\otimes u^*(\sh{L}_2)$ à isomorphisme canonique près (\hyperref[0.4.3.3]{$\mathbf{0}_{\mathrm I}$, 4.3.3}), on voit que le morphisme $u$ définit canoniquement un homomorphisme de groupes commutatifs
\[
\label{IV.21.3.2.4}
  \operatorname{Pic}(u):\operatorname{Pic}(X)\to\operatorname{Pic}(X').
  \tag{21.3.2.4}
\]
\end{env}

\begin{proposition}[21.3.3]
\label{IV.21.3.3}
\begin{enumerate}[label=(\roman*)]
  \item Le noyau de l'homomorphisme canonique $l:\operatorname{Div}(X)\to\operatorname{Pic}(X)$ est le sous-groupe $\operatorname{Div.princ}(X)$, autrement dit, pour que $\sh{O}_X(D)$ et $\sh{O}_X(D')$ soient isomorphes, il faut et il suffit que $D$ et $D'$ soient linéairement équivalents.
  On a donc un homomorphisme injectif canonique
  \[
  \label{IV.21.3.3.1}
    \operatorname{Div}(X)/\operatorname{Div.princ}(X)\to\operatorname{Pic}(X)
    \tag{21.3.3.1}
  \]
  déduit de $l$.
  \item Pour qu'un $\sh{O}_X$-Module inversible $\sh{L}$ soit tel que $\operatorname{cl}(\sh{L})$ soit de la forme $l(D)$, ou encore pour que $\sh{L}$ soit isomorphe à un $\sh{O}_X$-Module de la forme $\sh{O}_X(D)$, il faut et il suffit qu'il existe une section méromorphe régulière de $\sh{L}$.
\end{enumerate}
\end{proposition}

\begin{proof}
\oldpage[IV-4]{264}
La proposition résulte aussitôt des définitions et de \sref{IV.21.2.10}.
\end{proof}

\begin{proposition}[21.3.4]
\label{IV.21.3.4}
Soit $X$ un préschéma vérifiant l'une des deux hypothèses suivantes :
\begin{enumerate}[label=\emph{\alph*)}]
  \item $X$ est localement noethérien et $\operatorname{Ass}(\sh{O}_X)$ est contenu dans un ouvert affine de $X$.
  \item $X$ est réduit et l'ensemble de ses composantes irréductibles est localement fini.
\end{enumerate}
Alors l'homomorphisme canonique $l:\operatorname{Div}(X)\to\operatorname{Pic}(X)$ est surjectif, et donne par passage au quotient un isomorphisme
\[
  \operatorname{Div}(X)/\operatorname{Div.princ}(X)\xrightarrow{\sim}\operatorname{Pic}(X).
\]
\end{proposition}

\begin{proof}
Il suffit de montrer que tout $\sh{O}_X$-Module inversible $\sh{L}$ admet une section méromorphe régulière au-dessus de $X$ \sref{IV.21.3.3}.
Dans les deux cas, il suffira, grâce à \sref{IV.20.2.11},~(ii), de définir une section $s$ de $(\sh{M}_X(\sh{L}))^*$ au-dessus d'un ouvert schématiquement dense $U$ de $X$, ou encore dans le cas a), d'un ouvert $U$ contenant $\operatorname{Ass}(\sh{O}_X)$ \sref{IV.20.2.13},~(iv).
En effet, soit alors $V$ un ouvert quelconque de $X$ tel que $\sh{L}|V$ soit isomorphe à $\sh{O}_V$, de sorte qu'il existe un isomorphisme de $\sh{M}_X(\sh{L})|V$ sur $\sh{M}_X|V$, transformant $s|(U\cap V)$ en une section $f_V$ de $\sh{M}_X^*$ au-dessus de $U\cap V$.
Comme $U\cap V$ est schématiquement dense dans $V$, il résulte de \sref{IV.20.2.11} qu'il existe une fonction méromorphe régulière et une seule $g_V$ dans $V$ telle que $g_V|(U\cap V)=f_V$, et cette section correspond donc par l'isomorphisme considéré à une section méromorphe régulière $u_V$ de $\sh{L}$ au-dessus de $V$ telle que $u_V|(U\cap V)=s|(U\cap V)$.
En outre, si $V'$ est un second ouvert de $X$ tel que $\sh{L}|V'$ soit isomorphe à $\sh{O}_{V'}$, les restrictions de $u_V$ et $u_{V'}$ à $V\cap V'$ sont égales, car elles correspondent par isomorphisme à deux fonctions méromorphes qui coïncident dans un ouvert schématiquement dense $U\cap V\cap V'$, et on conclut encore par \sref{IV.20.2.11};
les $u_V$ sont donc bien les restrictions d'une section de $(\sh{M}_X(\sh{L}))^*$ au-dessus de $X$.

Cela étant, dans le cas b), on prend pour chacun des points maximaux $x_\lambda$ de $X$ un ouvert $U_\lambda$ contenant $x_\lambda$, ne rencontrant aucune composante irréductible de $X$ distincte de $\overline{\{x_\lambda\}}$ et tel que $\sh{L}|U_\lambda$ soit isomorphe à $\sh{O}_{U_\lambda}$;
on prendra pour $s$ la section telle que $s|U_\lambda$ soit la section de $\sh{L}|U_\lambda$ qui correspond par l'isomorphisme précédent à la section unité de $\sh{O}_{U_\lambda}$.

Dans le cas a), on peut prendre pour $U$ par hypothèse un ouvert affine (donc noethérien), autrement dit, supposer que $X=\operatorname{Spec}(A)$, où $A$ est noethérien, et $\sh{L}=\widetilde{P}$, où $P$ est un $A$-module projectif de rang $1$.
Si $S$ est l'ensemble des éléments réguliers de $A$, on a $\Gamma(X,\sh{M}_X)=S^{-1}A$ \sref{IV.20.2.12} et $\Gamma(X,\sh{M}_X(\sh{L}))=S^{-1}P$;
or $S$ est l'ensemble des éléments n'appartenant à aucun des idéaux associés à $A$, donc $S^{-1}A$ est un anneau semi-local dont les idéaux maximaux proviennent des éléments maximaux de $\operatorname{Ass}(A)$, et $S^{-1}P$ est un $S^{-1}A$-module projectif de rang $1$, donc ici libre de rang $1$ (Bourbaki, \emph{Alg. comm.}, chap.~II, \textsection~5, n\textsuperscript{o}~3, prop.~5);
un élément formant une base de ce $S^{-1}A$-module est par suite \sref{IV.20.1.8} une section méromorphe régulière de $\sh{L}$ au-dessus de $X$.
\end{proof}

\begin{corollary}[21.3.5]
\label{IV.21.3.5}
Si $X$ est un préschéma noethérien, tel qu'il existe un $\sh{O}_X$-Module inversible ample (\textbf{II}, 4.5.3) (par exemple un préschéma quasi-projectif sur un spectre d'anneau noethérien (\textbf{II}, 5.3.1 et 4.6.6)), l'homomorphisme canonique $l:\operatorname{Div}(X)\to\operatorname{Pic}(X)$ est surjectif.
\end{corollary}

\begin{proof}
\oldpage[IV-4]{265}
En effet (\textbf{II}, 4.5.4), il existe alors un voisinage ouvert affine de l'ensemble fini $\operatorname{Ass}(\sh{O}_X)$.
\end{proof}

\begin{remark}[21.3.6]
\label{IV.21.3.6}
Rappelons (\hyperref[0.5.4.7]{$\mathbf{0}_{\mathrm I}$, 5.4.7}) que l'on a un isomorphisme canonique $\pi:\mathrm{H}^1(X,\sh{O}_X^*)\xrightarrow{\sim}\operatorname{Pic}(X)$ défini de la façon suivante.
On part d'un $1$-cocycle $(c_{\alpha\beta})$ à valeurs dans $\sh{O}_X^*$ correspondant à un recouvrement ouvert $(U_\alpha)$ de $X$, $c_{\alpha\beta}$ étant un élément de $\Gamma(U_\alpha\cap U_\beta,\sh{O}_X^*)$, et on lui associe la classe du $\sh{O}_X$-Module inversible obtenu en recollant les $\sh{O}_{U_\alpha}$ suivant les isomorphismes $\sh{O}_{U_\alpha}|(U_\alpha\cap U_\beta)\xrightarrow{\sim}\sh{O}_{U_\beta}|(U_\alpha\cap U_\beta)$ définis par la multiplication par $c_{\alpha\beta}$.
D'autre part, on déduit de la suite exacte de faisceaux de groupes commutatifs
\[
  1\to\sh{O}_X^*\to\sh{M}_X^*\to\sh{D}iv_X\to0
\]
l'homomorphisme bord de la suite exacte de cohomologie
\[
\label{IV.21.3.6.1}
  \partial:\operatorname{Div}(X)\to\mathrm{H}^1(X,\sh{O}_X^*).
  \tag{21.3.6.1}
\]
Montrons que l'homomorphisme composé
\[
  \operatorname{Div}(X)\xrightarrow{\partial}\mathrm{H}^1(X,\sh{O}_X^*)\xrightarrow{\pi}\operatorname{Pic}(X)
\]
n'est autre que l'homomorphisme $l$ défini dans \sref{IV.21.3.2.2}.
En effet on doit partir d'un diviseur $D$ et d'un recouvrement ouvert $(U_\alpha)$ de $X$ tel que $D|U_\alpha=\operatorname{div}_{U_\alpha}(g_\alpha)$, où $g_\alpha$ est une fonction méromorphe régulière au-dessus de $U_\alpha$;
$\partial(D)$ est la classe de cohomologie du cocycle $(c_{\alpha\beta})$, où $c_{\alpha\beta}=g_{\alpha|\beta}g_{\beta|\alpha}^{-1}$, $g_{\alpha|\beta}$ désignant la restriction de $g_\alpha$ à $U_\alpha\cap U_\beta$.
Il est clair que l'image par $\pi$ de cette classe de cohomologie est la classe de l'Idéal fractionnaire inversible $\sh{L}$ tel que pour tout $\alpha$, $\sh{L}|U_\alpha=\sh{O}_{U_\alpha}.g_\alpha^{-1}$, qui n'est autre par définition que $\sh{O}_X(D)$ \sref{IV.21.2.8}.
\end{remark}

\subsection{Images réciproques de diviseurs}
\EGAReaderTarget{ega4.c04.subsec.21-4}{subsection:IV.21.4}
\label{IV.21.4}

\begin{env}[21.4.1]
\EGAReaderTarget{ega4.c04.statement.21-4-1}{IV.21.4.1}
Soit $f:X'\to X$ un morphisme d'espaces annelés;
proposons-nous de donner des conditions permettant d'associer à un diviseur $D$ sur $X$ un diviseur $D'$ sur $X'$, \emph{image réciproque} de $D$ par $f$.
Notons d'abord pour cela que pour toute section $t\in\Gamma(X,\sh{O}_X^*)$, l'image de $t$ par l'homomorphisme canonique $\Gamma(X,\sh{O}_X)\to\Gamma(X',\sh{O}_{X'})$ est encore inversible, autrement dit appartient à $\Gamma(X',\sh{O}_{X'}^*)$.
Considérons alors $D$ comme donné par la classe d'équivalence d'un couple $(\sh{L},s)$, où $\sh{L}$ est un $\sh{O}_X$-Module inversible et $s$ une section méromorphe régulière de $\sh{L}$ au-dessus de $X$ \sref{IV.21.2.11}.
Formons le $\sh{O}_{X'}$-Module inversible $f^*(\sh{L})=\sh{L}'$;
dire que l'image réciproque $s\circ f$ de $s$ par $f$ existe \sref{IV.20.1.11} et est une section méromorphe régulière de $\sh{L}'$ au-dessus de $X'$ revient à dire que les images réciproques par $f$ de $s$ et de $s^{\otimes(-1)}$ existent, autrement dit que $s\in\Gamma(X,\sh{M}_f(\sh{L}))$ et $s^{\otimes(-1)}\in\Gamma(X,\sh{M}_f(\sh{L}^{-1}))$;
la remarque faite ci-dessus montre alors que si $(\sh{L}_1,s_1)$ est un couple équivalent à $(\sh{L},s)$ au sens de \sref{IV.21.2.10}, l'image réciproque $s_1\circ f$ existe et est une section méromorphe régulière de $\sh{L}_1'=f^*(\sh{L}_1)$ au-dessus de $X'$, et les couples $(\sh{L}',s\circ f)$ et $(\sh{L}_1',s_1\circ f)$ sont équivalents.
On peut donc poser la définition suivante :
\end{env}

\begin{definition}[21.4.2]
\EGAReaderTarget{ega4.c04.def.21-4-2}{IV.21.4.2}
Étant donné un morphisme $f:X'\to X$ d'espaces annelés, on dit que l'\emph{image réciproque par $f$} d'un diviseur $D$ sur $X$ existe si l'on a $s_D\in\Gamma(X,\sh{M}_f(\sh{O}_X(D)))$ et $s_{-D}\in\Gamma(X,\sh{M}_f(\sh{O}_X(-D)))$ (cf. \sref{IV.20.1.11}).
On appelle alors \emph{image réciproque} de $D$ par $f$ et l'on note $f^*(D)$ le diviseur sur $X'$ qui correspond canoniquement \sref{IV.21.2.11} à la classe du couple $(f^*(\sh{O}_X(D)),s_D\circ f)$.
\end{definition}

Il résulte aussitôt de cette définition que si $D$ et $D'$ ont des images réciproques par $f$, il en est de même de $-D$ et de $D+D'$ (compte tenu de \sref{IV.21.2.9.4}) et que l'on a
\oldpage[IV-4]{266}
$f^*(D+D')=f^*(D)+f^*(D')$.
Autrement dit, l'ensemble $\operatorname{Div}^f(X)$ des diviseurs sur $X$ dont l'image réciproque par $f$ existe est un sous-groupe de $\operatorname{Div}(X)$, et l'application $D\mapsto f^*(D)$ est un homomorphisme croissant du sous-groupe ordonné $\operatorname{Div}^f(X)$ dans le groupe ordonné $\operatorname{Div}(X')$, rendant commutatif le diagramme
\[
\label{IV.21.4.2.1}
  \xymatrix{
    \operatorname{Div}^f(X) \ar[r]^{l_X} \ar[d]_{f^*} & \operatorname{Pic}(X) \ar[d]^{\operatorname{Pic}(f)} \\
    \operatorname{Div}(X') \ar[r]^{l_{X'}} & \operatorname{Pic}(X')
  }
  \tag{21.4.2.1}
\]

\begin{env}[21.4.3]
\label{IV.21.4.3}
La définition \sref{IV.21.4.2} montre aussitôt que, pour que $f^*(D)$ existe, il faut et il suffit que pour tout ouvert $U$ de $X$, l'image réciproque par $f^U:f^{-1}(U)\to U$ (restriction de $f$) de $D|U$ existe.
Or, si $D=\operatorname{div}(g)$, où $g$ est une fonction méromorphe régulière sur $X$, dire que l'image réciproque de $s_D$ par $f$ existe et est une section méromorphe régulière de $f^*(\sh{O}_X(D))$ signifie \sref{IV.21.2.9} que l'image réciproque de $g$ par $f$ existe et est une fonction méromorphe régulière sur $X'$.
On en déduit aussitôt une seconde description de $\operatorname{Div}^f(X)$ et de $f^*(D)$ : on considère le sous-faisceau de groupes de $\sh{M}_X^*$, noté $\sh{M}_f^*$, formé des germes de fonctions méromorphes régulières sur un ouvert de $X$ et dont l'image réciproque par $f$ existe et est régulière sur l'ouvert image réciproque \sref{IV.20.1.11}.
Alors si $f=(\psi,\theta)$, l'homomorphisme canonique \sref{IV.20.1.11} $\psi^*(\sh{M}_f)\to\sh{M}_{X'}$ donne par restriction un homomorphisme de faisceaux de groupes $\psi^*(\sh{M}_f^*)\to\sh{M}_{X'}^*$.
Posant $\sh{D}iv_X'=\sh{M}_f^*/\sh{O}_X^*$, on a $\operatorname{Div}^f(X)=\Gamma(X,\sh{D}iv_X')$, et l'application $D\mapsto f^*(D)$ correspond à l'homomorphisme de faisceaux de groupes $\psi^*(\sh{M}_f^*)/\psi^*(\sh{O}_X^*)\to\sh{M}_{X'}^*/\sh{O}_{X'}^*=\sh{D}iv_{X'}$ déduit du précédent par passage aux quotients.
\end{env}

\begin{env}[21.4.4]
\EGAReaderTarget{ega4.c04.statement.21-4-4}{IV.21.4.4}
Il résulte aussitôt des définitions précédentes que si $f':X''\to X'$ est un second morphisme d'espaces annelés, $D$ un diviseur sur $X$ tel que les images réciproques $f^*(D)$ et $f'^*(f^*(D))$ existent, alors $(f\circ f')^*(D)$ existe et est égal à $f'^*(f^*(D))$.
\end{env}

\begin{proposition}[21.4.5]
\label{IV.21.4.5}
Soit $f:X'\to X$ un morphisme d'espaces annelés.
Dans chacun des trois cas suivants, l'image réciproque par $f$ de tout diviseur sur $X$ est définie :
\begin{enumerate}[label=(\roman*)]
  \item $f$ est plat.
  \item $X$ et $X'$ sont des préschémas localement noethériens et l'on a $f(\operatorname{Ass}(\sh{O}_{X'}))\subset\operatorname{Ass}(\sh{O}_X)$.
  \item $X$ et $X'$ sont des préschémas, l'ensemble des composantes irréductibles de $X$ est localement fini, $X'$ est réduit et toute composante irréductible de $X'$ domine une composante irréductible de $X$.
\end{enumerate}
\end{proposition}

\begin{proof}
Il suffit en effet de montrer que dans les trois cas on a $\sh{M}_f=\sh{M}_X$;
dans le cas (i), cela résulte de \sref{IV.20.1.12}.
Dans le cas (iii), on peut se borner au cas où $X=\operatorname{Spec}(A)$ et $X'=\operatorname{Spec}(A')$ sont affines;
si $s\in A$ est régulier, il n'appartient à aucun idéal premier minimal de $A$ \sref{IV.20.1.3.1}, donc l'hypothèse entraîne que son image dans $A'$ n'appartient à aucun idéal premier minimal de $A'$, et est par suite un élément régulier de $A'$ \sref{IV.20.1.3.1}.
Dans le cas (ii) les fonctions méromorphes sur $X'$ s'identifient aux pseudo-fonctions sur $X'$ \sref{IV.20.2.11}, et l'hypothèse, jointe à \sref{IV.20.2.13},~(iv), assure que l'image réciproque
\oldpage[IV-4]{267}
par $f$ de tout ouvert schématiquement dense dans $X$ est un ouvert schématiquement dense dans $X'$;
on conclut donc par \sref{IV.20.3.12}.
\end{proof}

\begin{corollary}[21.4.6]
\label{IV.21.4.6}
Soit $X$ un préschéma ayant l'une des propriétés suivantes :
\begin{enumerate}[label=(\roman*)]
  \item $X$ est localement noethérien.
  \item $X$ est réduit et l'ensemble de ses composantes irréductibles est localement fini.
\end{enumerate}
Alors, pour tout $x\in X$, on a un isomorphisme canonique
\[
\label{IV.21.4.6.1}
  (\sh{D}iv_X)_x\xrightarrow{\sim}\operatorname{Div}(\sh{O}_{X,x}).
  \tag{21.4.6.1}
\]
\end{corollary}

\begin{proof}
Cela résulte en effet de \sref{IV.20.2.11}, \sref{IV.20.3.7} et de ce que $(\sh{O}_X^*)_x$ s'identifie au groupe des éléments inversibles de l'anneau $\sh{O}_{X,x}$.
\end{proof}

\begin{env}[21.4.7]
\label{IV.21.4.7}
Soient $X$, $X'$ deux préschémas, $f:X'\to X$ un morphisme.
Si $D$ est un diviseur \emph{positif} sur $X$ tel que l'image réciproque $f^*(D)$ soit définie \sref{IV.21.4.2}, alors le sous-préschéma fermé $Y(f^*(D))$ de $X'$ n'est autre que l'image réciproque $f^{-1}(Y(D))$;
cela résulte aussitôt des définitions \sref{IV.21.4.2} et \sref{IV.21.2.12}.
\end{env}

\begin{proposition}[21.4.8]
\EGAReaderTarget{ega4.c04.prop.21-4-8}{IV.21.4.8}
Soient $X$, $Y$ deux préschémas;
$f:X\to Y$ un morphisme fidèlement plat.
Alors, si un diviseur $D$ sur $Y$ est tel que $f^*(D)\geq0$ (l'existence de $f^*(D)$ résultant de \sref{IV.21.4.5}), on a $D\geq0$.
En particulier, l'application $D\mapsto f^*(D)$ de $\operatorname{Div}(Y)$ dans $\operatorname{Div}(X)$ est injective.
\end{proposition}

\begin{proof}
La question étant locale sur $Y$, on peut se borner au cas où $D=\operatorname{div}(w)$, avec $w=uv^{-1}$, $u$ et $v$ étant deux sections régulières de $\sh{O}_Y$ au-dessus de $Y$.
Par hypothèse on a $u\sh{O}_X\subset v\sh{O}_X$, donc, pour tout $x\in X$, si l'on pose $y=f(x)$, on a $u_y\sh{O}_{X,x}\subset v_y\sh{O}_{X,x}$;
on en conclut que $u_y\sh{O}_{Y,y}\subset v_y\sh{O}_{Y,y}$ en vertu de l'hypothèse que $\sh{O}_{X,x}$ est un $\sh{O}_{Y,y}$-module fidèlement plat et de Bourbaki, \emph{Alg. comm.}, chap.~I, \textsection~3, n\textsuperscript{o}~5, prop.~10,~(ii);
d'où $u\sh{O}_Y\subset v\sh{O}_Y$, puisque $f$ est surjectif, et par suite $D\geq0$.
\end{proof}

\subsection{Images directes de diviseurs}
\EGAReaderTarget{ega4.c04.subsec.21-5}{subsection:IV.21.5}
\label{IV.21.5}

\begin{env}[21.5.1]
\label{IV.21.5.1}
Soient $X$, $X'$ deux préschémas, $f:X'\to X$ un morphisme.
Nous allons, dans ce numéro, donner des conditions suffisantes pour qu'on puisse associer à tout diviseur $D'$ sur $X'$ un diviseur $D$ sur $X$, \emph{image directe} de $D'$ par $f$.
Nous nous bornerons au cas où $f$ est un morphisme \emph{fini} (pour des conditions plus générales, voir le chapitre de ce Traité consacré à la théorie des intersections).
\end{env}

\begin{lemma}[21.5.2]
\label{IV.21.5.2}
Soient $A$ un anneau, $E$ un $A$-module libre de rang fini.
Pour qu'un endomorphisme $u$ de $E$ soit injectif, il faut et il suffit que $\det(u)$ soit un élément régulier de $A$.
\end{lemma}

\begin{proof}
Cela est prouvé dans Bourbaki, \emph{Alg.}, chap.~III, 3\textsuperscript{e}~éd., \textsection~8, n\textsuperscript{o}~2, prop.~3.
\end{proof}

\begin{env}[21.5.3]
\EGAReaderTarget{ega4.c04.statement.21-5-3}{IV.21.5.3}
Supposons maintenant que le morphisme $f:X'\to X$ soit fini, et en outre que $f$ vérifie l'une des deux propriétés suivantes :
\begin{enumerate}[label=\Roman*)]
  \item $f$ est un morphisme fini \emph{localement libre}, autrement dit \sref{IV.18.2.7} le $\sh{O}_X$-Module quasi-cohérent de type fini $f_*(\sh{O}_{X'})$ est localement libre.
  \item $X$ est un préschéma localement noethérien réduit, le $\sh{R}(X)$-Module $f_*(\sh{O}_{X'})\otimes_{\sh{O}_X}\sh{R}(X)$ est localement libre, et pour toute section $s'\in\Gamma(f^{-1}(U),\sh{O}_{X'})$ ($U$ ouvert
  \oldpage[IV-4]{268}
  dans $X$), $\mathrm{N}_{f_*(\sh{O}_{X'})|\sh{O}_X}(s')$ est une section de $\sh{O}_X$ au-dessus de $U$ (cf. (\textbf{II}, 6.5.1)).
  (On rappelle que la condition (II) est vérifiée pour tout morphisme fini $f$ lorsque $X$ est un préschéma localement noethérien normal (\emph{loc. cit.}).)
\end{enumerate}

On sait alors (\textbf{II}, 6.5.5) que pour tout $\sh{O}_{X'}$-Module inversible $\sh{L}'$ on définit la norme $\sh{L}=\mathrm{N}_{X'/X}(\sh{L}')$ (que nous écrirons aussi $\mathrm{N}(\sh{L}')$), qui est un $\sh{O}_X$-Module inversible.
En outre, pour tout ouvert $U$ de $X$ et toute section régulière $s'\in\Gamma(f^{-1}(U),\sh{O}_{X'})$, la norme $\mathrm{N}_{X'/X}(s')=\mathrm{N}_{f_*(\sh{O}_{X'})|\sh{O}_X}(s')$ (que nous écrirons aussi $\mathrm{N}(s')$) est un élément régulier de $\Gamma(U,\sh{O}_X)$;
on est en effet ramené aussitôt au cas où $U=X$ est affine, et alors la conclusion résulte de \sref{IV.21.5.2} dans l'hypothèse (I);
d'autre part, dans l'hypothèse (II), le fait que $\sh{R}(X)$ soit un $\sh{O}_X$-Module plat entraîne que la section $s'\otimes1$ de $\Gamma(U,f_*(\sh{O}_{X'})\otimes_{\sh{O}_X}\sh{R}(X))$ est aussi régulière (\hyperref[0.6.1.2]{$\mathbf{0}_{\mathrm I}$, 6.1.2}), et la conclusion résulte encore de \sref{IV.21.5.2} appliqué à l'anneau $\Gamma(U,\sh{R}(X))$, compte tenu de la définition de la norme d'une section (\textbf{II}, 6.5.3).
Cela étant, soit $u'$ une section méromorphe de $\sh{L}'$ au-dessus de $X'$;
le morphisme $f$ étant affine, tout point $x\in X$ admet un voisinage ouvert $U$ tel que $u'|f^{-1}(U)$ soit de la forme $t'/s'$, où $t'\in\Gamma(f^{-1}(U),\sh{L}')$ et $s'$ est une section régulière dans $\Gamma(f^{-1}(U),\sh{O}_{X'})$;
l'élément $\mathrm{N}(t')/\mathrm{N}(s')$ (où $\mathrm{N}(t')$ est la section de $\sh{L}$ définie dans (\textbf{II}, 6.5.3)) est alors une section méromorphe de $\sh{L}$ au-dessus de $U$ en vertu de ce qui précède, et il résulte des propriétés multiplicatives de la norme (\textbf{II}, 6.5.3.1) que cette section ne dépend que de $u'|f^{-1}(U)$ et non de son écriture sous la forme $t'/s'$;
pour la même raison, lorsque $U$ varie, les sections méromorphes de $\sh{L}$ au-dessus de $U$ ainsi définies sont les restrictions d'une même section méromorphe de $\sh{L}$ au-dessus de $X$, que l'on appelle la \emph{norme} de $u'$ et que l'on note $\mathrm{N}_{X'/X}(u')$ (ou simplement $\mathrm{N}(u')$).
L'application ainsi définie
\[
\label{IV.21.5.3.1}
  \mathrm{N}_{X'/X}:\Gamma(X',\sh{M}_{X'}(\sh{L}'))\to\Gamma(X,\sh{M}_X(\mathrm{N}_{X'/X}(\sh{L}')))
  \tag{21.5.3.1}
\]
prolonge la norme définie dans (\textbf{II}, 6.5.5.3);
si $u'$ est une section méromorphe régulière de $\sh{L}'$ au-dessus de $X'$, il résulte aussitôt de ce qui précède que $\mathrm{N}(u')$ est une section méromorphe régulière de $\sh{L}$ au-dessus de $X$, car (avec les mêmes notations) $\mathrm{N}(t')$ est régulière si $t'$ l'est.
Enfin, si $\sh{L}_1'$, $\sh{L}_2'$ sont deux $\sh{O}_{X'}$-Modules inversibles, $s_1'$ (resp. $s_2'$) une section méromorphe régulière de $\sh{L}_1'$ (resp. $\sh{L}_2'$) au-dessus de $X'$, on a, en vertu de ce qui précède et de (\textbf{II}, 6.5.3.1)
\[
\label{IV.21.5.3.2}
  \mathrm{N}(s_1'\otimes s_2')=\mathrm{N}(s_1')\otimes\mathrm{N}(s_2').
  \tag{21.5.3.2}
\]
\end{env}

\begin{env}[21.5.4]
\label{IV.21.5.4}
Supposons toujours que $f$ vérifie l'une des hypothèses I), II) de \sref{IV.21.5.3}.
Si $(\sh{L}_1',s_1')$, $(\sh{L}_2',s_2')$ sont deux couples formés chacun d'un $\sh{O}_{X'}$-Module inversible et d'une section méromorphe régulière de ce Module au-dessus de $X'$, qui en outre sont équivalents au sens de \sref{IV.21.2.10}, alors les couples $(\mathrm{N}(\sh{L}_1'),\mathrm{N}(s_1'))$ et $(\mathrm{N}(\sh{L}_2'),\mathrm{N}(s_2'))$ sont aussi équivalents, car un isomorphisme de $\sh{O}_{X'}$-Modules inversibles a pour norme un isomorphisme de leurs normes (\textbf{II}, 6.5.3), et on a vu plus haut que $\mathrm{N}_{X'/X}$ transforme les sections de $\Gamma(X',\sh{O}_{X'}^*)$ en celles de $\Gamma(X,\sh{O}_X^*)$.
On peut donc poser la définition suivante :
\end{env}

\oldpage[IV-4]{269}

\begin{definition}[21.5.5]
\EGAReaderTarget{ega4.c04.def.21-5-5}{IV.21.5.5}
Étant donné un morphisme fini $f:X'\to X$ de préschémas, vérifiant l'une des conditions I), II) de \sref{IV.21.5.3}, on appelle \emph{image directe} (ou \emph{norme}) d'un diviseur $D'$ sur $X'$ par $f$, et l'on note $f_*(D')$ (ou $\mathrm{N}_{X'/X}(D')$), le diviseur sur $X$ qui correspond canoniquement \sref{IV.21.2.11} à la classe du couple $(\mathrm{N}_{X'/X}(\sh{O}_{X'}(D')),\mathrm{N}_{X'/X}(s_{D'}))$.
\end{definition}

Il résulte aussitôt de cette définition, compte tenu de \sref{IV.21.2.9.4} et de \sref{IV.21.5.3.2}, que si $D_1'$, $D_2'$, $D'$ sont des diviseurs sur $X'$, on a
\[
\label{IV.21.5.5.1}
  f_*(D_1'+D_2')=f_*(D_1')+f_*(D_2')
  \tag{21.5.5.1}
\]
et $D'\geq0$ entraîne $f_*(D')\geq0$, autrement dit, $D'\mapsto f_*(D')$ est un homomorphisme croissant du groupe ordonné $\operatorname{Div}(X')$ dans le groupe ordonné $\operatorname{Div}(X)$.
La définition \sref{IV.21.5.5} montre d'ailleurs aussitôt que pour tout ouvert $U$ de $X$, on a $f_*^U(D'|f^{-1}(U))=f_*(D')|U$ ($f^U$ étant la restriction $f^{-1}(U)\to U$ de $f$), et les homomorphismes $f_*^U$, pour $U$ variable, définissent donc un homomorphisme de faisceaux de groupes ordonnés
\[
\label{IV.21.5.5.2}
  \mathrm{N}_{X'/X}:f_*(\sh{D}iv_{X'})\to\sh{D}iv_X.
  \tag{21.5.5.2}
\]
En outre, pour tout ouvert $U$ de $X$, tout $\sh{O}_{X'}$-Module inversible $\sh{L}'$ et toute section méromorphe régulière $s'$ de $\sh{L}'$ au-dessus de $f^{-1}(U)$, on a, d'après les définitions précédentes et \sref{IV.21.1.4}
\[
\label{IV.21.5.5.3}
  \operatorname{div}_U(\mathrm{N}(s'))=f_*^U(\operatorname{div}_{f^{-1}(U)}(s')).
  \tag{21.5.5.3}
\]

\begin{proposition}[21.5.6]
\label{IV.21.5.6}
Soit $f:X'\to X$ un morphisme fini localement libre et supposons que $f_*(\sh{O}_{X'})$ soit de rang constant $n$.
Alors, pour tout diviseur $D$ sur $X$, $f_*(D)$ est défini et l'on a
\[
\label{IV.21.5.6.1}
  f_*(f^*(D))=nD.
  \tag{21.5.6.1}
\]
\end{proposition}

\begin{proof}
La première assertion résulte de ce que $f$ est plat \sref{IV.21.4.5}, et la seconde est conséquence immédiate des définitions et de (\textbf{II}, 6.5.3.2).
\end{proof}

\begin{proposition}[21.5.7]
\label{IV.21.5.7}
Soient $f:X'\to X$ un morphisme fini vérifiant l'une des hypothèses I), II) de \sref{IV.21.5.3}, $f':X''\to X'$ un morphisme fini localement libre de rang constant $n$.
Alors $f''=f\circ f':X''\to X$ vérifie la même hypothèse que $f$, et pour tout diviseur $D''$ sur $X''$, on a
\[
\label{IV.21.5.7.1}
  f_*''(D'')=f_*(f_*'(D'')).
  \tag{21.5.7.1}
\]
\end{proposition}

Compte tenu de la définition \sref{IV.21.5.5}, il suffit de prouver le résultat suivant :

\begin{lemma}[21.5.7.2]
\label{IV.21.5.7.2}
Sous les hypothèses de \sref{IV.21.5.7}, on a un isomorphisme fonctoriel
\[
\label{IV.21.5.7.3}
  \mathrm{N}_{X''/X}(\sh{L}'')\xrightarrow{\sim}\mathrm{N}_{X'/X}(\mathrm{N}_{X''/X'}(\sh{L}''))
  \tag{21.5.7.3}
\]
dans la catégorie des $\sh{O}_X$-Modules inversibles.
\end{lemma}

En effet, compte tenu de la définition de la norme d'une section de $\sh{L}''$ (\textbf{II}, 6.5.3) et de la définition \sref{IV.21.5.5}, on obtiendra aussitôt \sref{IV.21.5.7.1}.
Pour prouver \sref{IV.21.5.7.2}, il suffit, compte tenu des définitions de (\textbf{II}, 6.5.2 et 6.5.3), de prouver que pour toute section $s$ de $f_*''(\sh{O}_{X''})=f_*(f_*'(\sh{O}_{X''}))$ au-dessus d'un ouvert $U$ de $X$, on a
\[
\label{IV.21.5.7.4}
  \mathrm{N}_{f_*''(\sh{O}_{X''})|\sh{O}_X}(s)=\mathrm{N}_{f_*(\sh{O}_{X'})|\sh{O}_X}(\mathrm{N}_{f_*'(\sh{O}_{X''})|\sh{O}_{X'}}(s)).
  \tag{21.5.7.4}
\]

\oldpage[IV-4]{270}

La question est évidemment locale sur $X$, et on peut donc se borner au cas où $X=\operatorname{Spec}(A)$ est affine;
on a alors $X'=\operatorname{Spec}(A')$ et $X''=\operatorname{Spec}(A'')$, et on peut supposer que $A''$ est un $A'$-module projectif de rang $n$.
Lorsque $f$ est localement libre, on peut supposer que $A'$ est un $A$-module libre de rang $m$, et alors $A''$ est un $A$-module projectif de rang $mn$, et en restreignant $X$ à un ouvert convenable, on peut supposer que $A''$ est un $A$-module libre de rang $mn$;
la formule \sref{IV.21.5.7.4} résulte alors de la transitivité de la norme (Bourbaki, \emph{Alg.}, chap.~VIII, \textsection~12, n\textsuperscript{o}~2, prop.~7).
Lorsque $f$ vérifie l'hypothèse II), $A$ est noethérien réduit, et si $R$ est son anneau total de fractions, $A'\otimes_A R$ est un $R$-module libre de rang $m$, donc $A''\otimes_A R=A''\otimes_{A'}(A'\otimes_A R)$ est un $R$-module projectif de rang $mn$, et comme $R$ est alors un anneau semi-local, ce $R$-module est libre;
la proposition résulte alors encore de la transitivité des normes.

\begin{proposition}[21.5.8]
\label{IV.21.5.8}
Soient $f:X'\to X$ un morphisme fini, $g:Y\to X$ un morphisme;
on pose $Y'=X'\times_XY$, $f'=f_{(Y)}:Y'\to Y$, $g'=g_{(X')}:Y'\to X'$.
On suppose vérifiée l'une des conditions suivantes :
\begin{enumerate}[label=(\roman*)]
  \item $f$ est localement libre et $g$ est plat.
  \item $f$ est localement libre, $X$ et $Y$ sont localement noethériens, $g(\operatorname{Ass}(\sh{O}_Y))\subset\operatorname{Ass}(\sh{O}_X)$ et $g'(\operatorname{Ass}(\sh{O}_{Y'}))\subset\operatorname{Ass}(\sh{O}_{X'})$.
  \item $f$ vérifie l'hypothèse II) de \sref{IV.21.5.3}, $Y$ est localement noethérien, $Y$ et $Y'$ sont réduits et toute composante irréductible de $Y$ (resp. $Y'$) domine une composante irréductible de $X$ (resp. $X'$).
\end{enumerate}
Alors, pour tout diviseur $D'$ sur $X'$, $g'^*(D')$ est défini, $g^*(f_*(D'))$ est défini et l'on a
\[
\label{IV.21.5.8.1}
  g^*(f_*(D'))=f_*'(g'^*(D')).
  \tag{21.5.8.1}
\]
\end{proposition}

\begin{proof}
En effet, dans tous les cas, il résulte de (\textbf{II}, 6.5.8) que l'on a un isomorphisme fonctoriel
\[
  g^*(\mathrm{N}_{X'/X}(\sh{L}'))\xrightarrow{\sim}\mathrm{N}_{Y'/Y}(g'^*(\sh{L}'))
\]
dans la catégorie des $\sh{O}_Y$-Modules inversibles;
en outre (\textbf{II}, 6.5.4), si $s'$ est une section de $\sh{O}_{X'}$ au-dessus de $f^{-1}(U)$, $s''$ la section correspondante de $\sh{O}_{Y'}$ au-dessus de $g'^{-1}(f^{-1}(U))$ ($U$ ouvert de $X$), $\mathrm{N}_{Y'/Y}(s'')$ est la section de $\sh{O}_Y$ au-dessus de $g^{-1}(U)$ qui correspond à la section $\mathrm{N}_{X'/X}(s')$ de $\sh{O}_X$ au-dessus de $U$.
La formule \sref{IV.21.5.8.1} résultera donc des définitions si l'on prouve que $g'^*(D')$ et $g^*(D)$ sont définis, quels que soient les diviseurs $D'$ sur $X'$ et $D$ sur $X$.
En ce qui concerne $D$, cela résulte des hypothèses faites et de \sref{IV.21.4.5}.
En ce qui concerne $D'$, dans le cas (i) $g'$ est plat, donc dans tous les cas $g'^*(D')$ est défini en vertu de \sref{IV.21.4.5}.
\end{proof}

% FIN SOURCE ega4/ega4-21a-fr.tex

%
% DEBUT SOURCE ega4/ega4-21b-fr.tex
\subsection{Cycle 1-codimensionnel associé à un diviseur}
\EGAReaderTarget{ega4.c04.subsec.21-6}{subsection:IV.21.6}
\label{IV.21.6}

\begin{env}[21.6.1]
\EGAReaderTarget{ega4.c04.statement.21-6-1}{IV.21.6.1}
Soit $X$ un préschéma localement noethérien, et soit $\mathfrak{I}(X)$ l'ensemble des \emph{parties fermées irréductibles} de $X$ (qui est en correspondance biunivoque avec $X$ par l'application $x\mapsto\overline{\{x\}}$).
Dans le groupe produit $\bb{Z}^X$, considérons le sous-groupe $\mathfrak{Z}(X)$ des éléments $(n_x)_{x\in X}$ tels que l'ensemble des $\overline{\{x\}}\in\mathfrak{I}(X)$ tels que $n_x\neq0$ (ou, ce qui revient au même, l'ensemble des $x\in X$ tels que $n_x\neq0$) soit localement fini.
Il est clair que $\mathfrak{Z}(X)$ est un sous-groupe de $\bb{Z}^X$ qui contient le groupe somme directe $\bb{Z}^{(X)}$ (groupe libre ayant pour base $\mathfrak{I}(X)$), et lui est égal lorsque $X$ est noethérien.
Les éléments de $\mathfrak{Z}(X)$ sont appelés \emph{cycles} sur $X$ et ceux de $\mathfrak{I}(X)$ \emph{cycles premiers} (ils ne forment pas en général une base de $\mathfrak{Z}(X)$ lorsque $X$ n'est pas noethérien).
On considère toujours $\mathfrak{Z}(X)$ comme ordonné par l'ordre induit sur ce sous-groupe par l'ordre produit de $\bb{Z}^X$ et on notera $\mathfrak{Z}^+(X)$ l'ensemble des cycles $\geq0$.

\oldpage[IV-4]{271}

Pour tout cycle $Z\in\mathfrak{Z}(X)$, égal à $(n_x)_{x\in X}$, on écrira par abus de langage,
\[
  Z=\sum_{x\in X}n_x\cdot\overline{\{x\}}.
\]
On appelle \emph{multiplicité de $Z$ au point $x$} et on note $\operatorname{mult}_x(Z)$ l'entier $n_x$ (positif ou négatif).
Dire que $Z\geq0$ signifie que $\operatorname{mult}_x(Z)\geq0$ pour tout $x\in X$.
On appelle \emph{support de $Z$} et on note $\Supp(Z)$ la réunion des $\overline{\{x\}}$ tels que $\operatorname{mult}_x(Z)\neq0$;
par définition de $\mathfrak{Z}(X)$, c'est une partie fermée de $X$, comme réunion d'une famille localement finie de parties fermées.
On appelle \emph{dimension} (resp.\ \emph{codimension}) de $Z$ et on note $\dim(Z)$ (resp.\ $\codim(Z)$) la dimension (resp.\ codimension dans $X$) de $\Supp(Z)$.
\end{env}

\begin{env}[21.6.2]
\EGAReaderTarget{ega4.c04.statement.21-6-2}{IV.21.6.2}
On dit qu'une partie fermée $Y$ de $X$ est \emph{purement de codimension $p$} (dans $X$) si \emph{chacune} de ses composantes irréductibles est de codimension $p$ dans $X$.
On dit qu'un cycle est \emph{purement de codimension $p$}, ou (par abus de langage) est un \emph{cycle $p$-codimensionnel}, si son support est purement de codimension $p$.
On note $X^{(p)}$ l'ensemble des $x\in X$ tels que $\codim(\overline{\{x\}},X)=p$, ou, ce qui revient au même \sref{IV.5.1.2}, $\dim(\sh{O}_{X,x})=p$:
les cycles purement de codimension $p$ forment un sous-groupe $\mathfrak{Z}^p(X)$ de $\mathfrak{Z}(X)$, isomorphe au sous-groupe de $\bb{Z}^{X^{(p)}}$ formé des $(n_x)_{x\in X^{(p)}}$ tels que l'ensemble des $\overline{\{x\}}$ (ou l'ensemble des $x$) pour lesquels $n_x\neq0$ soit localement fini;
ce sous-groupe contient le groupe libre $\bb{Z}^{(X^{(p)})}$ et lui est identique lorsque $X$ est noethérien.
On note $\mathfrak{Z}^{p+}(X)$ l'ensemble des éléments $\geq0$ de $\mathfrak{Z}^p(X)$.
Il est clair que le groupe ordonné $\mathfrak{Z}(X)$ contient la somme directe des sous-groupes ordonnés $\mathfrak{Z}^p(X)$, et est identique à cette somme directe lorsque $X$ est noethérien;
dans ce dernier cas, on considère $\mathfrak{Z}(X)$ comme \emph{gradué} par les $\mathfrak{Z}^p(X)$ pour $p\in\bb{N}$.
\end{env}

\begin{env}[21.6.3]
\EGAReaderTarget{ega4.c04.statement.21-6-3}{IV.21.6.3}
Soient $Z=\sum_{x\in X}n_x\cdot\overline{\{x\}}$ un cycle sur $X$, $U$ un ouvert de $X$;
on appelle \emph{restriction} de $Z$ à $U$ et on note $Z|U$ le cycle $\sum_{x\in U}n_x\cdot(U\cap\overline{\{x\}})$ sur $U$;
on a $\Supp(Z|U)=\Supp(Z)\cap U$.
Il est clair que $Z\mapsto Z|U$ est un homomorphisme de groupes ordonnés de $\mathfrak{Z}(X)$ dans $\mathfrak{Z}(U)$ (resp.\ de $\mathfrak{Z}^p(X)$ dans $\mathfrak{Z}^p(U)$), et que si $V$ est un ouvert contenu dans $U$, on a $Z|V=(Z|U)|V$;
l'application $U\mapsto\mathfrak{Z}(U)$ (resp.\ $U\mapsto\mathfrak{Z}^p(U)$) est donc un \emph{préfaisceau} sur $X$ de groupes ordonnés commutatifs.
En fait ce préfaisceau est un \emph{faisceau}, somme directe des faisceaux $(i_x)_*(\bb{Z}_{\overline{\{x\}}})$, où $x$ parcourt $X$ (resp.\ $X^{(p)}$) et pour chaque $x\in X$, $i_x$ est l'injection canonique $\overline{\{x\}}\to X$ et $\bb{Z}_{\overline{\{x\}}}$ est le faisceau simple associé au faisceau constant $\bb{Z}$ sur l'espace $\overline{\{x\}}$:
cela résulte aussitôt de la description des sections d'une somme directe de faisceaux de groupes commutatifs (G, II, 2.7).
On note ce faisceau $\mathscr{Z}_X$ (resp.\ $\mathscr{Z}_X^p$).
On note $\mathscr{Z}_X^+$ (resp.\ $\mathscr{Z}_X^{p+}$) le sous-faisceau
\oldpage[IV-4]{272}
de monoïdes $U\mapsto\mathfrak{Z}^+(U)$ (resp.\ $U\mapsto\mathfrak{Z}^{p+}(U)$) de $\mathscr{Z}_X$ (resp.\ $\mathscr{Z}_X^p$).
Le faisceau $\mathscr{Z}_X$ est évidemment \emph{somme directe} des faisceaux $\mathscr{Z}_X^p$.

Notons enfin que les faisceaux $\mathscr{Z}_X^p$ (donc aussi $\mathscr{Z}_X$) sont \emph{flasques}.
Supposons en effet donnée une section $Z$ de $\mathscr{Z}_X^p$ au-dessus d'un ouvert $U$, de sorte que l'ensemble $S$ des $x\in U$ tels que $\operatorname{mult}_x(Z)\neq0$ est localement fini \emph{dans} $U$;
cet ensemble est aussi localement fini \emph{dans} $X$, car pour tout $z\in X$ et tout voisinage ouvert noethérien $V$ de $z$, $U\cap V$ est noethérien, donc ne contient qu'un nombre fini de points de $S$.
On définit donc une section $Z'$ de $\mathscr{Z}_X^p$ au-dessus de $X$ en posant $Z'=\sum_{x\in U}\operatorname{mult}_x(Z)\cdot\overline{\{x\}}$ et puisque l'adhérence de $x$ dans $U$ est $\overline{\{x\}}\cap U$, on a bien $Z'|U=Z$, d'où notre assertion.
\end{env}

\begin{env}[21.6.4]
\label{IV.21.6.4}
Nous nous proposons de définir un homomorphisme canonique
\[
\label{IV.21.6.4.1}
  c:\sh{D}iv_X\to\mathscr{Z}_X^1
  \tag{21.6.4.1}
\]
de faisceaux de groupes commutatifs.
Il suffit évidemment de définir un homomorphisme de $\sh{M}_X^*$ dans $\mathscr{Z}_X^1$, qui s'annule dans $\sh{O}_X^*$ \sref{IV.21.1.2};
or, par définition, $\sh{M}_X^*$ est le \emph{symétrisé} du faisceau de monoïdes $\sh{O}_X\cap\sh{M}_X^*$, et un homomorphisme $\sh{M}_X^*\to\mathscr{Z}_X^1$ est déterminé de façon unique par la donnée de sa restriction $\sh{S}(\sh{O}_X)\to\mathscr{Z}_X^1$ qui est un homomorphisme de faisceaux de \emph{monoïdes} soumis à la seule condition d'être nul dans $\sh{O}_X^*$;
il revient au même de dire que pour définir $c$, il suffit de définir un homomorphisme de \emph{faisceaux de monoïdes}
\[
\label{IV.21.6.4.2}
  c:\sh{D}iv_X^+\to\mathscr{Z}_X^{1+}.
  \tag{21.6.4.2}
\]
\end{env}

\begin{env}[21.6.5]
\label{IV.21.6.5}
Or, on a vu qu'à tout diviseur positif $D$ sur $X$ on fait correspondre le sous-préschéma fermé $Y(D)$ de $X$, défini par l'Idéal $\sh{J}_X(D)\subset\sh{O}_X$, et qui est régulièrement immergé et de codimension $1$.
En chacun des points maximaux $x$ de $Y(D)$, on a donc (\sref{IV.19.1.4} et \sref{IV.5.1.2}) $\codim(\overline{\{x\}},X)=\dim(\sh{O}_{X,x})=1$, autrement dit $x\in X^{(1)}$;
d'ailleurs l'ensemble de ces points maximaux est localement fini dans $X$ \sref{IV.3.1.6}, et $\sh{O}_{Y(D),x}$ est un anneau \emph{artinien}.
En tout point $x\in X^{(1)}$ qui n'est pas point maximal de $Y(D)$, on a nécessairement $x\notin Y(D)$, donc $\sh{O}_{Y(D),x}=0$.
Posons
\[
\label{IV.21.6.5.1}
  \operatorname{cyc}(D)=\sum_{x\in X^{(1)}}\operatorname{long}(\sh{O}_{Y(D),x})\cdot\overline{\{x\}}
  \tag{21.6.5.1}
\]
qui est donc un élément de $\mathfrak{Z}^1(X)$.
\end{env}

\begin{proposition}[21.6.6]
\label{IV.21.6.6}
L'application $D\mapsto\operatorname{cyc}(D)$ est un homomorphisme du monoïde $\operatorname{Div}^+(X)$ dans $\mathfrak{Z}^{1+}(X)$.
\end{proposition}

\begin{proof}
Tout revient à voir que pour deux diviseurs positifs $D$, $D'$, on a
\[
  \operatorname{cyc}(D+D')=\operatorname{cyc}(D)+\operatorname{cyc}(D').
\]
Or, on a \sref{IV.21.2.5} $\sh{J}_X(D+D')=\sh{J}_X(D)\sh{J}_X(D')$;
tout revient à voir que si $x\in X^{(1)}$, si l'on pose $A=\sh{O}_{X,x}$, et si $t$ et $t'$ sont deux éléments réguliers de $A$, on a $\operatorname{long}(A/tt'A)=\operatorname{long}(A/tA)+\operatorname{long}(A/t'A)$;
mais puisque $t$ est régulier, $tA/tt'A$ est isomorphe à $A/t'A$, d'où aussitôt la proposition.
\end{proof}

Il résulte aussitôt des définitions que pour tout ouvert $U\subset X$, on a
\[
  Y(D|U)=Y(D)\cap U,
\]
\oldpage[IV-4]{273}
donc $\operatorname{cyc}(D|U)=\operatorname{cyc}(D)|U$, et les homomorphismes $\operatorname{cyc}:\operatorname{Div}^+(U)\to\mathfrak{Z}^1(U)$ définissent donc un homomorphisme de faisceaux de monoïdes \sref{IV.21.6.4.2}, d'où l'homomorphisme cherché \sref{IV.21.6.4.1} de faisceaux de groupes.

On a $\Supp(\operatorname{cyc}(D))\subset\Supp(D)$ pour tout diviseur $D$ et
\begin{env}[21.6.6.2]
\label{IV.21.6.6.2}
\[
  \Supp(\operatorname{cyc}(D))=\Supp(D)\quad\text{lorsque}\quad D\geq0;
\]
\end{env}
on a déjà vu en effet la seconde relation \sref{IV.21.2.12};
lorsque $D$ est quelconque, la relation $D_x=0$ entraîne l'existence d'un voisinage ouvert $U$ de $x$ tel que $D|U=0$, d'où $\operatorname{cyc}(D)|U=0$, ce qui prouve notre assertion.

\begin{env}[21.6.7]
\EGAReaderTarget{ega4.c04.statement.21-6-7}{IV.21.6.7}
L'homomorphisme \sref{IV.21.6.4.1} donne, par passage aux groupes de sections au-dessus de $X$, un \emph{homomorphisme croissant de groupes ordonnés} $\operatorname{Div}(X)\to\mathfrak{Z}^1(X)$, que nous noterons encore $D\mapsto\operatorname{cyc}(D)$;
le nombre $\operatorname{mult}_x(\operatorname{cyc}(D))$ pour $x\in X^{(1)}$ se note encore $\operatorname{mult}_x(D)$ ou $\operatorname{mult}_{\overline{\{x\}}}(D)$ et s'appelle \emph{multiplicité du diviseur $D$ au point $x$}, ou \emph{multiplicité du cycle premier $\overline{\{x\}}$ dans $D$};
c'est un entier positif ou négatif, égal comme on l'a vu à $\operatorname{long}(\sh{O}_{Y(D),x})$ lorsque $D$ est un diviseur \emph{positif};
on a donc par définition
\[
\label{IV.21.6.7.1}
  \operatorname{cyc}(D)=\sum_{x\in X^{(1)}}\operatorname{mult}_x(D)\cdot\overline{\{x\}},
  \tag{21.6.7.1}
\]
et en vertu du fait que $D\mapsto\operatorname{cyc}(D)$ est un homomorphisme, on a
\[
\label{IV.21.6.7.2}
  \operatorname{mult}_x(-D)=-\operatorname{mult}_x(D),\qquad
  \operatorname{mult}_x(D+D')=\operatorname{mult}_x(D)+\operatorname{mult}_x(D')
  \tag{21.6.7.2}
\]
pour deux diviseurs quelconques $D$, $D'$.

Pour toute fonction méromorphe régulière $f$ sur $X$, on pose en particulier, pour $x\in X^{(1)}$
\[
\label{IV.21.6.7.3}
  o_x(f)=\operatorname{mult}_x(\operatorname{div}(f))
  \tag{21.6.7.3}
\]
et on dit que $o_x(f)$ (entier positif ou négatif) est \emph{l'ordre de $f$ au point} $x\in X^{(1)}$.
Si $f_x\in\sh{O}_{X,x}$ (élément régulier de cet anneau local par hypothèse), on a donc
\[
\label{IV.21.6.7.4}
  o_x(f)=\operatorname{long}(\sh{O}_{X,x}/f_x\sh{O}_{X,x});
  \tag{21.6.7.4}
\]
pour deux fonctions méromorphes régulières $f$, $f'$ sur $X$, on a
\[
\label{IV.21.6.7.5}
  o_x(ff')=o_x(f)+o_x(f'),\qquad o_x(f^{-1})=-o_x(f)
  \tag{21.6.7.5}
\]
pour tout $x\in X^{(1)}$.
Les cycles $1$-codimensionnels
\[
  Z^+(f)=\sum_{x\in X^{(1)}}(o_x(f))^+\cdot\overline{\{x\}},\qquad
  Z^-(f)=\sum_{x\in X^{(1)}}(o_x(f))^-\cdot\overline{\{x\}}
\]
sont appelés respectivement le \emph{cycle des zéros} et le \emph{cycle des pôles} (ou \emph{cycle polaire}) de $f$, et l'on a évidemment $\operatorname{cyc}(\operatorname{div}(f))=Z^+(f)-Z^-(f)$.
Les cycles $1$-codimensionnels de la forme $\operatorname{cyc}(\operatorname{div}(f))$ sont appelés \emph{principaux} (ou \emph{linéairement équivalents à zéro}) et forment un sous-groupe $\mathfrak{Z}^1_\mathrm{princ}(X)$ de $\mathfrak{Z}^1(X)$.
Les sections au-dessus de $X$ de l'image $c(\sh{D}iv_X)\subset\mathscr{Z}_X^1$ sont
\oldpage[IV-4]{274}
appelées \emph{cycles localement principaux};
un tel cycle $Z$ est donc caractérisé par le fait que, pour tout $y\in X$, il y a un voisinage ouvert $U$ de $y$ dans $X$ et une fonction méromorphe régulière $f$ sur $U$ telle que $Z|U=\sum_{x\in U\cap X^{(1)}}o_x(f)(\overline{\{x\}}\cap U)$.
Si $Z=\sum_{x\in X^{(1)}}n_x\cdot\overline{\{x\}}$, il revient au même de dire que, pour tout $y\in X$, si l'on pose $T_y=\Spec(\sh{O}_{X,y})$ et $Z_y=\sum_{x\in X^{(1)}\cap T_y}n_x(\overline{\{x\}}\cap T_y)$, $Z_y$ est \emph{principal}, autrement dit il existe une fonction méromorphe régulière $g$ sur $T_y$ telle que $Z_y=\operatorname{cyc}(\operatorname{div}(g))$.
Cela résulte aussitôt en effet de la définition précédente et de \sref{IV.20.3.7} qui établit une correspondance biunivoque entre les fonctions méromorphes régulières sur $T_y$ et les germes de fonctions méromorphes régulières sur les voisinages ouverts de $y$ dans $X$ lorsque $X$ est localement noethérien.
On exprime encore le fait que $Z_y$ est principal en disant que $Z$ est \emph{principal au point $y$}.
L'ensemble des points où $Z$ est principal est évidemment \emph{ouvert}, en vertu de ce qui précède.

On déduit de l'homomorphisme canonique $\operatorname{cyc}:\operatorname{Div}(X)\to\mathfrak{Z}^1(X)$ un homomorphisme $\operatorname{Div}(X)/\operatorname{Div.princ}(X)\to\mathfrak{Z}^1(X)/\mathfrak{Z}^1_\mathrm{princ}(X)$, en vertu de la définition de $\mathfrak{Z}^1_\mathrm{princ}(X)$.
On dit que le groupe $\mathfrak{Z}^1(X)/\mathfrak{Z}^1_\mathrm{princ}(X)$ est le \emph{groupe des classes de cycles $1$-codimensionnels} de $X$ et on le note $\mathfrak{Cl}(X)$.
Deux éléments de $\mathfrak{Z}^1(X)$ ayant même image dans $\mathfrak{Cl}(X)$ sont appelés \emph{linéairement équivalents}.
\end{env}

\begin{env}[21.6.8]
\label{IV.21.6.8}
Considérons en particulier le cas où $X=\Spec(A)$, où $A$ est un anneau \emph{noethérien intègre et intégralement clos}.
Alors $X^{(1)}$ est l'ensemble des \emph{idéaux premiers de hauteur $1$} de $A$, et $\mathfrak{Z}^1(X)$ s'identifie donc au \emph{groupe des diviseurs} (au sens de N.~Bourbaki) de l'anneau de Krull $A$ (Bourbaki, \emph{Alg.\ comm.}, chap.~VII, \textsection~1, n\textsuperscript{o}~3, cor.\ du th.~2 et n\textsuperscript{o}~6, th.~3).

Comme d'autre part, les fonctions méromorphes régulières sur $X$ s'identifient alors aux éléments non nuls du corps des fractions $K$ de $A$, l'application $f\mapsto\operatorname{cyc}(\operatorname{div}(f))$ de $M(X)$ dans $\mathfrak{Z}^1(X)$ s'identifie à l'application notée $f\mapsto\operatorname{div}(f)$ dans Bourbaki (\emph{loc.\ cit.}, \textsection~1, n\textsuperscript{o}~1);
$\mathfrak{Z}^1_\mathrm{princ}(X)$ s'identifie donc au \emph{groupe des diviseurs principaux} de $A$ au sens de Bourbaki, et $\mathfrak{Cl}(X)$ au \emph{groupe des classes de diviseurs} de $A$ au sens de Bourbaki (\emph{loc.\ cit.}, \textsection~1, n\textsuperscript{o}~2 et n\textsuperscript{o}~10).
\end{env}

\begin{theorem}[21.6.9]
\EGAReaderTarget{ega4.c04.thm.21-6-9}{IV.21.6.9}
Soit $X$ un préschéma localement noethérien normal.
\begin{enumerate}[label=(\roman*)]
  \item L'homomorphisme canonique $\operatorname{cyc}:\operatorname{Div}(X)\to\mathfrak{Z}^1(X)$ est injectif et son image est formée des cycles localement principaux.
  \item Les conditions suivantes sont équivalentes:
  \begin{enumerate}[label=\emph{\alph*)}]
    \item L'homomorphisme $\operatorname{cyc}:\operatorname{Div}(X)\to\mathfrak{Z}^1(X)$ est bijectif.
    \item Tout cycle $1$-codimensionnel est localement principal.
    \item Pour tout $x\in X$, l'anneau local $\sh{O}_{X,x}$ est factoriel.
  \end{enumerate}
  (On dit alors que $X$ est un préschéma \emph{localement factoriel}.)
\end{enumerate}
\end{theorem}

\begin{proof}
\begin{enumerate}[label=(\roman*)]
  \item Il suffit de démontrer que
\[
\label{IV.21.6.9.1}
  \operatorname{cyc}^{-1}(\mathfrak{Z}^{1+}(X))=\operatorname{Div}^+(X),
  \tag{21.6.9.1}
\]
puisque l'on a $\operatorname{Div}^+(X)\cap(-\operatorname{Div}^+(X))=0$ et $\mathfrak{Z}^{1+}(X)\cap(-\mathfrak{Z}^{1+}(X))=0$.
On est donc ramené à prouver que si $D$ est un diviseur sur $X$ tel que $\operatorname{mult}_x(D)\geq0$ pour tout $x\in X^{(1)}$, on a $D\geq0$.
Or, pour tout $x\notin X^{(1)}$, l'anneau local $\sh{O}_{X,x}$ est intègre et intégralement clos et
\oldpage[IV-4]{275}
de dimension $0$ ou $\geq2$, donc on a $\operatorname{prof}(\sh{O}_{X,x})=0$ ou $\operatorname{prof}(\sh{O}_{X,x})\geq2$ \sref[0]{0.16.5.1}.
D'autre part, aux points $x\in X^{(1)}$, l'anneau $\sh{O}_{X,x}$ est un anneau de valuation discrète (\sref[II]{II.7.1.6}), donc $\operatorname{prof}(\sh{O}_{X,x})=1$ \sref[0]{0.16.5.1};
les seuls points de $X$ tels que $\operatorname{prof}(\sh{O}_{X,x})=1$ sont donc les points de $X^{(1)}$, et la conclusion résulte par suite de \sref{IV.21.1.8}.

  \item L'équivalence de \emph{a)} et \emph{b)} est claire en vertu de (i).
D'après la caractérisation des cycles localement principaux donnée dans \sref{IV.21.6.7}, et la relation entre cycles $1$-codimensionnels sur $\Spec(A)$ et diviseurs (au sens de Bourbaki) de l'anneau $A$ lorsque $A$ est un anneau noethérien intégralement clos \sref{IV.21.6.8}, la condition \emph{b)} équivaut à dire que pour tout $x\in X$, tout diviseur de l'anneau $\sh{O}_{X,x}$ est principal, autrement dit que l'anneau $\sh{O}_{X,x}$ est factoriel (Bourbaki, \emph{Alg.\ comm.}, chap.~VII, \textsection~3, n\textsuperscript{o}~1), d'où l'équivalence de \emph{b)} et \emph{c)}.
\end{enumerate}
\end{proof}

\begin{env}[21.6.9.2]
\label{IV.21.6.9.2}
Lorsque $A$ est un anneau noethérien \emph{factoriel}, il est clair que $X=\Spec(A)$ est localement factoriel (Bourbaki, \emph{Alg.\ comm.}, chap.~VII, \textsection~3, n\textsuperscript{o}~4, prop.~3).
Si, dans ce cas, on écrit un élément $f\neq0$ du corps des fractions $K$ de $A$ sous la forme $r/s$, où $r$ et $s$ sont deux éléments \emph{étrangers} de $A$, dont les diviseurs sont bien déterminés (\emph{loc.\ cit.}, \textsection~3, n\textsuperscript{o}~3), ces diviseurs s'identifient respectivement au \emph{diviseur des zéros} et au \emph{diviseur des pôles} de $f$ \sref{IV.21.6.7}.
\end{env}

\begin{corollary}[21.6.10]
\label{IV.21.6.10}
Soit $X$ un préschéma localement noethérien normal.
\begin{enumerate}[label=(\roman*)]
  \item Il existe un homomorphisme canonique injectif
\[
\label{IV.21.6.10.1}
  \operatorname{Pic}(X)\to\mathfrak{Cl}(X).
  \tag{21.6.10.1}
\]

  \item Si $X$ est localement factoriel, l'homomorphisme \sref{IV.21.6.10.1} est bijectif, et réciproquement.
\end{enumerate}
\end{corollary}

\begin{proof}
On a vu \sref{IV.21.6.7} que l'image de $\operatorname{Div.princ}(X)$ par l'homomorphisme $\operatorname{cyc}:\operatorname{Div}(X)\to\mathfrak{Z}^1(X)$ est $\mathfrak{Z}^1_\mathrm{princ}(X)$;
on déduit donc de \sref{IV.21.6.9} que l'homomorphisme $\operatorname{Div}(X)/\operatorname{Div.princ}(X)\to\mathfrak{Z}^1(X)/\mathfrak{Z}^1_\mathrm{princ}(X)=\mathfrak{Cl}(X)$ déduit de $\operatorname{cyc}$ est injectif, et qu'il est bijectif si et seulement si $X$ est localement factoriel.
La conclusion résulte donc de ce que, lorsque $X$ est localement noethérien et réduit, l'homomorphisme canonique $\operatorname{Div}(X)/\operatorname{Div.princ}(X)\to\operatorname{Pic}(X)$ \sref{IV.21.3.3.1} est bijectif \sref{IV.21.3.4}, (ii).
\end{proof}

\begin{corollary}[21.6.11]
\label{IV.21.6.11}
Soit $X$ un préschéma localement noethérien et localement factoriel.
Alors le faisceau $\sh{D}iv_X$ est flasque, et pour tout ouvert $U$ de $X$, l'homomorphisme canonique $\operatorname{Pic}(X)\to\operatorname{Pic}(U)$ est surjectif.
\end{corollary}

\begin{proof}
Tenant compte de \sref{IV.21.6.9}, (ii), la première assertion équivaut à dire que le faisceau $\mathscr{Z}_X^1$ est flasque, ce qui a été vu plus haut \sref{IV.21.6.3}.
Pour tout ouvert $U$ de $X$, l'homomorphisme canonique $\mathfrak{Z}^1(X)\to\mathfrak{Z}^1(U)$ est donc surjectif;
comme en vertu de \sref{IV.21.6.10}, l'homomorphisme $\operatorname{Pic}(X)\to\operatorname{Pic}(U)$ s'identifie canoniquement à $\mathfrak{Z}^1(X)/\mathfrak{Z}^1_\mathrm{princ}(X)\to\mathfrak{Z}^1(U)/\mathfrak{Z}^1_\mathrm{princ}(U)$, il est aussi surjectif.
\end{proof}

\begin{proposition}[21.6.12]
\label{IV.21.6.12}
Soit $X$ un préschéma noethérien réduit.
Soit $(U_\lambda)_{\lambda\in L}$ une famille filtrante décroissante d'ouverts de $X$ vérifiant les conditions suivantes:
\begin{enumerate}
  \item[\rm 1\textsuperscript{o}] Si $Y_\lambda=X-U_\lambda$, on a $\codim(Y_\lambda,X)\geq2$ pour tout $\lambda\in L$.
  \item[\rm 2\textsuperscript{o}] Pour tout $x\in\bigcap_{\lambda\in L}U_\lambda$, l'anneau $\sh{O}_{X,x}$ est factoriel.
\end{enumerate}
Alors il existe des isomorphismes canoniques
\[
\label{IV.21.6.12.1}
  \varinjlim\operatorname{Div}(U_\lambda)\simto\mathfrak{Z}^1(X),\qquad
  \varinjlim\operatorname{Pic}(U_\lambda)\simto\mathfrak{Cl}(X)
  \tag{21.6.12.1}
\]
tels que, pour tout ouvert $V$ de $X$, les diagrammes
\[
\label{IV.21.6.12.2}
  \xymatrix{
    \varinjlim\operatorname{Div}(U_\lambda) \ar[r]^-\sim \ar[d] & \mathfrak{Z}^1(X) \ar[d] & & \varinjlim\operatorname{Pic}(U_\lambda) \ar[r]^-\sim \ar[d] & \mathfrak{Cl}(X) \ar[d]\\
    \varinjlim\operatorname{Div}(U_\lambda\cap V) \ar[r]^-\sim & \mathfrak{Z}^1(V) & & \varinjlim\operatorname{Pic}(U_\lambda\cap V) \ar[r]^-\sim & \mathfrak{Cl}(V)
  }
  \tag{21.6.12.2}
\]
soient commutatifs.
\end{proposition}

\begin{proof}
L'hypothèse 1\textsuperscript{o} sur $U_\lambda$ implique que pour tout $\lambda\in L$, on a $X^{(1)}\subset U_\lambda$ \sref{IV.5.1.3.1} donc l'homomorphisme de restriction $\mathfrak{Z}^1(X)\to\mathfrak{Z}^1(U_\lambda)$ est bijectif et par suite on a un isomorphisme canonique $\varinjlim\mathfrak{Z}^1(U_\lambda)\simto\mathfrak{Z}^1(X)$.
Les homomorphismes canoniques $\operatorname{cyc}:\operatorname{Div}(U_\lambda)\to\mathfrak{Z}^1(U_\lambda)$ définissent donc, par passage à la limite inductive, le premier des homomorphismes canoniques \sref{IV.21.6.12.1}, et le second s'en déduit par passage aux quotients.
En outre, il résulte de la condition 1\textsuperscript{o} que les $U_\lambda$ sont denses dans $X$, donc schématiquement denses puisque $X$ est réduit \sref{IV.11.10.4}, et par suite toute fonction méromorphe sur $X$ est entièrement déterminée par sa restriction à chaque $U_\lambda$;
on déduit aussitôt de là que dans l'isomorphisme $\mathfrak{Z}^1(X)\simto\mathfrak{Z}^1(U_\lambda)$, l'image de $\mathfrak{Z}^1_\mathrm{princ}(X)$ est $\mathfrak{Z}^1_\mathrm{princ}(U_\lambda)$, donc on a aussi un isomorphisme canonique $\varinjlim\mathfrak{Cl}(U_\lambda)\simto\mathfrak{Cl}(X)$.
Le second des homomorphismes canoniques \sref{IV.21.6.12.1} sera donc un isomorphisme si le premier l'est, et il reste donc à prouver ce dernier point, la commutativité des diagrammes \sref{IV.21.6.12.2} étant triviale.

\oldpage[IV-4]{276}
Montrons d'abord que l'homomorphisme $\varinjlim\operatorname{Div}(U_\lambda)\to\mathfrak{Z}^1(X)$ est \emph{injectif}.
Posons $T=\bigcap_\lambda U_\lambda$, et notons que les $U_\lambda$ forment un \emph{système fondamental de voisinages} de $T$.
En effet, pour tout ouvert $V\supset T$, $X-V$ est un sous-espace fermé de $X$, donc un espace noethérien dont toute partie fermée irréductible admet un point générique;
comme les ensembles $(X-V)\cap(X-U_\lambda)$ sont fermés dans $X-V$, forment une famille filtrante croissante et ont pour réunion $X-V$, notre assertion résulte de \sref[0]{0.9.2.4}.
Cela étant, il s'agit de voir que si $D\in\operatorname{Div}(U_\lambda)$ est tel que $\operatorname{cyc}(D)=0$ dans $\mathfrak{Z}^1(U_\lambda)$, alors il existe $\mu\geq\lambda$ tel que $D|U_\mu=0$.
En vertu de ce qui précède, il suffira de voir que pour tout $x\in T$, on a $D_x=0$;
en effet, il y aura alors un voisinage $W(x)$ de $x$ dans $X$ tel que $D|W(x)=0$, et la réunion des $W(x)$ pour $x\in T$ contient un $U_\mu$.
Compte tenu de \sref{IV.21.4.6}, on est donc ramené au cas où $X=\Spec(\sh{O}_{X,x})$ avec $x\in T$;
mais par hypothèse $\sh{O}_{X,x}$ est factoriel, donc intègre et intégralement clos, et $\Spec(\sh{O}_{X,x})$ est alors normal;
donc la conclusion résulte de \sref{IV.21.6.9}, (i).

\oldpage[IV-4]{277}
Pour prouver que l'homomorphisme $\varinjlim\operatorname{Div}(U_\lambda)\to\mathfrak{Z}^1(X)$ est bijectif, il suffira de prouver de même que pour tout $Z\in\mathfrak{Z}^1(X)$ et tout $x\in T$, il y a un voisinage $W(x)$ de $x$ dans $X$ et un diviseur $D^{(x)}$ sur $W(x)$ tel que $\operatorname{cyc}(D^{(x)})=Z|W(x)$.
En effet, on pourra alors recouvrir l'ensemble quasi-compact $T$ par un nombre fini de $W(x_i)$, et en vertu de la première partie de la démonstration, puisque les couples $(i,j)$ sont en nombre fini, il existera un indice $\lambda$ tel que dans chacun des $W(x_i)\cap W(x_j)\cap U_\lambda$, les restrictions de $D^{(x_i)}$ et $D^{(x_j)}$ coïncident;
comme on peut supposer en outre que $U_\lambda$ est contenu dans la réunion des $W(x_i)$, on voit que les restrictions $D^{(x_i)}|(W(x_i)\cap U_\lambda)$ sont les restrictions d'un même diviseur $D\in\operatorname{Div}(U_\lambda)$ qui sera tel que $\operatorname{cyc}(D)=Z|U_\lambda$.
On est donc ramené encore au cas où $X=\Spec(\sh{O}_{X,x})$ avec $x\in T$, mais comme $\sh{O}_{X,x}$ est factoriel, il en est de même de ses localisés (Bourbaki, \emph{Alg.\ comm.}, chap.~VII, \textsection~3, n\textsuperscript{o}~4, prop.~3), et il suffit d'appliquer \sref{IV.21.6.9}, (ii).
\end{proof}

\begin{corollary}[21.6.13]
\label{IV.21.6.13}
Soit $A$ un anneau local noethérien intègre et intégralement clos et de dimension $\geq2$.
Posons $X=\Spec(A)$, et soient $a$ le point fermé de $X$, $U=X-\{a\}$.
Pour que $A$ soit factoriel, il faut et il suffit que $U$ soit localement factoriel et que $\operatorname{Pic}(U)=0$.
\end{corollary}

\begin{proof}
En effet, dire que $A$ est factoriel équivaut à dire que $\mathfrak{Cl}(X)=0$ (Bourbaki, \emph{Alg.\ comm.}, chap.~VII, \textsection~1, n\textsuperscript{o}~4, cor.\ du th.~2 et \textsection~3, n\textsuperscript{o}~2, th.~1);
il suffit donc d'utiliser l'existence du second isomorphisme \sref{IV.21.6.12.1}, en prenant la famille $(U_\lambda)$ restreinte au seul ouvert $U$.
\end{proof}

\begin{corollary}[21.6.14]
\label{IV.21.6.14}
Soient $A$ un anneau local noethérien de dimension $\geq2$, $X=\Spec(A)$, $a$ le point fermé de $X$, $U=X-\{a\}$.
Les conditions suivantes sont équivalentes:
\begin{enumerate}[label=\emph{\alph*)}]
  \item $A$ est factoriel.
  \item $\operatorname{Pic}(U)=0$ et $\operatorname{prof}(A)\geq2$ (conditions que nous exprimerons plus loin \sref{IV.21.13} en disant que l'anneau $A$ est \emph{parafactoriel}), et $U$ est localement factoriel.
\end{enumerate}
\end{corollary}

\begin{proof}
Il est clair que si $A$ est factoriel, il est normal, donc $\operatorname{prof}(A)\geq2$ puisque $\dim(A)\geq2$ \sref[0]{0.16.5.1}, et \sref{IV.21.6.13} montre que \emph{a)} implique \emph{b)}.
Inversement, si \emph{b)} est vérifiée, il suffit de voir que $A$ est intégralement clos pour en conclure par \sref{IV.21.6.13} que \emph{b)} entraîne \emph{a)}.
En vertu du critère de Serre \sref{IV.5.8.6}, il suffit de vérifier pour $X$ les conditions $(\mathrm{R}_1)$ et $(\mathrm{S}_2)$.
Or, $U$ étant localement factoriel vérifie ces conditions, et l'hypothèse $\operatorname{prof}(A)\geq2$ entraîne que $X$ les vérifie également.
\end{proof}
\subsection{Interprétation des cycles positifs 1-codimensionnels en termes de sous-préschémas}
\EGAReaderTarget{ega4.c04.subsec.21-7}{subsection:IV.21.7}
\label{IV.21.7}

\begin{env}[21.7.1]
\label{IV.21.7.1}
Soient $X$ un préschéma localement noethérien, $C=\sum_{x\in X^{(1)}}n_x\cdot\overline{\{x\}}$ un cycle $1$-codimensionnel positif (de sorte que l'on a $n_x\geq0$ pour tout $x\in X^{(1)}$, et $n_x=0$ sauf en un ensemble localement fini de points).
Désignons par $Y(C)$ le sous-préschéma fermé de $X$, \emph{image fermée} (\sref[I]{I.9.5.3} et \sref{IV.9.5.1}) du préschéma \emph{somme} $Y'(C)=\coprod_{x\in X^{(1)}}\Spec(\sh{O}_{X,x}/\mathfrak{m}_x^{n_x})$ par le morphisme canonique, et par $\sh{J}_X(C)$ (ou $\sh{J}(C)$) l'Idéal de $\sh{O}_X$ définissant $Y(C)$.
\end{env}

\begin{proposition}[21.7.2]
\label{IV.21.7.2}
Soit $X$ un préschéma localement noethérien vérifiant la condition $(\mathrm{R}_1)$ \sref{IV.5.8.2}.
Pour qu'un sous-préschéma fermé $Y$ de $X$ soit de la forme $Y(C)$, où $C$ est un cycle $1$-codimensionnel positif, il faut et il suffit qu'il vérifie les deux conditions suivantes:
\begin{enumerate}[label=(\roman*)]
  \item $Y$ est purement de codimension $1$.
\oldpage[IV-4]{278}
  \item $Y$ vérifie la condition $(\mathrm{S}_1)$, autrement dit \sref{IV.5.7.5} n'a pas de cycle premier associé immergé.
\end{enumerate}
On a alors
\[
\label{IV.21.7.2.1}
  \operatorname{mult}_x(C)=\operatorname{long}\sh{O}_{Y,x}.
  \tag{21.7.2.1}
\]
L'application $C\mapsto Y(C)$ est une bijection de $\mathfrak{Z}^{1+}(X)$ sur l'ensemble des sous-préschémas fermés de $X$ vérifiant les conditions (i) et (ii).
\end{proposition}

\begin{proof}
Les conditions sont nécessaires (sans supposer que $X$ vérifie $(\mathrm{R}_1)$).
En effet, la question étant locale sur $X$, on peut supposer que $X=\Spec(A)$, où $A$ est un anneau noethérien;
on a alors $Y(C)=\Spec(A/\mathfrak{J})$, où $\mathfrak{J}$ est, par définition \sref[I]{I.9.5.1}, le noyau de l'homomorphisme canonique $A\to\bigoplus_i A_{\mathfrak{p}_i}/\mathfrak{p}_i^{n_i}A_{\mathfrak{p}_i}$, où les $\mathfrak{p}_i$ sont les idéaux premiers correspondant aux points $x\in X^{(1)}$ pour lesquels $n_x>0$, et où l'on a posé $n_i=n_x$.
L'image réciproque $\mathfrak{q}_i$ de $\mathfrak{p}_i^{n_i}A_{\mathfrak{p}_i}$ dans $A$ est un idéal $\mathfrak{p}_i$-primaire et on a $\mathfrak{J}=\bigcap_i\mathfrak{q}_i$.
D'ailleurs, comme les $x\in X^{(1)}$ sont tels que $\overline{\{x\}}$ est de codimension $1$, aucun point de $X^{(1)}$ ne peut être adhérent à un autre point de $X^{(1)}$.
Donc les $\mathfrak{p}_i$ sont les idéaux premiers minimaux de $A/\mathfrak{J}$ et l'ensemble des $\mathfrak{p}_i$ est égal à $\operatorname{Ass}(A/\mathfrak{J})$, ce qui prouve les conditions (i) et (ii).

Les conditions sont suffisantes.
Désignant par $\sh{J}$ l'Idéal de $\sh{O}_X$ définissant $Y$, l'hypothèse (ii) implique que $\operatorname{Ass}(\sh{O}_X/\sh{J})$ est identique à l'ensemble $F$ des points maximaux de $Y$, et l'hypothèse (i) implique que $F$ est contenu dans $X^{(1)}$;
donc \sref{IV.3.2.6} $\sh{O}_X/\sh{J}$ s'identifie à un sous-$\sh{O}_X$-Module de $\bigoplus_{x\in F}\mathscr{G}(x)$, où pour chaque $x\in F$, $\mathscr{G}(x)$ est un $\sh{O}_X$-Module irredondant tel que $\operatorname{Ass}(\mathscr{G}(x))=\{x\}$.
Or, puisque $X$ vérifie $(\mathrm{R}_1)$, pour chaque $x\in F$, $\sh{O}_{X,x}$ est un anneau de valuation discrète et par suite les idéaux primaires non nuls de cet anneau sont les puissances $\mathfrak{m}_x^{n_x}$ de l'idéal maximal;
supposant encore que $X=\Spec(A)$, on en conclut que $\sh{J}=\widetilde{\mathfrak{J}}$, où $\widetilde{\mathfrak{J}}=\bigcap_i\mathfrak{q}_i$, les $\mathfrak{q}_i$ étant les images réciproques dans $A$ d'idéaux $\mathfrak{p}_i^{n_i}A_{\mathfrak{p}_i}$, où les $\mathfrak{p}_i$ correspondent aux points de $F$.
On voit donc bien que $Y$ est de la forme $Y(C)$.
\end{proof}

\begin{corollary}[21.7.3]
\label{IV.21.7.3}
Soit $X$ un préschéma localement noethérien.
Pour tout diviseur positif $D$ sur $X$, tel que $X$ vérifie $(\mathrm{R}_1)$ aux points maximaux du sous-préschéma fermé $Y(D)$ \sref{IV.21.2.12}, $Y(D)$ majore le sous-préschéma fermé $Y(\operatorname{cyc}(D))$ \sref{IV.21.7.1}, et a même espace sous-jacent;
pour que ces deux sous-préschémas soient égaux, il faut et il suffit que $Y(D)$ vérifie la condition $(\mathrm{S}_1)$, ou encore que, pour tout $z\in Y(D)$ distinct d'un point maximal de $Y(D)$, on ait $\operatorname{prof}(\sh{O}_{X,z})\geq2$ (condition toujours vérifiée lorsque $X$ est \emph{normal}).
\end{corollary}

\begin{proof}
En effet, la question étant locale, on peut toujours supposer que $\sh{J}_X(D)=t\sh{O}_X$, $t$ étant une section régulière de $\sh{O}_X$ au-dessus de $X$;
en tout point maximal $x$ de $Y(D)$, appartenant nécessairement à $X^{(1)}$, $\sh{O}_{X,x}$ est par hypothèse un anneau de valuation discrète, donc $t_x\sh{O}_{X,x}=\mathfrak{m}_x^{n_x}$ où $n_x=\operatorname{mult}_x(D)$.
On peut supposer $X=\Spec(A)$, et alors, si $\sh{J}_X(\operatorname{cyc}(D))=\widetilde{\mathfrak{J}}$, il résulte de ce qui précède et de \sref{IV.21.7.1} que $\mathfrak{J}=tA$ et $\widetilde{\mathfrak{J}}$ ont mêmes \emph{idéaux primaires non immergés}, donc $\widetilde{\mathfrak{J}}\subset\mathfrak{J}$, puisque $\widetilde{\mathfrak{J}}$ est l'intersection
\oldpage[IV-4]{279}
de ces idéaux primaires \sref{IV.21.7.2}.
Cela prouve que $Y(D)$ majore $Y(\operatorname{cyc}(D))$ et que ces deux sous-préschémas sont égaux si et seulement si $Y(D)$ n'a pas de cycle premier associé immergé (autrement dit vérifie $(\mathrm{S}_1)$).
Comme pour tout $x\in Y(D)$, il y a par hypothèse un voisinage ouvert $U$ de $x$ dans $X$ et un élément régulier $t\in\Gamma(U,\sh{O}_X)$ tel que $\sh{O}_{Y(D)}|(U\cap Y(D))$ soit restriction à $Y(D)\cap U$ de $\sh{O}_U/t\sh{O}_U$, dire que $Y(D)$ vérifie $(\mathrm{S}_1)$ signifie aussi qu'en tout point $z\in Y(D)$ non maximal, on a $\operatorname{prof}(\sh{O}_{X,z}/t_z\sh{O}_{X,z})\geq1$, c'est-à-dire \sref[0]{0.16.4.6} $\operatorname{prof}(\sh{O}_{X,z})\geq2$.
L'assertion concernant le cas où $X$ est normal est alors immédiate puisque $X$ vérifie $(\mathrm{S}_2)$ et qu'aux points non maximaux $z$ de $Y(D)$ on a $\dim(\sh{O}_{X,z})\geq2$ \sref[0]{0.16.3.4}.
\end{proof}

\begin{env}[21.7.3.1]
\label{IV.21.7.3.1}
On voit donc que lorsque $X$ est \emph{normal}, $\operatorname{Div}^+(X)$ s'identifie canoniquement à l'ensemble des sous-préschémas fermés $Y$ de $X$ vérifiant les conditions (i) et (ii) de \sref{IV.21.7.2} et \emph{régulièrement immergés} dans $X$.
\end{env}

\begin{proposition}[21.7.4]
\label{IV.21.7.4}
Soit $X$ un préschéma localement noethérien, réduit en chacun de ses points isolés.
Les conditions suivantes sont équivalentes:
\begin{enumerate}
  \item[\emph{a)}] L'homomorphisme canonique $c:\sh{D}iv_X\to\mathscr{Z}_X^1$ \sref{IV.21.6.4} est un isomorphisme de faisceaux de groupes ordonnés.
  \item[\emph{a$'$)}] Tout cycle premier $1$-codimensionnel de $X$ est l'image par $\operatorname{cyc}$ d'un diviseur positif, et l'homomorphisme canonique $c:\sh{D}iv_X\to\mathscr{Z}_X^1$ est injectif.
  \item[\emph{a$''$)}] Pour tout sous-préschéma fermé intègre $Y$ de $X$, de codimension $1$, l'immersion canonique $Y\to X$ est régulière.
  \item[\emph{b)}] $X$ est normal et l'homomorphisme $\operatorname{cyc}:\operatorname{Div}(X)\to\mathfrak{Z}^1(X)$ est bijectif.
  \item[\emph{c)}] $X$ est localement factoriel.
\end{enumerate}
\end{proposition}

\begin{proof}
L'équivalence de \emph{b)} et \emph{c)} a été démontrée dans \sref{IV.21.6.9}, ainsi que le fait que \emph{c)} entraîne \emph{a)}.
De plus \emph{b)} entraîne \emph{a$''$)} en vertu de \sref{IV.21.7.3.1}, et \emph{a)} implique trivialement \emph{a$'$)}.
Il reste à voir que \emph{a$'$)} ou \emph{a$''$)} entraîne \emph{c)}.

Supposons donc vérifiée la condition \emph{a$'$)}, et montrons d'abord que $X$ est normal.
Notons en premier lieu que si $X$ vérifie \emph{a$'$)}, il en est de même de tout schéma local $\Spec(\sh{O}_{X,x})$.
Considérons alors $x\in X^{(1)}$, de sorte que $A=\sh{O}_{X,x}$ est un anneau local noethérien de dimension $1$.
Appliquant la condition \emph{a$'$)} à $\Spec(A)$ et au cycle premier $1$-codimensionnel formé du point fermé de $\Spec(A)$, on voit que dans $A$ l'idéal maximal est engendré par un seul élément régulier dans $A$;
donc \sref[0]{0.17.1.1} $A$ est un anneau régulier.
Comme les anneaux localisés $A_{\mathfrak p}$ sont aussi réguliers \sref[0]{0.17.3.2}, on voit que $X$ est régulier en tous ses points maximaux non isolés;
comme il est aussi réduit (donc régulier) en ses points isolés par hypothèse, on en conclut que $X$ vérifie la condition $(\mathrm{R}_1)$.
Montrons en outre que $X$ vérifie aussi $(\mathrm{S}_1)$, autrement dit que pour tout $x\in X$ tel que $\dim(\sh{O}_{X,x})\geq1$, on a $\mathfrak{m}_x\notin\operatorname{Ass}(\sh{O}_{X,x})$ \sref[0]{0.16.4.6}.
En effet, l'hypothèse \emph{a$'$)} appliquée à $X_1=\Spec(\sh{O}_{X,x})$ montre qu'il existe sur ce préschéma au moins un diviseur non nul, autrement dit que l'on a $\sh{M}_{X_1}^*\neq\sh{O}_{X_1}^*$, et il suffit d'appliquer \sref{IV.21.1.10}.
On déduit donc déjà de ces résultats que $X$ est réduit \sref{IV.5.8.5}.
Prouvons ensuite qu'il vérifie la condition $(\mathrm{S}_2)$ (ce qui établira que $X$ est normal, en vertu du critère de Serre \sref{IV.5.8.6}).
Raisonnons par l'absurde en supposant que l'ensemble $E$ des points $x\in X$ où $X$ ne vérifie pas $(\mathrm{S}_2)$ soit non vide, et soit $x\in E$ un point pour lequel $\dim(\sh{O}_{X,x})$ est la plus petite possible;
puisque $X$ vérifie $(\mathrm{S}_1)$, on a $\dim(\sh{O}_{X,x})\geq2$.
Dans $X_1=\Spec(\sh{O}_{X,x})$, l'ouvert $U=X_1-\{x\}$ vérifie $(\mathrm{S}_2)$;
montrons que $X_1$ vérifie $(\mathrm{S}_2)$, de sorte que l'on aura abouti à une contradiction.
Il suffit, en vertu de \sref{IV.5.10.5}, de montrer que toute section $f$ de $\sh{O}_{X_1}$ au-dessus de $U$ se prolonge en une section de $\sh{O}_{X_1}$ au-dessus de $X_1$.
Comme $X_1$ est réduit et $U$ dense dans $X_1$, $U$ est schématiquement dense dans $X_1$, et par suite $f$ est la restriction à $U$ d'une section méromorphe régulière $g\in M(X_1)$.
En outre, comme $\dim(\sh{O}_{X,x})\geq2$, on a $\operatorname{cyc}(\operatorname{div}(g))\geq0$ puisque $g|U=f$;
comme l'homomorphisme $\operatorname{cyc}$ est injectif, on a nécessairement $\operatorname{div}(g)\geq0$, donc \sref{IV.21.6.4} $g$ est une section de $\sh{O}_{X_1}$ au-dessus de $X_1$, ce qui établit notre assertion.

Pour montrer que \emph{a$'$)} implique \emph{c)}, remarquons alors que la condition \emph{a$'$)} implique que l'homomorphisme canonique $\operatorname{cyc}:\operatorname{Div}(X)\to\mathfrak{Z}^1(X)$ est surjectif;
puisque $X$ est normal, il suffit d'appliquer \sref{IV.21.6.9}.

Prouvons maintenant que \emph{a$''$)} entraîne \emph{c)}.
Il résulte de \sref{IV.21.6.5} que \emph{a$''$)} entraîne que tout cycle premier $1$-codimensionnel sur $X$ est l'image par $\operatorname{cyc}$ d'un diviseur positif;
la première partie du raisonnement fait ci-dessus prouve donc encore que $X$ vérifie $(\mathrm{R}_1)$ et $(\mathrm{S}_1)$.
Il reste encore à voir que $X$ est normal (la fin du raisonnement étant alors la même), c'est-à-dire que si $x\in X$ est tel que $\dim(\sh{O}_{X,x})\geq2$, on a $\operatorname{prof}(\sh{O}_{X,x})\geq2$.
Or, si $y$ est une générisation de $x$ telle que $\overline{\{y\}}$ soit de codimension $1$ dans $X$, le sous-préschéma réduit $Y$ de $X$ ayant $\overline{\{y\}}$ pour espace
\oldpage[IV-4]{280}
sous-jacent est intègre, donc régulièrement immergé dans $X$ par hypothèse.
Cela entraîne qu'il existe un élément régulier $t$ de $\sh{O}_{X,x}$ tel que l'idéal $t\sh{O}_{X,x}$ soit l'idéal premier définissant le préschéma $Y_1$ image réciproque de $Y$ dans $X_1=\Spec(\sh{O}_{X,x})$.
Comme $\sh{O}_{X,x}/t\sh{O}_{X,x}$ est intègre, cela prouve que $\operatorname{prof}(\sh{O}_{X,x})\geq2$ et termine la preuve de \sref{IV.21.7.4}.
\end{proof}

\begin{remark}[21.7.5]
\label{IV.21.7.5}
On ne peut dans \sref{IV.21.7.4} remplacer la condition \emph{a$'$)} par la condition plus faible que tout cycle premier $1$-codimensionnel de $X$ soit l'image par $\operatorname{cyc}$ d'un diviseur positif.
C'est ce que montre l'exemple où $X$ est le schéma affine défini dans \sref{IV.14.1.5} (« réunion d'un plan et d'une droite qui le coupe ») \emph{qui n'est pas normal};
avec les notations introduites dans \sref{IV.14.1.5}, les cycles premiers $1$-codimensionnels de $X$ sont ceux du plan $X_1$ et les points fermés de la droite $X_2$ distincts du point commun à $X_1$ et $X_2$;
si $t_1$, $t_2$, $t_3$ sont les images de $T_1$, $T_2$, $T_3$ dans $A/\mathfrak a$, on voit donc que les cycles premiers $1$-codimensionnels de $X$ sont définis par les idéaux premiers monogènes de générateurs $P(t_1,t_2)$ ($P$ irréductible dans $K[T_1,T_2]$) ou $t_3-a$, avec $a\neq0$ dans $K$;
ces éléments étant réguliers dans $A/\mathfrak a$, tout cycle $1$-codimensionnel est bien l'image canonique d'un diviseur positif.
\end{remark}

\begin{corollary}[21.7.6]
\label{IV.21.7.6}
Soit $X$ un préschéma localement noethérien, réduit en chacun de ses points isolés.
Les conditions suivantes sont équivalentes:
\begin{enumerate}
  \item[\emph{a)}] L'homomorphisme canonique $c:\sh{D}iv_X\to\mathscr{Z}_X^1$ est un isomorphisme de faisceaux de groupes ordonnés, et on a $\operatorname{Div}(X)=\operatorname{Div.princ}(X)$.
  \item[\emph{a$'$)}] L'homomorphisme canonique $c:\sh{D}iv_X\to\mathscr{Z}_X^1$ est injectif, et tout cycle premier $1$-codimensionnel est l'image par $\operatorname{cyc}$ d'un diviseur positif principal, autrement dit est de la forme $\operatorname{cyc}(\operatorname{div}(f))$, où $f$ est une section régulière de $\sh{O}_X$ au-dessus de $X$.
  \item[\emph{a$''$)}] Pour tout sous-préschéma fermé intègre $Y$ de $X$, de codimension $1$, il existe une section régulière $f$ de $\sh{O}_X$ au-dessus de $X$ telle que $Y$ soit défini par l'Idéal $f\sh{O}_X$.
  \item[\emph{b)}] $X$ est normal et tout cycle premier $1$-codimensionnel sur $X$ est principal.
  \item[\emph{c)}] $X$ est localement factoriel et $\operatorname{Pic}(X)=0$.
  \item[\emph{d)}] $X$ est normal, et pour tout ouvert $U$ de $X$, on a $\operatorname{Pic}(U)=0$.
\end{enumerate}
\end{corollary}

\begin{proof}
L'équivalence de \emph{a)}, \emph{a$'$)}, \emph{a$''$)}, \emph{b)} et \emph{c)} résulte aussitôt de \sref{IV.21.7.4} et de \sref{IV.21.6.11}.
En outre, il résulte aussitôt de \emph{a$''$)} que tout ouvert non vide $U$ de $X$ vérifie encore les mêmes conditions, autrement dit ces conditions impliquent \emph{d)}.
Reste à voir que \emph{d)} entraîne \emph{b)}.
Or, désignons par $U_\lambda$ les ouverts complémentaires des réunions finies d'ensembles de la forme $\overline{\{x\}}$, où $\dim(\sh{O}_{X,x})\geq2$;
il est clair que les $U_\lambda$ forment une famille filtrante décroissante et que pour tout $x\in\bigcap_\lambda U_\lambda$, on a $\dim(\sh{O}_{X,x})\leq1$, donc $\sh{O}_{X,x}$ est un anneau principal en vertu de l'hypothèse.
On peut donc appliquer à la famille $(U_\lambda)$ la proposition \sref{IV.21.6.12}, qui montre que $\mathfrak{Cl}(X)$ est isomorphe à $\varinjlim\operatorname{Pic}(U_\lambda)$, donc $\mathfrak{Cl}(X)=0$ en vertu de l'hypothèse, ce qui prouve \emph{b)} par définition.
\end{proof}

\begin{remark}[21.7.7]
\label{IV.21.7.7}
Lorsque $X=\Spec(A)$, où $A$ est un anneau noethérien intègre, les conditions équivalentes de \sref{IV.21.7.6} équivalent à dire que $A$ est un anneau \emph{factoriel}.
\end{remark}

\subsection{Diviseurs et normalisation}
\EGAReaderTarget{ega4.c04.subsec.21-8}{subsection:IV.21.8}

\begin{lemma}[21.8.1]
\label{IV.21.8.1}
Soit $f:X'\to X$ un morphisme entier de préschémas.
Pour tout $\sh{O}_{X'}$-Module localement libre $E'$ de rang constant $n$ et tout $x\in X$, il existe un voisinage ouvert $U$ de $x$ dans $X$ tel que, si l'on pose $U'=f^{-1}(U)$, $E'|U'$ soit isomorphe à $\sh{O}_{U'}^n$.
\end{lemma}

\begin{proof}
La question étant locale sur $X$, on peut se borner au cas où $X=\Spec(A)$ est affine, auquel cas on a $X'=\Spec(A')$, où $A'$ est une $A$-algèbre entière.
Alors $A'$ est limite inductive de ses sous-$A$-algèbres \emph{finies} $A'_\lambda$.
Posant $X'_\lambda=\Spec(A'_\lambda)$ et désignant par $p_\lambda:X'\to X'_\lambda$ le morphisme structural, il résulte de \sref{IV.8.5.2} et \sref{IV.8.5.5} qu'il existe un indice $\lambda$ et un $\sh{O}_{X'_\lambda}$-Module localement libre $E'_\lambda$ de rang $n$ tel que $E'=p_\lambda^*(E'_\lambda)$ et il suffira évidemment de prouver le lemme pour $X'_\lambda$ et $E'_\lambda$;
autrement dit, on peut se borner au cas où le morphisme $f$ est \emph{fini}.
Posons $B=\sh{O}_{X,x}$ et soit $B'$ l'anneau du schéma affine $X'_0=X'\times_X\Spec(\sh{O}_{X,x})$;
comme $B$ est un anneau local et que $B'$ est une $B$-algèbre finie, $B'$ est un anneau semi-local (Bourbaki, \emph{Alg.\ comm.}, chap.~V, \textsection~2, n\textsuperscript{o}~1, prop.~3);
on en conclut que le $\sh{O}_{X'_0}$-Module localement libre $E'_0=E'\otimes_{\sh{O}_{X'}}\sh{O}_{X'_0}$ est isomorphe à $\sh{O}_{X'_0}^n$ (Bourbaki, \emph{Alg.\ comm.}, chap.~II, \textsection~5, n\textsuperscript{o}~3, prop.~5).
Considérant $X'_0$ comme limite
\oldpage[IV-4]{281}
projective des préschémas induits par $X'$ sur les ouverts $f^{-1}(U)$, où $U$ parcourt l'ensemble filtrant des voisinages ouverts affines de $x$ dans $X$, suivant la méthode de \sref{IV.8.1.2}, \emph{a)}), et appliquant \sref{IV.8.5.2.5}, on obtient la conclusion cherchée.
\end{proof}

\begin{corollary}[21.8.2]
\label{IV.21.8.2}
Soit $f:X'\to X$ un morphisme entier de préschémas.
Alors:
\begin{enumerate}[label=(\roman*)]
  \item On a $\RR^1 f_*(\sh{O}^*_{X'})=0$ \sref[0]{0.12.2.1}.
  \item Le groupe $\operatorname{Pic}(X')$ est canoniquement isomorphe à $\HH^1(X,f_*(\sh{O}^*_{X'}))$.
\end{enumerate}
\end{corollary}

\begin{proof}
\begin{enumerate}[label=(\roman*)]
  \item $\RR^1 f_*(\sh{O}^*_{X'})$ est le faisceau associé au préfaisceau de groupes commutatifs $U\mapsto\HH^1(f^{-1}(U),\sh{O}^*_{X'})$ sur $X$ (\emph{loc.\ cit.});
la fibre de ce faisceau en un point $x$ peut donc \sref[0]{0.5.4.7} être identifiée au groupe commutatif $\varinjlim\operatorname{Pic}(f^{-1}(U))$, où $U$ parcourt l'ensemble des voisinages ouverts de $x$, les homomorphismes de transition $\operatorname{Pic}(f^{-1}(U'))\to\operatorname{Pic}(f^{-1}(U))$ pour deux ouverts $U\subset U'$ étant définis par \sref{IV.21.3.2.4}.
Or, pour tout $x\in X$, tout voisinage ouvert $U'$ de $x$ dans $X$, et tout élément $\zeta\in\operatorname{Pic}(f^{-1}(U'))$, il résulte de \sref{IV.21.8.1} qu'il y a un voisinage ouvert $U\subset U'$ de $x$ tel que l'image de $\zeta$ dans $\operatorname{Pic}(f^{-1}(U))$ soit nulle.
Donc la fibre de $\RR^1f_*(\sh{O}^*_{X'})$ au point $x$ est nulle.

  \item La suite spectrale de Leray pour le foncteur composé $\sh{F}'\mapsto\Gamma(X,f_*(\sh{F}'))$ de faisceaux de groupes commutatifs sur $X'$ (T, 2.4) donne la suite exacte de termes de bas degré
\[
  0\to\HH^1(X,f_*(\sh{O}^*_{X'}))\to\HH^1(X',\sh{O}^*_{X'})\to\HH^0(X,\RR^1f_*(\sh{O}^*_{X'}))
\]
et la conclusion résulte de (i) et de l'isomorphisme de $\operatorname{Pic}(X')$ et de $\HH^1(X',\sh{O}^*_{X'})$ \sref[0]{0.5.4.7}.
\end{enumerate}
\end{proof}

\begin{proposition}[21.8.3]
\label{IV.21.8.3}
Soit $f=(\psi,\theta):X'\to X$ un morphisme entier de préschémas;
supposons que, pour tout ouvert $U$ de $X$, l'homomorphisme $\Gamma(\theta):\Gamma(U,\sh{O}_X)\to\Gamma(U,f_*(\sh{O}_{X'}))$ transforme les éléments réguliers en éléments réguliers, ce qui permet de prolonger canoniquement l'homomorphisme $\theta:\sh{O}_X\to f_*(\sh{O}_{X'})$ en un homomorphisme de faisceaux d'anneaux $\theta':\sh{M}_X\to f_*(\sh{M}_{X'})$, d'où des homomorphismes de faisceaux de groupes multiplicatifs $\theta^*:\sh{O}^*_X\to f_*(\sh{O}^*_{X'})$ et $\theta'^*:\sh{M}^*_X\to f_*(\sh{M}^*_{X'})$, donnant par passage aux quotients un homomorphisme $\theta''^*:\sh{D}iv_X\to f_*(\sh{D}iv_{X'})$.
On a alors un diagramme commutatif
\[
\label{IV.21.8.3.1}
  \xymatrix{
    1 \ar[r] & \sh{O}^*_X \ar[r] \ar[d]^{\theta^*} & \sh{M}^*_X \ar[r] \ar[d]^{\theta'^*} & \sh{D}iv_X \ar[r] \ar[d]^{\theta''^*} & 0\\
    1 \ar[r] & f_*(\sh{O}^*_{X'}) \ar[r] & f_*(\sh{M}^*_{X'}) \ar[r] & f_*(\sh{D}iv_{X'}) \ar[r] & 0
  }
  \tag{21.8.3.1}
\]
dans lequel les deux lignes sont exactes.
\end{proposition}

\begin{proof}
Le seul point à vérifier est l'exactitude de la seconde ligne du diagramme, qui résulte de l'application de la suite exacte de cohomologie pour le foncteur $f_*$ à la suite exacte de faisceaux de groupes commutatifs sur $X'$
\[
  1\to\sh{O}^*_{X'}\to\sh{M}^*_{X'}\to\sh{D}iv_{X'}\to0
\]
et de \sref{IV.21.8.2}.
\end{proof}

\begin{corollary}[21.8.4]
\label{IV.21.8.4}
Si, en outre des hypothèses de \sref{IV.21.8.3}, l'homomorphisme $\theta'$ est un isomorphisme de faisceaux d'anneaux, alors $\theta^*:\sh{O}^*_X\to f_*(\sh{O}^*_{X'})$ est injectif, $\theta''^*:\sh{D}iv_X\to f_*(\sh{D}iv_{X'})$ est surjectif et $\operatorname{Ker}(\theta''^*)$ est isomorphe à $\operatorname{Coker}(\theta^*)$.
\end{corollary}

\begin{proof}
C'est une conséquence immédiate du diagramme du serpent (Bourbaki, \emph{Alg.\ comm.}, chap.~I, \textsection~2, n\textsuperscript{o}~4, prop.~2) appliqué au diagramme \sref{IV.21.8.3.1}.
\end{proof}

\oldpage[IV-4]{282}

\begin{proposition}[21.8.5]
\label{IV.21.8.5}
Soient $X$, $X'$ deux préschémas localement noethériens, $f:X'\to X$ un morphisme entier;
on suppose qu'il existe un ouvert schématiquement dense $U$ dans $X$ tel que $f^{-1}(U)$ soit schématiquement dense dans $X'$ (cf.\ \sref{IV.20.3.5}), et que le morphisme $f^{-1}(U)\to U$, restriction de $f$, soit un isomorphisme.
Alors:
\begin{enumerate}[label=(\roman*)]
  \item L'homomorphisme $\theta':\sh{M}_X\to f_*(\sh{M}_{X'})$ est défini et bijectif, l'homomorphisme
\[
  \theta^*:\sh{O}^*_X\to f_*(\sh{O}^*_{X'})
\]
est injectif, l'homomorphisme $\theta''^*:\sh{D}iv_X\to f_*(\sh{D}iv_{X'})$ est surjectif, $\operatorname{Ker}(\theta''^*)$ est isomorphe à $\sh{N}=\operatorname{Coker}(\theta^*)$ et le support du faisceau de groupes multiplicatifs $\sh{N}$ est contenu dans $S=X-U$.
  \item Supposons de plus l'ensemble $S$ discret et posons pour abréger $\sh{O}'_X=f_*(\sh{O}_{X'})$.
Alors on a un diagramme commutatif
\[
\label{IV.21.8.5.1}
  \xymatrix{
    1 \ar[r] & \prod_{s\in S}\sh{O}'^{*}_{X,s}/\sh{O}^*_{X,s} \ar[r]^-{j} \ar[d] & \operatorname{Div}(X) \ar[r] \ar[d]^{k_X} & \operatorname{Div}(X') \ar[r] \ar[d]^{k_{X'}} & 0\\
    1 \ar[r] & \bigl(\prod_{s\in S}\sh{O}'^{*}_{X,s}/\sh{O}^*_{X,s}\bigr)/\operatorname{Im}\Gamma(X',\sh{O}^*_{X'}) \ar[r]^-{j'} & \operatorname{Pic}(X) \ar[r] & \operatorname{Pic}(X') \ar[r] & 0
  }
  \tag{21.8.5.1}
\]
où les lignes sont exactes et où la flèche verticale de gauche est l'homomorphisme canonique.
\end{enumerate}
\end{proposition}

\begin{proof}
\begin{enumerate}[label=(\roman*)]
  \item L'hypothèse entraîne que l'on a un isomorphisme canonique $\sh{M}''_X\simto f_*(\sh{M}''_{X'})$ pour les faisceaux de germes de pseudo-fonctions \sref{IV.20.2.2}, d'où l'assertion relative à $\theta'$ vu l'existence des isomorphismes canoniques $\sh{M}_X\simto\sh{M}''_X$ et $\sh{M}_{X'}\simto\sh{M}''_{X'}$ \sref{IV.20.2.11}.
Le reste de l'assertion (i) découle de \sref{IV.21.8.4}, sauf ce qui concerne le support de $\sh{N}$, qui résulte directement de l'hypothèse sur $U$.

  \item Appliquant la suite exacte de cohomologie aux deux suites exactes de faisceaux de groupes commutatifs
\[
\label{IV.21.8.5.2}
\begin{gathered}
  1\to\sh{N}\to\sh{D}iv_X\to f_*(\sh{D}iv_{X'})\to0,\\
  1\to\sh{O}^*_X\to f_*(\sh{O}^*_{X'})\to\sh{N}\to1,
\end{gathered}
  \tag{21.8.5.2}
\]
il vient respectivement les suites exactes
\[
\label{IV.21.8.5.3}
\begin{gathered}
  1\to\Gamma(X,\sh{N})\xrightarrow{j}\operatorname{Div}(X)\to\operatorname{Div}(X')\to\HH^1(X,\sh{N}),\\
  \Gamma(X',\sh{O}^*_{X'})\to\Gamma(X,\sh{N})\xrightarrow{\partial}\operatorname{Pic}(X)\to\operatorname{Pic}(X')\to\HH^1(X,\sh{N}),
\end{gathered}
  \tag{21.8.5.3}
\]
\oldpage[IV-4]{283}
compte tenu de l'isomorphisme canonique $\operatorname{Pic}(X')\simto\HH^1(X,f_*(\sh{O}^*_{X'}))$ \sref{IV.21.8.2}.
Or, comme le support de $\sh{N}$ est contenu dans l'ensemble discret $S$, fermé dans $X$, on a
\[
  \HH^1(X,\sh{N})=\HH^1(S,\sh{N}|S)\quad\text{(G, II, 4.9.2)},\qquad
  \HH^1(S,\sh{N}|S)=\prod_{s\in S}\HH^1(\{s\},\sh{N}_s)=0
\]
par définition de la cohomologie (G, II, 4.4).
De même on a $\Gamma(X,\sh{N})=\prod_{s\in S}\sh{N}_s$ et $\sh{N}_s=\sh{O}'^{*}_{X,s}/\sh{O}^*_{X,s}$ en vertu de la seconde suite exacte \sref{IV.21.8.5.2}.
Reste à préciser les injections $j$ et $j'$.
Or, on peut décrire une section $t$ de $\sh{N}$ au-dessus de $X$ en prenant un recouvrement de $X$ formé de $U$ et d'ouverts $U_i$ tels que $U_i\cap S$ soit réduit à un seul point $s_i$, et que $U_i\cap U_j\subset U$ pour $i\neq j$, puis en considérant pour chaque $i$ une section $t_i$ de $\sh{N}$ au-dessus de $U_i$, nécessairement telle que $(t_i)_{s_i}\in\sh{O}'^{*}_{X,s_i}/\sh{O}^*_{X,s_i}$ et $(t_i)_x=0$ aux points $x\in U_i-\{s_i\}$.
Le germe $(t_i)_{s_i}$ provient d'une section $u_i$ de $\sh{O}^*_{X'}$ au-dessus de $f^{-1}(U_i)$, qu'on peut aussi considérer comme une section de $\sh{M}^*_{X'}$ au-dessus de $f^{-1}(U_i)$, donc une section de $\sh{M}^*_X$ au-dessus de $U_i$ (en vertu de l'hypothèse);
il correspond canoniquement à cette section une section $d_i\in\Gamma(U_i,\sh{D}iv_X)$ et comme dans $f^{-1}(U\cap U_i)$ la restriction de $u_i$ s'identifie à une section de $\sh{O}^*_X$ au-dessus de $U\cap U_i$ en vertu de l'hypothèse, les restrictions des $d_i$ à $U\cap U_i$ sont toutes nulles, donc les $d_i$ sont les restrictions d'un même diviseur de $\operatorname{Div}(X)$, qui est l'image de la section $t$ par $j$.
De même, l'image de $t$ par $\partial:\Gamma(X,\sh{N})\to\HH^1(X,\sh{O}^*_X)$ provient du cocycle égal à la restriction de $u_i$ dans $U\cap U_i$ pour chaque $i$, à $(u_i|(U_i\cap U_j))(u_j|(U_i\cap U_j))^{-1}$ dans $U_i\cap U_j$, d'où l'expression de $j'(t)$ et la commutativité du diagramme \sref{IV.21.8.5.1}.
\end{enumerate}
\end{proof}

\begin{remark}[21.8.6]
\label{IV.21.8.6}
\begin{enumerate}[label=(\roman*)]
  \item Les conditions d'application de \sref{IV.21.8.5}, (i) sont en particulier remplies lorsque $X$ est un préschéma localement noethérien réduit n'ayant qu'un nombre fini de composantes irréductibles, $X'$ son normalisé et que $X'$ est aussi localement noethérien (\sref[II]{II.6.3.8});
le $\sh{O}_X$-Module $\sh{D}iv_X$ est alors une \emph{extension} de $f_*(\sh{D}iv_{X'})$ par le conoyau $\sh{N}$ de $\sh{O}^{*}_X\to f_*(\sh{O}^*_{X'})$, que l'on peut considérer comme connu.
Si de plus $X$ est une courbe algébrique (réduite) sur un corps $k$, on est dans les conditions d'application de \sref{IV.21.8.5}, (ii).

  \item Lorsque les conditions d'application de \sref{IV.21.8.5}, (ii) sont remplies et en outre que l'homomorphisme canonique global $\Gamma(X,\sh{O}_X)\to\Gamma(X',\sh{O}_{X'})$ est \emph{bijectif}, on voit que les noyaux des homomorphismes surjectifs $\operatorname{Div}(X)\to\operatorname{Div}(X')$ et $\operatorname{Pic}(X)\to\operatorname{Pic}(X')$ sont \emph{isomorphes}.

  \item Lorsque les conditions d'application de \sref{IV.21.8.5}, (ii) sont remplies, on voit que l'homomorphisme $\operatorname{Div}(X)\to\operatorname{Div}(X')$ ne peut être injectif (auquel cas il est bijectif) que si $\sh{O}'^{*}_{X,s}=\sh{O}^*_{X,s}$ pour tout $s\in S$.
Pour un $s\in S$ tel que l'anneau $\sh{O}'_{X,s}$ soit \emph{local} (ce qui, compte tenu du fait que $X'=\Spec(\sh{O}'_X)$ \sref[II]{II.1.3}, signifie qu'il n'existe qu'un seul point $s'\in X'$ au-dessus de $s$) cela entraîne nécessairement que les corps résiduels $\boldsymbol{k}(s)$ et $\boldsymbol{k}(s')$ sont égaux et que $1+\mathfrak{m}_s=1+\mathfrak{m}_{s'}$, donc finalement équivaut à la relation $\sh{O}'_{X,s}=\sh{O}_{X,s}$ ou encore (compte tenu de l'hypothèse) au fait qu'il y a un voisinage ouvert $V$ de $s$ dans $X$ tel que le morphisme $f^{-1}(V)\to V$ soit un isomorphisme.
Dans le cas général, l'anneau $\sh{O}'_{X,s}$ est semi-local (c'est évident lorsque le morphisme $f$ est fini et dans le cas général cela résulte de \sref[0]{0.23.2.6}), l'homomorphisme canonique $\sh{O}^{*}_{X,s}\to\sh{O}'^{*}_{X,s}$ définit, par passage aux quotients, un homomorphisme du groupe multiplicatif $\boldsymbol{k}(s)^*$ dans le produit des groupes multiplicatifs $\boldsymbol{k}(s'_j)^*$, les $s'_j$ étant les points (en
\oldpage[IV-4]{284}
nombre $>1$) de $X'$ au-dessus de $s$.
Il est immédiat que cet homomorphisme ne peut être bijectif que si $\boldsymbol{k}(s)$ et tous les $\boldsymbol{k}(s'_j)$ sont égaux au corps $\bb{F}_2$ à $2$ éléments;
de plus, si cette condition est vérifiée, il faut en outre que le groupe multiplicatif $1+\mathfrak{m}_s$ ait pour image $1+\mathfrak{r}$, où $\mathfrak{r}$ est le radical de $\sh{O}'_{X,s}$, ou encore que $\mathfrak{m}_s=\mathfrak{r}$, ce qui entraîne que $\sh{O}'_{X,s}\otimes_{\sh{O}_{X,s}}\boldsymbol{k}(s)$ est composé direct de corps isomorphes à $\boldsymbol{k}(s)$ (ce qui, lorsque le morphisme $f$ est fini, entraîne qu'il est non ramifié aux points $s'_j$ \sref{IV.17.4.1}).
Si par exemple \emph{aucun} corps résiduel de $X$ n'est isomorphe à $\bb{F}_2$, l'homomorphisme canonique $\operatorname{Div}(X)\to\operatorname{Div}(X')$ ne peut être bijectif que si $f$ est un \emph{isomorphisme}.
Dans le cas où l'homomorphisme canonique $\Gamma(X,\sh{O}_X)\to\Gamma(X',\sh{O}_{X'})$ est bijectif, on conclut de ce qui précède et de (ii) que dans les considérations précédentes, on peut remplacer l'homomorphisme $\operatorname{Div}(X)\to\operatorname{Div}(X')$ par l'homomorphisme $\operatorname{Pic}(X)\to\operatorname{Pic}(X')$.
\end{enumerate}
\end{remark}

\subsection{Diviseurs sur les préschémas de dimension 1}
\EGAReaderTarget{ega4.c04.subsec.21-9}{subsection:IV.21.9}
\label{IV.21.9}

\begin{env}[21.9.1]
\label{IV.21.9.1}
Soient $X$ un espace topologique, $x$ un point de $X$, $i_x:\{x\}\to X$ l'injection canonique.
Si $A(x)$ est un groupe commutatif, on peut le considérer comme un faisceau de groupes commutatifs sur l'espace $\{x\}$ réduit à un seul point, et prendre son image directe $(i_x)_*(A(x))$ qui est un faisceau de groupes commutatifs sur $X$ \sref[0]{0.3.4.1};
il résulte aussitôt des définitions que pour tout ouvert $U$ de $X$, $\Gamma(U,(i_x)_*(A(x)))=A(x)$ si $x\in U$, et est réduit à $0$ dans le cas contraire;
donc, pour $y\in\overline{\{x\}}$ on a $((i_x)_*(A(x)))_y=A(x)$, et pour $y\notin\overline{\{x\}}$, $((i_x)_*(A(x)))_y=0$.

Si maintenant $\sh{F}$ est un faisceau de groupes commutatifs sur $X$ et $Y$ une partie de $X$, pour tout ouvert $U$ de $X$, on a un homomorphisme canonique
\[
\label{IV.21.9.1.1}
  \Gamma(U,\sh{F})\to\prod_{x\in U\cap Y}\sh{F}_x
  =\prod_{x\in Y}\Gamma(U,(i_x)_*(\sh{F}_x))
  \tag{21.9.1.1}
\]
et comme ces homomorphismes commutent aux opérateurs de restriction, ils définissent un homomorphisme canonique de \emph{faisceaux}
\[
\label{IV.21.9.1.2}
  j_Y:\sh{F}\to\prod_{x\in Y}(i_x)_*(\sh{F}_x).
  \tag{21.9.1.2}
\]
\end{env}

\begin{lemma}[21.9.2]
\label{IV.21.9.2}
Soient $X$ un espace topologique localement noethérien, $X_0$ l'ensemble de ses points fermés, $\sh{F}$ un faisceau de groupes commutatifs sur $X$.
Les conditions suivantes sont équivalentes:
\begin{enumerate}[label=\emph{\alph*)}]
  \item L'homomorphisme canonique $j_{X_0}:\sh{F}\to\prod_{x\in X_0}(i_x)_*(\sh{F}_x)$ est injectif et a pour image $\bigoplus_{x\in X_0}(i_x)_*(\sh{F}_x)$.
  \item[\rm a$'$)] Il existe une famille de groupes commutatifs $(A(x))_{x\in X_0}$ telle que $\sh{F}$ soit isomorphe à $\bigoplus_{x\in X_0}(i_x)_*(A(x))$.
  \item Toute section de $\sh{F}$ au-dessus d'un ouvert $U$ de $X$ a un support discret contenu dans $X_0$.
\end{enumerate}
Lorsqu'il en est ainsi, pour tout $x\in X_0$, le groupe $A(x)$ est déterminé à isomorphisme unique près et est isomorphe à $\sh{F}_x$.
En outre, le faisceau $\sh{F}$ est alors flasque.
\end{lemma}

\begin{proof}
\oldpage[IV-4]{285}
Comme les points de $X_0$ sont fermés dans $X$, on a $((i_x)_*(A(x)))_y=0$ pour tout $y\neq x\in X_0$;
cette remarque prouve l'assertion d'unicité relative aux groupes $A(x)$, et il est clair par ailleurs que \emph{a)} implique \emph{a$'$)}.
Pour voir que \emph{a$'$)} implique \emph{b)}, on peut se borner au cas où $U$ est noethérien;
alors (G, II, 3.10) on a
\[
  \Gamma(U,\bigoplus_{x\in X_0}(i_x)_*(A(x)))
  =\bigoplus_{x\in X_0}\Gamma(U,(i_x)_*(A(x))),
\]
donc toute section $s$ de $\sh{F}$ au-dessus de $U$ est somme directe d'un nombre \emph{fini} de sections $s_k\in\Gamma(U,(i_{x_k})_*(A(x_k)))$ ($1\leq k\leq n$) avec $x_k\in X_0$.
Mais puisque $x_k$ est fermé dans $X$, le support de $s_k$ est réduit à $\{x_k\}$, donc le support de $s$ est contenu dans l'ensemble fini des $x_k$, qui est évidemment discret puisque les points $x_k$ sont fermés dans $X$.
Montrons enfin que \emph{b)} entraîne \emph{a)};
pour tout ouvert noethérien $U$, tout support d'une section de $\sh{F}$ au-dessus de $U$ étant discret et quasi-compact, est \emph{fini}.
On déduit donc de l'hypothèse \emph{b)} que pour tout ouvert noethérien $U$, l'image de l'homomorphisme \sref{IV.21.9.1.1} (pour $Y=X_0$) est contenue dans $\bigoplus_{x\in X_0}\Gamma(U,(i_x)_*(\sh{F}_x))$ et que cet homomorphisme est injectif, ce qui établit \emph{a)}.

Enfin, pour montrer que $\sh{F}$ est flasque, considérons une section $s$ de $\sh{F}$ au-dessus d'un ouvert $U$ de $X$;
comme le support de $s$ est un sous-espace discret et fermé dans $X$, on prolonge $s$ en une section $s'$ de $\sh{F}$ au-dessus de $X$ en posant $s'_x=0$ dans $X-U$.
\end{proof}

\begin{remark}[21.9.3]
\label{IV.21.9.3}
\begin{enumerate}[label=(\roman*)]
  \item Dans la condition \emph{b)} de \sref{IV.21.9.2}, il ne suffit pas de supposer que le support de toute section de $\sh{F}$ au-dessus d'un ouvert quelconque de $X$ soit \emph{discret};
c'est ce que montre l'exemple où on prend pour $X$ un spectre d'anneau de valuation discrète, avec un point fermé $b$ et un point générique $a$, et pour $\sh{F}$ le faisceau de groupes commutatifs tel que $\sh{F}_b=0$ et $\sh{F}_a=\bb{Z}$.

  \item Les conditions de \sref{IV.21.9.2} étant supposées vérifiées, soit $E$ une partie discrète de $X_0$, et pour tout $x\in E$, soit $a(x)$ un élément de $A(x)$.
Il existe alors une section $t$ de $\bigoplus_{x\in X_0}(i_x)_*(A(x))$ et une seule au-dessus de $X$ telle que $t_x=a(x)$ pour tout $x\in E$ et que le support de $t$ soit contenu dans $E$.
En effet, pour tout ouvert noethérien $U$, $E\cap U$ est fini, et il suffit de prendre pour $t$ la section de $\bigoplus_{x\in E}(i_x)_*(A(x))$ dont la restriction à chaque ouvert noethérien $U$ est la somme des $a(x)$ pour $x\in E\cap U$.
\end{enumerate}
\end{remark}

\begin{proposition}[21.9.4]
\label{IV.21.9.4}
Soient $X$ un préschéma localement noethérien de dimension $\leq1$, $X^{(1)}$ l'ensemble des points $x\in X$ tels que $\dim(\sh{O}_{X,x})=1$.
On a alors un isomorphisme canonique
\[
\label{IV.21.9.4.1}
  \sh{D}iv_X\simto\bigoplus_{x\in X^{(1)}}(i_x)_*(\operatorname{Div}(\sh{O}_{X,x}))
  \tag{21.9.4.1}
\]
et $\sh{D}iv_X$ est flasque.
\end{proposition}

\begin{proof}
Compte tenu de l'isomorphisme \sref{IV.21.4.6.1}, l'homomorphisme \sref{IV.21.9.4.1} est défini par \sref{IV.21.9.1.2}:
prouvons que c'est une bijection.
Comme $\dim(X)\leq1$, les points de $X^{(1)}$ sont les \emph{points fermés non isolés} de $X$.
On est ramené à prouver que:
1\textsuperscript{o} $\sh{D}iv_X$ vérifie la condition \emph{b)} de \sref{IV.21.9.2};
2\textsuperscript{o} pour tout point isolé $x\in X$, on a $(\sh{D}iv_X)_x=0$.
Le second point résulte de ce qu'alors $\sh{O}_{X,x}$ est un anneau artinien et de ce que tout élément régulier de $\sh{O}_{X,x}$ est donc inversible.
\oldpage[IV-4]{286}
Pour le 1\textsuperscript{o}, il suffit de noter que pour tout ouvert $U$ de $X$ et tout diviseur $D\in\operatorname{Div}(U)$, les points maximaux du support $S$ de $D$ sont tels que $\operatorname{prof}(\sh{O}_{X,x})=1$ \sref{IV.21.1.9}, donc \emph{a fortiori} appartiennent à $X^{(1)}\cap U$;
puisque $\dim(X)\leq1$, l'ensemble de ces points est égal à $S$, et $S$ est donc discret, l'ensemble des composantes irréductibles de $S$ étant localement fini.
\end{proof}

\begin{corollary}[21.9.5]
\label{IV.21.9.5}
Soit $X$ un préschéma localement noethérien de dimension $\leq1$.
Pour toute partie discrète $E\subset X^{(1)}$ et toute famille $(D(x))_{x\in E}$ avec $D(x)\in\operatorname{Div}(\sh{O}_{X,x})$, il existe un diviseur $D$ sur $X$ et un seul tel que le support de $D$ soit contenu dans $E$ et que $D_x=D(x)$ pour tout $x\in E$.
\end{corollary}

\begin{proof}
Cela résulte de \sref{IV.21.9.4} et de \sref{IV.21.9.3}, (ii).
\end{proof}

\begin{corollary}[21.9.6]
\label{IV.21.9.6}
Soit $X$ un préschéma localement noethérien de dimension $\leq1$;
tout diviseur $D$ sur $X$ est de la forme $D'-D''$, où $D'$, $D''$ sont deux diviseurs $\geq0$ de supports contenus dans celui de $D$.
\end{corollary}

\begin{proof}
En vertu de \sref{IV.21.9.5}, on est aussitôt ramené au cas où $X=\Spec(A)$, où $A$ est un anneau local noethérien;
il suffit alors d'observer que $M(X)$ est l'anneau total des fractions de $A$, et qu'une section de $\sh{M}^*_X$ au-dessus de $X$ s'écrit par définition comme quotient $b/a$ de deux éléments réguliers de $A$.
\end{proof}

\begin{corollary}[21.9.7]
\label{IV.21.9.7}
Soient $X$ un préschéma localement noethérien de dimension $\leq1$, sans point isolé, et $U$ un ouvert dense dans $X$.
Il existe alors un diviseur $D\geq0$ sur $X$, de support contenu dans $U$ et rencontrant chacune des composantes irréductibles de $X$.
\end{corollary}

\begin{proof}
On peut supposer que $U$ est réunion d'ensembles ouverts disjoints non vides $U_\alpha$, dont chacun est contenu dans une seule composante irréductible de $X$;
il suffit alors de prendre dans chaque $U_\alpha$ un point $x_\alpha$ fermé dans $X$ (il en existe puisque l'unique point générique de $U_\alpha$ ne peut être isolé par hypothèse), et d'appliquer \sref{IV.21.9.5} à l'ensemble discret des $x_\alpha$.
\end{proof}

L'intérêt de ce corollaire est qu'il permettra de prouver qu'une courbe algébrique séparée $X$ sur un corps $k$ est quasi-projective, le diviseur $D\geq0$ défini dans \sref{IV.21.9.7} étant alors nécessairement \emph{ample} en vertu du théorème de Riemann--Roch pour les courbes (chap.~V).

\begin{env}[21.9.8]
\label{IV.21.9.8}
Pour les préschémas localement noethériens $X$ de dimension $\leq1$ la proposition \sref{IV.21.9.4} ramène la détermination de $\sh{D}iv_X$ au cas où $X=\Spec(A)$, $A$ étant un anneau local noethérien de dimension $1$.

Lorsque $A$ est un anneau local régulier de dimension $1$ (autrement dit, un anneau de valuation discrète), le groupe $\operatorname{Div}(A)$ s'identifie canoniquement à $\bb{Z}$ \sref{IV.21.6.8};
par suite, en vertu de \sref{IV.21.9.2}:
\end{env}

\begin{proposition}[21.9.9]
\label{IV.21.9.9}
Soit $X$ un préschéma localement noethérien régulier de dimension $\leq1$.
Alors le faisceau $\sh{D}iv_X$ est canoniquement isomorphe à $\bigoplus_{x\in X^{(1)}}(i_x)_*(\bb{Z}(x))$, où $\bb{Z}(x)$ est le groupe additif $\bb{Z}$ considéré comme faisceau de groupes sur l'espace $\{x\}$.
\end{proposition}

\begin{env}[21.9.10]
\label{IV.21.9.10}
Supposons seulement que l'anneau local noethérien $A$ de dimension $1$ soit \emph{réduit};
alors, si $A'$ est le \emph{normalisé} de $A$ (fermeture intégrale de $A$ dans son anneau total de fractions), $A'$ est noethérien en vertu du théorème de Krull--Akizuki (Bourbaki, \emph{Alg.\ comm.}, chap.~VII, \textsection~2, n\textsuperscript{o}~5, prop.~5), et on a vu dans \sref{IV.21.8.6} comment
\oldpage[IV-4]{287}
$\operatorname{Div}(A)$ peut s'obtenir comme \emph{extension} de $\operatorname{Div}(A')$ et ce dernier groupe est de la forme $\bb{Z}^r$.
\end{env}

\begin{proposition}[21.9.11]
\label{IV.21.9.11}
Soient $X$ un préschéma noethérien, $X_0$ un sous-préschéma fermé de $X$ ayant les propriétés suivantes:
\begin{enumerate}
  \item[\rm 1\textsuperscript{o}] $\dim(X_0)\leq1$.
  \item[\rm 2\textsuperscript{o}] Pour toute partie localement fermée $Y$ de $X$ telle que $Y\cap X_0$ soit discret, il existe une partie $Y'$ de $Y$, fermée dans $X$ et ouverte dans $Y$, contenant $Y\cap X_0$.
\end{enumerate}
Dans ces conditions:
\begin{enumerate}[label=(\roman*)]
  \item Soit $Z_0$ la réunion des ensembles $\overline{\{x\}}\cap X_0$ lorsque $x$ parcourt la partie de $\operatorname{Ass}(\sh{O}_X)$ formée des points tels que $\overline{\{x\}}\cap X_0$ soit fini.
Alors, pour tout diviseur $D_0$ sur $X_0$, dont le support ne rencontre pas $Z_0$, il existe un diviseur $D$ sur $X$ dont l'image réciproque par l'injection canonique $X_0\to X$ existe \sref{IV.21.4} et est égale à $D_0$;
si de plus $D_0\geq0$, on peut supposer $D\geq0$.

  \item Supposons en outre qu'il existe dans $X_0$ un ouvert affine $U_0$ contenant $\operatorname{Ass}(\sh{O}_{X_0})\cup Z_0$ (condition automatiquement vérifiée lorsqu'il existe un $\sh{O}_{X_0}$-Module ample \sref[II]{II.4.5.4}).
Alors l'homomorphisme canonique \sref{IV.21.3.2.4} $\operatorname{Pic}(X)\to\operatorname{Pic}(X_0)$ est surjectif.
\end{enumerate}
\end{proposition}

\begin{proof}
\begin{enumerate}[label=(\roman*)]
  \item Compte tenu de \sref{IV.21.9.6}, on peut se borner à prouver l'assertion relative aux diviseurs $D_0\geq0$;
en vertu de \sref{IV.21.9.4}, le support $T$ de $D_0$ est un ensemble discret et fermé dans $X_0$.
On peut supposer $D_0\neq0$, c'est-à-dire $T\neq\emp$, sans quoi il n'y a rien à démontrer.
Pour tout $x\in T$, il existe un élément $s_x\in\sh{O}_{X_0,x}$ qui est non diviseur de zéro dans cet anneau, appartient à son idéal maximal, et dont l'image dans $(\sh{D}iv_{X_0})_x$ est $(D_0)_x$.
Il existe un voisinage ouvert affine $U^{(x)}$ de $x$ dans $X$ et une section $g^{(x)}$ de $\sh{O}_X$ au-dessus de $U^{(x)}$ tels que $s_x$ soit l'image dans $\sh{O}_{X_0,x}$ du germe $(g^{(x)})_x\in\sh{O}_{X,x}$;
nous allons voir qu'en prenant $U^{(x)}$ assez petit, on peut faire en sorte que $g^{(x)}$ ne soit pas diviseur de zéro dans $\Gamma(U^{(x)},\sh{O}_X)$, autrement dit \sref{IV.3.1.9}, en désignant par $V^{(x)}$ la partie fermée de $U^{(x)}$ formée des $y$ tels que $g^{(x)}(y)=0$, que l'on ait $V^{(x)}\cap\operatorname{Ass}(\sh{O}_X)=\emp$.

Or, si $z\in\operatorname{Ass}(\sh{O}_X)$, l'hypothèse $x\notin Z_0$ entraîne que l'on a, ou bien $x\notin\overline{\{z\}}$, ou bien que $x$ n'est pas isolé dans $\overline{\{z\}}\cap X_0$.
En remplaçant $U^{(x)}$ par un ouvert plus petit, on peut supposer que c'est le second cas qui se produit pour un $z\in V^{(x)}\cap\operatorname{Ass}(\sh{O}_X)$.
$V^{(x)}$ contiendrait donc une composante irréductible de dimension $1$ de $X_0\cap U^{(x)}$, contenant $x$;
mais cela signifierait que l'image $\overline{g}^{(x)}$ de $g^{(x)}$ dans $\Gamma(U^{(x)}\cap X_0,\sh{O}_{X_0})$ serait dans le nilradical de cet anneau, et par suite $s_x$ appartiendrait au nilradical de $\sh{O}_{X_0,x}$, ce qui est absurde.
On a donc $x\notin\overline{\{z\}}$ pour tout point $z\in V^{(x)}\cap\operatorname{Ass}(\sh{O}_X)$, et comme cet ensemble est fini, on peut, en remplaçant $U^{(x)}$ par un voisinage ouvert plus petit de $x$, supposer que $V^{(x)}\cap\operatorname{Ass}(\sh{O}_X)=\emp$.

En vertu du choix de $U^{(x)}$, on peut définir un diviseur $D'^{(x)}$ sur $U^{(x)}$ par $D'^{(x)}=\operatorname{div}(g^{(x)})$;
d'ailleurs on a vu ci-dessus que $x$ est nécessairement isolé dans $V^{(x)}\cap X_0$, donc en remplaçant encore $U^{(x)}$ par un voisinage ouvert plus petit de $x$, on peut supposer que $V^{(x)}\cap X_0$ est réduit au point $x$.
Mais il existe alors, en vertu de la condition 2\textsuperscript{o}, une partie $W^{(x)}$ de $V^{(x)}$, fermée dans $X$ et ouverte dans $V^{(x)}$, telle que $W^{(x)}\cap X_0=\{x\}$.
Remplaçant encore $U^{(x)}$ par un voisinage ouvert plus petit de $x$, on peut donc supposer
\oldpage[IV-4]{288}
que $W^{(x)}=V^{(x)}$, autrement dit que $V^{(x)}$ est fermé dans $X$.
On peut alors définir un diviseur $D^{(x)}$ sur $X$ par la condition que $D^{(x)}|U^{(x)}=D'^{(x)}$, et $D^{(x)}|(X-V^{(x)})=0$, ce qui a un sens puisque dans $U^{(x)}\cap(X-V^{(x)})$ la restriction de $g^{(x)}$ est une section inversible, donc les restrictions à cet ouvert des deux diviseurs considérés sur $U^{(x)}$ et $X-V^{(x)}$ sont égales.
Il suffit alors, pour répondre à la question, de prendre $D=\sum_{x\in T}D^{(x)}$, ce qui a un sens puisque $T$ est fini.

  \item Compte tenu du diagramme commutatif \sref{IV.21.4.2.1}, la conclusion résultera de (i) si l'on prouve que tout $\sh{O}_{X_0}$-Module inversible $\sh{L}_0$ est isomorphe à un $\sh{O}_{X_0}$-Module de la forme $\sh{O}_{X_0}(D_0)$ \sref{IV.21.2.8}, où $D_0$ est un diviseur sur $X_0$ dont le support ne rencontre pas $Z_0$.
Or, comme $U_0$ est schématiquement dense dans $X_0$ \sref{IV.20.2.13}, (iv), il suffit pour cela de prouver qu'il existe une section $t\in\Gamma(U_0,\sh{L}_0)$ telle que $t(z_0)\neq0$ aux points de $\operatorname{Ass}(\sh{O}_{X_0})\cup Z_0$ \sref{IV.3.1.9}.
On peut donc se borner au cas où $X_0=\Spec(A_0)$ est affine;
mais alors $\sh{L}_0$ est ample \sref[II]{II.5.1.4} et comme l'ensemble $\operatorname{Ass}(\sh{O}_{X_0})\cup Z_0$ est fini, la conclusion résulte de \sref[II]{II.4.5.4}.
\end{enumerate}
\end{proof}

\begin{corollary}[21.9.12]
\label{IV.21.9.12}
Soient $A$ un anneau local hensélien \sref{IV.18.5.8}, $S=\Spec(A)$, $s_0$ le point fermé de $S$, $f:X\to S$ un morphisme séparé de présentation finie, et supposons que l'ensemble $X_0=f^{-1}(s_0)$ soit de dimension $\leq1$.
Alors, pour tout sous-schéma fermé $X'_0$ de $X$, ayant $X_0$ pour ensemble sous-jacent et de présentation finie sur $S$, l'homomorphisme canonique \sref{IV.21.3.2.4} $\operatorname{Pic}(X)\to\operatorname{Pic}(X'_0)$ est surjectif.
\end{corollary}

\begin{proof}
Montrons d'abord qu'il suffit de prouver le corollaire lorsque l'anneau local hensélien $A$ est de plus noethérien.
On sait en effet \sref{IV.18.6.15} que l'on peut écrire $A=\varinjlim A_\lambda$, où les $A_\lambda$ sont des anneaux locaux noethériens et henséliens, les homomorphismes $A_\lambda\to A$ étant locaux;
il existe en outre un indice $\alpha$ et un morphisme séparé de présentation finie $f_\alpha:X_\alpha\to S_\alpha=\Spec(A_\alpha)$ tels que $X$ et $f$ soient déduits de $X_\alpha$ et $f_\alpha$ par le changement de base $S\to S_\alpha$ \sref{IV.8.10.5}, (v);
avec les notations habituelles \sref{IV.8.8.1}, les morphismes $f_\lambda:X_\lambda\to S_\lambda$ seront séparés pour $\lambda\geq\alpha$ et on aura $X=X_\lambda\times_{S_\lambda}S$.
En outre, on peut supposer que, si $s_{0\alpha}$ est l'unique point fermé de $S_\alpha$, il y a un sous-schéma fermé $X'_{0\alpha}$ de $X_\alpha$, ayant même ensemble sous-jacent que $f_\alpha^{-1}(s_{0\alpha})$, tel que $X'_0=X'_{0\alpha}\times_{S_\alpha}S$ \sref{IV.8.6.3};
pour $\lambda\geq\alpha$, posons $X'_{0\lambda}=X'_{0\alpha}\times_{S_\alpha}S_\lambda$;
on a $\dim(X'_{0\lambda})=\dim(X'_0)\leq1$ par transitivité des fibres et \sref{IV.4.1.4}.
Il s'agit de voir que si l'on a prouvé que l'homomorphisme $\operatorname{Pic}(X_\lambda)\to\operatorname{Pic}(X'_{0\lambda})$ est surjectif pour tout $\lambda\geq\alpha$, il en est de même de $\operatorname{Pic}(X)\to\operatorname{Pic}(X'_0)$.
On a évidemment le diagramme commutatif d'homomorphismes canoniques
\[
  \xymatrix{
    \operatorname{Pic}(X_\lambda) \ar[r] \ar[d] & \operatorname{Pic}(X) \ar[d]\\
    \operatorname{Pic}(X'_{0\lambda}) \ar[r] & \operatorname{Pic}(X'_0)
  }
\]
Pour tout $\sh{O}_{X'_0}$-Module inversible $\sh{L}'_0$, il existe un $\lambda\geq\alpha$ et un $\sh{O}_{X'_{0\lambda}}$-Module inversible $\sh{L}'_{0\lambda}$ tel que $\sh{L}'_0$ s'en déduise par changement de base $S\to S_\lambda$ \sref{IV.8.5.2} et \sref{IV.8.5.5};
il suffit de considérer un $\sh{O}_{X_\lambda}$-Module inversible $\sh{L}_\lambda$ tel que $\operatorname{cl}(\sh{L}'_{0\lambda})$ soit l'image de $\operatorname{cl}(\sh{L}_\lambda)$ dans $\operatorname{Pic}(X'_{0\lambda})$, puis de prendre le $\sh{O}_X$-Module inversible $\sh{L}$ déduit de $\sh{L}_\lambda$ par changement de base;
en vertu de la commutativité du diagramme précédent, $\operatorname{cl}(\sh{L}'_0)$ sera l'image de $\operatorname{cl}(\sh{L})$.

Supposons donc $A$ noethérien, donc aussi $X$, et vérifions que $X$ et $X'_0$ satisfont aux conditions de \sref{IV.21.9.11}, (ii).
On a par hypothèse $\dim(X'_0)\leq1$;
d'autre part, pour vérifier la condition 2\textsuperscript{o} de \sref{IV.21.9.11}, considérons un sous-préschéma réduit $Y$ de $X$ ayant une partie localement fermée $Y$ pour ensemble sous-jacent;
le morphisme $g:Y\to S$, restriction de $f$, étant quasi-fini en chacun des points $x_i$ de $Y\cap X_0$ ($1\leq i\leq n$), on peut appliquer \sref{IV.18.5.11}, c), et on voit que $Y$ est somme des sous-préschémas ouverts $Y_i=\Spec(\sh{O}_{Y,x_i})$ ($1\leq i\leq n$) et d'un préschéma $Y''$ tel que $Y''\cap X_0=\emp$, et en outre que les injections canoniques $j_i:Y_i\to X$ sont telles que $f\circ j_i$ soit un morphisme fini.
Comme $f$ est séparé, $j_i$ est aussi un morphisme fini, donc $Y_i$ est fermé dans $X$;
on répond donc à la question en prenant $Y'=\bigcup_{1\leq i\leq n}Y_i$.

Il reste à vérifier l'hypothèse supplémentaire de \sref{IV.21.9.11}, (ii).
Or, $X'_0$ étant une courbe séparée sur le corps $\boldsymbol{k}(s_0)$, est un $\boldsymbol{k}(s_0)$-schéma quasi-projectif (chap.~V)\footnote{Le lecteur vérifiera que \sref{IV.21.9.12} n'est pas utilisé pour prouver cette propriété au chap.~V.}, donc il existe un $\sh{O}_{X'_0}$-Module ample \sref[II]{II.5.3.1}, et cela achève la démonstration.
\end{proof}

\begin{remark}[21.9.13]
\label{IV.21.9.13}
Sous les conditions de \sref{IV.21.9.12}, supposant en outre $f$ propre, le morphisme $f$ est même \emph{projectif}:
en effet, si $\sh{L}_0$ est un $\sh{O}_{X'_0}$-Module ample, il existe, en vertu de \sref{IV.21.9.12}, un $\sh{O}_X$-Module inversible $\sh{L}$ dont l'image réciproque dans $X'_0$ est isomorphe à $\sh{L}_0$, donc est ample.
Puisque tout voisinage de $s_0$ dans $S$ est nécessairement $S$ tout entier, on déduit alors de \sref{IV.9.6.4} que $\sh{L}$ est un $\sh{O}_X$-Module ample, d'où la conclusion (\sref[II]{II.5.3.1} et \sref[II]{II.5.5.3}).
\end{remark}

\subsection{Images réciproques et images directes de cycles 1-codimensionnels}
\EGAReaderTarget{ega4.c04.subsec.21-10}{subsection:IV.21.10}

Dans un chapitre ultérieur, consacré à la théorie des intersections, seront développées de façon systématique les notions d'image inverse et d'image directe de cycles.
Dans le présent numéro, nous nous contenterons de définir ces notions dans certains cas particuliers utiles, et pour les cycles $1$-codimensionnels, ces définitions étant choisies de telle sorte qu'elles soient compatibles avec les définitions correspondantes pour les diviseurs \sref{IV.21.4} et \sref{IV.21.5}, compte tenu de l'homomorphisme $\operatorname{cyc}$ défini en \sref{IV.21.6}.

\begin{env}[21.10.1]
\label{IV.21.10.1}
\oldpage[IV-4]{289}
Soient $X$, $X'$ deux préschémas localement noethériens, $f:X'\to X$ un morphisme, $T$ une partie de $X$;
supposons que l'image par $f$ de tout point maximal de $X'$ soit un point maximal de $X$, et que, pour tout $x'\in X'^{(1)}$ (i.e.\ tel que $\dim(\sh{O}_{X',x'})=1$), on ait l'une des trois conditions suivantes pour le point $x=f(x')$:
\begin{enumerate}[label=(\roman*)]
  \item $x\notin T$;
  \item $x\in X^{(1)}$ et $\sh{O}_{X',x'}$ est un $\sh{O}_{X,x}$-module plat;
  \item l'anneau $\sh{O}_{X,x}$ est factoriel et $\mathfrak{m}_{x'}\notin\operatorname{Ass}(\sh{O}_{X',x'})$.
\end{enumerate}
Sous ces conditions, nous nous proposons, pour tout cycle $1$-codimensionnel $Z$ sur $X$ de support contenu dans $T$, de définir un cycle $1$-codimensionnel $Z'=f^*(Z)$, de sorte que $Z\mapsto f^*(Z)$ soit un homomorphisme de groupes ordonnés du sous-groupe de $\mathfrak{Z}^1(X)$ formé des cycles de support contenu dans $T$, dans le groupe ordonné $\mathfrak{Z}^1(X')$.
Pour cela, soit
\[
\label{IV.21.10.1.1}
  Z=\sum_{x\in T\cap X^{(1)}}n_x\cdot\overline{\{x\}}
  \tag{21.10.1.1}
\]
où la famille des $x\in T\cap X^{(1)}$ tels que $n_x\neq0$ est localement finie.
Pour tout $x'\in X'^{(1)}$, définissons un entier $n_{x'}$ de la façon suivante, en posant $x=f(x')$:
\begin{enumerate}[label=\arabic*\textsuperscript{o}]
  \item si $x\notin T$, on prend $n_{x'}=0$;
  \item si $x\in X^{(1)}$ et si $\sh{O}_{X',x'}$ est un $\sh{O}_{X,x}$-module plat, on sait \sref{IV.6.1.1} que $\dim(\sh{O}_{X',x'}/\mathfrak{m}_x\sh{O}_{X',x'})=0$, autrement dit $\sh{O}_{X',x'}/\mathfrak{m}_x\sh{O}_{X',x'}$ est un $\sh{O}_{X',x'}$-module de longueur finie $\lambda_{x'}$, et on prend $n_{x'}=\lambda_{x'}n_x$;
  \item si $\sh{O}_{X,x}$ est factoriel et $\mathfrak{m}_{x'}\notin\operatorname{Ass}(\sh{O}_{X',x'})$, on sait \sref{IV.21.6.9} que l'homomorphisme canonique $\operatorname{cyc}:\operatorname{Div}(\sh{O}_{X,x})\to\mathfrak{Z}^1(\Spec(\sh{O}_{X,x}))$ est bijectif, et d'autre part comme $\dim(\sh{O}_{X',x'})=1$ et $\mathfrak{m}_{x'}\notin\operatorname{Ass}(\sh{O}_{X',x'})$, $\operatorname{Ass}(\sh{O}_{X',x'})$ se compose uniquement des points maximaux de $\Spec(\sh{O}_{X',x'})$, donc l'hypothèse sur $f$ implique que $f(\operatorname{Ass}(\sh{O}_{X',x'}))\subset\operatorname{Ass}(\sh{O}_{X,x})$ et il résulte de \sref{IV.21.4.5}, (ii) que l'homomorphisme $f^*:\operatorname{Div}(\sh{O}_{X,x})\to\operatorname{Div}(\sh{O}_{X',x'})$ est défini;
enfin, $x'$ étant l'unique point fermé de $\Spec(\sh{O}_{X',x'})$, $\mathfrak{Z}^1(\Spec(\sh{O}_{X',x'}))$ est canoniquement isomorphe à $\bb{Z}$.
\end{enumerate}
On a donc un homomorphisme canonique composé
\[
\label{IV.21.10.1.2}
  \mathfrak{Z}^1(\Spec(\sh{O}_{X,x}))
  \xrightarrow{\operatorname{cyc}^{-1}}\operatorname{Div}(\sh{O}_{X,x})
  \xrightarrow{f^*}\operatorname{Div}(\sh{O}_{X',x'})
  \xrightarrow{\operatorname{cyc}}\mathfrak{Z}^1(\Spec(\sh{O}_{X',x'}))
  \xrightarrow{\sim}\bb{Z}
  \tag{21.10.1.2}
\]
Si $Z_x$ est le cycle $\sum_{y\in\Spec(\sh{O}_{X,x})\cap X^{(1)}}n_y(\overline{\{y\}}\cap\Spec(\sh{O}_{X,x}))$ sur $\Spec(\sh{O}_{X,x})$, on prend alors pour $n_{x'}$ l'\emph{image} de $Z_x$ par l'homomorphisme \sref{IV.21.10.1.2}.
\end{env}

\oldpage[IV-4]{290}
\begin{env}[21.10.2]
\label{IV.21.10.2}
Nous nous proposons de montrer que:
\begin{enumerate}[label=\emph{\Alph*)}]
  \item Lorsque deux des conditions 1\textsuperscript{o}, 2\textsuperscript{o}, 3\textsuperscript{o} de \sref{IV.21.10.1} sont \emph{simultanément} satisfaites, les valeurs correspondantes de $n_{x'}$ \emph{coïncident}.
  \item L'ensemble des $x'\in X'^{(1)}$ tels que $n_{x'}\neq0$ est \emph{localement fini} dans $X'$.
\end{enumerate}

Pour démontrer A), supposons d'abord que $x\notin T$ et que $x$ vérifie l'une des conditions 2\textsuperscript{o} ou 3\textsuperscript{o} de \sref{IV.21.10.1};
alors $n_x=0$ et si l'on est dans le cas 2\textsuperscript{o}, on a $n_{x'}=0$;
si on est dans le cas 3\textsuperscript{o}, on a $Z_x=0$ puisque $\Supp(Z)\subset T$, donc encore $n_{x'}=0$.
Il reste à considérer le cas où l'on est à la fois dans le cas 2\textsuperscript{o} et dans le cas 3\textsuperscript{o};
alors, puisque $\dim(\sh{O}_{X,x})=1$, $\sh{O}_{X,x}$ est un anneau de valuation discrète;
si $t$ est une uniformisante de cet anneau, le diviseur correspondant à $Z_x$ est $\operatorname{div}(t^{n_x})$ dans $\operatorname{Div}(\sh{O}_{X,x})$, et son image dans $\operatorname{Div}(\sh{O}_{X',x'})$ est le diviseur de $t'^{n_x}$, où $t'$ est l'élément (régulier) de $\sh{O}_{X',x'}$ image de $t$.
On peut évidemment se borner au cas où $n_x=1$, et alors la définition \sref{IV.21.6.5.1} montre que l'image de $Z_x$ par \sref{IV.21.10.1.2} est le nombre $\lambda_{x'}$, ce qui achève de prouver A).

Démontrons maintenant B).
Posons $T_0=\Supp(Z)\subset T$, $T'_0=f^{-1}(T_0)$;
il suffit de prouver que la relation $n_{x'}\neq0$ implique que $x'$ appartient à l'ensemble des points maximaux de $T'_0$, ce dernier étant localement fini dans $X'$.
Il est immédiat que l'on a nécessairement $x'\in T'_0$;
si $x'$ n'était pas maximal dans l'ensemble fermé $T'_0$, il existerait
\oldpage[IV-4]{291}
une générisation $y'$ de $x'$ dans $T'_0$, distincte de $x'$, et puisque $x'\in X'^{(1)}$, $y'$ serait nécessairement un point maximal de $X'$;
par suite $y=f(y')$ serait point maximal de $X$ par hypothèse;
mais cela est absurde puisque $y\in T_0$ et que $T_0$ est purement de codimension $1$ dans $X$.
\end{env}

\begin{env}[21.10.3]
\EGAReaderTarget{ega4.c04.statement.21-10-3}{IV.21.10.3}
On peut maintenant poser
\[
  f^*(Z)=\sum_{x'\in X'^{(1)}}n_{x'}\overline{\{x'\}},
\]
la somme du second membre ayant un sens en vertu de ce qu'on a prouvé dans \sref{IV.21.10.2};
on dit que le cycle $1$-codimensionnel $f^*(Z)$ est l'\emph{image réciproque} de $Z$ par $f$.
Il est immédiat que l'application $f^*$ ainsi définie est un homomorphisme de groupes ordonnés.
En outre, si $U$ est un ouvert de $X$, $V$ un ouvert de $X'$ tel que $f(V)\subset U$ et $f':V\to U$ la restriction de $f$, il résulte aussitôt des définitions que l'on a
\[
\label{IV.21.10.3.1}
  f'^*(Z|U)=f^*(Z)|V.
  \tag{21.10.3.1}
\]

Désignons par $\Gamma_T(\mathscr{Z}^1_X)$ le plus grand sous-faisceau de groupes commutatifs de $\mathscr{Z}^1_X$ de support contenu dans $T$;
il résulte de la relation \sref{IV.21.10.3.1} que les applications $\Gamma(U,\Gamma_T(\mathscr{Z}^1_X))\to\Gamma(V,\mathscr{Z}^1_{X'})$ que l'on vient de définir, définissent un homomorphisme de faisceaux de groupes commutatifs ordonnés sur $X'$
\[
  \psi^*(\Gamma_T(\mathscr{Z}^1_X))\to\mathscr{Z}^1_{X'}
\]
où $\psi$ est l'application continue sous-jacente au morphisme $f$.
\end{env}

\begin{proposition}[21.10.4]
\label{IV.21.10.4}
Supposons vérifiées les conditions de \sref{IV.21.10.1}.
Alors, pour tout diviseur $D$ sur $X$ tel que $\Supp(D)\subset T$ et que $f^*(D)$ soit défini \sref{IV.21.4.2}, on a
\[
\label{IV.21.10.4.1}
  \operatorname{cyc}(f^*(D))=f^*(\operatorname{cyc}(D)).
  \tag{21.10.4.1}
\]
\end{proposition}

\begin{proof}
La question étant locale sur $X$, on peut se borner au cas où $X=\Spec(A)$ est affine, où $D=\operatorname{div}(t)$, où $t$ est un élément régulier non inversible de $A$, et où le sous-préschéma $Y(D)$ \sref{IV.21.2.12} n'a qu'un seul point maximal $y$, de sorte que $\operatorname{cyc}(D)=n_y\cdot\overline{\{y\}}$, où $n_y$ est la longueur de $\sh{O}_{X,y}/t_y\sh{O}_{X,y}$ \sref{IV.21.6.5.1}.
On a vu dans la démonstration de \sref{IV.21.10.2}, B) que les points $x'\in X'$ tels que $\operatorname{mult}_{x'}(f^*(\operatorname{cyc}(D)))\neq0$ sont des points maximaux de $f^{-1}(\overline{\{y\}})$.
Si le point $x'$ est dans le cas 3\textsuperscript{o} de \sref{IV.21.10.1}, l'égalité des multiplicités en $x'$ des deux membres de \sref{IV.21.10.4.1} résulte de la définition de $n_{x'}$ au moyen de l'homomorphisme \sref{IV.21.10.1.2}.
Supposons au contraire que $x'$ soit dans le cas 2\textsuperscript{o} de \sref{IV.21.10.1}, et soit $x=f(x')$;
puisque $x\in X^{(1)}\cap\overline{\{y\}}$, on a nécessairement $x=y$.
Remarquons maintenant que $f^*(D)=\operatorname{div}(t')$, où $t'$ est l'image de $t$ dans $\Gamma(X',\sh{O}_{X'})$, et $Y(f^*(D))=f^{-1}(Y(D))$;
la multiplicité de $f^*(D)$ au point $x'$ est donc la longueur $n'_{x'}$ du $\sh{O}_{X',x'}$-module
\[
  \sh{O}_{X',x'}/t'_{x'}\sh{O}_{X',x'}
  =(\sh{O}_{X,y}/t_y\sh{O}_{X,y})\otimes_{\sh{O}_{X,y}}\sh{O}_{X',x'};
\]
comme il résulte de \sref{IV.4.7.1} que l'on a $n'_{x'}=\lambda_{x'}n_y$, les multiplicités au point $x'$ des deux membres de \sref{IV.21.10.4.1} sont encore égales, ce qui achève la démonstration.
\end{proof}

\begin{env}[21.10.5]
\label{IV.21.10.5}
\oldpage[IV-4]{292}
Supposons maintenant que $f:X'\to X$ soit un morphisme de préschémas localement noethériens, transformant tout point maximal de $X'$ en point maximal de $X$, et supposons en outre que pour toute partie fermée rare $T$ de $X$, $f$ vérifie les conditions de \sref{IV.21.10.1};
cela signifie encore que pour tout $x'\in X'^{(1)}$, ou bien $x=f(x')$ est point maximal de $X$, ou bien $x'$ vérifie l'une des conditions (ii), (iii) de \sref{IV.21.10.1}.
Si on tient compte de ce que tout cycle $1$-codimensionnel a un support rare dans $X$, on voit que $f^*(Z)$ est défini pour tout cycle $1$-codimensionnel $Z$ sur $X$;
en vertu de \sref{IV.21.10.3.1}, on a donc défini ainsi un homomorphisme de faisceaux de groupes commutatifs ordonnés
\[
  \psi^*(\mathscr{Z}^1_X)\to\mathscr{Z}^1_{X'}.
\]
Si de plus, pour tout diviseur $D$ sur $X$, $f^*(D)$ est défini \sref{IV.21.4.5}, le fait que le support de $D$ soit rare dans $X$ \sref{IV.21.6.6} entraîne que \sref{IV.21.10.4} est applicable, et on a donc la formule \sref{IV.21.10.4.1} pour tout diviseur $D$ sur $X$.
En particulier:
\end{env}

\begin{proposition}[21.10.6]
\label{IV.21.10.6}
Soient $X$, $X'$ deux préschémas localement noethériens, $f:X'\to X$ un morphisme plat.
Alors $f^*(Z)$ est défini pour tout cycle $1$-codimensionnel $Z$ sur $X$, $f^*(D)$ est défini pour tout diviseur $D$ sur $X$ et l'on a la relation \sref{IV.21.10.4.1}.
\end{proposition}

\begin{proof}
En effet, si $x'\in X'^{(1)}$ est tel que $x=f(x')$ ne soit pas maximal, il résulte de \sref{IV.6.1.1} que l'on a nécessairement $x\in X^{(1)}$, donc on est dans le cas (ii) de \sref{IV.21.10.1}.
On peut donc appliquer \sref{IV.21.10.5}, compte tenu de \sref{IV.21.4.5} et de \sref{IV.2.3.4}.
\end{proof}

\begin{remark}[21.10.7]
\label{IV.21.10.7}
L'existence de $f^*(Z)$ pour tout cycle $1$-codimensionnel $Z$ sur $X$ résulte déjà de l'hypothèse que $f$ est \emph{plat aux points $x'$ de $X'$ de codimension $\leq1$ dans $X'$} (i.e.\ tels que $\dim\sh{O}_{X',x'}\leq1$);
en effet, pour tout point maximal $z'\in X'$, il résulte de \sref{IV.6.1.1} que $z=f(z')$ est point maximal de $X$ puisque $\dim\sh{O}_{X,z}\leq\dim\sh{O}_{X',z'}=0$.
De même, si $x'\in X'^{(1)}$, $x=f(x')$ est maximal ou appartient à $X^{(1)}$ par \sref{IV.6.1.1};
on peut donc encore appliquer \sref{IV.21.10.5}.
\end{remark}

\begin{proposition}[21.10.8]
\label{IV.21.10.8}
Soient $X$, $X'$, $X''$ trois préschémas localement noethériens, $f:X'\to X$, $g:X''\to X'$ deux morphismes;
on suppose que $f$ (resp.\ $g$) est plat en tout point de codimension $\leq1$ dans $X'$ (resp.\ $X''$).
Alors $f\circ g$ est plat en tout point de codimension $\leq1$ dans $X''$ et pour tout cycle $1$-codimensionnel $Z$ sur $X$, on a $(f\circ g)^*(Z)=g^*(f^*(Z))$.
\end{proposition}

\begin{proof}
La première assertion résulte de \sref{IV.6.1.1} et \sref{IV.2.1.6}.
La seconde résulte de ce que $f$ (resp.\ $g$, resp.\ $f\circ g$) transforme les points maximaux de $X'$ (resp.\ $X''$, resp.\ $X''$) en points maximaux de $X$ (resp.\ $X'$, resp.\ $X$) et de ce que, si $x''\in X''^{(1)}$ est tel que $x'=g(x'')\in X'^{(1)}$ et $x=f(x')\in X^{(1)}$, on a
\[
  \operatorname{long}(\sh{O}_{X'',x''}/\mathfrak{m}_x\sh{O}_{X'',x''})
  =\operatorname{long}(\sh{O}_{X'',x''}/\mathfrak{m}_{x'}\sh{O}_{X'',x''})
   \cdot\operatorname{long}(\sh{O}_{X',x'}/\mathfrak{m}_x\sh{O}_{X',x'}).
\]
En effet, si l'on pose $A=\sh{O}_{X,x}$, $A'=\sh{O}_{X',x'}$, $A''=\sh{O}_{X'',x''}$, $\mathfrak{m}=\mathfrak{m}_x$, $\mathfrak{m}'=\mathfrak{m}_{x'}$, on a $A''/\mathfrak{m}A''=(A'/\mathfrak{m}A')\otimes_{A'}A''$, et la formule
\[
  \operatorname{long}_{A''}(A''/\mathfrak{m}A'')
  =\operatorname{long}_{A'}(A'/\mathfrak{m}A')
   \cdot\operatorname{long}_{A''}(A''/\mathfrak{m}'A'')
\]
résulte de \sref{IV.4.7.1} et de l'hypothèse de platitude de $A''$ sur $A'$.
\end{proof}

\begin{env}[21.10.9]
\EGAReaderTarget{ega4.c04.statement.21-10-9}{IV.21.10.9}
Soient $X$ un préschéma localement noethérien, $\Lambda$ un groupe commutatif, noté additivement, et considéré donc comme $\bb{Z}$-module;
on notera encore par $\bb{Z}$ et $\Lambda$ les faisceaux simples sur $X$ associés aux préfaisceaux constants égaux respectivement à $\bb{Z}$ et $\Lambda$ \sref[0]{0.3.6}.
On appelle \emph{faisceau des germes de cycles $1$-codimensionnels à coefficients dans $\Lambda$} le faisceau de groupes commutatifs $\mathscr{Z}^1_X\otimes_{\bb{Z}}\Lambda$.
Si $\Lambda=\bb{Q}$, on dit que les sections
\oldpage[IV-4]{293}
de $\mathscr{Z}^1_X\otimes_{\bb{Z}}\bb{Q}$ au-dessus de $X$ sont les \emph{cycles $1$-codimensionnels à coefficients rationnels}.
Comme les fibres de $\mathscr{Z}^1_X$ sont des $\bb{Z}$-modules sans torsion \sref{IV.21.6.3}, l'homomorphisme canonique $\mathscr{Z}^1_X\to\mathscr{Z}^1_X\otimes_{\bb{Z}}\bb{Q}$ est \emph{injectif}, de sorte que les cycles $1$-codimensionnels s'identifient à des cycles $1$-codimensionnels à coefficients rationnels.
\end{env}

\begin{env}[21.10.10]
\label{IV.21.10.10}
Nous allons voir que sous certaines conditions, on peut élargir la définition de $f^*(Z)$ donnée dans \sref{IV.21.10.3} pour un cycle $1$-codimensionnel $Z$ sur $X$, mais à condition de prendre pour $f^*(Z)$ un cycle $1$-codimensionnel à coefficients rationnels sur $X'$.
Le cas plus général où nous nous plaçons est celui où $f$ transforme tout point maximal de $X'$ en un point maximal de $X$, et où, en tout point $x'\in X'^{(1)}$, on a l'une des conditions (i), (ii), (iii) de \sref{IV.21.10.1} \emph{ou une quatrième condition} (en posant $x=f(x')$):
\begin{enumerate}
  \item[(iv)] $x\in X^{(1)}$, $\mathfrak{m}_x\notin\operatorname{Ass}(\sh{O}_{X,x})$, et \emph{en outre, si l'on pose} $A=\widehat{\sh{O}}_{X,x}$, $A'=\widehat{\sh{O}}_{X',x'}$, \emph{et si $K$ désigne l'anneau total des fractions de $A$, alors $A'$ est une $A$-algèbre finie et $K'=A'\otimes_AK$ est un $K$-module} libre.
\end{enumerate}

Notons alors $r_{x'}$ le rang du $K$-module libre $K'$, $q_{x'}$ le degré de $\boldsymbol{k}(x')=\boldsymbol{k}(A')$ sur $\boldsymbol{k}(x)=\boldsymbol{k}(A)$, et posons $\mu_{x'}=r_{x'}/q_{x'}$.
Pour un cycle $1$-codimensionnel de support contenu dans $T$, donné par \sref{IV.21.10.1.1} et tout $x'\in X'^{(1)}$, on définit $c_{x'}\in\bb{Q}$, comme égal au nombre $n_{x'}\in\bb{Z}$ lorsque l'on est dans l'un des cas $1\textsuperscript{o}$, $2\textsuperscript{o}$, $3\textsuperscript{o}$ de \sref{IV.21.10.1}; mais il reste ici une quatrième possibilité:
\begin{enumerate}
  \item[$4\textsuperscript{o}$] si $x'$ vérifie la condition (iv) ci-dessus, on prend $c_{x'}=\mu_{x'}n_x\in\bb{Q}$.
\end{enumerate}
\end{env}

\begin{env}[21.10.11]
\EGAReaderTarget{ega4.c04.statement.21-10-11}{IV.21.10.11}
Il faut encore prouver que lorsque la condition $4\textsuperscript{o}$ de \sref{IV.21.10.10} est vérifiée \emph{en même temps} qu'une des conditions $1\textsuperscript{o}$, $2\textsuperscript{o}$, $3\textsuperscript{o}$ de \sref{IV.21.10.1}, on a $n_{x'}=c_{x'}$.
C'est évident si $x\notin T$ puisque alors $n_x=0$.
Pour étudier les deux autres cas, notons que $\mathfrak{m}_x\sh{O}_{X',x'}$ est fermé pour la topologie $\mathfrak{m}_{x'}$-préadique, donc le complété de $\sh{O}_{X',x'}/\mathfrak{m}_x\sh{O}_{X',x'}$ pour cette dernière topologie est $A'/\mathfrak{m}_xA'$.
Si l'on est à la fois dans le cas $2\textsuperscript{o}$ et le cas $4\textsuperscript{o}$, $\sh{O}_{X',x'}/\mathfrak{m}_x\sh{O}_{X',x'}$ est discret pour la topologie $\mathfrak{m}_{x'}$-préadique, donc isomorphe à $A'/\mathfrak{m}_xA'$.
Comme $A'$ est une $A$-algèbre finie et \emph{plate} $(\mathbf{0}_{\mathrm{III}},10.2.3)$, c'est un $A$-module \emph{libre} $(\mathbf{0}_{\mathrm{III}},10.1.3)$, et le rang de $A'/\mathfrak{m}_xA'$ sur $A/\mathfrak{m}_xA=\boldsymbol{k}(x)$ est égal au rang de $A'$ sur $A$, donc aussi à celui de $K'$ sur $K$.
D'autre part ce rang est aussi égal au produit de la longueur $\lambda_{x'}$ de $\sh{O}_{X',x'}/\mathfrak{m}_x\sh{O}_{X',x'}$ (en tant que $\sh{O}_{X',x'}$-module, ou de $\boldsymbol{k}(x')$-module) par le rang $[\boldsymbol{k}(x'): \boldsymbol{k}(x)]=q_{x'}$, ce qui prouve la relation $\mu_{x'}=r_{x'}/q_{x'}=\lambda_{x'}$ dans ce cas.

Supposons enfin que l'on soit à la fois dans le cas $3\textsuperscript{o}$ et le cas $4\textsuperscript{o}$.
Alors, puisque $\dim(\sh{O}_{X,x})=1$, $\sh{O}_{X,x}$ est un anneau de valuation discrète, donc régulier.
D'autre part, $\sh{O}_{X',x'}$ est de dimension $1$, et puisque $\mathfrak{m}_{x'}\notin\operatorname{Ass}(\sh{O}_{X',x'})$, $\sh{O}_{X',x'}$ est un anneau de Cohen--Macaulay $(\mathbf{0},16.4.6)$;
enfin, puisque $A'$ est un $A$-module de type fini, $A'/\mathfrak{m}_xA'$ est un $\boldsymbol{k}(x)$-espace vectoriel de rang fini;
puisque $\sh{O}_{X',x'}/\mathfrak{m}_x\sh{O}_{X',x'}$ est contenu dans $A'/\mathfrak{m}_xA'$, c'est aussi un $\boldsymbol{k}(x)$-espace vectoriel de rang fini, donc un anneau artinien.
L'application de \sref{IV.6.1.5} montre alors que $\sh{O}_{X',x'}$ est un $\sh{O}_{X,x}$-module \emph{plat}, donc on est aussi dans le cas $2\textsuperscript{o}$, et on conclut par ce qui a été vu plus haut.

Cela étant, dans le cas que l'on a considéré, on posera
\[
  f^*(Z)=\sum_{x'\in X'^{(1)}}c_{x'}\cdot\overline{\{x'\}}.
\]
\oldpage[IV-4]{294}
qui est donc un cycle $1$-codimensionnel sur $X'$ à coefficients \emph{rationnels}.
On a encore défini ainsi un homomorphisme $f^*$ de groupes ordonnés, vérifiant \sref{IV.21.10.3.1}, et par suite un homomorphisme de faisceaux de groupes commutatifs
\[
  \psi^*(\Gamma_T(\mathscr{Z}^1_X))\to\mathscr{Z}^1_{X'}\otimes_{\bb{Z}}\bb{Q}.
\]

Lorsque $f$ vérifie les conditions précédentes pour \emph{toute partie fermée $T$ rare} dans $X$, c'est-à-dire que pour tout $x'\in X'^{(1)}$, ou bien $x=f(x')$ est point maximal de $X$, ou bien $x'$ vérifie l'une des conditions (ii), (iii) ou (iv), $f^*(Z)$ est alors défini pour tout cycle $1$-codimensionnel $Z$ sur $X$, et on a défini un homomorphisme de faisceaux de groupes commutatifs ordonnés
\[
  \psi^*(\mathscr{Z}^1_X)\to\mathscr{Z}^1_{X'}\otimes_{\bb{Z}}\bb{Q},
\]
d'où, par tensorisation, un homomorphisme de faisceaux de $\bb{Q}$-espaces vectoriels ordonnés
\[
\label{IV.21.10.11.1}
  \psi^*(\mathscr{Z}^1_X\otimes_{\bb{Z}}\bb{Q})\to\mathscr{Z}^1_{X'}\otimes_{\bb{Z}}\bb{Q}.
  \tag{21.10.11.1}
\]
\end{env}

\begin{remark}[21.10.12]
\label{IV.21.10.12}
Lorsqu'on est dans la situation de \sref{IV.21.10.10}, on peut effectivement, pour des cycles $1$-codimensionnels $Z$ sur $X$, avoir pour $f^*(Z)$ des cycles $1$-codimensionnels à coefficients \emph{non entiers}; autrement dit les nombres $\mu_{x'}$ peuvent être non entiers.
On en a un exemple en prenant l'anneau intègre complet $A$ de $(\mathbf{0}_{\mathrm{IV}},6.15.11)$, (ii) et sa clôture intégrale $A'$: le point fermé $x'$ de $\Spec(A')$ vérifie la condition (iv) de \sref{IV.21.10.10} et l'on a $\mu_{x'}=\frac12$.
\end{remark}

\begin{lemma}[21.10.13]
\label{IV.21.10.13}
Soient $A$ un anneau local noethérien de dimension $1$, $t$ un élément régulier de $A$ appartenant à l'idéal maximal $\mathfrak{m}$ (ce qui implique que $\mathfrak{m}\notin\operatorname{Ass}(A)$).
\begin{enumerate}[label=(\roman*)]
  \item Pour tout $A$-module $M$ de type fini, le noyau $N_t(M)$ et le conoyau $P_t(M)$ de l'homothétie $t_M:M\to M$ de rapport $t$ sont de longueurs finies. On pose $d_t(M)=\operatorname{long}P_t(M)-\operatorname{long}N_t(M)$.
  \item Si $0\to M'\to M\to M''\to0$ est une suite exacte de $A$-modules de type fini, on a $d_t(M)=d_t(M')+d_t(M'')$.
  \item On a $d_t(M)\geq0$ pour tout $A$-module $M$ de type fini; pour que $d_t(M)=0$, il faut et il suffit que $M$ soit de longueur finie.
  \item Si $K$ est l'anneau total des fractions de $A$ et si $M\otimes_AK$ est un $K$-module libre de rang $n$, on a $d_t(M)=n\cdot d_t(A)=n\cdot\operatorname{long}(A/tA)$.
  \item Si $M$ vérifie les hypothèses de (iv) et si de plus $t$ est $M$-régulier, on a $\operatorname{long}(M/tM)=n\cdot\operatorname{long}(A/tA)$.
\end{enumerate}
\end{lemma}

\begin{proof}
(i) $\Spec(A)$ est formé du point $\mathfrak{m}$ et des idéaux premiers minimaux $\mathfrak{p}_i$;
comme par hypothèse $t\notin\mathfrak{p}_i$ pour tout $i$ (Bourbaki, \emph{Alg. comm.}, chap.~IV, \textsection~1, n\textsuperscript{o}~1, cor.~3 de la prop.~2), l'image de $t$ dans chacun des $A_{\mathfrak{p}_i}$ est inversible, et les supports des $A$-modules de type fini $N_t(M)$ et $P_t(M)$ sont donc vides ou réduits à $\mathfrak{m}$;
on en conclut $(\mathbf{0},16.1.10)$ que ces modules sont de longueur finie.

(ii) Puisque $t$ est régulier, on a une suite exacte
\[
  0\to A\xrightarrow{t_A}A\to A/tA\to0
\]
\oldpage[IV-4]{295}
et puisque $\operatorname{Tor}_i^A(M,A)=0$ pour $i\geq1$, la suite exacte des Tor donne d'une part la suite exacte
\[
  0\to\operatorname{Tor}_1^A(M,A/tA)\to M\xrightarrow{t_M}M\to M/tM\to0
\]
et d'autre part, pour $i\geq2$,
\[
  0=\operatorname{Tor}_i^A(M,A)\to\operatorname{Tor}_i^A(M,A/tA)\to\operatorname{Tor}_{i-1}^A(M,A)=0,
\]
donc on a $N_t(M)=\operatorname{Tor}_1^A(M,A/tA)$ et $\operatorname{Tor}_i^A(M,A/tA)=0$ pour $i\geq2$;
la suite exacte des Tor donne une suite exacte
\[
  0\to\operatorname{Tor}_1^A(M',A/tA)\to\operatorname{Tor}_1^A(M,A/tA)\to\operatorname{Tor}_1^A(M'',A/tA)\to
\]
\[
  M'\otimes_A(A/tA)\to M\otimes_A(A/tA)\to M''\otimes_A(A/tA)\to0
\]
et cette suite s'écrit, d'après ce qui précède,
\[
  0\to N_t(M')\to N_t(M)\to N_t(M'')\to P_t(M')\to P_t(M)\to P_t(M'')\to0,
\]
ce qui prouve (ii).

Pour prouver (iii), notons qu'il y a une suite de composition $(M_h)$ de $M$ dont les quotients $M_h/M_{h+1}$ sont isomorphes à $A/\mathfrak{m}$ ou à un des $A/\mathfrak{p}_i$ (Bourbaki, \emph{Alg. comm.}, chap.~IV, \textsection~1, n\textsuperscript{o}~4, th.~1), et pour que $M$ soit de longueur finie, il faut et il suffit que tous ces quotients soient isomorphes à $A/\mathfrak{m}$.
Tout revient donc (en vertu de (ii)) à prouver que $d_t(A/\mathfrak{m})=0$ et $d_t(A/\mathfrak{p}_i)>0$.
Or, l'image de $t$ dans $A/\mathfrak{m}$ étant $0$, on a $N_t(A/\mathfrak{m})=A/\mathfrak{m}$ et $P_t(A/\mathfrak{m})=A/\mathfrak{m}$, d'où la première assertion;
d'autre part, l'image de $t$ dans $A/\mathfrak{p}_i$ est régulière, donc $N_t(A/\mathfrak{p}_i)=0$ et $P_t(A/\mathfrak{p}_i)=A/(tA+\mathfrak{p}_i)$ qui n'est pas réduit à $0$, d'où la seconde assertion.

(iv) Il y a une base de $M\otimes_AK$ de la forme $(x_j/s)_{1\leq j\leq n}$, où $s$ est un élément régulier de $A$ et $x_j\in M$.
Considérons l'homomorphisme $u:A^n\to M$ qui transforme l'élément $e_j$ $(1\leq j\leq n)$ de la base canonique de $A^n$ en $x_j$ et montrons que $\operatorname{Ker}(u)$ et $\operatorname{Coker}(u)$ sont de \emph{longueur finie};
en effet, pour tout $i$, l'image de $s$ dans $A_{\mathfrak{p}_i}$ est inversible, et comme $K_{\mathfrak{p}_i}=A_{\mathfrak{p}_i}$, les images des $x_j/s$ dans $M_{\mathfrak{p}_i}$ forment une base de ce $A_{\mathfrak{p}_i}$-module;
donc $u_{\mathfrak{p}_i}:A_{\mathfrak{p}_i}^n\to M_{\mathfrak{p}_i}$ est bijectif.
Comme dans (i), on en conclut que les supports de $\operatorname{Ker}(u)$ et de $\operatorname{Coker}(u)$ sont contenus dans $\{\mathfrak{m}\}$, donc que ces modules sont de longueur finie.
Cela étant, il résulte de (ii) et (iii) que l'on a $d_t(M)=d_t(A^n)=n\cdot d_t(A)=n\cdot\operatorname{long}(A/tA)$ puisque $t$ est régulier.

Enfin, il est clair que (v) se déduit aussitôt de (iv), puisque alors $N_t(M)=0$.
\end{proof}

Ce lemme permet de généraliser \sref{IV.21.10.4}:

\begin{proposition}[21.10.13]
\label{IV.21.10.13bis}
Supposons que $f$ vérifie les conditions de \sref{IV.21.10.10}.
Alors, pour tout diviseur $D$ sur $X$ tel que $\Supp(D)\subset T$ et que $f^*(D)$ soit défini, on a
\[
\label{IV.21.10.13.1}
  \operatorname{cyc}(f^*(D))=f^*(\operatorname{cyc}(D)).
  \tag{21.10.13.1}
\]
\end{proposition}

\begin{proof}
Raisonnant comme dans \sref{IV.21.10.4}, tout revient à voir (avec les mêmes notations) que si $x'$ est dans le cas $4\textsuperscript{o}$ de \sref{IV.21.10.10} et si $n_y$ est la longueur de $\sh{O}_{X,y}/t_y\sh{O}_{X,y}$, alors la
\oldpage[IV-4]{296}
longueur $n_{x'}$ de $\sh{O}_{X',x'}/t'_{x'}\sh{O}_{X',x'}$ est égale à $\mu_{x'}n_y$, $\mu_{x'}$ étant le nombre rationnel défini dans \sref{IV.21.10.10}.
Comme $\sh{O}_{X',x'}/t'_{x'}\sh{O}_{X',x'}$ est de longueur finie, il a même longueur que son complété $\mathfrak{m}_{x'}$-préadique $A'/t'_{x'}A'$, qu'on peut aussi écrire $A'/t_yA'$;
d'ailleurs, comme $t'_{x'}$ est régulier par hypothèse dans $\sh{O}_{X',x'}$, il l'est aussi dans $A'$ par platitude $(\mathbf{0}_{\mathrm{I}},6.3.4)$, et lorsque $A'$ est considéré comme $A$-module, on peut aussi dire que $t_y$ est $A'$-régulier.
Comme $A$ est de dimension $1$ et que $A'$ est un $A$-module de type fini tel que $A'\otimes_AK$ soit un $K$-module libre de rang $r_{x'}$, on peut appliquer \sref{IV.21.10.13}, (v) à $M=A'$ et à $t_y$, et la longueur de $A'/t_yA'$ en tant que $A$-module est donc $r_{x'}n_y$.
Comme $\boldsymbol{k}(x')$ est un $\boldsymbol{k}(x)$-espace vectoriel de rang $q_{x'}$, la longueur de $A'/t_yA'$ en tant que $A'$-module est donc $r_{x'}n_y/q_{x'}=\mu_{x'}n_y$, ce qui achève la démonstration.
\end{proof}

\begin{env}[21.10.14]
\EGAReaderTarget{ega4.c04.statement.21-10-14}{IV.21.10.14}
Soient maintenant $X$, $X'$ deux préschémas localement noethériens, $f:X'\to X$ un morphisme ayant les \emph{deux propriétés} suivantes:
\begin{enumerate}[label=\emph{\alph*)}]
  \item $f$ est fini;
  \item l'image par $f$ de tout point maximal de $X'$ est un point maximal de $X$.
\end{enumerate}
Pour tout $x\in X^{(1)}$, les points $x'\in f^{-1}(x)$ appartiennent alors tous à $X'^{(1)}$, comme il résulte de l'hypothèse \emph{b)} et de l'inégalité $(\mathbf{0},16.3.9.1)$, la fibre $f^{-1}(x)$ étant discrète.
Soit alors
\[
  Z'=\sum_{x'\in X'^{(1)}}n_{x'}\cdot\overline{\{x'\}}
\]
un cycle $1$-codimensionnel sur $X'$.
Pour tout $x\in X^{(1)}$, définissons un entier $n_x$ par la formule
\[
  n_x=\sum_{x'\in f^{-1}(x)}n_{x'}\cdot[\boldsymbol{k}(x'): \boldsymbol{k}(x)],
\]
ce qui a un sens, les points de $f^{-1}(x)$ étant en nombre fini et $\boldsymbol{k}(x')$ étant un corps de degré fini sur $\boldsymbol{k}(x)$ \sref[I]{I.6.4.4}.
En outre, l'ensemble des $x\in X^{(1)}$ tels que $n_x\neq0$ est \emph{localement fini} dans $X$, car il est contenu dans l'image par $f$ de l'ensemble des $x'\in X'^{(1)}$ tels que $n_{x'}\neq0$, et la conclusion résulte de ce que le morphisme $f$ est quasi-compact.
On peut donc définir un cycle $1$-codimensionnel sur $X$ en posant
\[
\label{IV.21.10.14.1}
  f_*(Z')=\sum_{x\in X^{(1)}}n_x\cdot\overline{\{x\}}
  \tag{21.10.14.1}
\]
et on dit que $f_*(Z')$ est l'\emph{image directe} de $Z'$ par $f$.
Il est clair que l'application $f_*:\mathfrak{Z}^1(X')\to\mathfrak{Z}^1(X)$ ainsi définie est un homomorphisme de groupes ordonnés.
En outre, si $U$ est un ouvert de $X$, $f_U:f^{-1}(U)\to U$ la restriction de $f$, il résulte aussitôt des définitions que l'on a
\[
\label{IV.21.10.14.2}
  (f_U)_*(Z'|f^{-1}(U))=f_*(Z')|U,
  \tag{21.10.14.2}
\]
donc, en désignant par $\psi$ l'application continue sous-jacente au morphisme $f$, les applications $\Gamma(f^{-1}(U),\mathscr{Z}^1_{X'})\to\Gamma(U,\mathscr{Z}^1_X)$ que l'on vient de définir définissent un homomorphisme de faisceaux de groupes commutatifs ordonnés sur $X$
\[
  \psi_*(\mathscr{Z}^1_{X'})\to\mathscr{Z}^1_X.
\]
\end{env}

\begin{proposition}[21.10.15]
\label{IV.21.10.15}
\oldpage[IV-4]{297}
Soient $X$, $X'$, $X''$ trois préschémas localement noethériens, $f:X'\to X$, $g:X''\to X'$ deux morphismes vérifiant les conditions \emph{a)} et \emph{b)} de \sref{IV.21.10.14}.
Alors $f\circ g$ vérifie les mêmes conditions, et pour tout cycle $1$-codimensionnel $Z''$ sur $X''$, on a
\[
  (f\circ g)_*(Z'')=f_*(g_*(Z'')).
\]
\end{proposition}

\begin{proof}
Cela résulte aussitôt des définitions.
\end{proof}

\begin{proposition}[21.10.16]
\label{IV.21.10.16}
Soient $X$, $X'$, $X_1$ trois préschémas localement noethériens, $f:X'\to X$ un morphisme vérifiant les conditions \emph{a)} et \emph{b)} de \sref{IV.21.10.14}, $g:X_1\to X$ un morphisme plat.
On pose $X'_1=X'\times_XX_1$ (de sorte que $X'_1$ est localement noethérien) et on note $f_1:X'_1\to X_1$ et $g_1:X'_1\to X'$ les projections canoniques.
Alors $f_1$ vérifie les conditions \emph{a)} et \emph{b)} de \sref{IV.21.10.14}, et pour tout $1$-cycle codimensionnel $Z'$ sur $X'$, on a
\[
\label{IV.21.10.16.1}
  g^*(f_*(Z'))=(f_1)_*(g_1^*(Z')).
  \tag{21.10.16.1}
\]
\end{proposition}

\begin{proof}
Il est clair que $f_1$ est fini, et il vérifie la condition \emph{b)} de \sref{IV.21.10.14} en vertu de \sref{IV.2.3.7}.
Pour prouver \sref{IV.21.10.16.1}, on est aussitôt ramené au cas où $X$, $X'$ et $X_1$ sont des spectres d'anneaux locaux noethériens de dimension $1$, $A$, $A'$ et $A_1$, $A'$ étant un $A$-module fini et $A_1$ un $A$-module plat.
Désignant par $x$, $x'$ et $x_1$ les points fermés de $X$, $X'$ et $X_1$ respectivement, il s'agit de voir que l'on a
\[
\label{IV.21.10.16.2}
  \sum_{x'_1}\lambda_{x'_1}[\boldsymbol{k}(x'_1):\boldsymbol{k}(x_1)]
  =\lambda_{x_1}[\boldsymbol{k}(x'):\boldsymbol{k}(x)]
  \tag{21.10.16.2}
\]
où $x'_1$ parcourt l'ensemble des points fermés de $X'_1$ (i.e. l'ensemble des points à la fois au-dessus de $x'$ et de $x_1$) et où $\lambda_{x'_1}=\operatorname{long}(\sh{O}_{X'_1,x'_1}/\mathfrak{m}_{x'}\sh{O}_{X'_1,x'_1})$ et $\lambda_{x_1}=\operatorname{long}(A_1/\mathfrak{m}_xA_1)$.
Comme on a
\[
  [\boldsymbol{k}(x'_1):\boldsymbol{k}(x')]\cdot[\boldsymbol{k}(x'):\boldsymbol{k}(x)]
  =[\boldsymbol{k}(x'_1):\boldsymbol{k}(x_1)]\cdot[\boldsymbol{k}(x_1):\boldsymbol{k}(x)],
\]
le premier membre de \sref{IV.21.10.16.2} s'écrit aussi
\[
  \frac{[\boldsymbol{k}(x'): \boldsymbol{k}(x)]}{[\boldsymbol{k}(x_1): \boldsymbol{k}(x)]}
  \operatorname{long}_{A'}(A'_1/\mathfrak{m}_{x'}A'_1)
  =\operatorname{long}_{A_1}(A'_1/\mathfrak{m}_{x'}A'_1),
\]
où on a posé $A'_1=A'\otimes_AA_1$.
On a donc $A'_1/\mathfrak{m}_{x'}A'_1=(A'/\mathfrak{m}_{x'}A')\otimes_AA_1$, et comme $A_1$ est un $A$-module plat, on a par \sref{IV.4.7.1}
\[
  \operatorname{long}_{A_1}(A'_1/\mathfrak{m}_{x'}A'_1)
  =\operatorname{long}_{A}(A'/\mathfrak{m}_{x'}A')\operatorname{long}_{A_1}(A_1/\mathfrak{m}_xA_1)
  =[\boldsymbol{k}(x'): \boldsymbol{k}(x)]\cdot\lambda_{x_1},
\]
ce qui achève la démonstration.
\end{proof}

\begin{proposition}[21.10.17]
\label{IV.21.10.17}
Soient $X$, $X'$ deux préschémas localement noethériens, $f:X'\to X$ un morphisme fini localement libre.
Alors $f$ vérifie la condition \emph{b)} de \sref{IV.21.10.14}, et pour tout diviseur $D'$ sur $X'$, on a
\[
\label{IV.21.10.17.1}
  f_*(\operatorname{cyc}(D'))=\operatorname{cyc}(f_*(D')).
  \tag{21.10.17.1}
\]
\end{proposition}

\begin{proof}
Comme $f$ est plat et fini, la condition \emph{b)} de \sref{IV.21.10.14} et la relation $f(X'^{(1)})\subset X^{(1)}$ résultent de \sref{IV.6.1.1}.
La définition \sref{IV.21.10.14.1} montre qu'on peut se ramener au cas où $X=\Spec(A)=\Spec(\sh{O}_{X,x})$ pour un $x\in X^{(1)}$;
alors $X'=\Spec(A')$, où $A'$ est une
\oldpage[IV-4]{298}
$A$-algèbre qui est un $A$-module libre de rang fini;
en outre, on peut supposer que $D'=\operatorname{div}(t')$, où $t'$ est un élément régulier de $A'$;
on a alors $f_*(D')=\operatorname{div}(t)$, où $t=\mathrm{N}_{A'/A}(t')$ est un élément régulier de $A$ \sref{IV.21.5.2}.
On peut se borner au cas où $t'$ n'est pas inversible dans $A'$, ce qui équivaut au fait que $t$ ne l'est pas dans $A$;
l'anneau $A/tA$ est alors de longueur finie et $\neq0$, et $A'/tA'$ est composé direct d'anneaux locaux artiniens $A'_i$ $(1\leq i\leq r)$, le corps résiduel de $A'_i$ étant $\boldsymbol{k}(x'_i)$, où les $x'_i$ $(1\leq i\leq r)$ sont les points de $X'$ au-dessus de $x$.
Si $t'_i$ $(1\leq i\leq r)$ est l'image de $t'$ dans $A'_i$, $A'/t'A'$ est composé direct des $A'_i/t'_iA'_i$;
comme le produit $\operatorname{long}_{A'_i}(A'_i/t'_iA'_i)\cdot[\boldsymbol{k}(x'_i):\boldsymbol{k}(x)]$ est égal à $\operatorname{long}_A(A'_i/t'_iA'_i)$, on voit que la multiplicité au point $x$ du premier membre de \sref{IV.21.10.17.1} est $\operatorname{long}_A(A'/t'A')$, si bien que la formule à démontrer se réduit à
\[
\label{IV.21.10.17.2}
  \operatorname{long}_A(A'/t'A')=\operatorname{long}_A(A/tA).
  \tag{21.10.17.2}
\]

Cette relation découle du lemme plus général suivant.
\end{proof}

\begin{lemma}[21.10.17.3]
\label{IV.21.10.17.3}
Soient $A$ un anneau local noethérien de dimension $1$, $M$ un $A$-module libre de rang fini, $u$ un endomorphisme injectif de $M$.
Alors on a
\[
\label{IV.21.10.17.4}
  \operatorname{long}_A(\operatorname{Coker}u)=\operatorname{long}_A(A/(\det u)A).
  \tag{21.10.17.4}
\]
\end{lemma}

\begin{proof}
Distinguons plusieurs cas.

I) $A$ est un anneau de valuation discrète.
Soit en effet $\pi$ une uniformisante de $A$, et remarquons que $\operatorname{long}_A(A/\pi^kA)=k$;
si $n$ est le rang de $M$ et si $\pi^{m_i}$ $(1\leq i\leq n)$ sont les facteurs invariants de $u$, $\operatorname{Coker}(u)$ est somme directe des $A$-modules $A/\pi^{m_i}A$, donc a pour longueur $m=\sum_{i=1}^nm_i$ et $\det(u)$ est produit d'un élément inversible et de $\pi^m$, d'où la conclusion dans ce cas (Bourbaki, \emph{Alg.}, chap.~VII, \textsection~4, n\textsuperscript{o}~5, cor.~1 de la prop.~4).

II) $A$ est un anneau complet intègre (de dimension $1$).
On sait alors $(\mathbf{0},19.8.8)$, (ii) qu'il y a un sous-anneau $B$ de $A$, qui est un anneau de valuation discrète, tel que $B\to A$ soit un homomorphisme local faisant de $A$ un $B$-module de type fini;
comme ce $B$-module est évidemment sans torsion, il est libre (Bourbaki, \emph{Alg. comm.}, chap.~VI, \textsection~3, n\textsuperscript{o}~6, lemme~1).
Notons $M'$ l'ensemble $M$ muni de sa structure de $B$-module (libre), par $u'$ l'endomorphisme $u$ considéré comme $B$-endomorphisme.
Il résulte de I) que l'on a
\[
\label{IV.21.10.17.5}
  \operatorname{long}_B(\operatorname{Coker}u')=\operatorname{long}_B(B/(\det u')B).
  \tag{21.10.17.5}
\]
Mais on a (Bourbaki, \emph{Alg.}, chap.~VIII, \textsection~12, n\textsuperscript{o}~2, prop.~7)
\[
  \det u'=\mathrm{N}_{A/B}(\det u),
\]
donc, en appliquant \sref{IV.21.10.17.4} à l'homothétie $x\mapsto(\det u)x$ du $B$-module libre $A$, il vient
\[
  \operatorname{long}_B(B/(\det u')B)=\operatorname{long}_B(A/(\det u)A),
\]
d'où, en portant dans \sref{IV.21.10.17.5} et divisant par $[\boldsymbol{k}(A):\boldsymbol{k}(B)]$, longueur du $B$-module $\boldsymbol{k}(A)$, on obtient \sref{IV.21.10.17.4} dans le cas envisagé.

\oldpage[IV-4]{299}
III) $A$ est un anneau complet.
Notons qu'on peut supposer en outre que $\mathfrak{m}\notin\operatorname{Ass}(A)$;
en effet, puisque $\det(u)$ est un élément régulier de $A$, si l'on avait $\mathfrak{m}\in\operatorname{Ass}(A)$, on en déduirait $\det(u)\notin\mathfrak{m}$, donc $\det(u)$ serait inversible, $u$ un automorphisme de $M$, et la formule \sref{IV.21.10.17.4} devient alors triviale, les deux membres étant nuls.

Dans ce qui suit, pour un endomorphisme $v$ d'un module $N$ sur un anneau $R$, tel que $\operatorname{Ker}v$ et $\operatorname{Coker}v$ soient de longueur finie, on posera
\[
  \chi(N,v)=\operatorname{long}_R(\operatorname{Ker}(v))-\operatorname{long}_R(\operatorname{Coker}(v)).
\]
On notera que l'hypothèse sur $v$ revient à dire que le complexe
\[
  K^0:0\to N\xrightarrow{v}N\to0
\]
a ses modules de cohomologie de longueur finie et que $\chi(N,v)=\chi(H^0(K^0))$ $(\mathbf{0}_{\mathrm{III}},11.10)$.
On en déduit que si $N'$, $N$, $N''$ sont trois $R$-modules, si l'on a un diagramme commutatif
\[
  \xymatrix{
    0 \ar[r] & N' \ar[r]^r \ar[d]^{v'} & N \ar[r]^s \ar[d]^v & N'' \ar[r] \ar[d]^{v''} & 0\\
    0 \ar[r] & N' \ar[r]_r & N \ar[r]_s & N'' \ar[r] & 0
  }
\]
dont les lignes sont exactes, et si deux des trois nombres $\chi(N,v)$, $\chi(N',v')$ et $\chi(N'',v'')$ sont définis, alors il en est de même du troisième et on a
\[
\label{IV.21.10.17.6}
  \chi(N,v)=\chi(N',v')+\chi(N'',v'').
  \tag{21.10.17.6}
\]
Cela résulte en effet aussitôt de la suite exacte de cohomologie.

Enfin, si $N$ est un $R$-module de longueur finie, on a $\chi(N,v)=0$.

Avec ces notations, on a le lemme suivant.
\end{proof}

\begin{lemma}[21.10.17.7]
\label{IV.21.10.17.7}
Soient $A$ un anneau local noethérien de dimension $1$ dont l'idéal maximal $\mathfrak{m}$ est tel que $\mathfrak{m}\notin\operatorname{Ass}(A)$;
soient $\mathfrak{p}_i$ $(1\leq i\leq n)$ les idéaux premiers minimaux de $A$.
Soient $M$ un $A$-module libre de type fini, $u$ un endomorphisme de $M$ tel que $\chi(M,u)$ soit défini;
pour chaque $i$, posons $M_i=M\otimes_A(A/\mathfrak{p}_i)$ et soit $u_i$ l'endomorphisme $u\otimes1_{A/\mathfrak{p}_i}$ de $M_i$;
alors si $\chi(M_i,u_i)$ est défini pour chaque $i$, on a
\[
\label{IV.21.10.17.8}
  \chi(M,u)=\sum_{i=1}^n\operatorname{long}(A_{\mathfrak{p}_i})\cdot\chi(M_i,u_i).
  \tag{21.10.17.8}
\]
\end{lemma}

\begin{proof}
Puisque $\mathfrak{m}\notin\operatorname{Ass}(A)$, on a une décomposition primaire réduite unique $(0)=\bigcap_{1\leq i\leq n}\mathfrak{q}_i$, où l'idéal $\mathfrak{q}_i$ est $\mathfrak{p}_i$-primaire pour $1\leq i\leq n$.
Si l'on pose $M'_i=M\otimes_A(A/\mathfrak{q}_i)$, on a donc une suite exacte
\[
  0\to M\to\bigoplus_iM'_i\to M''\to0
\]
de $A$-modules, où $M''$ est de longueur finie:
en effet, en localisant la suite exacte précédente en chacun des $\mathfrak{p}_i$, on obtient $M''_{\mathfrak{p}_i}=0$, car $(\mathfrak{q}_i)_{\mathfrak{p}_i}=0$ et $(\mathfrak{q}_j)_{\mathfrak{p}_i}=A_{\mathfrak{p}_i}$ pour $j\neq i$ (Bourbaki, \emph{Alg. comm.}, chap.~IV, \textsection~2, n\textsuperscript{o}~4, prop.~6), donc $(M'_i)_{\mathfrak{p}_i}=M_{\mathfrak{p}_i}$ et $(M'_j)_{\mathfrak{p}_i}=0$;
\oldpage[IV-4]{300}
le support de $M''$ étant donc réduit à $\mathfrak{m}$, $M''$ est de longueur finie $(\mathbf{0},16.1.10)$.
Si l'on pose $u'_i=u\otimes1_{A/\mathfrak{q}_i}$, et si $u''$ est l'endomorphisme de $M''$ déduit de $\bigoplus u'_i$ par passage aux quotients, on déduit de \sref{IV.21.10.17.6} que $\chi(M,u)+\chi(M'',u'')=\sum_i\chi(M'_i,u'_i)$, et puisque $M''$ est de longueur finie, $\chi(M'',u'')=0$.
Pour prouver \sref{IV.21.10.17.8}, on peut donc se ramener au cas où $n=1$.
On désignera alors par $\mathfrak{p}$ l'unique idéal premier minimal, qui est le nilradical de $A$;
si $A_0=A/\mathfrak{p}$, $M_0=M\otimes_AA_0$, $u_0=u\otimes1_{A_0}$, il s'agit de prouver que si $\chi(M_0,u_0)$ est défini, on a
\[
\label{IV.21.10.17.9}
  \chi(M,u)=\operatorname{long}A_{\mathfrak{p}}\cdot\chi(M_0,u_0).
  \tag{21.10.17.9}
\]

Soit $\mathfrak{n}_j$ $(0\leq j\leq r)$ la « puissance symbolique $j$-ème » de $\mathfrak{p}$, image réciproque dans $A$ de l'idéal $(\mathfrak{p}A_{\mathfrak{p}})^j$ de $A_{\mathfrak{p}}$ $(1\leq j\leq r)$, avec $\mathfrak{n}_0=A$ et $\mathfrak{n}_r=0$;
posons
\[
  M_j=\mathfrak{n}_jM/\mathfrak{n}_{j+1}M\qquad(0\leq j\leq r-1),
\]
et notons $v_j$ l'endomorphisme de $M_j$ déduit de $u$ par restriction à $\mathfrak{n}_jM$ et passage aux quotients.
Nous allons d'abord montrer que chacun des nombres $\chi(M_j,v_j)$ est défini et que l'on a
\[
\label{IV.21.10.17.10}
  \chi(M,u)=\sum_j\chi(M_j,v_j).
  \tag{21.10.17.10}
\]
La première assertion entraînera la seconde, en appliquant \sref{IV.21.10.17.6} à chacune des suites exactes
\[
  0\to\mathfrak{n}_jM/\mathfrak{n}_{j+1}M\to M/\mathfrak{n}_{j+1}M\to M/\mathfrak{n}_jM\to0.
\]

Pour prouver la première assertion, on note que si $m$ est le rang du $A$-module libre $M$, $M_j$ est isomorphe à $(\mathfrak{n}_j/\mathfrak{n}_{j+1})^m$, ou encore (puisque $\mathfrak{p}$ annule chacun des quotients $\mathfrak{n}_j/\mathfrak{n}_{j+1}$), $M_j$ est un $A_0$-module isomorphe à $M_0\otimes_{A_0}(\mathfrak{n}_j/\mathfrak{n}_{j+1})$.
Désignons par $l_j$ le rang du $A_0$-module $\mathfrak{n}_j/\mathfrak{n}_{j+1}$;
comme le corps des fractions $K_0$ de $A_0$ est le corps résiduel de $A_{\mathfrak{p}}$, $l_j$ est aussi la longueur du $A_{\mathfrak{p}}$-module $(\mathfrak{p}A_{\mathfrak{p}})^j/(\mathfrak{p}A_{\mathfrak{p}})^{j+1}$.
Il y a un système de générateurs du $A_0$-module $M_j$ qui contient une base de $M_j\otimes_{A_0}K_0$;
donc il y a un $A_0$-homomorphisme
\[
  w_j:M_0^{l_j}\to M_j
\]
dont le localisé en l'idéal $(0)$ de $A_0$ est un isomorphisme, de sorte que le support de $\operatorname{Ker}(w_j)$ et de $\operatorname{Coker}(w_j)$ est réduit à l'idéal maximal $\mathfrak{m}/\mathfrak{p}$ de $A_0$;
$\operatorname{Ker}(w_j)$ et $\operatorname{Coker}(w_j)$ sont donc des $A_0$-modules de longueur finie $(\mathbf{0},16.1.10)$.
Comme par hypothèse $\chi(M_0,u_0)$ est défini, il en est de même de $\chi(M_0^{l_j},u_0^{l_j})=l_j\chi(M_0,u_0)$ et en vertu de \sref{IV.21.10.17.6} et du fait que $\operatorname{Ker}(w_j)$ et $\operatorname{Coker}(w_j)$ sont de longueur finie, on voit que $\chi(M_j,v_j)$ est défini et égal à $l_j\chi(M_0,u_0)$;
la relation \sref{IV.21.10.17.10} donne alors
\[
  \chi(M,u)=\left(\sum_jl_j\right)\chi(M_0,u_0),
\]
et en vertu d'une remarque précédente, $\sum_jl_j$ n'est autre que la longueur de $A_{\mathfrak{p}}$, ce qui achève de prouver le lemme \sref{IV.21.10.17.7}.
\end{proof}

\oldpage[IV-4]{301}
Pour appliquer ce lemme lorsque $A$ est un anneau complet et $\mathfrak{m}\notin\operatorname{Ass}(A)$, on observe que si $u$ est injectif, il en est de même des $u_i$ (avec les notations du lemme):
en effet, $\det u$ est alors un élément régulier de $A$, donc n'appartient à aucun des $\mathfrak{p}_i$, et son image $\det u_i$ dans $A/\mathfrak{p}_i$ est par suite un élément $\neq0$ de cet anneau intègre, ce qui prouve que $u_i$ est injectif.
Comme $\operatorname{Coker}(u_i)$ est image de $\operatorname{Coker}(u)$, il est aussi de longueur finie et $\chi(M_i,u_i)$ est donc défini pour tout $i$;
on a par suite la formule \sref{IV.21.10.17.8}.
D'autre part, puisque $\det(u)$ est un élément régulier de $A$, il n'est contenu dans aucun des $\mathfrak{p}_i$;
l'idéal $(\det u)A$ est donc $\mathfrak{m}$-primaire et le quotient $A/(\det u)A$ de longueur finie.
Appliquant le même raisonnement que ci-dessus à l'homothétie injective $t:\xi\mapsto(\det u)\xi$ de $A$ et à ses images $t_i=\det u_i$ dans les $A/\mathfrak{p}_i$, il vient
\[
  \chi(A,t)=\sum_i\operatorname{long}(A_{\mathfrak{p}_i})\cdot\chi(A/\mathfrak{p}_i,t_i).
\]
Mais les anneaux $A/\mathfrak{p}_i$ sont intègres et complets, et en appliquant le résultat de II) à chacun d'eux, on obtient encore la relation \sref{IV.21.10.17.4} pour $M$ et $u$.

IV) Cas général.
Posons $A'=\widehat{A}$, $M'=M\otimes_AA'$, $u'=u\otimes1_{A'}$;
on a $\det(u')=\det(u)$, et par platitude, $\operatorname{Coker}(u')=(\operatorname{Coker}(u))_{(A')}$ et $A'/(\det u)A'=(A/(\det u)A)_{(A')}$;
comme la formule \sref{IV.21.10.17.4} est vraie pour $A'$ et $u'$ en vertu de III), elle est aussi vraie pour $A$ et $u$ en vertu de \sref{IV.4.7.1}.

Ceci termine donc la démonstration de \sref{IV.21.10.17}.

\begin{proposition}[21.10.18]
\label{IV.21.10.18}
Soient $X$, $X'$ deux préschémas localement noethériens, $f:X'\to X$ un morphisme fini, transformant tout point maximal de $X'$ en un point maximal de $X$, et vérifiant pour tout $x'\in X'$ l'une des conditions (ii), (iii), (iv) de \sref{IV.21.10.10}.
On suppose en outre qu'il existe un entier $n$ tel que, pour tout point maximal $x$ de $X$, $(f_*(\sh{O}_{X'}))_x$ soit un $\sh{O}_{X,x}$-module libre de rang $n$.
Alors, pour tout cycle $1$-codimensionnel $Z$ sur $X$, on a
\[
\label{IV.21.10.18.1}
  f_*(f^*(Z))=n\cdot Z
  \tag{21.10.18.1}
\]
(« formule de projection »).
\end{proposition}

\begin{proof}
En vertu des définitions, on est aussitôt ramené au cas où $X$ est spectre d'un anneau local noethérien $A$ de dimension $1$, de point fermé $x$, et où $Z=\overline{\{x\}}$;
posons $X'=\Spec(A')$ où $A'$ est une $A$-algèbre finie, et, pour tout idéal premier minimal $\mathfrak{p}_i$ de $A$, $A'_{\mathfrak{p}_i}$ est un $A_{\mathfrak{p}_i}$-module libre de rang $n$.
Montrons qu'on peut se borner en outre au cas où $A$ est complet.
Faisons en effet le changement de base $h:Y\to X$, où $Y=\Spec(B)$, avec $B=\widehat{A}$, et posons
\[
  Y'=X'\times_XY=\Spec(B\otimes_AA'),\qquad g=f_{\{Y\}}:Y'\to Y;
\]
le morphisme $g$ est alors fini, et comme $h$ est plat, les points maximaux de $Y$ sont au-dessus de ceux de $X$;
au-dessus de chacun des $\mathfrak{p}_i$, il n'y a qu'un nombre fini d'idéaux minimaux $\mathfrak{p}_{ij}$ de $B$, et $(B\otimes_AA')_{\mathfrak{p}_{ij}}$ est un $B_{\mathfrak{p}_{ij}}$-module libre de rang $n$.
Enfin, si $f$ vérifie l'une des conditions (ii), (iii) ou (iv) de \sref{IV.21.10.10} en chacun des points $x'$ de $f^{-1}(x)$, $g$ vérifie la condition correspondante en l'unique point $y'$ de $g^{-1}(y)$ au-dessus de $x'$ (en désignant par $y$ le point fermé de $Y$);
c'est immédiat pour les conditions (ii) et (iv);
pour la condition (iii), elle implique que $A$ est un anneau de valuation discrète, donc il en est de même de $B$, et la condition
\oldpage[IV-4]{302}
$\mathfrak{m}_{y'}\notin\operatorname{Ass}(\sh{O}_{Y',y'})$ résulte, par platitude, de la condition $\mathfrak{m}_{x'}\notin\operatorname{Ass}(\sh{O}_{X',x'})$ $(\mathbf{0},3.3.1)$.
Le morphisme $g$ vérifie donc les mêmes conditions que $f$;
si l'on prouve la formule \sref{IV.21.10.18.1} pour $g$ et $\overline{\{y\}}$, la même formule sera valable pour $f$ et $\overline{\{x\}}$, en vertu de \sref{IV.21.10.16.1}.

On peut donc supposer que $A$ est \emph{complet};
alors il en est de même de $A'$ qui est donc composé direct d'anneaux locaux complets;
on peut par suite se borner au cas où $A'$ est un anneau \emph{local}, et il reste à vérifier \sref{IV.21.10.18.1} dans chacun des cas (ii), (iii), (iv) pour le point fermé $x'$ de $X'$.
Dans le cas (ii), $A'$ étant un $A$-module plat de type fini, est un $A$-module \emph{libre} $(\mathbf{0}_{\mathrm{III}},10.1.3)$, de rang $n$ en vertu de l'hypothèse.
Or, on a par définition \sref{IV.21.10.1} et \sref{IV.21.10.3}
\[
  f^*(Z)=\lambda_{x'}\cdot\overline{\{x'\}},
\]
où $\lambda_{x'}$ est la longueur du $A'$-module $A'/\mathfrak{m}_xA'$, puis
\[
  f_*(f^*(Z))=(\lambda_{x'}[\boldsymbol{k}(x'): \boldsymbol{k}(x)])\cdot\overline{\{x\}};
\]
mais $\lambda_{x'}[\boldsymbol{k}(x'): \boldsymbol{k}(x)]$ est la longueur de $A'/\mathfrak{m}_xA'=A'\otimes_A\boldsymbol{k}(x)$ en tant que $A$-module, ou encore son rang en tant que $\boldsymbol{k}(x)$-espace vectoriel, donc est égal à $n$.

Dans le cas (iii), $A$ est un anneau de valuation discrète, donc régulier, et l'hypothèse $\mathfrak{m}_{x'}\notin\operatorname{Ass}(\sh{O}_{X',x'})$ entraîne que $A'$ est un anneau de Cohen--Macaulay $(\mathbf{0},16.4.6)$;
comme $\dim(A')=\dim(A)=1$ et que $A'/\mathfrak{m}_xA'$ est un anneau artinien, il résulte de \sref{IV.6.1.5} que $A'$ est un $A$-module plat et on est ramené au cas (ii).

Dans le cas (iv), si $K$ est l'anneau total des fractions de $A$, $K'=A'\otimes_AK$ est par hypothèse un $K$-module libre de rang $n$ et par définition \sref{IV.21.10.10}, on a
\[
  f^*(Z)=\frac{n}{[\boldsymbol{k}(x'): \boldsymbol{k}(x)]}\cdot\overline{\{x'\}},
\]
d'où $f_*(f^*(Z))=n\cdot\overline{\{x\}}$.
\end{proof}

On notera que la formule \sref{IV.21.10.18.1} est applicable en particulier lorsque le morphisme $f$ est \emph{fini} et \emph{plat} et tel que pour tout point maximal $x$ de $X$, $(f_*(\sh{O}_{X'}))_x$ soit un $\sh{O}_{X,x}$-module libre de rang $n$.

\begin{corollary}[21.10.19]
\label{IV.21.10.19}
Sous les hypothèses de \sref{IV.21.10.18}, soit $D$ un diviseur sur $X$ tel que $f^*(D)$ soit défini \sref{IV.21.4.5}; alors on a
\[
\label{IV.21.10.19.1}
  f_*(\operatorname{cyc}(f^*(D)))=n\cdot\operatorname{cyc}(D).
  \tag{21.10.19.1}
\]
\end{corollary}

\begin{proof}
Cela résulte de \sref{IV.21.10.18} et \sref{IV.21.10.13bis}.
\end{proof}

% FIN SOURCE ega4/ega4-21b-fr.tex

%
% DEBUT SOURCE ega4/ega4-21c-fr.tex
\subsection{Factorialité des anneaux locaux réguliers}
\EGAReaderTarget{ega4.c04.subsec.21-11}{subsection:IV.21.11}
\label{IV.21.11}

\begin{theorem}[21.11.1]
\label{IV.21.11.1}
\textup{(Auslander-Buchsbaum)} ---
Un anneau local noethérien régulier est factoriel.
\end{theorem}

\begin{proof}
La démonstration qui suit est due à I.~Kaplansky.

Soit $A$ un anneau local noethérien régulier de dimension $n$; nous allons raisonner
par récurrence sur $n$. Pour $n=0$, $A$ est un corps et pour $n=1$, un anneau de
valuation discrète, donc principal (et \emph{a fortiori} factoriel). Supposons donc
$n\geq2$ et le théorème démontré pour des anneaux réguliers de dimension $<n$.
Posons $X=\operatorname{Spec}(A)$, désignons par $a$ le point fermé de $X$ et posons
$U=X-\{a\}$. En tout point $y\in U$, on a
$\dim(\mathcal O_{X,y})\leq n-1$, et puisque $A$ est régulier, les anneaux
$\mathcal O_{X,y}$ le sont aussi $(\mathbf 0, 17.3.2)$, donc l'hypothèse de
récurrence entraîne qu'ils sont factoriels. De plus on a
$\operatorname{prof}(A)=\dim(A)\geq2$ puisque $A$ est régulier, donc de
Cohen--Macaulay $(\mathbf 0, 17.1.3)$.
\oldpage[IV-4]{303}
Utilisant (21.6.14), on est ramené à prouver que $\operatorname{Pic}(U)=0$.
Considérons donc un $\mathcal O_U$-Module inversible $\mathscr L$, et prouvons qu'il
est isomorphe à $\mathcal O_U$. Il résulte de $(\mathbf I, 9.4.5)$ qu'il existe un
$\mathcal O_X$-Module \emph{cohérent} $\mathscr F$ tel que
$\mathscr F|U=\mathscr L$. Comme $A$ est régulier, donc de dimension cohomologique
finie $(\mathbf 0, 17.3.1)$, il existe une résolution gauche \emph{finie} de
$\mathscr F$ :
\[
  0\leftarrow\mathscr F\leftarrow\mathcal O_X^{n_1}
  \leftarrow\mathcal O_X^{n_2}\leftarrow\cdots
  \leftarrow\mathcal O_X^{n_h}\leftarrow0
\]
$(\mathbf 0, 17.2.8$ et $\mathbf 0, 17.2.2$, (iii)). Par restriction à $U$, on a
donc une résolution finie
\[
\label{IV.21.11.1.1}
  0\leftarrow\mathscr L\leftarrow\mathcal O_U^{n_1}
  \leftarrow\mathcal O_U^{n_2}\leftarrow\cdots
  \leftarrow\mathcal O_U^{n_h}\leftarrow0.
  \tag{21.11.1.1}
\]
\end{proof}

Le théorème résultera alors des considérations générales suivantes. Sur un espace
annelé $X$, soit $\mathscr E$ un $\mathcal O_X$-Module localement libre de rang fini;
on désignera par $\Lambda^{\max}\mathscr E$ le $\mathcal O_X$-Module inversible qui,
au voisinage de chaque point de $X$, est égal à $\Lambda^p\mathscr E$, en désignant
par $p$ le rang de $\mathscr E$ dans ce voisinage (qui peut varier avec la composante
connexe de $X$). Avec cette notation, on a le

\begin{lemma}[21.11.1.2]
\label{IV.21.11.1.2}
Soient $X$ un espace annelé en anneaux locaux, et
\[
  0\leftarrow\mathscr E_0\xleftarrow{u_0}\mathscr E_1
  \leftarrow\cdots\leftarrow\mathscr E_h\leftarrow0
\]
une suite exacte de $\mathcal O_X$-Modules localement libres de rang fini; alors le
$\mathcal O_X$-Module inversible
\[
  \bigotimes_{0\leq i\leq h}
  (\Lambda^{\max}\mathscr E_i)^{\otimes(-1)^i}
\]
est isomorphe à $\mathcal O_X$.
\end{lemma}

\begin{proof}
Montrons comment ce lemme permettra de conclure la preuve de (21.11.1). Il suffit de
noter pour cela que, pour tout entier $n$,
$\Lambda^{\max}(\mathcal O_U^n)=\Lambda^n\mathcal O_U^n$ est isomorphe à
$\mathcal O_U$, ainsi que $\mathcal O_U^{\otimes(-1)}$. Comme d'autre part
$\Lambda^{\max}\mathscr L=\mathscr L$ pour tout $\mathcal O_U$-Module inversible
$\mathscr L$, le lemme (21.11.1.2), appliqué à la suite exacte (21.11.1.1) montre
que $\mathscr L$ est isomorphe à $\mathcal O_U$.

Reste à prouver (21.11.1.2); on procède par récurrence sur $h$, le lemme étant trivial
pour $h=1$. Pour $h>1$, $\mathscr N=\operatorname{Ker}(u_0)$ est un
$\mathcal O_X$-Module localement libre de rang fini $(\mathbf 0_{\mathrm I}, 5.5.5)$,
et on a les deux suites exactes
\[
  0\leftarrow\mathscr E_0\leftarrow\mathscr E_1
  \leftarrow\mathscr N\leftarrow0
\]
\[
  0\leftarrow\mathscr N\leftarrow\mathscr E_2
  \leftarrow\cdots\leftarrow\mathscr E_h\leftarrow0.
\]
En vertu de l'hypothèse de récurrence,
$(\Lambda^{\max}\mathscr N)\otimes
(\bigotimes_{2\leq i\leq h}(\Lambda^{\max}\mathscr E_i)^{\otimes(-1)^{i-1}})$
est isomorphe à $\mathcal O_X$; il suffira donc de définir un isomorphisme canonique
\[
  (\Lambda^{\max}\mathscr N)\otimes(\Lambda^{\max}\mathscr E_0)
  \xrightarrow{\sim}\Lambda^{\max}\mathscr E_1
\]
pour achever la démonstration. Or, il existe un recouvrement ouvert $(U_\alpha)$ de
$X$ tel que dans tout $U_\alpha$, $\mathscr E_1|U_\alpha$ soit somme directe de
$\mathscr N|U_\alpha$ et d'un $\mathcal O_{U_\alpha}$-Module localement libre
$\mathscr M_\alpha$ $(\mathbf 0_{\mathrm I}, 5.5.5)$, d'où un isomorphisme canonique
$v_\alpha:\mathscr M_\alpha\xrightarrow{\sim}\mathscr E_0|U_\alpha$.
Comme on a un isomorphisme canonique
\[
  r_\alpha:(\Lambda^{\max}\mathscr N|U_\alpha)
  \otimes(\Lambda^{\max}\mathscr M_\alpha)
  \xrightarrow{\sim}(\Lambda^{\max}\mathscr E_1)|U_\alpha,
\]
on en déduit au moyen de $v_\alpha$ un isomorphisme
\[
  u_\alpha:(\Lambda^{\max}\mathscr N|U_\alpha)
  \otimes(\Lambda^{\max}\mathscr E_0|U_\alpha)
  \xrightarrow{\sim}\Lambda^{\max}\mathscr E_1|U_\alpha.
\]
\oldpage[IV-4]{304}
Et il s'agit de montrer que $u_\alpha$ et $u_\beta$ coïncident dans
$U_\alpha\cap U_\beta$ pour deux indices quelconques $\alpha$, $\beta$. Or, si
$v'_\alpha$ et $v'_\beta$ sont les restrictions à $U_\alpha\cap U_\beta$ de
$v_\alpha$ et $v_\beta$ respectivement, on a
$v'_\alpha=v'_\beta\circ w_{\beta\alpha}$, où
\[
  w_{\beta\alpha}:\mathscr M_\alpha|(U_\alpha\cap U_\beta)
  \xrightarrow{\sim}\mathscr M_\beta|(U_\alpha\cap U_\beta)
\]
est la « projection parallèle à $\mathscr N$ » telle que pour toute section
$s\in\Gamma(U_\alpha\cap U_\beta,\mathscr M_\alpha)$,
$w_{\beta\alpha}(s)=s+t$ avec
$t\in\Gamma(U_\alpha\cap U_\beta,\mathscr N)$; l'identité de $u_\alpha$ et de
$u_\beta$ résulte aussitôt de ce fait et de la définition de l'isomorphisme canonique
$r_\alpha$ (Bourbaki, \emph{Alg.}, chap.~III, 3\textsuperscript{e} éd.).
\end{proof}

\subsection{Le théorème de pureté de Van der Waerden pour l'ensemble de
ramification d'un morphisme birationnel}
\EGAReaderTarget{ega4.c04.subsec.21-12}{subsection:IV.21.12}
\label{IV.21.12}

\begin{env}[21.12.1]
\EGAReaderTarget{ega4.c04.statement.21-12-1}{IV.21.12.1}
Soient $X$ et $U$ deux préschémas, $f:U\to X$ un morphisme quasi-compact et
quasi-séparé, de sorte que $f_*(\mathcal O_U)$ est une $\mathcal O_X$-Algèbre
quasi-cohérente (1.7.4). On appelle \emph{enveloppe affine} du $X$-préschéma $U$ le
$X$-préschéma affine sur $X$
\[
  U^0=\operatorname{Aff}(U/X)=\operatorname{Spec}(f_*(\mathcal O_U))
  =\operatorname{Spec}(\mathscr A(U))\qquad(\mathbf{II}, 1.1.1).
\]
Si $f^0:U^0\to X$ est le morphisme structural, on a donc par définition
\[
  \mathscr A(U^0)=f^0_*(\mathcal O_{U^0})=f_*(\mathcal O_U)=\mathscr A(U),
\]
et à l'isomorphisme identique de $f_*(\mathcal O_U)$ il correspond par
$(\mathbf{II}, 1.2.7)$ un $X$-morphisme canonique
\[
\label{IV.21.12.1.1}
  i_U:U\to U^0.
  \tag{21.12.1.1}
\]

Pour que $i_U$ soit un \emph{isomorphisme}, il faut et il suffit évidemment que le
morphisme $f:U\to X$ soit \emph{affine}. Pour tout $X$-préschéma $V$ \emph{affine sur
$X$}, l'application $u\mapsto u\circ i_U$ est une \emph{bijection}
\[
  \operatorname{Hom}_X(U^0,V)\xrightarrow{\sim}\operatorname{Hom}_X(U,V)
\]
fonctorielle en $V$: cela résulte de l'existence des bijections canoniques
\[
  \operatorname{Hom}_X(U,V)\xrightarrow{\sim}
  \operatorname{Hom}_{\mathcal O_X}(\mathscr A(V),\mathscr A(U))
\]
et
$\operatorname{Hom}_X(U^0,V)\xrightarrow{\sim}
\operatorname{Hom}_{\mathcal O_X}(\mathscr A(V),\mathscr A(U^0))$
$(\mathbf{II}, 1.2.7)$.

On peut donc dire que $U^0$ \emph{représente} le foncteur covariant
$V\mapsto\operatorname{Hom}_X(U,V)$ dans la catégorie des $X$-préschémas affines sur
$X$ $(\mathbf 0_{\mathrm{III}}, 8.1.11)$. On en déduit
$(\mathbf 0_{\mathrm{III}}, 8.1.7)$ que pour $X$ fixé,
$U\mapsto\operatorname{Aff}(U/X)$ est un foncteur covariant de la catégorie des
$X$-préschémas quasi-compacts et quasi-séparés sur $X$, dans la catégorie des
$X$-préschémas affines sur $X$; de façon plus précise, si $U_1$, $U_2$ sont deux
$X$-préschémas quasi-compacts et quasi-séparés sur $X$, à tout $X$-morphisme
$g:U_1\to U_2$ correspond l'unique $X$-morphisme $g^0:U_1^0\to U_2^0$ rendant
commutatif le diagramme
\[
  \xymatrix{
    U_1 \ar[r]^g \ar[d]_{i_{U_1}} & U_2 \ar[d]^{i_{U_2}} \\
    U_1^0 \ar[r]_{g^0} & U_2^0.
  }
\]
\end{env}

\oldpage[IV-4]{305}
Plus généralement, considérons un diagramme commutatif
\[
  \xymatrix{
    U \ar[r]^v \ar[d]_f & U' \ar[d]^{f'} \\
    X \ar[r]_u & X'
  }
\]
où les morphismes $f$, $f'$ sont quasi-compacts et quasi-séparés et le morphisme $u$
\emph{affine}. Si l'on pose $h=u\circ f$, on a
$h_*(\mathcal O_U)=u_*(f_*(\mathcal O_U))=u_*(f^0_*(\mathcal O_{U^0}))$ et
$u\circ f^0$ est un morphisme affine; on a par suite
$U^0=\operatorname{Aff}(U/X')$ (relativement au morphisme $h$), d'où un unique
$X'$-morphisme $v^0:U^0\to U'^0$ rendant commutatif le diagramme
\[
  \xymatrix{
    U \ar[r]^v \ar[d]_{i_U} & U' \ar[d]^{i_{U'}} \\
    U^0 \ar[r]_{v^0} & U'^0.
  }
\]

\begin{proposition}[21.12.2]
\label{IV.21.12.2}
Soient $f:U\to X$ un morphisme quasi-compact et quasi-séparé,
$h:X'\to X$ un morphisme plat; on pose $U'=U\times_X X'$,
$f'=f_{\{X'\}}:U'\to X'$. Alors on a un $X'$-isomorphisme canonique
\[
\label{IV.21.12.2.1}
  \operatorname{Aff}(U'/X')\xrightarrow{\sim}
  \operatorname{Aff}(U/X)\times_X X'.
  \tag{21.12.2.1}
\]
\end{proposition}

\begin{proof}
En effet, on a
$\operatorname{Aff}(U'/X')=\operatorname{Spec}(f'_*(\mathcal O_{U'}))$, et
$\operatorname{Aff}(U/X)\times_X X'=\operatorname{Spec}(h^*(f_*(\mathcal O_U)))$
$(\mathbf{II}, 1.5.1)$; l'isomorphisme de l'énoncé provient de l'isomorphisme canonique
$h^*(f_*(\mathcal O_U))\xrightarrow{\sim}f'_*(\mathcal O_{U'})$ (2.3.1).
\end{proof}

\begin{corollary}[21.12.3]
\label{IV.21.12.3}
Pour tout morphisme quasi-compact et quasi-séparé $f:U\to X$ et tout $x\in X$, on a,
à un isomorphisme canonique près
\[
  U^0\times_X\operatorname{Spec}(\mathcal O_{X,x})
  =(U\times_X\operatorname{Spec}(\mathcal O_{X,x}))^0.
\]
\end{corollary}

Il résulte aussi de (21.12.2) que l'on a, à un isomorphisme canonique près, pour tout
ouvert $V$ de $X$
\[
\label{IV.21.12.4}
  (f^{-1}(V))^0=(f^0)^{-1}(V).
  \tag{21.12.4}
\]

\begin{env}[21.12.5]
\EGAReaderTarget{ega4.c04.statement.21-12-5}{IV.21.12.5}
Nous allons considérer en particulier le cas où $f:U\to X$ est une \emph{immersion
ouverte}, $U$ s'identifiant donc à un préschéma induit sur un ouvert de $X$. Comme le
morphisme $f^{-1}(U)\to U$ est l'identité, il résulte de (21.12.4) que le morphisme
$(f^0)^{-1}(U)\to U$ restriction de $f^0$ est un isomorphisme, $i_U:U\to U^0$ étant
donc une \emph{immersion ouverte} permettant d'identifier $U$ à un préschéma induit
sur un ouvert de $U^0$.

Plus généralement, pour un ouvert $V$ de $X$, la restriction $(f^0)^{-1}(V)\to V$ de
$f^0$ est un isomorphisme de $(f^0)^{-1}(V)$ sur le préschéma induit sur l'ouvert
$V\cap U$ si et seulement si l'immersion ouverte $U\cap V\to V$ est un morphisme
\emph{affine}. Il est clair $(\mathbf{II}, 1.2.1)$
\oldpage[IV-4]{306}
que la réunion de ces ouverts $V$ est le \emph{plus grand} d'entre eux, $U_1$, qui
contient $U$; en vertu de ce qui précède, $U_1$ est aussi le plus grand ouvert ne
rencontrant pas l'ensemble
\[
  \operatorname{Daf}(U/X)=f^0(U^0)-U
\]
(« défaut d'affinité » de l'ouvert $U$ relativement à $X$, qui est vide si et seulement
si $U$ est affine sur $X$); autrement dit, l'ensemble fermé $Z=X-U_1$ est
l'\emph{adhérence} de l'ensemble $\operatorname{Daf}(U/X)$.

On notera que pour tout morphisme \emph{plat} $h:X'\to X$, si l'on pose
$U'=h^{-1}(U)$, on a
\[
\label{IV.21.12.5.1}
  \operatorname{Daf}(U'/X')=h^{-1}(\operatorname{Daf}(U/X))
  \tag{21.12.5.1}
\]
comme il résulte aussitôt de (21.12.2.1) et de $(\mathbf I, 3.4.8)$. En particulier,
pour tout ouvert $V$ de $X$, on a
\[
\label{IV.21.12.5.2}
  \operatorname{Daf}((U\cap V)/V)=\operatorname{Daf}(U/X)\cap V
  \tag{21.12.5.2}
\]
et pour tout $x\in X$,
\[
\label{IV.21.12.5.3}
  \operatorname{Daf}((U\cap\operatorname{Spec}(\mathcal O_{X,x}))
  /\operatorname{Spec}(\mathcal O_{X,x}))
  =\operatorname{Daf}(U/X)\cap\operatorname{Spec}(\mathcal O_{X,x}).
  \tag{21.12.5.3}
\]
\end{env}

Nous allons, lorsque $X$ est localement noethérien, donner des précisions sur la nature
de l'ensemble $\operatorname{Daf}(U/X)$, qui montreront par exemple que lorsque $U$ est
partout dense dans $X$, $\operatorname{Daf}(U/X)$ n'est pas un ensemble fermé rare
quelconque :

\begin{theorem}[21.12.6]
\label{IV.21.12.6}
Soient $X$ un préschéma localement noethérien, $U$ un ouvert non vide de $X$,
$f:U\to X$ l'injection canonique. Alors :
\begin{enumerate}[label=(\roman*)]
  \item L'adhérence $Z=X-U_1$ de $\operatorname{Daf}(U/X)$ est de codimension
  $\geq2$ dans $X$.
  \item Si $T=X-U\supset Z$ est de codimension $\geq2$, le morphisme
  $f^0:U^0\to X$ est surjectif.
\end{enumerate}
\end{theorem}

\begin{proof}
(i) Montrons d'abord que pour tout point $x\in\operatorname{Daf}(U/X)$ on a
nécessairement $\dim(\mathcal O_{X,x})>1$. En effet, $x$ ne peut évidemment être point
maximal de $X$, étant contenu dans $\overline U-U$; il s'agit donc de voir que l'on ne
peut avoir $\dim(\mathcal O_{X,x})=1$. Or, cette relation entraînerait, par (21.12.5.3),
$x\in\operatorname{Daf}(U^{(x)}/\operatorname{Spec}(\mathcal O_{X,x}))$, où
$U^{(x)}=U\cap\operatorname{Spec}(\mathcal O_{X,x})$. Mais les seuls ouverts de
$\operatorname{Spec}(\mathcal O_{X,x})$ sont $\operatorname{Spec}(\mathcal O_{X,x})$
lui-même et les parties de l'ensemble (fini) des points maximaux de
$\operatorname{Spec}(\mathcal O_{X,x})$. Or, l'ensemble ouvert réduit à un point maximal
de $\operatorname{Spec}(\mathcal O_{X,x})$ est affine, par définition des préschémas; on
en conclut que \emph{tous} les ensembles ouverts dans
$\operatorname{Spec}(\mathcal O_{X,x})$ sont affines, donc
$\operatorname{Daf}(U^{(x)}/\operatorname{Spec}(\mathcal O_{X,x}))=\varnothing$
contrairement à l'hypothèse faite.

Pour prouver (i) il faut montrer davantage, savoir que si $x\in X$ est tel que
$\dim(\mathcal O_{X,x})=1$, il existe un voisinage ouvert $W$ de $x$ dans $X$ tel que
$W$ ne rencontre pas $\operatorname{Daf}(U/X)$, c'est-à-dire tel que l'injection
canonique $f_W:U\cap W\to W$ soit affine. Mais, avec les mêmes notations que ci-dessus,
on vient de voir que l'injection canonique
$f^{(x)}:U^{(x)}\to\operatorname{Spec}(\mathcal O_{X,x})$ est affine. On peut
évidemment se borner au cas où $X$ est noethérien; comme le morphisme $f$ est de
présentation finie, la conclusion résulte de (8.10.5, (viii)) appliqué suivant la
méthode de (8.1.2, a)).

\oldpage[IV-4]{307}
(ii) Il s'agit de prouver que pour tout point $x\in T$, on a $x\in f^0(U^0)$; nous
allons montrer d'abord qu'on peut se réduire au cas où $X=\operatorname{Spec}(A)$, où
$A$ est un anneau local noethérien \emph{complet}, et $x$ le point fermé de $X$. Pour
cela, il suffit de faire le changement de base
$h:X'=\operatorname{Spec}(\widehat{\mathcal O}_{X,x})\to X$; si l'on pose
$U'=h^{-1}(U)$, $f'=f_{\{X'\}}$ est l'injection canonique $U'\to X'$, et puisque le
morphisme $h$ est plat, il résulte de (21.12.2) que si l'on prouve que $x$ appartient à
$f^{\prime 0}(U^{\prime 0})$, on en déduira que $x\in f^0(U^0)$. En vertu de (6.1.1), on a
bien opéré la réduction cherchée.

Soit alors $X_1$ un sous-préschéma fermé réduit de $X$ ayant pour espace sous-jacent
une composante irréductible de $X$, de dimension \emph{maximale}, parmi celles qui
contiennent une composante irréductible de $T$, et posons
$U_1=U\cap X_1$, $T_1=T\cap X_1=X_1-U_1$; on a donc
$\operatorname{codim}(T_1,X_1)\geq2$ $(\mathbf 0, 14.2.1)$, et le couple $(U_1,X_1)$
vérifie donc les mêmes hypothèses que le couple $(U,X)$ ($X_1$ étant spectre d'un
anneau quotient de $A$, donc local noethérien et complet). On a par ailleurs un
diagramme commutatif
\[
  \xymatrix{
    U_1 \ar[r]^j \ar[d]_{f_1} & U \ar[d]^f \\
    X_1 \ar[r]_i & X
  }
\]
où $i$ et $j$ sont les injections canoniques, $i$ étant donc un morphisme affine; on en
déduit (21.12.1) l'existence d'un morphisme $j^0:U_1^0\to U^0$ rendant commutatif le
diagramme
\[
  \xymatrix{
    U_1^0 \ar[r]^{j^0} \ar[d]_{f_1^0} & U^0 \ar[d]^{f^0} \\
    X_1 \ar[r]_i & X
  }
\]
et par suite, pour prouver que $x\in f^0(U^0)$, il suffit de prouver que
$x\in f_1^0(U_1^0)$. On peut donc remplacer $X$ par $X_1$, et on peut par suite
supposer que l'anneau $A$ est en outre \emph{intègre}. Mais $A$ est quotient d'un
anneau noethérien régulier en vertu du théorème de Cohen $(\mathbf 0, 19.8.8$, (i)) et
puisqu'il est intègre, on peut appliquer (5.11.1) avec la famille $(x_\alpha)$ réduite au
seul point maximal de $X$; comme $\operatorname{codim}(T,X)\geq2$ par hypothèse, on
voit que $f_*(\mathcal O_U)$ est un $\mathcal O_X$-Module \emph{cohérent}, donc le
morphisme $f^0:U^0\to X$ est fini; comme ce morphisme est dominant ($U$ étant partout
dense), il est surjectif $(\mathbf{II}, 6.1.10)$, et par suite $x\in f^0(U^0)$. C.Q.F.D.
\end{proof}

\begin{corollary}[21.12.7]
\label{IV.21.12.7}
Soient $X$ un préschéma localement noethérien, $U$ un sous-préschéma induit sur un
ouvert de $X$, $j:U\to X$ l'immersion canonique; supposons que $j$ soit un morphisme
affine. Alors toute composante irréductible de $T=X-U$ est de codimension $\leq1$
(et par suite de codimension $1$ si $U$ est partout dense).
\end{corollary}

\begin{proof}
Supposons en effet qu'une des composantes irréductibles $T_1$ de $T$ soit de
codimension $\geq2$ dans $X$. Remplaçant au besoin $X$ par un voisinage ouvert du
point générique de $T_1$, on peut supposer $T$ irréductible et de codimension $\geq2$.
Mais alors l'hypothèse
\oldpage[IV-4]{308}
que $j:U\to X$ est affine implique que $U^0$ s'identifie à $U$ et $j^0$ à $j$,
autrement dit que $j^0(U^0)=U$; or, cela contredit la conclusion de (21.12.6) qui,
sous l'hypothèse de dimension faite implique que $j^0$ doit être surjectif.
\end{proof}

\begin{env}[21.12.8]
\label{IV.21.12.8}
Soient $A$ un anneau local noethérien, $Y=\operatorname{Spec}(A)$, $y$ l'unique point
fermé de $Y$, $Y'=Y-\{y\}$. Considérons la condition suivante :
\begin{itemize}
  \item[\textup{(W)}] \emph{Pour tout ouvert $U$ de $Y$ contenu dans $Y'$, ne
  contenant aucune composante irréductible de $Y'$ et tel que l'immersion canonique
  $U\to Y'$ soit affine, $U$ est lui-même un ouvert affine.}
\end{itemize}
Cette condition est entraînée par la suivante :
\begin{itemize}
  \item[\textup{(W')}] \emph{Pour toute partie fermée $T$ de $Y$ dont toute composante
  irréductible est de codimension $1$ dans $Y$, et tel que pour toute composante
  irréductible $Y_i$ de $Y$, on ait $Y_i\cap T\ne\{y\}$, l'ouvert $U=Y-T$ est affine.}
\end{itemize}
En effet, si (W') est vérifiée et si l'ouvert $U$ de $Y$ vérifie l'hypothèse de (W), il
résulte de (21.12.7) qu'aucune composante irréductible de $T=Y-U$ ne peut être de
codimension $\geq2$; comme par ailleurs $U$ ne contient aucune des composantes
irréductibles $Y_i-\{y\}$ de $Y'$, la condition (W') montre que $U$ est affine.

On notera que la condition (W') se simplifie lorsque $Y$ est \emph{irréductible} et est
alors équivalente à la suivante :
\begin{itemize}
  \item[\textup{(W'')}] \emph{Pour toute partie fermée irréductible $T$ de $Y$, de
  codimension $1$ dans $Y$, $Y-T$ est un ouvert affine.}
\end{itemize}
En effet, il est clair que (W') entraîne (W'') lorsque $Y$ est irréductible, et la
réciproque est vraie en considérant les composantes irréductibles de $T$ et en notant
que l'intersection d'un nombre fini d'ouverts affines est un ouvert affine
$(\mathbf I, 5.5.6)$.
\end{env}

\begin{env}[21.12.9]
\label{IV.21.12.9}
\emph{Exemples}. ---
Si $A$ est un anneau local noethérien \emph{factoriel}, il vérifie (W') (et
\emph{a fortiori} (W)), puisque tout idéal premier de hauteur $1$ est principal
$(\mathbf I, 1.3.6)$. Mais il y a des anneaux locaux noethériens non factoriels
vérifiant (W''), par exemple ceux de dimension $\leq1$: en effet, on a remarqué dans la
démonstration de (21.12.6, (i)) que \emph{tous} les ouverts de $Y$ sont alors affines.
On peut d'autre part prouver en utilisant la théorie de la dualité locale (chap.~III,
3\textsuperscript{e} partie) que tout anneau local noethérien de dimension $2$ vérifie
(W').
\end{env}

L'intérêt de la condition (W) réside dans le résultat suivant :

\begin{proposition}[21.12.10]
\label{IV.21.12.10}
Soient $X$, $Y$ deux préschémas localement noethériens, $g:X\to Y$ un morphisme,
$y$ un point fermé de $Y$, $Y'=Y-\{y\}$, $X'=g^{-1}(Y')$; supposons que le morphisme
$g':X'\to Y'$ restriction de $g$ soit une immersion ouverte, et que l'anneau local
$\mathcal O_{Y,y}$ vérifie la condition (W) (21.12.8). Alors, pour toute composante
irréductible $Z$ de $X_y=g^{-1}(y)$, ou bien $Z$ est de codimension $\leq1$ dans $X$,
ou bien son point générique est isolé dans $Z$. Si $Z$ est localement de type fini sur
$k(y)$, la seconde alternative implique que $Z$ est réduit à un seul point.
\end{proposition}

\begin{proof}
La dernière assertion de l'énoncé résulte de ce que $Z$ est alors un préschéma de
Jacobson (10.4.7), et comme l'ensemble des points fermés de $Z$ est alors dense dans
$Z$, le point générique de $Z$ ne peut être isolé dans $Z$ que si $Z$ est réduit à un
seul point.

\oldpage[IV-4]{309}
Supposons qu'il existe dans $X_y$ une composante irréductible $Z$ dont le point
générique $z$ n'est pas isolé dans $Z$, et telle que $\operatorname{codim}(Z,X)\geq2$.
La question étant locale sur $X$ et $Y$ puisque $z$ est non isolé dans $Z$, on peut, en
remplaçant $X$ par un voisinage ouvert de $z$ dans $X$, supposer $X$ et $Y$ affines,
$X_y=Z$ irréductible, non réduit à un point et telle que
$\operatorname{codim}(X_y,X)\geq2$. L'image $g(X-X_y)=U$ est un ouvert de $Y'$
isomorphe à $X-X_y$ par hypothèse; montrons qu'en remplaçant au besoin $X$ par un
voisinage ouvert de $z$ dans $X$, on peut supposer que $U$ ne contient aucune des
composantes irréductibles de $Y'$ dont l'adhérence contient $y$.

En effet, on peut tout d'abord supposer que tous les points maximaux de $X$ sont des
générisations de $z$, l'ensemble de ces points étant fini; donc l'ensemble des points
maximaux de $U$ est l'ensemble des images $y_i$ par $g$ des points maximaux $x_i$ de
$X$ (aucun point maximal de $X$ ne pouvant par hypothèse être contenu dans $X_y$,
puisque $\dim(\mathcal O_{X,z})\geq2$). Posons $X_i=\overline{\{x_i\}}$. Par
hypothèse, on a $z\in X_i$ et puisque $z$ n'est pas isolé dans $X_y$, il existe dans
$X_y$ un point $x\ne z$. Comme $X_i\ne Z$, il existe dans $X_i$ un point $t_i$ tel
que si l'on pose $T_i=\overline{\{t_i\}}$, on ait $\dim(\mathcal O_{T_i,x})=1$
(10.5.9); cela entraîne par hypothèse $t_i\notin Z$, donc $t_i$ n'est pas une
générisation de $z$. Remplaçant $X$ par $X'=X-\bigcup_iT_i$, on voit que l'image par
$g$ de $X'-X_y$ ne contient pas les images $g(t_i)$, donc ne contient aucune des
composantes irréductibles de $Y'$ dont l'adhérence contient $y$.

Cela étant, comme $X$ et $Y$ sont affines, le morphisme $g:X\to Y$ est affine, donc il
en est de même de la restriction $g':X'\to Y'$; puisque $g'$ est une immersion ouverte,
l'immersion canonique $U\to Y'$ est affine. Posons
$Y_1=\operatorname{Spec}(\mathcal O_{Y,y})$, $Y'_1=Y_1-\{y\}=Y'\cap Y_1$,
$U_1=U\cap Y_1$. Ce qui précède prouve que l'immersion canonique $U_1\to Y'_1$ est
affine, donc $U_1$ est un ouvert affine dans $Y_1$ en vertu de (W), en d'autres termes
l'immersion canonique $U_1\to Y_1$ est affine. Mais puisque cette immersion est de
présentation finie (1.6.2), il résulte de (8.10.5, (viii)) appliqué suivant la méthode de
(8.1.2, a)) qu'en restreignant au besoin $Y$ à un voisinage ouvert affine de $y$, on peut
supposer que l'immersion $U\to Y$ est affine. On en conclurait que l'ouvert $X-X_y$ de
$X$, isomorphe à $U$, serait affine, ce qui contredirait (21.12.7); la proposition est
ainsi démontrée.
\end{proof}

\begin{corollary}[21.12.11]
\label{IV.21.12.11}
Soient $X$, $Y$ deux préschémas localement noethériens, $X$ étant supposé
irréductible; soit $g:X\to Y$ un morphisme localement de type fini, et soit $V$ le plus
grand ouvert de $X$ tel que la restriction $g|V:V\to Y$ soit un isomorphisme local.
Supposons que pour tout point $y\in Y$, l'anneau local $\mathcal O_{Y,y}$ vérifie la
condition (W) (21.12.8). Alors toute composante irréductible de $T=X-V$ est, ou bien
de codimension $\leq1$, ou bien telle que son point générique $z$ soit isolé dans
$g^{-1}(g(z))$.
\end{corollary}

\begin{proof}
Posons $g(z)=y$. La question ne dépendant que de la fibre $g^{-1}(y)$ et de l'anneau
$\mathcal O_{Y,y}$, on peut par changement de base se borner au cas où
$Y=\operatorname{Spec}(\mathcal O_{Y,y})$ et où $X$ est affine
$((\mathbf I, 3.6.5)$ et $(\mathbf I, 4.5.5))$; remplaçant éventuellement $X$ par un
voisinage ouvert de $z$ dans $X$, on peut supposer que $T$ est irréductible; en outre on
peut se borner au cas où $V$ est non vide. Le morphisme $g$ peut donc être supposé
séparé et $y$ fermé dans $Y$; puisque la restriction $g|V:V\to Y$ est un isomorphisme
local et que l'ouvert non vide $V$ de $X$ est irréductible, il résulte de
$(\mathbf I, 8.2.8)$ que $g|V:V\to Y$ est une immersion ouverte. La restriction
$g|(V\cap X_y):V\cap X_y\to\{y\}=\operatorname{Spec}(k(y))$ est donc aussi une
immersion ouverte, ce qui montre que $V\cap X_y$ est, soit vide, soit réduit à un point
$x$ rationnel
\oldpage[IV-4]{310}
sur $k(y)$, donc fermé dans $X_y$ $(\mathbf I, 6.4.2)$. Remplaçant encore
éventuellement $X$ par un voisinage ouvert plus petit de $z$ dans $X$, on peut donc se
borner au cas où $V\cap X_y=\varnothing$, autrement dit $T\supset X_y$; comme d'autre
part $z\in X_y$ est le point générique de $T$, on a
$g(T)\subset\overline{\{y\}}=\{y\}$, autrement dit $T=X_y$. On est alors exactement
dans les conditions d'application de (21.12.10), d'où la conclusion.
\end{proof}

\begin{theorem}[21.12.12]
\label{IV.21.12.12}
\textup{(van der Waerden)}. ---
Soient $X$, $Y$ deux préschémas localement noethériens intègres, $g:X\to Y$ un
morphisme birationnel et localement de type fini. On suppose en outre que $Y$ est
normal et que pour tout $y\in Y$, l'anneau $\mathcal O_{Y,y}$ vérifie la condition (W)
de (21.12.8) (conditions remplies en particulier lorsque $Y$ est localement factoriel
(21.12.9)). Si $V$ est le plus grand ouvert de $X$ tel que la restriction
$g|V:V\to Y$ soit un isomorphisme local, toutes les composantes irréductibles de
$T=X-V$ sont de codimension $1$ dans $X$.
\end{theorem}

\begin{proof}
On notera que puisque $g$ est birationnel, l'ensemble ouvert $V$ n'est pas vide; il
suffit donc de prouver qu'un point maximal $z$ de $T$ ne peut être isolé dans $X_y$,
où $y=g(z)$. Mais, quitte à se restreindre à un voisinage ouvert de $z$, on peut se
borner au cas où $T=X_y$; alors toutes les fibres $g^{-1}(y')$ ($y'\in Y$) seraient
vides ou réduites à un point et il résulterait du « Main theorem » (8.12.10) que $g$
serait un isomorphisme local, contrairement à l'hypothèse $z\in T$.
\end{proof}

\begin{corollary}[21.12.13]
\label{IV.21.12.13}
Supposons vérifiées les hypothèses de (21.12.12) et en outre, supposons que $g$ soit
quasi-fini en chacun des points de $X^{(1)}$ (on rappelle que ce sont les points où
$\dim(\mathcal O_{X,x})=1$); alors $g$ est un isomorphisme local, et si de plus $g$ est
séparé, $g$ est une immersion ouverte.
\end{corollary}

\begin{proof}
Il suffit de prouver la première assertion, la seconde étant une conséquence de la
première et de $(\mathbf I, 8.2.8)$. Tout revient à prouver, avec les notations de
(21.12.12), que $T=\varnothing$; dans le cas contraire, un point générique $z$ d'une
composante irréductible de $T$ appartiendrait à $X^{(1)}$ en vertu de (21.12.12), et par
hypothèse il serait isolé dans $X_y$ avec $y=g(z)$; mais on a vu dans la démonstration
de (21.12.12) que cela n'est pas possible.
\end{proof}

\begin{remarks}[21.12.14]
\label{IV.21.12.14}
(i) La conclusion de (21.12.12) s'applique lorsque $Y$ est régulier et intègre et $X$
intègre, car en vertu du théorème d'Auslander--Buchsbaum (21.11.1), $Y$ est alors
localement factoriel. Par contre, on connaît des exemples où $X$ et $Y$ sont des schémas
algébriques normaux de dimension $3$, sur un corps algébriquement clos de
caractéristique quelconque, et où la conclusion de (21.12.12) n'est plus valable.

(ii) L'ensemble $T$ de (21.12.12) est aussi l'ensemble des points où $g$ est
\emph{ramifié}. En effet, si $g$ est non ramifié en un point $x\in X$, il est aussi non
ramifié dans un voisinage ouvert affine $U$ de $x$ (17.3.7), donc $g|U$ est un morphisme
séparé, quasi-fini et birationnel (17.4.3). Puisque $Y$ est normal, on conclut du « Main
theorem » $(\mathbf{III}, 4.4.9)$ que $g$ est un isomorphisme local au point $x$, donc
$x\in X-T$. Réciproquement, $g$ est évidemment non ramifié en tout point où il est un
isomorphisme local. Ceci justifie le titre donné à cette section.

\oldpage[IV-4]{311}
(iii) Sans supposer que $Y$ soit localement factoriel, mais en supposant par contre que
$X$ est intersection complète relativement à $Y$ (chap.~V), nous verrons au chap.~V
qu'on a un résultat plus précis que (21.12.12), en exprimant $T$ comme cycle
1-codimensionnel d'une section d'un $\mathcal O_X$-Module inversible, canoniquement
associée à $g$. Ceci s'appliquera notamment quand $X$ et $Y$ sont tous deux réguliers,
ou lorsque $g$ est un $S$-morphisme, $X$ et $Y$ étant des préschémas lisses sur $S$.

(iv) On peut se demander si (21.12.10) admet une réciproque; autrement dit, pour un
anneau local noethérien donné $A$, si l'on pose $Y=\operatorname{Spec}(A)$, $y$
désignant le point fermé de $Y$, et si, pour tout préschéma localement noethérien $X$ et
tout morphisme $g:X\to Y$ tel que $g':X'\to Y'$ soit une immersion ouverte, toute
composante irréductible de $X_y$ est de codimension $\leq1$ dans $X$ ou est réduite à
un point, alors est-il vrai que $A$ vérifie la condition (W) de (21.12.8)?

(v) Soient $Y$ un préschéma intègre \emph{régulier}, $X$ un préschéma localement
noethérien \emph{normal}, $g:X\to Y$ un morphisme localement de type fini; supposons
en outre que, si $\eta$ est le point générique de $Y$, la fibre $X_\eta$ soit \emph{étale}
sur $k(\eta)$, et soit $T'$ le complémentaire du plus grand ouvert $V'$ de $X$ tel que
$g|V':V'\to Y$ soit étale; est-il vrai alors que toutes les composantes irréductibles de
$T'$ soient de codimension $1$? Il en est bien ainsi lorsqu'on suppose en outre que $g$
est \emph{localement quasi-fini} (« théorème de pureté » de Zariski--Nagata). On peut
montrer que la conjecture précédente serait conséquence de la suivante (et lui serait
même équivalente si la conjecture de (iv) était vérifiée): pour un anneau local régulier
$A$ contenu dans un anneau local intègre et intégralement clos $B$, tel que $B$ soit un
$A$-module fini, l'ouvert $U$ des points de $\operatorname{Spec}(B)$ en lesquels
$\operatorname{Spec}(B)$ est non ramifié sur $\operatorname{Spec}(A)$ (ou, ce qui
revient au même par (18.10.1), étale sur $\operatorname{Spec}(A)$) est affine.
\end{remarks}

\begin{proposition}[21.12.15]
\label{IV.21.12.15}
Soient $S$ un préschéma, $g:X\to S$ un morphisme plat et localement de présentation
finie, dont les fibres $X_s=g^{-1}(s)$ sont géométriquement irréductibles (4.5.2),
$h:Y\to S$ un morphisme lisse; on note $Y_s=h^{-1}(s)$ les fibres de ce morphisme.
Soit $f:X\to Y$ un $S$-morphisme propre tel que, pour tout point maximal $\eta$ de
$S$, le morphisme $f_\eta:X_\eta\to Y_\eta$ soit un isomorphisme. Alors $f$ est un
isomorphisme.
\end{proposition}

\begin{proof}
Puisque $h$ est plat, tout point maximal de $Y$ est au-dessus d'un point maximal de $S$
(2.3.4), donc la réunion des $Y_\eta$, lorsque $\eta$ parcourt l'ensemble des points
maximaux de $S$, est dense dans $Y$; comme les $f_\eta$ sont des isomorphismes, $f$
est dominant et par suite \emph{surjectif} puisqu'il est propre. Il suffit donc de montrer
que $f$ est une \emph{immersion ouverte}. Compte tenu de (17.9.5), il suffit de prouver
que pour tout $s\in S$, $f_s:X_s\to Y_s$ est une immersion ouverte.

Comme tout $s\in S$ est générisation d'un point maximal $\eta$, on peut déjà, en
faisant le changement de base $\operatorname{Spec}(\mathcal O_{S',s})\to S$, où $S'$
est le sous-préschéma réduit de $S$ ayant $\overline{\{\eta\}}$ comme espace
sous-jacent, se ramener au cas où $S$ est un schéma local et intègre de point fermé $s$:
les fibres aux points $\eta$ et $s$ ne sont en effet pas changées et les propriétés des
morphismes $g$, $h$ et $f$ se conservent par changement de base. En outre, la question
est locale sur $Y$; remplaçant $Y$ par un voisinage ouvert affine $U$ dans $Y$ d'un point
de $Y_s$, et $X$ par $f^{-1}(U)$, qui est quasi-compact puisque $f$ est propre, on voit
qu'on peut supposer en outre que $X$ et $Y$ sont de présentation finie sur $S$.

Il existe alors un sous-anneau local noethérien $A_0$ de
$A=\mathcal O_{S,s}$ tel que $A_0\to A$ soit un homomorphisme local, deux morphismes
de type fini $g_0:X_0\to S_0=\operatorname{Spec}(A_0)$,
$h_0:Y_0\to S_0$ et un $S_0$-morphisme
\oldpage[IV-4]{312}
$f_0:X_0\to Y_0$ tels que $g$, $h$ et $f$ se déduisent de $g_0$, $h_0$, $f_0$ par le
changement de base $S\to S_0$ ((8.9.1) et (5.13.3)); on peut en outre supposer $g_0$
plat (11.2.7), $h_0$ lisse (17.7.9) et $f_0$ propre (8.10.5, (xii)). D'autre part, la
projection sur $S_0$ du point générique $\eta$ de $S$ est le point générique $\eta_0$ de
$S_0$ et il résulte de (2.7.1, (viii)) que le morphisme
$(f_0)_{\eta_0}:(X_0)_{\eta_0}\to(Y_0)_{\eta_0}$ est un isomorphisme; enfin le point
fermé $s$ de $S$ est au-dessus du point fermé $s_0$ de $S_0$, donc la fibre
$(X_0)_{s_0}$ de $g_0$ est géométriquement irréductible (4.5.6).

On voit donc qu'on peut se borner au cas où $S$ est le spectre d'un anneau local intègre
noethérien, de point fermé $s$, affaiblir l'hypothèse sur les fibres en supposant seulement
que $X_s$ et $X_\eta$ sont géométriquement irréductibles, et prouver que sous ces
hypothèses $f_s$ est une immersion ouverte. Il y a alors un anneau de valuation discrète
$A'$ et un morphisme $S'=\operatorname{Spec}(A')\to S$ qui transforme le point fermé
$a$ de $S'$ en $s$ et le point générique $b$ de $S'$ en $\eta$ $(\mathbf{II}, 7.1.9)$;
soient $g':X'\to S'$, $h':Y'\to S'$, $f':X'\to Y'$ les morphismes déduits de $g$,
$h$, $f$ par le changement de base $S'\to S$, qui vérifient de nouveau les hypothèses
vérifiées par $g$, $h$, et $f$ dans l'énoncé (21.12.15); si l'on prouve que
$f'_a:X'_a\to Y'_a$ est une immersion ouverte, il résultera de (2.7.1, (x)) qu'il en
sera de même de $f_s:X_s\to Y_s$. On peut donc finalement se borner au cas où $S$ est
le \emph{spectre d'un anneau de valuation discrète}.

Alors, puisque $h:Y\to S$ est lisse, $Y$ est régulier (17.5.8), et puisque la question
est locale sur $Y$, on peut se borner au cas où $Y$ est intègre. Comme $g:X\to S$ est
plat, les points maximaux de $X$ sont contenus dans $X_\eta$ (2.3.4) et puisque
$f_\eta$ est un isomorphisme, $X$ est irréductible et l'anneau local en son point
générique est un corps; de plus, par platitude (3.3.2), on voit que $X$ n'a pas de cycles
premiers immergés, donc $X$ est réduit (3.2.1) et par suite intègre et $f$ est un
morphisme birationnel et séparé. Pour prouver que $f$ est une immersion ouverte, on peut
donc appliquer le critère (21.12.13); pour voir que $f$ est quasi-fini aux points de
$X^{(1)}$, on remarque que l'assertion est évidente en ceux de ces points qui
appartiennent à $X_\eta$; le seul point de $X^{(1)}$ appartenant à $X_s$ est le point
générique $x$ de $X_s$, en vertu de (6.1.1). Or, il résulte de (2.4.6) et de (14.2.2)
que les morphismes $g$ et $h$ sont équidimensionnels; comme $X_\eta$ et $Y_\eta$ sont
isomorphes, les composantes irréductibles de $X_s$ et de $Y_s$ ont toutes même
dimension. Mais par hypothèse $X_s$ est irréductible et on a vu que $f$ est surjectif,
donc il en est de même de $f_s:X_s\to Y_s$, ce qui entraîne que $Y_s$ est aussi
irréductible; si $y$ est le point générique de $Y_s$, on a donc $f(x)=y$, et la relation
$\dim(X_s)=\dim(Y_s)$ entraîne alors, en vertu de (5.6.6), que
$\dim(f_s^{-1}(y))=0$, autrement dit $f_s$ (et par suite aussi $f$) est bien quasi-fini
au point $x$, ce qui achève la démonstration.
\end{proof}

\begin{corollary}[21.12.16]
\label{IV.21.12.16}
Soient $g:X\to S$ un morphisme propre, plat et de présentation finie, $h:Y\to S$ un
morphisme lisse, propre et de présentation finie, $f:X\to Y$ un $S$-morphisme. On
suppose que les fibres $X_s=g^{-1}(s)$ de $g$ sont géométriquement irréductibles. Soit
$U$ l'ensemble des $s\in S$ tels que $f_s:X_s\to Y_s$ soit un isomorphisme. Alors
$U$ est ouvert et fermé dans $S$, et le morphisme $g^{-1}(U)\to h^{-1}(U)$,
restriction de $f$, est un isomorphisme.
\end{corollary}

\begin{proof}
La dernière assertion résultera de la première et de (17.9.5). On sait déjà (9.6.1,
(xi)) que $U$ est localement constructible dans $S$. En vertu de (1.10.1), il suffit de
prouver que $U$ est stable par \emph{spécialisation} et par \emph{générisation}. Pour
démontrer ces assertions on
\oldpage[IV-4]{313}
peut, par un changement de base $\operatorname{Spec}(\mathcal O_{S,s})\to S$, se
ramener au cas où $S$ est un schéma local de point fermé $s$ et de point générique
$\eta$, et il faut montrer que, pour que $f_s$ soit un isomorphisme, il faut et il suffit
que $f_\eta$ le soit. Or, la suffisance de cette condition résulte de (21.12.15). Pour
montrer qu'elle est nécessaire, on raisonne comme dans la démonstration de (21.12.15)
(en remarquant, avec les mêmes notations, que la projection du point fermé de $S$ est le
point fermé de $S_0$) pour se ramener au cas où en outre $S$ est \emph{noethérien};
mais alors la conclusion résulte de $(\mathbf{III}, 4.6.7$, (ii)).
\end{proof}

\begin{remarks}[21.12.17]
\label{IV.21.12.17}
(i) La conclusion de la proposition (21.12.15) n'est plus valable si on élimine
l'hypothèse que les fibres $X_s$ sont irréductibles. Prenons en effet
$S=\operatorname{Spec}(A)$, où $A$ est un anneau de valuation discrète,
$Y=\mathbb P^1_S$, qui est propre et lisse sur $S$ (17.3.9). Notons encore $s$ et
$\eta$ le point fermé et le point générique de $S$; soit $z$ un point fermé de $Y_s$,
par exemple un point rationnel sur $k=k(s)$, et soit $X$ le $Y$-préschéma obtenu en
faisant éclater le point $z$. Comme l'anneau de polynômes $A[T]$ est régulier
$(\mathbf 0, 17.3.7)$ et de dimension $2$, on voit comme dans la démonstration de
(15.1.1.6) que, si $f:X\to Y$ est le morphisme structural, $f^{-1}(z)$ est isomorphe
à $\mathbb P^1_k$, et d'autre part, le complémentaire de $f^{-1}(z)$ dans $X_s$ est
isomorphe à $Y_s-\{z\}$, donc $Z_1=f^{-1}(z)$ et $Z_2$, adhérence dans $X_s$ du
complémentaire de $Z_1$, sont les deux composantes irréductibles de $X_s$. Par
ailleurs, $f$ est évidemment propre et $g=h\circ f$ est plat, car $X$ est intègre
$(\mathbf{II}, 8.1.4)$ et pour tout ouvert affine $U$ de $X$, l'homomorphisme
$A\to\Gamma(U,\mathcal O_X)$ est injectif $(\mathbf I, 1.2.7)$, donc
$\Gamma(U,\mathcal O_X)$ est un $A$-module sans torsion, et par suite plat puisque $A$
est un anneau de valuation discrète $(\mathbf 0_{\mathrm I}, 6.3.4)$. Il est clair que
$f_\eta:X_\eta\to Y_\eta$ est un isomorphisme, bien que
$f_s:X_s\to Y_s$ n'en soit pas un.

(ii) La conclusion de (21.12.15) tombe aussi en défaut si, dans les hypothèses, on
supprime l'hypothèse que $f$ est \emph{propre}. Il suffit pour le voir de définir $S$ et
$Y$ comme dans (i), mais de remplacer $X$ par le complémentaire $X'$ de $Z_2$ dans le
préschéma $X$ défini dans (i), $f$ par la restriction $f':X'\to Y$ de $f$; le morphisme
$g':X'\to S$, restriction de $g$, est encore plat, et cette fois $X'_s$ est
géométriquement irréductible; $X'_s$ est d'ailleurs isomorphe au complémentaire d'un
point fermé dans $\mathbb P^1_k$, donc est un schéma affine isomorphe à $\mathbb A^1_k$;
comme son image dans $Y_s$ est réduite au point fermé $z$, $f'$ n'est pas propre
$(\mathbf{III}, 4.4.2)$ et $f'_s$ n'est pas un isomorphisme, bien que $f'_\eta$ en soit
un.

(iii) Il se peut que l'énoncé de la proposition (21.12.15) reste valable quand on y
remplace le mot « isomorphisme » par « morphisme étale » (cf. (21.12.14, (v))). Il en
sera alors encore de même de (21.12.16).
\end{remarks}
\subsection{Couples parafactoriels. Anneaux locaux parafactoriels}
\EGAReaderTarget{ega4.c04.subsec.21-13}{subsection:IV.21.13}
\label{IV.21.13}

\begin{definition}[21.13.1]
\EGAReaderTarget{ega4.c04.def.21-13-1}{IV.21.13.1}
Soient $X$ un espace annelé, $Y$ une partie fermée de $X$, $U=X-Y$.
On dit que le couple $(X,Y)$ est \emph{parafactoriel} si, pour tout ouvert $V$ de $X$,
le foncteur restriction
$\mathscr L\mapsto\mathscr L|(U\cap V)$, de la catégorie des
$\mathcal O_V$-Modules inversibles dans celle des $\mathcal O_{U\cap V}$-Modules
inversibles, est une \emph{équivalence de catégories}.
\end{definition}

Dire que le couple $(X,Y)$ est parafactoriel signifie donc que, pour tout ouvert $V$
de $X$, les conditions suivantes sont vérifiées :
\begin{enumerate}[label=\arabic*\textsuperscript{o}]
  \item le foncteur $\mathscr L\mapsto\mathscr L|(U\cap V)$ est \emph{pleinement
  fidèle}, autrement dit, pour deux $\mathcal O_V$-Modules inversibles $\mathscr L$,
  $\mathscr L'$, l'application de restriction
  \[
    \operatorname{Hom}_{\mathcal O_V}(\mathscr L,\mathscr L')
    \longrightarrow
    \operatorname{Hom}_{\mathcal O_{U\cap V}}
    (\mathscr L|(U\cap V),\mathscr L'|(U\cap V))
  \]
  est \emph{bijective};
  \item le foncteur $\mathscr L\mapsto\mathscr L|(U\cap V)$ est
  \emph{essentiellement surjectif}, autrement dit, pour tout
  $\mathcal O_{U\cap V}$-Module inversible $\mathscr L_0$, il existe un
  $\mathcal O_V$-Module inversible $\mathscr L$ tel que $\mathscr L_0$ soit
  \oldpage[IV-4]{314}
  isomorphe à $\mathscr L|(U\cap V)$; cela s'exprime encore en disant que
  l'homomorphisme canonique (21.3.2.4)
  \[
    \operatorname{Pic}(V)\longrightarrow\operatorname{Pic}(U\cap V)
  \]
  est \emph{surjectif}.
\end{enumerate}

\begin{lemma}[21.13.2]
\label{IV.21.13.2}
Soit $f:X'\to X$ un morphisme d'espaces annelés; pour tout ouvert $V$ de $X$, on
note $f_V:f^{-1}(V)\to V$ la restriction de $f$. Les conditions suivantes sont
équivalentes :
\begin{enumerate}[label=\emph{\alph*)}]
  \item Pour tout ouvert $V$ de $X$, le foncteur
  $\mathscr E\mapsto f_V^*(\mathscr E)$ de la catégorie des $\mathcal O_V$-Modules
  localement libres de rang fini dans la catégorie des
  $\mathcal O_{f^{-1}(V)}$-Modules localement libres de rang fini, est fidèle
  (resp.\ pleinement fidèle).
  \item[\emph{a')}] Pour tout ouvert $V$ de $X$, le foncteur
  $\mathscr L\mapsto f_V^*(\mathscr L)$ de la catégorie des $\mathcal O_V$-Modules
  localement libres de rang $1$ dans la catégorie des
  $\mathcal O_{f^{-1}(V)}$-Modules localement libres de rang $1$ est fidèle
  (resp.\ pleinement fidèle).
  \item Pour tout ouvert $V$ de $X$, l'homomorphisme canonique
  $(\mathbf 0_{\mathrm I}, 4.4.3.2)$
  $\mathscr E\to(f_V)_*(f_V^*(\mathscr E))$ est un monomorphisme
  (resp.\ un isomorphisme) pour tout $\mathcal O_V$-Module localement libre
  $\mathscr E$ de rang fini.
  \item[\emph{b')}] L'homomorphisme canonique
  $\mathcal O_X\to f_*(\mathcal O_{X'})$ est un monomorphisme
  (resp.\ un isomorphisme).
\end{enumerate}

Supposons que l'homomorphisme canonique $\mathcal O_X\to f_*(\mathcal O_{X'})$ soit
bijectif; alors, pour qu'un $\mathcal O_{X'}$-Module localement libre de rang fini
$\mathscr E'$ soit isomorphe à un $\mathcal O_{X'}$-Module de la forme
$f^*(\mathscr E)$, où $\mathscr E$ est un $\mathcal O_X$-Module localement libre de
rang fini, il faut et il suffit que les deux conditions suivantes soient remplies :
\begin{enumerate}[label=(\roman*)]
  \item $f_*(\mathscr E')$ est un $\mathcal O_X$-Module localement libre;
  \item l'homomorphisme canonique $(\mathbf 0_{\mathrm I}, 4.4.3.3)$
  $f^*(f_*(\mathscr E'))\to\mathscr E'$ est un isomorphisme.
\end{enumerate}
Lorsque ces deux conditions sont remplies, $\mathscr E$ est isomorphe à
$f_*(\mathscr E')$.
\end{lemma}

\begin{proof}
Dire que le foncteur $\mathscr E\mapsto f_V^*(\mathscr E)$ est fidèle
(resp.\ pleinement fidèle) signifie que pour deux $\mathcal O_V$-Modules localement
libres de rang fini $\mathscr E_1$, $\mathscr E_2$, l'application
\[
  \operatorname{Hom}_{\mathcal O_V}(\mathscr E_1,\mathscr E_2)
  \longrightarrow
  \operatorname{Hom}_{\mathcal O_{f^{-1}(V)}}
  (f_V^*(\mathscr E_1),f_V^*(\mathscr E_2))
\]
est injective (resp.\ bijective); comme ceci doit avoir lieu en remplaçant $V$ par un
ouvert $W\subset V$ et $\mathscr E_1$, $\mathscr E_2$ par $\mathscr E_1|W$,
$\mathscr E_2|W$, et que
\[
  \operatorname{Hom}_{\mathcal O_W}(\mathscr E_1|W,\mathscr E_2|W)
  =\Gamma(W,\mathscr Hom_{\mathcal O_V}(\mathscr E_1,\mathscr E_2)),
\]
la condition signifie encore que l'homomorphisme canonique de faisceaux
\[
\label{IV.21.13.2.1}
  \mathscr Hom_{\mathcal O_V}(\mathscr E_1,\mathscr E_2)
  \longrightarrow
  (f_V)_*(f_V^*(\mathscr Hom_{\mathcal O_V}(\mathscr E_1,\mathscr E_2)))
  \tag{21.13.2.1}
\]
est injectif (resp.\ bijectif). Mais puisque $\mathscr E_1$ et $\mathscr E_2$ sont
localement libres de rang fini,
$\mathscr Hom_{\mathcal O_V}(\mathscr E_1,\mathscr E_2)$, isomorphe à
$\mathscr E_1^\vee\otimes_{\mathcal O_V}\mathscr E_2$
$(\mathbf 0_{\mathrm I}, 5.4.2)$, est aussi localement libre de rang fini; cela montre
déjà que $b)$ entraîne $a)$, et inversement $b)$ est un cas particulier de $a)$,
compte tenu de l'isomorphisme
$\mathscr E\xrightarrow{\sim}\mathscr Hom_{\mathcal O_V}(\mathcal O_V,\mathscr E)$.
Il est clair que $a')$ est un cas particulier de $a)$ et $b')$ un cas particulier de
$a')$. Inversement, comme $b)$ est une propriété locale, on peut, pour la vérifier,
se borner au cas où $\mathscr E=\mathcal O_V^n$, et cela montre que $b')$ entraîne
$b)$.

Si l'homomorphisme canonique $\mathcal O_X\to f_*(\mathcal O_{X'})$ est bijectif,
et si les conditions (i) et (ii) sont remplies, il est clair que
$\mathscr E'=f^*(\mathscr E)$ avec $\mathscr E=f_*(\mathscr E')$, à un isomorphisme
près.
\oldpage[IV-4]{315}
Inversement, supposons que $\mathscr E'=f^*(\mathscr E)$, avec $\mathscr E$ localement
libre; la question étant locale sur $X$, on peut supposer que
$\mathscr E=\mathcal O_X^n$, d'où $\mathscr E'=\mathcal O_{X'}^n$, et les conditions
(i) et (ii) résultent de l'hypothèse que
$\mathcal O_X\to f_*(\mathcal O_{X'})$ est un isomorphisme.
\end{proof}

Dans ce numéro, nous appliquerons le lemme précédent à une injection canonique
$j:U\to X$, où $U$ est l'espace annelé induit par $X$ sur un \emph{ouvert} de $X$.
Avec ces notations, (21.13.2) se traduit en :

\begin{corollary}[21.13.3]
\label{IV.21.13.3}
Soient $X$ un espace annelé, $Y$ une partie fermée de $X$, $U=X-Y$,
$j:U\to X$ l'injection canonique. Les conditions suivantes sont équivalentes :
\begin{enumerate}[label=\emph{\alph*)}]
  \item Pour tout ouvert $V$ de $X$, le foncteur restriction
  $\mathscr E\mapsto\mathscr E|(U\cap V)$ de la catégorie des
  $\mathcal O_V$-Modules localement libres de rang fini dans la catégorie des
  $\mathcal O_{U\cap V}$-Modules localement libres de rang fini est fidèle
  (resp.\ pleinement fidèle).
  \item[\emph{a')}] Pour tout ouvert $V$ de $X$, le foncteur restriction
  $\mathscr L\mapsto\mathscr L|(U\cap V)$ de la catégorie des
  $\mathcal O_V$-Modules localement libres de rang $1$ dans la catégorie des
  $\mathcal O_{U\cap V}$-Modules localement libres de rang $1$ est fidèle
  (resp.\ pleinement fidèle).
  \item Pour tout ouvert $V$ de $X$, l'homomorphisme canonique
  $\mathscr E\to(j_V)_*(\mathscr E|(U\cap V))$ est injectif (resp.\ bijectif) pour
  tout $\mathcal O_V$-Module localement libre de rang fini $\mathscr E$.
  \item[\emph{b')}] L'homomorphisme canonique
  $\mathcal O_X\to j_*(\mathcal O_U)$ est injectif (resp.\ bijectif).
\end{enumerate}

Supposons que l'homomorphisme canonique $\mathcal O_X\to j_*(\mathcal O_U)$ soit
bijectif. Alors, pour qu'un $\mathcal O_U$-Module localement libre de rang fini
$\mathscr F$ soit de la forme $\mathscr E|U$, où $\mathscr E$ est un
$\mathcal O_X$-Module localement libre de rang fini, il faut et il suffit que
$j_*(\mathscr F)$ soit un $\mathcal O_X$-Module localement libre, et dans ce cas, on
peut prendre $\mathscr E=j_*(\mathscr F)$.
\end{corollary}

\begin{lemma}[21.13.4]
\label{IV.21.13.4}
Soient $X$ un préschéma localement noethérien, $Y$ une partie fermée de $X$,
$U=X-Y$, $j:U\to X$ l'injection canonique. Pour que l'homomorphisme canonique
$\mathcal O_X\to j_*(\mathcal O_U)$ soit injectif (resp.\ bijectif), il faut et il
suffit que $\operatorname{prof}_Y(\mathcal O_X)\geq1$
(resp.\ $\operatorname{prof}_Y(\mathcal O_X)\geq2$) (autrement dit, que
$\operatorname{prof}(\mathcal O_{X,x})\geq1$
(resp.\ $\operatorname{prof}(\mathcal O_{X,x})\geq2$) pour tout $x\in Y$).
\end{lemma}

\begin{proof}
Ces assertions sont en effet des cas particuliers de (5.10.2) et (5.10.5).
\end{proof}

\begin{proposition}[21.13.5]
\label{IV.21.13.5}
Soient $X$ un espace annelé, $Y$ une partie fermée de $X$, $U=X-Y$,
$j:U\to X$ l'injection canonique. Pour que le couple $(X,Y)$ soit parafactoriel, il
faut et il suffit que l'homomorphisme canonique $\mathcal O_X\to j_*(\mathcal O_U)$
soit bijectif, et que pour tout ouvert $V$ de $X$ et tout
$\mathcal O_{U\cap V}$-Module inversible $\mathscr L$, $(j_V)_*(\mathscr L)$ soit un
$\mathcal O_V$-Module inversible (notation de (21.13.3)).
\end{proposition}

\begin{proof}
C'est une conséquence immédiate de la définition (21.13.1) et de (21.13.3).
\end{proof}

\begin{corollary}[21.13.6]
\label{IV.21.13.6}
Soient $X$ un espace annelé, $Y$ une partie fermée de $X$.
\begin{enumerate}[label=(\roman*)]
  \item Si le couple $(X,Y)$ est parafactoriel, il en est de même de
  $(W,Y\cap W)$ pour tout ouvert $W$ de $X$. Inversement, si $(W_\alpha)$ est un
  recouvrement ouvert de $X$ tel que chacun des couples
  $(W_\alpha,Y\cap W_\alpha)$ soit parafactoriel, alors le couple $(X,Y)$ est
  parafactoriel.
  \item Si le couple $(X,Y)$ est parafactoriel, il en est de même de $(X,Y')$ pour
  toute partie fermée $Y'$ de $Y$.
  \item Supposons que $X$ soit un préschéma, et soient $X'$ un préschéma,
  $f:X'\to X$ un morphisme fidèlement plat et quasi-compact; posons
  $Y'=f^{-1}(Y)$. Supposons que le couple $(X',Y')$ soit parafactoriel et que
  l'ouvert $U=X-Y$ soit rétrocompact dans $X$ $(\mathbf 0_{\mathrm{III}}, 9.1.1)$.
  Alors le couple $(X,Y)$ est parafactoriel.
\end{enumerate}
\end{corollary}

\begin{proof}
\oldpage[IV-4]{316}
(i) Comme le fait que l'homomorphisme canonique
$\mathcal O_X\to j_*(\mathcal O_U)$ soit bijectif est une propriété locale sur $X$,
tout revient à voir (en vertu de (21.13.5)), en désignant par
$j_\alpha:U\cap W_\alpha\to W_\alpha$ l'injection canonique, que l'on a la propriété
suivante : si, pour tout $\alpha$ et pour tout $\mathcal O_U$-Module inversible
$\mathscr L$, $(j_\alpha)_*(\mathscr L|(U\cap W_\alpha))$ est un
$\mathcal O_{W_\alpha}$-Module inversible, alors $j_*(\mathscr L)$ est un
$\mathcal O_X$-Module inversible. Mais la propriété d'être un $\mathcal O_X$-Module
inversible est locale sur $X$, et $(W_\alpha)$ est un recouvrement ouvert de $X$;
comme on a
\[
  j_*(\mathscr L)|W_\alpha
  =(j_\alpha)_*(\mathscr L|(U\cap W_\alpha)),
\]
cela prouve notre assertion.

(ii) Posons $U'=X-Y'$, et soit $j':U'\to X$ l'injection canonique. Dire que
l'homomorphisme $\mathcal O_X\to j'_*(\mathcal O_{U'})$ est bijectif revient à dire
que pour tout ouvert $V$ de $X$, l'homomorphisme
$\Gamma(V,\mathcal O_X)\to\Gamma(V\cap U',\mathcal O_X)$ est bijectif; mais
l'homomorphisme composé
\[
  \Gamma(V,\mathcal O_X)\longrightarrow\Gamma(V\cap U',\mathcal O_X)
  \longrightarrow\Gamma(V\cap U,\mathcal O_X)
\]
est bijectif par hypothèse (21.13.5), et en remplaçant $V$ par $V\cap U'$, on voit que
$\Gamma(V\cap U',\mathcal O_X)\to\Gamma(V\cap U,\mathcal O_X)$ est bijectif, donc
$\Gamma(V,\mathcal O_X)\to\Gamma(V\cap U',\mathcal O_X)$ est bijectif. Notons ensuite
que si $j'':U\to U'$ est l'injection canonique et si $\mathscr L'$ est un
$\mathcal O_{U'}$-Module inversible, $\mathscr L'|U$ est un $\mathcal O_U$-Module
inversible, donc
$j_*(\mathscr L'|U)=j'_*(j''_*(\mathscr L'|U))$ est par hypothèse un
$\mathcal O_X$-Module inversible. Comme le couple $(U',U'\cap Y)$ est parafactoriel
en vertu de (i), on a $j''_*(\mathscr L'|U)\cong\mathscr L'$, donc
$j'_*(\mathscr L')$ est un $\mathcal O_X$-Module inversible. Il suffit alors pour
conclure de remplacer dans le raisonnement précédent $X$, $U$ et $U'$ par $V$,
$V\cap U$ et $V\cap U'$, où $V$ est un ouvert de $X$.

(iii) Posons $U'=X'-Y'=f^{-1}(U)$ et notons que puisqu'il s'agit de préschémas, on
peut écrire $U'=U\times_X X'$; soit $f_U:f^{-1}(U)\to U$ la restriction de $f$ et
soit $j':U'\to X'$ l'injection canonique, qui s'écrit aussi $j_{(X')}$. Montrons
d'abord que l'homomorphisme canonique
$\rho:\mathcal O_X\to j_*(\mathcal O_U)$ est bijectif; pour cela, notons que par
hypothèse l'homomorphisme canonique
$\mathcal O_{X'}\to j'_*(\mathcal O_{U'})$ est bijectif; mais puisque le morphisme
$j$ est quasi-compact et séparé par hypothèse, et le morphisme $f$ plat,
$j'_*(\mathcal O_{U'})=j'_*(f_U^*(\mathcal O_U))$ s'identifie canoniquement à
$f^*(j_*(\mathcal O_U))$ en vertu de (2.3.1); comme
$\mathcal O_{X'}=f^*(\mathcal O_X)$, on voit que l'homomorphisme
$f^*(\rho):f^*(\mathcal O_X)\to f^*(j_*(\mathcal O_U))$ est bijectif; on en conclut
que $\rho$ est bijectif puisque $f$ est fidèlement plat (2.2.7). Considérons ensuite
un $\mathcal O_U$-Module inversible $\mathscr L$, et soit
$\mathscr L'=f_U^*(\mathscr L)$, qui est un $\mathcal O_{U'}$-Module inversible. Par
hypothèse $j'_*(f_U^*(\mathscr L))$ est un $\mathcal O_{X'}$-Module inversible, qui
est isomorphe à $f^*(j_*(\mathscr L))$ par (2.3.1) comme ci-dessus. Mais puisque $f$
est fidèlement plat et quasi-compact, cela entraîne que $j_*(\mathscr L)$ est un
$\mathcal O_X$-Module localement libre (2.5.2). Pour achever la démonstration, il
suffit de remplacer $X$ par un ouvert $V$ et $X'$ par $f^{-1}(V)$ dans ce qui
précède.
\end{proof}

\begin{definition}[21.13.7]
\EGAReaderTarget{ega4.c04.def.21-13-7}{IV.21.13.7}
On dit qu'un anneau local $A$ est \emph{parafactoriel} si, posant
$X=\operatorname{Spec}(A)$ et désignant par $a$ le point fermé de $X$, le couple
$(X,\{a\})$ est parafactoriel.
\end{definition}

\begin{proposition}[21.13.8]
\label{IV.21.13.8}
Les notations étant celles de (21.13.7), on pose $U=X-\{a\}$. Pour que $A$ soit
parafactoriel, il faut et il suffit qu'il vérifie les conditions suivantes :
\begin{enumerate}[label=(\roman*)]
  \item L'homomorphisme canonique
  $A=\Gamma(X,\mathcal O_X)\to\Gamma(U,\mathcal O_X)$ est bijectif.
  \item $\operatorname{Pic}(U)=0$.
\end{enumerate}
\oldpage[IV-4]{317}
Si en outre $A$ est noethérien, la condition (i) équivaut à
\begin{enumerate}
  \item[\textup{(i bis)}] $\operatorname{prof}(A)\geq2$.
\end{enumerate}
\end{proposition}

\begin{proof}
En effet, le seul ouvert de $X$ contenant $a$ est $X$ lui-même, et par suite tout
$\mathcal O_X$-Module inversible est isomorphe à $\mathcal O_X$, autrement dit
$\operatorname{Pic}(X)=0$; la première assertion résulte donc de la définition
(21.13.1) et de (21.13.3). La seconde assertion n'est autre qu'un cas particulier de
(21.13.4).
\end{proof}

\begin{env}[21.13.9]
\label{IV.21.13.9}
\emph{Exemples}. ---

(i) Un anneau local noethérien parafactoriel est nécessairement de dimension
$\geq2$ en vertu de (21.13.8); en d'autres termes un anneau local noethérien de
dimension $\leq1$ n'est pas parafactoriel.

(ii) Un anneau local noethérien \emph{factoriel} $A$ de dimension $\geq2$ est
parafactoriel, comme il résulte de (21.13.8) et de (21.6.14).

(iii) Si $A$ est un anneau local noethérien de dimension $\geq3$ et parafactoriel, il
n'est pas nécessairement normal ni même réduit. Considérons un anneau local régulier
$B$ de dimension $\geq3$, et soit $A=D_B(B)$ $(\mathbf 0, 18.2.3)$, isomorphe à
$B[T]/(T^2)$; on voit aussitôt que l'on a
$\operatorname{prof}(A)=\operatorname{prof}(B)=\dim(B)\geq3$. Pour montrer que $A$
est parafactoriel, il suffit donc, avec les notations de (21.13.8), de prouver que
$\operatorname{Pic}(U)=0$. Soit $\mathfrak J$ le noyau de l'augmentation $A\to B$,
qui est tel que $\mathfrak J^2=0$ et qui, en tant que $B$-module, est isomorphe à $B$.
Comme $B=A/\mathfrak J$, $X=\operatorname{Spec}(A)$ et
$X_0=\operatorname{Spec}(B)$ ont même espace sous-jacent; si
$\mathscr J=\widetilde{\mathfrak J}$, $X_0$ est le sous-schéma défini par
$\mathscr J$; nous noterons $U_0$ le sous-préschéma induit par $X_0$ sur l'ouvert
$U$. Pour tout $z\in\mathfrak J$, posons $\varphi(z)=1+z$; puisque
$(1+z)(1-z)=1$, $1+z$ est inversible dans $A$ et $\varphi(\mathfrak J)$ est le noyau
de l'homomorphisme surjectif canonique de groupes multiplicatifs $A^*\to B^*$; en
d'autres termes, on a une suite exacte de groupes commutatifs
\[
  0\longrightarrow\mathfrak J\xrightarrow{\varphi}A^*\longrightarrow B^*
  \longrightarrow1
\]
(les trois derniers groupes étant multiplicatifs, les deux premiers additifs). Pour la
même raison, pour tout $t\in A$, on a la suite exacte
\[
  0\longrightarrow\mathfrak J_t\xrightarrow{\varphi_t}(A_t)^*
  \longrightarrow(B_{t_0})^*\longrightarrow1
\]
en désignant par $t_0$ l'image de $t$ dans $B$, puisque
$B_{t_0}=A_t/\mathfrak J_t$; autrement dit, on a une suite exacte de faisceaux de
groupes commutatifs sur l'espace topologique $X$
\[
  0\longrightarrow\mathscr J\xrightarrow{u}\mathcal O_X^*
  \longrightarrow\mathcal O_{X_0}^*\longrightarrow1
\]
d'où par restriction à l'ouvert $U$, une suite exacte
\[
  0\longrightarrow\mathscr J|U\longrightarrow\mathcal O_U^*
  \longrightarrow\mathcal O_{U_0}^*\longrightarrow1.
\]
Par la suite exacte de cohomologie, on en déduit une suite exacte
\[
\label{IV.21.13.9.1}
  \mathrm H^1(U,\mathscr J)\longrightarrow\mathrm H^1(U,\mathcal O_U^*)
  \xrightarrow{u_1}\mathrm H^1(U,\mathcal O_{U_0}^*)
  \tag{21.13.9.1}
\]
Mais puisque $\mathfrak J$ est un $B$-module isomorphe à $B$, on a
$\mathrm H^1(U,\mathscr J)=\mathrm H^1(U_0,\mathcal O_{U_0})$, et il résulte du
chap.~III, 3\textsuperscript{e} partie (voir aussi [41, III, Exemple III-1]) que l'on a
\oldpage[IV-4]{318}
$\mathrm H^1(U_0,\mathcal O_{U_0})=0$ en raison de la relation
$\operatorname{prof}(B)\geq3$. Par ailleurs, comme $B$ est factoriel, on a
$\mathrm H^1(U,\mathcal O_{U_0}^*)=\operatorname{Pic}(U_0)=0$; on tire donc bien de
la suite exacte précédente que
$\operatorname{Pic}(U)=\mathrm H^1(U,\mathcal O_U^*)=0$.

(iv) Il existe aussi des anneaux locaux noethériens de dimension $3$ qui sont
intègres, intégralement clos et parafactoriels, mais non factoriels. Soit en effet $B$
un anneau local noethérien, complet, intègre et intégralement clos, de dimension
$\geq2$ et non factoriel (par exemple le complété de l'anneau localisé de l'algèbre
intègre $k[U,V,W]/(W^2-UV)$ en l'idéal image de $(U)+(V)+(W)$). Alors il résulte de
ce que l'on verra plus bas (21.14.2) que l'anneau local $A=B[[T]]$ des séries formelles
en $T$ sur $B$ est parafactoriel, mais il n'est pas factoriel, sans quoi $B$ le serait
en vertu de (21.13.12) ci-dessous.

(v) On peut montrer qu'un anneau local d'intersection complète absolue (19.3.1) de
dimension $\geq4$, est parafactoriel (cf. [41, XI, 3.13 (i)]).

(vi) On a vu (Remarque (ii)) que tout anneau local noethérien factoriel $A$ de
dimension $2$ est parafactoriel. Mais il y a des anneaux locaux noethériens de
dimension $2$ qui sont parafactoriels et non factoriels. On peut montrer que, pour
qu'un anneau local noethérien $A$ de dimension $2$ soit parafactoriel et non
factoriel, il faut et il suffit qu'il satisfasse aux trois conditions suivantes :
\begin{enumerate}[label=\arabic*\textsuperscript{o}]
  \item $A$ est un anneau de Cohen--Macaulay (autrement dit
  $\operatorname{prof}(A)=2$).
  \item $A$ est intègre et si $A'$ est sa clôture intégrale, $A'$ est factoriel et est
  une $A$-algèbre finie.
  \item Soit $\mathfrak J$ le conducteur de $A$ dans $A'$ (annulateur du $A$-module
  $A'/A$, ou encore le plus grand idéal de $A'$ contenu dans $A$); posons
  $B=A/\mathfrak J$, $B'=A'/\mathfrak J$; alors $\dim(B)=1$ (ce qui implique
  $A'\neq A$, autrement dit $A$ n'est pas intégralement clos), et l'application
  canonique $\operatorname{Div}(B)\to\operatorname{Div}(B')$ (21.4.5) est
  \emph{surjective}.
\end{enumerate}
On peut montrer en outre que ces conditions entraînent la propriété suivante :

$4^{\mathrm o}$ L'anneau $B'$ (et \emph{a fortiori} $B\subset B'$) est réduit, et le
morphisme $\operatorname{Spec}(B')\to\operatorname{Spec}(B)$ est bijectif.

Si l'on pose $X=\operatorname{Spec}(A)$, $X'=\operatorname{Spec}(A')$,
$Y=\operatorname{Spec}(B)$, $Y'=\operatorname{Spec}(B')$, de sorte que $Y$
(resp.\ $Y'$) est défini par l'idéal $\mathfrak J$ de $A$ (resp.\ $A'$), le morphisme
structural $f:X'\to X$ est un isomorphisme de $X'-Y'$ sur $X-Y$ (Bourbaki,
\emph{Alg. comm.}, chap.~V, \S1, n\textsuperscript{o}~5, cor.~5 de la prop.~16). On
voit donc en vertu de $4^{\mathrm o}$ que $f$ est un morphisme \emph{bijectif}; en
d'autres termes, $X$ est un préschéma \emph{unibranche} (6.15.1) (et en particulier
$A'$ est un anneau local noethérien); en général $X$ n'est pas géométriquement
unibranche. L'espace $Y$, de dimension $1$, est constitué par le point fermé $x$ de
$X$ et les points maximaux $y_i$ ($1\leq i\leq r$) de $Y$, et l'unique point $y_i'$
de $Y'$ au-dessus de $y_i$ est aussi un point maximal de $Y'$; comme $B$ et $B'$ sont
réduits, on en déduit que l'on a
\[
  \mathfrak m_{y_i}\mathcal O_{X',y_i'}=\mathfrak m_{y_i'}.
\]
Pour $z=y_i$ et $z'=y_i'$, cette relation est aussi vérifiée de façon évidente pour
$z\notin Y$ et $z'$ l'unique point de $f^{-1}(z)$, donc pour tout
$z\in U=X-\{x\}$ et $z'$ l'unique point de $f^{-1}(z)$. En particulier, si l'anneau
$A$ est de caractéristique $0$ (0, 21.1.1) on voit que $U'$ est \emph{non ramifié} sur
$U$ (mais non étale en général).

Nous nous bornerons ici à démontrer que les conditions $1^{\mathrm o}$,
$2^{\mathrm o}$, $3^{\mathrm o}$ ci-dessus sont \emph{suffisantes} pour que $A$ soit
parafactoriel. Or, en vertu de la condition $1^{\mathrm o}$ et de (21.13.8), il suffit
de démontrer que $\operatorname{Pic}(U)=0$. Puisque $A'$ est factoriel en vertu de
$2^{\mathrm o}$, on a $\operatorname{Pic}(U')=0$, et on déduit de (21.8.5, (ii)) une
suite exacte
\[
  1\longrightarrow
  \left(\prod_{i=1}^r
  \mathcal O_{U',y_i'}^*/\mathcal O_{U,y_i}^*\right)
  \big/\operatorname{Im}\Gamma(U',\mathcal O_{U'}^*)
  \longrightarrow\operatorname{Pic}(U)\longrightarrow\operatorname{Pic}(U')
  \longrightarrow0
\]
et il s'agit donc de montrer que le second terme de cette suite est réduit à l'élément
neutre, en prenant pour $S$ l'ensemble des $y_i$ ($1\leq i\leq r$) et en notant
qu'ici $f_*(\mathcal O_{X'})$ est la $\mathcal O_X$-Algèbre $\widetilde{A'}$. Soient
$\mathfrak p_i$, $\mathfrak p_i'$ les idéaux premiers de $A$ et $A'$ correspondant
aux points $y_i$, $y_i'$, et remarquons que $\mathfrak p_i'\cap A=\mathfrak p_i$;
donnons-nous pour chaque $i$ un élément inversible $a_i'/s_i$ de
$A'_{\mathfrak p_i'}$, avec $a_i'\in A'$, $s_i\notin\mathfrak p_i'$; comme nous
n'avons à examiner que les groupes quotients
$(A'_{\mathfrak p_i'})^*/(A_{\mathfrak p_i})^*$, on peut supposer $s_i=1$ pour tout
$i$; soit $b_i'\in B'=A'/\mathfrak J$ l'image canonique de $a_i'$, de sorte que
$b_i'/1$ est un élément non nul de $k(y_i')=\mathcal O_{Y',y_i'}$. L'ensemble
$U\cap Y$ formé par les $y_i$ est un ouvert affine de la forme $D(t)$ dans $Y$, avec
$t\in\mathfrak m$, idéal maximal de $A$; il existe donc un élément inversible
$b'\in B'_t$ dont les $b_i'/1$ sont les images canoniques, et comme $t$ est régulier
dans l'anneau intègre $A'$, on peut supposer $b'$ de la forme $b''/t^m$, où
$b''\in B'$ est inversible. Soit $a''$ un élément de $A'$ dans la classe $b''$, qui
est donc nécessairement inversible; $a'=a''/t^m$ est inversible dans $A'_t$ et pour
tout $i$, on a
\[
  u_i=a''-a_i't^m\in\mathfrak J,
  \qquad
  t^m a_i'=a''-u_i=a''(1-a''^{-1}u_i);
\]
mais $a''^{-1}u_i\in\mathfrak J\subset\mathfrak m$, donc $1-a''^{-1}u_i$ est un
élément inversible de $A$, et les classes de $a'$ et $a_i'$ dans
$(A'_{\mathfrak p_i'})^*/(A_{\mathfrak p_i})^*$ sont les mêmes, ce qui achève la
démonstration.

\oldpage[IV-4]{319}
Pour avoir un exemple explicite d'anneau parafactoriel de dimension $2$ obtenu de
cette manière et \emph{non factoriel}, considérons l'anneau
\[
  E=\mathbb R[[U,V]]/(U^2+V^2),
\]
dont la clôture intégrale $E'$ s'identifie à $\mathbb C[[U]]$ (6.15.11). Si l'on pose
$A=E[[T]]$, la clôture intégrale de $A$ est l'anneau $A'=E'[[T]]$ (Bourbaki,
\emph{Alg. comm.}, chap.~V, \S1, n\textsuperscript{o}~4, prop.~14); on vérifie
aussitôt que le conducteur de $E$ dans $E'$ est l'idéal maximal $\mathfrak n$ de
$E$, donc le conducteur $\mathfrak J$ de $A$ dans $A'$ est $\mathfrak n[[T]]$, et on
a $B=A/\mathfrak J\cong\mathbb R[[T]]$, $B'=A'/\mathfrak J\cong\mathbb C[[T]]$.
Il est alors immédiat que les conditions $1^{\mathrm o}$, $2^{\mathrm o}$ et
$3^{\mathrm o}$ énoncées ci-dessus sont bien vérifiées, mais $A$ n'est même pas
intégralement clos.

On peut varier cet exemple, et le lecteur verra sans peine que si $k$ est un corps
algébriquement clos, l'anneau $A$, localisé de l'anneau
$k[U,V,W]/(U^2-WV^2)$ en l'idéal maximal engendré par les images de $U$, $V$ et $W$,
vérifie aussi les conditions $1^{\mathrm o}$, $2^{\mathrm o}$ et $3^{\mathrm o}$
ci-dessus.

Il se pourrait que ces trois conditions impliquent même que chacun des anneaux
$B'/\mathfrak p_i'$ soit un anneau de valuation discrète, ce qui entraînerait que
$A'$ est même un anneau \emph{régulier}.
\end{env}

\begin{proposition}[21.13.10]
\label{IV.21.13.10}
Soient $X$ un préschéma localement noethérien, $Y$ une partie fermée de $X$. Pour que
le couple $(X,Y)$ soit parafactoriel il faut et il suffit que, pour tout $y\in Y$,
l'anneau local $\mathcal O_{X,y}$ soit parafactoriel.
\end{proposition}

\begin{proof}
Posons $U=X-Y$ et soit $j:U\to X$ l'injection canonique.

En vertu de (21.13.4), dire que l'homomorphisme canonique
$\mathcal O_X\to j_*(\mathcal O_U)$ soit bijectif équivaut à dire que chacun des
anneaux locaux $\mathcal O_{X,y}$ pour $y\in Y$ satisfait à la condition (i bis) de
(21.13.8).

Montrons d'abord que si le couple $(X,Y)$ est parafactoriel, chacun des anneaux
$\mathcal O_{X,y}$ ($y\in Y$) est parafactoriel. Compte tenu de la remarque
précédente, tout revient à voir que, si l'on pose
$T_y=\operatorname{Spec}(\mathcal O_{X,y})$ et $U_y=T_y-\{y\}$, tout
$\mathcal O_{U_y}$-Module inversible $\mathscr L_0$ est isomorphe à
$\mathcal O_{U_y}$. Or, lorsque $V$ parcourt l'ensemble des voisinages ouverts
affines de $y$ dans $X$, il résulte de (8.2.13) et $(\mathbf I, 2.4.2)$ que le
préschéma $U_y$ est limite projective des préschémas induits par $X$ sur les ouverts
$V-(V\cap\overline{\{y\}})$, qui sont quasi-compacts puisque $X$ est localement
noethérien. Comme les $V-(V\cap\overline{\{y\}})$ sont séparés, il résulte alors de
(8.5.2, (ii)) et (8.5.5) qu'il y a un voisinage ouvert affine $V$ de $y$ dans $X$ et
un $\mathcal O_{U_V}$-Module inversible $\mathscr L'_V$ tel que $\mathscr L_0$ soit le
faisceau induit par $\mathscr L'_V$ sur $U_y$. Mais l'hypothèse entraîne, en vertu de
(21.13.6, (i) et (iii)), que le couple $(V,V\cap\overline{\{y\}})$ est parafactoriel;
on en conclut par définition (21.13.1) qu'il existe un $\mathcal O_V$-Module inversible
$\mathscr L_V$ induisant $\mathscr L'_V$ sur $U_V$. Remplaçant au besoin $V$ par un
voisinage ouvert plus petit de $y$, on peut supposer que $\mathscr L_V$ est isomorphe
à $\mathcal O_V$, d'où on conclut que $\mathscr L_0$ est isomorphe à
$\mathcal O_{U_y}$.

Inversement, prouvons que si tous les $\mathcal O_{X,y}$ ($y\in Y$) sont
parafactoriels le couple $(X,Y)$ est parafactoriel. On est évidemment ramené, compte
tenu de la remarque du début, à montrer que pour tout $\mathcal O_U$-Module inversible
$\mathscr L$, $j_*(\mathscr L)$ est un $\mathcal O_X$-Module inversible (21.13.5).
La question étant locale sur $X$, on peut supposer $X$ noethérien. L'ensemble $V$ des
$x\in X$ tels que la restriction de $j_*(\mathscr L)$ à un voisinage ouvert de $x$
dans $X$ soit inversible est évidemment ouvert dans $X$; il s'agit de montrer que
$V=X$. Pour cela, supposons le contraire et posons $Z=X-V\neq\varnothing$. Soit $y$
un point \emph{maximal} de $Z$; puisque $Z\subset Y$ par définition,
$\mathcal O_{X,y}$ est par hypothèse un anneau parafactoriel. Remplaçant au besoin
$X$ par un voisinage ouvert de $y$ dans $X$, on peut supposer que la restriction de
$j_*(\mathscr L)$
\oldpage[IV-4]{320}
à $V-(V\cap\overline{\{y\}})$ est inversible; donc, avec les notations introduites
ci-dessus, le $\mathcal O_{U_y}$-Module $\mathscr L_0$ induit par $j_*(\mathscr L)$
sur $U_y$ est inversible. Puisque $\mathcal O_{X,y}$ est parafactoriel,
$\mathscr L_0$ est induit par un $\mathcal O_{T_y}$-Module inversible $\mathscr L'$;
mais comme $T_y$ est limite projective des voisinages ouverts $W$ de $y$ dans $X$,
on peut de nouveau appliquer (8.5.2, (ii)) et (8.5.5), prouvant l'existence d'un tel
voisinage $W$ et d'un $\mathcal O_W$-Module inversible $\mathscr L''$ induisant
$\mathscr L'$ sur $T_y$, donc induisant $\mathscr L_0$ sur $U_y$; appliquant cette
fois (8.5.2.5) et (8.5.2, (i)), on en déduit qu'en remplaçant au besoin $W$ par un
voisinage ouvert plus petit de $y$, on peut supposer que les restrictions de
$j_*(\mathscr L)$ et de $\mathscr L''$ à
$W-(W\cap\overline{\{y\}})$ soient égales. Si l'on pose alors $V_1=V\cup W$ et si
l'on désigne par $\mathscr L_1$ le $\mathcal O_{V_1}$-Module inversible égal à
$\mathscr L''$ dans $W$, à $j_*(\mathscr L)|V$ dans $V$, on a $U\subset V_1$ et
$\mathscr L_1|U=\mathscr L$. On conclut alors du fait que l'homomorphisme
$\mathcal O_X\to j_*(\mathcal O_U)$ est bijectif et de (21.13.2, $b$)) que
$j_*(\mathscr L)|V_1$ est isomorphe à $\mathscr L_1$, donc inversible. Mais cela
contredit la définition de $V$, et conclut donc la démonstration de (21.13.10).
\end{proof}

\begin{corollary}[21.13.11]
\label{IV.21.13.11}
Soient $X$ un préschéma localement noethérien, $Y$ une partie fermée de $X$,
$U=X-Y$. Supposons que le couple $(X,Y)$ soit parafactoriel et que $U$ soit
localement factoriel; en d'autres termes (21.13.10), pour tout $x\in X$, l'anneau
$\mathcal O_{X,x}$ est : $1^{\mathrm o}$ parafactoriel si $x\in Y$;
$2^{\mathrm o}$ factoriel si $x\notin Y$. Alors $X$ est localement factoriel
(autrement dit, $\mathcal O_{X,x}$ est en fait factoriel pour tout $x\in X$).
\end{corollary}

\begin{proof}
Supposons en effet le contraire, et soit $y$ un point de $Y$ tel que
$\mathcal O_{X,y}$ ne soit pas factoriel et ait une dimension \emph{minima} parmi
tous les points de $Y$ ayant cette propriété. Alors, si l'on pose
$T_y=\operatorname{Spec}(\mathcal O_{X,y})$, $U_y=T_y-\{y\}$, l'hypothèse que
$\mathcal O_{X,y}$ est parafactoriel entraîne $\operatorname{Pic}(U_y)=0$ et
$\operatorname{prof}(\mathcal O_{X,y})\geq2$ (21.13.8); de plus le choix de $y$
entraîne que $U_y$ est localement factoriel. Mais alors (21.6.14) prouve que
$\mathcal O_{X,y}$ est factoriel, contrairement à l'hypothèse.
\end{proof}

\begin{proposition}[21.13.12]
\label{IV.21.13.12}
Soient $A$, $B$ deux anneaux locaux noethériens,
$\rho:A\to B$ un homomorphisme local faisant de $B$ un $A$-module plat. Alors, si
$B$ est factoriel, il en est de même de $A$.
\end{proposition}

\begin{proof}
On sait déjà que $A$ est intègre et intégralement clos (6.5.4), donc on peut se borner
au cas où $\dim(A)\geq2$, et raisonner par récurrence sur $n=\dim(A)$. Soient
$X=\operatorname{Spec}(A)$, $a$ le point fermé de $X$,
$X'=\operatorname{Spec}(B)$, $f:X'\to X$ le morphisme structural,
$U=X-\{a\}$, $U'=f^{-1}(U)$.

Les anneaux locaux $\mathcal O_{X',x'}$ aux points $x'\in f^{-1}(a)$ sont factoriels
et de dimension $\geq2$ (6.1.2), donc parafactoriels (21.13.9, (ii)); on en conclut
(21.13.10) que le couple $(X',f^{-1}(a))$ est parafactoriel. En vertu de (21.13.6,
(iii)), cela montre que l'anneau $A$ est parafactoriel, donc
$\operatorname{prof}(A)\geq2$ et $\operatorname{Pic}(U)=0$ (21.13.8). Enfin, les
anneaux locaux de $U$ étant de dimension $<n$, l'hypothèse de récurrence prouve que
$U$ est \emph{localement factoriel}; la conclusion résulte alors de (21.6.14).
\end{proof}

\begin{env}[21.13.12.1]
\label{IV.21.13.12.1}
On notera que l'énoncé (21.13.12) où on remplace « factoriel » par « parafactoriel »
n'est plus exact, comme le montre l'exemple où $A=k$ est un corps et $B$ l'anneau
local factoriel $k[[T_1,T_2]]$. Cependant, il résulte de (21.13.6, (iii)) que, sous
les conditions de (21.13.12), si $\mathfrak m$ désigne l'idéal maximal de $A$ et si
$\mathfrak mB$ est un \emph{idéal}
\oldpage[IV-4]{321}
\emph{de définition} de $B$ (autrement dit, si $\dim(B/\mathfrak mB)=0$), alors,
lorsque $B$ est parafactoriel, il en est de même de $A$. En particulier, si le
complété $\widehat A$ d'un anneau local noethérien est parafactoriel, il en est de
même de $A$.
\end{env}

\begin{remark}[21.13.13]
\label{IV.21.13.13}
Une partie des définitions et résultats précédents se rattache à des considérations
plus générales sur la cohomologie des faisceaux de groupes commutatifs sur un espace
topologique, développées au chap.~III, 3\textsuperscript{e} partie, et nous n'en avons
donné ici un exposé indépendant que pour la commodité du lecteur. En effet, si $X$
est un espace annelé, on a une \emph{équivalence} canonique entre la catégorie des
$\mathcal O_X$-Modules inversibles et celle des \emph{faisceaux principaux homogènes
sous le faisceau de groupes} $\mathcal O_X^*$ (16.5.15). Cette équivalence se définit
en faisant correspondre à tout $\mathcal O_X$-Module inversible $\mathscr L$ le
faisceau $\mathscr L^*=\mathscr Isom_{\mathcal O_X}(\mathcal O_X,\mathscr L)$ des
germes d'isomorphismes de $\mathcal O_X$ sur $\mathscr L$; il est immédiat que
$\mathscr L^*$ est canoniquement muni d'une structure de faisceau principal homogène
sous $\mathcal O_X^*\cong\mathscr Isom_{\mathcal O_X}(\mathcal O_X,\mathcal O_X)$.
La vérification du fait que le foncteur $\mathscr L\mapsto\mathscr L^*$ est une
équivalence de catégories est immédiate.

Cela étant, considérons de façon générale un espace topologique $X$ et un faisceau de
groupes $\mathscr G$ sur $X$; soit $Y$ une partie fermée de $X$, posons $U=X-Y$, et
soit $i$ un entier tel que $0\leq i\leq2$; on dira que le couple $(X,Y)$ est
\emph{$i$-pur pour $\mathscr G$} si, pour tout ouvert $V$ de $X$, le foncteur
restriction $\mathscr P\mapsto\mathscr P|(U\cap V)$, de la catégorie de faisceaux
principaux homogènes sous le faisceau de groupes $\mathscr G|V$, dans la catégorie
analogue sous le faisceau de groupes $\mathscr G|(U\cap V)$, est \emph{fidèle}
($i=0$), resp.\ \emph{pleinement fidèle} ($i=1$), resp.\ une \emph{équivalence de
catégories} ($i=2$). Dans le cas où $X$ est un espace annelé et
$\mathscr G=\mathcal O_X^*$, dire que $(X,Y)$ est $i$-pur signifie donc pour $i=0$,
que l'homomorphisme $\mathcal O_X\to j_*(\mathcal O_U)$ est \emph{injectif}; pour
$i=1$, que cet homomorphisme est \emph{bijectif}; et enfin, pour $i=2$, que le couple
$(X,Y)$ est \emph{parafactoriel} (21.13.3).

Revenant au cas général, rappelons que les morphismes de la catégorie des faisceaux
principaux sous $\mathscr G$ sont par définition les isomorphismes. Dire que le couple
$(X,Y)$ est $0$-pur signifie donc que pour tout ouvert $V$ de $X$, l'homomorphisme
canonique $\mathrm H^0(V,\mathscr G)\to\mathrm H^0(U\cap V,\mathscr G)$ est injectif;
dire que le couple $(X,Y)$ est $1$-pur signifie que l'homomorphisme canonique
$\mathrm H^0(V,\mathscr G)\to\mathrm H^0(U\cap V,\mathscr G)$ est bijectif (et alors
l'homomorphisme canonique
$\mathrm H^1(V,\mathscr G)\to\mathrm H^1(U\cap V,\mathscr G)$ est injectif); enfin,
on peut montrer que pour que le couple $(X,Y)$ soit $2$-pur, il faut et il suffit que
les homomorphismes
$\mathrm H^i(V,\mathscr G)\to\mathrm H^i(U\cap V,\mathscr G)$ soient bijectifs pour
$i=0$ et $i=1$. Lorsque $\mathscr G$ est un faisceau de groupes commutatifs, en
introduisant les faisceaux de cohomologie $\mathscr H_Y^p(\mathscr G)$ définis au
chap.~III, 3\textsuperscript{e} partie, dire que $(X,Y)$ est $i$-pur pour $i\leq2$
signifie que $\mathscr H_Y^p(\mathscr G)=0$ pour $p\leq i$; sous cette forme, la
notion se généralise donc aussitôt pour tout entier $i$.
\end{remark}

\begin{proposition}[21.13.14]
\label{IV.21.13.14}
Soient $X$ un préschéma localement noethérien \emph{normal},
$(U_\lambda)_{\lambda\in L}$ une famille filtrante décroissante d'ouverts de $X$. Les
conditions suivantes sont équivalentes :
\begin{enumerate}[label=\emph{\alph*)}]
  \item Tout cycle $1$-codimensionnel $Z$ sur $X$ tel qu'il existe un indice
  $\lambda$ pour lequel $Z|U_\lambda$ soit localement principal (21.6.7), est
  localement principal.
  \item Pour tout $x\in X$ tel que $\dim(\mathcal O_{X,x})\geq2$ et que $x$
  n'appartienne pas à $\bigcap_{\lambda\in L}U_\lambda$, l'anneau
  $\mathcal O_{X,x}$ est parafactoriel.
\end{enumerate}
\end{proposition}

\begin{proof}
\oldpage[IV-4]{322}
La condition $b)$ peut encore s'exprimer de la façon suivante :

\emph{$b')$ Pour toute partie fermée $Y$ de $X$ telle que
$\operatorname{codim}(Y,X)\geq2$ et telle qu'il existe $\lambda$ pour lequel
$Y\subset X-U_\lambda$, le couple $(X,Y)$ est parafactoriel.}

La propriété $a)$ peut encore s'exprimer de la façon suivante : si l'ensemble fermé
$N$ des points $x\in X$ où $Z$ n'est pas principal est contenu dans l'un des
$X-U_\lambda$, alors $Z$ est localement principal. Comme $X$ est normal, la condition
$\dim(\mathcal O_{X,x})\leq1$ entraîne que $\mathcal O_{X,x}$ est un corps ou un
anneau de valuation discrète, donc $Z_x$ est principal; en d'autres termes, on a
nécessairement $\operatorname{codim}(N,X)\geq2$ (5.1.3). La propriété $a)$ est donc
équivalente à la suivante :

\emph{$a')$ S'il existe un ensemble fermé $Y$ contenu dans un des $X-U_\lambda$, tel
que $\operatorname{codim}(Y,X)\geq2$ et que $Z|(X-Y)$ soit localement principal,
alors $Z$ est localement principal.}

Notons d'autre part que $b)$ implique $b')$ en vertu de (21.13.10); inversement, si
$b')$ est vérifiée et si $x\notin U_\lambda$, alors $Y=\{x\}$ est contenue dans
$X-U_\lambda$ et on a $\operatorname{codim}(Y,X)\geq2$, donc il résulte encore de
(21.13.10) que $\mathcal O_{X,x}$ est parafactoriel, ce qui prouve l'équivalence de
$b)$ et $b')$. On est ainsi ramené à prouver l'équivalence de $a')$ et de $b')$.
Notons que puisque $X$ est normal, on a, pour toute partie fermée $Y$ de $X$ telle que
$\operatorname{codim}(Y,X)\geq2$, $\operatorname{prof}_Y(\mathcal O_X)\geq2$
(5.8.6); si l'on pose $U=X-Y$, l'homomorphisme
$\mathcal O_X\to j_*(\mathcal O_U)$ est donc bijectif (21.13.4); comme les conditions
$a')$ et $b')$ sont locales, on voit qu'il suffit de montrer l'équivalence des
conditions suivantes, lorsque $X$ est noethérien et normal et $Y$ fermé dans $X$ et
tel que $\operatorname{codim}(Y,X)\geq2$ :

\emph{$a'')$ Pour tout cycle $1$-codimensionnel $Z$ sur $X$, l'hypothèse que $Z|U$
(où $U=X-Y$) est localement principal entraîne que $Z$ lui-même est localement
principal.}

\emph{$b'')$ L'homomorphisme canonique
$\operatorname{Pic}(X)\to\operatorname{Pic}(U)$ est surjectif.}

Prouvons d'abord que $a'')$ entraîne $b'')$; un élément
$\zeta_0\in\operatorname{Pic}(U)$ a pour image dans $\mathfrak{Cl}(U)$ (21.6.11) la
classe d'un cycle $1$-codimensionnel $Z_0$ sur $U$, qui est localement principal.
L'hypothèse $\operatorname{codim}(Y,X)\geq2$ entraîne que l'homomorphisme de
restriction $\mathscr Z^1(X)\to\mathscr Z^1(U)$ est bijectif, donc $Z_0=Z|U$, où $Z$
est un cycle $1$-codimensionnel sur $X$; puisque $Z_0$ est localement principal, il en
est de même de $Z$ en vertu de $a'')$; il résulte de (21.6.11) que l'image de $Z$ dans
$\mathfrak{Cl}(X)$ est l'image d'un unique élément $\zeta\in\operatorname{Pic}(X)$,
et il est clair alors que $\zeta_0$ est l'image de $\zeta$.

Inversement, prouvons que $b'')$ entraîne $a'')$. Soit donc $Z$ un cycle
$1$-codimensionnel sur $X$ tel que $Z|U$ soit localement principal; l'image de $Z|U$
dans $\mathfrak{Cl}(U)$ est donc l'image d'un unique élément
$\zeta_0\in\operatorname{Pic}(U)$ (21.6.11). Par hypothèse il existe
$\zeta\in\operatorname{Pic}(X)$ dont l'image dans $\operatorname{Pic}(U)$ est
$\zeta_0$; l'image de $\zeta$ dans $\mathfrak{Cl}(X)$ est donc la classe d'un cycle
$1$-codimensionnel $Z'$ sur $X$ tel que $Z|U$ et $Z'|U$ soient linéairement
équivalents. Mais puisque $U$ est schématiquement dense dans $X$ l'image de
$\mathscr Z^1_{\mathrm{princ}}(X)$ par l'isomorphisme
$\mathscr Z^1(X)\to\mathscr Z^1(U)$ est $\mathscr Z^1_{\mathrm{princ}}(U)$, donc $Z$
et $Z'$ sont linéairement équivalents, et puisque $Z'$ est localement principal, il
en est de même de $Z$.
\end{proof}

\begin{corollary}[21.13.15]
\label{IV.21.13.15}
Soient $X$ un préschéma localement noethérien et normal, $S$ une partie de $X$. Les
conditions suivantes sont équivalentes :
\begin{enumerate}[label=\emph{\alph*)}]
  \item Tout cycle $1$-codimensionnel sur $X$ qui est principal aux points de $S$ est
  localement principal.
  \oldpage[IV-4]{323}
  \item Pour tout $x\in X$ qui n'est pas générisation d'un point de $S$ (autrement
  dit, tel que $\overline{\{x\}}\cap S=\varnothing$) et qui est tel que
  $\dim(\mathcal O_{X,x})\geq2$, l'anneau $\mathcal O_{X,x}$ est parafactoriel.
\end{enumerate}
\end{corollary}

\begin{proof}
Il est clair que l'ensemble $S'$ des points de $X$ générisations de points de $S$ est
l'intersection des voisinages ouverts de $S$. On peut se borner au cas où $X$ est
noethérien (les propriétés $a)$ et $b)$ étant locales sur $X$); comme tout cycle
$1$-codimensionnel sur $X$ qui est principal aux points de $S$ l'est aussi aux points
d'un voisinage ouvert de $S$, on voit que la condition $a)$ de (21.13.15) équivaut à
la condition $a)$ de (21.13.14) appliquée à la famille $(U_\lambda)$ des voisinages
ouverts de $S$. Le corollaire résulte alors de l'équivalence de $a)$ et $b)$ dans
(21.13.14).
\end{proof}

\begin{remark}[21.13.16]
\label{IV.21.13.16}
Sous les conditions générales de (21.13.14), supposons de plus que l'on ait
$\varinjlim\operatorname{Pic}(U_\lambda)=0$. Alors la condition $a)$ de (21.13.14)
est aussi équivalente à la suivante :
\begin{enumerate}
  \item[c)] \emph{Tout cycle $1$-codimensionnel sur $X$ dont le support est contenu
  dans un des ensembles $X-U_\lambda$ est localement principal.}
\end{enumerate}

On peut se borner au cas où $X$ est irréductible et tous les $U_\lambda$ sont non
vides. Il est clair que $a)$ entraîne $c)$, car si le cycle $1$-codimensionnel $Z$ sur
$X$ a son support dans $X-U_\lambda$, $Z|U_\lambda=0$, donc $Z|U_\lambda$ est
localement principal. Inversement $c)$ entraîne $a)$ : soit en effet $Z$ un cycle
$1$-codimensionnel sur $X$ tel que $Z|U_\lambda$ soit localement principal; comme $X$
est normal, $Z|U_\lambda=\operatorname{cyc}(D_\lambda)$, où $D_\lambda$ est un
diviseur sur $U_\lambda$ (21.6.10, (ii)), et l'hypothèse implique (en vertu de
(21.3.4)) qu'il y a un ensemble $U_\mu\subset U_\lambda$ tel que
$D_\mu=D_\lambda|U_\mu$ soit équivalent à $0$. Si
$D_\mu=\operatorname{div}(f_\mu)$, où $f_\mu$ est une fonction rationnelle régulière
sur $U_\mu$, $f_\mu$ peut être considérée comme une fonction rationnelle régulière
sur $X$. Si $Z'=\operatorname{cyc}(\operatorname{div}(f_\mu))$, on voit donc que
$Z''=Z-Z'$ a son support dans $X-U_\mu$; en vertu de l'hypothèse, $Z''$ est
localement principal, d'où la conclusion.

La condition $\varinjlim\operatorname{Pic}(U_\lambda)=0$ sera en particulier remplie
si $(U_\lambda)$ est la famille des voisinages ouverts d'un \emph{point} $z\in X$,
car pour tout $\mathcal O_{U_\lambda}$-Module inversible $\mathscr L_\lambda$, il
existe par définition un $U_\mu\subset U_\lambda$ tel que
$\mathscr L_\lambda|U_\mu$ soit isomorphe à $\mathcal O_{U_\mu}$. On voit donc que,
dans l'énoncé de (21.13.15), si $S=\{z\}$, les conditions $a)$ et $b)$ équivalent
aussi à la suivante :
\begin{enumerate}
  \item[c)] \emph{Tout cycle $1$-codimensionnel sur $X$ dont le support ne contient
  pas $z$ est localement principal.}
\end{enumerate}
Si de plus $\operatorname{Pic}(X)=0$, on en conclura que cette condition entraîne que
tout cycle $1$-codimensionnel dont le support ne contient pas $z$ est principal.
\end{remark}

\subsection{Le théorème de Ramanujam--Samuel}
\EGAReaderTarget{ega4.c04.subsec.21-14}{subsection:IV.21.14}
\label{IV.21.14}

\begin{theorem}[21.14.1]
\label{IV.21.14.1}
\textup{(Ramanujam--Samuel)}. ---
Soit $A$ un anneau local noethérien d'idéal maximal $\mathfrak m$, tel que son complété
$\widehat A$ soit intègre et intégralement clos. Soient $B$ un anneau local noethérien
tel que $\dim(B)\geq\dim(A)$, $\rho:A\to B$ un homomorphisme local faisant de $B$
une $A$-algèbre formellement lisse (pour les topologies préadiques
$(\mathbf 0, 19.3.1)$) et tel que le corps résiduel de $B$ soit fini sur celui de $A$.
Alors tout cycle $1$-codimensionnel sur $\operatorname{Spec}(B)$ qui est principal au
point $\mathfrak p=\mathfrak mB$, est un cycle $1$-codimensionnel principal.
\end{theorem}

\begin{proof}
Si $k=A/\mathfrak m$ est le corps résiduel de $A$, $B/\mathfrak mB$ est une
$k$-algèbre formellement lisse (pour sa topologie préadique) $(\mathbf 0, 19.3.5)$,
\oldpage[IV-4]{324}
donc régulière, et en particulier intègre, autrement dit
$\mathfrak p=\mathfrak mB$ est bien un idéal premier de $B$, ce qui justifie
l'énoncé. Tout revient évidemment à prouver que tout idéal premier $\mathfrak q$ de
$B$ non contenu dans $\mathfrak p$ est \emph{principal}.

Soient $\widehat A$, $\widehat B$ les complétés de $A$ et $B$ respectivement, de
sorte que l'idéal maximal de $\widehat A$ est $\mathfrak m\widehat A$; on sait
$(\mathbf 0, 19.3.6)$ que $\widehat B$ est une $\widehat A$-algèbre formellement
lisse pour les topologies adiques. Soit $k'$ le corps résiduel de $\widehat B$,
extension finie de $k$; il existe un homomorphisme local $\widehat A\to C$, où $C$
est un anneau local noethérien qui est un $\widehat A$-module \emph{fini} et
\emph{plat} (donc libre) et est tel que $C/\mathfrak mC$ soit isomorphe à $k'$
$(\mathbf 0_{\mathrm{III}}, 10.3.1)$; on en déduit que $C$ est \emph{complet}, et il
résulte alors de (7.5.1), (7.5.3) et (6.5.4, (ii)) que $C$ est \emph{intègre et
intégralement clos}. En outre,
\[
  D=\widehat B\widehat\otimes_{\widehat A}C
\]
est un anneau semi-local complet, composé direct d'anneaux locaux complets dont l'un,
$D_0$, a pour corps résiduel $k'$ (puisque $k'\otimes_k k'$ est composé direct
d'anneaux locaux dont l'un est isomorphe à $k'$). Comme $D$ est formellement lisse
sur $C$, il en est de même de $D_0$; par suite $D_0/\mathfrak mD_0$ est une
$k'$-algèbre formellement lisse, de corps résiduel $k'$, ce qui entraîne qu'elle est
$k'$-isomorphe à une algèbre de séries formelles
$k'[[T_1,\ldots,T_n]]$ $(\mathbf 0, 19.6.4)$; on en conclut, par
$(\mathbf 0, 19.7.1.5)$, que $D_0$ est $C$-isomorphe à
$C[[T_1,\ldots,T_n]]$, et par suite \emph{intègre et intégralement clos} (Bourbaki,
\emph{Alg. comm.}, chap.~V, \S1, n\textsuperscript{o}~4, prop.~14). Comme les
morphismes $\operatorname{Spec}(D_0)\to\operatorname{Spec}(\widehat B)$ et
$\operatorname{Spec}(D_0)\to\operatorname{Spec}(B)$ sont fidèlement plats, on en
déduit que $B$ et $\widehat B$ sont aussi \emph{intègres et intégralement clos}
(2.1.13). Cela prouve que l'idéal $\mathfrak qD_0$ est divisoriel (Bourbaki,
\emph{Alg. comm.}, chap.~VII, \S1, n\textsuperscript{o}~10, prop.~15) et non contenu
dans $\mathfrak mD_0$, sans quoi on aurait
\[
  \mathfrak q=(\mathfrak qD_0)\cap B
  \subset(\mathfrak mD_0)\cap B=\mathfrak mB=\mathfrak p
\]
par fidèle platitude $(\mathbf 0_{\mathrm I}, 6.5.1)$. On en conclut qu'aucun des
idéaux premiers $\mathfrak r_h$ de hauteur $1$ dans $D_0$ qui contiennent
$\mathfrak qD_0$ ne peut être contenu dans $\mathfrak mD_0$; si l'on prouve que ces
idéaux sont principaux, il en résultera que $\mathfrak qD_0$ est principal, les
diviseurs (au sens de Bourbaki) de $\mathfrak qD_0$ et d'un produit de puissances des
$\mathfrak r_h$ étant égaux (Bourbaki, \emph{Alg. comm.}, chap.~VII, \S1,
n\textsuperscript{o}~4, prop.~5). Comme $\mathfrak qD_0\cap B=\mathfrak q$, on en
déduit par fidèle platitude que $\mathfrak q$ est principal (2.5.2), en utilisant le
fait que $B$ est un anneau local.

On est ainsi ramené à prouver le théorème dans le cas particulier où $A$ est
\emph{complet} et $B=A[[T_1,\ldots,T_n]]$. Montrons d'abord qu'on peut se ramener au
cas où $n=1$. En effet, procédons par récurrence sur $n$, et soit
$f\in\mathfrak q$ un élément n'appartenant pas à $\mathfrak p=\mathfrak mB$; il
suffira de montrer qu'il existe un $A$-automorphisme $\sigma$ de $B$ tel que
$\sigma(f)$ n'appartienne pas à l'idéal
\[
  \mathfrak p'=\mathfrak p+BT_1+\cdots+BT_{n-1};
\]
en effet, si $B'=A[[T_1,\ldots,T_{n-1}]]$, $B'$ est complet et a pour idéal maximal
$\mathfrak mB'+B'T_1+\cdots+B'T_{n-1}=\mathfrak m'$. Comme on a
$B=B'[[T_n]]$ et $\mathfrak p'=\mathfrak m'B$, il résultera de l'hypothèse de
récurrence (compte tenu de ce que $B'$ est intègre et intégralement clos) que l'idéal
$\sigma(\mathfrak q)$ est principal dans $B$, donc aussi $\mathfrak q$. Or, si
$\overline f\in k[[T_1,\ldots,T_n]]$ est l'image de $f$ (« série réduite » de $f$),
on sait (Bourbaki, \emph{Alg. comm.}, chap.~VII, \S3, n\textsuperscript{o}~7,
lemme~3) qu'il y a un $A$-automorphisme $\sigma$ de $B$ tel que
$(\overline{\sigma(f)})(0,\ldots,0,T_n)\neq0$, ce qui implique évidemment
$\sigma(f)\notin\mathfrak p'$.

Supposons donc $n=1$ et écrivons $T$ au lieu de $T_1$, de sorte que $B=A[[T]]$. Il
suffit de montrer que dans l'anneau de polynômes $A[T]$, l'idéal
$\mathfrak q\cap A[T]$ est principal;
\oldpage[IV-4]{325}
en effet, soit $f_0$ un élément de $\mathfrak q$ n'appartenant pas à $\mathfrak mB$;
c'est une série formelle dont la série réduite $\overline f_0$ n'est pas nulle; donc,
en vertu du théorème de préparation (Bourbaki, \emph{Alg. comm.}, chap.~VII, \S3,
n\textsuperscript{o}~8, prop.~5), pour tout $f\in\mathfrak q$, il existe $g\in B$ et
un polynôme $r\in A[T]$ tels que $f=gf_0+r$ et on a donc
$r\in\mathfrak q\cap A[T]$; d'autre part (\emph{loc. cit.}, prop.~6) il existe un
polynôme distingué non constant $F_0\in A[T]$ et un élément inversible $u\in B$ tels
que $f_0=uF_0$, donc on a aussi $F_0\in\mathfrak q\cap A[T]$, ce qui prouve que
$\mathfrak q$ est engendré par $\mathfrak q\cap A[T]=\mathfrak q_1$. Puisque $B$
est plat sur $A[T]$ $(\mathbf 0_{\mathrm I}, 7.3.3)$, il résulte de Bourbaki,
\emph{Alg. comm.}, chap.~VII, \S1, n\textsuperscript{o}~10, prop.~15, que
$\mathfrak q_1$ est un idéal premier de hauteur $1$ dans $A[T]$.

En outre, on a nécessairement $\mathfrak q_1\cap A=0$; sinon,
$\mathfrak q_1\cap A$ serait nécessairement de hauteur $\geq1$, et il résulterait de
(5.5.3) que l'on aurait $\mathfrak q_1=(\mathfrak q_1\cap A)A[T]$. Mais alors,
puisque $\mathfrak q_1\cap A\subset\mathfrak m$, on aurait
$\mathfrak q_1\subset\mathfrak mA[T]$ contrairement à l'hypothèse sur
$\mathfrak q$. Si $K$ est le corps des fractions de $A$, $\mathfrak q_1K[T]$ est
donc un idéal premier distinct de $0$ et de $K[T]$ dans $K[T]$, donc de la forme
$hK[T]$, où
\[
  h(T)=T^m+a_1T^{m-1}+\cdots+a_m,
  \qquad m\geq1,
\]
avec les $a_i\in K$, $h$ étant irréductible dans $K[T]$. Mais on a vu plus haut qu'il
existe dans $\mathfrak q_1$ un polynôme distingué non constant $F_0\in A[T]$. Si $t$
est la classe de $T$ dans $K[T]/\mathfrak q_1K[T]$, $t$ est donc racine du polynôme
$F_0$ dans une extension de $K$ et par suite $h$ divise $F_0$ dans $K[T]$; mais
puisque $h$ et $F_0$ sont unitaires, cela entraîne que les coefficients $a_i$ de $h$
sont entiers sur $A$ (Bourbaki, \emph{Alg. comm.}, chap.~V, \S1,
n\textsuperscript{o}~3, prop.~11), donc appartiennent à $A$ puisque $A$ est
intégralement clos. Autrement dit, on a
\[
  h\in A[T]\cap\mathfrak q_1K[T]=\mathfrak q_1
\]
(Bourbaki, \emph{Alg. comm.}, chap.~II, \S1, n\textsuperscript{o}~5, prop.~11);
comme tout polynôme $g\in\mathfrak q_1$ est divisible par $h$ dans $K[T]$ et que $h$
est unitaire, les coefficients de $g/h$ appartiennent à $A$, donc
$\mathfrak q_1=hA[T]$. C.Q.F.D.
\end{proof}

L'énoncé (21.14.1) est \emph{équivalent} au suivant :

\begin{corollary}[21.14.2]
\label{IV.21.14.2}
Sous les hypothèses de (21.14.1) sur $A$, $B$ et $\rho$ pour tout idéal premier
$\mathfrak q$ de $B$ non contenu dans $\mathfrak p=\mathfrak mB$ et tel que
$\dim(B_\mathfrak q)\geq2$, l'anneau $B_\mathfrak q$ est parafactoriel. En
particulier, si $\dim(B\otimes_A k)>0$ (\emph{i.e.} $\dim B>\dim A$
$((\mathbf 0, 19.7.1)$ et $(\mathbf 0, 6.1.1)))$, l'anneau $B$ est parafactoriel.
\end{corollary}

\begin{proof}
En effet, on a vu dans la démonstration de (21.14.1) que les conditions de l'énoncé
entraînent que $B$ est intègre et intégralement clos; l'équivalence des énoncés
(21.14.1) et (21.14.2) résulte alors de (21.13.15) appliqué à
$X=\operatorname{Spec}(B)$ et $S=\{\mathfrak p\}$.
\end{proof}

\begin{proposition}[21.14.3]
\label{IV.21.14.3}
Soient $S$ un préschéma normal, $f:X\to S$ un morphisme lisse.
\begin{enumerate}[label=(\roman*)]
  \item Si $S$ est localement noethérien (donc aussi $X$), alors, pour tout $x\in X$
  tel que $\dim(\mathcal O_{X,x})\geq2$ et que $x$ ne soit pas point maximal de sa
  fibre $f^{-1}(f(x))$, l'anneau $\mathcal O_{X,x}$ est parafactoriel. Tout cycle
  $1$-codimensionnel $Z$ sur $X$ tel que $\operatorname{Supp}(Z)$ ne contienne aucune
  composante irréductible d'une fibre $f^{-1}(s)=X_s$, est localement principal.
  \item Soit $Y$ une partie fermée de $X$ ne contenant aucune composante irréductible
  d'une fibre $X_s$, et telle que, pour tout point maximal $\eta$ de $S$,
  $Y_\eta=Y\cap X_\eta$ soit de codimension $\geq2$ dans $X_\eta$. Alors le couple
  $(X,Y)$ est parafactoriel.
\end{enumerate}
\end{proposition}

\begin{proof}
On notera que les conditions de (ii) sont remplies si $Y$ est une partie fermée de $X$
telle que, pour tout $s\in S$, $Y_s=Y\cap X_s$ soit de codimension $\geq2$ dans
$X_s$.

Pour prouver (21.14.3), notons d'abord que sous les hypothèses de (i), si l'on pose
$Y=\overline{\{x\}}$, il existe un voisinage ouvert $V$ de $x$ dans $X$ tel que
$Y\cap V$ et $V$ vérifient les conditions de (ii) : en effet, il résulte de l'hypothèse
et de (9.5.3) que l'on peut prendre $V$ tel que $V\cap Y$ ne contienne aucune
composante irréductible d'un $X_s$; d'autre part, pour tout point maximal $\eta$ de
$S$ tel que $Y\cap X_\eta\neq\varnothing$, les points de $Y\cap X_\eta$ sont des
spécialisations de $x$, donc en un tel point $z$ on a
$\dim(\mathcal O_{X,z})\geq2$, et comme
$\mathcal O_{X,z}=\mathcal O_{X_\eta,z}$ ($S$ étant réduit au point $\eta$ en vertu
de (2.1.13)), $Y\cap X_\eta$ est bien de codimension $\geq2$ dans $X_\eta$. Comme le
couple $(V,Y\cap V)$ est alors parafactoriel en vertu de (ii), on conclut de (21.13.10)
que l'anneau $\mathcal O_{X,x}$ est parafactoriel, c'est-à-dire la première assertion
de (i). Pour prouver la seconde, on peut se borner au cas où
$Z=\overline{\{z\}}$, où $z\in X^{(1)}$; comme $\mathcal O_{X,z}$ est un anneau local
intégralement clos et de dimension $1$, c'est un anneau de valuation discrète et $Z$
est donc principal au point $z$; pour tout autre point $x$ de
$\operatorname{Supp}(Z)$, on a donc $\dim(\mathcal O_{X,x})\geq2$ et $x$ n'est pas
point maximal de sa fibre par hypothèse, donc $\mathcal O_{X,x}$ est parafactoriel.
Appliquons alors (21.13.15) en prenant pour $S$ l'ensemble formé de $z$ et des points
maximaux des fibres de $f$; il est clair que $Z$ est principal aux points de $S$
puisque $z$ est le seul point de $S$ appartenant à $\operatorname{Supp}(Z)$; comme
d'autre part la condition $b)$ de (21.13.15) est évidemment remplie, on en conclut que
$Z$ est localement principal.

On est donc ramené à prouver (ii). Posons $U=X-Y$ et soit $j:U\to X$ l'injection
canonique.
\oldpage[IV-4]{326}
Comme les hypothèses et la conclusion sont locales sur $S$ et sur $X$ en vertu de
(21.13.6, (i)), on peut se borner au cas où $S$ et $X$ sont affines, $f$ étant donc de
présentation finie. Puisque les fibres de $f$ sont régulières, l'hypothèse sur $Y$
entraîne que, pour tout point $x\in Y$, on a
\[
  \operatorname{prof}(\mathcal O_{X_{f(x)},x})
  =\dim(\mathcal O_{X_{f(x)},x})\geq1;
\]
il résulte de (19.9.8) que l'homomorphisme canonique
$\mathcal O_X\to j_*(\mathcal O_U)$ est \emph{injectif}. De plus, remplaçant $Y$ par
une partie fermée plus grande suivant la méthode de la démonstration de (19.9.8), et
utilisant (21.13.6, (ii)), on peut supposer que $Y$ est défini par un idéal \emph{de
type fini} de l'anneau de $X$, donc l'ouvert $U=X-Y$ est \emph{quasi-compact et
quasi-séparé} et l'ensemble fermé $Y$ est \emph{constructible}.

Pour prouver (ii), il suffit, pour tout $\mathcal O_U$-Module inversible $\mathscr L$
et tout point $x\in Y$ donnés, d'établir l'existence d'un voisinage ouvert $V$ de $x$
dans $X$ tel que : $1^{\mathrm o}$ l'homomorphisme canonique
$\mathcal O_V\to(j_V)_*(\mathcal O_{U\cap V})$ est surjectif;
$2^{\mathrm o}$ il existe un $\mathcal O_V$-Module inversible $\mathscr L'_V$ tel que
$\mathscr L'_V|(U\cap V)=\mathscr L|(U\cap V)$ ((21.13.5) et (21.13.3)).

Posons $s=f(x)$; nous allons voir qu'on peut se borner au cas où
$S=S_0=\operatorname{Spec}(\mathcal O_{S,s})$. Soit en effet $(S_\nu)$ un système
fondamental de voisinages ouverts affines de $s$ dans $S$, et posons
$X_\nu=f^{-1}(S_\nu)$, $Y_\nu=Y\cap X_\nu$, $U_\nu=X_\nu-Y_\nu$,
$j_\nu:U_\nu\to X_\nu$ étant l'injection canonique; $X_0=X\times_S S_0$ est alors
limite projective des $X_\nu$ (8.1.2, $a$)); supposons la proposition vraie pour le
morphisme $f_0:X_0\to S_0$ et pour $Y_0=Y\cap X_0$, et posons $U_0=X_0-Y_0$,
$j_0:U_0\to X_0$ étant l'injection canonique. La projection
$p_\nu:X_0\to X_\nu$ étant un morphisme affine, on a
\[
  (j_0)_*(\mathcal O_{U_0})
  =p_\nu^*((j_\nu)_*(\mathcal O_{U_\nu}))
  \qquad(\mathbf{II}, 1.5.2)
\]
et l'homomorphisme canonique
$\rho_0:\mathcal O_{X_0}\to(j_0)_*(\mathcal O_{U_0})$ n'est autre que
$p_\nu^*(\rho_\nu)$, où
$\rho_\nu:\mathcal O_{X_\nu}\to(j_\nu)_*(\mathcal O_{U_\nu})$ est l'homomorphisme
canonique; comme par hypothèse $\rho_0$ est surjectif, il en est de même de $\rho_\nu$
pour $\nu$ assez grand (8.5.7). Soit d'autre part $\mathscr L_0$ le
$\mathcal O_{U_0}$-Module inversible restriction de $\mathscr L$, et notons que les
$X_\nu$ sont affines, donc quasi-compacts et quasi-séparés; puisque par hypothèse il
existe un $\mathcal O_{X_0}$-Module inversible $\mathscr L'_0$ tel que
$\mathscr L'_0|U_0=\mathscr L_0$, il existe un $\nu$ assez grand et un
$\mathcal O_{X_\nu}$-Module quasi-cohérent $\mathscr L'_\nu$ de présentation finie
tels que $\mathscr L'_0$ soit isomorphe à $p_\nu^*(\mathscr L'_\nu)$ (8.5.2, (ii));
de plus (8.5.5) on peut supposer que $\mathscr L'_\nu$ est inversible. Enfin, comme
les $U_\nu$ sont quasi-compacts et quasi-séparés et que
$\mathscr L'_0|U_0=\mathscr L_0$, il existe $\nu$ assez grand tel que
$\mathscr L'_\nu|U_\nu$ et $\mathscr L|U_\nu$ soient isomorphes (8.5.2.5).

\oldpage[IV-4]{327}
On est ainsi ramené au cas où $S=\operatorname{Spec}(A)$,
$X=\operatorname{Spec}(B)$, où $A$ est un anneau \emph{local}; puisque $X$ est
normal et $f$ lisse, $S$ est normal (17.5.7), donc $A$ est \emph{intègre et
intégralement clos}. Considérons alors $A$ comme limite inductive de ses
\emph{sous-$\mathbb Z$-algèbres de type fini}; comme la clôture intégrale d'une telle
sous-$\mathbb Z$-algèbre est un sous-anneau de $A$ et est aussi une
$\mathbb Z$-algèbre de type fini (7.8.3, (ii), (iii) et (vi)), $A$ est limite inductive
filtrante des sous-$\mathbb Z$-algèbres de type fini $A_\lambda$ de $A$ qui sont des
anneaux \emph{intégralement clos}. En vertu de (1.8.4.2), il existe un indice
$\lambda$ et une $A_\lambda$-algèbre de type fini $B_\lambda$ tels que
\[
  B=B_\lambda\otimes_{A_\lambda}A
\]
à un $A$-isomorphisme près. Posons $S_\lambda=\operatorname{Spec}(A_\lambda)$,
$X_\lambda=\operatorname{Spec}(B_\lambda)$; si $p_\lambda:X\to X_\lambda$ est la
projection canonique, on peut en outre supposer (puisque $Y$ est
\emph{constructible}) que $Y=p_\lambda^{-1}(Y_\lambda)$, où $Y_\lambda$ est une
partie fermée de $X_\lambda$ (8.3.11). Soit $f_\lambda$ le morphisme
$X_\lambda\to S_\lambda$; en vertu de la transitivité des fibres et de (4.2.6),
$Y_\lambda$ ne contient aucune composante irréductible des fibres
$f_\lambda^{-1}(s_\lambda)$ pour aucun $s_\lambda\in S_\lambda$. Puisque $f$ est
lisse, on peut aussi supposer que $f_\lambda$ est lisse (17.7.8); enfin, l'image du
point générique $\eta$ de $S$ dans $S_\lambda$ est le point générique $\eta_\lambda$
de $S_\lambda$, et on a
\[
  \operatorname{codim}((Y_\lambda)_{\eta_\lambda},
  (X_\lambda)_{\eta_\lambda})
  =\operatorname{codim}(Y_\eta,X_\eta)\geq2
\]
en vertu de (6.1.4). On voit donc que $S_\lambda$, $X_\lambda$ et $Y_\lambda$
vérifient toutes les hypothèses de (ii), et en posant
$X_\mu=X_\lambda\times_{S_\lambda}S_\mu$ pour $\lambda\leq\mu$, il en est de même de
$S_\mu$, $X_\mu$ et $Y_\mu$. Montrons que si l'on prouve que le couple
$(X_\mu,Y_\mu)$ est parafactoriel pour tout $\mu\gg\lambda$, il en est de même de
$(X,Y)$. En effet, soient $U=X-Y$, $U_\mu=X_\mu-Y_\mu$, $j:U\to X$,
$j_\mu:U_\mu\to X_\mu$ les injections canoniques; la projection
$p_\mu:X\to X_\mu$ étant un morphisme affine, on a
\[
  j_*(\mathcal O_U)=p_\mu^*((j_\mu)_*(\mathcal O_{U_\mu}))
  \qquad(\mathbf{II}, 1.5.2),
\]
et par suite l'homomorphisme canonique
$\rho:\mathcal O_X\to j_*(\mathcal O_U)$ n'est autre que $p_\mu^*(\rho_\mu)$, où
$\rho_\mu:\mathcal O_{X_\mu}\to(j_\mu)_*(\mathcal O_{U_\mu})$ est l'homomorphisme
canonique; ce dernier étant par hypothèse bijectif, il en est de même de $\rho$.
D'autre part, pour tout $\mathcal O_U$-Module inversible $\mathscr L$, il existe un
$\mu\gg\lambda$ et un $\mathcal O_{U_\mu}$-Module inversible $\mathscr L_\mu$ tels
que $\mathscr L=p_\mu'^*(\mathscr L_\mu)$ (en désignant par
$p_\mu':U\to U_\mu$ la restriction de $p_\mu$) (8.5.2 et 8.5.5), $U_\lambda$ étant
quasi-compact puisque $X_\lambda$ est noethérien. Par hypothèse, il existe un
$\mathcal O_{X_\mu}$-Module inversible $\mathscr L'_\mu$ tel que
$\mathscr L'_\mu|U_\mu$ soit isomorphe à $\mathscr L_\mu$;
$\mathscr L'=p_\mu^*(\mathscr L'_\mu)$ est alors un $\mathcal O_X$-Module
inversible tel que $\mathscr L'|U$ soit isomorphe à $\mathscr L$, ce qui prouve notre
assertion.

On est ainsi ramené à prouver (ii) lorsque l'anneau $A$ est une
$\mathbb Z$-algèbre de type fini intégralement close; comme les anneaux locaux
$\mathcal O_{S,s}$ de $S$ sont alors des anneaux excellents intégralement clos, leurs
\emph{complétés} sont aussi intégralement clos (7.8.3, (ii), (iii) et (vii)). Pour
prouver que le couple $(X,Y)$ est parafactoriel, il suffit de montrer que pour tout
$x\in Y$, l'anneau $\mathcal O_{X,x}$ est parafactoriel (21.13.10). Soit $y$ un point
fermé de $Y_{f(x)}$, qui soit spécialisation de $x$ (5.1.11); si l'on pose
$s=f(x)=f(y)$, $\mathcal O_{S,s}$ a un
\oldpage[IV-4]{328}
complété intégralement clos, et $\mathcal O_{X,y}$ est une
$\mathcal O_{S,s}$-algèbre formellement lisse pour les topologies préadiques
(17.5.3); puisque $Y$ ne contient aucune composante irréductible de $X_s$, on a
$\dim(\mathcal O_{X,y})>\dim(\mathcal O_{S,s})$ (17.5.8). Si $s\neq\eta$, on a
$\dim(\mathcal O_{S,s})\geq1$; si au contraire $s=\eta$, on a par hypothèse
$\dim(\mathcal O_{X,x})\geq2$; donc dans tous les cas
$\dim(\mathcal O_{X,y})\geq2$. En outre, l'idéal premier de $\mathcal O_{X,y}$
correspondant au point $x$ n'est pas contenu dans
$\mathfrak m_s\mathcal O_{X,y}$, puisque $Y$ ne contient aucune composante
irréductible de $X_s$. Enfin, puisque $y$ est fermé dans $X_s$, $k(y)$ est extension
finie de $k(s)$ $(\mathbf I, 6.4.2)$. Dans tous les cas, on peut appliquer à
$A=\mathcal O_{S,s}$, $B=\mathcal O_{X,y}$ et
$B_\mathfrak q=\mathcal O_{X,x}$ le résultat de (21.14.2), ce qui achève la
démonstration.
\end{proof}

\begin{env}[21.14.4]
\label{IV.21.14.4}
\emph{Remarques}. ---

(i) Il se peut que les énoncés (21.14.1) et (21.14.2) restent valables \emph{sans
hypothèse} sur le corps résiduel de $B$. Énoncé avec cette généralité, le résultat
équivaudrait, en vertu de (21.13.15) et (21.13.12.1) au suivant : Soient $A$ un
anneau local noethérien complet, intègre et intégralement clos, $B$ un anneau local
noethérien qui est une $A$-algèbre formellement lisse, tel que
$\dim(B)>\dim(A)$; alors $B$ est parafactoriel.

(ii) Nous verrons au chap.~VI, en employant une technique de « descente finie »
appliquée au morphisme $Y'\to Y=\operatorname{Spec}(A)$, où $Y'$ est le normalisé de
$Y$, que la conclusion de (21.14.1) (ou de (21.14.2)) reste valable en remplaçant
l'hypothèse que $\widehat A$ est intégralement clos par l'hypothèse que
$\widehat A$ est \emph{réduit}, pourvu que $\dim B\geq\dim A+2$. Si on remplace cette
dernière condition par $\dim B\geq\dim A+3$, on peut même supprimer l'hypothèse que
$\widehat A$ est réduit. De même, la conclusion de (21.14.3) reste valable en
remplaçant l'hypothèse que $X$ est normal par l'hypothèse que $X$ est réduit, pourvu
que $\dim(\mathcal O_{X,x})\geq\dim(\mathcal O_{S,f(x)})+2$.

(iii) Au chap.~III, 3\textsuperscript{e} partie, on prouve que si $f:X\to S$ est un
morphisme lisse, $Y$ une partie fermée localement constructible de $X$ telle que,
pour tout $s\in S$, on ait (avec les notations de (21.14.3))
$\operatorname{codim}(Y_s,X_s)\geq3$, alors le couple $(X,Y)$ est parafactoriel
([41], XII, 4.8). Cette conclusion n'est plus valable lorsqu'on suppose seulement que
$\operatorname{codim}(Y_s,X_s)\geq2$ pour tout $s\in S$ et lorsqu'on ne suppose plus
$X$ réduit. Par exemple, soient $k$ un corps, $A=k[T]/(T^2)$, algèbre des nombres
duaux sur $k$, $S=\operatorname{Spec}(A)$,
$X=\operatorname{Spec}(A[T_1,T_2])$ ($T_1$, $T_2$ indéterminées), de sorte que
$X=\mathbb V_S^2$, $Y$ étant la « section nulle » de ce fibré; si $s$ est l'unique
point fermé de $S$, on a $X_s=\mathbb V_k^2$ et $Y_s$ est la « section nulle » de
$X_s$. Pour voir que le couple $(X,Y)$ n'est pas parafactoriel, il suffit de montrer
que l'anneau $B=A[T_1,T_2]_\mathfrak m$, où $\mathfrak m$ est l'idéal
$(T_1)+(T_2)$ qui définit $Y$, n'est pas parafactoriel (21.13.10). Soit
$B_0=B_{\mathrm{red}}$, et notons $U$ et $U_0$ les complémentaires du point fermé
dans $\operatorname{Spec}(B)$ et $\operatorname{Spec}(B_0)$; en raisonnant comme
dans (21.13.9), on a la suite exacte, extension de (21.13.9.1)
\[
  \Gamma(U,\mathcal O_U)^*\longrightarrow\Gamma(U_0,\mathcal O_{U_0})^*
  \longrightarrow\mathrm H^1(U_0,\mathcal O_{U_0})
  \longrightarrow\operatorname{Pic}(U)\longrightarrow\operatorname{Pic}(U_0).
\]
Or, on a $\Gamma(U,\mathcal O_U)=B$, $\Gamma(U_0,\mathcal O_{U_0})=B_0$ puisque
$\operatorname{prof}(B)=\operatorname{prof}(B_0)=2$ (5.10.5). Par ailleurs
$\operatorname{Pic}(U_0)=0$ puisque $B_0$ est factoriel, et
$\mathrm H^1(U_0,\mathcal O_{U_0})\neq0$ [41, 3, Exemple III-1]; comme
l'homomorphisme $B^*\to B_0^*$ est surjectif, on en conclut que
$\operatorname{Pic}(U)\neq0$, donc que $B$ n'est pas parafactoriel.
\oldpage[IV-4]{329}
(iv) Le résultat (21.14.3) a d'abord été démontré par C.~Seshadri [44] dans le cas
particulier où $S$ est un schéma algébrique normal sur un corps algébriquement clos
$k$ et $X=S\times_k T$, où $T$ est un $k$-préschéma algébrique lisse sur $k$. La
démonstration de Seshadri [44, p.~188--189] est de nature globale et utilise la
théorie des schémas de Picard. Elle donne en outre (\emph{loc. cit.}) des résultats
tels que le suivant (pour lequel on ne possède pas à l'heure actuelle de démonstration
par voie locale). Soient $S$, $T$ deux préschémas localement de type fini sur un
corps $k$, $X=S\times_k T$, $Z$ un cycle $1$-codimensionnel sur $X$ (considéré comme
$S$-préschéma); supposons vérifiées les conditions suivantes :
\begin{enumerate}[label=\arabic*\textsuperscript{o}]
  \item $S$ et $T$ sont géométriquement normaux sur $k$ (6.7.6);
  \item Pour tout point maximal $\eta$ de $S$, le cycle $1$-codimensionnel $Z_\eta$
  sur la fibre $X_\eta$, ayant même multiplicité que $Z$ en tout point de
  $X^{(1)}\cap X_\eta$, est localement principal (autrement dit, est l'image d'un
  diviseur de $X_\eta$, puisque $X_\eta$ est normal);
  \item Pour tout $s\in S$, $Z$ est principal aux points maximaux de la fibre $X_s$.
\end{enumerate}
Alors $Z$ est localement principal. En d'autres termes, $X$ étant normal (6.14.1),
pour tout $x\in X$ qui n'est pas maximal dans sa fibre $X_s$ et qui n'appartient à
aucune des « fibres génériques » $X_\eta$ (ce qui implique
$\dim(\mathcal O_{X,x})\geq2$ en vertu de (6.1.1)), l'anneau local
$\mathcal O_{X,x}$ est parafactoriel, en vertu de (21.12.15).
\footnote{En fait, dans l'article précité, Seshadri suppose que $k$ est algébriquement
clos, $T$ séparé et « semi-complet » (i.e.\ tel que
$\Gamma(T,\mathcal O_T)$ soit $k$-isomorphe à $k$) et remplace l'hypothèse
$3^{\mathrm o}$ par l'hypothèse plus forte que $\operatorname{Supp}(Z)$ ne contient
aucune des fibres $X_s$ pour $s\in S$. Mais comme l'énoncé est local sur $S$, on en
conclut aussitôt qu'il suffit de faire l'hypothèse $3^{\mathrm o}$, et cela prouve que
la conclusion (interprétée comme ci-dessus en termes de propriété de parafactorialité
des anneaux $\mathcal O_{X,x}$) est locale sur $S$ et sur $T$, ce qui permet
d'éliminer complètement l'hypothèse que $T$ est « semi-complet » et que $k$ est
algébriquement clos, puisque (quitte à passer d'abord à la clôture algébrique de $k$)
on peut supposer d'abord $T$ affine, ce qui permet de le plonger comme ouvert d'un
schéma projectif et normal sur $k$, auquel le résultat de Seshadri s'applique. Notons
aussi que, grâce à cette réduction, il suffit de faire la démonstration de Seshadri dans
le cas où $T$ est \emph{projectif} (et non seulement « semi-complet »), cas où la
théorie du schéma de Picard utilisée par Seshadri est contenue dans celle qui sera
développée au chap.~VI de notre Traité.}
\end{env}

\subsection{Diviseurs relatifs}
\EGAReaderTarget{ega4.c04.subsec.21-15}{subsection:IV.21.15}
\label{IV.21.15}

\begin{env}[21.15.1]
\label{IV.21.15.1}
Soient $S$ un préschéma, $f:X\to S$ un morphisme \emph{plat et localement de
présentation finie}. On a défini dans (20.6.1) le faisceau d'anneaux
$\mathscr M_{X/S}$ des \emph{germes de fonctions méromorphes sur $X$ relatives à
$S$}, sous-faisceau de $\mathscr M_X$; il est clair que l'injection canonique
$\mathcal O_X\to\mathscr M_X$ (20.1.4.1) applique $\mathcal O_X$ sur un
sous-faisceau de $\mathscr M_{X/S}$, avec lequel on l'identifie. Soit
$\mathscr M_{X/S}^*$ le faisceau (en groupes multiplicatifs) des germes de sections
inversibles de $\mathscr M_{X/S}$ qui est donc un sous-faisceau de $\mathscr M_X^*$
et contient $\mathcal O_X^*$ comme sous-faisceau.
\end{env}

\begin{definition}[21.15.2]
\EGAReaderTarget{ega4.c04.def.21-15-2}{IV.21.15.2}
On appelle \emph{faisceau des diviseurs sur $X$ relativement à $S$}, ou
\emph{faisceau des diviseurs sur $X$ transversaux à $f$}, et on note
$\mathscr{D}iv_{X/S}$ le faisceau quotient (de groupes commutatifs)
$\mathscr M_{X/S}^*/\mathcal O_X^*$; les sections de ce faisceau au-dessus de $X$ se
nomment \emph{diviseurs sur $X$ relatifs à $S$}, ou \emph{diviseurs sur $X$
transversaux à $f$}; ils forment un groupe commutatif noté $\operatorname{Div}(X/S)$.
\end{definition}

Il est clair que $\mathscr{D}iv_{X/S}$ s'identifie canoniquement à un sous-faisceau de
$\mathscr{D}iv_X$, et par suite $\operatorname{Div}(X/S)$ à un sous-groupe de
$\operatorname{Div}(X)$, que l'on note encore \emph{additivement}. Pour
\oldpage[IV-4]{330}
tout ouvert $U$ de $X$, on a
$\mathscr M_{X/S}^*|U=\mathscr M_{U/S}^*$, donc
$\mathscr{D}iv_{X/S}|U=\mathscr{D}iv_{U/S}$, et le faisceau
$\mathscr{D}iv_{X/S}$ est donc égal au préfaisceau
$U\mapsto\operatorname{Div}(U/S)$.

Puisque $\mathscr M_{X/S}^*$ est un sous-faisceau de $\mathscr M_X^*$, les
définitions, notations et formules relatives aux diviseurs des sections de
$\mathscr M_X^*$ au-dessus de $X$ (21.1.3) s'appliquent sans changement aux sections
de $\mathscr M_{X/S}^*$ au-dessus de $X$.

\begin{env}[21.15.3]
\label{IV.21.15.3}
La structure de faisceau de groupes ordonnés sur $\mathscr{D}iv_X$ (21.1.6) induit
sur le sous-faisceau $\mathscr{D}iv_{X/S}$ une structure de faisceau de groupes
ordonnés, pour laquelle le faisceau de monoïdes des germes de sections positives est
$\mathscr{D}iv_X^+\cap\mathscr{D}iv_{X/S}$, que l'on note
$\mathscr{D}iv_{X/S}^+$. On a
$\Gamma(X,\mathscr{D}iv_{X/S}^+)=\operatorname{Div}^+(X)\cap
\operatorname{Div}(X/S)$; on note ce sous-monoïde de $\operatorname{Div}(X/S)$ par
$\operatorname{Div}^+(X/S)$, et il est formé des éléments $\geq0$ pour une structure
de groupe ordonné sur $\operatorname{Div}(X/S)$. Il résulte de (21.1.5.1) que
$\mathscr{D}iv_{X/S}^+$ est l'image dans $\mathscr{D}iv_{X/S}$ du sous-faisceau de
monoïdes
\[
\label{IV.21.15.3.1}
  \mathcal O_{X/S}^{\mathrm{reg}}=\mathcal O_X\cap\mathscr M_{X/S}^*.
  \tag{21.15.3.1}
\]
\end{env}

Pour tout ouvert $U$ de $X$, $\Gamma(U,\mathcal O_{X/S}^{\mathrm{reg}})$ est
l'ensemble des sections $t$ de $\mathcal O_X$ au-dessus de $U$ telles que $t$ soit
régulière et que $1/t$ appartienne à $\Gamma(U,\mathscr M_{X/S})$, ce qui signifie,
avec les notations de (20.6.1), que $t\in\mathscr T_{X/S}(U)$, de sorte que le faisceau
$\mathcal O_{X/S}^{\mathrm{reg}}$ n'est autre que le faisceau noté $\mathscr T_{X/S}$
dans (20.6.1). On peut donc écrire
\[
\label{IV.21.15.3.2}
  \mathscr{D}iv_{X/S}^+=\mathcal O_{X/S}^{\mathrm{reg}}/\mathcal O_X^*
  \tag{21.15.3.2}
\]
à un isomorphisme canonique près.

\begin{env}[21.15.3.3]
\label{IV.21.15.3.3}
Soit $D\in\operatorname{Div}^+(X/S)$ et considérons le sous-préschéma fermé $Y(D)$
de $X$ défini par l'Idéal $\mathscr J_X(D)$ de $\mathcal O_X$ (21.6.5); en vertu de
ce qui précède, pour tout $x\in Y(D)$, il y a un voisinage ouvert $U$ de $x$ dans $X$
et une section $t\in\mathscr T_{X/S}(U)$ telle que
$\mathscr J_X(D)|U=\mathcal O_U\cdot t$; puisque $x\in Y(D)$, l'image $t_{f(x)}$ de
$t$ dans $\Gamma(U\cap X_{f(x)},\mathcal O_{X_{f(x)}})$ appartient à l'idéal maximal
de $\mathcal O_{X_{f(x)},x}$ et en outre, par définition, pour tout $s\in S$, l'image
$t_s$ de $t$ dans $\Gamma(U\cap X_s,\mathcal O_{X_s})$ est une section
\emph{régulière}. On déduit donc de (11.3.8) et (19.2.4) que l'immersion canonique
$Y(D)\to X$ est \emph{transversalement régulière relativement à $S$ et de
codimension $1$} au point $x$. La réciproque étant immédiate, on voit qu'on peut
\emph{identifier canoniquement} les diviseurs positifs sur $X$ relatifs à $S$ aux
sous-préschémas fermés $Y$ de $X$ tels que l'injection canonique $Y\to X$ soit une
immersion transversalement régulière relativement à $S$ et de codimension $1$. Nous
ferons d'ordinaire cette identification.
\end{env}

\begin{proposition}[21.15.4]
\label{IV.21.15.4}
Soient $D$ un diviseur sur $X$, $\mathscr J_X(D)$ l'Idéal fractionnaire inversible
correspondant (21.2.5). Pour que $D\in\operatorname{Div}(X/S)$, il faut et il suffit
que, pour tout $x\in X$, on ait
$(\mathscr J_X(D)\cap\mathscr M_{X/S}^*)_x\ne0$ (ou, ce qui revient au même, que
l'on ait $\mathscr J_X(D)\subset\mathscr M_{X/S}$ et
$\mathscr J_X(D)^{-1}\subset\mathscr M_{X/S}$).
\end{proposition}

\begin{proof}
En effet, dire que $D\in\operatorname{Div}(X/S)$ signifie que pour tout $x\in X$, il
existe un voisinage ouvert $U$ de $x$ et une section
$f\in\Gamma(U,\mathscr M_{X/S}^*)$ tels que $D|U=\operatorname{div}_U(f)$; comme
$\mathscr J_X(D)|U$ est l'Idéal fractionnaire inversible $\mathcal O_Uf$, la
proposition en résulte aussitôt.
\end{proof}

\oldpage[IV-4]{331}
\begin{proposition}[21.15.5]
\label{IV.21.15.5}
Soient $D$ un diviseur sur $X$, $\mathcal O_X(D)$ l'Idéal fractionnaire inversible et
$s_D$ la section méromorphe régulière de $\mathcal O_X(D)$ définis canoniquement par
$D$ (21.2.8 et 21.2.9). Pour que $D\in\operatorname{Div}(X/S)$, il faut et il suffit
que $s_D\in\Gamma(X,\mathscr M_{X/S}(\mathcal O_X(D)))$.
\end{proposition}

\begin{proof}
En effet, si $U$ est un ouvert de $X$ tel que $D|U=\operatorname{div}_U(f)$, où
$f\in\Gamma(U,\mathscr M_X^*)$, dire que
$s_D\in\Gamma(U,\mathscr M_{X/S}(\mathcal O_X(D)))$ signifie, en vertu des
définitions (20.6.2), que $f^{-1}\in\Gamma(U,\mathscr M_{X/S})$, d'où la
proposition.
\end{proof}

L'interprétation des diviseurs sur $X$ à l'aide des classes
$\operatorname{cl}(\mathscr L,s)$ (21.2.11) permet donc d'interpréter les éléments de
$\operatorname{Div}(X/S)$ comme les couples (à isomorphisme près) $(\mathscr L,s)$
tels que $\mathscr L$ soit un $\mathcal O_X$-Module inversible et que $s$ soit une
section méromorphe de $\mathscr L$ au-dessus de $X$, \emph{régulière relativement à
$S$} (20.6.5, (iii)).

\begin{proposition}[21.15.6]
\label{IV.21.15.6}
Soit $D$ un diviseur sur $X$ relatif à $S$, et supposons que pour tout $x\in X$ tel
que $\operatorname{prof}(\mathcal O_{X_{f(x)},x})=1$, on ait $D_x\geq0$
(resp. $D_x=0$). Alors on a $D\geq0$ (resp. $D=0$).
\end{proposition}

\begin{proof}
Comme dans (21.1.8), on peut se borner au cas où $D=\operatorname{div}(\varphi)$, où
$\varphi$ est une fonction méromorphe régulière relative à $S$, et tout revient à voir
que si $D_x\geq0$ en tout point $x\in X$ tel que
$\operatorname{prof}(\mathcal O_{X_{f(x)},x})=1$, $\varphi$ est partout définie dans
$X$. Mais cette hypothèse signifie que, si $T=X-\operatorname{dom}(\varphi)$, on a
$\operatorname{prof}(\mathcal O_{X_{f(x)},x})\geq2$ pour tout $x\in T$, et il suffit
pour conclure d'appliquer (20.6.6).
\end{proof}

\begin{env}[21.15.7]
\label{IV.21.15.7}
Soit $X'$ un second $S$-préschéma plat et localement de présentation finie sur $S$, et
soit $g:X'\to X$ un $S$-morphisme. Si le $S$-morphisme $g$ est \emph{plat}, on sait
(21.4.5) que l'image réciproque par $g$ de tout diviseur sur $X$ est définie; si de
plus $D\in\operatorname{Div}(X/S)$, il résulte de la définition (21.15.2) et de
(20.6.8) que l'on a alors $g^*(D)\in\operatorname{Div}(X'/S)$.
\end{env}

\begin{env}[21.15.8]
\label{IV.21.15.8}
Soit $X'$ un second $S$-préschéma plat et localement de présentation finie sur $S$, et
soit $g:X'\to X$ un $S$-morphisme \emph{fini et plat}. Notons que $g$ est alors
nécessairement de présentation finie (1.4.3, 1.4.6 et 1.6.3), donc
$g_*(\mathcal O_{X'})$ est un $\mathcal O_X$-Module plat et de présentation finie, et
par suite \emph{localement libre} (2.1.12), autrement dit $g$ est un morphisme
localement libre (18.2.7); pour tout $s\in S$, le morphisme correspondant
$g_s:X'_s\to X_s$ est donc aussi fini et localement libre. On déduit alors de (21.5.2)
et (20.6.1) que pour tout ouvert $U$ de $X$ et toute section
$t'\in\Gamma(g^{-1}(U),\mathscr T_{X'/S})$, la norme $N_{X'/X}(t')$ appartient à
$\Gamma(U,\mathscr T_{X/S})$; le raisonnement de (21.5.3) prouve ensuite que pour tout
$\mathcal O_{X'}$-Module inversible $\mathscr L'$ et toute section méromorphe $u'$ de
$\mathscr L'$ au-dessus de $X'$, \emph{régulière relativement à $S$}, la norme
$N_{X'/X}(u')$ est une section méromorphe de
$\mathscr L=N_{X'/X}(\mathscr L')$ au-dessus de $X$, régulière relativement à $S$.
L'interprétation des diviseurs relatifs à $S$ donnée dans (21.15.5) et la définition de
l'image directe d'un diviseur (21.5.5) prouvent alors que \emph{pour tout diviseur}
$D'\in\operatorname{Div}(X'/S)$, on a $g_*(D')\in\operatorname{Div}(X/S)$.
\end{env}

\begin{env}[21.15.9]
\EGAReaderTarget{ega4.c04.statement.21-15-9}{IV.21.15.9}
Considérons enfin un morphisme quelconque $S'\to S$, et (sous les hypothèses de
(21.15.1)) posons $X'=X\times_S S'$, qui est plat et localement de présentation finie
sur $S'$; si $p:X'\to X$ est la projection canonique, on a vu (20.6.9) qu'on a un
homomorphisme canonique
$p^*(\mathscr M_{X/S})\to\mathscr M_{X'/S'}$, qui transforme évidemment toute
section de $\mathscr M_{X/S}$ au-dessus d'un ouvert $U$, \emph{régulière relativement
à $S$}, en une section de $\mathscr M_{X'/S'}$ au-dessus de $p^{-1}(U)$, régulière
relativement à $S$ (20.6.5, (iii)); on conclut alors de la
\oldpage[IV-4]{332}
définition (21.15.2) et de l'exactitude à droite du foncteur $p^*$, que
l'homomorphisme précédent définit par passage aux quotients un homomorphisme canonique
\[
\label{IV.21.15.9.1}
  \operatorname{Div}(X/S)\longrightarrow\operatorname{Div}(X'/S')
  \tag{21.15.9.1}
\]
qui transforme évidemment les éléments de $\operatorname{Div}^+(X/S)$ en éléments de
$\operatorname{Div}^+(X'/S')$. On pose encore
$\operatorname{Div}(X'/S')=\operatorname{Div}_{X/S}(S')$ (resp.
$\operatorname{Div}^+(X'/S')=\operatorname{Div}_{X/S}^+(S')$), et on voit aussitôt
qu'on a ainsi défini deux foncteurs contravariants
\[
  \operatorname{Div}_{X/S}:\mathbf{Sch}_{/S}\longrightarrow\mathbf{Ab},\qquad
  \operatorname{Div}_{X/S}^+:\mathbf{Sch}_{/S}\longrightarrow\mathbf{Ens}
\]
de la catégorie des $S$-préschémas dans celle des groupes commutatifs (resp. des
ensembles). On verra plus loin (chap.~VI) des cas importants où le foncteur
$\operatorname{Div}_{X/S}^+$ est \emph{représentable} $(\mathbf 0_{\mathrm{III}},
8.1.8)$.

Pour tout diviseur $D\in\operatorname{Div}(X/S)$, l'image de $D$ par
l'homomorphisme (21.15.9.1) n'est autre d'ailleurs que l'\emph{image réciproque}
$p^*(D)$ (au sens de (21.4.2)) : l'existence de cette image réciproque et l'assertion
précédente sont en effet des conséquences immédiates de (20.6.5, (iii)) et (20.6.9).

Les éléments de $\operatorname{Div}(X'/S')$ s'appellent souvent, par abus de langage,
« \emph{familles de diviseurs sur $X$ relatifs à $S$, paramétrées par le
$S$-préschéma $S'$} »; on utilise surtout cette terminologie lorsqu'il s'agit de
diviseurs \emph{positifs}.
\end{env}

% FIN SOURCE ega4/ega4-21c-fr.tex

%
% DEBUT SOURCE ega4/ega4-backmatter-bibliography-fr.tex
% Transcription diplomatique française de EGA IV, quatrième partie, p. 333.
\clearpage
\oldpage[IV]{333}
\phantomsection
\addcontentsline{toc}{section}{Bibliographie}
\section*{BIBLIOGRAPHIE}

\begin{thebibliography}{44}

\bibitem[41]{IV-41}
A.~Grothendieck, \emph{Séminaire de Géométrie algébrique, 1962,
Cohomologie locale des faisceaux cohérents, et théorèmes de Lefschetz locaux
et globaux} (Institut des Hautes Études scientifiques, Bures-sur-Yvette).

\bibitem[42]{IV-42}
M.~Demazure et A.~Grothendieck, \emph{Séminaire de Géométrie algébrique,
1963--1964, Schémas en groupes} (Institut des Hautes Études scientifiques,
Bures-sur-Yvette).

\bibitem[43]{IV-43}
M.~Artin et A.~Grothendieck, \emph{Séminaire de Géométrie algébrique
1963--1964, Cohomologie étale des schémas} (Institut des Hautes Études
scientifiques, Bures-sur-Yvette).

\bibitem[44]{IV-44}
C.~S.~Seshadri, \emph{Quotient space by an abelian variety},
\emph{Mathematische Annalen}, t.~152 (1962), p.~185--194.

\end{thebibliography}

% FIN SOURCE ega4/ega4-backmatter-bibliography-fr.tex

%
% DEBUT SOURCE ega4/ega4-backmatter-notation-fr.tex
% Transcription diplomatique française de EGA IV, quatrième partie, pp. 334--335.
\clearpage
\oldpage[IV]{334}
\phantomsection
\addcontentsline{toc}{section}{Index des notations}
\section*{INDEX DES NOTATIONS}
\markboth{INDEX DES NOTATIONS}{INDEX DES NOTATIONS}

\begingroup
\small
\setlength{\parindent}{0pt}
\setlength{\parskip}{1.5pt}

$Y_f^{(n)}$, $Y^{(n)}$, $\mathcal{G}r_{\bullet}(f)$,
$\mathcal{G}r_n(f)$, $\mathcal{N}_{Y/X}$
($f:Y\to X$ morphisme d'espaces annelés) : 16.1.2.

$\mathcal{O}_{Y^{(\infty)}}$ : 16.1.11.

$\mathcal{P}^{n}_{f}$, $\mathcal{P}^{n}_{X/S}$,
$\mathcal{P}^{\infty}_{f}$, $\mathcal{P}^{\infty}_{X/S}$,
$\Omega^{1}_{f}$, $\Omega^{1}_{X/S}$
($f:X\to S$ morphisme de préschémas) : 16.3.1.

$d^{n}_{f}$, $d^{n}_{X/S}$, $d^{\infty}_{f}$,
$d^{\infty}_{X/S}$, $d^{n}$, $d^{\infty}$, $dt$, $d_{X/S}(t)$
($f:X\to S$ morphisme de préschémas) : 16.3.6.

$P^{n}(u)$, $\Gr_{\bullet}(u)$, $\Gr_{1}(u)$
($u:X'\to X$ morphisme de préschémas) : 16.4.3.

$\operatorname{Der}_{S}(\mathcal{O}_{X},\mathcal{F})$ : 16.5.1.

$\mathcal{D}er_{S}(\mathcal{O}_{X},\mathcal{F})$ : 16.5.4.

$\mathcal{T}_{X/S}$ : 16.5.7.

$D_i$, $\partial/\partial s_i$ : 16.5.8.

$T_{X/S}$ : 16.5.12.

$T_{X/S}(x)$, $T_Y(s)$ : 16.5.13.

$\Omega^p_{X/S}$, $\Omega^\bullet_{X/S}$, $\Omega^p_{B/A}$ : 16.6.1.

$d$, $d_{X/S}$ : 16.6.3.

$\mathcal{P}^n_{X/S}(\mathcal{F})$ : 16.7.2.

$d^n_{X/S,\mathcal{F}}$, $d^n_{X/S}$ : 16.7.5.

$\mathcal{P}^{\infty}_{X/S}(\mathcal{F})$,
$d^{\infty}_{X/S,\mathcal{F}}$ : 16.7.7.

$\operatorname{Diff}^{n}_{X/S}(\mathcal{F},\mathcal{G})$,
$\operatorname{Diff}^{n}_{X/S}$,
$\mathcal{D}iff^{n}_{X/S}(\mathcal{F},\mathcal{G})$ : 16.8.3.

$\mathcal{D}iff^{n}_{X/S}$, $\langle t,D\rangle$ : 16.8.4.

$aD$ ($D$ opérateur différentiel, $a\in\Gamma(U,\mathcal{O}_X)$) : 16.8.5.

$\operatorname{Diff}_{X/S}(\mathcal{F},\mathcal{G})$,
$\mathcal{D}iff_{X/S}(\mathcal{F},\mathcal{G})$ : 16.8.7.

$\mathcal{D}iff_{X/S}$ : 16.8.10.

$|\mathbf p|$, $\mathbf p!$, $\binom{\mathbf p}{\mathbf q}$,
$\mathbf z^{\mathbf p}$, $\boldsymbol\zeta^{\mathbf p}$
($\mathbf p,\mathbf q$ multi-indices) : 16.11.1.

$\dim_x(f)$ ($f:X\to Y$ morphisme localement de type fini, $x\in X$) : 17.10.1.

$\operatorname{Tr}_{B/A}$, $\operatorname{astr}_{B/A}$
($B$ $A$-algèbre finie et libre) : 18.2.1.

$\operatorname{Tr}_{\mathcal B/\mathcal O_X,U}$,
$\operatorname{astr}_{\mathcal B/\mathcal O_X,U}$,
$\operatorname{Tr}_{\mathcal B/\mathcal O_X}$,
$\operatorname{astr}_{\mathcal B/\mathcal O_X}$
($X$ espace annelé, $\mathcal B$ $\mathcal O_X$-Algèbre localement libre de rang fini) : 18.2.2.

$\operatorname{Tr}_f$ ($f$ morphisme fini et étale) : 18.2.6.

$d_{X/Y}$, $\mathcal D_{X/Y}$ : 18.2.7.

${}^{t}V_{\mathcal E}$ : 18.5.1.

$Al(X,\mathcal E,\mathcal J)$, $Al(\mathcal E,\mathcal J)$ : 18.5.2.

$Of(S)$, $Of(X)$, $Id(\mathcal B')$ : 18.5.3.

$\operatorname{Hom.loc}(A,B)$ : 18.6.1.

${}^{h}A$ : 18.6.5 et 18.6.7.

$\operatorname{Hom.sloc}(A,B)$ : 18.6.7.

${}^{hs}A_{(i)}$, ${}^{hs}A$ : 18.8.7 et 18.8.9.

$\operatorname{codim}^{t}_{y}(Y,X)$ : 19.1.3.

$\operatorname{prof}_{T,t}(\mathcal F)$ : 19.9.1.

$\mathcal O_X[\mathcal S^{-1}]$, $\mathcal F[\mathcal S^{-1}]$ : 20.1.1.

$\mathcal S(\mathcal O_X)$, $\mathcal M_X$, $M(X)$,
$\mathcal M_X(\mathcal F)$, $M(X,\mathcal F)$ : 20.1.3.

$\operatorname{dom}(\varphi)$ (fonction méromorphe) : 20.1.4.

$\operatorname{dom}(u)$ ($u$ section méromorphe de $\mathcal F$) : 20.1.7.

$(\mathcal M_X(\mathcal L))^{*}$
($\mathcal L$ $\mathcal O_X$-Module inversible) : 20.1.8.

$s^{-1}$ ($s$ section méromorphe d'un $\mathcal O_X$-Module inversible) : 20.1.10.

$\mathcal S_f(U)$, $\mathcal S_f$, $\mathcal M_f$
($f$ morphisme d'espaces annelés), $\varphi\circ f$ ($\varphi$ fonction méromorphe),
$u\circ f$ ($u$ section méromorphe de $\mathcal F$), $\mathcal M_f(\mathcal F)$ : 20.1.11.

\oldpage[IV]{335}
$\operatorname{Ps.hom}(X,Y)$, $\operatorname{Ps.hom}_S(X,Y)$ : 20.2.1.

$\omega|U$ ($\omega$ pseudo-morphisme),
$\mathcal P\!s.\operatorname{hom}(X,Y)$,
$\mathcal P\!s.\operatorname{hom}_S(X,Y)$ : 20.2.2.

$\operatorname{dom}_S(\omega)$, $\operatorname{dom}(\omega)$
($\omega$ pseudo-morphisme) : 20.2.3.

$\mathcal M'_X$, $M'(X)$ : 20.2.8.

$f\circ\omega$ ($f$ morphisme, $\omega$ pseudo-morphisme) : 20.3.1.

$\omega\circ f$ ($f$ morphisme, $\omega$ pseudo-morphisme) : 20.3.2.

$\operatorname{Ps.hom}_S(X,Y)^f$,
$\mathcal P\!s.\operatorname{hom}_S(X,Y)^f$ : 20.3.2.

$\Gamma_\omega$ ($\omega$ pseudo-morphisme) : 20.4.1.

$p^{-1}$ ($p:\Gamma_\omega\to X$ restriction de la première projection) : 20.4.2.

$\omega(M)$ ($\omega$ pseudo-morphisme, $M$ partie de $X$) : 20.4.2.

$\operatorname{Ps.hom}_{X/S}(X,Y)$ : 20.5.1.

$\mathcal P\!s.\operatorname{hom}_{X/S}(X,Y)$ : 20.5.3.

$\mathcal M'_{X/S}$, $M'(X/S)$ : 20.5.4.

$\mathcal S_{X/S}(U)$, $\mathcal S_{X/S}$, $\mathcal M_{X/S}$,
$M(X/S)$ : 20.6.1.

$\mathcal M_{X/S}(\mathcal F)$, $M(X/S,\mathcal F)$ : 20.6.2.

$\mathcal S_f(U)$, $\mathcal S_f$, $\mathcal M_{X/S,f}$ : 20.6.7.

$\mathcal{D}iv_X$, $\operatorname{Div}(X)$, $\operatorname{div}(f)$,
$\operatorname{div}_X(f)$ ($X$ espace annelé) : 21.1.2.

$\operatorname{Supp}(D)$ ($D$ diviseur) : 21.1.2.

$\operatorname{Div}(A)$ ($A$ anneau) : 21.2.1.

$\operatorname{div}(s)$ ($s$ section d'un Module inversible) : 21.1.4.

$\mathcal{D}iv_X^+$, $\operatorname{Div}^+(X)$, $D\leq D'$
($D,D'$ diviseurs) : 21.1.6.

$\operatorname{Id.inv}(X)$, $\sh{J}d.inv_X$ : 21.2.4.

$\mathcal J_U$, $\mathcal J_X(D)$ : 21.2.5.

$\operatorname{div}(\mathcal I)$ ($\mathcal I$ Idéal fractionnaire) : 21.2.6.

$\mathcal O_X(D)$ ($D$ diviseur) : 21.2.8.

$s_D$ ($D$ diviseur) : 21.2.9.

$D(X)$, $\operatorname{cl}(\mathcal L,s)$ : 21.2.10.

$Y(D)$ ($D$ diviseur positif) : 21.2.12.

$\operatorname{Div.princ}(X)$ : 21.3.1.

$\operatorname{Pic}(X)$, $\operatorname{cl}(\mathcal L)$, $l$, $l_X$,
$\operatorname{Pic}(u)$ : 21.3.2.

$f^*(D)$ ($f$ morphisme d'espaces annelés, $D$ diviseur),
$\operatorname{Div}^f(X)$ : 21.4.2.

$\mathcal M_f^{**}$, $\mathcal{D}iv_X^f$ : 21.4.3.

$N_{X'/X}(u')$, $N(u')$
($u'$ section méromorphe d'un $\mathcal O_{X'}$-Module inversible) : 21.5.3.

$f_*(D')$, $N_{X'/X}(D')$
($f$ morphisme fini, $D'$ diviseur sur $X'$) : 21.5.5.

$\mathfrak I(X)$, $\mathfrak Z(X)$, $\mathfrak Z^+(X)$
($X$ préschéma localement noethérien) : 21.6.1.

$\operatorname{mult}_x(Z)$, $\operatorname{Supp}(Z)$,
$\dim(Z)$, $\operatorname{codim}(Z)$ ($Z$ cycle) : 21.6.1.

$X^{(p)}$, $\mathfrak Z^p(X)$, $\mathfrak Z^{p+}(X)$ : 21.6.2.

$Z|U$ ($Z$ cycle), $\mathscr Z_X$, $\mathscr Z_X^p$,
$\mathscr Z_X^+$, $\mathscr Z_X^{p+}$ : 21.6.3.

$c$ : 21.6.4.

$\operatorname{cyc}(D)$ ($D$ diviseur) : 21.6.5 et 21.6.7.

$\operatorname{mult}_x(D)$, $\operatorname{mult}_{\{x\}}(D)$
($D$ diviseur), $o_x(f)$ ($f$ fonction méromorphe régulière) : 21.6.7.

$\mathrm Z^+(f)$, $\mathrm Z^-(f)$, $\mathfrak Z^1_{\mathrm{princ}}(X)$,
$\mathfrak{Cl}(X)$ : 21.6.7.

$Y(C)$, $Y'(C)$, $\sh{J}_X(C)$, $\sh{J}(C)$
($C$ cycle 1-codimensionnel positif) : 21.7.1.

$f^*(Z)$ ($Z$ cycle 1-codimensionnel, $f$ morphisme) : 21.10.3 et 21.10.11.

$\Gamma_T(\mathscr Z_X^1)$ : 21.10.3.

$f_*(Z')$ ($Z'$ cycle 1-codimensionnel, $f$ morphisme fini) : 21.10.14.

$\Lambda^{\max}\mathcal E$ ($\mathcal E$ module localement libre) : 21.11.1.

$U'$, $\operatorname{Aff}(U/X)$ : 21.12.1.

$\operatorname{Daf}(U/X)$ : 21.12.5.

$\mathcal{D}iv_{X/S}$, $\operatorname{Div}(X/S)$ : 21.15.2.

$\mathcal{D}iv_{X/S}^+$, $\operatorname{Div}^+(X/S)$ : 21.15.3.

$\mathcal O_{X/S}^{\mathrm{reg}}$ : 21.15.3.

$\operatorname{Div}_{X/S}(S')$, $\operatorname{Div}_{X/S}^+(S')$
($S'$ $S$-préschéma) : 21.15.9.

\endgroup

% FIN SOURCE ega4/ega4-backmatter-notation-fr.tex

%
% DEBUT SOURCE ega4/ega4-backmatter-terminology-fr.tex
% Transcription diplomatique française de EGA IV, quatrième partie, pp. 336--339.
% Les entrées et leurs renvois reproduisent l'index imprimé.
\clearpage
\oldpage[IV]{336}
\phantomsection
\addcontentsline{toc}{section}{Index terminologique}
\section*{INDEX TERMINOLOGIQUE}
\markboth{INDEX TERMINOLOGIQUE}{INDEX TERMINOLOGIQUE}

\begingroup
\small
\setlength{\parindent}{0pt}
\setlength{\parskip}{2pt}
\newcommand{\EGAIVterm}[2]{\textbf{#1}: #2.\par}
\newcommand{\EGAIVidxref}[1]{\hyperref[IV.#1]{#1}}

\EGAIVterm{Algèbre étale}{{\hyperlink{ega4.c04.subsec.17-3-2}{17.3.2}}}
\EGAIVterm{Algèbre essentiellement étale}{{\hyperlink{ega4.c04.statement.18-6-1}{18.6.1}}}
\EGAIVterm{$\sh{O}_X$-Algèbre finie et étale}{{\hyperlink{ega4.c04.prop.18-2-3}{18.2.3}}}
\EGAIVterm{Algèbre lisse}{{\hyperlink{ega4.c04.subsec.17-3-2}{17.3.2}}}
\EGAIVterm{Algèbre non ramifiée}{{\hyperlink{ega4.c04.subsec.17-3-2}{17.3.2}}}
\EGAIVterm{$K$-algèbre non ramifiée sur $Y$, sur $A$}{{\hyperlink{ega4.c04.rem.18-10-11}{18.10.11}}}
\EGAIVterm{Algèbre strictement essentiellement étale}{{\hyperlink{ega4.c04.lem.18-6-2}{18.6.2}}}
\EGAIVterm{Anneau d'intersection complète, anneau d'intersection complète absolue}{{\hyperlink{ega4.c04.def.19-3-1}{19.3.1}}}
\EGAIVterm{Anneau hensélien}{{\hyperlink{ega4.c04.def.18-5-8}{18.5.8}}}
\EGAIVterm{Anneau local parafactoriel}{{\hyperlink{ega4.c04.def.21-13-7}{21.13.7}}}
\EGAIVterm{Anneau strictement local}{{\hyperlink{ega4.c04.def.18-8-2}{18.8.2}}}
\EGAIVterm{Application linéaire tangente}{{\hyperlink{ega4.c04.statement.16-5-13}{16.5.13}}}
\EGAIVterm{Application rationnelle stricte}{{\hyperlink{ega4.c04.statement.20-2-1}{20.2.1}}}
\EGAIVterm{Codimension d'un cycle}{{\hyperlink{ega4.c04.statement.21-6-1}{21.6.1}}}
\EGAIVterm{Codimension transversale en un point d'un sous-préschéma}{{\hyperlink{ega4.c04.def.19-1-3}{19.1.3}}}
\EGAIVterm{Codimension transversale d'un morphisme d'immersion en un point}{{\hyperlink{ega4.c04.def.19-1-3}{19.1.3}}}
\EGAIVterm{Couple de morphismes transversaux en un point}{{\hyperlink{ega4.c04.cor.17-13-6}{17.13.6}}}
\EGAIVterm{Couple hensélien}{{\hyperlink{ega4.c04.def.18-5-5}{18.5.5}}}
\EGAIVterm{Couple parafactoriel}{{\hyperlink{ega4.c04.def.21-13-1}{21.13.1}}}
\EGAIVterm{Cycle}{{\hyperlink{ega4.c04.statement.21-6-1}{21.6.1}}}
\EGAIVterm{Cycle 1-codimensionnel à coefficients rationnels}{{\hyperlink{ega4.c04.statement.21-10-9}{21.10.9}}}
\EGAIVterm{Cycle $p$-codimensionnel}{{\hyperlink{ega4.c04.statement.21-6-2}{21.6.2}}}
\EGAIVterm{Cycle des pôles, cycle des zéros d'une fonction méromorphe}{{\hyperlink{ega4.c04.statement.21-6-7}{21.6.7}}}
\EGAIVterm{Cycle linéairement équivalent à 0}{{\hyperlink{ega4.c04.statement.21-6-7}{21.6.7}}}
\EGAIVterm{Cycle localement principal}{{\hyperlink{ega4.c04.statement.21-6-7}{21.6.7}}}
\EGAIVterm{Cycle polaire d'une fonction méromorphe}{{\hyperlink{ega4.c04.statement.21-6-7}{21.6.7}}}
\EGAIVterm{Cycle premier}{{\hyperlink{ega4.c04.statement.21-6-1}{21.6.1}}}
\EGAIVterm{Cycle principal}{{\hyperlink{ega4.c04.statement.21-6-7}{21.6.7}}}
\EGAIVterm{Cycle principal en un point}{{\hyperlink{ega4.c04.statement.21-6-7}{21.6.7}}}
\EGAIVterm{Cycle purement de codimension $p$}{{\hyperlink{ega4.c04.statement.21-6-2}{21.6.2}}}
\EGAIVterm{Défaut d'affinité d'un ouvert dans un préschéma}{{\hyperlink{ega4.c04.statement.21-12-5}{21.12.5}}}
\EGAIVterm{$S$-dérivation, $(X/S)$-dérivation, $f$-dérivation}{{\hyperlink{ega4.c04.statement.16-5-1}{16.5.1}}}
\EGAIVterm{Différentielle d'une section de $\sh{O}_X$}{{\hyperlink{ega4.c04.def.16-3-6}{16.3.6}}}
\EGAIVterm{Différentielle extérieure}{{\hyperlink{ega4.c04.def.16-6-3}{16.6.3}}}
\EGAIVterm{Dimension d'un cycle}{{\hyperlink{ega4.c04.statement.21-6-1}{21.6.1}}}
\EGAIVterm{Dimension relative d'un morphisme en un point, d'un $Y$-préschéma en un point}{{\hyperlink{ega4.c04.def.17-10-1}{17.10.1}}}
\EGAIVterm{Discriminant d'un $Y$-préschéma fini et localement libre}{{\hyperlink{ega4.c04.item.18-2-7}{18.2.7}}}
\EGAIVterm{Diviseur sur un espace annelé}{{\hyperlink{ega4.c04.statement.21-1-2}{21.1.2}}}
\EGAIVterm{Diviseur d'une fonction méromorphe}{{\hyperlink{ega4.c04.statement.21-1-2}{21.1.2}}}
\EGAIVterm{Diviseur d'une section méromorphe d'un Module inversible}{{\hyperlink{ega4.c04.statement.21-1-4}{21.1.4}}}
\EGAIVterm{Diviseur positif}{{\hyperlink{ega4.c04.def.21-1-6}{21.1.6}}}
\EGAIVterm{Diviseur principal}{{\hyperlink{ega4.c04.statement.21-3-1}{21.3.1}}}
\EGAIVterm{Diviseur relativement à $S$, diviseur transversal à $f$}{{\hyperlink{ega4.c04.def.21-15-2}{21.15.2}}}
\EGAIVterm{Diviseurs linéairement équivalents}{{\hyperlink{ega4.c04.statement.21-3-1}{21.3.1}}}

\oldpage[IV]{337}
\EGAIVterm{Domaine de définition d'une fonction méromorphe}{{\hyperlink{ega4.c04.statement.20-1-4}{20.1.4}}}
\EGAIVterm{Domaine de définition d'un pseudo-morphisme}{{\hyperlink{ega4.c04.statement.20-2-3}{20.2.3}}}
\EGAIVterm{Domaine de définition d'une section méromorphe d'un Module}{{\hyperlink{ega4.c04.statement.20-1-7}{20.1.7}}}
\EGAIVterm{Enveloppe affine d'un $X$-préschéma}{{\hyperlink{ega4.c04.statement.21-12-1}{21.12.1}}}
\EGAIVterm{Espace tangent à un préschéma en un point}{{\hyperlink{ega4.c04.statement.16-5-13}{16.5.13}}}
\EGAIVterm{Faisceau conormal}{{\hyperlink{ega4.c04.def.16-1-2}{16.1.2}}}
\EGAIVterm{Faisceau d'anneaux de fractions à dénominateurs dans $\sh{S}$}{{\hyperlink{ega4.c04.statement.20-1-1}{20.1.1}}}
\EGAIVterm{Faisceau d'anneaux gradués associé à un morphisme}{{\hyperlink{ega4.c04.def.16-1-2}{16.1.2}}}
\EGAIVterm{Faisceau de modules de fractions à dénominateurs dans $\sh{S}$}{{\hyperlink{ega4.c04.statement.20-1-2}{20.1.2}}}
\EGAIVterm{Faisceau des $S$-dérivations de $\sh{O}_X$}{{\hyperlink{ega4.c04.statement.16-5-7}{16.5.7}}}
\EGAIVterm{Faisceau des $S$-dérivations de $\sh{O}_X$ dans $\sh{F}$}{{\hyperlink{ega4.c04.statement.16-5-4}{16.5.4}}}
\EGAIVterm{Faisceau des $p$-différentielles de $X$ relativement à $S$}{{\hyperlink{ega4.c04.statement.16-6-1}{16.6.1}}}
\EGAIVterm{Faisceau des diviseurs sur un espace annelé}{{\hyperlink{ega4.c04.statement.21-1-2}{21.1.2}}}
\EGAIVterm{Faisceau des diviseurs relativement à $S$, ou transversaux à $f$}{{\hyperlink{ega4.c04.def.21-15-2}{21.15.2}}}
\EGAIVterm{Faisceau des germes de cycles 1-codimensionnels à coefficients dans $\Lambda$}{{\hyperlink{ega4.c04.statement.21-10-9}{21.10.9}}}
\EGAIVterm{Faisceau des germes de fonctions méromorphes sur $X$}{{\hyperlink{ega4.c04.statement.20-1-4}{20.1.4}}}
\EGAIVterm{Faisceau des germes de fonctions méromorphes relativement à $S$}{{\hyperlink{ega4.c04.statement.20-6-1}{20.6.1}}}
\EGAIVterm{Faisceau des germes de sections méromorphes d'un Module}{{\hyperlink{ega4.c04.statement.20-1-3}{20.1.3}}}
\EGAIVterm{Faisceau des parties principales d'ordre $n$ du $S$-préschéma $X$}{{\hyperlink{ega4.c04.def.16-3-1}{16.3.1}}}
\EGAIVterm{Faisceau formellement principal homogène (ou pseudo-torseur) sous un faisceau de groupes}{{\hyperlink{ega4.c04.statement.16-5-15}{16.5.15}}}
\EGAIVterm{Faisceau principal homogène (ou torseur) sous un faisceau de groupes}{{\hyperlink{ega4.c04.statement.16-5-15}{16.5.15}}}
\EGAIVterm{Faisceau principal homogène (ou torseur) trivial sous un faisceau de groupes}{{\hyperlink{ega4.c04.statement.16-5-15}{16.5.15}}}
\EGAIVterm{Faisceau tangent de $X$ relativement à $S$}{{\hyperlink{ega4.c04.statement.16-5-7}{16.5.7}}}
\EGAIVterm{Famille de diviseurs sur $X$ relativement à $S$, paramétrés par $S'$}{{\hyperlink{ega4.c04.statement.21-15-9}{21.15.9}}}
\EGAIVterm{Famille de morphismes transversaux en un point}{{\hyperlink{ega4.c04.cor.17-13-10}{17.13.10}}}
\EGAIVterm{Famille de sous-préschémas se coupant transversalement en un point}{{\hyperlink{ega4.c04.statement.17-13-11}{17.13.11}}}
\EGAIVterm{Fibré tangent de $X$ relativement à $S$}{{\hyperlink{ega4.c04.statement.16-5-12}{16.5.12}}}
\EGAIVterm{Fonction méromorphe sur un espace annelé}{{\hyperlink{ega4.c04.statement.20-1-3}{20.1.3}}}
\EGAIVterm{Fonction méromorphe définie en un point}{{\hyperlink{ega4.c04.statement.20-1-4}{20.1.4}}}
\EGAIVterm{Fonction méromorphe régulière}{{\hyperlink{ega4.c04.statement.20-1-8}{20.1.8}}}
\EGAIVterm{Fonction méromorphe régulière relativement à $S$}{{\hyperlink{ega4.c04.statement.20-6-5}{20.6.5}}}
\EGAIVterm{Fonction méromorphe relativement à $S$}{{\hyperlink{ega4.c04.statement.20-6-1}{20.6.1}}}
\EGAIVterm{Forme trace}{{\hyperlink{ega4.c04.statement.18-2-1}{18.2.1}}}
\EGAIVterm{Graphe d'un pseudo-$S$-morphisme}{{\hyperlink{ega4.c04.statement.20-4-1}{20.4.1}}}
\EGAIVterm{Groupe des classes de cycles 1-codimensionnels}{{\hyperlink{ega4.c04.statement.21-6-7}{21.6.7}}}
\EGAIVterm{Hensélisé d'un anneau local}{{\hyperlink{ega4.c04.def.18-6-5}{18.6.5}}}
\EGAIVterm{Hensélisé d'un anneau semi-local}{{\hyperlink{ega4.c04.statement.18-6-7}{18.6.7}}}
\EGAIVterm{Hensélisé strict d'un anneau local}{{\hyperlink{ega4.c04.def.18-8-7}{18.8.7}}}
\EGAIVterm{Hensélisé strict d'un anneau semi-local}{{\hyperlink{ega4.c04.statement.18-8-9}{18.8.9}}}
\EGAIVterm{Homomorphisme associé à la trace}{{\hyperlink{ega4.c04.statement.18-2-2}{18.2.2}}}
\EGAIVterm{Homomorphisme régulier de $\sh{O}_X$-Modules}{{\hyperlink{ega4.c04.prop.19-4-11}{19.4.11}}}
\EGAIVterm{Homomorphisme semi-local}{{\hyperlink{ega4.c04.statement.18-6-7}{18.6.7}}}
\EGAIVterm{Homomorphisme trace}{{\hyperlink{ega4.c04.statement.18-2-2}{18.2.2}}}
\EGAIVterm{Idéal des dénominateurs d'une section méromorphe}{{\hyperlink{ega4.c04.statement.20-2-14}{20.2.14}}}
\EGAIVterm{Idéal discriminant d'un $Y$-préschéma fini et localement libre}{{\hyperlink{ega4.c04.item.18-2-7}{18.2.7}}}
\EGAIVterm{Idéal entier}{{\hyperlink{ega4.c04.statement.21-2-7}{21.2.7}}}
\EGAIVterm{Idéal fractionnaire}{{\hyperlink{ega4.c04.statement.21-2-1}{21.2.1}}}
\EGAIVterm{Idéal fractionnaire inversible}{{\hyperlink{ega4.c04.statement.21-2-1}{21.2.1}}}
\EGAIVterm{Idéal quasi-régulier d'un faisceau d'anneaux}{{\hyperlink{ega4.c04.def.16-9-1}{16.9.1}}}
\EGAIVterm{Idéal quasi-régulier d'un anneau}{{\hyperlink{ega4.c04.statement.16-9-7}{16.9.7}}}
\EGAIVterm{Idéal régulier d'un faisceau d'anneaux}{{\hyperlink{ega4.c04.def.16-9-1}{16.9.1}}}
\EGAIVterm{Idéal régulier d'un anneau}{{\hyperlink{ega4.c04.statement.16-9-7}{16.9.7}}}
\EGAIVterm{Idéal transversalement régulier en un point d'un préschéma}{{\hyperlink{ega4.c04.def.19-2-1}{19.2.1}}}

\oldpage[IV]{338}
\EGAIVterm{Idéal transversalement régulier d'une algèbre}{{\hyperlink{ega4.c04.def.19-2-1}{19.2.1}}}
\EGAIVterm{Image directe d'un cycle 1-codimensionnel}{{\hyperlink{ega4.c04.statement.21-10-14}{21.10.14}}}
\EGAIVterm{Image directe d'un diviseur}{{\hyperlink{ega4.c04.def.21-5-5}{21.5.5}}}
\EGAIVterm{Image réciproque d'un cycle 1-codimensionnel}{{\hyperlink{ega4.c04.statement.21-10-3}{21.10.3}} et {\hyperlink{ega4.c04.statement.21-10-11}{21.10.11}}}
\EGAIVterm{Image réciproque d'un diviseur}{{\hyperlink{ega4.c04.def.21-4-2}{21.4.2}}}
\EGAIVterm{Image réciproque d'une fonction méromorphe}{{\hyperlink{ega4.c04.statement.20-1-11}{20.1.11}}}
\EGAIVterm{Image réciproque d'une fonction méromorphe relative à $S$}{{\hyperlink{ega4.c04.statement.20-6-7}{20.6.7}}}
\EGAIVterm{Immersion quasi-régulière}{{\hyperlink{ega4.c04.def.16-9-2}{16.9.2}}}
\EGAIVterm{Immersion régulière}{{\hyperlink{ega4.c04.def.16-9-2}{16.9.2}}}
\EGAIVterm{Immersion régulière en un point}{{\hyperlink{ega4.c04.prop.16-9-10}{16.9.10}}}
\EGAIVterm{Immersion transversalement régulière en un point relativement à $S$, immersion transversalement régulière relativement à $S$}{{\hyperlink{ega4.c04.def.19-2-2}{19.2.2}}}
\EGAIVterm{Intersection complète en un point relativement à $S$, intersection complète relativement à $S$ (pour un préschéma plat et localement de présentation finie)}{{\hyperlink{ega4.c04.def.19-3-6}{19.3.6}}}
\EGAIVterm{$n$-ème invariant normal d'un morphisme}{{\hyperlink{ega4.c04.def.16-1-2}{16.1.2}}}
\EGAIVterm{Invariant normal d'ordre infini d'un morphisme}{{\hyperlink{ega4.c04.rem.16-1-11}{16.1.11}}}
\EGAIVterm{$\sh{O}_X$-Module des 1-différentielles de $f$ (ou de $X$ par rapport à $S$, ou du $S$-préschéma $X$)}{{\hyperlink{ega4.c04.def.16-3-1}{16.3.1}}}
\EGAIVterm{$\sh{O}_X$-Module sans torsion relativement à $S$}{{\hyperlink{ega4.c04.statement.20-6-2}{20.6.2}}}
\EGAIVterm{$\sh{O}_X$-Module sans torsion}{{\hyperlink{ega4.c04.statement.20-1-5}{20.1.5}}}
\EGAIVterm{Morphisme différentiellement lisse}{{\hyperlink{ega4.c04.def.16-10-1}{16.10.1}}}
\EGAIVterm{Morphisme différentiellement lisse jusqu'à l'ordre $r$}{{\hyperlink{ega4.c04.statement.16-11-3}{16.11.3}}}
\EGAIVterm{Morphisme (plat) d'intersection complète}{{\hyperlink{ega4.c04.def.19-3-6}{19.3.6}}}
\EGAIVterm{Morphisme essentiellement propre}{{\hyperlink{ega4.c04.rem.18-10-20}{18.10.20}}}
\EGAIVterm{Morphisme étale}{{\hyperlink{ega4.c04.def.17-3-1}{17.3.1}}}
\EGAIVterm{Morphisme étale en un point}{{\hyperlink{ega4.c04.def.17-3-7}{17.3.7}}}
\EGAIVterm{Morphisme fini localement libre de rang $n$, morphisme fini localement libre}{{\hyperlink{ega4.c04.item.18-2-7}{18.2.7}}}
\EGAIVterm{Morphisme formellement étale, formellement lisse, formellement net, formellement non ramifié}{{\hyperlink{ega4.c04.def.17-1-1}{17.1.1}}}
\EGAIVterm{Morphisme lisse}{{\hyperlink{ega4.c04.def.17-3-1}{17.3.1}}}
\EGAIVterm{Morphisme lisse en un point}{{\hyperlink{ega4.c04.def.17-3-7}{17.3.7}}}
\EGAIVterm{Morphisme non ramifié, morphisme net}{{\hyperlink{ega4.c04.def.17-3-1}{17.3.1}}}
\EGAIVterm{Morphisme non ramifié en un point}{{\hyperlink{ega4.c04.def.17-3-7}{17.3.7}}}
\EGAIVterm{Morphisme transversal à $Y$ au point $x$ relativement à $S$}{{\hyperlink{ega4.c04.def.17-13-3}{17.13.3}}}
\EGAIVterm{Morphismes équivalents}{{\hyperlink{ega4.c04.statement.20-2-1}{20.2.1}}}
\EGAIVterm{Multiplicité d'un cycle en un point}{{\hyperlink{ega4.c04.statement.21-6-1}{21.6.1}}}
\EGAIVterm{Norme d'un diviseur}{{\hyperlink{ega4.c04.def.21-5-5}{21.5.5}}}
\EGAIVterm{Norme d'une section méromorphe}{{\hyperlink{ega4.c04.statement.21-5-3}{21.5.3}}}
\EGAIVterm{Opérateur différentiel relativement à $S$}{{\hyperlink{ega4.c04.def.16-8-1}{16.8.1}}}
\EGAIVterm{Ordre d'un opérateur différentiel}{{\hyperlink{ega4.c04.def.16-8-1}{16.8.1}}}
\EGAIVterm{Ordre d'une fonction méromorphe en un point $x\in X^{(1)}$}{{\hyperlink{ega4.c04.statement.21-6-7}{21.6.7}}}
\EGAIVterm{Partie principale d'ordre $n$, d'ordre infini, d'une section de $\sh{O}_X$}{{\hyperlink{ega4.c04.def.16-3-6}{16.3.6}}}
\EGAIVterm{Partie fermée purement de codimension $p$}{{\hyperlink{ega4.c04.statement.21-6-2}{21.6.2}}}
\EGAIVterm{Polynôme séparable sur un anneau}{{\hyperlink{ega4.c04.cor.18-4-3}{18.4.3}}}
\EGAIVterm{Préschéma de caractéristique $p$}{{\hyperlink{ega4.c04.statement.16-12-1}{16.12.1}}}
\EGAIVterm{Préschéma des coïncidences de deux morphismes}{{\hyperlink{ega4.c04.statement.17-4-5}{17.4.5}}}
\EGAIVterm{Préschéma différentiellement lisse sur $S$}{{\hyperlink{ega4.c04.def.16-10-1}{16.10.1}}}
\EGAIVterm{Préschéma étale sur $S$}{{\hyperlink{ega4.c04.def.17-3-1}{17.3.1}}}
\EGAIVterm{Préschéma étale sur $S$ en un point}{{\hyperlink{ega4.c04.def.17-3-7}{17.3.7}}}
\EGAIVterm{Préschéma formellement étale, formellement lisse, formellement non ramifié sur $S$}{{\hyperlink{ega4.c04.def.17-1-1}{17.1.1}}}
\EGAIVterm{Préschéma intersection complète absolue en un point}{{\hyperlink{ega4.c04.def.19-3-1}{19.3.1}}}
\EGAIVterm{Préschéma lisse sur $S$}{{\hyperlink{ega4.c04.def.17-3-1}{17.3.1}}}
\EGAIVterm{Préschéma lisse sur $S$ en un point}{{\hyperlink{ega4.c04.def.17-3-7}{17.3.7}}}
\EGAIVterm{Préschéma localement factoriel}{{\hyperlink{ega4.c04.thm.21-6-9}{21.6.9}}}
\EGAIVterm{Préschéma net sur $S$, non ramifié sur $S$}{{\hyperlink{ega4.c04.def.17-3-1}{17.3.1}}}

\oldpage[IV]{339}
\EGAIVterm{Préschéma net sur $S$ en un point, non ramifié sur $S$ en un point}{{\hyperlink{ega4.c04.def.17-3-7}{17.3.7}}}
\EGAIVterm{$T$-profondeur d'un $\sh{O}_X$-Module cohérent en un point}{{\hyperlink{ega4.c04.statement.19-9-1}{19.9.1}}}
\EGAIVterm{Pseudo-fonction}{{\hyperlink{ega4.c04.statement.20-2-8}{20.2.8}}}
\EGAIVterm{Pseudo-fonction relativement à $S$}{{\hyperlink{ega4.c04.statement.20-5-4}{20.5.4}}}
\EGAIVterm{Pseudo-morphisme, pseudo-$S$-morphisme}{{\hyperlink{ega4.c04.statement.20-2-1}{20.2.1}}}
\EGAIVterm{Pseudo-morphisme composé d'un morphisme et d'un pseudo-morphisme}{{\hyperlink{ega4.c04.statement.20-3-1}{20.3.1}}}
\EGAIVterm{Pseudo-morphisme composé d'un pseudo-morphisme et d'un morphisme}{{\hyperlink{ega4.c04.statement.20-3-2}{20.3.2}}}
\EGAIVterm{Pseudo-morphisme défini en un point}{{\hyperlink{ega4.c04.statement.20-2-3}{20.2.3}}}
\EGAIVterm{Pseudo-morphisme relativement à $S$}{{\hyperlink{ega4.c04.statement.20-5-1}{20.5.1}}}
\EGAIVterm{Pseudo-torseur}{{\hyperlink{ega4.c04.statement.16-5-15}{16.5.15}}}
\EGAIVterm{Restriction d'un cycle à un ouvert}{{\hyperlink{ega4.c04.statement.21-6-3}{21.6.3}}}
\EGAIVterm{Restriction d'un pseudo-morphisme à un ouvert}{{\hyperlink{ega4.c04.statement.20-2-2}{20.2.2}}}
\EGAIVterm{Restriction d'un pseudo-morphisme à un schéma local}{{\hyperlink{ega4.c04.statement.20-3-6}{20.3.6}}}
\EGAIVterm{Revêtement étale, revêtement étale localement trivial, revêtement étale trivial}{{\hyperlink{ega4.c04.item.18-2-7}{18.2.7}}}
\EGAIVterm{Schéma local hensélien}{{\hyperlink{ega4.c04.def.18-5-8}{18.5.8}}}
\EGAIVterm{Schéma strictement local}{{\hyperlink{ega4.c04.def.18-8-2}{18.8.2}}}
\EGAIVterm{Section méromorphe d'un $\sh{O}_X$-Module}{{\hyperlink{ega4.c04.statement.20-1-3}{20.1.3}}}
\EGAIVterm{Section méromorphe d'un $\sh{O}_X$-Module définie en un point}{{\hyperlink{ega4.c04.statement.20-1-7}{20.1.7}}}
\EGAIVterm{Section méromorphe régulière d'un $\sh{O}_X$-Module inversible}{{\hyperlink{ega4.c04.statement.20-1-8}{20.1.8}}}
\EGAIVterm{Section méromorphe d'un $\sh{O}_X$-Module relativement à $S$}{{\hyperlink{ega4.c04.statement.20-6-2}{20.6.2}}}
\EGAIVterm{Section méromorphe d'un $\sh{O}_X$-Module inversible régulière relativement à $S$}{{\hyperlink{ega4.c04.statement.20-6-5}{20.6.5}}}
\EGAIVterm{Sous-préschéma quasi-régulièrement immergé, régulièrement immergé}{{\hyperlink{ega4.c04.def.16-9-2}{16.9.2}}}
\EGAIVterm{Sous-préschémas se coupant transversalement en un point relativement à $S$}{{\hyperlink{ega4.c04.statement.17-13-7}{17.13.7}}}
\EGAIVterm{Support d'un cycle}{{\hyperlink{ega4.c04.statement.21-6-1}{21.6.1}}}
\EGAIVterm{Support d'un diviseur}{{\hyperlink{ega4.c04.statement.21-1-2}{21.1.2}}}
\EGAIVterm{Suite de sections transversalement régulière relativement à $S$}{{\hyperlink{ega4.c04.cor.19-1-2}{19.1.2}}}
\EGAIVterm{Symétrie canonique de $X\times_S X$}{{\hyperlink{ega4.c04.statement.16-3-3}{16.3.3}}}
\EGAIVterm{Symétrie canonique de $\sh{P}^{n}_{X/S}$, de $\sh{P}^{\infty}_{X/S}$}{{\hyperlink{ega4.c04.prop.16-3-4}{16.3.4}}}
\EGAIVterm{Système quasi-régulier (régulier) de générateurs d'un Idéal}{{\hyperlink{ega4.c04.def.16-9-1}{16.9.1}}}
\EGAIVterm{Torseur}{{\hyperlink{ega4.c04.statement.16-5-15}{16.5.15}}}
\EGAIVterm{Voisinage formel d'un sous-préschéma}{{\hyperlink{ega4.c04.rem.16-1-11}{16.1.11}}}
\EGAIVterm{$n$-ème voisinage infinitésimal}{{\hyperlink{ega4.c04.def.16-1-2}{16.1.2}}}

\endgroup

% FIN SOURCE ega4/ega4-backmatter-terminology-fr.tex

%
% DEBUT SOURCE ega4/ega4-backmatter-contents-fr.tex
% Transcription diplomatique française de EGA IV, quatrième partie, pp. 341--343.
\clearpage
\oldpage[IV]{341}
\phantomsection
\addcontentsline{toc}{section}{Table des matières}
\section*{TABLE DES MATIÈRES}
\markboth{TABLE DES MATIÈRES}{TABLE DES MATIÈRES}

\begingroup
\small
\newcommand{\EGAIVcontentsline}[3]{%
  \par\begingroup
  \rightskip=2.5em\parfillskip=-\rightskip
  \noindent\textbf{#1}\quad #2\nobreak\hspace{0.5em}%
  \dotfill\nobreak\hbox to 2.5em{\hss #3}\par
  \endgroup}

\begin{center}\textbf{CHAPITRE IV. --- Étude locale des schémas et des morphismes de schémas (fin)}\end{center}

\EGAIVcontentsline{\S16}{Invariants différentiels. Morphismes différentiellement lisses}{5}
\EGAIVcontentsline{{\texorpdfstring{\protect\hyperlink{ega4.c04.subsec.16-1}{16.1}}{16.1}}}{Invariants normaux d'une immersion}{5}
\EGAIVcontentsline{{\texorpdfstring{\protect\hyperlink{ega4.c04.subsec.16-2}{16.2}}{16.2}}}{Propriétés fonctorielles des invariants normaux d'une immersion}{9}
\EGAIVcontentsline{{\texorpdfstring{\protect\hyperlink{ega4.c04.subsec.16-3}{16.3}}{16.3}}}{Invariants différentiels fondamentaux d'un morphisme de préschémas}{14}
\EGAIVcontentsline{{\texorpdfstring{\protect\hyperlink{ega4.c04.subsec.16-4}{16.4}}{16.4}}}{Propriétés fonctorielles des invariants différentiels}{16}
\EGAIVcontentsline{{\texorpdfstring{\protect\hyperlink{ega4.c04.subsec.16-5}{16.5}}{16.5}}}{Faisceaux et fibrés tangents relatifs; dérivations}{27}
\EGAIVcontentsline{{\texorpdfstring{\protect\hyperlink{ega4.c04.subsec.16-6}{16.6}}{16.6}}}{Faisceaux de $p$-différentielles et différentielle extérieure}{34}
\EGAIVcontentsline{{\texorpdfstring{\protect\hyperlink{ega4.c04.subsec.16-7}{16.7}}{16.7}}}{Les $\sh{P}^{n}_{X/S}(\sh{F})$}{36}
\EGAIVcontentsline{{\texorpdfstring{\protect\hyperlink{ega4.c04.subsec.16-8}{16.8}}{16.8}}}{Opérateurs différentiels}{39}
\EGAIVcontentsline{{\texorpdfstring{\protect\hyperlink{ega4.c04.subsec.16-9}{16.9}}{16.9}}}{Immersions régulières et quasi-régulières}{46}
\EGAIVcontentsline{{\texorpdfstring{\protect\hyperlink{ega4.c04.subsec.16-10}{16.10}}{16.10}}}{Morphismes différentiellement lisses}{51}
\EGAIVcontentsline{{\texorpdfstring{\protect\hyperlink{ega4.c04.subsec.16-11}{16.11}}{16.11}}}{Opérateurs différentiels sur un $S$-préschéma différentiellement lisse}{53}
\EGAIVcontentsline{{\texorpdfstring{\protect\hyperlink{ega4.c04.subsec.16-12}{16.12}}{16.12}}}{Cas de la caractéristique nulle : critère jacobien pour les morphismes différentiellement lisses}{55}

\EGAIVcontentsline{\S17}{Morphismes lisses, morphismes non ramifiés, morphismes étales}{56}
\EGAIVcontentsline{{\texorpdfstring{\protect\hyperlink{ega4.c04.subsec.17-1}{17.1}}{17.1}}}{Morphismes formellement lisses, morphismes formellement non ramifiés, morphismes formellement étales}{56}
\EGAIVcontentsline{{\texorpdfstring{\protect\hyperlink{ega4.c04.subsec.17-2}{17.2}}{17.2}}}{Propriétés différentielles générales}{59}
\EGAIVcontentsline{{\texorpdfstring{\protect\hyperlink{ega4.c04.subsec.17-3}{17.3}}{17.3}}}{Morphismes lisses, morphismes non ramifiés, morphismes étales}{61}
\EGAIVcontentsline{{\texorpdfstring{\protect\hyperlink{ega4.c04.subsec.17-4}{17.4}}{17.4}}}{Caractérisations des morphismes non ramifiés}{65}
\EGAIVcontentsline{{\texorpdfstring{\protect\hyperlink{ega4.c04.subsec.17-5}{17.5}}{17.5}}}{Caractérisations des morphismes lisses}{67}
\EGAIVcontentsline{{\texorpdfstring{\protect\hyperlink{ega4.c04.subsec.17-6}{17.6}}{17.6}}}{Caractérisations des morphismes étales}{70}
\EGAIVcontentsline{{\texorpdfstring{\protect\hyperlink{ega4.c04.subsec.17-7}{17.7}}{17.7}}}{Propriétés de descente et de passage à la limite}{72}
\EGAIVcontentsline{{\texorpdfstring{\protect\hyperlink{ega4.c04.subsec.17-8}{17.8}}{17.8}}}{Critères de lissité et de non ramification par fibres}{79}
\EGAIVcontentsline{{\texorpdfstring{\protect\hyperlink{ega4.c04.subsec.17-9}{17.9}}{17.9}}}{Morphismes étales et immersions ouvertes}{79}
\EGAIVcontentsline{{\texorpdfstring{\protect\hyperlink{ega4.c04.subsec.17-10}{17.10}}{17.10}}}{Dimension relative d'un préschéma lisse sur un autre}{81}
\EGAIVcontentsline{{\texorpdfstring{\protect\hyperlink{ega4.c04.subsec.17-11}{17.11}}{17.11}}}{Morphismes lisses de préschémas lisses}{82}
\EGAIVcontentsline{{\texorpdfstring{\protect\hyperlink{ega4.c04.subsec.17-12}{17.12}}{17.12}}}{Sous-préschémas lisses d'un préschéma lisse. Morphismes lisses et morphismes différentiellement lisses}{85}

\oldpage[IV]{342}
\EGAIVcontentsline{{\texorpdfstring{\protect\hyperlink{ega4.c04.subsec.17-13}{17.13}}{17.13}}}{Morphismes transversaux}{89}
\EGAIVcontentsline{{\texorpdfstring{\protect\hyperlink{ega4.c04.subsec.17-14}{17.14}}{17.14}}}{Caractérisations locales et infinitésimales des morphismes lisses, des morphismes non ramifiés et des morphismes étales}{98}
\EGAIVcontentsline{{\texorpdfstring{\protect\hyperlink{ega4.c04.subsec.17-15}{17.15}}{17.15}}}{Cas des préschémas sur un corps de base}{99}
\EGAIVcontentsline{{\texorpdfstring{\protect\hyperlink{ega4.c04.subsec.17-16}{17.16}}{17.16}}}{Quasi-sections de morphismes plats ou lisses}{105}

\EGAIVcontentsline{\S18}{Compléments sur les morphismes étales. Anneaux locaux henséliens et anneaux strictement locaux}{109}
\EGAIVcontentsline{{\texorpdfstring{\protect\hyperlink{ega4.c04.subsec.18-1}{18.1}}{18.1}}}{Une équivalence remarquable de catégories}{109}
\EGAIVcontentsline{{\texorpdfstring{\protect\hyperlink{ega4.c04.subsec.18-2}{18.2}}{18.2}}}{Revêtements étales}{111}
\EGAIVcontentsline{{\texorpdfstring{\protect\hyperlink{ega4.c04.subsec.18-3}{18.3}}{18.3}}}{Algèbres finies et étales}{114}
\EGAIVcontentsline{{\texorpdfstring{\protect\hyperlink{ega4.c04.subsec.18-4}{18.4}}{18.4}}}{Structure locale des morphismes non ramifiés et des morphismes étales}{118}
\EGAIVcontentsline{{\texorpdfstring{\protect\hyperlink{ega4.c04.subsec.18-5}{18.5}}{18.5}}}{Anneaux locaux henséliens}{125}
\EGAIVcontentsline{{\texorpdfstring{\protect\hyperlink{ega4.c04.subsec.18-6}{18.6}}{18.6}}}{Hensélisation}{135}
\EGAIVcontentsline{{\texorpdfstring{\protect\hyperlink{ega4.c04.subsec.18-7}{18.7}}{18.7}}}{Hensélisation et anneaux excellents}{142}
\EGAIVcontentsline{{\texorpdfstring{\protect\hyperlink{ega4.c04.subsec.18-8}{18.8}}{18.8}}}{Anneaux strictement locaux et hensélisation stricte}{144}
\EGAIVcontentsline{{\texorpdfstring{\protect\hyperlink{ega4.c04.subsec.18-9}{18.9}}{18.9}}}{Fibres formelles des anneaux noethériens henséliens}{150}
\EGAIVcontentsline{{\texorpdfstring{\protect\hyperlink{ega4.c04.subsec.18-10}{18.10}}{18.10}}}{Préschémas étales sur un préschéma géométriquement unibranche ou normal}{157}
\EGAIVcontentsline{{\texorpdfstring{\protect\hyperlink{ega4.c04.subsec.18-11}{18.11}}{18.11}}}{Application aux algèbres locales noethériennes complètes sur un corps}{169}
\EGAIVcontentsline{{\texorpdfstring{\protect\hyperlink{ega4.c04.subsec.18-12}{18.12}}{18.12}}}{Applications de la localisation étale aux morphismes quasi-finis (généralisations de résultats antérieurs)}{181}

\EGAIVcontentsline{\S19}{Immersions régulières et platitude normale}{185}
\EGAIVcontentsline{{\texorpdfstring{\protect\hyperlink{ega4.c04.subsec.19-1}{19.1}}{19.1}}}{Propriétés des immersions régulières}{185}
\EGAIVcontentsline{{\texorpdfstring{\protect\hyperlink{ega4.c04.subsec.19-2}{19.2}}{19.2}}}{Immersions transversalement régulières}{190}
\EGAIVcontentsline{{\texorpdfstring{\protect\hyperlink{ega4.c04.subsec.19-3}{19.3}}{19.3}}}{Intersections complètes relatives (cas plat)}{194}
\EGAIVcontentsline{{\texorpdfstring{\protect\hyperlink{ega4.c04.subsec.19-4}{19.4}}{19.4}}}{Application : critères de régularité et de lissité pour les préschémas éclatés}{198}
\EGAIVcontentsline{{\texorpdfstring{\protect\hyperlink{ega4.c04.subsec.19-5}{19.5}}{19.5}}}{Critères de $M$-régularité}{204}
\EGAIVcontentsline{{\texorpdfstring{\protect\hyperlink{ega4.c04.subsec.19-6}{19.6}}{19.6}}}{Suites régulières relativement à un module filtré quotient}{209}
\EGAIVcontentsline{{\texorpdfstring{\protect\hyperlink{ega4.c04.subsec.19-7}{19.7}}{19.7}}}{Critère de platitude normale de Hironaka}{212}
\EGAIVcontentsline{{\texorpdfstring{\protect\hyperlink{ega4.c04.subsec.19-8}{19.8}}{19.8}}}{Propriétés de passage à la limite projective}{219}
\EGAIVcontentsline{{\texorpdfstring{\protect\hyperlink{ega4.c04.subsec.19-9}{19.9}}{19.9}}}{Suites $\sh{F}$-régulières et profondeur}{222}

\EGAIVcontentsline{\S20}{Fonctions méromorphes et pseudo-morphismes}{223}
\EGAIVcontentsline{{\texorpdfstring{\protect\hyperlink{ega4.c04.subsec.20-0}{20.0}}{20.0}}}{Introduction}{223}
\EGAIVcontentsline{{\texorpdfstring{\protect\hyperlink{ega4.c04.subsec.20-1}{20.1}}{20.1}}}{Fonctions méromorphes}{224}
\EGAIVcontentsline{{\texorpdfstring{\protect\hyperlink{ega4.c04.subsec.20-2}{20.2}}{20.2}}}{Pseudo-morphismes et pseudo-fonctions}{231}

\oldpage[IV]{343}
\EGAIVcontentsline{{\texorpdfstring{\protect\hyperlink{ega4.c04.subsec.20-3}{20.3}}{20.3}}}{Composition des pseudo-morphismes}{237}
\EGAIVcontentsline{{\texorpdfstring{\protect\hyperlink{ega4.c04.subsec.20-4}{20.4}}{20.4}}}{Propriétés des domaines de définition des fonctions rationnelles}{244}
\EGAIVcontentsline{{\texorpdfstring{\protect\hyperlink{ega4.c04.subsec.20-5}{20.5}}{20.5}}}{Pseudo-morphismes relatifs}{249}
\EGAIVcontentsline{{\texorpdfstring{\protect\hyperlink{ega4.c04.subsec.20-6}{20.6}}{20.6}}}{Fonctions méromorphes relatives}{252}

\EGAIVcontentsline{\S21}{Diviseurs}{255}
\EGAIVcontentsline{{\texorpdfstring{\protect\hyperlink{ega4.c04.subsec.21-1}{21.1}}{21.1}}}{Diviseurs sur un espace annelé}{255}
\EGAIVcontentsline{{\texorpdfstring{\protect\hyperlink{ega4.c04.subsec.21-2}{21.2}}{21.2}}}{Diviseurs et Idéaux fractionnaires inversibles}{258}
\EGAIVcontentsline{{\texorpdfstring{\protect\hyperlink{ega4.c04.subsec.21-3}{21.3}}{21.3}}}{Équivalence linéaire des diviseurs}{263}
\EGAIVcontentsline{{\texorpdfstring{\protect\hyperlink{ega4.c04.subsec.21-4}{21.4}}{21.4}}}{Images réciproques de diviseurs}{265}
\EGAIVcontentsline{{\texorpdfstring{\protect\hyperlink{ega4.c04.subsec.21-5}{21.5}}{21.5}}}{Images directes de diviseurs}{267}
\EGAIVcontentsline{{\texorpdfstring{\protect\hyperlink{ega4.c04.subsec.21-6}{21.6}}{21.6}}}{Cycle 1-codimensionnel associé à un diviseur}{270}
\EGAIVcontentsline{{\texorpdfstring{\protect\hyperlink{ega4.c04.subsec.21-7}{21.7}}{21.7}}}{Interprétation des cycles positifs 1-codimensionnels en termes de sous-préschémas}{277}
\EGAIVcontentsline{{\texorpdfstring{\protect\hyperlink{ega4.c04.subsec.21-8}{21.8}}{21.8}}}{Diviseurs et normalisation}{280}
\EGAIVcontentsline{{\texorpdfstring{\protect\hyperlink{ega4.c04.subsec.21-9}{21.9}}{21.9}}}{Diviseurs sur les préschémas de dimension 1}{284}
\EGAIVcontentsline{{\texorpdfstring{\protect\hyperlink{ega4.c04.subsec.21-10}{21.10}}{21.10}}}{Images réciproques et images directes de cycles 1-codimensionnels}{289}
\EGAIVcontentsline{{\texorpdfstring{\protect\hyperlink{ega4.c04.subsec.21-11}{21.11}}{21.11}}}{Factorialité des anneaux réguliers}{302}
\EGAIVcontentsline{{\texorpdfstring{\protect\hyperlink{ega4.c04.subsec.21-12}{21.12}}{21.12}}}{Le théorème de pureté de Van der Waerden pour l'ensemble de ramification d'un morphisme birationnel}{304}
\EGAIVcontentsline{{\texorpdfstring{\protect\hyperlink{ega4.c04.subsec.21-13}{21.13}}{21.13}}}{Couples parafactoriels. Anneaux locaux parafactoriels}{313}
\EGAIVcontentsline{{\texorpdfstring{\protect\hyperlink{ega4.c04.subsec.21-14}{21.14}}{21.14}}}{Le théorème de Ramanujam--Samuel}{323}
\EGAIVcontentsline{{\texorpdfstring{\protect\hyperlink{ega4.c04.subsec.21-15}{21.15}}{21.15}}}{Diviseurs relatifs}{329}
\EGAIVcontentsline{}{BIBLIOGRAPHIE}{333}
\EGAIVcontentsline{}{INDEX DES NOTATIONS}{334}
\EGAIVcontentsline{}{INDEX TERMINOLOGIQUE}{336}

\endgroup

% FIN SOURCE ega4/ega4-backmatter-contents-fr.tex

%
% DEBUT SOURCE ega4/ega4-errata-addenda-list3-fr.tex
% Transcription diplomatique française de EGA IV, quatrième partie, pp. 345--361.
\clearpage
\oldpage[IV]{345}
\phantomsection
\addcontentsline{toc}{section}{Errata et Addenda (Liste 3)}
\section*{ERRATA ET ADDENDA}
\markboth{ERRATA ET ADDENDA, LISTE 3}{ERRATA ET ADDENDA, LISTE 3}
\begin{center}\large (Liste 3)\end{center}

\subsection*{A) Erreurs typographiques}

\paragraph{(II, 3.3.2).}
Ligne 1 de la p.~58, remplacer $\sh{S}$ par $\sh{M}$. Ligne 6 de la
p.~58, remplacer $S(n)$ par $M(n)$.

\paragraph{(II, 4.1.1).}
Ligne 8 de la p.~71, remplacer $\sh{O}$-$P$-Module par
$\sh{O}_P$-Module.

\paragraph{(II, 5.4.5).}
Ligne 6 de la p.~102, remplacer $u$ par $f\circ u$.

\paragraph{(II, 5.4.8).}
Avant la ligne 4 du bas de la p.~102, remplacer $f\times1_{S'}$ par
$f_{(Y')}$ dans le diagramme. Ligne 4 du bas de la p.~102, remplacer
$f\times1_{S'}$ par $f_{(Y')}$.

\paragraph{(II, 6.4.4).}
Ligne 1 du haut de la p.~122, remplacer normal par intégralement clos.

\paragraph{(II, 6.4.7).}
Ligne 14 du bas de la p.~122, remplacer dimension par rang.

\paragraph{(II, 6.5.2).}
Ligne 5 de la p.~127, remplacer $(\sh{B}|U_\lambda)$ par
$(B|U_\lambda)$. Ligne 15 de la p.~127, supprimer le premier
$(N_{\sh{B}/\sh{A}}\omega_{\lambda\mu})$.

\paragraph{($0_{\mathrm{III}}$, 8.1.11).}
Ligne 8 du bas de la p.~8, remplacer $h_X(u)$ par $h'_X(u)$.

\paragraph{($0_{\mathrm{III}}$, 9.1.2).}
Ligne 10 du bas de la p.~12, remplacer $\sh{F}$ par $X$.

\paragraph{($0_{\mathrm{III}}$, 10.2.1).}
Ligne 6 du bas de la p.~18, remplacer : « en outre si » par : « si en
outre ».

\paragraph{(III, 4.4.12).}
Ligne 10 de la p.~138, remplacer : chap.~V par : chap.~IV (8.11.2).

\paragraph{(III, 5.2.5).}
Ligne 12 de la p.~153, remplacer (I, 5.5.4) par (II, 5.5.4).

\paragraph{(III, 7.4.7).}
Ligne 15 de la p.~57, remplacer : tout $A$-module par : tout
$A$-module $M$.

\paragraph{(III, 7.8.5).}
Ligne 18 du bas de la p.~74, remplacer $F$ par $\sh{F}$.

\paragraph{(III, 7.8.9).}
Ligne 9 du bas de la p.~75, remplacer $M$ par $\sh{M}$.

\paragraph{(III, 7.9.11).}
Ligne 17 du bas de la p.~79, remplacer chap.~V par chap.~VI.

\paragraph{($0_{\mathrm{IV}}$, 16.2.2.2).}
Ligne 13 de la p.~26, remplacer $M'\cap\mathfrak q^nM'$ par
$M'\cap\mathfrak q^nM$.

\paragraph{($0_{\mathrm{IV}}$, 17.3.7).}
Ligne 16 du bas de la p.~50, remplacer (17.1.10) par (17.3.3).

\paragraph{($0_{\mathrm{IV}}$, 18.1.3).}
Dans le diagramme (18.1.3.1), $p_2$ doit être à côté de la flèche
$G\to F$.

\paragraph{($0_{\mathrm{IV}}$, 19.8.7).}
Lignes 7, 8 et 9 du bas de la p.~111, remplacer (21.5.5) par (21.5.3).

\paragraph{($0_{\mathrm{IV}}$, 19.8.8).}
Ligne 2 du bas de la p.~112, remplacer $A$-module par $B$-module.

\paragraph{($0_{\mathrm{IV}}$, 20.5.4).}
Ligne 5 du bas de la p.~130, remplacer $\Omega_{B'/A'}$ par
$\Omega^1_{B'/A'}$.

\paragraph{($0_{\mathrm{IV}}$, 20.7.8).}
Ligne 9 du bas de la p.~149, remplacer (18.4.3) par (18.5.4). Ligne 4
du bas, remplacer (18.4.1) par (18.5.1).

\paragraph{($0_{\mathrm{IV}}$, 21.5.3).}
Ligne 1 du bas de la p.~166, remplacer ramifiée par ramifié.

\paragraph{($0_{\mathrm{IV}}$, 21.7.1).}
Ligne 4 de la p.~170, remplacer (20.6.15.1) par (20.6.17.1).

\paragraph{($0_{\mathrm{IV}}$, 22.1.1).}
Ligne 12 de la p.~183, remplacer (20.6.20) par (20.6.10).

\oldpage[IV]{346}
\paragraph{($0_{\mathrm{IV}}$, 22.3.3).}
Ligne 13 du bas de la p.~192, remplacer $\widehat B_{\mathfrak p}$ par
$\widehat B_{\mathfrak q}$.

\paragraph{($0_{\mathrm{IV}}$, 22.7.3).}
Ligne 21 du bas de la p.~212, remplacer $\operatorname{Ker}(i)$ par
$\operatorname{Ker}(i_0)$. Ligne 12 du bas de la p.~212, remplacer
$(\mathfrak p/\mathfrak p^2)\otimes_L L$ par
$(\mathfrak p/\mathfrak p^2)\otimes_A L$.

\paragraph{($0_{\mathrm{IV}}$, 23.1.3).}
Ligne 4 de la p.~215, remplacer $x'^q\in K'$ par $x'\in K'$. Ligne 9
de la p.~215, remplacer $x'^q\in A'$ par $x'^q\in A$.

\paragraph{({\texorpdfstring{\protect\hyperref[canon-index.ega4.c04.statement.2-1-1]{IV, 2.1.1}}{IV, 2.1.1}}).}
Ligne 4 du bas de la p.~5, remplacer (III, 1.4.15.1) par
(III, 1.4.15.5).

\paragraph{({\texorpdfstring{\protect\hyperref[canon-index.ega4.c04.prop.2-4-3]{IV, 2.4.3}}{IV, 2.4.3}}).}
Ligne 12 de la p.~20, remplacer
$(X_{\mathrm{red}}\times_{Y_{\mathrm{red}}}Y_{\mathrm{red}})_{\mathrm{red}}$
par
$(X_{\mathrm{red}}\times_{Y_{\mathrm{red}}}Y'_{\mathrm{red}})_{\mathrm{red}}$.

\paragraph{({\texorpdfstring{\protect\hyperref[canon-index.ega4.c04.prop.2-6-1]{IV, 2.6.1}}{IV, 2.6.1}}).}
Ligne 10 de la p.~27, remplacer : $X'\to X_{(S')}$ par
$X'=X_{(S')}$. 

\paragraph{({\texorpdfstring{\protect\hyperref[canon-index.ega4.c04.prop.4-2-1]{IV, 4.2.1}}{IV, 4.2.1}}).}
Ligne 20 du bas de la p.~55, remplacer
$\operatorname{deg.tr.}_kE$ par $\operatorname{deg.tr.}_LE$ et
$\operatorname{deg.tr.}_kL(t)$ par $\operatorname{deg.tr.}_LL(t)$;
ligne 19 du bas, échanger « première » et « seconde ».

\paragraph{({\texorpdfstring{\protect\hyperref[canon-index.ega4.c04.lem.4-4-2]{IV, 4.4.2}}{IV, 4.4.2}}).}
Ligne 1 du bas de la p.~59, remplacer $X$ par $X'$.

\paragraph{({\texorpdfstring{\protect\hyperref[canon-index.ega4.c04.prop.4-7-3]{IV, 4.7.3}}{IV, 4.7.3}}).}
Ligne 10 du bas de la p.~75, remplacer préschéma par $k$-préschéma.

\paragraph{({\texorpdfstring{\protect\hyperref[canon-index.ega4.c04.prop.5-6-1]{IV, 5.6.1}}{IV, 5.6.1}}).}
Ligne 16 de la p.~98, remplacer (en indice) $k(\mathfrak p($ par
$k(\mathfrak p)$.

\paragraph{({\texorpdfstring{\protect\hyperref[canon-index.ega4.c04.prop.5-6-5]{IV, 5.6.5}}{IV, 5.6.5}}).}
Ligne 20 de la p.~99, supprimer la seconde parenthèse dans $k(\xi))$.

\paragraph{({\texorpdfstring{\protect\hyperref[canon-index.ega4.c04.lem.5-6-11-1]{IV, 5.6.11.1}}{IV, 5.6.11.1}}).}
Ligne 16 de la p.~102, remplacer $K$ par $\mathfrak K$.

\paragraph{({\texorpdfstring{\protect\hyperref[canon-index.ega4.c04.cor.5-11-2]{IV, 5.11.2}}{IV, 5.11.2}}).}
Note de bas de page de la p.~123, remplacer 5.11.2) par (5.11.2).

\paragraph{({\texorpdfstring{\protect\hyperref[canon-index.ega4.c04.prop.6-1-5]{IV, 6.1.5}}{IV, 6.1.5}}).}
Ligne 2 du bas de la p.~136, remplacer (5.5.1.2) par
($0$, 16.3.9.2).

\paragraph{({\texorpdfstring{\protect\hyperref[canon-index.ega4.c04.prop.6-3-1]{IV, 6.3.1}}{IV, 6.3.1}}).}
Ligne 9 de la p.~139, remplacer $y\in\mathfrak m$ par
$y\in\mathfrak n$.

\paragraph{({\texorpdfstring{\protect\hyperref[canon-index.ega4.c04.cor.6-7-2]{IV, 6.7.2}}{IV, 6.7.2}}).}
Ligne 18 de la p.~146, remplacer $\sh{O}'_{x'}$ par $\sh{O}_{x'}$.

\paragraph{({\texorpdfstring{\protect\hyperref[canon-index.ega4.c04.statement.6-10-1]{IV, 6.10.1}}{IV, 6.10.1}}).}
Ligne 13 de la p.~155, remplacer
$\sh{F}\otimes_{\sh{O}_Y}\sh{O}_Y$ par
$\sh{F}\otimes_{\sh{O}_X}\sh{O}_Y$.

\paragraph{({\texorpdfstring{\protect\hyperref[canon-index.ega4.c04.prop.6-14-1]{IV, 6.14.1}}{IV, 6.14.1}}).}
Ligne 21 de la p.~170, remplacer : n\textsuperscript{o}~3, prop.~12
par : n\textsuperscript{o}~2, prop.~8.

\paragraph{({\texorpdfstring{\protect\hyperref[canon-index.ega4.c04.statement.6-15-2]{IV, 6.15.2}}{IV, 6.15.2}}).}
Ligne 13 de la p.~177, remplacer $t(t-v)$ par $t^2(t-v)^2$.

\paragraph{({\texorpdfstring{\protect\hyperref[canon-index.ega4.c04.prop.6-15-7]{IV, 6.15.7}}{IV, 6.15.7}}).}
Ligne 6 du bas de la p.~179, remplacer $X\otimes_Ak'$ par
$X\otimes_kk'$.

\paragraph{({\texorpdfstring{\protect\hyperref[canon-index.ega4.c04.item.7-5-4]{IV, 7.5.4}}{IV, 7.5.4}}).}
Ligne 14 du bas de la p.~206, remplacer
$\operatorname{Spec}(\widehat B\otimes k)$ par
$\operatorname{Spec}(\widehat B\otimes_{\widehat A}k)$.

\paragraph{({\texorpdfstring{\protect\hyperref[canon-index.ega4.c04.item.7-8-4]{IV, 7.8.4}}{IV, 7.8.4}}).}
Ligne 10 du bas de la p.~216, remplacer « vérifiant »
({\hyperref[canon-index.ega4.c04.def.7-8-2]{7.8.2}}, (ii) et (iii)) par
« vérifiant ({\hyperref[canon-index.ega4.c04.def.7-8-2]{7.8.2}}, (ii) et (iii)) ».

\paragraph{(IV, 2\textsuperscript{e} Partie, Bibliographie).}
Ligne 6 de la p.~224, remplacer : \emph{Amer. J. of Math.}, par
\emph{Annals of Math.}

\paragraph{({\texorpdfstring{\protect\hyperref[canon-index.ega4.c04.cor.8-5-8]{IV, 8.5.8}}{IV, 8.5.8}}).}
Ligne 6 de la p.~24, remplacer $\operatorname{Ker}g_\mu$ à la fin de
la ligne par $\operatorname{Ker}g_\lambda$.

\paragraph{({\texorpdfstring{\protect\hyperref[canon-index.ega4.c04.prop.8-10-1]{IV, 8.10.1}}{IV, 8.10.1}}).}
Ligne 13 du bas de la p.~36, remplacer (2.3.11) par
{\hyperref[canon-index.ega4.c04.cor.8-2-11]{(8.2.11)}}.

\paragraph{({\texorpdfstring{\protect\hyperref[canon-index.ega4.c04.prop.10-2-1]{IV, 10.2.1}}{IV, 10.2.1}}).}
Ligne 14 du bas de la p.~98, remplacer $y$ par $Y$.

\paragraph{({\texorpdfstring{\protect\hyperref[canon-index.ega4.c04.statement.10-4-11]{IV, 10.4.11}}{IV, 10.4.11}}).}
Ligne 10 du bas de la p.~104, remplacer (17.9.7) par
{\hyperlink{ega4.c04.prop.17-9-6}{(17.9.6)}}.

\paragraph{({\texorpdfstring{\protect\hyperref[canon-index.ega4.c04.thm.11-2-9]{IV, 11.2.9}}{IV, 11.2.9}}).}
Ligne 3 du bas de la p.~126, remplacer $A_\lambda$ par $A_\mu$.

\paragraph{({\texorpdfstring{\protect\hyperref[canon-index.ega4.c04.prop.11-3-4]{IV, 11.3.4}}{IV, 11.3.4}}).}
Ligne 16 de la p.~134, remplacer $F$ par $\sh{F}$.

\paragraph{({\texorpdfstring{\protect\hyperref[canon-index.ega4.c04.thm.11-3-10]{IV, 11.3.10}}{IV, 11.3.10}}).}
Ligne 9 du bas de la p.~139, remplacer (11.3.7) par
{\hyperref[canon-index.ega4.c04.thm.11-3-10]{(11.3.10)}}.

\paragraph{({\texorpdfstring{\protect\hyperref[canon-index.ega4.c04.statement.11-7-5]{IV, 11.7.5}}{IV, 11.7.5}}).}
Ligne 11 de la p.~159, remplacer (6.5.11) par (6.15.11).

\paragraph{({\texorpdfstring{\protect\hyperref[canon-index.ega4.c04.thm.12-1-1]{IV, 12.1.1}}{IV, 12.1.1}}).}
Ligne 13 de la p.~174, ajouter à la fin : pour $k\geq1$.

\paragraph{({\texorpdfstring{\protect\hyperref[canon-index.ega4.c04.thm.14-2-1]{IV, 14.2.1}}{IV, 14.2.1}}).}
Ligne 11 de la p.~202, remplacer (6.1.2) par (6.1.1).

\oldpage[IV]{347}
\paragraph{({\texorpdfstring{\protect\hyperref[canon-index.ega4.c04.prop.14-3-8]{IV, 14.3.8}}{IV, 14.3.8}}).}
Ligne 10 du bas de la p.~206, remplacer $(f^{-1}y)$ par $f^{-1}(y)$.

\paragraph{({\texorpdfstring{\protect\hyperref[canon-index.ega4.c04.rem.14-5-7]{IV, 14.5.7}}{IV, 14.5.7}}).}
Ligne 16 de la p.~218, remplacer (14.5.5) par (14.5.6).

\paragraph{({\texorpdfstring{\protect\hyperref[canon-index.ega4.c04.cor.14-5-11]{IV, 14.5.11}}{IV, 14.5.11}}).}
Dans le diagramme du haut de la page 223, les flèches horizontales
doivent toutes aller de droite à gauche.

\paragraph{({\texorpdfstring{\protect\hyperref[canon-index.ega4.c04.prop.15-4-2]{IV, 15.4.2}}{IV, 15.4.2}}).}
Ligne 23 de la p.~230, remplacer (6.1.4) par (6.1.5).

\paragraph{({\texorpdfstring{\protect\hyperref[canon-index.ega4.c04.cor.15-5-2]{IV, 15.5.2}}{IV, 15.5.2}}).}
Ligne 16 de la p.~233, remplacer $f^{-1}(y)$ par $f^{-1}(y')$.

\paragraph{({\texorpdfstring{\protect\hyperref[canon-index.ega4.c04.item.15-6-10]{IV, 15.6.10}}{IV, 15.6.10}}).}
Ligne 6 de la p.~242, remplacer ${}^{-1}(y')$ par $f^{-1}(y')$.
Ligne 11 de la p.~242, remplacer : $\alpha$) et $\beta$), par :
$\alpha$) ou $\beta$).

\subsection*{B) Modifications de texte}

\paragraph{(Err\textsubscript{IV}, 1).}
Dans ($0_{\mathrm{III}}$, 11.1.6), l'énoncé (reproduisant un énoncé de
(G, I, 4.4.1)) est incorrect lorsque $q(n)$ est un entier quelconque;
il n'est correct que si l'on a $q(n)=0$ ou $q(n)=n$. En effet, dans le
premier cas, le noyau de $d_r^{p,0}$ est alors $E_r^{p,0}$, puisque
l'image de $d_r^{p,0}$ est $E_r^{p+r,-r+1}=0$ par hypothèse; de même
l'image de $d_r^{p-r,r-1}$ est nulle puisque
$E_r^{p-r,r-1}=0$; on a donc bien $E_{r+1}^{p,0}=E_r^{p,0}$. On
raisonne de même lorsque $q(n)=n$ pour tout $n$.

\paragraph{(Err\textsubscript{IV}, 2).}
Dans (III, 3.4), supprimer les lignes 2 et 3 de la p.~119.

\oldpage[IV]{348}
\paragraph{(Err\textsubscript{IV}, 3).}
Dans (III, 4.4.9), ligne 15 du bas, p.~137, remplacer : finies par :
finies en tant qu'ensembles.

\paragraph{(Err\textsubscript{IV}, 4).}
Dans (III, 7.8.10), ligne 9 de la p.~76, remplacer simple par lisse.

\paragraph{(Err\textsubscript{IV}, 5).}
Dans ($0_{\mathrm{IV}}$, 16.1.4), ligne 8 du bas de la p.~23, remplacer :
Dire que $A$ est caténaire, par : Dire qu'un anneau intègre $A$ est
caténaire.

\paragraph{(Err\textsubscript{IV}, 6).}
Dans ($0_{\mathrm{IV}}$, 16.1.9), on peut supprimer l'hypothèse que
$A$ est noethérien. En effet, pour tout élément $u\in B_M$, $A[u]$
est une sous-$A$-algèbre de $C$, donc $u$ est entier sur $A$ (Bourbaki,
\emph{Alg. comm.}, chap.~V, \S1, n\textsuperscript{o}~1, cor. de la
prop.~1), et par suite aussi sur $A_M$; on conclut encore par
{\hyperlink{ega4.c04.prop.16-1-5}{(16.1.5)}}.

\paragraph{(Err\textsubscript{IV}, 7).}
Dans ($0_{\mathrm{IV}}$, 16.5.12), on notera que la relation
(16.5.11.1) équivaut à la même relation où l'on remplace $A$ par les
anneaux quotients $A/\mathfrak r$, $\mathfrak r$ parcourant l'ensemble
des idéaux premiers minimaux de $A$. On peut donc dans la preuve de
{\hyperlink{ega4.c04.statement.16-5-12}{(16.5.12)}} se borner au cas
où $A$ est intègre, et appliquer (16.1.4.2).

\paragraph{(Err\textsubscript{IV}, 8).}
Supprimer ($0_{\mathrm{IV}}$, 19.3.12) dont la démonstration est erronée
(les ensembles $W_\lambda$, $U_\lambda$ et $V_\lambda$ ne sont pas des
variétés linéaires).

\paragraph{(Err\textsubscript{IV}, 9).}
Dans ($0_{\mathrm{IV}}$, 19.8.3), lignes 6 et 8 du bas de la p.~110,
remplacer homomorphisme par homomorphisme local; de même dans
($0_{\mathrm{IV}}$, 19.8.4), ligne 1 de la p.~111.

\paragraph{(Err\textsubscript{IV}, 10).}
Dans ($0_{\mathrm{IV}}$, 20.1.2), ligne 2 du bas de la p.~117,
remplacer $a\in A$ par $b\in B$ et $aD$ par $bD$; ligne 1 du bas de la
p.~117, remplacer $A$-module par $B$-module.

\paragraph{(Err\textsubscript{IV}, 11).}
Dans ($0_{\mathrm{IV}}$, 20.2.1), ligne 8 de la p.~119, remplacer
$A$-modules par $B$-modules. Lignes 21 et 22 de la p.~119, remplacer
« di-homomorphismes de modules (relatif à $u$) » par : homomorphisme
de $B$-modules.

\paragraph{(Err\textsubscript{IV}, 12).}
Dans ($0_{\mathrm{IV}}$, 20.4.6.1), remplacer $p_2-p_1$ par $p_1-p_2$;
ligne 10 du bas de la p.~125, remplacer $x\otimes1-1\otimes x$ par
$1\otimes x-x\otimes1$.

\paragraph{(Err\textsubscript{IV}, 13).}
Dans ($0_{\mathrm{IV}}$, 20.4.8), ligne 14 de la p.~126, remplacer
$p_2=p_1-d_{B/A}$ par $p_1=p_2-d_{B/A}$.

\paragraph{(Err\textsubscript{IV}, 14).}
Dans ($0_{\mathrm{IV}}$, 20.4.13), remplacer les lignes 6 à 11 du bas
de la p.~127 par le texte suivant :

\medskip
\noindent
(iv) Soient $A$, $B$ deux anneaux intègres tels que $A\subset B$, $B$
étant entier sur $A$, et que le corps des fractions de $B$ soit une
extension séparable de celui de $A$. Alors le $B$-module
$\Omega^1_{B/A}$ est un module de torsion. En effet, pour tout $x\in B$,
soit $f$ le polynôme minimal de $x$ par rapport au corps des fractions
$K$ de $A$; il existe $a\in A$ non nul tel que $af[T]\in A[T]$; comme
$x$ est séparable sur $K$, on a $f'(x)\neq0$, et d'autre part de la
relation $af(x)=0$ on déduit $af'(x)d_{B/A}x=0$, d'où notre assertion
en vertu de {\hyperlink{ega4.c04.cor.20-4-7}{(20.4.7)}}.

\paragraph{(Err\textsubscript{IV}, 15).}
Dans ($0_{\mathrm{IV}}$, 22.2.7), ligne 5 de la p.~189, au lieu de
« formellement lisses », il faut « formellement lisses pour les
topologies préadiques ».

\paragraph{(Err\textsubscript{IV}, 16).}
Dans le Sommaire du chap.~IV, ligne 4 du bas de la p.~222, ajouter :
et anneaux strictement locaux. Lignes 1 à 3 du bas, remplacer par :
\begin{quote}
\S19. Immersions régulières et platitude normale.\\
\S20. Fonctions méromorphes et pseudo-morphismes.\\
\S21. Diviseurs.
\end{quote}
Ligne 16 de la p.~223, remplacer « de présentation finie » par
« localement de présentation finie »; ligne 19 de la p.~223, remplacer
« comme » par « comme étant ».

\paragraph{(Err\textsubscript{IV}, 17).}
Dans (IV, 3.1.10), ajouter après la ligne 16 de la p.~39 :

\begin{remark}[3.1.10.2]
\label{IV.3.1.10.2}
Soit $f:X\to Y$ un morphisme quasi-compact et quasi-séparé, et soit
$\sh{F}$ un $\sh{O}_X$-Module quasi-cohérent, de sorte que
$f_*(\sh{F})$ est un $\sh{O}_Y$-Module quasi-cohérent
{\hyperref[canon-index.ega4.c04.statement.1-7-4]{(1.7.4)}}. Alors on a
\begin{align}
 \operatorname{Ass}(f_*(\sh{F}))&\subset f(X),
 \label{IV.3.1.10.3}\tag{3.1.10.3}\\
 \operatorname{Ass}(f_*(\sh{F}))&\subset f(\operatorname{Ass}(\sh{F}))
 \quad\text{si $X$ est localement noethérien}.
 \label{IV.3.1.10.4}\tag{3.1.10.4}
\end{align}

Soit en effet $y\in\operatorname{Ass}(f_*(\sh{F}))$; en vertu de
{\hyperref[canon-index.ega4.c04.lem.2-3-1]{(2.3.1)}} et du fait que le morphisme
$\operatorname{Spec}(\sh{O}_{Y,y})\to Y$ est plat,
$(f_*(\sh{F}))_y$ est égal à $(f'_*(\sh{F}'))_y$, où
\[
 f':X\times_Y\operatorname{Spec}(\sh{O}_{Y,y})
 \longrightarrow\operatorname{Spec}(\sh{O}_{Y,y})
\]
est déduit de $f$ par changement de base et
$\sh{F}'=\sh{F}\otimes_{\sh{O}_Y}\sh{O}_{Y,y}$. On est donc ramené à
prouver (3.1.10.3) et (3.1.10.4) en y remplaçant $Y$ par
$\operatorname{Spec}(\sh{O}_{Y,y})$, $f$ par $f'$ et $\sh{F}$ par
$\sh{F}'$, compte tenu de (I, 3.6.5) et de la définition de
$\operatorname{Ass}(\sh{F}')$.

Par définition, il y a un élément
$t_y\in(f_*(\sh{F}))_y$ non nul et dont $\mathfrak m_y$ est
l'annulateur, et comme $(f_*(\sh{F}))_y$ s'identifie à
$\Gamma(Y,f_*(\sh{F}))$ puisque $Y$ est local, $t_y$ s'identifie aussi
à une section $s\in\Gamma(X,\sh{F})$. L'Idéal de $\sh{O}_X$ annulateur
de $s$ contient donc $\mathfrak m_y\sh{O}_X$, et il résulte de
($0_{\mathrm I}$, 1.7.4) que le support de $s$ est contenu dans
$f^{-1}(y)$; si ce dernier était vide, on aurait $s=0$, ce qui est
absurde. Lorsque de plus $X$ est localement noethérien, alors le
sous-$\sh{O}_X$-Module cohérent $\sh{O}_Xs$ de $\sh{F}$ n'est pas nul,
donc $\operatorname{Ass}(\sh{O}_Xs)\neq\varnothing$
{\hyperref[canon-index.ega4.c04.cor.3-1-5]{(3.1.5)}}; mais
\[
 \operatorname{Ass}(\sh{O}_Xs)\subset
 \operatorname{Supp}(\sh{O}_Xs)\subset f^{-1}(y),\qquad
 \operatorname{Ass}(\sh{O}_Xs)\subset\operatorname{Ass}(\sh{F})
\]
{\hyperref[canon-index.ega4.c04.prop.3-1-7]{(3.1.7)}}, donc
$y\in f(\operatorname{Ass}(\sh{F}))$.
\end{remark}

\oldpage[IV]{349}
\begin{remark}[3.1.10.5]
\label{IV.3.1.10.5}
On notera que même lorsque $f$ est un morphisme propre, on n'a pas
nécessairement
\[
 f(\operatorname{Ass}(\sh{F}))\subset
 \operatorname{Ass}(f_*(\sh{F})).
\]
On le voit en prenant pour $Y$ le spectre d'un corps et pour $X$ un
espace projectif sur $k$, par exemple $\mathbf P^1_k$; alors
$\sh{F}=\sh{O}_X(-1)$ est un $\sh{O}_X$-Module cohérent $\neq0$ tel
que $\Gamma(X,\sh{F})=0$ (III, 2.1.13), donc
$\operatorname{Ass}(f_*(\sh{F}))=\varnothing$ et
$\operatorname{Ass}(\sh{F})\neq\varnothing$.
\end{remark}

\paragraph{(Err\textsubscript{IV}, 18).}
Après (IV, 3.2.8), insérer après la ligne 17 de la p.~43 :

\begin{remark}[3.2.9]
\label{IV.3.2.9}
Pour tout point $s$ d'un préschéma $X$, notons
$i_s:\operatorname{Spec}(k(s))\to X$ le morphisme canonique, qui est
un monomorphisme quasi-compact. Pour tout $k(s)$-module $G(s)$,
considéré comme Module sur $\operatorname{Spec}(k(s))$,
$(i_s)_*(G(s))$ est donc un $\sh{O}_X$-Module quasi-cohérent
{\hyperref[canon-index.ega4.c04.statement.1-7-4]{(1.7.4)}}. On vérifie aussitôt
que la fibre de $(i_s)_*(G(s))$ est nulle en tout point
$x\notin\overline{\{s\}}$ et isomorphe à $G(s)$ dans le cas contraire.
Si $X$ est localement noethérien et $G(s)$ de longueur finie, il est
clair que $\operatorname{Ass}((i_s)_*(G(s)))$ est réduit au seul point
$s$, autrement dit $(i_s)_*(G(s))$ est irredondant.

Cela étant, supposons $X$ localement noethérien, et soit $\sh{F}$ un
$\sh{O}_X$-Module cohérent. Soit d'autre part $S$ une partie discrète
de $X$, et pour tout $s\in S$, posons
$\sh{F}(s)=i_s^*(\sh{F})$, qui est un $k(s)$-module; considérons pour
chaque $s\in S$ un module quotient $G(s)$ de $\sh{F}(s)$ de longueur
finie. On a donc un homomorphisme canonique
\[
 u:\sh{F}\longrightarrow
 \bigoplus_{s\in S}(i_s)_*(\sh{F}(s))\longrightarrow
 \bigoplus_{s\in S}(i_s)_*(G(s))
\]
de $\sh{O}_X$-Modules cohérents (puisque $S$, étant discret, est
localement fini dans $X$). Si $\sh{H}=\operatorname{Ker}u$,
$\sh{G}=\sh{F}/\sh{H}$ est un $\sh{O}_X$-Module cohérent et on déduit
de $u$ un homomorphisme injectif
\[
 v:\sh{G}\longrightarrow\bigoplus_{s\in S}(i_s)_*(G(s)).
\]
Si $\sh{G}^{(s)}$ est l'image de $\sh{G}$ dans $(i_s)_*(G(s))$,
$(\sh{G}^{(s)})_{s\in S}$ est une décomposition irredondante, qui est
réduite lorsque les $G(s)$ sont tous $\neq0$, et telle que
$\operatorname{Ass}(\sh{G})=S$ ($\sh{G}$ étant donc sans cycle premier
associé immergé). On a ainsi obtenu une bijection
$(G(s))_{s\in S}\leftrightarrow\sh{G}$, où $S$ parcourt l'ensemble des
parties discrètes de $X$ pour lesquelles $\sh{F}(s)\neq0$ pour tout
$s\in S$, et, pour chaque $S$, $(G(s))_{s\in S}$ parcourt l'ensemble
des quotients $\neq0$ et de longueur finie des $\sh{F}(s)$; d'autre
part $\sh{G}$ parcourt l'ensemble des $\sh{O}_X$-Modules quotients
cohérents de $\sh{F}$, sans cycle premier immergé. La bijection
réciproque associe à chaque $\sh{G}$ sa décomposition irredondante
réduite.
\end{remark}

\paragraph{(Err\textsubscript{IV}, 19).}
Dans (IV, 4.2.1), avant la ligne 17 du bas de la p.~55, insérer :

\begin{remark}[4.2.1.4]
\label{IV.4.2.1.4}
Sous les hypothèses de {\hyperref[canon-index.ega4.c04.prop.4-2-1]{(4.2.1)}}, on a
en outre
\begin{equation}
 \dim(K\otimes_kL)=
 \inf(\operatorname{deg.tr.}_kK,\operatorname{deg.tr.}_kL).
 \label{IV.4.2.1.5}\tag{4.2.1.5}
\end{equation}
où il doit être entendu que si les deux corps $K$, $L$ sont de degré de
transcendance infini (qui peuvent être des cardinaux quelconques), le
premier membre doit être pris égal à $+\infty$ (la dimension considérée
ici est celle définie dans ($0$, 16.1.1)).

Supposons en effet par exemple que
$\operatorname{deg.tr.}_kK\leq\operatorname{deg.tr.}_kL$ et en premier
lieu que $\operatorname{deg.tr.}_kK=n<+\infty$. On a vu dans la preuve
de {\hyperref[canon-index.ega4.c04.prop.4-2-1]{(4.2.1)}} que $K\otimes_kL$ est une
extension algébrique de $K'\otimes_kL$, et a donc même dimension
($0$, 16.1.5). On peut donc se borner au cas où
$K=k[T_1,\ldots,T_n]$; alors $K\otimes_kL$ est un anneau de fractions
de
\[
 k[T_1,\ldots,T_n]\otimes_kL=L[T_1,\ldots,T_n],
\]
qui est de dimension $n$ {\hyperref[canon-index.ega4.c04.cor.5-5-4]{(5.5.4)}}, et
par suite $\dim(K\otimes_kL)\leq n$ ($0$, 16.1.3.1). Pour prouver que
$\dim(K\otimes_kL)=n$, il suffit évidemment de montrer qu'il y a dans
$\operatorname{Spec}(L[T_1,\ldots,T_n])$ un point rationnel sur $L$
qui appartient à $\operatorname{Spec}(K\otimes_kL)$
{\hyperref[canon-index.ega4.c04.cor.5-2-3]{(5.2.3)}}, ou encore (I, 1.7.4) qu'il
existe un $L$-homomorphisme $u:K\otimes_kL\to L$. Mais par définition
du produit tensoriel, cela revient à prouver l'existence d'un
$k$-homomorphisme $v:K\to L$, ce qui est évident en vertu de
l'hypothèse
$\operatorname{deg.tr.}_kK\leq\operatorname{deg.tr.}_kL$.

Si $\operatorname{deg.tr.}_kK\leq\operatorname{deg.tr.}_kL$ et si ces
deux cardinaux sont infinis, on voit de la même façon qu'il y a un point
rationnel dans $\operatorname{Spec}(L[(t_\alpha)_{\alpha\in I}])$ qui
appartient à $\operatorname{Spec}(K\otimes_kL)$; mais le noyau de tout
$L$-homomorphisme de $L[(t_\alpha)_{\alpha\in I}]$ dans $L$ est un
idéal maximal engendré par une famille de la forme
$(t_\alpha-a_\alpha)_{\alpha\in I}$ avec $a_\alpha\in L$. Il est clair
qu'un tel idéal maximal contient une suite infinie décroissante
d'idéaux premiers distincts, donc l'anneau local de
$L[(t_\alpha)_{\alpha\in I}]$ correspondant est non noethérien et de
dimension infinie; comme c'est un anneau local de $K\otimes_kL$, on
voit que ce dernier est non noethérien et de dimension infinie.
\end{remark}

\oldpage[IV]{350}
\paragraph{(Err\textsubscript{IV}, 20).}
Après (IV, 4.5.19), insérer, avant la ligne 12 du bas de la p.~67 :

\begin{corollary}[4.5.19.3]
\label{IV.4.5.19.3}
Soient $k$ un corps, $X$ un $k$-préschéma de type fini sur $k$, qui est
irréductible mais non géométriquement irréductible. Alors il existe une
partie fermée $Y$ de $X$, rare dans $X$, contenant tous les points
$x\in X$ tels que $k(x)$ soit radiciel sur $k$ (et en particulier tous
les points de $X$ rationnels sur $k$).
\end{corollary}

Soient $k'$ une clôture algébrique de $k$, $X'=X\otimes_kk'$,
$p:X'\to X$ la projection canonique. Si $x\in X$ est tel que $k(x)$
soit radiciel sur $k$, il n'existe qu'un seul point $x'\in X'$ au-dessus
de $x$ (I, 3.4.9); ce point appartient nécessairement à deux composantes
irréductibles distinctes de $X'$, sans quoi $X$ serait géométriquement
irréductible en vertu de
{\hyperref[canon-index.ega4.c04.prop.4-5-19]{(4.5.19)}}, contrairement à
l'hypothèse. Soient alors $X'_i$ ($1\leq i\leq n$) les composantes
irréductibles du préschéma noethérien $X'$, et pour
$1\leq i<j\leq n$, posons $X'_{ij}=X'_i\cap X'_j$; l'ensemble
\[
 Y=p\left(\bigcup_{i<j}X'_{ij}\right)
\]
répond à la question. En effet, comme $p$ est un morphisme entier et
que les $X'_{ij}$ sont fermés dans $X'$, $Y$ est fermé dans $X$.
D'autre part, si $\xi$ est le point générique de $X$, $p^{-1}(\xi)$
est discret et formé des points génériques des $X'_i$
{\hyperref[canon-index.ega4.c04.cor.4-5-11]{(4.5.11)}}, donc ne rencontre aucun
des $X'_{ij}$; l'ensemble fermé $Y$ est par suite rare.

\paragraph{(Err\textsubscript{IV}, 21).}
Dans (IV, 5.6.11.1), ligne 7 du bas de la p.~102, après « intègre »,
ajouter : « de dimension 2 ».

\paragraph{(Err\textsubscript{IV}, 22).}
Après (IV, 6.2.3), après la ligne 18 de la p.~138, insérer :

\begin{proposition}[6.2.4]
\label{IV.6.2.4}
Sous les hypothèses de {\hyperref[canon-index.ega4.c04.prop.6-2-3]{(6.2.3)}}, soit
de plus $M$ un $A$-module de type fini et supposons $M\neq0$ et
$N\neq0$. Alors on a
\begin{equation}
 \operatorname{dim.proj}_B(M\otimes_AN)=
 \operatorname{dim.proj}_A(M)+
 \operatorname{dim.proj}_{B\otimes_Ak}(N\otimes_Ak).
 \label{IV.6.2.4.1}\tag{6.2.4.1}
\end{equation}
\end{proposition}

Considérons la résolution gauche $L_\bullet$ de $N$ introduite dans
{\hyperref[canon-index.ega4.c04.prop.6-2-3]{(6.2.3)}}, et une résolution gauche
$K_\bullet$ de $M$ par des $A$-modules libres de type fini :
\[
 \cdots\longrightarrow K_i\longrightarrow K_{i-1}\longrightarrow
 \cdots\longrightarrow K_0\longrightarrow M\longrightarrow0.
\]
On a vu que les $L_i$, $N$, les $Z_i(L_\bullet)$ sont des $A$-modules
plats; il en est de même des $B_i(L_\bullet)$ qui sont égaux aux
$Z_i(L_\bullet)$ sauf pour $i=0$, et dans ce dernier cas
$B_0(L_\bullet)$ est

\oldpage[IV]{351}
plat par ($0_{\mathrm I}$, 6.1.2); enfin, les $H_i(L_\bullet)$ sont nuls
par définition sauf $H_0(L_\bullet)=N$, qui est un $A$-module plat. On
sait que, dans ces conditions,
\[
 H_n(K_\bullet\otimes_AL_\bullet)\simeq
 \bigoplus_{p+q=n}
 H_p(K_\bullet)\otimes_AH_q(L_\bullet)
\]
(G, I, 5.5), donc $K_\bullet\otimes_AL_\bullet$ (gradué comme
d'ordinaire) est une résolution de $M\otimes_AN$ par des $B$-modules
libres. Posons
\[
 r=\operatorname{dim.proj}_A(M),\qquad
 s=\operatorname{dim.proj}_B(N)=
 \operatorname{dim.proj}_{B\otimes_Ak}(N\otimes_Ak).
\]
Il y a donc une résolution libre $K_\bullet$ de longueur $r$ et une
résolution libre $L_\bullet$ de longueur $s$, et
$K_\bullet\otimes_AL_\bullet$ est alors une résolution libre de longueur
$r+s$, ce qui prouve déjà que
$\operatorname{dim.proj}_B(M\otimes_AN)\leq r+s$. Il reste à prouver
($0$, 17.2.6) que l'on a
\[
 \operatorname{Tor}^{B}_{r+s}(M\otimes_AN,k(B))\neq0.
\]
Utilisant la résolution précédente $K_\bullet\otimes_AL_\bullet$ pour
calculer ce foncteur, on obtient
\[
 \operatorname{Tor}^{B}_{r+s}(M\otimes_AN,k(B))=
 H_{r+s}((K_\bullet\otimes_AL_\bullet)\otimes_Bk(B)).
\]
Mais on peut écrire
\[
 (K_\bullet\otimes_AL_\bullet)\otimes_Bk(B)=
 (K_\bullet\otimes_Ak)\otimes_k
 (L_\bullet\otimes_Bk(B)),
\]
qui est donc un produit tensoriel de deux complexes d'espaces
vectoriels sur $k$; tous les $k$-modules étant plats, la formule de
Künneth donne
\[
 H_{r+s}((K_\bullet\otimes_Ak)\otimes_k
 (L_\bullet\otimes_Bk(B)))\simeq
 H_r(K_\bullet\otimes_Ak)\otimes_k
 H_s(L_\bullet\otimes_Bk(B)),
\]
l'homologie du premier (resp. second) complexe étant nulle en degrés
$>r$ (resp. $>s$). On a donc
\[
 \operatorname{Tor}^{B}_{r+s}(M\otimes_AN,k(B))=
 \operatorname{Tor}^{A}_{r}(M,k)\otimes_k
 \operatorname{Tor}^{B}_{s}(N,k(B))
\]
à isomorphie près, et comme chacun des facteurs est $\neq0$ par
hypothèse, il en est de même de leur produit tensoriel sur $k$.

\paragraph{(Err\textsubscript{IV}, 23).}
Dans (IV, 6.10.1), ligne 8 de la p.~155, après : canonique, ajouter :
$\sh{F}$ un $\sh{O}_X$-Module cohérent.

\paragraph{(Err\textsubscript{IV}, 24).}
Dans ({\hyperref[canon-index.ega4.c04.prop.6-14-4]{IV, 6.14.4}}), la référence
donnée ligne 12 du bas de la p.~175 est insuffisante; il faut raisonner
de la façon suivante. Tout d'abord, si $M$ est une base de transcendance
de $K'$ sur $A$, $A(M)$ est algébriquement fermé dans $B(A(M))$
(Bourbaki, \emph{Alg.}, chap.~V, \S6, exerc.~8 c)). On peut donc se
borner au cas où $K'$ est une extension algébrique séparable de $A$.
Soit $p$ l'exposant caractéristique de $A$; il résulte de Bourbaki,
\emph{loc. cit.}, \S9, exerc.~2, que la fermeture algébrique de $K'$
dans $B(K')$ est radicielle sur $K'$; montrons que tout élément
$x\in B(K')$ tel que $x^{p^r}\in K'$ appartient déjà à $K'$. Pour cela,
considérons une base $(e_\lambda)$ de $K'$ sur $A$; puisque $K'$ est une
extension algébrique séparable de $A$, les $e_\lambda^{p^r}$ forment
aussi une base de $K'$ sur $A$ (Bourbaki, \emph{loc. cit.}, \S8,
n\textsuperscript{o}~3, cor. de la prop.~4). On a
$x=\sum_\lambda e_\lambda b_\lambda$, avec $b_\lambda\in B$ et, par
hypothèse, il existe des éléments $a_\lambda\in A$ tels que l'on ait
\[
 \sum_\lambda e_\lambda^{p^r}b_\lambda^{p^r}=
 \sum_\lambda e_\lambda^{p^r}a_\lambda;
\]
mais cela entraîne $b_\lambda^{p^r}=a_\lambda$ pour tout $\lambda$;
comme par hypothèse $A$ est algébriquement fermé dans $B$, on a donc
$b_\lambda\in A$ pour tout $\lambda$, et par suite $x\in K'$.

\paragraph{(Err\textsubscript{IV}, 25).}
Dans (IV, 7.4.8, B)), p.~203, la question a été résolue affirmativement
par R.~Kiehl, \emph{Ausgezeichnete Ringe der nichtarchimedischen
analytischen Geometrie}, à paraître dans \emph{Inventiones
Mathematicae}.

\paragraph{(Err\textsubscript{IV}, 26).}
Dans (IV, 7.8.4), ligne 9 du bas de la p.~217, remplacer :
\begin{quote}
Nous verrons plus loin (\S18) que si $k$ est un corps valué complet non
discret\ldots
\end{quote}
par :
\begin{quote}
Il a été démontré par R.~Kiehl (dans l'article cité en
(Err\textsubscript{IV}, 25)) que si $k$ est un corps valué complet non
discret\ldots
\end{quote}

\oldpage[IV]{352}
\paragraph{(Err\textsubscript{IV}, 27).}
Dans (IV, 7.8.4), ajouter, avant la ligne 6 du bas de la p.~217 :

\smallskip
\noindent
\textup{(vi)} Soit $A$ un anneau local noethérien intègre de dimension
1; il est alors universellement caténaire
{\hyperref[canon-index.ega4.c04.cor.7-2-9]{(7.2.9)}} et, pour qu'il soit excellent,
il faut et il suffit donc, en vertu de
{\hyperref[canon-index.ega4.c04.statement.7-3-19]{(7.3.19)}}, que ses fibres
formelles soient géométriquement réduites, ou encore que son complété
$\widehat A$ soit réduit et que les corps composants de son anneau total
de fractions soient des extensions séparables du corps des fractions
$K$ de $A$. En vertu de {\hyperref[canon-index.ega4.c04.thm.7-6-4]{(7.6.4)}} et
{\hyperref[canon-index.ega4.c04.thm.7-7-2]{(7.7.2)}}, il revient aussi au même de
dire que $A$ est universellement japonais. Lorsque $A$ est un anneau de
valuation discrète, il en est de même de $\widehat A$; dire que $A$ est
excellent signifie alors que le corps des fractions de $\widehat A$
est une extension séparable de $K$, condition toujours satisfaite si
$K$ est de caractéristique 0.

\paragraph{(Err\textsubscript{IV}, 28).}
Après (IV, 8.10.3), après la ligne 14 de la p.~37, insérer :

\begin{remark}[8.10.3.1]
\label{IV.8.10.3.1}
On ne peut dans {\hyperref[canon-index.ega4.c04.prop.8-10-3]{(8.10.3)}} supprimer
l'hypothèse que les $\overline{f_\alpha(X_\alpha)}$ soient
constructibles. Prenons par exemple $S=\operatorname{Spec}(A)$, où $A$
est un anneau tel que $S$ soit compact, totalement discontinu et sans
point isolé, tout idéal premier $\mathfrak p$ de $A$ étant maximal et
$A_\mathfrak p$ étant un corps (Bourbaki, \emph{Alg. comm.}, chap.~II,
\S4, exerc.~16 et 17). Soit $s\in S$, et prenons pour
$Y_\alpha=S_\alpha$ les voisinages ouverts affines de $s$ dans $S$; on
prend $X_\alpha=\operatorname{Spec}(k(s))$ pour tout $\alpha$,
$f_\alpha$ et $f$ étant les morphismes canoniques; il est clair que
$f$, qui est l'identité, est dominant, mais aucun des $f_\alpha$ ne
l'est.
\end{remark}

\paragraph{(Err\textsubscript{IV}, 29).}
Dans (IV, 8.12.8), ligne 7 de la p.~47, remplacer « tel qu'il existe un
$\sh{O}_Y$-Module ample » par : « et quasi-séparé ». Ligne 12 de la
p.~47, remplacer « L'hypothèse » par : « Il résulte de
{\hyperref[canon-index.ega4.c04.prop.8-12-3]{(8.12.3)}} que l'existence de la
factorisation de l'énoncé est une propriété locale sur $Y$; on peut
donc supposer $Y$ affine, ce qui ».

\paragraph{(Err\textsubscript{IV}, 30).}
Après (IV, 8.12.11), ajouter à la fin de la p.~48 :

\begin{corollary}[8.12.12]
\EGAReaderTarget{ega4.c04.cor.8-12-12}{IV.8.12.12}
Soient $X$, $Y$ deux préschémas intègres,
$f:X\to Y$ un morphisme séparé, birationnel et localement de type fini.
Soit $y\in Y$ un point normal de $Y$ et supposons qu'un point
$x\in f^{-1}(y)$ soit isolé dans $f^{-1}(y)$. Alors il existe un
voisinage ouvert $U$ de $y$ dans $Y$ tel que la restriction
$f^{-1}(U)\to U$ de $f$ soit un isomorphisme.
\end{corollary}

\begin{proof}
Il suffit de prouver qu'il existe un voisinage ouvert $U$ de $y$ dans
$Y$ et une $U$-section $g:U\to f^{-1}(U)$ de $f^{-1}(U)$ telle que
$g(y)=x$. En effet, comme $U$ contient le point générique $\eta$ de
$Y$ et $f^{-1}(U)$ le point générique $\xi$ de $X$, on a
nécessairement $g(\eta)=\xi$ puisque $f$ est birationnel
{\hyperref[canon-index.ega4.c04.statement.6-15-4]{(6.15.4)}}; comme $g$ est une
immersion, $g(U)$ est localement fermé dans $f^{-1}(U)$, donc ouvert
dans son adhérence dans $f^{-1}(U)$, et comme cette dernière est égale
à $f^{-1}(U)$ d'après ce qui précède, $g(U)$ est ouvert dans
$f^{-1}(U)$; comme le sous-préschéma de $f^{-1}(U)$ associé à $g$ est
isomorphe à $U$, donc réduit, et que $f^{-1}(U)$ est réduit, il résulte
de (I, 5.2.1) que $g$ est une immersion ouverte. D'autre part, puisque
$f$ est séparé, $g$ est une immersion fermée (I, 5.4.6), donc $g(U)$
est un ensemble à la fois ouvert et fermé dans $f^{-1}(U)$. Mais comme
$f^{-1}(U)$ est intègre, cela entraînera $g(U)=f^{-1}(U)$, d'où la
conclusion.

Considérons alors deux cas :

\textup{I)} Supposons d'abord $Y$ normal, et soit $V$ l'ensemble des
points de $X$ en lesquels $f$ est quasi-fini; c'est un ensemble ouvert
dans $X$ {\hyperref[canon-index.ega4.c04.cor.13-1-4]{(13.1.4)}}.\footnote{Le
lecteur vérifiera que ni {\hyperlink{ega4.c04.cor.8-12-12}{(8.12.12)}}
ni aucune de ses conséquences n'est utilisée dans la preuve de
{\hyperref[canon-index.ega4.c04.cor.13-1-4]{(13.1.4)}}.} Il résulte alors de
{\hyperref[canon-index.ega4.c04.cor.8-12-10]{(8.12.10)}}

\oldpage[IV]{353}
que la restriction $f|V:V\to Y$ est une immersion ouverte; donc il y a
un voisinage ouvert $U$ de $y\in Y$ tel que $f|V$ soit un isomorphisme
de $V$ sur le sous-préschéma induit par $Y$ sur $U$; l'isomorphisme
réciproque $g$ est alors une $U$-section de $f^{-1}(U)$ telle que
$g(y)=x$.

\textup{II)} L'obtention d'une $U$-section $g$ telle que $g(y)=x$ étant
un problème local sur $X$ et sur $Y$, on peut supposer que $X$ et $Y$
sont affines et que $f$ est de type fini, donc que $X$ est un
sous-préschéma fermé d'un $Y$-préschéma $Z$ de présentation finie sur
$Y$ (par exemple de la forme $Y[T_1,\ldots,T_n]$). Posons
\[
 Y'=\operatorname{Spec}(\sh{O}_{Y,y}),\qquad
 X'=X_{(Y')},\qquad Z'=Z_{(Y')},
\]
de sorte que $X'$ est un sous-préschéma fermé de $Z'$. En outre, si
$f'=f_{(Y')}$, on a $f'^{-1}(y)=f^{-1}(y)$ (I, 3.6.5) et $x$ est donc
isolé dans $f'^{-1}(y)$. Puisque $Y'$ est normal, on peut appliquer le
résultat de I) à $f'$, et puisque $Y'$ est le seul voisinage de $y$ dans
$Y'$, il existe une $Y'$-section $g'$ de $X'$ telle que $g'(y)=x$, que
l'on peut évidemment aussi considérer comme une $Y'$-section de $Z'$.
Puisque le morphisme structural $h:Z\to Y$ est de présentation finie,
on peut appliquer ({\hyperref[canon-index.ega4.c04.thm.8-8-2]{8.8.2}}, (i) et
(ii)) suivant la méthode de
({\hyperref[canon-index.ega4.c04.statement.8-1-2]{8.1.2}}, a)), d'où l'on conclut
à l'existence d'un voisinage ouvert $U$ de $y$ dans $Y$ et d'une
$U$-section $g$ de $h^{-1}(U)$ telle que $g(y)=x$. Il reste à voir
qu'en fait $g$ est une $U$-section de $f^{-1}(U)$. Mais puisque
$f^{-1}(U)$ est un sous-préschéma fermé de $h^{-1}(U)$ et que $U$ est
réduit, il suffit pour cela (I, 5.2.2) de voir que
$g(U)\subset f^{-1}(U)$. Or, puisque $U$ est intègre et $f^{-1}(U)$
fermé dans $h^{-1}(U)$, il suffit de prouver que si $\eta$ est le point
générique de $Y$ (et de $U$), on a $g(\eta)\in f^{-1}(U)$. Or, cela
résulte de ce que la restriction de $g$ à $Y'$ est une $Y'$-section de
$X'\subset f^{-1}(U)$ et de ce que $\eta\in Y'$.
\end{proof}

\paragraph{(Err\textsubscript{IV}, 31).}
Dans ({\hyperref[canon-index.ega4.c04.subsec.10-4]{IV, 10.4}}), ajouter, après la
proposition {\hyperref[canon-index.ega4.c04.statement.10-4-11]{(10.4.11)}} :

\begin{corollary}[10.4.12]
\EGAReaderTarget{ega4.c04.cor.10-4-12}{IV.10.4.12}
Soient $S$ un préschéma, $X$ un $S$-préschéma de présentation finie,
$f$ un $S$-endomorphisme radiciel de $X$.
\begin{enumerate}[label=\textup{(\roman*)}]
\item L'endomorphisme $f$ est bijectif et fini (donc un homéomorphisme
universel {\hyperref[canon-index.ega4.c04.prop.2-4-5]{(2.4.5)}}).
\item Supposons de plus que $X$ soit réduit, que les anneaux locaux de
$S$ soient noethériens et universellement japonais, et que $f$ soit
birationnel (condition automatiquement vérifiée si pour tout point
maximal $x$ de $X$, $k(x)$ est un corps parfait). Alors $f$ est un
$S$-automorphisme.
\end{enumerate}
\end{corollary}

\begin{proof}
\textup{(i)} La question étant locale sur $S$, on peut se borner au cas
où $S=\operatorname{Spec}(A)$ est affine. Il existe alors une
sous-$\mathbb Z$-algèbre de type fini $A_0$ de $A$ telle qu'en posant
$S_0=\operatorname{Spec}(A_0)$, on ait
$X=X_0\times_{S_0}S$, où $X_0$ est un préschéma de type fini sur $S_0$,
et que $f$ se déduise d'un $S_0$-endomorphisme $f_0$ de $X_0$ par le
changement de base $S\to S_0$
{\hyperref[canon-index.ega4.c04.prop.8-9-1]{(8.9.1)}}; on peut en outre supposer
$f_0$ radiciel ({\hyperref[canon-index.ega4.c04.thm.8-10-5]{8.10.5}}, (vii)). Il
suffira de prouver que $f_0$ est bijectif et fini, et on peut donc
supposer désormais que $A$ est une $\mathbb Z$-algèbre de type fini, ce
qui entraîne que $S$ et $X$ sont des préschémas excellents
({\hyperref[canon-index.ega4.c04.statement.7-8-3]{(7.8.3)}} et
{\hyperref[canon-index.ega4.c04.prop.7-8-6]{(7.8.6)}}) et en particulier que $X$
est un préschéma universellement japonais
({\hyperref[canon-index.ega4.c04.cor.5-11-4]{(5.11.4)}} et
({\hyperref[canon-index.ega4.c04.statement.7-8-3]{7.8.3}}, (vi))). On peut en outre
supposer $X$ réduit; en effet,
$f_{\mathrm{red}}:X_{\mathrm{red}}\to X_{\mathrm{red}}$ est un
$S_{\mathrm{red}}$-endomorphisme radiciel (I, 5.1.6); si l'on prouve que
$f_{\mathrm{red}}$ est fini, il en sera de même de $f$ : en effet,
puisque $f$ est radiciel, il est séparé
{\hyperref[canon-index.ega4.c04.lem.1-8-7-1]{(1.8.7.1)}}; comme
$f_{\mathrm{red}}$ est propre il en est de même de $f$ (II, 5.4.6);
mais puisque $f$ est évidemment quasi-fini, il est fini (III, 4.4.2).

\oldpage[IV]{354}
Soit alors $X'$ le normalisé de $X$ (II, 6.3.8); si $X_i$
($1\leq i\leq n$) sont les sous-préschémas réduits de $X$ ayant pour
espaces sous-jacents les composantes irréductibles de $X$, $X'$ est
somme des normalisés $X'_i$ des préschémas intègres $X_i$; nous
désignerons par $x_i$ (resp. $x'_i$) le point générique de $X_i$
(resp. $X'_i$). Le morphisme canonique $p:X'\to X$ est fini puisque
$X$ est universellement japonais et on a
$p^{-1}(x_i)=\{x'_i\}$, l'homomorphisme $k(x_i)\to k(x'_i)$ étant
bijectif.

Soit $f=(\psi,\theta)$; on sait déjà
{\hyperref[canon-index.ega4.c04.statement.10-4-11]{(10.4.11)}} que $\psi$ est une
bijection continue de l'espace $X$ sur lui-même; si l'on pose
$y_i=\psi(x_i)$, $y_i$ est aussi un point maximal de $X$, car si $y'_i$
était une générisation de $y_i$ distincte de $y_i$, $x_i$ serait une
générisation $\neq x_i$ de $\psi^{-1}(y'_i)$. Puisque l'ensemble $M$
des points maximaux de $X$ est fini, la restriction $\psi|M$ est une
bijection de $M$ sur lui-même, et on peut écrire
$\psi(x_i)=x_{\sigma(i)}$, où $\sigma$ est une permutation de
$[1,n]$. Alors, pour tout $i$, la restriction $f|X_i$ de $f$ se
factorise en
\[
 X_i\xrightarrow{f_i}X_{\sigma(i)}\longrightarrow X
\]
(I, 5.2.2), où la seconde flèche est l'injection canonique. Comme les
$X_i$ sont intègres, il existe (II, 6.3.9) un unique morphisme
$f'_i:X'_i\to X'_{\sigma(i)}$ rendant commutatif le diagramme
\[
\begin{array}{ccc}
X'_i&\xrightarrow{f'_i}&X'_{\sigma(i)}\\
\downarrow p_i&&\downarrow p_{\sigma(i)}\\
X_i&\xrightarrow{f_i}&X_{\sigma(i)}
\end{array}
\]
et si $f'$ coïncide avec $f'_i$ dans chacun des $X'_i$, on a donc un
diagramme commutatif
\phantomsection\EGAReaderTarget{ega4.c04.sec.10-4-12-1}{IV.10.4.12.1}
\begin{equation}
\label{IV.10.4.12.1.errata}
\begin{array}{ccc}
X'&\xrightarrow{f'}&X'\\
\downarrow p&&\downarrow p\\
X&\xrightarrow{f}&X.
\end{array}
\tag{10.4.12.1}
\end{equation}

D'autre part, comme $f$ est séparé, il en est de même de
$f\circ p=p\circ f'$, donc $f'$ est séparé (I, 5.5.1); enfin, puisque
$p$ est fini et $f$ bijectif, il résulte de
{\hyperlink{ega4.c04.sec.10-4-12-1}{(10.4.12.1)}} que $f'$ est
quasi-fini. Puisque les $X'_i$ sont intègres et normaux, il résulte de
{\hyperref[canon-index.ega4.c04.cor.8-12-11]{(8.12.11)}} que $f'_i$ se factorise en
\[
 X'_i\xrightarrow{f''_i}X''_{\sigma(i)}
 \xrightarrow{u_{\sigma(i)}}X'_{\sigma(i)},
\]
où $X''_{\sigma(i)}$ est la fermeture intégrale de
$X'_{\sigma(i)}$ relative à l'extension $k(x'_i)$ de
$k(x'_{\sigma(i)})$, et $f''_i$ une immersion ouverte. Mais comme
$k(x'_i)$ est une extension radicielle de $k(x'_{\sigma(i)})$,
$u_{\sigma(i)}$ est un morphisme radiciel (Bourbaki,
\emph{Alg. comm.}, chap.~V, \S2, n\textsuperscript{o}~3, lemme 4);
donc $f'$ est radiciel, et il résulte alors de
{\hyperref[canon-index.ega4.c04.statement.10-4-11]{(10.4.11)}} que $f'$ est
surjectif; on en conclut (puisque les $u_{\sigma(i)}$ sont injectifs)
que $f''_i$ est bijectif, donc un isomorphisme; on en déduit enfin que
le morphisme $f'$ est fini (puisqu'il en est ainsi de $u_{\sigma(i)}$,
$X'$ étant universellement japonais
{\hyperref[canon-index.ega4.c04.statement.7-8-3]{(7.8.3)}}). Cela étant, comme
$p\circ f'=f\circ p$ est alors fini et que $p$ est surjectif et $f$
séparé, $f$ est propre (II, 5.4.3) et quasi-fini, donc fini
{\hyperref[canon-index.ega4.c04.thm.8-11-1]{(8.11.1)}}.

\oldpage[IV]{355}
\textup{(ii)} Compte tenu de
({\hyperref[canon-index.ega4.c04.thm.8-10-5]{8.10.5}}, (i)) appliqué suivant la
méthode de ({\hyperref[canon-index.ega4.c04.statement.8-1-2]{8.1.2}}, a)), on peut
se borner au cas où $S$ est un schéma local. Les réductions préalables
faites dans la preuve de (i) sont donc ici inutiles, et avec les mêmes
notations on voit que l'on a un diagramme commutatif
{\hyperlink{ega4.c04.sec.10-4-12-1}{(10.4.12.1)}} dans lequel $X$ est
un $S$-préschéma de type fini, $p$ est fini, $f'$ est un isomorphisme et
l'homomorphisme canonique $\sh{O}_X\to p_*(\sh{O}_{X'})$ est injectif.
Mais on montrera au chap.~VI, dans la théorie des quotients des
préschémas par des relations d'équivalence, que sous ces conditions $f$
est lui-même un isomorphisme (le lecteur vérifiera qu'on n'utilisera pas
{\hyperlink{ega4.c04.cor.10-4-12}{(10.4.12)}} dans la théorie des
relations d'équivalence du chap.~VI).
\end{proof}

\begin{corollary}[10.4.13]
\EGAReaderTarget{ega4.c04.cor.10-4-13}{IV.10.4.13}
Soient $k$ un corps, $X$ un préschéma algébrique sur $k$. Alors tout
$k$-endomorphisme radiciel $f$ de $X$ est fini et bijectif, donc un
homéomorphisme universel. Si de plus $X$ est réduit et $f$ birationnel
(par exemple si $k$ est de caractéristique 0), $f$ est un
$k$-automorphisme.
\end{corollary}

\begin{remark}[10.4.14]
\label{IV.10.4.14}
\begin{enumerate}[label=\textup{(\roman*)}]
\item Dans {\hyperlink{ega4.c04.cor.10-4-12}{(10.4.12)}} et
{\hyperlink{ega4.c04.cor.10-4-13}{(10.4.13)}} on ne peut supprimer
l'hypothèse que $X$ est réduit pour la seconde conclusion; c'est ce que
montre l'exemple du $B$-endomorphisme correspondant à l'augmentation
$D_B(L)\to B$ de l'anneau $D_B(L)=A$ défini dans ($0$, 18.2.3), où
$B$ est identifié à un sous-anneau de $A$.

\item Soit $X$ un préschéma noethérien. Nous ignorons si tout
endomorphisme radiciel $f$ de $X$ est nécessairement surjectif et fini.
Si $f$ est supposé surjectif et si de plus $X$ est un préschéma
japonais, le raisonnement de
{\hyperlink{ega4.c04.cor.10-4-12}{(10.4.12)}} montre que $f$ est fini;
si de plus $X$ est réduit et si $f$ est birationnel, par exemple si,
pour tout point maximal $x\in X$, $k(x)$ est parfait, le même
raisonnement montre que $f$ est un automorphisme.

\item On peut donner des exemples d'endomorphismes radiciels non
surjectifs de schémas affines intègres non noethériens : si $k$ est un
anneau, $A=k[(T_\lambda)_{\lambda\in L}]$ un anneau de polynômes où $L$
est infini dénombrable, et $\mathfrak J$ l'idéal engendré dans $A$ par
une sous-famille $(T_\lambda)_{\lambda\in H}$, où $L-H$ est infini et
$H\neq\varnothing$, on remarque que $A/\mathfrak J$ est isomorphe à
$k[(T_\lambda)_{\lambda\in L-H}]$, donc à $A$; l'endomorphisme de
$\operatorname{Spec}(A)$ correspondant à l'endomorphisme d'anneaux
composé
\[
 A\xrightarrow{p}A/\mathfrak J\xrightarrow{u}A
\]
où $p$ est la projection canonique et $u$ un isomorphisme, répond à la
question, son image étant le sous-préschéma fermé de
$\operatorname{Spec}(A)$ défini par $\mathfrak J$.
\end{enumerate}
\end{remark}

\paragraph{(Err\textsubscript{IV}, 32).}
Dans (IV, 11.3.14), ligne 7 de la p.~141, avant « Soient », insérer :
(i). Après la ligne 9 de la p.~141, ajouter :

\smallskip
\noindent
\textup{(ii)} Inversement, soit $f:X\to Y$ un morphisme plat. Si $X$
est unibranche (resp. géométriquement unibranche) en un point $x$, $Y$
est unibranche (resp. géométriquement unibranche) au point $y=f(x)$.

Au début de la ligne 10 de la p.~141, insérer : (i). Après la ligne 20
de la p.~141, ajouter :

\smallskip
\noindent
\textup{(ii)} Notons que si $Y_1=Y_{\mathrm{red}}$,
$X_1=X\times_Y Y_1$ a pour espace sous-jacent celui de $X$ et le
morphisme $f_1:X_1\to Y_1$ déduit de $f$ par changement de base est
plat. On peut donc se borner au cas où $Y=\operatorname{Spec}(A)$,
$A$ étant un anneau local réduit, et $y$ le point fermé de $Y$. Posons
$B=\sh{O}_{X,x}$; par hypothèse $B_{\mathrm{red}}$ est unibranche, donc
$\operatorname{Spec}(B)$ n'a qu'un seul point maximal, et par platitude
{\hyperref[canon-index.ega4.c04.prop.2-3-4]{(2.3.4)}} on en déduit qu'il en est de
même

\oldpage[IV]{356}
de $\operatorname{Spec}(A)$; puisque $A$ est réduit, il est donc
intègre. Soient $K$ le corps des fractions de $A$, $A'$ la clôture
intégrale de $A$; par platitude, $B\otimes_AA'$ s'identifie à un
sous-anneau de $B\otimes_AK$ contenant $B$, et comme les éléments non
nuls de $A$ sont $B$-réguliers par platitude, $B\otimes_AK$ s'identifie
à un sous-anneau de l'anneau total des fractions $R$ de $B$. Pour
montrer que $A$ est unibranche, il suffit de voir que $A'$ ne peut avoir
deux idéaux maximaux distincts (nécessairement au-dessus de $y$
(Bourbaki, \emph{Alg. comm.}, chap.~V, \S2,
n\textsuperscript{o}~1, prop.~1)); s'il en était ainsi,
$B'=B\otimes_AA'$ aurait deux idéaux premiers distincts au-dessus de
$x$ (I, 3.4.7), et il suffit de voir que ce n'est pas possible.

Soit $\mathfrak N$ le nilradical de $B$; si $S$ est l'ensemble des
éléments réguliers de $B$, on a $R=S^{-1}B$, et
$S^{-1}\mathfrak N=R\mathfrak N$ est le nilradical de $R$, donc
\[
 R/R\mathfrak N=S^{-1}B/S^{-1}\mathfrak N=S^{-1}(B/\mathfrak N)
\]
est contenu dans le corps des fractions $L$ de l'anneau local intègre
$B/\mathfrak N=B_{\mathrm{red}}$; il en est donc de même de
\[
 B'_{\mathrm{red}}=B'/(B'\cap R\mathfrak N)
\]
qui est entier sur $B_{\mathrm{red}}$; par hypothèse
$B'_{\mathrm{red}}$, étant contenu dans la clôture intégrale de
$B_{\mathrm{red}}$, est un anneau local, donc il en est de même de
$B'$, ce qui établit notre assertion.

Si l'on suppose en outre $B$ géométriquement unibranche, le corps
résiduel de $B'_{\mathrm{red}}$, étant contenu dans le corps résiduel de
la clôture intégrale de $B_{\mathrm{red}}$, est une extension radicielle
du corps résiduel $k(x)$ de $B$. Notons maintenant que le corps résiduel
de $B'_{\mathrm{red}}$ est aussi celui de $B'=B\otimes_AA'$, donc est
égal à $k(x)\otimes_{k(y)}k'$, où $k'$ est le corps résiduel de $A'$;
on déduit alors de {\hyperref[canon-index.ega4.c04.prop.2-6-1]{(2.6.1)}} que $k'$
est une extension radicielle de $k(y)$, donc $A$ est aussi
géométriquement unibranche.

\paragraph{(Err\textsubscript{IV}, 33).}
Dans ({\hyperref[canon-index.ega4.c04.item.12-1-2]{IV, 12.1.2}}), l'exemple donné
en (i) n'est pas correct, le morphisme considéré étant ouvert. On
construit un exemple d'un morphisme plat $f:X\to Y$ dont la restriction
à une composante irréductible de $X$ n'est pas un morphisme ouvert de
la façon suivante. Soit $Z=\operatorname{Spec}(\mathbb C[U,V])$ le
plan affine sur le corps des nombres complexes, et considérons le
schéma $Y$ obtenu en recollant deux points fermés distincts $a$, $b$ de
$Z$; on a $Y=\operatorname{Spec}(B)$, où $B$ est le sous-anneau de
$\mathbb C[U,V]$ formé des polynômes prenant la même valeur aux points
$a$ et $b$ (l'ensemble des points fermés de $Z$ étant identifié à
$\mathbb C^2$); il est immédiat que $B$ est un sous-$\mathbb C$-espace
vectoriel de $\mathbb C[U,V]$ de codimension 1, donc $\mathbb C[U,V]$
est un $B$-module de type fini, et par suite le morphisme $Z\to Y$ est
fini, l'image réciproque du point $c$ de $Y$ obtenu par recollement de
$a$ et de $b$ étant l'ensemble $\{a,b\}$.

Considérons alors le schéma $X$ obtenu en recollant deux exemplaires de
$Z$, soient $Z_1$, $Z_2$, de sorte que si $a_i$, $b_i$ ($i=1,2$) sont
les points de $Z_i$ correspondant à $a$ et $b$, $a_1$ soit recollé à
$b_2$ et $b_1$ à $a_2$ (cf. ({\hyperref[canon-index.ega4.c04.item.6-5-5]{6.5.5}},
(ii))). On voit aisément (\emph{loc. cit.}) que le morphisme $X\to Y$
ainsi obtenu est plat et fini (et même étale) et que $Z_1$ et $Z_2$
s'identifient aux deux composantes irréductibles de $X$. Cependant, si
$D$ est une droite dans $Z_1$ passant par $a_1$ mais non par $b_1$,
l'image $V$ dans $Y$ de l'ouvert $Z_1-D$ de $Z_1$ contient l'image de
$b_1$, c'est-à-dire le point $c$, mais ne contient pas l'image du point
générique de $D$, donc $V$ n'est pas un voisinage de $c$.

\paragraph{(Err\textsubscript{IV}, 34).}
Dans le titre de (IV, 13.2), ligne 5 du bas de la p.~190, remplacer :
« équidimensionnels » par « localement équidimensionnels ». De même
dans (IV, 13.2.2), ligne 4 de la p.~191, remplacer deux fois
« équidimensionnel » par « localement équidimensionnel ».

\oldpage[IV]{357}
Ligne 5 de la p.~191, ajouter : \emph{Si de plus les fibres
$f^{-1}(f(x))$ de $f$ sont équidimensionnelles ($0_{\mathrm{IV}}$,
14.1.3), on dit que $f$ est équidimensionnel.} Ligne 6 du bas de la
p.~194, remplacer « équidimensionnelles » par « localement
équidimensionnelles ».

\paragraph{(Err\textsubscript{IV}, 35).}
Dans le titre de (IV, 13.3), ligne 4 du bas de la p.~94, remplacer
« équidimensionnels » par « localement équidimensionnels ». Dans
(IV, 13.3.2), ligne 6 du bas de la p.~196, remplacer deux fois
« équidimensionnel » par « localement équidimensionnel ». Ligne 5 du bas
de la p.~196, ajouter : \emph{Si de plus les fibres $f^{-1}(f(x))$ de
$f$ sont équidimensionnelles ($0_{\mathrm{IV}}$, 14.1.3), on dit que
$f$ est équidimensionnel.} Dans (IV, 13.3.7), ligne 14 de la p.~198,
remplacer « équidimensionnel » par « localement équidimensionnel ».
Dans ({\hyperref[canon-index.ega4.c04.statement.13-3-9]{IV, 13.3.9}}), ligne 3 du
bas de la p.~198, remplacer « équidimensionalité » par
« équidimensionalité locale ». Lignes 10, 12, 20 et 24 de la p.~199,
remplacer « équidimensionnel » par « localement équidimensionnel ».

\paragraph{(Err\textsubscript{IV}, 36).}
Dans l'introduction du \S14, ligne 6 du bas de la p.~199, remplacer
« équidimensionnel » par « localement équidimensionnel ».

\paragraph{(Err\textsubscript{IV}, 37).}
Dans (IV, 14.4.1), ligne 5 du bas de la p.~209, remplacer « localement
de type fini » par « localement de présentation finie ». L'énoncé
actuel de ({\hyperref[canon-index.ega4.c04.thm.14-4-1]{IV, 14.4.1}}) est inexact,
même lorsque $f:X\to Y$ est une immersion fermée, $X$ étant réduit à
un seul point $y$ (fermé) de $Y$, et $\sh{O}_{Y,y}$ étant un corps. Il
y a en effet des anneaux $A$ tels que $Y=\operatorname{Spec}(A)$ soit
compact, sans point isolé et totalement discontinu, tout idéal premier
$\mathfrak p$ de $A$ étant maximal et $A_\mathfrak p$ étant un corps
(Bourbaki, \emph{Alg. comm.}, chap.~II, \S4, exerc.~16 et 17); il est
clair que l'immersion fermée $\{y\}\to Y$ n'est pas un morphisme
ouvert, mais vérifie les conditions b) et b') de
({\hyperref[canon-index.ega4.c04.thm.14-4-1]{IV, 14.4.1}}).

\paragraph{(Err\textsubscript{IV}, 38).}
Dans (IV, 14.4.1.3), ligne 14 de la p.~211, remplacer « localement
quasi-fini et dominant » par « localement quasi-fini, localement de
présentation finie et dominant ».

\paragraph{(Err\textsubscript{IV}, 39).}
Dans (IV, 14.4.2), ligne 10 du bas de la p.~211, remplacer « localement
de type fini » par « localement de présentation finie ».

\paragraph{(Err\textsubscript{IV}, 40).}
Dans (IV, 14.4.3), ligne 18 de la p.~212, remplacer « localement de
type fini » par « localement de présentation finie ».

\paragraph{(Err\textsubscript{IV}, 41).}
Dans ({\hyperref[canon-index.ega4.c04.cor.14-4-4]{IV, 14.4.4}}), lignes 15 et 16 du
bas de la p.~212, remplacer « localement de type fini » par
« localement de présentation finie ». Lignes 10 et 11 du bas, remplacer
« équidimensionnel » par « localement équidimensionnel ».

\paragraph{(Err\textsubscript{IV}, 42).}
Dans (IV, 14.4.8), après la ligne 1 de la p.~214, ajouter :
\begin{quote}
\emph{Si de plus l'ensemble des composantes irréductibles de $Y$ est
localement fini, ces deux conditions sont équivalentes aux suivantes :}
\end{quote}

\paragraph{(Err\textsubscript{IV}, 43).}
Dans (IV, 14.4.10), ligne 15 de la p.~215, remplacer : quelconque, par :
dont l'ensemble des composantes irréductibles est localement fini.
Ligne 16 de la p.~215, remplacer : de type fini, par : de présentation
finie.

\paragraph{(Err\textsubscript{IV}, 44).}
Dans (IV, 15.5.6.1), ligne 15 du bas de la p.~234, après la première
phrase, ajouter :

\begin{quote}
On peut se limiter au cas où dans $X$ toute partie fermée irréductible
n'admet qu'un seul point générique. En effet, dans le cas contraire on
considère l'espace $Y=X/R$,

\oldpage[IV]{358}
quotient de $X$ par la relation d'équivalence
$\overline{\{x\}}=\overline{\{x'\}}$. Il résulte aussitôt de cette
définition que tout ouvert de $X$ est saturé pour $R$, donc la topologie
de $X$ est l'image réciproque de celle de $Y$ par l'application
canonique $\pi:X\to X/R$. Un ouvert dense dans $X$ est de la forme
$\pi^{-1}(U)$, où $U$ est un ouvert dense dans $Y$, donc un fermé rare
$Z$ dans $X$ est de la forme $\pi^{-1}(S)$, où $S=\pi(Z)$ est fermé
rare dans $Y$, et il revient au même de dire que $X-Z$ est connexe ou
que $Y-S$ est connexe. De même $\pi(T_x)$ est l'ensemble des
générisations de $\pi(x)$, et il est équivalent de dire que
$T_x-\{x\}$ est connexe ou que $\pi(T_x)-\{\pi(x)\}$ l'est. Enfin, par
construction $Y$ est un espace de Kolmogoroff localement noethérien,
dont toute partie fermée admet donc un seul point générique. On voit
alors qu'il revient au même de prouver le lemme pour $X$ ou pour $Y$,
et on peut donc supposer que toute partie fermée irréductible de $X$ a
un point générique unique.
\end{quote}

\paragraph{(Err\textsubscript{IV}, 45).}
Dans (IV, 15.6.1), ligne 18 du bas de la p.~236, supprimer « localement
noethérien » (hypothèse non utilisée dans la démonstration).

\paragraph{(Err\textsubscript{IV}, 46).}
Dans (IV, 15.6.3), ligne 11 de la p.~237, au lieu de « supposons que
pour tout $y\in Y$ », lire : « supposons que $f$ soit localement de
présentation finie et que pour tout $y\in Y$ ». Remplacer la
démonstration (lignes 14 à 24 de la p.~237) par la suivante :
\begin{quote}
En vertu de {\hyperref[canon-index.ega4.c04.prop.1-10-3]{(1.10.3)}}, il suffit de
prouver que pour tout $x\in X_y^0$ et toute générisation $y'$ de $y$,
il existe une générisation $x'$ de $x$ telle que $f(x')=y'$. Or il
suffit de prendre pour $x'$ le point générique de $X_{y'}^0$, en vertu
de {\hyperref[canon-index.ega4.c04.cor.15-6-2]{(15.6.2)}}.
\end{quote}

\paragraph{(Err\textsubscript{IV}, 47).}
Dans (IV, 15.6.7), ligne 2 de la p.~241, après « noethérien », ajouter :
« ou si la condition $\beta$) de
({\hyperref[canon-index.ega4.c04.prop.15-6-6]{15.6.6}}, (iii)) est vérifiée ».
Après la ligne 9 de la p.~241, ajouter :
\begin{quote}
Lorsque la condition $\beta$) de
({\hyperref[canon-index.ega4.c04.prop.15-6-6]{15.6.6}}, (iii)) est satisfaite, le
fait que a) entraîne b') résulte de
({\hyperref[canon-index.ega4.c04.prop.15-6-6]{15.6.6}}, (ii)), et le fait que b')
entraîne a) résulte de
({\hyperref[canon-index.ega4.c04.prop.15-6-6]{15.6.6}}, (iii)).
\end{quote}

\paragraph{(Err\textsubscript{IV}, 48).}
Dans (IV, 20.1.5), ligne 13 de la p.~228, (IV, 20.1.6), ligne 17 de la
p.~228, (IV, 20.1.7), ligne 13 du bas de la p.~228, supprimer
« strictement ».

\paragraph{(Err\textsubscript{IV}, 49).}
Dans (IV, 20.2.13), ajouter après la ligne 6 de la p.~236 : En
particulier si $X$ est réduit et n'a qu'un nombre fini de composantes
irréductibles, la notion de $\sh{O}_X$-Module sans torsion définie dans
({\hyperlink{ega4.c04.statement.20-1-5}{IV, 20.1.5}}) coïncide avec
celle de (I, 7.4.1).

\paragraph{(Err\textsubscript{IV}, 50).}
Dans (IV, 20.2.14), lignes 18 et 25 de la p.~236, supprimer
« strictement ».

\paragraph{(Err\textsubscript{IV}, 51).}
Dans (IV, 20.6.2), ligne 13 du bas de la p.~252, supprimer
« strictement ».

\paragraph{(Err\textsubscript{IV}, 52).}
Dans (IV, 20.6.4), lignes 23 et 22 de la p.~253, supprimer
« strictement ».

\paragraph{(Err\textsubscript{IV}, 53).}
Dans (IV, 21.4), après la ligne 23 de la p.~267, ajouter :

La proposition suivante et ses corollaires nous ont été signalés par
M.~Raynaud.

\begin{proposition}[21.4.9]
\EGAReaderTarget{ega4.c04.prop.21-4-9}{IV.21.4.9}
Soient $X$, $Y$ deux préschémas localement noethériens,
$f:X\to Y$ un morphisme fidèlement plat; on suppose en outre, soit que
$f$ est quasi-compact, soit que $f$ est ouvert. Pour qu'un diviseur
$D'$ sur $X$ soit de la forme $f^*(D)$, où $D$ est un diviseur sur
$Y$, il faut et il suffit que pour tout $y\in Y$ tel que
$\operatorname{prof}(\sh{O}_{Y,y})\leq1$, en posant
\[
 Y_1=\operatorname{Spec}(\sh{O}_{Y,y}),\qquad
 X_1=X\times_YY_1,\qquad
 f_1=f_{(Y_1)}:X_1\to Y_1,
\]
l'image réciproque $D'_1$ de $D'$ sur $X_1$ soit de la forme
$f_1^*(D_1)$, où $D_1\in\operatorname{Div}(Y_1)$.
\end{proposition}

\oldpage[IV]{359}
\begin{proof}
La nécessité de la condition est évidente par transitivité des images
réciproques de diviseurs, tous les morphismes considérés étant plats
{\hyperlink{ega4.c04.statement.21-4-4}{(21.4.4)}}. Pour tout ouvert
$U$ de $Y$, notons $f_U$ la restriction $f^{-1}(U)\to U$ de $f$; en
vertu de {\hyperlink{ega4.c04.prop.21-4-8}{(21.4.8)}} l'application
\[
 D_U\longmapsto f_U^*(D_U)
\]
de $\operatorname{Div}(U)$ dans $\operatorname{Div}(f^{-1}(U))$ est
injective; donc, si $U_1$, $U_2$ sont deux ouverts de $Y$, et si
\[
 D'|f^{-1}(U_1)=f_{U_1}^*(D_1),\qquad
 D'|f^{-1}(U_2)=f_{U_2}^*(D_2),
\]
on a nécessairement
$D_1|U_1\cap U_2=D_2|U_1\cap U_2$. Ceci montre aussitôt qu'il existe
un plus grand ouvert $U\subset Y$ tel que
\[
 D'|f^{-1}(U)=f_U^*(D_U),
 \qquad D_U\in\operatorname{Div}(U);
\]
il s'agit de voir que $U=Y$. Supposons le contraire, et soit $y$ un
point maximal de $Y-U$; on va montrer qu'il existe un voisinage ouvert
$V$ de $y$ dans $Y$ tel que
$D'|f^{-1}(V)=f_V^*(D_V)$ pour un $D_V\in\operatorname{Div}(V)$, ce qui
contredira le caractère maximal de $U$.

\textup{I) Cas où $f$ est quasi-compact.} Appliquant le fait que $f$
est quasi-compact et {\hyperlink{ega4.c04.prop.20-3-8}{(20.3.8)}}
suivant la méthode de ({\hyperref[canon-index.ega4.c04.statement.8-1-2]{8.1.2}},
a)), on est ramené à prouver qu'avec les notations de l'énoncé (mais
sans supposer d'abord de limitation sur
$\operatorname{prof}(\sh{O}_{Y,y})$), $D'_1$ est de la forme
$f_1^*(D_1)$. Or, si $\operatorname{prof}(\sh{O}_{Y,y})\leq1$, cette
assertion n'est autre que l'hypothèse. Supposons donc
$\operatorname{prof}(\sh{O}_{Y,y})\geq2$; puisque $y$ est point maximal
de $Y-U$, l'ouvert $W=Y_1\cap U$ de $Y_1$ est égal à $Y_1-\{y\}$, et
on peut se borner au cas où $Y=Y_1$, $U=Y-\{y\}$. Par hypothèse, il
existe un diviseur $D_U\in\operatorname{Div}(U)$ tel que
$f_U^*(D_U)=D'|f^{-1}(U)$; $D_U$ est classe d'équivalence d'un couple
$(\sh{L}_U,s)$, où $\sh{L}_U$ est un $\sh{O}_U$-Module inversible et
$s$ une section méromorphe régulière de $\sh{L}_U$ au-dessus de $U$
{\hyperlink{ega4.c04.prop.21-2-11}{(21.2.11)}}; par suite
{\hyperlink{ega4.c04.statement.21-4-1}{(21.4.1)}},
$D'|f^{-1}(U)$ est classe d'équivalence du couple
$(\sh{L}'_U,s')$, où
\[
 \sh{L}'_U=f_U^*(\sh{L}_U),\qquad s'=s\circ f_U.
\]
Notons $i:U\to Y$, $j:f^{-1}(U)\to X$ les injections canoniques, et
posons
\[
 \sh{L}=i_*(\sh{L}_U),\qquad
 \sh{L}'=j_*(\sh{L}'_U).
\]
Comme $f$ est plat, $f^*(\sh{L})$ est canoniquement isomorphe à
$\sh{L}'$ {\hyperref[canon-index.ega4.c04.lem.2-3-1]{(2.3.1)}}; d'autre part,
puisque $f$ est plat, la relation
$\operatorname{prof}(\sh{O}_{Y,y})\geq2$ entraîne
$\operatorname{prof}(\sh{O}_{X,x})\geq2$ pour tout $x\in f^{-1}(y)$
{\hyperref[canon-index.ega4.c04.prop.6-3-1]{(6.3.1)}}.

Or, $D'$ est classe d'équivalence d'un couple $(\sh{L}'',s'')$, où
$\sh{L}''$ est un $\sh{O}_X$-Module inversible et $s''$ une section
méromorphe régulière de $\sh{L}''$ au-dessus de $X$; $\sh{L}'_U$ est
par suite isomorphe à $\sh{L}''|f^{-1}(U)$, et il résulte alors de ce
qui précède et de {\hyperref[canon-index.ega4.c04.prop.5-9-8]{(5.9.8)}} et
{\hyperref[canon-index.ega4.c04.thm.5-10-5]{(5.10.5)}} que $\sh{L}'$ est un
$\sh{O}_X$-Module inversible isomorphe à $\sh{L}''$. Puisque $f$ est
fidèlement plat et quasi-compact, on en conclut que $\sh{L}$ est un
$\sh{O}_Y$-Module inversible
{\hyperref[canon-index.ega4.c04.prop.2-5-2]{(2.5.2)}} et par suite isomorphe à
$\sh{O}_Y$, puisque $Y$ est un schéma local. D'ailleurs, puisque
$\operatorname{prof}(\sh{O}_{Y,y})\geq1$, $U=Y-\{y\}$ est
schématiquement dense dans $Y$
({\hyperlink{ega4.c04.statement.20-2-13}{20.2.13}}, (iv)) et par suite
la section méromorphe régulière $s$ de $\sh{L}_U$ au-dessus de $U$ se
prolonge en une section méromorphe régulière $s_0$ de $\sh{L}$
au-dessus de $Y$ {\hyperlink{ega4.c04.prop.20-2-11}{(20.2.11)}}; la
classe d'équivalence de $(\sh{L},s_0)$ est donc un diviseur $D$ sur
$Y$; et il est clair que $f^*(D)$ et $D'$ induisent le même diviseur sur
$f^{-1}(U)=X-f^{-1}(y)$; mais comme
$\operatorname{prof}(\sh{O}_{X,x})\geq2$ pour tout $x\in f^{-1}(y)$,
on a $D'=f^*(D)$ {\hyperlink{ega4.c04.prop.21-1-8}{(21.1.8)}}, ce qui
achève la démonstration dans ce cas.

\textup{II) Cas où $f$ est ouvert.} Comme il s'agit de prouver
l'existence d'un voisinage ouvert $V$ de $y$ dans $Y$ tel que
$D'|f^{-1}(V)=f_V^*(D_V)$ pour un $D_V\in\operatorname{Div}(V)$, on
peut déjà supposer $Y$ affine. Puisque $f$ est ouvert, les images
$f(W)$, où $W$ parcourt les ouverts affines de $X$, forment un
recouvrement ouvert de $Y$, donc, puisque $Y$ est quasi-compact, il
existe un ouvert quasi-compact $T\subset X$ tel que $f(T)=Y$, et le
morphisme $f|T$ est quasi-compact (I, 1.1.1); \emph{a fortiori}, pour
tout ouvert quasi-compact $T_\alpha\supset T$ de $X$,

\oldpage[IV]{360}
$f(T_\alpha)=Y$ et $f|T_\alpha$ est quasi-compact. On peut par suite
appliquer à $f|T_\alpha$ le résultat de I), et il existe un diviseur
$D_\alpha$ sur $Y$ tel que
\[
 (f|T_\alpha)^*(D_\alpha)=D'|T_\alpha;
\]
en vertu de la transitivité des images réciproques de diviseurs par
morphismes plats {\hyperlink{ega4.c04.statement.21-4-4}{(21.4.4)}}, on
a aussi $(f|T)^*(D_\alpha)=D'|T$. On en conclut que pour deux indices
$\alpha$, $\beta$ distincts, on a
\[
 (f|T)^*(D_\alpha)=(f|T)^*(D_\beta),
\]
d'où (par {\hyperlink{ega4.c04.prop.21-4-8}{(21.4.8)}}),
$D_\alpha=D_\beta$. Si $D$ est la valeur commune des $D_\alpha$, on a
donc $f^*(D)|T_\alpha=D'|T_\alpha$ pour tout $\alpha$, et par suite
$f^*(D)=D'$.
\end{proof}

\begin{corollary}[21.4.10]
\EGAReaderTarget{ega4.c04.cor.21-4-10}{IV.21.4.10}
Supposons vérifiées les hypothèses générales de
{\hyperlink{ega4.c04.prop.21-4-9}{(21.4.9)}}, et en outre supposons que
$Y$ soit normal et que pour tout $y\in Y$ tel que
$\dim(\sh{O}_{Y,y})=1$, la fibre $X_y=f^{-1}(y)$ n'ait pas de cycle
premier associé immergé. Pour qu'un diviseur $D'$ sur $X$ soit de la
forme $f^*(D)$, où $D\in\operatorname{Div}(Y)$, il faut et il suffit
que les deux conditions suivantes soient vérifiées :
\begin{enumerate}[label=\textup{\arabic*\textdegree}]
\item pour tout point maximal $y$ de $Y$, l'image réciproque de $D'$
dans $\operatorname{Div}(X_y)$ soit nulle;
\item pour tout $y\in Y$ tel que $\dim(\sh{O}_{Y,y})=1$, il existe un
entier $n_y$ tel que, pour tout point maximal $x$ de $X_y$, l'image
réciproque de $D'$ dans
$\operatorname{Div}(\operatorname{Spec}(\sh{O}_{X,x}))$ soit égale à
$\operatorname{div}(t_y^{n_y})$, où $t_y$ est l'image par
l'homomorphisme $\sh{O}_{Y,y}\to\sh{O}_{X,x}$ d'une uniformisante de
l'anneau de valuation discrète $\sh{O}_{Y,y}$.
\end{enumerate}
\end{corollary}

\begin{proof}
Appliquons le critère {\hyperlink{ega4.c04.prop.21-4-9}{(21.4.9)}}.
Aux points maximaux $y$ de $Y$, on a $\operatorname{Div}(Y_1)=0$
puisque $\sh{O}_{Y,y}$ est artinien, et la condition 1\textdegree{} de
(21.4.10) est équivalente au critère de (21.4.9) en un tel point. Il
faut voir que les conditions 1\textdegree{} et 2\textdegree{} impliquent
de même la condition de (21.4.9) en un point $y\in Y$ où
$\operatorname{prof}(\sh{O}_{Y,y})=1$; comme $Y$ est normal, ces points
sont aussi ceux où $\dim(\sh{O}_{Y,y})=1$, ou encore où
$\sh{O}_{Y,y}$ est un anneau de valuation discrète
{\hyperref[canon-index.ega4.c04.thm.5-8-6]{(5.8.6)}}. On peut donc se borner au
cas où $Y=\operatorname{Spec}(\sh{O}_{Y,y})$. Si $t$ est une
uniformisante de $\sh{O}_{Y,y}$, l'application
$n\mapsto\operatorname{div}(t^n)$ est un isomorphisme de $\mathbb Z$
sur $\operatorname{Div}(Y)$; la condition de (21.4.9) signifie que
$D'$ est de la forme $\operatorname{div}(t^n\circ f)$ pour un
$n\in\mathbb Z$. Cela montre aussitôt que cette condition entraîne la
condition 2\textdegree{} de l'énoncé. Inversement, supposons cette
dernière vérifiée (ainsi que la condition 1\textdegree{}); posant pour
simplifier $n=n_y$, tout revient à voir que l'on a alors
$D'=\operatorname{div}(t^n\circ f)$. Si $\eta$ est le point générique
de $Y$, les germes de $D'$ et de $\operatorname{div}(t^n\circ f)$ en
un point quelconque de $X_\eta$ sont tous deux nuls en vertu de
1\textdegree{}. De même, les germes de $D'$ et
$\operatorname{div}(t^n\circ f)$ en un point maximal de $X_y$ sont
égaux par la condition 2\textdegree{}. Mais en un point $x\in X_y$ qui
n'est pas maximal, on a
$\operatorname{prof}(\sh{O}_{X_y,x})\geq1$ puisque $X_y$ n'a pas de
cycle premier associé immergé, donc aussi
$\operatorname{prof}(\sh{O}_{X,x})\geq2$ par
{\hyperref[canon-index.ega4.c04.prop.6-3-1]{(6.3.1)}}. On a donc
$D'=\operatorname{div}(t^n\circ f)$ en vertu de
{\hyperlink{ega4.c04.prop.21-1-8}{(21.1.8)}}.
\end{proof}

\begin{corollary}[21.4.11]
\EGAReaderTarget{ega4.c04.cor.21-4-11}{IV.21.4.11}
Supposons vérifiées les hypothèses générales de
{\hyperlink{ega4.c04.prop.21-4-9}{(21.4.9)}}, et supposons en outre que
$Y$ soit normal et que pour tout $y\in Y$ tel que
$\dim(\sh{O}_{Y,y})=1$, la fibre $X_y$ soit intègre. Alors, pour qu'un
diviseur $D'$ soit de la forme $f^*(D)$, où
$D\in\operatorname{Div}(Y)$, il faut et il suffit que, pour tout point
maximal $y$ de $Y$, l'image réciproque de $D'$ dans
$\operatorname{Div}(X_y)$ soit nulle.
\end{corollary}

\begin{proof}
Il suffit de vérifier que la condition 2\textdegree{} de
{\hyperlink{ega4.c04.cor.21-4-10}{(21.4.10)}} est alors
automatiquement vérifiée. Or, si $\dim(\sh{O}_{Y,y})=1$ et si $X_y$ est
intègre, on a, en l'unique point maximal $x$ de $X_y$,
$\dim(\sh{O}_{X,x})=1$ en vertu de
{\hyperref[canon-index.ega4.c04.cor.6-1-3]{(6.1.3)}}; en outre, comme
$\sh{O}_{Y,y}$ est un anneau régulier et $\sh{O}_{X_y,x}$ un corps,
$\sh{O}_{X,x}$ est régulier
{\hyperref[canon-index.ega4.c04.prop.6-5-1]{(6.5.1)}}, donc un anneau de valuation
discrète, et si $t$ est une uniformisante de $\sh{O}_{Y,y}$, son image
$t_y$ dans $\sh{O}_{X,x}$ est une uniformisante de $\sh{O}_{X,x}$
($0$, 17.3.3); l'image réciproque de $D'$ dans
$\operatorname{Div}(\operatorname{Spec}(\sh{O}_{X,x}))$ est donc bien
de la forme

\oldpage[IV]{361}
$\operatorname{div}(t_y^n)$, ce qui n'est autre que la condition
2\textdegree{} de {\hyperlink{ega4.c04.cor.21-4-10}{(21.4.10)}},
puisque $X_y$ n'a qu'un seul point maximal.
\end{proof}

\begin{remark}[21.4.12]
\label{IV.21.4.12}
Si dans {\hyperlink{ega4.c04.cor.21-4-11}{(21.4.11)}} on suppose
seulement que les $X_y$ sont réduits, il faut ajouter à la condition de
l'énoncé que les images réciproques de $D'$ dans chacun des ensembles
$\operatorname{Div}(\operatorname{Spec}(\sh{O}_{X,x}))$, où $x$
parcourt l'ensemble des points maximaux de $X_y$, sont toutes de la
forme $\operatorname{div}((t'_x)^n)$ où $t'_x$ est une uniformisante
de $\sh{O}_{X,x}$, l'entier $n$ étant le même pour tous ces points.
\end{remark}

\begin{corollary}[21.4.13]
\label{IV.21.4.13}
Supposons vérifiées les conditions de
{\hyperlink{ega4.c04.cor.21-4-11}{(21.4.11)}}. Pour tout
$\sh{O}_X$-Module inversible $\sh{L}$ tel que, pour tout point maximal
$\eta$ de $Y$, l'image réciproque
\[
 \sh{L}_\eta=\sh{L}\otimes_{\sh{O}_X}\sh{O}_{X_\eta}
\]
soit isomorphe à $\sh{O}_{X_\eta}$, il existe un
$\sh{O}_Y$-Module inversible $\sh{M}$ et un isomorphisme de
$f^*(\sh{M})$ sur $\sh{L}$.
\end{corollary}

\begin{proof}
Il suffit de montrer qu'il existe une section méromorphe régulière $s$
de $\sh{L}$ au-dessus de $X$ telle que le diviseur $D'$ correspondant
à $(\sh{L},s)$ soit tel que $D'_x=0$ en tout point $x$ d'une fibre
$X_\eta$ d'un point maximal $\eta$ de $Y$. On appliquera alors
{\hyperlink{ega4.c04.cor.21-4-11}{(21.4.11)}} à $D'$, et la conclusion
résultera de {\hyperlink{ega4.c04.def.21-4-2}{(21.4.2)}}.

Notons que les points de $\operatorname{Ass}(\sh{O}_X)$ sont au-dessus
des points maximaux de $Y$ {\hyperref[canon-index.ega4.c04.prop.3-3-1]{(3.3.1)}};
comme $Y$ est réduit et localement noethérien, il y a un ouvert dense
$V$ de $Y$, réunion d'une famille disjointe d'ouverts affines contenant
chacun un seul point maximal de $Y$, et $f^{-1}(V)$ est
schématiquement dense dans $X$
{\hyperref[canon-index.ega4.c04.prop.5-10-2]{(5.10.2)}}; on peut donc se borner au
cas où $Y$ est affine et intègre. Puisque, pour l'unique point maximal
$\eta$ de $Y$, $X_\eta$ est alors intègre,
$\operatorname{Ass}(\sh{O}_X)$ est réduit à un seul point $z$, qui est
le point générique de $X$. Soit $s_\eta$ la section de $\sh{L}_\eta$
au-dessus de $X_\eta$ qui correspond par isomorphie à la section unité
de $\sh{O}_{X_\eta}$, et soient $U$ un voisinage ouvert affine d'un
point $x\in X_\eta$ dans $X$ et $s$ une section de $\sh{L}$ au-dessus
de $U$ ayant même germe que $s_\eta$ au point $x$. Comme $U$ est
schématiquement dense dans $X$
({\hyperlink{ega4.c04.statement.20-2-13}{20.2.13}}, (iv)), $s$
s'identifie à une section méromorphe régulière bien déterminée de
$\sh{L}$ au-dessus de $X$, et il est clair
({\hyperlink{ega4.c04.prop.20-2-11}{20.2.11}} et
{\hyperlink{ega4.c04.prop.20-3-5}{20.3.5}}) que $s$ induit sur $X_\eta$
une section méromorphe régulière de $\sh{L}_\eta$ égale à $s_\eta$, et
répond donc à la question.
\end{proof}

\bigskip
\begin{center}
\small
1967. --- Imprimerie des Presses Universitaires de France. --- Vendôme
(France)\\
ÉDIT. N\textsuperscript{o} 29 470\qquad
IMPRIMÉ EN FRANCE\qquad
IMP. N\textsuperscript{o} 20 298
\end{center}

% FIN SOURCE ega4/ega4-errata-addenda-list3-fr.tex

\end{document}
